% YM4 proof sketch; revised source-supported construction arguments.
\documentclass[11pt]{amsart}
\usepackage[margin=1.05in]{geometry}
\usepackage{amsmath,amssymb,mathrsfs,longtable}
\usepackage{hyperref}
\hypersetup{hidelinks,pdftitle={Four-dimensional SU(2) Yang-Mills theory: an outline of the proof},pdfauthor={}}
\allowdisplaybreaks
\numberwithin{equation}{section}
\newtheorem{theorem}{Theorem}[section]
\newtheorem{proposition}[theorem]{Proposition}
\newtheorem{lemma}[theorem]{Lemma}
\newtheorem{corollary}[theorem]{Corollary}
\theoremstyle{definition}
\newtheorem{definition}[theorem]{Definition}
\theoremstyle{remark}
\newtheorem{remark}[theorem]{Remark}
\newcommand{\R}{\mathbb R}
\newcommand{\C}{\mathbb C}
\newcommand{\Z}{\mathbb Z}
\newcommand{\N}{\mathbb N}
\newcommand{\E}{\mathbb E}
\newcommand{\Sspace}{\mathcal S}
\newcommand{\cS}{\mathcal S}
\newcommand{\dd}{\mathrm d}
\DeclareMathOperator{\supp}{supp}
\DeclareMathOperator{\tr}{tr}
\DeclareMathOperator{\diag}{diag}

% Helpers retained for the outline's existing displays and argument codes.
\usepackage{array,graphicx}
\DeclareOldFontCommand{\cal}{\relax}{\mathcal}
\setlength{\emergencystretch}{3em}
\newsavebox{\MathBox}
\newcommand{\FitMath}[1]{\sbox{\MathBox}{$\displaystyle #1$}\ifdim\wd\MathBox>\linewidth\typeout{MATHFIT line=\the\inputlineno\space width=\the\wd\MathBox}\resizebox{\linewidth}{!}{\usebox{\MathBox}}\else\usebox{\MathBox}\fi}
\newcommand{\formulacode}[1]{\par\nopagebreak\noindent\hfill{\footnotesize\textup{(#1)}}\par}
\newcommand{\candidatebanner}[1]{{\small\noindent\textbf{Proof outline.} The supporting analytic results used in this section, including the obligations that we must still prove, are listed under #1 in the dependency catalogue. Their hypotheses and precise forms are retained at the points where the argument uses them.\par}\medskip}
\theoremstyle{definition}
\newtheorem*{maintheorem}{Theorem 0}
\title[Four-dimensional SU(2) Yang--Mills theory]{Four-dimensional $\mathrm{SU}(2)$ Yang--Mills theory\newline An outline of the proof for the $\operatorname{tr}(F^2)$ observable}
\author{}
\date{}

\begin{document}
\pagenumbering{Alph}

% BEGIN MODULE chapters/main_theorem_page.tex
\begin{abstract}
This manuscript outlines the proof of the continuum construction of the
renormalized $\operatorname{tr}(F^2)$ observable in four-dimensional Wilson
$\mathrm{SU}(2)$ theory. The argument connects multiscale estimates,
Euclidean correlation limits, and scalar reconstruction.
\end{abstract}
\maketitle
\setcounter{tocdepth}{1}
\tableofcontents
\clearpage

Here $\operatorname{tr}(F^2)$ denotes the renormalized scalar
curvature-squared observable, with the Wilson normalization in
Section~\ref{ch:foundations}. At each lattice cutoff, we use
$\mathcal O(f)-\mathbb E\mathcal O(f)$: the expectation of the smeared
observable is subtracted before taking the continuum limit.

\phantomsection\label{page:main-theorem}
\addcontentsline{toc}{section}{Main theorem}
\begin{maintheorem}
For sufficiently small reference coupling, the correlations of
$\operatorname{tr}(F^2)$ in four-dimensional Wilson $\mathrm{SU}(2)$
theory have nontrivial volume-first subsequential continuum limits.
The full moment functionals $G_{\xi,n}$ of every canonical component,
normalized by $G_{\xi,0}=1$, satisfy the following
Osterwalder--Schrader axioms.

Let $\mathcal S_{\ne,n}$ be the Schwartz functions on
$(\mathbb R^4)^n$ whose every derivative vanishes whenever two
arguments coincide. Let $\mathcal S_{+,n}$ be the subspace also
flat outside the region where all time coordinates are positive.
Set $\mathcal S_{\ne,0}=\mathcal S_{+,0}=\mathbb C$, and write
\[
 |\varphi|_r=\max_{|\alpha|\le r}\sup_x
       (1+|x|^2)^{r/2}|\partial^\alpha\varphi(x)|.
\]
\begin{enumerate}
\renewcommand{\labelenumi}{(\roman{enumi})}
\item \textbf{Regularity and growth.}
Each $G_{\xi,n}$ is continuous on $\mathcal S_{\ne,n}$.
There are constants $A,B>0$, $p\ge0$ and an integer $k\ge1$,
independent of $n$ and $\xi$, such that
\[
 |G_{\xi,n}(\varphi)|\le A B^n(n!)^p|\varphi|_{kn},
 \qquad n\ge1.
\]

\item \textbf{Euclidean invariance.}
The functionals are invariant under simultaneous translations of
all arguments and every rotation or reflection in $\mathrm O(4)$.

\item \textbf{Reflection positivity.}
Let $\theta(x_1,x_2,x_3,x_4)=(x_1,x_2,x_3,-x_4)$ and put
\[
 (\Theta\overline\varphi)(x_1,\ldots,x_n)
       =\overline{\varphi(\theta x_1,\ldots,\theta x_n)}.
\]
For every finite sequence of tests
$\varphi_n\in\mathcal S_{+,n}$, $0\le n\le N$,
\[
 \sum_{m,n=0}^{N}
 G_{\xi,m+n}(\Theta\overline\varphi_m\otimes\varphi_n)\ge0.
\]

\item \textbf{Symmetry and reality.}
The functionals are invariant under permutations of the arguments,
and
\[
 G_{\xi,n}(\overline\varphi)=\overline{G_{\xi,n}(\varphi)}.
\]

\item \textbf{Clustering.}
Two groups of observables factorize when their spatial separation
tends to infinity. Explicitly, for
$\varphi\in\mathcal S_{+,m}$ and $\psi\in\mathcal S_{+,n}$,
\[
 G_{\xi,m+n}(\Theta\overline\varphi\otimes\tau_{(v,0)}\psi)
 \longrightarrow
 \overline{G_{\xi,m}(\varphi)}\,G_{\xi,n}(\psi)
 \qquad (|v|\longrightarrow\infty),
\]
where $v\in\mathbb R^3$ and $\tau_{(v,0)}$ translates all arguments
of $\psi$ by the same spatial vector.
\end{enumerate}
\end{maintheorem}

The precise construction and component statements are Theorems~0
and~0' in Section~\ref{candidate:orientation}; the test-space conventions
are given in Section~\ref{ch:precise-reconstruction}.

Here volume-first means taking the infinite-volume limit at each fixed
lattice spacing before sending the spacing to zero along the stated
admissible weak-coupling sequence. A canonical component is one phase
of the limiting translation-invariant state. It retains the reference
subtraction used at the cutoffs, so its one-point function need not
vanish. The conclusion concerns the scalar sector generated by
$\operatorname{tr}(F^2)$.

\paragraph*{Proof sketch.}
The multiscale construction supplies source bounds uniform in cutoff
and volume. Compactness and quantitative extension give the Euclidean
correlations and their factorial growth on diagonal-flat Schwartz
tests. The mirror, angular and phase arguments establish reflection
positivity, Euclidean invariance and clustering in each component;
symmetry and reality pass from the lattice correlations.

% END MODULE chapters/main_theorem_page.tex


% BEGIN MODULE chapters/frontmatter.tex
\section*{Plan of the proof}

This manuscript outlines the proof of the continuum construction of the
renormalized $\operatorname{tr}(F^2)$ observable in four-dimensional Wilson
$\mathrm{SU}(2)$ theory. Its lattice mean is subtracted before taking the
continuum limit. The resulting theory is the scalar sector generated by
this observable. Theorems 0 and 0' give the full statements.

The outline presents the main constructions, estimates, and connections
between the steps. It develops the central mechanisms in detail and gives
shorter arguments for supporting results; full derivations are reserved
for the fuller proof. Definitions and detailed technical statements appear
at the points where they enter the argument.

\subsection*{The multiscale construction}

Block averaging provides the coordinates for the proof. Local pivot charts
straighten the average near regular configurations and give the coarea and
determinant formulas. The classical flat problem supplies the constrained
gap, curl-energy comparison, average lift, and exponential kernels used in
the many-level background construction.

The regular integration combines local Gaussian comparison with a polymer
expansion. The pointed Hessian and nonlinear window estimates control the
extracted coefficients and the remainder. The large-field stage sums
histories while retaining the Wilson weight and the reserve needed for
source estimates. Together these steps lead to ordinary source stability
and the partition floor on the localization geometry of the construction.

\subsection*{Sources, limits, and short distance}

Mixed source differentiation selects connected contributions meeting every
source support. Cauchy estimates and geometric counting give bounds
uniform in cutoff and volume. Compactness then extracts continuum
correlation forms on separated compact tests. The quantitative extension
passes to Schwartz tests flat on all coincidence diagonals.

For the volume-first construction, the infinite-volume limit is taken at
each fixed lattice spacing through compatible periodic tori. The spacing
then tends to zero along the stated admissible weak-coupling sequence.
The two-point extraction and running normalization give the controlled
short-distance law, whose lower bound establishes nontriviality.

\subsection*{Euclidean structure and reconstruction}

Lattice mirror inequalities give reflection positivity. Uniform insertion
estimates, plate densities, angular kernels, and ultraviolet harmonic
elimination yield Euclidean covariance. The phase decomposition identifies
the canonical components and their clustering properties.

The growth estimates, covariance, positivity, and clustering enter the
analytic-continuation step on the diagonal-flat test domain. Reconstruction
gives, for each canonical component, a local Hermitian scalar Wightman
theory on a separable Hilbert space, with Poincar\'e covariance,
forward-cone spectrum, and a unique vacuum. Its cyclic sector is generated
by the renormalized $\operatorname{tr}(F^2)$ observable.

\subsection*{Reading the outline}

Codes such as ``4.8'', ``6c.2'', and ``8a.1.11'' track arguments across
sections. The argument index locates their statements and uses; the
notation appendix fixes the multiscale objects. The dependency catalogue
collects the detailed analytic statements, and
Section~\ref{ch:precise-reconstruction} gives the reconstruction step.

% END MODULE chapters/frontmatter.tex

\clearpage
\pagenumbering{arabic}


% BEGIN MODULE chapters/foundations.tex
\section{Finite-cutoff foundations}
\label{ch:foundations}

This section establishes the finite-dimensional identities and compactness
arguments used throughout the construction. None of its statements assumes
a continuum Yang--Mills measure. A constant depending on the number of
plaquettes is not a bound uniform in the cutoff. We retain this distinction
even when every integral under discussion is finite.

\subsection{The configuration space and its probability measure}

Let $G=\mathrm{SU}(2)$, the compact group of complex $2\times2$ matrices
$U$ with $U^*U=1$ and $\det U=1$. Its normalized Haar probability measure
is denoted by $dg$. Fix a finite oriented graph obtained from a cubic lattice
on a four-dimensional torus. Choose one orientation of each geometric bond;
write $B_+$ for this set and impose $U_{\bar b}=U_b^{-1}$ for the reverse
orientation. The configuration space is the compact space
$\mathcal V=G^{B_+}$, with product probability measure $dU$.

A path is an ordered list of consecutive oriented bonds. Its transport is
the ordered product of their matrices. For each geometrical elementary
square choose one boundary orientation and one starting vertex, and denote
its transport by $U_p$. Define

\begin{gather*}
\FitMath{a_p(U)=1-\tfrac12\operatorname{Re}\operatorname{tr}U_p,} \\[3pt]
\FitMath{A(U)=\sum_{p\in P}a_p(U),} \\[3pt]
\FitMath{N_P=|P|.}
\end{gather*}

Changing the starting vertex conjugates $U_p$, and reversing orientation
inverts it. Both changes preserve $a_p$.

\begin{proposition}[Finite-cutoff existence and elementary bounds]
For every real $\beta\ge0$,

\begin{gather*}
\FitMath{e^{-2\beta N_P}\le Z(\beta):=\int e^{-\beta A(U)}dU\le1.}
\end{gather*}

Consequently $Z(\beta)^{-1}e^{-\beta A(U)}dU$ is a well-defined probability
measure. Every bounded Borel observable has a finite expectation.
\end{proposition}
\begin{proof}
The eigenvalues of an element of $G$ are $e^{it},e^{-it}$ for a real $t$.
Thus $\frac12\operatorname{Re}\operatorname{tr}U_p=\cos t$ lies in $[-1,1]$,
so $0\le a_p\le2$ and $0\le A\le2N_P$. Exponentiate these inequalities and
integrate against a probability measure. The lower bound is strictly
positive and supplies the normalization.
\end{proof}

The proof uses only compactness and positivity. Its constants degenerate
when $N_P\to\infty$. In particular it supplies neither tightness of
renormalized local fields nor a useful estimate for a continuum limit.

A gauge transformation is a family $(h_x)$ of group elements indexed by
vertices. It acts by $U_b\mapsto h_{b_-}U_bh_{b_+}^{-1}$. Successive factors
cancel along a path, so a closed transport based at $x$ is conjugated by
$h_x$. Haar measure on each bond is unchanged by its independent left and
right multiplications. This proves gauge invariance of the action and of
the probability measure. A permutation of vertices and bonds coming from
a lattice isometry also preserves the product measure and permutes the
plaquettes. Therefore the finite torus measure has every isometry of that
torus. This conclusion concerns these finite torus measures; it is not a
proof that every infinite-volume periodic Gibbs state has every reflection.

\subsection{The continuum normalization is a local expansion}

Represent the Lie algebra by Hermitian traceless matrices and set
$|F|^2=\frac12\operatorname{tr}(F^2)$. For a fixed such matrix,

\begin{gather*}
\FitMath{1-\tfrac12\operatorname{Re}\operatorname{tr}e^{itF}
   =\tfrac{t^2}{4}\operatorname{tr}(F^2)+O(t^4|F|^4)
   =\tfrac12t^2|F|^2+O(t^4|F|^4).}
\end{gather*}

The expansion follows by the matrix exponential series: the linear trace
vanishes and the odd powers have zero real trace. If a smooth connection
has elementary holonomy $U_p=\exp(i a^2F_{\mu\nu}(x_p)+O(a^3))$, with all
required derivatives uniformly bounded, it follows that

\begin{gather*}
\FitMath{a_p=\tfrac12a^4|F_{\mu\nu}(x_p)|^2+O(a^5).}
\end{gather*}

Hence, for a fixed continuous test function $f$ supported in a compact region,
summing $f(x_p)a_p$ over plaquettes gives the corresponding Riemann integral of
$\frac12 f\sum_{\mu<\nu}|F_{\mu\nu}|^2$ as $a\downarrow0$.

This is a deterministic smooth-field calculation. Typical configurations
under a quantum measure have not been shown to satisfy its hypotheses.
The calculation fixes the intended normalization of the observable; it
does not prove convergence of that observable under the lattice measures.

\subsection{Sources, cumulants, and centering}

For bounded functions $f_i$ on the plaquette barycentres put
$O_i=\sum_p f_i(x_p)a_p$ and

\begin{gather*}
\FitMath{Z(z)=\int\exp\{-\beta A(U)+\sum_{i=1}^n z_iO_i(U)\}\,dU,} \\[3pt]
\FitMath{z\in\mathbb C^n.}
\end{gather*}

The estimate $|O_i|\le2\sum_p|f_i(x_p)|$ gives an integrable uniform majorant
on every compact source set. Termwise differentiation of the exponential
series is therefore valid, and $Z$ is entire. Since $Z(0)>0$, continuity
gives a sufficiently small polydisc about zero where $Z$ has no zeros. On that polydisc
there is a unique holomorphic logarithm $W$ with $W(0)=\log Z(0)$.

The joint cumulant is

\begin{gather*}
\FitMath{\kappa(O_1,\ldots,O_n)=
 \left.\partial_{z_1}\cdots\partial_{z_n}W(z)\right|_{z=0}.}
\end{gather*}

Its expression in moments is

\begin{gather*}
\FitMath{\kappa(O_1,\ldots,O_n)=
 \sum_{\pi\in\operatorname{Part}(n)}
 (|\pi|-1)!(-1)^{|\pi|-1}
 \prod_{J\in\pi}\mathbb E\prod_{i\in J}O_i.}
\end{gather*}

Indeed, write $Z(z)/Z(0)=1+u(z)$ and use
$\log(1+u)=\sum_{m\ge1}(-1)^{m-1}u^m/m$ near zero. The coefficient of
$z_1\cdots z_n$ in $u^m$ sums over ordered partitions into $m$ nonempty
blocks. Dividing by $m$ turns their $m!$ orderings into $(m-1)!$.
Conversely, expanding $e^{W(z)-W(0)}$ proves the moment--cumulant relation

\begin{gather*}
\FitMath{\mathbb E\prod_iO_i=\sum_{\pi\in\operatorname{Part}(n)}
 \prod_{J\in\pi}\kappa((O_i)_{i\in J}).}
\end{gather*}


Replacing $O_i$ by $O_i-c_i$ changes the logarithm only by
$-\sum_i z_ic_i$. Thus every cumulant of order at least two is unchanged.
If $c_i=\mathbb E O_i$, every singleton cumulant vanishes. The centered
moment functional is consequently

\begin{gather*}
\FitMath{\mathbb E\prod_i(O_i-\mathbb E O_i)
 =\sum_{\pi\in\operatorname{Part}_{\ge2}(n)}
 \prod_{J\in\pi}\kappa((O_i)_{i\in J}).}
\end{gather*}

There are finitely many partitions at fixed $n$. Convergence of every
separated cumulant of order at least two therefore implies convergence of
the centered moments on the same tuples. Positivity must be applied to
these moments, not to the cumulants separately.

The source is exactly a local coupling shift:

\begin{gather*}
\FitMath{-\beta A+\sum_i z_iO_i=-\sum_p(\beta-w(x_p))a_p,} \\[3pt]
\FitMath{w=\sum_i z_if_i.}
\end{gather*}

For a constant source $c$, $Z_\beta(c+w)=Z_{\beta-c}(w)$. This algebraic
identity holds before any renormalization procedure is chosen.

\subsection{What a mass-preserving transformation preserves}

Let $T$ be a fixed linear map on finite measures on a measurable space,
with $(T\nu)(1)=\nu(1)$ for all measures in its domain. If
$\nu_z=e^{\sum_i z_iO_i}\nu$ belongs to that domain, then

\begin{gather*}
\FitMath{(T\nu_z)(1)=\nu_z(1).}
\end{gather*}

This proves exact equality of sourced partition functions, and hence of
their derivatives when the analytic dependence can be justified. It does
not identify $T\nu$ as the distribution of a predetermined block variable.
For that stronger conclusion one needs a measurable map $M$ and the
identity $(T\nu)(F)=\nu(F\circ M)$ for every bounded measurable $F$.

For example, on the two-point set $\{0,1\}$ let
$T\nu=\nu(\{0,1\})\delta_0$. This positive linear map preserves total mass
but changes the expectation of the coordinate from $\nu(\{1\})$ to zero
when $\nu$ is a probability measure. Thus mass preservation alone contains
no statement about observable expectations. Sources must be transported
through the exact integral, with the choices of the transformation fixed
throughout the source family.

\subsection{A fixed-spacing infinite-volume limit}

Consider the countable cubic lattice at a fixed positive spacing. Its
configuration space is $\mathcal V_\infty=G^{B_+}$ with the product
topology. This is a compact metrizable space: enumerate the bonds and use,
for example, the sum of $2^{-j}$ times bounded compatible distances on
their group coordinates. The space of Borel probability measures is
sequentially compact in the weak topology.

For a finite set of bonds $D$, define the conditional specification

\begin{gather*}
\FitMath{(\gamma_DF)(U)=
 \frac{\int F(V_D,U_{D^c})\exp[-\beta A_D(V_D,U_{D^c})]\,dV_D}
 {\int\exp[-\beta A_D(V_D,U_{D^c})]\,dV_D},} \\[3pt]
\FitMath{A_D=\sum_{p:p\cap D\ne\varnothing}a_p.}
\end{gather*}

The sum is finite and the denominator is bounded below by
$e^{-2\beta N(D)}$, where $N(D)$ is the number of plaquettes meeting $D$.
For continuous cylinder $F$, numerator and denominator depend continuously
on only finitely many boundary bonds; dominated convergence proves that
$\gamma_DF$ is another continuous cylinder function.

\begin{proposition}[DLR existence at fixed lattice parameters]
Take torus Wilson measures at fixed spacing and fixed $\beta$, and extend
their bond configurations periodically to the infinite lattice. Every
sequence of torus side lengths tending to infinity has a weakly convergent
subsequence. Each limit $\nu$ satisfies $\nu(\gamma_DF)=\nu(F)$ for every
finite $D$ and bounded Borel cylinder function $F$. It is invariant under
all lattice translations and all lattice isometries preserved by the
finite torus measures.
\end{proposition}
\begin{proof}
Compactness gives the subsequence. Fix $D$ and a continuous cylinder $F$.
For sufficiently large tori the finite set consisting of $D$, all adjacent
plaquettes and the support of $F$ embeds without identifications. The
finite-volume conditional expectation identity then reads
$\mu_\ell(\gamma_DF)=\mu_\ell(F)$ after periodic extension. Both integrands
are continuous cylinder functions, so the identity passes to the weak
limit. Continuous functions on the finite-dimensional compact coordinate
spaces determine measures. The identity therefore extends first to bounded
Borel cylinder functions, and, if desired, to all bounded Borel functions
by the monotone class theorem. Invariance passes to the limit by the same
continuous-cylinder argument, since a fixed lattice isometry acts
continuously on the product space.
\end{proof}

This proposition does not change the spacing or the coupling. It provides
no bounds uniform as the spacing tends to zero and gives no uniqueness or
extremality of the limiting state. A theorem about all periodic Gibbs
states is stronger than a theorem about these torus limit points.

\subsection{Compactness of separated correlation forms}

Let $\mathbb E$ be $\mathbb R^4$ or a fixed torus. A separated tuple consists
of finitely many compactly supported $C^1$ functions with pairwise disjoint
supports. Their compactness makes the minimum distance between nonempty
supports positive. Fix the norm
$\|f\|_1=\max(\|f\|_\infty,\|\nabla f\|_\infty)$.

\begin{proposition}[A diagonal extraction principle]
Suppose $T_{n,j}$ are symmetric multilinear forms on separated tuples and
for every fixed $n$, $R<\infty$ and $d>0$ there is a finite $C(n,R,d)$,
independent of all sufficiently large $j$, such that

\begin{gather*}
\FitMath{|T_{n,j}(f_1,\ldots,f_n)|\le C(n,R,d)\prod_i\|f_i\|_1}
\end{gather*}

whenever the supports lie in $B(0,R)$ and are separated by at least $d$.
Then a subsequence converges at every separated tuple, simultaneously for
all $n$. The limits obey the same bounds.
\end{proposition}
\begin{proof}
Choose all finite families of finite unions of rational closed boxes with positive mutual
distance. The $C^1$ functions supported in each union form a separable normed
space, being a subspace of a finite product of spaces of continuous
functions on a compact set. Choose countable dense subsets. There are
countably many finite tuples drawn from these subsets. Each sequence of
values is eventually bounded by the assumed estimate. Bolzano--Weierstrass
and a successive diagonal extraction give a common subsequence on which
all these values converge.

Given any separated tuple, enclose its supports in pairwise separated
finite unions of rational boxes, still inside a fixed bounded region.
Approximate each function in $C^1$ by elements of the chosen dense sets.
The telescoping multilinear identity

\begin{gather*}
\FitMath{T(f_1,\ldots,f_n)-T(g_1,\ldots,g_n)
 =\sum_{i=1}^n T(g_1,\ldots,g_{i-1},f_i-g_i,f_{i+1},\ldots,f_n)}
\end{gather*}

and the uniform bound make the approximation error uniformly small along
the extracted subsequence. This proves the Cauchy property for the
original tuple. Multilinearity, symmetry and the bound pass to the limit
on each product of spaces with separated supports.
\end{proof}

The difficult Yang--Mills step, which we must prove, is the estimate
supplying $C(n,R,d)$.
Finite-cutoff boundedness with a constant involving the number of
plaquettes does not meet it. An exponentially large dependence on the
physical torus volume also does not supply a bound uniform on
$\mathbb R^4$. These are different compactness problems.

\subsection{Reflection positivity and closed limits}

Let $\theta$ be a reflection. Real reflection-invariant centering preserves
reflection positivity: replacing each $O_j(f)$ by $O_j(f)-m_j(f)$ maps every
positive-half polynomial to another such polynomial and commutes with reflection
when $m_j(\theta f)=m_j(f)$. Suppose for each cutoff a centered moment
functional $\sigma_j$ satisfies

\begin{gather*}
\FitMath{\sum_{r,s=1}^m\overline{c_r}c_s\,
 \sigma_j(\theta\overline{X_r}\cdot X_s)\ge0}
\end{gather*}

for every finite collection of complex test tuples $X_r$ in the positive
half-space, where $\overline{X_r}$ conjugates every test and $\theta$
acts by pullback. Within each $X_r$ the supports are separated; the reflected negative tuple
and the positive tuple are automatically separated from one another.
Overlap between supports from $X_r$ and $X_s$ on the same side causes no
problem, because these supports do not occur together in the same
concatenated pair $\theta\overline{X_r}\cdot X_s$.

If all the entries converge, the finite sum converges and remains
nonnegative. Thus the limiting moment functional is reflection positive
on this vector space of separated monomials. This elementary closedness statement is the
precise sense in which positivity passes to limits. Passing from that
domain to a Schwartz test space requires a density result and a continuous
bound for the limiting distributions; it cannot be justified by closedness
alone.

The algebraic form $\langle F,G\rangle=\sigma(\theta\overline F\cdot G)$
is positive semidefinite. Positivity of
$\langle F+zG,F+zG\rangle$ for every complex $z$ gives
$|\langle F,G\rangle|^2\le\langle F,F\rangle\langle G,G\rangle$.
Consequently a null vector is orthogonal to every vector. The quotient by
the null space is an inner-product space, and its completion is a Hilbert
space. Defining translations on this quotient additionally requires that
they preserve the null space and satisfy the appropriate bounds. Strong
continuity and a nonnegative generator are later conclusions with their
own hypotheses.

\subsection{Two limits of what scalar correlations prove}

Let $X$ be a real standard Gaussian variable and put $Y=X^2-1$. For $t<1/2$
near zero,

\begin{gather*}
\FitMath{\log\mathbb E e^{tY}=-t-\tfrac12\log(1-2t)
  =\sum_{n\ge2}\frac{2^{n-1}}n t^n.}
\end{gather*}

It follows that $\kappa_n(Y)=2^{n-1}(n-1)!$ for every $n\ge2$.
In particular a nonzero connected four-point function of a quadratic
observable is compatible with a Gaussian underlying theory. For distinct
centered jointly Gaussian variables the same conclusion is expressed by
the connected cycle contractions of their Wick squares.

Accordingly a non-Gaussian action-density field is an important property
of that field, but is not by itself a characterization of an interacting
Yang--Mills gauge theory. Reconstruction from the action-density
correlations produces the cyclic sector generated by that observable.
Identifying a richer gauge-invariant algebra requires additional
observables and compatibility arguments.

A polynomial cluster estimate likewise gives no positive mass gap. In a
Hilbert space with a nonnegative self-adjoint $H$, a connected time
correlation has the spectral form
$\langle v,e^{-tH}v\rangle=\int_0^\infty e^{-tE}\,d\mu_v(E)$ for a vector
orthogonal to the vacuum. If $d\mu_v(E)=E e^{-E}\,dE$, the integral is
$(1+t)^{-2}$. It decays, yet its spectral measure meets every interval
$(0,\varepsilon)$. The support condition required for a gap would instead
exclude all such intervals below one fixed positive threshold.

These examples clarify the scope of the construction and the additional
estimates needed for richer observable algebras and a positive spectral
gap. The ultraviolet estimates enter with their own stated hypotheses.

% END MODULE chapters/foundations.tex


% BEGIN MODULE chapters/block_normalization_v2.tex
\section{Block normalization and a uniform scalar gap}
\subsection{Block means, adjoints, and a uniform scalar Poincar\'e inequality}

An adjoint depends on the inner products of both its domain and its range.
This matters when a block mean maps a fine lattice to a coarse lattice. The
following finite-dimensional calculation fixes the convention and gives the
scalar estimate used in a gauge-fixed Gaussian comparison. It applies to real
or complex scalar functions; a finite number of color components are handled
by summing the same estimate component by component.

Fix an integer $m\ge1$ and put $\eta=m^{-1}$. A block consists of the sites

\begin{gather*}
\FitMath{B=\{0,\eta,\ldots,(m-1)\eta\}^{4}.}
\end{gather*}

Its associated physical cell has side $m\eta=1$ and volume one. Translation of
the block, including translation by half a fine spacing to use cell centers,
does not change the calculation. Let $\mathcal B$ be a finite disjoint family
of such blocks, for example the blocks of a periodic torus whose side is an
integer number of coarse cells. Write $\Lambda_f$ for their union and
$\Lambda_c$ for the set indexing the blocks. A coarse function has one value
per block. We count every fine bond once, with positive orientation. The block
Neumann form uses the intrinsic open rectangular graph. In a chosen lift
$B=x_B+\eta\{0,\ldots,m-1\}^{4}$, include exactly the bonds
$x_B+\eta j\longrightarrow x_B+\eta(j+e_\mu)$ with
$j\in\{0,\ldots,m-1\}^{4}$ and $j_\mu\le m-2$.
Thus there are no internal bonds if $m=1$. Exclude opposite-face periodic
wrap bonds even when their two torus endpoints belong to the same block.
The ambient periodic $d_\eta$ retains all fine bonds, with their usual
positive-orientation multiplicities.

The physical inner products are

\begin{gather*}
\FitMath{\langle f,g\rangle_f=\eta^4\sum_{x\in\Lambda_f}\overline{f(x)}g(x),} \\[3pt]
\FitMath{\langle u,v\rangle_c=\sum_{B\in\mathcal B}\overline{u(B)}v(B),} \\[3pt]
\FitMath{\langle a,b\rangle_e=\eta^4\sum_{e}\overline{a(e)}b(e).}
\end{gather*}

The last sum can mean all fine bonds of the periodic lattice or just the
internal bonds of the blocks; the operator being considered determines which
set is used. The same weight is used on fine sites and fine bonds. The physical
difference is
$d_\eta f(x,x+\eta e_\mu)=\eta^{-1}[f(x+\eta e_\mu)-f(x)]$.
The unscaled combinatorial difference is $d_0=\eta d_\eta$. These two operators
must not be interchanged in an estimate with a mesh-independent constant.

\begin{proposition}[The mean, its adjoint, and its projection]
Define the block mean $Q:H_f\to H_c$ by

\begin{gather*}
\FitMath{(Qf)(B)=\eta^4\sum_{x\in B}f(x)=m^{-4}\sum_{x\in B}f(x).}
\end{gather*}

For the physical inner products just specified, its adjoint is
$(Q^*u)(x)=u(B)$ for $x\in B$. Consequently $QQ^*=1$, and
$P=Q^*Q$ is the orthogonal projection onto the functions constant in each
block. In particular,

\begin{gather*}
\FitMath{\|f\|_f^2=\|(1-P)f\|_f^2+\|Qf\|_c^2.}
\end{gather*}

\end{proposition}

\begin{proof}
Directly from the definitions,

\begin{gather*}
\FitMath{\langle Qf,u\rangle_c
 =\sum_B\eta^4\sum_{x\in B}\overline{f(x)}u(B)
 =\langle f,Q^*u\rangle_f.}
\end{gather*}

Since a block contains $m^4$ sites and $\eta^4m^4=1$, its mean applied to a
constant is that constant. Thus $QQ^*=1$. It follows that $P^*=P$ and
$P^2=Q^*(QQ^*)Q=P$. Its range is precisely the block-constant functions, and
the displayed norm identity is the orthogonal decomposition by that projection.
No limiting argument is involved.
\end{proof}

\begin{proposition}[A block Neumann estimate]
Let $d_{\eta,N}$ denote the physical difference restricted to internal bonds
of each block. If $m\ge2$, then

\begin{gather*}
\FitMath{\|(1-P)f\|_f^2
 \le \frac{\eta^2}{4\sin^2(\pi\eta/2)}
                 \|d_{\eta,N}f\|_e^2
 \le \tfrac14\|d_{\eta,N}f\|_e^2.}
\end{gather*}

For $m=1$ the left side is zero. If $d_\eta$ includes bonds between blocks,
the same bound holds with $d_\eta$ on the right. With the combinatorial
difference the exact first bound is instead

\begin{gather*}
\FitMath{\|(1-P)f\|_f^2
 \le \frac{1}{4\sin^2(\pi/(2m))}\|d_{0,N}f\|_e^2.}
\end{gather*}

\end{proposition}

\begin{proof}
First consider a one-dimensional path with $m$ vertices indexed by
$j=0,\ldots,m-1$. Its Neumann graph Laplacian is the matrix associated with
the quadratic form $\sum_{j=0}^{m-2}|f_{j+1}-f_j|^2$. For an interior
vertex it acts as $2f_j-f_{j-1}-f_{j+1}$; at the two endpoints it acts as
$f_0-f_1$ and $f_{m-1}-f_{m-2}$. For $k=0,\ldots,m-1$, set

\begin{gather*}
\FitMath{v_k(j)=\cos\bigl(\pi k(j+\tfrac12)/m\bigr).}
\end{gather*}

The cosine addition formula gives, in the interior,

\begin{gather*}
\FitMath{2v_k(j)-v_k(j-1)-v_k(j+1)
       =4\sin^2(\pi k/(2m))v_k(j).}
\end{gather*}

At the left endpoint,

\begin{gather*}
\FitMath{\cos(\theta/2)-\cos(3\theta/2)
 =4\sin^2(\theta/2)\cos(\theta/2),} \\[3pt]
\FitMath{\theta=\pi k/m.}
\end{gather*}

The right endpoint follows by reflection. Thus these are eigenvectors with
eigenvalues $4\sin^2(\pi k/(2m))$. The eigenvalues are distinct because sine
is strictly increasing on $[0,\pi/2)$; the nonzero eigenvectors therefore form
an orthogonal basis. The $k=0$ vector is constant.

On the four-dimensional block the graph Laplacian is the sum of the four
one-dimensional matrices acting on separate tensor factors. Tensor products
of the $v_k$ therefore form an eigenbasis. The eigenvalues of the physical
Neumann Laplacian $d_{\eta,N}^*d_{\eta,N}$ are

\begin{gather*}
\FitMath{4\eta^{-2}\sum_{\mu=1}^4\sin^2(\pi k_\mu/(2m)),} \\[3pt]
\FitMath{0\le k_\mu<m.}
\end{gather*}

Its only zero mode is the constant vector, and the smallest positive eigenvalue
is $\lambda_N=4\eta^{-2}\sin^2(\pi/(2m))$. Multiplication of the site and
bond inner products by the common constant $\eta^4$ does not change this
matrix or these eigenvalues. Expanding $f-Pf$ in the eigenbasis gives
$\lambda_N\|f-Pf\|_f^2\le\|d_{\eta,N}f\|_e^2$ on one block and hence on
their disjoint union.

Finally, concavity of sine on $[0,\pi/2]$ gives
$\sin t\ge2t/\pi$ there. At $t=\pi/(2m)$ this yields
$\sin(\pi/(2m))\ge1/m=\eta$, so $\lambda_N\ge4$.
Adding the nonnegative contributions of the other ambient lattice bonds,
including periodic wraps whose endpoints may lie in the same block, can only
increase the Dirichlet form. Substitution of $d_{0,N}=\eta d_{\eta,N}$ gives the
combinatorial version.
\end{proof}

These estimates immediately give a uniform scalar gap on a periodic fine
lattice. If $\mathcal L_\eta=d_\eta^*d_\eta$ and
$\mathcal A=\mathcal L_\eta+Q^*Q$, then

\begin{gather*}
\FitMath{\langle f,\mathcal A f\rangle_f
 =\|d_\eta f\|_e^2+\|Qf\|_c^2
 \ge\|(1-P)f\|_f^2+\|Qf\|_c^2=\|f\|_f^2.}
\end{gather*}

Thus $\mathcal A$ is invertible and $\|\mathcal A^{-1}\|\le1$, independently
of the mesh and number of blocks. This proves a scalar block-constrained gap.
It supplies neither a vector-field coercivity theorem nor a mass gap for a
physical Hamiltonian; those are different assertions.

\subsubsection*{A uniform inverse for a coarse Gram matrix}

The same convention also makes a useful projection formula precise. On the
periodic lattice put $G'=\mathcal A^{-1}$, with $\mathcal A$ as above. There
is a numerical constant $c>0$, independent of $m$ and the number of blocks,
such that

\begin{gather*}
\FitMath{QG'^2Q^*\ge c1.}
\end{gather*}

Here is a finite-dimensional proof, including a choice of the test functions.
For $m\ge2$ define on a block indexed by $j=(j_1,\ldots,j_4)$

\begin{gather*}
\FitMath{\phi_m(j)=16\prod_{\mu=1}^4
             \sin^2\bigl(\pi(j_\mu+\tfrac12)/m\bigr).}
\end{gather*}

The sum of the $m$ roots $e^{2\pi i(j+1/2)/m}$ is zero. Taking real parts
in $\sin^2 t=(1-\cos 2t)/2$ shows that the one-dimensional arithmetic
mean of each sine-square factor is $1/2$. Thus the mean of $\phi_m$ on a
block is one. For a coarse function $u$, define $Tu(x)=u(B)\phi_m(x)$ on
the block $B$. Then $QT=1$ and $\|T\|\le16$.

There is also a bound $\|\mathcal L_\eta Tu\|_f\le C\|u\|_c$ uniform in
$m$. To see it, the second physical difference of
$\sin^2(\pi t)$ at interior sites is bounded by the supremum $2\pi^2$ of
its second derivative. At an end site, its value is
$\sin^2(\pi\eta/2)\le\pi^2\eta^2/4$. A second difference that crosses a
block boundary therefore contributes at most a constant times the absolute
coarse values in the two neighboring blocks: the factor $\eta^{-2}$ is
canceled by this endpoint bound. The other three sine-square factors are
bounded by one. After summing the four coordinate differences, a pointwise
bound is $C\sum_{B'\sim B}|u(B')|$, where the sum includes $B$ and its
face neighbors. Squaring, summing with weight $\eta^4$, and using the bounded
number of neighbors proves the norm estimate. For $m=1$, set $T=1$; the
unit-spacing Laplacian has norm at most $16$, so the same conclusion holds.
Consequently $\|\mathcal A Tu\|_f\le C\|u\|_c$ for all $m$.

Self-adjointness and $G'\mathcal A=1$ now give

\begin{gather*}
\FitMath{\|u\|_c^2=\langle Q^*u,Tu\rangle_f
          =\langle G'Q^*u,\mathcal A Tu\rangle_f
          \le C\|G'Q^*u\|_f\|u\|_c.}
\end{gather*}

This proves the claimed positive lower bound for $QG'^2Q^*$.

The map $V=G'Q^*$ is injective, its Gram matrix is $V^*V=QG'^2Q^*$,
and elementary linear algebra shows that

\begin{gather*}
\FitMath{R=1-G'Q^*(QG'^2Q^*)^{-1}QG'}
\end{gather*}

is the orthogonal projection onto $\ker(QG')$. This kernel equals
$\mathcal L_\eta\ker Q$. Indeed, if $Qv=0$ then
$\mathcal A v=\mathcal L_\eta v$, so $QG'\mathcal L_\eta v=0$.
Conversely, if $QG'h=0$, put $v=G'h$. Then $Qv=0$ and
$h=\mathcal A v=\mathcal L_\eta v$. The expression for $R$ is therefore
well defined with a uniformly invertible coarse Gram matrix. Every adjoint in
this argument uses the physical inner products fixed above.

\subsubsection*{Translation to plain Euclidean coordinates}

There is an equivalent convention with unweighted coordinate inner products,
but the averaging matrix has to change. Let $S_f f=\eta^2 f$ on fine sites,
$S_e a=\eta^2 a$ on fine bonds, and $S_cu=u$ on coarse sites. These maps are
unitaries from the physical spaces just defined to the spaces with plain
coordinate sums. The correctly conjugated averaging matrix is

\begin{gather*}
\FitMath{\widetilde Q=S_cQS_f^{-1},} \\[3pt]
\FitMath{(\widetilde Q v)(B)=\eta^2\sum_{x\in B}v(x),} \\[3pt]
\FitMath{(\widetilde Q^*u)(x)=\eta^2u(B).}
\end{gather*}

Each row has $m^4$ entries of size $\eta^2=m^{-2}$, so its squared Euclidean
length is one. Rows have disjoint supports, proving
$\widetilde Q\widetilde Q^*=1$. The conjugated projection is
$\widetilde P=\widetilde Q^*\widetilde Q=S_fPS_f^{-1}$. Because sites and
bonds have the same physical weight,
$S_e d_\eta S_f^{-1}$ is the same physical difference matrix
$\eta^{-1}(\text{shift}-1)$. Therefore the entire Poincar\'e estimate survives
unitary conjugation with its exact constant.

In contrast, suppose that one retains the arithmetic mean matrix
$\overline Qv=\eta^4\sum_{x\in B}v(x)$ while using plain inner products on
both lattices. Its plain adjoint is
$(\overline Q^*u)(x)=\eta^4u(B)$, and

\begin{gather*}
\FitMath{\overline Q\overline Q^*=\eta^4 1,} \\[3pt]
\FitMath{P=\eta^{-4}\overline Q^*\overline Q.}
\end{gather*}

Thus $\overline Q^*\overline Q$ itself is not the projection. For a nonzero
constant fine function and $m>1$, the expression
$\|f-\overline Q^*\overline Q f\|^2$ is positive whereas $d_\eta f=0$.
This gives an immediate check against accidentally combining the arithmetic
mean with the wrong adjoint. Either the physical weighted convention or the
conjugated unweighted convention is consistent; mixing them is not.

\subsubsection*{Consequences for cochains and Gaussian formulas}

For the centered exact-cochain Gaussian that is standard on the exact-form
subspace in the plain inner product (Section~\ref{ch:lattice-curvature-comparison}),
its covariance is the orthogonal projection onto that subspace.
After changing to physical coordinates, transform the random vector and the
covariance together. For a linear random vector $F$ with covariance $C$ in
plain coordinates and a linear coordinate change $T$, the covariance of $TF$
is $TCT^*$. An orthogonal projection is carried to an orthogonal projection
when the coordinate change is interpreted as the unitary between the two
specified Hilbert spaces; an arbitrary rescaling of its coordinate entries
without the corresponding inner product change has a different covariance.

For a Hilbert-space operator $B:H_1\to H_2$, unitary changes of coordinates
$S_1,S_2$ give

\begin{gather*}
\FitMath{\widetilde B=S_2BS_1^{-1},} \\[3pt]
\FitMath{\widetilde B^*=S_1B^*S_2^{-1}.}
\end{gather*}

In particular products, orthogonal projections, positive inverses and
Moore--Penrose inverses commute with this conjugation. A conditional Gaussian
covariance written as
$C-CB^*(BCB^*)^+BC$ remains the same Hilbert-space operator after every factor
is conjugated. This observation permits a proof to begin in plain finite
cochain coordinates and then use physical block coordinates. The change must
be announced before the first adjoint, constraint, or Poincar\'e estimate that
uses the latter convention. Every discrete derivative must also say whether
it is $d_0$ or $d_\eta$. These declarations prevent a missing power of the mesh
from being hidden inside an adjoint symbol.

% END MODULE chapters/block_normalization_v2.tex


% BEGIN MODULE chapters/block_decay_v2.tex
\subsection{A scalar block resolvent with uniform off-diagonal decay}

This finite-dimensional lemma supplies a scalar part of the propagator input.
All distances are physical distances. The argument uses a relative form bound
and a Neumann series; the exponentially conjugated operator is generally not
self-adjoint.

Let a four-dimensional periodic lattice of spacing $\eta=1/m$, $m\ge1$, be
partitioned into blocks of $m^4$ sites and physical side one. Its torus has an
integer number of blocks in each direction. On site functions use
$\langle u,v\rangle=\eta^4\sum_x\overline{u_x}v_x$. On positively oriented
fine bonds use the same weight. Put

\begin{gather*}
\FitMath{(d_\eta u)(x,x+\eta e_\mu)
       =\eta^{-1}(u_{x+\eta e_\mu}-u_x),} \\[3pt]
\FitMath{\mathcal L_\eta=d_\eta^*d_\eta.}
\end{gather*}

The block mean $Q$ takes the arithmetic mean of the $m^4$ fine values. Its
coarse inner product is the unweighted sum, so $Q^*$ embeds a coarse value as
a constant on its block. Thus $P=Q^*Q$ is an orthogonal projection and its
matrix entries, as an operator acting on coordinate vectors, are
$P_{xy}=\eta^4$ when $x,y$ are in the same block and zero otherwise.

Write $\operatorname{dist}_\eta$ for shortest-path distance on the periodic
fine graph, with length $\eta$ per bond. The distance includes wrapping paths.
The diameter of a single block in this metric is at most $4$.
The block Poincar\'e inequality proves

\begin{gather*}
\FitMath{A_\eta:=\mathcal L_\eta+P\ge1,} \\[3pt]
\FitMath{\|u\|\le\|A_\eta^{1/2}u\|,} \\[3pt]
\FitMath{\|d_\eta u\|\le\|A_\eta^{1/2}u\|.}
\end{gather*}

These statements also hold when $m=1$, since then $P=1$.

\begin{proposition}[A mesh-independent Combes--Thomas bound]
Let $w$ be a real function with $|w_x-w_y|\le\eta$ on every neighboring pair.
For real $t$ let $M_t$ multiply by $e^{tw_x}$ and set
$A_t=M_tA_\eta M_t^{-1}=A_\eta+E_t$. There is a numerical constant $C$,
independent of the mesh, volume and function $w$, such that, for $|t|\le1$,

\begin{gather*}
\FitMath{\|A_\eta^{-1/2}E_t A_\eta^{-1/2}\|
                  \le C(|t|+t^2).}
\end{gather*}

Consequently one may choose a fixed $t_0>0$, independent of those parameters,
so that $\|A_t^{-1}\|\le2$ for $|t|\le t_0$. For arbitrary nonempty sets of
sites $X,Y$ this gives

\begin{gather*}
\FitMath{\|1_X A_\eta^{-1}1_Y\|
 \le2\exp[-t_0\operatorname{dist}_\eta(X,Y)].}
\end{gather*}

\end{proposition}

\begin{proof}
Set $A=A_\eta$. First consider the contribution $E_t^{\mathcal L}$ from the
Laplacian. On a positively oriented bond with endpoints $x,y$, put
$\delta=w_x-w_y$, $a=e^{t\delta}-1$ and $b=e^{-t\delta}-1$. The diagonal
entries of the conjugated edge Laplacian do not change. Its off-diagonal
contribution to $\langle u,E_t^{\mathcal L}v\rangle$ is therefore

\begin{gather*}
\FitMath{-\eta^4\eta^{-2}
             (a\overline{u_x}v_y+b\overline{u_y}v_x).}
\end{gather*}

The exact identity

\begin{gather*}
\FitMath{a\overline{u_x}v_y+b\overline{u_y}v_x
 =a(\overline{u_x}-\overline{u_y})v_y
   +b\overline{u_y}(v_x-v_y)+(a+b)\overline{u_y}v_y}
\end{gather*}

exposes a difference of $u$ or $v$ in the terms linear in $t$. Since
$|\delta|\le\eta\le1$ and $|t|\le1$,

\begin{gather*}
\FitMath{|a|,|b|\le e|t|\eta,} \\[3pt]
\FitMath{|a+b|=2|\cosh(t\delta)-1|\le e t^2\eta^2.}
\end{gather*}

Each vertex occurs as the head of at most four positively oriented bonds.
Cauchy--Schwarz applied to the first term, summed over the bonds, consequently
gives $2e|t|\|d_\eta u\|\|v\|$. The second term is bounded by
$2e|t|\|u\|\|d_\eta v\|$. The last term is bounded by
$4e t^2\|u\|\|v\|$. In particular,

\begin{gather*}
\FitMath{|\langle u,E_t^{\mathcal L}v\rangle|
 \le 2e|t|(\|d_\eta u\|\|v\|+\|u\|\|d_\eta v\|)
                         +4e t^2\|u\|\|v\|.}
\end{gather*}

The two powers of $\eta^{-1}$ in the Laplacian have been canceled by the
Lipschitz difference and the difference of a test vector. Estimating the
off-diagonal matrix entries directly would miss this cancellation.

For the block term, the neighbor condition implies
$|w_x-w_y|\le4$ for sites in one block. Its conjugation error has entries

\begin{gather*}
\FitMath{(E_t^P)_{xy}
 =\eta^4 1_{\{x,y\text{ in one block}\}}
                              (e^{t(w_x-w_y)}-1).}
\end{gather*}

Their absolute row sums and column sums are at most $4e^4|t|$, because a
block has $m^4$ sites and $m^4\eta^4=1$. The Schur test therefore yields
$\|E_t^P\|\le4e^4|t|$. The common site weight $\eta^4$ does not alter this
Schur estimate. Using the two energy bounds preceding the proposition gives

\begin{gather*}
\FitMath{|\langle u,E_t v\rangle|
 \le C(|t|+t^2)\|A^{1/2}u\|\|A^{1/2}v\|,}
\end{gather*}

where, for example, $C=8e+4e^4$ suffices. Substitute $u=A^{-1/2}f$ and
$v=A^{-1/2}g$ and take the supremum over unit vectors $f,g$. This proves
the asserted relative operator bound.

Choose $t_0=\min\{1,(4C)^{-1}\}$. Then
$B_t=A^{-1/2}E_t A^{-1/2}$ has norm at most $1/2$ when $|t|\le t_0$.
The exact factorization and its inverse are

\begin{gather*}
\FitMath{A_t=A^{1/2}(1+B_t)A^{1/2},} \\[3pt]
\FitMath{A_t^{-1}=A^{-1/2}(1+B_t)^{-1}A^{-1/2}.}
\end{gather*}

The Neumann series gives $\|(1+B_t)^{-1}\|\le2$ and hence
$\|A_t^{-1}\|\le2$. This step requires no assertion that $A_t$ is
self-adjoint or ordered below by a positive scalar.

Finally take $w_x=\operatorname{dist}_\eta(x,Y)$ and $t=t_0$. The triangle
inequality makes $w$ admissible, and $w=0$ on $Y$. If
$D=\operatorname{dist}_\eta(X,Y)$, then $w\ge D$ on $X$. Similarity of the
two inverses gives

\begin{gather*}
\FitMath{1_XA^{-1}1_Y=1_X M_t^{-1}A_t^{-1}M_t1_Y.}
\end{gather*}

The rightmost multiplication is the identity on $Y$, while the norm of the
leftmost multiplication is at most $e^{-t_0D}$. Combining these statements
with $\|A_t^{-1}\|\le2$ proves the final estimate.
\end{proof}

The estimate concerns operator norms between sets in the weighted lattice
Hilbert space. It does not assert a mesh-independent bound for a Green
function kernel normalized relative to the physical counting measure
$\eta^4\sum_x$. Converting a coordinate matrix to that kernel introduces its
own power of $\eta$. The scalar lemma likewise does not establish the
non-Abelian vector-field or multiscale background estimates. Those require
their additional gauge constraints and response bounds.

\subsubsection*{Dirichlet restrictions and the distance being estimated}

The same estimate holds for a Dirichlet compression of this particular
operator, with constants independent of the shape of the region. To state it
without ambiguity, let $D$ be any set of sites of the periodic fine lattice.
Let $H_D$ consist of the functions supported in $D$, with the inherited weighted
inner product, and let $J_D:H_D\to H$ be extension by zero. Define

\begin{gather*}
\FitMath{A_D=J_D^*A_\eta J_D.}
\end{gather*}

In particular the block term is the compression of the \emph{original}
projection $P$. On a block intersected by the boundary it is not a new mean
taken over just the surviving sites. Replacing it by that new mean would
define a different operator and is not justified by this corollary.

For $u\in H_D$, the quadratic form is

\begin{gather*}
\FitMath{\langle u,A_Du\rangle
 =\|d_\eta J_Du\|^2+\|QJ_Du\|_c^2\ge\|u\|^2.}
\end{gather*}

It controls the physical gradient of the zero extension, including boundary
bonds. Restrict an admissible ambient weight $w$ to $D$ and conjugate by its
exponential on $H_D$. Because multiplication by this exponential commutes with
the site indicator $1_D$, the new error form is exactly the restriction of
$E_t$ to zero-extended vectors. The estimates in the proof therefore give

\begin{gather*}
\FitMath{\|A_D^{-1/2}(M_tA_DM_t^{-1}-A_D)A_D^{-1/2}\|
 \le C(|t|+t^2).}
\end{gather*}

The same Neumann argument implies, for $X,Y\subset D$,

\begin{gather*}
\FitMath{\|1_XA_D^{-1}1_Y\|
 \le2e^{-t_0\operatorname{dist}_\eta(X,Y)}.}
\end{gather*}

Here the distance is computed in the ambient periodic lattice. It may use a
path outside $D$. No estimate in the possibly much larger intrinsic distance
of a region with holes has been asserted.

For comparison with continuum notation, the shortest-path distance is the
periodic physical Manhattan distance, and hence is no smaller than the
periodic Euclidean distance. The displayed bound therefore implies the same
bound with Euclidean distance in its exponent. One can alternatively take
$w$ to be Euclidean distance to $Y$ directly, since its change across a fine
bond is at most $\eta$. Both formulations concern physical distances; using
the number of bonds without the factor $\eta$ would incorrectly strengthen
the assertion as the mesh is refined.

The calculation explains the uniform mechanism: the scalar energy controls
one physical difference, the weight varies by only $O(t\eta)$ across an
edge, and its genuinely quadratic error is $O(t^2\eta^2)$. The block projection
has range at most one physical cell and row mass one. Consequently the
conjugation cost depends on the dimension and physical block size, with no
factor proportional to the number of fine sites in a block or the torus.

% END MODULE chapters/block_decay_v2.tex


% BEGIN MODULE chapters/polymer_theorem_v2.tex
\section{A polymer expansion with a complete convergence proof}
\label{ch:polymer-exact}

The small-field integration must produce activities satisfying a summability
criterion. Once that estimate is available, the passage from activities to a
local logarithm is a combinatorial theorem. We prove that theorem here, so its
convergence mechanism is independent of the particular gauge-field
coordinates. The estimate on the activities remains a separate analytic task.

\subsection{Finite polymer systems}

A finite polymer system consists of a finite set $\mathcal P$ and a symmetric
incompatibility relation $\not\sim$ on it. Every polymer is incompatible with
itself. Give each $X\in\mathcal P$ a complex activity $z_X$. A compatible
family contains no two incompatible polymers. Its partition function is

\begin{gather*}
\FitMath{Z(z)=\sum_{\substack{\mathcal F\subset\mathcal P\\
                      \mathcal F\,\mathrm{compatible}}}
                              \prod_{X\in\mathcal F}z_X.}
\end{gather*}

The empty family contributes one. There are no questions of convergence in
this polynomial. The problem is to construct its logarithm and bound it in
terms of local activities without a loss exponential in $|\mathcal P|$.

For a list $\mathbf X=(X_1,\ldots,X_n)$, let $H(\mathbf X)$ be the simple
graph on the labeled vertices $1,\ldots,n$ whose edges are the pairs with
$X_i\not\sim X_j$. Repeated polymers are permitted in lists; their distinct
vertices are adjacent by self-incompatibility. Define

\begin{gather*}
\FitMath{\phi(\mathbf X)=
  \sum_{\substack{F\subset E(H(\mathbf X))\\([n],F)\,\mathrm{connected}}}
                       (-1)^{|F|},} \\[3pt]
\FitMath{n\ge2,} \\[3pt]
\FitMath{\phi(X_1)=1.}
\end{gather*}

If $H(\mathbf X)$ is disconnected the sum is zero.

\begin{lemma}[The connected graph identity]
As a formal power series at $z=0$,

\begin{gather*}
\FitMath{\log Z(z)=\sum_{n\ge1}\frac1{n!}
                  \sum_{X_1,\ldots,X_n\in\mathcal P}
                      \phi(\mathbf X)\prod_{i=1}^nz_{X_i}.}
\end{gather*}

\end{lemma}
\begin{proof}
The ordered-list expression for $Z$ is

\begin{gather*}
\FitMath{Z(z)=\sum_{n\ge0}\frac1{n!}\sum_{\mathbf X\in\mathcal P^n}
  \prod_i z_{X_i}\prod_{i<j}
                  (1-\mathbf1_{\{X_i\not\sim X_j\}}).}
\end{gather*}

The product is one on pairwise compatible lists and zero otherwise. Such
lists contain no repetition, and their $n!$ orderings give exactly the
original polynomial. Expand the product as a sum over edge sets. Each graph
decomposes uniquely into its connected components. The weight factors over
the components, including the activity at each vertex. Expanding the
exponential of the proposed connected-graph series enumerates a partition
of the labeled vertices into these components. The factor for ordering the
components is canceled by the exponential's factorial, and the factorials
within the components cancel their vertex orderings. This reconstructs the
ordered-list expression coefficient by coefficient. Its constant term is
one, so its formal logarithm is the connected series.
\end{proof}

\subsection{The tree bound, including the cancellation}

\begin{lemma}[A hard-core tree inequality]
For a finite simple graph $H$ on $n\ge1$ labeled vertices,

\begin{gather*}
\FitMath{\left|\sum_{\substack{F\subset E(H)\\([n],F)\,\mathrm{connected}}}
                  (-1)^{|F|}\right|
       \le \#\{\text{spanning trees of }H\}.}
\end{gather*}

For $n=1$ the empty tree is counted once.
\end{lemma}
\begin{proof}
Fix a total order of the edges of $H$. Call an edge $e$ eligible for an edge
set $F$ if the endpoints of $e$ can be joined by a path in $F$ consisting
only of edges smaller than $e$. This definition does not require $e\in F$.
If at least one edge is eligible, choose the smallest eligible edge $e$ and
toggle its membership in $F$.

This is an involution. Eligibility of $e$ depends only on edges smaller than
$e$, which were not changed; the same holds for every smaller edge, so the
smallest eligible edge stays the same. Toggling $e$ preserves the connected
components, because its endpoints already have a path using smaller edges.
It reverses the parity of $|F|$. Thus all connected edge sets with an
eligible edge cancel in pairs in the alternating sum.

A surviving edge set has no cycle: if it contained one, the largest edge of
that cycle would be eligible. Consequently every surviving connected edge
set is a spanning tree. Each contributes the same sign $(-1)^{n-1}$, and
there can be no more survivors than spanning trees. This proves the bound.
The same argument also proves that the alternating sum has sign
$(-1)^{n-1}$ when it is nonzero.
\end{proof}

In particular,

\begin{gather*}
\FitMath{|\phi(X_1,\ldots,X_n)|\le
 \sum_{T\,\mathrm{tree\,on}\,[n]}
                       \prod_{\{i,j\}\in E(T)}
                             \mathbf1_{\{X_i\not\sim X_j\}}.}
\end{gather*}

The cancellation in this inequality is essential. Taking absolute values
before summing over connected graphs would count all connected subgraphs,
which need not admit the rooted-tree bound used next.

\subsection{The rooted-tree estimate}

Let $a_X\ge0$ be real numbers and suppose
\begin{equation}
 \sum_{Y:Y\not\sim X}|z_Y|e^{a_Y}\le a_X
                      \quad\text{for every }X\in\mathcal P.
 \label{eq:kp-explicit}
\end{equation}
This is the polymer criterion used here \cite{KP86}.
For a fixed root polymer $X$, let $\mathcal T_X$ be the total nonnegative
weight of all rooted labeled trees, with the root fixed to have polymer
$X$, every nonroot vertex of polymer $Y$ carrying weight $|z_Y|$, and every
edge requiring incompatibility of its endpoint polymers. A tree on $m$
vertices has the factor $1/(m-1)!$ for the nonroot vertex labels.

One can define this sum without assuming its convergence. Start with
$T_X^{(0)}=1$ and set

\begin{gather*}
\FitMath{T_X^{(h+1)}=\exp\left(\sum_{Y:Y\not\sim X}|z_Y|T_Y^{(h)}\right).}
\end{gather*}

The quantity $T_X^{(h)}$ sums trees of height at most $h$. To see this,
delete the root. Each child becomes the root of a tree of height at most
$h-1$, and an unordered collection of $r$ child trees contributes the
factor $1/r!$ in the exponential. Expanding those exponentials and assigning
the nonroot labels gives exactly $1/(m-1)!$: the multinomial coefficient
for assigning labels to the child subtrees cancels their factorials and
the ordering of the $r$ children. This also verifies the formula for every
fixed degree before any infinite sum is taken.

The iterates increase with $h$. Moreover $T_X^{(h)}\le e^{a_X}$ for every
$h$: this is true at $h=0$, and the induction step is precisely
\eqref{eq:kp-explicit}. Monotone convergence proves
\begin{equation}
 \mathcal T_X=\lim_{h\to\infty}T_X^{(h)}\le e^{a_X}.
 \label{eq:rooted-tree-bound}
\end{equation}

\begin{theorem}[Absolute convergence and a nonzero partition function]
Under \eqref{eq:kp-explicit} the connected series converges absolutely and
defines a number $L(z)$ with $e^{L(z)}=Z(z)$. In particular $Z(z)\ne0$.
For each root polymer $X$,

\begin{gather*}
\FitMath{\sum_{n\ge1}\frac1{(n-1)!}
 \sum_{X_2,\ldots,X_n}
 |\phi(X,X_2,\ldots,X_n)|\prod_{i=2}^n|z_{X_i}|
                       \le e^{a_X}.}
\end{gather*}

Also $|L(z)|\le\sum_X|z_X|e^{a_X}$.
If the criterion holds throughout a complex activity domain with a uniform
majorant, the series converges locally normally there and gives the branch
of $\log Z$ that equals zero at zero.
\end{theorem}
\begin{proof}
The tree inequality bounds the root series by $\mathcal T_X$, proving its
estimate by \eqref{eq:rooted-tree-bound}. Rooting each nonempty connected
list at its first vertex bounds the absolute unrooted series by
$\sum_X|z_X|e^{a_X}$: replacing $1/n!$ by $1/(n-1)!$ only enlarges it.

For $|t|<1$ the same majorant applies to the activities $tz_X$, and the
connected series is analytic in $t$ there. Near zero its exponential is
$Z(tz)$ by the formal identity and absolute convergence. The analytic
identity theorem gives $e^{L(tz)}=Z(tz)$ on the entire unit disc. The
absolute majorant is summable also at $|t|=1$, so the connected series is
continuous as $t$ tends radially to one. Taking that limit proves
$Z(z)=e^{L(z)}\ne0$.

On a domain with a uniform summable majorant, the Weierstrass criterion
gives normal convergence on smaller domains and hence holomorphy. The
branch statement follows by agreement near zero and continuation within
the connected component of that domain containing zero. This statement
does not assert that an arbitrary disconnected activity domain has a
single preferred logarithm.
\end{proof}

For a countable polymer set the same statements hold in a finite region
whenever $\sum_X|z_X|e^{a_X}<\infty$ there, by exhaustion and the displayed
absolute majorant. In an infinite lattice the total unrooted sum is
generally extensive and infinite. The rooted bound remains meaningful if
the incompatibility sums satisfy the criterion, and it is the form used
to define local connected expansions. No finite infinite-volume partition
function is asserted.

\subsection{Application to fluctuation activities}

If a polymer activity depends only on field variables in an enlargement
of that polymer, each connected term depends only on the union of those
enlargements. This is a direct consequence of its product formula. Extra
exponential weights can be incorporated in the activity majorant before
the criterion is checked. The root estimate then retains those weights.
Holomorphic dependence on external fields or sources passes to the
connected expansion on every domain where the same absolute majorant is
valid.

For gauge-field fluctuation activities, apply the theorem after
verifying \eqref{eq:kp-explicit} with the covariance, localization-domain
count, source radii, boundary costs and required derivatives included
in its majorant. Large-field histories use their separate weighted
label estimate. These are the analytic conditions that we must verify in
the multiscale application of the combinatorial convergence theorem.

% END MODULE chapters/polymer_theorem_v2.tex


% BEGIN MODULE chapters/free_curvature_v2.tex
\section{The free curvature model: an exact calculation}
\label{ch:free-curvature-exact}

This section computes the noncoincident correlations of the action density
of a free curvature in four Euclidean dimensions. The calculation fixes a
normalization and proves the cycle identities used in the short-distance
comparison. It does not identify an interacting lattice measure with this
model. That comparison, including its error, remains a separate estimate.

\subsection{A Gaussian field defined without a pointwise potential}

Fix a positive number $g$ and an integer $d_G\ge1$; for the adjoint color
space of $\mathrm{SU}(2)$, $d_G=3$. Coordinates in $\mathbb R^4$ are denoted
by $x_1,\ldots,x_4$. Our Fourier transform is

\begin{gather*}
\FitMath{\widehat f(k)=\int_{\mathbb R^4}e^{-ik\cdot x}f(x)\,dx,} \\[3pt]
\FitMath{f(x)=(2\pi)^{-4}\int e^{ik\cdot x}\widehat f(k)\,dk.}
\end{gather*}

Let $G$ be the tempered distribution with Fourier transform $|k|^{-2}$.
In four dimensions this multiplier is locally integrable at zero. Its
inverse transform, away from zero, is

\begin{gather*}
\FitMath{G(x)=\frac1{4\pi^2|x|^2}.}
\end{gather*}

One direct verification uses
$|k|^{-2}=\int_0^\infty e^{-t|k|^2}\,dt$ and the product Gaussian integral.
The inverse transform of $e^{-t|k|^2}$ is
$(4\pi t)^{-2}e^{-|x|^2/(4t)}$. Integrating in $t$, with
$u=|x|^2/(4t)$, gives $1/(4\pi^2|x|^2)$. The same identity holds as
tempered distributions by Fubini: the integral
$\int |\varphi(x)|\,|x|^{-2}dx$ is finite for every Schwartz function
$\varphi$, controlling the absolute integral of the heat-kernel formula.

For real Schwartz functions define

\begin{gather*}
\FitMath{(f,h)_G=(2\pi)^{-4}\int_{\mathbb R^4}
              \frac{\overline{\widehat f(k)}\widehat h(k)}{|k|^2}\,dk.}
\end{gather*}

This is a positive real bilinear form. The integral is finite near zero
because $\int_0^1r^3r^{-2}dr<\infty$, and is finite at infinity by Schwartz
decay. Quotient by its null space and complete to a real Hilbert space
$\mathcal H_G$. It is separable: the map $f\mapsto|k|^{-1}\widehat f(k)$
embeds it into a separable real form of $L^2(\mathbb R^4)$.

Choose an orthonormal basis $(e_j)$ and independent real standard normal
variables $(N_j)$ on their countable product probability space. For
$f\in\mathcal H_G$, the series

\begin{gather*}
\FitMath{W(f)=\sum_j(f,e_j)_G N_j}
\end{gather*}

converges in $L^2$ and defines a centered jointly Gaussian family with
$\mathbb E[W(f)W(h)]=(f,h)_G$. Finite partial sums prove the covariance
formula and joint Gaussianity; convergence of their characteristic
functions proves it for the limit. Independent copies give families
$A_\mu^a(f)=gW_\mu^a(f)$ for $1\le\mu\le4$, $1\le a\le d_G$.
Their covariance is

\begin{gather*}
\FitMath{\mathbb E[A_\mu^a(f)A_\nu^b(h)]
       =g^2\delta_{ab}\delta_{\mu\nu}(f,h)_G.}
\end{gather*}

No value $A_\mu^a(x)$ at a point is used in this definition.

Define the curvature as the Gaussian distribution

\begin{gather*}
\FitMath{F_{\mu\nu}^a(f)=-A_\nu^a(\partial_\mu f)
                         +A_\mu^a(\partial_\nu f),} \\[3pt]
\FitMath{\mu<\nu.}
\end{gather*}

There are six independent coordinate labels $P=(\mu,\nu)$, although the
six curvature components need not be independent random variables.
Differentiating the covariance as a distribution gives

\begin{gather*}
\FitMath{\mathbb E[F_{\mu\nu}^a(x)F_{\rho\sigma}^b(y)]
       =g^2\delta_{ab}K_{\mu\nu,\rho\sigma}(x-y),}
\end{gather*}

where the notation on the left denotes a distributional kernel and

\begin{gather*}
\FitMath{K_{\mu\nu,\rho\sigma}(x)=
 -\delta_{\nu\sigma}\partial_\mu\partial_\rho G(x)
 +\delta_{\nu\rho}\partial_\mu\partial_\sigma G(x)
 +\delta_{\mu\sigma}\partial_\nu\partial_\rho G(x)
 -\delta_{\mu\rho}\partial_\nu\partial_\sigma G(x).}
\end{gather*}

The signs follow from differentiating one argument at $x$ and the other
at $y$: $\partial_{y_\rho}G(x-y)=-\partial_{x_\rho}G(x-y)$.
At nonzero separation this is an ordinary smooth matrix-valued function.

\subsection{The six-dimensional curvature matrix}

Give $\Lambda^2\mathbb R^4$ the Euclidean inner product for which
$(e_\mu\wedge e_\nu)_{\mu<\nu}$ is an orthonormal basis. Put $r=|x|$,
$u=x/r$, and $R_u=I-2u\otimes u$. This $R_u$ is the orthogonal reflection
normal to $u$. Write $\Lambda^2R_u$ for its induced operator on two-forms.

\begin{proposition}[The curvature coefficient]
For $x\ne0$,

\begin{gather*}
\FitMath{K(x)=\frac1{\pi^2r^4}\Lambda^2R_u,} \\[3pt]
\FitMath{K(x)^\mathsf T K(x)=\frac1{\pi^4r^8}I_6,} \\[3pt]
\FitMath{\operatorname{tr}(K(x)^\mathsf T K(x))
                  =\frac6{\pi^4r^8}.}
\end{gather*}

\end{proposition}
\begin{proof}
Twice differentiating $G$ gives

\begin{gather*}
\FitMath{\partial_\mu\partial_\rho G(x)
   =\frac1{2\pi^2r^4}(4u_\mu u_\rho-\delta_{\mu\rho}).}
\end{gather*}

At $u=e_4$ this Hessian is diagonal with entries
$(-1,-1,-1,3)/(2\pi^2r^4)$. Inserting it in the component formula for $K$
shows that $K$ is diagonal: it has eigenvalue $1/(\pi^2r^4)$ on
$e_1\wedge e_2,e_1\wedge e_3,e_2\wedge e_3$, and eigenvalue
$-1/(\pi^2r^4)$ on $e_1\wedge e_4,e_2\wedge e_4,e_3\wedge e_4$.
This is exactly the action of $\Lambda^2R_{e_4}$ times the stated scalar.
Both the differentiated radial kernel and the claimed expression transform
by conjugation under orthogonal changes of coordinates, proving the first
identity for every $u$. Since $\Lambda^2R_u$ is an orthogonal involution
on a six-dimensional space, its squared norm matrix is $I_6$, giving the
other two identities.
\end{proof}

Fix the standard orientation and let $J=*$ be the Hodge operator on
two-forms. It satisfies $J^2=I$. For an orthogonal map $R$,

\begin{gather*}
\FitMath{J(\Lambda^2R)=(\det R)(\Lambda^2R)J.}
\end{gather*}

To check the formula, use the defining identity
$\alpha\wedge *\beta=\langle\alpha,\beta\rangle\,dx_1\wedge\cdots\wedge dx_4$
and the fact that the top exterior power of $R$ multiplies the oriented
volume by $\det R$. Since $\det R_u=-1$, $JK(x)=-K(x)J$.
Every product of an odd number of matrices $K(x_i)$ therefore anticommutes
with $J$ and has trace zero: conjugation by $J$ preserves its trace while
changing its sign.

\subsection{Regularized densities and connected cycle contractions}

Take a real smooth even function $\rho$ supported in the unit ball with
$\int\rho=1$ and put $\rho_\varepsilon(x)=\varepsilon^{-4}\rho(x/\varepsilon)$.
Smear every curvature component with its translate to define
$F_{\varepsilon,P}^a(x)$. This is a Gaussian variable with a smooth
covariance kernel. Its centered quadratic density is

\begin{gather*}
\FitMath{q_\varepsilon(x)=\frac12\sum_{a=1}^{d_G}\sum_{P=1}^6
  \big(F_{\varepsilon,P}^a(x)^2
           -\mathbb E[F_{\varepsilon,P}^a(x)^2]\big).}
\end{gather*}

On any compact set its covariance is continuous and bounded. The integral
$q_\varepsilon(f)=\int f(x)q_\varepsilon(x)dx$ for compactly supported
smooth $f$ is defined by $L^2$ limits of Riemann sums. The continuous
covariance makes those sums Cauchy. The Gaussian moment formula gives
convergence of their moments of each fixed order as well: all the
integrands are finite sums of products of bounded continuous covariances
on a fixed compact set.

For clarity, first fix distinct points $x_1,\ldots,x_n$ at a fixed
$\varepsilon$. For one color let $C_{ij}$ be the $6\times6$ covariance block
of $F_\varepsilon(x_i)$ and $F_\varepsilon(x_j)$, including the factor $g^2$.
Put $Z=\operatorname{diag}(z_1I_6,\ldots,z_nI_6)$ and let $C$ be the full
block covariance matrix. The Gaussian integral yields, near $z=0$,

\begin{gather*}
\FitMath{\log\mathbb E\exp\big(\textstyle\sum_i z_iq_\varepsilon(x_i)\big)
   =-\frac{d_G}{2}\log\det(I-CZ)
                         -\frac{d_G}{2}\sum_i z_i\operatorname{tr}C_{ii}.}
\end{gather*}

Here $C$ is positive semidefinite; write the Gaussian vector as
$C^{1/2}N$ and diagonalize the real symmetric quadratic form for small
real $z$. The one-dimensional Gaussian integral proves the determinant
identity there. Holomorphic continuation proves it in a complex
neighborhood of zero. Sylvester's determinant identity identifies
$\det(I-C^{1/2}ZC^{1/2})$ with $\det(I-CZ)$, including singular $C$.

The last term only subtracts the first cumulants. Expand the logarithm of
the determinant by its convergent matrix power series. For $n\ge2$ the
joint cumulant is consequently

\begin{gather*}
\FitMath{\kappa_n(q_\varepsilon(x_1),\ldots,q_\varepsilon(x_n))
  =\frac{d_G}{2n}\sum_{\sigma\in S_n}
  \operatorname{tr}\big(C_{\sigma(1)\sigma(2)}\cdots
                      C_{\sigma(n)\sigma(1)}\big).}
\end{gather*}

Indeed only the term of degree $n$ in
$\frac{d_G}{2}\sum_{m\ge1}\operatorname{tr}((CZ)^m)/m$ contributes to the
mixed derivative. Each permutation records the order of the $n$ distinct
vertices of the cycle. The factor $1/n$ removes its cyclic repetitions;
the factor $1/2$ remains even when $n=2$.

\subsection{The noncoincident limit and its consequences}

Let the test functions $f_1,\ldots,f_n$ have pairwise disjoint compact
supports, and let $d>0$ be their minimum separation. The regularized
off-diagonal kernel is
$C_{ij}=g^2(K*\rho_\varepsilon*\rho_\varepsilon)(x_i-x_j)$.
For $\varepsilon<d/4$ the convolution samples $K$ only away from zero.
It converges uniformly, with every fixed derivative, on the compact set
of relevant differences. The cycle formula and ordinary dominated
convergence therefore give a finite limit of the smeared connected
correlations. Its kernel is

\begin{gather*}
\FitMath{C_n(x_1,\ldots,x_n)=\frac{d_Gg^{2n}}{2n}\sum_{\sigma\in S_n}
 \operatorname{tr}\big(K(x_{\sigma(1)}-x_{\sigma(2)})\cdots
 K(x_{\sigma(n)}-x_{\sigma(1)})\big),} \\[3pt]
\FitMath{x_i\ne x_j.}
\end{gather*}

This defines separated correlations. It does not assert convergence of
$q_\varepsilon(f)$ in $L^2$ when the same test is used twice; coincidence
singularities require a different distributional treatment.

\begin{corollary}[Two points, odd cycles, and a positive four-point kernel]
The limiting connected two-point kernel is

\begin{gather*}
\FitMath{C_2(x,y)=\frac{3d_Gg^4}{\pi^4|x-y|^8}.}
\end{gather*}

Every odd connected kernel $C_n$ with $n\ge3$ vanishes at noncoincident
arguments. If four distinct points are collinear, then

\begin{gather*}
\FitMath{C_4(x_1,x_2,x_3,x_4)=\frac{6d_Gg^8}{\pi^8}
 \sum_{\text{three unoriented cycles}}
       \prod_{(i,j)\text{ in the cycle}}|x_i-x_j|^{-4}>0.}
\end{gather*}

For four sufficiently small disjoint neighborhoods of such points the
four-point kernel remains strictly positive throughout their product.
\end{corollary}
\begin{proof}
For two points, the permutation formula is
$d_Gg^4\operatorname{tr}(K(x-y)K(y-x))/2$. The kernel is even and symmetric,
so the curvature coefficient proposition proves the stated value. For odd
$n$, every trace in the cycle sum vanishes by the Hodge anticommutation
identity. For collinear points all difference directions give the same
reflection $R_u=R_{-u}$. The product of four curvature matrices is the
positive scalar product of their radial factors times $(\Lambda^2R_u)^4=I$.
Every trace is therefore six times that scalar. The 24 permutations divide
into three unoriented cycles, each represented eight times, canceling
the factor $2n=8$. Positivity on small neighborhoods follows from
continuity and the strictly positive value at the chosen configuration.
\end{proof}

For nonnegative nonzero separated tests this gives

\begin{gather*}
\FitMath{C_2(f,h)=\frac{3d_Gg^4}{\pi^4}
                 \iint\frac{f(x)h(y)}{|x-y|^8}\,dx\,dy>0.}
\end{gather*}

Thus the coefficient $3(N^2-1)/\pi^4$ used for $\mathrm{SU}(N)$ is the
exact free coefficient in this normalization. For $\mathrm{SU}(2)$ it
is $9/\pi^4$. Recovering this coefficient from the interacting lattice
theory still requires its ultraviolet comparison and uniform remainder
estimate. The four-point calculation also shows why non-Gaussianity of
the scalar density alone cannot certify an interacting underlying gauge
theory: this density already has a positive connected four-point kernel
in an explicitly Gaussian model.

% END MODULE chapters/free_curvature_v2.tex


% BEGIN MODULE chapters/lattice_curvature_comparison_v2.tex
\section{A uniform comparison for the free lattice curvature}
\label{ch:lattice-curvature-comparison}

This section proves a free-field input of the construction. The centered Gaussian exact lattice two-form that is standard on the exact-form subspace has covariance equal to the orthogonal projection onto that subspace. After multiplication by the physical factor $a^{-4}$, its entries converge away from coincidence to the continuum curvature kernel. We prove an error estimate uniform in the lattice spacing and the torus side, and compute the resulting separated quadratic-density coefficient. No interacting measure is compared with this Gaussian model here.

\subsection{Cochains, Fourier conventions, and the statement}

Let $a>0$, $\ell>0$, and $N=\ell/a\in\mathbb N$. The lattice is
$\Lambda_{a,\ell}=a\mathbb Z^4/\ell\mathbb Z^4$. A one-form has coordinates $A_\mu(x)$, one per positive bond based at $x$. A two-form has coordinates $F_{\mu\nu}(x)$ for $\mu<\nu$, based at the lower vertex of its positively oriented plaquette. Both cochain spaces use the plain coordinate inner products. Define

\begin{gather*}
\FitMath{(d_1A)_{\mu\nu}(x)
 =a^{-1}\{A_\nu(x+ae_\mu)-A_\nu(x)
          -A_\mu(x+ae_\nu)+A_\mu(x)\}.}
\end{gather*}

The factor $a^{-1}$ does not affect the range. Let $P_{\mathrm{ex}}$ be the orthogonal projection onto $\operatorname{Ran}d_1$. It is the covariance of a centered Gaussian vector which is standard on that subspace. Potential gauge directions need not be inverted to define this curvature covariance. The constant two-forms are excluded because they are not in $\operatorname{Ran}d_1$.

Write $P=(\mu,\nu)$, $\mu<\nu$, and put $c_P=(e_\mu+e_\nu)/2$. The centre of the plaquette $(x,P)$ is $x+ac_P$. The physical lattice kernel is

\begin{gather*}
\FitMath{\mathsf K_{a,\ell}^{PQ}(x,y)
    =a^{-4}P_{\mathrm{ex}}((x,P),(y,Q)).}
\end{gather*}

For a displacement on the torus, its length means distance to $\ell\mathbb Z^4$. Representatives are chosen in $[-\ell/2,\ell/2]^4$ when needed.

The continuum Fourier transform on $\mathbb R^4$ is

\begin{gather*}
\FitMath{\widehat f(k)=\int e^{-ik\cdot x}f(x)\,dx,} \\[3pt]
\FitMath{f(x)=(2\pi)^{-4}\int e^{ik\cdot x}\widehat f(k)\,dk.}
\end{gather*}

Let $G(x)=1/(4\pi^2|x|^2)$ and define the continuum curvature distribution by
\begin{align*}
 \mathsf K_{\mu\nu,\rho\sigma}(x)
 ={}&-\delta_{\nu\sigma}\partial_\mu\partial_\rho G(x)
      +\delta_{\nu\rho}\partial_\mu\partial_\sigma G(x)\\
    &+\delta_{\mu\sigma}\partial_\nu\partial_\rho G(x)
      -\delta_{\mu\rho}\partial_\nu\partial_\sigma G(x).
\end{align*}
At $x\ne0$ this is an ordinary smooth matrix. Its contact terms at zero are retained in its Fourier definition and are not squared in the statements below.

\begin{theorem}[Free lattice curvature comparison]
There is a numerical constant $C$ such that, for every $a,\ell$ as above and every pair of plaquettes with centre displacement

\begin{gather*}
\FitMath{u=x-y+a(c_P-c_Q),} \\[3pt]
\FitMath{10a\le r=\operatorname{dist}(u,\ell\mathbb Z^4)\le\ell/10,}
\end{gather*}

one has
\begin{equation}
 \left|\mathsf K_{a,\ell}^{PQ}(x,y)-\mathsf K_{PQ}(u)\right|
 \le Cr^{-4}\left\{(a/r)^{1/2}+(r/\ell)^4\right\}.
 \label{eq:lcc-main}
\end{equation}
The continuum kernel on the right is evaluated at the short representative of $u$. The constants do not depend on the orientation labels, the lattice spacing, the number of sites, or the torus side.
\end{theorem}

The exponent $1/2$ is sufficient for convergence and is not claimed optimal. Keeping the barycentre phases in a sharper symbol expansion can improve the lattice power. The proof here deliberately keeps a periodic raw symbol and moves to barycentres only after the Fourier estimate.

\subsection{The exact projection symbol}

The torus frequencies are $k\in\delta\mathbb Z^4$, where $\delta=2\pi/\ell$. Choose one Brillouin representative of each frequency modulo $(2\pi/a)\mathbb Z^4$ and denote this set by $\mathcal Q_a$. Its coordinates lie between $-\pi/a$ and $\pi/a$, with one endpoint omitted when necessary. Define

\begin{gather*}
\FitMath{t_{a,\mu}(k)=a^{-1}(e^{iak_\mu}-1),} \\[3pt]
\FitMath{\omega_a(k)=\sum_\mu|t_{a,\mu}(k)|^2
        =4a^{-2}\sum_\mu\sin^2(ak_\mu/2).}
\end{gather*}

Let $B_a(k):\mathbb C^4\to\Lambda^2\mathbb C^4$ be

\begin{gather*}
\FitMath{(B_a(k)v)_{\mu\nu}=t_{a,\mu}(k)v_\nu-t_{a,\nu}(k)v_\mu.}
\end{gather*}

This is exactly the Fourier symbol of $d_1$. Direct multiplication gives

\begin{gather*}
\FitMath{B_a^*B_a=\omega_a I-t_a t_a^*,} \\[3pt]
\FitMath{B_at_a=0.}
\end{gather*}

Consequently $(B_aB_a^*)^2=\omega_a B_aB_a^*$. For $k\ne0$ in the Brillouin zone,

\begin{gather*}
\FitMath{P_a(k)=B_a(k)B_a(k)^*/\omega_a(k)}
\end{gather*}

is the orthogonal projection onto $\operatorname{Ran}B_a(k)$. Set $P_a(0)=0$. Thus $\|P_a(k)\|\le1$ and its rank at nonzero frequency is three. At frequency zero the exact-form subspace is zero.

The unitary finite Fourier transform and its inverse are

\begin{gather*}
\FitMath{\widehat F(k)=N^{-2}\sum_{x\in\Lambda_{a,\ell}}e^{-ik\cdot x}F(x),} \\[3pt]
\FitMath{F(x)=N^{-2}\sum_{k\in\mathcal Q_a}e^{ik\cdot x}\widehat F(k).}
\end{gather*}

Therefore
\begin{equation}
 a^{-4}P_{\mathrm{ex}}((x,P),(y,Q))
  =\ell^{-4}\sum_{k\in\mathcal Q_a}
          e^{ik\cdot(x-y)}[P_a(k)]_{PQ}.
 \label{eq:lcc-fourier-lattice}
\end{equation}
Indeed the projection entry before multiplication by $a^{-4}$ has prefactor $N^{-4}$, and $a^{-4}N^{-4}=\ell^{-4}$.

For the continuum symbol replace $t_a$ by $t_0(k)=ik$. Put

\begin{gather*}
\FitMath{P_0(k)=B_0(k)B_0(k)^*/|k|^2} \\[3pt]
\FitMath{(k\ne0),} \\[3pt]
\FitMath{P_0(0)=0.}
\end{gather*}

It is again a norm-one rank-three projection. With the specified Fourier convention its inverse transform is precisely the distribution $\mathsf K$ defined above. On the torus define the distribution
\begin{equation}
 \mathsf K_\ell(z)=\ell^{-4}
        \sum_{k\in\delta\mathbb Z^4}e^{ik\cdot z}P_0(k).
 \label{eq:lcc-fourier-torus}
\end{equation}
The bounded Fourier coefficients define a periodic distribution. Below we show that it is smooth away from the origin and identify its finite-volume correction.

\subsection{Uniform derivative and counting estimates}

For each multi-index $\alpha$ there is a constant $C_\alpha$ such that, in the Brillouin zone away from zero,
\begin{equation}
 \|\partial^\alpha P_a(k)\|\le C_\alpha |k|^{-|\alpha|}.
 \label{eq:lcc-symbol-derivatives}
\end{equation}
The same estimate holds for $P_0$. To prove it, use
$2|\sin(t/2)|\ge2|t|/\pi$ for $|t|\le\pi$. Hence
$c|k|^2\le\omega_a(k)\le|k|^2$ throughout the Brillouin zone. Differentiating the exponential gives
$|\partial^\alpha t_a(k)|\le C_\alpha |k|^{1-|\alpha|}$; for positive derivative order this uses $a|k|\le2\pi$. The product rule gives the analogous bounds of order $2-|\alpha|$ for $\omega_a$ and $B_aB_a^*$. Induction in the identity $\omega_a\omega_a^{-1}=1$ bounds derivatives of the reciprocal. The quotient rule proves \eqref{eq:lcc-symbol-derivatives}. All functions involved are periodic; near a Brillouin face the same bounds hold in periodic coordinates, where $|k|$ is comparable to $a^{-1}$.

For $a|k|\le c_0$, with a fixed small $c_0>0$, Taylor's formula gives
$t_a(k)=ik+O(a|k|^2)$. The preceding denominator bounds imply
\begin{equation}
 \|P_a(k)-P_0(k)\|\le Ca|k|.
 \label{eq:lcc-low-symbol}
\end{equation}
For example, the numerator difference is $O(a|k|^3)$ and the denominator difference is $O(a|k|^3)$, while both denominators are comparable to $|k|^2$. Both symbols have been assigned zero at $k=0$, so that frequency contributes no error.

We will use two elementary grid bounds. If $R\ge C\delta$ and $M>4$, then
\begin{align}
 \ell^{-4}\#\{k\in\delta\mathbb Z^4:|k|\le2R\}
       &\le C R^4,\label{eq:lcc-count-low}\\
 \ell^{-4}\sum_{k\in\delta\mathbb Z^4}(R+|k|)^{-M}
       &\le C_M R^{4-M}.\label{eq:lcc-count-high}
\end{align}
For a proof, place a cube of side $\delta$ about each grid point. The cubes are disjoint up to boundaries. If $R\ge C\delta$, the values of $R+|k|$ and $R+|q|$ on such a cube differ by bounded factors. The second sum is bounded by a constant times $(2\pi)^{-4}\int(R+|q|)^{-M}dq$, which is $C_MR^{4-M}$ after scaling. The first bound follows by enclosing the cubes in a ball of radius $3R$.

\subsection{Summation by parts on the frequency grid}

Fix a smooth cutoff $\chi$ equal to one on $|k|\le1$ and zero on $|k|\ge2$. Suppose

\begin{gather*}
\FitMath{R\ge2M\delta,} \\[3pt]
\FitMath{2R<\pi/(2a).}
\end{gather*}

Extend $\chi(k/R)$ periodically on the frequency torus of side $2\pi/a$. It is smooth because its support stays away from the Brillouin faces. The high-frequency lattice symbol

\begin{gather*}
\FitMath{H_a(k)=(1-\chi(k/R))P_a(k)}
\end{gather*}

is smooth on that frequency torus, including at zero. Write

\begin{gather*}
\FitMath{\rho_a(k)=\operatorname{dist}(k,(2\pi/a)\mathbb Z^4).}
\end{gather*}

The derivative estimates and the cutoff product rule give

\begin{gather*}
\FitMath{\|\partial_\mu^M H_a(k)\|
       \le C_M(R+\rho_a(k))^{-M}.}
\end{gather*}

Let $\Delta_\mu H(k)=H(k+\delta e_\mu)-H(k)$. The iterated fundamental theorem of calculus yields

\begin{gather*}
\FitMath{\Delta_\mu^M H(k)=\delta^M
 \int_{[0,1]^M}\partial_\mu^M H
       (k+\delta(t_1+\cdots+t_M)e_\mu)\,dt.}
\end{gather*}

The torus distance $\rho_a$ changes by at most $M\delta$ along this segment. Since $R\ge2M\delta$, $R+\rho_a$ remains comparable to its value at $k$. Thus
\begin{equation}
 \|\Delta_\mu^M H_a(k)\|
       \le C_M\delta^M(R+\rho_a(k))^{-M}.
 \label{eq:lcc-differences}
\end{equation}

Now let $z\in a\mathbb Z^4/\ell\mathbb Z^4$, $z\ne0$, and choose a short representative. Its length $r_z$ is at most $\ell$, and some coordinate satisfies $|z_\mu|\ge r_z/2$. On $|z_\mu|\le\ell/2$,

\begin{gather*}
\FitMath{|e^{-i\delta z_\mu}-1|
 =2|\sin(\pi z_\mu/\ell)|\ge4|z_\mu|/\ell
 \ge2r_z/\ell.}
\end{gather*}

The phase $e^{ik\cdot z}$ is periodic on the frequency grid because $z/a$ is integral. Finite summation by parts therefore gives the exact identity

\begin{gather*}
\FitMath{(e^{-i\delta z_\mu}-1)^M
   \sum_{k\in\mathcal Q_a}e^{ik\cdot z}H_a(k)
 =\sum_{k\in\mathcal Q_a}e^{ik\cdot z}\Delta_\mu^M H_a(k).}
\end{gather*}

There is no Brillouin boundary term. Apply \eqref{eq:lcc-differences} and \eqref{eq:lcc-count-high}, using $\rho_a(k)=|k|$ on the chosen representatives. The result is
\begin{equation}
 \left\|\ell^{-4}\sum_{k\in\mathcal Q_a}
                 e^{ik\cdot z}H_a(k)\right\|
 \le C_M r_z^{-M}R^{4-M}.
 \label{eq:lcc-high-lattice}
\end{equation}
The powers of the torus side cancel: the inverse difference factor contributes $(\ell/r_z)^M$ and the frequency difference contributes $\delta^M=(2\pi/\ell)^M$.

The continuum torus tail obeys the same estimate. Here
$H_0(k)=(1-\chi(k/R))P_0(k)$ on the infinite reciprocal grid. Its derivatives satisfy the same bounds with $|k|$ in place of $\rho_a$. The $M$th differences are absolutely summable for $M>4$. Multiplication of a Fourier distribution by $(e^{-i\delta z_\mu}-1)^M$ shifts its Fourier coefficients and hence gives the same identity distributionally. The right side is an absolutely convergent continuous Fourier series. On the open set where the multiplying factor is nonzero, division by that smooth factor identifies the tail with an ordinary continuous function and proves
\begin{equation}
 \|\mathsf K_\ell^{\mathrm{high}}(z)\|
        \le C_M r_z^{-M}R^{4-M}.
 \label{eq:lcc-high-continuum}
\end{equation}
This distributional argument avoids assuming convergence of an undifferentiated order-zero Fourier series.

\subsection{Comparison with the continuum torus kernel}

On the common low-frequency support, the finite and infinite reciprocal grids agree. Equations \eqref{eq:lcc-low-symbol} and \eqref{eq:lcc-count-low} give

\begin{gather*}
\FitMath{\left\|\ell^{-4}\sum_k e^{ik\cdot z}\chi(k/R)
                  (P_a(k)-P_0(k))\right\|\le CaR^5.}
\end{gather*}

Combining this with the two high-frequency bounds yields, for $M>4$,
\begin{equation}
 \|\mathsf K_{a,\ell}^{\mathrm{raw}}(z)-\mathsf K_\ell(z)\|
       \le C_M\{aR^5+r_z^{-M}R^{4-M}\}.
 \label{eq:lcc-split}
\end{equation}
Here $\mathsf K_{a,\ell}^{\mathrm{raw}}$ is the matrix Fourier sum in \eqref{eq:lcc-fourier-lattice} at the common plaquette-origin displacement $z$.

Take $M=9$ and choose

\begin{gather*}
\FitMath{R=\frac D{r_z}(r_z/a)^{1/10},}
\end{gather*}

where $D$ is a fixed sufficiently large numerical constant. Because $r_z\le\ell$, this ensures $R\ge18\delta$. For $r_z/a$ larger than a fixed numerical threshold $C_0$, it also ensures $2aR<c_0$ and $2R<\pi/(2a)$. Both terms in \eqref{eq:lcc-split} are then at most a constant times

\begin{gather*}
\FitMath{r_z^{-4}(a/r_z)^{1/2}.}
\end{gather*}

Indeed $aR^5=D^5r_z^{-4}(a/r_z)^{1/2}$ and
$r_z^{-9}R^{-5}=D^{-5}r_z^{-4}(a/r_z)^{1/2}$.

For $1\le r_z/a\le C_0$, the lattice Fourier sum has norm at most $a^{-4}$ because $\|P_a\|\le1$. The continuum bound
$\|\mathsf K_\ell(z)\|\le C r_z^{-4}$, established in the next section, shows that their difference is at most $C'a^{-4}$. Increasing the numerical constant by a factor depending only on $C_0$ bounds this by
$Cr_z^{-4}(a/r_z)^{1/2}$. Consequently, for every nonzero raw lattice displacement,
\begin{equation}
 \|\mathsf K_{a,\ell}^{\mathrm{raw}}(z)-\mathsf K_\ell(z)\|
       \le Cr_z^{-4}(a/r_z)^{1/2}.
 \label{eq:lcc-torus-comparison}
\end{equation}
All thresholds used here are fixed numbers. They impose no relation between the physical volume and the ultraviolet cutoff beyond $\ell/a\in\mathbb N$.

\subsection{The finite-volume correction and contact terms}

The continuum torus heat kernel is

\begin{gather*}
\FitMath{p_t^\ell(x)=\sum_{m\in\mathbb Z^4}
       (4\pi t)^{-2}e^{-|x+m\ell|^2/(4t)}
   =\ell^{-4}\sum_{k\in\delta\mathbb Z^4}e^{-t|k|^2}e^{ik\cdot x}.}
\end{gather*}

For completeness, the first expression is an absolutely convergent periodic sum for $t>0$. Its Fourier coefficient at $k$ is the ordinary Gaussian Fourier integral $\ell^{-4}e^{-t|k|^2}$, which proves the second expression. Define the mean-zero torus Green distribution by

\begin{gather*}
\FitMath{G_\ell(x)=\int_0^\infty(p_t^\ell(x)-\ell^{-4})\,dt.}
\end{gather*}

Its nonzero Fourier coefficient is $\ell^{-4}|k|^{-2}$ and its zero coefficient is zero. Applying the same second-derivative combination as in the definition of $\mathsf K$ gives \eqref{eq:lcc-fourier-torus}.

We record the uniform local correction. At $\ell=1$ and $|x|\le1/4$,
\begin{align*}
 G_1(x)-G(x)
 ={}&\int_0^1\sum_{m\ne0}(4\pi t)^{-2}
                 e^{-|x+m|^2/(4t)}\,dt-1\\
 &+\int_1^\infty(p_t^1(x)-1)\,dt
       -\int_1^\infty(4\pi t)^{-2}e^{-|x|^2/(4t)}\,dt.
\end{align*}
Every fixed derivative of the first integral is uniformly bounded: differentiating a Gaussian gives a polynomial times a Gaussian, while $|x+m|\ge c|m|$ and
$\sum_{m\ne0}e^{-c|m|^2/t}\le C e^{-c'/t}$ for $0<t\le1$. The resulting power of $t^{-1}$ is integrable against this last exponential. For the second integral, its differentiated Fourier series is bounded by a constant times $e^{-ct}$ for $t\ge1$. Derivatives of the last integral are bounded by $C_\alpha t^{-2-|\alpha|/2}$ on this ball and are integrable. Thus $G_1-G$ extends smoothly across zero with bounded derivatives of every fixed order.

Scaling gives $G_\ell(x)=\ell^{-2}G_1(x/\ell)$ and the corresponding homogeneity for $G$. In particular,
\begin{equation}
 \|\mathsf K_\ell(z)-\mathsf K(z)\|\le C\ell^{-4}
       \quad(0<|z|\le\ell/4).
 \label{eq:lcc-volume-error}
\end{equation}
The same argument gives an extra factor $\ell^{-1}$ for each additional derivative of the difference. On the remainder of the unit torus, away from zero, the heat-kernel formula and its derivatives converge uniformly on compact sets. Scaling and compactness therefore also give the global bounds

\begin{gather*}
\FitMath{\|\partial^\alpha\mathsf K_\ell(z)\|
       \le C_\alpha r_z^{-4-|\alpha|},} \\[3pt]
\FitMath{z\ne0.}
\end{gather*}

These justify the estimate used in the bounded $r_z/a$ regime above.

All these statements concern the kernels away from coincidence. The symbols have not discarded their order-zero distributional terms. For example, the trace of $P_0(k)$ is three for $k\ne0$. Accordingly the trace of the torus kernel as a distribution is
$3(\delta_{\mathrm{per}}-\ell^{-4})$; away from zero it is $-3\ell^{-4}$. The trace of the free noncoincident kernel is zero. This is consistent with \eqref{eq:lcc-volume-error}; it would be inconsistent to replace the full distribution by its noncoincident formula before taking Fourier coefficients.

\subsection{Moving to barycentres and proving the theorem}

Let $u$ be the centre displacement in \eqref{eq:lcc-main} and let
$z=u-a(c_P-c_Q)$ be its raw-origin displacement. The shift has length at most $a$. Choose representatives so $|u|=r\le\ell/10$ and $|z-u|\le a$. Then

\begin{gather*}
\FitMath{0.9r\le|z|\le1.1r<\ell/4.}
\end{gather*}

Use \eqref{eq:lcc-torus-comparison} for the raw kernel and \eqref{eq:lcc-volume-error} for the continuum torus kernel at $z$. The free kernel is homogeneous of degree $-4$ and smooth away from zero, so $\|\nabla\mathsf K(v)\|\le C|v|^{-5}$. The segment from $z$ to $u$ stays at distance at least $0.9r$ from zero. Hence

\begin{gather*}
\FitMath{|\mathsf K_{PQ}(z)-\mathsf K_{PQ}(u)|\le Ca r^{-5}.}
\end{gather*}

Combining the three differences proves

\begin{gather*}
\FitMath{|\mathsf K_{a,\ell}^{PQ}(x,y)-\mathsf K_{PQ}(u)|
 \le C\{r^{-4}(a/r)^{1/2}+\ell^{-4}+ar^{-5}\}.}
\end{gather*}

Since $a/r\le1$, the last term is absorbed by the first, and
$\ell^{-4}=r^{-4}(r/\ell)^4$. This proves \eqref{eq:lcc-main}.

The argument also proves convergence to $\mathsf K_\ell$ uniformly on any fixed separated subset of the torus, when plaquette centres are used in both arguments. In this statement use the global derivative bound for $\mathsf K_\ell$ to move from origins to centres. No assumption that a chosen separation is small relative to $\ell$ is needed for this fixed-torus convergence.

\subsection{The separated quadratic-density coefficient}

Let $\xi_p^b$, $1\le b\le d_G$, be independent color copies of the Gaussian exact two-form, each with covariance $g^2P_{\mathrm{ex}}$. For a test $f$ on the torus define the finite-dimensional centered density

\begin{gather*}
\FitMath{Q_a(f)=\frac12\sum_{b=1}^{d_G}\sum_p f(x_p)
             \{(\xi_p^b)^2-\mathbb E(\xi_p^b)^2\}.}
\end{gather*}

Here $x_p$ is the plaquette barycentre. Equivalently, $a^{-2}\xi_p^b$ is the physical curvature and the displayed sum includes its $a^4$ Riemann weight. For two tests of positive mutual separation, the Gaussian identity
$\operatorname{Cov}(X^2,Y^2)=2\operatorname{Cov}(X,Y)^2$ gives

\begin{gather*}
\FitMath{\operatorname{Cov}(Q_a(f),Q_a(h))
 =\frac{d_Gg^4}{2}\,a^8\sum_{p,q}
         f(x_p)h(x_q)\,|\mathsf K_{a,\ell}(p,q)|^2.}
\end{gather*}

For real kernels the square here is the ordinary real square of the scalar entry. The uniform separated convergence just proved, followed by ordinary shifted-grid Riemann sums, gives

\begin{gather*}
\FitMath{\lim_{a\downarrow0}\operatorname{Cov}(Q_a(f),Q_a(h))
 =\frac{d_Gg^4}{2}\iint f(x)h(y)
             \sum_{P,Q}\mathsf K_{\ell,PQ}(x-y)^2\,dx\,dy.}
\end{gather*}

For continuous compactly supported separated tests the Riemann-sum integrand is continuous after setting it to zero in a neighbourhood of the diagonal; for $C^1$ tests it is $C^1$ there. There is therefore no use of an integral of $|x-y|^{-8}$ through coincidence.

Now keep the tests in a fixed compact subset of $\mathbb R^4$ and let $\ell\to\infty$. The correction estimate makes $\mathsf K_\ell\to\mathsf K$ uniformly on the relevant separated differences. Direct differentiation of $G$ shows

\begin{gather*}
\FitMath{\mathsf K(x)=\frac1{\pi^2|x|^4}\Lambda^2(I-2\widehat x\otimes\widehat x).}
\end{gather*}

One can check the coefficient at $\widehat x=e_4$: the six eigenvalues are three copies of $1/(\pi^2|x|^4)$ and three copies of its negative. Orthogonal covariance proves the identity in every direction. The induced two-form reflection is orthogonal, and therefore

\begin{gather*}
\FitMath{\sum_{P,Q}\mathsf K_{PQ}(x)^2=\frac6{\pi^4|x|^8}.}
\end{gather*}

It follows that the infinite-volume separated covariance is exactly
\begin{equation}
 \frac{3d_Gg^4}{\pi^4}
       \iint f(x)h(y)|x-y|^{-8}\,dx\,dy.
 \label{eq:lcc-density-coefficient}
\end{equation}
For $\mathrm{SU}(2)$, $d_G=3$, so the coefficient is $9g^4/\pi^4$ in this convention. Higher separated Gaussian cumulants can likewise be passed through their finite cycle formulas using the same kernel convergence.

There is also a relative finite-volume statement for nonnegative pairs. Suppose every cross-distance between the supports lies between $d$ and $C_{\mathrm{sh}}d$, and $C_{\mathrm{sh}}d\le\ell/10$. The correction estimate and the Frobenius norm of the free kernel give pointwise

\begin{gather*}
\FitMath{\left|\sum_{P,Q}\mathsf K_{\ell,PQ}(v)^2
              -\frac6{\pi^4|v|^8}\right|
 \le C\left(\frac{|v|}{\ell}\right)^4|v|^{-8}.}
\end{gather*}

Indeed the difference of the squares is bounded by
$C|v|^{-4}\ell^{-4}+C\ell^{-8}$. Multiplication by $f(x)h(y)\ge0$ and integration show that the limiting torus covariance differs from \eqref{eq:lcc-density-coefficient} by a relative error at most $C_{\mathrm{sh}}'(d/\ell)^4$, provided both tests are nonzero. The mesh error tends to zero for each such fixed test pair by the uniform kernel estimate and the Riemann-sum argument above. A rate relative to a varying family of tests would require their norms and nondegeneracy to be specified.

This proves the free covariance comparison and the free density coefficient. It asserts neither an $L^2$ limit of $Q_a(f)$ with the same test repeated nor an interacting short-distance asymptotic. Establishing the latter requires a comparison of the interacting measure with this Gaussian model and a remainder smaller than the leading term.

% END MODULE chapters/lattice_curvature_comparison_v2.tex


% BEGIN MODULE chapters/gibbs_symmetry_v2.tex
\section{Periodic Gibbs states: exact layer elimination and link reflections}
\label{ch:periodic-link-symmetry}

This section concerns one fixed lattice and one fixed positive coupling.
No continuum limit or renormalization estimate is used. The conclusion
quantifies over every periodic DLR state. In particular, the state is not
assumed to arise as a limit of torus measures. The finite-dimensional
transfer calculation below, together with a local elimination identity,
is what permits that quantifier.

Set the lattice spacing equal to one, let $G=\mathrm{SU}(2)$, and use
normalized Haar measure. On one orientation of the bonds of
$\mathbb Z^4$ assign variables $U_b\in G$; reverse orientation means
$U_{\bar b}=U_b^{-1}$. For an elementary plaquette set

\begin{gather*}
\FitMath{a_p=1-\tfrac12\operatorname{tr}U_p,} \\[3pt]
\FitMath{0\le a_p\le2.}
\end{gather*}

The trace is real for $G=\mathrm{SU}(2)$. Fix $\beta>0$. A DLR state
is a probability measure whose conditional law on any finite bond set
$D$, given all other bonds, has density proportional to
$\exp[-\beta\sum_{p:p\cap D\ne\varnothing}a_p]$. A state is periodic
when it is invariant under a subgroup of finite index of
$\mathbb Z^4$. A link reflection is reflection in a plane
$\{x_\mu=i+\frac12\}$, with $i\in\mathbb Z$.

\begin{theorem}[Link symmetry of every periodic Wilson state]
Every periodic DLR state of the four-dimensional $\mathrm{SU}(2)$
Wilson interaction at $\beta>0$ is invariant under every link
reflection. Consequently every such state is invariant under
$2\mathbb Z^4$.
\end{theorem}

The proof uses positivity and strict positivity of different objects.
The one-bond convolution operator is positive as an operator, with
strictly positive Peter--Weyl eigenvalues. Its original kernel is
strictly positive pointwise. An auxiliary operator inserted between
two half-kernels will be a contraction, and need not be self-adjoint
or positive. Confusing these conditions would invalidate the
variational step.

\subsubsection*{A lower tangent valid for an arbitrary periodic state}

Let $\phi$ be a real continuous local observable and let
$m\mathbb Z^4$ be a subgroup of periods of a DLR state $\nu$.
For torus sides $s$ divisible by $m$, write

\begin{gather*}
\FitMath{W_s=\sum_{v\in m\mathbb Z^4/s\mathbb Z^4}\tau_v\phi,} \\[3pt]
\FitMath{Z_s(\phi)=\int e^{W_s-\beta\sum_pa_p}\,dU.}
\end{gather*}

Here local observables are evaluated on the periodic extension of
the torus configuration. Suppose

\begin{gather*}
\FitMath{Z_s(\phi)\le Z_s(0)} \\[3pt]
\FitMath{\hbox{for all sufficiently large multiples }s\hbox{ of }m.}\tag{GL1}
\end{gather*}

Then

\begin{gather*}
\FitMath{\nu(\phi)\le0.}\tag{GL2}
\end{gather*}


To prove this implication, take a cube $D_s$ of sites of side $s$,
and let $B_s$ be the positively oriented bonds whose initial sites
belong to the cube. Choose the representatives of the source
anchors inside $D_s$, and denote their sum, evaluated with an
arbitrary outside configuration, by $\widetilde W_s$. This is a
finite local function; terms near the boundary can involve outside
bonds. The number of anchors is exactly $M_s=(s/m)^4$.
The DLR equation and Jensen's inequality give

\begin{gather*}
\FitMath{\log\gamma_{B_s}(e^{\widetilde W_s})(U')
       \ge\gamma_{B_s}(\widetilde W_s)(U').}
\end{gather*}

The conditional partition-function ratio on the left can be
compared with $Z_s(\phi)/Z_s(0)$ uniformly in $U'$. Interior terms
are identical under the identification of $B_s$ with the torus
bond coordinates. Only plaquettes and translates of $\phi$ in a
fixed-width boundary layer can differ. Their number is at most
$C_{\phi,m}s^3$, and each difference is bounded by
$C(\beta+\|\phi\|_\infty)$. Thus

\begin{gather*}
\FitMath{\log\frac{Z_s(\phi)}{Z_s(0)}
   \ge \nu(\widetilde W_s)-C_{\phi,m,\beta}s^3
   = M_s\nu(\phi)-C_{\phi,m,\beta}s^3.}
\end{gather*}

In the first inequality we integrated the conditional inequality
against $\nu$ and used the DLR equation. The equality uses only
$m\mathbb Z^4$ invariance. Divide by $M_s$ and let $s\to\infty$
through multiples of $m$. Assumption (GL1) proves (GL2).

This argument needs neither differentiability of an infinite-volume
pressure nor representation of $\nu$ by finite torus measures.
The local source is fixed during the limit $s\to\infty$. If the
source later depends on a parameter tending to zero, that parameter
is kept fixed until the surface error has disappeared. Constant
parts of $\phi$ cause no difficulty with the anchored definition of
$\widetilde W_s$: they occur with exactly the same multiplicity in
the box and torus source sums.

\subsubsection*{The one-bond convolution and its smooth conjugations}

Define

\begin{gather*}
\FitMath{c(g)=e^{-\beta(1-\frac12\operatorname{tr}g)},} \\[3pt]
\FitMath{(CF)(w)=\int_G c(wv^{-1})F(v)\,dv.}
\end{gather*}

The kernel is real, symmetric, and bounded below by $e^{-2\beta}$.
Let $\Pi_l$ be the Peter--Weyl projection corresponding to the
irreducible representation of spin $l\in\frac12\mathbb N_0$ and
dimension $d_l=2l+1$. Schur's lemma and the Weyl integration
formula give

\begin{gather*}
\FitMath{C=\sum_l c_l\Pi_l,} \\[3pt]
\FitMath{c_l=\frac{2e^{-\beta}}{\beta}I_{2l+1}(\beta)>0.}\tag{GL3}
\end{gather*}

For completeness, write $g$ with eigenvalues $e^{it},e^{-it}$ and
use $\chi_l(g)=\sin((2l+1)t)/\sin t$. The character coefficient is

\begin{gather*}
\FitMath{d_lc_l=\frac{2e^{-\beta}}\pi
  \int_0^\pi e^{\beta\cos t}\sin t\sin(d_lt)\,dt
  =e^{-\beta}(I_{2l}-I_{2l+2})(\beta).}
\end{gather*}

The Bessel recurrence gives (GL3). Its series

\begin{gather*}
\FitMath{I_n(\beta)=\sum_{k\ge0}
       \frac{(\beta/2)^{n+2k}}{k!(n+k)!}}
\end{gather*}

has positive terms and implies factorial decay in $n$.
Moreover,

\begin{gather*}
\FitMath{\frac{I_{n+1}}{I_n}\le\frac{\beta}{2(n+1)},} \\[3pt]
\FitMath{\frac{I_n}{I_{n+1}}
  =\frac{2(n+1)}\beta+\frac{I_{n+2}}{I_{n+1}}
  \le C_\beta(1+n).}\tag{GL4}
\end{gather*}

Consequently both ratios between neighboring $c_l$ grow at most
polynomially in $l$. Factorial decay also shows that $C$ is trace
class and that $C^{1/2}$ has a real symmetric smooth kernel.

For a finite block $\mathcal B$ of spatial bonds put

\begin{gather*}
\FitMath{\mathcal C=C_{\mathcal B}=\bigotimes_{b\in\mathcal B}C,} \\[3pt]
\FitMath{c_{\boldsymbol l}=\prod_{b\in\mathcal B}c_{l_b},} \\[3pt]
\FitMath{\Pi_{\boldsymbol l}=\bigotimes_{b\in\mathcal B}\Pi_{l_b}.}
\end{gather*}

Its kernel is
$\mathcal C(u,v)=\prod_{b\in\mathcal B}c(u_bv_b^{-1})$.
In particular it is strictly positive on the compact product
space. It commutes with the simultaneous vertex gauge action on
its two arguments.

Let $f$ be a polynomial in matrix coordinates and their conjugates
on $G^{\mathcal B}$. Define on $C^\infty(G^{\mathcal B})$

\begin{gather*}
\FitMath{\widehat f F=
   \sum_{\boldsymbol l,\boldsymbol l'}
    \frac{c_{\boldsymbol l}}{c_{\boldsymbol l'}}
       \Pi_{\boldsymbol l'}M_f\Pi_{\boldsymbol l}F,}\tag{GL5}
\end{gather*}

where $M_f$ denotes multiplication by $f$. This is a continuous
map of the smooth Frechet space into itself. Indeed, if the degree
in coordinate $b$ is $d_b$, the Clebsch--Gordan rule makes the
summand zero unless $|l_b-l_b'|\le d_b/2$. By (GL4), the remaining
ratios are bounded by a fixed polynomial in the indices. Smooth
functions have Peter--Weyl coefficients decreasing faster than
every polynomial, as follows by applying powers of the Casimir.
The finite-dimensional reproducing-kernel estimate bounds the
supremum norm on a block by a polynomial times its $L^2$ norm.
The same holds after any fixed number of derivatives. These facts
give convergence of (GL5) with all derivatives and the asserted
continuity.

The notation $\widehat f=\mathcal C^{-1}M_f\mathcal C$ is therefore
legitimate on this smooth domain when interpreted by (GL5).
It does not assert that $\mathcal C^{-1}$ acts on every smooth
function or that $\widehat f$ is bounded on $L^2$. The formal
$L^2$ adjoint on smooth functions is given by the same finite-band
series with the coefficient ratios reversed and $f$ replaced by
$\overline f$. Those ratios also have polynomial growth, so this
formal adjoint maps smooth functions continuously to smooth
functions. In particular no unbounded inverse is applied outside
its justified spectral domain in the following identity:

\begin{gather*}
\FitMath{\mathcal C\widehat fF=M_f\mathcal CF.}\tag{GL6}
\end{gather*}

It follows by multiplying (GL5) by $\mathcal C$ and summing the
projections. In kernel form, for every $w$,

\begin{gather*}
\FitMath{f(w)\int\mathcal C(v,w)F(v)\,dv
       =\int\mathcal C(v,w)(\widehat fF)(v)\,dv.}\tag{GL7}
\end{gather*}

The equality is pointwise, since the series and both integrals
define smooth functions. All these assertions also hold for
smooth families depending on extra variables on a finite compact
product: differentiation in those variables commutes with the
series. The operator $\widehat f$ commutes with multiplication by
functions of variables on which it does not act, and operators
acting on disjoint variable sets commute.

\subsubsection*{A transfer contraction inequality}

Fix a direction, called time, and take a cubic torus of side $s$.
Let $\mathcal E$ be its spatial bonds in one layer. On
$L^2(G^{\mathcal E})$, let $P$ be vertex gauge averaging and let
$V$ be $\beta$ times the spatial plaquette action in that layer.
The operator

\begin{gather*}
\FitMath{L=P C_{\mathcal E}^{1/2}e^{-V}C_{\mathcal E}^{1/2}P}
\end{gather*}

is positive and trace class. The convolution, its square root,
and $e^{-V}$ all commute with $P$. Integrating the time-like bonds
in the torus integral inserts $P$ between successive layers.
Using cyclicity of the trace gives

\begin{gather*}
\FitMath{Z_s(0)=\operatorname{Tr}L^s.}\tag{GL8}
\end{gather*}

Here normalized Haar measure makes the gauge average an
orthogonal projection without an additional volume factor.

Let $\Xi$ be a real bounded operator on $L^2(G^{\mathcal B})$
commuting with the block vertex gauge action, and suppose
$\|\Xi\|\le1$. Suppose the kernel

\begin{gather*}
\FitMath{k_\Xi=\text{kernel of }\mathcal C^{1/2}\Xi\mathcal C^{1/2}}
\end{gather*}

is continuous and strictly positive pointwise. In two adjacent
layers, use transported spatial variables $u,v$ and set

\begin{gather*}
\FitMath{\phi_\Xi(u,v)=\log\frac{k_\Xi(u,v)}{\mathcal C(u,v)}.}\tag{GL9}
\end{gather*}

This is a real continuous gauge-invariant local observable.

Place translates of the block so far apart that blocks in a
single time step have disjoint bond sets. Multiplication by
$e^{\phi_\Xi}$ replaces the corresponding factor $\mathcal C$
by $\mathcal C^{1/2}\Xi\mathcal C^{1/2}$. Thus the full kernel in
step $j$ is
$C_{\mathcal E}^{1/2}\Xi_j C_{\mathcal E}^{1/2}$, where $\Xi_j$
is a tensor product of copies of $\Xi$ and identities and has
norm at most one. Gauge averaging the time-like bonds gives
the trace of the product of
$e^{-V}C_{\mathcal E}^{1/2}\Xi_j C_{\mathcal E}^{1/2}P$.
Moving the first half-kernel and layer weight cyclically to the
end, and using commutation with $P$, rewrites this trace as

\begin{gather*}
\FitMath{\operatorname{Tr}(\Xi_1L\Xi_2L\cdots\Xi_sL).}
\end{gather*}

The unperturbed steps have $\Xi_j=1$. Schatten Holder gives
\begin{equation}
\FitMath{\left|\operatorname{Tr}(\Xi_1L\cdots\Xi_sL)\right|
 \le\prod_{j=1}^s\|\Xi_jL\|_s
 \le\|L\|_s^s=\operatorname{Tr}L^s.}
 \tag{GL10}
\end{equation}
The last equality uses $L\ge0$. The modified integral is
positive because every replaced kernel is positive pointwise.
Therefore (GL10) is precisely the partition-function inequality
(GL1) for a sufficiently sparse periodic array of $\phi_\Xi$.

Given an arbitrary periodic state $\nu$, choose $m$ large enough
that $m\mathbb Z^4$ is a subgroup of its periods and the copies
are disjoint. Equations (GL1)--(GL2) imply

\begin{gather*}
\FitMath{\nu(\phi_\Xi)\le0.}\tag{GL11}
\end{gather*}

The norm contraction and the lower tangent are the only
ingredients of this conclusion. No self-adjointness of $\Xi$
and no differentiability of a limiting pressure has been used.

\subsubsection*{Two adjacent layers cannot carry an antisymmetric mean}

Let $Y(u,v)$ be a real continuous function on
$G^{\mathcal B}\times G^{\mathcal B}$ that is invariant under
simultaneous vertex gauge transformations and antisymmetric:
$Y(v,u)=-Y(u,v)$. Then every periodic DLR state satisfies

\begin{gather*}
\FitMath{\nu(Y(U_i,U_{i+1}))=0.}\tag{GL12}
\end{gather*}

To prove this, first suppose $K(u,v)=Y(u,v)\mathcal C(u,v)$ is
a polynomial. Its integral operator has finite Peter--Weyl
support. Antisymmetry makes $K$ a real skew-adjoint operator.
On that finite-dimensional support define

\begin{gather*}
\FitMath{D=\mathcal C^{-1/2}K\mathcal C^{-1/2}.}
\end{gather*}

Every eigenvalue $c_{\boldsymbol l}$ is strictly positive, so
the definition is valid. Extend $D$ by zero on the orthogonal
complement. It is a finite-rank real skew-adjoint operator,
commutes with gauge transformations, and
$\Xi_\varepsilon=e^{\varepsilon D}$ is a real orthogonal
operator for real $\varepsilon$ of either sign. In particular
$\|\Xi_\varepsilon\|=1$.

The kernel difference from $\mathcal C$ has finite
Peter--Weyl support and admits, uniformly in $(u,v)$,

\begin{gather*}
\FitMath{k_{\Xi_\varepsilon}(u,v)
   =\mathcal C(u,v)+\varepsilon K(u,v)+O(\varepsilon^2).}
\end{gather*}

Since $\mathcal C\ge e^{-2\beta|\mathcal B|}$, this kernel
is strictly positive for sufficiently small $|\varepsilon|$.
Consequently
$\phi_{\Xi_\varepsilon}=\varepsilon Y+O(\varepsilon^2)$
uniformly. Apply (GL11), first at each fixed nonzero
$\varepsilon$ after taking the thermodynamic limit in the
lower-tangent argument. For both signs it gives
$\varepsilon\nu(Y)\le C\varepsilon^2$. Letting
$\varepsilon\to0$ proves (GL12) in this case.

For general $Y$, approximate $Y\mathcal C$ uniformly by
polynomials on the compact product group. Antisymmetrize each
approximant and average it over the simultaneous gauge group.
These operations preserve polynomiality, because a finite
matrix-coordinate degree spans a finite-dimensional
representation, and preserve the uniform approximation.
Divide by the strictly positive kernel $\mathcal C$.
The resulting functions $Y_n$ converge uniformly to $Y$ and
have the already proved property $\nu(Y_n)=0$. This proves
(GL12). The argument also specifies why injectivity of
$\mathcal C$, rather than merely its operator positivity,
is needed.

\subsubsection*{The exact one-layer DLR elimination identity}

Fix a finite stack of layers $[a,b]$. Write $T_{j,y}$ for
the time-like bond from $(j,y)$ to $(j+1,y)$ and
$\mathring U_{j,e}$ for a spatial bond in layer $j$. Define

\begin{gather*}
\FitMath{g_{j,y}=T_{a,y}\cdots T_{j-1,y},} \\[3pt]
\FitMath{U_{j,\langle y,y'\rangle}
  =g_{j,y}\mathring U_{j,\langle y,y'\rangle}g_{j,y'}^{-1}.}\tag{GL13}
\end{gather*}

These are temporal-gauge coordinates on the stack, with the
empty product at $j=a$. A gauge-invariant local function of
the stack is exactly a function of the $U_j$ invariant under
one common vertex gauge action in all layers. This follows
by setting the time-like bonds to the identity successively.
For any fixed finite observable only finitely many vertices
are involved, so only finite-support gauge transformations
are needed.

Let $a<j<b$ and resample the spatial bonds
$\{j\}\times\mathcal B$, leaving every other original bond
conditioned. The matrices $g_{j,y}$ use only time-like bonds
and hence are fixed during this resampling. The substitution
$\mathring U_j\mapsto v=U_j|_{\mathcal B}$ preserves Haar
measure. The two time-like plaquettes attached to each
resampled spatial bond produce the factors below; all
spatial plaquettes that meet the block produce the finite
potential $V_j(v;\eta)$, where $\eta$ denotes the other
spatial variables of layer $j$. Thus the conditional density
is exactly

\begin{gather*}
\FitMath{\frac1{z(w,w',\eta)}
 \mathcal C(w,v)e^{-V_j(v;\eta)}\mathcal C(v,w')\,dv,} \\[3pt]
\FitMath{w=U_{j-1}|_{\mathcal B},} \\[3pt]
\FitMath{w'=U_{j+1}|_{\mathcal B}.}\tag{GL14}
\end{gather*}

No plaquette meeting the resampled bonds is omitted: it is
either spatial in layer $j$ or one of the two time-like
plaquettes of that bond. The potential can involve spatial
neighbors outside $\mathcal B$, but they lie in the same
layer and in a finite enlargement of the block.

Let $f$ be a real polynomial on $G^{\mathcal B}$ and let
$\Phi$ be smooth in finitely many transported variables
from layers $j$ and above, and independent of the spatial
variables of layer $j-1$. Define

\begin{gather*}
\FitMath{\Phi^\flat
  =\frac{e^{V_j(v;\eta)}}{\mathcal C(v,w')}
     \widehat f_v\!\left[
        e^{-V_j(v;\eta)}\mathcal C(v,w')\Phi\right].}\tag{GL15}
\end{gather*}

All factors are smooth on a finite compact product, and
the denominator is strictly positive. Equation (GL5)
therefore proves that $\Phi^\flat$ is smooth. Its new
dependence is confined to layer $j$ and the already present
higher layers; it acquires no spatial variable in layer
$j-1$.

Every DLR state, periodic or not, satisfies the exact identity
\begin{equation}
\FitMath{\nu\big(f(U_{j-1}|_{\mathcal B})\Phi\big)
       =\nu(\Phi^\flat).}
 \tag{GL16}
\end{equation}
Indeed, condition on all bonds except the specified block
in layer $j$. The conditional numerator on the left is

\begin{gather*}
\FitMath{f(w)\int\mathcal C(v,w)
       e^{-V_j(v;\eta)}\mathcal C(v,w')\Phi\,dv.}
\end{gather*}

Apply (GL7) to the smooth bracketed function. Substitution
of (GL15) identifies the result with the numerator of the
conditional expectation on the right. The same normalizer
$z(w,w',\eta)$ occurs on both sides; it is not
differentiated or canceled through an operator. Integration
of this conditional identity gives (GL16). Reflection of
the construction yields the corresponding identity that
eliminates an upper outer layer.

\subsubsection*{Removing two outer layers and preserving reflection}

Let $\theta$ be reflection about the middle of a stack with
$2k+2$ layers, numbered $i-k,\ldots,i+k+1$. Let $X$ be a
real gauge-invariant polynomial observable in that stack.
In temporal coordinates, a finite sum of products writes
it as

\begin{gather*}
\FitMath{X=\sum_\alpha f_\alpha(w)h_\alpha(z)\Phi_\alpha,}
\end{gather*}

where $w=U_{i-k}|_{\mathcal B}$ and
$z=U_{i+k+1}|_{\mathcal B}$ are the outer layers and
$\Phi_\alpha$ depends on the inner layers. Take one finite
block containing all the variables that occur in this
decomposition. The potentials used during elimination can
enlarge this block spatially, but only finitely and only in
the inner layers.

Use (GL16) to remove $f_\alpha(w)$ through the first inner
layer $u=U_{i-k+1}|_{\mathcal B}$, and then its reflected
version to remove $h_\alpha(z)$ through the last inner
layer $v=U_{i+k}|_{\mathcal B}$. The result is

\begin{gather*}
\FitMath{\nu(X)=\nu(EX),} \\[3pt]
\FitMath{EX=\sum_\alpha \Upsilon^{-1}
       \widehat f_{\alpha,u}\widehat h_{\alpha,v}
          (\Upsilon\Phi_\alpha).}\tag{GL17}
\end{gather*}

When the two inner layers are not adjacent, $\Upsilon$ is
the product of their two factors $e^{-V}$ and their bond
kernels toward the remaining interior. When $k=1$ they
are adjacent and
\begin{equation}
\FitMath{\Upsilon=e^{-V(u;\eta)}e^{-V(v;\eta')}
                    \mathcal C(u,v).}
 \tag{GL18}
\end{equation}
The common bond kernel occurs once, not twice. To see the
cancellation directly, the first elimination introduces
$e^{V(u)}/\mathcal C(u,v)$ outside its conjugated
operator. In the upper elimination the conditional
numerator contains $e^{-V(v)}\mathcal C(u,v)$, canceling
that denominator before the second operator acts. The
operators act on distinct variables and commute. Moving
the remaining independent factors through them gives
exactly (GL18). This calculation is also a check that
normalizing factors from different conditional laws have
not been interchanged.

Let $R$ denote literal reversal of the array of transported
layer coordinates. An actual reflection of the original
bond configuration induces $R$ followed by a common vertex
gauge transformation arising from the change of temporal-gauge
base. The spatial potential is the same function of its own
layer variables at both ends, and $\mathcal C(u,v)$ is
symmetric. Reflection interchanges $u$ and $v$, including
all their spatial neighbors, and preserves $\Upsilon$.
It also sends $\widehat f_u$ to $\widehat f_v$.
Consequently coordinate elimination commutes with literal
reversal:

\begin{gather*}
\FitMath{E(RX)=R(EX).}\tag{GL19}
\end{gather*}

The output need not be a polynomial, but it is a smooth
function of finitely many inner-layer variables.

Every DLR state is invariant under finite-support vertex
gauge transformations. This follows by applying its
conditional law to the finite set of bonds incident to
the transformed vertices and changing variables using
Haar invariance and plaquette conjugacy. We may therefore
replace $EX$ in expectation by its average under the
common temporal-gauge action; denote it by $\mathcal A$.
This averaging commutes with $R$. Once averaged, the choice of the
lowest layer in (GL13) disappears: changing the base from
$i-k$ to $i-k+1$ conjugates all remaining transported
variables by one common vertex gauge transformation.
The averaged function is invariant under that
transformation, even when its values are the actual
time-like transport matrices. This last statement is a
pointwise identity of invariant functions, and does not
invoke invariance of the state under a field-dependent
change of variables. Thus the averaged output is a
continuous gauge-invariant observable of the smaller
stack of $2k$ layers. Since $X$ is gauge invariant, its actual
reflection is represented by $RX$ in temporal coordinates.
After averaging the common gauge factor on the output also
disappears. Thus the identity needed on the original bond
space is $\mathcal A E(\theta X)=\theta\mathcal A EX$.

\subsubsection*{Proof of the link-reflection theorem}

Fix the reflection $\theta$ in $x_\mu=i+\frac12$.
We prove by induction on $k$ that every real
gauge-invariant polynomial in the stack
$[i-k+1,i+k]$ has equal expectation to its reflection.

For $k=1$, temporal gauge writes the observable as
$\Psi(U_i,U_{i+1})$. Reflection interchanges the two
arguments, up to a common gauge transformation, which
does not affect $\Psi$. The difference from its
reflection is the function
$Y(u,v)=\Psi(u,v)-\Psi(v,u)$, to which (GL12)
applies. Hence the expectations agree.

Assume the assertion for $k$. For a polynomial $X$ in
the stack with $2k+2$ layers, (GL17)--(GL19) and gauge
averaging reduce the comparison of $\nu(X)$ and
$\nu(\theta X)$ to the same comparison for a
continuous gauge-invariant function of the smaller
stack. Invariant polynomials are uniformly dense in
the invariant continuous functions on a finite product
of compact matrix groups: first apply Stone--Weierstrass
to matrix coordinates and their conjugates, then
average the approximants over the finite vertex gauge
group. The induction hypothesis therefore applies to
this continuous output. It gives
$\nu(\mathcal A EX)=\nu(\theta\mathcal A EX)$ and proves
the next step.

Every gauge-invariant continuous local function fits
in a finite symmetric stack and is uniformly
approximated as above. We have proved equality of
its expectation and that of its reflection.
For a general continuous local $F$, gauge average
over the finite set of vertices relevant to its
support. Gauge invariance of both $\nu$ and its
reflected state shows that both expectations of
$F$ equal the expectations of that gauge average.
Thus $\nu(F)=\nu(\theta F)$ for all continuous
local $F$, which determine the probability measure
on the compact product configuration space. The
state is reflection invariant. The direction $\mu$
and integer $i$ were arbitrary. Composing
reflections in two successive link planes gives
translation by $2e_\mu$, proving the stated
$2\mathbb Z^4$ invariance.\hfill$\square$

\subsubsection*{Passage to all mirrors}

The common period is the input to the independent fold-and-separation
proof in Section~\ref{ch:all-mirror-periodic}. That section supplies
adapted diagonal split tori, finite crossing kernels, the full
chessboard inequality, and the arbitrary-state separation argument.
The link theorem above does not require those later ingredients.

% END MODULE chapters/gibbs_symmetry_v2.tex


% BEGIN MODULE chapters/all_mirror_rp_audit_v3.tex
\section{All lattice mirrors for every periodic Wilson state}
\label{ch:all-mirror-periodic}

The result in this section concerns the complete algebra of bonds on a
closed side of a mirror. Its proof does not identify an arbitrary Gibbs
state with a limit of unperturbed torus measures. Instead, a pressure
inequality produces, for each separating observable, a suitable limit of
\emph{perturbed} torus measures. The common period supplied by link
reflection invariance is what makes the separation argument work.

We use the following result proved in the preceding link-symmetry
argument: every periodic DLR state of the four-dimensional
$\mathrm{SU}(2)$ Wilson specification, at $\beta>0$, is invariant under
all coordinate link reflections and consequently under
$P=2\mathbb Z^4$. No uniqueness or extremality is assumed here.

Let $\Omega=\mathrm{SU}(2)^{\mathcal B}$, with one group variable for each
positively oriented nearest-neighbour bond; reversing a bond inverts its
variable. Put

\begin{gather*}
\FitMath{a_p(U)=1-\tfrac12\operatorname{tr}U_p,} \\[3pt]
\FitMath{c(g)=\exp[-\beta(1-\tfrac12\operatorname{tr}g)].}
\end{gather*}

The trace is real for $\mathrm{SU}(2)$, and $0\le a_p\le2$.
The product topology makes $\Omega$ compact and metrizable. A state is
periodic if it is invariant under some finite-index subgroup of
$\mathbb Z^4$. Denote by $\mathcal G_\beta^P$ the compact convex set of
Wilson DLR states invariant under $P$.

The normals considered below are

\begin{gather*}
\FitMath{\mathcal N=\{e_\mu:1\le\mu\le4\}
 \cup\{e_\mu+e_\lambda,e_\mu-e_\lambda:1\le\mu<\lambda\le4\}.}
\end{gather*}

For $n\in\mathcal N$ write $q=|n|^2$ and

\begin{gather*}
\FitMath{\theta_c x=x-2(n\cdot x-c)n/q.}
\end{gather*}

The admissible positions are $c\in\frac12\mathbb Z$ when $q=1$ and
$c\in\mathbb Z$ when $q=2$. These and only these positions in the stated
families give lattice reflections. Let $\mathcal A_c^+$ be the bounded
Borel functions of bonds both of whose endpoints satisfy $n\cdot x\ge c$,
and define $\mathcal A_c^-$ similarly. The antilinear reflection is
$\Theta_cF=\overline{F\circ\theta_c^*}$.

\begin{theorem}[All-mirror positivity of periodic states]
\label{thm:amrp-all-periodic}
Every periodic Wilson DLR state at $\beta>0$ is invariant under every
admissible reflection $\theta_c$ in the sixteen families above. For every
such mirror,

\begin{gather*}
\FitMath{\nu(\Theta_cF\,F)\ge0} \\[3pt]
\FitMath{(F\in\mathcal A_c^+\text{ or }F\in\mathcal A_c^-).}\tag{AM1}
\end{gather*}

In particular, the conclusion applies to every bounded observable on
any finite stack contained in either closed half-space, regardless of
the number of layers in that stack.
\end{theorem}

Composing the reflections in $x_\mu=0$ and $x_\mu=1/2$ gives
translation by one lattice unit. Coordinate reflections give sign
changes, and reflection in $x_\mu-x_\lambda=0$ exchanges two
coordinates. Thus every such state is invariant under all lattice
translations and the full signed-permutation group. At $\beta=0$,
the DLR equations directly force every finite-bond marginal to be
product Haar; that unique product state has the same properties.

\subsection*{Positive crossing kernels on split tori}

Write an element of $\mathrm{SU}(2)$ as a unit quaternion
$s(g)\in S^3\subset\mathbb R^4$. Then

\begin{gather*}
\FitMath{\tfrac12\operatorname{tr}(gh^{-1})=s(g)\cdot s(h).}
\end{gather*}

Consequently

\begin{gather*}
\FitMath{c(gh^{-1})=e^{-\beta}\sum_{m=0}^\infty\frac{\beta^m}{m!}
 \sum_{i_1,\ldots,i_m=1}^4
 \prod_{j=1}^m s_{i_j}(g)\prod_{j=1}^m s_{i_j}(h).}\tag{AM2}
\end{gather*}

Every coefficient is nonnegative. Each feature has absolute value at
most one, and the sum of the coefficients is $e^{3\beta}$, so the
series converges absolutely and uniformly. Tensor products of finitely
many such expansions are again uniformly convergent positive Gram
expansions. The coarse total coefficient bound is immaterial; each
torus used below is finite.

A torus is called split for the mirror normal $n$ if its period lattice
has the form

\begin{gather*}
\FitMath{\Gamma=\Gamma_\parallel\oplus\mathbb Z(2z n/q),} \\[3pt]
\FitMath{\Gamma_\parallel\subset n^\perp,}\tag{AM3}
\end{gather*}

where $\Gamma_\parallel$ has rank three, $2zn/q$ is a lattice vector,
and $z$ is a sufficiently large integer. The reflection has two fixed
cuts, at heights $c$ and $c+z$, on this torus. The two intervening closed
halves are exchanged. All tori below have arbitrarily large shortest
nonzero period, so that no small-quotient ambiguity affects plaquettes.

We verify reflection positivity on these tori for all bounded
half-space functions. Partition the bond variables into the interiors
of the two halves, the bonds lying on the cuts, and, when present, the
bonds crossing a cut. A coordinate site cut has no plaquette with
vertices strictly on both sides. Conditional on its boundary bonds,
the action is a sum of a positive-side action, its reflected action,
and an action on the boundary. Its reflection form is thus an integral
of absolute squares.

At a diagonal cut, signed permutations reduce the geometry to
$n=e_1+e_2$, $c\in\mathbb Z$. A bond changes the height $x_1+x_2$ by
$0$ or $1$ in absolute value, so no bond crosses from a strictly
negative height to a strictly positive one. The only plaquettes meeting
both strict sides are in the $(1,2)$ plane, with heights
$c-1,c,c,c+1$. Each is cut into two two-bond paths between the two
vertices at height $c$. Orienting both paths in the same direction
expresses its Wilson factor as $c(P_+P_-^{-1})$. Reflection interchanges
the paths, and (AM2) supplies the required Gram expansion. Plaquettes
in all other planes lie in a closed half. The bonds on the mirror, in
directions $3$ and $4$, are fixed pointwise. Conditioning on them and
multiplying the crossing expansions again gives an integral of squares.
The same argument applies at the second cut.

For a coordinate link cut there are crossing bonds. Integrate them
before evaluating a half-space function, which cannot depend on them.
Let $X$ and $Y$ be the collections of bonds parallel to the cut in its
two adjacent layers. The product crossing kernel before integration is

\begin{gather*}
\FitMath{C(X,Y)=\prod_b c(X_bY_b^{-1}).}
\end{gather*}

It is a positive kernel by (AM2). If $T=(T_x)$ is a collection of group
elements at the vertices of one boundary layer, set
$(T\cdot X)_{xy}=T_xX_{xy}T_y^{-1}$. Let $P_0$ be Haar averaging under
this action. Simultaneous gauge covariance gives
$C(T\cdot X,T\cdot Y)=C(X,Y)$, since each plaquette trace changes by
conjugation. Hence the associated integral operator commutes with the
orthogonal projection $P_0$. Integration of the crossing bonds gives
the kernel of

\begin{gather*}
\FitMath{P_0C=CP_0=P_0CP_0\ge0.}\tag{AM4}
\end{gather*}

Equivalently, in each Gram feature of the product expansion one may
Haar-average both its $X$ argument and its $Y$ argument. Absolute
uniform convergence justifies this operation. The crossing bonds at
the two cuts are disjoint, so the two constructions multiply. This
argument places no gauge-invariance condition on the test function.

In all three geometries the resulting reflection form has the form

\begin{gather*}
\FitMath{\int\Theta_cF' F\,e^{-\beta\sum_pa_p}\,dU
   =\sum_\gamma b_\gamma\int
       \overline{I_\gamma(F';Y)}I_\gamma(F;Y)\,d\lambda(Y),} \\[3pt]
\FitMath{b_\gamma\ge0,}\tag{AM5}
\end{gather*}

where $Y$ denotes the fixed boundary variables; in the link case $Y$
is absent. The functions $I_\gamma$ integrate over one half, with its
Boltzmann factor and the corresponding Gram feature. The boundary
Boltzmann factor is included in the positive measure $\lambda$.
Fubini's theorem follows from the coefficient summability above.
This proves the Cauchy--Schwarz inequality for the reflection form and
reflection positivity for all bounded half-space functions. Haar
invariance, orientation reversal, and invariance of the plaquette trace
also prove reflection invariance of the finite measure.

If $E$ is a bounded real function on one closed half, multiplying the
density by $e^{E+\theta_cE}$ preserves both properties: insert $Fe^E$
in (AM5). This observation will apply to folded perturbations.

\subsection*{A finite one-direction chessboard proof}

Fix parallel slab faces at heights $c+kH$, where $H\ge2$ is an integer.
Use a split torus with normal height period $2MH$, $M\ge1$. Its $2M$
slabs are indexed by $\mathbb Z/(2M)$. The group generated by reflection
in the slab faces acts simply transitively on the slabs: its
translations change the index by an even integer and its reflections
send $k$ to $2j-1-k$. Thus a function on any one slab has a well-defined
copy on every other slab obtained using this group. Boundary bonds may
belong to two closed slabs; this causes no problem, since the products
are functions of the same boundary variables and (AM5) permits those
variables in both half-algebras.

For positive bounded slab functions $f_0,\ldots,f_{2M-1}$ with strictly
positive infima, denote by $Z^{\rm dis}(f_j)$ the partition function
obtained by copying $f_j$ into every slab. We claim

\begin{gather*}
\FitMath{\int\prod_{j=0}^{2M-1}f_j\,e^{-\beta\sum_pa_p}\,dU
 \le\prod_{j=0}^{2M-1}Z^{\rm dis}(f_j)^{1/(2M)}.}\tag{AM6}
\end{gather*}

Here is a proof valid for every $M$, without a power-of-two restriction.
Transport the $2M$ functions to one reference slab and retain their
labels. For any filling $\sigma:\mathbb Z/(2M)\to\{0,\ldots,2M-1\}$,
let $Z(\sigma)$ be the partition function obtained by placing profile
$\sigma(k)$ in slab $k$. Put

\begin{gather*}
\FitMath{w_j=Z^{\rm dis}(f_j)^{1/(2M)},} \\[3pt]
\FitMath{T(\sigma)=Z(\sigma)\big/\prod_kw_{\sigma(k)}.}
\end{gather*}

All these numbers are positive, and every constant filling has $T=1$.
Reflection Cauchy--Schwarz at any pair of opposite slab faces gives

\begin{gather*}
\FitMath{T(\sigma)^2\le T(\sigma_+)T(\sigma_-),}\tag{AM7}
\end{gather*}

where $\sigma_+$ retains one semicircle of labels and reflects it into
the other, and $\sigma_-$ does the reverse. Indeed, the unnormalized
inequality is (AM5); the two denominators together contain each
original denominator factor twice.

There are finitely many fillings. If $\sigma$ maximizes $T$, (AM7)
forces both reflected fillings to maximize it as well. Take a cyclic
constant run of length $l$ in a maximizing filling and reflect at the
face at the start of that run. If $l\le M$, retaining the half beginning
there gives a maximizing filling with a run of length $2l$. If
$l\ge M$, retain a semicircle inside the run and reflect it; the result
is constant. Iterating therefore produces a constant maximizing
filling. The maximum is one. Applying this to the original filling
proves (AM6).

This argument uses reflection positivity only for finite tori. It does
not infer the positivity of an arbitrary DLR state from the positivity
of one adjacent-layer kernel.

\subsection*{A lower tangent for an arbitrary periodic DLR state}

Let $Q$ be a fixed finite-index period subgroup, let $A$ be a real
continuous local observable, and form the locally finite interaction
array $\mathbb A=(\tau_vA)_{v\in Q}$. Write
$\kappa=[\mathbb Z^4:Q]^{-1}$. Let $\Gamma_i\subset Q$ be rectangular
adapted period lattices with every side tending to infinity. On the
corresponding tori put

\begin{gather*}
\FitMath{Z_i(t\mathbb A)=\int
   \exp\left[t\sum_{v\in Q/\Gamma_i}(\tau_vA)^{\Gamma_i}
                    -\beta\sum_pa_p\right]dU.}
\end{gather*}

For every $Q$-invariant Wilson DLR state $\nu$ and every fixed real $t$,

\begin{gather*}
\FitMath{\frac1{|D_i|}\log\frac{Z_i(t\mathbb A)}{Z_i(0)}
 \ge t\kappa\nu(A)-o(1).}\tag{AM8}
\end{gather*}

The error may depend on $A,Q,\beta,t$, but is uniform in $\nu$.

For completeness, choose a fundamental parallelepiped $D_i$ and
integrate the positively oriented bonds anchored in it, using arbitrary
external bonds as boundary conditions. Let $W_i$ sum the source copies
whose entire support is interior to this bond set. Conditional Jensen
and the DLR equation give

\begin{gather*}
\FitMath{\log\frac{Z_{D_i}^{t}(\eta)}{Z_{D_i}^{0}(\eta)}
 \ge t\gamma_{D_i}W_i(\eta),} \\[3pt]
\FitMath{\int\gamma_{D_i}W_i\,d\nu=\nu(W_i).}
\end{gather*}

The open and torus Hamiltonians differ only in a boundary layer of
fixed thickness, determined by the range of $A$ and the plaquettes.
Each energy term is bounded; each bond is met by only a fixed finite
number of array terms. Thus their logarithmic partition ratios differ
by at most
$C_A(\beta+|t|\|A\|)|\partial_R D_i|$, uniformly in $\eta$. The number
of omitted array anchors is at most $C_A|\partial_R D_i|$. Since the
remaining $Q$-translates all have expectation $\nu(A)$, integrating the
Jensen bound proves (AM8), with an error bounded by a constant times
$|\partial_R D_i|/|D_i|$. This ratio tends to zero for the adapted
rectangular tori, including the diagonal ones. No differentiability of
infinite-volume pressure is used.

\subsection*{Fold domination with a fixed common period}

Fix one admissible mirror $(n,c)$. Choose an even positive integer $h$
and set

\begin{gather*}
\FitMath{w=hn\in P,} \\[3pt]
\FitMath{H=n\cdot w=hq,} \\[3pt]
\FitMath{P_\parallel=P\cap n^\perp,} \\[3pt]
\FitMath{Q=P_\parallel\oplus\mathbb Zw,} \\[3pt]
\FitMath{Q_f=P_\parallel\oplus\mathbb Z(2w).}\tag{AM9}
\end{gather*}

Both $Q$ and $Q_f$ are finite-index subgroups of $P$ and
$[Q:Q_f]=2$. The tangent lattice $P_\parallel$ has an orthogonal
integer basis: for a coordinate normal use the three other vectors
$2e_j$; for $e_1\pm e_2$ use $2(e_1\mp e_2),2e_3,2e_4$, with the
corresponding signed permutations for the other cases. Thus

\begin{gather*}
\FitMath{\Gamma_{M,J}=J P_\parallel\oplus\mathbb Z(2Mw)}\tag{AM10}
\end{gather*}

is an adapted rectangular split torus. Its shortest period tends to
infinity as $M,J\to\infty$.

Assume for the moment that $A$ is supported in the slab
$c\le n\cdot x\le c+H$. The unfolded array contains all its translates
by $Q$. The folded array contains translates of $A$ in the even slabs
and the reflected copies $\theta_{c+H}A$ in the odd slabs, with period
$Q_f$. It is invariant under reflection at every face $c+kH$.
In (AM6) take

\begin{gather*}
\FitMath{f_k=\exp\left[t\sum_{v\in P_\parallel/JP_\parallel}
                    (\tau_{v+kw}A)^{\Gamma_{M,J}}\right],} \\[3pt]
\FitMath{t>0.}
\end{gather*}

Their product is the unfolded perturbation. Their disseminated
partition functions are identical: normal translation by $kw$ carries
one dissemination to another, permutes the slab faces, and preserves
the unperturbed Wilson measure. Each is the folded partition function.
Hence

\begin{gather*}
\FitMath{Z_{M,J}(t\mathbb A^{\rm u})
 \le Z_{M,J}(t\mathbb A^{\rm f}).}\tag{AM11}
\end{gather*}


Let $\mu_{M,J,t}$ be the normalized folded measure, and define

\begin{gather*}
\FitMath{p_{M,J}(t)=|D_{M,J}|^{-1}
       \log[Z_{M,J}(t\mathbb A^{\rm f})/Z_{M,J}(0)].}
\end{gather*}

It is convex, $p_{M,J}(0)=0$, and therefore
$p_{M,J}(t)\le t p'_{M,J}(t)$ for $t>0$. The folded measure is invariant
under $Q_f$ and under $\theta_{c+H}$. All its source terms consequently
have the same expectation as $A$. There are $\kappa|D_{M,J}|$ terms,
where $\kappa=[\mathbb Z^4:Q]^{-1}$, so

\begin{gather*}
\FitMath{p'_{M,J}(t)=\kappa\mu_{M,J,t}(A).}
\end{gather*}

Combining this identity, (AM8), and (AM11) gives, for every
$\nu\in\mathcal G_\beta^P$,

\begin{gather*}
\FitMath{t\kappa\nu(A)-o(1)
 \le p_{M,J}(t)\le t\kappa\mu_{M,J,t}(A).}\tag{AM12}
\end{gather*}


The finite folded measure is invariant and reflection positive at
every face $c+kH$. Indeed, its perturbation is $E+\theta_{c+kH}E$,
where $E$ is the sum of its copies lying in one closed half of the
torus. Copies on a boundary can be assigned half weight to each side.
The observation following (AM5) applies.

First fix $t>0$ and let $M,J\to\infty$. Compactness gives a limit
$\rho_t$. It is a DLR state for the folded, finite-range perturbed
interaction, is $Q_f$-periodic, is invariant and reflection positive at
all its slab faces, and satisfies $\nu(A)\le\rho_t(A)$. Here DLR passage
is local and exact: any fixed collection of bonds, its interaction
neighbourhood, and the support of a cylinder test embed injectively in
all sufficiently large tori. The finite torus conditional equation is
then exactly the required local DLR equation on that test. Its kernel
is continuous, so the equation passes to a weak limit. A finite test
on one half-space eventually lies in a single torus half; hence its
reflection inequality passes to the limit as well.

Next let $t\downarrow0$ along a subsequence for which $\rho_t$ converges
to $\rho_0$. On any fixed finite set, only finitely many source copies
meet that set. Their bounded total contribution is $O(t)$ uniformly in
the exterior configuration. The perturbed conditional densities thus
converge uniformly to the Wilson conditional density. It follows that
$\rho_0$ is a Wilson DLR state. It retains $Q_f$ periodicity, reflection
invariance and positivity at the slab faces, and
$\nu(A)\le\rho_0(A)$.

Now apply the previously proved link-symmetry theorem to $\rho_0$.
Its finite-index period $Q_f$ is enough for that theorem; the result
upgrades its period to $P$. We have proved the following domination:

\begin{gather*}
\FitMath{\begin{gathered}
 A\text{ supported in }[c,c+H],\quad \nu\in\mathcal G_\beta^P
 \quad\Longrightarrow\quad
 \exists\rho_0\in\mathcal G_\beta^P:
 \\
 \rho_0\text{ is invariant and positive at }(n,c),
 \qquad \nu(A)\le\rho_0(A).
 \end{gathered}}\tag{AM13}
\end{gather*}

The order of limits is essential: the surface error in (AM12) is
divided by $t\kappa$, so volume tends to infinity before $t$ tends to
zero.

\subsection*{Separation closes the full half-space statement}

Let $K=\mathcal G_\beta^P$ and let $K_c\subset K$ consist of the states
invariant under $\theta_c$ and positive on every continuous cylinder
function in the positive half-space. The set $K$ is compact and convex:
the DLR kernels preserve continuous cylinder functions, and all its
defining equations and period equations are closed affine conditions.
The set $K_c$ is also closed and convex. It is nonempty, for example
by the folded construction with $A=0$.

For any real continuous local $A$, choose $v\in P$ so that $\tau_vA$
is supported strictly above the fixed mirror $c$. Then choose the even
$h$ large enough that this support lies below $c+H$. Such a $v$ exists
because $n\cdot(2ke_j)=2k$ for a coordinate $j$ with $n_j=1$.
Apply (AM13) to $\tau_vA$. Both $\nu$ and the resulting $\rho_0$ are
$P$-invariant, so

\begin{gather*}
\FitMath{\nu(A)=\nu(\tau_vA)\le\rho_0(\tau_vA)=\rho_0(A),} \\[3pt]
\FitMath{\rho_0\in K_c.}\tag{AM14}
\end{gather*}

If some $\nu\in K$ were outside $K_c$, separation of a point from a
compact convex set of measures would supply a real continuous function
$A$ with $\nu(A)>\sup_{\rho\in K_c}\rho(A)$. Continuous cylinder
functions are uniformly dense in $C(\Omega)$, so $A$ could be chosen
local while retaining a strict inequality. This contradicts (AM14).
Therefore $K=K_c$.

This step explains why the common period cannot be omitted. A state
produced by folding initially has only the normal period $2w$. Such a
large period need not move a test straddling the mirror into one slab.
The link-symmetry theorem supplies $2\mathbb Z^4$ for both states in
(AM14), independently of the observable and slab width. No claim that
periodic DLR states are torus-selected is needed.

The mirror was arbitrary, so every periodic state is invariant and
positive at all the stated mirrors on continuous local half-space
functions. Finally, on the compact metrizable product of the positive
half-space bonds, continuous cylinder functions are dense in
$L^2(\nu)$. For bounded Borel $F$ choose such $F_j$ with
$\|F_j-F\|_2\to0$. Reflection invariance makes $\Theta_c$ an isometry
of $L^2(\nu)$, and Cauchy--Schwarz gives

\begin{gather*}
\FitMath{\left|\nu(\Theta_cF_jF_j)-\nu(\Theta_cFF)\right|
 \le\|F_j-F\|_2(\|F_j\|_2+\|F\|_2)\longrightarrow0.}
\end{gather*}

This proves (AM1) on the complete positive-half algebra. Replacing a
negative-half function by its reflected conjugate proves the other
half. In particular, arbitrary finite stacks and their bounded Borel
tests are covered.\hfill$\square$

As a further consequence, every periodic state is invariant under all
unit lattice translations: in direction $e_\mu$ the product of the
site reflection at $0$ and the link reflection at $1/2$ is translation
by $e_\mu$. It is also invariant under the hypercubic group. The
coordinate reflections at zero generate coordinate sign changes, and
the diagonal reflections with normal $e_\mu-e_\lambda$ interchange
coordinates. These transformations generate all signed permutations.

\subsection*{Scope of the implication}

Positive semidefinite contractions of an adjacent-layer kernel can
produce a one-layer positivity test, just as skew contractions produce
an antisymmetric invariance test. That observation alone supplies no
identity reducing arbitrary odd or even stacks to a single layer.
The proof above uses a different route: exact finite-torus positivity
on the full half-algebra, a finite chessboard inequality, a lower
tangent for each arbitrary periodic DLR state, and separation after
the common period is established. Thus the all-stack and diagonal
conclusions use no layer-elimination positivity assumption. The result is at a fixed lattice spacing and coupling; it makes
no continuum-limit or interacting renormalization-group assertion.

% END MODULE chapters/all_mirror_rp_audit_v3.tex


% BEGIN MODULE chapters/phase_plateau_audit_v3.tex
\section{A uniform lattice plateau estimate from reflection positivity}
\label{ch:plateau-proof}

The estimate in this section is a consequence of translation invariance and
reflection positivity in coordinate and diagonal directions. No infrared bound,
mass gap, or differentiability of a pressure is assumed. The Gibbs equations
are not needed in the proof. Their role in the Wilson application is to place
every periodic Wilson Gibbs state in the reflection-positive class by
Theorem~\ref{thm:amrp-all-periodic}, proved independently in
Section~\ref{ch:all-mirror-periodic}. The state here is a measure on the
infinite lattice, rather than a finite-torus law lifted periodically.

By dilation it suffices initially to take lattice spacing one. Let $\nu$ be
a translation-invariant probability measure on lattice configurations. A local
function means a bounded Borel function of finitely many bonds. The algebra
$\mathfrak A(S)$ contains the functions of bonds with both endpoints in $S$.
Assume reflection positivity on the closed side of every coordinate mirror
$x_\mu=c$, $c\in\frac12\mathbb Z$, and every diagonal mirror
$x_\mu\pm x_\lambda=c$, $c\in\mathbb Z$. These are sixteen families.
Reflections act linearly on functions by pullback, without complex conjugation.
Write $\theta_\mu$ for reflection in $x_\mu=0$.

For $r\ge1$ an integer put $Q_r=[-r,r]^4$. For $A\in\mathfrak A(Q_r)$
define

\begin{gather*}
\FitMath{\mathcal N_{\nu,r}(A)
 =\max_{\mu,\gamma}
   \nu\bigl(\overline{\theta_\mu\gamma A}
                     \tau_{2re_\mu}\gamma A\bigr),}\tag{PL1}
\end{gather*}

where $\gamma$ ranges over the sixteen coordinate sign changes. The terms
in this maximum are nonnegative squared reflection norms. Let $\mathcal I$
be the completed invariant $\sigma$-field for all lattice translations and set

\begin{gather*}
\FitMath{\operatorname{Pl}_\nu(A,B)
   =\nu\bigl(\overline{E_\nu[A\mid\mathcal I]}
                           E_\nu[B\mid\mathcal I]\bigr).}\tag{PL2}
\end{gather*}


\begin{theorem}[Uniform plateau bound]
\label{thm:lattice-plateau}
There is a numerical constant $C$ such that, for the measures just specified,
$A,B\in\mathfrak A(Q_r)$, and $y\in\mathbb Z^4$ with
$n=|y|_\infty-2r\ge1$,

\begin{gather*}
\FitMath{\left|\nu(\overline A\tau_yB)-\operatorname{Pl}_\nu(A,B)\right|
 \le C\left[(r/n)^2+(r/n)^3\right]
          \mathcal N_{\nu,r}(A)^{1/2}
          \mathcal N_{\nu,r}(B)^{1/2}.}\tag{PL3}
\end{gather*}

The constant is independent of the measure, the support size, and any
coupling. At lattice spacing $a$, the same statement holds with
$Q_r=a[-r,r]^4$ and $n=|y|_\infty/a-2r$.
\end{theorem}

We use the spectral theorem for commuting bounded normal operators and the
mean ergodic theorem for a unitary operator. All estimates supplying the
power in (PL3), including the momentum-density estimate, are proved below.

\subsection{Reflection Hilbert spaces and their spectral carriers}

Reflection positivity and translation invariance first imply invariance under
each of the indicated reflections. On one closed side, positivity makes
$\langle F,G\rangle_\theta=\nu(\overline{\theta F}G)$ Hermitian.
Testing it against $1$ shows $\nu(F)=\nu(\theta F)$ for functions supported
on that side. Any local function can be translated into the side by a
lattice vector normal to the mirror. Reflecting reverses that translation;
translation invariance therefore proves the equality for the original
function. The reflections consequently act unitarily on $L^2(\nu)$.

In the coordinate picture $\mu$, complete the quotient of local functions on
$x_\mu\ge0$ for the positive semidefinite form

\begin{gather*}
\FitMath{\langle F,G\rangle_\mu=\nu(\overline{\theta_\mu F}G).}
\end{gather*}

Denote the class of $F$ by $[F]_\mu$. Its norm is at most
$\|F\|_\infty$. Translation by $e_\mu$ induces a positive self-adjoint
contraction $T_\mu$, and tangential coordinate translations induce commuting
unitaries $U_{\mu,\lambda}$, $\lambda\ne\mu$.

Here are the details of this elementary reconstruction. The normal translation
is symmetric on the pre-Hilbert space by translation invariance and
$\theta_\mu\tau_{e_\mu}=\tau_{-e_\mu}\theta_\mu$. Its quadratic form is
nonnegative by reflection positivity at the mirror $x_\mu=-1/2$.
Repeated Cauchy--Schwarz gives

\begin{gather*}
\FitMath{\|T_\mu[F]\|\le
       \|[F]\|^{1-2^{-k}}\|T_\mu^{2^k}[F]\|^{2^{-k}}.}
\end{gather*}

The last norm is at most $\|F\|_\infty$. Letting $k\to\infty$ proves
contraction and preservation of null vectors. The bounded symmetric extension
is self-adjoint and positive. Tangential translations preserve the form and
have inverses of the same kind, so are unitaries. Their commutations on local
functions pass to the completion. In particular, for integers $s\ge0$ and tangential
$w\in\mathbb Z^4$,

\begin{gather*}
\FitMath{\langle[F]_\mu,T_\mu^s U_\mu^w[G]_\mu\rangle
       =\nu(\overline{\theta_\mu F}\tau_{se_\mu+w}G).}\tag{PL4}
\end{gather*}


There is a joint spectral measure with variables
$(E,\eta)\in[0,\infty]\times(-\pi,\pi]^3$ for which
$T_\mu=e^{-E}$ and $U_{\mu,\lambda}=e^{i\eta_\lambda}$.
The value $E=\infty$ represents a possible kernel of $T_\mu$.
The following carrier bound is essential:

\begin{gather*}
\FitMath{\cosh E+\cos\eta_\lambda\ge2} \\[3pt]
\FitMath{(\lambda\ne\mu),} \\[3pt]
\FitMath{E\ge\kappa|\eta|_\infty,} \\[3pt]
\FitMath{\kappa=\frac2\pi\log(1+\sqrt2).}\tag{PL5}
\end{gather*}

We now derive it from the diagonal mirrors.

\paragraph{A polynomial-measure lemma.}
Suppose a finite positive measure $m$ on the closed unit disc satisfies

\begin{gather*}
\FitMath{\int|p(z)|^2\,dm(z)
   \le C_0\sup_{[0,1]}|p|\sup_{|z|=1}|p|}\tag{PL6}
\end{gather*}

for every polynomial. Then it gives zero mass to
$\{|1+z|^2<4|z|\}$.

To prove this without invoking an extremal-function theorem, set

\begin{gather*}
\FitMath{L(\zeta)=\operatorname{Log}\frac{1+\zeta}{1-\zeta},} \\[3pt]
\FitMath{f_s(z)=\frac{\cos(sL(\sqrt z))}{\cosh(\pi s/2)},} \\[3pt]
\FitMath{s>0.}
\end{gather*}

The Cayley transform has positive real part in the disc. Thus $L$ is
holomorphic, odd, and satisfies $|\operatorname{Im}L|<\pi/2$.
The cosine is even in $\sqrt z$, so $f_s$ is a single-valued holomorphic
function of $z$ in the disc. Put
$\varphi(z)=|\operatorname{Im}L(\sqrt z)|$.
The identity $|\cos(u+iv)|^2=\cos^2u+\sinh^2v$ gives

\begin{gather*}
\FitMath{|f_s|\le1,} \\[3pt]
\FitMath{\sup_{[0,1)}|f_s|\le\cosh(\pi s/2)^{-1},} \\[3pt]
\FitMath{|f_s(z)|\ge\frac{\sinh(s\varphi(z))}{\cosh(\pi s/2)}.}
\end{gather*}

Apply (PL6) to Taylor approximants of $f_s(qz)$, $0<q<1$,
which converge uniformly on the closed disc. The resulting inequality,
followed by Fatou as $q\uparrow1$, is

\begin{gather*}
\FitMath{\int\sinh^2(s\widetilde\varphi(z))\,dm(z)
       \le C_0\cosh(\pi s/2).}
\end{gather*}

Here $\widetilde\varphi=\varphi$ inside the disc, equals $\pi/2$ on
the boundary except at $1$, and may be set to zero at $1$.
Indeed the radial limits of $\varphi$ at the other boundary points
equal $\pi/2$. For each $\epsilon>0$ this inequality forces
$m(\{\widetilde\varphi\ge\pi/4+\epsilon\})=0$ by letting
$s\to\infty$. Finally, if $z=\zeta^2$ inside the disc,

\begin{gather*}
\FitMath{\tan\varphi(z)=\frac{2|\operatorname{Im}\zeta|}{1-|\zeta|^2},} \\[3pt]
\FitMath{|1+\zeta^2|^2-4|\zeta|^2
       =(1-|\zeta|^2)^2-4(\operatorname{Im}\zeta)^2.}
\end{gather*}

Thus $\varphi>\pi/4$ is precisely the excluded set. The same
assertion holds on the boundary, where only $z=1$ is not excluded.
This proves the lemma.

Fix $\lambda\ne\mu$ and a local $A$ on $x_\mu\ge0$. Choose an
integer $R\ge1$ enclosing its support in
$0\le x_\mu\le R$, $|x_\lambda|\le R$. Let $\theta'$ be reflection
in $x_\mu+x_\lambda=R$, put
$B=\theta'\theta_\mu A$, $A_0=\tau_{2Re_\mu}A$, and
$d_\pm=e_\mu\pm e_\lambda$. Both $B$ and $A_0$ lie on the positive
side of this diagonal mirror. Reconstruction at that mirror gives a
positive contraction $T_+$ for $d_+$ and a unitary $V_-$ for $d_-$.
The preceding reconstruction proof applies: positivity of $T_+$ uses
the admissible diagonal mirror one height unit behind, and its
contraction uses the bounded local representative. For integers $k,l\ge0$,

\begin{gather*}
\FitMath{\langle[A]_\mu,T_\mu^{2R+k+l}U_{\mu,\lambda}^{k-l}[A]_\mu\rangle
       =\langle[B]',T_+^kV_-^l[A_0]'\rangle.}\tag{PL7}
\end{gather*}

Both sides are the same lattice correlation: their displacement is
$2Re_\mu+kd_++ld_-$. No analytic continuation is used here.

Push the spectral measure of $[A]_\mu$, multiplied by $t^{2R}$,
forward by $z=t\zeta$, where $t=e^{-E}$ and
$\zeta=e^{i\eta_\lambda}$. Expanding a polynomial and its conjugate
in (PL7) gives

\begin{gather*}
\FitMath{\int|p|^2\,dm
  =\langle[B]',p(T_+)\overline p(V_-)[A_0]'\rangle,} \\[3pt]
\FitMath{\overline p(z)=\sum_k\overline{p_k}z^k.}
\end{gather*}

The spectral theorem bounds this by the right side of (PL6), with
$C_0=\|[B]'\|\|[A_0]'\|$. The lemma excludes
$|1+t\zeta|^2<4t$. Since $t^{2R}>0$ for $t>0$, the joint spectral
measure of $[A]_\mu$ itself vanishes there. The local classes are
dense, so the operator spectral measure vanishes there. At $t=0$
the equivalent operator inequality is automatically satisfied.
We have proved

\begin{gather*}
\FitMath{(1-T_\mu)^2\ge T_\mu(2-U_{\mu,\lambda}-U_{\mu,\lambda}^*).}
\end{gather*}

For $0<t\le1$ this is equivalent to the first assertion of (PL5).
The curve is $E\ge2\operatorname{arcsinh}|\sin(\eta_\lambda/2)|$.
Concavity of $x\mapsto\operatorname{arcsinh}\sin x$ on $[0,\pi/2]$
puts it above its endpoint chord, giving the constant $\kappa$.
Intersect the three carriers to obtain (PL5).

\subsection{The zero-energy plateau is the invariant conditional field}

Write $\Pi_\mu=1_{\{E=0\}}$; spectral dominated convergence gives
$T_\mu^s\to\Pi_\mu$ strongly as $s\to\infty$. By (PL5), every
tangential unitary equals the identity on this zero-energy space.
Let $\mathcal I_\mu$ denote the invariant field for the single
translation $\tau_{e_\mu}$. The mean ergodic theorem on $L^2(\nu)$
and (PL4) give

\begin{gather*}
\FitMath{\langle[F]_\mu,\Pi_\mu[G]_\mu\rangle
       =\nu(\overline{\theta_\mu F}E_\nu[G\mid\mathcal I_\mu]).}\tag{PL8}
\end{gather*}


Reflection acts as a unitary involution on
$L^2(\mathcal I_\mu,\nu)$. That involution is positive.
Indeed, approximate an invariant function $J$ in $L^2$ by
bounded local $G_k$ and translate each approximant into
$x_\mu\ge0$. The error is unchanged because $J$ is invariant.
Reflection positivity then passes to the $L^2$ limit and gives
$\langle\theta_\mu J,J\rangle\ge0$. A positive self-adjoint
involution has only the eigenvalue $+1$, so reflection fixes $J$.
Consequently (PL8) equals

\begin{gather*}
\FitMath{\nu\bigl(\overline{E_\nu[F\mid\mathcal I_\mu]}
                       E_\nu[G\mid\mathcal I_\mu]\bigr).}
\end{gather*}

The map $\Pi_\mu[F]_\mu\mapsto E_\nu[F\mid\mathcal I_\mu]$
is therefore an isometry onto a dense subspace of
$L^2(\mathcal I_\mu,\nu)$. Density follows by the same shifted
local approximations. It intertwines the tangential translations.
Since those translations are the identity on $\Pi_\mu$, every
$\mathcal I_\mu$-measurable function is invariant under all four
coordinate translations. Hence $\mathcal I_\mu=\mathcal I$ modulo
$\nu$, and

\begin{gather*}
\FitMath{\langle[F]_\mu,\Pi_\mu[G]_\mu\rangle
       =\nu\bigl(\overline{E_\nu[F\mid\mathcal I]}
                           E_\nu[G\mid\mathcal I]\bigr).}\tag{PL9}
\end{gather*}

This also proves that $E_\nu[\theta_\mu F\mid\mathcal I]
=E_\nu[F\mid\mathcal I]$. The possible $T_\mu$ kernel has
energy infinity and does not affect any of these zero-energy identities.

\subsection{A positive time filter and a bounded momentum density}

For $A\in\mathfrak A(Q_r)$ define its cube vector
$\Phi_A^\mu=[\tau_{re_\mu}A]_\mu$. Translation invariance shows
$\|\Phi_A^\mu\|^2=\nu(\overline{\theta_\mu A}\tau_{2re_\mu}A)$.
Thus (PL1) is the maximum of the corresponding squared norms of
$\gamma A$. In particular $\mathcal N^{1/2}$ is a seminorm,
is invariant under the sign-change group, and is at most the
supremum norm.

Let $b$ be the nonnegative periodic triangle function supported on
$[-\pi/4,\pi/4]$, normalized to have mean one. Its Fourier
coefficients are

\begin{gather*}
\FitMath{b_m=\left(\frac{\sin(m\pi/8)}{m\pi/8}\right)^2,} \\[3pt]
\FitMath{b_0=1.}
\end{gather*}

This follows by writing the triangle as the convolution of two
interval indicators and normalizing their zeroth coefficients.
Let $z_m$ be the even period-four sequence $(0,1,-2,1)$ and set

\begin{gather*}
\FitMath{\widehat q(u)=\sum_{m\in\mathbb Z}z_mb_m e^{imu},} \\[3pt]
\FitMath{\Psi(q)=2\sum_{m\ge1}z_mb_mq^{m-1},} \\[3pt]
\FitMath{0\le q\le1.}\tag{PL10}
\end{gather*}

Both series converge absolutely. The identity
$z_m=\frac12i^m-(-1)^m+\frac12(-i)^m$ gives

\begin{gather*}
\FitMath{\widehat q(u)=\tfrac12b(u+\pi/2)-b(u+\pi)
                               +\tfrac12b(u-\pi/2).}
\end{gather*}

In particular $\widehat q$ vanishes whenever
$\operatorname{dist}(u,2\pi\mathbb Z)\le\pi/4$.
The crucial energy property is

\begin{gather*}
\FitMath{0\le\Psi\le C,} \\[3pt]
\FitMath{\Psi(1)=0,} \\[3pt]
\FitMath{\Psi(e^{-x})\ge c\tanh x} \\[3pt]
\FitMath{(x>0).}\tag{PL11}
\end{gather*}


To verify its nontrivial lower bound, use the Poisson kernel

\begin{gather*}
\FitMath{P_x(u)=\sum_me^{-|m|x}e^{imu}=\frac{\sinh x}{\cosh x-\cos u}.}
\end{gather*}

Since $z_0=0$, the Poisson integral yields

\begin{gather*}
\FitMath{\Psi(e^{-x})=\frac{e^x}{2\pi}\int_{|v|\le\pi/4}b(v)
  [\tfrac12P_x(v-\pi/2)+\tfrac12P_x(v+\pi/2)-P_x(v-\pi)]\,dv.}
\end{gather*}

Writing $C_x=\cosh x$, the bracket is

\begin{gather*}
\FitMath{\sinh x\left[
  \frac{C_x}{C_x^2-\sin^2v}-\frac1{C_x+\cos v}\right]
 \ge\frac{\tanh x}{\sqrt2(C_x+2^{-1/2})}.}
\end{gather*}

For this inequality first replace the first fraction by $1/C_x$,
then use $\cos v\ge2^{-1/2}$ and monotonicity of
$q/(C_x+q)$. The factor $e^x/(C_x+2^{-1/2})$ has a positive
lower bound, and $b$ has mean one. This proves the lower bound.
The upper bound follows from absolute convergence of (PL10), and
$\Psi(1)=\widehat q(0)=0$. These statements extend continuously
to $q=0$, where $\Psi(0)=2b_1>0$.

First assume $\theta_\mu A=\epsilon A$, $\epsilon\in\{1,-1\}$,
and put $\Phi=\Phi_A^\mu$. For $s\in\mathbb Z$ and tangential
$w\in\mathbb Z^3$, define

\begin{gather*}
\FitMath{g_A(s,w)=\epsilon\nu(\overline A\tau_{se_\mu+w}A).}
\end{gather*}

For $s\ge2r$, the cube-crossing identity from (PL4) gives

\begin{gather*}
\FitMath{g_A(s,w)=\langle\Phi,T_\mu^{s-2r}U_\mu^w\Phi\rangle,} \\[3pt]
\FitMath{|g_A(s,w)|\le\mathcal N(A).}\tag{PL12}
\end{gather*}

Reflection in $x_\mu=0$ and translation invariance give
$g_A(-s,w)=g_A(s,w)$ and
$g_A(s,-w)=\overline{g_A(s,w)}$. Define

\begin{gather*}
\FitMath{K_A(w)=\sum_{m\in\mathbb Z}z_mb_mg_A(2rm,w).}
\end{gather*}

The term at $m=0$ vanishes. Every other time separation has
absolute value at least $2r$, so (PL12) bounds this sum by
$C\mathcal N(A)$, uniformly in $w$. By the same identities,

\begin{gather*}
\FitMath{K_A(w)=\langle\Phi,\Psi(T_\mu^{2r})U_\mu^w\Phi\rangle
       =\int e^{i\eta\cdot w}\,dM_A(\eta),}\tag{PL13}
\end{gather*}

where $M_A$ is the momentum pushforward of the positive measure

\begin{gather*}
\FitMath{\Psi(e^{-2rE})\,d\langle\Phi,\mathsf E_\mu(E,\eta)\Phi\rangle.}
\end{gather*}

Thus $M_A$ is positive; the signs in $z_m$ do not compromise
positivity because (PL11) controls their spectral sum.

For $|w|_\infty\ge2r$, choose $\lambda\ne\mu$ at a maximal
coordinate. Replace $w$ by $-w$ if necessary, which only conjugates
$K_A(w)$, and suppose $w_\lambda\ge2r$. In the coordinate picture
$\lambda$, for every integer $s$ the same correlation is

\begin{gather*}
\FitMath{g_A(s,w)=\epsilon\left\langle
 \Phi_{\theta_\lambda A}^\lambda,
 T_\lambda^{w_\lambda-2r}U_{\lambda,\mu}^s
       \prod_{\lambda'\ne\lambda,\mu}U_{\lambda,\lambda'}^{w_{\lambda'}}
 \Phi_A^\lambda\right\rangle.}
\end{gather*}

Sum this identity at $s=2rm$ against $z_mb_m$. In this picture
the filter becomes the bounded spectral function
$\widehat q(2r\eta'_\mu)$ of a unitary translation. Whenever it
is nonzero, $|\eta'_\mu|>\pi/(8r)$, since the gap around zero
modulo $2\pi$ includes the ordinary interval about zero. Aliases
at other multiples of $2\pi$ do not invalidate this implication.
The carrier (PL5) consequently gives
$E'\ge\alpha=\kappa\pi/(8r)$ on the filter support. Both cube
vectors have norm at most $\mathcal N(A)^{1/2}$. Hence

\begin{gather*}
\FitMath{|K_A(w)|\le C\mathcal N(A)
       e^{-\alpha(|w|_\infty-2r)}} \\[3pt]
\FitMath{(|w|_\infty\ge2r).}\tag{PL14}
\end{gather*}


There are $O(r^3)$ points with $|w|_\infty<2r$ and $O(k^2)$
points on the shell $|w|_\infty=k$. Since $\alpha$ is a fixed
constant times $r^{-1}$, (PL14) and the previous bound give

\begin{gather*}
\FitMath{\sum_w|K_A(w)|
 \le C\mathcal N(A)
       \left[r^3+\sum_{j\ge0}(r^2+j^2)e^{-\alpha j}\right]
 \le Cr^3\mathcal N(A).}
\end{gather*}

Therefore $h_A(\eta)=(2\pi)^{-3}\sum_wK_A(w)e^{-i\eta\cdot w}$
is continuous and bounded by $Cr^3\mathcal N(A)$.
It has the same Fourier coefficients as $M_A$. Uniqueness of a
finite measure from its Fourier coefficients follows by density
of trigonometric polynomials in the continuous functions on the
torus. Consequently

\begin{gather*}
\FitMath{M_A=h_A(\eta)\,d\eta,} \\[3pt]
\FitMath{0\le h_A\le Cr^3\mathcal N(A).}\tag{PL15}
\end{gather*}

This is a derived bound for an energy-weighted momentum measure.
It is not an assumed bound for the unweighted spectral measure.

\subsection{Removing the filter and proving the plateau estimate}

Let $m_\Phi=\langle\Phi,\mathsf E_\mu(\cdot)\Phi\rangle$.
Since $M_A$ has a density, it has no atom at $\eta=0$.
By (PL11), its weight is strictly positive at every finite
$E>0$. Thus $m_\Phi$ gives zero mass to
$\{0<E<\infty,\eta=0\}$. The energy-zero contribution is precisely
the plateau and was not removed by this inference: its weight is zero.
The point $E=\infty$ contributes zero to $T_\mu^n$ for $n\ge1$.

For $E>0$ the function $e^{-nE}/\tanh(2rE)$ decreases in $E$.
Using (PL11), then (PL5), and then (PL15), gives
\begin{align*}
 \langle\Phi,T_\mu^n(1-\Pi_\mu)\Phi\rangle
 &\le C\int_{\mathbb T^3}
       \frac{e^{-n\kappa|\eta|_\infty}}
            {\tanh(2r\kappa|\eta|_\infty)}h_A(\eta)\,d\eta\\
 &\le Cr^3\mathcal N(A)\int_0^\infty
       u^2e^{-n\kappa u}\left(1+\frac1{2r\kappa u}\right)du\\
 &\le C\mathcal N(A)\bigl[(r/n)^3+(r/n)^2\bigr].
 \tag{PL16}
\end{align*}
In the second line we used $1/\tanh v\le1+1/v$ and the
three-dimensional cube shell volume $24u^2du$. The apparent singularity
at $u=0$ is integrable. The pushforward step can include $E=\infty$
and other unused spectral mass, since the integrand is nonnegative;
it does not add an atom at $\eta=0$ by (PL15).

For an arbitrary complex $A$, write
$A=A_++A_-$ with $A_\pm=(A\pm\theta_\mu A)/2$.
The seminorm property and sign-change invariance imply
$\mathcal N(A_\pm)\le\mathcal N(A)$.
The operator $D_n=T_\mu^n(1-\Pi_\mu)$ is positive, so its
quadratic form at a sum is at most twice the sum of its two
quadratic forms. Applying (PL16) to $A_+$ and $A_-$ proves
(PL16), with an enlarged numerical constant, for every $A$.

Finally suppose $y_\mu=|y|_\infty>0$. The cube-crossing identity is

\begin{gather*}
\FitMath{\nu(\overline A\tau_yB)=
   \langle\Phi_{\theta_\mu A}^\mu,
        T_\mu^n U_\mu^{y_\perp}\Phi_B^\mu\rangle.}
\end{gather*}

The $\Pi_\mu$ term is (PL2) by (PL9), conditional invariance under
translations and reflection, and $U_\mu^{y_\perp}\Pi_\mu=\Pi_\mu$.
For the remaining term, Cauchy--Schwarz in the positive form of
$D_n$ bounds the modulus by the geometric mean of the two quantities
in (PL16). The tangential unitary commutes with $D_n$, so does not
change that form. Also
$\mathcal N(\theta_\mu A)=\mathcal N(A)$. This proves (PL3).
For a negative maximal coordinate use
$\nu(\overline A\tau_{-y}B)
=\overline{\nu(\overline B\tau_yA)}$ and conjugate symmetry of
the plateau. Dilation restores spacing $a$ and replaces distances
by their numbers of lattice steps, proving the stated scaling.
\hfill$\square$

For an ergodic state the invariant conditional expectations are constants,
so (PL3) is clustering. For a mixture its generally nonzero invariant plateau
is retained. When local continuum representatives have large supremum norms,
the reflection sizes in (PL1) can still be bounded by separated reflected
moments. This is the reason (PL3), rather than a supremum-norm substitute,
is the useful estimate for a continuum phase construction. No spectral gap
above the invariant subspace has been proved: the filter weight vanishes
linearly near zero energy, and removing it is exactly what produces the
inverse-square term.

\begin{corollary}[Closedness of the ergodic reflection-positive class]
At fixed spacing, the translation-invariant measures satisfying these
mirror-positivity hypotheses form a weakly closed set, and so does its
subset of ergodic measures.
\end{corollary}
\begin{proof}
Translation invariance and each positivity inequality on continuous
local functions pass to a weak limit. Reflection invariance passes
as well, and approximation in $L^2$ extends positivity to bounded
Borel local functions. Thus the full reflection-positive class is
closed.

Suppose now that $\nu_j$ are ergodic in this class and converge weakly
to $\nu$. For fixed continuous local $A,B$ in a common cube,
their plateaux under $\nu_j$ are
$\overline{\nu_j(A)}\nu_j(B)$. Passing to the limit in (PL3) gives

\begin{gather*}
\FitMath{|\nu(\overline A\tau_yB)-\overline{\nu(A)}\nu(B)|
 \le C[(r/n)^2+(r/n)^3]\|A\|_\infty\|B\|_\infty.}
\end{gather*}

In particular the correlation converges to the product along
$y=ke_\mu$, $k\to\infty$. Its Cesaro means, with $A=B$, identify
the squared $L^2$ norm of $E_\nu[A\mid\mathcal I_\mu]$ by the mean
ergodic theorem. That norm is $|\nu(A)|^2$, so the conditional
expectation is constant. Continuous local functions are dense in
$L^2(\nu)$; hence $\mathcal I_\mu$, and therefore $\mathcal I$, is
trivial. This is ergodicity of $\nu$.
\end{proof}

% END MODULE chapters/phase_plateau_audit_v3.tex


% BEGIN MODULE chapters/dependency_catalogue.tex
\section{Supporting analytic results and proof dependencies}
\label{ch:dependency-catalogue}
\subsection*{Supporting results and their use in the proof}

This catalogue fixes the estimates and their domains used by the
construction. The labels identify proof tasks. The regular Gaussian
matrices and response, the hereditary substitution room, the weighted
history sum and the finite ordinary-source induction are developed below
for the actual construction. The interacting extraction then uses its
clean conditional law and the uniform finite-reference comparison.
The one-loop calculation is developed in
Section~\ref{ym4:new-one-loop-reduction}.

The main scalar route uses D1 in its actual rooted geometry, G1 and
F2--F3. The proof outlines supplying these inputs are in
Sections~\ref{sec:substitution-room},
\ref{sec:actual-history-closure},
\ref{sec:actual-ordinary-source-proof},
\ref{ym4:new-clean-box}--\ref{ym4:new-gaussian-segment} and
\ref{sec:finite-reference-normalizer}. The later bridge sections derive
compactness, D2, H1, Hilbert evolution and G2 from those inputs.
Existence of an arbitrarily fine admissible sequence uses the
determining coupling law. Weighted and matched-run comparisons and typed Ward
estimates are required only for the stronger statements invoking them.
Every use of a package retains its stated uniformity and scope.

Here are conventions shared by the packages. For numerical coefficients written in terms of $N$, the group is $\mathrm{SU}(N)$ and here $N=2$. Fix an odd integer $L>1$, a compact
simple matrix group $G$, and dimension four. The principal example is
$G=\mathrm{SU}(2)$. The lattice spacing is $a_K=L^{-K}$, and a finite torus has
physical side $\ell$. There is one group variable $U_b$ per positively oriented
bond, $U_{b^{-1}}=U_b^{-1}$, and normalized product Haar measure $dU$.
Write

\begin{gather*}
\FitMath{a_p(U)=1-\tfrac12\operatorname{Re}\operatorname{tr}U_p,} \\[3pt]
\FitMath{A(w,U)=\sum_p w(x(p))a_p(U),} \\[3pt]
\FitMath{Z_{K,\ell}(w)=\int e^{-\beta A(1,U)+A(w,U)}dU.}
\end{gather*}

The trace convention in this display is for $\mathrm{SU}(2)$; another group
requires its own normalization. Let $\mathcal O_K(f)=A(f,U)$. The connected
functions are the mixed derivatives at zero of
$\log Z_{K,\ell}(\sum_i z_if_i)$. Test functions used together have pairwise
disjoint compact supports, separation at least $d>0$, and individual supports in
balls of radius $R$. Set
$\|f\|_\sharp=\max(\|f\|_\infty,c_M\|\nabla f\|_\infty)$, with $c_M$ fixed.
All uniform constants below must be independent of $K$, the torus, the history,
the positions of the tests, and the relevant Gibbs state. Dependence on
$g,n,d,R$ must be displayed. A lower bound on admissible $\ell$ may depend on
$g$ and on constants fixed earlier. Replacing this uniformity by a bound at one
fixed cutoff does not discharge an obligation.

\subsubsection*{A. Background fields and propagators: four obligations}

\noindent\textbf{Obligation A1: the constrained critical orbit.}
We must prove that, for every allowed nested family of small-field regions
and every regular coarse datum $V$, there exists a regular critical orbit
$[U(V)]$ for the Wilson action with the prescribed block averages, unique
among regular critical orbits in a specified larger neighborhood, and that
a chosen gauge gives a representative which is Borel in all real data and
holomorphic in a complex neighborhood of each regular real datum. Both
neighborhoods have stated radii on the local scale;
their radii and constants are uniform in the number of scales. Existence of an
unrestricted global minimizer is insufficient: the needed assertion concerns
the constrained regular orbit and its gauge representative. Compatibility with
restriction to subregions, with seam constraints, and with changing the torus
must be part of the statement. This is the background-orbit obligation behind the
local assertions numbered 1a.3.6, 1a.3.8 and the background theorems.
For the pointed window deduction, the actual restrictions and substitutions
must remain in their prescribed pair cores. These cores contain the usable
Taylor chart ball and obey the curvature and own tent-action bounds on the
whole admitted core, including the finitely many young ages.
Section~\ref{ch:block-average-pivot-chart} proves a local chart for
the actual one-step average near a pure-gauge double block, including
the pivot derivative, exact straightening and measure factors. Its
explicit chart radius and nonwrapping premise do not imply existence,
uniformity or all-substitution room for the many-level critical orbit.

\medskip\noindent\textbf{Obligation A2: coercivity on a fluctuation fiber.}
Let $Q$ be the derivative of the chosen block average, and let $\mathcal E$ be
the gauge-fixed Hessian on the fluctuation variables at a regular background.
The chosen coordinates must remove all residual gauge directions. On each
conditioning set and with its specified boundary conditions we must prove

\begin{gather*}
\FitMath{\operatorname{Re}\langle B,\mathcal E B\rangle
 \ge c\|B\|_2^2,} \\[3pt]
\FitMath{Q B=0,}
\end{gather*}

with the norm in the units of that fluctuation scale. The same estimate, with
$c/2$, must hold throughout the claimed complex neighborhoods. Zero modes,
torus holonomies, and node constraints must either be absent from this space or
treated separately. A mass for this constrained operator is not a mass gap for
the physical continuum Hamiltonian. This is the missing quantitative content of
the conditioning and covariance assertions numbered 1c.6 and related claims.
For the actual unweighted classical flat reference,
Sections~\ref{ch:flat-contour-precision-gap}
and~\ref{ch:flat-whitney-comparison} prove
\(c(L)I\le\mathsf A_k^{\rm flat}\le CL^8I\), uniformly in the
number of previous classical blocking steps and in compatible tori,
including real flat holonomies. They use the actual plain constraint
and the soft internal-contour gauge. The curl-energy comparison is
\(\|da\|^2\le E_k(a)\le256\|da\|^2\). This supplies the flat
classical case, not the non-Abelian, nonzero-current, signed-profile,
layered or enlarged-complex-domain versions of A2.

\medskip\noindent\textbf{Obligation A3: local inverse and response bounds.}
We must prove that the inverses and response operators actually used,
including conditional covariances and derivatives of $U(V)$, have kernels
bounded by
$C e^{-c\operatorname{dist}(x,y)}$ in scale units. The bound must also control the
stated first differences and the stated H\"older differences throughout the
full complex datum polydisc, not only at the identity or at real data.
For the pointed window this is the product bond-log domain (PW1), with
the scale-weighted potential, gradient and current response in (PW2)--(PW3).
For a region with
several local scales, use a precisely defined weighted path distance and the
corresponding powers of the local spacings. Holomorphic derivatives in each
decoupling parameter must satisfy Cauchy bounds on a fixed enlarged complex
disc. A derivative joining distant sets must retain the decay needed by the
polymer norm. The mesh-uniform scalar estimate and its Dirichlet compressions are proved
in the block-resolvent section. The corresponding non-Abelian background and
multiscale response estimates remain to be proved. This includes the decoupled response assertion numbered 1c.10;
decay of one real inverse alone does not establish it.
Section~\ref{ch:flat-exponential-kernels} proves exponential kernels
for the actual unweighted classical flat precision and its constrained
covariance, uniformly in depth and compatible volume. It derives the
exact alias symbol, cancels its apparent gauge denominator on the full
complex zero cone, proves a uniform strip, and then inverts using the
flat gap. Periodization and unitary flat holonomies preserve the bounds.
This establishes the stated flat reference kernels, not the nonlinear
response or decoupling bounds required on the full interacting domains.
The finite-box comparison also requires compatible actual linearized
equations and boundary rows, including exponentially local gauge rows
when present. Cutting off an affine solution must leave polynomially
bounded boundary forcing, or an exponentially decaying boundary tail.
The one-pillar inverse bound is in the natural force norms, with the
fine-cell weights included. These properties, used in (PW12), do not
follow from bounds on unrelated inverse kernels.

\medskip\noindent\textbf{Obligation A4: averaging schemes and locality.}
We must prove that any transfer from one averaging scheme to another
provides a uniform intertwining map, including its coarse gauge factor. On fluctuation fibers the
transferred Hessians, inverses and constraints must obey A1--A3 with the same
allowed losses, for every number of blocking steps. Restrictions to a region
must depend only on an explicitly bounded enlargement of that region. Equality
of two linearized averages only at the identity is insufficient. A proof must
bound the accumulation of their differences at nonflat backgrounds and must
state the gauge in which the comparison is made.

\subsubsection*{B. Cluster expansion and power counting: four requirements}

\noindent\textbf{Requirement B1: a polymer majorant.}
Let polymers be finite connected sets of scale cubes, with incompatibility
meaning intersecting supports. A specified positive size $a(X)$ and activities
$z(X)$ must satisfy, uniformly in all external data,

\begin{gather*}
\FitMath{\sum_{Y:Y\not\sim X}|z(Y)|e^{a(Y)}\le a(X).}
\end{gather*}

The same inequality, with strict room, must hold for complex sources and
decoupling parameters on the domains used for differentiation. Require the
absolutely convergent connected-cluster formula for the logarithm, including
termwise differentiation and a uniform majorant for each derivative used later.
The finite-volume identity must be exact before taking limits. This is the
content needed from the convergence statement numbered 1d.4.1; a formal graph
expansion is not sufficient.
Section~\ref{ch:unmarked-polymer-majorant} derives the local unmarked
activity and union-filed logarithm bounds from its actual enlarged
Gaussian matrix/core hypotheses, sharp-cutoff structure and local
potential budgets. The Mayer-parameter discs are established at the
Gaussian integral, before zero-freeness is inferred. Gaussian volume
costs, failed bonds, actual neighbor-layer closures and the spare
potential slot are all paid for in the rate allocation. This also
supplies SI5 under those primitive hypotheses. We must still verify
them for the actual complex interacting matrices and all allowed cores;
convergence of an unrelated model gas would not supply this.

\medskip\noindent\textbf{Requirement B2: pointed marginal extraction.}
Sections~\ref{ch:pointed-hessian-cancellation}
and~\ref{ch:pointed-window-remainder} prove the following implication
for an actual regular member with no stored color spectators.
Its pointed total and differentiated sums must converge normally;
it is locally gauge invariant, and the individual member's pointed
totals are covariant under unit translations and the full hypercubic
group. Use the actual free response, with gauge covariance and affine
constant-curvature reproduction modulo gauge, with zero linear current.
The usable pair core, localized response and entropy gap must give
the absolute Hessian moments in \eqref{ph:exponential-moment}.
Then \eqref{ph:extractor} defines its coefficient and
\eqref{ph:own-counterterm} proves cancellation by that same member's
own midpoint tent action. Different members may be chosen at different
sites; their coefficients cannot be interchanged.

On physical radius \(a\), the readout retains
\(|b(E)|+\|T[E]\|\le Ca^{-2}\|E\|\).
Section~\ref{ch:extraction-radius-normalization} supplies its stated
near-zero inclusion and normalization. For the nonlinear remainder
also require A1--A3 throughout the complex datum domain, the actual
core inclusions, spatial H\"older regularity, compatible boundary
equations and the whole-core bound on the own tent action.
Put \(\sigma=L^{-(n-j)}\), \(\rho=\max(1,a_{\rm out}/a)\) and
\(0<\alpha_H\le1\). Under the window domain conditions the
near-domain sum at one base point, after its vacuum and own
counterterm subtraction, has bound

\begin{gather*}
\FitMath{C\|E\|\left[
 \rho^2\sigma^{4+\alpha_H}+\rho^3\sigma^6
 +(1+\rho^2\sigma^4)e^{-c/\sigma}\right].}
\end{gather*}

The final exponential term includes remote domains without a
nonwrapping chart. No factor proportional to the actual curvature
is asserted for that tail. Under the ordinary age bound
\(\rho^2\le C(1+n-j)\), the finite young-age core fallback and the
\(O(\sigma^{-4})\) site count give the summable weight
\(C\|E\|(1+n-j)^{3/2}L^{-\alpha_H(n-j)}\).

These deductions use the full member and field-domain hypotheses above.
In the actual RG application we must prove those hypotheses.
Torus-sensitive or field-sensitive tails, source derivatives of the
formation map, and the fluctuation disintegration laws in F2 use their
separate estimates. Stronger assertions in the arguments coded
1e.1.13, 1e.1.15 or 1e.1.16 need their corresponding estimates, which
we must also prove.

\medskip\noindent\textbf{Obligation B3: the one-loop law with a summable lattice correction and controlled nonlinear errors.}
For the chosen normalization let
$c_0=11N^2\log L/(12\pi^2)$. We must prove that the Gaussian contribution
$b_j^0$ is independent of the bare coupling and obeys
$|b_j^0-c_0|\le C L^{-\sigma j}$, or a stated summable substitute, and that
the full increment has a decomposition

\begin{gather*}
\FitMath{b_j=b_j^0+f_j+q_j,} \\[3pt]
\FitMath{|f_j|\le Cx_{j-1}^{-1/2}(\log x_{j-1})^p,} \\[3pt]
\FitMath{|q_j|\le Cx_{j-1}^{-(\kappa_0-4)/2},} \\[3pt]
\FitMath{\kappa_0>4.}
\end{gather*}

If a later deduction needs a stronger error, that stronger estimate must be
proved separately. We must obtain $c_0$ from an actual lattice trace
calculation with the stated zero-mode convention. A continuum perturbative
coefficient quoted without this comparison does not discharge this
obligation.

\medskip\noindent\textbf{Obligation B4: the ordinary step and its finer derivatives.}
The ordinary source box allows a separate Euclidean regular member at each base
point, its own scalar counterterm, a variable Wilson weight, and scale-small
noncovariant defects. We must prove that the step is holomorphic on this
product box with a majorant independent of the number of labels. It is not necessary that the
whole spatial array be translation invariant. The ordinary regular-output
estimate does not by itself prove every marked or large-field derivative.

For the stronger comparisons actually invoked, we must also prove the potential
derivative bound \eqref{eq:wext-potential-test} and the one-zone complex-mass bound
\eqref{eq:wext-zone-test}, on the domains stated in that section, with constants
independent of the levels and source positions. The regular irrelevance weight,
squared-radius weights for defects, and large-field survival weights have
different roles. In particular the survival factor belongs in the bound on a
zone's output; it does not enlarge its input radius by its reciprocal. Joint
marked derivatives and typed tensor counter-insertions must satisfy their
corresponding estimates. We must prove these finer estimates, with their stated complex
domains and uniform bounds, for the stronger comparison and Ward routes.
The ordinary-source insertion and angular proofs below do not obtain their
operator estimates by assuming these typed Ward derivatives.
Section~\ref{ch:ordinary-source-formation-audit} proves ordinary
old-level summations, scalar action-load and reserve implications.
The finite Gaussian, insertion and renewed-pointing sections derive
the local first response under their local Gaussian and bond-split
hypotheses; the unmarked section derives the needed majorant from
the stronger primitive matrix/core and potential package. Section~\ref{ch:window-potential-formation}
derives the ordinary formation bound using one near window member
at each old base point before its bracket is applied. The linear
bootstrap retains the source norm and gradient gain under these
inputs, the constant-source slopes and the old-remainder bounds.
Actual nonlinear background/core inclusions and history formation
remain to be verified; these implications do not supply them.

\subsubsection*{C. Large-field histories and closed induction: four obligations}

\noindent\textbf{Obligation C1: exact measurable operations.}
We must construct every move as a specified Borel kernel, change of
variables, or partition of unity on an explicitly defined space, and prove
the following. The tests are total, including at configurations where
logarithms or minimizers fail to exist. Every stored variable is integrated
exactly once or retained exactly once. For every bounded gauge-invariant
insertion $F$, the sourced construction preserves the total integral of
$F e^{-\beta A}$. An operation preserving only the integral without
$F$ cannot be used to transport that observable. For real inputs every summand
used as a probability weight must be nonnegative.

\medskip\noindent\textbf{Obligation C2: geometric margins and determination.}
We must give the allowed histories a constructive definition and prove the
following. At level $j$, tests,
local backgrounds and retained terms depend on at most a stated enlargement of
the determining regions. Every enlargement lies inside the margin supplied by
the next region. Healing, restriction, joining and continuation must preserve
these properties, including near holes, seams, and waiting windows. On an
admissible torus no object in these operations wraps around the torus.
The number of histories with a specified size must have a proved entropy bound,
with the same size function used in the rarity estimate.

\medskip\noindent\textbf{Obligation C3: rarity with entropy and insertions paid for.}
For each created large-field component $X$ at level $j$, we must prove a mass bound
$\mathfrak M_j(X)\le \exp[-c p_0(g_j)-\kappa |X|_j]$ with the stated
supremum over external complex data inside its internal integral when
component majorants are multiplied.
Here $p_0(g)=(\log g^{-2})^p$, and $\kappa$ exceeds the entropy rate of C2.
The bound must include the Wilson factor actually available to pay for that
component, all stored Jacobians and gauge factors, and the designated insertion
weights. A failure of existence of a background, as opposed to a quantitative
large-field inequality, must also be covered. Drop-outs without their own
Wilson factor need a separate bound after attachment to their carrying term.

\medskip\noindent\textbf{Obligation C4: invariance of the inductive class.}
We must construct one nonempty class $\mathfrak C_j$ of complete
presentations, with the same constants at every level, containing the bare
measure, and prove that under the increment and tube conditions the exact
step maps $\mathfrak C_j$ into $\mathfrak C_{j+1}$. Its conclusions include the analytic
polymer bounds, C1--C3, boundary bounds, covariance, and the required stability
floor. All constants must admit one simultaneous order of choice, without a
parameter needing to be both smaller and larger than quantities chosen later.
This is the full obligation represented by the step assertion numbered 2.11;
listing its clauses does not discharge it.

\subsubsection*{D. Sources and local correlation bounds: four requirements}

\noindent\textbf{Obligation D1: the ordinary localized source bound and stability.}
We must prove the precise torus conditions (ST5)--(ST6) of
Section~\ref{ch:all-state-source-transfer}: holomorphic point counterterms
and localized source energies, a strict entropy gap, source dependence only
on the specified regions, and independence of the representation of a
source function by a list of tests. On the ordinary source ball, we must
also prove the normalized upper stability bound and unsourced
partition-function floor there, with constants uniform in cutoff and allowed torus. The
source class must include the global mirror-periodized Lipschitz profiles
used in the chessboard construction. Fixed-cutoff bounds for point
counterterms may depend on the cutoff but must be uniform in the torus.

For the graded deduction, this obligation also specifies the actual domain
indexing of Section~\ref{ch:graded-rooted-source}. Take $M=L^{m_0}$ with
fixed integer $m_0\ge0$, and the cofinal compatible tori $\ell=2L^m$,
$m\ge m_0$. At scale $a_s=L^{-s}$ the localization cubes have side
$Ma_s$; their counts in every direction are even integers. Every fixed
power-of-$L$ chessboard tiles all sufficiently large members of this family.
All actual scales $0\le s\le K$ are retained. A domain $X$ is indexed by
its union of these cubes, with $m(X)$ the minimum number of edges of a
finite connected cubical edge tree meeting every cube. Its source
dependence is contained in a fixed number of neighboring cube layers.
The source-energy bound uses $e^{-\kappa m(X)}$, with $\kappa$ above the
rooted entropy rate for this actual indexing. If extra domain types are
present, their multiplicity must be included in that stated entropy gap.

The closed-walk count in (GS1), the ball-root count in (GS2), and the global
weighted count (GS3) are then proved geometric consequences. In particular,
the sharp local count is $C(1+u/a_s)^4e^{-ct}$, and the global count is
$C\ell^4a_s^{-4}$. The weaker abstract (ST6) estimate alone does not imply
these sharper bounds and is insufficient for the graded conclusion below.
Checking that the localized RG energies use this specified geometry is
part of D1, not supplied by writing an abstract domain sum.

These are raw partition-function bounds. A zero-free logarithm on the
whole global source ball is not part of this obligation. The fixed-volume normalized
zero-free lemma gives a suitable neighborhood when needed; the all-state
transfer theorem instead constructs common holomorphic cell normalizers
and obtains the derivatives at zero directly. The elementary transfer does
not prove the ordinary RG source estimates in this obligation.
Section~\ref{ch:ordinary-history-stability} proves explicit sufficient
criteria for the upper bound and floor from the normalized positive
label/history representation: total-variation bounds with a weighted
label sum, and a regular all-small-field label with its own action
bound. We must still verify those criteria, including the reserve and
allocation of Wilson factors. A pointwise
polymer logarithm cannot replace that history representation without
an exact common integral formula.

% Revision4 promotion follows the completed independent QF review.
\medskip\noindent\textbf{Requirement D2: graded all-state bounds derived from D1 and G1.}
Theorem~\ref{thm:graded-rooted-source} uses D1 in its actual rooted setting
above and the proved chessboard domination from G1. It yields one fixed
finite integer $k$ and $B(R)=C(R+1)^4$ such that, for periodic Gibbs states
$\rho,\nu$ at the same datum,

\begin{gather*}
\FitMath{\left|\rho\prod_{i=1}^n
    [\mathcal O_K(f_i)-\nu(\mathcal O_K(f_i))]\right|
 \le B(R)^n(n!)^k\max\{1,d^{-5n}\}\prod_i\|f_i\|_\sharp.}
\end{gather*}

For $n\ge2$ the cumulants under $\rho$ satisfy the same estimate with
$[2B(R)]^n$ in place of $B(R)^n$. Each support lies in its own radius-$R$
ball, placed anywhere, and different supports have distance at least $d$.
The first-order case uses the one-point difference bound. All constants
are uniform in cutoff, states and support positions.

The proof gives a polynomial cell-normalizer loss $Q_i\le Cn^p$ and
$k=\lceil p\rceil+5$. It uses a genuine product source polydisc from
separate chessboard factors and common univariate normalizers. The
separation singularity is explicitly $d^{-5n}$, not a single $d^{-5}$
prefactor independent of $n$. Thus D1, its specified geometry and G1 imply D2. We must
retain the actual D1 construction of
Section~\ref{sec:actual-ordinary-source-proof}; the present counting
argument does not supply it independently.

\medskip\noindent\textbf{Obligation D3: age-resolved uniform rarity.}
We must prove that the probabilities or positive weights of live regions
persisting from a given level to a later level have a summable bound uniform
in volume and in the location of a test. Witness selection, replication by reflections, and pattern
counting must preserve the insertion and source weights. State the resulting
tail function $r(m)$ with $\sum_m r(m)<\infty$, together with every additional
moment in $m$ used later. A rarity estimate for one newly created component is
insufficient for this persistence assertion. This obligation represents the age and
survival arguments wherever actual history persistence or weighted comparisons
are used. The ordinary graded all-state estimate D2 is obtained by its GS
route above and does not require this separate history-persistence obligation.

For deductions from a fixed-bump Riemann sum, the cutoff is required to be sufficiently fine.

\medskip\noindent\textbf{Requirement D4: one-point transfer derived from D1.}
Theorem~\ref{thm:all-state-source-transfer} proves, under D1, that periodic
Gibbs states $\nu,\rho$ at the same datum satisfy

\begin{gather*}
\FitMath{|\nu(\mathcal O_K(f))-\rho(\mathcal O_K(f))|
       \le C_1(R)\|f\|_\sharp.}
\end{gather*}

If both states are invariant under all plaquette isometries, the bound implies
$|\nu(a_p)-\rho(a_p)|\le C a_K^4$ by testing a fixed nonnegative bump of
nonzero integral. The uniform bound on differences follows from a common holomorphic
normalizer; neither
individual uncentered mean is required to remain bounded. This is essential
when a pure component is centered with respect to a mixture.

\subsubsection*{E. Flow, tuning and comparison of cutoffs: four obligations}

\noindent\textbf{Obligation E1: a consistent flow.}
We must prove that the recursion $x_j=\beta-\sum_{i\le j}b_i$ is defined up
to first passage, with $b_i\le B_\sharp$ and
$\lambda(j-i)-C\le x_i-x_j\le B_\sharp(j-i)$, $\lambda>0$, and that its
increments are independent of the admissible torus and of the unused finer
cutoff. The elementary first-passage argument is applied to a recursion
whose analytic formation and domain conditions we must establish through
every preceding step of the run. Together with B3, these properties imply
the leading logarithmic relation between $\beta$ and the cutoff.

\medskip\noindent\textbf{Obligation E2: a frozen calibration and tuning map.}
For auxiliary choices fixed independently of $\beta$, we must prove that the
terminal map $\beta\mapsto x_K^{\natural}(\beta)$ is holomorphic on a stated
complex neighborhood of a real interval, that
$0<c\le\partial_\beta x_K^{\natural}\le C$ on that interval, with endpoint
values bracketing the prescribed renormalized inverse coupling, and that all
intermediate values remain in the allowed tubes. This then supplies existence
and local uniqueness of a tuned bare coupling by elementary analysis. The holomorphy assertion is not
valid merely because the original finite-lattice partition function is entire.
The flow section proves the slope implication both from the stronger
independent coupling polydiscs and from the weaker marked bare-disc
route. In the latter, the marked forcing is of size
\(C\epsilon'_i/\varrho\), not \(C\epsilon'_i/x_{i-1}\), and
is retained in the discrete Gronwall estimate. Its summability needs
the stated exponent restriction. Canonical tuning uses the admitted
real interval; broader presentation gluing additionally needs an
admitted real overlap. Overlapping geometric complex discs alone
do not establish that admission.

\medskip\noindent\textbf{Obligation E3: a summable matched-cutoff comparison.}
At matched terminal coupling, we must prove for each fixed torus and separated tuple

\begin{gather*}
\FitMath{|S^T_{n;K+1}(X)-S^T_{n;K}(X)|
     \le C(n,d,g,\ell)\epsilon_K\prod_i\|f_i\|_\sharp,} \\[3pt]
\FitMath{\sum_K\epsilon_K<\infty.}
\end{gather*}

The floor in $K$ may depend on the fixed torus. To compare two nearby terminal
couplings, add a uniform Lipschitz term proportional to their difference.
Neither assertion implies infinite-volume uniqueness. This is the actual
analytic obligation behind the matched-run claims numbered 8a and the tuned Cauchy
sequence; a bounded difference at each cutoff does not suffice.

\medskip\noindent\textbf{Obligation E4: compatibility of calibrations.}
When moving between a self-calibrated first-passage run and a frozen run, we
must prove that the maps represent the same bare integral and satisfy a stated
uniform bound on their coupling discrepancy, small enough for every comparison
domain and source radius used above. Exact dilation identifies the
same finite lattice in two units, but supplies no estimate on this calibration
discrepancy. Uniqueness statements must state which terminal coupling and which
calibration they use.

\subsubsection*{F. Short-distance normalization and nontriviality: four requirements}

\noindent\textbf{Requirement F1: the proved free-kernel comparison.}
The exact continuum coefficient is proved in Section~\ref{ch:free-curvature-exact}.
For the remaining lattice comparison, let $a=\ell/n$, with integer $n$, and use
plain coordinate inner products on site, bond and positively oriented plaquette
cochains of the periodic lattice. Let $d_1$ be its unscaled circulation matrix
and $P_{\rm ex}$ the orthogonal projection onto $\operatorname{ran}d_1$.
A standard Gaussian on that range has covariance $P_{\rm ex}$. Its physical
curvature kernel is the $6\times6$ matrix of components

\begin{gather*}
\FitMath{K_{a,\ell;P,Q}(p,q)=a^{-4}P_{\rm ex}(p_P,q_Q),}
\end{gather*}

with each component sampled at its own plaquette barycentre. Section~\ref{ch:lattice-curvature-comparison} proves the following
comparison with the continuum matrix $K$, with $\gamma=1/2$:

\begin{gather*}
\FitMath{|K_{a,\ell;P,Q}(p,q)-K_{P,Q}(x(p)-x(q))|
 \le C r^{-4}\big[(a/r)^\gamma+(r/\ell)^4\big],} \\[3pt]
\FitMath{a\le r/10,} \\[3pt]
\FitMath{r\le\ell/10,} \\[3pt]
\FitMath{r=|x(p)-x(q)|,}
\end{gather*}

for some fixed $\gamma>0$, uniformly in $a,\ell$, plaquette positions and
orientations. Periodic differences are represented by the short displacement
in this range. The same section proves separated smearing convergence and the resulting
quadratic-density coefficient. This closes the free exact-cochain Gaussian
requirement; the interacting extraction F2 must still be proved separately.

The projection convention in this requirement is distinct from physical weighted
block means. When passing to unit physical blocks use the adjoints and physical
difference $d_\eta$ established in the block-normalization section. Under plain
fine and coarse inner products an arithmetic block mean does not have an
isometric adjoint; no projection identity may silently interchange these
conventions. Neither normalization proves the interacting extraction F2.

\medskip\noindent\textbf{Obligation F2: extraction of the leading two-point term.}
For the prescribed short-distance test pairs we must prove the finite-cutoff identity

\begin{gather*}
\FitMath{S^T_{2;K,\ell}(f,h)=C_G\zeta_K(d)^2g_K(d)^4
                      \mathcal K(f,h)(1+r_{K,\ell}(f,h)),}
\end{gather*}

where $C_G>0$, the real-data error satisfies $|r|\le1/2$, and all lower
cutoff requirements are explicit. For the angular deduction this obligation
must cover the full class $\mathfrak T(d;\vartheta_*)$ in (TU9)--(TU13)
of Section~\ref{ch:angular-two-point-elimination}, with
$\vartheta_*=1/12000$, or a broader class. These are nonnegative $C^1$
tests in radius-$d$ balls, with support distance between $d$ and $6d$,
the stated lower integral and upper Lipschitz bounds, and arbitrary
positions. A law proved only for a narrower shape class is insufficient
without a further inclusion argument.

We must prove the same bound for $S^T_2[\rho]$ in every periodic Gibbs state
used to approximate the limit or any component. At fixed $d$ its
eventual cutoff threshold and constants must be uniform in those states
and support positions. This uniformity in the state and support positions is part of F2. On each fixed torus any
asserted $r\to0$ needs a separate uniform estimate; it must not be silently
promoted to a statement uniform in volume.
This is the substantive Gaussian-extraction obligation behind the short-distance
law, including all interaction and large-field remainders.
The finite source insertion and renewed-pointing proofs retain the
\(\alpha^2s\) potential norm under the actual unmarked-majorant,
leg-path and complex-domain hypotheses. The window formation and
linear bootstrap propagate it under their explicit ordinary inputs;
the resulting near-field estimate requires the complete field to
have size \(\delta\), not merely a \(\delta\)-sized increment on an
arbitrary \(\alpha\)-background. Remote and holonomy tails retain
their separate bounds. These deductions do not prove the background
and history obligations for the interacting construction.

The actual clean-history normal form must still control every old
activity's localized Hessian, the conditional exterior force and
pointwise deletion with the same exterior field and correct flag mass.
The actual nonlinear averaging-fibre comparison is an additional
obligation for that particular route; gauge fixing the actual run
provides a route that avoids it. That route has its own normal-form, exterior-force and
pointwise clean-class conditions, which we must prove.
Section~\ref{ch:two-point-extraction-tests} proves the covariance
arithmetic and common-normalizer transfer once their inputs hold.
Section~\ref{ch:old-activity-hessian} derives the old-activity
localized Hessian from the actual all-substitution window room,
window-first value bounds and holomorphy on a Gaussian-coordinate
ball of radius \(c\alpha/g\). The resulting scale-small bound has
factor \((g/\alpha)^2\). The room mechanism is in Section~\ref{sec:substitution-room}; the
conditional exterior-force and pointwise clean-class laws are
outlined in Sections~\ref{ym4:new-clean-box}--\ref{ym4:new-gaussian-segment}.
The local one-step chart or a value bound without a complex domain
does not supply those conditions.

The transfer of a Gaussian reference law has two different scopes.
D1's common cell normalizer compares any two periodic infinite-volume
states; an infinite reference obtained with a Gaussian estimate uniform
along arbitrarily large tori therefore needs no finite-volume vacuum
comparison. A law proved only on one prescribed finite reference torus
instead requires a comparison with those states. Compatible local
counterterms, the floor and vacuum-density comparison give one
sufficient route. The local companion criterion (ET8)--(ET10) proves
that vacuum comparison from explicit nonwrapping and normalization
data; it is not required in the infinite-reference route. Per-torus
estimates whose distance thresholds decrease with the volume do not
establish the latter reference law. Section~\ref{sec:finite-reference-normalizer}
supplies the finite-reference route directly: upper stability is
used on the reference torus, and all larger tori contribute only
their unsourced floors and the common vacuum-density bound.

\medskip\noindent\textbf{Obligation F3: the running normalization at the witness scale.}
The witness depth must be defined without reference to a desired conclusion,
and its lattice separation must lie within the permitted comparison range.
We must prove

\begin{gather*}
\FitMath{g_K(d)^{-2}=c_*\log(1/d)(1+o(1)),} \\[3pt]
\FitMath{\zeta_K(d)=1+o(1),} \\[3pt]
\FitMath{c_*>0,}
\end{gather*}

with the precise uniformity needed by F2. A derivative of a running coupling
with respect to a constant source must be shown to equal this normalization;
it is not an independent choice of field renormalization. F2--F3 give a
nonzero two-point limit by a closed lower bound.
Appendix~\ref{ch:comparison-capsule-v7} proves summability of the
normalization slopes from the actual increment estimates on
coupling-sized independent complex polydiscs and control of the
large-field remainder loss. The flow section preserves the alternative
with an actual marked remainder on a fixed bare-coupling disc and its
separate summable forcing. A sharp \(O(g\log^p g^{-2})\) slope rate
on this weaker route needs \(p_R\ge3/2\); summability alone needs
\(p_R>1\). Analyticity on an ordinary source disc alone does not
provide the required marked size and history bounds.

Here is their sufficient ultraviolet consequence, with its size estimate
included. Use

\begin{gather*}
\FitMath{\mathcal K(f,h)=\iint f(x)h(y)|x-y|^{-8}\,dx\,dy .}
\end{gather*}

If the nonnegative supports lie in radius-$d$ balls at distance at least
$d$, and $v_4$ is the volume of the unit ball in $\mathbb R^4$, then

\begin{gather*}
\FitMath{0\le\mathcal K(f,h)\le d^{-8}\|f\|_{L^1}\|h\|_{L^1}
       \le v_4^2\|f\|_\infty\|h\|_\infty .}
\end{gather*}

The uniform F3 bounds imply $|\zeta_K(d)|\le Z$ and
$g_K(d)^4\le C[\log(1/d)]^{-2}$ for small $d$ and all permitted sufficiently
fine cutoffs. Thus F2 gives

\begin{gather*}
\FitMath{0\le S^T_2[\rho](f,h)
     \le C'[\log(1/d)]^{-2}\|f\|_\infty\|h\|_\infty .}
\end{gather*}

The lower inequality uses $1+r\ge1/2$; the upper inequality uses
$1+r\le3/2$. At each fixed pair these closed inequalities pass to every
limit and to the all-state phase realizations. Taking
$u(d)=C'[\log(1/d)]^{-2}$ proves precisely (TU10), with $u(d)\to0$.
This deduction does not verify the state-uniform extraction or running
normalization required in F2--F3.

\medskip\noindent\textbf{Obligation F4: interpretation of the observable sector.}
The conclusions of F1--F3 concern the scalar curvature-density sector. To
identify a reconstructed theory as the full Yang--Mills theory, we must
specify a sufficient gauge-invariant observable family, its consistency
relations, renormalization, and its relation to this scalar sector, and
prove the needed continuum gauge identities or an equivalent identification.
A nonzero scalar Wightman field alone does not provide this identification.
Until this is proved, the conclusion must be stated as existence of the
conditional scalar observable theory, with full gauge-theory identification
left open.

\subsubsection*{G. Symmetries, rotations and phases: four requirements}

\noindent\textbf{Requirement G1: proved finite-lattice symmetries and positivity.}
Section~\ref{ch:periodic-link-symmetry} proves link-reflection invariance
and the common period $2\mathbb Z^4$ for every periodic $\mathrm{SU}(2)$
Wilson DLR state at fixed $\beta>0$. Section~\ref{ch:all-mirror-periodic}
then proves invariance and reflection positivity at every coordinate site,
coordinate link, and diagonal lattice mirror, on the full bounded Borel
half-space bond algebra. No restriction to torus-selected states is made.
The compositions give all unit translations and signed permutations.
At $\beta=0$ the unique product Haar state has the same properties.
Adapted split tori are used for diagonal mirrors; an invalid diagonal
split of an ordinary cubic torus is not assumed.

\medskip\noindent\textbf{Requirement G2: scalar rotations derived from ordinary sources and the ultraviolet law.}
Consider limits through periodic Gibbs states, including the canonical
components with their all-state realizations. G1 and the compactness
inheritance argument provide their coordinate and diagonal reflection
positivity; Section~\ref{ch:euclidean-hilbert-bridge} provides the exact
coordinate cube cone. A component with mean $c\int f$ is first centered
at its own mean. Its self-centered cutoff moments converge to this
centered functional by the finite subset expansion and the convergent
one-point differences; the same expansion restores the invariant mean
after the proof. Under D1 in its actual rooted geometry,
Theorem~\ref{thm:uniform-one-field-insertion} proves a Hilbert insertion
bound uniform in the number of surrounding fields. This operator estimate
does not follow from D2 alone.

Theorem~\ref{thm:plate-density-from-insertion} converts that bound to an
analytic $L^1$ plate density, with the explicit loss of five powers of the
time margin. Section~\ref{ch:angular-kernel-construction} constructs the
holomorphic product kernels and vectors, glues their frames, and proves the
finite angular band and continued Gram positivity.
Theorem~\ref{thm:angular-two-point-elimination} uses the exact (TU10)
shape law to remove every nonconstant two-point harmonic, including the
critical frequency; its diagonal measure and cone support are constructed
in the proof. Theorem~\ref{thm:all-orders-angular-closure} proves polygon
divisibility, integer convexity, removal of exceptional distributional
support by common cuts, and full $\mathrm O(4)$ covariance. QF transfers
this covariance to the diagonal-flat Schwartz domain.

Together with the preceding deduction from F2--F3 to (TU10), this is a
conditional proof of G2 for this state class. It assumes neither typed
Ward counter-insertions nor vanishing distant-shell flux. It does not
verify D1, F2 or F3. An arbitrary growing ordinary cubic-torus limit must
separately satisfy the inherited mirror and cone interface before this
route applies; being a separated-point limit alone supplies no such
identification.

The retained stronger Ward route still has its own contract: an exact
lattice identity must pass to
$\delta_\omega S_n(X)=-\lim_{R\to\infty}\lim_j
\langle\mathfrak B_{\omega,R};X\rangle_j$, with every bulk and
discretization error controlled and the claimed boundary limit existing
and vanishing. Wherever that route is invoked, we must prove these limits
in that order, with the typed estimates they need. The angular
proof does not establish those separate comparison estimates.

\medskip
\noindent\textbf{Requirement G3: the proved finite-spacing plateau law.}
For a spacing $a>0$, let $\Omega_a=\mathrm{SU}(2)^{\mathcal B_a}$, where
$\mathcal B_a$ is the countable set of positively oriented bonds of
$a\mathbb Z^4$. Give this space its product topology. Let $\mathcal G_{a,\beta}$
be the translation-invariant probability measures satisfying the Wilson DLR
equations at inverse coupling $\beta$. The finite-bond conditional kernel
integrates the bonds in $W$ against the normalized density
$\exp[-\beta\sum_{p:p\cap W\ne\varnothing}a_p]$ and holds the other bonds
fixed. These kernels take continuous local functions to continuous local
functions. Thus $\mathcal G_{a,\beta}$ is a compact convex set. Its reflection invariance and reflection positivity are proved under G1.

Let $\mathcal I$ be the completed $\sigma$-field of events invariant, modulo
the measure, under all translations in $a\mathbb Z^4$. For bounded local
functions $A,B$ define

\begin{gather*}
\FitMath{\operatorname{Pl}_\nu(A,B)
 =\nu\bigl(\overline{E_\nu[A\mid\mathcal I]}
                         E_\nu[B\mid\mathcal I]\bigr).}
\end{gather*}

Theorem~\ref{thm:lattice-plateau} proves the following law, with one numerical
constant $C$ independent of $a,\beta,\nu,r,n,A,B$. If $r\ge1$ is an integer,
$A,B$ depend only on bonds with both endpoints in $a[-r,r]^4$, and
$y\in a\mathbb Z^4$ has $n=|y|_\infty/a-2r\ge1$, then

\begin{gather*}
\FitMath{|\nu(\overline A\,\tau_yB)-\operatorname{Pl}_\nu(A,B)|
 \le C\left[(r/n)^2+(r/n)^3\right]
                \mathcal N_{\nu,r}(A)^{1/2}
                \mathcal N_{\nu,r}(B)^{1/2},}
\end{gather*}

where

\begin{gather*}
\FitMath{\mathcal N_{\nu,r}(A)
 =\max_{\mu=1,\ldots,4;\,\gamma\in\Gamma_{16}}
 \nu\bigl(\overline{\theta_\mu\gamma A}
                         \tau_{2rae_\mu}\gamma A\bigr)
 \le\|A\|_\infty^2.}
\end{gather*}

Here $\theta_\mu$ reflects $x_\mu$ to $-x_\mu$, and $\Gamma_{16}$ is the group
of coordinate sign changes. Nonnegativity of these reflection sizes follows from G1. The law holds for every $\nu\in\mathcal G_{a,\beta}$, including
mixtures. It is a statement before any continuum component is constructed.

Theorem~\ref{thm:compact-conditional-law} constructs the required conditional
probabilities on this compact metric space, including Borel versions and
simultaneous probability-valued uniqueness. Applied to $\mathcal I$, this
theorem, the DLR equations, and the law above yield a unique decomposition
into ergodic DLR states. The argument and the required closedness of the pure
state set are proved in the phase-interface section. Compactness of that pure
set and existence of continuum components are conclusions, not assumptions.

\medskip\noindent\textbf{Requirement G4: all-state moments derived from ordinary sources.}
Choose $a_j=L^{-K_j}\downarrow0$, couplings $\beta_j$, and centering states
$\nu_j\in\mathcal G_{a_j,\beta_j}$. Write
$\mathcal O_j(f)=\sum_p f(x(p))a_p$. Once the precise ordinary source obligation D1 is
proved, Theorem~\ref{thm:all-state-source-transfer} gives for every
$\rho\in\mathcal G_{a_j,\beta_j}$ that

\begin{gather*}
\FitMath{|\rho(\mathcal O_j(f))-\nu_j(\mathcal O_j(f))|
       \le C_1(R)\|f\|_\sharp,} \\[3pt]
\FitMath{|\kappa^\rho_n(\mathcal O_j(f_1),\ldots,\mathcal O_j(f_n))|
       \le C_n(d;R)\prod_i\|f_i\|_\sharp} \\[3pt]
\FitMath{(n\ge2).}
\end{gather*}

The tests are real $C^1$ functions; each support is contained in a ball of
radius $R$, and different supports have distance at least $d>0$ in the second
bound. Here $\|f\|_\sharp=\max(\|f\|_\infty,c_M\|\nabla f\|_\infty)$.
The constants are independent of $j,\rho$ and support positions. They also
apply to the separated reflected copies used in G3. For the chosen sequence,
we take the moments centered at $\nu_j$ to converge on every separated tuple,
passing to a subsequence by the preceding compactness theorem if necessary.
The per-order transfer alone would not establish factorial growth. In the
actual rooted setting of D1, Theorem~\ref{thm:graded-rooted-source} also
provides D2, and Theorem~\ref{thm:quantitative-flat-extension} gives H1.
The bounds displayed here and G3 suffice for the canonical decomposition constructed below.
They do not assume that decomposition or clustering of its future components.


\subsubsection*{H. Test spaces and reconstruction: four requirements}

\noindent\textbf{Requirement H1: quantitative diagonal-flat extension derived from D2.}
Let $\mathcal S_{\ne}(\mathbb R^{4n})$ be the Schwartz functions flat, to all
orders, whenever two arguments coincide. Under the D2 moment bound, passed
to a separated-tuple limit, Theorem~\ref{thm:quantitative-flat-extension}
constructs its unique continuous extension to this space and proves

\begin{gather*}
\FitMath{|\mathsf S_n(\varphi)|\le C_{\rm E}^n(n!)^{\kappa_{\rm E}}
 \sup_{|\alpha|\le k_{\rm E} n,x}(1+|x|^2)^{k_{\rm E}n/2}
                          |\partial^\alpha\varphi(x)|,} \\[3pt]
\FitMath{k_{\rm E}=20,} \\[3pt]
\FitMath{\kappa_{\rm E}=k+32.}
\end{gather*}

The constants are independent of $n$. The same uniform statement applies
to component limits obtained from the all-state bounds. Only the
position-uniform radius-one moment bound is used; a unit partition already
gives an effective $B(R)\le CB(1)(1+R)^4$. The explicit dyadic/product-box
proof includes tensor-relation independence, all diagonal cutoffs,
infinity cutoffs, and the positive-time boundary. Its constrained tensor
density transfers reflection positivity to finite positive-time Schwartz
sequences. Thus D1, its actual GS geometry and G1 imply H1.

A distribution only on the complement of the diagonals with unspecified
constants would not establish this result. Diagonal-flat functions are
not dense in the full Schwartz space; an extension across diagonals or
contact-term identities still need a separate renormalization statement.

\medskip\noindent\textbf{Requirement H2: positivity and continuous Euclidean evolution.}
Section~\ref{ch:euclidean-hilbert-bridge} proves the compact-tuple Hilbert
construction from invariant separated moments, full reflection positivity
and position-independent per-ball bounds. It proves strong continuity, time
contraction and the cube spectral cone for the stated lattice approximation.
For the full positive-time test sequences in the precise reconstruction
theorem, the required finite quadratic inequalities are
$\sum_{p,q}\mathsf S_{p+q}(\Theta\overline\varphi_p\otimes\varphi_q)\ge0$.
All tensors in this expression must belong to the declared domain. Under H1 and the constrained tensor-density lemma, those full sequence
inequalities and the same Hilbert evolution follow from the compact-tuple
construction. The proof bounds all translated matrix elements before
iterating Cauchy--Schwarz. Its position-independent per-ball hypotheses
are supplied by D1 for the ordinary phase moment bounds. The compact
argument does not by itself give quantitative Schwartz growth; that
additional conclusion is now the GS/QF deduction recorded as H1 above.

\medskip\noindent\textbf{Imported theorem H3: the reconstruction theorem in its applicable form.}
The reconstruction step uses an Osterwalder--Schrader theorem whose test
spaces and growth condition match H1--H2. Its other hypotheses
are Euclidean invariance, permutation symmetry, and the required cluster
property for a unique vacuum. G2 supplies the invariance for the periodic-state
limits and components once the ordinary-source and ultraviolet obligations
D1 and F2--F3 are proved. The reconstruction theorem itself remains external.
Its conclusion is a scalar Wightman theory with
tempered Minkowski distributions, locality, a positive Hilbert metric, and the
forward spectral cone. Here it is used as an imported theorem. Contact terms invisible on diagonal-flat Euclidean tests
must not be assigned arbitrary values and then treated as consequences of the
reconstruction.

\noindent\textbf{H4: conclusions and scope.}
To complete the proof on its declared route for periodic-state limits,
the actual D1 and F2--F3 proof outlines must be completed and a
sequence of admissible data constructed. The supporting propositions then supply
the downstream scalar conditions, and H3 is applied in its matching
test spaces and growth class. D1 and the interacting short-distance law
are thus obligations of this route, not assumptions.
The conclusions that use weighted, matched-cutoff or Ward comparisons
require those comparisons to be proved as well. Identification with the
full gauge theory additionally requires F4. A mass-gap conclusion requires
a separate spectral estimate, for example one placing the spectrum
outside zero in $[m,\infty)$ for some $m>0$.
The conclusions are therefore taken in stages: a finite lattice measure,
a subsequential separated-point limit, an axiomatic scalar theory, and
its Yang--Mills identification, each with its own obligations.

% END MODULE chapters/dependency_catalogue.tex


% BEGIN MODULE chapters/rg_audit_chapter.tex
\section{From local source bounds to continuum limits}
\label{ch:conditional-source-bridge}
\subsection{From local source estimates to separated correlation limits}
\label{sec:conditional-source-bridge}

This section shows that the local analytic energies and residual
finite-volume integral bounds stated below, which we must prove, yield
the source-disc, cumulant and limit conclusions of this section. Exact
measure transport and geometric forgetting of older scales enter through
their own stated conditions, which we must also prove.

\subsubsection{Finite-dimensional starting point}

Let $\mathcal P$ be a finite set of plaquettes, let $G=\mathrm{SU}(2)$, and assign one variable $U_b\in G$ to each positively oriented bond. Opposite orientations carry $U_{-b}=U_b^{-1}$. The product of normalized Haar measures is denoted by $dU$. Write

\begin{gather*}
\FitMath{a_p(U)=1-\tfrac12\operatorname{Re}\operatorname{tr}U_p,} \\[3pt]
\FitMath{A(U)=\sum_{p\in\mathcal P}a_p(U).}
\end{gather*}

Since $0\le a_p\le2$, the probability measure

\begin{gather*}
\FitMath{d\mu_\beta(U)=Z_\beta^{-1}e^{-\beta A(U)}dU,} \\[3pt]
\FitMath{Z_\beta=\int e^{-\beta A(U)}dU}
\end{gather*}

is well defined for every real $\beta$. Given bounded real functions $f_i$ on plaquette barycentres, set $X_i=\sum_pf_i(x_p)a_p(U)$ and

\begin{gather*}
\FitMath{\Gamma(z)=\int e^{\sum_i z_iX_i}d\mu_\beta.}
\end{gather*}

The bound $|X_i|\le2\sum_p|f_i(x_p)|$ implies uniform domination on every compact subset of $\mathbb C^n$. Differentiation under the integral therefore proves that $\Gamma$ is entire. Since $\Gamma(0)=1$, it has a holomorphic logarithm in a neighbourhood of the origin, and the joint cumulant is defined by

\begin{gather*}
\FitMath{\langle X_1;\ldots;X_n\rangle
 =\partial_{z_1}\cdots\partial_{z_n}\log\Gamma(0).}
\end{gather*}

The neighbourhood obtained this way can depend on the number of plaquettes.

If $T$ is a fixed linear operation on finite measures which preserves their total mass, then $T(e^{\sum z_iX_i}\mu_\beta)$ has total mass $\Gamma(z)$. For a finite sequence of such operations the same conclusion follows by induction. If $T$ is specified only on gauge-invariant densities, all insertions in this assertion must be gauge invariant. Neither linearity nor mass preservation supplies a uniform analytic bound on the resulting density.

\subsubsection{A zero-free disc from a normalized bound}

\paragraph{Lemma.}
Let $c_i>0$ and put $\|z\|_c=\sum_i c_i|z_i|$. Suppose $h$ is holomorphic on $\|z\|_c<1$, $h(0)=1$, and $|h|\le Q$, where $Q\ge1$. Then

\begin{gather*}
\FitMath{|h(z)-1|\le(Q+1)\|z\|_c.}
\end{gather*}

In particular $h$ is nonzero on $\|z\|_c<(Q+1)^{-1}$. On the smaller ball $\|z\|_c<[2(Q+1)]^{-1}$ its principal logarithm satisfies $|\operatorname{Log}h|\le\log2$.

\begin{proof}
For $z\ne0$ write $z=tv$, where $t=\|z\|_c$ and $\|v\|_c=1$. The function $(h(\zeta v)-1)/(Q+1)$ maps the unit disc into its closed unit disc and vanishes at zero. The Schwarz lemma gives the first estimate. Strictness on the stated open balls follows by using slightly smaller compact discs and taking limits. For $|u|\le1/2$, the logarithmic series gives

\begin{gather*}
\FitMath{|\operatorname{Log}(1+u)|\le\sum_{k\ge1}\frac{|u|^k}{k}
 =-\log(1-|u|)\le\log2.}
\end{gather*}

The same series defines a single holomorphic branch on that ball.
\end{proof}

\paragraph{Corollary.}
Suppose $\Gamma=e^{\mathcal E}h$ near zero, where $\mathcal E(0)=0$ and the lemma applies to $h$. Then

\begin{gather*}
\FitMath{\left|\partial_{z_1}\cdots\partial_{z_n}
       (\log\Gamma-\mathcal E)(0)\right|
 \le\log2\,[2(Q+1)n]^n\prod_i c_i.}
\end{gather*}


\begin{proof}
Use Cauchy's integral formula on a polydisc with radii strictly smaller than $[2(Q+1)nc_i]^{-1}$. The sum of the weighted radii is less than $[2(Q+1)]^{-1}$. Each variable is differentiated once, so no factorial occurs in this application of the formula. Let the radii increase to their stated values.
\end{proof}

The dependence on $Q$ is material. For example, $\cosh(Vz)$ is the moment-generating function of a random variable taking the values $V$ and $-V$ with equal probabilities. It is at most $e^V$ on the unit disc, but has a zero at $i\pi/(2V)$. Thus an exponential-in-volume estimate alone cannot give a volume-independent zero-free disc.

\subsubsection{A precise local-energy hypothesis}

Fix $L>1$, a finite physical volume $V$, and a scale index $s\ge0$, with physical scale $a_s=L^{-s}$. A local domain $D$ has an integer size $m(D)\ge0$ and a dependence region $D^+$. Assume there are constants $C_0,C_1,\kappa_{\mathrm{ent}}$ such that

\begin{gather*}
\FitMath{\operatorname{diam}D^+\le C_0a_s(m(D)+1),}
\end{gather*}

and, for every fixed compact test support $B$, the number of domains of size $m$ whose dependence region meets $B$ is at most

\begin{gather*}
\FitMath{C_1V a_s^{-4}e^{\kappa_{\mathrm{ent}}m}.}
\end{gather*}

The factor $V$ is a finite-volume allowance; a more local counting estimate can replace it when available.

For cutoff $K$ let $e_{K,s,D}(w)$, $0\le s<K$, be holomorphic on the source ball

\begin{gather*}
\FitMath{\sum_i|z_i|\|f_i\|_*/r<1,} \\[3pt]
\FitMath{w=\sum_i z_if_i,}
\end{gather*}

where $r>0$ is independent of $K$. Suppose

\begin{gather*}
\FitMath{e_{K,s,D}(0)=0,} \\[3pt]
\FitMath{|e_{K,s,D}(w)|\le A e^{-\kappa m(D)},} \\[3pt]
\FitMath{\kappa>\kappa_{\mathrm{ent}}.}
\end{gather*}

Finally suppose $e_{K,s,D}$ depends on $w$ only through its values in $D^+$. These statements are the local-energy hypothesis. They include analyticity, an entropy gap, and geometric control of the dependence region.

\paragraph{Proposition.}
Let $n\ge2$ and suppose the supports of $f_1,\ldots,f_n$ have pairwise distance at least $d>0$. Under the local-energy hypothesis,

\begin{gather*}
\FitMath{\left|\partial_{z_1}\cdots\partial_{z_n}
 \sum_{s<K}\sum_D e_{K,s,D}(w)\bigg|_{z=0}\right|
 \le CAV\max\{1,d^{-4}\}
       (n/r)^n\prod_i\|f_i\|_*.}
\end{gather*}

The constant is independent of $K$; it can depend on $L,C_0,C_1$ and the entropy gap.

\begin{proof}
If $D^+$ misses one support, the energy is independent of the corresponding variable, and the mixed derivative vanishes. Every contributing domain therefore meets both the first and second supports. Consequently

\begin{gather*}
\FitMath{m(D)\ge d/(C_0a_s)-1.}
\end{gather*}

Cauchy's formula with radii approaching $r/(n\|f_i\|_*)$ bounds each derivative by

\begin{gather*}
\FitMath{A(n/r)^n\prod_i\|f_i\|_*\,e^{-\kappa m(D)}.}
\end{gather*}

A zero test function makes the assertion immediate, so these radii may be assumed finite. The entropy estimate and a geometric series in $m$ bound the sum at scale $s$ by a constant times

\begin{gather*}
\FitMath{AVa_s^{-4}\exp(-c d/a_s),}
\end{gather*}

with a larger harmless constant when $d/a_s$ is bounded. It remains to estimate $\sum_{s\ge0}L^{4s}e^{-cdL^s}$.

For $0<d\le1$, let $s_*$ be the largest integer with $dL^{s_*}\le1$. The sum over $s\le s_*$ is at most $CL^{4s_*}\le Cd^{-4}$. Writing $s=s_*+j$ in the remaining sum bounds it by

\begin{gather*}
\FitMath{Cd^{-4}\sum_{j\ge1}L^{4j}e^{-c'L^j},}
\end{gather*}

which converges by the ratio test. For $d>1$, replace $d$ by $1$ in the exponential. This proves the stated uniform bound. At each fixed cutoff the domain sums converge normally on smaller source balls, permitting termwise differentiation there. Only the separated mixed-derivative estimate is asserted uniformly as the number of scales grows; the undifferentiated energy sums need not converge in that limit.
\end{proof}

If in addition $\Gamma_K=e^{\mathcal E_K}h_K$, where $\mathcal E_K=\sum_{s,D}e_{K,s,D}$ and $|h_K|\le Q(V)$ with $h_K(0)=1$, the corollary above gives a cutoff-independent bound on every separated cumulant. This conclusion remains a finite-volume conclusion when $Q$ depends on $V$.

\subsubsection{What a comparison of cutoffs must prove}

Uniform boundedness gives compactness only after equicontinuity and a suitable countable test class have been supplied. Convergence along a prescribed whole sequence requires more. One sufficient condition is a summable comparison: for every fixed separated tuple,

\begin{gather*}
\FitMath{|S_{K+1}(f_1,\ldots,f_n)-S_K(f_1,\ldots,f_n)|
 \le C(f_1,\ldots,f_n)\epsilon_K,} \\[3pt]
\FitMath{\sum_K\epsilon_K<\infty.}
\end{gather*}

The triangle inequality then makes the sequence Cauchy. Its tail is bounded by $C\sum_{k\ge K}\epsilon_k$.

A local proof of this comparison can split the scales at $K/2$. At $s\le K/2$, it is enough to compare paired local energies with a bound

\begin{gather*}
\FitMath{|e_{K+1,s,D}-e_{K,s,D}|
 \le C(1+s)\epsilon_K e^{-\kappa m(D)}.}
\end{gather*}

At $s>K/2$, the individual energy bounds already give a tail controlled by $\sum_{s>K/2}L^{4s}e^{-cdL^s}$. This tail is summable in $K$. The factor $1+s$ in the paired bound is harmless because $\sum_s(1+s)L^{4s}e^{-cdL^s}<\infty$ for each $d>0$. The residual integrals must also be compared on a common zero-free source ball. If both residuals are within $1/2$ of $1$, then the logarithmic difference is bounded by twice the difference of the residuals, using the line segment between their values and the derivative $1/z$ of the logarithm. Cauchy's formula converts this bound into a cumulant comparison.

These are consequences of an actual paired estimate. A bound for two covariant constant-source lists does not automatically prove the estimate for position-dependent arrays. Similarly, ordinary holomorphy in an unweighted norm does not create geometric factors for the age of an inherited activity. Those factors must be present in a separately established weighted estimate or in the domain and bound of a weighted holomorphic extension.

\subsubsection{A conditional tuning lemma}

Let $F_K$ be a real differentiable renormalized-coupling coordinate on an interval $I_K$, and suppose

\begin{gather*}
\FitMath{2/3\le F_K'(\beta)\le4/3} \\[3pt]
\FitMath{(\beta\in I_K).}
\end{gather*}

If $F_K$ takes values on both sides of $x_*$, the intermediate value theorem gives a unique $\beta_K\in I_K$ with $F_K(\beta_K)=x_*$. Moreover,

\begin{gather*}
\FitMath{\tfrac23|\beta-\beta'|
 \le|F_K(\beta)-F_K(\beta')|
 \le\tfrac43|\beta-\beta'|.}
\end{gather*}

Both inequalities follow from the mean value theorem on the interval between the two points. The inverse coordinate is therefore Lipschitz with constant $3/2$ on its image. If $F_K(\beta)=\beta-\sum_{j\le K}b_j(\beta)$, the derivative hypothesis follows from $\sum_j\sup_{I_K}|b_j'|\le1/3$. We must establish the existence of the interval, the crossing of the target, and the derivative estimate through the construction.

The conclusions of this section concern separated correlation functionals. They do not, without additional estimates, provide extensions across coincident points, reflection-positive distributions on the full test-function algebra, rotation covariance, or a characterization of the resulting scalar observable sector as a Yang--Mills theory. Those are additional mathematical steps with their own hypotheses.

% END MODULE chapters/rg_audit_chapter.tex


% BEGIN MODULE chapters/weighted_extension_audit_v2.tex
\section{Weighted analytic domains and a local test for scale forgetting}
\label{ch:weighted-extension-audit}

A multiscale argument needs several analyticity statements with different meanings. Holomorphy of a finite-dimensional partition function is the weakest of them. Holomorphy of a renormalization step on a class of effective interactions is stronger. Geometric forgetting of old interactions is stronger still. The latter can follow from an analytic extension with appropriate domains and bounds, but the weights in that extension must themselves be justified.

There is an additional distinction for large-field regions. A small probability of retaining an old region can appear as a small factor in the bound on an analytic function. It does not in general permit the analytic input radius to grow by the reciprocal of that probability. This section makes these distinctions precise, reconstructs a suitable input space, and isolates local inequalities that can be tested before undertaking a full continuum reconstruction. It does not assert that the required inequalities are false for the Yang--Mills step.

\subsection{The ordinary input space}

Fix a cutoff, a reference sequence of inverse couplings $\theta_0,\ldots,\theta_k$, a complete set of labels and tests, and hereditary complex domains $\mathcal U_j(X;r_{0,j},r_{1,j})$. A point of such a domain is a pair $(\mathbf U,\mathbf J)$; the current $\mathbf J$ is an independent complex variable. A regular background gives the special pair $(U,J(U))$. An identity verified at these special pairs is not automatically an identity at all pairs in the complex domain.

One convenient Banach space of pointed local families has the norm

\begin{gather*}
\FitMath{\|E\|_{j,\kappa}
 =\sup_{X,z}\,e^{\kappa d_j(X)}
    \sup_{\mathbf p\in\mathcal U_j(X;r_{0,j},r_{1,j})}
       |E(X,\mathbf p,z)|.}
\end{gather*}

When families have infinite-lattice companions, include their norms in the maximum. Gauge invariance, the prescribed translation covariance, and compatibility with restriction are closed linear conditions. Holomorphy is preserved under convergence in this supremum norm on the declared domains. These facts specify a closed space of analytic families once the domains themselves have been fixed. We must still prove the existence and uniform size of those domains.

A regular Euclidean member comprises a family $G$, its companion, and its scalar quadratic coefficient $b(G)$. It has covariant pointed totals on the prescribed regular backgrounds and carries its own counterterm. A spatially varying input can assign a different such member $G^{(j)}_z$ to every base point $z$. This does not require the array $z\mapsto G^{(j)}_z$ to be translation invariant. Its norm is $\sup_z\|G^{(j)}_z\|_{j,\kappa}$. Each assigned member must retain its own coefficient and the required pointed covariance; covariance of the sum of all base points is a weaker condition.

The complete input also contains a bare plaquette weight $u$, defect families $D^{(j)}$, remainders $H^{(j)}$, boundary terms $B$, and excess exponents $\Psi$ in stored large-field operations. A useful product box has inequalities of the following form:

\begin{gather*}
\FitMath{\begin{split}
 &\|u\|_\infty+
   \sum_{j\le k}\sup_z|b(G^{(j)}_z)-b_j^{\mathrm{ctr}}|\le C_u,
 \qquad \sup_z\|G^{(j)}_z\|_{j,\kappa}\le E_0,\\
 &\|D^{(j)}\|_{j,\kappa}\le L^{-(m-j)},
 \qquad\|H^{(j)}\|_{j,\kappa}\le C_H\theta_j^{-\kappa_0/2},
 \qquad m\in\{k,k+1\}.
\end{split}}
\end{gather*}

Here $b_j^{\mathrm{ctr}}$ denotes the coefficient of the chosen real centre. Boundary terms are measured against specified positive majorants $\overline B_{j,h}(X)$:

\begin{gather*}
\FitMath{\|B\|_{\mathrm{bdry}}
  =\sup_{j,h,X}|B_{j,h}(X)|/\overline B_{j,h}(X).}
\end{gather*}

For a stored history $\omega$ let $N(\omega)$ count its stored sites, each at its own scale. Excess exponents require both a bound on $|\Psi_\omega|/N(\omega)$ and a carrier-by-carrier oscillation bound. If $\Psi_\omega=\Psi_\omega^0+\sum_f\psi_f$, the latter can be expressed as

\begin{gather*}
\FitMath{\sup_{\omega,R}\overline M_{\omega}(R)^{-1}
       \sum_{f:R_f=R}\operatorname{osc}\psi_f<\infty.}
\end{gather*}

The carrier majorant and the history size are part of the definition, not consequences of holomorphy.

A source trajectory fits this format by freezing the source separately at each base point. Schematically,

\begin{gather*}
\FitMath{G^{(j)}_z=E^{(j),\zeta}\big|_{\zeta=w(p_z)},} \\[3pt]
\FitMath{D^{(j)}(X,\cdot,z)
   =E^{(j),w}(X,\cdot,z)-E^{(j),w(p_z)}(X,\cdot,z).}
\end{gather*}

For a Lipschitz source, a local freezing estimate can give
$\|D^{(j)}\|\le C L^{-(K-j)}\|w\|_{\mathrm{Lip}}$.
This smallness is a property of the source trajectory. An arbitrary defect in an ordinary box has only the size stipulated by that box.

This organization resolves a possible misconception. Global noncovariance does not by itself invalidate a construction with sources: Euclidean members can be frozen site by site. The remaining task is to prove the asserted bounds for the full array, including its noncovariant defects, its variable Wilson weight, and its stored-history operations. Scalar marginal cancellation for a globally covariant list alone does not complete that task.

\subsection{Three kinds of weights}

Let $e_j$ measure a difference of regular families and let $r_j$ measure a difference of remainders. The regular irrelevance weights have the form

\begin{gather*}
\FitMath{\omega_\nu\le C\nu^4L^{-\sigma(\nu-1)},} \\[3pt]
\FitMath{\sigma>0.}
\end{gather*}

They represent decay after the vacuum and the applicable marginal quadratic terms have been removed. In contrast, an unrenormalized remainder can enter with a squared ratio of analyticity radii,

\begin{gather*}
\FitMath{w_{n,j}=\left(\frac{r_{0,n}}{r_{0,j}}\right)^2
 \asymp\frac{\theta_j}{\theta_n}
    \left(\frac{\log\theta_n}{\log\theta_j}\right)^{2q}.}
\end{gather*}

These weights are not geometric forgetting factors. A small size of the remainder, rather than an extra symmetry cancellation, pays for their contribution.

For a typed tangent list, let $z_j$ and $b_j$ be nonnegative bounds on the absolute marginal amplitude and increment of the direction. Here $b_j$ is a size bound, distinct from the signed running increment in the flow; the weights and the size entries $e_j,r_j$ in the load below are also nonnegative. Define

\begin{gather*}
\FitMath{z_n^\sharp=\sum_{j\le n}L^{-(n-j)}(z_j+z_{j-1}),} \\[3pt]
\FitMath{z_{-1}=0,}
\end{gather*}

and a load

\begin{gather*}
\FitMath{p_n(\delta)=\sum_{j\le n}\omega_{n+1-j}e_j
  +\sum_{j\in\mathcal J_n}w_{n,j}r_j
  +\gamma_n\left(z_n^\sharp+
          \sum_{j\le n}L^{-c(n-j)}b_j\right),} \\[3pt]
\FitMath{\gamma_n\asymp\theta_n^{-1}(\log\theta_n)^C.}
\end{gather*}

The set $\mathcal J_n$ consists of remainder slots still present at level $n$ in the chosen filing convention. If a remainder is re-filed into a regular family, it must be removed from this set and counted in that family instead. A recent-level window can have a separate term without geometric age decay. This bookkeeping cannot be suppressed by changing notation.

The survival age of a large-field zone has a different role. For $0<q_a<1$ the desired integrated bound involves

\begin{gather*}
\FitMath{\Sigma_k(\delta)=\sum_{n\le k}(k-n+1)
        q_a^{(k-n-1)_+}p_n(\delta).}
\end{gather*}

The factor $k-n+1$ counts the steps at which the zone can be used. The survival factor $q_a^{(k-n-1)_+}$ belongs to an output estimate on the relevant histories. It need not, and generally cannot, be put into a single enlarged analytic input radius.

\subsection{What weighted holomorphy actually proves}

\begin{proposition}[A segment estimate]
Let $X,Y$ be complex Banach spaces, let $C\subset X$ be convex, and let $p$ be a continuous seminorm. Suppose $F$ is holomorphic on

\begin{gather*}
\FitMath{\mathcal T=\{a+h:a\in C,\,p(h)<\epsilon\},} \\[3pt]
\FitMath{\epsilon>0,}
\end{gather*}

and, for every $a\in C$ and $p(h)<\epsilon$, one has
$\|F(a+h)-F(a)\|_Y\le B$. Then, for $a,b\in C$,

\begin{gather*}
\FitMath{\|F(b)-F(a)\|_Y\le (B/\epsilon)p(b-a).}
\end{gather*}

\end{proposition}
\begin{proof}
Put $h=b-a$. At every $a_t=a+th$, $0\le t\le1$, the map
$z\mapsto F(a_t+zh)-F(a_t)$ is holomorphic on
$|z|<\epsilon/p(h)$ and has norm at most $B$. The Banach-valued Cauchy formula gives
$\|DF(a_t)h\|\le Bp(h)/\epsilon$. Integrating over $t$ proves the inequality. If $p(h)=0$, the line is entire and bounded in oscillation. Every bounded linear functional on $Y$ makes it a bounded scalar entire function, so Liouville's theorem and separation by the dual imply that it is constant.
\end{proof}

The seminorm in this proposition is an assumption. Applying it with
$p(h)=\sum_j\omega_{k+1-j}\|h_j\|$ produces the displayed weights because they already occur in the domain. It does not derive them from an ordinary product box.

A differential bound is a weaker sufficient input. If
$\|DF(a)h\|\le C p(h)$ holds on the segment joining two actual trajectories, integrating this inequality proves the same comparison without constructing a full weighted tube around every off-trajectory point. This is useful when formulating a minimal proof obligation.

\begin{proposition}[A marked slot]
Suppose $F(a,m)$ is jointly holomorphic when
$p(a-a_0)<r_a$ and $\|m\|<r_m$, is bounded there by $B$, and satisfies $F(a,0)=0$. For $p(a-a_0)\le r_a/2$ and $\|m\|\le r_m/2$,

\begin{gather*}
\FitMath{\|D_aF(a,m)h\|\le
       \frac{CB}{r_a r_m}\,p(h)\|m\|.}
\end{gather*}

\end{proposition}
\begin{proof}
Cauchy's formula first in the complex line $a+zh$, and then in the line $tm$, bounds the mixed derivative by
$CBp(h)\|m\|/(r_a r_m)$ on the reduced domains. Integrate that derivative from $m=0$ to $m$. The derivative in $a$ vanishes at $m=0$ because $F(a,0)$ is identically zero. The argument uses two complex lines and introduces no factor equal to the dimension of either input space.
\end{proof}

This explains how a small marked-slot size can multiply the sensitivity to other data. It does not establish the required radii, the marked-slot vanishing identity, or a uniform bound on the enlarged domain.

\subsection{Exact examples separating the assertions}

\begin{proposition}[Ordinary analyticity does not imply forgetting]
Fix $0<q<1$ and any positive numbers $s_k$. There are entire scalar functions $F_k$ bounded by $s_k$ on the unit polydisc of $\mathbb C^{k+1}$ for which no constant $C$ independent of $k$ satisfies

\begin{gather*}
\FitMath{|F_k(x)-F_k(y)|\le Cs_k
      \sum_{j=0}^k q^{k-j}|x_j-y_j|}
\end{gather*}

on the polydisc of radius $1/4$.
\end{proposition}
\begin{proof}
Take $F_k(x)=s_kx_0$, $x=0$, and $y=(1/4,0,\ldots,0)$. The inequality would give $1\le Cq^k$. This is impossible as $k\to\infty$. The same example is holomorphic on every weighted tube, but its supremum on a tube of fixed weighted radius grows like $s_kq^{-k}$. Thus the missing issue can be a uniform bound, even when the extension itself is entire.
\end{proof}

This example is a logical separation, not a model satisfying the gauge-theory power-counting identities. It shows exactly why analyticity and a small one-step factor alone do not establish those identities.

\begin{proposition}[Rarity does not enlarge the complex radius]
Let $0<q<1$, let $a$ be a nonnegative integer, and let an event have probability $q^a$. For $\delta>0$ its excess generating function can be

\begin{gather*}
\FitMath{R_a(\lambda)=q^a(e^{\lambda\delta}-1).}
\end{gather*}

On $|\lambda|\delta<c$ it has a bound $C_cq^a$, and $R_a'(0)=q^a\delta$. There is no age-independent bound $C_cq^a$ on the larger disc $|\lambda|q^a\delta<c$.
\end{proposition}
\begin{proof}
The first bound follows from $|e^z-1|\le e^{|z|}+1$. At the positive point $\lambda=cq^{-a}/(2\delta)$ the ratio $R_a(\lambda)/q^a$ equals $e^{cq^{-a}/2}-1$, which tends to infinity. The event model is a genuine two-point probability space, and its generating function is entire.
\end{proof}

The correct large-field strategy is therefore a bounded per-zone radius together with a rare-history factor in the bound. A demand for a bounded extension on a single tube measured only by the aged load would be stronger than what the argument needs and can fail even in this elementary example.

There is a separate reason for the Euclidean-member and marginal-slot conditions. Let

\begin{gather*}
\FitMath{H=\begin{pmatrix}0&1\\1&0\end{pmatrix},} \\[3pt]
\FitMath{M_\varepsilon=I+\varepsilon H.}
\end{gather*}

For real $|\varepsilon|<1$, take independent centered Gaussian vectors $X_1,\ldots,X_m$ with precision matrix $M_\varepsilon$. The normalized block average $m^{-1/2}\sum_iX_i$ has exactly the same covariance $M_\varepsilon^{-1}$, hence exactly the same precision matrix. The traceless quadratic perturbation $\varepsilon x_1x_2$ survives this blocking unchanged, although its scalar trace coefficient is zero. The Gaussian integrals are holomorphic and uniformly coercive on each fixed complex disc $|\varepsilon|<r<1$.

This example explains why a scalar counterterm cannot remove an arbitrary anisotropic quadratic perturbation. It does not contradict the use of a scalar counterterm in the Euclidean class. It also explains why a typed extension must supply all of its prescribed tensor counter-insertions before claiming that the remaining part is irrelevant.

Finally, identities on backgrounds do not determine an arbitrary extension to independent currents. The polynomial $P(z,j)=j-z^2$ vanishes on every pair $(z,z^2)$ and is nonzero elsewhere. An analogous cancellation known only on a graph of background pairs needs an additional argument before it can be used throughout a domain where the current is independent.

\subsection{Two local inequalities to establish}

The first test concerns the small-field potentials before exponentiation. Let $\dot V_{k,j}$ be the variation of the actual localized fluctuation potential caused by one regular family at level $j$. Include its own scalar counterterm, or the complete tensor counter-insertions for a typed family. A suitable test is
\begin{equation}
 \|\dot V_{k,j}[\delta E]\|_{\mathrm{pot}}
 \le C\varkappa_k\omega_{k+1-j}\|\delta E\|_{j,\kappa}.
 \label{eq:wext-potential-test}
\end{equation}
The supremum in the potential norm must range over the complex pairs and decoupling parameters actually used inside the fluctuation integral, not just real backgrounds. One possible norm is the least $M$ such that

\begin{gather*}
\FitMath{|\dot V(Y;\mathbf p,B)|\le
 M e^{-\lambda d_k(Y)}
   \left(1+M_0^{-4}\sum_{b\subset Y}|B(b)|^2\right).}
\end{gather*}

Here $M_0$ is the fixed block-size parameter. The complex field domains, the allowed fluctuation cut, and the enlarged decoupling discs must be specified. Defects and remainders have their own test with squared-radius weights in place of $\omega$. A variable Wilson weight is a marginal slot; it does not receive an irrelevance weight merely by being written as a local family.

For ordinary site-frozen Euclidean members, a regular-output dependence estimate already has this intended structure after the cluster expansion. The additional typed claim requires the same cancellation at the complex pairs in the integrand, with the full tensor removed. A background extraction formula alone is insufficient for that extension. Proving \eqref{eq:wext-potential-test} for these typed inputs is a concrete local task.

The second test concerns a single zone in a large-field operation. Fix a zone $n$ and a step $k'$ at which it is used, where $n\le k'\le k$. Let $\lambda_{n,k'}$ multiply only that use of a direction of load $p_n(\delta)$. Isolate the normally convergent subseries of histories affected by this dial, and denote the resulting output by $R_{k;n,k'}(\lambda)$. The required axis estimate has the form
\begin{equation}
 \sup_{|\lambda|p_n(\delta)<c}
   \|R_{k;n,k'}(\lambda)\|_{k,\kappa}
 \le C\epsilon_k q_a^{(k-n-1)_+}.
 \label{eq:wext-zone-test}
\end{equation}
The constants $c,C$ are independent of the levels, the history, and the positions of the sources. The source is allowed to be complex on the declared source ball. The factor $\epsilon_k$ is the small mass assigned to the output under consideration. For an unintegrated boundary term it must be replaced by that term's actual majorant; its residual dependence is subsequently paid under the relevant mass, rather than being assigned a false termwise survival bound.

Cauchy's formula applied to \eqref{eq:wext-zone-test} gives

\begin{gather*}
\FitMath{\|\partial_\lambda R_{k;n,k'}(0)\|_{k,\kappa}
 \le (C/c)\epsilon_k q_a^{(k-n-1)_+}p_n(\delta).}
\end{gather*}

Summing over the $k-n+1$ uses of that zone, and then over $n$, gives $(C/c)\epsilon_k\Sigma_k(\delta)$. This proves the differentiated age estimate from a precise bounded complex-mass statement. The proof needs no age-dependent increase of the analytic radius.

Use \eqref{eq:wext-zone-test} under the moduli, suprema and reserve
weights of the actual operation. Together with the needed first and
mixed derivative estimates on its trajectories, this controls the
perturbed integrand on rare histories. These direct axis estimates
provide the required hypothesis; a full off-trajectory analytic tube with the required uniform bounds
is an alternative sufficient input.

\subsection{Ordinary sources and the stronger comparison routes}

The ordinary product-box theorem, a weighted dependence estimate for its regular output, and the finer marked and zone-wise statements are different inputs. Coefficient readout uses the physical radii of the chain, as proved in Section~\ref{ch:extraction-radius-normalization}: fixed radii give a uniform norm, and running radii carry the inverse-square factor already recorded by the radius weights. Near-zero core inclusion is supplied by the stated tower bounds. The finer assertions concern actual marked and typed evaluations in complex windows, with their proper radius normalization. Contraction estimates derived from a linearized class retain all its hypotheses. The weighted and zone-wise estimates are additional inputs to the ordinary complex-box theorem.

For a prescribed whole-sequence continuum limit, the smallest useful replacement is an actual comparison of the two tuned trajectories: a step inequality, domination of the observable evaluations by the comparison distance, and a terminal-integral inequality with summable defects. These three estimates can be proved directly. Establishing a large weighted analytic class is one route to them, not a logical necessity.

Finite-cutoff source identities do not require this finer statement. Fixed-spacing DLR compactness and closedness of reflection positivity do not require it either. Subsequential separated-correlation compactness follows from its own cutoff- and volume-uniform source bounds. Those bounds remain substantial analytic requirements; their volume-free proof can use undifferentiated age-resolved rarity without proving the differentiated zone estimates above. Similarly, tuning at one fixed cutoff can follow from its own holomorphic slope and target-crossing estimates without establishing a comparison of consecutive cutoffs.

% Keep the closing statement and argument together.
\clearpage
The summable matched-cutoff comparison and the resulting whole-sequence tuned convergence use the stronger comparison route. The typed tangent contractions, scale-uniform memory estimates, and differentiated anisotropy method for rotations also require their corresponding finer inputs. Estimates for several insertions at nonzero source require joint bounds beyond a single first derivative at zero. An independent vanishing-flux theorem would provide a different route to rotations and should be assessed on its own hypotheses.

The potential test and the one-zone complex-mass test specify the uniform estimates on the actual inputs used by these stronger routes. Their stated domains, reserve weights and derivative bounds are retained in each application.

% END MODULE chapters/weighted_extension_audit_v2.tex


% BEGIN MODULE chapters/zone_reserve_test_v3.tex
\section{A reserve test for a single large-field zone}
\label{ch:zone-reserve-test}

The one-zone estimate uses the reserve lemma and the local
evaluation-map norm bound stated below. The reserve argument
and the direct Wilson-weight channel apply on regular zones.
For the formed-copy application, use the evaluation-map bound
for every actual formed copy of the direction, with the stated
complex perturbation and reserve weights.

\subsection{The integral before taking a logarithm}

Keep the reference sequence, tests, history labels, changes of variables and reference kernels fixed. A local operation is then represented on a disjoint union of stored-variable spaces by

\begin{gather*}
\FitMath{I(\lambda,p)=\sum_\omega\int_{S_\omega}
       e^{\Phi_\omega(\lambda,\xi;p)}\,\mu_\omega(d\xi),}
\end{gather*}

where each $\mu_\omega$ is positive and independent of $\lambda$. A harmless finite frame average can be included in these measures. The parameter $p$ ranges over the external complex-pair domain. A single zone-use parameter has an index $(n,k')$: it multiplies the chosen direction where the stored zone $n$ is used at step $k'$. Only histories containing that use are affected.

Before re-exponentiation, the substitution of an input list by a line of lists is affine in the exponent when its evaluation maps have been frozen. Its perturbation has the concrete form
\begin{equation}
 L_{\omega,n,k'}[\delta](\xi,p)
   =\sum_{e\in\mathcal E(\omega,n,k')}\epsilon_e\,
         \delta F_e\bigl(\kappa_{\omega,e}(\xi,p)\bigr)
       +\delta c_\omega(\xi).
 \label{eq:zr-local-evaluation-map}
\end{equation}
The index $e$ records the entry, its counter-insertions, its carrier, and its particular occurrence. The signs or fixed coefficients are $\epsilon_e$. The formed argument $\kappa_{\omega,e}$ must lie in that entry's original analytic domain. The last term records scalar parts, including the values subtracted at the trivial field and their contributions to the vacuum bookkeeping. Constants cannot be omitted from this map merely because the observable ultimately uses a connected derivative.

For this affine use one has
$\Phi_\omega(\lambda)=\Phi_\omega(0)+\lambda L_{\omega,n,k'}[\delta]$.
After local histories or polymer clusters have been summed, the resulting effective families are generally nonlinear in $\lambda$. They should be controlled by a convergent marked expansion, rather than by pretending that their logarithms remain affine. If the formed arguments themselves are changed with the input, their derivatives are additional terms in \eqref{eq:zr-local-evaluation-map}; they require their own domain and norm bounds.

\subsection{What an exponential reserve proves}

\begin{proposition}[A one-zone reserve lemma]
Let the fixed affected event $E_a$ consist of histories of age at least $a$
in a measurable disjoint union of history spaces with positive reference
measure. Assume that the integrands are measurable in the stored variables
and holomorphic in $\lambda$ for almost every stored configuration.
Suppose $W_\omega(\xi)\ge0$ is measurable, and assume a measurable majorant $M_\omega$ satisfies

\begin{gather*}
\FitMath{\sup_p e^{\operatorname{Re}\Phi_\omega(0,\xi;p)}\le M_\omega(\xi),} \\[3pt]
\FitMath{\sum_\omega\int_{E_a}M_\omega e^{cW_\omega}\,d\mu_\omega
       \le B\varepsilon q^a,}
\end{gather*}

where $c,C>0$, $0<q<1$ and the fixed load satisfies $p_n\ge0$. Put $R=c/(2Cp_n)$, interpreted as infinity
when $p_n=0$. Suppose the perturbation is defined and holomorphic on
$|\lambda|<R$, and throughout that disc
\begin{equation}
 \sup_p|\Phi_\omega(\lambda,\xi;p)-\Phi_\omega(0,\xi;p)|
       \le C|\lambda|p_n W_\omega(\xi).
 \label{eq:zr-phase-hypothesis}
\end{equation}
Then on $|\lambda|p_n<c/(2C)$ the integral over $E_a$ is bounded in modulus by $B\varepsilon q^a$, and is holomorphic. Its difference from its value at zero is bounded by $2B\varepsilon q^a$. Its first derivative at zero is at most

\begin{gather*}
\FitMath{(2CB/c)\varepsilon q^a p_n.}
\end{gather*}

The same conclusions hold with an additional insertion which is measurable
in the stored variables, holomorphic in $\lambda$ on the same disc, and
uniformly bounded there, multiplied by that uniform bound. A fixed bounded
insertion is a special case.
\end{proposition}
\begin{proof}
The modulus of the perturbed integrand is at most
$M_\omega e^{cW_\omega/2}$. The assumed positive majorant is integrable over the disjoint union of histories. Dominated integration, or Morera's theorem and Fubini, proves holomorphy on the open disc. The mass estimate follows from the same domination. The difference is bounded by the sum of the two masses. Cauchy's formula on circles of radius approaching $c/(2Cp_n)$ gives the derivative bound. If $p_n=0$, the hypothesis makes the perturbation independent of $\lambda$, and the derivative is zero.
\end{proof}

Finite upper integrals can replace the measurable majorants in this proposition: choose measurable majorants with integrals arbitrarily close to their upper integrals, apply the argument, and take the infimum. The supremum over external pairs must remain inside the history integral.

A reserve can also produce the age factor. If $W\ge a$ on $E_a$ and

\begin{gather*}
\FitMath{\sum_\omega\int M_\omega e^{(c_{\rm pert}+c_{\rm age})W}\,d\mu_\omega
       \le B\varepsilon,}
\end{gather*}

then a perturbation paid by $c_{\rm pert}W$ leaves the estimate
$B\varepsilon e^{-c_{\rm age}a}$. Thus a per-zone radius independent of the age is compatible with geometric forgetting. The perturbation budget and age budget must be allocated before the coupling ceiling is selected.

This argument is effective only after \eqref{eq:zr-phase-hypothesis} has been proved. An oscillation bound for the unperturbed factors does not establish it.

\subsection{A common reserve for many zone uses}

One-zone bounds do not by themselves give a common polydisc for an
arbitrarily large number of zone uses. The reserve also needs a capacity
bound. The following precise version distinguishes the two issues.

\begin{proposition}[A capacity criterion for joint parameters]
Let $\mathcal Z$ be a finite set of zone-use indices and let $i$ label
finitely many directions at each index. Their nonnegative loads are
$p_{z,i}$. Use the measurable stored spaces and fixed reference measures
of the preceding proposition, and assume that all phase coefficients are
measurable in the stored variables. Fix nonnegative measurable history charges $W_{\omega,z}$,
a finite family of nonnegative measurable action pieces $S_{\omega,b}$, and nonnegative
coefficients $a_{\omega,z,b}$, with

\begin{gather*}
\FitMath{\sum_{z\in\mathcal Z}W_{\omega,z}\le W_\omega,} \\[3pt]
\FitMath{\sum_{z\in\mathcal Z}a_{\omega,z,b}\le1} \\[3pt]
\FitMath{\hbox{for every }b.}
\end{gather*}

Suppose the affine perturbation is

\begin{gather*}
\FitMath{\Phi_\omega(\vec\lambda)-\Phi_\omega(0)
    =\sum_{z,i}\lambda_{z,i}L_{\omega,z,i},} \\[3pt]
\FitMath{|L_{\omega,z,i}|
 \le p_{z,i}\left(CW_{\omega,z}
                    +\sum_ba_{\omega,z,b}S_{\omega,b}\right).}
\end{gather*}

Assume a majorant at zero obeys, on a fixed event $E$,

\begin{gather*}
\FitMath{\sum_\omega\int_E M_\omega
       \exp\left(cW_\omega+\sigma\sum_bS_{\omega,b}\right)d\mu_\omega
     \le B_E,} \\[3pt]
\FitMath{\sup_p e^{\operatorname{Re}\Phi_\omega(0,\xi;p)}\le M_\omega,}
\end{gather*}

where $c,\sigma>0$. Choose
$0<r<\min\{c/(2C),\sigma/2\}$. Then the integral over $E$ is jointly
holomorphic and bounded by $B_E$ on the product of weighted balls

\begin{gather*}
\FitMath{\mathcal D_r=\left\{\vec\lambda:
       \sum_i|\lambda_{z,i}|p_{z,i}<r
                  \text{ for every }z\right\}.}
\end{gather*}

All statements use fixed reference measures, fixed events and fixed
evaluation maps; if additional parameter dependence is present, the
integrands must instead be given holomorphically on this domain with
the same phase bound.
\end{proposition}
\begin{proof}
For $\vec\lambda\in\mathcal D_r$, the triangle inequality and the two
capacity inequalities give

\begin{gather*}
\FitMath{|\Phi_\omega(\vec\lambda)-\Phi_\omega(0)|
 \le Cr\sum_zW_{\omega,z}
       +r\sum_b\left(\sum_za_{\omega,z,b}\right)S_{\omega,b}
 \le (c/2)W_\omega+(\sigma/2)\sum_bS_{\omega,b}.}
\end{gather*}

Thus a single integrable majorant works on the entire parameter domain,
independently of $|\mathcal Z|$ and the number of directions. Dominated
integration proves separate holomorphy and local boundedness, and the
finite-dimensional Hartogs theorem gives joint holomorphy. Equivalently,
one can apply Morera's theorem successively on each compact polydisc.
\end{proof}

If a designated parameter affects only an event $E_a$ whose majorant
has $B_{E_a}\le B\varepsilon q^a$, setting that parameter to zero
changes the whole integral by at most $2B\varepsilon q^a$, uniformly
in the other parameters on $\mathcal D_r$. At their common zero,
Cauchy's formula gives the first derivative bound
$B\varepsilon q^a p_{z,i}/r$. To see this without subtracting a large
unaffected part, differentiate only the integral on $E_a$; outside
$E_a$ the integrand is independent of the designated parameter.
For $m$ distinct directions at the same index, Cauchy's formula on
circles with radii tending to $r/(mp_{z,i})$ gives

\begin{gather*}
\FitMath{\left|\partial_{\lambda_{z,1}}\cdots
             \partial_{\lambda_{z,m}} I_E(0)\right|
   \le B_E(m/r)^m\prod_{i=1}^m p_{z,i}.}
\end{gather*}

Here zero loads mean that the corresponding affine perturbations vanish,
and those derivatives are zero. Higher mixed bounds require an appropriate
majorant for the history event on which all marked uses occur; they do not
acquire a product of rarity factors from holomorphy alone.

The capacity assumptions have content. Without them, let a single
history have $W=1$ and let each of $N$ indices contribute
$L_z=p_z=1$. Every separate perturbation is bounded by $|\lambda_z|W$,
but choosing all $\lambda_z=r/2$ gives a total exponent change $Nr/2$.
A fixed reserve in $W$ cannot pay this uniformly in $N$. A proof for the
actual formation must therefore assign all simultaneous uses to a common
site and action budget, in addition to bounding each separate direction.

\subsection{A channel which can be bounded directly}

For $U\in\mathrm{SL}(2,\mathbb C)$, the determinant identity

\begin{gather*}
\FitMath{\det(U-I)=\det U-\operatorname{tr}U+1
          =2-\operatorname{tr}U}
\end{gather*}

gives $a(U)=\frac12\det(U-I)$. With the operator matrix norm,
\begin{equation}
 |a(U)|\le\tfrac12\|U-I\|^2.
 \label{eq:zr-plaquette-square}
\end{equation}
For real $U\in\mathrm{SU}(2)$ this is an equality, since the two singular values of $U-I$ coincide.

Consider a regular zone expressed in its own units, with bare mesh $\eta_n$, containing $N_n$ unit cells. Suppose the configurations used in the actual substitutions satisfy
$\|U_p-I\|\le\alpha_n\eta_n^2$ on that zone. There are at most $C\eta_n^{-4}$ plaquettes per unit cell. Hence
\begin{equation}
 |A(\delta c,U)|
    \le C\alpha_n^2\|\delta c\|_\infty N_n.
 \label{eq:zr-regular-weight}
\end{equation}
Consequently, if the zone load includes
$p_n\ge\alpha_n^2\|\delta c\|_\infty$, this direct Wilson-weight contribution to \eqref{eq:zr-local-evaluation-map} is bounded by $Cp_nN_n$. If the reserve volume dominates the number of such stored cells, it satisfies \eqref{eq:zr-phase-hypothesis}.

The usual scale $\alpha_n^2\asymp\theta_n^{-1}(\log\theta_n)^{2q}$ has precisely the form of the coefficient multiplying a marginal amplitude in the load. This calculation uses only the stated curvature regularity at the actual formed arguments. It does not prove that every substitution has that regularity; that is a separate inclusion requirement.

\subsection{An unrestricted bare zone needs action reserve}

There is a genuine limitation to a reserve measured only by stored sites. On a bare region where the real field is unrestricted, a plaquette can have $U_p=-I$, and then $a_p=2$. A direction which changes a local Wilson weight by one has an exponent variation of order one at such a configuration. If its proposed load is

\begin{gather*}
\FitMath{p_0=\beta^{-1}(\log\beta)^C}
\end{gather*}

and the site reserve is of order $W=N\log\beta$ on a fixed finite region, then

\begin{gather*}
\FitMath{\frac{|\delta\Phi|}{p_0W}
     \ge c\,\frac{\beta}{(\log\beta)^{C+1}}}
\end{gather*}

for an appropriate fixed configuration and fixed $N$. This diverges as $\beta\to\infty$. Such configurations are obtained, for example, by setting one bond to $-I$ and the other bonds of a small patch to $I$. The number of altered plaquettes remains bounded. Thus \eqref{eq:zr-phase-hypothesis} with this site reserve cannot hold pointwise on the entire unrestricted bare-field domain.

This does not refute the complex mass estimate. The Wilson factor strongly suppresses those configurations. For a bare insertion $\mathfrak b$ satisfying a pointwise action bound

\begin{gather*}
\FitMath{|\mathfrak b|\le C_b\|k\|_\infty A_{\rm available}(U),}
\end{gather*}

one instead has

\begin{gather*}
\FitMath{|\lambda\mathfrak b|
  \le\frac{C_b|\lambda|\|k\|_\infty}{\beta}
           S_{\rm available}(U),} \\[3pt]
\FitMath{S_{\rm available}=\beta A_{\rm available}.}
\end{gather*}

At a radius $|\lambda|p_0<c$, with
$p_0\ge C_b\beta^{-1}(\log\beta)^C\|k\|_\infty$, this spends at most
$c(\log\beta)^{-C}$ of the available Wilson action. It is a fractional action estimate, not a site-volume estimate.

A correct reserve lemma can include both budgets. If the majorant has the stronger reserve
$e^{cV+\sigma S_{\rm available}}$ and

\begin{gather*}
\FitMath{|\Phi(\lambda)-\Phi(0)|
       \le C|\lambda|p_nV+\eta(\lambda)S_{\rm available},} \\[3pt]
\FitMath{C|\lambda|p_n\le c/2,} \\[3pt]
\FitMath{|\eta(\lambda)|\le\sigma/2,}
\end{gather*}

the mass proof above applies on the intersection of these parameter
domains. A derivative estimate proportional to $p_n$ requires this
intersection to contain a disc of radius $c' /p_n$ with a fixed $c'>0$.
One sufficient condition is $|\eta(\lambda)|\le C_{\rm act}|\lambda|p_n$;
otherwise the Cauchy radius is the smaller action-controlled radius.
To use this statement one must identify action genuinely left available
after the failed-test payments, gauge fixing, and other insertions.
It cannot be spent twice. For several active zone-use parameters, their
fractions must have a bounded sum on each shared action piece. The existence
of a site-volume reserve alone proves neither of these action-allocation facts.

\subsection{What the carrier bounds do not control}

Write a standing exponent as

\begin{gather*}
\FitMath{\Phi_\omega(\xi;p)=\Phi_\omega^0(\xi)
                 +\sum_f\phi_{\omega,f}(\xi;p|_{R_f}).}
\end{gather*}

The carrier condition bounds the oscillation of each factor in its external pair argument. Choose a permitted anchor $p_\omega^0(\xi)$ and split

\begin{gather*}
\FitMath{\phi_{\omega,f}(\xi;p)=
  \phi_{\omega,f}(\xi;p_\omega^0)
  +\{\phi_{\omega,f}(\xi;p)-\phi_{\omega,f}(\xi;p_\omega^0)\}.}
\end{gather*}

Oscillation controls the braced expression. It controls neither the scalar anchor value nor the dependence of either expression on an input-list parameter. Cauchy's estimate in the pair variables is not Cauchy's estimate in that new input parameter.

A simple exact example separates these statements. Let $\pi_a\downarrow0$ be a sequence of putative loads, let one stored history have volume $W_a\ge a$, and let its formed exponent be

\begin{gather*}
\FitMath{\Phi_a(x;p)=\log\varepsilon-(c+s+1)W_a+xW_a,} \\[3pt]
\FitMath{s>0.}
\end{gather*}

It is entire, independent of $p$, and its carrier oscillation is zero. At $x=0$ it has the prepaid mass bound
$e^{\Phi_a(0)+cW_a}\le\varepsilon e^{-sa}$. The same bound, with room, holds throughout a fixed ordinary disc $|x|\le1/2$. It fails on a disc $|x|\pi_a<c'$: at $x=c'/(2\pi_a)$ the positive exponent eventually overwhelms every fixed reserve coefficient.

This is a counterexample to deriving the enlarged-radius estimate from carrier oscillation and an ordinary reserve alone. It is not a counterexample to the Wilson operation. Even if the incoming direction has size $\pi_a$, the example can be produced by an abstract local formation map which multiplies that direction by $\pi_a^{-1}$. All input sizes are then small, while the output exponent has the displayed variation. Locality and holomorphy do not exclude this map; a uniform norm bound on formation must exclude it. Exact preservation of a represented integral, by itself, also supplies no such bound on a particular decomposition into exponents and vacuum factors.

\subsection{A sufficient structural bound on the actual formation map}

The needed additional statement is precise. For the actual index set and formed arguments in \eqref{eq:zr-local-evaluation-map}, prove
\begin{equation}
 |L_{\omega,n,k'}[\delta](\xi,p)|
      \le C p_n(\delta)V_\omega(\xi)
       \quad\text{for every admissible }(\xi,p),
 \label{eq:zr-formation-test}
\end{equation}
with a separate fractional-action term when the bare zone is unrestricted. The norm must include scalar anchors, not only oscillations. The constant must be independent of the age, the level, and the history. If a parameter changes the arguments as well, prove the corresponding derivative bound throughout its proposed disc and integrate it to obtain \eqref{eq:zr-phase-hypothesis}.

One verifiable route to \eqref{eq:zr-formation-test} is a carrier incidence estimate.
Let the charged stored cells be a finite set $\mathcal C_\omega(\xi)$
with $|\mathcal C_\omega(\xi)|\le V_\omega(\xi)$, and decompose the
scalar part as $\delta c_\omega=\sum_{c\in\mathcal C_\omega}\delta c_{\omega,c}$.
Assign each formed occurrence to these cells and give nonnegative weights
$w_{e,c}$ with $\sum_cw_{e,c}=1$. Suppose the actual evaluations stay in their original domains and

\begin{gather*}
\FitMath{\sum_e |\epsilon_e|w_{e,c}\,
       \|\delta F_e\|_{\mathrm{evaluation}}
       +|\delta c_{\omega,c}|\le Cp_n(\delta)}
\end{gather*}

at every stored cell $c$. Summing over the cells proves
\eqref{eq:zr-formation-test}. The evaluation norm here is the norm after the specified unit change and counter-insertions. Using a norm at an earlier level without those transformations can lose exactly the scale weights under investigation.

A condition on the sum of sizes of the original direction terms in a unit cell is helpful, but it is not yet this incidence statement. The proof must show that substitutions and repeated occurrences preserve that bound. The formation rules distinguish copies still anchored in a live component, copies whose last anchor is removed now, and copies inside older already formed terms. The first kind is not a small new remainder. The last kind must be summed with its own norm and history mass. Counting its occurrence does not permit its bound to be reused without payment. In particular a factor obtained only after summation of histories is not a pointwise bound on the unsummed exponent.


If ``size in a unit cell'' is defined to mean the sum of the evaluated sizes of all formed occurrences, including their scalar parts and the indicated unit changes, then the incidence estimate is already part of that definition. Under that stronger interpretation the reserve implication is proved. The substantive remaining task is to show that directions generated by the actual renormalization step meet that definition with the proposed load. A condition on the unformed entries cannot be silently replaced by this stronger condition.

\subsection{Passing the bound through a logarithmic sum}

Suppose that a finite polymer system has activities $z_X(\lambda)$
with a common Koteck\'y--Preiss majorant $r_X$, and that the corresponding
pinned bound is $e^{a_X}$. Choose the logarithm defined by the connected
expansion, equal to zero at the zero activity vector. Then a useful marked estimate is

\begin{gather*}
\FitMath{|\log Z(z(\lambda))-\log Z(z(0))|
    \le\sum_Xe^{a_X}|z_X(\lambda)-z_X(0)|,}
\end{gather*}

To prove it, join the two activity vectors by their straight segment. Every point on that segment has modulus at most $r_X$ coordinatewise. The absolutely convergent connected expansion gives
$|\partial_{z_X}\log Z|\le e^{a_X}$ by the pinned bound. Differentiate along the segment and integrate. Thus a summable one-zone age bound on the changed activities passes to the logarithm. No logarithm of an individual possibly zero history integral is required.

For a countably infinite system, apply the same estimate first to finite
restrictions of the same activity and incompatibility system.
If $\sum_X e^{a_X}|z_X(\lambda)-z_X(0)|<\infty$, the integrated
pinned series defines a relative cluster logarithm and satisfies this
bound, independently of the finite exhaustion by dominated convergence
of the rooted expansions. Writing two separate infinite-volume logarithms
requires convergence of both unmarked global series; a common local
majorant alone does not ensure that. This deduction has two inputs: the
activity majorant must remain valid on the complex parameter disc, and
the changed activities must satisfy the marked bound after the actual
formation and copy bookkeeping. A convergence theorem for polymers does
not prove those inputs. The exponential reserve can pay a controlled
perturbation of an integrand; it cannot manufacture control of the
perturbation map.

A recursive treatment of old copied terms can be sufficient instead of a single global phase estimate.
Fix $0<q_0<q<1$ and $C>0$. For example, suppose an already justified local operation estimate gives nonnegative bounds, for $j\ge n+1$,

\begin{gather*}
\FitMath{A_{j,n}\le Cq_0^{j-n}
       +C\sum_{n<e<j}\epsilon_e q_0^{j-e}A_{e,n},} \\[3pt]
\FitMath{0\le\epsilon_e\le\epsilon.}
\end{gather*}

Divide by $q_0^{j-n}$. A discrete Gronwall induction gives

\begin{gather*}
\FitMath{A_{j,n}\le Cq_0^{j-n}
       \prod_{n<e<j}(1+C\epsilon_e)
 \le C\{q_0(1+C\epsilon)\}^{j-n}.}
\end{gather*}

The induction follows by expanding the product, or using
$\prod_{e<j}(1+C\epsilon_e)=1+\sum_{e<j}C\epsilon_e\prod_{r<e}(1+C\epsilon_r)$.
Thus choosing $q_0\le q/(1+C\epsilon)$ preserves the desired rate $q$. No small sum of the $\epsilon_e$ is needed. This pays growth of output bounds after local integration. It does not pay an exponentially amplified perturbation inside an unsummed exponent. To apply this alternative one must establish the local parameterized operation bound at each new loss of an anchor and the stated copying recurrence. The sufficient formation test above is not claimed to be the only possible proof route.

The present test therefore proves the reserve implication and the direct regular-zone Wilson channel. It shows why an unrestricted bare channel requires available action, and isolates the additional uniform formation estimate for the remaining entries. The specified carrier oscillation, ordinary presentation bounds, and end-mark definitions do not alone establish that estimate. A proof of it for the actual local substitutions, with their scalar parts and copy incidence, would complete the proposed one-zone route rather than assume its conclusion.

% END MODULE chapters/zone_reserve_test_v3.tex


% BEGIN MODULE chapters/rg_objects_v2.tex
\section{The objects carried by the renormalisation argument}
\label{ch:rg-objects-v2}

This section fixes the types, domains and norms in the conditional
renormalisation argument. A presentation contains functions and measures,
as well as numerical estimates. In particular, an array of numbers
satisfying the estimates below does not constitute a presentation.
Statements about the existence of backgrounds, the inclusion of substituted
fields in complex domains, and the reproduction of the presentation by a
step must be proved separately.
No conclusion about those statements follows from the definitions here.

\subsection{Lattices, currents, and auxiliary geometric data}

The normalization in this section is that of $G=\mathrm{SU}(2)$.
Its Lie algebra is represented by traceless Hermitian matrices, with
$\langle A,B\rangle=\tfrac12\operatorname{tr}(AB)$ and complexification
$\mathfrak g^{\mathbb C}$. Fix an odd integer $L>1$. The level-$j$
lattice, denoted $\mathbb T^{(j)}$, consists of the sites retained after
$j$ blockings of a periodic bare lattice. Its bare spacing, expressed in
level-$j$ units, is $\eta_j=L^{-j}$. A bond configuration satisfies
$U(\bar b)=U(b)^{-1}$. Products along oriented paths are ordered in path
order. Positive coordinate bonds suffice to specify a configuration.
All suprema over bonds below use these bonds. In the left multiplication
chart, reversal is $J(\bar b)=-\operatorname{Ad}(U_b^{-1})J(b)$.
The complex configuration group is
$\mathrm{SL}(2,\mathbb C)$.

For a positively oriented plaquette $p$, based at its chosen corner,
write $U_p$ for its holonomy and set

\begin{gather*}
\FitMath{a_p(U)=1-\tfrac12\operatorname{tr}U_p,} \\[3pt]
\FitMath{A(U)=\sum_pa_p(U),} \\[3pt]
\FitMath{A(c,U)=\sum_pc(x(p))a_p(U).}
\end{gather*}

Here $x(p)$ is the plaquette midpoint. On real configurations this is
the real Wilson action; the same formula is its holomorphic extension.
For definiteness the current on a lattice of spacing $\eta$ is specified
by its variational normalization:

\begin{gather*}
\FitMath{\left.\frac{d}{dt}A(e^{i\eta tB}U)\right|_{t=0}
   =\eta^4\sum_b\langle J_\eta(U)(b),B(b)\rangle.}\tag{R1}
\end{gather*}

The sum is over positive bonds and $B$ is arbitrary. This identity
defines every component of $J_\eta$, also for complex $U$ by analytic
continuation. It fixes the transports and factors implicit in the
notation $D_U^*F_U$. One must use the plaquette divergence associated
with (R1) in that notation; the adjoint of the site gradient alone has
a different domain. At $U=\mathbf1$, differentiation in the potential coordinate
$U(t)=e^{i\eta tB}$ gives
$J'_\eta(\mathbf1)B=d_\eta^*d_\eta B$, where $d_\eta$ here is the
oriented plaquette difference on one-forms divided by $\eta$.
If the numerical mesh changes from $\eta$ to $\eta'=L^m\eta$ while the group
configuration stays fixed, potentials are multiplied by $L^{-m}$ and currents
by $L^{-3m}$. This follows directly from (R1) and the chart $e^{i\eta A}$,
and specifies the direction of the unit change.
The current variable of a complex pair, introduced below, is independent
of this particular function of $U$.

Fix powers of $L$ called $\mathfrak M,M_1',M$, a gap
$\rho=R_1\mathfrak M$, and constants $C_{\rm box},C_{\rm pot}$.
They are part of the auxiliary data, selected before domains or norms
are formed. The first two sizes concern backgrounds and walks; $M$
is the localization cube size. They are not interchangeable. The
partition $\pi_j$ consists of the centered, half-open $M$-cubes in
level-$j$ units. Boundary conventions are fixed once on every lattice.
For a cubical region $Y$, $Y^{\sim n}$ means addition of $n$ layers
in the partition currently named; $Y^{\sim-n}$ means removal of those
layers. A hull is the smallest union of cubes of the named partition
containing the set. These operations include their periodic versions.

The geometric data include two erosion assignments. Given a level-$p$
cubical subdomain $Z$ and an integer $0\le n\le p$, each assignment
returns a finite nested sequence of cubical site sets, cut at depth
$n$. In the first assignment the sets are unions of
$\mathfrak M$-cubes in their own units, with consecutive sets separated
by a gap greater than $\rho$ in the outer unit. In the second they are
unions of $M_1'$-cubes, with gap at least $M_1'$ in the inner unit.
An assignment specifies the actual sets, not just these inequalities.
It must be a single translation-compatible rule on all tori, commuting
with common lifts of non-wrapping regions. If more than one erosion
has these properties, either a particular assignment is supplied or the
core clauses below are imposed on all admissible assignments. A choice
cannot be recovered from a gap inequality alone.

Here is the determining set of a nested sequence. Work temporarily in
level-$k$ units, so the spacing of level $j$ is $\xi_j=L^{j-k}$.
Let $\Omega_0\supset\cdots\supset\Omega_k$ and
$\Omega_{k+1}=\varnothing$. The set $\Gamma_j$ consists of bonds of
$\mathbb T^{(j)}$ whose endpoints lie in $\Omega_0$, at least one
endpoint lies in $\Omega_j$, and neither lies in $\Omega_{j+1}$.
Thus inward bonds crossing into the next layer are excluded. Put

\begin{gather*}
\FitMath{\mathbf B=\bigcup_j\Gamma_j,} \\[3pt]
\FitMath{M_{\mathbf B}(U)
       =(\mathcal M^jU|_{\Gamma_j})_j.}\tag{R2}
\end{gather*}

The endpoints of these bonds are the nodes. They are the tops of the
blocks, or pillars, partitioning the determining region. A residual
gauge function equals $\mathbf1$ at all nodes. In particular, imposing
additional inward-bond constraints changes this problem.

The averaging map $\mathcal M$ used in (R2) is the symmetric contour
average. For group elements $V_i$ of small diameter their mean is the
unique nearby $V$ satisfying $\sum_i\log(V_iV^{-1})=0$. The contour
transport $U(y,x)$ is the mean over the shortest paths from $y$ to $x$
in its centered block that traverse the coordinate directions in one
of their orders. Set $U(x,y)=U(y,x)^{-1}$. For a coarse bond $c$,
$x'=x+(c_+-c_-)$, and straight transports $U[c]$ and $U([x,x'])$,

\begin{gather*}
\FitMath{\mathcal M(U)(c)=
 \exp\!\left(L^{-4}\sum_{x\in B(c_-)}
 \log\{U(c_-,x)U([x,x'])U(x',c_+)U[c]^{-1}\}\right)U[c].}\tag{R3}
\end{gather*}

All logarithms in regular neighborhoods are their near-identity
branches. Away from those neighborhoods a measurable extension of
the contour means and logarithms is auxiliary input. Exact identities
using that extension require its gauge covariance; small-field
analyticity concerns only the displayed branches. Existence of a
globally symmetric extension with every desired property is not
asserted merely by writing (R3).

For every determining sequence and real regular datum $\mathbf V$,
the background input is a regular critical orbit for
$\sum_{p\subset\Omega_0}a_p$ on
$\{U:M_{\mathbf B}(U)=\mathbf V\}$, together with a chosen
analytic continuation $\Phi_{\mathbf V}(B)$ at datum
$e^{iB}\mathbf V$. The representatives, gauge, analytic domain and
transport under local charts are included in this input. Uniqueness,
existence and uniform estimates for these maps are separate background
theorems. The word ``background'' does not substitute for them.

\subsection{Hereditary domains and their shifts}

A triple of level $j$ on $X$ is defined on the bare mesh $\eta_j$
when lengths are expressed in level-$j$ units. It is
$P=(U,A',\mathbf J)$, with $U$ real and $A',\mathbf J$ complex
Lie-algebra-valued bond forms. Its pair is
$\mathbf p(P)=(\mathbf U,\mathbf J)$, where
$\mathbf U=e^{i\eta_jA'}U$. A triple is a splitting of its pair.
Fix positive radii $a_0,a_1,a_2$. The simple conditions are

\begin{gather*}
\FitMath{\begin{aligned}
 &|\partial U-\mathbf1|,|\partial\mathbf U-\mathbf1|
       <a_0\eta_j^2,\qquad
   |A'|,|\nabla_UA'|<a_1,\qquad |\mathbf J|<a_2;\\
 &|A|,|\nabla A|<C_{\rm pot}Ma_0
       \quad\text{in the real gauges described below.}
 \end{aligned}}\tag{R4}
\end{gather*}

For each cubical box contained in $X$ with side at most
$C_{\rm box}LM$, such a real gauge must exist, with
$U=e^{i\eta_jA}$. The norms are maximum matrix norms in the named local units.
For a one-form component $A_\mu$ based at its initial site, fix the forward convention

\begin{gather*}
\FitMath{(\nabla_\nu^U A)_\mu(x)=\eta_j^{-1}
 \{\operatorname{Ad}(U_\nu(x))A_\mu(x+\eta_je_\nu)-A_\mu(x)\}.}
\end{gather*}

The maximum of these component norms defines $|\nabla_U A|$.
In the usual two-radius notation $a_2=a_0$.

For every subdomain $Z\subset X$ at every level $p\le j$, for both
erosion assignments, and every truncation depth $1\le n\le p$, form
its determining set $\mathbf B'$. Define its tower configuration by

\begin{gather*}
\FitMath{U_n(Z;P)=\Phi_{M_{\mathbf B'}(U)}
 \left(\frac1i\log\{M_{\mathbf B'}(\mathbf U)
                           M_{\mathbf B'}(U)^{-1}\}\right).}\tag{R5}
\end{gather*}

This expression is required to exist in the chosen continuation
domain. On $Z$ with two layers of the $M$-cubes of level $p-1$
removed, require

\begin{gather*}
\FitMath{\begin{aligned}
 |\partial U_n(Z;P)-\mathbf1|
      &<a_0L^{-(j-p)}\eta_p^2,\\
 |J(U_n(Z;P))|
      &<a_0L^{-(j-p)}L^{2(n-p)}.
 \end{aligned}}\tag{R6}
\end{gather*}

The current in the second line is formed in level-$n$ units.
Impose the same conditions with $\mathbf U$ replaced by $U$.
These are quantifiers over every allowed $Z,p,n$, not just the domain
on which an activity is eventually evaluated. The hereditary core
$\mathcal U_j^\circ(X;a_0,a_1,a_2)$ is the set of triples satisfying
(R4)--(R6). Its pair image is the analytic pair domain; holomorphy
means holomorphy in $(A',\mathbf J)$ on every open slice with $U$
fixed. The inequalities are strict. Requiring the tower maps to exist
in their open continuation domains is part of membership.

An activity on this core must have the tie: its values at any two
splittings of the same complex pair agree. Gauge invariance means
invariance under simultaneous transformation of both entries of the
pair; the current at a bond is based at its initial site and transforms
by conjugation there. The core itself is not saturated by complex
gauge transformations. Extension of a function to the union of complex
orbits requires equality on all intersections and holomorphic descent;
it is an additional statement when used.

Restriction $\rho_X$ discards bonds and variables outside $X$ without
extending them across the boundary. If an evaluation changes its numerical
mesh from $\eta_j$ to $\eta'=L^m\eta_j$, the coordinate change is explicitly

\begin{gather*}
\FitMath{(U,A',\mathbf J)\longmapsto
 (U,L^{-m}A',L^{-3m}\mathbf J),} \\[3pt]
\FitMath{m\ge0.}\tag{R7}
\end{gather*}

For a shift $(\mathbf H,\mathsf j)$ on the same lattice define

\begin{gather*}
\FitMath{P\dotplus(\mathbf H,\mathsf j)
  =(U,A'\boxplus\mathbf H,\mathbf J+\mathsf j),} \\[3pt]
\FitMath{e^{i\eta_j(A'\boxplus\mathbf H)}
       =e^{i\eta_j\mathbf H}e^{i\eta_jA'}.}\tag{R8}
\end{gather*}

The logarithm defining $\boxplus$ is the small branch; inclusion in
its domain is required. The independent current is shifted additively.
Replacing it by $J(e^{i\eta_j\mathbf H}\mathbf U)$ at every complex
argument would change the construction.

A decoupled response input on a cubical region $Y$ is a pair of maps
$\mathfrak h_Y(s;P,B')=(\mathbf H,\mathsf j)$, holomorphic in
complex cube parameters $|s(\Delta)|<e^{\bar\kappa}$, the complex
triple variables, and the datum $B'$. It vanishes at $B'=0$.
Its first component is the solution of the actual background equations
with the propagators decoupled by these parameters:

\begin{gather*}
\FitMath{\begin{gathered}
 \mathbf H=\mathbf a-\mathsf H(s)\mathsf D,
 \qquad \mathsf D=C_k(\mathbf a-\mathsf H(s)\mathsf D),\\
 \mathbf a=\mathsf H_1(s)B'+\mathbf A_1,
 \qquad \mathbf A_1+\mathfrak G(s)\mathsf W_s(\mathbf a)=0.
 \end{gathered}}\tag{R9}
\end{gather*}

Here $C_k$ is the nonlinear part of the straightened constraint map;
$\mathsf H,\mathsf H_1$ are the constrained linear minimizers;
$\mathfrak G$ is the inverse on the simultaneous Landau and
zero-average slice; and $\mathsf W_s$ is the nonlinear gradient
remainder. As data for (R9), these are specified linear maps between
the finite-dimensional bond and constraint spaces, and specified
holomorphic maps $C_k,\mathsf W_s$ on open balls of those spaces.
They include their walk coefficients and convergence domains.
The reduced current is a second component of this input, obtained
from the constraint multiplier. At one change of units its value is
$L^{-3}\mathbf J+\mathsf j$; the multiplier construction and its
uniform bounds are required independently. Bounds on first
derivatives of $\mathbf H$ alone do not define this component.

The response hypothesis requires realness, gauge covariance,
compatibility under restriction to wall-components, and transport by
common lifts. A wall-component is a component of cubes sharing
three-dimensional faces. When $s$ vanishes on separating cubes, the
equations split over the remaining components, and the response on a
component with no datum is zero. In local first-order norms its size
is at most $C e^{C\bar\kappa}\sup|B'|$; a response viewed $k-j$
levels below has the additional potential factor $L^{-(k-j)}$.
At $s=1$ and real background data its shifted pair is the old
background and its own current, up to the recorded node gauge factor.
Existence, uniqueness and these properties of (R9) are substantive
response inputs.

For a term $F$ on $X$ let $X_0$ be its hull in the cubes of the
decoupling scale. The interpolation is

\begin{gather*}
\FitMath{\Phi_Y[F](s,\nu)=F\big((\rho_XP)
                          \dotplus\nu\mathfrak h_Y(s;P,B')\big).}\tag{R10}
\end{gather*}

Every occurrence of (R10) requires a domain-inclusion hypothesis:
for all the indicated $s$ and $|\nu|<\varrho_F$, its triple lies in
the core of $F$. One sufficient quantitative hypothesis, with a fixed
constant $c_F>0$, is inclusion throughout

\begin{gather*}
\FitMath{|\nu|<c_F L^{k-j}a_{0,j}\delta(\theta_k)^{-1}
                   e^{-C_F\bar\kappa}=:\varrho_F>1.}\tag{R11}
\end{gather*}

In inner and removed substitutions separate inclusion bounds and
their own radii are required. Equation (R11) is not a consequence of
the formal shift law (R8).

For disjoint finite cube sets $\mathcal K,\mathcal N$ define

\begin{gather*}
\FitMath{F[\mathcal K\mid\mathcal N]
 =\sum_{\mathcal K'\subset\mathcal K}
   (-1)^{|\mathcal K|-|\mathcal K'|}
      F(\mathbf1_{\mathcal K'\cup\mathcal N}).}\tag{R12}
\end{gather*}

This equals the integral of $\prod_{\Delta\in\mathcal K}
\partial_{s(\Delta)}F(s+\mathbf1_{\mathcal N})$ over the unit cube,
by iterating the one-variable fundamental theorem of calculus.
Summing (R12) over subsets gives its exact resummation identity.
For the rooting assertion, use the wall-components of
$\mathcal K\cup\mathcal N\cup X_0$. If such a component contains a
cube of the differenced set $\mathcal K$ and misses the base $X_0$,
then locality and splitting make the function independent of that component's
parameters, so the alternating bracket is zero. An additional component made
only of fixed $\mathcal N$ cubes need not make a bracket zero. For
$\Phi_Y[F](s,\nu)-\Phi_Y[F](s,0)$ the expression vanishes if the
component of the base has no datum for the shifted response, since that
response is then zero on the base. These conclusions require the stated
splitting and support properties of the actual response. If $|\Phi_Y[F]|\le M_F$ on the stated polydisc,
Schwarz's lemma gives $|\Phi_Y[F](s,1)-\Phi_Y[F](s,0)|
\le2M_F/\varrho_F$. Cauchy's formula on the $s$ discs then bounds
the bracket of this excess by
$2M_F\varrho_F^{-1}(e^{\bar\kappa}-1)^{-|\mathcal K|}$.
For $\bar\kappa\ge1$ this is at most
$2M_F\varrho_F^{-1}e^{-(\bar\kappa-1)|\mathcal K|}$.
When $\mathcal K$ is nonempty, this also bounds the bracket of
$\Phi_Y[F](s,1)$ itself: its value at $\nu=0$ is the $s$-independent
constant $F(\rho_XP)$, which the bracket annihilates. For an empty differenced
set, an un-subtracted constant does not acquire $\varrho_F^{-1}$.
Thus the algebraic bracket bound is proved once the response and
inclusion hypotheses have been supplied.

\subsection{Pointed families, marginal coefficients, and list norms}

Let $\mathcal D_j$ be the connected unions of cubes of $\pi_j$.
Set $d_j(X)=M^{-1}$ times the length of a shortest tree in $X$
meeting each cube; it is zero for one cube. A pointed family is the
collection $E(X,\mathbf p,z)$ for $X\in\mathcal D_j$ and unit sites
$z\in X\cap\mathbb T^{(j)}$, with each term analytic on its core,
gauge invariant and tied. Its norm is

\begin{gather*}
\FitMath{\|E\|_{j,\kappa}=
 \sup_{\text{admissible tori and their companion}}
 \sup_{X,z\in X} e^{\kappa d_j(X)}
 \sup_{P\in\mathcal U_j^\circ(X)}|E(X,\mathbf p(P),z)|.}\tag{R13}
\end{gather*}

The supremum in $z$ is not a sum over base points. For an unpointed
family omit $z$ from this formula. Sums over domains require the
separate lattice-animal summability bound at the chosen decay rate.

A companion is a single family $\widehat E$ on $\mathbb Z^4$,
agreeing with every torus term on every common lift of a non-wrapping
domain with the necessary collar. For a fixed level, fixed frame prefix and fixed initial weight, the same
companion serves every cutoff on which those input data coincide. Its
existence is compatibility between tori, not a property
that can be tested on just one torus. Individual terms are required
to be covariant under the translations preserving $\pi_j$.
The pointed total has the stronger property

\begin{gather*}
\FitMath{\mathcal P_z^{\mathsf T}(W)=
    \sum_{X\ni z}E(X,\mathbf p_W,z),} \\[3pt]
\FitMath{\mathcal P_{rz}^{\mathsf T}(rW)=\mathcal P_z^{\mathsf T}(W)}\tag{R14}
\end{gather*}

for every lattice isometry $r$, including unit translations. Here

\begin{gather*}
\FitMath{\mathbf p_W=(U_j(W),J(U_j(W))).}
\end{gather*}

The covariance class consists
of the $W$ for which these backgrounds exist, every restriction lies
in the required core, and

\begin{gather*}
\FitMath{L^{2j}\max_p|\partial U_j(W)(p)-\mathbf1|
       <c\min\{\varepsilon(\theta_{j-1}),
                         \varepsilon(\theta_j)\}.}\tag{R15}
\end{gather*}

Both $W$ and $rW$ must be in this class. Invariance of the sum over
$z$ alone is weaker than (R14).

The site tent is
$h_z(x)=\prod_{\mu=1}^4\max\{1-|x_\mu-z_\mu|,0\}$.
Evaluate it at bond and plaquette midpoints. Then
$\sum_zh_z(x)=1$ and $h_{rz}(rx)=h_z(x)$. A bond-rooted term is
pointed by $E(X,\cdot,z)=\sum_bh_z(x_b)E(X,b;\cdot)$; its collar
must contain the supports of these tents. The pointed Wilson family
has terms $A(h_z,\mathbf U)$ on the cube of $z$ enlarged by one
layer, and zero elsewhere. This convention distributes a coupling
counterterm by the same base points as the regular family.

To define the coefficient, use the chart
$\mathsf e_X(A,J)=(e^{i\eta_jA},J)$ at $(\mathbf1,0)$.
Directions are pairs of complex bond forms, with norm
$\max\{\sup|A|,\sup|\nabla A|,\sup|J|\}$.
Let $\mathsf H_\infty$ be the derivative of the infinite-lattice
background map at its trivial datum, in its fixed gauge, and put

\begin{gather*}
\FitMath{\mathfrak h_\infty B=
       (\mathsf H_\infty B,d^*d\mathsf H_\infty B).}
\end{gather*}

This linear response is an input from the background theory; it must
include its domain and decay estimates. Choose an orthonormal color
basis $t_1,t_2,t_3$, and let $\delta_c^a$ be $t_a$ on the unit bond
$c$ and zero elsewhere. With Hessians taken in the displayed chart,

\begin{gather*}
\FitMath{\begin{aligned}
 \Pi^E(c,c';z)&=\tfrac13\sum_{a=1}^3\sum_{\widehat X\ni z}
   \operatorname{Hess}\widehat E(\widehat X,\cdot,z)
       [\mathfrak h_\infty\delta_c^a,
        \mathfrak h_\infty\delta_{c'}^a],\\
 \Pi^E(c,c')&=\sum_z\Pi^E(c,c';z)
                    =:\Pi^E_{\mu\nu}(c'_--c_-),\\
 b[\widehat E]&=\sum_{u\in\mathbb Z^4}
                   \Pi^E_{\mu\nu}(u)u_\mu u_\nu,
                    \qquad\mu\ne\nu.
 \end{aligned}}\tag{R16}
\end{gather*}

The coefficient is defined only when all these sums converge
absolutely, the middle expression depends just on directions and
displacement, and the last value is independent of the distinct
directions chosen. Gauge independence, the Ward identities and the
reduction of the second-moment tensor to this scalar require proof.
They are not incorporated by declaring an arbitrary Hessian to have
a coefficient. In particular, a claim that the Wilson tent family has
coefficient one requires the linear-response identities with the
normalization (R1).

A member is a triple $(E,\widehat E,e)$ satisfying (R13)--(R16),
with $e=b[\widehat E]$. Activity bounds and coefficient bounds use their
declared physical radii, including the inverse-square readout factor of
Section~\ref{ch:extraction-radius-normalization}. A list at step
$k+1$ contains such members for $j\le k$ and unpointed slot
families $H^{(j)}$ with no coupling counterterm. Its action is

\begin{gather*}
\FitMath{\begin{aligned}
 S[\mathfrak f](U)=&\sum_{j\le k}\left[
   \sum_z\sum_{X\ni z}
       \{E^{(j)}(X,U,z)-E^{(j)}(X,\mathbf1,z)\}
                         -e_jA(U)\right]\\
 &+\sum_{j,X}\{H^{(j)}(X,U)-H^{(j)}(X,\mathbf1)\}.
 \end{aligned}}\tag{R17}
\end{gather*}

Every evaluation uses the background and units appropriate to its
level. Slot families have companions, term covariance under the
localization translations, and invariant totals on the stated
background class, but no requirement of a pointed coefficient.
For $0<\alpha_{\rm H}<1$ their list size is

\begin{gather*}
\FitMath{N_{\mathfrak a}(\mathfrak f)=
 \sum_{j\le k}(k+2-j)^{3/2}L^{-\alpha_{\rm H}(k-j)}
                      \|E^{(j)}\|_{j,\kappa}
       +\sum_j\varrho_j^2\|H^{(j)}\|_{j,\kappa},} \\[3pt]
\FitMath{\varrho_j=\max_\iota
       \frac{(1+\beta_{\rm s})a_{\iota,k+1}}{a_{\iota,j}}.}\tag{R18}
\end{gather*}

All three radii, including the independent-current radius, enter this
maximum. Admissible radius systems require
$CMa_{0,j}<a_{1,j}$, $a_{1,j}^2\le Ca_{0,j}$,
$a_{\iota,j}\le(1+\beta_{\rm s})a_{\iota,k+1}
\le C(k+2-j)^{1/2}a_{\iota,j}$, and a fixed small upper bound on
the absolute radii. They also require
$\mathsf s_{\mathfrak a}:=\delta(\theta_k)/a_{0,k+1}\le c$.
The constants and the admissible systems are specified before taking
(R18). For the running system use
$a_{\iota,j}=\check\alpha_{\iota,j}
:=\min\{\alpha_{\iota,j-1},\alpha_{\iota,j}\}$ at $j\le k$ and
$a_{\iota,k+1}=\alpha_{\iota,k}$. For the absolute system use fixed
$\alpha_\iota$ at $j\le k$ and
$a_{\iota,k+1}=\alpha_\iota/(1+\beta_{\rm s})$.

\subsection{Frames, histories, structures, and stored variables}

A frame is a real sequence $\theta_0,\ldots,\theta_k\ge\gamma^{-2}$,
with $0<\gamma<e^{-1/2}$, such that, for $i\le j$,

\begin{gather*}
\FitMath{\underline\lambda(j-i)-c_5\le\theta_i-\theta_j
                \le\bar B(j-i),} \\[3pt]
\FitMath{\check R_j=R(\min_{i\le j}\theta_i),} \\[3pt]
\FitMath{\check R_j\le L\check R_{j+1},}\tag{R19}
\end{gather*}

where $R(y)$ is the smallest power of $L$ not below $(\log y)^r$.
Set $\bar g_j=\theta_j^{-1/2}$,
$p_0(g)=A_0(\log g^{-2})^{\mathfrak p_0}$,
$p_1(g)=(\log g^{-2})^{\mathfrak p_1}$,
$\varepsilon_j=\bar g_jp_0(\bar g_j)$,
$\delta_j=(A_1/A_0)\varepsilon_j$ and
$\alpha_{\iota,j}=C_\iota\theta_j^{-1/2}(\log\theta_j)^q$.
Here $\mathfrak p_0>\mathfrak p_1$, $A_1<A_0$, and all constants
are fixed. The waiting depth is $\mathsf N_j=\check R_j$.
The barred couplings belong only to the auxiliary frame. The actual running
couplings $x_j^{-1/2}$ are separate. The weights $x_j$ are separate entries, with
$x_0=\beta$ and $x_j=x_{j-1}-b_j$. A tube requires

\begin{gather*}
\FitMath{|\theta_j-x_j|\le\mathsf t_j,} \\[3pt]
\FitMath{\mathsf t_j\le c\theta_{j+1}
              (\log\theta_{j+1})^{-2q_{\rm t}},} \\[3pt]
\FitMath{q_{\rm t}\ge q,} \\[3pt]
\FitMath{2q_{\rm t}>q+\mathfrak p_0.}\tag{R20}
\end{gather*}

Its radii are non-increasing. At a terminal frame with no next
member the radius must instead be supplied explicitly; for a
one-member frame require
$\mathsf t_0\le2c\theta_0(\log\theta_0)^{-2q_{\rm t}}$.
All tests, regions and kernels are formed from $\theta$, not from
the unknown entries $x_j$.

A history is a finite record of every outcome of the tests of the
steps and of the operations inside large-field components, including
coordinate frames in the signed-permutation group $W_4$. The
formation data assign to each record its regions, stored bonds and
integration variables. In particular they give total tests
$\chi,1-\chi$ as Borel functions on the actual variable space,
with each value in $\{0,1\}$. A plaquette test is
$\mathbf1\{|\partial U_{j,\square}-\mathbf1|
<\varepsilon_j\eta_j^2\}$; a fluctuation test is
$\mathbf1\{|\mathcal A_j|<\delta_j\}$.
The local background map, the neighborhoods of each test, and the
collar and rim tests with their sharper thresholds are part of the
formation data. Listing the names of tests does not define them.

The label is $\mathcal S=(\Omega_j,\Lambda_j,S_j)_{j\le k}$ with
$T_\ell\supset\Omega_1\supset\Lambda_1\supset\cdots
\supset\Lambda_k$; $S_j\subset\Omega_j\setminus\Lambda_j$
records strips and failed fluctuation regions. Write
$Z_j=\Lambda_j^c$. The step creates unions of
$M\check R_j$-cubes, with gaps at least $cM\check R_j$ in local
units. Nests inside an operation may have gaps of only one $M$-cube
of the later level. The formation data must state the actual
enlargements and test neighborhoods; bounds described by a fixed
number of layers leave that number as a required datum.

An operation at level $t$ acts on the prescribed class of components
which are small, quiet for $\mathsf N_t$ steps, and satisfy the
waiting and quarantine conditions. It first integrates stored fields
in $\mathsf N_t$ stages, then tests extension of the exterior field,
integrates the inner box in an axial gauge with frame in $W_4$,
restores block variables through normalized reference laws, localizes
the exponent, and performs its last test. A healed component leaves
the label; an island passes the earlier tests and fails the last;
a drop-out fails after $n\le\mathsf N_t$ stages and keeps those
stages. These alternatives and every operation kernel must be given
as measurable maps and positive measures. Their exactness is a
separate equality of represented measures.

An operation can use its waiting window only when
$t_\flat=t-\mathsf N_t\ge0$. Before then it is not eligible; no negative-level
history is introduced. A drop-out $C$ records its former
regions $\Omega_l^\circ\cap C$ and its stage boxes $Z''_l(C)$.
For each integer $0\le\nu(C)\le n(C)$ define

\begin{gather*}
\FitMath{\Omega_l^{\boldsymbol\nu}\cap C=
 \begin{cases}
 C\setminus Z''_l(C),&t_\flat<l\le t_\flat+\nu(C),\\
 \Omega_l^\circ\cap C,&t_\flat+\nu(C)<l
                                      \le t_\flat+n(C).
 \end{cases}}\tag{R21}
\end{gather*}

Outside the affected levels retain the label regions. The choice
$\boldsymbol\nu=n$ is the determining sequence; it governs
backgrounds. The choice $\boldsymbol\nu=0$ is the action sequence;
it governs the sums and the coupling profiles. Standing marks record
the analogous $\Lambda_l,S_l$ data. Their waiting windows are
required to be disjoint at a point, and old marks disappear when
swallowed by a healed component or an island. These geometric
facts require verification for the actual operation.

A structure $\sigma=(\mathcal S,\boldsymbol\nu;t,\phi)$ also
records a phase. Its phase weight is $0$ at output, $1$ at the
step input, and $1-n/(2\mathsf N_t)$ after $n$ stages. A structure
precedes another if it counts no more completed stages and has
lower level, or the same level with larger phase weight. At an index
$0\le m\le t$ the radius coefficient is

\begin{gather*}
\FitMath{\begin{gathered}
 c_\sigma(x\mid m)=1-\beta_{\rm s}(1-2^{-(t-m)})
       +\varpi_\phi\mathsf g_t(m)
       -\beta_{\rm R}\sum_{l\in\operatorname{Re}_\sigma(x)}
                                      2^{-|m-l|},\\
 \mathsf g_t(m)=
 \begin{cases}\theta_{\rm s}\beta_{\rm s}2^{-(t-m)},&m<t,\\
 \beta_{\rm s},&m=t.
 \end{cases}
 \end{gathered}}\tag{R22}
\end{gather*}

Here $0<\theta_{\rm s}<1$ and
$\beta_{\rm R}=c_{\rm R}\beta_{\rm s}$ are fixed fractions;
$\operatorname{Re}_\sigma(x)$ consists exactly of the levels of
the first $\nu(C)$ stages of each drop-out containing $x$.
The layer-radius formation data give the original regularity layer
and the layer after each stage. Their nonnegative difference is
$e_\sigma(x)$. The lift is $L_\flat^{2e_\sigma(x)}$ for a fixed
$L_\flat\ge1$. The layer-$m$ radii are multiplied by this lift and,
except for the real-potential gauge bounds, by (R22). Unfinished structures
here have $t\ge1$: at layer zero use level-one radii; at the unfinished
top use level-$t-1$ radii. The initial $t=0$ base structure is separate
input with its own specified radii. The data must ensure $c\le c_\sigma\le1+\beta_{\rm s}$.

The space $\widetilde U^c_\sigma(X)$ imposes the simple, tower
and small-field-core clauses with those layer radii, with one
splitting, on every subdomain and every preceding structure. The
layer assignments, erosion rules, continuation maps and recorded
remade levels are thus indispensable parts of its definition.
The claims that restriction or a stage substitution maps one such
space into another are analytic inclusion hypotheses. They require
the specified radii and constants; a formal partial order by itself
does not prove them.

Stored fluctuation variables are real bond forms
$\mathcal A_l^{\rm s}(b)$ belonging to the kernels outside the
exponent. They transform by $\operatorname{Ad}u(b_-)$.
Near bonds are the bonds in collars, rims and the nonanalytic cutoff
zone; all other stored bonds are far. Fix
$\lambda_l^{\rm st}=(\log\theta_l)^{6r+1}>1$.
A stored-field function is Borel in the near variables and is the
value at $\mu_b=1$ of a function jointly analytic in the pair, the
far variables, and switches $|\mu_b|<\lambda_l^{\rm st}$.
It is independent of a near variable when its switch is zero, and
independent of that switch when the variable is zero. For far
variables the analytic ball has radius $\lambda_l^{\rm st}$ times
their specified real range. That range is part of the kernel data.
The pair tie is required at every switch value.

\subsection{Presentations, standing operations, and their bounds}

A real presentation consists of the frame and formation data above,
weights $x_j$, regular members
$(E^{(j),\mathrm u},E^{(j),\mathrm m},\widehat E^{(j)},
b_j^{\mathrm u},b_j^{\mathrm m})$, remainder families
$(R^{(j)},\widehat R^{(j)})$, boundary functions
$B_{\mathcal S}(j,X,o)$, and the standing-operation exponents below.
Here $E=E^{\mathrm u}+E^{\mathrm m}$,
$b_j=b_j^{\mathrm u}+b_j^{\mathrm m}=x_{j-1}-x_j
=b[\widehat E^{(j)}]$. Boundary origins $o$ range over five
types: the step, the last layer of a regular family, a step remainder
beside an island, a stage, and exponentiation. A slot is the tuple
$(j,X,o)$ together with the structure-domain on which the function
was formed. It retains that domain at later steps.

For a finite exempt family $\mathsf H$ of nonempty regions, define
$d^\natural_{j,\mathsf H}(X)$ by minimizing $M^{-1}$ times tree
length also over subfamilies $\mathsf F\subset\mathsf H$ contained
in $X$: the tree meets each cube outside $\bigcup\mathsf F$ and
each member of $\mathsf F$, but need not fill those members.
For a label use the level-$j$ hulls of its standing islands as
$\mathsf H$. Denote the result by $d_j^\natural(\mathcal S;X)$.
It is at most $d_j(X)$ and asserts no decay across an exempt island.
The real class requires

\begin{gather*}
\FitMath{\|E^{(j)}\|_{j,\kappa}\le E_0,} \\[3pt]
\FitMath{\|R^{(j)}\|_{j,\kappa}\le \bar g_j^{\kappa_0},} \\[3pt]
\FitMath{|B_{\mathcal S}(j,X,o)|\le\mathsf B
                e^{-\kappa d_j^\natural(\mathcal S;X)}.}\tag{R23}
\end{gather*}

Boundary functions have the stored-field property. All entries are
real on real arguments and invariant under the fixed simultaneous
action on pairs and stored data. For regular families (R23) is
evaluated at $\check\alpha_{\iota,j}$; for the unmarked family
also require a member of size at most $E_0$ on fixed absolute radii.
The whole-lattice step is additional formed data: its positive tame
kernel is $\mathbf T^{\rm sf}_{j-1}[\theta_{j-1}]$, with its
non-pivot bond coordinates and sharp fluctuation box. For a list
action $S$ its step functional is

\begin{gather*}
\FitMath{\begin{split}
 \mathcal E^{(j)}[\theta;x;S](W)
   ={}&\log\big(\mathbf T^{\rm sf}_{j-1}[\theta]
                   e^{-xA(U_{j-1})+S(U_{j-1})}\big)(W)\\
    &+xA(U_j(W))-S(U_j(W)).
 \end{split}}
\end{gather*}

The kernel, the fibre coordinates, the normalization and the
near-real logarithm branch must be given for this definition.
The localization rule producing a pointed family from this
functional is supplied once, for every torus and list; it is not
determined by its total. The identification clause requires
$E^{(j)}$ to be exactly this localized output from the earlier
entries. Its unmarked version is the same output using only their
unmarked entries.

The formed data also contain pointed Gaussian families
$\mathbf N_j^\circ$ and $\mathbf G_{j,\tau}$ on those absolute
cores, the first localizing the Gaussian logarithmic normalizer
and its trivial-field local numbers, and the second localizing
the change of its quadratic form when $x/\theta=\tau$.
Their finite-dimensional covariance matrices, localization maps and
logarithm branches are typed input data, so the following conditions
have definite meanings. Set
$b_j^0=b[\mathbf N_j^\circ]$ and
$\tau_{j-1}=x_{j-1}/\theta_{j-1}$.
Require, for fixed constants $C_{\rm m},C_{\rm u},C_b$,

\begin{gather*}
\FitMath{\begin{split}
 \|E^{(j),\mathrm u}-\mathbf N_j^\circ
                 -\mathbf G_{j,\tau_{j-1}}\|_{j,\kappa}
       &\le C_{\rm u}\bar g_{j-1}(\log\theta_{j-1})^{\mathfrak p_0},\\
 \|E^{(j),\mathrm m}\|_{j,\kappa}
       &\le C_{\rm m}\bar g_{j-1}^{\kappa_0-C_{\rm m}},\\
 |b_j^{\mathrm m}|&\le C_b \bar g_{j-1}^{\kappa_0-4},\qquad
 |b_j^{\mathrm u}-b_j^0|
       \le C_b \bar g_{j-1}(\log\theta_{j-1})^{\mathfrak p_0}.
 \end{split}}
\end{gather*}

No universality formula or asymptotic-freedom estimate for $b_j^0$
is implied by calling these families Gaussian. We must prove their
covariance construction and coefficient estimates separately.

A standing operation at an island $(t,Y)$ has a countable or finite
local-history index $\omega$, a Borel subset $S_\omega$ of a finite
product of group variables, real bond variables and finite frame
sets, and a positive Borel measure $\mu_\omega$ on that set.
These sets and measures are formed independently of the entries.
Its exponents are Borel in $\xi\in S_\omega$, holomorphic in the
pair, real at real arguments, and have the form

\begin{gather*}
\FitMath{\mathbf T'_t(Y)F=\frac1{|W_4|}\sum_\omega
       \int_{S_\omega}F(\xi)
            e^{\Phi_\omega(\xi;\mathbf p)}\mu_\omega(d\xi),} \\[3pt]
\FitMath{\Phi_\omega=\Phi_\omega^\circ(\xi)
          +\sum_f\Phi_{\omega,f}(\xi;\mathbf p|_{R_f}).}\tag{R24}
\end{gather*}

Each carrier $R_f\subset Y^\sim$ supports both dependencies of
its factor. The finite factor index set, the measurable coordinate
projections onto stored variables in each carrier, and any nested
island variables are part of the data. No variable is integrated
twice: swallowed-operation variables retain their earlier labels
inside the product assigned to the swallowing component.

If $(s,D)$ ranges over the forest of nested stored components, set
$v(\omega)=\sum_{(s,D)}\log\theta_s\,\#_s(D)$, where
$\#_s$ counts sites in their own level. For a fixed creation-time
complex domain $\mathcal K_{t,Y}$ require

\begin{gather*}
\FitMath{\begin{aligned}
 &\frac1{|W_4|}\sum_\omega e^{c_{\rm r}v(\omega)}
    \int_{S_\omega}^{*}\sup_{\mathbf p\in\mathcal K_{t,Y}}
       |e^{\Phi_\omega(\xi;\mathbf p)}|\mu_\omega(d\xi)
       \le C e^{-2p_0(\bar g_t)/(1+\beta_0)},\\
 &\sum_{f:R_f=R}\operatorname{osc}\Phi_{\omega,f}
       \le\mathsf B_\Phi(\log\theta_t)^{2\mathfrak p_0}
                  e^{-\kappa d_t^\natural(\mathcal S;R)}.
 \end{aligned}}\tag{R25}
\end{gather*}

The first integral is an upper integral, allowing a nonmeasurable
supremum. The oscillation is the supremum of
$|\Phi_{\omega,f}(\xi;p)-\Phi_{\omega,f}(\xi;p')|$ over $\xi$
and both pairs in that fixed domain. The domain is supplied from
the data at creation and must contain every later evaluation;
this last inclusion is an independent requirement.
$c_{\rm r}\ge0$ is fixed before the coupling ceiling, and
$\beta_0>0$ satisfies
$p_0(\bar g_m)\ge(1+\beta_0)^{-1}p_0(\bar g_n)$ for $m<n$.
The supremum remains inside the integral and sum. There is no
asserted small factor in the size of $Y$. A drop-out has no new free
entry of type (R24): it keeps its partially formed operation, and
its carrying Wilson factor is included in its separate majorant.
Such a majorant must be supplied to use a drop-out estimate.

The effective exponent additionally requires finite index maps and
vacuum numbers, which are formed data. At label $\mathcal S$ the
map $\mathcal I_j^E(\mathcal S)$ is a subset of
$\{(z,X):z\in\Lambda_j\cap\mathbb T^{(j)},\ X\in\mathcal D_j,
z\in X\}$; it gives the actual enlarged-region rule for the regular
terms. The enlargement width and treatment of boundary-crossing
domains must be included in that map. The remainder index map is
$\mathcal I_j^R(\mathcal S)=
\{X\in\mathcal D_j:X\subset\Lambda_j^{\sim-1}\}$.
The boundary index map selects the slots meeting the one-layer
enlargement of $Z_j$. These maps use the action sequence in (R21).
Profiles $\varphi_j$ are supplied real functions on plaquette
midpoints, equal to zero outside their prescribed small-field
regions and to one on their prescribed interior regions. The
interpolation rule in the intervening layers is part of the data.
With $x_k(p)=\beta-\sum_{j\le k}b_j\varphi_j(p)$, the exponent is

\begin{gather*}
\FitMath{\begin{split}
 A_{k,\mathcal S}(U)={}&-A(x_k(\cdot),U)
    +\sum_{j,(z,X)\in\mathcal I_j^E(\mathcal S)}
       \{E^{(j)}(X,U,z)-E^{(j)}(X,\mathbf1,z)\}\\
 &-\sum_{j\le k}b_jA(\varphi_j,U)
    +\sum_{j,X\in\mathcal I_j^R(\mathcal S)}
       \{R^{(j)}(X,U)-R^{(j)}(X,\mathbf1)\}\\
 &+\sum_{(j,X,o)\text{ in boundary index map}}
          B_{\mathcal S}(j,X,o;U,\text{stored data})
                  -\mathcal E_{k,\mathcal S}.
 \end{split}}
\end{gather*}

The real scalar $\mathcal E_{k,\mathcal S}$ is supplied with its
vacuum bookkeeping rule, including the stage numbers retained by
drop-outs. The pointwise identity
$x_k(p)+\sum_{j\le k}b_j\varphi_j(p)=\beta$ is exact. It does not
authorize changing the index maps or the vacuum rule. Thus these
objects have explicit mathematical types, even where their actual
formation and their uniform estimates remain upstream input.

Erase every healed component together with its ancestral cluster
from a history $h$; the result is its standing part $h^\flat$.
Erasure being an admissible history with the same label, and the
hard-core decomposition used to sum healed clusters, are structural
requirements on the formation data. They are not assumed merely
from the word ``erase.'' The never-large standing part has
$\Omega_j=\Lambda_j=T_\ell$. The presentation gives finite positive
measures $\widehat\tau_{h^\flat}$ on $V_k$, stored fields and
coordinate frames, with marginals $\tau_{h^\flat}$ on $V_k$.
It also retains the individual-history measures $\widehat\tau_h$
produced by the exact kernels. Their identities are

\begin{gather*}
\FitMath{\rho_k(V_k)dV_k=\sum_{h^\flat}\tau_{h^\flat}
             =\sum_h\tau_h,} \\[3pt]
\FitMath{\sum_h\tau_h(\mathcal V_k)=Z.}\tag{R26}
\end{gather*}

The first equality uses the summed healed operations; the exponent
of an individual $h$ need not have the effective-action form.

The reading at $X$ records the standing regions, outcomes, marks,
stored test bonds and frames within the hull of $X$ enlarged by one
level-$l$ cell at every scale $l$. A chart is a translation by
top-level cells, a common lift between tori, or their composite.
Coherence requires a chart matching readings to carry the terms,
carriers, Borel sets and measures to their counterparts. Factorwise,

\begin{gather*}
\FitMath{\Phi_{\iota\omega,\iota f}(\iota\xi;\iota p)
                       =\Phi_{\omega,f}(\xi;p).}\tag{R27}
\end{gather*}

Unit translations not preserving the partitions are covered by the
pointed-total requirement (R14), not by this chart definition.
The simultaneous lattice symmetries preserving level-$j$ cells
must also transport the measures with their variables.

For a datum at $a$ just below a pillar top $y$, let
$\tau[U](a)=(\mathcal M^{m-1}U)(y,a)$ be its contour transport.
An adapted background satisfies $\tau[U](a)=\mathbf1$ at every
stored datum site. To define $\widehat\tau$ require that all
backgrounds exist on its support, every argument belongs to its
entry's domain, and the exponential has one value at all adapted
representatives. The last clause is independent of simultaneous
gauge invariance: an orbit transformation need not fix every datum
site. These are part of class membership, not formal consequences
of the norm bounds.

Finally let $\mathcal Z_k(h)$ be the live region after healing.
For a live component $X$, define its lineage downwards by
$\mathsf Z_k(X)=X$ and by taking at level $j$ all live components
contained in $\mathsf Z_{j+1}(X)$. Write $|C|_j$ for its number of
level-$j$ cells and $d'_j(W)$ for tree length in those cells.
The created sets in a lineage are the components $W$ of the step's
creation set there; operation marks do not enter its deposit.
With one fixed deposit constant $c>0$, set

\begin{gather*}
\FitMath{\begin{aligned}
 \operatorname{Vol}_k(X;h)&=
       \sum_{(j,C)\text{ in its lineage}}
                       2M^4\check R_j^5|C|_j,\\
 \operatorname{Dep}_k(X;h)&=
       \sum_{(j,W)\text{ created in its lineage}}
                       c p_0(\bar g_j)^2(d'_j(W)+1),\\
 \Theta^{(k)}(h)&=\sum_{X\text{ live at }k}
       \{w\operatorname{Dep}_k(X;h)
                          +c_a\operatorname{Vol}_k(X;h)\},
                 \quad 0\le w<1,\quad c_a\ge0.
 \end{aligned}}\tag{R28}
\end{gather*}

The end mark is $\Theta^{(K)}$. Label marks and this numerical
weight are different objects. The prepaid-stability input is
$\sum_he^{-E}\tau_h(\mathcal V_k)e^{\Theta^{(k)}(h)}
\le e^{E_+\ell^4}$, with the stipulated normalization and dials.
Its use requires the creation-deposit and operation-majorant bounds;
(R28) defines the numerical weight.

Thus the class $\mathfrak C_k(\theta)$ consists of the complete
typed presentations satisfying (R13)--(R28), their structural and
positivity clauses, and the specified analytic domains. The induction
uses the step-closure input: formation of the next presentation,
preservation of the same constants and margins, and equality of the
represented measures under its kernels. The background, response,
substitution, island-mass, drop-out-majorant and closure estimates
are the supporting analytic estimates used at this step, which we must
prove.

% END MODULE chapters/rg_objects_v2.tex


% BEGIN MODULE chapters/extraction_radius_audit_v3.tex
\section{Field coordinates and the radius of quadratic extraction}
\label{ch:extraction-radius-normalization}

The quadratic coefficient of an effective interaction is defined in the
physical block field. The field used to normalize a fluctuation Gaussian
is a different coordinate. Keeping these coordinates distinct explains
both the fixed radii of the unmarked chain and the squared radius ratios
in the marked chain. The required inclusion near the trivial field follows
from the background and tower estimates already used to define the
analytic field domains. It is not an additional loss caused by the number
of blocking steps.

All constants below may depend on the fixed blocking and localization
parameters. They are independent of the level, torus size and number of
earlier levels. The non-Abelian background and tower estimates are
explicit hypotheses in the statements that use them.

\subsection{Physical fields, Gaussian fields and the coefficient}

At level $j$, write the configuration in the trivial-field chart as

\begin{gather*}
\FitMath{\mathbf U=e^{i\eta_j A'},} \\[3pt]
\FitMath{v=(A',J),} \\[3pt]
\FitMath{\|v\|_{1,j}=\max\{\|A'\|_\infty,
                    \|\nabla_j A'\|_\infty,\|J\|_\infty\}.}\tag{ER1}
\end{gather*}

Here $0<\eta_j\le1$ is the bare spacing in level-$j$ units and $\nabla_j$
is a difference quotient in those units. The current $J$ is an
independent analytic variable. In particular, $A'$ in this chart is a
physical potential; no coupling multiplies it by definition.

The block datum of the background is $V=e^{iB}$. Let $p_j(B)$ be its
configuration-current pair in a fixed gauge near zero. Derivatives of its
configuration entry are taken in the chart (ER1). Suppose
\begin{equation}
\FitMath{p_j(0)=(\mathbf1,0),\qquad
 Dp_j(0)B=\mathfrak h_jB
       =(H_jB,D^*D H_jB).}
 \tag{ER2}
\end{equation}
The datum $B$ is again physical. Its continuation radius $r_{\rm B}$
is a radius in $e^{iB}$, chosen independently of the coupling.
For a gauge-invariant local interaction $E$, the first derivative at
the trivial pair vanishes: invariance under constant gauge
transformations makes it an invariant linear functional on the
semisimple Lie algebra, and that functional is zero. Consequently

\begin{gather*}
\FitMath{D^2(E\circ p_j)(0)[B_1,B_2]
 =\operatorname{Hess}E(0)[\mathfrak h_jB_1,
                         \mathfrak h_jB_2].}\tag{ER3}
\end{gather*}

The possible second derivative of $p_j$ contributes nothing because it
is multiplied by $DE(0)$.

The polarisation kernel is the kernel of (ER3). Its second spatial
moments define the curvature tensor $T[E]$, and, in the invariant
scalar channel, the action coefficient $b(E)$. These definitions use
derivatives with respect to $B$. For the Wilson functional of the block
datum itself, use the normalization

\begin{gather*}
\FitMath{b(A_{\rm block})=1,}\tag{ER4}
\end{gather*}

whose quadratic part is $\frac12\|dB\|^2$. For a Wilson family evaluated
after the background map, the analogous normalization requires the
linear-response identities in that map. Likewise
$T[\mathsf A[c]]=c$ for a constant allowed curvature tensor $c$ is
the prescribed counter-insertion normalization. Neither identity for
pulled-back families is proved merely by the coordinate change below.

In a fluctuation integral one instead sets
$\mathcal A_j=g_j\mathcal A_j^{\rm s}$. If
$\widetilde E_j(B^{\rm s})=E(p_j(g_jB^{\rm s}))$, the chain rule gives

\begin{gather*}
\FitMath{D^2\widetilde E_j(0)=g_j^2D^2(E\circ p_j)(0),} \\[3pt]
\FitMath{T_{\rm s}[E]=g_j^2T[E],} \\[3pt]
\FitMath{b_{\rm s}(E)=g_j^2b(E).}\tag{ER5}
\end{gather*}

Thus a fixed probe in the Gaussian variable may fit inside a shrinking
physical domain. It extracts the rescaled coefficient in (ER5).
Recovering the coefficient of the physical action requires division by
$g_j^2$.

For example, impose the response-decay and localization hypotheses of
the amplitude proposition below, and suppose that the usable physical
radius is $a_j$ and the family norm is at most $M_j$ there. Its radius
in $B^{\rm s}$ is comparable to $a_j/g_j$. The Cauchy radius factor
in the amplitude bound then transforms as
\begin{equation}
\FitMath{|g_j^2 b(E)|\le C M_j(g_j/a_j)^2,
 \qquad |b(E)|\le C M_j a_j^{-2}.}
 \tag{ER6}
\end{equation}
The coupling cancels from the coordinate conversion. The same
calculation holds if the physical radius is written $g_j^\gamma$:
the canonical radius is $g_j^{\gamma-1}$, and the physical coefficient
bound remains $CM_jg_j^{-2\gamma}$. The change of variables creates no
extra loss and removes no loss already present in the physical radius.

\subsection{A near-zero ball lies in the hereditary core}

We state precisely the background information needed for this local
inclusion. A core at radii $(a_0,a_1)$ is described by the following
simple conditions and their tower counterparts. A triple is
$(U,A',J)$, with $U$ real and $\mathbf U=e^{i\eta_jA'}U$.
The real and complex curvatures are smaller than $a_0\eta_j^2$;
the real factor has the prescribed local regular gauges;
$|A'|,|\nabla_UA'|<a_1$; and $|J|<a_0$.
The tower counterparts require the corresponding curvature and current
bounds for every allowed restricted tower inside the domain.

Assume the following uniform tower estimate. If the simple conditions
hold at sufficiently small radii $(\epsilon_0,\epsilon_1)$, the tower
configurations exist holomorphically. Their curvature and current,
in the local units of their layer, satisfy the defining tower bounds
with radius

\begin{gather*}
\FitMath{K_{\rm T}(\epsilon_0+\epsilon_1^2).}\tag{ER7}
\end{gather*}

The estimate includes the prescribed change-of-unit factors and holds
for both tower rules, every restricted domain and every depth. This is
the substantive background estimate used here. It is stronger than
existence of a background for one real datum.

\begin{proposition}[Near-zero core inclusion]
Under (ER7), for sufficiently small $a_0,a_1$ there is a constant
$c>0$, independent of $j$, such that

\begin{gather*}
\FitMath{\|v\|_{1,j}<r_j,} \\[3pt]
\FitMath{r_j=c\min\{a_0,a_1,1\}} \\[3pt]
\FitMath{\Longrightarrow} \\[3pt]
\FitMath{(\mathbf1,A',J)\text{ belongs to the core at }(a_0,a_1).}\tag{ER8}
\end{gather*}

\end{proposition}
\begin{proof}
The plaquette product in the exponential chart obeys

\begin{gather*}
\FitMath{|\partial e^{i\eta_jA'}-\mathbf1|
 \le C\eta_j^2
       (\|\nabla_jA'\|_\infty+\|A'\|_\infty^2)
 \le C\eta_j^2(r_j+r_j^2).}\tag{ER9}
\end{gather*}

This follows by expanding a plaquette product: the linear term is the
discrete curl and each nonlinear term contains at least two bond
potentials. On a fixed small chart the remaining exponential series is
bounded by a uniform multiple of the quadratic term. In particular,
there is no inverse power of $\eta_j$ in (ER9). This is the reason
for including the first difference quotient in (ER1).

Choose

\begin{gather*}
\FitMath{\epsilon_0=\frac{a_0}{4K_{\rm T}},} \\[3pt]
\FitMath{\epsilon_1=\min\{a_1/2,\sqrt{a_0/(4K_{\rm T})}\}.}
\end{gather*}

Taking $c$ smaller if necessary ensures
$C(r_j+r_j^2)<\epsilon_0$ and $r_j<\min\{\epsilon_0,\epsilon_1\}$.
The real configuration is $\mathbf1$, so its curvature and local
gauge potential vanish. All simple conditions therefore hold at these
reduced radii. Moreover,
$K_{\rm T}(\epsilon_0+\epsilon_1^2)\le a_0/2$.
The tower estimate gives every hereditary clause at the desired
radii, and $|J|<r_j<a_0$. This proves (ER8).
\end{proof}

A bound only on $\|A'\|_\infty$ would be different: arbitrary
oscillation on the bare lattice can give
$\|\nabla_jA'\|_\infty$ of order
$\eta_j^{-1}\|A'\|_\infty$. The ball relevant to (ER8) is the
$C^1$ ball, not that weaker ball.

For completeness, suppose the continued background in (ER2) is
holomorphic for $\|B\|_\infty<r_{\rm B}$ and satisfies

\begin{gather*}
\FitMath{\|\mathcal H_j(B)\|_{C^1}
       +\|J(e^{i\eta_j\mathcal H_j(B)})\|_\infty
 \le C_{\rm B}\|B\|_\infty.}\tag{ER10}
\end{gather*}

Then its triple belongs to the core whenever

\begin{gather*}
\FitMath{\|B\|_\infty<
 \min\{r_{\rm B}/2,r_j/(2C_{\rm B})\}.}\tag{ER11}
\end{gather*}

The same conclusion for a first-order response uses the derivative of
(ER10). Equations (ER8)--(ER11) give a fixed datum radius when
$(a_0,a_1)$ are fixed, and a radius comparable to
$\min\{a_0,a_1\}$ when they run. The hypotheses are exactly the
background continuation, its uniform first-order bounds and the
hereditary tower estimate. No analyticity radius in the normalized
Gaussian variable is substituted for this physical datum radius.

\subsection{Cauchy's estimate, response decay and entropy}

Work first with the infinite-lattice companion. Its torus versions use
the uniform response bounds and the lift convergence supplied by the
background hypotheses. Consider pointed local families satisfying

\begin{gather*}
\FitMath{|E(X,p,z)|\le M e^{-\kappa d(X)}}
\end{gather*}

on the core, and assume that the near-zero ball there has usable
radius $a$. Assume also the uniform response decay

\begin{gather*}
\FitMath{\|\mathfrak h_j\delta_c\|_{1,j;X}
 \le C e^{-\delta\operatorname{dist}(c,X)}.}\tag{ER12}
\end{gather*}

The distances in this section are measured in the unit of level $j$.
For a fixed root cube, suppose the number of domains with tree length
between $m$ and $m+1$ is at most $C e^{\kappa_{\rm D}m}$, while the
number of pointed sites in such a domain is at most $C(m+1)^4$.
These are the usual localization entropy and volume bounds.

\begin{proposition}[The amplitude bound with its radius]
If $\kappa>\kappa_{\rm D}$ with a fixed positive margin,
the polarisation sums converge absolutely and

\begin{gather*}
\FitMath{|\Pi_{\mu\nu}(u)|\le C a^{-2}M e^{-\delta_1|u|},} \\[3pt]
\FitMath{|b(E)|+\|T[E]\|\le C a^{-2}M,}\tag{ER13}
\end{gather*}

whenever the scalar coefficient or the stated summed tensor is defined
by the prescribed covariance conditions. The constants are uniform in
the level and torus volume. A tensor at a fixed base point satisfies
the same bound before any isotropy assertion is made.
\end{proposition}
\begin{proof}
The two-variable Cauchy formula on a bidisc inside the ball gives

\begin{gather*}
\FitMath{|\operatorname{Hess}E(X,z)[v,w]|
 \le Ca^{-2}M e^{-\kappa d(X)}\|v\|_{1,j;X}\|w\|_{1,j;X}.}
\end{gather*}

Combining this with (ER12) bounds the contribution of one term by

\begin{gather*}
\FitMath{Ca^{-2}M e^{-\kappa d(X)}
 e^{-\delta(\operatorname{dist}(c,X)+\operatorname{dist}(c',X))}.}\tag{ER14}
\end{gather*}

If the localization cubes have fixed side $M_0$, then

\begin{gather*}
\FitMath{|c_--c'_-|\le\operatorname{dist}(c,X)
          +\operatorname{dist}(c',X)+C M_0(d(X)+1).}
\end{gather*}

A sufficiently small fixed fraction of the exponents in (ER14) supplies
$e^{-\delta_1|c_--c'_-|}$ for a fixed
$\delta_1>0$. Choose that fraction so that the remaining localization rate stays
strictly above $\kappa_{\rm D}$. The remainder is summable. Indeed, root a domain
at a cube closest to $c$. Summing its pointed sites costs at most a
polynomial in $m+1$, and summing its shapes costs
$e^{\kappa_{\rm D}m}$. These costs are absorbed by the remaining
$e^{-\kappa_{\rm rem}m}$, where $\kappa_{\rm rem}>\kappa_{\rm D}$.
Summing the root location is bounded by the
remaining exponential in its distance from $c$. This proves the first
inequality in (ER13) and absolute convergence.

The second moment defining $b(E)$ is bounded by
$Ca^{-2}M\sum_u|u|^2e^{-\delta_1|u|}$, a finite constant times
$a^{-2}M$. For a tensor at a fixed base point, the two moment factors
are bounded by polynomials in the domain diameter and the two
response distances. Those polynomials are absorbed by a further
fixed fraction of the same exponentials. This proves the tensor
bound. Identifying an invariant tensor with the scalar action
coefficient uses the separate Ward and lattice-symmetry identities;
it is not needed for the amplitude estimate.
\end{proof}

In particular, the extraction norm is uniform on fixed physical
radii. On the running radii

\begin{gather*}
\FitMath{a_j\asymp\theta_j^{-1/2}(\log\theta_j)^q,} \\[3pt]
\FitMath{q\ge0,}
\end{gather*}

the correct general bound has the factor
$a_j^{-2}\asymp\theta_j(\log\theta_j)^{-2q}$.
One can express uniformity instead as
$a_j^2|b(E)|\le C\|E\|$, provided that this normalization is retained
in all later estimates.

\subsection{What the squared radius ratios pay for}

Let a remainder born at level $j$ have size $r_j$ on physical radius
$a_j$ and have its vacuum value removed. There is no linear term.
Assume a gauge-invariant window estimate in which the quadratic curvature
factors each scale by $L^{-2(n-j)}$, while a unit volume at level $n$ contains
$L^{4(n-j)}$ sites of level $j$. After these geometric factors cancel,
evaluating the quadratic part on a background of radius $a_n$ costs

\begin{gather*}
\FitMath{w_{n,j}r_j,} \\[3pt]
\FitMath{w_{n,j}=(a_n/a_j)^2.}\tag{ER15}
\end{gather*}

This is the relative quadratic evaluation weight. Extracting the
physical coefficient from the resulting family at level $n$ costs
$a_n^{-2}$, so

\begin{gather*}
\FitMath{a_n^{-2}w_{n,j}r_j=a_j^{-2}r_j.}\tag{ER16}
\end{gather*}

Charging another factor $a_j^{-2}$ on top of the left side would
count the same shrinking input radius twice. Omitting the displayed
readout factor $a_n^{-2}$ would instead confuse an evaluated
potential with its physical quadratic coefficient.
The full nonlinear remainder estimate also requires its actual
backgrounds to lie in the original analytic domains and the separate
control of large domains. Equation (ER15) records the radius factor
in that estimate; it provides no continuation outside those domains.

For example, let $r_j\le C\theta_j^{-\kappa_0/2}$,
$\kappa_0>4$, and suppose
$\theta_j\ge\theta_n+\lambda(n-j)-C_0$ with $\lambda>0$ and a
sufficiently large common lower bound on $\theta_n$. Integral
comparison gives

\begin{gather*}
\FitMath{\begin{aligned}
 \sum_{j\le n}w_{n,j}r_j
 &\le C a_n^2
    \sum_{j\le n}\theta_j^{1-\kappa_0/2}
                         (\log\theta_j)^{-2q}
 \le C\theta_n^{1-\kappa_0/2},\\
 a_n^{-2}\sum_{j\le n}w_{n,j}r_j
 &\le C\theta_n^{2-\kappa_0/2}
                       (\log\theta_n)^{-2q}.
 \end{aligned}}\tag{ER17}
\end{gather*}

The first loss of one power of $\theta$ comes from summing the older
levels; the second comes from quadratic extraction. This is the
usual $g^{\kappa_0-4}$ bound, with its logarithmic improvement.
The two losses are already present in (ER17).

% Keep the closing scope paragraph together.
\clearpage
The unmarked chain uses fixed physical radii, so (ER13) preserves
its small interaction bound. The marked chain uses running radii,
and its stronger smallness pays the readout loss in (ER17). These
are the radius estimates used under the stated background, response
and entropy hypotheses. The complex-window, tensor-marginal,
weighted-sensitivity and large-field-zone derivative propositions
use their additional stated estimates.

% END MODULE chapters/extraction_radius_audit_v3.tex



% BEGIN MODULE chapters/01_orientation.tex
\section{The construction and the plan of proof}
\label{candidate:orientation}
\candidatebanner{A--H}

\noindent A dagger marks auxiliary assertions outside the central target
statements. Their hypotheses and dependencies are read in the same way as
those of the main argument.\par\medskip
\subsection*{0. Objects and proof strategy}

This section states the two theorems and explains the structure of the
proof. It identifies the methods adapted from Ba{\l}aban, the classical
tools, and the additional arguments, and shows how the multiscale
construction addresses the difficulties of a direct approach. Each later
step contributes to a clause of (0a).

\textbf{Objects.} We take $G=\mathrm{SU}(2)$ and $N=2$. On the torus $T_\ell$, with $\ell\in\mathfrak{L}$, we consider the following objects: the lattice $\Lambda_{K,\ell}$, whose spacing is $a_K=L^{-K}$; the configurations $U$ on it; the quantity $a_p(U)=1-\tfrac12\operatorname{Re}\operatorname{tr}U_p$; Wilson's measure; the smeared curvature density; and the connected correlation functions of that density:

\begin{gather*}
\FitMath{\mu_{K,\ell,\beta}(dU)=Z_{K,\ell,\beta}^{-1}\,e^{-\beta\sum_pa_p(U)}\,dU,} \\[3pt]
\FitMath{\mathcal{O}(f)=\sum_pf(x(p))\,a_p(U),} \\[3pt]
\FitMath{S^T_{n;K,\ell,\beta}(f_1,\dots,f_n)=\langle\mathcal{O}(f_1);\dots;\mathcal{O}(f_n)\rangle_{K,\ell,\beta}.}
\end{gather*}

In view of $a_p\approx\tfrac12a_K^4|F_U(p)|^2$, the sum $\mathcal{O}(f)$, which carries no prefactor involving the spacing, is a Riemann-sum approximation of $\tfrac12\int f\sum_{\mu<\nu}|F_{\mu\nu}|^2$. It is therefore the action density $\operatorname{tr}F\cdot F$. The test functions are restricted to separated tuples $(f_i)\in\mathfrak{S}_n(\mathbb{E})$, which are tuples of real $C^1_c$ functions with pairwise disjoint supports. The reason is given under ``The frame'' in (0a).

\textbf{The cutoff is tied to the coupling.} The renormalisation group of steps 1 and 2 assigns to $\beta=g_0^{-2}$ a sequence of running inverse couplings $x_0=\beta$, $x_k=x_{k-1}-b_k$, whose increments satisfy $b_k\le B_\sharp$. These increments are independent of the torus (3b), and also of the cutoff $K'\ge k$ used to carry out the steps (3d). We fix a small reference coupling $g$. As the unit of length we take the scale at which the flow reaches, for the first time, the reference window $g^{-2}\le x\le g^{-2}+B_\sharp$ belonging to $g$:

\begin{gather*}
\FitMath{K(\beta,g)=\min\{k:\,x_k(\beta)\le g^{-2}+B_\sharp\},} \\[3pt]
\FitMath{\mathcal{W}_K(g)=\{\beta\ge g^{-2}+B_\sharp:\,K(\beta,g)=K\};}
\end{gather*}

a pair $(K,\beta)$ such that $\beta\in\mathcal{W}_K(g)$ is called a first-passage datum.

\ensuremath{\dagger} The remaining pairs with $\beta\ge g^{-2}+B_\sharp$ are accounted for as follows. Suppose that $K\ge K'=K(\beta,g)$, put $b=K-K'$, and let $f^{[b]}(x)=f(L^{-b}x)$ on $T_{L^b\ell}$. Then

\begin{gather*}
\FitMath{\langle\mathcal{O}(f_1);\dots;\mathcal{O}(f_n)\rangle_{K,\ell,\beta}=\langle\mathcal{O}(f_1^{[b]});\dots;\mathcal{O}(f_n^{[b]})\rangle_{K',\,L^b\ell,\,\beta}}
\end{gather*}

holds exactly, for all bounded $f_i$. To see this, note that the dilation $x\mapsto L^bx$ maps $\Lambda_{K,\ell}$ onto $\Lambda_{K',L^b\ell}$, plaquettes to plaquettes and barycentres to barycentres. Wilson's weight depends only on the plaquette variables and on Haar measures, and no power of the spacing enters $\mathcal{O}$. The same identity holds between a periodic Gibbs state at $\beta$ on $\mathbb{L}_K$ and its image on $\mathbb{L}_{K'}$ (3d). Consequently a pair with $K>K(\beta,g)$ amounts to the first-passage datum of $\beta$ on a dilated torus, on which $\|f^{[b]}\|_\sharp\le\|f\|_\sharp$ and every distance is multiplied by $L^b$. The restriction to first passage therefore loses none of these lattice measures; its only effect is to fix the unit of length. (When $K<K(\beta,g)$ the identity is used in the opposite direction, on a contracted torus. This requires a lower bound on $\ell$, and the norms grow by at most the factor $L^{K(\beta,g)-K}$ (3d).)

The final value $x_K(\beta)$ lies in the window but is not fixed at $g^{-2}$. This is why (a)(ii) passes to subsequences, and it is the purpose of clause (d). On each torus, provided the conditions of (d)(1) are satisfied, that clause dispenses with the subsequences as soon as the renormalised coupling, defined below, converges along the sequence. On $\mathbb{R}^4$, in (a)(iii), a subsequence is taken in every case, including the case in which the renormalised coupling converges: the fixed-torus statement does not assert uniqueness of the infinite-volume limit. The renormalised inverse coupling $x^\natural_K(\beta)$ is the final value of a particular run of $K$ steps starting at $\beta$. In that run, the auxiliary choices of level $k$ that depend on the coupling are made with $\theta_k=g^{-2}+c_0(K-k)$ instead of $x_k$, where $c_0=\tfrac{11N^2}{12\pi^2}\log L$. Here and in all steps, $\log$ denotes the natural logarithm. A second kind of run is needed for the following reason. Continuity and holomorphy in the bare coupling are established by setting the constructions at two bare couplings side by side and comparing them entry by entry, and by letting a single fixed linear operation act on a whole family of bare densities. Each of these two devices needs the choices to coincide at the two couplings. In a run made with its own couplings, however, the choices vary with $\beta$. Some of them, namely the cube sizes $\check{R}_k$, change through integer values (3c). Freezing the choices leaves $Z_{K,\ell,\beta}$ and $S^T_n$ unchanged, since these are quantities of the measure. At first-passage data, and under the conditions of (d)(1), the shift that freezing produces in the couplings is at most $g$ times a power of $\log g^{-2}$ (3c), (8b). By a continuum limit at $g$ we mean a family $S=(S_n)_{n\ge2}$ with

\begin{gather*}
\FitMath{S_n(X)=\lim_{j\to\infty}S^T_{n;K_j,\ell_j,\beta_j}(X),} \\[3pt]
\FitMath{\beta_j\in\mathcal{W}_{K_j}(g),} \\[3pt]
\FitMath{K_j\to\infty,}
\end{gather*}

taken either on $T_\ell$ (with all $\ell_j=\ell$) or on $\mathbb{R}^4$ (with $\ell_j\to\infty$). Limits on a torus along the tuned couplings of (a)(i) are also counted as continuum limits. A volume-first limit at $g$ is a limit $\lim_jS^T_n[\nu_j]$, taken along first-passage data $(K_j,\beta_j)$ with $K_j\to\infty$. In it each $\nu_j\in\mathcal{G}^{\mathrm{tor}}_{K_j,\beta_j}$ is a limit point of the periodically extended torus measures as $\ell\to\infty$, with spacing and coupling held fixed. Thus the volume is removed first and the spacing only afterwards.

\textbf{The two theorems in brief.} The exponents are fixed first. Then $L$ is taken large, then the remaining constants are chosen in their order, and finally $g$ is taken small. Theorem 0 states the following. The two assertions below hold uniformly in the first-passage datum, and the left-hand one holds also uniformly in the torus as soon as $\ell\ge Cn(R+1)$:

\begin{gather*}
\FitMath{\big|S^T_{n;K,\ell,\beta}(f_1,\dots,f_n)\big|\le C^\infty_n(\mathsf{d};R)\prod_i\|f_i\|_\sharp,} \\[3pt]
\FitMath{S_2(f,h)=\frac{3(N^2-1)}{\pi^4}\Big(\frac{12\pi^2}{11N^2}\Big)^2\frac{\mathcal{K}_{\mathrm{free}}(f,h)}{(\log(1/\mathsf{d}))^2}\,(1+\eta)>0.}
\end{gather*}

In the left-hand bound, $n\ge2$, and the supports lie in balls of radius $R$ and are pairwise at distance $\ge\mathsf{d}$. In the right-hand formula, $S$ is a limit along first-passage data, $f,h\ge0$ are bumps whose size and separation are of the order of a small $\mathsf{d}$, and $\mathcal{K}_{\mathrm{free}}(f,h)=\iint f(x)h(y)|x-y|^{-8}\,dx\,dy$. The factor $1+\eta$ is bounded above and below by positive constants, and on each torus $\eta\to0$ as $\mathsf{d}\to0$; on $\mathbb{R}^4$ the clause consists of the bounds by the two constants and nothing more. The five clauses of this theorem are the following:
\begin{itemize}
\item limits exist on the tori and on $\mathbb{R}^4$ (a);
\item they inherit translation and hypercubic invariance, reflection positivity, symmetry and regularity; the factorial growth assertion in (b) is supplied for limits formed from periodic Wilson DLR states;
\item they are non-trivial and have, at short distance, the normalisation of asymptotic freedom (c);
\item on a torus they depend only on the renormalised coupling (d);
\item on $\mathbb{R}^4$ they are rotation invariant if and only if a flux through far shells vanishes (e).
\end{itemize}

Theorem 0' states that a volume-first limit $S$ is $\mathrm{O}(4)$-invariant. Its centred functional $\sigma_S=\sum_{\pi\in\operatorname{Part}_{\ge2}}\prod_{J\in\pi}S_{|J|}$ is precisely the system of moments obtained when the one-point function is set to zero, which is to say the moments of the centred density. This functional admits a canonical decomposition, which consists of a probability measure $\lambda$, carried by a compact set $\mathsf{X}=\operatorname{supp}\lambda$, together with a continuous injective map $\xi\mapsto\sigma_\xi$ such that

\begin{gather*}
\FitMath{\int_{\mathsf{X}}\prod_{i=1}^q\sigma_\xi(X^i)\,\lambda(d\xi)=\lim_{|v_q|\to\infty}\cdots\lim_{|v_2|\to\infty}\sigma_S\big(X^1\cdot\tau_{v_2}X^2\cdots\tau_{v_q}X^q\big)} \\[3pt]
\FitMath{(q\ge1),} \\[3pt]
\FitMath{\text{so that}} \\[3pt]
\FitMath{\sigma_S=\int_{\mathsf{X}}\sigma_\xi\,\lambda(d\xi).}
\end{gather*}

Once the stated analytic obligations are proved, each component $\sigma_\xi$ has the Euclidean properties listed in Theorem 0'. The reconstruction implication additionally uses the quantitative growth condition and the external analytic-continuation theorem specified in Section~\ref{ch:precise-reconstruction}.


\subsubsection*{(0a) The theorem, stated clause by clause}

There exist a shape constant $\vartheta_0\in(0,1]$ and, for each source radius $\varrho>0$ (step 4) and each shape $\vartheta\in(0,1]$, thresholds $0<\gamma_{U'}\le\gamma_U\le\gamma_J$ and $0<\gamma_W(\vartheta)\le\gamma_J$ such that Theorem 0' and the clauses below hold. Throughout, $0<g\le\gamma_J$, $\ell\in\mathfrak{L}(g)$, and $(K,\beta)$ is a first-passage datum. Apart from dependences that are displayed, the constants do not vary with $K,\beta,\ell$ or with positions. The ceilings $\gamma_{U'},\gamma_U,\gamma_J$ depend only on the constants fixed before them, among which is $\varrho$, and $\gamma_W(\vartheta)$ depends only on these and on $\vartheta$. Each constant, threshold and function introduced by a clause of the two theorems may in addition depend on $g$ and on the quantities displayed with it. The brackets after a clause give the steps that yield that clause, together with some of what those steps take from other steps.

\textbf{The frame is adapted to the tools.} When the side is $\ell=2L^m$ with $m\in\mathbb{N}$, each axis carries an even number $2L^{m+K}$ of sites. Each mirror pair $\{y_\mu=c\}\cup\{y_\mu=c+\ell/2\}$ therefore cuts the torus into two halves, and the chessboard cells come in pairs (4d). Because $L$ is odd, the mirrors that bisect bonds at successive cutoffs are nested: $a_K(\mathbb{Z}+\tfrac12)\subset a_{K+1}(\mathbb{Z}+\tfrac12)$. Because barycentres are used, the identity $\mathcal{O}(f)\circ\mathsf{r}^*=\mathcal{O}(\mathsf{r}f)$ is exact for each lattice isometry $\mathsf{r}$, where $\mathsf{r}^*$ denotes its action on configurations. For the same reason $\mathcal{O}(f)$ stays on the same side of such a mirror as $f$. In the exact-cochain Gaussian reference of Section~\ref{ch:lattice-curvature-comparison}, the mean of the quadratic approximation to $\mathcal O(f)$ is $\tfrac{N^2-1}{4\beta}(1-V_K^{-1})\sum_p f(x(p))$, where $V_K=(\ell/a_K)^4$ is the number of sites. For a fixed smooth $f$ of nonzero integral on a fixed torus, its leading Riemann-sum behavior is $\tfrac{3(N^2-1)}{2\beta}a_K^{-4}\int f$, which diverges in magnitude along the stated first-passage scaling. Moreover $|x-y|^{-8}$ is not locally integrable. For these reasons only $n\ge2$ and only separated tuples are considered, and positivity is stated for the centred $\sigma_S$. This is the diagonal-flat test domain used in Section~\ref{ch:precise-reconstruction}. Flatness removes sensitivity to contact terms; it does not supply the growth bound, continuity or analytic-continuation hypotheses needed for reconstruction. By a mirror $\Sigma$ of $\mathbb{R}^4$ we mean a hyperplane $\{u\cdot x=c\}$ with either $u=e_\mu$, the unit vector of the $\mu$-th coordinate (coordinate), or $u=(e_\mu\pm e_\nu)/\sqrt2$ for some $1\le\mu<\nu\le4$ (diagonal). Its reflection is $\theta_\Sigma$ and its open side is $\mathbb{E}_+(\Sigma)$. A torus has only the coordinate pairs as mirrors, and there the open side is a slab.

\textbf{Theorem 0 (a) [existence; 3, 5, 8].} Each run started from $\beta\ge g^{-2}+B_\sharp$ reaches first passage. For such a run, $c(\beta-g^{-2})-C\le K(\beta,g)\le C(\beta-g^{-2})+C$, and since $b_k\le B_\sharp$ also $g^{-2}<x_K(\beta)\le g^{-2}+B_\sharp$. Hence $K(\beta,g)\to\infty$ as $\beta\to\infty$, and so $\mathcal{W}_K(g)\ne\varnothing$ for infinitely many $K$. The values $x_k(\beta)$ are the same whatever torus the steps are carried out on, and $\beta/\log(1/a_K)\to\tfrac{11N^2}{12\pi^2}$ uniformly on $\mathcal{W}_K(g)$. (i) Let $g\le\gamma_U$. Then for each $K$ there is a tuned $\beta^\natural_K$, that is, one with $x^\natural_K(\beta^\natural_K)=g^{-2}$, and it satisfies $\tfrac12c_0K-c_5\le\beta^\natural_K-g^{-2}\le\tfrac32c_0K+c_5$. Moreover $S^T_{n;K,\ell,\beta^\natural_K}\to S^{\natural,\ell}_n$ as $K\to\infty$ along the whole sequence. This convergence holds on each torus of the family and on tuples of Lipschitz functions whose supports are pairwise disjoint. (ii) On a fixed torus, from every sequence of first-passage data with $K_j\to\infty$ one can extract a subsequence that converges on $\mathfrak{S}(T_\ell)$. (iii) The same holds on $\mathfrak{S}(\mathbb{R}^4)$ when in addition $\ell_j\to\infty$. The limits obtained in this way along tori, with defining sequence $(K_j,\ell_j,\beta_j)$, are what the phrase ``a limit of (a)(iii)'' means. The same also holds for $S^T_n[\nu_j]$ with $\nu_j\in\mathcal{G}^{\mathrm{tor}}_{K_j,\beta_j}$, so volume-first limits exist.

\textbf{(b) [inheritance; 6; the bounds from (4c)--(4e), those of $S^{\natural,\ell}$ through (8a), the last sentence from (5c) and (10a)].} Each continuum limit has the following properties.
(i) It is invariant under $W_4$ and under the translations of $\mathbb{E}$.
(ii) It is reflection positive at each coordinate mirror: $\sum_{i,j}\bar c_ic_j\,\sigma_S(\theta_\Sigma X_i\cdot X_j)\ge0$ whenever $X_i\in\mathfrak{S}(\mathbb{E}_+(\Sigma))$.
(iii) It is symmetric.
(iv) It is regular: there is a unique distribution of locally finite order on $\mathbb{E}^n_{\ne}$ of which $S_n$ is the restriction. On $C^1$ functions, and with the same constants, it satisfies (c)(1) on $T_\ell$. The per-order bound labelled (c)(2) passes to limits of (a)(ii)--(iii), with its displayed volume condition retained in the fixed-torus case. The factorial bound displayed after it passes to limits formed from periodic Wilson DLR states.
(v) On $\mathbb{R}^4$, a limit formed from periodic Wilson DLR states satisfies the growth condition

\begin{gather*}
\FitMath{|\mathsf{S}^{\sigma_S}_n(\varphi)|\le c_{\mathrm{E}}C_{\mathrm{E}}^n(n!)^{\varkappa}|\varphi|_{k_{\mathrm{E}}n}}
\end{gather*}

 on $\mathcal{S}_{\ne}(\mathbb{R}^{4n})$. For every $\mathbb{R}^4$ limit, take the form $\langle F,G\rangle=\sigma_S(\theta\bar F\cdot G)$ on $\mathbb{C}\mathfrak{S}(\{x_4>0\})$. The separable space $\mathcal{H}(S)$ then carries three structures:
\begin{itemize}
\item $\mathcal{T}_t[F]=[\tau_{te_4}F]$, $\mathcal{T}_t=e^{-tH}$, which form a strongly continuous semigroup of self-adjoint contractions, where $H\ge0$ and $H\Omega=0$;
\item a strongly continuous unitary representation $e^{iv\cdot P}$ of $\mathbb{R}^3$ that leaves $\Omega$ fixed, the $P_j$ commuting strongly with $H$;
\item a unitary representation of $W_3$ that leaves $\Omega$ fixed and commutes with $H$.
\end{itemize}

Already for $g\le\gamma_J$, each volume-first limit is a limit of (a)(iii) and is, in addition, reflection positive at each diagonal mirror.

\textbf{(c) [bounds, non-triviality; 4, 7; the envelope of $g(\mathsf{d})$ from (3a), the periodic Gibbs states through (10a)].} Let $n\ge2$, let the $f_i$ be Lipschitz with supports pairwise at distance $\ge\mathsf{d}>0$, and in (2) let the supports also lie in balls of radius $R$. Then

\begin{gather*}
\FitMath{(1)\,|S^T_n|\le C_n(\mathsf{d})\,e^{nc_V\ell^4}\prod_i\|f_i\|_\sharp;} \\[3pt]
\FitMath{(2)\,|S^T_n|\le C^\infty_n(\mathsf{d};R)\prod_i\|f_i\|_\sharp\,\,(\ell\ge Cn(R+1)),} \\[3pt]
\FitMath{|S^T_n[\nu]|\le [2B(R)]^n(n!)^k\max\{1,\mathsf{d}^{-5n}\}\prod_i\|f_i\|_\sharp;}
\end{gather*}

Here $C^\infty_n(\mathsf{d};R)$ does not increase as $\mathsf{d}$ grows.
The last bound is for periodic Wilson DLR states $\nu$ and real $C^1$
tests, under D1 in its actual rooted geometry and on its compatible
source tori, together with G1. One may take $B(R)=C(R+1)^4$ and one
fixed finite integer $k$; see Theorem~\ref{thm:graded-rooted-source}.
A finite-torus factorial bound, or a single polynomial separation
prefactor of degree independent of $n$, requires an additional estimate. The bounds (2) also hold, with no condition on a volume, for $S^T_n[\nu]$, where $\nu$ is an arbitrary periodic Gibbs state on $\mathbb{L}_K$ of Wilson's weight at $\beta$, and where $(K,\beta)$ is a first-passage datum.

(3) The class $\mathfrak{T}_{\mathbb{E}}(\mathsf{d};\vartheta)$ consists of the pairs of Lipschitz functions $f,h\ge0$, neither of which is identically zero, that satisfy the following conditions. Each is supported in a ball of radius $\mathsf{d}$, and $\mathsf{d}\le\operatorname{dist}(\operatorname{supp}f,\operatorname{supp}h)\le6\mathsf{d}$. In addition, $\int f\ge\vartheta\|f\|_\infty\mathsf{d}^4$ and $\operatorname{Lip}f\le\|f\|_\infty/(\vartheta\mathsf{d})$, and the same two conditions hold for $h$. Assume $g\le\gamma_W(\vartheta)$, $\mathsf{d}\le\mathsf{d}_0$ and $(f,h)\in\mathfrak{T}_{T_\ell}(\mathsf{d};\vartheta)$. Assume also $K-s(\mathsf{d})\ge n_0$, so that at least $n_0$ steps lie between the scale $\mathsf{d}$ and the cutoff. Then

\begin{gather*}
\FitMath{\langle\mathcal{O}(f);\mathcal{O}(h)\rangle_{K,\ell,\beta}=\zeta_K(\mathsf{d})^2\,\tfrac{3(N^2-1)}{\pi^4}\,g(\mathsf{d})^4\,\mathcal{K}_{\mathrm{free}}(f,h)\,(1+\varepsilon),} \\[3pt]
\FitMath{|\varepsilon|\le\min\{\tfrac12,\bar\varepsilon_\ell(\mathsf{d})\}.}
\end{gather*}

In this formula $g(\mathsf{d})=g_{K-s(\mathsf{d})}(\beta)$, where $L^{-s(\mathsf{d})}\approx\mathsf{d}$ up to logarithmic factors. Also $\underline x_s\le g(\mathsf{d})^{-2}\le\bar x_s$, and both bounds are affine in $s=s(\mathsf{d})$. The factor $\zeta_K(\mathsf{d})=\zeta'_{K-s(\mathsf{d})}$ is the running normalisation, that is, it measures how $x_{K-s(\mathsf{d})}$ responds when $\beta$ is shifted by a constant (4b). One has $|\zeta_K(\mathsf{d})-1|\le\bar\eta(\mathsf{d})$ and $|\tfrac{11N^2}{12\pi^2}g(\mathsf{d})^2\log\tfrac1{\mathsf{d}}-1|\le\bar\eta(\mathsf{d})$. Here $\mathsf{d}_0\in(0,1)$, and the two functions $\bar\eta:(0,\mathsf{d}_0]\to[0,\frac13]$ and $\bar\varepsilon_\ell:(0,\mathsf{d}_0]\to[0,\frac12]$ are non-decreasing; moreover $\bar\eta,\bar\varepsilon_\ell\to0$ as $\mathsf{d}\to0$. Hence, on each torus, $\varepsilon\to0$ uniformly with respect to the cutoff, to the bare coupling as it ranges over its window, to the pair, and to where the pair is placed. For periodic Gibbs states the identity remains true, but only with $|\varepsilon|\le\tfrac12$.

(4) Consequently the limits of (a)(ii)--(iii) obey the law stated at the opening of the text, for those pairs in $\mathfrak{T}_{\mathbb{E}}(\mathsf{d};\vartheta)$ in which $f,h$ are of class $C^1$; the tuned $S^{\natural,\ell}$ are not included. The error satisfies $(1-\bar\varepsilon)(1-\bar\eta)^4\le1+\eta\le(1+\bar\varepsilon)(1+\bar\eta)^4$, with $\bar\varepsilon=\bar\varepsilon_\ell(\mathsf{d})$ on $T_\ell$ and $\bar\varepsilon=\tfrac12$ on $\mathbb{R}^4$, where in addition $\dim\mathcal{H}(S)\ge2$.

(5) The class $\mathfrak{T}^{(4)}_{\mathbb{E}}(\mathsf{d};\vartheta)$ consists of the quadruples whose pairs all belong to $\mathfrak{T}_{\mathbb{E}}(\mathsf{d};\vartheta)$ and for which $\operatorname{supp}f_i\subset B(c_i,\mathsf{d}/16)$, the four centres $c_i$ lying on one line. On the tori there exists $\vartheta_\square\in(0,1]$ with the following property. Let $\vartheta\le\vartheta_\square$ and $g\le\gamma_W(\vartheta)$. Then there are constants $0<c_\square\le C_\square$, which may depend on $g$ and $\vartheta$, such that the bound below holds whenever $K-s(\mathsf{d})\ge n_4$, where $n_4\ge n_0$ is an integer, $\mathsf{d}$ is smaller than a threshold, which may depend on $\ell$, and the quadruple belongs to $\mathfrak{T}^{(4)}_{T_\ell}(\mathsf{d};\vartheta)$:

\begin{gather*}
\FitMath{c_\square(N^2-1)\,g(\mathsf{d})^8\prod_i\|f_i\|_\infty\le S^T_{4;K,\ell,\beta}(f_1,\dots,f_4)\le C_\square(N^2-1)\,g(\mathsf{d})^8\prod_i\|f_i\|_\infty:}
\end{gather*}

on these quadruples the connected four-point function is therefore strictly positive. On rescaled quadruples it decreases as the fourth power of the reciprocal of the logarithm of the scale.

\textbf{(d) [uniqueness, $g\le\gamma_U$; 8].} Take $(\beta^\natural_K)$ and $S^{\natural,\ell}$ as in (a)(i).
(1) Let $\varrho$ be large and $g\le\gamma_{U'}$. Then the first-passage windows lie close to the tuned couplings: $\mathcal{W}_K(g)\subset[\beta^\natural_K-\tfrac14\varrho,\ \beta^\natural_K+\tfrac14\varrho]$ for each $K$.
(2) For $|\beta-\beta^\natural_K|\le\tfrac14\varrho$ the reference runs of length $K$ are well defined, and $\beta\mapsto x^\natural_K(\beta)$ is bi-Lipschitz uniformly in $K$. As $K\to\infty$, the functions $t\mapsto x^\natural_K(\beta^\natural_K-t)$ converge uniformly on $|t|\le\tfrac14\varrho$. The limit function has a compact interval $I$ as its range, and this interval contains $[g^{-2}-c'\varrho,g^{-2}+c'\varrho]$.
(3) For each torus of the family, one has a family $(S^{(\ell,x)})_{x\in I}$. Each of its members is a family of connected functions, defined on tuples of Lipschitz functions whose supports are pairwise disjoint. It satisfies $S^{(\ell,g^{-2})}=S^{\natural,\ell}$, and $x\mapsto S^{(\ell,x)}_n$ is continuous on $I$. Moreover, whenever $K_j\to\infty$, $|\beta_j-\beta^\natural_{K_j}|\le\tfrac14\varrho$ and $x^\natural_{K_j}(\beta_j)\to x$, it follows that $x\in I$ and $S^T_{n;K_j,\ell,\beta_j}\to S^{(\ell,x)}_n$, with no passage to a subsequence.

Consequently, take two sequences of bare data close to the tuned couplings whose renormalised couplings converge to the same value. They have the same limit on $T_\ell$. Under the conditions of (1) this applies to every sequence of first-passage data, and every limit of (a)(ii) is equal to some $S^{(\ell,x)}$. On each torus, the renormalised coupling by itself determines the limit theory.

\textbf{(e) [rotations; 9].} Let $\omega\in\mathfrak{so}(4)$, and let $\chi$ be a cutoff whose support lies in $(-\ell/2+2a_K,\ell/2-2a_K)^4$. Consider the rotational Ward identity satisfied by $\mu_{K,\ell,\beta}$ with respect to the deformation field $\chi(y)\,\omega y$. The shell flux is the boundary term of this identity:

\begin{gather*}
\FitMath{\mathfrak{B}_{\omega,\chi}=\beta\sum_pa_p\,(\nabla^{\mathrm{s}}_\omega\chi)(x(p))-\beta\sum_p\mathfrak{h}_{\omega,\chi}(p)-\tfrac32\sum_pa_p\sum_{\rho\in\{\mu,\nu\}}(\omega x(p))_\rho\,(\nabla^{\mathrm{c}}_\rho\chi)(x(p)),}
\end{gather*}

where $\mu<\nu$ are the two directions spanning $p$ and $\nabla^{\mathrm{s}}_\omega,\nabla^{\mathrm{c}}_\rho$ are lattice differences, and where 

\begin{gather*}
\FitMath{\mathfrak{h}_{\omega,\chi}(p)=a_K^{-1}\sum_{b\subset\partial p}\sum_\rho[\chi(x_b)-\chi(x(p))]\,(\omega x_b)_\rho\,\mathfrak{d}_{b,p,\rho}}
\end{gather*}

, the $\mathfrak{d}_{b,p,\rho}$ being bounded, gauge-invariant local densities. Each term contains a difference of $\chi$. Therefore $\mathfrak{B}_{\omega,\mathcal{R}}:=\mathfrak{B}_{\omega,\chi_{\mathcal{R}}}$ is supported on the shell $\{\mathcal{R}-a\le|x(p)|\le2\mathcal{R}+a\}$, $a=a_K$, which surrounds the shell $\mathcal{R}\le|y|\le2\mathcal{R}$ where $\chi_{\mathcal{R}}$ falls from $1$ to $0$. A limit of (a)(iii) is $\mathrm{O}(4)$-invariant if and only if the following holds for one $\omega\ne0$ (equivalently, for all $\omega$), for all $n\ge2$ and for all separated tuples with smooth entries:

\begin{gather*}
\FitMath{\liminf_{\mathcal{R}\to\infty}\limsup_{j\to\infty}\big|\langle\mathfrak{B}_{\omega,\mathcal{R}};\mathcal{O}(f_1);\dots;\mathcal{O}(f_n)\rangle_{K_j,\ell_j,\beta_j}\big|=0.}
\end{gather*}

In that case $\mathcal{H}(S)$ carries a strongly continuous unitary representation of $\mathrm{O}(3)$. This representation leaves $\Omega$ fixed, commutes with $H$, and extends the representation of $W_3$. This equivalence singles out the one infrared input, and for volume-first limits Theorem 0' settles it.

\emph{What the equivalence uses.} When $\chi=1$ near the supports, the identity takes the following form (Theorem 9.3 [to be proved; Package G], in (9a)):

\begin{gather*}
\FitMath{\sum_{i=1}^n\big\langle\mathcal{O}(f_1);\dots;\mathcal{O}(\nabla^{\mathrm{s}}_\omega f_i);\dots;\mathcal{O}(f_n)\big\rangle_{K,\ell,\beta}=\mathfrak{U}_{\omega,\chi}(f_1,\dots,f_n)} \\[3pt]
\FitMath{+\big\langle\mathfrak{B}_{\omega,\chi};\mathcal{O}(f_1);\dots;\mathcal{O}(f_n)\big\rangle_{K,\ell,\beta}:}
\end{gather*}

besides the flux, it contains a bulk term $\mathfrak{U}_{\omega,\chi}$. This term is the correlation between the $\mathcal{O}(f_i)$ and lattice artefacts of dimension six, distributed over $\operatorname{supp}\chi$. Once $\mathfrak{U}_{\omega,\chi_{\mathcal{R}}}\to0$ along every defining sequence, the rotation defect of $S$, namely the derivative $\tfrac{d}{dt}S_n(e^{t\omega}f_1,\dots,e^{t\omega}f_n)|_{t=0}$, equals minus the limit of the correlation between the flux and the $\mathcal{O}(f_i)$, and clause (e) follows.

This vanishing (Theorem 9.1 [to be proved; Package G]) is a consequence of three propositions about an insertion that does not have the form of the action, followed to first order through the $K$ steps (Propositions 9.7 [to be proved; Package G], 9.8, 9.9), together with a fourth. The first three yield the vanishing for slowly growing volumes as well as on a fixed torus. The fourth is a moment bound in which several insertions occur simultaneously (Proposition 9.11 [to be proved; Package G]); by means of it, the transfer in the volume (Proposition 9.10 [to be proved; Package G]) extends the vanishing to every sequence with $\ell_j\to\infty$. All of these propositions are in (9b). Along either route, the vanishing depends on a single number.

One marginal operator remains in $\mathfrak{U}$: a traceless diagonal anisotropy of Wilson's weight. Its coefficient is of the size $g_0^{-1}\approx K^{1/2}$, and it disappears only through its own decay $K^{-c_1}$ along the flow. What is needed is $c_1>\tfrac12$.

This number is expressed through one-loop data. Put the weight $1+\mathsf{e}_{\pi(p)}$ on each $a_p$, where $\pi(p)$ is the coordinate plane of $p$. Differentiating the one-loop coefficient of step $l$ plane by plane at $\mathsf{e}=0$ gives a derivative that commutes with $W_4$. Let $\mathsf{m}(l)$ denote the eigenvalue of this derivative on the traceless diagonal weights $\mathsf{e}_{\{\mu,\nu\}}=\mathsf{e}_\mu+\mathsf{e}_\nu$, $\sum_\mu\mathsf{e}_\mu=0$. Then $c_1\in(\tfrac12,1]$ is an exponent such that $\sum_{i<l\le j}\mathsf{m}(l)\ge c_1c_0(j-i)-C$. The one-loop law in this channel, $|\sum_{l\le k}\mathsf{m}(l)-c_0k|\le C$, gives $c_1=1$.

The law comes from a telescoping sum. By the law of total covariance, the sum of the one-loop responses of the steps $l\le k$ is a single trace on the bare lattice. The logarithmic part of this trace is $\int_1^{L^{2k}}d\sigma/\sigma=2k\log L$, in proper time $\sigma$, multiplied by the heat coefficient on a constant abelian background. It equals $kc_0$, with the same $c_0$ as the coupling itself. In the continuum, an anisotropy of the weight amounts, to first order, to a change of coordinates, and it does not run (9b). Rotation invariance in Theorem 0' (1) is obtained along a different route, without any Ward identity (10b), and does not depend on $c_1$.

\clearpage
\textbf{Theorem 0' [10]} (Let 

\begin{gather*}
\FitMath{g\le\min\{\gamma_W(\vartheta),\gamma_W(\vartheta_0)\}}
\end{gather*}

, and let $S$ be a volume-first limit with defining sequence $(K_j,\beta_j,\nu_j)$.)
(1) Every $\nu_j$ is a periodic Gibbs state and is invariant under all translations of the lattice; moreover $S$ is a limit of (a)(iii). In addition to (b) and (c)(4), it satisfies $H\ge|P_j|$ for $j=1,2,3$, is $\mathrm{O}(4)$-invariant, and is reflection positive at the diagonal mirrors.
(2) Its canonical decomposition exists and is unique up to isomorphism. Here an isomorphism is a homeomorphism of the spaces $\mathsf{X}$ that carries $\lambda$ and the $\sigma_\xi$ along with it. Let $P_0$ be the projection onto $\ker H$. The map $P_0[F]\mapsto(\xi\mapsto\sigma_\xi(F))$ extends to a unitary map $\ker H\to L^2(\mathsf{X},\lambda)$ with $\Omega\mapsto1$, so that $\dim\ker H=\#\mathsf X$ when $\mathsf X$ is finite, while $\ker H$ is infinite-dimensional when $\mathsf X$ is infinite. The components are permuted, and $\lambda$ is preserved, by the isometries that leave $S$ invariant (translations, $W_4$, rotations).
(3) Each $\sigma_\xi$ has the following properties.
(i) It has the same invariances and positivity properties, and it is symmetric and regular.
(ii) Its one-point function is $S^{\sigma_\xi}_1(f)=\mathfrak{c}_\xi\int f$, where $\mathfrak{c}_\xi$ is real and continuous in $\xi$ and $\int\mathfrak{c}_\xi\,\lambda(d\xi)=0$. With the same constants, it satisfies the growth condition and the bounds (c)(2); it also satisfies the law (c)(4) in its form on $\mathbb{R}^4$.
(iii) On $\mathcal{H}(\sigma_\xi)$ it has the operators of (b)(v) and a unique vacuum, $\ker H_{\sigma_\xi}=\mathbb{C}\Omega_{\sigma_\xi}$. Its joint spectrum lies in the closed forward cone $\{E\ge|q|\}$. It also has a strongly continuous unitary representation of $\mathrm{O}(3)$ that leaves $\Omega_{\sigma_\xi}$ fixed and commutes with $H_{\sigma_\xi}$.
(iv) It has the cluster property:

\begin{gather*}
\FitMath{|\sigma_\xi(X\cdot\tau_vY)-\sigma_\xi(X)\,\sigma_\xi(Y)|\le C(X,Y)\,|v|^{-2}} \\[3pt]
\FitMath{(|v|\ge r(X,Y)),}
\end{gather*}

where $C(X,Y)$ and $r(X,Y)$ are the same for all $\xi$ and for all volume-first limits at $g$.
(4) Under the analytic hypotheses of this theorem, $\sigma_\xi$ has the stated growth, invariance under $\mathbb{R}^4\rtimes\mathrm{O}(4)$, reflection positivity, symmetry and clustering. Applying the external reconstruction theorem with the quantitative hypotheses of Section~\ref{ch:precise-reconstruction} then identifies these Schwinger functions with those of a Wightman theory. This theory has a single Hermitian scalar field on a separable Hilbert space and a unique vacuum. It is covariant under the proper orthochronous Poincar{\'e} group, its spectrum lies in the closed forward cone, and it is local.
% END MODULE chapters/01_orientation.tex


% BEGIN MODULE chapters/02_averaging.tex
\section{Block averages and gauge fixing}
\label{candidate:averaging}
\candidatebanner{A and C}

\subsection*{(1a) Block averages, gauge fixing, and the renormalisation group of this construction}

\textbf{Setting and objects.} \emph{Requirements on an average.} We work in level-$k$ units, with $L$ odd and centred blocks $B(y)=\{x:|x-y|_\infty\le\frac12(L-1)\}$. We put $\varpi(W)=\sup_p|\partial W(p)-\mathbf{1}|$, and $\varpi_c$ denotes the same supremum taken over the double block $B(c_-)\cup B(c_+)$. For kernels we use the norms $\|K\|_{\mathrm{row}}=\sup_c\sum_b|K(c,b)|$ and $\|K\|_{\mathrm{col}}=\sup_b\sum_c|K(c,b)|$. Elements of the group cannot be added, and transports joining different pairs of points cannot even be compared. A block average $\mathcal{M}:\mathcal{V}_k\to\mathcal{V}_{k+1}$ is required to have the following properties: gauge covariance; locality; linearity at first order; analyticity near flat fields, on $G^c$ as well; symmetry under the isometries of the coarse lattice; and Borel measurability on the whole of $\mathcal{V}_k$, because $\mathbf{T}$ integrates over all configurations. For the mean of group elements Ba{\l}aban gives axioms (109 (1987), (0.5)--(0.9)). Such a mean is any analytic $G^c$-valued function $M$ on families of small diameter that satisfies

\begin{gather*}
\FitMath{M(\{\mathbf{U}_j^{-1}\})=M(\{\mathbf{U}_j\})^{-1},} \\[3pt]
\FitMath{M(\{u\mathbf{U}_jv\})=u\,M(\{\mathbf{U}_j\})\,v,} \\[3pt]
\FitMath{M(\pi\{\mathbf{U}_j\})=M(\{\mathbf{U}_j\}),}
\end{gather*}


\begin{gather*}
\FitMath{\tfrac1i\log M(\{e^{iA_j}\})=\tfrac1n\sum_jA_j+(\text{higher orders}),} \\[3pt]
\FitMath{\mathbf{U}_j\in G\text{ for all }j\,\Rightarrow\,M(\{\mathbf{U}_j\})\in G.}
\end{gather*}

According to Ba{\l}aban, the results of his 1987 paper and of the earlier papers hold for every mean of this kind. For the map $\mathcal{M}$, the second of the two symmetric block averages that he introduces, he states that every result of the paper on averaging operations (Comm. Math. Phys. 98) holds, with proofs that are for the most part unchanged. This statement concerns that single paper; the broader sentence, which refers to ``this, and previous papers'', is the one about means. Of the two symmetric averages together he states that they are equally good, and that many other averages can also be used. He also states that the block average admits an axiomatisation as well. Already in 1985 he states that the analysis of the paper on averaging operations, together with its propositions, extends easily to other averages that satisfy the two requirements set out in its introduction, namely gauge covariance and agreement with the linear block average at first order. No proof accompanies any of these statements. For the theorems of 1984--85 on averaging, gauge fixing, propagators and backgrounds, the proofs are carried out for an average that is not symmetric. Proposition 1a.1.13 states which properties of an average the proofs use, and it compares the linear parts (1a.1).

\emph{The mean, the contours, the map.} The \emph{mean} $\operatorname{av}(V_\bullet)$ (Federbush) is the solution, close to the family, of $\sum_i\log(V_iV^{-1})=0$. By $\mathcal{G}(y,x)$ we denote the family of shortest paths from $y$ to $x\in B(y)$ that complete the displacement in one direction before passing to the next. There is one such path for each ordering of the directions, so the family has at most $4!$ members and does not contain every shortest path. The \emph{comb} $\Gamma_{y,x}$ is the member with the order $4,3,2,1$. The averaged contour variable is $U(y,x):=\operatorname{av}(U(\Gamma))_{\Gamma\in\mathcal{G}(y,x)}$, and we set $U(x,y)=U(y,x)^{-1}$ and $U(y,y)=\mathbf{1}$. With $x'=x+(c_+-c_-)$, one defines

\begin{gather*}
\FitMath{W_x=U(c_-,x)\,U([x,x'])\,U(x',c_+)\,U[c]^{-1},} \\[3pt]
\FitMath{\mathcal{M}(U)(c)=\exp\Big(L^{-4}\!\!\sum_{x\in B(c_-)}\!\!\log W_x\Big)U[c]} \\[3pt]
\FitMath{\text{on }\mathrm{Dom}(c)=\{\varpi_c<cL^{-2}\}:}
\end{gather*}

the $W_x$ are closed transports based at $c_-$ and lying near $\mathbf{1}$, so their logarithms can be added. Where $U(y,x)=\mathbf{1}$, only straight segments are left, so the form is adapted to the block axial gauge. Off $\mathrm{Dom}(c)$ the map is given by two cruder rules, both covariant. When none of the elements $Y_{x,\Gamma,\Gamma'}=U(\Gamma)U([x,x'])U(\Gamma')^{-1}U[c]^{-1}$ equals $-\mathbf{1}$, one replaces $\log W_x$ by the arithmetic mean, over $(\Gamma,\Gamma')$, of the intrinsic logarithms of the $Y_{x,\Gamma,\Gamma'}$; this is Ba{\l}aban's first symmetrised average. In the remaining case one sets $\mathcal{M}(U)(c)=U[c]$. This third rule is required because Ba{\l}aban defines his average on all configurations by means of the spectral logarithm, which satisfies $\log(-\mathbf{1})=i\pi\mathbf{1}\notin i\mathfrak{g}$ and so leads out of $\mathrm{SU}(2)$ (1a.1). By $\mathcal{M}^k$ we mean the iterate only, and not a mean taken in a single stroke over $B^k$. In Ba{\l}aban's paper of 1987 the lattices consist of centres of cubes and have coordinates in $a_j(\mathbb{Z}+\frac12)$. The single translation $x\mapsto x+\frac12L^m\sum_\mu e_\mu$ of $T_\ell$, $\ell=2L^m$, maps them onto the $\Lambda_{j,\ell}$ at every level simultaneously, since $\frac12L^m=\frac12L^{m+j}a_j$ and $L^{m+j}$ is odd. This is the second place where the oddness of $L$ is used, the first being the centring of the blocks (1a.1). Ba{\l}aban also remarks that all the results were obtained first for the two-dimensional chiral models, in which the fields sit on the sites and the action is $\sum_b[1-\frac12\operatorname{Re}\operatorname{tr}U(b_-)U(b_+)^{-1}]$. In those models block averages are means of group elements, and neither contours nor a gauge group enter.

\emph{Linear averages} of bond functions $\mathsf{a}$, which enter in the form $e^{i\mathsf{a}}W$ without any length factor, are defined as follows, where $\bar{\mathsf{a}}(y,x)$ denotes the arithmetic mean of $\mathsf{a}(\Gamma)$ over $\mathcal{G}(y,x)$:

\begin{gather*}
\FitMath{(\mathsf{Q}^{\mathrm{s}}\mathsf{a})(c)=L^{-4}\!\!\sum_{x\in B(c_-)}\!\!\mathsf{a}([x,x']),} \\[3pt]
\FitMath{\sigma(y)=L^{-4}\!\!\sum_{x\in B(y)}\!\!\bar{\mathsf{a}}(y,x),} \\[3pt]
\FitMath{(\mathsf{Q}^{\mathrm{pl}}\mathsf{a})(c)=\sigma(c_-)+(\mathsf{Q}^{\mathrm{s}}\mathsf{a})(c)-\sigma(c_+),} \\[3pt]
\FitMath{(q\lambda)(y)=L^{-4}\!\!\sum_{x\in B(y)}\!\!\lambda(x);}
\end{gather*}

here $\mathsf{Q}^{\mathrm{pl}}$ is the derivative of $\mathcal{M}$ at $\mathbf{1}$. The operator $\mathsf{S}(W_0)$ is $\mathsf{Q}^{\mathrm{s}}$ formed with parallel transport to $c_-$, namely 

\begin{gather*}
\FitMath{(\mathsf{S}(W_0)\mathsf{a})(c)=L^{-4}\sum_{x\in B(c_-)}\sum_{b\in[x,x']}\operatorname{Ad}\big(W_0(c_-,x)W_0([x,b_-])\big)\mathsf{a}(b)}
\end{gather*}

.

\emph{The straightened average} of $\mathbf{W}=e^{i\mathsf{a}}W_0$ relative to $W_0$ and its \emph{block factor} $\mathsf{u}$ (written $\mathsf{n}_{\mathbf{W}}$ in 1a.3) are given by

\begin{gather*}
\FitMath{\mathsf{u}(y)=\exp\Big(L^{-4}\!\!\sum_{x\in B(y)}\!\!\log\big[\mathbf{W}(y,x)W_0(y,x)^{-1}\big]\Big),} \\[3pt]
\FitMath{e^{i\mathcal{Q}^1(W_0,\mathsf{a})(c)}\mathcal{M}(W_0)(c)=\mathsf{u}(c_-)^{-1}\mathcal{M}(\mathbf{W})(c)\,\mathsf{u}(c_+).}
\end{gather*}

The right-hand side equals $\mathcal{M}(\mathbf{W}^v)$, where $v$ is the gauge transformation with $\mathbf{W}^v(y,x)=W_0(y,x)$ (the block axial gauge relative to $W_0$) and with trivial average $\bar{R}_{W_0}v=\mathbf{1}$. Together these two conditions imply $v=\mathsf{u}^{-1}$ at the coarse sites. Ba{\l}aban gives two reasons for the straightening. First, the variational problems of (1b) involve composing the average with the passage from the axial gauge to the Landau gauge. Second, only the straight linear average admits the explicit representations of the propagators (1c); this average is compatible with composition, has an explicit Fourier symbol and does not change under $q\lambda=0$. The change of gauge has to preserve this form, and the non-linear counterpart of $q\lambda=0$ is $\bar{R}_{W_0}v=\mathbf{1}$. A third reason appears in the proof of Theorem 1a.1.6.

\emph{Order $k$; the maps denoted by $Q$.} Along the sequence $W_{0,n}=\mathcal{M}^nU_0$ we set $\mathsf{a}_0=\eta_kA$, $\mathsf{a}_{n+1}=\mathcal{Q}^1(W_{0,n},\mathsf{a}_n)$, $v_0=\mathbf{1}$, $v_{n+1}=v_n\,\mathsf{u}(W_{0,n},\mathsf{a}_n)$, and define

\begin{gather*}
\FitMath{\mathcal{Q}^{\sharp}_k(U_0;A):=\mathsf{a}_k=Q^{\sharp}_k(U_0)A+C_k(U_0,A),} \\[3pt]
\FitMath{\mathcal{Q}_k(U_0;A)=\tfrac1i\log\big[\mathcal{M}^k(e^{i\eta_kA}U_0)\mathcal{M}^k(U_0)^{-1}\big]=Q_k(U_0)A+C^{\mathrm{pl}}_k(U_0,A).}
\end{gather*}

A sharp indicates the straightened objects, which are the ones that enter the propagators; Ba{\l}aban's $Q_k(U_0,\eta A)$ is the same object built with combs. Without the sharp, $\mathcal{Q}_k$ and $Q_k$ refer to the plain average, whose remainder is denoted $C^{\mathrm{pl}}_k$, while $C_k$ denotes the remainder of the straightened average. The product $v_k=v_k[A]$ of the $k$ block factors is the \emph{node factor}, and $\dot{v}_k$ denotes the linear part of $\frac1i\log v_k$. The letter $Q$ is used for maps of four kinds (1a.1). (1) The non-linear constraint map $\mathcal{Q}_k$ of the plain iterate $\mathcal{M}^k$, and next to it its straightened form $\mathcal{Q}^{\sharp}_k$. (2) $Q_k(U)=\partial_A\mathcal{Q}_k(U;A)|_{A=0}$, together with $Q=Q_1(\mathbf{1})$ and $Q_V=Q_1(V)$; this is the operator in the delta function of a step. (3) $Q^{\sharp}_k(U)$, the operator that enters the propagators and the Landau gauge. Its leading term is the \emph{straight} average with parallel transport, $Q^{\mathrm{st}}_k(U)=\eta_k\mathsf{S}(W_{0,k-1})\cdots\mathsf{S}(W_{0,0})$. The two differ by $O(\alpha_0)$ in row norm, and they are equal at $U=\mathbf{1}$. There, after the translation of Proposition 1a.1.13, each of them is Ba{\l}aban's $Q_k$ of 1984. (4) $Q'_k(U)$, which averages site functions. In a background, the linear average in Ba{\l}aban's papers (99 (1985) 389, (3.15)) is $Q^{\mathrm{B}}_j(U):=Q(\bar{U}^{j-1})\cdots Q(\bar{U})Q(U)$. It is the linear part of his gauge-corrected average built with combs, and each of its factors equals the straight covariant average up to $O(L^2\alpha_0)$. Theorem 1a.1.6 and Proposition 1a.1.13 provide two comparisons, that of (2) with (3) and that of $Q^{\mathrm{B}}_k$ with $Q^{\sharp}_k$, and (1c) relies on both.

\emph{Gauge transformations} are averaged by means of the same contours,

\begin{gather*}
\FitMath{(\bar{R}_Wv)(y)=v(y)\exp\Big(L^{-4}\!\!\sum_{x\in B(y)}\!\!\log\big[v(y)^{-1}\operatorname{Ad}(W(y,x))v(x)\big]\Big),}
\end{gather*}

and the iterates of this map give $\bar{R}^{(j)}_U$. The linear part, $Q'_j(U)\lambda$, is the mean over $B^j(y)$ of $\lambda$ after transport to $y$.

\emph{Isometries on $1$-forms.} An isometry $\mathsf{r}$ acts by $(\mathsf{r}U)(b)=U(\mathsf{r}^{-1}b)$. Bond functions are specified on the positively oriented bonds, whereas a reflection reverses orientation. Since $(e^{i\mathsf{a}}W_0)(\bar{b})=\exp\big(-i\operatorname{Ad}(W_0(b)^{-1})\mathsf{a}(b)\big)W_0(\bar{b})$, the action relative to $W_0$ must include a twist:

\begin{gather*}
\FitMath{(\mathsf{r}_{W_0}\mathsf{a})(b')=\mathsf{a}(b)\,\text{ if }\mathsf{r}^{-1}b'=b,} \\[3pt]
\FitMath{(\mathsf{r}_{W_0}\mathsf{a})(b')=-\operatorname{Ad}\big(W_0(b)^{-1}\big)\mathsf{a}(b)\,\text{ if }\mathsf{r}^{-1}b'=\bar{b};}
\end{gather*}

on pairs the action is $\mathsf{r}(\mathbf{U},\mathbf{J})=(\mathsf{r}\mathbf{U},\mathsf{r}_{\mathbf{U}}\mathbf{J})$ (1a.1).

\emph{Lift, regions, determining sets.} The \emph{lift} is defined by $(Q^{s*}_kV)(b)=V(c)$ when $b_\pm\in B^k(c_\pm)$; on every bond contained in a single block of order $k$ its value is $\mathbf{1}$. For a bond function $B$ on $\mathbb{T}^{(k)}$, its linear lift is the $1$-form that takes the value $\eta_k^{-1}B(c)$ on these crossing bonds and the value $0$ on bonds lying within a block. The value $\mathcal{M}^q(W)(c)$, and also which rule applies, depend on $W$ only in $B^q(c_-)\cup B^q(c_+)$. Averages are therefore defined for a field given on part of a lattice; they remain covariant and obey the semigroup law $\mathcal{M}^{q+1}W(c')=\mathcal{M}(\mathcal{M}^qW)(c')$. Understood in this way, $M_{\mathbf{B}}(U):=(\mathcal{M}^jU|_{\Gamma_j})_j$ is defined on configurations given on a region, and the problems that define the localised backgrounds $U_{k,\square}$ are set up near a cube without reference to the torus (1a.1).

\emph{The transformation.} Here $\delta$ denotes the Dirac distribution at $\mathbf{1}\in G$ with respect to Haar measure. It is a class function, and it is the limit $\mathfrak{a}\to\infty$ of the normalised weights $\propto e^{-\mathfrak{a}(1-\frac12\operatorname{Re}\operatorname{tr}W)}$. The transformation sends a complex Borel measure $\nu$ to its image $\mathbf{T}\nu=\mathcal{M}_*\nu$, which is written with the kernel $\delta(\mathcal{M}(U)V^{-1})$. Thus a coarse configuration receives the total weight of the fine configurations whose average it is. Wilson and Kadanoff introduced the block-spin transformations $(\mathbf{T}_t\rho)(V)=\int dU\,t(V,U)\rho(U)$, with $t\ge0$, $\int dV\,t(V,U)=1$ and $t(V^{u|_{\mathrm{coarse}}},U^u)=t(V,U)$; the transformation here is the member of this family with a hard constraint. Decompositions of unity and the factor $\mathfrak{G}_k$ are inserted into the integral. They are built with the auxiliary choices $\varepsilon_k,\delta_k,\alpha_{i,k},\check{R}_k$ and with the weight of $\mathfrak{G}_k$. All of these depend on a single number $\mathsf{t}$, which is $x_k$ in the run and $\theta_k$ in a reference run. Only $\check{R}_k$ is formed with the minimum of these numbers taken over all levels $i\le k$, which makes it non-increasing in $k$. A history $h$ records outcomes, and these yield the term $\tau_h$, or the term $\tau^{\mathfrak{w}}_h$ when $\rho_0$ is multiplied by a bounded complex weight $\mathfrak{w}(U)$. In addition to $\chi_k$ and $\chi^{(k)}$ there is the approximate fluctuation test $\chi'_k(\square)=\mathbf{1}\{|V_k(V^{\mathrm{bg}}_{k,\square})^{-1}-\mathbf{1}|<2\delta_k\text{ near }\square\}$, which is built with the local averaged background $V^{\mathrm{bg}}_{k,\square}=\mathcal{M}^kU_{k+1,\square}$. Braces denote a difference between backgrounds: $\{F\}:=F(U_k(V_k))-F(U_{k+1})$.

\emph{The block algebra.} The block field of depth $s$, $V^{(s)}:=\mathcal{M}^{K-s}U$, is defined on $\Lambda_{s,\ell}$, a lattice that does not depend on the cutoff $K\ge s$. The \emph{block algebra} $\mathfrak{A}$ is the set of gauge-invariant functions $F$, bounded and continuous, that depend on only finitely many block variables $V^{(s)}(b)$. The composite $F\circ\mathcal{M}^{K-s}$ is defined for every $K\ge s$, so these observables exist at all cutoffs simultaneously (1a.2). The block algebra is a notion separate from Theorem 0 and Theorem 0', whose observables are the smeared curvature evaluated at the bare scale; neither theorem refers to $\mathfrak{A}$.

\emph{Two operations} (1a.2). Let $\Delta=\mathsf{C}^{\mathsf{T}}\Delta^{(k)}\mathsf{C}$, let $z^{(k)}$ be a product of Jacobians, and let $\chi(S_{k+1})$ denote the tests on $\Omega_{k+1}\setminus\Lambda_{k+1}$. To obtain $\zeta(\Omega^c_{k+1})$, form for each outcome that determines this same $\Omega_{k+1}$ the product of the tests and gauge-fixing factors localised in $\Omega_k\setminus\Omega_{k+1}$, and add these products over all such outcomes. The one-step operation collects everything that step $k+1$ does not carry out:

\begin{gather*}
\FitMath{\mathbf{T}^{(k)}(Z_{k+1})F=\int dV_k\big|_{\Omega_{k+1}^c}\delta(\mathcal{M}(V_k)V_{k+1}^{-1})\,\zeta(\Omega^c_{k+1})\;z^{(k)}\!\int d\mathcal{A}^{\mathrm{s}}_k\big|_{\Omega_{k+1}\setminus\Lambda_{k+1}}\chi(S_{k+1})\,e^{-\frac12\langle\mathcal{A}^{\mathrm{s}}_k,[\Delta-\Delta C^{(k)}(\Lambda_{k+1})\Delta]\mathcal{A}^{\mathrm{s}}_k\rangle}F,}
\end{gather*}

where, more precisely, the Gaussian exponent of the operation is a part of this quadratic form, namely a part that couples no two components of $\Omega_{k+1}\setminus\Lambda_{k+1}$; the remainder of the form is assigned to the boundary terms $B^{(k+1)}$ (2b). As long as $\mathbf{R}$ does not intervene, $\mathbf{T}'_{k+1}=\mathbf{T}^{(k)}(Z_{k+1})\mathbf{T}'_k$. For $x\in\mathbb{C}$ and a function $S$ of the bare fields, we define the pure small-field integral and then the regular part of step $k+1$ by

\begin{gather*}
\FitMath{(\mathbf{T}^{\mathrm{sf}}_k[\mathsf{t}]F)(W)=\int dV\,\delta(\mathcal{M}(V)W^{-1})\,\mathfrak{G}_k(V)\,\chi^{(k)}(T_\ell)\,F(V),}
\end{gather*}


\begin{gather*}
\FitMath{\mathfrak{E}_{k+1}[\mathsf{t};x;S](W)=\log\big(\mathbf{T}^{\mathrm{sf}}_k[\mathsf{t}]\,e^{-xA(U_k(\cdot))+S(U_k(\cdot))}\big)(W)+xA(U_{k+1}(W))-S(U_{k+1}(W)).}
\end{gather*}


\emph{Regular configurations.} We consider sequences $\Omega_0\supset\dots\supset\Omega_k$, with $\Omega_{k+1}=\emptyset$, in which each $\Omega_j$ is a union of $M_1$-cubes of level $j$ and $\operatorname{dist}(\Omega_j^c,\Omega_{j+1})\ge CM_1L^j\eta_k$. The small-field case is the constant sequence $\Omega_j=T_\ell$. Layer $j$ is the set $\Omega_j\setminus\Omega_{j+1}$. It has the local scale $\xi_j=L^j\eta_k$ and the tops $\mathbb{Y}_j=\Lambda_j(\mathbf{B})$, which are its sites in $\mathbb{T}^{(j)}$. Its pillars are the $B^j(y)$, $y\in\mathbb{Y}_j$, and over all $j\le k$ the pillars partition the bare sites of $\Omega_0$. We define

\begin{gather*}
\FitMath{\mathfrak{U}_k(\Omega_\bullet,\alpha)=\big\{U:\,|\partial U-\mathbf{1}|<\alpha\,\eta_k^2\xi_j^{-2},\,\,|J(U)|<\alpha\,\xi_j^{-3}\,\text{ on layer }j\big\},} \\[3pt]
\FitMath{|f|_{(-s)}=\sup_j\sup_{\text{layer }j}\xi_j^s|f|,}
\end{gather*}

together with 

\begin{gather*}
\FitMath{\|A\|_1=\max\{|A|_{(-1)},|\nabla_{U_0}A|_{(-2)}\}}
\end{gather*}

 and 

\begin{gather*}
\FitMath{\|\lambda\|_0=\max\{\sup|\lambda|,|D_{U_0}\lambda|_{(-1)}\}}
\end{gather*}

. Both conditions are invariant under gauge transformations. Combined with the Bianchi identity, the second condition is an elliptic system for the field strength $F_U$, and all bounds on derivatives come from this system; minimisers satisfy the second condition. The complex neighbourhood of radius $\alpha_1$ is the set of $\mathbf{U}=e^{i\eta_kA'}U$ with $U$ real and regular and $\|A'\|_1<\alpha_1$. Since $\|A'\|_1$ bounds only first differences, it does not control the current of such a $\mathbf{U}$, which is a second difference. For this reason the current enters the domains $U^c_j(X,\alpha_0,\alpha_1)$ as an independent variable $\mathbf{J}$ with $|\mathbf{J}|<\alpha_0$. The real centre $U$ is required to have small curvature and small local potentials, but not a small current. The backgrounds of the averages of $\mathbf{U}$ are required to satisfy what the condition on the current provided, and $U^c_j$ is the set of all such pairs. This set is not saturated with respect to the $G^c$-valued $u$, which act by 

\begin{gather*}
\FitMath{(\mathbf{U},\mathbf{J})\to(\mathbf{U}^u,\operatorname{Ad}(u)\mathbf{J})}
\end{gather*}

. The extent to which a term is invariant under these transformations is given in 1a.3.10 and in (1d).

\emph{Gauges.} Because $G$ is compact, the gauge is not fixed in order to obtain finiteness. The reason for fixing it is that the integrand of $\mathbf{T}$ does not change under the block group $\mathcal{G}_{\mathrm{bl}}=\{u:u=\mathbf{1}\text{ on }\mathbb{T}^{(k+1)}\}$. Consequently, in each block there are $3(L^4-1)$ directions $D_V\lambda$ along which the Hessian at the minimum vanishes, and there is then no Gaussian measure, no covariance and no decay. Similarly, the residual group $\mathcal{G}^{\mathrm{res}}$ leaves the constraints $\mathcal{M}^jU=V_j$ on $\Gamma_j$ unchanged. This group consists of the bare $u$ that equal $\mathbf{1}$ at the \emph{nodes}, which are the end points of the bonds of $\mathbf{B}[\Omega_\bullet]$. Hence a background is an orbit and not a point. No gauge has all of the following properties at once: locality (on a torus, no global gauge turns small curvature into a small potential), compatibility with the averages, covariance under the remaining gauge group, symmetry under the hypercubic group $W_4$, and regularity uniform in $k$. Axial gauges are algebraic and strictly local and have Jacobian one; they appear in the integrals and serve to select representatives. \emph{Box axial:} the condition $V=\mathbf{1}$ on the tree of combs that issues from the base point $y_P$ of a box. \emph{Block axial:} the condition $V(y,x)=\mathbf{1}$ on the averaged contour variables. It is imposed softly on a set $P''$ of blocks, at those pairs $(y,x)$ for which the mean $V(y,x)$ is defined (no factor is inserted at the other pairs), by means of

\begin{gather*}
\FitMath{\mathfrak{G}_k(V;P'')=\prod_{y\in P''}\prod_{x\in B(y)\setminus y}z_{\mathfrak{G}}^{-1}\,\mathbf{1}\{|V(y,x)-\mathbf{1}|<\varepsilon_k\}\,e^{-x_k(1-\frac12\operatorname{Re}\operatorname{tr}V(y,x))},} \\[3pt]
\FitMath{\mathrm{GF}_k(V)=\sum_{y,x}\big[1-\tfrac12\operatorname{Re}\operatorname{tr}V(y,x)\big]\ge0.}
\end{gather*}

\emph{Tower axial} relative to $U_0$, written $\operatorname{Ax}(\Omega_\bullet,U_0)$: in every pillar and at every level $n<j$, the contour variables of $\mathcal{M}^nU$ coincide with those of $\mathcal{M}^nU_0$. When the tops are the only nodes, solving this condition from the top down selects one representative in each orbit of $\mathcal{G}^{\mathrm{res}}$. An example, for $U_0=\mathbf{1}$, is the representative of the orbit $U_k(\mathbf{V})$ of the regular critical configurations of (1b). Suppose that an inward bond of layer $j-1$ is retained as an additional constraint (see (1b), \emph{Inward bonds}; by the Notation, \ensuremath{\S}4, the determining set contains no such bond). Its end point then lies inside a pillar of layer $j$. That end point is a node but not a top, and it is called a \emph{junction node}; in (1b) it is the seam node. The transport from the top to a junction node is invariant under $\mathcal{G}^{\mathrm{res}}$, so the condition cannot be achieved there. Leaving out these conditions gives the gauge $\operatorname{Ax}^\circ(\Omega_\bullet,U_0)$, which meets in exactly one point every orbit where the contour variables are defined. Representatives are covariant. Given a gauge transformation $v$ on the nodes, its lift is the bare transformation $\tilde{v}$ whose value at a bare site is the value of $v$ at the first node reached when one ascends from that site through the centres of its blocks. The map $W\mapsto W^{\tilde{v}}$ then sends the constraint set of $\mathbf{V}$ onto the constraint set of $\mathbf{V}^v$; it also sends orbits to orbits and maps $\operatorname{Ax}^\circ(\Omega_\bullet,\mathbf{1})$ onto itself. Hence, for the representative taken at the reference $U_0=\mathbf{1}$, $U_k(\mathbf{V}^v)=U_k(\mathbf{V})^{\tilde{v}}$. For nests of boxes whose lattices are of different levels (2b), the tree consists of shell trees, connected by one entrance bond in each box. Trees depend on the choice of a frame in $W_4$, and averaging over the $2^4\,4!$ frames restores the symmetry (1a.3).

The estimates are carried out in the elliptic Landau gauge, and this gauge comes with a determinant. We put $Q^{\sharp}=(Q^{\sharp}_j|_{\Gamma_j})_j$, with each component normalised to total weight $1+O(\alpha_0)$, so that $\mathcal{Q}^{\sharp}_j(U_0;A)=\xi_jQ^{\sharp}_jA+C_j(A)$. Further, $Q'=(Q'_j|_{\mathbb{Y}_j})_j$, and its null space is denoted $N(Q')$. On layer $j$ we set $\xi=\xi_j$ and $a=\xi_j^{-2}$. We also put $G'=(\Delta_{U_0}+Q'^*aQ')^{-1}$, and $R$ is the orthogonal projection onto $\Delta_{U_0}N(Q')$. Relative to $U_0$, a \emph{Landau pair} $(A,u)$ of $U$ is a pair satisfying

\begin{gather*}
\FitMath{U^{u^{-1}}=e^{i\eta_kA}U_0,} \\[3pt]
\FitMath{R\,D^*_{U_0}A=0,} \\[3pt]
\FitMath{\bar{R}^{(j)}_{U_0}u=\mathbf{1}\,\text{ on }\mathbb{Y}_j,}
\end{gather*}

in which conditions on averages take the place of the condition $u=\mathbf{1}$ at the nodes, since that condition cannot be handled analytically. This has a cost. Transformations $u$ of this kind do not form a group: $\bar{R}u_1=\bar{R}u_2=\mathbf{1}$ does not imply $\bar{R}(u_1u_2)=\mathbf{1}$. The word ``orbit'' therefore has no meaning, and the set that such transformations generate depends on its base point. For this reason the operations come in a fixed order. First, the tower axial gauge $\operatorname{Ax}^\circ$ selects a base point in each orbit of $\mathcal{G}^{\mathrm{res}}$. Next, the set is generated from that base point. Finally, the set is intersected with the slice. The average that belongs to this gauge is the \emph{reduced average} 

\begin{gather*}
\FitMath{\mathcal{M}_{\mathrm{red}}(W;U_0)=\mathcal{M}(W)^{\mathsf{n}_W^{-1}}\mathcal{M}(U_0)^{-1}=e^{i\mathcal{Q}^1(U_0,\mathsf{a})}}
\end{gather*}

; here $\mathsf{n}_W$ stands for the block factor $\mathsf{u}(U_0,\mathsf{a})$ taken at $W=e^{i\mathsf{a}}U_0$. The reduced average is the straightened average, under the name it bears in the theory of gauges. The contour variable is required to be the same in all four objects $\mathcal{M}$, $\mathsf{n}_W$, $\bar{R}$ and $\operatorname{Ax}$: it must be the mean in each of them, or the comb in each of them.

\textbf{Statements.} \textbf{Lemma 1a.1.1 (the mean).} Consider families $V_\bullet=(V_1,\dots,V_n)$ in $G^c$ of small diameter $\operatorname{diam}V_\bullet:=\max_{i,i'}|V_iV_{i'}^{-1}-\mathbf{1}|$, where the radii do not depend on $n$. For such families the mean near the family is unique. It is holomorphic, with bounds independent of $n$ (by the contraction principle), and it is symmetric. On families in $G$ it takes values in $G$, and there it is the only mean lying in $G$. Here a mean of $V_\bullet$ is an element $V$ with $\max_i|V_iV^{-1}-\mathbf{1}|<1$ and $\sum_i\log(V_iV^{-1})=0$; the logarithm is the series logarithm, which has a meaning because of the first condition. Next, $\operatorname{av}(TV_\bullet T')=T\operatorname{av}(V_\bullet)T'$ holds for $T,T'\in G$. It also holds for $T\in G^c$ as long as $\operatorname{cn}(T)\operatorname{diam}V_\bullet$ remains below the radius, where $\operatorname{cn}(T):=|T|\,|T^{-1}|$ is the condition number. The identity $\operatorname{av}(V_\bullet^{-1})=\operatorname{av}(V_\bullet)^{-1}$ holds in $G$. It holds in $G^c$ when $\operatorname{diam}V_\bullet$ is below half the radius and $\operatorname{cn}(V_i)\le2$ for all $i$. Finally, $\frac1i\log\operatorname{av}(e^{iA_\bullet})=\frac1n\sum_jA_j+O(\max_j|A_j|^2)$.

\textbf{Lemma 1a.1.2 (small curvature: nearby transports, comb gauge).} Consider a box that contains at most $s$ sites along each direction and has plaquette size $\varpi$. Suppose the $m$ steps of a path are rearranged so that the rearranged paths stay inside the box. Then the transport along the path changes by at most $(1+\varpi)^{m(m-1)/2}-1$. Moreover, with $h(x)=W(\Gamma_{y,x})$,

\begin{gather*}
\FitMath{|W^h(b)-\mathbf{1}|\le|b_--y|_1\,\varpi,} \\[3pt]
\FitMath{W^h=e^{i\mathsf{a}^0},} \\[3pt]
\FitMath{|\mathsf{a}^0|\le Cs\varpi.}
\end{gather*}

Consequently the contour families on $\mathrm{Dom}(c)$ are of diameter $\le CL^2\varpi_c$, and one has $\mathcal{M}(W)(c)=e^{iY}W[c]$ where $|Y|\le CL^2\varpi_c$.

\textbf{Proposition 1a.1.3 (the averaging map).} Let $U$, $u$ be arbitrary, and let $\mathsf{r}$ be any isometry that preserves $\mathbb{T}^{(k+1)}$ (translations, $W_4$, coordinate and diagonal reflections). Then

\begin{gather*}
\FitMath{\mathcal{M}(U^u)=\mathcal{M}(U)^{u|_{\mathrm{coarse}}},} \\[3pt]
\FitMath{\mathcal{M}(\mathsf{r}U)=\mathsf{r}\,\mathcal{M}(U);}
\end{gather*}

and $\mathcal{M}^n(U)(c)$ depends on $U$ only in $B^n(c_-)\cup B^n(c_+)$. The averaging map has four further properties.
\begin{itemize}
\item It is Borel. Each of the three rules applies on a Borel set, and $\mathcal{M}$ is continuous on each of these sets but not across their boundaries. Outside $\mathrm{Dom}(c)$ no further regularity is required.
\item It is holomorphic at $e^{i\mathsf{a}}W_0$ for $\varpi_c(W_0),|\mathsf{a}|<c_L$.
\item It is onto: $\mathcal{M}^k(Q^{s*}_kV)=V$.
\item On small fields of one colour it is linear: $\mathcal{M}(e^{ifT})=e^{i(\mathsf{Q}^{\mathrm{pl}}f)T}$ whenever $T\in\mathfrak{g}$, $|T|=1$, and $f$ is real with $\sup|f|<c_L$.
\end{itemize}

For a $G^c$-valued $g$, covariance is obtained at the expense of room in the domain, and the amount of room lost is $\operatorname{cn}(g)$. Precisely, consider the graded domains $\mathrm{Dom}^c_\theta(c)$, $0<\theta\le1$, on which the contour families and the $W_x$ lie within $\theta$ times the radii. If $\operatorname{cn}(g(c_\pm))\le\varkappa$ and $\varkappa\theta\le1$, then $\mathbf{W}^g\in\mathrm{Dom}^c_{\varkappa\theta}(c)$ and

\begin{gather*}
\FitMath{\mathcal{M}(\mathbf{W}^g)(c)=g(c_-)\,\mathcal{M}(\mathbf{W})(c)\,g(c_+)^{-1}.}
\end{gather*}


\textbf{Proposition 1a.1.4 (one averaging multiplies curvature by $L^2$).}

\begin{gather*}
\FitMath{|\partial\mathcal{M}(W)(p')-\mathbf{1}|\le L^2\varpi+C(L^2\varpi)^2;} \\[3pt]
\FitMath{|\partial U-\mathbf{1}|\le\alpha_0\eta_k^2,\,\alpha_0\le c} \\[3pt]
\FitMath{\Rightarrow} \\[3pt]
\FitMath{\varpi(\mathcal{M}^nU)\le2\alpha_0(L^n\eta_k)^2} \\[3pt]
\FitMath{(n\le k),}
\end{gather*}

and $\mathcal{M}^nU(c)$ differs from the straight bare transport by at most $C\alpha_0(L^n\eta_k)^2$. Consequently the first rule is applicable on every level, and on $\{|\partial U-\mathbf{1}|<c\,\eta_k^2\}$ the map $\mathcal{M}^k$ is real-analytic.

\textbf{Proposition 1a.1.5 (one step relative to a reference).} Let $\varpi=\varpi(W_0)\le c_L$, and consider complex $|\mathsf{a}|<c_L$. Then $\mathcal{Q}^1$ and $\frac1i\log\mathsf{u}$ depend holomorphically on $\mathsf{a}$, are covariant and do not exceed $C|\mathsf{a}|$ in size; moreover

\begin{gather*}
\FitMath{\mathcal{Q}^1(W_0,\mathsf{a})=\mathsf{Q}(W_0)\mathsf{a}+C^1(W_0,\mathsf{a}),} \\[3pt]
\FitMath{\mathsf{Q}(\mathbf{1})=\mathsf{Q}^{\mathrm{s}},} \\[3pt]
\FitMath{\|\mathsf{Q}(W_0)-\mathsf{S}(W_0)\|_{\mathrm{row},\,\mathrm{col}}\le C\varpi,}
\end{gather*}


\begin{gather*}
\FitMath{\|\mathsf{S}(W_0)\|_{\mathrm{row}}\le L,} \\[3pt]
\FitMath{\|\mathsf{S}(W_0)\|_{\mathrm{col}}\le L^{-3},} \\[3pt]
\FitMath{|C^1|\le C|\mathsf{a}|^2.}
\end{gather*}


\textbf{Theorem 1a.1.6 ($k$ steps, uniformly in $k$).} Assume $|\partial U_0-\mathbf{1}|\le\alpha_0\eta_k^2$ with $\alpha_0\le c$, and let $A$ be complex with $\sup|A|<c$. Here $c,C$ depend only on $L$. The following hold holomorphically in $A$ and in the complex shifts $\mathbf{U}=e^{i\eta_kA'}U_0$ of $U_0$ with $\sup|A'|<\alpha_1\le c$. The bounds are uniform in $A'$; in the bounds on the row norm and on the functional derivative, $\alpha_0+\alpha_1$ then replaces $\alpha_0$.

\begin{gather*}
\FitMath{\mathcal{M}^n(e^{i\eta_kA}U_0)=\big(e^{i\mathsf{a}_n}\mathcal{M}^nU_0\big)^{v_n},} \\[3pt]
\FitMath{|\mathsf{a}_n|\le2L^n\eta_k\sup|A|,} \\[3pt]
\FitMath{|v_n-\mathbf{1}|\le CL^n\eta_k\sup|A|,}
\end{gather*}


\begin{gather*}
\FitMath{Q^{\sharp}_k(U_0)=\eta_k\,\mathsf{Q}(W_{0,k-1})\cdots\mathsf{Q}(W_{0,0}),} \\[3pt]
\FitMath{\|Q^{\sharp}_k(U_0)\|_{\mathrm{row}}\le1+C\alpha_0,} \\[3pt]
\FitMath{|C_k|\le C\sup|A|^2,}
\end{gather*}


\begin{gather*}
\FitMath{Q_k(U_0)A=Q^{\sharp}_k(U_0)A-D_{\mathcal{M}^kU_0}\dot{v}_k(A),} \\[3pt]
\FitMath{|Q_kA|\le C\sup|A|,} \\[3pt]
\FitMath{|C^{\mathrm{pl}}_k|\le C\sup|A|^2:}
\end{gather*}

thus a coarse gradient separates the two linear parts, and this gradient does not become small as the curvature does. More generally, let $h$ be $G^c$-valued with $\operatorname{cn}(h)\operatorname{cn}(v_n)\le2$. Then $(e^{i\eta_kA}U_0)^h$ has well-defined averages, and $\mathcal{M}^n\big((e^{i\eta_kA}U_0)^h\big)=\big(e^{i\mathsf{a}_n}\mathcal{M}^nU_0\big)^{hv_n}$. This ball is much larger than the regular class, because the plaquettes of $e^{i\eta_kA}U_0$ are $O(\eta_k)$ rather than $O(\eta_k^2)$. Accordingly, the validity of the first rule on every level and the analyticity are consequences of the factorisation, not of 1a.1.4. The remainder decomposes as $C_k=C_k^{(2)}+C_k^{(3)}+\cdots$, where

\begin{gather*}
\FitMath{C_k^{(2)}(U_0,A)=\sum_{n<k}\mathsf{Q}(W_{0,k-1})\cdots\mathsf{Q}(W_{0,n+1})\,C^{1,(2)}\big(W_{0,n},\mathsf{a}^{\mathrm{lin}}_n\big),} \\[3pt]
\FitMath{|n\text{-th term}|\le CL^{n-k}\sup|A|^2:}
\end{gather*}

the non-linearity consists of commutators and is concentrated at the coarsest of the $k$ scales. The derivative $K_A(c,b):=\partial\mathcal{Q}^{\sharp}_k(c)/\partial A(b)$ satisfies $\|K_A\|_{\mathrm{col}}\le2\eta_k^4$ and the Lipschitz bound $\|K_A-K_{A'}\|_{\mathrm{col}}\le C\eta_k^4\sup|A-A'|$. Consequently $1+C\alpha_0+C|A|$ is a bound for the functional derivative $\eta_k^{-4}K_A$, the adjoint $K_A^*$ is a bounded operator for the sup norm, and $K_A$ acts boundedly from $L^2$ to $\ell^2$. The Lagrange equations of (1b) involve $K_A^*$. Complex regular fields $\mathbf{U}=e^{i\eta_kA'}U_0$ with $\sup|A'|<\alpha_1\le c$ have averages that are regular, each on its own scale: 

\begin{gather*}
\FitMath{|\partial\mathbf{U}-\mathbf{1}|\le\alpha_0'\eta_k^2,\,\alpha_0'\le c\,\Rightarrow|\partial\mathcal{M}^n\mathbf{U}-\mathbf{1}|\le2\alpha_0'(L^n\eta_k)^2}
\end{gather*}

. From $|\nabla_{U_0}A'|<\alpha_1$ one obtains only $\alpha_0'=\alpha_0+C\alpha_1$. That is why the domains $U^c_j$ require, as a separate condition, that the complex plaquettes be small; averaging reproduces exactly this condition.

All of this holds on determining sets, layer by layer. At the bare bonds of layer $j$ the local scale is $\xi(b)=\xi_j=L^j\eta_k$; it is the local scale of layer $j$ in the sense of ``Regular configurations''. On determining sets the ceilings on $\alpha_0$ and on $\sup_b\xi(b)|A(b)|$ are lowered by the factors $L^{-2}$ and $L^{-1}$. At a bond of $\Gamma_j$ that lies astride a seam, one factor is lost. Writing $K_A(c,b)=\partial\mathsf{a}_j(c)/\partial A(b)$ for $c\in\Gamma_j$, one has

\begin{gather*}
\FitMath{|\partial\mathcal{M}^jU_0-\mathbf{1}|\le2L^2\alpha_0,} \\[3pt]
\FitMath{|\mathsf{a}_j|\le2L\sup_b\xi(b)|A(b)|,} \\[3pt]
\FitMath{\sup_c\sum_b\frac{|K_0(c,b)|}{\xi(b)}\le2L,} \\[3pt]
\FitMath{\sum_j\sum_{c\in\Gamma_j}|K_A(c,b)|\le\frac{C\eta_k^4}{\xi(b)^3},}
\end{gather*}

and this factor is lost once, without compounding; it cannot be removed (1a.1.14). The last bound is a sum over the levels at which a given bare bond enters.

\textbf{Propositions 1a.1.7, 1a.1.8 (linear averages; gauge transformations).} The operator 

\begin{gather*}
\FitMath{Q^{\mathrm{st}}_k:=Q^{\mathrm{st}}_k(\mathbf{1})=Q^{\sharp}_k(\mathbf{1})=\eta_k(\mathsf{Q}^{\mathrm{s}})^k}
\end{gather*}

 is the straight average taken over the blocks of order $k$, and

\begin{gather*}
\FitMath{(Q^{\mathrm{st}}_kA)(c)=\eta_k^5\sum_{x\in B^k(c_-)}A([x,x']),} \\[3pt]
\FitMath{u_k(p)=\prod_\mu\frac{\eta_k(e^{ip_\mu}-1)}{e^{i\eta_kp_\mu}-1},} \\[3pt]
\FitMath{c\le Q^{\mathrm{st}}_kQ^{\mathrm{st}*}_k\le1,}
\end{gather*}

where $u_k$ coincides with the Fourier symbol of $Q'_k$ up to a phase factor. The normalisation is as follows. Constants are mapped to the same constants, and $\|Q^{\mathrm{st}}_k\|_{\mathrm{row}}=1$. Composing $Q^{\mathrm{st}}_k$ with the linear lift gives the identity, and $Q'_kQ_k'^*=1$. The adjoint does not coincide with the lift. The operator $Q^{\mathrm{st}}_kQ^{\mathrm{st}*}_k$ is not a multiple of the identity: it is a convolution that couples each coarse bond to its two neighbours on the same line. This is why the display gives only a two-sided bound. In a background one has $|Q^{\sharp}_k(U_0)A|\le Q^{\mathrm{st}}_k|A|+C\alpha_0Q''_k|A|$, where $Q''_k$ takes the mean of $|A|$ over every bond in the double block of order $k$. With gradients $(d\lambda)(b)=\lambda(b_+)-\lambda(b_-)$, the two averages intertwine gradients in different ways:

\begin{gather*}
\FitMath{\mathsf{Q}^{\mathrm{pl}}d\lambda=d(\lambda|_{\mathrm{coarse}}),} \\[3pt]
\FitMath{\mathsf{Q}^{\mathrm{s}}d\lambda=d(q\lambda),} \\[3pt]
\FitMath{Q^{\sharp}_k(U)D_U\lambda=D_{\mathcal{M}^kU}\,Q'_k(U)\lambda,}
\end{gather*}

the last of these relations being exact and valid for the straightened average only. In a background the plain average satisfies, likewise exactly, $Q_k(U)D_U\lambda=D_{\mathcal{M}^kU}(\lambda|_{\mathbb{T}^{(k)}})$, where the restriction stands in place of $Q'_k$. For the main term $Q^{\mathrm{st}}_k(U)$, by contrast, the relation involving $Q'_k$ is valid only up to $O(\alpha_0)$.

It follows that there are two residual invariances. The constraint $\mathcal{M}^kU=V$ is preserved by every transformation with $u=\mathbf{1}$ on $\mathbb{T}^{(k)}$. The straightened constraint is preserved by the $w$ with $r_k=\mathbf{1}$, where $r_k$ is defined below. In linearised form this condition is $Q'_k\lambda=0$. When $A=0$, and likewise at a field tower axial relative to $U_0$, these transformations are precisely the $w$ whose averages are trivial. Of the two sets, only the first forms a group. The set generated by the second depends on its base point, and the tower axial gauge selects that base point in each orbit of the first. The two invariances correspond to each other through $v_k[A]$.

Now let $A$ be real with $\sup|A|\le c$. Put $w=e^{i\lambda}$ and $e^{i\eta_kA^w}U_0:=(e^{i\eta_kA}U_0)^w$, and for a site function assume $\|\lambda\|_1:=\max\{\sup|\lambda|,\sup|D_{U_0}\lambda|\}\le c$. This norm is taken in level-$k$ units, without local scales, and it differs from the norm $\|\cdot\|_1$ of $1$-forms. It equals $\|\lambda\|_0$ for the constant sequence $\Omega_j=T_\ell$, whose only layer is $k$, with $\xi_k=1$. The straightened averages change by $r_n$, where $r_0=w$ and $r_{n+1}=\mathsf{u}(W_{0,n},\mathsf{a}_n(A^w))^{-1}\,r_n\,\mathsf{u}(W_{0,n},\mathsf{a}_n(A))$:

\begin{gather*}
\FitMath{e^{i\mathsf{a}_n(A^w)}W_{0,n}=\big(e^{i\mathsf{a}_n(A)}W_{0,n}\big)^{r_n},} \\[3pt]
\FitMath{w|_{\mathbb{T}^{(n)}}\,v_n[A]=v_n[A^w]\,r_n,} \\[3pt]
\FitMath{r_n|_{A=0}=\bar{R}^{(n)}_{U_0}w,} \\[3pt]
\FitMath{\tfrac1i\log r_k=Q'_k(U_0)\lambda+O\big((\sup|A|+\|\lambda\|_1)\|\lambda\|_1\big).}
\end{gather*}

We say that $u$ is of class $\alpha$ when each $\bar{R}^{(j)}u$ is within $\alpha$ of $\mathbf{1}$ and its variation inside any block is less than $\alpha L^{j+1}\eta_k$. If $u_1,u_2$ are of class $\alpha\le c$, the averages are multiplicative up to second order,

\begin{gather*}
\FitMath{\bar{R}^{(j)}(u_1u_2)=\bar{R}^{(j)}u_1\,\bar{R}^{(j)}u_2\,e^{i\mathsf{z}_j},} \\[3pt]
\FitMath{|\mathsf{z}_j|\le C(\alpha L^j\eta_k)^2,}
\end{gather*}

and $u_1u_2$ is of class $2\alpha+C\alpha^2$. The averages of products depend holomorphically on complex $\lambda$. Precisely, let $u'=e^{i\lambda}$ with $\|\lambda\|_1\le\alpha'\le c$, and let $u_1$ be of class $\alpha$. Then the quotients $\tilde{u}^{(j)}:=\bar{R}^{(j)}(u'u_1)(\bar{R}^{(j)}u_1)^{-1}$ have four properties: they are holomorphic in $\lambda$; they lie within $C\alpha'$ of $\mathbf{1}$; their covariant differences are $\le C\alpha'L^j\eta_k$; and $\frac1i\log\tilde{u}^{(k)}=Q'_k(U_0)\lambda+O((\alpha+\alpha')\alpha')$. 1a.3 uses this to show that the Landau gauge depends analytically on the fields.

\textbf{Lemma 1a.1.9 (rigidity of one step).} Suppose that $W,W_0$ have plaquette size $\le\varpi$, that their contour variables agree up to a defect $\delta$, and that $|\mathcal{M}(W)\mathcal{M}(W_0)^{-1}-\mathbf{1}|\le\alpha_1$. Then

\begin{gather*}
\FitMath{|W(b)W_0(b)^{-1}-\mathbf{1}|\le\alpha_1+C\varpi+2\delta}
\end{gather*}

at every bond of the double block. Put differently, modulo the curvature, each bond variable is a function of the contour variables together with the average:

\begin{gather*}
\FitMath{\big|W(b)-W(y,b_-)^{-1}W(y,b_+)\big|\le C\varpi,} \\[3pt]
\FitMath{\big|W(b')-W(c_-,b'_-)^{-1}\,\mathcal{M}(W)(c)\,W(c_+,b'_+)\big|\le C\varpi,}
\end{gather*}

for $b$ inside $B(y)$ with $y=c_\pm$, and for $b'$ in the corridor of $c$, that is, among the bonds parallel to $c$ joining $B(c_-)$ to $B(c_+)$. The defect accounts for two things. One is what the soft gauge fixing leaves, which gives $|V(y,x)-\mathbf{1}|<\varepsilon_k$ rather than equality. The other is what a seam leaves (1a.3.11).

\textbf{Proposition 1a.1.10 (the fibre is curved; the current enters).} Let $|\partial U-\mathbf{1}|\le\alpha_0\eta_k^2$ with $\alpha_0\le c$ (the range of Theorem 1a.1.6). Let moreover $U\in\mathfrak{U}_k(\Omega_\bullet,\alpha_0)$ at the constant sequence $\Omega_j=T_\ell$; this membership brings in the further condition $|J(U)|<\alpha_0$. Suppose that $M_1$, $M\alpha_0$ and the number of sites of $\mathbb{T}^{(k)}$ satisfy the hypotheses of Proposition 1a.3.5 [to be proved; Package A], which provides $G$ and $\mathsf{H}_k$. Finally, let $A$ be small. On $Q^{\sharp}_k(U)A+C_k(U,A)=0$, in the Landau gauge, one has $A=A_1-\mathsf{H}_kC_k^{(2)}(A_1)+\dots$ with $Q^{\sharp}_k(U)A_1=0$, and

\begin{gather*}
\FitMath{A(e^{i\eta_kA}U)=A(U)+\langle J(U),A_1\rangle+\tfrac12\langle A_1,\Delta^{\mathrm{W}}_UA_1\rangle-\langle\mathsf{H}_kC_k^{(2)}(A_1),J(U)\rangle+O(|A_1|^3).}
\end{gather*}

In this expansion $\mathsf{H}_k$ denotes the right inverse $\mathsf{H}=GQ^*(QGQ^*)^{-1}$ of $Q=Q^{\sharp}_k(U)$, built from the $G$ of (1c); with this choice the relation $RD^*GQ^*=0$ is exact. For the derivative of $V\mapsto U_k(V)$ the corresponding objects are $G_1$, $\mathsf{H}_1$. The term $O(|A_1|^3)$ collects all terms of order $\ge3$ in $A_1$. Indeed, the constraint is solved up to second order, and the solution is substituted into the Taylor expansion of the action.

\textbf{Proposition 1a.1.11 (isometries and linearised objects).} Suppose that $\mathsf{r}$ preserves every lattice involved. Then $\mathsf{r}(e^{i\mathsf{a}}W_0)=e^{i\mathsf{r}_{W_0}\mathsf{a}}\,\mathsf{r}W_0$, and

\begin{gather*}
\FitMath{\mathcal{Q}^1(\mathsf{r}W_0,\mathsf{r}_{W_0}\mathsf{a})=\mathsf{r}_{\mathcal{M}(W_0)}\mathcal{Q}^1(W_0,\mathsf{a}),} \\[3pt]
\FitMath{Q^{\sharp}_k(\mathsf{r}U)\,\mathsf{r}_U=\mathsf{r}_{\mathcal{M}^kU}\,Q^{\sharp}_k(U),} \\[3pt]
\FitMath{J(\mathsf{r}U)=\mathsf{r}_UJ(U),} \\[3pt]
\FitMath{\Delta^{\mathrm{W}}_{\mathsf{r}U}\mathsf{r}_U=\mathsf{r}_U\Delta^{\mathrm{W}}_U,}
\end{gather*}

and the corresponding relations hold for $Q_k$, $Q'_k$, $\bar{R}$, $\mathsf{u}$. For real $W_0$, the map $\mathsf{r}_{W_0}$ is an isometry both for the sup norm and for $\langle\cdot,\cdot\rangle$. These maps compose according to $(\mathsf{r}\mathsf{r}')_{W_0}=\mathsf{r}_{\mathsf{r}'W_0}\mathsf{r}'_{W_0}$. If the twisted action is replaced by the plain one (the action at $W_0=\mathbf{1}$), these relations in general fail for a reflection reversing bonds. The plain and the twisted action are equal where $\operatorname{Ad}(W_0(b))$ acts as the identity on the reversed bonds $b$. Consequently, propagators, $\mathsf{H}_k$, $\Delta^{(k)}$ and $C^{(k)}$ are covariant under all of $W_4$ only in the twisted sense, and (1c), (1d) and (1e) use the symmetry in this form.

\textbf{Proposition 1a.1.12 (the fundamental equality; the transformation to the tower axial gauge).} Let $U_0$, $A$ be as in 1a.1.6. There exists exactly one $u$, defined on the bare sites, for which $(e^{i\eta_kA}U_0)^u$ is tower axial relative to $U_0$ and $\bar{R}^{(k)}_{U_0}u=\mathbf{1}$. Its averages are given by $\bar{R}^{(n)}_{U_0}u=u\,v_n$, so that $u=v_k^{-1}$ on $\mathbb{T}^{(k)}$, and

\begin{gather*}
\FitMath{\mathcal{M}^k\big((e^{i\eta_kA}U_0)^u\big)=e^{i\mathcal{Q}^{\sharp}_k(U_0;A)}\,\mathcal{M}^kU_0:}
\end{gather*}

once the gauge has been changed in this way, averaging of order $k$ amounts to composing the straightened one-step averages $k$ times. The transformation can be written explicitly. Let $x=x_0,\dots,x_k=y$ be the tower of block centres. Let $\mathsf{w}^{\mathrm{rel}}_j(z,x)=(e^{i\mathsf{a}_j}W_{0,j})(z,x)W_{0,j}(z,x)^{-1}$ be the relative contour variables of the straightened averages, let $\mathsf{u}_j$ be the block factors, and let $\mathsf{T}_{U_0}$ be the transport down the tower. Then

\begin{gather*}
\FitMath{u(x)=\prod_{j=k-1}^{0}\operatorname{Ad}\big(\mathsf{T}_{U_0}(x_{j+1},x)\big)^{-1}\Big[\mathsf{u}_j(x_{j+1})^{-1}\,\mathsf{w}^{\mathrm{rel}}_j(x_{j+1},x_j)\Big],} \\[3pt]
\FitMath{|u-\mathbf{1}|\le C\sup|A|.}
\end{gather*}

Consequently $u$ depends locally and holomorphically on $A$ and is of class $C\sup|A|$. At the bare scale, however, it fails to be regular: it jumps by $O(L^j\eta_k\sup|A|)$ between neighbouring bare sites belonging to different blocks of order $j$.

\textbf{Proposition 1a.1.13 (two schemes; the linear parts compared; what carries over).} Ba{\l}aban's theorems of 1984--85 are proved for blocks that have their site at a corner and for single combs; this is scheme $\mathsf{c}$. The map $\mathcal{M}$ uses centred blocks and means; this is scheme $\mathsf{m}$. Since $L$ is odd, the levelwise translations by $\mathsf{e}_j=\frac12(L^j-1)a_K(1,1,1,1)$ carry the cornered blocks of level $j$ onto centred blocks. Compatibility of the operators across levels is part of the comparison below. Of the contour variable, the proofs of 1a.1.2--1a.1.9 use only the following properties:
\begin{itemize}
\item there is one contour variable for each pair $(y,x)$, and it is covariant;
\item its paths stay in the block and make fewer than $L$ steps in each direction;
\item coarse translations carry the contour variables into one another;
\item it is $G$-valued on $G$;
\item it is holomorphic near $\mathbf{1}$, with a mean of the $\mathsf{a}(\Gamma)$ as its linear part;
\item where the bond variables commute with one another, it is the exponential of that mean.
\end{itemize}

This list is one way of carrying out the axiomatisation that Ba{\l}aban mentions. The results of 1a.1.2--1a.1.9 therefore hold in both schemes, with two exceptions. The first exception is covariance under isometries other than coarse translations. It fails for combs, since under an element of $W_4$ a comb goes over into a comb of a different order, and a lowest corner into a different corner; this failure is precisely why the mean is used. The second exception is the three rules, which are formulated for $\mathcal{M}$ only. When the background is trivial, the two schemes have the same straightened linear part, $\mathsf{Q}_\star(\mathbf{1})=\mathsf{Q}^{\mathrm{s}}$, but different plain linear parts. In a background, let $t$ denote a transport of $U_0$ joining the two base sites of a block. Then

\begin{gather*}
\FitMath{Q^{\sharp}_k(U_0)=\operatorname{Ad}(t)\big(Q^{\mathrm{B}}_k(U_0)+\mathsf{D}_k\big),} \\[3pt]
\FitMath{\|\mathsf{D}_k\|_{\mathrm{row}}\le C\alpha_0,} \\[3pt]
\FitMath{\|\mathsf{D}_k\|_{\mathrm{col}}\le C\alpha_0\eta_k^4,}
\end{gather*}

and the same holds for $Q'_k$. Moreover, $\mathsf{D}_k(c,b)=0$ unless $b$ lies under $c$. In each scheme, the linear part and the straight average of that scheme differ by a kernel of row norm $C\alpha_0$. In this leading part, a bare bond $b$ parallel to $c$ receives as its weight $\eta_k^5$ multiplied by the count of segments $[x,x']$, $x\in B^k(c_-)$, passing through $b$ (a tent along the direction of $c$). Every such segment enters with the transport of $U_0$ taken down a chain of contours and segments. The exact identity $Q_k=Q^{\sharp}_k-D\dot{v}_k$ is supplied in the present text. The identity involving $\mathsf{D}_k$ takes the place of the assertion of 1987 (109, p. 254) that the new average satisfies all the results of the paper on averaging operations. For real $U_0$ the factor $\operatorname{Ad}(t)$ is orthogonal, so it cancels from $Q^{\sharp}_k(U_0)^*\,a\,Q^{\sharp}_k(U_0)$. Consequently, in the two schemes the forms defining $G'$, $R$, $\mathsf{G}_k$, $\mathsf{H}_k$ differ from each other by local forms whose relative size is $C\alpha_0$. Through this comparison, the operators built from $Q^{\mathrm{B}}_k(U_0)$ are used in (1b). The operators are built afresh from $Q^{\sharp}_k(U_0)$, and their bounds are derived from the case $U_0=\mathbf{1}$, in which, after the blocks have been aligned, the two linear parts are equal (1c).

\textbf{Proposition 1a.1.14 (sharpness: fields of one colour).} Let $T\in\mathfrak{g}$ with $|T|=1$. For $U=e^{ifT}$, with $f$ real and small (1a.1.3), everything is linear and explicitly computable. A constant curvature $\varphi$ gives $\partial\mathcal{M}U(p')=e^{iL^2\varphi T}$, so the factor $L^2$ of 1a.1.4 is attained.

Now consider a seam. Let $\Omega_1$ be, locally, a slab $\{x_1\ge\text{const}\}$. Let $z\in\mathbb{T}^{(1)}$ be the last coarse site outside it, and let $p'$ be the coarse plaquette with corners $z,z+Le_1,z+Le_1+Le_2,z+Le_2$. Set $f(\langle x,x+e_2\rangle)=\mathsf{f}^\Sigma(x_1)$ and $f=0$ on the other bonds, where $\mathsf{f}^\Sigma(t):=\sum_{t'<t}\mathsf{f}(t')$. Here $\mathsf{f}=\varphi$ on the $L-1$ bonds in direction $1$ that lie in the block of $z$; on the $L-1$ bonds preceding them the value is $-\varphi$, and everywhere else it is $0$. The curvature has size at most $\varphi$ and vanishes in $\Omega_1$, while $\partial\mathcal{M}U(p')=e^{i\Theta T}$ with $\Theta=\frac12L(L-1)\varphi$. Thus no bound $C\alpha_0$ with $C$ independent of $L$ can take the place of the bound $2L^2\alpha_0$ of 1a.1.6.

At $U_0=\mathbf{1}$, take a bond $c\in\Gamma_j$ astride a seam with outer end $c_-$. Such a bond has one end, its outer end, outside $\Omega_j$; in (1b) it is called an outward bond. The weighted row sum equals $\sum_b\frac{|K_0(c,b)|}{\xi(b)}=\tfrac12L(1-L^{-j})+\tfrac12(1+L^{-j})$, and its limit as $j\to\infty$ is $\frac12(L+1)$. Indeed, along a segment $[x,x']$, a bare bond contained in $B^j(c_-)$ has $\xi(b)=L^{j-1}\eta_k$, while the remaining ones extend into $B^j(c_+)$, where the scale equals $L^j\eta_k$.

For a bond of $\Gamma_j$ astride a seam, an additional block factor arises at its outer end. Moreover, the datum entering the straightened constraint (the reduced datum of (1b)) is not the same as the plain datum. The factor is genuinely present. At $U_0=\mathbf{1}$, take $\eta_kA=tT$ on the single bare bond $\langle z,z+e_\nu\rangle$, where $z$ is the outer end of a bond of $\Gamma_1$, and $A=0$ elsewhere. Then $\mathsf{u}(z)=\exp\big(it\frac{L^4-1}{8L^4}T\big)\ne\mathbf{1}$. Finally, $\|\mathsf{Q}^{\mathrm{pl}}\|_{\mathrm{row}}>2L$.

\textbf{Proposition 1a.2.1 (the exact identities).} The transformation $\mathbf{T}$ preserves mass, positivity ($|\mathbf{T}\nu|\le\mathbf{T}|\nu|$) and gauge invariance. It commutes with the isometries that map $\mathbb{T}^{(k+1)}$ onto itself, and

\begin{gather*}
\FitMath{\mathbf{T}^n\nu=(\mathcal{M}^n)_*\nu,} \\[3pt]
\FitMath{\langle F(\mathcal{M}^kU)\rangle_{K,\ell,\beta}=\int F\,d(\mathbf{T}^k\rho_0)\Big/\int d(\mathbf{T}^k\rho_0).}
\end{gather*}

Consequently, whenever $F\in\mathfrak{A}$ has largest depth $s$, each of the numbers $\langle F\circ\mathcal{M}^{K-s}\rangle_{K,\ell,\beta}$, $K\ge s$, is the integral of one fixed function with respect to the normalised measures $\mathbf{T}^{K-s}\rho_0$, all on the same compact space $\mathcal{V}(\Lambda_{s,\ell})$; this idea stands alongside the two theorems.

\textbf{Proposition 1a.2.2 (soft and hard constraints).} Each kernel $t$ of the kind described above preserves mass, positivity and gauge invariance. In the case of linear site fields and of block means $Q$, the Gaussian kernels $t_{\mathfrak{a}}\propto e^{-\frac{\mathfrak{a}}2|\Psi-Q\Phi|^2}$, taken after the canonical rescaling by $L^{(d-2)/2}$, form a family that is closed under composition. In these kernels $\Phi$, $\Psi$ are site functions in dimension $d$, $Q$ is the rescaled block mean, and $Q_k$ is the mean over $B^k$, rescaled in the same way; the composition law reads

\begin{gather*}
\FitMath{t_{\mathfrak{a}}*t_{\mathfrak{a}_k}=t_{\mathfrak{a}_{k+1}},} \\[3pt]
\FitMath{\mathfrak{a}_{k+1}^{-1}=\mathfrak{a}^{-1}+L^{-2}\mathfrak{a}_k^{-1},} \\[3pt]
\FitMath{\mathfrak{a}_k\to\mathfrak{a}(1-L^{-2}),}
\end{gather*}

and these kernels give rise to the propagators $(\Delta+\mathfrak{a}_kQ_k^*Q_k)^{-1}$, where $\Delta=D^*D$ acts on site functions. For the delta kernel the fibres are nested, and over regular data, up to gauge,

\begin{gather*}
\FitMath{\{\mathcal{M}^{k+1}U=W\}=\bigsqcup_{\mathcal{M}(V)=W}\{\mathcal{M}^kU=V\}} \\[3pt]
\FitMath{\Rightarrow} \\[3pt]
\FitMath{U_k(V^{\mathrm{bg}}_k)=U_{k+1}(W),} \\[3pt]
\FitMath{V^{\mathrm{bg}}_k=\mathcal{M}^kU_{k+1}(W),} \\[3pt]
\FitMath{U_j(\mathcal{M}^jU_k(V))=U_k(V)\,\,(j\le k):}
\end{gather*}

every background therefore coincides with the background of its own averages, and $V^{\mathrm{bg}}_k$, with $\mathcal{M}^kU$ taken in the block axial gauge, is in that gauge the critical point of $V\mapsto A(U_k(V))$ on $\{\mathcal{M}(V)=W\}$. Suppose that $\Delta\ge0$ is positive on $\ker Q$ and that $Q$ is surjective. Then, under the constraint $QA=B$, the form $\frac12\langle A,\Delta A\rangle$ attains its minimum at $\mathsf{H}B$, where $\mathsf{H}=\mathsf{G}Q^*(Q\mathsf{G}Q^*)^{-1}$ and $\mathsf{G}=(\Delta+aQ^*Q)^{-1}$. This holds for every $a>0$. \emph{Remark.} Let $W$ be real and lie in a neighbourhood of $\mathbf{1}$; this neighbourhood may depend on $k$ and on the torus. Then $V^{\mathrm{bg}}_k$ is, in that gauge, the unique point of minimum of $V\mapsto A(U_k(V))$ on $\{\mathcal{M}(V)=W\}$.

\textbf{Proposition 1a.2.3 (the steps with their insertions are exact).} Each operation of $\mathbf{T}$ and of $\mathbf{R}$ consists of finitely many positive kernels, one for each outcome of the operation, whose sum is a probability kernel. These kernels depend on the auxiliary choices. Once the choices are fixed, the kernels depend neither on $\beta$ nor on the density. When the run makes its own choices, the operations themselves vary with $\beta$. For a bounded, gauge-invariant, complex-valued $\mathfrak{w}$,

\begin{gather*}
\FitMath{\tau^{\mathfrak{w}}_h(dV)=\int dU\,\mathfrak{w}(U)\rho_0(U)\,\mathsf{K}_h(U,dV),} \\[3pt]
\FitMath{\sum_h\mathsf{K}_h(U,\mathcal{V}_k)=1.}
\end{gather*}

Consequently $\tau_h\ge0$ and $|\tau^{\mathfrak{w}}_h|\le\sup|\mathfrak{w}|\,\tau_h$. Moreover $\int\rho_k=Z_{K,\ell,\beta}$ and $\int\rho^w_k=Z(w)$ at every level, on a torus that is large compared with the cubes of side $M\check{R}_j$ of every level. This holds irrespective of the gauges and representatives chosen. It also holds for every real sequence of numbers $\mathsf{t}$ among those admitted in (2e). These are the sequences whose terms are all $\ge\gamma^{-2}$, whose differences obey two-sided bounds of the same type as the flow inequalities of 1a.2.7, and for which $\check{R}_j\le L\check{R}_{j+1}$. The sum over outcomes of the operations of $\mathbf{T}$ is $\mathcal{M}_*$. By contrast, $\mathbf{R}$ preserves only those integrals of $F$ that do not depend on the components on which it acts; in general, therefore, $\rho_k\ne\mathbf{T}^k\rho_0$. \emph{Exact insertion:} for every gauge-invariant $F$,

\begin{gather*}
\FitMath{\langle F(\mathcal{M}^kU)\rangle_{K,\ell,\beta}=\sum_h\int F^{(h)}d\tau_h\Big/\sum_h\int d\tau_h,} \\[3pt]
\FitMath{|F^{(h)}|\le\sup|F|,}
\end{gather*}

where $F^{(h)}$ denotes the conditional expectation, under the joint law constructed from the kernels, of $F(\mathcal{M}^kU)$ given $(h,V_k)$. If $F$ depends only on fields away from those components, then $F^{(h)}=F$. \emph{Covariance:} $\mathsf{r}_*\tau_h=\tau_{\mathsf{r}h}$ for every isometry that preserves all the partitions involved. An individual $E^{(j)}(X,\cdot)$ is covariant under the translations that preserve $\pi_j$. An individual $R^{(j)}(X,\cdot)$ is covariant under the translations that preserve the $M\check{R}_j$-cubes. Covariance under all isometries $\mathsf{r}$ of $\mathbb{T}^{(j)}$ holds only for the unrestricted pointed totals $\sum_{X\ni z}E^{(j)}(X,U_j,z)$. These totals are taken at $U_j=U_j(V)$, where $V$ and $\mathsf{r}V$ both belong to a class of real block fields of level $j$ lying near the regular ones (1d). Wherever the covariance of pointed totals is invoked below, it has this meaning.

\textbf{Lemmas 1a.2.4--1a.2.6 (units; the delta function resolved; regrouping).} When the plaquette variables are close to $\mathbf{1}$, $\xi^{d-4}a_p=\frac12\xi^d\langle F_U,F_U\rangle(1+O(\xi^4|F_U|^2))$. Hence, for $d=4$, Wilson's action is a Riemann sum for $\frac12\int|F|^2$ whatever the unit of length. A change of unit leaves $A$, $\rho_k$ and $x_k$ unchanged. On small fields, $\mathcal{M}(e^{i\mathcal{A}}V^{\mathrm{bg}}_k)V_{k+1}^{-1}=\exp i\mathcal{Q}_1(L\mathcal{A})$ and $\mathcal{Q}_1(L\mathcal{A})=LQ\mathcal{A}+C^{\mathrm{pl},(2)}(\mathcal{A})+\dots$. Here $\mathcal{A}$ carries no factor of length. Up to the scaling of its argument by $L$, the map is $\mathcal{Q}_1(\cdot)=\mathcal{Q}_1(V^{\mathrm{bg}}_k;\cdot)$, and this scaling accounts for the factor $L$ in front of its linear part. The average here is the plain $Q=Q_{V^{\mathrm{bg}}_k}$, and $C^{\mathrm{pl},(2)}$ is the second-order part of the plain remainder $C^{\mathrm{pl}}_1(V^{\mathrm{bg}}_k,L\mathcal{A})$. The symbol $C^{(2)}_k$ is reserved for the straightened average.

An analytic substitution $\mathcal{A}=\psi^{\mathrm{lin}}(\mathcal{A}')=\mathcal{A}'-\mathsf{h}\tilde{D}(\mathcal{A}')$ changes only the variables on the pivot bonds. It turns the delta function into the linear one, $\delta(LQ\mathcal{A}')$, which is solved for a single pivot bond $b_0(c)$ for each coarse bond. The substitution comes from the implicit function theorem. In it, $\mathsf{h}$ inverts $LQ$ on the pivot bonds, $LQ\mathsf{h}=1$. The linearising correction $\tilde{D}$ is of second order and is written $\widetilde{\mathsf{D}}$ in (1d). It is the unique small solution of $\mathcal{Q}_1\big(L(\mathcal{A}'-\mathsf{h}\tilde{D}(\mathcal{A}'))\big)=LQ\mathcal{A}'$, obtained as the fixed point of a contraction. For $F$ supported in the analytic chart just constructed, the local coarea identity is

\begin{gather*}
\FitMath{\int_{\rm chart} dV_k\,\delta(\mathcal{M}(V_k)V_{k+1}^{-1})F(V_k)=z^{(k)}\int\prod_{b\ne b_0(c)}d\mathcal{A}'(b)\;e^{\operatorname{Tr}\log(1-\mathsf{h}\,\partial\tilde{D})+\sum_b\log\sigma_{\mathrm{H}}(\mathcal{A}(b))}\,F\big(e^{i\psi^{\mathrm{lin}}(\mathsf{C}\mathcal{A}')}V^{\mathrm{bg}}_k\big),}
\end{gather*}

where the right-hand integral is over the preimage of that chart in the nonpivot coordinates, and $\sigma_{\mathrm{H}}$ is the density of Haar measure in exponential coordinates. A decomposition into charts or the remaining large-field terms is required for the unrestricted integral. The explicit normalization of this local identity is given in Section~\ref{ch:block-average-pivot-chart}. Once this density is divided by its value $\sigma_{\mathrm{H}}(0)$ at $\mathbf{1}$, the Jacobians are $1+O(\delta_k)$ for each bond. They are local and analytic, and they are included in $\mathbb{P}^{(k)}$. The factor $\sigma_{\mathrm{H}}(0)$, taken once for each non-pivot bond, is absorbed into the constant of the step. For every number $b$,

\begin{gather*}
\FitMath{-x_kA+\sum_XE^{(k+1)}(X,\cdot)=-(x_k-b)A+\Big(\sum_X\big[E^{(k+1)}(X,\cdot)-E^{(k+1)}(X,\mathbf{1})\big]-bA\Big)+\sum_XE^{(k+1)}(X,\mathbf{1}),}
\end{gather*}

and, summed over the levels, $-A(x_k(\cdot),U)-\sum_{j\le k}b_jA(\varphi_j,U)=-\beta A(U)$. The constant of a step depends on the regions of that step. Let $|\Omega^{(k)*}|$ be the number of non-pivot bonds of $\mathbb{T}^{(k)}$ in $\Omega$, and $\#_{k+1}(\Omega)$ the corresponding number of sites of $\mathbb{T}^{(k+1)}$. Let $z_{\mathfrak{G}}=\mathfrak{z}(x_k,\varepsilon_k)$ be the normaliser of $\mathfrak{G}_k$ for each block site (1a.3), and $\nu_{\mathfrak{g}}=\log\int_{\mathfrak{g}}e^{-|Y|^2/2}dY$. Then

\begin{gather*}
\FitMath{\mathsf{E}^{(k)}(\Omega,\Lambda)=\big(\log\sigma_{\mathrm{H}}(0)-\tfrac12(N^2-1)\log x_k\big)|\Omega^{(k)*}|-(L^4-1)\log z_{\mathfrak{G}}\,\#_{k+1}(\Omega)+\nu_{\mathfrak{g}}|\Lambda^{(k)*}|+\sum_{X\subset\Lambda}\big[E^{(k+1)}+R^{(k+1)}\big](X,\mathbf{1}),}
\end{gather*}

together with $-\mathcal{E}_k=\sum_{j<k}\mathsf{E}^{(j)}(\Omega_{j+1},\Lambda_{j+1})$ and $E=\sum_{j<K}\mathsf{E}^{(j)}(T_\ell,T_\ell)$. Here $\mathcal{E}_k$ depends on the label, and at the all-small-field label the vacuum constant $E$ equals $-\mathcal{E}_K$.

\textbf{Theorem 1a.2.7 (the invariant form: the programme).} Let $\mathfrak{F}_k$ be the class of measures of the form below. In it the $E^{(j)}$ are pointed families, the pointed totals of which are covariant with respect to every isometry of $\mathbb{T}^{(j)}$, and $b_j$ is their coefficient (1e). A step does not preserve the form and the bounds alone (2e). The measures of the class, with their bounds, are

\begin{gather*}
\FitMath{\sum_{\mathcal{S}}\mathbf{T}'_k(\mathcal{S})\big[\chi_{k,\mathcal{S}}\,e^{A_{k,\mathcal{S}}(U_k(\mathbf{V}))}\big],} \\[3pt]
\FitMath{\|E^{(j)}\|_{j,\kappa}\le E_0,} \\[3pt]
\FitMath{\|R^{(j)}\|_{j,\kappa}\le g_j^{\kappa_0},} \\[3pt]
\FitMath{|B^{(j)}(X,\cdot)|\le\mathsf{B}\,e^{-\kappa d^\natural_j(\mathcal{S};X)},} \\[3pt]
\FitMath{\mathsf{M}(i,X)\le e^{-(2-c')\,p_0(g_i)},}
\end{gather*}

where $\mathsf{B}$ is a constant chosen before $\gamma$ ((2d), (2e)), and where the last inequality, the \emph{mass bound}, concerns those components $X$ that $\mathbf{R}$ has acted on at some level $i\le k$. Such a component was healed, or left standing as an island, or else it dropped out at a failed test, either of a stage or of the main step of $\mathbf{R}$. The proposed induction must derive the drop-out bound from the lower-level island bounds while retaining the Wilson weight of the carrying term over $X$. Until that derivation is verified for the actual operations, the drop-out majorant remains a required input, as in Section~\ref{ch:rg-objects-v2}. Let the constants be chosen in order, and let the torus be large compared with the $M\check{R}_j$-cubes at each level. If, in addition, $g_j\le\gamma$ and

\begin{gather*}
\FitMath{\lambda_\sharp(j-i)-c_5\le x_i-x_j\le B_\sharp(j-i)} \\[3pt]
\FitMath{(i\le j\le k+1),}
\end{gather*}

hold, then $\rho_0\in\mathfrak{F}_0$ and $\mathbf{R}\mathbf{T}\mathfrak{F}_k\subset\mathfrak{F}_{k+1}$ \emph{with the same constants}. No family of a level $j\le k$ is altered; only the background at which they are evaluated changes, $U_k\to U_{k+1}$. The new families come out of $\mathbf{T}$ with a surplus, and $\mathbf{R}$ uses it up: radii $(1+\beta_{\mathrm{s}})\alpha_{i,k}\ge\alpha_{i,k+1}$ and rate $(1+4\beta_{\mathrm{s}})\kappa$, where $\beta_{\mathrm{s}}$ is a fixed fraction. For given $(x_0,\dots,x_j)$, the activities $E^{(j)}(X,\cdot)$, $R^{(j)}(X,\cdot)$ do not depend on the label. When $X$ does not wrap around the torus, they depend neither on the torus nor on the cutoff. Further, $x_{k+1}=x_k-b_{k+1}(x_0,\dots,x_k)$ with $|b_{k+1}|\le B_\sharp$. When the label has no large fields, $|\mathbf{E}_k(U_k)|+|\mathbf{R}_k(U_k)|\le C|\mathbb{T}^{(k)}|$. For a general label, every layer is measured in its own scale, and for each $j\le k$ one adds the sites of level $j$ that lie in a rim around the large-field region $Z_j$ ((1e), (3a)). Finally, stability holds if the flow inequalities hold for all $i\le j\le K$:

\begin{gather*}
\FitMath{\chi_k\,e^{-x_kA(U_k)-E_-|\mathbb{T}^{(k)}|}\le\bar\rho_k,} \\[3pt]
\FitMath{\int\bar\rho_k\le e^{E_+|\mathbb{T}^{(k)}|},} \\[3pt]
\FitMath{E_+\le\tfrac32(N^2-1)\log x_k+O(1).}
\end{gather*}

\emph{The large-field factor} (clause (f) of 1a.2.7; (2c)). The operation $\mathbf{T}_k(Z)$ contains delta functions, tests and conditional Gaussian weights. Apart from the core factors of the islands in $Z$, it contains no Wilson weight. Hence, for a component produced by failed plaquette tests, $\mathbf{T}_k(Z)1$ carries no small factor. If no $\mathbf{R}$ has acted on $Z$, then under the tests of the term it obeys only the bound $\le \exp[C\sum_{t\le k}\#_t(Z_t\cap Z)]$. The small factor lies in three places. They depend on pairwise disjoint sets of plaquettes, factors and bonds, so no part of the small factor is used twice. The first is the Wilson weight: near a failed plaquette test of step $j$, $x_jA\ge c\,p_0(g_j)^2$. The second is the kept gauge-fixing exponents; together with the Wilson sum, they pay for a failed approximate fluctuation test. The third is the Gaussian forms of the fluctuation fields, which pay for a failed fluctuation test.

(1) Fix a history, and a component $X$ of $Z_k$ that has not been acted on by any $\mathbf{R}$. Let $(j_i,\mathfrak{c}_i)$ be the creation cells of the history in $X$, namely those cells of the creation set of step $j_i$, which are $M\check{R}_{j_i}$-cubes, that meet a failed test (2c). Then, for all values of the exterior variables, and with the norm taken as the total variation of that part which the term of the history carries,

\begin{gather*}
\FitMath{\big\|\mathbf{T}_k(X)\big[\chi_{k,\mathcal{S}}\,e^{-A(x_k(\cdot),U_k)}\big]\big\|\le\prod_ie^{-c\,p_0(g_{j_i})^2}\,\exp\Big[C\sum_{t\le k}\#_t(Z_t\cap X)\Big]:}
\end{gather*}

there is one factor for each creation, taken at the coupling of its own step, and one constant for each site.

(2) Bounds of this kind do not multiply over the components, because the background couples the components to one another.

(3) What does multiply, piece by piece, is the pointwise majorant on which (1) rests. A piece $(t,D)$ consists of a component $D$ of the large-field region formed at step $t$, whose enlargement $D^\sim$ shares no cube with any component that $\mathbf{R}$ has acted on at a level in $[t,k]$, together with the stored variables $\xi_{t,D}$ belonging to it and the weight $\mathsf{w}_{t,D}$ of those variables. Throughout the support of the history,

\begin{gather*}
\FitMath{A(x_k(\cdot),U_k)\ge A_{\mathrm{cr}},} \\[3pt]
\FitMath{e^{-A_{\mathrm{cr}}}\prod_{(t,D)}\mathsf{w}_{t,D}\le\prod_{\mathrm{creations}}e^{-c\,p_0(g_t)^2}\prod_{(t,D)}\bar{\mathsf{w}}_{t,D}(\xi_{t,D}),} \\[3pt]
\FitMath{\int\bar{\mathsf{w}}_{t,D}\le e^{C\#_t(D)},}
\end{gather*}

where $A_{\mathrm{cr}}$ denotes the Wilson sums near the creations in these pieces. Only those creations are counted on the right-hand side. Each $\bar{\mathsf{w}}_{t,D}$ is a function of the variables of its own piece only.

(4) After integration, take a whole label together with the list of those outermost components $(i,X)$ that $\mathbf{R}$ has acted on and that still stand in the label. For some history having this label and this list, the bound carries one factor for each of its creation cells lying outside the components that $\mathbf{R}$ has acted on and that still stand in the label (islands and drop-outs; the creations inside such a component are spent in its mass), and one mass for each outermost one among these components:

\begin{gather*}
\FitMath{\int\mathbf{T}'_k(\mathcal{S})\big[\chi_{k,\mathcal{S}}\,e^{-A(x_k(\cdot),U_k)}\big]dV_k\le\exp\Big[C\sum_j(1+\log x_j)\#_j(Z_j)\Big]\prod_{\text{creation cells }(t,\mathfrak{c})}e^{-c\,p_0(g_t)}\prod_{(i,X)}\mathsf{M}(i,X),}
\end{gather*}


\begin{gather*}
\FitMath{\mathsf{M}(i,X)=\frac1{|W_4|}\sum_\omega e^{c\,\mathsf{v}(\omega)}\int\sup_{\mathbf{p}}\big|e^{\Phi_\omega(\xi;\mathbf{p})}\big|\,\mu_\omega(d\xi),}
\end{gather*}

where $\omega$ ranges over the local histories of $X$, each equipped with a hypercubic frame in $W_4$ of the gauge tree, and $\xi$ denotes the stored variables. The variable $\mathbf{p}=(\mathbf{U},\mathbf{J})$ ranges over the complex pairs on $X^\sim$ at which a term evaluates its factor $e^{\Phi_\omega}$. The exponent has the form $\Phi_\omega=\Phi^\circ_\omega(\xi)+\sum_f\Phi_{\omega,f}(\xi;\mathbf{p}|_{R_f})$ with $R_f\subset X^\sim$, and $\mathsf{v}(\omega)$ is a space-time volume. The supremum is taken inside the integral and the sum. This is the form of the mass for a healed component or an island. In the case of a drop-out, the Wilson weight of the carrying term over $X$ appears inside the integral, and the supremum extends over the evaluations at $\xi$. These are the real configurations that are compatible with the stored variables and regular layer by layer, and at which every test listed by $\omega$ takes its recorded value, with the exception of the test whose failure made $X$ a drop-out. They form a set fibred over $\xi$ (2c). The induction carries the mass bound at the coupling of level $i$ of the action. In this bound the $c$ of $2-c$ is small, whatever the $c$ in the weight $e^{c\,\mathsf{v}(\omega)}$. The mass evaluated at a later level does not exceed the mass at level $i$.

(5) The integral operation $\mathbf{T}'_k(X)$ on a component on which $\mathbf{R}$ acts contains, except in the case of a drop-out, a Wilson factor of its own, namely that of the core. When the component is healed or remains as an island, the total mass that is bounded is that of this operation, not that of $\mathbf{T}_k(Z)$:

\begin{gather*}
\FitMath{\mathbf{T}'_k(X)\,1\le\sup\,\exp\Big\{\sum_{j\le k}CM^4\check{R}_j^5\,d'_j(Z_j)\Big\}\prod_{j\le k}\prod_i\exp\Big(-c\,p_0(g_j)^2\big(d'_j(Z^{(i)}_j)+1\big)-2p_0(g_j)\Big),}
\end{gather*}

where the supremum extends over every admissible choice of the domains $Z_j$, $Z^{(i)}_j$ (Ba{\l}aban, Comm. Math. Phys. 122 (1989) 355, (1.79)), and the sums over the admissible sequences are counted among the costs. The inner product is taken over the regions $Z^{(i)}_j$ created at step $j$. For a component $Z$ of $Z_j$ which has been neither healed nor marked, the factor is written $\exp(-\kappa_j(Z)-2p_0(g_{j(Z)}))$; here $j(Z)$ is the level at which the first creation in $Z$ occurs. The \emph{credit} satisfies $\kappa_j(Z)\ge\sum_{n=j+1}^{j+\mathsf{k}}CM^4\check{R}_n^5\,d'_n(S^{n-j}Z)$, where the sum runs over the $\mathsf{k}$ further steps at which nothing is created. Here $S$ denotes the hull in the next partition with a margin added. The number of $M\check{R}_n$-cubes that meet the region takes the place of the linear size $d'_n$, because a single cube has $d'=0$ and still carries a cost. The credit covers the volume costs of every step until $\mathbf{R}$ acts on $Z$. Once no further step is owed, $\mathbf{T}'_k(X)1\le e^{-(2-c)\,p_0(g_k)}$. Moreover, in the weighted sum over the labels, every component $C$ of $Z_k$ keeps a factor $\exp[-c\,p_0(g_k)-\kappa_1|C|_k]$, where $|C|_k$ is its number of cells.

\textbf{Proposition 1a.2.8 (the small-field step as a map on actions).} Suppose that the tests and $\mathfrak{G}_k$ are formed at $\mathsf{t}$, where $\mathsf{t}$ is large, and let $x\in\mathbb{C}$ with $|x-\mathsf{t}|\le c\,\mathsf{t}$. Suppose that the analyticity radii $(\alpha_0,\alpha_1)$ of the levels $j\le k+1$ are small, with $CM\alpha_0<\alpha_1$ and $\alpha_1^2\le C\alpha_0$ at each level. Suppose also that they are comparable: the quotient of the radii of level $k+1$ by those of any level $j\le k$ lies between $(1+\beta_{\mathrm{s}})^{-1}$ and $(1+\beta_{\mathrm{s}})(k+2-j)^{1/2}$. These are the admissible systems of (1d). Let $S$ be local, analytic and \emph{renormalised} in the same way as $\mathbf{E}_k+\mathbf{R}_k$:

\begin{gather*}
\FitMath{S=\sum_{j\le k}\Big[\sum_z\sum_{X\ni z}\big(G^{(j)}(X,\cdot,z)-G^{(j)}(X,\mathbf{1},z)\big)-e_jA\Big]+\sum_{j\in\mathcal{J}}\sum_X\big(H^{(j)}(X,\cdot)-H^{(j)}(X,\mathbf{1})\big),}
\end{gather*}


\begin{gather*}
\FitMath{N(S)=\sum_{j\le k}\omega_{k+1-j}\|G^{(j)}\|_{j,\kappa}+\sum_{j\in\mathcal{J}}\mathsf{q}_j^2\|H^{(j)}\|_{j,\kappa},} \\[3pt]
\FitMath{\omega_n\le Cn^CL^{-\alpha_{\mathrm{H}}n},}
\end{gather*}

where the pointed totals of $G^{(j)}$ are covariant under all isometries of $\mathbb{T}^{(j)}$, and $e_j$ is the coefficient (1e) of $G^{(j)}$. The totals of $H^{(j)}$ are invariant under the isometries of $\mathbb{T}^{(k+1)}$. Finally, $\mathsf{q}_j$ is the ratio of the analyticity radii of level $k+1$ to those of level $j$. Geometric weights occur only in the first sum. Both symmetries are part of the definition of ``renormalised''. The covariance of the pointed totals is what makes the bracket left after the subtraction of $e_jA$ irrelevant (1e). Consider $W$ in a class of real block fields, namely those at which every background formed in the step exists and yields pairs in the domains of the terms. Then, on the enlarged radii and uniformly in $k$, 

\begin{gather*}
\FitMath{\mathfrak{E}_{k+1}[\mathsf{t};x;S]=\log Z^{(k)}+\sum_{X\in\mathcal{D}_{k+1}}\mathsf{F}(X,U_{k+1})+\mathrm{const}}
\end{gather*}

, where

\begin{gather*}
\FitMath{|\mathsf{F}(X,\cdot)|\le C\big(|1-x/\mathsf{t}|+\delta(\mathsf{t})+\mathsf{s}\,N(S)+e^{-c\,\mathsf{t}\delta(\mathsf{t})^2}\big)\,e^{-(1+4\beta_{\mathrm{s}})\kappa d_{k+1}(X)},}
\end{gather*}

with holomorphic dependence on $x$ and on the parameters of $S$. The map $S\mapsto\mathsf{F}$ satisfies a Lipschitz bound with constant $C\mathsf{s}$; here $\mathsf{s}\ll1$ denotes the quotient of the fluctuation threshold by the analyticity radius. Everything above holds under the condition $\mathsf{s}\,(1+N(S))\le c$. Along the run, $G^{(j)}=E^{(j)}$ and $e_j=b_j$, while the $H^{(j)}$ are the aged $R^{(j)}$ of 1a.2.9.

\textbf{Proposition 1a.2.9 (anatomy of step $k+1$).} The insertions are made in the following order. First, the plaquettes of the local backgrounds of $V_{k+1}$ are tested. Next, on the blocks $P''$ that lie away from failures and from $Z_k$, where the $V(y,x)$ are defined, the factor $\mathfrak{G}_k(\cdot;P'')$ is inserted. Then the approximate fluctuation $V_k(V^{\mathrm{bg}}_{k,\square})^{-1}$ is tested, and this determines $\Omega_{k+1}$. On that region $V^{\mathrm{bg}}_k$ is defined and $V_k=e^{i\mathcal{A}_k}V^{\mathrm{bg}}_k$ is set. Finally, $\mathcal{A}_k$ is tested, and this determines $\Lambda_{k+1}$. The insertions and tests just described, together with the regions that they determine, are those of Ba{\l}aban's form, whereas the sharper tests on collar and rim, together with their outcomes, are given in (2a). As a formula,

\begin{gather*}
\FitMath{(\mathbf{T}\rho_k)(V_{k+1})=\sum_{P,Q,R}\int dV_k\,\delta(\mathcal{M}(V_k)V_{k+1}^{-1})\;\chi_{k+1}\chi^c_{k+1}(P)\;\mathfrak{G}_k(V_k)\;\chi'_k\chi'^c_k(Q)\;\chi^{(k)}\chi^{(k)c}(R)\;\rho_k(V_k).}
\end{gather*}

For the node-free small-field step, apply 1a.2.5 and scale $\mathcal{A}_k=g_kB$. The Taylor expansion about the critical point then has the following form. On a label with seams, this is the interior expansion only: the additional node and seam contributions, including their constant and linear terms, remain in the boundary exponent and require its separate estimates.

\begin{gather*}
\FitMath{x_k\big[A(U_k(V_k))+\mathrm{GF}_k(V_k)\big]=x_kA(U_{k+1})+\tfrac12\langle B,\Delta^{(k)}B\rangle+\sum_bO(g_k|B|^3).}
\end{gather*}

Completing the square in $B|_{\Lambda_{k+1}}$ produces $C^{(k)}(\Lambda_{k+1})$, together with the conditional Gaussian that is kept in $\mathbf{T}^{(k)}(Z_{k+1})$:

\begin{gather*}
\FitMath{(\mathbf{T}\rho_k)(V_{k+1})=\sum_{\mathcal{S}'}\chi_{k+1}(\Omega_{k+1})\,\mathbf{T}^{(k)}(Z_{k+1})\mathbf{T}'_k(\mathcal{S})\Big[e^{A_{k,\mathcal{S}}(U_{k+1})+E^{(k)}_0(\Omega_{k+1})}\,\mathsf{I}\Big],} \\[3pt]
\FitMath{\mathsf{I}=Z^{(k)}(\Lambda_{k+1})\!\int d\mu_{C^{(k)}(\Lambda_{k+1})}(B)\,\chi^{(k)}(\Lambda_{k+1})\,e^{\mathbb{P}^{(k)}+\{A_{k,\mathcal{S}}\}'+\mathsf{W}^{(k)}},}
\end{gather*}

where $E^{(k)}_0(\Omega)$ consists of the first two terms of $\mathsf{E}^{(k)}(\Omega,\Lambda)$. The measure $\mu_{C^{(k)}(\Lambda_{k+1})}$ is the Gaussian measure conditioned on the exterior fluctuation field, and its mean is linear in that field. The prime indicates that the quadratic part of $-x_kA$ has been absorbed into the measure. The term $\mathsf{W}^{(k)}$ is left over from a change of variables on the strip, that is, on the part of $S_{k+1}$ lying along $\partial\Lambda_{k+1}$. Before any expansion, the logarithm is divided into three parts by interpolation:

\begin{gather*}
\FitMath{\log\mathsf{I}=\int_0^1\!dt_1\,\big\langle\{\mathbf{B}_k-A(x_k(\cdot)-x_k,\cdot)\}+\mathsf{W}^{(k)}\big\rangle_{t_1,1}+\int_0^1\!dt_2\,\big\langle\{\mathbf{R}''_k\}\big\rangle_{0,t_2}+\log\mathsf{I}\big|_{t_1=t_2=0},}
\end{gather*}

where $\mathbf{R}''_k$ contains the young $R^{(j)}$, namely those whose cubes are larger than a block. The logarithm is expanded (1d) and regrouped (1a.2.6, with $b_{k+1}$, $\varphi_{k+1}$ and $\mathsf{E}^{(k)}(\Omega_{k+1},\Lambda_{k+1})$). This gives the form of level $k+1$:

\begin{gather*}
\FitMath{(\mathbf{T}\rho_k)(V_{k+1})=\sum_{\mathcal{S}'}\chi_{k+1}\,\mathbf{T}'_{k+1}(\mathcal{S}')\exp\Big[A_{k,\mathcal{S}}(U_{k+1})+\!\!\sum_{X\subset\Lambda_{k+1}}\!\!\big[E^{(k+1)}+R^{(k+1)}\big](X,U_{k+1})} \\[3pt]
\FitMath{+\!\!\sum_{X\cap Z^\sim_{k+1}\ne\emptyset}\!\!B^{(k+1)}(X;U_{k+1},\underline{\mathfrak{f}})+E^{(k)}_0(\Omega_{k+1})+\nu_{\mathfrak{g}}|\Lambda_{k+1}^{(k)*}|\Big]}
\end{gather*}


\begin{gather*}
\FitMath{=\sum_{\mathcal{S}'}\chi_{k+1}\,\mathbf{T}'_{k+1}(\mathcal{S}')\exp A_{k+1,\mathcal{S}'}(U_{k+1});}
\end{gather*}

after this, $\mathbf{R}$ acts. The terms with $X\subset\Lambda_{k+1}$ coincide with the terms of the pure small-field integral (among the $R^{(k+1)}$, here and in $\mathsf{E}^{(k)}$, a label keeps those whose $X$ lies one cell inside $\Lambda_{k+1}$, and the others are counted among the $B^{(k+1)}$); this coincidence is the locality stated in 1a.2.7.

\textbf{Proposition 1a.2.10 (a model case: the free field; the renormalisation group as nested orthogonal projections).} Let the \emph{free curvature} $F$ be the standard Gaussian vector in the Euclidean space $\Omega^2_{\mathrm{ex}}=\operatorname{ran}d$ of exact $2$-cochains on the bare lattice. Here $d$ is the coboundary taking real $1$-cochains to real $2$-cochains on the bare lattice, and $d_{(j)}$ is the coboundary of $\mathbb{T}^{(j)}$. This vector stands for one colour of the gauge field, taken in the quadratic approximation and up to closed forms. The plain linear average and the straight linear average induce the same average of curvatures, called the \emph{plaquette average}:

\begin{gather*}
\FitMath{d_{(j)}(\mathsf{Q}^{\mathrm{pl}})^j=d_{(j)}(\mathsf{Q}^{\mathrm{s}})^j=\mathsf{A}^{(j)}d,} \\[3pt]
\FitMath{(\mathsf{A}^{(j)}F)(P)=L^{-4j}\sum_{z\in B^j(y)}\,\sum_{p\in\mathrm{Sq}^{(j)}(z;P)}F(p),}
\end{gather*}

where $\mathrm{Sq}^{(j)}(z;P)$ is the set of $L^{2j}$ bare plaquettes which together cover the plaquette $P$ of $\mathbb{T}^{(j)}$ with corner $y$, once it is translated to $z$. Put $\mathsf{W}_j:=\Omega^2_{\mathrm{ex}}\ominus(\Omega^2_{\mathrm{ex}}\cap\ker\mathsf{A}^{(j)})$, and let $P_{\mathsf{W}}$ denote the orthogonal projections. Then the spaces are nested, $\mathsf{W}_0\supset\mathsf{W}_1\supset\dots\supset\mathsf{W}_K$. Given the block field of level $k$, the conditional expectation of $F$ is $P_{\mathsf{W}_k}F$. This is the curvature of the linear background, namely the exact cochain of smallest norm that has the prescribed plaquette averages. The covariances of the background and of the fluctuation are two complementary orthogonal projections, $\mathbb{E}\big[P_{\mathsf{W}_k}F(p)\,P_{\mathsf{W}_k}F(q)\big]=P_{\mathsf{W}_k}(p,q)$ and $P_{\mathrm{ex}}=P_{\mathsf{W}_k}+P_{\mathsf{W}_k^\perp}$, where $P_{\mathrm{ex}}$ is the projection onto $\Omega^2_{\mathrm{ex}}$. Finally,

\begin{gather*}
\FitMath{F=P_{\mathsf{W}_K}F+\sum_{j=1}^K\mathcal{C}_{(j)}F,} \\[3pt]
\FitMath{\mathcal{C}_{(j)}=P_{\mathsf{W}_{j-1}}-P_{\mathsf{W}_j}=\mathcal{H}_{j-1}C^{\mathrm{tr}}_j\mathcal{H}_{j-1}^*,} \\[3pt]
\FitMath{|\mathcal{C}_{(j)}(p,q)|\le C\eta_{j-1}^4e^{-\delta|x(p)-x(q)|/a_{K-j+1}},}
\end{gather*}

where the summands are independent and the covariance of each summand is its own projection. Here, for a block field $B$ of level $j-1$, $\mathcal{H}_{j-1}B$ denotes the curvature of its linear background, namely the exact $2$-cochain on the bare lattice of smallest norm satisfying $\mathsf{A}^{(j-1)}\mathcal{H}_{j-1}B=d_{(j-1)}B$; and $C^{\mathrm{tr}}_j$ is the fluctuation covariance, in a tree gauge, of step $j$. Thus step $j$ integrates out precisely $\mathcal{C}_{(j)}F$, and the amplitude and range of each slice are those of its own scale. In this model the semigroup law is the nesting of the spaces, and the delta constraint is conditioning. The nesting of backgrounds takes the form $P_{\mathsf{W}_j}P_{\mathsf{W}_k}=P_{\mathsf{W}_k}$ ($j\le k$). Three features of the full theory are absent from the model: the fibre is flat, no coupling runs, and no field is large. The current of a linear background is not zero: $d^*P_{\mathsf{W}_k}F=((\mathsf{Q}^{\mathrm{s}})^k)^*d_{(k)}^*\phi$, which is the image of a Lagrange multiplier. What is missing is the term $\langle\mathsf{H}_kC^{(2)}_k,J\rangle$ of 1a.1.10, because $C^{(2)}_k=0$ on a flat fibre. In the non-linear theory the current and the multiplier do vanish, at $U=\mathbf{1}$. For this reason the linearisation at that point of the minimiser of $A$ subject to $\mathcal{M}^kU=V$ is the minimiser of $\frac12\|d\mathsf{a}\|^2$ subject to the linearised constraint, and the second-order part of the averages makes no contribution. Hence the model is the linearisation at $\mathbf{1}$. It serves as the template for (7a), (7b).

\textbf{Lemmas 1a.3.1, 1a.3.2 (axial gauges).} Consider a box of side $D$ on which $|\partial V-\mathbf{1}|\le\epsilon$. Among the representatives $V^{\psi'}$ with $\psi'(y_P)=\mathbf{1}$, the configuration $V^{\mathrm{ax}}=V^\psi$, $\psi(w)=V(\Gamma_{y_P,w})$, is the only one that equals $\mathbf{1}$ on the tree. The non-abelian Stokes theorem expresses $V^{\mathrm{ax}}(b)$ as a product of conjugated plaquette variables, at most $|b_--y_P|_1$ in number:

\begin{gather*}
\FitMath{V^{\mathrm{ax}}(b)=\prod_{i=1}^m\psi(w_i)\,\partial V(p_i)^{\pm1}\psi(w_i)^{-1},} \\[3pt]
\FitMath{|V^{\mathrm{ax}}(b)-\mathbf{1}|\le|b_--y_P|_1\,\epsilon\le4D\epsilon;}
\end{gather*}

Suppose that $D\epsilon\le c$ and that $F=\frac1i\log\partial V$ is H{\"o}lder along combs; it is distinct from the field strength $F_V$ introduced in the Notation, \ensuremath{\S}2, being the logarithm of each plaquette variable. Then $\frac1i\log V^{\mathrm{ax}}_\mu(x)=\sum_{\rho<\mu}(x_\rho-y_\rho)F_{\rho\mu}(y_P)+O(|x-y_P|^2)\times(\text{H{\"o}lder constant}+\epsilon^2)$, which is the expansion used in (1e) and (2b). Now suppose that $U_0,U\in\mathfrak{U}_k(\Omega_\bullet,\alpha_0)$, that $U$ is tower axial, and that the averages at the tops differ by an amount $<\alpha_1$. Then $|\mathcal{M}^jU(\mathcal{M}^jU_0)^{-1}-\mathbf{1}|<C\alpha_0+\alpha_1$ at every level, uniformly in $k$. This bound holds bond by bond and is coarser than what scaling requires. On a single level (compare the rigidity Lemma 1a.1.9), $\mathcal{A}_k$ is defined, and is $O(\varepsilon_k)$, only after the gauge has been fixed.

\textbf{Proposition 1a.3.3 (Faddeev--Popov identities with Jacobian one).} Let a compact group $H$ act, and let $\mathsf{k}\ge0$ satisfy $\int_H\mathsf{k}(h\cdot x)\,dh=1$. Then for every finite complex measure $\nu$ and every bounded invariant $F$,

\begin{gather*}
\FitMath{\int\nu(dx)\int_Hdh\,\mathsf{k}(h\cdot x)\,F(h\cdot x)=\int F\,d\nu;}
\end{gather*}

The measure obtained in this way is constructed as follows: one chooses $h$ with law $\mathsf{k}(h\cdot x)dh$ and replaces $x$ by $h\cdot x$. This measure equals $\mathsf{k}\,\bar\nu$, where $\bar\nu$ is the gauge average of $\nu$. When $\nu$ is invariant, it is the product $\mathsf{k}\,\nu$. If $u=\mathbf{1}$ on the coarse sites, then $V^u(y,x)=V(y,x)u(x)^{-1}$. Where $V(y,x)$ is undefined, no factor is inserted at the pair $(y,x)$, and with this convention

\begin{gather*}
\FitMath{\int\prod_xdu(x)\,\mathfrak{G}_k(V^u;P'')=1,} \\[3pt]
\FitMath{(\mathbf{T}\rho_k)(V_{k+1})=\int dV\,\delta(\mathcal{M}(V)V_{k+1}^{-1})\,\mathfrak{G}_k(V;P'')\,\rho_k(V).}
\end{gather*}

In addition, $\mathfrak{G}_k(V^v)=\mathfrak{G}_k(V)$ whenever $v$ is constant on blocks. The same holds for delta functions on spanning trees. Let $T$ be a tree whose root lies outside the sites of the group, and let $V_0$ be any configuration on $T$. Then a unique $u$ satisfies $(u\cdot V)|_T=V_0$, and the measure $dV$ splits into factors along the tree:

\begin{gather*}
\FitMath{\int f\,dV=\int du\int\prod_{b\notin T}dS(b)\,f\big(u\cdot(S\sqcup V_0)\big),} \\[3pt]
\FitMath{\text{so}} \\[3pt]
\FitMath{\int dV\,\delta_T(VV_0^{-1})\,\rho F=\int dV\,\rho F}
\end{gather*}

for invariant $\rho,F$. Finally, the delta function can be removed: the gauge average of $\Psi\,\delta_T(VV_0^{-1})dV$ is $(\Psi\circ\pi_T)\,dV$, where $\pi_T$ is the projection onto the slice along the orbits.

\textbf{Lemma 1a.3.4, Proposition 1a.3.5 [to be proved; Package A] (the Landau gauge).} Let $U$ be tower axial and let $W=U^{u^{-1}}$. Then $\mathsf{n}_W=u^{-1}\bar{R}_{U_0}u$, so a $u$ with trivial averages leaves the straightened averages exactly unchanged. In linearised form this reads $Q^{\sharp}_j(U_0)D\lambda=D_{\mathcal{M}^jU_0}Q'_j(U_0)\lambda$. Hence $Q^{\sharp}A$ does not change under $A\to A-D\lambda$, $\lambda\in N(Q')$. Up to terms that contain $J(U_0)$, gradients are null directions of $\Delta^{\mathrm{W}}$. Indeed $D^*J(U_0)=0$, and, writing $\lambda_+$ for $\lambda(b_+)$ transported to $b_-$,

\begin{gather*}
\FitMath{\langle A-D\lambda,\Delta^{\mathrm{W}}(A-D\lambda)\rangle=\langle A,\Delta^{\mathrm{W}}A\rangle-\big\langle i[\lambda,A]-i[A,\lambda_+]-i[\lambda,D\lambda],\,J(U_0)\big\rangle.}
\end{gather*}

The projection is given by

\begin{gather*}
\FitMath{R=1-G'Q'^*(Q'G'^2Q'^*)^{-1}Q'G',} \\[3pt]
\FitMath{\operatorname{ran}R=\{f:Q'G'f=0\},} \\[3pt]
\FitMath{\inf_{\lambda\in N(Q')}\|D^*A-\Delta\lambda\|^2=\|(1-R)D^*A\|^2,}
\end{gather*}

and $R\Delta\lambda=\Delta\lambda=G'^{-1}\lambda$ for $\lambda\in N(Q')$. The map $\lambda\mapsto(R\Delta\lambda,Q'\lambda)$ is a bijection onto $\operatorname{ran}R\times\{\text{functions on tops}\}$. Its inverse is $(f,\nu)\mapsto G'f+H'\nu$, where $\nu$ is a function on the tops and $H'$ is a right inverse of $Q'$, constructed from $G'$, with $R\Delta H'=0$. Let $U\in\mathfrak{U}_k(\Omega_\bullet,\alpha_0)$ and let $M_1$ be large. Assume that $M\alpha_0$ (and therefore $M_1\alpha_0$) lies below a number that depends only on $L$, and that $\mathbb{T}^{(k)}$ has at least $CM_1$ sites in each coordinate direction. Then the form

\begin{gather*}
\FitMath{\langle A,\Delta^{\mathrm{W}}_UA\rangle+\|RD^*_UA\|^2+\langle Q^{\sharp}A,aQ^{\sharp}A\rangle}
\end{gather*}

is positive definite. Its inverse $\mathsf{G}_k(U)$ maps $|\cdot|_{(-3)}$ into $\|\cdot\|_1$, with a norm that depends only on $L$. On the set $\{RD^*A=0,\ \xi Q^{\sharp}A=B\}$ the form differs from the Hessian by a constant. Consequently $\mathsf{H}_kB=\mathsf{G}Q^{\sharp*}(Q^{\sharp}\mathsf{G}Q^{\sharp*})^{-1}\xi^{-1}B$ minimises the Hessian itself under the constraint, and $RD^*\mathsf{G}Q^{\sharp*}=0$. Here $\mathsf{G}$ stands for Ba{\l}aban's $G$, and not for the inverse of the displayed form, which is his $G_0$. The formula for $\mathsf{H}_kB$ and the identity just stated hold exactly only when the Hessian in the form is replaced by its invariant extension $\Delta_\pi$. This extension is defined by $\langle A,\Delta_\pi A\rangle=\langle A-D\lambda,\Delta^{\mathrm{W}}(A-D\lambda)\rangle$, where $\lambda\in N(Q')$ is the one with $RD^*(A-D\lambda)=0$. The replacement is a perturbation of size $CM\alpha_0$, which is small together with the current because $\alpha_0$ is chosen after $M$. The bounds persist under this perturbation, and it vanishes where $J=0$. In the symbols of the Notation, \ensuremath{\S}5, this operator equals $\mathsf{H}\,\xi^{-1}$. There $\mathsf{H}$ denotes the linear minimiser constructed from $G$, and it satisfies $Q^{\sharp}\mathsf{H}=1$. Where $J=0$ and $\xi=1$, the operator coincides with the $\mathsf{H}_k$ of the Notation. It carries the normalisation for the constraint $\xi Q^{\sharp}A=B$. With his $Q^{\mathrm{B}}$, Ba{\l}aban uses both normalisations: the minimiser under $QA=B$ and the minimiser under $L^j\eta\,Q_jA=B$. On layer $j$ the two differ by the factor $\xi^{-1}=(L^j\eta_k)^{-1}$. The normalisation with $\xi$ is the one taken in his paper on the variational problem.

\textbf{Theorem 1a.3.6 (from the axial to the Landau gauge).} Work in the scheme $\mathsf{c}$, in which Ba{\l}aban's average of 1985 replaces $\mathcal{M}$ in every object, and let $c_1,B_1$ depend only on $L$. Assume the hypotheses of Lemma 1a.3.2 [to be proved; Package A], with the full gauge $\operatorname{Ax}(\Omega_\bullet,U_0)$ and not $\operatorname{Ax}^\circ$, and assume $\alpha_0+\alpha_1\le c_1$. Then $U$ has exactly one Landau pair that satisfies

\begin{gather*}
\FitMath{|A|_{(-1)},\,|\nabla_{U_0}A|_{(-2)},\,|D^*DA|_{(-3)},\,|\Delta_{U_0}A|_{(-3)}\le B_1(\alpha_0+\alpha_1),} \\[3pt]
\FitMath{\sup|u-\mathbf{1}|\le CB_1(\alpha_0+\alpha_1),}
\end{gather*}

that satisfies H{\"o}lder bounds on $\nabla A$ (these bounds give rise to the $\alpha_{\mathrm{H}}$ below), and that has the straightened averages of $U$,

\begin{gather*}
\FitMath{\mathcal{Q}^{\sharp}_j(U_0;A)=B\,\text{ on the bonds at }\mathbb{Y}_j,} \\[3pt]
\FitMath{B=\tfrac1i\log\big[\mathcal{M}^jU(\mathcal{M}^jU_0)^{-1}\big],} \\[3pt]
\FitMath{|B|\le C\alpha_1}
\end{gather*}

(with an explicit correction entering on bonds that join two layers). For $\mathcal{M}$ itself, under the same hypotheses with $\alpha_0+\alpha_1\le c$, there exists a pair $(A,u)$ in the slice of (1b) that has these averages and satisfies $\|A\|_1,\ |D^*DA|_{(-3)}\le C(\alpha_0+\alpha_1)$. That slice has the same form as the Landau gauge but a different projection, whose site averages run along combs from the block corners.

\textbf{Proposition 1a.3.7 (a priori bounds on the slice).} Assume that $U_0,e^{i\eta_kA}U_0\in\mathfrak{U}_k(\Omega_\bullet,\alpha_0)$, that $RD^*A=0$, and that the averages are as above. Assume further that the only bound given is $|A|_{(-1)}<\alpha_2$, with $\alpha_0,\alpha_1,\alpha_2\le c$ ($c$ depending only on $L$) and $C\alpha_2^2\le\alpha_0+\alpha_1$. Then, in the scheme $\mathsf{c}$, all the bounds of the theorem on $A$ hold, with $B_1=CL\times$ the norms of $\mathsf{G}_k(U_0)$, $\mathsf{H}_k(U_0)$. For $\mathcal{M}$, the bounds on $\|A\|_1$ and $|D^*DA|_{(-3)}$ hold, in the slice of (1b).

\textbf{Proposition 1a.3.8 (regular means small potential, locally).} Let $U\in\mathfrak{U}_k(\Omega_\bullet,\alpha_0)$, take a cube of side $M\xi_j$ in $\Omega_j$, and assume $CM\alpha_0\le c_1$. Then for some $u$

\begin{gather*}
\FitMath{U^{u^{-1}}=e^{i\eta_kA},} \\[3pt]
\FitMath{\xi_j|A|,\,\xi_j^2|\nabla A|,\,\xi_j^3|D_{\mathbf{1}}^*D_{\mathbf{1}}A|,\,\xi_j^3|\Delta A|\le CM\alpha_0,}
\end{gather*}

and, in the scheme $\mathsf{c}$, the box axial potential of the averaged field gives the averages of $A$: $\mathcal{Q}^{\sharp}_j(\mathbf{1};A)=\frac1i\log(\mathcal{M}^jU)^{\mathrm{ax}}$. Conversely, let $U_0\in\mathfrak{U}_k(\Omega_\bullet,\alpha_0)$ and $\|A\|_1,|D^*DA|_{(-3)}<\alpha_2$, with $\alpha_0,\alpha_2$ small. Then $e^{i\eta_kA}U_0\in\mathfrak{U}_k(\Omega_\bullet,\alpha_0+3\alpha_2)$. Now take a complex $\mathbf{U}=e^{i\eta_kA'}U$ with $\xi_j|A'|<\alpha_1$ whose own plaquettes are small, $|\partial\mathbf{U}-\mathbf{1}|<\alpha_0\eta_k^2\xi_j^{-2}$ on the cube. (From $\|A'\|_1<\alpha_1$ this follows only with $\alpha_0+C\alpha_1$ in place of $\alpha_0$, 1a.1.6.) Assume also $CM(\alpha_0+\alpha_1)\le c$. Then the averages are $V^v$. Here $V$ is what the averages of $\mathbf{U}$ become in the comb gauge of the cube, and $v=v'v_{\mathrm{r}}$ is a $G^c$-valued function on the nodes. The factor $v_{\mathrm{r}}$ is $G$-valued and does not depend on $A'$, the factor $v'$ is holomorphic in $A'$, and

\begin{gather*}
\FitMath{|V-\mathbf{1}|\le CM\alpha_0,} \\[3pt]
\FitMath{\operatorname{cn}(v)=|v|\,|v^{-1}|\le e^{CM\alpha_1},} \\[3pt]
\FitMath{|v'-\mathbf{1}|\le CM\alpha_1.}
\end{gather*}

The third bound of Ba{\l}aban, $|v-\mathbf{1}|\le CM\alpha_1$, holds for $v'$. For $v$ it holds only at pairs normalised by $U(\Gamma_{y_\square,x})=\mathbf{1}$. Indeed, for $U=\mathbf{1}^g$, $A'=0$ one has $v=g(\cdot)g(y_\square)^{-1}$, which is arbitrary, whereas $\operatorname{cn}(v)=1$.

\textbf{Propositions 1a.3.9, 1a.3.10 (transversality, locality, node factor; the choices disappear).} The slice is transversal. Let $\tilde{A}$ be real with $\|\tilde{A}\|_1\le c$, and let a $G$-valued $\mathsf{y}$ near $\mathbf{1}$ be prescribed on the tops. Then some $w=e^{i\lambda}$ brings $\tilde{A}$ into the slice with $r_j=\mathsf{y}$ on $\mathbb{Y}_j$, and $\|\lambda\|_0\le C(|D^*\tilde{A}|_{(-2)}+\sup|\mathsf{y}-\mathbf{1}|)$. Here $r_j$ is defined as in 1a.1.8 and formed at $\tilde{A}$; at $\tilde{A}=0$ it is $\bar{R}^{(j)}_{U_0}w$. The pair $(A,u)$ depends locally on the data $(J(U),B)$, with tails $Ce^{-\delta d(y,y')}$ and $\delta<\delta_0$. Distances are measured in local scales, $J(U)$ in $|\cdot|_{(-3)}$, $B$ in the sup norm and $A$ in $\|\cdot\|_1$. Along differentiable families of $U$, this locality follows from the relation

\begin{gather*}
\FitMath{\delta A=\mathsf{G}_k(U_0)\big[\operatorname{Ad}(u^{-1})\delta J(U)+\mathcal{L}\delta A\big]+\mathsf{G}_k(U_0)Q^{\sharp*}\xi^{-3}\big[\delta B-\partial_AC(A)\delta A\big],}
\end{gather*}

and, for backgrounds over a real datum $\mathbf{V}$ that is regular in the sense of (1b) (its test loops lie within a small $\varepsilon$ of $\mathbf{1}$, and, where inward bonds are kept, so do the two-scale loops at the seams), the pair depends holomorphically on the datum within a ball of radius $c$, the representative at $e^{iB}\mathbf{V}$ being $(e^{i\eta_kA(B)}U_k(\mathbf{V}))^{h(B)}$. None of these statements depends on $k$. The plain averages of $e^{i\eta_kA}U_0$ are the datum transformed by the \emph{node factor}: $\mathcal{M}^j(e^{i\eta_kA}U_0)(c)=\mathsf{n}(c_-)V_j(c)\mathsf{n}(c_+)^{-1}$, where $\mathsf{n}=u^{-1}|_{\mathrm{nodes}}$. At the tops the node factor is the iterated $\mathsf{n}_W$, namely $v_j[A]$. At a junction node it contains one further factor, the quotient of the datum by the reduced average on one bond ending at that node. Every ingredient of a step is gauge covariant, and the map $(A,u)\to(\operatorname{Ad}(v)A,vuv^{-1})$ takes the set of Landau pairs of $U$ with respect to $U_0$ onto the set of Landau pairs of $U^v$ with respect to $U_0^v$. Consequently $E^{(k+1)}(X,\cdot)$ is invariant under 

\begin{gather*}
\FitMath{(\mathbf{U},\mathbf{J})\to(\mathbf{U}^u,\operatorname{Ad}(u)\mathbf{J})}
\end{gather*}

 for $G$-valued $u$. It is therefore also invariant for the $G^c$-valued $u$ that are connected to these within the domain, and it can be evaluated in any gauge. A function $F$ constructed on a slice is invariant under the subgroup that preserves the slice, and it has a unique invariant extension, $F^\natural(\mathfrak{z})=F(\mathsf{c}[\mathfrak{z}]\cdot\mathfrak{z})$. Here $\mathsf{c}[\mathfrak{z}]$ is the transformation that takes $\mathfrak{z}=(U,A',\mathbf{J})$ to the slice. The extension retains locality, the bounds, and holomorphy in $(A',\mathbf{J})$ at fixed real $U$.

\textbf{Lemma 1a.3.11 (the block axial gauge outside the nodes).} Near a seam, the block of a coarse site contains nodes, and every element of $\mathcal{G}^{\mathrm{res}}$ is pinned at these nodes. In that case the gauge condition is $W(y,x)=\mathbf{1}$ for $y\in P''$ and for those $x\in B(y)\setminus y$ other than nodes. For a pair of nodes, the value $(\mathcal{M}^iU)(y,x)$ is invariant on the orbit, and the datum does not contain it. If $P''$ consists of nodes, the gauge is met once by every orbit of $\mathcal{G}^{\mathrm{res}}$, meaning that the averages $\mathcal{M}^iU$ on the bonds within $B(P'')$ and the nodes depend only on the orbit. Hence $V^{\mathrm{bg}}_k=\mathcal{M}^kU$, with $U$ in the orbit $U_{k+1}(\mathbf{V})$ and in the gauge, is a function of the datum; this is the case $i=k$, with $k+1$ in place of $k$. In this situation

\begin{gather*}
\FitMath{\mathrm{GF}_i(\mathcal{M}^iU;P'')=\sum_{y\in P''}\,\sum_{x\in B(y)\setminus y,\,x\text{ a node}}\big[1-\tfrac12\operatorname{Re}\operatorname{tr}(\mathcal{M}^iU)(y,x)\big],}
\end{gather*}

which is constant on the orbit. Near a seam, this is what is left of the identity $\mathrm{GF}_k(V^{\mathrm{bg}}_k)=0$.

\textbf{Proposition 1a.3.12 (a generalisation of Theorem 1a.3.6, in the scheme $\mathsf{c}$).} Fix $\gamma'>0$. Where the constant $B_1$ of the theorem enters, $\gamma'$ is not large; $\gamma'=1$ is an example. Assume the hypotheses of the theorem, with $c_1$ depending on $\gamma'$ as well; for arbitrary $\gamma'$, the constant $B_1$ must also depend on $\gamma'$. Let $f\in\operatorname{ran}R$ with $|f|_{(-2)}<\gamma'(\alpha_0+\alpha_1)$. Then the contraction in $\|\lambda\|_0$ from the proof of the theorem produces, for $\lambda$ in a ball of fixed radius, a unique $u$ whose averages are trivial and for which $R\,D^*_{U_0}A=f$. Nothing more than this enters the consequence below. Ba{\l}aban's Theorem 8 [to be proved; Package E] also asserts, without proof, all the bounds of the theorem, with the same $B_1$. The second-order and H{\"o}lder bounds require, in addition, a bound on $Df$ in $|\cdot|_{(-3)}$. As a consequence, the set $\{A:\mathcal{Q}^{\sharp}(U_0;A)=B,\ RD^*A=f\text{ for some small }f\}$ consists of pieces of gauge orbits, and each piece meets the slice $f=0$ in exactly one point. On each piece the action takes a single value. Hence a point of the slice that is critical there is critical on the thickened set as well, so that the Lagrange equations no longer involve the gauge condition. Ba{\l}aban makes use of this a single time, to obtain the second-order bounds on the minimiser. In (1c) the corresponding step, which passes from the slice to a multiplier attached to the averages only, is carried out with gauge directions the images of which under $RD^*$ fill all of $\operatorname{ran}R$. The domains $U^c_j$ have real centres that are regular, though not nearly critical, and at these centres the construction demands small local potentials outright.

\textbf{Proposition 1a.3.13 (a seam in dimension one).} Take all definitions on $\mathbb{Z}\supset L\mathbb{Z}$, and perform one step on the weight $1$ on two blocks, the second of which is centred at $y$; the third block, $B(y')$, remains fine. The seam bond $s$ is kept in the datum, $\mathsf{p}$ is the transport from $s$ to $y'$, and $f_{\mathfrak{G}}$ is the one-site factor of $\mathfrak{G}_k$. Let $V^+(c)$ be the coarse variable on $c=\langle y,y'\rangle$. Then the integral of the step can be computed exactly:

\begin{gather*}
\FitMath{\Phi=f_{\mathfrak{G}}\big(V^+(c)\,\mathsf{p}^{-1}V(s)^{-1}\big).}
\end{gather*}

Under the node-wise gauge group of the datum, this integral fails to be invariant. Let a transformation be $g$ at the centre $y$ and $\mathbf{1}$ at all remaining end points of bonds of the datum, the end of $s$ among them. It changes the integral into $f_{\mathfrak{G}}(gV^+(c)\mathsf{p}^{-1}V(s)^{-1})$. The integral is invariant under the action in which the end of $s$ moves with the centre of its block, and with respect to this action the tree insertion is a valid identity. Under the node-wise action, $\Phi$ has gauge average equal to the constant $1$. The procedure of 1a.3.3 therefore produces the bare slice measure, not $\Phi$ times that measure, and in this case the insertion is not an identity. If $s$ is integrated together with its block, then $\int dV(s)\Phi=1$, and nothing remains to vary.

\textbf{Proofs.} \emph{Averaging.} Every covariance property follows from the identity $\log(TXT^{-1})=T(\log X)T^{-1}$. Each expression under a logarithm is conjugated at a single site ($c_-$ for $\log W_x$). The norms are gauge invariant, and so therefore is the choice of rule. For 1a.1.1, in the chart $V=e^YV_1$ the means are exactly the fixed points of the map $Y\mapsto Y+\frac1n\sum_i\log(V_iV_1^{-1}e^{-Y})$. This map is a contraction at rate $O(\operatorname{diam}+r)$, with bounds independent of $n$. Under a conjugation that is not unitary, diameters change by at most $\operatorname{cn}(T)$. The isometry $\mathsf{r}$ permutes the centred blocks and the families $\mathcal{G}$, but it carries a comb to a comb in a different order. This is why Ba{\l}aban's first average, built on cornered blocks and combs, needed to be symmetrised. Lemma 1a.1.2 (lattice homotopy): each interchange of two steps in a path costs one plaquette. The contours sweep out $O(L^2)$ plaquettes, and this determines $\mathrm{Dom}(c)$. In the comb gauge the bound increases with the size of the box, so small curvature yields a small potential only on boxes of bounded size. In the comb gauge one has $\mathcal{M}(e^{i\mathsf{a}})=\exp(i\mathsf{Q}^{\mathrm{pl}}\mathsf{a}+O(|\mathsf{a}|^2))$. Since the contour terms $\sigma$ cancel at the corners of $p'$, the lattice Stokes formula yields

\begin{gather*}
\FitMath{\sum_{c\in\partial p'}\pm(\mathsf{Q}^{\mathrm{pl}}\mathsf{a})(c)=L^{-4}\sum_{x\in B(y)}\sum_{p\subset S_x}\mathsf{a}(\partial p),} \\[3pt]
\FitMath{S_x=p'\text{, with corner }y\text{, translated to }x,\,L^2\text{ plaquettes}.}
\end{gather*}

Consequently $\hat\varpi_n=\varpi(\mathcal{M}^nU)/\alpha_0(L^n\eta_k)^2$ satisfies $\hat\varpi_{n+1}\le\hat\varpi_n(1+C\alpha_0L^{-2(k-n)})$. Measured in its own scale, the curvature of level $n$ is summable, and this outweighs the accumulation of $k$ errors.

A single step (Proposition 1a.1.5) is proved by universality and Cauchy bounds, whereas Ba{\l}aban computes $\mathsf{Q}(W_0)$ in closed form from the Campbell--Hausdorff formula. In the comb gauge everything reduces to finitely many holomorphic functions that depend on $L$ alone, and hence $\mathcal{Q}^1=\mathsf{Q}(W_0)\mathsf{a}+O(|\mathsf{a}|^2)$. At $W_0=\mathbf{1}$ one has $\frac1i\log\mathsf{u}=\sigma+O(|\mathsf{a}|^2)$. Thus the block factors take away precisely the contour terms of $\mathsf{Q}^{\mathrm{pl}}$, and $\mathsf{Q}(\mathbf{1})=\mathsf{Q}^{\mathrm{s}}$. The kernels of $\mathsf{Q}(W_0)$ and $\mathsf{S}(W_0)$ are holomorphic and coincide at $\mathbf{1}$, so they differ by $O(\varpi)$. The norms of $\mathsf{S}$ reduce to counting: there are $L^4\cdot L$ pairs $(x,b)$ for each coarse bond, each fine bond is contained in $L$ segments, and every pair carries the weight $L^{-4}$. For $k$ steps, the factorisation follows from covariance by induction. At a complex $1$-form $A$, covariance under the $G^c$-valued $v_n$ uses up part of the room in the domain, and it is applied with $\operatorname{cn}(v_n)=|v_n|\,|v_n^{-1}|\le2$. Setting $t_n=\sup|\mathsf{a}_n|/L^n\eta_k$ and $\varpi_n=\varpi(W_{0,n})$, one has

\begin{gather*}
\FitMath{t_{n+1}\le t_n\big(1+C\varpi_n+CL^n\eta_kt_n\big),} \\[3pt]
\FitMath{\sum_n\varpi_n\le C\alpha_0,} \\[3pt]
\FitMath{\sum_{n<k}L^n\eta_k\le2L^{-1}.}
\end{gather*}

The recursion closes because $\|\mathsf{S}(W_{0,n})\|_{\mathrm{row}}$ does not exceed $L$, the ratio of the spacings, with no loss. This is the third gain of straightening. Since $\|\mathsf{Q}^{\mathrm{pl}}\|_{\mathrm{row}}\ge L+\frac{3(L^2-1)}{2L}>2L$, a product of one-step bounds for the plain average grows geometrically in $k$. The derivative is bounded through column sums and Schur's test. When the plaquettes are complex, the real comb gauge does not work: at every level its error is as large as the bound itself, which gives only $1+Ck\alpha_1$. The $G^c$-valued comb gauge of $\mathcal{M}^n\mathbf{U}$, whose condition number is $1+C\alpha_1L^n\eta_k$, gives $\varpi^c_{n+1}\le L^2\varpi^c_n(1+C\alpha_1L^n\eta_k)(1+CL^2\varpi^c_n)$. The price of the non-unitary conjugation is therefore one factor per level, and these factors are summable. The determining sets are obtained by folding with reflections. The double block of order $j$ beneath a bond of $\Gamma_j$ lies in $\Omega_{j-1}$. There the fields are, in the worst case, regular only in the next finer scale, and this costs one factor, a single time.

Intertwining: at a pure gauge $e^{i\mathsf{a}^w}W_0=W_0^w$ one has $\mathsf{u}(W_0,\mathsf{a}^w)=w\,(\bar{R}_{W_0}w)^{-1}$. It is for the sake of this identity that $\bar{R}$ is defined as it is. 1a.1.8 rests on a coupled recursion for the size of $r_n$ and for its covariant difference, which the same geometric sums close. For products one has $|\mathsf{z}_{j+1}|\le|\mathsf{z}_j|+C(|\mathsf{z}_j|+\alpha L^{j+1}\eta_k)^2$, because the defect added at each step is quadratic in the variation within a block. This is the origin of the second condition in the definition of a class. 1a.1.9: up to an error $C\varpi$, a bond inside a block is the quotient of two contour transports, and the average then determines the crossing bonds. The comparison is made factor by factor. 1a.1.10: the constraint is solved to second order by means of the right inverse $\mathsf{H}_k$. At a critical point under the constraint, $J(U)=Q^{\sharp}_k(U)^*\omega$ is in general not zero. It is the image of a Lagrange multiplier $\omega$, whose size is that of the curvature of $V$. The extra term does not occur in abelian theories, and it gives rise to $G_1$, $\mathsf{H}_1$ in (1c). 1a.1.11 follows from the first identity together with $\mathcal{M}(\mathsf{r}U)=\mathsf{r}\mathcal{M}(U)$ and $(\mathsf{r}U)(\mathsf{r}y,\mathsf{r}x)=U(y,x)$. 1a.1.12 is purely algebraic. Writing $\hat{u}_n=u\,v_n$, the axial condition at level $n$ takes the form 

\begin{gather*}
\FitMath{\hat{u}_n(x)=\operatorname{Ad}(W_{0,n}(z,x))^{-1}[\hat{u}_{n+1}(z)\mathsf{u}_n(z)^{-1}\mathsf{w}^{\mathrm{rel}}_n(z,x)]}
\end{gather*}

. This fixes $u$ level by level, starting at the top and going down, and it yields $\bar{R}_{W_{0,n}}\hat{u}_n=\hat{u}_{n+1}$. 1a.1.13: two chains of transports with common end points differ by the curvature that they enclose. At level $n$ this is $O(\alpha_0L^{-2(k-n)})$, so the differences form a geometric series. The transport $t$ is unavoidable, because in each block the two schemes see a gauge transformation at two different sites. 1a.1.14 follows from linearity and the Stokes formula, which give $\Theta=L^{-4}\sum_{x\in B(z)}L[\mathsf{f}^\Sigma(x_1+L)-\mathsf{f}^\Sigma(x_1)]$, together with a count.

\emph{The transformation.} 1a.2.1 combines the definition of an image measure with a property of $\mathcal{M}$. The name ``Group'' is inaccurate. All that holds is the semigroup law $\mathbf{T}^n\nu=(\mathcal{M}^n)_*\nu$, which relates different spaces. In physical units the symmetry group becomes smaller as $k$ grows; in the units of the level it is the symmetry group of the unit lattice. Finite range does not survive, since $\mathbf{T}\rho_0$ couples all bonds. Reflection positivity is required of the bare measure only (4d), (6b), (10a). 1a.2.2: the variances are additive (taking a block mean divides by $L^d$, the rescaling multiplies by $L^{d-2}$, and the new kernel contributes $\mathfrak{a}^{-1}$). Nested backgrounds: backgrounds exist only over regular data, so minimising in two stages is not enough. The argument is instead as follows. A point that is critical on the larger fibre is critical on the smaller one, regularity improves in finer units, and uniqueness (1b) completes the proof. Because the argument involves critical points alone, it extends to complex data. The minimiser $\mathsf{H}B$ is the one given by the method of Lagrange multipliers. The delta kernel secures four things:
\begin{itemize}
\item closure under composition. Since $\mathcal{M}$ is non-linear, the composite of two soft kernels on $G$ is not a function of $W$ and $\mathcal{M}^2U$ alone.
\item a single large term. The critical point does not depend on the coupling, the leading term is exactly $x_kA(U_{k+1})$, and terms produced at $U_j$ can be evaluated at $U_k$ at every later level.
\item exact averages. Since $V_k=\mathcal{M}^kU$, a test on the block fields is a test on the bare field.
\item no analytic loss. The mass $aQ^{\sharp*}Q^{\sharp}$ remains in $\mathsf{G}_k$ as a device, with $a>0$ arbitrary.
\end{itemize}

The cost is that the fibre is curved. For the gauge fixing the opposite choice is made: it enters through a soft weight.

1a.2.3 (stochastic kernels, balance of masses). The mass, summed over the outcomes, is preserved by operations of three kinds. \emph{Tests:} at any given state exactly one outcome takes the value $1$. A test is counted as failed whenever the object it examines (a mean, a logarithm, a minimal orbit, a minimiser) is undefined, so every test is defined on all configurations. The actual test maps, including the selections of representatives and the failure sets of partially defined objects, must be Borel. Definability can establish this when it is checked for those maps and their selections; it is not supplied by the convention that an undefined object fails its test. \emph{Draws:} the kernel $\kappa_{\mathfrak{G}}(V,\cdot)=\int du\,\mathfrak{G}_k(V^u)\delta_{V^u}$. \emph{Maps:} adjoining $\mathcal{M}(V_k)$, and forgetting a variable. No smallness is involved in any of this. The operation $\mathbf{R}$ draws the variables of a component afresh, and after a Haar draw a plaquette trace in that component has integral $0$. Because of the partitions, single terms have only restricted symmetry. Only the totals at a base point are covariant under the full group $W_4$, which (1e) requires. Base points are required because the sum $\sum_{X\subset\Lambda_j}$ is not invariant and the cancellation against $b_jA(\varphi_j,\cdot)$ must be made at each point separately. When $z$ lies one layer of $M\check{R}_j$-cubes inside $\Lambda_j$, the restricted pointed total differs from the unrestricted one by $O(e^{-\frac12\kappa\check{R}_j})$.

1a.2.4: a compact group admits no dilation. Hence $x_k$ \emph{is} the normalisation of the kinetic term. Gauge invariance links the quadratic, cubic and quartic parts of the marginal term, so the renormalisations of the field strength and of the coupling are given by a single number. The substitution $\mathcal{A}_k=g_k\mathcal{A}^{\mathrm{s}}_k$ takes the place of the rescaling. Consequently there is no scaling transformation, no fixed point and no eigen-direction. The word ``irrelevant'' can then mean only a bound $O(\eta_{k-j}^{4+\alpha_{\mathrm{H}}})$ at regular backgrounds, where $\alpha_{\mathrm{H}}<1$ is a H{\"o}lder exponent (1e). Level $k$ is a problem on a unit lattice that carries a fine lattice of spacing $\eta_k\to0$. Every estimate in (1b), (1c) must therefore hold uniformly in $\eta_k$. 1a.2.6: the density itself does not change. Its exponent is always $-\beta A(U_k)$ together with the terms in their unrenormalised form. The point is which grouping of the terms admits uniform bounds:
\begin{itemize}
\item without grouping, the bound $\|E^{(j)}\|\le E_0$ yields no more than $CE_0\eta_k^{-4}|\mathbb{T}^{(k)}|$;
\item the number $b=b_{k+1}$ by itself leaves a geometric sum over $j$;
\item if the marginal terms are left among the corrections, they add up to $c_0kA$ to leading order;
\item once absorbed into $x_k$, they fix the next thresholds and $g_{k+1}$, the width of the following Gaussian.
\end{itemize}

The profiles $\varphi_j$ keep track of the place at which step $j$ was performed. In the constant, the contribution $-\frac12(N^2-1)\log x_k$ for each free bond is the Jacobian of $\mathcal{A}_k=g_kB$. Laplace's method gives $z_{\mathfrak{G}}\approx Cx_k^{-(N^2-1)/2}$. With four bonds per site, less one gauge direction, this gives $\mathsf{E}^{(j)}(T_\ell,T_\ell)=-(\frac32(N^2-1)\log x_j+O(1))(|\mathbb{T}^{(j)}|-|\mathbb{T}^{(j+1)}|)$. The term $\log x_k$ in $E_+$ arises from the constants $\mathsf{E}^{(j)}(T_\ell,T_\ell)$, $k\le j<K$, which belong to the steps $k+1,\dots,K$ and remain owed in $\bar\rho_k=e^{-E}\rho_k$ (when $k=K$ no constant remains owed, and $E_+=O(1)$). Over a large-field region, the constants differ by a volume factor $e^{C\log x_j\#_j(Z_{j+1})}$, and the small factors of (2c) must outweigh it.

1a.2.7 (the division of the work among the parts: 1a.2). Each part supplies the following:
\begin{itemize}
\item (1b): the backgrounds;
\item (1c): the propagators;
\item (1d): the new $E^{(k+1)}$, by a polymer expansion under a criterion of Koteck{\'y}--Preiss type;
\item (1e): the coupling and the bounds on the sums;
\item (2a): tests, regions and histories;
\item (2b): $\mathbf{R}$;
\item (2c): the large-field factor and stability;
\item (2d): the terms that $\mathbf{R}$ leaves behind;
\item (2e): the order of the constants and the inductive form;
\item (3a): the flow;
\item (3b): independence from the label, the torus and the cutoff.
\end{itemize}

The last item rests on locality. A term localised in $X$ is assembled from walks, potentials and tests within $X$. It therefore does not depend on the label, and, when $X$ does not wrap around the torus, it depends on neither the torus nor the cutoff. The small-field half requires $g_j\le\gamma$. It also requires $-C\le x_i-x_j\le C(j-i)$ for $i\le j\le k$, because the analyticity radii of level $j$ run with $g_j$ and the radii of two levels must be comparable. For radii that do not run, no relation between the couplings is required; Ba{\l}aban's remark that asymptotic freedom is not assumed refers to this case. For the complete model, what is used is the lower half of asymptotic freedom. At the levels of creation $j\le k$ the small factors were deposited, and since $x_j\ge x_k-C$ one has $p_0(g_j)\ge c\,p_0(g_k)$. The terms $R^{(j)}$ carry the vacuum subtraction but no counterterm for the coupling, and they are marginal. On the support of the tests, the following bound holds for every $X\in\mathcal{D}_j$ that contains a site $z$ of $\Omega_k$ and remains within $cML^{k-j}$ of $z$ in level-$j$ units (a longer $X$ satisfies $d_j(X)\ge cL^{k-j}-C$ and has only the bound $2g_j^{\kappa_0}e^{-\kappa d_j(X)}$; around a site of $\Omega_n\setminus\Omega_{n+1}$, $j\le n<k$, the right-hand side holds with $n$ in place of $k$):

\begin{gather*}
\FitMath{|R^{(j)}(X,U_k)-R^{(j)}(X,\mathbf{1})|\le C(d_j(X)+3)^2(x_j/x_k)\,\eta_{k-j}^4\,g_j^{\kappa_0}e^{-\kappa d_j(X)},} \\[3pt]
\FitMath{\sum_{j\le k}(x_j/x_k)\,g_j^{\kappa_0}\le Cg_k^{\kappa_0-2}:}
\end{gather*}

these terms do not decrease as $k-j$ grows. Only the linear growth of $x_j$ makes them summable, and it does so for $\kappa_0>4$. Here $x_j/x_k$ is, up to a bounded logarithmic factor, the squared ratio of the analyticity radii of levels $k$ and $j$, that is, the relative quadratic evaluation weight of (ER15). If one assumes no more than $x_j\ge x_k-C$, the sum can grow to order $kg_k^{\kappa_0}$. It is this summability that makes a counterterm for them unnecessary. It is also why aged $R$-terms change $b_j$ by no more than $C\gamma^{\kappa_0-C}$, so that $\lambda_\sharp>0$ is preserved (3a). The upper half ensures that the thresholds of adjacent levels are comparable: $\varepsilon_{k+1}\le C\varepsilon_k$.

\emph{The large-field factor: the reasons for these forms.} The fields at which a plaquette test fails do not form a small set for the measures contained in $\mathbf{T}_k(Z)$. They do form a small set for the Wilson weight, since $e^{-x_jA}\le e^{-c\,x_j\varepsilon_j^2}=e^{-c\,p_0(g_j)^2}$. Therefore, if no $\mathbf{R}$ has acted on a component $X$, the small factor that a failed plaquette test carries there is obtained from $\mathbf{T}_k(X)$ applied to $\chi\,e^{-A(x_k(\cdot),U_k)}$, and never from $\mathbf{T}_k(X)1$. This weight is evaluated at the background, which depends on the variables of all the components simultaneously. For this reason, bounds that are uniform in the exterior variables do not combine multiplicatively. The phenomenon already occurs in measure theory. Take Lebesgue measure on $[0,\mathsf{l}]$ in each of two variables, $0<\epsilon\ll\mathsf{l}$, and $F=\mathbf{1}\{|\xi_1-\xi_2|<\epsilon\}$. Then

\begin{gather*}
\FitMath{\sup_{\xi_2}\int F\,d\xi_1=2\epsilon,} \\[3pt]
\FitMath{\iint F\,d\xi_1d\xi_2=2\epsilon\,\mathsf{l}-\epsilon^2\gg(2\epsilon)^2.}
\end{gather*}

The remedy is a pointwise bound by a product of functions, each depending only on its own group of variables, the groups being disjoint, combined with Tonelli's theorem. The small factors come out as numbers from the Wilson sums near the creations, joined with a fixed fraction both of the Gaussian forms and of the kept gauge-fixing exponents. What remains of the weights is bounded by functions that no longer depend on the background. The supremum in $\mathsf{M}$ is placed under the integral for the same reason. The factor of $X$ then becomes a function of $\xi$ only, and the masses multiply. When a component has dropped out, the supremum is taken over the evaluations at $\xi$. These are the real configurations that are compatible with the stored variables and regular layer by layer, and at which every test listed by $\omega$ takes its recorded value, with one exception: the test whose failure turned $X$ into a drop-out. They form a set fibred over $\xi$. A product set contains a flat $\mathbf{U}$, at which $e^{-A}=1$, so taking the supremum over it eliminates the Wilson weight. For a live component, a bound $e^{-c\,p_0(g_k)}$ does not reproduce itself, since a step without a creation costs $e^{+CM^4\check{R}_k^5(d'_k(Z)+1)}$. What is carried along is the credit. As long as nothing is created, linear sizes counted in cells decrease geometrically. At a drop of $\check{R}$ the two partitions coincide and nothing shrinks, and no two consecutive steps are both drops. The waiting lasts at most $\check{R}_j$ steps, unless two components join or a mark is set. In that case the waiting begins anew. In the count of (2c), every joining is charged to a separate created cell belonging to it alone, and a mark sets one factor $e^{-p_0(g_t)}$ of its mass against the same $\check{R}_t^6$. Hence the costs still owed add up to at most $CM^4L^5\check{R}_j^6(d'_j(Z)+1)$. The $p_0(g_j)^2$ per cube deposited at the creation covers $\check{R}_j^6$ per cube, and the powers of $\log g^{-2}$ are chosen in this order. After $\mathbf{R}$ has acted, the mass bound is what is carried forward.

Three losses are made good:
\begin{itemize}
\item The expansions use up part of $\kappa$ and of the radii. Cubes that are $L$ times larger give back a factor of order $L$ in the rate, and the curvature decreases by $L^{-2}$ at each level. $\mathbf{T}$ reserves a surplus of this gain for $\mathbf{R}$, whose re-expansions consume that surplus and restore exactly $\kappa$, $\alpha_i$.
\item The marginal term is absorbed into $x_k$.
\item $e^{-p_0(g_j)}$ is smaller than any power of $g_j$, but it is not a power of the spacing when $d=4$, so it cannot pay for all later steps. This is the reason for $\mathbf{R}$.
\end{itemize}

In the presence of large fields, $\log\rho_k$ is not a sum of local terms, and this is why the construction uses densities and sums over labels. Analyticity serves as a device. On a smaller complex domain, a supremum bound supplies derivative bounds by Cauchy's formula. Its transport through substitution, integration and logarithms uses the stated domain inclusions, integrable majorants and zero-free cluster estimates.

1a.2.8 is (1d) carried out with parameters. The size of $\mathbb{P}^{(k)}$ is $g|B|^3$, and under the test $g|B|<\delta(\mathsf{t})$ this is $\le\delta(\mathsf{t})|B|^2$. The bracket of level $j$ in $\{S\}$ is $O(\mathsf{q}_j^3\,\eta_{k-j}^{4+\alpha_{\mathrm{H}}})$, because every Cauchy estimate taken at a later background loses a power of $\mathsf{q}_j$, up to the third (this loss does not occur on radii that do not run). Applied to a shift of size $\delta$, Cauchy's estimate produces $\mathsf{s}$, and a cube of level $k$ contains $\eta_{k-j}^{-4}$ cubes of $\pi_j$. What is left at the age $n=k+1-j$ is $\mathsf{q}_j^3\,\eta_{k-j}^{\alpha_{\mathrm{H}}}\le C(n+1)^{3/2}L^{-\alpha_{\mathrm{H}}(n-1)}$, and these are the weights $\omega_n$. The $H^{(j)}$ are marginal, and their sum is controlled only through their norms. If the raw local terms are taken without the counterterms $b_jA$, then $\{S\}$ contains $\sum_{j\le k}b_j\{A\}\approx c_0k\{A\}$. In that case no bound holds uniformly in $k$, and the statement does not hold. Keeping $\mathsf{t}$ (which governs the tests, the gauge-fixing weight and the scaling) separate from $x$, the coefficient of $A$, comes at the price of a quadratic term proportional to $1-x/\mathsf{t}$. In return it allows frozen choices (3c) and complex weights (4a), (8a). Once the choices are fixed, the steps are fixed operations that are linear in the measure, and $\beta$ enters only through the weight of the action.

1a.2.9: on blocks with no additional node, the block axial representative satisfies $\mathrm{GF}_k(V^{\mathrm{bg}}_k)=0$. At a seam, the node terms described in 1a.3.11 can remain constant on the orbit and must be retained with the boundary terms. The following expansion uses the gauge-fixing factors on the blocks where the vanishing holds. The linear term $x_k\langle\mathsf{H}_k\mathcal{A}',J(U_{k+1})\rangle$ vanishes on $Q_{V^{\mathrm{bg}}_k}\mathcal{A}'=0$, since the current is the image of a Lagrange multiplier. Because the fibre is curved, the term $\langle\mathsf{H}_k\mathsf{h}C^{\mathrm{pl},(2)}(B),J\rangle$ remains, and it is included in $\Delta^{(k)}$. In this node-free interior expansion, $x_kA(U_{k+1})$ is the only term with a negative power of $g_k$. This assertion does not discard the seam contribution $x_k\mathrm{GF}_k(V^{\mathrm{bg}}_k)$ or its fluctuation variations. They must be retained in the boundary exponent, whose estimates we must prove separately. The choices are made in the following order. $b_{k+1}$ is obtained on the infinite lattice from the choices of level $k$, so $x_{k+1}$, $\varepsilon_{k+1}$, $\check{R}_{k+1}$ are fixed in advance of testing $V_{k+1}$. \emph{The new tests imply the old:} on $\Omega_{k+1}$, wherever the new functions are supported, the old ones take the value $1$. By the nesting, $U_{k,\square}=(e^{i\eta_k\mathbf{H}}U_{k+1,\square'})^u$. Here $\mathbf{H}$ denotes the response of (1b) when the datum is shifted, the shift being $O(\delta_k)$ everywhere but in a boundary layer at a distance of order $M_2\check{R}_k$. Therefore

\begin{gather*}
\FitMath{|\mathbf{H}|,|\nabla\mathbf{H}|\le C\big(\delta_k+e^{-\check{R}_k}\varepsilon_k\big)\le c\,\varepsilon_k,} \\[3pt]
\FitMath{|\partial U_{k,\square}-\mathbf{1}|\le2\varepsilon_{k+1}\eta_{k+1}^2+c\,\varepsilon_k\eta_k^2<\varepsilon_k\eta_k^2:}
\end{gather*}

so $\chi_k$ passes into $\chi_{k+1}$. This holds because the new threshold is smaller by the factor $L^{-2}\varepsilon_{k+1}/\varepsilon_k\le CL^{-2}$ (by the upper half of the flow inequalities, thresholds of neighbouring levels remain comparable), and because $\delta_k=(A_1/A_0)\varepsilon_k\ll\varepsilon_k$ (both carry the same power of the logarithm, and their ratio is small and fixed). Likewise every indicator in $\mathfrak{G}_k$ is equal to $1$, and $\chi'_k$ follows from $\chi^{(k)}$. Since the local backgrounds are so close to the global ones, the tests can be carried out on local objects, which factorise, while the expansions are made around global ones. \emph{The split:} an aged $R^{(j)}$ has each of its cubes inside a single block, its total transforms covariantly under the isometries of $\mathbb{T}^{(k+1)}$, and it is counted together with $\mathbf{E}_k$. For a young one this is not so, and a young one may not be admitted to the family from which $b_{k+1}$ is computed. The $t_1$-integral, which comprises boundary terms, the part of the coupling that depends on position, and the strip term $\mathsf{W}^{(k)}$, contributes only to $B^{(k+1)}$. The $t_2$-integral takes the expectation of the young terms in a measure which itself contains them. Consequently at least one factor $g_j^{\kappa_0}$ appears in every term of its expansion, and this integral contributes to both $R^{(k+1)}$ and $B^{(k+1)}$. The third term contains only regular terms. A further interpolation makes it independent of the exterior fluctuation field, and the resulting contribution goes to $B^{(k+1)}$. Its cluster expansion then produces, for $X\subset\Lambda_{k+1}$, the pure small-field integral's terms, $E^{(k+1)}$, and in addition boundary terms. $b_{k+1}$ is computed from this symmetric remainder. Thus a contribution to $R^{(k+1)}$ comes from $\mathbf{T}$ already, ahead of the one that $\mathbf{R}$ adds. 1a.2.10: iterating the Stokes formula of 1a.1.4 gives the first display. When a standard Gaussian vector is conditioned on the value of a linear map, its conditional expectation projects orthogonally onto the complement of the kernel of that map, and its components in mutually orthogonal subspaces are independent. The slice is a Schur complement, and $\mathcal{H}_{j-1}$ transports it. The only property of it that has to be analysed is its decay, and this is the analysis of (1c) at $\mathbf{1}$.

\emph{Gauges.} 1a.3.1 is Stokes' formula on a tree. Because of the factor $|b_--y_P|$, boxes have side $O(M)$ and $M\varepsilon_k$ has to be small; that is, the coupling is taken small after $M$. 1a.3.2 is a descent through the levels, with increments $C\alpha_0\sum_iL^{-2i}$. 1a.3.3 follows from Fubini's theorem and the invariance of Haar measure. No determinant arises, because the gauge function is a group variable translated by $u(x)$. The mean is used because single combs break $W_4$. The constraint is taken soft because fixing the mean by a delta function imposes a non-linear restriction which couples all the bonds of the paths to one another and has no explicit solution. The weight is $x_k$. Let $P^{\prime\prime}$ be a family of blocks with no additional node, on which every background contour variable is $\mathbf1$. Restricting the gauge sum to these blocks gives

\begin{gather*}
\FitMath{x_k\,\mathrm{GF}_k(e^{ig_kB}V^{\mathrm{bg}}_k;P^{\prime\prime})=\tfrac12\sum_{y\in P^{\prime\prime}}\sum_{x\in B(y)\setminus\{y\}}\big[|B(y,x)|^2+O(g_k|B|^3)\big].}
\end{gather*}

The quadratic coefficient is of order one in these Gaussian variables. The omitted seam-node factors are retained, with their constant and linear contributions, in the conditional boundary estimates; this formula provides no bound for them. Here $g_kB(y,x)$ denotes the linear part in $B$ of $\frac1i\log$ applied to the mean $(e^{ig_kB}V^{\mathrm{bg}}_k)(y,x)$, and when $V^{\mathrm{bg}}_k=\mathbf{1}$, $B(y,x)$ coincides with the mean of $B(\Gamma)$ taken over $\mathcal{G}(y,x)$. This quadratic gauge form is strictly positive on each nonzero residual gauge direction $D\lambda$ supported in the blocks $P^{\prime\prime}$, with $\lambda$ zero at their coarse sites: there $|B(y,x)|=|\lambda(x)|$. This local calculation does not identify the full layered null space. Its removal and the coercivity of the complete fluctuation form remain the separate requirements of A2 in the dependency catalogue. The weight itself is auxiliary: the first identity of 1a.3.3 yields an identity for every choice of weight, so the weight may be frozen (3c) and need not depend on the source (4a). \emph{Procedure and product:} the insertion, taken as a draw, is exact on any measure, and on invariant measures it agrees with the written product. A term with stored fields is invariant. Indeed, the distance from $P''$ to $Z_k$ is $\ge cM\check{R}_k$. Hence every transformation over $P''$ can be lifted to a bare transformation which is constant on pillars and equal to $\mathbf{1}$ near $Z_k$, the region where all stored quantities enter. In the case of the trees of (2b), the stored fields are acted on by the group via the tops of their pillars. The exact model 1a.3.13 shows that this verification is not idle: the group that acts is decided by whether or not the seam bonds lie in the determining set. The removal identity shows that dropping $\delta_T$ from the density gives an equivalent density, not an equal one: the two coincide when integrated against invariant functions.

1a.3.4: in the pillars, $U$ and $U_0$ have the same contour variables. Hence $W(y,x)U_0(y,x)^{-1}=u(y)^{-1}\operatorname{Ad}(U_0(y,x))u(x)$, and the block factor $\mathsf{n}_W$ becomes the definition of $u^{-1}\bar{R}_{U_0}u$. This is Ba{\l}aban's deeper reason for imposing the axial gauge first and for keeping a single contour variable throughout (with two of them, circulations of $A$ survive in the linear part). \emph{Origin of $R$:} the term sought is Feynman's $\|\partial^*A\|^2$, but only the $\lambda\in N(Q')$ leave $\delta(B-Q^{\sharp}A)$ unchanged. Over these, the Faddeev--Popov quotient with gauge parameter $\alpha_{\mathrm{g}}$ is

\begin{gather*}
\FitMath{e^{-\frac1{2\alpha_{\mathrm{g}}}\|\partial^*A\|^2}\Big/\int d\lambda\,\delta(Q'\lambda)\,e^{-\frac1{2\alpha_{\mathrm{g}}}\|\partial^*A-\Delta\lambda\|^2}=\mathrm{const}\cdot e^{-\frac1{2\alpha_{\mathrm{g}}}\|R\,\partial^*A\|^2}.}
\end{gather*}

In a background the constants turn into determinants that depend on $U_0$:

\begin{gather*}
\FitMath{e^{-\frac12\langle A,\Delta_\pi A\rangle}=\big|\det\big(\Delta_{U_0}|_{N(Q')}\big)\big|\int d\lambda\,\delta(Q'\lambda)\,\delta_R(RD^*A-\Delta\lambda)\,e^{-\frac12\langle A-D\lambda,\Delta^{\mathrm{W}}(A-D\lambda)\rangle}:}
\end{gather*}

the determinant here is the Landau gauge's Faddeev--Popov determinant, and axial gauges have no such determinant. In normalised expressions such as $\mathsf{H}_kB$ it cancels; in Gaussian normalisers it does not. In the one-loop coefficient (1e), (3a) it enters together with the determinant of the gauge-fixed Hessian. Positivity (1a.3.5): consider $U=\mathbf{1}$ (the remaining cases are perturbations of it). A gradient $\partial\lambda$ ($\partial=D_{\mathbf{1}}$) is detected by $Q^{\sharp}$ if $\partial Q'\lambda\ne0$, and by $R\,\partial^*$ otherwise, since in that case $\lambda\in N(Q')$ once a constant is subtracted. The full $DD^*$ is positive as well, but admissible transformations cannot reach $D^*A=0$, and with the full operator $\mathsf{G}Q^{\sharp*}(Q^{\sharp}\mathsf{G}Q^{\sharp*})^{-1}$ is no longer the minimiser of the Hessian. $RD^*\mathsf{G}Q^{\sharp*}=0$ requires the invariant extension $\Delta_\pi$, which is a perturbation of size $CM\alpha_0$. (Ba{\l}aban's $G_0$ is defined by the form of the statement, and his $G$ goes with the identity. Adding the term of 1a.1.10 gives $G_1$. The three satisfy the same bounds, with the exception of the one with the covariant Laplacian, and they coincide where $J=0$.)

1a.3.7 (elliptic regularity through $\mathsf{G}_k(U_0)$): both currents, each smaller than $\alpha_0$, give $D^*DA=\mathcal{J}$ with $|\mathcal{J}|_{(-3)}\le2\alpha_0+C\alpha_2|\nabla A|_{(-2)}+\dots$. The averages give $\xi Q^{\sharp}A=B-C(A)=:\mathsf{b}$ with $|C(A)|\le C\alpha_2^2$ (1a.1.6 leaves a remainder that is of second order), and this accounts for the $\alpha_1$ in the result and for the condition on $\alpha_2^2$. On the slice,

\begin{gather*}
\FitMath{A=\mathsf{G}_k(U_0)\mathcal{J}+\mathsf{G}_k(U_0)Q^{\sharp*}\xi^{-3}\mathsf{b},} \\[3pt]
\FitMath{\|A\|_1\le B_0\big(|\mathcal{J}|_{(-3)}+\sup|\mathsf{b}|\big),}
\end{gather*}

and absorption of the gradient term takes place as soon as $B_0C\alpha_2\le\frac12$. 1a.3.6 (Ba{\l}aban's argument, written out for the scheme $\mathsf{c}$) is proved by induction on $k$. On $\Omega_k$ the pair $(A,u_1)$ of depth $k-1$ is a factor $L$ worse; here $u_1$ denotes the explicit product from 1a.1.12. One seeks $u=u_1e^{i\lambda}$. By Campbell--Hausdorff, once the non-linear average has been straightened to $Q'\lambda=0$,

\begin{gather*}
\FitMath{\lambda=-G'RD^*A+G'R\,\mathfrak{N}_4(\lambda,D\lambda,A,D^*A),} \\[3pt]
\FitMath{|R\mathfrak{N}_4|_{(-2)}\le CB_1(\alpha_0+\alpha_1)\|\lambda\|_0,}
\end{gather*}

and this is a contraction on complex $\lambda$ in $\|\lambda\|_0$. (The bound is valid on the ball of radius $CB_1(\alpha_0+\alpha_1)$, which contains the solution; within a ball of fixed radius the solution is unique.) Curvature and mismatch thus have to beat only norms that depend on $L$. The new constant is of order $LB_1$, and 1a.3.7 resets it to $B_1$; without this reset its growth in $k$ is geometric. Uniqueness holds among small $A$, so no Gribov ambiguity arises. In (1b) the passage to the slice is redone for $\mathcal{M}$. There Brouwer's theorem replaces the contraction. It provides a pair, though not its uniqueness, and with it the bounds on $\|A\|_1$ and $|D^*DA|_{(-3)}$. For a sequence with layers, the induction is replaced by a single preliminary gauge followed by an ascent through the truncations $(\Omega_0,\dots,\Omega_j)$. 1a.3.8: the cube is surrounded by a tower that descends to the bare scale. After passing to tower and box axial gauges, one applies the theorem to $(\mathbf{1},U)$ within the scheme $\mathsf{c}$. Only propagators at $\mathbf{1}$ enter, and this breaks the circularity with (1c). At complex fields we have $v=w^{-1}$ and $w(x)=\mathbf{U}(\Gamma_{y_\square,x})=\Pi\,v_{\mathrm{r}}^{-1}$. Here $\Pi$ is a product of factors $\exp(i\eta_k\,tA'(b)t^{-1})$, where each $t$ is a unitary transport of $U$, so $\operatorname{cn}(v)=\operatorname{cn}(\Pi)$. The potential of the background inside a large-field box is obtained in (2b) by the same sequence of steps: the box axial gauge, used for the average, beneath it the tower axial gauge, and then Landau. Every change of its gauge is compensated on the stored variables. Local gauges $\hat{k}_y$ on overlapping boxes are glued together by $k(x)=\mathcal{N}\big(\sum_y\chi_y(x)\hat{k}_y(x)\big)$, $\mathcal{N}(X)=X/\sqrt{\det X}$, using real weights $\chi_y(x)\ge0$ with $\sum_y\chi_y(x)=1$. A convex combination of sufficiently close elements of $\mathrm{SU}(2)$ is nonzero and is a positive scalar times a group element. On a complex neighborhood of these combinations, $\mathcal{N}$ is holomorphic for the local square-root branch continuing the positive real root.

1a.3.9: with weighted norms, or on complex balls, the induction incurs a loss at every level, and there is no reset. Instead one differentiates the identity of 1a.3.7, with $\mathsf{G}_k$ bounded in the sup norms that carry the weight $e^{\delta d(\cdot,y')}$. This is a single linear estimate and needs no induction. That the averages are trivial does not force $u$ to be trivial at the nodes, and this gives rise to the node factor. Invariant functions do not detect it. At sites that carry stored fields but have ceased to be tops, $\mathsf{n}^{-1}$ is continued into each pillar by means of a smooth bump, and in this way the stored data are not acted on. 1a.3.10: the objects are words in transports and logarithms, each typed at a site, and fixed points are covariant by uniqueness. 1a.3.6 gives the uniqueness of its Landau pairs only in the scheme $\mathsf{c}$; the set of these pairs is nevertheless covariant, since every one of the three defining conditions of a pair is covariant. The case of complex $u$ follows by the identity theorem, because the $G$-valued $u$ form a totally real submanifold of full dimension. 1a.3.11 is a matter of covariance only: $(\mathcal{M}^iU^w)(y,x)=w(y)(\mathcal{M}^iU)(y,x)w(x)^{-1}$ with $w=\mathbf{1}$ at the nodes. 1a.3.12: the equation becomes $\lambda=G'(f-RD^*A)+G'R\,\mathfrak{N}_4$, whose inhomogeneity is of the size of the other terms. Ba{\l}aban states that the remaining proofs carry over with only minor changes. However, $\Delta_{U_0}A-D^*DA$ equals $Df+D(1-R)D^*A$ up to curvature terms, so the bounds of the theorem force $|Df|_{(-3)}\le C(\alpha_0+\alpha_1)$. On the other hand, $|Df|$ is of order $\eta_k^{-1}|f|$ for $f=\Delta\lambda$ with $\lambda\in N(Q')$ oscillating at the bare spacing. 1a.3.13: the contour variables of a block are coordinates that preserve Haar measure. All the integrals except one give $\int f_{\mathfrak{G}}=1$, and the remaining one, $\int dt\,f_{\mathfrak{G}}(t)\,\delta(tV(s)\mathsf{p}V^+(c)^{-1})$, is the display. In four dimensions the same occurs at each coarse site whose block contains an end of a seam bond.

\emph{Routine and delicate.} Routine are the algebra of means, of logarithms and of covariance, the exact identities, Stokes' formula and the Haar identities. The delicate points are the following:
\begin{itemize}
\item uniformity in $k$, which is reached solely via the straightened average, and only because curvature and potentials are measured on the scale of each level;
\item complex arguments, where each use of covariance costs room in the domain;
\item the two averages, together with the invariances of the constraints they impose, which have to be kept apart;
\item normalisation by averages instead of pointwise normalisation, the group property that is lost, and the node factor that is left over;
\item the junctions of layers;
\item the invariance of whatever an insertion encounters after fields have been stored.
\end{itemize}

For the construction as a whole, the difficulty lies in the words ``with the same constants''.

% END MODULE chapters/02_averaging.tex


% BEGIN MODULE chapters/block_average_pivot_chart_v8.tex
\section{The nonlinear block average and its pivot chart}
\label{ch:block-average-pivot-chart}

The finite calculation here concerns the actual centered one-step
average.  There are two means in its definition, of different kinds.
An implicit group mean defines each internal contour transport.
An explicit exponential of an arithmetic mean of closed logarithms
then defines the block bond.  We retain both operations.  The
calculation establishes the local pivot and quadratic-remainder
estimates. The many-level background and exterior-force estimates enter
through their separate analytic propositions.

\subsection*{The exact finite map}

Work on a unit-spacing lattice in dimension four, with odd \(L\ge3\).
The block at a coarse site \(y\) is
\(B(y)=\{x:|x-y|_\infty\le(L-1)/2\}\).
For a positive coarse bond \(c=(y,y+Le_\mu)\), put
\(y'=y+Le_\mu\) and \(D_c=B(y)\cup B(y')\).
On a torus we require that this double block and its one-bond
neighborhood have an embedded rectangular lift preserving their
unit bond graph, with no periodic edge joining opposite faces of
the patch.  A sufficient condition for a cubic torus is that its
period exceeds \(4L+4\).  All bonds used below have both endpoints
in \(D_c\), with the rectangular bond graph of this lift.

For \(x\in B(y)\), let \(\mathcal G(y,x)\) be the shortest paths
which complete one coordinate displacement before starting the next.
There are at most \(24\) such paths, each of length at most \(2L\).
Repeated paths, if retained, are counted with their prescribed equal
weights.  The transport along an oriented path is the ordered product
of its bond matrices; reversed bonds use the inverse matrix.
For nearby matrices \(T_1,\ldots,T_n\), define their contour mean
as the solution near them of
\[
 \frac1n\sum_{j=1}^n\log(T_jT^{-1})=0.
 \tag{QC1}
\]
All logarithms below are the local analytic logarithm near the
identity.  Set
\[
 U(y,x)=\operatorname{av}\big(U(\Gamma):\Gamma\in\mathcal G(y,x)\big),
 \qquad U(x,y)=U(y,x)^{-1}.
 \tag{QC2}
\]
With \(x'=x+Le_\mu\), the actual block map is
\[
 \begin{split}
 Z_x(U)&=U(y,x)U([x,x'])U(x',y')U[c]^{-1},\\
 \mathcal M(U)(c)
   &=\exp\left[L^{-4}\sum_{x\in B(y)}\log Z_x(U)\right]U[c].
 \end{split}\tag{QC3}
\]
Here \([x,x']\) and \([c]\) are straight segments of length \(L\).
In particular, the outer operation in (QC3) is not the implicit
mean (QC1).  The cruder Borel definitions of the block map away
from this analytic neighborhood do not enter the calculation.

For a background \(V\), let \(W=\mathcal M(V)\).  In physical bond
logarithms, with no scale factor in the exponential, put
\[
 q_{V,c}(a)=\frac1i\log\left[
       \mathcal M(e^{ia}V)(c)\mathcal M(V)(c)^{-1}\right].
 \tag{QC4}
\]
The matrices \(a(b)\) are complex traceless \(2\) by \(2\) matrices,
and their norm is the supremum of the operator norms over positive
bonds.  On real variables they are Hermitian.  All constants denoted
\(C,c\) below are universal; their values may be enlarged or decreased
at successive occurrences.

\subsection*{The contour mean, with quantitative bounds}

The elementary matrix inequalities
\[
 \|e^X-I-X\|\le\tfrac12e^{\|X\|}\|X\|^2,
 \qquad
 \|\log(I+H)-H\|\le\frac{\|H\|^2}{2(1-\|H\|)}
 \tag{QC5}
\]
hold for \(\|H\|<1\), by the power series and submultiplicativity.
They show, on a fixed small ball, that
\[
 \begin{aligned}
 \log(e^Xe^{-Z})&=X-Z+E(X,Z),\\
 \|E(X,Z)\|&\le C(\|X\|+\|Z\|)^2,\\
 \|D_ZE(X,Z)\|&\le C(\|X\|+\|Z\|).
 \end{aligned}
 \tag{QC6}
\]
For the derivative estimate one can differentiate the convergent
logarithm series term by term and use
\(D e^Z[H]=\int_0^1e^{(1-t)Z}He^{tZ}\,dt\).
The derivative of each remainder in (QC5) is bounded by a constant
times the radius.  Thus the constants in (QC6) depend on no number
of contour paths.

Write \(T_j=e^{X_j}\), with \(\max_j\|X_j\|\le t\), and seek
\(T=e^Z\).  Equation (QC1) is equivalent to
\[
 Z=\overline X+\frac1n\sum_jE(X_j,Z),
 \qquad \overline X=\frac1n\sum_jX_j.
 \tag{QC7}
\]
For a fixed sufficiently small \(t\), the right side maps
\(\|Z\|\le2t\) into itself and has Lipschitz constant at most
\(1/2\).  Iteration from \(\overline X\) proves existence and
uniqueness there.  It also proves
\(Z=\overline X+O(t^2)\), uniformly in \(n\).
The iterates are holomorphic in all \(X_j\), and convergence is
locally uniform on a slightly smaller fixed ball.  Their limit is
therefore holomorphic.  Trace zero is preserved in (QC7).
For anti-Hermitian \(X_j\), anti-Hermiticity is preserved as well,
so the mean lies in \(\mathrm{SU}(2)\) on real group inputs.

Uniqueness proves permutation symmetry and left-right covariance.
Indeed under \(T_j\mapsto uT_jv\), with \(T\mapsto uTv\), every
logarithm in (QC1) is conjugated by \(u\), whenever the transformed
families remain in their analytic domains.  The inverse property
follows from
\(\log(T_j^{-1}T)=-\operatorname{Ad}(T^{-1})\log(T_jT^{-1})\).
Thus the locally constructed mean is the mean in (QC2), including
families centered at a unitary matrix other than \(I\).

\subsection*{Local gauges do not constrain coarse holonomies}

For a real background \(V_0\) on \(D_c\), put
\(\varpi_c=\sup_{p\subset D_c}\|\partial V_0(p)-I\|\).
There is a unitary gauge \(s_c\) on this rectangle such that
\[
 \sup_{b\subset D_c}\|V_0^{s_c}(b)-I\|
                                      \le C L\varpi_c.
 \tag{QC8}
\]
Here \(V^s(b)=s(b_-)V(b)s(b_+)^{-1}\).
To prove (QC8), fix a coordinate-order comb from \(y\) to every
site of the rectangle and set \(s_c(x)=V_0(\Gamma_{y,x})\).
For any bond \(b\), the product
\(s_c(b_-)V_0(b)s_c(b_+)^{-1}\) is its comb loop.  Slide the
additional bond past the later coordinate segments of the comb.
Each slide exchanges two adjacent perpendicular steps and inserts
one plaquette, conjugated by a partial transport.  There are at
most \(CL\) slides, since every coordinate interval in the rectangle
has length at most \(2L\).  All these conjugations are unitary.
The norm of a product of such plaquette factors minus \(I\) is
at most the sum of their deviations, proving (QC8).

The same gauge, fixed from the real background, works for
\(V=e^{iE}V_0\): its distance from \(I\) is at most
\(C(L\varpi_c+\|E\|_\infty)\), when these quantities are small.
More generally we can start directly with a unitary local gauge
in which
\[
 \sup_{b\subset D_c}\|V^{s_c}(b)-I\|\le\epsilon .
 \tag{QC9}
\]
Different double blocks may use different gauges.  We never require
one gauge close to \(I\) on an entire observation window.

The group mean and (QC3) are gauge covariant.  The internal transport
transforms with its two endpoint factors; each \(Z_x\) transforms
by conjugation at \(y\).  Consequently
\[
 q_{V^{s_c},c}(a^{s_c})
    =\operatorname{Ad}(s_c(y))q_{V,c}(a),
 \qquad a^{s_c}(b)=\operatorname{Ad}(s_c(b_-))a(b).
 \tag{QC10}
\]
All these adjoint actions have norm one.  A coarse bond value
\(\mathcal M(V)(c)\) in the original gauge may be anywhere in
\(\mathrm{SU}(2)\).  It is the ratio in (QC4), rather than that
coarse bond itself, whose logarithm is taken near \(I\).
Nontrivial coarse holonomies therefore cause no loss in these
local estimates.

\subsection*{Uniform Taylor estimates for the actual average}

We first work in the gauge (QC9), so \(V=e^{iv}\) with
\(\|v\|_\infty\le C\epsilon\).  If
\(\|v\|_\infty,\|a\|_\infty<c_0/L\), every path product entering
(QC2) stays in a fixed small neighborhood of \(I\).
For example, a product of \(N\) factors, each within \(t\) of \(I\),
is within \(e^{CNt}-1\) of \(I\).  Its reversed factors satisfy
the same bound.  Here the internal paths have \(N\le2L\).
The contour-mean estimate just proved applies, and the closed
products in (QC3), of total path length at most \(6L\), are also
in a fixed logarithm neighborhood.  The outer logarithm average
is normalized by \(L^{-4}\); its size is bounded by its largest
summand, without a factor \(L^4\).

It follows that
\[
 F_c(v,a)=\frac1i\log\left[
   \mathcal M(e^{ia}e^{iv})(c)\mathcal M(e^{iv})(c)^{-1}\right]
 \tag{QC11}
\]
is holomorphic and bounded by a universal constant on the joint
supremum balls of radius \(c_0/L\).  It vanishes at \(a=0\).
This establishes the domain before any derivative estimate.
On balls a fixed fraction smaller, two scalar Cauchy formulas
give
\[
 \|D_a F_c\|_{\infty\to\mathrm{op}}\le CL,
 \qquad
 \|D_a^2F_c\|+\|D_vD_aF_c\|\le CL^2.
 \tag{QC12}
\]
For example, for unit supremum directions \(u,w\), apply Cauchy
to \(F_c(v,a+zu+tw)\), with both scalar circles of radius
\(c/L\).  This proves the mixed bilinear bound without dividing
the radius by the number of bond coordinates.

Let \(K_V=Dq_V(0)\) and \(R_V(a)=q_V(a)-K_Va\).
Taylor's formula with integral remainder gives
\[
 \begin{aligned}
 \|K_V\|_{\infty\to\infty}&\le CL,\\
 \|R_V(a)\|_\infty&\le CL^2\|a\|_\infty^2,\\
 \|DR_V(a)\|_{\infty\to\infty}&\le CL^2\|a\|_\infty .
 \end{aligned}
 \tag{QC13}\label{qc:quadratic-residual}
\]
for \(\|a\|_\infty<c/L\) and \(\epsilon<c/L\).
The estimates transfer to the original frame by (QC10).
They are uniform over all coarse bonds, hence over the volume.
In addition, varying the background in (QC12) gives
\[
 \|Dq_{e^{iv}}(0)-Dq_I(0)\|_{\infty\to\infty}
                                      \le CL^2\|v\|_\infty.
 \tag{QC14}
\]
These are estimates for the explicit nonlinear average (QC3), not
for a model linear average.

\subsection*{The pivot coefficient and finite elimination}

Write \(h=(L-1)/2\).  The pivot of \(c=(y,y+Le_\mu)\) is
\(b_0(c)=(y+he_\mu,y+(h+1)e_\mu)\).
It crosses from \(B(y)\) to \(B(y')\).  No internal contour path
contains it.  Among all coarse pivots, only \(b_0(c)\) occurs
in the support of \(q_{V,c}\).  Indeed internal contours stay
inside one block; each straight segment in (QC3) crosses only
the face between these two blocks and is parallel to \(e_\mu\).
The only coarse pivot on that face is \(b_0(c)\).  This support
assertion, rather than a claim about all bonds joining two
vertices of a periodic double block, is what is used below.

At the identity background define
\(\overline a(y,x)\) to be the arithmetic mean of the oriented
path sums of \(a\) over \(\mathcal G(y,x)\), and
\(\sigma_a(y)=L^{-4}\sum_{x\in B(y)}\overline a(y,x)\).
The first derivative of the contour mean is this arithmetic mean.
Differentiating (QC3), the negative derivative of \(U[c]\) in
each closed logarithm cancels its final positive derivative.
Thus
\[
 (K_Ia)(c)=\sigma_a(y)
       +L^{-4}\sum_{x\in B(y)}a([x,x'])-\sigma_a(y').
 \tag{QC15}
\]
The two \(\sigma_a\) terms contain no pivot.  Exactly \(L\)
segments in the remaining sum contain \(b_0(c)\): their transverse
coordinates equal those of \(y\), and their longitudinal starting
coordinate can be any of the \(L\) coordinates in the block.
All occurrences have positive orientation.  Therefore
\[
 J_c(I):=D_{a(b_0(c))}q_{I,c}(0)=L^{-3}I_3.
 \tag{QC16}
\]
For a pure gauge \(V^s=I\), the coefficient in its original frame is
\[
 J_c(V)=L^{-3}\operatorname{Ad}(s(y)^{-1}s(b_0(c)_-)).
 \tag{QC17}
\]
The unitary matrix in this formula is the prefix transport from
\(y\) to the start of the pivot.  Replacing it by \(I\) in the
original gauge would be incorrect.

From (QC14), after the unitary local change of frame, we have
\(\|J_c(V^{s_c})-L^{-3}I_3\|\le CL^2\epsilon\).
Consequently, under the conservative condition
\[
 \epsilon\le cL^{-5},
 \qquad\text{in particular if }\quad
       \|E\|_\infty+L\varpi_c\le cL^{-5},
 \tag{QC18}
\]
every \(J_c\) is invertible and \(\|J_c^{-1}\|\le2L^3\).
There is no off-diagonal pivot coefficient.  Set
\[
 (h_V\nu)(b_0(c))=J_c^{-1}\nu(c),\qquad
 (h_V\nu)(b)=0\ (b\text{ nonpivot}),
 \qquad
 \mathsf C_V=\iota-h_VK_V\iota,
 \tag{QC19}
\]
where \(\iota\) extends the nonpivot coordinates by zero.
Then
\[
 K_Vh_V=1,\quad K_V\mathsf C_V=0,\quad
 (\mathsf C_Vu)|_{\rm free}=u,
 \qquad
 \|h_V\|\le2L^3,\quad\|\mathsf C_V\|\le CL^4 .
 \tag{QC20}\label{qc:pivot-linear-maps}
\]
These are supremum operator norms.  The same powers bound column
norms.  Indeed a fine bond belongs to at most a fixed number of
double blocks, and (QC12) bounds the coefficient in each of their
rows.  The maps are local: a pivot row uses only its own double
block.  They are holomorphic in the background, by the Neumann
series for \(J_c^{-1}\).

For comparison with a convention which inserts \(a/L\) in the
one-step exponential, its linear average is \(Q_V=K_V/L\), and
its pivot coefficient is
\[
 \mathsf m_c=J_c/L,\qquad
 \mathsf m_c(I)=L^{-4}I_3,
 \qquad \det J_c=L^3\det\mathsf m_c.
 \tag{QC21}
\]
This accounts for the factor \(L^{-3}\) in a real elimination
Jacobian; it is not a change in the number of color coordinates.

\subsection*{The nonlinear graph, constructed rather than assumed}

For a full physical bond variable \(t\), solve
\[
 d=R_V(t-h_Vd),\qquad \Psi_V(t)=t-h_Vd(t).
 \tag{QC22}
\]
On \(\|t\|_\infty<r_1=cL^{-5}\), (QC13) maps the ball
\(\|d\|_\infty\le CL^2\|t\|_\infty^2\) into itself, after
choosing its fixed prefactor large enough.  Its Lipschitz constant
in \(d\) is at most \(CL^5r_1<1/2\).  Iteration proves the unique
solution, its holomorphic dependence, and
\[
 \begin{aligned}
 q_V(\Psi_V(t))&=K_Vt,&
 \|d(t)\|&\le CL^2\|t\|^2,\\
 \|D_td(t)\|&\le CL^2\|t\|,&
 \|D\Psi_V(t)-1\|&\le CL^5\|t\|.
 \end{aligned}
 \tag{QC23}
\]
The derivative estimate follows by differentiating the contraction
and inverting \(1+DR_Vh_V\) by its convergent Neumann series.
Only the pivot values are changed.  There is even an explicit
inverse, on a smaller neighborhood:
\(\Psi_V^{-1}(a)=a+h_VR_V(a)\).  Substitution in (QC22) verifies
this formula by uniqueness.

Each equation for \(d(c)\) uses only its own pivot value and
the nonpivot data in \(D_c\).  The equations for different
coarse bonds are independent once the nonpivot data are fixed.
Thus the graph is exactly finite range, not merely exponentially
local.  In particular
\[
 a(u)=\Psi_V(\mathsf C_Vu),\quad q_V(a(u))=0,\quad
 a(u)|_{\rm free}=u,
 \qquad \|u\|_\infty<cL^{-9}
 \tag{QC24}
\]
is the actual small fibre chart.  It obeys
\(\|a(u)\|_\infty+\|D_ua(u)\|\|u\|_\infty
 \le CL^4\|u\|_\infty\).
Conversely, any sufficiently small solution of \(q_V(a)=0\)
has this form: \(t=a+h_VR_V(a)\) satisfies \(K_Vt=0\), hence
\(t=\mathsf C_V(a|_{\rm free})\), and (QC23) gives uniqueness.

For the canonical Gaussian coordinate \(u=gy\), the physical
chart therefore contains \(\|y\|_\infty<cL^{-9}/g\).
At weak coupling it contains any radius \(c_*\alpha/g\) with
fixed \(c_*\) and sufficiently small \(\alpha\).
This establishes the one-step primitive radius; inclusion in
all the older background domains remains a different estimate.

\subsection*{Real coarea normalization and complex continuation}

We record the normalization only on the real positive chart first.
Let \(\mathcal E,\mathcal C,\mathcal F\) be the fine positive bonds,
coarse positive bonds and nonpivots, respectively.  Thus
\(|\mathcal F|=|\mathcal E|-|\mathcal C|\).
In real Lie algebra coordinates Haar measure is
\(d(e^{ia})=\sigma(a)d^3a\), where
\(\sigma(a)=\sigma_0(\sin|a|/|a|)^2\) and \(\sigma_0>0\).
Group delta in the same chart satisfies
\(\delta_G(e^{iq})=\sigma_0^{-1}\delta_{\mathbb R^3}(q)\).

In the linear variables \(t=\mathsf C_Vu+h_Vv\), free coordinates
are \(u\) and \(K_Vt=v\).  Its determinant is
\(\prod_c\det J_c^{-1}\).  The full straightening \(a=\Psi_V(t)\)
has derivative close to the identity by (QC23).  At real
backgrounds, the determinants \(\det J_c\) are positive: in the
local unitary frame \(J_c\) is within half of the positive scalar
matrix \(L^{-3}I_3\), and the real adjoint frame transformations
have determinant one.  Also \(\det D\Psi_V>0\), by continuity
from its value one at zero and the uniform derivative bound.
The ordinary real change of variables therefore gives, on the fibre,
\[
 \begin{split}
 &\prod_{b\in\mathcal E}dU_b
       \prod_{c\in\mathcal C}
            \delta_G(\mathcal M(U)(c)W(c)^{-1})\\
 &\quad=\sigma_0^{|\mathcal F|}
       \prod_c(\det J_c)^{-1}
       \det D\Psi_V(\mathsf C_Vu)
       \prod_{b\in\mathcal E}\frac{\sigma(a_b(u))}{\sigma_0}
       \prod_{b\in\mathcal F}d^3u_b .
 \end{split}\tag{QC25}
\]
This is an identity of local fibre measures, with the group
delta interpreted by this change of coordinates.
The straightening derivative is block triangular: its free block
is the identity, and its pivot diagonal has one three by three
block per coarse bond.  Hence its log determinant is a sum of
local trace logarithms.  Substituting \(u=gy\) in (QC25) supplies
\((\sigma_0g^3)^{|\mathcal F|}\), while
\((\det J_c)^{-1}=L^{-3}(\det\mathsf m_c)^{-1}\).

For complex backgrounds the coefficient functions in (QC25)
have their coherent holomorphic continuations on the chart just
constructed.  In particular \((\det J_c)^{-1}\) and the branch
of \(\log\det J_c\) continuing its real value are used, not an
absolute value of a complex determinant.  The Haar density is
continued by the entire power series
\(\sigma_0[\sin\sqrt z/\sqrt z]^2\),
\(z=\tfrac12\operatorname{tr}a^2\), rather than by a complex
operator norm.  No complex measure or complex delta function is
being asserted by this continuation.

\subsection*{Dependence on the remaining scales}

The average, its raw Taylor constant \(CL^2\), and the local
pivot construction contain no \(M\), observation radius or
number of free variables.  If the supplied local near-pure-gauge
bound is only \(\|V^s-I\|\le C_bMa\), (QC18) requires
\(a\le c/(C_bML^5)\).  Alternatively the direct real condition
\(\varpi_c\le cL^{-6}\), with a small complex perturbation, suffices.
The input-to-chart constants are conservative fixed powers of \(L\);
no claim of optimal powers is needed.

If only a complex gauge with condition numbers at most \(K\ge1\)
is available, conjugation charges one factor \(K\) per input
direction and one for the output.  Thus safe bounds are
\(C K^3L^2\) for the quadratic residual,
\(2K^2L^3\) for \(h_V\), and \(CK^4L^4\) for \(\mathsf C_V\).
The conservative physical nonpivot radius becomes
\(cK^{-9}L^{-9}\).  A condition number growing with a larger
window must therefore be charged; it cannot be discarded by
calling the one-step chart local.

The calculation discharges the one-step pivot and Taylor primitives
in these actual physical coordinates.  It does not bound the
many-level response or a further coordinate transformation on a
large run window.  In particular, a polynomial loss in that window
radius in the background/core inclusion must still be retained
in a later Cauchy radius.  The free precision gap, the all-window
core and tower estimates, the conditional exterior forces and the
clean-class law remain separate assertions.

% END MODULE chapters/block_average_pivot_chart_v8.tex


% BEGIN MODULE chapters/flat_contour_precision_gap_v8.tex
\section{A gap for the classical flat fluctuation precision}
\label{ch:flat-contour-precision-gap}

The gauge form in this section is the actual soft contour gauge.
It is not a divergence projector.  At the flat field its quadratic
form is the sum of the squares of the arithmetic internal contour
path sums.  The constraint is the plain derivative of the nonlinear
block average, rather than the straight average alone.  With these
operators the classical flat precision has a gap uniform in the
torus volume and in the number of finer averaging levels.

The statement concerns the unweighted Wilson action at a flat
background, with its prescribed positive soft gauge coefficient.
Signed counterterms, nonconstant layer weights, nonzero-current
constraint corrections and complex decoupling matrices are not
included in the conclusion.

\subsection*{The actual quadratic forms and their normalization}

Use raw bond logarithms on a unit-spacing four-dimensional lattice.
For a scalar site function \(f\) and a positively oriented bond
field \(a\), set
\[
 \begin{aligned}
 (df)_\mu(x)&=f(x+e_\mu)-f(x),\\
 (da)_{\mu\nu}(x)&=a_\mu(x)+a_\nu(x+e_\mu)
                         -a_\mu(x+e_\nu)-a_\nu(x).
 \end{aligned}
 \tag{CF1}
\]
The same convention, with step \(n\), is written \(d_n\) on a
lattice of spacing \(n\).  Norms are counting \(\ell^2\) norms of
sites, positive bonds or plaquettes, with each unordered direction
pair counted once.  One color suffices for the proof; the flat
\(\mathrm{SU}(2)\) forms are three identical copies.

The Wilson convention
\(A(U)=\sum_p[1-\tfrac12\operatorname{Re}\operatorname{tr}U_p]\)
gives
\(A(e^{ita})=\tfrac12t^2\|da\|^2+O(t^3)\).
For the actual soft gauge,
\[
 \mathrm{GF}(V)=\sum_y\sum_{x\in B(y)\setminus\{y\}}
          [1-\tfrac12\operatorname{Re}\operatorname{tr}V(y,x)].
 \tag{CF2}
\]
Here \(V(y,x)\) is the implicit mean of the transports along the
coordinate-order internal paths.  Its first derivative at the
identity is the arithmetic mean of the corresponding oriented
path sums.  Denote this linear functional by
\(\gamma_{y,x}(a)\), and put \((Ga)_{y,x}=\gamma_{y,x}(a)\).
Consequently
\[
 \mathrm{GF}(e^{ita})=\tfrac12t^2\|Ga\|^2+O(t^3).
 \tag{CF3}
\]
Internal paths contain no crossing pivot.  Thus the gauge form
on nonpivot variables is local to each block and has no pivot
coordinate.  For a gradient,
\(\gamma_{y,x}(df)=f(x)-f(y)\).

The physical step uses the scale \(g=\theta^{-1/2}\) and gauge
weight \(\theta\).  For the real unweighted step the action
coefficient is also \(\theta\).  Hence \(\theta g^2=1\), and
the two quadratic forms above have coefficient one in the
canonical Gaussian variables.  Normalizing constants and the
Haar Jacobian belong to other parts of the step; they do not
change this definition of its free precision.

\subsection*{Plain and straight averaging}

For odd \(n\), let \(B_n(y)\) be the centered block of \(n^4\)
sites.  The raw straight bond average is
\[
 (S_na)_\mu(y)=n^{-4}\sum_{x\in B_n(y)}
                         \sum_{t=0}^{n-1}a_\mu(x+te_\mu).
 \tag{CF4}
\]
At one step define
\(\sigma(a)(y)=L^{-4}\sum_{x\in B_L(y)}\gamma_{y,x}(a)\).
Direct differentiation of the actual outer exponential of closed
logarithms gives the plain derivative
\[
 P_La=S_La-d_L\sigma(a).
 \tag{CF5}
\]
The derivative of its reference straight segment cancels its
negative occurrence in every closed logarithm.  This explains
why no extra reference-segment term is present in (CF5).

Let \(q_Lf\) be the arithmetic site average over \(B_L(y)\),
and \(r_Lf=f|_{\text{block centers}}\).  Telescoping along paths
proves
\[
 S_Ldf=d_Lq_Lf,\qquad
 \sigma(df)=q_Lf-r_Lf,\qquad
 P_Ldf=d_Lr_Lf.
 \tag{CF6}
\]
In particular it is restriction at the centers, not a site
average, that intertwines with the plain constraint.

Write \(P^{[k]}\) for the derivative at the identity of the
plain \(k\)-fold nonlinear average, with raw logarithms at both
ends.  For \(n=L^k\),
\[
 P^{[k]}a=S_na+d_n\psi_k(a)
 \tag{CF7}
\]
for a linear site functional \(\psi_k\).
Here is a direct verification.  Composition of the straight
averages equals \(S_{nL}\): in each coordinate the centered
digits \(nr+u\), with
\(r\in\{-(L-1)/2,\ldots,(L-1)/2\}\) and
\(u\in\{-(n-1)/2,\ldots,(n-1)/2\}\),
tile the larger centered block, and the segment indices
\(nt+j\), \(0\le t<L\), \(0\le j<n\), tile the longer straight
segment.  All weights multiply
to \((nL)^{-4}\).  Applying the next plain step to the gradient
in (CF7) gives another coarse gradient by (CF6).  Its application
to \(S_na\) differs from the next straight average by a gradient,
by (CF5).  This proves (CF7) inductively, starting with (CF5).

\subsection*{Curvature averaging is a contraction in dimension four}

Expanding the boundary of an \(n\) by \(n\) square gives
\[
 (d_nS_na)_{\mu\nu}(y)
    =n^{-4}\sum_{x\in B_n(y)}\sum_{r,s=0}^{n-1}
             (da)_{\mu\nu}(x+re_\mu+se_\nu).
 \tag{CF8}
\]
Every interior bond cancels in this expansion; its four remaining
sides are exactly the four straight averages in \(d_nS_n\).
Call the positive averaging operator on plaquettes on the right
\(T_n\).  Each row has sum \(n^2\).  Each fine plaquette has
total column weight \(n^{-2}\): for each pair \((r,s)\) its
starting point determines exactly one block \(B_n(y)\), since
those blocks partition the fine sites.  There are \(n^2\) such
pairs, each with weight \(n^{-4}\).

For nonnegative weights \(t_{Pp}\), weighted Cauchy gives
\[
 \sum_P\left|\sum_p t_{Pp}F_p\right|^2
 \le\sum_P\left(\sum_pt_{Pp}\right)
                   \left(\sum_pt_{Pp}|F_p|^2\right)
 \le n^2n^{-2}\sum_p|F_p|^2.
 \tag{CF9}
\]
Thus \(\|T_n\|_{2\to2}\le1\).  Counting with multiplicities
makes the same proof valid on a periodic lattice.  The product
of row and column sums in dimension \(d\) would be \(n^{4-d}\);
its cancellation here uses dimension four.
Taking curls of (CF7) therefore proves the exact uniform bound
\[
 \|d_nP^{[k]}w\|\le\|dw\|.
 \tag{CF10}\label{cf:curvature-contraction}
\]
No bound on a background propagator was used.

\subsection*{The classical constrained pullback}

For current bond data \(a\), on the level-\(k\) lattice, define
the classical flat quadratic value
\[
 E_k(a)=\min\{\|dw\|^2:P^{[k]}w=a\}.
 \tag{CF11}
\]
The affine set is nonempty.  For instance place the value of a
coarse bond on every fine bond crossing its block face and put
zero on fine bonds internal to a block of side \(n=L^k\).
The nonlinear group lift with internal bonds \(I\) has exactly
the prescribed average: all parallel segments cross the same
face value and all closed logarithms vanish.  Applying this
observation successively through the centered subdivisions gives
the same assertion at \(k\) levels.  Its derivative is a right
inverse for \(P^{[k]}\).

Existence of the minimum in (CF11) needs no uniform inverse
theorem.  The image under \(d\) of the finite-dimensional affine
constraint set is a closed affine subspace; its point nearest
zero exists.  Equivalently, quotient
\(\ker P^{[k]}\cap\ker d\), restrict to its orthogonal complement,
and solve the positive finite-dimensional normal equations there.
This also shows that the minimum is a quadratic form in \(a\).
Its construction at this stage has no claimed uniform inverse bound.
Equation (CF10), applied to every admissible lift, gives
\[
 E_k(a)\ge\|d_na\|^2.
 \tag{CF12}
\]
We henceforth express the current lattice in units of its own
spacing, so the curl on the right is written \(da\).

This is the classical action part of the actual flat precision.
Indeed the first variation of the bare action at the identity is
zero.  A twice differentiable exact background lift has a tangent \(w\)
with \(P^{[k]}w=a\); its action Hessian is \(\|dw\|^2\), because
its second derivative is multiplied by that zero first variation.
Linearized constrained stationarity makes this tangent minimize
the quadratic form in (CF11), modulo its gauge null space.
The same observation eliminates the second-order one-step
constraint correction at a flat point.  In particular, no
nonzero-current correction has been dropped in establishing
(CF11) at the flat point.  The finite Gaussian form (CF11) is
defined without assuming a nonlinear background theorem; when
that background is used, this argument identifies its flat
quadratic pullback.  Even an arbitrary exact differentiable lift
would suffice for the lower bound (CF12).

Let \(C\) be the actual flat pivot inclusion, with identity on
nonpivot coordinates.  The physical flat free precision is
characterized by
\[
 \langle u,A_k^{\rm flat}u\rangle
       =E_k(Cu)+\|GCu\|^2
       \ge\|dCu\|^2+\|GCu\|^2.
 \tag{CF13}
\]
This formula is for the actual contour gauge and the actual
plain fibre.  It is not obtained by substituting a Landau
projection for \(G\).

\subsection*{A fixed-cell Fourier calculation}

Group the current lattice into blocks of side \(L\).  In one
cell let \(R=\{-(L-1)/2,\ldots,(L-1)/2\}^4\).
A bond is indexed by its starting site \(r\in R\) and direction.
There are \(4L^4\) bond coordinates and four pivot coordinates.
The free space consequently has dimension \(4L^4-4\), independently
of the volume.  For a coarse Bloch parameter
\(p\in(\mathbb R/2\pi\mathbb Z)^4\), extend a cell field by
\[
 a(Lm+r,\mu)=e^{ip\cdot m}a(r,\mu).
 \tag{CF14}
\]
The curl matrix \(d_p\) has entries which are constants or
\(e^{\pm ip_\mu}\).  The internal-contour matrix \(G\) uses
only edges within its cell, hence has no \(p\) dependence.
The plain constraint matrix \(P(p)\), obtained from (CF5), is
also a finite matrix with such phase entries.

The middle pivot of direction \(\mu\) is the crossing bond
from \(he_\mu\) to \((h+1)e_\mu\), \(h=(L-1)/2\).
At the identity its coefficient in the raw constraint is
\(L^{-3}\), since exactly \(L\) of the \(L^4\) parallel
segments contain it.  There is no other pivot in that constraint.
It is the own pivot in the cell of the coarse bond's start,
so this coefficient has no Bloch phase.  Thus the pivot block
of \(P(p)\) is \(L^{-3}I_4\) for every \(p\).
If \(h_0\) inserts \(L^3\) on those pivots and \(\iota\)
inserts free coordinates with zero pivots, then
\[
 C(p)=\iota-h_0P(p)\iota,
 \quad P(p)C(p)=0,\quad C(p)|_{\rm free}=1,
 \quad \operatorname{Ran}C(p)=\ker P(p).
 \tag{CF15}
\]
The last equality follows by solving each pivot equation.
It also proves the rank \(4L^4-4\) for all \(p\).
All entries of \(C(p)\) are continuous Laurent polynomials in
the Bloch phases.  This is the same finite-range pivot map as
in the physical chart, evaluated at the flat field.

Define the Hermitian finite matrix
\[
 B(p)=C(p)^*\big(d_p^*d_p+G^*G\big)C(p).
 \tag{CF16}
\]
We prove that it has no zero vector for every \(p\), including
\(p=0\).  Suppose \(a=C(p)u\) and
\(d_pa=0\), \(Ga=0\).  Extend it by (CF14).
Within a block, two coordinate-order paths to the same site
differ by plaquette swaps.  Zero curl makes their sums equal.
The same swaps express a comb loop with one additional internal
bond as a sum of internal plaquette curls.  Consequently zero
\(G\) forces every internal path sum, and then every internal
bond, to vanish.

Across a face perpendicular to \(e_\mu\), a plaquette with
one transverse step compares two neighboring crossing values.
Its other two bonds are internal to the two blocks, hence zero.
The crossing values are therefore equal.  The transverse face
grid is connected, so all crossing bonds on that face have a
common value \(\beta_{y,\mu}\).  Every straight segment from
\(x\in B(y)\) to \(x+Le_\mu\) then has sum
\(\beta_{y,\mu}\), and all internal contour sums are zero.
The actual constraint gives
\((P_La)(y,\mu)=\beta_{y,\mu}\).
It vanishes by (CF15), killing every crossing bond as well.
Thus \(a=0\), and \(u=0\) by its free restriction.
The argument is valid for every real Bloch parameter, with complex
fiber vectors.

One can also describe the possible closed fields explicitly.
On the cell, the \(L^4\) functions
\(L^{-2}\exp(i\xi_\ell\cdot r)\),
\(\xi_\ell=(p+2\pi\ell)/L\), form an orthonormal basis,
where each component of \(\ell\) runs from \(0\) to \(L-1\).
Orthogonality is the finite geometric-series identity.
Fine differences multiply a mode by
\(t_\mu=e^{i\xi_{\ell,\mu}}-1\).
The closedness equations are
\(t_\mu a_\nu-t_\nu a_\mu=0\).  If one \(t_\mu\ne0\),
they say exactly \(a=t\phi\), a gradient.
The only remaining case is \(p=0,\ell=0\), giving a constant
harmonic vector.  The local argument above kills that case
as well; no zero-frequency mode has been silently omitted.

The entries of \(B(p)\) are continuous on the compact Bloch
torus.  Its least Rayleigh quotient therefore has a positive
uniform minimum \(c(L)>0\).  Explicitly, if the infimum were
zero, a sequence of parameters and unit free vectors would
have a convergent subsequence.  Their limit would be a unit
zero vector for the limiting \(B(p)\), contradicting the
previous paragraph.  Thus
\[
 B(p)\ge c(L)1\qquad\hbox{for every }p.
 \tag{CF17}
\]

On a torus with \(N_\mu\) coarse blocks in each direction,
the allowed parameters are \(p_\mu=2\pi j_\mu/N_\mu\).
The normalized discrete Fourier transform in the block index
is unitary, again by the same geometric-series identity.
It converts the comparison form on the right of (CF13) into
the sum of (CF16) over these parameters.  Parseval and (CF17)
therefore prove
\[
 A_k^{\rm flat}\ge c(L)1
 \tag{CF18}\label{cf:classical-flat-gap}
\]
on the actual free nonpivot coordinates, uniformly in the torus
volume and in \(k\).  Local formulas can first be used on the
admissible large tori; no small-torus identification of different
pivots is needed.  The three color copies have the same constant.

\subsection*{Flat holonomies and the limit of the conclusion}

The same conclusion holds for a real flat background in its
prescribed block axial representative.  Flatness makes internal
transports path independent.  Since their means in this
representative are \(I\), every internal bond is \(I\).
Flat plaquettes make the crossing value constant on each block
face.  The background is therefore the lift of a flat coarse
connection.  A gauge constant on each block flattens it on
the universal cover and preserves the soft contour gauge norm.
Only twisted periodicity remains.

The four torus holonomies commute: their commutator is the
product of the plaquettes of a fundamental rectangle, all equal
to \(I\).  Commuting unitary matrices have a common orthonormal
eigenbasis; choose an eigenspace of any noncentral holonomy and
use its invariance under the others, or any basis if all are
central.  Their adjoint actions are consequently commuting
unitaries on the complexified color space.  Decomposing into
their common eigenspaces gives characters \(e^{i\chi_\mu}\).
The allowed block momenta in that character are
\((2\pi j_\mu+\chi_\mu)/N_\mu\), up to the sign convention
for the boundary twist.  The local matrices and the curvature
averaging proof are unchanged: the boundary phases have modulus
one, so the absolute row and column counts in (CF9) still apply.
Since (CF17) holds for every
real Bloch parameter, all these shifted sets satisfy the same
gap.  No smallness of global holonomies is required.

The result proves the classical flat gap for the actual contour
gauge and plain constraint.  It is a lower bound, and does not
prove a uniform upper bound or a decay estimate for a dressed
precision.  Nonzero-current terms from curved constraints,
signed layer profiles, counterterm arrays, non-Abelian
perturbations and enlarged complex decoupling parameters still
need their own quantitative comparison with this positive
flat form.  No full background or interacting Gaussian-law
assertion follows from (CF18) alone.

% END MODULE chapters/flat_contour_precision_gap_v8.tex


% BEGIN MODULE chapters/flat_whitney_comparison_v8.tex
\section{A constructive upper bound for the flat classical action}
\label{ch:flat-whitney-comparison}

This section concerns the same unweighted classical Gaussian value
as the flat-gap calculation.  Let \(n=L^k\), with \(L\) odd, and
let \(P^{[k]}\) be the raw derivative of the actual plain iterated
average.  For a coarse bond field \(a\), put
\[
 E_k(a)=\min\{\|dw\|^2:P^{[k]}w=a\}.
 \tag{WI1}
\]
We construct an exact lift with energy at most
\(256\|d_na\|^2\), uniformly in \(k\) and in the volume.
No bound on a non-Abelian background response is used.

\subsection*{Discrete interpolation on the actual centered grids}

Fine sites are integers and coarse sites are \(nm\),
\(m\in\mathbb Z^4\).  Write
\[
 \varphi_n(t)=\max\{1-|t|/n,0\},\qquad
 \beta_n(t)=\mathbf1_{\{0\le t<n\}},\quad t\in\mathbb Z.
 \tag{WI2}
\]
Periodize these functions on a finite torus.  The identities
\(\sum_m\varphi_n(x-nm)=1\),
\(\sum_m\beta_n(x-nm)=1\), and
\(\sum_{x\in\mathbb Z}\varphi_n(x)=n\) follow by summing the
two linear pieces.  Also
\[
 \varphi_n(t+1)-\varphi_n(t)
       =n^{-1}\big[\beta_n(t+n)-\beta_n(t)\big].
 \tag{WI3}
\]
Define interpolation of sites, bonds and plaquettes by
\[
 \begin{aligned}
 (I_0f)(x)&=\sum_m f(m)\prod_\lambda\varphi_n(x_\lambda-nm_\lambda),\\
 (I_1a)_\mu(x)&=n^{-1}\sum_m a_\mu(m)
       \beta_n(x_\mu-nm_\mu)
       \prod_{\lambda\ne\mu}\varphi_n(x_\lambda-nm_\lambda),\\
 (I_2F)_{\mu\nu}(x)&=n^{-2}\sum_m F_{\mu\nu}(m)
       \beta_n(x_\mu-nm_\mu)\beta_n(x_\nu-nm_\nu)\\
   &\qquad\times\prod_{\lambda\notin\{\mu,\nu\}}
                          \varphi_n(x_\lambda-nm_\lambda).
 \end{aligned}\tag{WI4}
\]
The coarse arrays in these formulas are indexed by their starting
node \(m\).  Their values are raw bond or plaquette values; no
length factor has been included in the coarse difference \(d_n\).
Equation (WI3), followed by reindexing \(m\), proves
\[
 dI_0=I_1d_n,\qquad dI_1=I_2d_n.
 \tag{WI5}
\]
For example, differentiating a transverse nodal hat in \(I_1\)
replaces it by the difference of two interval indicators divided
by \(n\); their coefficients are precisely the two terms of
the corresponding coarse plaquette boundary.  This proves the
second identity with the same orientation as
\((da)_{\mu\nu}=a_\mu+\nabla_\mu a_\nu-a_\mu(\cdot+e_\nu)\).

For a fixed plaquette direction, \(I_2\) has nonnegative weights.
Their row sum is \(n^{-2}\), by the two partition identities.
Their column sum is \(n^2\): each interval contributes \(n\)
fine positions, each of the remaining two nodal hats has sum
\(n\), and the prefactor is \(n^{-2}\).
Weighted Cauchy therefore gives
\[
 \|I_2F\|^2\le n^{-2}n^2\|F\|^2=\|F\|^2.
 \tag{WI6}
\]
These counting norms are those of the four-dimensional Wilson
quadratic energy.  As for curvature averaging, dimension four
is used in the cancellation of the powers of \(n\).

\subsection*{The exact averaging symbol}

The raw straight average is
\[
 (S_nw)_\mu(y)=n^{-4}\sum_{x\in B_n(y)}
\sum_{j=0}^{n-1}w_\mu(x+je_\mu),
\]
with centered blocks.
We compute \(S_nI_1\) exactly, rather than using an approximate
right inverse.  Put \(h=(n-1)/2\).  In a transverse direction
the coefficient of the neighboring coarse index \(m\) is
\(n^{-1}\sum_{x=-h}^h\varphi_n(x-nm)\).
In the component direction it is
\[
 n^{-2}\#\{(x,j):-h\le x\le h,\ 0\le j<n,
\ nm\le x+j<n(m+1)\}.
\]
Both expressions are zero unless \(m=0,\pm1\), and both have
the same three values:
\[
 a_0=1-\frac{h(h+1)}{n^2}=\frac34+\frac1{4n^2},
 \qquad
 a_1=\frac{h(h+1)}{2n^2}=\frac{1-n^{-2}}8
       \quad(m=\pm1).
 \tag{WI7}
\]
For the longitudinal count, the negative and positive overshoots
each number \(1+\cdots+h=h(h+1)/2\); the remaining
\(n^2-h(h+1)\) pairs give the central coefficient.
For the transverse hat, its positive-side and negative-side
linear sums give the same two triangular numbers.

Let \(T_\lambda\) translate a coarse array by one node in
direction \(\lambda\).  The tensor product in (WI4) and the
identical one-dimensional coefficients just computed show
\[
 S_nI_1=\mathcal A_n,
 \qquad \mathcal A_n=\prod_{\lambda=1}^4
       \big[a_0I+a_1(T_\lambda+T_\lambda^{-1})\big].
 \tag{WI8}
\]
Crucially, this is the same scalar convolution on all four
bond components.  Each one-dimensional symbol is
\[
 a_n(t)=\frac34+\frac1{4n^2}
                   +\frac{1-n^{-2}}4\cos t,
 \qquad \frac12\le a_n(t)\le1.
 \tag{WI9}
\]
Fourier transform on the coarse torus, or the quadratic-form
bound \(a_0-2a_1\ge1/2\) for a unitary shift, proves
\[
 \|\mathcal A_n^{-1}\|_{2\to2}\le16.
 \tag{WI10}
\]
All shifts commute.  Their scalar convolution commutes with
every coarse difference, on sites, bonds and plaquettes.
Its inverse has the same property.

\subsection*{An exact lift for the plain constraint}

For prescribed coarse \(a\), take
\(w_0=I_1\mathcal A_n^{-1}a\).
Equations (WI5), (WI6), (WI8) and (WI10) imply
\[
 S_nw_0=a,\qquad
 \|dw_0\|\le\|d_n\mathcal A_n^{-1}a\|
                       \le16\|d_na\|.
 \tag{WI11}
\]
This is an exact straight lift.  The actual constraint is plain.
For completeness its two relevant identities are
\[
 P^{[k]}w=S_nw+d_n\psi_k(w),\qquad
 P^{[k]}df=d_n(r_nf),
 \tag{WI12}
\]
where \(r_n\) restricts a fine site function to the centers
\(n\mathbb Z^4\).  At one step these follow by differentiating
the actual closed-logarithm formula: \(P_L=S_L-d_L\sigma\),
and on gradients an internal contour is
\(f(x)-f(y)\), so \(P_Ldf=d_Lr_Lf\).
Iterating proves the second identity.  Composition of straight
averages is \(S_n\); applying a plain average to a gradient
again gives a gradient, proving the first by induction.

Nodal interpolation satisfies \(r_nI_0f=f\).  Set
\[
 f=-\psi_k(w_0),\qquad w=w_0+dI_0f.
 \tag{WI13}
\]
Then (WI11)--(WI12) give \(P^{[k]}w=a\) exactly, while
\(dw=dw_0\).  The gradient correction has no classical
curvature cost.  The energy estimate itself needs no amplitude
bound; such a bound is proved below.  We have proved
\[
 E_k(a)\le256\|d_na\|^2.
 \tag{WI14}
\]
The complementary lower bound follows from the exact curvature
average: \(d_nP^{[k]}w=T_ndw\), where \(T_n\) has row sum
\(n^2\) and total column weight \(n^{-2}\).  The same weighted
Cauchy argument gives \(\|T_n\|\le1\).  Thus
\[
 \boxed{\ \|d_na\|^2\le E_k(a)\le256\|d_na\|^2\ }.
 \tag{WI15}\label{wi:energy-comparison}
\]
All statements include \(n=1\), where the interpolations and
\(\mathcal A_1\) are identities.  No zero-frequency inverse
is taken: \(\mathcal A_n\) itself is strictly positive even
on constants.

\subsection*{Consequences and scope}

For the actual flat free nonpivot precision,
\(\langle u,A_k^{\rm flat}u\rangle
=E_k(Cu)+\|GCu\|^2\).
The local pivot basis has counting \(\ell^2\) norm at most
\(CL^4\), from its row and column bounds.  Also
\(\|d\|^2\le24\) on one-forms in four dimensions.
The internal contour matrix has row sum at most \(CL\)
and column sum at most \(CL^4\), hence \(\|G\|^2\le CL^5\).
It contains no pivot, so \(GC=G\iota\), where \(\iota\)
inserts free coordinates.  Equations (WI14) and these bounds
give the conservative uniform upper bound
\[
 A_k^{\rm flat}\le CL^8 I.
 \tag{WI16}
\]
Together with the separate flat-gap theorem this bounds the
same classical precision above and below independently of
the number of finer levels and the volume.

\subsection*{Supremum norms of the constructed linear lift}

The same scalar inverse is bounded in the supremum norm.
For one direction the off-diagonal part of
\(a_0I+a_1(T+T^{-1})\), divided by \(a_0\), has supremum
operator norm at most \(2a_1/a_0<1\).  Its inverse Neumann
series therefore has norm at most
\((a_0-2a_1)^{-1}\le2\).  Multiplying the four commuting
direction inverses gives
\[
 \|\mathcal A_n^{-1}\|_{\infty\to\infty}\le16,
 \qquad \|w_0\|_\infty\le\frac{16}{n}\|a\|_\infty .
 \tag{WI17}
\]
The second inequality uses the row sum \(1/n\) of \(I_1\).

There is an explicit choice of the scalar in (WI12).
Let \(\psi_0=0\).  When the current spacing is \(L^j\),
apply \(P_L=S_L-d_L\sigma_L\) to
\(P^{[j]}w_0=S_{L^j}w_0+d_{L^j}\psi_j\).
The next plain step sends the last gradient to the gradient
of the center restriction.  Thus
\[
 \psi_{j+1}=r_L\psi_j-\sigma_L(S_{L^j}w_0).
 \tag{WI18}
\]
All fields in this recurrence are raw bond or site values
on their current lattices.  There is no missing spacing
factor.  Internal paths have at most \(2(L-1)\) steps,
so \(\|\sigma_L\|_{\infty\to\infty}\le2(L-1)\).
Also \(\|S_{L^j}\|_{\infty\to\infty}=L^j\), whereas
restriction has norm one.  The geometric sum yields
\(\|\psi_k\|_\infty
 \le2(L-1)\sum_{j<k}L^j\|w_0\|_\infty
 =2(n-1)\|w_0\|_\infty\le32\|a\|_\infty\).

Since \(dI_0f=I_1d_nf\), the correcting gradient in (WI13)
has norm at most \(2\|f\|_\infty/n\).  Consequently the
particular exact lift constructed above obeys
\[
 \|dI_0f\|_\infty\le\frac{64}{n}\|a\|_\infty,
 \qquad \|w\|_\infty\le\frac{80}{n}\|a\|_\infty .
 \tag{WI19}\label{wi:linear-lift-sup}
\]
The scalar convolution and its inverse commute with curl.
The row sum \(n^{-2}\) of \(I_2\) and (WI17) also give
\[
 \|dw\|_\infty=\|dw_0\|_\infty
              \le\frac{16}{n^2}\|d_na\|_\infty.
 \tag{WI20}
\]
These are norms of one explicitly defined linear right
inverse.  They are not estimates for a minimizing response.

The interpolation identities also hold with commuting unitary
boundary twists.  Work on the universal cover and periodize
with the corresponding character; absolute row and column
sums are unchanged.  Twisted shifts remain unitary and
commuting, so (WI9)--(WI10) retain their lower bound for every
shifted Bloch parameter.  Thus the same energy comparison
applies to real flat holonomies in the block axial gauge.

The construction supplies energy and supremum estimates for the
specified linear lift of the unweighted classical flat form.
The exponential flat-precision kernels are supplied by
Section~\ref{ch:flat-exponential-kernels}. Applications to a minimizing
or nonlinear background, complex cores and tower substitutions,
signed layer profiles, nonzero currents, counterterms and complex
decoupling arrays use the corresponding actual-operator comparison
propositions.

% END MODULE chapters/flat_whitney_comparison_v8.tex


% BEGIN MODULE chapters/flat_exponential_kernels_v8.tex
\section{Exponential kernels of the classical flat precision}
\label{ch:flat-exponential-kernels}

We compute the classical flat quadratic form of the actual plain
iterated average, and then the covariance on its next-step nonpivot
fiber with the actual soft contour gauge.  A spectral gap by itself
does not imply kernel decay.  The proof below supplies a uniform
complex strip for the exact symbols before using the flat gap.
All fields are in raw bond coordinates, indexed by their starting
sites.  No midpoint phase is suppressed.

Let \(n=L^k\), with \(L\) odd.  The current lattice is in units
of its own spacing; its finer lattice has \(n\) subdivisions per
direction.  We first use one scalar color.  At the identity the
three \(\mathrm{SU}(2)\) colors give identical copies.  The
normalization of the action is
\(A(e^{ita})=\tfrac12t^2\|da\|^2+O(t^3)\).
After the canonical scaling \(g=\theta^{-1/2}\), the real
unweighted step has \(\theta g^2=1\).  Thus the matrices calculated
here are precisely its classical quadratic precision matrices.

\subsection*{The two constrained minima are equal}

Write \(P^{[k]}\) for the raw derivative of the actual plain
\(k\)-fold average.  The straight average is
\[
 (S_nw)_\mu(ny)=n^{-4}\sum_{r\in B_n(0)}
                         \sum_{j=0}^{n-1}w_\mu(ny+r+je_\mu).
 \tag{FX1}
\]
Here \(B_n(0)=\{-(n-1)/2,\ldots,(n-1)/2\}^4\).
The actual one-step map is that of
Section~\ref{ch:block-average-pivot-chart}.  Let
\(\mathcal G_L(y,x)\) be its coordinate-order shortest-path
multiset from \(y\) to \(x\in B_L(y)\), with its prescribed
multiplicities and equal weights.  In positive bond coordinates set
\[
 \begin{aligned}
 \gamma_{y,x}(a)&=\frac1{\#\mathcal G_L(y,x)}
       \sum_{\Gamma\in\mathcal G_L(y,x)}
                \sum_{b\in\Gamma}\varepsilon_\Gamma(b)a(b),\\
 (Ga)_{y,x}&=\gamma_{y,x}(a)\quad(x\ne y),\qquad
 \sigma(a)(y)=L^{-4}\sum_{x\in B_L(y)}\gamma_{y,x}(a).
 \end{aligned}\tag{FX1a}
\]
Here the signs record the path orientation and
\(\gamma_{y,y}=0\).  Differentiation gives
\(P_L=S_L-d_L\sigma\).  On gradients the plain map intertwines
with restriction at the centers, whereas the straight map
intertwines with the site average.  Iteration gives
\[
 P^{[k]}w=S_nw+d_n\psi_k(w),\qquad
 P^{[k]}df=d_n r_nf,\qquad S_ndf=d_nq_nf.
 \tag{FX2}
\]
Both site maps are onto: for \(r_n\) assign the values at the
centers, and for \(q_n\) use a function constant on each centered
block.  These are statements about finite linear cochains; no
regularity bound for those corrections is required.

Suppose \(P^{[k]}w=a\).  Then \(S_nw=a-d_n\psi_k(w)\).
Choose \(q_nf=\psi_k(w)\).  The field \(w+df\) has straight
average \(a\) and the same curl.  Conversely, for \(S_nw=a\),
choose \(r_nf=-\psi_k(w)\) to correct the plain average.
Consequently the actual classical quadratic value is
\[
 E_n(a):=\min_{P^{[k]}w=a}\|dw\|^2
          =\min_{S_nw=a}\|dw\|^2.
 \tag{FX3}
\]
Nonemptiness also follows from the filter calculation below:
its central alias is nonzero in every component at \(p\ne0\),
and a constant fine field lifts a constant coarse field.
The two gradient corrections transfer feasibility between the
constraints.  Finite dimensional projection after quotienting
the curl-null directions then gives both minima.
The equality does not change the actual soft gauge, which is
added later on the current one-step fiber.

\subsection*{Fourier normalization and the auxiliary covariance}

Take a coarse torus with \(\mathsf s\) sites per direction and
a fine torus with \(n\mathsf s\).  Use the unitary transforms
\[
 \widehat w(\xi)=(n\mathsf s)^{-2}\sum_xe^{-i\xi\cdot x}w(x),
 \qquad
 \widehat a(p)=\mathsf s^{-2}\sum_ye^{-ip\cdot y}a(ny).
 \tag{FX4}
\]
Fine aliases of a coarse momentum \(p\) are
\(\xi_m=(p+2\pi m)/n\), with
\(m\in I_n^4\), \(I_n=\{-(n-1)/2,\ldots,(n-1)/2\}\).
Put
\[
 b_n(z)=\frac{\sin(nz/2)}{n\sin(z/2)},\qquad
 B_n(\xi)=\prod_\nu b_n(\xi_\nu),\qquad h=(n-1)/2.
 \tag{FX5}
\]
The centered block sum is \(B_n\), and the straight segment
sum is \(ne^{ih\xi_\mu}b_n(\xi_\mu)\).  Thus the exact filter
of component \(\mu\) in (FX1) is
\[
 s_{\mu,n}(\xi)=ne^{ih\xi_\mu}B_n(\xi)b_n(\xi_\mu),
 \qquad
 \widehat a_\mu(p)=n^{-2}\sum_{m\in I_n^4}
                      s_{\mu,n}(\xi_m)\widehat w_\mu(\xi_m).
 \tag{FX6}
\]
The factor \(n^{-2}\) is the square root of the ratio of the
coarse and fine numbers of sites.  In particular its square,
not an unnormalized alias sum, enters a covariance.

The auxiliary fine Feynman form is
\(\|dw\|^2+\|d^*w\|^2\).  Its symbol is component-diagonal,
with
\(\lambda(\xi)=4\sum_\nu\sin^2(\xi_\nu/2)\).
For \(p\ne0\), every alias has positive \(\lambda\).
The coarse covariance of this auxiliary form is diagonal:
\[
 D_{\mu,n}(p)=n^{-4}\sum_m
                  \frac{|s_{\mu,n}(\xi_m)|^2}{\lambda(\xi_m)}.
 \tag{FX7}
\]
This formula can equally be obtained by minimizing the positive
fine Feynman form under its linear constraint; no probability
measure on the fine constant mode is needed.

For complex continuation define
\[
 \begin{aligned}
 q_{\nu,m}(p)&=
 \frac{\sin((p_\nu+2\pi m_\nu)/2)}
      {n\sin((p_\nu+2\pi m_\nu)/(2n))},\\
 f_{\mu,m}(p)&=\left(\prod_\nu q_{\nu,m}(p)^2\right)
                                              q_{\mu,m}(p)^2,\\
 s_{n,m}(p)&=4n^2\sum_\nu
                    \sin^2((p_\nu+2\pi m_\nu)/(2n)).
 \end{aligned}\tag{FX8}
\]
Removable ratios are filled in at zero.  On real momenta the
phase \(e^{ih\xi_\mu}\) in (FX6) cancels with its conjugate.
Its analytic version is the product of the filters at \(\xi\)
and \(-\xi\), giving exactly
\[
 D_{\mu,n}(p)=\sum_{m\in I_n^4}
                           \frac{f_{\mu,m}(p)}{s_{n,m}(p)}.
 \tag{FX9}\label{fx:alias-covariance}
\]
Indeed \(s_{\mu,n}(\xi)s_{\mu,n}(-\xi)
=n^2B_n(\xi)^2b_n(\xi_\mu)^2\), and
\(s_{n,m}=n^2\lambda(\xi_m)\).  This checks every power of
\(n\).  Squared absolute values in (FX7) are not used as
holomorphic functions.

\subsection*{Recovering the gauge-invariant energy}

For real \(p\), let
\(t_\mu^+(p)=e^{ip_\mu}-1\) and
\(t_\mu^-(p)=e^{-ip_\mu}-1\).
The latter is the conjugate of the former only on real momenta.
In analytic formulas below \((t^-)^{\mathsf T}\) is a holomorphic
row, not a conjugation of a complex variable.

At nonzero fine frequency, the elementary Hodge identity is
\(\|d\widehat w\|^2+\|d^*\widehat w\|^2
=\lambda\|\widehat w\|^2\).
Subtracting a fine gradient makes a field divergence free
without changing its curl.  Its straight average changes by
a coarse gradient, by (FX2).  Conversely any coarse gradient
can be corrected by a fine gradient, since \(q_n\) is onto.
It follows that the curl minimum in (FX3) equals the minimum
of the Feynman quadratic value over coarse gauge changes:
\[
 E_n(p;a)=\min_{z\in\mathbb C}
            (a+t^+z)^*D_n(p)^{-1}(a+t^+z).
 \tag{FX10}
\]
For clarity, the two inequalities use different simple steps.
A divergence-free representative of a curl-minimizing field
has Feynman energy equal to its curl energy, with coarse value
\(a+t^+z\).  In the other direction a Feynman minimizer at
that coarse value can be gradient-corrected back to \(a\),
and its curl energy is no greater than its Feynman energy.

Solving the scalar minimum gives the precision symbol
\[
 \mathcal E_n=D_n^{-1}
   -\frac{D_n^{-1}t^+(t^-)^{\mathsf T}D_n^{-1}}
                 {(t^-)^{\mathsf T}D_n^{-1}t^+}.
 \tag{FX11}
\]
Here \(D_n=\operatorname{diag}(D_{\mu,n})\).
At the real zero momentum the classical value is zero on every
constant coarse one-form: its constant fine lift is \(a/n\).
We accordingly set \(\mathcal E_n(0)=0\).  This zero mode is
retained, not removed by a mass term.  We next prove that this
definition is the value of a uniformly holomorphic symbol.

\subsection*{Uniform noncentral alias bounds}

Fix a small \(\rho>0\) and allow
\(|\operatorname{Re}p_\nu|<\pi+\rho\),
\(|\operatorname{Im}p_\nu|<\rho\).
For \(m_\nu\ne0\),
\[
 |n\sin((p_\nu+2\pi m_\nu)/(2n))|
                      \ge c(1+|m_\nu|).
 \tag{FX12}
\]
To see this, the real part of the sine argument stays in a
fixed subinterval of \((-\pi,\pi)\), and its absolute value
is at least \(c|m_\nu|/n\).  On that interval
\(|\sin x|\ge c|x|\), and
\(|\sin(x+iy)|^2=\sin^2x+\sinh^2y\).
The numerator in (FX8) is bounded.  For \(m_\nu=0\) the ratio
is a quotient of two sinc functions on a fixed compact
zero-free neighborhood.  Hence, uniformly in \(n\),
\[
 |q_{\nu,m}(p)|\le\frac C{1+|m_\nu|}.
 \tag{FX13}
\]

For \(m\ne0\) the denominator in (FX8) has a uniform positive
real part.  In fact
\(\operatorname{Re}\sin^2(x+iy)
\ge\sin^2x-\sinh^2y\).
After multiplication by \(4n^2\), the imaginary displacement
costs at most \(C\rho^2\), whereas the real sum is at least
\(c|m|^2\).  Decrease \(\rho\) to obtain
\[
 \operatorname{Re}s_{n,m}(p)\ge c(1+|m|^2)
                                \qquad(m\ne0).
 \tag{FX14}
\]
The constants do not grow as the alias set grows.

Write \(s_n=s_{n,0}\), \(q_\nu=q_{\nu,0}\),
\(Q_n=\prod_\nu q_\nu^2\), and \(f_\mu=Q_nq_\mu^2\).
Then
\[
 D_{\mu,n}=\frac{f_\mu}{s_n}+R_{\mu,n},\qquad
 R_{\mu,n}=\sum_{m\ne0}\frac{f_{\mu,m}}{s_{n,m}}.
 \tag{FX15}
\]
Each summand of \(R\) is bounded by the absolutely summable
majorant
\[
 \frac{C}{1+|m|^2}
          \prod_{\nu=1}^4(1+|m_\nu|)^{-2}.
 \tag{FX16}
\]
The extra factor \((1+|m_\mu|)^{-2}\) may be discarded.
Thus \(R\) is holomorphic and uniformly bounded on the stated
box, independently of \(n\).  Cauchy on a smaller box bounds
all its fixed-order derivatives uniformly as well.  For \(n=1\)
the sum is empty.  Every noncentral summand vanishes at \(p=0\),
so \(R_{\mu,n}(0)=0\); only boundedness is needed below.

\subsection*{Cancellation of the apparent gauge denominator}

Near \(p=0\), uniformly in \(n\),
\(q_\nu=1+O(|p|^2)\), \(s_n=O(|p|^2)\), and
\(f_\mu=1+O(|p|^2)\).  Define
\[
 h_\mu=(f_\mu+s_nR_{\mu,n})^{-1},\qquad
 H_n=\operatorname{diag}(h_\mu).
 \tag{FX17}
\]
On a sufficiently small fixed complex polydisc these are
holomorphic and uniformly bounded.  The exact filter identity
\[
 \sum_\mu\frac{t_\mu^+t_\mu^-}{f_\mu}
       =\frac{4n^2\sum_\mu\sin^2(p_\mu/(2n))}{Q_n}
       =\frac{s_n}{Q_n}
 \tag{FX18}
\]
is essential.  It follows directly by canceling the numerator
\(\sin^2(p_\mu/2)\) against \(q_\mu^2\), including removable
values.  It is not an approximation to a gauge denominator.

Since \(h_\mu=f_\mu^{-1}-s_nR_{\mu,n}h_\mu/f_\mu\),
equation (FX18) gives
\[
 \begin{aligned}
 F_n&:=\sum_\mu t_\mu^+t_\mu^-h_\mu=s_nJ_n,\\
 J_n&=Q_n^{-1}
       -\sum_\mu\frac{t_\mu^+t_\mu^-R_{\mu,n}h_\mu}{f_\mu}.
 \end{aligned}\tag{FX19}
\]
Here \(J_n(0)=1\) and \(|J_n-1|\le C|p|^2\) uniformly.
It is therefore nonzero on a smaller fixed polydisc.
On its meromorphic domain \(D_n^{-1}=s_nH_n\), and the
denominator in (FX11) equals \(s_nF_n=s_n^2J_n\).
The actual precision consequently has the holomorphic form
\[
 \mathcal E_n(p)=s_nH_n
       -\frac{H_nt^+(t^-)^{\mathsf T}H_n}{J_n}.
 \tag{FX20}\label{fx:origin-cancellation}
\]
This removes the whole complex zero cone of \(s_n\), not merely
the value at the real origin.  The expression is uniformly
bounded there, is \(O(|p|^2)\), and gives
\(\mathcal E_n(0)=0\).  It also obeys the exact Ward identity
\(\mathcal E_nt^+=0\), by (FX19).
For \(n=1\), it reduces to
\(\lambda(p)I-t^+(t^-)^{\mathsf T}\), the ordinary curl
quadratic form.  This checks the normalization and phases.

\subsection*{A uniform strip and its Fourier kernels}

We supply the argument away from zero rather than infer a
complex domain from real positivity alone.  On the real
Brillouin cube the central \(q_\nu\) are positive and bounded
above and below uniformly in \(n\).  Also
\(c|p|^2\le s_n(p)\le C|p|^2\).
For real \(|p|\ge r>0\), the central term in (FX15) and the
nonnegative alias terms give
\[
 c_D\le D_{\mu,n}(p)\le C_r,
 \qquad
 (t^-)^{\mathsf T}D_n^{-1}t^+\ge c_r>0.
 \tag{FX21}
\]
The second inequality uses
\(\sum_\mu|t_\mu^+|^2\ge c|p|^2\) and the upper bound on \(D\).

On a sufficiently thin complex neighborhood of this real
compact set, \(\operatorname{Re}s_n\ge cr^2-C|\operatorname{Im}p|^2\).
Thus the central term is holomorphic and bounded there;
(FX12)--(FX16) do the same for the aliases.  Cauchy bounds
the first derivatives uniformly in \(n\).  Shrinking the
neighborhood by a fixed amount depending on \(r\), (FX21)
therefore keeps every \(D_\mu\) and its scalar gauge denominator
nonzero.  Formula (FX11) is holomorphic and uniformly bounded
on that neighborhood.  This is a quantitative perturbation
of uniformly separated real denominators.

Choose \(r\) inside the origin polydisc of (FX20), then reduce
the common width once.  The two formulas agree wherever their
denominators are nonzero, by their exact algebraic derivation,
and hence on their overlaps by continuity.  Together they give
constants \(\eta>0,M_E<\infty\), independent of \(n\), such that
\[
 \mathcal E_n\text{ is holomorphic and }
 \sup_{|\operatorname{Im}p_\nu|<\eta}\|\mathcal E_n(p)\|\le M_E.
 \tag{FX22}
\]
The symbol is \(2\pi\)-periodic.  In (FX9), changing a coarse
momentum by \(2\pi\) reindexes the aliases; an endpoint alias
changes its fine frequency by \(2\pi\), under which the actual
filter and denominator are periodic.  The same is true of
(FX11), and its analytic continuation has that periodicity.

Define its infinite-lattice kernel by
\[
 \mathcal K_n(z)=(2\pi)^{-4}\int_{[-\pi,\pi]^4}
                         e^{ip\cdot z}\mathcal E_n(p)\,dp.
 \tag{FX23}
\]
For \(0<\eta'<\eta\), shift the \(\nu\)-th contour to
\(\operatorname{Im}p_\nu=\eta'\operatorname{sgn}z_\nu\).
Periodicity cancels the vertical faces.  The exponential then
gives
\[
 \|\mathcal K_n(z)\|\le M_Ee^{-\eta'|z|_1}.
 \tag{FX24}
\]
All diagonal terms are retained in this Fourier definition.
The kernel and its weighted absolute row sums are bounded
uniformly in \(k\).  Its zero-frequency value and Ward identity
are those already proved; no contact term is added afterward.

\subsection*{The actual next-step precision and its inverse}

Now group the current lattice into the fixed blocks of side
\(L\) used by the next RG step.  For block Bloch momentum
\(P\), let \(r\in B_L(0)\), \(j\in I_L^4\), and
\(\zeta_j=(P+2\pi j)/L\).  The finite cell transform and its
analytic inverse are
\[
 U(P)_{rj}=L^{-2}e^{i\zeta_j\cdot r},\qquad
 V(P)_{jr}=L^{-2}e^{-i\zeta_j\cdot r}.
 \tag{FX25}
\]
Finite geometric-series orthogonality proves \(VU=UV=1\).
For real \(P\), \(V=U^*\).  For complex \(P\) the displayed
inverse, not literal conjugation, is used.  Their norms are
bounded on a fixed strip, since their extra diagonal factors
are \(e^{\pm iP\cdot r/L}\) and \(|r_\nu|/L<1/2\).
The folded classical precision is
\[
 \mathcal E_n^{\rm cell}(P)
     =U(P)\operatorname{diag}_j(\mathcal E_n(\zeta_j))V(P),
 \tag{FX26}
\]
with the identity on component indices understood.  It is
periodic and holomorphic on a fixed block strip, uniformly
in \(n\).  Periodicity again follows by the alias permutation.

Let \(C(P)\) be the actual flat middle-pivot basis for the
plain one-step constraint.  Its free restriction is the
identity and its pivot block solves the coefficient
\(L^{-3}I_4\); hence its entries are fixed Laurent polynomials
in the block phases.  Let \(G\) be the internal contour
matrix defined in (FX1a), with rows \((y,x)\), \(x\ne y\).
Define
\(C^\sharp(P)=C(-P)^{\mathsf T}\); on real \(P\) this is \(C(P)^*\).
The exact classical free nonpivot precision is
\[
 \mathsf A_n(P)=C^\sharp(P)
        [\mathcal E_n^{\rm cell}(P)+G^{\mathsf T}G]C(P).
 \tag{FX27}
\]
This is the action pullback plus the actual soft gauge.  The
Feynman form was used only to compute the gauge-invariant
minimum in (FX3); it has not replaced the soft gauge in (FX27).

For clarity the uniform real gap used next can be obtained
independently of this symbol calculation.  Expanding the curl
of (FX1) gives an average of \(n\) by \(n\) fine plaquette
sheets.  Its nonnegative row weights sum to \(n^2\) and its
total column weights to \(n^{-2}\).  Weighted Cauchy gives
\(\|d_nS_nw\|\le\|dw\|\), hence
\(E_n(a)\ge\|da\|^2\) in current units.
Consequently (FX27) dominates the fixed-cell form
\(C(P)^*[d_P^*d_P+G^{\mathsf T}G]C(P)\).
This form has no zero vector: zero curl makes internal path
sums independent of the path, so zero \(G\) kills every
internal bond.  A plaquette crossing a block face then makes
its crossing values equal tangentially.  The plain constraint
kills that common face value.  The free restriction of \(C\)
is the identity, so the original free vector is zero.
The matrices in this comparison are continuous on the compact
block Bloch torus, with fixed dimension and fixed rank.  Their
least Rayleigh quotient on the unit sphere therefore has a
positive minimum \(c(L)\).  This proves
\(\mathsf A_n(P)\ge c(L)1\) uniformly in \(n\), including \(P=0\).
Equation (FX22), the finite transform and the fixed polynomial
matrices in (FX27) give a uniform bound for \(\mathsf A_n\)
on a complex strip.  Cauchy on a smaller strip bounds its
first derivatives by a constant \(M_A(L)\).
Choose its width so that
\[
 \|\mathsf A_n(P+iY)-\mathsf A_n(P)\|
          \le4M_A(L)|Y|_\infty<c(L)/2.
 \tag{FX28}
\]
The Neumann series around the real positive matrix then proves
\[
 \mathsf A_n^{-1}\text{ is holomorphic},\qquad
          \|\mathsf A_n(P+iY)^{-1}\|\le2/c(L)
 \tag{FX29}
\]
on one strip independent of \(n\).  This step uses both the
proved analytic bounds and the gap.

Contour shifting as in (FX23)--(FX24) now bounds the Fourier
coefficients of the block precision and its inverse by
\(C(L)e^{-c(L)|m|_1}\).  The cell has fixed size, so converting
its indices back to current nonpivot bonds only changes the
constants.  We obtain
\[
 |\mathsf A_k^{\rm flat}(b,b')|
       +|\mathsf C_k^{\rm flat}(b,b')|
       \le C_L e^{-\mu_L d(b,b')},
 \qquad \mathsf C_k^{\rm flat}=(\mathsf A_k^{\rm flat})^{-1}.
 \tag{FX30}\label{fx:flat-kernel-bounds}
\]
Here \(d(b,b')\) is bond distance in current lattice units;
using starting-site distance or midpoint distance only changes
the fixed constants.
The symbols \(\mathsf C_k^{\rm flat}\) in this conclusion are
covariances; the undecorated pivot map \(C(P)\) is a different
rectangular matrix.  Exponential weighted row and column
sums follow by summing the fixed-dimensional lattice shells.

\subsection*{Finite tori, holonomies, and scope}

For a periodic lattice the exact quadratic minimization is
diagonal on its allowed Fourier momenta.  Its symbols are
the samples calculated above, with no volume correction.
If \(K_n(m)\) is any infinite block kernel just obtained and
the block period is \(\mathsf s_{\rm b}\), discrete Fourier
orthogonality and absolute convergence give
\[
 K_{n,\mathsf s_{\rm b}}(m)
          =\sum_{v\in\mathbb Z^4}K_n(m+\mathsf s_{\rm b}v).
 \tag{FX31}
\]
The image sum preserves an exponential bound in periodic
distance, uniformly over the admissible periods.  It identifies
the infinite-lattice kernel determined by the same symbol,
with exponentially small image errors on fixed nonwrapping
neighborhoods.  It does not assert exact equality of a full
torus kernel and an infinite kernel on a lift.  The same
argument applies to the inverse symbol, since its sampled
matrix is the inverse of each sampled positive precision.

For real flat holonomies in the prescribed block axial
representative, a unitary gauge constant on each block
flattens the connection on the universal cover and preserves
the contour gauge norm.  The commuting adjoint holonomies
split into unitary characters.  Each character shifts the
allowed block momenta by a real amount; the infinite symbols
are unchanged.  Equivalently the image sum (FX31) receives
unit-modulus factors.  All absolute kernel bounds remain
valid, without requiring small global holonomies.

The conclusion is the exponential precision and covariance
bound for the actual unweighted classical flat quadratic
problem.  The auxiliary \(D_n\) itself is massless and is not
asserted to have an exponential kernel.  The gap and decay
in (FX30) concern the next-step constrained nonpivot variables
with the soft gauge, not a physical Yang--Mills mass gap.
No signed profile, counterterm, nonzero-current, non-Abelian
or enlarged complex-decoupling comparison has been proved
by this calculation.  Those require their actual operators
and uniform perturbation estimates.

% END MODULE chapters/flat_exponential_kernels_v8.tex


% BEGIN MODULE chapters/03_background.tex
\section{Constrained backgrounds and their regularity}
\label{candidate:background}
\candidatebanner{A}

\subsection*{(1b) The background field: the regular critical orbit with prescribed averages}

A renormalisation step expands about the background field. This field is the regular critical point of the Wilson action among bare fields whose block averages are prescribed. At each level the regions of large field are left behind (Step 2 of the outline). The averages are therefore prescribed on a lattice whose scale varies: on the $j$-th layer of a decreasing sequence of regions they are prescribed at level $j$, and every constant is required to be independent of the regions and of the number of levels.

\textbf{(i) The setting and the objects involved.}

\emph{Sequence, scales, distance.} The \emph{depth} of a sequence is the index $k\le K$ of its last member. Every depth in this part is the depth of a sequence: the depth $K-k$ of a level (Notation \ensuremath{\S}1) does not occur here. We work in level-$k$ units: $\eta=L^{-k}$, and the spacing of $\mathbb T^{(j)}$ is $\xi_j:=L^{j-k}$. Take regions $\Omega_0\supset\dots\supset\Omega_k$, with $\Omega_{k+1}=\emptyset$ and each $\Omega_j$ a union of cubes of side $\mathfrak M\xi_j$. They are said to be \emph{separated} when

\begin{gather*}
\FitMath{\operatorname{dist}_\infty(\Omega_j^c,\Omega_{j+1})>\rho\,\xi_j .}
\end{gather*}

Both the \emph{cube size} $\mathfrak M$ and the \emph{gap} $\rho=R_1\mathfrak M$ are large powers of $L$; here $R_1$ is the gap number of the Notation. The cube size is the quantity $M_1$, which gives the size of the walks both in (1c) and in the Notation. The cube size $\mathfrak M$ and the side $M$ of the localisation cubes, which is fixed at a later stage, are different quantities, and so are the gap $\rho$ and the densities $\rho_k$. A member may be $T_\ell$ or empty, and no condition is imposed on the shape or the connectivity of the members. The case $\Omega_j\equiv T_\ell$ is the problem on a single lattice. By the \emph{seams} we mean the boundaries $\partial\Omega_{j+1}$ that separate consecutive layers. The \emph{tops} are the sites $y$ of $\mathbb T^{(j)}$ that lie in the \emph{layer} $\Omega_j\setminus\Omega_{j+1}$, and for such a site we set $\xi_y:=\xi_j$. The blocks $B^j(y)$ of the tops are called \emph{pillars}, and they form a partition of $\Omega_0$. Each site $x$ of a pillar other than its top has a \emph{level} $n(x)<j$, given by $x\in\mathbb T^{(n)}\setminus\mathbb T^{(n+1)}$. It also has a \emph{parent} $z(x)\in\mathbb T^{(n(x)+1)}$, characterised by $x\in B(z(x))$, and its \emph{ascent} $x,z(x),z(z(x)),\dots$ rises to the top. For a bare bond or plaquette $\mathfrak e$, $\mathsf L(\mathfrak e)$ is the largest $\xi_j$ attached to its corners; the bond or plaquette is called \emph{inner} when all its corners lie in $\Omega_0$. Each layer is measured in its own unit, by the distance

\begin{gather*}
\FitMath{d_{\mathbf B}(x,x'):=\inf_{\Gamma\ni x,x'}\Big[\sum_{j\ge1}\xi_j^{-1}\operatorname{length}\big(\Gamma\cap(\Omega_j\setminus\Omega_{j+1})\big)+\eta^{-1}\operatorname{length}(\Gamma\setminus\Omega_1)\Big]}
\end{gather*}

(the infimum runs over connected paths $\Gamma$ in the torus, and the first member of the sequence does not enter; the index $\mathbf B$ refers to the determining set of the sequence, defined below).

\emph{Variational problem.} The set $\Gamma_j$ consists of those bonds of $\mathbb T^{(j)}$ in $\Omega_0$ that have an end point in $\Omega_j$ and no end point in $\Omega_{j+1}$. Such a bond is either \emph{interior}, joining two tops of level $j$, or \emph{outward}. For an outward bond, the end $z\notin\Omega_j$, called the \emph{outer end}, is a top of level $j-1$, and every bond lying within $B(z)$ belongs to $\Gamma_{j-1}$ as one of its interior bonds. The \emph{nodes} are the end points of the bonds of the \emph{determining set} $\mathbf B=\bigcup_j\Gamma_j$, and they are the tops. No constraint is placed on the \emph{inward} bonds, which lead from the layer $j$ into $\Omega_{j+1}$; the variant in which some of them are kept is treated in the next paragraph. In terms of these sets we define

\begin{gather*}
\FitMath{M_{\mathbf B}(U)=(\mathcal M^jU|_{\Gamma_j})_j,} \\[3pt]
\FitMath{C(\mathbf B,\mathbf V)=\{U\text{ on }\Omega_0:M_{\mathbf B}(U)=\mathbf V\},} \\[3pt]
\FitMath{A_{\Omega_0}=\sum_{p\subset\Omega_0}a_p,} \\[3pt]
\FitMath{\mathcal G^{\rm res}=\{u:u=\mathbf 1\text{ at the nodes}\}.}
\end{gather*}

A configuration $U$ is called \emph{critical} if $A_{\Omega_0}$ is stationary along the curves in $C(\mathbf B,\mathbf V)$ that pass through $U$. It is called \emph{regular at the radius $\varepsilon$} if, in the local units, its curvature and its current are of size $\varepsilon$:

\begin{gather*}
\FitMath{|\partial U(p)-\mathbf 1|<\varepsilon\,\eta^2\mathsf L(p)^{-2},} \\[3pt]
\FitMath{|J(U)(b)|<\varepsilon\,\mathsf L(b)^{-3}.}
\end{gather*}

A datum $\mathbf V$ is called \emph{$\varepsilon_1$-regular} if the values $\partial\mathbf V(\tau)$ at its \emph{test loops} $\tau$ (also called \emph{tests}) lie within $\varepsilon_1$ of $\mathbf 1$. These values are the plaquette products $\partial\mathbf V(p')$ of $V_j$, in which $\mathcal M(V_{j-1})$ is used on those bonds of $p'$ that lie in the finer layer.

\emph{Inward bonds.} If an arbitrary set of inward bonds is kept as additional constraints, Theorems 1, 2, 3 and the first half of Theorem 4 remain valid, though Theorem 4 then holds with a rate specific to that case. The problem without inward bonds is the setting of the complex pairs of Theorem 4 and of Theorem 5. It is also the setting of Theorem 6, although its restriction identity also holds when they are kept. Every inward bond of level $j$ that is kept has as its end point a \emph{seam node} $x\in\Omega_{j+1}$, that is, a node that is not a top. In that case the tests also contain two-scale loops at the seams, namely the \emph{pinch holonomies} $\mathsf f_{\mathbf V}(b)\mathsf f_{\mathbf V}(\tilde b)^{-1}$, where $b,\tilde b$ is a pair of kept inward bonds of level $j$ ending at a common site $x\in\Omega_{j+1}$, such that the plaquette they span again has its fourth corner in $\Omega_{j+1}$. The \emph{frame} $\mathsf f_{\mathbf V}(b)=V_{j+1}(c_b)\,V_j(y_b,b_-)\,V_j(b)$ transports from $y$, the centre of the block containing $x$, to $x$ itself, along a route that stays outside $\Omega_{j+1}$. In it, $y_b$ denotes the centre of the block of $b_-$, $c_b=\langle y,y_b\rangle$, and $V_j(y_b,b_-)$ is the averaged contour variable (defined in the paragraph \emph{Chart}). Ba{\l}aban's ``bonds of $\Lambda_j$'', where his $\Lambda_j$ is the site set $\Lambda_j(\mathbf B)$ of the Notation, \ensuremath{\S}4, admit two readings, which differ by the inward bonds. They are either the bonds having at least one end point in that site set, or the bond set of $\Omega_j$ minus the bond set of $\Omega_{j+1}$. In Commun. Math. Phys. 96 (1984) 223, (2.3), the difference is explicitly taken on bonds, and with this reading no inward bond occurs at $j\ge1$. At level $0$, the set $\Lambda_0=\Omega_1^c$ of that paper and the set $\Gamma_0=\Omega_1^c$ of Commun. Math. Phys. 119 (1988) 243, (2.2), contain the inward bonds. The set $\Omega_0^{(0)}\setminus\Omega_1^{(0)}$ of Commun. Math. Phys. 99 (1985) 75, (1.5), understood as a difference, does not contain them. By convention, the present construction omits the inward bonds at every level. The reason for this convention is an incompatibility that arises at a single level. Take a sequence of depth $k+1$ with $\emptyset\ne\Omega_{k+1}\ne T_\ell$ on which the inward bonds of level $k$ are kept. Then for some regular data, every element of the critical orbit (Theorem 1) has a $k$-th average that fails to be in block axial gauge somewhere on $\mathbb T^{(k+1)}\cap\Omega_{k+1}$. Hence imposing constraints on these bonds is incompatible with imposing the gauge by which Theorem 6 defines $V^{\rm bg}_k$. There are two remedies, and both of them are used. A step discards these bonds. The strip of Step 2 of the outline, in which nodes lie inside the blocks for a different reason, retains its set and imposes the axial gauge except at the nodes.

\emph{Chart.} Fix a real $U_0$ with $\mathbf V_0=M_{\mathbf B}(U_0)$ and curvature radius $\alpha_0$, meaning that $|\partial U_0(p)-\mathbf 1|\le\alpha_0\eta^2\mathsf L(p)^{-2}$. About it we write $U=e^{i\eta A}U_0$ and $F(A)=A_{\Omega_0}(e^{i\eta A}U_0)$. The gradient is $F'(A)=F'(0)+\Delta^{\rm W}_{U_0}A+\mathcal V(A)$, where $\Delta^{\rm W}_{U_0}=D^*D+\Delta'$. The operator $\Delta'$ is local and of the size of the curvature of $U_0$, and $\mathcal V$ is of second order. Moreover

\begin{gather*}
\FitMath{\mathcal M^j(e^{i\eta A}U_0)(c)=v_j(c_-)\,e^{i\mathcal Q^{\rm st}(A)(c)}\,\mathbf V_0(c)\,v_j(c_+)^{-1},} \\[3pt]
\FitMath{K_A:=(\mathcal Q^{\rm st})'(A),}
\end{gather*}

where $\mathcal Q^{\rm st}=\mathcal Q^{\rm st}_j(U_0;\cdot)$ on $\Gamma_j$ is the \emph{straightened average}. The site factors $v_j=v_j[A]=v_{j-1}\mathsf v_{j-1}$ are products of one-step \emph{block factors} $\mathsf v_m(y)$, each of which is the mean over a block of the change in the contour variables of $\mathcal M^mU_0$. For $x\in B(y)$, the \emph{averaged contour variable} $W(y,x)$ is the mean of the transports of $W$ along the contour family of $(y,x)$; it is the variable that enters the definition of $\mathcal M$ (Notation, \ensuremath{\S}3). The \emph{comb} $\Gamma_{y,x}$ is the member of this family that goes through the coordinate directions one after another, always in the same order. To each single bond $c$ of $\mathbf B$ belongs a functional of $A$, its \emph{row}. When $c\in\Gamma_j$, it depends only on the values of $A$ in the \emph{double block} $B^j(c_-)\cup B^j(c_+)$. The superscript st abbreviates ``straightened''; the \emph{straight} linear average $Q^{\rm st}_k(U)$ of the Notation, \ensuremath{\S}3, carries the same superscript but denotes a different operator. The Notation's plain map $\mathcal Q_j(U_0;A)=\frac1i\log[\mathcal M^j(e^{i\eta A}U_0)\mathbf V_0^{-1}]$ is expressed in level-$j$ units, and the site factors account for the difference between it and $\mathcal Q^{\rm st}$. Its derivative $\xi_jQ_j(U_0)$, in which the length $\xi_j$ of the coarse bond appears, differs from $K_0$ by a gauge term on the coarse lattice. This term does not become small with the curvature of $U_0$. The Landau gauge is the condition $RD^*A=0$, where $R$ is the orthogonal projection onto $\Delta_{U_0}\{\lambda:Q'\lambda=0\}$. In it, $(Q'\lambda)(y)=L^{-4j}\sum_{x\in B^j(y)}\operatorname{Ad}(\mathsf T(y,x))\lambda(x)$ is the \emph{pillar average} at a top $y$ of level $j$. The transport $\mathsf T(y,x)$, a function of two sites, is obtained by multiplying, down the pillar from $y$ to $x$, the averaged contour variables of the successive $\mathcal M^nU_0$; the superscript ${}^{\mathsf T}$, by contrast, marks a transpose. In the norms, each derivative carries one power of the local scale. The first norm is $\|A\|_X=\sup\max\{\mathsf L|A|,\mathsf L^2|\nabla_{U_0}A|\}$, the norm of $1$-forms in the gauge slice, also written $\|A\|_1$ in (1a); the subscript $X$ in it has nothing to do with the domain $X$ of the paragraph \emph{Complex fields}. The second is $|f|_{(-s)}=\sup\mathsf L^s|f|$. For complex fields every pairing is bilinear and every adjoint is taken to be a transpose ${}^{\mathsf T}$, so that each identity is complex-linear.

\emph{Ba{\l}aban's scheme.} Ba{\l}aban's own average (Commun. Math. Phys. 98 (1985) 17) takes the lowest corner of each block as its base site and uses one comb contour for each site. Moved onto the centred blocks by a translation through half a block, it has linear rows $\mathsf Q$ on the bonds of $\mathbf B$ and pillar averages $\mathsf Q'$, the latter formed from comb transports that issue from the block corners. On $\Gamma_j$, once the length $\xi_j$ of the coarse bond is divided out, $\xi_j^{-1}\mathsf Q$ is the operator written $Q^{\rm B}_j(U_0)$ in the Notation, \ensuremath{\S}3. The projection $R^{\mathsf c}$ is built in the same way as $R$, with $\mathsf Q'$ replacing $Q'$, and the condition $R^{\mathsf c}D^*A=0$ defines the \emph{comb slice}. Ba{\l}aban's propagators are built from these operators.

\emph{Reduced datum.} An outer end is a node shared by two levels. For small $B:\mathbf B\to\mathfrak g^c$ we set $\mathsf b(B)=B$ at the interior bonds ($\mathsf b(B)$ and the coefficient $\mathsf b[\Phi]$ of the Notation, \ensuremath{\S}6, are different objects, distinguished by their arguments). At an outward bond $c\in\Gamma_j$ whose outer end is $z=c_+$ we put the following (the formula for $z=c_-$ is the mirror image):

\begin{gather*}
\FitMath{e^{i\mathsf b(B)(c)}=e^{iB(c)}\,\operatorname{Ad}(\mathbf V_0(c))\,\mathsf v_z(B),}
\end{gather*}

where $\mathsf v_z(B)$ is the one-step block factor at $z$. It is built from $\mathcal M^{j-1}U_0$ and from the values of $B$ on the bonds inside the block $B(z)$. The argument of the map $B\mapsto\mathsf b(B)$ is a family on $\mathbf B$; the coefficient $\mathsf b(f)$ of (1e), like $\mathsf b[\Phi]$ of the Notation, \ensuremath{\S}6, has an argument of another kind. The map is holomorphic and local. It is triangular, because $\mathsf v_z$ depends only on interior bonds, and hence invertible. It satisfies $\sup|\mathsf b(B)|\le C\sup|B|$.

\emph{Complex fields.} On a domain $X$ at level $j$, a \emph{triple} $(U,A',\mathbf J)$ consists of a real $U$, the complex configuration $\mathbf U=e^{i\eta_jA'}U$ built from it, and an independent current $\mathbf J$. A \emph{tower} of $Z\subset X$ is the separated sequence obtained by eroding $Z$ level by level according to one of two rules. The first uses cubes of side $\mathfrak M$ in the unit of their own level, with a gap exceeding $\rho$ in the unit of the outer member. The second uses cubes of side $M_1'$ and requires a gap no smaller than $M_1'$, measured in the unit of the inner member. The \emph{tower configuration} $U_n(Z;\mathbf U)$ of the tower is obtained by cutting the tower off at depth $n$, taking the background of the truncated tower, whose determining set is $\mathbf B'$, and continuing it analytically to $M_{\mathbf B'}(\mathbf U)$, which is a complex datum. Explicitly, it is $\Phi_{M_{\mathbf B'}(U)}(B)$ of Theorem 3 evaluated at $B=\frac1i\log[M_{\mathbf B'}(\mathbf U)M_{\mathbf B'}(U)^{-1}]$. By this definition the tower configuration is a function of the decomposition $(U,A')$; the definition does not make it a function of the product $\mathbf U$ alone. The \emph{core} $\mathcal U^\circ_j(X;\alpha_0,\alpha_1)$ is defined by two sets of requirements. The first are the \emph{simple conditions}:
\begin{itemize}
\item $|\partial U-\mathbf 1|,|\partial\mathbf U-\mathbf 1|<\alpha_0\eta_j^2$;
\item on cubes of side $O(LM)$ there exist real gauges in which $|A|,|\nabla A|<CM\alpha_0$;
\item $|A'|,|\nabla_UA'|<\alpha_1$;
\item $|\mathbf J|<\alpha_0$.
\end{itemize}

The second are the \emph{tower clauses}. Every tower configuration exists, for $Z\subset X$ of every level $p\le j$, for every $n\le p$ and for both rules. Let $Z$ have two layers of $M$-cubes removed (the cubes of $\pi_{p-1}$, Notation \ensuremath{\S}1; $1\le n\le p$). On that set, $|\partial U_n(Z;\mathbf U)-\mathbf 1|<\alpha_0L^{-(j-p)}\eta_p^2$ and $|J(U_n(Z;\mathbf U))|<\alpha_0L^{-(j-p)}L^{2(n-p)}$, with the current formed in level-$n$ units. The same holds with $U$ in place of $\mathbf U$. A term extends to the union of the complex gauge orbits of the pairs $(\mathbf U,\mathbf J)$ of the core to the extent that it takes one value at any two pairs of the core lying on one orbit. This holds at two radii (1d). Complex gauge orbits do not saturate the core itself. Of the domain of analyticity $U^c_j(X,\alpha_0,\alpha_1)$ of the Notation, the core is the hereditary form: the regularity that this domain demands of the backgrounds of the averages of $\mathbf U$, the core demands on every $Z\subset X$ of every level $p\le j$. On a domain meeting several layers, this space takes the form of the \emph{layered space} $\tilde U^c_j$, whose radii are assigned layer by layer. Radii written $(\alpha_0',\alpha_1')$ are radii of the simple conditions only, and the $\alpha_{i,j}$ are the running radii of the Notation, \ensuremath{\S}7.

In the sub-parts some objects carry other symbols as well. In 1b.1, 1b.2, 1b.4 and 1b.6 the depth is denoted $k$; in 1b.3, 1b.5 and 1b.7--1b.9 it is denoted $n$. The gap is denoted $\mathsf g$ in 1b.1 and $\rho_{\rm g}$ in 1b.4; in 1b.6, 1b.7 and 1b.9 it is denoted $R_1\mathfrak M$, $\rho_\Omega$ and $\rho_{\rm s}$ respectively. In 1b.4 and 1b.6 the letter $\rho$ without decoration denotes a length on which the thresholds depend, and $L^2\rho$ does not exceed the gap. The straightened average keeps the symbol $\mathcal Q^{\rm st}$ in 1b.1--1b.4. In 1b.5 it is denoted $a(A)$, in 1b.7 $\mathcal Q$, and in 1b.6 $\mathcal Q^\natural$. In 1b.3 and 1b.4 the same symbol also denotes the family obtained by adjoining the rows of the kept inward bonds. In 1b.2 the reduced datum $\mathsf b(B)$ is denoted $\tilde B$; it is denoted $B^\natural[B]$ in 1b.4, and $B^\natural$, the row datum, in 1b.6. In 1b.4, however, the words ``reduced datum'' refer to $\mathbf V^\flat$: this is what remains of a datum with inward bonds after the gauge group at the seam nodes has been divided out. In 1b.6 and 1b.7 the symbol $\mathscr H$ stands for $\mathcal H$, the Landau potential of Theorem 3. In 1b.3--1b.6 the projection $R^{\mathsf c}$ is denoted $\mathsf R$ or $R$. In 1b.6, $A$ stands for the potential $A^{\mathsf c}$ of the comb slice. In 1b.8 and 1b.9 the radii of a core are denoted $b_0,b_1$.

\textbf{(ii) The principal statements.} The constants $C$ and the rates depend on $L$ alone, and in H{\"o}lder members also on $\alpha_{\mathrm H}$. There are two exceptions. The first is $K^\sharp$ and the Lipschitz constant of (b) (Theorem 5); the dependence of this Lipschitz constant is on $L$ together with the radii $\bar\varepsilon$, $r_{\rm B}$ from Theorems 1 and 3. The second is the constant and the threshold of Theorem 6 (c), and their dependence is shown in that theorem. Thresholds and radii depend in addition on the lengths chosen before them (end of (iii)), and also on $M$ where this is shown. Above three floors on the torus, there is no dependence on $K$, $\ell$, $k$, on the regions, on the number of layers or on the datum. The first floor requires at least $O(\rho)$ sites of $\mathbb T^{(k)}$ in every direction. It ensures that the double blocks beneath the rows, and likewise the cubes on which gauges are chosen in the proofs, are shorter than half the torus. The first half of Theorem 2 and Theorems 1, 3, 4, 6 (b), (c) require no other floor than this one. The second floor is that each cube on which a statement chooses a gauge stays shorter than half the torus after it is enlarged by $O(\rho)$ of its units. In the second half of Theorem 2 these cubes have side $m$; in Theorems 5 and 6 (d) they have side $O(LM)$, measured at the level of the space. The third floor is that a tower or a localisation is formed only down to depths $p$ at which $M_1'$ spacings of $\mathbb T^{(p)}$ do not exceed one half of the side length of the torus. For the tori $\ell\in\mathfrak L(g)$ introduced in the Notation, these depths are all $p\le K$.

\textbf{Theorem 1 (one regular critical orbit, linear bound).} Let $\mathbf V$ be real and $\varepsilon_1$-regular, with $\varepsilon_1\le\bar\varepsilon$. The critical configurations of $C(\mathbf B,\mathbf V)$ that are regular at a fixed small radius then form precisely one orbit of $\mathcal G^{\rm res}$, and this orbit is non-empty. It is the background $U_k(\mathbf V)$, also written $U_{\mathbf B}(\mathbf V)$ when more than one sequence is involved, and

\begin{gather*}
\FitMath{|\partial U_k(\mathbf V)(p)-\mathbf 1|<C\varepsilon_1\eta^2\mathsf L(p)^{-2},} \\[3pt]
\FitMath{|J(U_k(\mathbf V))(b)|<C\varepsilon_1\mathsf L(b)^{-3}.}
\end{gather*}

Exactly one element $U^\circ(\mathbf V)$ of the orbit lies in the \emph{axial gauge outside the nodes}. This gauge is defined as follows. Let $x$ be a bare site that is not a node, let $n$ be its level ($x\in\mathbb T^{(n)}\setminus\mathbb T^{(n+1)}$), and let $z\in\mathbb T^{(n+1)}$ be the centre of its block; then $(\mathcal M^nU)(\Gamma_{z,x})=\mathbf 1$. Hence, all the way down each pillar, the comb transports of the successive averages equal $\mathbf 1$. When $\mathbf V=\mathbf 1$ the orbit is $\{\mathbf 1^u\}$. Conversely, if $U$ is any configuration whose plaquettes are regular at a small radius $\varepsilon\le cL^{-2}$, then the tests of $M_{\mathbf B}(U)$ are within $2\varepsilon$ of $\mathbf 1$ in the interior of a layer and within $CL^2\varepsilon$ across a seam.

\emph{Remark (the gauge with the contour means).} Replacing the comb transport by the averaged contour variable $(\mathcal M^nU)(z,x)$ gives a second gauge of the same kind, and each orbit again has exactly one element in it. The proof uses only three properties of the variable:
\begin{itemize}
\item it depends on the bonds inside $B(z)$;
\item it is covariant, $W^h(z,x)=h(z)W(z,x)h(x)^{-1}$;
\item $W(z,z)=\mathbf 1$ and $\mathbf 1(z,x)=\mathbf 1$.
\end{itemize}

The two representatives are related by an element of $\mathcal G^{\rm res}$, and both equal $\mathbf 1$ over $\mathbf V=\mathbf 1$. In 1b.1 the gauges are developed with the means. Theorems 3 and 5 concern the comb representative, and 1b.8 uses only properties that the two representatives share. In both cases the block axial gauge of a single level, through which $V^{\rm bg}_k$ is defined in Theorem 6, is formed with the mean.

\textbf{Theorem 2 (local regularity).} Let $U$ be critical and regular at a fixed small radius. Then at each $\mathfrak e$ with corner $x$,

\begin{gather*}
\FitMath{\eta^{-2}\mathsf L^2|\partial U-\mathbf 1|,\,\,\mathsf L^3|J(U)|\,\le\,C\max_\tau e^{-\delta\,d_{\mathbf B}(x,\tau)}\,|\partial\mathbf V(\tau)-\mathbf 1| ,}
\end{gather*}

and consequently, for thresholds $\varepsilon^{(q)}$ given layer by layer, the same quantities are $\le C\sup_q\varepsilon^{(q)}e^{-\delta d_{\mathbf B}(x,\text{layer }q)}$. Now let $\square$ be a cube of side $m\xi_j$, where $m$ is a multiple of $\rho$. Assume that its enlargement $\tilde\square$ by $O(\rho)\xi_j$ is contained in $\Omega_j$ and is shorter than half the torus. Take \emph{any} $U$ that is regular at the radius $\varepsilon$ on $\tilde\square$, where $m\varepsilon\le c$. Then near $\square$ there is a gauge $U^u=e^{i\eta A}$ in which, on $\square$,

\begin{gather*}
\FitMath{\xi_j|A|,\,\xi_j^2|\nabla A|,\,\xi_j^3|\partial^*\partial A|,\,\xi_j^3|\Delta A|\le Cm\varepsilon,} \\[3pt]
\FitMath{|\nabla A(x)-\nabla A(x')|\le Cm\varepsilon\,\xi_j^{-2-\alpha_{\mathrm H}}|x-x'|^{\alpha_{\mathrm H}}\,\,(\alpha_{\mathrm H}<1),}
\end{gather*}

where the last $C$ depends on $\alpha_{\mathrm H}$ as well and comes with no bound as $\alpha_{\mathrm H}\to1$. The differences are taken with respect to the configuration $\mathbf 1$, and $\Delta=\sum_\nu\nabla_\nu^{\mathsf T}\nabla_\nu$ acts componentwise on $A$.

\textbf{Theorem 3 (continuation in the datum; Landau pair).} Let $\mathbf V_0$ be real and $\varepsilon_1$-regular. Let $\varepsilon_1$ and $r_{\rm B}$ be small enough that $e^{iB}\mathbf V_0$ is $\bar\varepsilon$-regular for every real $B$ with $\sup|B|<r_{\rm B}$. Put $U_0=U^\circ(\mathbf V_0)$ and $\mathcal B=\{B:\mathbf B\to\mathfrak g^c,\ \sup|B|<r_{\rm B}\}$.

There is exactly one holomorphic map $\Phi=\Phi_{\mathbf V_0}$ on $\mathcal B$ that agrees with $U^\circ(e^{iB}\mathbf V_0)$ for real $B$. Throughout $\mathcal B$ its complex averages are defined and $M_{\mathbf B}\Phi(B)=e^{iB}\mathbf V_0$, and $\Phi(B)$ is regular at the radius $C(\varepsilon_1+\sup|B|)$.

Assume in addition that $\bar\varepsilon\le c(L)/\mathfrak M$; this is a ceiling on the radius, chosen after $\mathfrak M$. Then exactly one holomorphic pair $(\mathcal H,u)$ on $\mathcal B$ takes the value $(0,\mathbf 1)$ at $B=0$ and satisfies

\begin{gather*}
\FitMath{\Phi(B)=\big(e^{i\eta\mathcal H(B)}U_0\big)^{u_B},} \\[3pt]
\FitMath{RD^{\mathsf T}\mathcal H(B)=0,} \\[3pt]
\FitMath{u_B\,v_j[\mathcal H(B)]=\mathbf 1\,\text{at the tops of level }j.}
\end{gather*}

For real $B$ this pair is real, and there is a holomorphic multiplier $\omega(B)$ on $\mathbf B$ such that

\begin{gather*}
\FitMath{F'(\mathcal H(B))=K^{\mathsf T}_{\mathcal H(B)}\,\omega(B),} \\[3pt]
\FitMath{\mathcal Q^{\rm st}(\mathcal H(B))=\mathsf b(B),}
\end{gather*}


\begin{gather*}
\FitMath{\|\mathcal H(B)\|_X,\,|D^{\mathsf T}D\mathcal H(B)|_{(-3)}\le C\sup|B|,} \\[3pt]
\FitMath{|u_B-\mathbf 1|\le C\sup|B|\,\text{at the tops}.}
\end{gather*}

Let $v$ be $G$-valued on the nodes, and let $\tilde v$ be its extension that is constant on pillars (when inward bonds are kept, the extension is constant along the ascent of each site as far as the first node). Then $\Phi_{\mathbf V_0^v}(\operatorname{Ad}(v)B)=\Phi_{\mathbf V_0}(B)^{\tilde v}$.

Every one of the continued maps is real for the conjugation ${}^{\mathfrak c}$ of the Notation, \ensuremath{\S}11. This conjugation is $A^{\mathfrak c}=A^*$ on $\mathfrak g^c$ and $g^{\mathfrak c}=(g^*)^{-1}$ on $G^c$, applied value by value. Thus $\Phi(B^{\mathfrak c})=\Phi(B)^{\mathfrak c}$, $\mathcal H(B^{\mathfrak c})=\mathcal H(B)^{\mathfrak c}$ and $u_{B^{\mathfrak c}}=(u_B)^{\mathfrak c}$. At $B=B^{\mathfrak c}$ this is reality at real data.

\emph{Remark (the comb slice).} The comb slice carries a second holomorphic pair $(A^{\mathsf c},h)$ on $\mathcal B$. It equals $(0,\mathbf 1)$ at $B=0$, is real in the same sense, and satisfies

\begin{gather*}
\FitMath{\Phi(B)=\big(e^{i\eta A^{\mathsf c}(B)}U_0\big)^{h(B)},} \\[3pt]
\FitMath{R^{\mathsf c}D^{\mathsf T}A^{\mathsf c}(B)=0,} \\[3pt]
\FitMath{h(B)\,v_j[A^{\mathsf c}(B)]=\mathbf 1\,\text{at the tops of level }j.}
\end{gather*}

This pair satisfies the two equations $F'=K^{\mathsf T}\omega$, $\mathcal Q^{\rm st}=\mathsf b(B)$ as well, and its potential obeys the same bounds. In addition, $\sup_x|h(B)(x)-\mathbf 1|\le C\sup|B|$ at \emph{all} bare sites.

$R^{\mathsf c}$ is the projection that belongs with Ba{\l}aban's propagators, and $\Phi$ is constructed by means of this pair (1b.6). For $R$, one has $K_0D\lambda=0$ whenever the pillar averages of $\lambda$ vanish, and $(\mathcal H,u)$ is constructed on the basis of $\Phi$ (1b.7). The two potentials are gauge transforms of each other: $e^{i\eta A^{\mathsf c}(B)}U_0=(e^{i\eta\mathcal H(B)}U_0)^{w_B}$, where $w_B=h(B)^{-1}u_B$ is holomorphic. The two pairs are equal exactly when $RD^{\mathsf T}A^{\mathsf c}(B)=0$ near $0$. This happens whenever the restriction of $U_0$ to each pillar is a pure gauge, and so in particular at $\mathbf V_0=\mathbf 1$. Uniqueness as a pair concerns $(\mathcal H,u)$, while the bound at all bare sites concerns $h$.

\textbf{Theorem 4 (exponential decay of the response).} Let $\bar\varepsilon\le c(L)/\mathfrak M$, with $\bar\varepsilon$ chosen after $\mathfrak M$. Put $\mathsf h=\partial_s\mathcal H(B+se_c)|_{s=0}$, where $e\in\mathfrak g^c$ sits at the bond $c$ and is placed at its top $y_c$. Let $\delta J$ be the same derivative of $J(e^{i\eta\mathcal H}U_0)$. Then on the whole polydisc, over the top $y$,

\begin{gather*}
\FitMath{\xi_y|\mathsf h|,\,\xi_y^2|\nabla\mathsf h|,\,\xi_y^3|D^{\mathsf T}D\mathsf h|,\,\xi_y^3|\Delta_{U_0}\mathsf h|,\,\xi_y^3|\delta J|\,\le\,Ce^{-c\delta_0d_{\mathbf B}(y,y_c)}|e| ,}
\end{gather*}

together with a H{\"o}lder member $\xi_y^{-2-\alpha_{\mathrm H}}$ for $\nabla\mathsf h$. For these variables the bounds used are the value estimates $|\dot\omega(c')|\le C e^{-c\delta_0d_{\mathbf B}(y_{c'},y_c)}|e|$ and $|\dot u(y)|\le C e^{-c\delta_0d_{\mathbf B}(y,y_c)}|e|$, respectively. Here $\dot\omega(c')=\partial_s\omega(B+se_c)(c')|_{s=0}$ for a determining bond $c'$, with assigned top $y_{c'}$, and $\dot u(y)=\partial_su_{B+se_c}(y)|_{s=0}$ for a top $y$; both are measured in the unweighted colour or matrix value norm. No bare $1$-form difference or Laplacian is applied to these coarse-row or node variables. Here $\Delta_{U_0}=\sum_\nu\nabla_\nu^{\mathsf T}\nabla_\nu$ denotes the covariant Laplacian acting on $1$-forms. The member containing this Laplacian holds at every bond whose neighbours all belong to $\Omega_0$, and the member containing the current holds at every bond where the current is defined.

At a regular complex pair $(\mathbf U,\mathbf J)$, critical or not, the corresponding system has a single small solution, and this solution is analytic in the pair. Its response satisfies the bounds on $\mathsf h$, $\nabla\mathsf h$ and $\dot\omega$, and the H{\"o}lder member over distances $\le\xi_y$. The second-order members are then $(D^{\mathsf T}D+DRD^{\mathsf T})\mathsf h$, $D(1-R)D^{\mathsf T}\mathsf h$ and $\Delta_{\mathbf U}\mathsf h$, because away from criticality the slice condition is not available.

At $\mathbf V_0=\mathbf 1$ on a single lattice, $\mathsf H_k=\partial_B\mathcal H(0)$ is the periodisation of a $\mathbb Z^4$-invariant kernel. At the lifts of two bonds lying within half the torus of each other, it differs from that kernel by $Ce^{-c\delta_0\ell_k/2}$, with $\ell_k$ denoting the side length of the torus measured in level-$k$ units.

\textbf{Theorem 5 (spaces of complex fields).} (a) Suppose that the simple conditions of level $p$ hold on $Z$ at small radii $(\alpha_0',\alpha_1')$. Then the tower configurations exist and depend holomorphically on $A'$; well inside $Z$ they satisfy

\begin{gather*}
\FitMath{|\partial U_n(Z;\mathbf U)-\mathbf 1|<K^\sharp(\alpha_0'+\alpha_1'^2)\,\eta_p^2,} \\[3pt]
\FitMath{|J(U_n(Z;\mathbf U))|<K^\sharp(\alpha_0'+\alpha_1'^2)\,L^{2(n-p)}\,\,(\text{level-}n\text{ units}).}
\end{gather*}

(Here $K^\sharp$ depends on $L$ and $\rho$.) Consequently, once $K^\sharp(\alpha_0'+\alpha_1'^2)\le\alpha_0$, a triple of level $j$ that satisfies the simple conditions belongs to the core $\mathcal U^\circ_j(X;\alpha_0,\alpha_1)$, $\alpha_1\ge\alpha_1'$. The weight $L^{-(j-p)}$ in the clauses comes from the change of unit. The triple $(\mathbf 1,0,0)$ belongs to every core.

(b) Suppose that the tower lies, member by member, in an ambient separated sequence, which may be deeper. Let $l\ge\xi'$ be the local scales of the ambient sequence and the tower, and assume that $\eta^{-2}l^2|\partial U-\mathbf 1|$, $l|A'|$, $l^2|\nabla_UA'|$ stay below a threshold $a_{\rm T}$. Then the tower configuration exists, and it depends holomorphically on $A'$. Fix $U$, and at an \emph{element} at $x$ take the tower configuration's curvature and current. Such an element is a bare plaquette, or a bond together with its six plaquettes. Take these quantities in the scales $\xi'$ and multiply them by $(l/\xi')^2$, that is, use the factors $\eta^{-2}l^2$ and $l^2\xi'$. At the \emph{settled} elements, those at which the tower has reached the greatest height available to it, these quantities are Lipschitz functions of $A'$ in first-order norms, with a weight $e^{-c\delta\,d^\flat(x,x_b)}$. Here $d^\flat$ is the weighted distance of the ambient sequence. The rate $\delta$ is that of Theorem 2; it is the smaller of the two rates involved, the other being the rate of Theorem 4.

(c) The change of unit $(U,L^{-m}A',L^{-3m}\mathbf J)$, $m=j'-j\ge0$, sends each core of level $j'$ into a core of level $j$, the radii becoming $L^{-m}\alpha_i$. The \emph{outer substitution} is $\mathbf U\mapsto e^{i\eta_j\mathbf A}\mathbf U$, with $|\mathbf A|,|\nabla_{\mathbf U}\mathbf A|<c\alpha_0$ and $c\sim1/K^\sharp$. For radii with $\alpha_0\le\alpha_1$, $\alpha_1^2\le C\alpha_0$, it takes half the radii to the full radii.

By the \emph{inner substitution} we mean the background over $\mathbf 1$, at the datum $\varsigma B+B'$ with $\varsigma\in[0,1]$, of a box $\square_0$ made up of $O(1)$ cubes. Its ingredients are as follows:
\begin{itemize}
\item $e^{iB}$ are the averages of $e^{i\eta_k\mathbf A}e^{i\eta_{k+1}\mathbf H_{k+1}}$;
\item $\mathbf H_{k+1}$ denotes the Landau potential belonging to the box background of $\mathbf U$;
\item $\mathbf A$ is arbitrary subject to $|\mathbf A|<c\,\alpha_{0,k}$, measured in level-$k$ units; in the application one takes for it the cut-off response to the fluctuation field;
\item $B'$ is arbitrary subject to $\sup|B'|<c\min\{\alpha_{0,j-1},\alpha_{0,j}\}$, where $j$ is the level to be reached.
\end{itemize}

Suppose that $(U,A')$ satisfies, on some region containing $\square_0$ and at radii larger than $\alpha_{i,k}$ by a fixed factor, both the simple conditions of level $k+1$ and, at that same level, the tower clause of $\square_0$. Suppose further that the couplings $g_q=x_q^{-1/2}$, $q\le k$, are small and that $-C\le x_q-x_{q'}\le C(q'-q)$ for $q\le q'$. Then, on the domains well inside $\square_0$, the inner substitution in Landau form takes values in the cores of every level $1\le j\le k$, with the running radii $\alpha_{i,j}$. In detail:
\begin{itemize}
\item once the tower of $\square_0$ (first rule) is cut off at depth $j$, the background takes the Landau form $\Phi(\varsigma B+B')=(e^{i\eta_j\mathbf H})^{u}$, with $\mathbf H=\mathcal H(\varsigma B+B')$, where $(\mathcal H,u)$ denotes the pair of Theorem 3 over $\mathbf 1$;
\item the triple $(\mathbf 1,\mathbf H,J(e^{i\eta_j\mathbf H}))$ belongs to the core of level $j$;
\item the pair formed by $\Phi(\varsigma B+B')$ and its current, expressed in level-$j$ units, lies in the complex gauge orbit of the pair given by that triple.
\end{itemize}

(d) Consider the tower of $\square_0$, with determining set $\mathbf B'$, and assume the simple conditions at $(\alpha_0,\alpha_1)$. Then $M_{\mathbf B'}(\mathbf U)=V^v$ with $|V-\mathbf 1|<CM\alpha_0$ and $|v|\,|v^{-1}|<e^{CM\alpha_1}$, and both factors depend only on $\mathbf U|_{\square_0}$.

\textbf{Theorem 6 (the sequences of the construction; restriction; comparison).} (a) Each of the following sequences is separated: the sequences $(\Omega_j)$ of the labels (Notation, \ensuremath{\S}9), together with the sequences of labels to which the domain of the next step has been added; the healed sequences (which arise once a component has been absorbed into $\mathbf R$); the prepared sequences and the stage sequences of $\mathbf R$, which belong respectively to its preparatory steps and to its stages; the towers; the joined sequences of collared regions (for their definition see (c)), including those of the strip and those of the glued sets of $\mathbf R$; and the localisations, which are cube-wise meets of the tower of a region with a separated sequence. (b) Let a separated sequence with $\Omega'_q\subset\Omega_q$ be given, together with a real $\varepsilon_1$-regular $\mathbf V$ such that $CL^2\varepsilon_1\le\bar\varepsilon$ (by the bound and the converse in Theorem 1, the datum of the restriction is then $\bar\varepsilon$-regular). Then

\begin{gather*}
\FitMath{U_{\mathbf B'}\big(M_{\mathbf B'}(U_{\mathbf B}(\mathbf V)|_{\Omega'_0})\big)=U_{\mathbf B}(\mathbf V)|_{\Omega'_0}} \\[3pt]
\FitMath{\text{modulo }\mathcal G^{\rm res}(\mathbf B').}
\end{gather*}

For a sequence of depth $k+1$, no block $B(y)$, $y\in\mathbb T^{(k+1)}\cap\Omega_{k+1}$, contains a node other than $y$. Hence the orbit $U_{k+1}(\mathbf V)$ contains some $U$ for which $\mathcal M^kU$ is block axial there. Moreover, $V^{\rm bg}_k:=\mathcal M^kU$, restricted to the bonds of $\mathbb T^{(k)}$ that have an end point in $\Omega_{k+1}$, does not depend on which such $U$ is chosen, as a count of nodes shows. Applied to the truncation, the restriction identity then gives $U_k(V^{\rm bg}_k)=U_{k+1}(V_{k+1})$, that is, $U_k(\mathbf V')=U_{k+1}(\mathbf V)$, up to gauge. Here $\mathbf V'=(V',(V_q)_{q<k})$, where $V'=V^{\rm bg}_k$ on these bonds and $V_k$ on the remaining bonds of level $k$. (c) Let $\mathbf V$ be real, and let $\mu\ge\rho$ be a power of $L$. A \emph{$\mu$-collared region} $W$ is a cube-wise meet of \emph{parts} $A_a$. In statements on regions the letter $W$ denotes a set of points, and in statements on blocks and boxes it denotes a configuration; its arguments distinguish the two. Each part is a union of cubes of side $\mu\xi_{j_a}$, and two conditions hold. First, the $O(1)$ outermost layers of cubes of $A_a$, which form its \emph{collar}, lie within the single layer $j_a$ of the sequence. Second, the boundaries of distinct parts $A_a$, $A_{a'}$ are separated by at least $O(1)$ cubes of side $\mu\max\{\xi_{j_a},\xi_{j_{a'}}\}$. Its \emph{joined sequence} coincides with the given sequence well inside $W$ and with the tower of $A_a$ in the collar. The \emph{joined datum} equals $\mathbf V$ on the old bonds. On the new bonds, the \emph{collar bonds}, whose levels are $q<j_a$, it is the \emph{lift} of $V_{j_a}$. The lift takes the value $\mathbf 1$ at a bond lying inside a block of level $j_a$, and the value $V_{j_a}(c)$ at a bond leading from the block of $c_-$ into the block of $c_+$. The \emph{localised background} $U^\circ_W$ is defined as the representative of the background attached to this problem. Take one part on one lattice, with $\mu=M_1'$. The problem is then the one posed on the tower of the part, with values taken from the lift of $V$. This is the problem, set on a neighbourhood of a cube $\square$, that defines the localised background $U_{k,\square}(V)$ of the Notation, \ensuremath{\S}4. If $\varepsilon_1\le\mathsf a(L,\rho,\mu)$, then $U|_W=(e^{i\eta f}U^\circ_W)^w$, where

\begin{gather*}
\FitMath{\xi_y|f|,\,\xi_y^2|\nabla f|,\,\xi_y^3|D^*Df|\le C_\mu\,\varepsilon_1e^{-\delta_Wd_{\mathbf B}(y,W^c)},}
\end{gather*}

over each top $y$ of the joined sequence. Here $\xi_y$ is its scale in that sequence: well inside $W$ it is the scale of the given sequence, and in the collar it is the scale of the tower, which can be finer than $\xi_{j_a}$. The rate $\delta_W$ is a fixed fraction of a single rate chosen below those of Theorems 2 and 4. The constant $C_\mu$ depends on $L$, $\rho$ and $\mu$, and $C_\mu\sim e^{O(\mu)\delta_0}$. The reason is that the mismatch is located on the collar bonds, which lie up to $O(\mu)$ from $W^c$ in $d_{\mathbf B}$, so that the decay sets in from the inner edge of the collar onward. The gauge transformation $w$ is determined only at the tops lying away from the collars, where $w=v_q[f]^{-1}$; at every other point it is subject to no restriction. Two backgrounds whose sequences and whose data coincide on $W$ differ there to this extent. (d) The \emph{profile} of the datum at the rate $\delta$ of Theorem 2 is the function $x\mapsto\max_\tau e^{-\delta\,d_{\mathbf B}(x,\tau)}\,|\partial\mathbf V(\tau)-\mathbf 1|$, that is, the right-hand side of the first inequality of that theorem without the constant $C$. Let $e_X$ be its supremum over the neighbourhood of $X$ whose width, in level-$j$ units, is $O(\rho)$, and suppose that $Me_X$ is small. If this neighbourhood is contained in $\Omega_j$, then $(U,0,J(U))|_X$ belongs to the core of $X$ at level $j$ for all radii $\alpha_0\ge Ce_X$, $\alpha_0,\alpha_1>0$. In the opposite case three further conditions are imposed. First, the gaps are wider by a fixed factor, $\operatorname{dist}_\infty(\Omega_q^c,\Omega_{q+1})>C\rho\,\xi_q$ for every $q$; this holds for the sequences of (a) listed before the towers. Second, $X$ is a union of cubes of side $M_1'\xi_j$. Third, the profile multiplied by $M$ is small also near each cube of side $O(M\xi_r)$ in $\Omega_r$, $r\le j$, meeting $X$. Under these conditions the restriction belongs to the layered space of $X$, with radii that, layer by layer, dominate $C$ times the profile.

The later parts rely on three properties of the orbit: criticality, uniqueness among regular configurations, and regularity. Over $\mathbf 1$ the orbit consists of the absolute minima. At other data ``minimiser'' stands for ``regular critical point'', which is the notion that extends to complex data.

\textbf{(iii) Plan of the proofs.}

\emph{Chain.} 1b.1 contains the geometry of seams, the gauges and the definitions given above. 1b.2 begins with the mean of group elements that lie close to one another. It then treats $\mathcal M$ on its domain (analytic, local, covariant) together with its holomorphic extension, and the plaquettes of averages. Next come $\mathcal Q^{\rm st}$, $v_j$, the gauge functional $\mathcal R$ ($w\,v[A]=v[A^w]\mathcal R$), rigidity and the section of one step, all on a single lattice. Finally, everything except the section is extended by folding onto boxes, one layer after another. The folding lemma belongs to 1b.2; 1b.3 adopts it and states it once more so as to apply the statements of 1b.2 on a double block, and 1b.4 constructs a section of one step at a seam. 1b.3 contains the linear operators of mechanisms 2 and 4, the calculus of mechanism 6 and the gauge slice. 1b.4 (Theorem 1) treats first a single lattice and then sequences, carrying the datum down through the truncations and the configurations upward. Its induction is closed by an improvement of the radius that is specific to it: the global halving $\varepsilon_0\to\max\{C\varepsilon_1,\frac12\varepsilon_0\}$, following Proposition 8 [cited from Ba{\l}aban; Package E] of Commun. Math. Phys. 102 (1985) 277--309, with operators taken at a reference configuration cut off to $\mathbf 1$ near a site. Over $\mathbf 1$, flatness follows from a separate energy inequality in scale-invariant $\ell^2$ norms. 1b.5 (Theorem 2) gives that halving in local form. It concerns an arbitrary regular critical configuration and uses only operators at $\mathbf 1$. It uses neither existence nor uniqueness, and it yields the global improvement as a corollary. 1b.6 (Theorem 3: the map $\Phi$, constructed by way of the comb slice) gives the complexification of the system of 1b.4. 1b.7 (Theorem 4; the pair $(\mathcal H,u)$ and its uniqueness) carries out the argument once more in weighted norms, with $R$ and with $K_0$, which on $\Gamma_j$ equals $\xi_j$ times the $Q^\sharp_j(U_0)$ of the Notation, \ensuremath{\S}3, that is, the linear part of the straightened average. 1b.8 (Theorem 5) uses Theorems 1, 2 at real points only, and Theorems 3, 4 on the whole polydisc. From the relation $\Phi_{\mathbf V_0}=(e^{i\eta\mathcal H}U_0)^u$ it takes only the real $B$ near $0$, and it carries out the continuation itself, with an explicit gauge factor. 1b.9 (Theorem 6) combines geometry of cubes, exact algebra of averages and Theorems 1 (one orbit, which drives the restriction), 2 (profile, local gauges), 3, 4 and 5 (cores).

\emph{Interplay with (1c).} (1b) draws on four groups of results of (1c), and (1c) draws on Theorems 3 and 4. Three of these groups depend only on 1b.1--1b.2. They are: (a) on separated sequences, the flat operators at $\mathbf 1$; (b) at complex configurations, the linearised averages together with their gauge identities; and (c) at \emph{regular pairs}, the propagators, together with the rows of $\mathcal M$. The regular pairs form a class defined by inequalities (curvature, gauges on cubes, complex part, independent current) that do not mention any critical point. The fourth group, (d), concerns $G'$, $(Q'G'^2Q'^*)^{-1}$ and $G$ at a real regular $U_0$, with Ba{\l}aban's rows $\mathsf Q$, $\mathsf Q'$. On a single lattice, (d) applies the theorems of Commun. Math. Phys. 99 (1985) 389--434 to the sequence in which every member is the torus. In that paper the first member carries Dirichlet conditions outside it, and members equal to the torus are explicitly admitted in the companion papers (volume 96 after (2.4), volume 99, 75--102, after (1.6)). This group uses these theorems for his own average, which is moved onto the centred blocks by a translation by half a block. For a separated sequence with an arbitrary first member $\Omega_0$, the group derives the same bounds and the positivity by a separate argument. Where that paper uses an expansion whose cover covers the whole torus, this argument uses a parametrix made from one flat operator on the torus, dressed with local gauges. Group (d) depends on those theorems, on (a), on 1b.2, and on the lemmas of 1b.3 that make no use of a bound on a propagator: the algebra of $\mathsf Q',\mathsf Q,R^{\mathsf c}$, pillars and rows, folding, sums over tops, Combes--Thomas, and the second variation of a single plaquette. Its operators are different from those of (c), and the two kinds of rows coincide at $\mathbf 1$. The proposed proof must establish these results in the following order: 1b.1, 1b.2; (a); the geometric and algebraic half of 1b.3; (d); the remainder of 1b.3 and 1b.4--1b.6, which depend on (d); and then 1b.7. 1b.7 uses $G_0,G_1,\mathsf H_1$ of (c) at the pair $(U_0,F'(0))$, which is a regular pair of small radius as soon as $\mathfrak M\bar\varepsilon\le c(L)$. This condition is a ceiling on the radius $\bar\varepsilon$, which is chosen after $\mathfrak M$; it places no upper bound on $\mathfrak M$.

On a sequence, 1b.3 and 1b.4 need less than (d) provides. They use only the positivity and the four bounds that involve no distance: $|G'f|\le C\xi^2\sup|f|$, $|DG'f|\le C\xi\sup|f|$ and the corresponding bounds for $G$, $\nabla G$, where $\xi$ is the local scale. These are needed at an auxiliary first member that closely surrounds $\Omega_1$, namely the one described in that paper on p. 394, which is $\Omega_1$ together with a layer of big blocks around it. Theorems 3.1 and 3.3 of that paper are formulated for a first member of this kind, and they also state these bounds; the paper proves them by an expansion that is set up for a cover covering the whole torus. In the following form, which includes the decay, the positivity and these bounds hold as well. Take a separated sequence with $\Omega_1\ne\emptyset$ and an arbitrary first member. Its cube size $\mathfrak M$ exceeds a threshold that depends on $L$ alone, and its gap is at least $L\mathfrak M$ and exceeds a second threshold of the same kind. Take also a real configuration with small gauges on cubes and with curvature small in the local units. Then the three operators are positive definite. The third of them, and the bounds on $G$ and $\nabla G$, require in addition that the curvature be small in the local units at those bare plaquettes which protrude from $\Omega_0$ while having a corner in it. For $f$ supported over the top $y'$, one has $|G'f|_y\le C\xi_y^2e^{-\delta_0d_{\mathbf B}(y,y')}\sup|f|$. The same bound holds with $\xi_y$ for $DG'$, with $\xi_y^{-4}$ for $(Q'G'^2Q'^*)^{-1}$, and with $\xi_y^2$, $\xi_y$ for $G$, $\nabla G$. Here $C$ and $\delta_0$ depend on $L$ alone, and the proof uses the parametrix built from a single flat operator on the torus and dressed with local gauges. The auxiliary sequence is a sequence of this kind. Its cube size and gap number, both functions of $L$, replace $\mathfrak M$ and the gap number of the given sequence. Three consequences follow (1b.3, 1b.4): the bound on $(Q'G'^2Q'^*)^{-1}$, the decay in $d_{\mathbf B}$ (mechanism 6), and the passage to the given $\Omega_0$. The last of these is obtained by a partition of unity in layer $0$, where the rows act as a mass $\eta^{-2}$. The required smallness is the slope of the cut-offs, and the gap pays for it.

A lemma of 1b.7 states the following. Let $U_0$ be real and regular at the radius $\varepsilon$, and let $\mathfrak M\varepsilon\le c$. Then $(U_0,F'(0))$ is a regular pair of radius $C\mathfrak M\varepsilon$, since the potential is $\sim$ side $\times$ curvature. Only the gauges on cubes require an argument. In layer $0$ the comb gauge is used. Above it, the gauge on a cube is given by Proposition 6 of Commun. Math. Phys. 99 (1985) 75--102; the second half of Theorem 2 above is the counterpart of that proposition for sequences. From that proposition the lemma needs the conclusions (1.135)--(1.136) and the condition (1.7)--(1.9), the latter with (side of the cube) $\times\,\alpha_0$ small. Neither of these involves an average; (1.137)--(1.138) involve one and are not used, while the paper proves the proposition by means of his average. In the paper the proposition is formulated for configurations defined on the whole torus and regular at all that touches $\Omega_0$. In the lemma, the configuration is given on $\Omega_0$ only, and it is regular at the inner plaquettes and at the bonds where the current is defined. When $\Omega_0$ is not the torus, the lemma extends the configuration by $\mathbf 1$. It then applies the proposition to a sequence truncated at the level of the cube. The first member of that sequence is $\Omega_0$ with the layer of cubes of side $\mathfrak M\eta$ that touch $\Omega_0^c$ removed, so the first gap is $\rho-\mathfrak M$. Everything that touches this member, together with the plaquettes at it, lies inside $\Omega_0$, so (1.7)--(1.9) hold at every place where the paper checks them. For the cubes, which are of the levels $j\ge1$, the enlargements are contained both in $\Omega_{j-1}$ and in the smaller first member, and their length is less than half that of the torus. The blocks of the paper are brought into line with the centred ones by the translation by half a block.

The decoupled response of (1c) invokes Theorems 3 and 4 only once 1b.7 is complete, and it does so in order to identify its solution, at the value $1$ of all parameters, with $\mathcal H$; thus this dependency order does not by itself introduce a cycle. The uniform estimates and the applicability of each comparison still have to be proved. The pairs of (1c) are taken from the spaces of Theorem 5 and from the sequences of Theorem 6. These spaces and sequences depend on (1c) only through Theorems 1--4, which is to say only through those results of (1c) that come before these theorems in the order fixed in (1c) (``The order of (1b) and (1c): no circle'').

\emph{The principal mechanisms.}

\begin{enumerate}
\item[1.] \textbf{Straightening, and geometric series in the level.} Once the block factors are removed, the average of a perturbation becomes a perturbation of the average. For a single step the row sums are $\le L$ and the column sums are $\le L^{-3}$, up to relative errors $\mathsf e_m\le C(\alpha_0L^{-2(j-m)}+L^{m-j}\xi_j\sup|A|)$, and these errors are summable. Consequently, when $\alpha_0$ and $|A|_{(-1)}$ are small, and for complex $A$ as well, the rows $c\in\Gamma_j$ satisfy

\begin{gather*}
\FitMath{\sum_b|K_A(c,b)|\le2\xi_j,} \\[3pt]
\FitMath{\sum_c|K_A(c,b)|\le2\eta^4\xi_j^{-3},} \\[3pt]
\FitMath{v_j=\prod_me^{iY_m},\,\,\sum_m|Y_m|\le C|A|_{(-1)} .}
\end{gather*}

It follows that $K_A^*$ is bounded from $\sup$ into $|\cdot|_{(-3)}$, and that the gauge factors obey bounds uniform in $k$. Non-unitary gauge transformations do not preserve the domain of the complex average. The identity $\mathcal M(\mathbf W^g)(c)=g(c_-)\mathcal M(\mathbf W)(c)g(c_+)^{-1}$ remains valid provided that the condition number $|g|\,|g^{-1}|$ at the two coarse end points, multiplied by the field size in the double block, remains small. In the opposite direction, one application of the average multiplies plaquettes by exactly $L^2$:

\begin{gather*}
\FitMath{|\partial\mathcal M(W)(p')-\mathbf 1|\le L^2\varpi+C\varpi^2} \\[3pt]
\FitMath{(\varpi:\text{ the largest }|\partial W-\mathbf 1|\text{ in the four blocks at }p'),}
\end{gather*}

because around $p'$ the contour terms cancel at the corners (lattice Stokes). The factor $L^2$ is absorbed by the change of unit, and the errors add up as a geometric series. Hence, if $|\partial U-\mathbf 1|\le\alpha_0\eta^2$ with $\alpha_0$ small, then $|\partial\mathcal M^nU-\mathbf 1|\le2\alpha_0(L^n\eta)^2$ for every $n\le k$. Beneath a corner of $p'$ that lies outside $\Omega_j$, the curvature is small only when it is measured in the unit of level $j-1$, and this is the loss of $L^2$ across a seam.

\item[2.] \textbf{A single system with a multiplier.}

\begin{gather*}
\FitMath{F'(A)+DRD^*A=K_A^*\omega,} \\[3pt]
\FitMath{\mathcal Q^{\rm st}(A)=\mathsf b(B).}
\end{gather*}

The constraint stays non-linear and is imposed through the multiplier $\omega$, which takes the place of an analytic change of variables bringing the constraint to the linear form $QA=B$. We use Ba{\l}aban's rows $\mathsf Q$ of (i), from which the propagators are constructed, and accordingly $R^{\mathsf c}$ in place of $R$. This holds in mechanisms 2--4, while mechanism 7 uses the form with $K_0$ and $R$. Up to a rotation of the rows, which is omitted from the notation, $\mathcal Q^{\rm st}=\mathsf Q+\mathsf C$. Besides its quadratic part, $\mathsf C$ has a \emph{linear} part of relative size $\alpha_0$, the radius of $U_0$. With $\mathsf G_0=(\Delta^{\rm W}_{U_0}+DRD^*+\mathsf Q^*\mathsf Q)^{-1}$, $\mathsf H_0=\mathsf G_0\mathsf Q^*(\mathsf Q\mathsf G_0\mathsf Q^*)^{-1}$ and $\tilde{\mathsf G}_0=\mathsf G_0-\mathsf H_0\mathsf Q\mathsf G_0$, the system takes the form of a fixed-point equation $(A,\omega)=\Psi(A,\omega)$:

\begin{gather*}
\FitMath{A=\tilde{\mathsf G}_0\mathsf N+\mathsf H_0\big(\mathsf b(B)-\mathsf C(A)\big),} \\[3pt]
\FitMath{\omega=\mathsf H_0^*\big[F'(A)+DRD^*A-\mathsf C'(A)^*\omega\big],} \\[3pt]
\FitMath{\mathsf N=-F'(0)-\mathcal V(A)+\mathsf C'(A)^*\omega .}
\end{gather*}

The current $F'(0)$ of $U_0$ acts as an inhomogeneous term alongside $B$, and every bound is linear in $|F'(0)|_{(-3)}+\sup|B|$. When $U_0$ is critical, the current's only effect is to produce the multiplier $\omega_0$ of $U_0$. After summation by parts, $\mathsf H_0^*D^*DA=(D\mathsf H_0)^*DA$, only first differences of $A$ remain. In terms of size the Lipschitz matrix is triangular. Its entries are $O(\alpha_0+r)$, where $r$ is the radius of the ball, except for one entry $O(1)$, and the norm $\max\{\|A\|_X,\Theta^{-1}\sup|\omega|\}$ overcomes that entry. (The system with the multiplier is solved in this way on one lattice, in 1b.6 and in 1b.7. For sequences the multiplier is instead eliminated, in 1b.4. Since $1+\mathsf H_0^*\mathsf C'(A)^*$ is invertible, the second equation by itself forces $\omega=\mathsf w(A)$. The contraction is then carried out in $\|\cdot\|_X$ alone: the Lipschitz constant of $\mathsf w$ is not small, but it enters only through the small operator $\mathsf C'(A)^*$. Uniqueness then imposes no condition on $\omega$.) Two balls are used: a fixed ball for uniqueness, and a ball of radius $C(|F'(0)|_{(-3)}+\sup|B|)$ for the linear bound. Pairing with the gauge directions that fix $\mathcal Q^{\rm st}$ yields $(RD^*A,RD^*\delta_\lambda A)=0$, and transversality then yields the Landau condition a posteriori.

\item[3.] \textbf{Each regular configuration has a representative in the slice; the three resets.} For uniqueness in a ball to yield \emph{one orbit}, it is necessary that every regular $U$ whose averages lie within $\alpha_1$ of $\mathbf V_0$ can be written as $U^g=e^{i\eta A}U_0$ with $RD^*A=0$, with $g=v_k[A]$ at the nodes, and with $A$ in that ball. Within this mechanism $\alpha_1$ denotes this distance. On one lattice the proof is an induction on the level. A gauge transformation brings the averaged contour variables of $\mathcal M^{k-1}U$ into agreement with those of $\mathcal M^{k-1}U_0$. The rigidity of one step then transfers the closeness to the level below. Rigidity states that if the contour variables are equal, the averages lie within $\alpha_1$ and the plaquettes are $\le\varpi$, then the bonds lie within $\alpha_1+C\varpi$, with coefficient $1$. Rescaling the representative obtained there loses a factor $L^2$. A gauge correction brings back the slice and the normalisation, the last factor of $v_k$ being prescribed at the nodes; it is obtained by Brouwer's theorem and so comes with bounds only. The loss is recovered by an a priori estimate in the slice. This estimate requires no criticality. It requires only that $U_0$ and $e^{i\eta A}U_0$ be regular at a small radius $\alpha_0$, and that $|A|_{(-1)}$ lie below a fixed threshold, which the constant $\mathsf B_1$ below does not depend on:

\begin{gather*}
\FitMath{\mathsf G_0^{-1}A=D^*DA+\Delta'A+\mathsf Q^*\big(\mathcal Q^{\rm st}(A)-\mathsf C(A)\big),} \\[3pt]
\FitMath{\|A\|_X,\,|D^*DA|_{(-3)}\le\mathsf B_1\big(\alpha_0+\sup|\mathcal Q^{\rm st}(A)|\big);}
\end{gather*}

here the currents of $U_0$ and $U$ bound $D^*DA$, and the constant $\mathsf B_1$, which is a multiple of the norm of $\mathsf G_0$, does not change from level to level. On a sequence, the statement that every such regular configuration has a representative in the slice takes the same form, now with the weighted norms and with one gauge condition at each top, but the induction that proves it is a different one. The induction above rests on the gain $\alpha_0\to L^{-2}\alpha_0$ in passing one level down. With a single radius this gain is absent whenever an outer layer is non-empty, since the gain is located on $\Omega_k$ alone. Instead, a single \emph{preliminary gauge} $\hat g$, which is $\mathbf 1$ at the tops, brings the contour variables of every average into line at once. Rigidity \emph{with a defect} then carries the closeness downwards. In this form one adds the mismatch between the contour variables at the two end points of a bond, while the averages keep the coefficient $1$. Let $j$ be the largest index for which both end points of the bond lie in $\Omega_j$. The increments $C\alpha_0L^{2(n-j)}$ are then summable because of the position of the bond, while the factor $2L-1$ is incurred a single time, at a seam, the gauge being frozen there in the outer block:

\begin{gather*}
\FitMath{|\mathcal M^n(U^{\hat g})(b)\,\mathcal M^nU_0(b)^{-1}-\mathbf 1|\le(2L-1)\alpha_1+C\alpha_0} \\[3pt]
\FitMath{(n\le k,\,b\text{ a bond of }\mathbb T^{(n)}\text{ with an end in }\Omega_n),}
\end{gather*}

and this family of bounds passes to truncations. The representative is then constructed by ascending along the truncations $(\Omega_0,\dots,\Omega_j)$, $j=0,\dots,k$. Each stage consists of a rescaling (with loss $L^2$), a gauge correction on the tops and the estimate in the slice, and the output of the estimate does not retain the size of its input. It therefore suffices that, after the loss, the potential remains under the fixed thresholds required by the gauge correction and by the estimate; this condition determines how small $\alpha_0+\alpha_1$ must be. Given two critical configurations, take one of them as $U_0$. Then both $(A,\omega)$ and $(0,\omega_0)$ solve the system with $B=0$, hence $A=0$ and $g=\mathbf 1$ at the nodes. Of the three resets, each of which overcomes a loss that grows with the level, the a priori estimate is the first. The improvement of the radius (mechanism 5) is the second reset. These two devices go back to Ba{\l}aban (see (iv)). The third prevents the descent of the datum from compounding its losses. Away from the seam, the section of a single step is exact; near $\partial\Omega_n$ it loses a factor $C_\downarrow$, multiplied by the tests in a collar, and the gap keeps this collar apart from whatever the earlier sections have spoilt. In a class with two radii, $\varepsilon$ in the collars and $C_\downarrow\varepsilon$ elsewhere, at each level the datum descends with $C_\downarrow\varepsilon$, and not with $C_\downarrow^k\varepsilon$.

\item[4.] \textbf{A lower bound uniform in the levels.} The bound $\mathsf Q\mathsf G_0\mathsf Q^*\ge\gamma_0$ follows from the variational formula for this form, together with a right inverse of all the rows built from bumps at the local scales. The Dirichlet energy of such a bump, $\xi^4(\xi^{-2})^2=1$, is scale invariant, and no bound on the kernel of $\mathsf G_0$ enters. On one lattice the bumps are placed across the face shared by the two blocks of a row, and it suffices to use diagonal dominance. At a seam, rows of two neighbouring levels overlap. There each row carries a single bump lying well inside \emph{one} pillar, and the profiles are biorthogonal, exactly, to both chain counts that neighbouring rows assign as weights to a bare bond. The coupling between neighbouring levels is not assumed small. The interlevel part of the bump argument uses two inputs: the nilpotent structure of the coupling and a bound on its finite elimination polynomial that is uniform in the depth. The first makes the elimination finite; the second yields a depth-independent $\gamma_0$. The bound also holds for rows that the mass term of $\mathsf G_0^{-1}$ does not contain, namely the inward rows and their differences. Conjugation in the manner of Combes--Thomas yields the decay of the inverse.

\item[5.] \textbf{A critical configuration has the regularity of its datum.} One trivialises on a cube of side $\sigma$, applies a cutoff, aligns the Landau slice with a comb gauge of the datum and splits into near and far parts. This gives

\begin{gather*}
\FitMath{\mathsf r(x_0)\le C\,\mathsf v_R(x_0)+\big(C\sigma^2e^{-c\delta_*\sigma}+C\sigma^4\mathsf r_R(x_0)\big)\mathsf r_R(x_0),} \\[3pt]
\FitMath{C,c\text{ independent of }\sigma}
\end{gather*}

(here $\mathsf r,\mathsf v$ denote the local radii of $U$ and the sizes of the tests, and the index $R$ indicates the maximum over a $d_{\mathbf B}$-ball of radius $R\sim L\sigma$). The decay dominates the linear growth of the comb, and for this reason the constant $C$ in front of $\mathsf v_R$ does not depend on $\sigma$. Taking $\sigma$ large then makes $e^{-c\delta_*\sigma}$ dominate $\sigma^2$, the side of the cube times the length of a comb. After that, a regularity radius $\lesssim\sigma^{-4}$ compensates $\sigma^4$. Theorem 2 is then obtained from a maximum principle whose weight varies slowly, at rate $\delta\sim(L\sigma)^{-1}$. In 1b.4 the inequality, in its global form, closes at every level the cycle of radii $C\varepsilon\to L^3C\varepsilon\to C'L^3C\varepsilon\to C\varepsilon$, made up of the change of level, the solution and the improvement. The radius that comes out is therefore the radius that went in.

\item[6.] \textbf{Each crossing of a layer costs $\rho$.} For every rate $a$ that occurs, with $\rho$ chosen so that $a\rho\ge C\log L$,

\begin{gather*}
\FitMath{d_{\mathbf B}(y,y')\ge\rho(|j-j'|-1)_+,} \\[3pt]
\FitMath{\sum_{y'}e^{-a\,d_{\mathbf B}(y,y')}\le\mathsf S(a) :}
\end{gather*}

the factor $e^{-a\rho}$ outweighs the entropy, $L^4$ for each level, and a fraction of the rate pays for powers of $\xi_y/\xi_{y'}$. The distance $d_{\mathbf B}$ itself is too rough to serve as a conjugation weight for the propagators. In 1b.3 and 1b.4 one therefore works with a \emph{regularised distance} $\mathsf f=\mathsf f_{y'}$, whose first and second differences are controlled at the local scale $\xi(x)$:

\begin{gather*}
\FitMath{|\mathsf f-d_{\mathbf B}(\cdot,y')|\le C,} \\[3pt]
\FitMath{|\mathsf f(x')-\mathsf f(x)|\le C\frac{\eta}{\xi(x)},} \\[3pt]
\FitMath{|\mathsf f(x+\eta e_\mu)+\mathsf f(x-\eta e_\mu)-2\mathsf f(x)|\le C\frac{\eta^2}{\xi(x)^2}.}
\end{gather*}

Estimated term by term, the scale-normalised Laplacian perturbation from conjugation by $e^{\theta\mathsf f}$ has an apparent factor $\theta\,\xi/\eta$. Pairing the directions $\pm e_\mu$ and separating a field from its neighbouring increments gives, for small $\theta$,
\[
 |(e^{\theta\mathsf f}\Delta e^{-\theta\mathsf f}-\Delta)v|
 \le C\theta\big(\xi^{-2}|v|+\xi^{-1}|\nabla v|\big),
\]
with the field and gradient norms taken on the local stencil. Thus the no-distance bounds on $G$ and $\nabla G$ stated above bound the composed perturbation $(e^{\theta\mathsf f}\Delta e^{-\theta\mathsf f}-\Delta)G$ by $C\theta$ in the corresponding force supremum norms. This is a relative inverse estimate, not a uniform bound for the bare Laplacian perturbation on arbitrary bounded fields. Together with the corresponding relative bounds for the other operator terms, it is the estimate used by the Combes--Thomas argument to obtain decay in $d_{\mathbf B}$.

\item[7.] \textbf{Decay from the same contraction.} In 1b.7 the map $\Psi$ of mechanism 2 is formed with $K_0$ itself as the linear rows, with the projection $R$ of the Landau gauge, and with $\mathsf H=G_1K_0^{\mathsf T}(K_0G_1K_0^{\mathsf T})^{-1}$ and $\tilde{\mathsf G}=G_1-\mathsf HK_0G_1$. Since $K_0$ carries the length $\xi_c$ of its row $c$, this $\mathsf H$ is the $\mathsf H_1$ of the Notation, \ensuremath{\S}5, multiplied by $\xi_c^{-1}$. At the critical $U_0$, the operator $G_1^{-1}$ contains the Hessian, the gauge term and the mass term of the rows, and in addition the local term $-\mathsf W_0$. Here $\mathsf W_0\delta=\frac{d}{dt}K_{t\delta}^{\mathsf T}\omega_0\big|_{t=0}$ is obtained by pairing the second-order part of the constraint with the multiplier $\omega_0$, and accordingly $\mathsf N$ contains here the further term $-\mathsf W_0A$. Because the constraint is kept non-linear, every non-linear term is local. The non-locality is confined to $\mathsf H$, $\tilde{\mathsf G}$, $\mathsf H^{\mathsf T}$ and $1-R$, and the last of these decays like the others. The projection $R$ enters through the term $\mathsf H^{\mathsf T}DRD^{\mathsf T}A$ of the equation for $\omega$, and at the critical $U_0$ this term vanishes, because $R^{\mathsf T}=R$ and $RD^{\mathsf T}G_1K_0^{\mathsf T}=0$. Consequently $\|\partial\Psi\,Z\|_{\theta,y_0}\le\frac12\|Z\|_{\theta,y_0}$ in the norms weighted by $e^{\theta d_{\mathbf B}(\cdot,y_0)}$. This holds for every centre $y_0$, for every $\theta$ from $0$ up to a fixed fraction of the decay rate of $\mathsf H$, $\tilde{\mathsf G}$, and at every point of a fixed small ball around $(0,\omega_0)$. Denote by $\mathbf A(\mathsf b)$ the solution of the system in which the right-hand side $\mathsf b$ of the constraint is left free. Its derivative with respect to $\mathsf b(c)$, together with the derivative of the multiplier, satisfies

\begin{gather*}
\FitMath{(1-\partial\Psi)(\mathsf h,\dot\omega)=(\mathsf He_c,0).}
\end{gather*}

Then $\mathcal H=\mathbf A\circ\mathsf b$, and Theorem 4 follows by the chain rule through $\mathsf b$, since each bond enters $O(1)$ rows within $O(L)$. The second-order members are obtained from the differentiated equation, not from a kernel.

\item[8.] \textbf{Passing from real points to the polydisc; the covariant curl is small.} The iterates are holomorphic, and by Weierstrass's theorem their uniform limit is holomorphic as well. The identities that hold at real $B$ (the slice, criticality, covariance) extend to complex arguments by the identity theorem, and the Schwarz lemma makes the bounds linear. Let $\phi$ be holomorphic and bounded on a ball of radius $r$, and let $w=w_{\rm R}+iw_{\rm I}$, where $w_{\rm R},w_{\rm I}$ are real and $\|w_{\rm R}\|,\|w_{\rm I}\|\le\|w\|\le r/2$. Then

\begin{gather*}
\FitMath{|\phi(w)|\le C\big(|\phi(0)|+|\phi(w_{\rm R})|+|\phi(w_{\rm I})|\big)+C\sup|\phi|\,\|w\|^2/r^2 :}
\end{gather*}

the first-order Taylor term is complex linear, so the crude supremum appears only in the quadratic remainder. We apply this with $\phi$ the rescaled plaquette or the rescaled current of a tower configuration, and with $w=A'=A'_{\rm R}+iA'_{\rm I}$. If one uses only $|A'|,|\nabla_UA'|<\alpha_1'$, the curvature of the real fields $e^{i\eta A'_{\rm R}}U$, $e^{i\eta A'_{\rm I}}U$ is no better than $O(\alpha_0'+\alpha_1')$. However, the two bounds $|\partial U-\mathbf 1|,|\partial\mathbf U-\mathbf 1|<\alpha_0'\eta^2$, together with

\begin{gather*}
\FitMath{\partial(e^{i\eta A'}U)=\big(\mathbf 1+i\eta^2D_UA'+O(\eta^2|A'|^2+\eta|A'|\,|\partial U-\mathbf 1|)\big)\,\partial U}
\end{gather*}

yield $|D_UA'|\le C(\alpha_0'+\alpha_1'^2)$. Since $\operatorname{Ad}(U)$ commutes with taking adjoints, the real part and the imaginary part satisfy the same bound, and therefore both real fields have curvature $C(\alpha_0'+\alpha_1'^2)$. Theorem 1 is linear in the datum, and it bounds $\phi$ at the three real points. This gives the bound $K^\sharp(\alpha_0'+\alpha_1'^2)$ without any bound on $\nabla A'$. The square is essential because of homogeneity. Measured in units of level $n$, the radii equal $L^{-2(p-n)}\alpha_0'$ and $L^{-(p-n)}\alpha_1'$. Hence $\alpha_1'^2$ has the same scaling as $\alpha_0'$ and as the two bounds of the clause, which in level-$n$ units are $\eta_p^2=L^{-2(p-n)}\eta_n^2$ and $L^{2(n-p)}$ (Theorem 5 (a); the weight $L^{-(j-p)}$ enters through the radius and is independent of $n$). A bound that is linear in these two radii, with a constant that does not depend on the level (the Schwarz lemma provides such a bound on the ball), costs a factor $L^{p-n}$. In these units $|\nabla_UA'|$ is bounded by $L^{-2(p-n)}\alpha_1'$, which is smaller than the radius by a factor $L^{-(p-n)}$. Consequently, when $\alpha_0'+\alpha_1'\le c$, the inequality for $\phi$ together with Theorem 1, applied to the two real fields, whose curvature is $O(\alpha_0'+\alpha_1')$, yields $C(\alpha_0'+\alpha_1')\eta_p^2$ and $C(\alpha_0'+\alpha_1')L^{2(n-p)}$, and these bounds do not incur that factor. Compared with this bound there is a second gain. The condition $K^\sharp(\alpha_0'+\alpha_1'^2)\le\alpha_0$ allows $\alpha_1'\gg M\alpha_0$, and therefore allows $\alpha_1'\ge c\alpha_1$ together with $M\alpha_0\ll\alpha_1$, the relation that holds between the two radii in the application. A condition linear in $\alpha_1'$ instead forces $\alpha_1'\le\alpha_0/K^\sharp$.

\item[9.] \textbf{The coarser constraints are frozen by the finer ones.} Let $U'\in C(\mathbf B',\mathbf V')$ be extended by $U$. Induction on the level then gives $M_{\mathbf B}(\widetilde{U'})=\mathbf V$ and $A_{\Omega_0}(\widetilde{U'})-A_{\Omega_0}(U)=A_W(U')-A_W(U|_W)$, so restriction preserves criticality. By mechanism 1 the restricted datum is $CL^2\varepsilon$-regular. It therefore has an orbit when $CL^2\varepsilon\le\bar\varepsilon$, and the restriction identity follows from uniqueness. For the comparison one performs a gauge alignment on the collar; it turns the mismatch into a real $B_1$ with $|B_1|\le C\varepsilon_1$, supported on the bonds of the collar, and $f=\mathcal H(B_1)$ decays by Theorem 4. The alignment relies on 1b.9, the lemma on the \emph{defect of lifts}, whose step from one level to the next is the reconstruction bound of 1b.2. After a gauge transformation fixing the sites of level $j$, each lower average of $U$ agrees with the lift of the level-$j$ average up to the total, over the intermediate averages, of their plaquette sizes, the coefficient being $1$ at each passage from one level to the next. When the collar lies inside one layer $j$, this sum is the geometric series $\sum_{q<j}2\varepsilon L^{2(q-j)}\le C\varepsilon$, uniformly in $j$. It is here that the argument needs the boundary to run within one layer.
\end{enumerate}

\emph{Constants} are chosen in the following order, each one large or small enough in terms of those chosen before it: first $L$; then the one-step constants of $\mathcal M$; then $\delta_0$ together with the propagator constants; then $\gamma_0$, with the derived rate $\delta_*$; then the length $\sigma$; next, the constant $\delta\sim(L\sigma)^{-1}\ll\delta_0$ that appears in Theorem 2 and, at the same step, the regularity radius $\lesssim\sigma^{-4}$; then the cube size $\mathfrak M$, together with the gap $\rho=R_1\mathfrak M\gg\sigma$, subject to (rate)$\,\times\rho\ge C\log L$ for each rate in use; then the remaining radii ($\bar\varepsilon$, $r_{\rm B}$), with $\mathfrak M\bar\varepsilon$ required to be small; then $M_1'$ and $M$; then the threshold of Theorem 6 (c), chosen after $\mu$ (in the application $\mu=M_1'$); finally the field strengths, with $M\alpha$ small. In the order of choice fixed in the Notation, everything listed here before $M_1'$ comes after $L$ and before $M_1'$, and $\mathfrak M$, the cube size, stands where the Notation has $M_1$. For Theorem 1, and for Theorem 3 in its first half (the map $\Phi_{\mathbf V_0}$, together with the pair of the comb slice), the thresholds involve the gap only through one length. That length may be fixed at its least admissible value, which is a function of $L$; with this choice the thresholds depend on $L$ alone, and so do the constants that the Schwarz lemma derives from them (crude bound divided by radius). The pair $(\mathcal H,u)$ of Theorem 3 requires in addition $\mathfrak M\bar\varepsilon\le c(L)$, and so, along with it, does Theorem 4; for both, the threshold on the datum therefore depends also on $\mathfrak M$. No upper bound is placed on $\mathfrak M$, $\rho$, $M$.

\subsection*{The reduced datum has an explicit local inverse}

We now detail the two finite inversions used in the construction.
In (BGD1)--(BGD4), $B$ and $b$ are raw row logarithms.
Their per-length versions are $\ell^{-1}B$ and $\ell^{-1}b$,
where $\ell=\operatorname{diag}(\xi_c)$ is the row-length matrix.
Fix a real reference field \(U_0\) whose used averages lie on the
regular branches of the actual finite mean. For an interior row
the reduced datum is \(b(B)(c)=B(c)\). On an outward row oriented
toward its outer endpoint \(z\), it is
\[
 e^{ib(B)(c)}=e^{iB(c)}Z_c(B_{\rm int}),\qquad
 Z_c(B_{\rm int})
       =\operatorname{Ad}(V_0(c))\,v_z(B_{\rm int}).
 \tag{BGD1}\label{ym4:exp500-bg:mb:reduced-datum}
\]
The factor \(v_z\) is the actual one-step block factor made from
the interior data in the finer block centered at \(z\).
The determining-set geometry above shows that these are prescribed
interior rows, not another set of unknown outward rows.
For the reverse orientation use the inverse group identity.

There is therefore a direct inverse. Recover \(B_{\rm int}=b_{\rm int}\),
then set
\[
 B(c)=\frac1i\log\big(e^{ib(c)}
                         Z_c(b_{\rm int})^{-1}\big)
 \quad\text{on outward rows}.
 \tag{BGD2}\label{ym4:exp500-bg:mb:reduced-inverse}
\]
All logarithms are on the common small datum chart. Conjugation by
the real \(V_0(c)\) has norm one; it imposes no smallness of that
coarse edge or of a global holonomy. The finite mean construction
gives a radius \(r_0(L)>0\) on its admitted local reference class.
On a fixed smaller ball, its block factors and their first two
derivatives are bounded by constants depending on that finite
construction. Thus (BGD1)--(BGD2) are admitted on balls
\(\|B\|_\infty<r_0/C_L\), respectively
\(\|b\|_\infty<r_0/C'_L\), and satisfy
\[
 \|b(B)\|_\infty\le C_L\|B\|_\infty,\qquad
 \|b^{-1}(b)\|_\infty\le C'_L\|b\|_\infty.
 \tag{BGD3}\label{ym4:exp500-bg:mb:reduced-bounds}
\]
The uniform statement is over local reference patches on which that
same finite chart and its radius are admitted. Establishing those
patches for a multilevel background is part of its regular-domain
estimate; it is not inferred just from the existence of \(U_0\).

It is useful to see exactly why the inverse does not contain a
long product of triangular matrices. Define the right logarithmic
differential
\(\mathcal E_X(H)=(1/i)(d/dt)e^{i(X+tH)}|_{t=0}e^{-iX}\).
Differentiating (BGD1) gives the outward diagonal
\(D_c=\mathcal E_{b(c)}^{-1}\mathcal E_{B(c)}\), not the identity
in general. Let \(D\) be this block diagonal matrix, with identity
on interior rows, and let \(T\) contain all derivatives of block
factors. Its only nonzero blocks have interior columns and outward
rows. Consequently
\[
 Db=D+T,\qquad TD^{-1}T=0,\qquad
 (Db)^{-1}=D^{-1}-D^{-1}TD^{-1}.
 \tag{BGD4}\label{ym4:exp500-bg:mb:seam-datum-nilpotence}
\]
Multiplying the last two matrices verifies both inverse identities.
This nilpotence is a property of the actual reduced datum. An
arbitrary lower triangular coupling between successive levels would
not have it.

For the row count, one outward factor reads the \(O(L^4)\) interior
data of its single finer block. For the column count, an interior
bond belongs to one assigned finer block, and that block can be
the outer endpoint block of at most eight coarse incident rows.
Thus both row and column counts are bounded at fixed \(L\),
independently of the number of layers. After inserting the actual
cell weights and amplitudes, an interaction changes scale only
from \(j-1\) to \(j\); its weight ratios are fixed powers of \(L\).
Equations (BGD3)--(BGD4) therefore also give two-sided weighted
Schur bounds in those declared norms. Higher derivatives follow
from the same finite group formula and its fixed datum disc.

This local datum calculation is different from constructing a
right inverse of all the field constraint rows. The latter must
also lift fields and control their energy and gauge slice at a
seam. That primitive linear task is retained in the precision
section. The exact derivative \(Db\), including its nonidentity
outward diagonal, is the datum forcing in the response equations.


\subsection*{Eliminating the multiplier}

We next write the multiplier as a function of the potential.
Here we use the per-length constraint
$q(A)=\ell^{-1}\mathcal Q^{\rm st}(A)$ and put $Q_0=Dq(0)$.
The field and row pairings are the declared cell-weighted pairings,
and $\sharp$ denotes their transpose. With these pairings fixed,
the corresponding multiplier is
$\omega_{\rm len}=\ell\omega_{\rm raw}$; we write it as $\omega$
in (BGD5)--(BGD8). The field space
\(X\) carries the norm \(\|A\|_X=\sup_y\max\{\xi_y|A|,\xi_y^2|\nabla A|\}\),
and \(Z\) is the multiplier space with its row weights. Here \(X^*\) denotes the designated physical force space, with norm \(\|f\|_{X^*}=\sup_y\xi_y^3|f|\), paired with fields using the declared cell measures; it is not the Banach dual of \(X\).
No a priori multiplier bound enters the uniqueness argument. Fix
\(Q_0E_0=1\), and write
\(Dq(A)=Q_0+C'(A)\). At a critical
field, apply the transpose of this right inverse to the
multiplier equation. This gives
\[
 [1+(C'(A)E_0)^\sharp]\omega=E_0^\sharp DF(A).
 \tag{BGD5}\label{ym4:exp500-bg:mb:eliminate-multiplier}
\]
If \(\|(C'(A)E_0)^\sharp\|_{Z\to Z}<1/2\),
the bracket has its convergent Neumann inverse,
and the multiplier is the explicit function
\[
 \omega(A)=
    [1+(C'(A)E_0)^\sharp]^{-1}E_0^\sharp DF(A).
 \tag{BGD6}\label{ym4:exp500-bg:mb:multiplier-function}
\]
The required transpose bound is a consequence of
the column as well as row estimate of the
actual lift, with the declared force and
multiplier amplitudes. It does not follow
from an unweighted supremum bound on \(E_0\).

Use the prescribed propagator lift $E_0$ and put
$F_0(A)=E_0^\sharp DF(A)$. Here $\gamma$ controls the paired
force throughout the chosen field ball, including its variation away
from $A=0$. The bounds needed on the field ball are
$\|F_0(A)\|_Z\le C_E\gamma$ and
$\|F_0(A)-F_0(A')\|_Z\le C_EC_F\|A-A'\|_X$.
They concern the paired force: in the quadratic part,
$E_0^\sharp D^{\mathsf T}DA=(DE_0)^{\mathsf T}DA$,
so the weighted derivative bounds on the lift leave only first
differences of $A$. No pointwise second-difference bound on all
fields of the ball is needed. Let the bracket's perturbation
have Lipschitz constant $C_N$.
The inverse difference identity gives
\[
 \begin{split}
 \|\omega(A)\|_Z&\le2C_E\gamma,\\
 \|\omega(A)-\omega(A')\|_Z
       &\le(2C_EC_F+4C_E\gamma C_N)\|A-A'\|_X .
 \end{split}
 \tag{BGD7}\label{ym4:exp500-bg:mb:multiplier-bounds}
\]
The second constant need not be small.
However a field fixed-point equation reads
it through the small factor \(C'(A)^\sharp\).
If \(\|C'(A)^\sharp\|\le cr\) and its
Lipschitz constant is \(c\), then
\[
 \|C'(A)^\sharp\omega(A)
                 -C'(A')^\sharp\omega(A')\|_{X^*}
 \le \{cr\,{\rm Lip}(\omega)+2cC_E\gamma\}
                                      \|A-A'\|_X .
 \tag{BGD8}\label{ym4:exp500-bg:mb:small-multiplier-entry}
\]
The coefficient is small with \(r+\gamma\).
This is the field-only form of the coupled
contraction. It shows why a non-small
derivative of the multiplier is not an
obstruction.

Alternatively suppose the two-component
iteration has a Lipschitz matrix bounded by
\(\left(\begin{smallmatrix}\epsilon&\epsilon\\C&\epsilon\end{smallmatrix}\right)\).
In the norm
\(\max(\|A\|_X,\Theta^{-1}\|\omega\|_Z)\)
its row sums are at most
\(\epsilon+\epsilon\Theta\) and
\(C/\Theta+\epsilon\).
Choose a fixed \(\Theta\ge4C\) first,
then \(\epsilon(1+\Theta)\le1/4\).
Both row sums are less than one half.
The admitted multiplier radius must be
at least \(\Theta\) times the field
radius used by this ball. These choices
are made with the actual physical
norms, not by suppressing the multiplier
dimension in notation.

% END MODULE chapters/03_background.tex


% BEGIN MODULE chapters/04_propagators.tex
\section{Constrained propagators and local decay}
\label{candidate:propagators}
\candidatebanner{A}

\subsection*{(1c) Propagators in a background field: exponential decay produced by the averages}

Every object produced by a renormalisation step is assembled from inverses of ``Hessian of the action + gauge term + mass made of block averages''. These inverses are taken around a background on a decreasing sequence of regions, in which each layer has a scale of its own. The construction uses five properties of these inverses: they exist; they depend analytically on the background; in each layer they behave like massive Green's functions measured in the unit of that layer; they decay exponentially; and they are sums of terms that are local both in the field and in the regions. A single requirement governs the whole construction. Since the regions are supplied by the large-field analysis, none of the constants may vary with their shape or their number, nor with the level, the cutoff or the torus. Two numbers are exceptions, and neither of them enters any localised bound, rate, weight or sum of weights. (Ba{\l}aban's papers are referred to by the volume of Commun. Math. Phys. in bold. Where two papers must be told apart, the first page is added: \textbf{99}, 75 and \textbf{99}, 389; \textbf{102}, 277; \textbf{122}, 175. The references are collected in the bibliography. The principal statements supplied in (ii) are numbered 0--5. The mechanism descriptions also refer to auxiliary hypotheses 6--8, whose required operator data are listed before (iii). The mechanisms in (iii) carry the letters (a)--(k).)

\textbf{(i) The setting and the objects.}

\emph{Geometry.} The integer $L$ is odd and large. The cube size of the walk expansions, $M_1$, and the gap number, $R_1$, are both powers of $L$, with $R_1\ge L$. We work in level-$k$ units, in which $\mathbb T^{(j)}$ has spacing $\xi_j:=L^{j-k}$; moreover $\eta=\eta_k=\xi_0$. A sequence $\Omega_0\supset\dots\supset\Omega_k$ is called \emph{separated} when every $\Omega_j$ is a union of $M_1$-cubes at level $j$ and

\begin{gather*}
\FitMath{\operatorname{dist}_\infty(\Omega_j^c,\Omega_{j+1})>R_1M_1\xi_j .}
\end{gather*}

No further condition is placed on the regions. In particular $\Omega_0$ is arbitrary, and Dirichlet conditions are imposed outside it. \emph{Layers} are the sets $\mathcal L_j=\overline{\Omega_j\setminus\Omega_{j+1}}$. The \emph{tops} $y$ are the sites of $\mathbb T^{(j)}$ lying in $\mathcal L_j$. Such a top has local scale $\xi_y=\xi_j$ and block $\Delta(y)$, which is the cube of side $\xi_y$ with centre $y$. The \emph{rows} $c$ are the bonds of $\mathbb T^{(j)}$ that have one end point in $\mathcal L_j$ and the other end point either in $\mathcal L_j$ (the \emph{interior} rows) or, when $j\ge1$, outside $\Omega_j$ (the \emph{outward} rows); for them $\xi_c=\xi_j$. The \emph{inward} bonds, those that enter $\Omega_{j+1}$, carry no row. The distance

\begin{gather*}
\FitMath{d_{\mathbf B}(y,y')=\inf_{\Gamma\ni y,y'}|\Gamma|_{\mathbf B},} \\[3pt]
\FitMath{|\Gamma|_{\mathbf B}=\sum_{j\ge1}\xi_j^{-1}\operatorname{length}(\Gamma\cap\mathcal L_j)+\xi_0^{-1}\operatorname{length}(\Gamma\setminus\Omega_1)} \\[3pt]
\FitMath{(\Gamma\text{ connected})}
\end{gather*}

measures each layer in the unit of that layer. Outside $\Omega_1$, and hence also outside $\Omega_0$, length is counted in the bare unit. Therefore $d_{\mathbf B}$ is independent of $\Omega_0$, and a path that leaves $\Omega_0$ finds no cheaper unit of length out there. This distance differs from the distance (2.46) of \textbf{96}, whose contours are made up, in each layer, of bonds of the lattice of that layer. A larger class of paths is allowed here, and so $d_{\mathbf B}$ is the smaller of the two.

\emph{Regular pairs} $(\mathbf U,\mathbf J)$ consist of $\mathbf U=e^{i\eta A'}U$, with $U$ real, together with an \textbf{independent} $\mathfrak g^c$-valued $1$-form $\mathbf J$ that represents the current. Such a pair is required to satisfy, on $\mathcal L_j$,

\begin{gather*}
\FitMath{|\partial U-\mathbf 1|,|\partial\mathbf U-\mathbf 1|<\alpha_0\eta^2\xi_j^{-2},} \\[3pt]
\FitMath{\xi_j|A'|,\,\xi_j^2|\nabla_UA'|<\alpha_1,} \\[3pt]
\FitMath{|\mathbf J|<\gamma_0\xi_j^{-3},}
\end{gather*}

and, on every cube of side $O(M_1\xi_j)$ contained in $\mathcal L_j\cup\mathcal L_{j+1}$, to admit a real gauge $U=e^{i\eta A}$ with $\xi_j|A|,\xi_j^2|\nabla A|\le\alpha_0$. In this condition the bound is $\alpha_0$ on the entire cube, and not $O(M_1)\alpha_0$. Pairs in which the current is a separate variable, with a radius $\gamma_0$ of its own, are due to Ba{\l}aban (\textbf{109}, p. 262, (1.10)--(1.14)). In the present setting they are formed on a sequence, with the local scales. Throughout, $J(\mathbf U)$ is the first variation of the action of the \emph{inner} plaquettes, namely those inside $\Omega_0$. At every bond whose six plaquettes are all inner, it coincides with the current $D^*_{\mathbf U}F_{\mathbf U}$; when $\Omega_0=T_\ell$, every bond is of this kind. \emph{On shell}: the pair satisfies $\mathbf J=J(\mathbf U)$ at the bonds that touch $\Omega_1$. \emph{Real}: $A'=0$, and $\mathbf J$ is real. \emph{Local class on $W$}: the same conditions, required inside $W$.

\emph{Operators.} Relative to $\mathbf U$, the \emph{straightened} $j$-fold average of $e^{i\eta A}\mathbf U$ takes, at a row $c$ of level $j$, the value $\exp i\mathcal Q(A)(c)$. Here $\mathcal Q=\mathsf m\,(Q+C^{(2)}+\dots)$, and $\mathsf m$ is multiplication by $\xi_c$. Thus $Q$, which maps $1$-forms to rows, is the linear part per unit length, and $C^{(2)}$ is the quadratic part per unit length. When the level has to be displayed, the straightened average is written $\mathcal Q^{\rm st}_j(\mathbf U;A)$. In the units of level $j$, its linear part is the operator $Q^\sharp_j(\mathbf U)$ of the Notation, \ensuremath{\S}3, which is built from the means over the contour families of $\mathcal M$. The index st stands for ``straightened''; the \emph{straight} linear average $Q^{\rm st}_j(\mathbf U)$ of the Notation is a different operator. Consequently, in this part the unadorned letters $\mathcal Q$ and $Q$ differ from those of \ensuremath{\S}3 of the Notation: they denote neither the logarithm $\mathcal Q_j(\mathbf U;A)$ of $\mathcal M^j(e^{i\eta A}\mathbf U)\,\mathcal M^j(\mathbf U)^{-1}$ nor its linear part $Q_j(\mathbf U)$. The straightened and the plain averages are related by \emph{site factors} $v_n$. These are group-valued functions on the sites of $\mathbb T^{(n)}$, and they accumulate from one step to the next (statement 3: $\mathcal M^n(e^{i\eta A}\mathbf U)=\big(e^{i\mathcal Q^{\rm st}_n(A)}\mathcal M^n\mathbf U\big)^{v_n}$). The linear parts of the straightened and of the plain average already differ at $\mathbf 1$, by a gauge term on the coarse lattice. For the plain averages, the identity $QD=\bar DQ'$ fails, and so does the bound $1+C(\dots)$ on the derivative in statement 3. The operator $Q'$ maps site functions to tops. It is the linearisation of the average of a gauge transformation, taken with the same contour variables. The symbol ${}^{\mathsf T}$ denotes the transpose with respect to bilinear pairings, so that everything remains holomorphic. The operator $a$ acts on rows as multiplication by $a_0\xi_c^{-2}$, where $a_0$ is a fixed number (equal to $1$ in the sub-parts). On tops it acts as multiplication by a comparable weight $\approx\xi_y^{-2}$. We write $D=D_{\mathbf U}$. Finally, $\Delta^{\rm W}_{\mathbf U}$ is the Hessian of the action in $\Omega_0$. In statements 3--8 it is the Hessian of the action of the inner plaquettes, which is an exactly invariant function. In statements 1, 2 it is obtained by taking the Hessian on the entire lattice and compressing it to the bonds of $\Omega_0$.

\begin{gather*}
\FitMath{G'=(\Delta_{\mathbf U}+Q'^{\mathsf T}aQ')^{-1},} \\[3pt]
\FitMath{R=1-G'Q'^{\mathsf T}(Q'G'^2Q'^{\mathsf T})^{-1}Q'G',} \\[3pt]
\FitMath{\mathsf S=1-DG'RD^{\mathsf T},}
\end{gather*}


\begin{gather*}
\FitMath{\mathsf G_0=(\Delta^{\rm W}_{\mathbf U}+DRD^{\mathsf T}+Q^{\mathsf T}aQ)^{-1},} \\[3pt]
\FitMath{\mathsf G=(\mathsf G_0^{-1}-\Delta'_\pi[\mathbf J])^{-1},} \\[3pt]
\FitMath{\mathsf G_1=(\mathsf G^{-1}-\mathsf S^{\mathsf T}\Delta^{(2)}[\mathsf H^{\mathsf T}\mathbf J]\mathsf S)^{-1},} \\[3pt]
\FitMath{\mathsf H_{(1)}=\mathsf G_{(1)}Q^{\mathsf T}(Q\mathsf G_{(1)}Q^{\mathsf T})^{-1}}
\end{gather*}

(where a bracketed index occurs, the formula holds both with that index and without it). The operators are defined in this order: $\mathsf G$, then $\mathsf H$ from $\mathsf G$, then $\mathsf G_1$, then $\mathsf H_1$ from $\mathsf G_1$. The operator $\Delta'_\pi[\mathbf J]$ is linear in $\mathbf J$; at the actual current it restores the gauge invariance of the Hessian, as in statement 3. Moreover, $(A,\Delta^{(2)}[\mathsf j]A)=2(C^{(2)}(A),\mathsf j)$. The \emph{six propagators} are 

\begin{gather*}
\FitMath{\mathsf G,\mathsf G_1,\mathsf H,\mathsf H_1,(Q\mathsf GQ^{\mathsf T})^{-1},(Q\mathsf G_1Q^{\mathsf T})^{-1}}
\end{gather*}

.

The current correction is defined by the gauge identity of
\cite[(3.117)--(3.120)]{B85c}, with the current treated as an
independent variable as in \cite[(1.8)--(1.14)]{B87}.
All operators are formed at the same reference $\mathbf U$, with
the site map, slice and pairings fixed above. Put
\[
 \lambda_A=G'RD^{\mathsf T}A,\qquad
 \lambda_A^+(b)=\operatorname{Ad}(\mathbf U(b))\lambda_A(b_+),
 \qquad \mathsf SA=A-D\lambda_A.
\]
In every bracket below, $\lambda_A$ without a superscript is
evaluated at $b_-$, so all factors are based at the negative
endpoint. The symmetric bilinear form is
\begin{equation}
\begin{split}
 \langle A,\Delta'_\pi[\mathbf J]C\rangle_{\rm f}
 =\tfrac12\big\langle
  &i[\lambda_A,C]+i[\lambda_C,A]
   -i[A,\lambda_C^+]-i[C,\lambda_A^+]\\
  &-i[\lambda_A,D\lambda_C]-i[\lambda_C,D\lambda_A],
                              \mathbf J\big\rangle_{\rm f}.
\end{split}
\label{ym4:explicit-current-continuation}
\end{equation}
The fine pairing includes its cell volume $\eta^4$. For complex
fields the transpose and colour pairing are bilinear. On the physical
current, for the invariant inner-plaquette action, the second variation
of gauge invariance gives
\[
 \Delta^{\rm W}-\Delta'_\pi[J(\mathbf U)]
                    =\mathsf S^{\mathsf T}\Delta^{\rm W}\mathsf S.
\]
Equation~\eqref{ym4:explicit-current-continuation} specifies the
continuation away from that physical current. The physical identity
alone would not specify it. A compressed whole-lattice Hessian retains
the boundary comparison described below.

In addition we set $\mathfrak G=\mathsf G_1-\mathsf G_1DRD^{\mathsf T}\mathsf G_1-\mathsf H_1Q\mathsf G_1$. For step $k+1$ we work on the bonds of $\mathbb T^{(k)}$ that touch $\Omega_{k+1}$, with $V=\mathcal M^k\mathbf U$. Let $\mathsf C_V$ be the elimination of the pivot bonds, and let $\mathsf G^{(2)}_V$ be the quadratic part of the soft axial gauge fixing. Let $\mathsf E^{(2)}$ be the quadratic part of the one-step constraint, which is imposed through the plain $\mathcal M$ at $V$ and whose linear part is the $Q_V$ of the Notation, \ensuremath{\S}3. We set

\begin{gather*}
\FitMath{\Delta_k=(Q\mathsf G_1Q^{\mathsf T})^{-1}-a-\mathsf E^{(2)}[\mathsf H_1^{\mathsf T}\mathbf J],} \\[3pt]
\FitMath{\mathsf A=\mathsf C_V^{\mathsf T}(\mathsf G^{(2)}_V+\Delta_k)\mathsf C_V,} \\[3pt]
\FitMath{C^{(k)}(Y)=(\chi_Y\mathsf A\chi_Y)^{-1}.}
\end{gather*}

At the pair $(U_{k+1},J(U_{k+1}))$ of the step, one has $\mathsf A=\mathsf C^{\mathsf T}\Delta^{(k)}\mathsf C$. In this identification (1c.1), the elimination and the gauge-fixing term are obtained directly from the definitions. The only step of substance is the quadratic term of the action of the background, regarded as a function of its datum. This is Ba{\l}aban's expansion \textbf{102}, 277, (174), together with \textbf{109}, (2.4)--(2.8), (2.11). There the derivation is carried out on the whole lattice, using Ba{\l}aban's own rows and his averages. For $\mathcal M$ and the straightened rows, on the sequence of a label, the expansion is obtained from the system with a multiplier of (1b). That multiplier coincides with the gradient of the action of the background with respect to its datum, as one sees by differentiating the constraint and pairing the result with the Euler--Lagrange equation. A single differentiation of the system yields the derivative of the multiplier. For the critical pair, the operator $\mathsf G_1^{-1}$ equals the sum of the Hessian of the Lagrangian, the gauge term, and the mass term carried by the rows. Hence this derivative is obtained by applying $(Q\mathsf G_1Q^{\mathsf T})^{-1}-a$ to the datum. When the datum moves along a curve contained in the fibre of the one-step average, it picks up a second-order term, and this term is determined by the quadratic part of the one-step constraint. Paired with the gradient, this term yields $-\mathsf E^{(2)}[\mathsf H_1^{\mathsf T}\mathbf J]$. Ba{\l}aban also gives a second expression for the first two terms of $\Delta_k$, namely $\mathsf H_1^{\mathsf T}[\mathsf G_1^{-1}-DRD^{\mathsf T}-Q^{\mathsf T}aQ]\mathsf H_1$. It differs from them by $-\mathsf H_1^{\mathsf T}DRD^{\mathsf T}\mathsf H_1$, a term that vanishes on shell. The bounds and the positivity in statement 6 concern the form displayed above.

\emph{Walks.} A \emph{cell} is a union of $O(1)$ cubes of side $O(M_1\xi_j)$ lying in two consecutive layers. In a \emph{walk} $\omega=(X_0,\dots,X_m)$, any two consecutive cells overlap. In a \emph{tree walk}, each $X_i$ overlaps an earlier cell, called its parent, and the cells are enumerated depth first. With arbitrary parent maps, one sequence of length $m$ is counted as many as $m!$ times. The \emph{hull} $[\![\omega]\!]$ enlarges the walk by $O(M_1)$ local units and is not cut off at $\Omega_0$. The \emph{carrier} is $[\omega]=[\![\omega]\!]\cap\Omega_0$. Finally, $d_{\mathbf B}(\omega;y,y')$ is the smallest weighted length of a tree that passes through $y,y'$ and meets the cells $X_1,\dots,X_m$. The tree is not required to meet $X_0$.

Several letters have a different meaning in the sub-parts 1c.1--1c.10, as follows. In 1c.1--1c.5, in all sections of 1c.8, and in the section of 1c.9 that expands $\mathsf A$ over sets of cubes, the cube size of the walks is denoted $\mathfrak M$. In those places $M_1$ denotes $M_1'\ge R_1M_1$, which is the second cube size introduced in the Notation, the larger one. Within 1c.5 and 1c.9, the product $R_1M_1$ appears as $\rho$ and as $\mathsf w_0$. The cubes $\square$ that carry parameters (statement 7) have side $\ell_{\mathsf P}$. This side is denoted $w$ in 1c.9, and in 1c.10 the cubes themselves are denoted $\Delta$. In 1c.2, 1c.4 and 1c.10, $\Gamma_j$ is the set of rows of level $j$. In 1c.6, 1c.7 it also contains the inward bonds, and the rows are its bonds that are not inward. In 1c.10, $Q$ stands for $\mathsf mQ$, and slanted $H,H_1$ stand for $\mathsf H\mathsf m^{-1},\mathsf H_1\mathsf m^{-1}$. In the section that treats the averaging operators at complex configurations, the operator $Q^{\rm st}_k(\mathbf U)=\eta\,\mathbf Q(W_{k-1})\cdots\mathbf Q(W_0)$, where $\mathbf Q(W)$ denotes the linear part of a single step, stands everywhere for the linear part of the straightened average, namely the $Q^\sharp_k(\mathbf U)$ of \ensuremath{\S}3 of the Notation; it is not the straight average defined in that section of the Notation.

In comparison with the Notation:
\begin{itemize}
\item $R_1$ denotes a number (as in \textbf{116}, p. 3) and not a region $R_t$.
\item A cell means a cell of a walk, as opposed to a level-$k$ cell of \ensuremath{\S}9.
\item The hull of a walk differs from the hull in a coarser partition.
\item $\mathfrak G$ denotes an operator, not the gauge-fixing factor $\mathfrak G_k$.
\item $\mathsf A$ denotes the fluctuation operator; in \ensuremath{\S}5 of the Notation the same operator is written $\mathsf K$.
\item $\Delta(y)$ is the block $B^j(y)$, regarded as a cube.
\item $s(\square)$ is a parameter, not a depth.
\item $\bar\kappa$ denotes, as in \textbf{116}, the exponent that enters the radius of the parameters, and not the rate at which the large-field factor decays. In that paper the exponent is written $\kappa_1$; it is the $\bar\kappa$ that appears in the Notation's order of the constants.
\item $\gamma,\gamma^{\rm pos}$ are lower bounds, not ceilings on couplings.
\item $\gamma_0$ is the radius of the current, which is also written $\alpha_0$ in \ensuremath{\S}11.
\item $\lambda$ denotes a site function (an infinitesimal gauge transformation); in statement 8 and in mechanisms (j), (k), however, it denotes a multiplier, that is, a function on rows.
\item $X$ denotes a localisation domain in statement 0, $X_i$ is a cell, $X$ in (I) is a function on rows, and $X_\chi$ is a gauge direction.
\item $A'$ denotes the complex factor of the triple $(U,A',\mathbf J)$; in the system of clause 1 of the theorem of 1c.10, however, it denotes the unknown $H_1(s)B'+\mathbf A$.
\item $S$ denotes, in statement 6, a set of coarse cubes and, in (II), an operator; $\mathsf S$ denotes the operator $1-DG'RD^{\mathsf T}$.
\end{itemize}

In 1c.10 the following further conventions hold:
\begin{itemize}
\item $\theta_0$ denotes the decay rate with respect to $d_{\mathbf B}$; it is not one of the members of the sequence $\theta$ of \ensuremath{\S}6.
\item $\mathsf r=v[\mathbf H]$ is a function on the nodes with values in $G$, and not an isometry.
\item $\beta$ and $\beta_\natural$ denote a small real datum, not the bare inverse coupling.
\item $\omega(\cdot)$ is the map that sends a field to its multiplier; it is not a walk.
\item $\mathsf E_1=G_1DRD^{\mathsf T}G_1$ denotes an operator rather than an event.
\item $P_1(s)$ is the operator built from the walk data of $D(1-R)D^{\mathsf T}$, and neither a region nor a momentum.
\end{itemize}

\textbf{(ii) The principal statements.} All constants depend only on $L$ (and on H{\"o}lder exponents). There are two exceptions, and neither enters any localised bound, rate, weight or sum of weights. The first is the bound of statement 7, which serves only for convergence. The second is the entrywise bound of mechanism (i), which makes the family normal. The numbers $\gamma$, $\gamma^{\rm pos}$, $\delta_0$, $\delta_*$, $\delta'$ are positive constants of this kind.

In statement 2 the side of the torus, measured in top units, is at least a small fixed multiple of $R_1M_1$; in the sentence of that statement which concerns a single lattice, it is at least $C$. From statement 4 onwards the side is at least $CM_1$. In 5, the single assignment of terms valid for all tori needs room for three cubes of side $O(M_1)$ in each direction. Statements 7, 8 need room for three cubes of side $\ell_{\mathsf P}$. These are lower bounds on the volume, and no constant depends on the volume except the bound of statement 7 just mentioned.

In what follows:
\begin{itemize}
\item $x\in\Delta(y)$, and the sources $f$ lie near $\Delta(y')$;
\item $\mathsf e=e^{-\delta_0d_{\mathbf B}(y,y')}$, up to a fixed fraction of the rate;
\item $\|\cdot\|_\varepsilon$ denotes a H{\"o}lder seminorm;
\item kernels are taken with respect to the counting measure;
\item on $1$-forms, the covariant Laplacian is $\Delta_{\mathbf U}=\sum_\nu\nabla_\nu^{\mathsf T}\nabla_\nu$, where $\nabla_\nu$ are the covariant differences of the components;
\item $\Delta=\Delta_{\mathbf 1}$ is the full vector Laplacian $\partial^*\partial+\partial\partial^*$, with $\partial=D_{\mathbf 1}$. It acts componentwise and is not the $\Delta=\Delta^{\rm W}_U$ of the Notation, \ensuremath{\S}5.
\end{itemize}

\textbf{0. The domains of analyticity are contained in the classes.} Consider a pair in the domain of an activity on $X\in\mathcal D_j$. Such a pair satisfies $|\partial U-\mathbf 1|,|\partial\mathbf U-\mathbf 1|<\alpha_0\eta^2$, $|A'|,|\nabla_UA'|<\alpha_1$ and $|\mathbf J|<\gamma_0$, and on cubes of side $O(LM)$ it admits real gauges with $|A|,|\nabla A|<O(LM)\alpha_0$. It is a regular pair for the tower of regions obtained from $X$ by erosion (1b.8), and the radii with which it is regular are

\begin{gather*}
\FitMath{\big(O(LM)\,\alpha_0,\,\alpha_1,\,\gamma_0\big).}
\end{gather*}

Let a label be produced by steps in which each region $\Omega_n$, $n\ge1$, is a union of $LM$-cubes of level $n$. By a \emph{layered domain} we mean a domain defined by the layered bounds of (2a), \emph{Structures}, whose radii do not exceed the running radii of the level of each layer. Only the following of these bounds enter. On the layer $\mathcal L_n$ of the label, $n\ge1$, one has $|\partial U-\mathbf 1|,|\partial\mathbf U-\mathbf 1|<\alpha_{0,n}\eta^2\xi_n^{-2}$, $|\mathbf J|<\alpha_{0,n}\xi_n^{-3}$ and $\xi_n|A'|,\xi_n^2|\nabla_UA'|<\alpha_{1,n}$. The cubes of level $n$ of the domain are the closed cubes of side $CM\xi_n$ that are contained in $\mathcal L_n\cup\mathcal L_{n+1}$ and meet $\Omega_{n+1}^c$; on the part in $X$ of each of them there is a real gauge with $\xi_n|A|,\xi_n^2|\nabla A|<CM\alpha_{0,n}$. Then on $X\cap\Omega_1$ a pair of a layered domain belongs to the local class of the \emph{label sequence}, which is the sequence formed by the $\Omega_n$ of the label. The radii are $O(M)\max_n\alpha_{0,n}$ and $\max_n\alpha_{1,n}$, together with $\max_n\alpha_{0,n}$ for the current. In that class the gauge clause is derived for those cubes of level-$n$ size which meet $\mathcal L_n$. For $n=0$, on the cubes inside $\mathcal L_1$, it is derived from the gauge of a cube of level $1$, which can be used at the index $0$ since $\xi_0<\xi_1$. The same remark applies to a cube of level-$n$ size inside $\mathcal L_{n+1}$, because some cube of level $n+1$ of the domain contains it. This smaller cube is contained in $\Omega_{n+1}$, intersects $\Omega_{n+2}^c$, and its side is below $M\xi_{n+1}$. The derivation of the clause on the cubes meeting $\mathcal L_n$ uses no other property of a cube; hence repeating the same steps with the grids of level $n+1$ produces the larger cube. The factors $O(LM)$, $O(M)$ are the cost of the gauge clause, since on a cube of side $O(M)$ a gauge has $|A|\le O(M)\alpha_0$. The field strengths are fixed after $M$.

\textbf{1. The flat theorem.} Let $\mathbf U=\mathbf 1$, $\mathbf J=0$, and let $Q$, $Q'$ be the straight averages taken over the centred blocks. Then on every separated sequence

\begin{gather*}
\FitMath{|\mathsf G_0f|,|\nabla\mathsf G_0f|,|\mathsf G_0\nabla^*f|,|\Delta\mathsf G_0f|\,(x)\le C[\xi_y^2,\xi_y,\xi_y,1]\,\mathsf e\,|f|,} \\[3pt]
\FitMath{|\nabla\mathsf G_0\nabla^*f(x)|\le C_\varepsilon\,\mathsf e\,\big(\xi_{y'}^\varepsilon\|f\|_\varepsilon+|f|\big),}
\end{gather*}

together with matching bounds in H{\"o}lder and $L^2$ norms. The same bounds hold for $G'$, and

\begin{gather*}
\FitMath{|(Q'G'^2Q'^*)^{-1}(y,y')|\le C\xi_y^{-4}\mathsf e,} \\[3pt]
\FitMath{|(Q\mathsf G_0Q^*)^{-1}(c,c')|\le C\xi_c^{-2}\mathsf e,} \\[3pt]
\FitMath{|\mathsf H(b,c)|,|\nabla\mathsf H(b,c)|\le C[1,\xi_y^{-1}]\mathsf e,}
\end{gather*}


\begin{gather*}
\FitMath{Q\mathsf G_0Q^*\ge\gamma\,\xi_c^2,} \\[3pt]
\FitMath{\big\langle B,\big[(Q\mathsf G_0Q^*)^{-1}-a\big]B\big\rangle\,\ge\,c\,L^{-5}\|B\|^2}
\end{gather*}

for every $B$ with three properties: it lives on the bonds of $\mathbb T^{(k)}$ that touch a union of $L$-blocks lying well inside $\Omega_k$; its straight one-step average vanishes; and it vanishes on the axial trees of the blocks. Each block has one such tree. It consists of the contours joining the site $y\in\mathbb T^{(k+1)}$ of the block to each of its points $x$, along which the coordinates are changed successively, always in the same fixed order. The gauge so defined is the strict axial gauge, which is a hard gauge fixing. It is not the block axial gauge (Notation, \ensuremath{\S}3), where the means $V(y,x)$ over contour families are all equal to $\mathbf 1$. For the plain one-step constraint, in that gauge or on trees, the lower bound is proved again in 1c.8. If the gauge condition is dropped, the form has null directions.

\textbf{2. Real backgrounds with Ba{\l}aban's average.} Let $Q,Q'$ be formed with one contour per site and moved to centred blocks by a translation. Let the Hessian be the compressed one. Let $U$ be real, with local gauges and with small plaquettes; for $\mathsf G_0$ this includes the plaquettes that stick out of $\Omega_0$. Then $\Delta_U+Q'^*aQ'$, $Q'G'^2Q'^*$ and $\mathsf G_0^{-1}$ are positive. The operators $G',\nabla G',\mathsf G_0,\nabla\mathsf G_0,(Q'G'^2Q'^*)^{-1}$ satisfy the bounds of 1 for every $\Omega_0$, including a collar of $\Omega_1$ and a concentric cube. On a single lattice, and for $U$ with $|\partial U-\mathbf 1|,|D_U^*\partial U|<\alpha\eta^2$, where $\alpha$ is small, also $Q\mathsf G_0Q^*\ge\gamma>0$, and $\mathsf G_0Q^*(Q\mathsf G_0Q^*)^{-1}$ is a right inverse of the linearised constraint; this right inverse and its gradient decay like $e^{-\delta_*|y-y'|}$. The rows of $\mathcal M$ agree with these rows up to a transport and an error $O(\alpha_0)$ in the row and column norms, uniformly in the level. Statements 3--8 concern different operators, built from the rows of $\mathcal M$ and the Hessian of the inner plaquettes. At $\mathbf 1$ the two kinds of rows coincide.

\textbf{3. The averages and the gauge identities.} The operators $Q,Q'$ are local and holomorphic in $\mathbf U$. Their row sums are $\le1+C(\alpha_0+\alpha_1)$ and their entries are $\le C(\eta/\xi_c)^4$, independently of the number of levels. At each row the non-linear average is holomorphic in $(A',A)$ on a polydisc $\sup|A'|,\sup|A|<r_L$, where $r_L$ is determined by $L$ and measured in units of the row's level. For the functional-derivative bounds, work on a single lattice of depth $k$ in its top units. Let $b$ be a bare bond and $c$ a row of level $k$, so $\xi_c=1$, and write $\delta/\delta A_b=\eta^{-4}\partial/\partial A(b)$. In the first two inequalities below, $\mathcal Q=\mathcal Q^{\rm st}_k$. Let $\mathcal Q^{\rm st}_n$ be the $n$-fold straightened average, and let $v_n$ be its accumulated site factor. Then on this polydisc

\begin{gather*}
\FitMath{\Big|\frac{\delta\mathcal Q(A)(c)}{\delta A_b}\Big|\le1+C(\alpha_0+\sup|A'|+\sup|A|),} \\[3pt]
\FitMath{\Big|\frac{\delta(\mathcal Q-\mathsf mQ)(A)(c)}{\delta A_b}\Big|\le C\sup|A|,} \\[3pt]
\FitMath{\mathcal M^n(e^{i\eta A}\mathbf U)=\big(e^{i\mathcal Q^{\rm st}_n(A)}\mathcal M^n\mathbf U\big)^{v_n}.}
\end{gather*}

Let $\bar D$ be the covariant difference of $\mathcal M^j\mathbf U$ on $\mathbb T^{(j)}$. At a row of level $j$, let $Q'\lambda$ denote the $j$-fold average taken at both ends of the row. For an outward row, the outer end is a top of level $j-1$, not a top of level $j$. Then

\begin{gather*}
\FitMath{QD\lambda=\bar DQ'\lambda,} \\[3pt]
\FitMath{\Delta^{\rm W}_{\mathbf U}-\Delta'_\pi[J(\mathbf U)]=\mathsf S^{\mathsf T}\Delta^{\rm W}_{\mathbf U}\mathsf S,} \\[3pt]
\FitMath{\mathsf G_{(1)}DR=DG'R+\mathsf G_{(1)}\Delta'_\pi[\mathbf J-J(\mathbf U)]DG'R.}
\end{gather*}

On shell, the defect is zero, $RD^{\mathsf T}\mathsf H_{(1)}=0$, and $\mathsf HB$ is a critical point of $\frac12(A,\Delta^{\rm W}_{\mathbf U}A)$ on $\{QA=B,\ RD^{\mathsf T}A=0\}$, with critical value $\frac12(B,[(Q\mathsf GQ^{\mathsf T})^{-1}-a]B)$.

\textbf{4. The six propagators at regular pairs.} The ceiling on the radii $\alpha_0,\alpha_1,\gamma_0$ is fixed before $M_1$ and is \textbf{independent of $M_1$}. At every regular pair all the inverses exist, they are holomorphic in $(\mathbf U,\mathbf J)$, and they are real at real pairs. At real pairs $\mathsf G_0,\mathsf G,\mathsf G_1>0$, both on shell and off shell. The six propagators and $(Q'G'^2Q'^{\mathsf T})^{-1}$ satisfy the upper bounds of 1 with covariant differences, and $\mathsf G$ and $\mathsf G_1$ satisfy those of $\mathsf G_0$. The exception is the bound on $\Delta_{\mathbf U}\mathsf G_{(1)}$, and the H{\"o}lder quotients are taken only over $|x-x'|\le\xi_y$. For $G'$ at complex pairs the bounds reduce to the supremum bounds on $G'f$ and $DG'f$. Moreover $|D(1-R)D^{\mathsf T}(b,b')|\le C(\eta/\xi_{y'})^4\xi_y^{-2}\mathsf e$.

\textbf{5. Walk expansions.} Let $\mathsf O$ be any of the following operators:
\begin{itemize}
\item one of the six propagators;
\item $1-R$ or one of its derivatives;
\item $\mathfrak G$;
\item Ba{\l}aban's $H_0,\tilde G$, namely $\mathsf H\mathsf m^{-1}$ and $\mathsf G-\mathsf HQ\mathsf G$ at $\mathbf J=0$;
\item $C^{(k)}(Y)$ or $\chi_Y\mathsf A\chi_Y$, these two taken at the pairs that are regular for the longer sequence $(\Omega_j)_{j\le k+1}$, where $M_1$ is larger and the ceiling on the radii is smaller than in 4.
\end{itemize}

For $\tilde G$ the indices are the walks together with the pairs of walks in which the second walk starts at a cell that a row joins to some cell of the first walk, its cells being placed after those of the first. In the display, $|\omega|$ is then the total length of the two walks, and the step that joins them contributes no small factor. Each such operator is then a sum $\sum_\omega\mathsf O_\omega$, where, in the norms in which $\mathsf O$ is bounded,

\begin{gather*}
\FitMath{\|\mathsf O_\omega\|\le(\text{bound of }\mathsf O)\cdot(CM_1^{-1/2})^{|\omega|}M_1^{-|\omega|/2},} \\[3pt]
\FitMath{d_{\mathbf B}(y,y')\,\text{replaced by}\,d_{\mathbf B}(\omega;y,y').}
\end{gather*}

The first factor makes the sum over the walks converge. The second factor is left over, and it absorbs the factor $|\omega|+1$ of mechanism (i). The bounds of $\mathsf O$ are as follows:
\begin{itemize}
\item For $\mathfrak G$ and $\tilde G$ there are exactly four: the bounds on $|\mathsf Of|$, $|\nabla\mathsf Of|$, $|\mathsf O\nabla^{\mathsf T}f|$ and on the H{\"o}lder seminorm of $\nabla\mathsf Of$ that 4 provides for $\mathsf G$, $\mathsf G_1$. For $\tilde G$ they hold with a smaller fraction of $\delta_0$.
\item For $1-R$ and for its derivatives $D(1-R)$, $(1-R)D^{\mathsf T}$, $D(1-R)D^{\mathsf T}$, the relevant bounds are the kernel bounds $C(\eta/\xi_{y'})^4[1,\xi_y^{-1},\xi_y^{-1},\xi_y^{-2}]\,\mathsf e$. The last of these is displayed in 4.
\item For $\chi_Y\mathsf A\chi_Y$, the bound is the kernel bound that 6 states for $C^{(k)}(Y)$.
\end{itemize}

Each $\mathsf O_\omega$ has the following properties:
\begin{itemize}
\item it is holomorphic on the local class on $[\omega]$;
\item its dependence on the regions is through their traces on $[\![\omega]\!]$ alone;
\item it is gauge covariant and real at real pairs;
\item it is \textbf{the same function} for any two sequences, tori or cutoffs that a translation in $M_1\xi_k\mathbb Z^4$ relates on the hull. For $C^{(k)}(Y)$, $\chi_Y\mathsf A\chi_Y$ the translation is in $M_1\xi_{k+1}\mathbb Z^4$ and must also carry $\Omega_{k+1}$ and $Y$.
\end{itemize}

The operators are covariant under every isometry that preserves $\mathbb T^{(k)}$ and the regions, with the twist $-\operatorname{Ad}\mathbf U(b)^{-1}A(b)$ at reversed bonds. For $C^{(k)}(Y)$, $\chi_Y\mathsf A\chi_Y$ the isometry must also preserve $\mathbb T^{(k+1)}$ and $\Omega_{k+1}$, and $Y$ is carried along. Pivot bonds go to pivot bonds, since each pivot bond is the middle bond of its path. For single terms the covariance holds only under those isometries which in addition preserve the $M_1$-cubes of the top level. That level is $k+1$ for the covariance, and in the strict axial gauge the isometry must also preserve the axial trees. This covariance holds at the pairs whose image again lies in the class.


\paragraph{The additional operator estimates.}
The fluctuation matrix, its conditional inverses, the independent
complex decoupling and the nonlinear response are the additional
statements numbered 6--8. Their estimates are obtained from the walk
data and the core bounds. The response construction is given below;
the actual Gaussian comparison is in
Section~\ref{ch:localized-gaussian-comparison}. The substitutions use the
hereditary room estimate in Section~\ref{sec:substitution-room}.
The layered and nonzero-current versions retain the same weighted
operator bounds and the assigned determining structures.

\subsection*{The response and its independent current}

The decoupled response is formed from the operators of the walk
expansions. Put $\widehat Q=\mathsf mQ$, where $\mathsf m$ is
the length of a row, and write the nonlinear straightened
average as $\mathcal Q(A)=\widehat QA+\mathcal C(A)$.
System (R9) is then
\begin{gather*}
 X=\mathcal C(A_\bullet-H(s)X),\qquad
 A_\bullet=H_1(s)B'+\mathbf A,\\
 \mathbf A+\mathfrak G(s)\mathsf W_s(A_\bullet)=0,
 \qquad \mathbf H=A_\bullet-H(s)X,
\end{gather*}
where $H=\mathsf H\mathsf m^{-1}$ and
$H_1=\mathsf H_1\mathsf m^{-1}$. The reduced expression
$\mathsf W_s$ is obtained by cancelling the quadratic part
before the parameters are inserted. At arbitrary $s$ it is
this expression that defines the map; identities between the
undecoupled operators are not imposed on their decoupled sums.

For a positive weight $\varpi$ with
$\varpi(y)\le e^{\theta d_{\mathbf B}(y,y')}\varpi(y')$,
use the norms
\[
 \|A\|_\varpi=\sup_y\varpi(y)
       \max\{\xi_y|A|,\xi_y^2|\nabla_{\mathbf U}A|\},
 \qquad
 |f|^\sharp_\varpi=\sup_y\varpi(y)\xi_y^3|f|.
\]
Each decoupled operator has norm at most
$\widehat b=Ce^{16\bar\kappa}$ in these norms. On a ball
$\|A_i\|\le r<c/\widehat b$, the nonlinear expression obeys
\[
 |\mathsf W_s(A_1)-\mathsf W_s(A_2)|^\sharp_\varpi
       \le C\widehat b^2r\|A_1-A_2\|_\varpi.
\]
The fixed-point Lipschitz constant is bounded by
$C\widehat b^4\sup|B'|$. Thus $\sup|B'|\le c\widehat b^{-4}$
gives the same contraction in every such weight. Its solution
is holomorphic in the parameters, the datum, the complex
potential and the independent current, and
\[
 \xi_y|\mathbf H|+\xi_y^2|\nabla_{\mathbf U}\mathbf H|
       \le C\widehat b\sup|B'|.
\]
Choosing an exponential weight gives the first-response decay.
The same contraction, by uniqueness, gives consistency and
splitting over the wall-components. It also transports the
solution whenever the assigned walk data are transported.

The current is a separate entry of the pair. Its increment is
formed from the multiplier, so that the second-order part of
the gradient enters only through
$(D^{\mathsf T}DA,He_c)=(DA,DHe_c)$. It is therefore bounded
in $|\cdot|^\sharp_\varpi$ by the same first-order response
norms. No level-independent bound on the full Laplacian of an
arbitrary decoupled response is needed.

For completeness, a second system provides the other
second-order bound at physical parameters. Let
\[
 S=(\widehat Q\mathsf G\widehat Q^{\mathsf T})^{-1},
 \quad H_0=\mathsf G_0\widehat Q^{\mathsf T}
       (\widehat Q\mathsf G_0\widehat Q^{\mathsf T})^{-1},
 \quad\widetilde G=\mathsf G_0-H_0\widehat Q\mathsf G_0,
\]
and let $\mathsf a$ satisfy
$\widehat Q^{\mathsf T}\mathsf a\widehat Q=Q^{\mathsf T}aQ$.
In the action chart set
$\mathcal N(T)=F'(T)-F'(0)-\Delta^{\rm W}_{\mathbf U}T$.
The reduced system for $(T,\lambda)$ is
\begin{align*}
 T={}&H_0(s)(B'-\mathcal C(T))
       -\widetilde G(s)
       [\mathcal N(T)-\mathcal C_T'{}^{\mathsf T}\lambda],\\
 \lambda={}&H(s)^{\mathsf T}\mathbf J
       +(S(s)-\mathsf a)(B'-\mathcal C(T))\\
 &\quad+H(s)^{\mathsf T}
       [\mathcal N(T)-\mathcal C_T'{}^{\mathsf T}\lambda].
\end{align*}
The terms quadratic in the unknowns are contractions for
$\widehat b^2\sup|B'|\ll1$. The terms linear in $T$ that
contain the current require the additional ceiling
$\widehat b^2\gamma_{\rm cur}\ll1$, where $\gamma_{\rm cur}$ is the
ceiling on the independent current in this response problem. With
these choices the system has
one small holomorphic solution in
$(T,\lambda-H(s)^{\mathsf T}\mathbf J)$, with the same weighted
Lipschitz estimates. This second system is used on its stated
global domains and, in local form, on strip sequences. It is
not substituted for the first system on arbitrary local cores.

At $s=1$ over a critical real background, the first solution
satisfies
\[
 \mathcal Q(\mathbf H)=B',\qquad
 RD^{\mathsf T}\mathbf H=0,\qquad
 F'(\mathbf H)=K_{\mathbf H}^{\mathsf T}\lambda.
\]
For real data near zero, invariance gives all missing gauge
directions transverse to the slice; hence the slice equation
has a multiplier for the averages alone. Holomorphy extends
the identities to the full polydisc. They characterize the
background response there. On the common physical polydisc,
the second system follows from these equations and uniqueness
identifies its field with $\mathbf H$. The first system then
gives the bound on $DD^{\mathsf T}\mathbf H$, and the second
gives the bound on $D^{\mathsf T}D\mathbf H$, both in local
units. Away from critical physical pairs the two systems
define separate functions. In particular no constraint
identity is asserted at arbitrary complex decoupling parameters.

\textbf{(iii) The strategy.}

\emph{The chain.} The parts have the following content.
\begin{itemize}
\item 1c.1 (the operators and the statements about them) fixes the operators, the classes and the form of the statements. It counts the walks and derives statement 0. It also shows that the identification of $\mathsf A$ at the step reduces to Ba{\l}aban's expansion of the background action.
\item 1c.2 (the flat case, I) states the flat theorem and compares it clause by clause with the published statements. It also contains the exact algebra at $\mathbf 1$: existence, identities with boundary terms, and the independence of $R$ and $\mathsf H$ from $\Omega_0$.
\item In 1c.3 (the flat case, II) the one-scale operator is treated uniformly in the spacing; then come $G'$, $(Q'G'^2Q'^*)^{-1}$ and $R$.
\item In 1c.4 (the flat case, III) the local two-scale operators are treated, followed by $\mathsf G_0$, $(Q\mathsf G_0Q^*)^{-1}$ and $\mathsf H$. With this, statement 1 is complete.
\item 1c.5 (Ba{\l}aban's theorems in a background) yields statement 2. On the whole lattice it comes from the published theorems; on sequences it comes from the flat case, by a shortened form of mechanism (f).
\item 1c.6 (the averaging operators at complex pairs; the algebraic identities) yields statement 3.
\item 1c.7 (the six propagators) uses mechanisms (f), (g) to derive, from the flat bounds and 1c.6, statement 4 and the expansions of 5.
\item In 1c.8 (the covariance; positivity) these two mechanisms are applied to a larger system. This yields statement 6, except for the expansion of $\mathsf A$ over sets of coarse cubes.
\item In 1c.9 (walk expansions as data) an expansion is regarded as a datum enjoying the short list of properties of 7. Statement 7 is deduced from these properties together with the geometry. Moreover, 1c.9 \emph{makes} terms that are real, covariant and identical on all tori with room; these properties do not follow from the list.
\item 1c.10 (the decoupled response) takes the operators as data and solves the systems (I), (II). Next it inserts the local maps and operators of 1c.6--1c.9; for $\mathfrak G$ and $D(1-R)D^{\mathsf T}$ the inserted data are those named at the end of 7. Last, it treats the case $s\equiv1$.
\end{itemize}

\emph{The order of (1b) and (1c), which involves no circle.} Statements 1--7 refer to no minimiser, since the classes are defined by inequalities. The results are obtained in the following order:
\begin{itemize}
\item the geometry and the real calculus of $\mathcal M$ (1b.1--1b.2);
\item the flat case;
\item the lemmas of 1b.3 that are geometric or algebraic (folding, cells and rows, the algebra of $Q',Q,R$, sums over tops, weights for a general metric), and none of them uses any bound on a propagator;
\item 1c.5;
\item the rest of 1b.3; then existence, regularity (which provides the backgrounds of (1b) with their local gauges) and continuation (1b.4--1b.6);
\item 1c.6--1c.9. Their inputs are the flat case, the alignment of the blocks in 1c.5, and 1b.1--1b.2, together with, besides these, geometric facts involving neither propagators nor minimisers (meets of sequences; the covariance of $\mathcal M$ with respect to the isometries of the lattice);
\item the decay of the response, and the membership in the classes of the backgrounds of (1b) (1b.7--1b.9);
\item finally, the identification at $s\equiv1$.
\end{itemize}

\emph{Mechanisms.}

(a) \textbf{The metric matches the mass.} The form satisfies $\langle\lambda,(\Delta+Q'^*aQ')\lambda\rangle\ge\gamma\sum_j\xi_j^{-2}\|\lambda\|^2_{\mathcal L_j}$, and on each layer $d_{\mathbf B}$ has slope $\xi_j^{-1}$. A single conjugation by $e^{\varsigma\varphi}$, with $\varphi$ Lipschitz with respect to $d_{\mathbf B}$ (Combes--Thomas), therefore serves for all layers at once. It costs $O(\varsigma L^2)$, and this cost stays below $\gamma$ under a ceiling on the rate, $\varsigma\le c\gamma L^{-2}$, where $L$ is chosen first. Every layer that has to be traversed in full costs the gap, $d_{\mathbf B}(y,y')\ge R_1M_1(|j-j'|-1)_+$. Let a rate $\tau>0$ and a power $q\ge0$ be given, and choose the product $R_1M_1$ after them so that $\tau R_1M_1\ge C(1+q)\log L$. Then

\begin{gather*}
\FitMath{\sum_{y'}e^{-\tau d_{\mathbf B}(y,y')}\le C(\tau,L),} \\[3pt]
\FitMath{(\xi_{y'}/\xi_y)^q\le L^qe^{\tau d_{\mathbf B}(y,y')}:}
\end{gather*}

the factor $e^{-c\tau R_1M_1}$ for each level outweighs the growth $L^4$ in the number of sites. Conjugation and walk expansions do different work. Conjugation yields decay of the inverse as a whole, with a rate given by a fraction of the gap, and it requires no cube size to be large. A walk expansion gives, in addition, a sum of terms, each of which depends on the background and on the regions only near its own walk. This property is what statements 5--8, the comparison of tori and (1d) make use of. Both methods are used. Conjugation serves for the scalar, one-scale and local operators and for the inverses on tops and rows. Walks serve for the assembly, namely of $\mathsf G_0$ in (d) and of the propagators in a background in (f). In \textbf{95}, p. 36, Ba{\l}aban describes conjugation as ``probably the simplest proof'' of the decay; he builds a walk representation instead, since ``such representations will be our basic tool in investigation of more complicated operators''.

(b) \textbf{No Dirichlet condition on a lattice finer than the blocks.} Because of the gap, $\partial\Omega_0$ lies in the bare layer $\mathcal L_0$, and the boundary cubes and their cut-off collars are chosen inside its all-bare collar. In the local model on such a cube $Q_\square=Q'_\square=1$, hence $R_\square=0$; the global $R$ need not vanish throughout $\mathcal L_0$ near its interface with $\Omega_1$. The local operator $G_\square$, which inverts this all-bare form (see (d)), obeys $a_0\le\eta^2G_\square^{-1}\le a_0+16$. Finite range together with a mass comparable with the norm gives decay for an arbitrary set of bonds. The following examples show that this placement of the boundary matters. Let $G$ be a one-scale operator (spacing $\eta$, unit blocks) whose order is $\ge1$, with Dirichlet conditions imposed outside a union of \emph{whole} blocks, so that the lattice is finer than the blocks and no block is cut. Two examples, a flat face and a re-entrant edge, rule out constants uniform in the order, by powers of $\eta^{-1}$. This applies to the $L^2$ bounds on second differences, to the bounds on $\Delta G$, to those on $\nabla G$ and to its H{\"o}lder norm. It applies equally to the pointwise bound and to the H{\"o}lder bound involving $\nabla G\nabla^*$, taken in the form of statement 1, the powers being those that occur for $\nabla G$ and for its H{\"o}lder norm. Indeed, on the region every bounded $f$ supported in a block can be written as $\nabla^*F$ with $F=\nabla\Delta^{-1}f$, where $\Delta^{-1}$ denotes the free inverse on $\eta\mathbb Z^4$ (on a torus, the inverse on functions of mean zero, after $f$ has been modified on a block lying outside the region); moreover $|F|,\|F\|_\varepsilon\le C_\varepsilon\sup|f|$ uniformly in $\eta$, and the sum over the blocks is controlled by the exponential factor in the bound. For the $L^2$ bound involving $\nabla G\nabla^*$, by contrast, the two examples rule out no uniform constant. On a separated sequence no operator of this kind arises.

(c) \textbf{One scale, two scales.} Uniformity in $\eta$ comes from two sources. The first is the representation $(\Delta+1)^{-1}=\int_0^\infty e^{-t}e^{-t\Delta}dt$, combined with heat-kernel bounds that hold both below and above $t=\eta^2$. The second is a factorisation of $\partial(1-R)\partial^*$ through the unit lattice. At a seam one makes an exact change of gauge $A=A_2+H_jB-\partial H'_j\vartheta$. Here $A_2\in\ker Q_j$, and $B$ is a $1$-form on the coarse lattice whose values on the axial trees of the $L$-blocks of the coarse region are zero. Further, $\vartheta$ is a site function on the coarse lattice, supported in the coarse region, whose mean over every one of those blocks is zero; for these two classes the representation is unique. Finally, $H_j,H'_j$ are the one-scale minimisers, with $H'_j$ the minimiser of $\|\Delta\lambda\|^2$ rather than of $\|\partial\lambda\|^2$. This change of gauge splits the form orthogonally into three parts: a form in $A_2$, a form in $B$ and a form in $\vartheta-\vartheta_1$, where $\vartheta_1$ is linear in $A_2$ and $B$. If the coarse region is all of the torus, the two-scale representation so obtained becomes a recursion $G_j\to G_{j+1}$, where $G_j$ denotes the one-scale propagator of order $j$. It is applied for a single step and is never iterated. In the scalar theory, recursions of this kind add up to a series over the scales, and they served as the method of proof of regularity (\textbf{89}, (2.34)). Here $G_j$ occurs multiplied by other operators. Iteration then leads to long products whose estimation is very difficult, since the gauge conditions are not the same at different scales (a remark of Ba{\l}aban, \textbf{96}, pp. 246--247). For gauge propagators, accordingly, no series over the scales is employed; instead the one-scale $G_j$ has to be known simultaneously for all $j$ (1c.3). Three gaps on the unit lattice take the place of a small parameter. The sharp one is stated for $B$ on a unit lattice that vanishes on the axial trees. In it, $Q_1$ and $\partial_1$ are the straight one-step average and the exterior derivative; the $Q_1(\mathbf 1)$ of the Notation, \ensuremath{\S}3, also contains site terms. The sharp gap reads

\begin{gather*}
\FitMath{L^{2}\sum_c|(Q_1B)(c)|^2+\sum_p|(\partial_1B)(p)|^2\ge cL^{-5}\|B\|^2.}
\end{gather*}


(d) \textbf{Many scales at $\mathbf 1$.} No expansion of $G'$ is required. For $u=G'(f+\nabla^*F)$ the conjugation of (a) yields 

\begin{gather*}
\FitMath{\|\xi^{-1}e^{\varsigma\varphi}u\|+\|e^{\varsigma\varphi}\nabla u\|\le C\gamma^{-1}(\|\xi e^{\varsigma\varphi}f\|+\|e^{\varsigma\varphi}F\|)}
\end{gather*}

, where $\xi$ denotes multiplication by the local scale. From this one passes to local $L^2$ norms, and then to interior estimates at the scale $\xi_y$, with the mass term treated as a source. The walk expansions are obtained afterwards, as data. For $\mathsf G_0$ we take $\sum_\square h_\square^2=1$, with $\zeta_\square=1$ near $\operatorname{supp}h_\square$, together with the $G_\square$ of (c) on periodised cubes. For the cubes of level $0$ that extend beyond $\Omega_0$, the operators of (b) are taken instead. For a cube of level $\ge1$ the periodic two-scale model is retained, irrespective of where $\Omega_0^c$ and the finer layers lie; this is legitimate since across the two seams to $\mathcal L_j\cup\mathcal L_{j+1}$ the function $h_\square$ is cut off. With $P=1-R$ and $\mathcal P=\partial P\partial^*$ one has

\begin{gather*}
\FitMath{\mathsf G_0^{-1}\sum_\square h_\square G_\square h_\square=1-\mathcal R,} \\[3pt]
\FitMath{\mathcal Ph-h\mathcal P_\square=\zeta[\mathcal P_\square,h]+\zeta(\mathcal P-\mathcal P_\square)h+(1-\zeta)\mathcal Ph:}
\end{gather*}

of all these operators, only the non-local $\mathcal P$ cannot be transplanted. Each commutator costs $M_1^{-1}$, because $|\nabla h|\le C(M_1\xi_j)^{-1}$ stands against $|\nabla G_\square|\le C\xi_j$. The change of domain costs $e^{-cM_1}$; this follows from resolvent identities on the two domains, recombined by means of $1-P=R$. The tail is likewise exponentially small, and it is turned into a step $\square\to\square'$ by means of $\sum h_{\square'}^2=1$. The factor $M_1^{-1}$ for each step outweighs the number $C(\tau,L)$ of neighbours.

Lower bounds independent of the number of levels underlie the inverses on tops and rows. The bound $Q'G'^2Q'^*\ge\gamma\xi_y^4$ follows by duality with test functions $\psi_y$ on the blocks, with $Q'\psi_y=\delta_y$, whose second differences are $\le C\xi_y^{-2}$, across the faces as well. The bound $Q\mathsf G_0Q^*\ge\gamma\xi_c^2$, with $\gamma$ independent of $L$, follows from a right inverse $\mathsf E$ for all rows of all levels at once. This right inverse has scale-invariant energy, $\|\nabla\mathsf EB\|^2\le C\|\mathsf m^{-1}B\|^2$. (Choosing on each level test forms biorthogonal to the rows of that level, one obtains $Q\mathsf E_0=1+\mathsf N$. The operator $\mathsf N$ goes up by one level only, and it goes in one direction: from rows that do not cross a seam to rows, one level higher, that do cross one. Hence $\mathsf N^2=0$ and $\mathsf E=\mathsf E_0(1-\mathsf N)$. The weights of $\mathsf N$ are each $\le L^{-4}$ and have sum $\le1$. Set against $\xi_{j+1}^{2}=L^{2}\xi_j^{2}$, they give $\|\mathsf m^{-1}\mathsf N\mathsf m\|\le CL^{-1}$, and from this the constant is independent of $L$.) In 2, where there is a single lattice, the test forms are placed on the face separating the two blocks of a bond; there the share $\theta_0>\frac12$ of its own chains outweighs $1-\theta_0$. Neither matrix is bounded below by a number independent of the scales $\xi$. They are therefore rescaled, to $\Sigma=\xi^{-2}Q'G'^2Q'^*\xi^{-2}$ or $\xi^{-1}Q\mathsf G_0Q^*\xi^{-1}$. A matrix $\Sigma\ge\gamma$ with $|\Sigma|\le Ce^{-\tau d_{\mathbf B}}$ then satisfies $|\Sigma^{-1}|\le2\gamma^{-1}e^{-\tau'd_{\mathbf B}}$, where $\tau'$ is small compared with the spectral gap $\gamma$ (Combes--Thomas, Schur). Thus both inverses decay without any expansion, and $\mathsf H$ decays column by column. A sequence is compared with the same sequence cut down around a cube. The cut-down sequence is again separated and coincides with the given one on the doubled cube, and its $Q\mathsf G_0Q^*$ provides the local operators that enter the walk expansion of $(Q\mathsf G_0Q^*)^{-1}$. The comparison uses Green's identity with a cut-off, by way of $(\partial^*\partial+Q^*aQ)\,\mathsf G_0Q^*=Q^*$, and the gauge term vanishes on the columns.

(e) \textbf{Gauge algebra without Gaussian integrals.} On a pure gauge, take the derivative of the covariance of the straightened average; the site factors then convert point values into block means: $QD=\bar DQ'$. Consequently the rows annihilate $D\ker Q'$. Combine this with the invariance of the action to second order, which is exact for the inner plaquettes; this accounts for that Hessian. The result is $\mathsf G^{-1}D\lambda=D\Delta_{\mathbf U}\lambda-\Delta'_\pi[\mathbf J-J(\mathbf U)]D\lambda$ on $\ker Q'$. The propagators are taken at pairs whose current is an independent variable carrying a bound of its own, as in the spaces of \textbf{109}, p. 262. The reason is that bounds on the plaquettes, on the potential and on the first differences do not control $J(U)$ uniformly in the level: a staggered field satisfies all of them with $|J(U)|\ge\alpha\eta^{-1}$. In \textbf{99}, 389, by contrast, $J=J(U)$, and a condition on second differences is imposed in addition ((3.36), (3.38)). The price of a free current is the defect in 3. Uniformity in the level holds for the following reason. Row norms and column norms are both submultiplicative, and for a single step they lie near $L$ and $L^{-3}$. The relative errors $O(\alpha_0(L^n\eta)^2+\alpha_1L^n\eta)$ of the steps form geometric series that are dominated by the top level.

(f) \textbf{Unfolding and parametrix.} With auxiliary unknowns 

\begin{gather*}
\FitMath{\Psi=(A;v,w,\sigma;p,q,r;\nu)}
\end{gather*}

, the operators $R$, $\Delta'_\pi$, $\mathsf S^{\mathsf T}\Delta^{(2)}\mathsf S$ and the constraint are combined into a single system $\mathbb K$ of finite range. The inverse of this system has the propagators as blocks. Let $\varkappa$ be an $M_1$-cube; these cubes have side $M_1$ in the local unit and tile $\Omega_0$, whereas the $\square$ of (b), (d) are the cubes of the flat expansions. On such a cube, in the gauge of the class, $\mathbb K_\varkappa$ equals the flat system plus $O(\alpha_0+\alpha_1+\gamma_0)$, and this perturbation need only be small compared with flat constants. Let $\theta_\varkappa$ be a partition of unity, let the cut-offs $\tilde h_\varkappa=1$ near $\operatorname{supp}\theta_\varkappa$ be such that $\mathbb K=\mathbb K_\varkappa$ on $\operatorname{supp}\tilde h_\varkappa$, and let $\delta$ denote the rate of the local inverses. Then

\begin{gather*}
\FitMath{\mathbb K\sum_\varkappa\tilde h_\varkappa\mathbb K_\varkappa^{-1}\theta_\varkappa=1-\sum_\varkappa\mathbb R_\varkappa,} \\[3pt]
\FitMath{\mathbb R_\varkappa=-[\mathbb K_\varkappa,\tilde h_\varkappa]\mathbb K_\varkappa^{-1}\theta_\varkappa=O\big(M_1^{-1}e^{-c\delta M_1/L}\big),}
\end{gather*}

where the power comes from the commutator, and the exponential from how far $\operatorname{supp}\theta_\varkappa$ lies from $\operatorname{supp}\nabla\tilde h_\varkappa$. This bound outweighs the $O(1)$ neighbours. Three features keep the perturbation and the commutators from losing a derivative: scaling weights, one component of the source that admits a divergence, and the absence of any derivative on $p$.

(g) \textbf{Trees.} The system $\mathbb K$ depends affinely on the non-local coefficient $\mathsf H^{\mathsf T}\mathbf J$. The part of walk length zero stays in the local operators. The remainder, of size $O(\gamma_0M_1^{-1}e^{-cM_1/L})$, forms a vertex that carries a walk of $\mathsf H$. Thus vertices are small because of $M_1$ and not because of the field, and each term is indexed by a trunk with branches. When the cells are listed depth first, the number of such trees grows geometrically: for each first cell, there are $C^m$ tree walks of length $m$ starting from it. Both this count and the size of $M_1$ needed for the sum over trees to converge coincide with the ones in 1c.9. In 1c.7 and 1c.8 a branch is listed as soon as its cut is reached, and the trunk is resumed afterwards. For the data of 1c.9 the listing first continues at every cell along the trunk (within a branch, along the branch). So the whole trunk comes first, and then the branches come in the reverse order of their cuts, beginning with the last cut and ending with the first. Each of the two orders is a depth-first order. The order of listing changes neither the cells nor their parents, and so it changes neither the hull, nor the tree length, nor the bounds.

(h) \textbf{Multiplier and positivity.} Adding the equation ``$QA=$ fluctuation field'' brings in $(Q\mathsf G_1Q^{\mathsf T})^{-1}$ as a Lagrange multiplier. Consequently $C^{(k)}(Y)$ occurs as a block in the inverse of a system of finite range. In the flat case, positivity follows from $\sum_{p'}|\bar\partial QA(p')|^2\le\eta^4\sum_p|\partial A(p)|^2$, with $\bar\partial$ taken on $\mathbb T^{(k)}$, together with a definite form on double blocks. In a background, a homotopy to $\mathbf 1$ preserves $C_\varkappa>0$ for the local inverses. The eigenvalues $z$ of $\chi_Y\mathsf A\chi_Y\sum h_\varkappa C_\varkappa h_\varkappa$ satisfy $|1-z|\le CM_1^{-1}$.

(i) \textbf{Length against parameters.} A tree of length $t$ comes close to at most $2^4\max\{1,Ct/\ell_{\mathsf P}\}$ cubes $\square$ of the partition of side $\ell_{\mathsf P}$ in statement 7; these are not the cubes $\square$ of (b), (d). Hence $\bar\kappa\#\mathsf P(\omega)\le2^4\bar\kappa+(C\bar\kappa/\ell_{\mathsf P})|\Gamma|_{\mathbf B}$ for every tree $\Gamma$ admitted in $d_{\mathbf B}(\omega;y,y')$, provided that $X_0$, which $\Gamma$ need not meet, lies within $O(1)$ top units of $y$ (no other term enters a functional near $y$, since the left argument of the kernel of $\mathsf O_\omega$ touches $X_0$). Consequently, for $|s(\square)|\le e^{\bar\kappa}$ and $\delta\ge C\bar\kappa/\ell_{\mathsf P}$,

\begin{gather*}
\FitMath{|s^\omega|\,e^{-\delta d_{\mathbf B}(\omega;y,y')}\le e^{2^4\bar\kappa}\,e^{-\bar\kappa\#(\mathsf P(\omega)\setminus\sigma_0)}\,e^{-(\delta-C\bar\kappa/\ell_{\mathsf P})d_{\mathbf B}(y,y')}.}
\end{gather*}

Each cube beyond the first $2^4$ adds to the tree a length $\ge c\ell_{\mathsf P}$. Since $\ell_{\mathsf P}$ is chosen last, $e^{-c\delta\ell_{\mathsf P}}$ outweighs $e^{\bar\kappa}$. The factor $e^{2^4\bar\kappa}$ is left over and cannot be avoided, as a cell at a corner of the grid shows. The middle factor amounts to $e^{-\bar\kappa}$ for each cube that is crossed without a parameter. The estimate is Ba{\l}aban's (\textbf{116}, p. 5, after (1.11)); in his setting as well, the exponential by itself does the paying, and $e^{16\bar\kappa}$ remains. He keeps the factor $M_1^{-|\omega|/2}$ in reserve together with the exponential (p. 3, after (1.7)). This factor is not used up on the parameters, and here it is still at our disposal for the factor $|\omega|+1$ that occurs below. Since the uncut hull, which is his localisation domain, is connected, $\mathsf P(\omega)$ is connected through walls; this yields the splitting. The carrier is needed because the fields are given only on $\Omega_0$, and it need not be connected when $\Omega_0\ne T_\ell$.

The properties of single terms are produced by construction; they are not deduced. The sum is itself covariant and real. Therefore the Haar mean of the terms, made symmetric under the conjugation fixing the real pairs, is once more an expansion of the sum. By convexity it satisfies the same bounds, provided that the family of bounds is invariant under the real gauge transformations and under this conjugation; the family of localised bounds has this invariance. A single assignment for all tori is obtained in three steps: a mean over translations; a push-forward from covering tori; and a passage to the limit by Montel's theorem, along a subsequence fixed once, at the given top level and units. In the push-forward, a walk wrapping around the small torus gains at most $|\omega|+1$ blocks; the weight therefore loses a factor $|\omega|+1$, and nothing else is lost. The family is made normal by an entrywise bound, which is one and the same number for all tori with room. This number is allowed to depend on the level, since the entries of $(Q\mathsf GQ^{\mathsf T})^{-1}$ are $O(\eta^{-2})$, and none of the norms, rates, weights or sums of weights involves it. If a family of terms is already carried by all shifts and compatible with all coverings, it is not changed.

(j) \textbf{Operators as data.} The principle goes back to Ba{\l}aban, who observes that in the equations for the response ``we can replace the propagators by arbitrary operators having the same regularity properties and satisfying the same bounds'' (\textbf{116}, p. 3, (1.5)). In its published form, however, the second equation relies on the identities $RD^{\mathsf T}\mathsf H=0$ and $\mathsf G^{-1}\mathsf H=Q^{\mathsf T}(Q\mathsf GQ^{\mathsf T})^{-1}$, and therefore on their consequences $Q\mathsf H=1$ and $\Delta_\pi\mathsf H=Q^{\mathsf T}[(Q\mathsf GQ^{\mathsf T})^{-1}-a]$ ($\Delta_\pi=\Delta^{\rm W}_{\mathbf U}-\Delta'_\pi$). None of these is available for the sums $\sum_\omega s^\omega\mathsf O_\omega$: on the cubes around a row where $s$ vanishes, every term is absent, so that $Q\mathsf H(s)$ equals $0$ there and not $1$. The first identity, moreover, holds only for the pairs on shell, even at $s\equiv1$. Consequently no identity enters the argument. The identities are applied earlier, where they hold. With their help the gradient of the non-quadratic part of the action is put into a shape from which its second-order term $\Delta_\pi\mathsf H$ is absent, and this expression \emph{defines} $\mathsf W_s$ for every $s$ and every pair (it is a gradient only at $s\equiv1$ and on shell). From then on only bounds are used, in the norms

\begin{gather*}
\FitMath{\|A\|_\varpi=\sup_y\varpi(y)\max\{\xi_y|A|,\,\xi_y^2|\nabla_{\mathbf U}A|\},} \\[3pt]
\FitMath{|f|^\sharp_\varpi=\sup_y\varpi(y)\,\xi_y^3|f|,} \\[3pt]
\FitMath{\varpi(y)\le e^{\theta d_{\mathbf B}(y,y')}\varpi(y'):}
\end{gather*}

the first measures a potential together with its first derivative, the second a force. The operator $\mathfrak G(s)$ takes the second norm back to the first, with norm $\hat{\mathsf b}\sim e^{2^4\bar\kappa}$, which is the size of every one of the operators; this closes the fixed-point argument. Whenever $\|A_i\|\le r<c/\hat{\mathsf b}$, one has, in every weight whose $\theta$ is at most a fixed fraction of $\delta_0$,

\begin{gather*}
\FitMath{|\mathsf W_s(A_1)-\mathsf W_s(A_2)|^\sharp_\varpi\le C\hat{\mathsf b}^2r\,\|A_1-A_2\|_\varpi;}
\end{gather*}

in the only term containing $\mathbf J$, the factor $r$ is supplied by the Schwarz lemma, applied to the full ball of radius $c/\hat{\mathsf b}$. The Lipschitz constant of the fixed-point map is $\hat{\mathsf b}\cdot\hat{\mathsf b}^2\cdot\hat{\mathsf b}\sup|B'|$, with contributions from $\mathfrak G$, $\mathsf W_s$ and $\mathsf H_1$, and this is why the radius is $\hat{\mathsf b}^{-4}$. Decay amounts to choosing the weight $\varpi=e^{\theta d_{\mathbf B}(\cdot,y')}$. Locality, splitting, covariance and transport follow from exhibiting a candidate and from uniqueness in the balls. When $s\not\equiv1$, bounds of first order on the operators do not yield a bound on $\Delta_{\mathbf U}\mathbf H$ independent of the level; an oscillating kernel belonging to the class of data is an example. The reduced current takes its place. It is built from the multiplier, and in it the second-order part of the gradient occurs only in paired form, $(D^{\mathsf T}DA,\mathsf He_c)=(DA,D\mathsf He_c)$, where $e_c$ is an element of $\mathfrak g^c$ placed at the row $c$. As a result, only first-order norms occur.

System (II) is handled by the same method, but a single map replaces the two nested ones. In each of its two components one operator of norm $\hat{\mathsf b}$ acts, on a ball of radius $2\hat{\mathsf b}r$, on terms vanishing to second order ($\mathcal C$, $\mathcal N$, $\mathcal C'^{\,\mathsf T}(\lambda-H(s)^{\mathsf T}\mathbf J)$). Hence $\hat{\mathsf b}^2r$ must be small, and this gives the radius $\hat{\mathsf b}^{-2}$ in $B'$. Two of the terms are of first order in the unknown. One is $\tilde G(s)\,\mathcal C_T'^{\,\mathsf T}H(s)^{\mathsf T}\mathbf J$ in the first component, whose linear part at $s\equiv1$ is $\tilde G\Delta^{(2)}T$. The other is $-H(s)^{\mathsf T}\mathcal C_T'^{\,\mathsf T}H(s)^{\mathsf T}\mathbf J$ in the second component. Each of them involves two operator norms and can be made small only through the current, which requires $\hat{\mathsf b}^2\gamma_0$ to be small. Under these conditions the map is a contraction with constant $\frac12$ in all the weighted norms $\max\{\|T\|_\varpi,|\lambda|_\varpi\}$ simultaneously. Since the member $-\mathsf a$ is diagonal and is carried by no proper set, the splitting is built around this member.

(k) \textbf{From the slice to a multiplier.} Let $s\equiv1$, and fix a critical real $U_0$ such that $(U_0,J(U_0))$ is a regular pair whose radii lie below the ceiling. Then (I) is equivalent to the constraint, the Landau condition and $\mathsf z:=F'(\mathbf H)-K_{\mathbf H}^{\mathsf T}\lambda\in DR(\text{site functions})$, where $\lambda$ is the multiplier formed from $\mathbf H$. In other words, (I) amounts to stationarity on the gauge slice alone. If $B'$ is real and close to $0$, then $\mathbf H(B')$ is real and lies in a neighbourhood of $0$ whose size is not uniform. On that neighbourhood the invariance of the action and the factorisation in 3 provide gauge directions $X_\chi$ with $(F'(\mathbf H),X_\chi)=0$ and $K_{\mathbf H}X_\chi=0$, and the $RD^{\mathsf T}X_\chi$ fill out the range of $R$. Therefore $\mathsf z=0$ on this neighbourhood, and by the identity theorem also on the whole polydisc, so the neighbourhood enters no constant. The factorisation of the averages holds only up to site factors and, at rows that leave their layer, up to seam factors; this is the reason for the row datum. For (II) the corresponding passage is organised by means of two defects. Take arbitrary small $T$ and $\lambda$, and set $B'=\mathcal Q(T)$, $\mathsf z=F'(T)-K_T^{\mathsf T}\lambda$ and $\phi=RD^{\mathsf T}T$. Then both equations of (II) have residuals that depend linearly on $(\mathsf z,\phi)$, the residual of the first being $\tilde G(\mathsf z+D\phi)$. Consequently (II) follows from the three equations and is tied to them.

\emph{Constants} are chosen in the following order, each large or small enough relative to those chosen before it. First comes $L$, and with it, as a single group, the numbers determined by $L$ alone: the flat constants and the flat rate, the lower thresholds that the flat case imposes on $M_1$ and on $R_1M_1$, the radius of the local perturbation theory, the rates $\delta_0$, $\delta_*$, $\delta'$, and the ceiling on $\alpha_0,\alpha_1,\gamma_0$. Next come $M_1$ and $R_1M_1$; the ceiling does not depend on them, and they do not depend on the ceiling. Then comes $\bar\kappa$. It does not depend on $M_1,R_1$, and they do not depend on it, so it may just as well be chosen before them. After that come $\ell_{\mathsf P}$, the radius in $B'$, and the second ceiling on $\gamma_0$, which concerns (II). The final choices are $M$ and the field strengths of the construction, taken so small that $O(LM)\alpha_0$ is below the ceiling. Two numbers are exceptions: among the bounds of statement 7, the one bound that is used for convergence alone; and, in mechanism (i), the entrywise bound by which the family is normal. Neither of them appears in a localised bound, a rate, a weight or a sum of weights. With the exception of these two, every constant, rate and ceiling, and the radius in $B'$, is independent of $k$, $K$, $\ell$, of the sequence (in particular of $\Omega_0$), of $Y$ and of the pair, and also of $M_1,R_1$ once these are large enough. The sizes $\ell_{\mathsf P}$ and $M$, and the side of the coarse cubes of 6, are each at least a multiple of $M_1$. The field strengths are small compared with $(LM)^{-1}$. Where statement 8 makes the identification with the response of (1b), the datum of the background is small compared with $M_1^{-1}$.

\subsection{Layered row maps and physical Hessians}
\label{ym4:exp500-rows}

We make explicit the endpoint readouts, the three successive Hessians, and the row extensions used by the propagator construction. The lattices, separated layers and determining rows are those fixed above. In this subsection every transpose uses the physical bilinear pairing, with fine volume $\eta^4$ and row volume $\xi_c^4$; real lower bounds use the corresponding Hermitian form. The row map $Q$ is per unit length. Its linearization is taken after the accumulated node factors of the average have been removed.

Write $\xi_b$ for the larger pillar scale at the endpoints of a fine bond. We use\[
 \|u\|_{\rm mass}^2=\sum_b\eta^4\xi_b^{-2}|u(b)|^2,
 \qquad \|\psi\|_{\rm row}^2=\sum_c\xi_c^2|\psi(c)|^2.
\]Thus $\|\psi\|_{\rm row}=\|\Xi^{-1}\psi\|_{2,{\rm rows}}$ when $\Xi$ multiplies row $c$ by $\xi_c$. All constants in this calculation may depend on the fixed block factor $L$. Their uniformity is in the depth, fine spacing, volume, and admitted separated regions.

Write $d=d_{\mathbf B}$ for the row distance inherited from the layer metric. For a complex variation put $U_E=e^{i\eta E}U$; $h$ bounds the physical first-order quantities $\xi|E|$ and $\xi^2|\nabla_UE|$ on the determining neighborhoods in the regular-pair domain. The real curvature and local real-gauge bounds have a common ceiling $\epsilon$. The Wilson Hessian at $U_E$ means its holomorphic continuation, with the bilinear transpose convention retained.

On site functions let $Q'=Q'_{\rm pill}$ be the own-pillar mean and set\[
 G'=(D_U^{\mathsf T}D_U+Q'^{\mathsf T}aQ')^{-1},\quad
 R=1-G'Q'^{\mathsf T}(Q'G'^2Q'^{\mathsf T})^{-1}Q'G',\quad
 \mathsf S=1-D_UG'RD_U^{\mathsf T}.
\]Here $R$ is the Landau projection, as in the propagator definitions above. The site identities $Q'G'R=0$ and $\Delta_U G'R=R$ imply that $\mathsf S$ projects to the slice $RD_U^{\mathsf T}u=0$. The field-row intertwining proved below gives $Q\mathsf S=Q$. On the physical graph the completed Hessian is $\Delta_\pi=\mathsf S^{\mathsf T}\Delta_U^{\rm W}\mathsf S$. We use the positive difference convention $\eta D_U\lambda=\operatorname{Ad}(U_b)\lambda(b_+)-\lambda(b_-)$.

\subsubsection{Cell norms retain the ultraviolet scale}

Work on a finite compatible torus, with a separated aligned sequence and the spacings of Section~\ref{candidate:background}. Let \(\Delta(y)\) be the pillar of a top \(y\), of side \(\xi_y\). Outside the action region one can add bare cells and extend the local scale to \(\eta\); a compressed operator is zero on that added space. The point scale \(\xi(x)\) is the scale of its pillar. Use the exact path metric\[
 d(x,z)=\inf_{\gamma:x\to z}
                   \int_\gamma\frac{ds_\infty}{\xi(s)} .
 \tag{PX1}\label{px:layer-metric}
\]The half-open convention only assigns a measure-zero cell boundary; paths and their infima can be approximated on either side. The corresponding lattice metric is comparable after the fixed collar constants have been inserted.

For each field species \(s\), declare a positive amplitude \(a_s(y)\) and use \(\|f\|_{s,y}=a_s(y)^{-1} \|f\|_{\infty,\Delta(y)}\). For example a physical potential has amplitude \(\xi_y^{-1}\), a physical force density \(\xi_y^{-3}\), and a raw datum logarithm amplitude one. The site gauge parameter also has amplitude one. A raw multiplier uses counting measure and amplitude one; its per-length density has amplitude \(\xi_y^{-3}\). The conversion is the weighted dual conversion of the domain section.

If \(P_y\) is restriction to a cell, define the normalized cell entry of an operator \(K:t\to s\) by\[
 k_{s t}(y,z)=
       \|P_yKP_z\|_{\,\ell^\infty(t,z)\to\ell^\infty(s,y)} .
 \tag{PX2}\label{px:cell-entry}
\]All physical cell measures in the action of \(K\) are included before taking this norm. It is a block operator norm, not a pointwise bound on the continuum kernel. A Green kernel with its fine Coulomb singularity can have a bounded cell norm. Replacing it by a pointwise kernel bound independent of \(\eta\) would be a stronger and generally false assertion.

For compatible species, the cell partition and triangle inequality give\[
 (k\ell)_{s u}(y,z)
       \le\sum_w k_{s t}(y,w)\ell_{t u}(w,z).
 \tag{PX3}\label{px:cell-composition}
\]This follows first for finite lattices by inserting \(\sum_wP_w=1\). In an infinite-volume application the same formula holds when the row sums below give absolute convergence. Thus a two-sided exponential cell norm is an algebra norm, without discarding the fine singularity inside a block.

\subsubsection{The growth of a ball in the layer metric}

The layer metric is not assumed to have polynomial volume growth. A small-cell region can contain many more pillars than a nearby large-cell region. Separation controls that growth quantitatively.

Suppose \(x\) has layer index \(j\) and \(z\) has index \(i\). Any path between them crosses each complete intervening layer. For a layer \(r\), a crossing from \(\Omega_{r+1}\) to \(\Omega_r^c\) has an arc whose interior lies in \(\Omega_r\setminus\Omega_{r+1}\). The separation makes its physical length greater than \(\rho\xi_r\). On that arc the metric density is \(\xi_r^{-1}\), so its length in \eqref{px:layer-metric} is at least \(\rho\). Arcs belonging to distinct layers have disjoint interiors. Consequently\[
 d(x,z)\ge\rho(|i-j|-1)_+ .
 \tag{PX4}\label{px:level-barrier}
\]One layer is lost because an endpoint can already lie next to its neighboring layer. The same proof applies to an intermediate point of a path, so a path of metric length at most \(R+1\) starting in layer \(j\) never uses a scale larger than\[
 \xi_{\max}\le
            L\xi_j L^{(R+1)/\rho}.
 \tag{PX5}\label{px:path-max-scale}
\]If the top index would exceed the depth, this is simply an overestimate.

Its physical length is at most \((R+1)\xi_{\max}\). Every pillar center in the metric \(R\)-ball therefore has a lift in the physical sup-norm ball of radius\[
 H_R=L(R+1)\xi_j L^{(R+1)/\rho}.
 \tag{PX6}\label{px:physical-ball}
\]For each fixed level \(i\) the centers lie on its grid of spacing \(\xi_i\). The number of grid points in that physical ball is at most \(C(1+H_R/\xi_i)^4\). Equation \eqref{px:level-barrier} restricts the possible indices to \(|i-j|\le1+(R+1)/\rho\). Using the smallest such \(\xi_i\) in the grid count and then summing the possible levels gives\[
 \#\{z:d(y,z)\le R\}
   \le C_L(1+R)^5
                  \exp\!\left(\frac{8\log L}{\rho}R\right).
 \tag{PX7}\label{px:metric-ball-count}
\]Indeed \(H_R/\xi_i\le C_L(R+1)L^{2(R+1)/\rho}\); its fourth power gives the exponential in \eqref{px:metric-ball-count}, and the number of levels costs one more polynomial power. Since \(\rho\ge1\), the harmless factors \(L^{8/\rho}\) and the bound on the number of levels can be included in \(C_L\) uniformly in \(\rho\).

The argument works on a torus without a nonwrapping-ball assumption. For every reached center choose the lift reached by its path. Distinct periodic centers give distinct lifted grid points. Counting all grid points in the cover ball is therefore an upper bound, even if that ball crosses many fundamental domains. Restricting the allowed regions only decreases the number counted.

An entry bound of the form\[
 k_{s t}(y,z)\le
   C\left(\frac{\xi_y}{\xi_z}\right)^p e^{-m d(y,z)}
 \tag{PX8}\label{px:entry-with-scale}
\]now implies a weighted row and column sum, provided the rate pays for scale growth. By \eqref{px:level-barrier} the scale ratio on an annulus of metric radius \(R\) is at most \(L^{|p|}e^{|p|(\log L)(R+1)/\rho}\). Break the space into unit annuli and use \eqref{px:metric-ball-count}. Thus\[
 \sup_y\sum_z e^{\delta d(y,z)}k_{s t}(y,z)
 +\sup_z\sum_y e^{\delta d(y,z)}k_{s t}(y,z)
       \le C'
 \tag{PX9}\label{px:entry-to-schur}
\]whenever\[
 m-\delta>\frac{(8+|p|)\log L}{\rho}.
 \tag{PX10}\label{px:summability-rate}
\]For the column sum use the same ball argument centered at \(z\); the absolute scale exponent is unchanged. The convergent scalar sum is bounded by \(C\sum_{n\ge0}(n+2)^5 e^{-[m-\delta-(8+|p|)\log L/\rho]n}\). This proves the claim with its precise uniformity.

If the amplitude normalization has already removed the scale ratio, take \(p=0\). For several species use the largest of their finitely many exponents. An exponential bound with an arbitrarily small \(m\) does not automatically imply \eqref{px:entry-to-schur}. The geometry and rate must satisfy \eqref{px:summability-rate}, or a stronger actual row estimate must be proved directly. This is a numerical condition in the order of choice of separation and kernel rates.


\subsubsection{The site readout at an outward row}

The site average used in the Landau projector has one value at each top, taken at that top's assigned level. The endpoint readout of a row is a different list when the row is outward. We now give their exact relation.

Let \(\mathsf T_j(y,x)\) be the product of the admitted internal mean transports from a level-\(j\) center \(y\) down its pillar to a bare site \(x\). For a top of level \(j\) define\[
 (Q'_{\rm pill}\lambda)(y)
   =L^{-4j}\sum_{x\in B^j(y)}
                 \operatorname{Ad}(\mathsf T_j(y,x))\lambda(x).
 \tag{PX31}\label{px:own-pillar-site-average}
\]These pillars are disjoint. There is therefore an explicit right inverse on node fields:\[
 (E_{\rm pill}\psi)(x)
      =\operatorname{Ad}(\mathsf T_j(y,x))^{-1}\psi(y)
       \quad (x\in B^j(y)),\qquad
                 Q'_{\rm pill}E_{\rm pill}=1 .
 \tag{PX32}\label{px:site-right-inverse}
\]This proves rank without an elliptic estimate. For real transports the right inverse has supremum norm one. With fine site measure \(\eta^4\) and top measure \(\xi_y^4\), it is also an isometry in \(L^2\). Indeed a pillar contains \(L^{4j}\) bare sites, so \(\eta^4L^{4j}=\xi_y^4\). The invariant bilinear color pairing gives\[
 Q'_{\rm pill}=E_{\rm pill}^{\mathsf T},\qquad
                  Q'_{\rm pill}(Q'_{\rm pill})^{\mathsf T}=1.
 \tag{PX33}\label{px:site-weighted-adjoint}
\]The algebraic identities hold for admitted complex transports as well. Their Hermitian norm bounds then include the actual condition numbers of those transports. A counting kernel must include the ratio \(\xi_y^4/\eta^4\) in the transpose; an ordinary unweighted array transpose would not be \eqref{px:site-weighted-adjoint}.

For any level-\(j\) center \(z\), write \(Q'_j\lambda(z)\) for its full \(j\)-fold transported site average, whether or not its assigned top level is \(j\). If \(z\) is the outer endpoint of a determining level-\(j\) row, its assigned level is \(j-1\). The entire block \(B^j(z)\) lies in that outer layer. Decompose it into its \(L^4\) level-\((j-1)\) pillars. The product definition of the transport gives\[
 Q'_j\lambda(z)
     =L^{-4}\sum_{w\in B(z)\cap\mathbb T^{(j-1)}}
        \operatorname{Ad}(T_{z,w}(\mathcal M^{j-1}U))
                               (Q'_{\rm pill}\lambda)(w).
 \tag{PX34}\label{px:outward-site-extension}
\]To verify the coefficient, each smaller average has weight \(L^{-4(j-1)}\); multiplying it by \(L^{-4}\) gives \(L^{-4j}\). The two transport factors multiply in their path order. This proves \eqref{px:outward-site-extension} exactly, without an asymptotic or small-field approximation.

Define \(\mathcal E_{j,z}\psi\) to be \(\psi(z)\) at an assigned level-\(j\) top, and to be the finite average on the right of \eqref{px:outward-site-extension} at an outer endpoint. Then define the extended coarse difference by\[
 (\bar D_{\rm ext}\psi)(c)
   =\xi_c^{-1}\{
          \operatorname{Ad}((\mathcal M^jU)(c))\mathcal E_{j,c_+}\psi-\mathcal E_{j,c_-}\psi\},
                         \qquad c\in\Gamma_j .
 \tag{PX35}\label{px:extended-coarse-difference}
\]For an interior row it is the ordinary covariant difference of its two top values. For an outward row it first averages the outer finer tops. There is no need to duplicate that site as a new independent top in the site Gram matrix.

Here is a direct verification of the field intertwining. Under an infinitesimal gauge \(e^{-it\lambda}\) the derivative of the level-\(j\) node factor is\[
 \left.\frac1i\frac d{dt}v_j[tD_U\lambda](y)\right|_{t=0}
                         =Q'_j\lambda(y)-\lambda(y).
 \tag{PX36}\label{px:node-factor-gauge-derivative}
\]At one step the change in each relative contour has derivative \(\operatorname{Ad}(T_{y,x})\lambda(x)-\lambda(y)\). The implicit mean at the identity has differential the arithmetic mean of its inputs, so this gives the one-step formula. Iterating and using the product rule for the accumulated node factors gives \eqref{px:node-factor-gauge-derivative}, with the hierarchical site average \eqref{px:own-pillar-site-average}.

Differentiate the actual straightening identity \(\mathcal M^j(e^{i\eta A}U) =v_j[A]\,e^{iq_j^{\rm st}(A)} \mathcal M^jU\,v_j[A]^{-1}\) on a row. The plain coarse variation is \(\operatorname{Ad}(\mathcal M^jU(c))\lambda(c_+)-\lambda(c_-)\). Subtract the two node derivatives \eqref{px:node-factor-gauge-derivative}, with their endpoint transports. The remaining variation is\[
 Dq_j^{\rm st}(0)\,D_U\lambda
     =\operatorname{Ad}(\mathcal M^jU(c))Q'_j\lambda(c_+)-Q'_j\lambda(c_-).
 \tag{PX37}\label{px:raw-straightened-intertwining}
\]Divide by \(\xi_c\) to obtain the per-length row \(Q\). Using \eqref{px:outward-site-extension} proves the intertwining identity on interior and outward rows\[
                 QD_U
                          =\bar D_{\rm ext}Q'_{\rm pill}.
 \tag{PX38}\label{px:typed-seam-intertwining}
\]This supplies the intertwining needed in the slice identities above. There \(\bar D\) means this extended row difference and \(Q'\) means the own-pillar site map. Using a single unextended value of the outer top would give the wrong formula.

Each endpoint extension reads at most \(L^4\) tops in one adjacent layer. Its real row sum is at most one, since the transports are unitary and the weights add to one. A finer top occurs in one such outer block and at most eight incident coarse rows. Thus its two-sided normalized cell counts are bounded at fixed \(L\), without a factor depending on the depth. For complex pairs the same statement uses the actual bounded transport condition numbers. The latter are bounds on the local transports in the regular-pair domain; the algebraic intertwining does not prove them.


\subsubsection{Finite real existence of the site operators}

There is also a real finite positivity statement independent of a quantitative inverse bound. Assume the actual mean transports are admitted and unitary, and that every active fine site lies in a connected pillar with its top. Let \(a\) be strictly positive on the top space. Then\[
 L_{\rm site}=D_U^{\mathsf T} D_U+
                    (Q'_{\rm pill})^{\mathsf T} aQ'_{\rm pill}
                         \quad\text{is strictly positive}.
 \tag{PX39}\label{px:finite-site-positivity}
\]To prove it, a vector of zero Hermitian energy must obey \(D_U\lambda=0\) and \(Q'_{\rm pill}\lambda=0\). Split into Hermitian real and imaginary parts; it suffices to treat real \(\lambda\). The first identity means that \(e^{it\lambda(x)}\) leaves every active bond unchanged:\[
 e^{it\lambda(b_-)}U_b
                        e^{-it\lambda(b_+)}=U_b.
 \tag{PX40}\label{px:parallel-gauge-stabilizer}
\]Gauge covariance of each actual mean transport therefore gives \(\operatorname{Ad}(\mathsf T_j(y,x))\lambda(x) =\lambda(y)\). Average over the pillar: \(Q'_{\rm pill}\lambda(y)=\lambda(y)\). Its zero value makes \(\lambda(y)=0\). Parallel transport along the active graph from the top to every site makes \(\lambda=0\). This proves \eqref{px:finite-site-positivity}.

It follows that \(G'=L_{\rm site}^{-1}\) exists on every finite real problem of this type. The site Gram matrix\[
 Q'_{\rm pill}G'^2(Q'_{\rm pill})^{\mathsf T}
 \tag{PX41}\label{px:finite-site-gram}
\]is strictly positive as well. For a nonzero node vector \(\psi\), surjectivity makes \((Q'_{\rm pill})^{\mathsf T}\psi\ne0\); invertibility of \(G'\) then makes its squared norm positive. Thus the formulas for \(R\) and \(\mathsf S\) are genuine finite real matrices before their uniform bounds are established.

The quantitative site gap is proved next.

\subsubsection{Transport control on one pillar}

The site gap has a quantitative proof from the actual averages. Let a pillar have side \(\xi=n\eta\), where \(n=L^j\), and assume\[
 |\partial U(p)-1|\le
                 \epsilon\,\eta^2\xi^{-2}
 \tag{PX42}\label{px:pillar-curvature}
\]on its determining neighborhood. The fixed path neighborhood has an embedded rectangular lift. If an extra fine collar is needed, the regularity bounds in the adjacent layer give the same inequality with a fixed \(L\)-dependent factor. We include that factor in the constants below.

Take the fine coordinate comb from the pillar center and put \(g(x)=U(\Gamma_{y,x})\). For adjacent \(x,x+\eta e_\mu\), the loop made by the two root combs and that bond is filled by strips of width one fine cell. Commuting the extra \(\mu\)-step through the later coordinate segments uses at most \(Cn\) plaquettes, not \(Cn^2\). The unitary Stokes product inequality therefore gives\[
 |U^g(b)-1|\le C\epsilon/n
                  \quad\text{on internal fine bonds}.
 \tag{PX43}\label{px:fine-comb-strip}
\]Here \(g(y)=1\). Every curvature norm is unchanged. The gauge need not be close to one, and its derivatives are not estimated in this step.

Write \(\xi_r=L^r\eta\), and let \(a_r\) be the supremum of \(|\mathcal M^r(U^g)(c)-1|\) on the level-\(r\) internal edges used in the hierarchical transport. Let \(\varpi_r\) be the corresponding coarse plaquette bound. The determining neighborhoods remain in the rectangular pillar neighborhood: each parent block is a union of its child blocks, and each contour stays in its assigned block. The real curvature iteration proved in the background section gives\[
 \varpi_r\le C_L\epsilon(\xi_r/\xi)^2 .
 \tag{PX44}\label{px:pillar-coarse-curvature}
\]The same finite comparison shows that one actual averaged edge differs from its central straight transport by at most \(C_L\varpi_r\). Indeed every closed logarithm in its definition is a comparison loop with that central transport. The unitary strip estimate bounds those logarithms, and then their exponential average. The central route has \(L\) edges, so\[
       a_{r+1}\le L a_r+C_L\varpi_r .
 \tag{PX45}\label{px:pillar-open-recurrence}
\]The coefficient of \(a_r\) is exactly \(L\), by the unitary product inequality.

Iterating \eqref{px:pillar-open-recurrence}, using \(a_0\le C\epsilon\eta/\xi\), gives\[
 \begin{split}
 a_r&\le L^ra_0+
       C_L\epsilon\sum_{s<r}
                 L^{r-1-s}(\xi_s/\xi)^2\\
     &\le C_L\epsilon\,\xi_r/\xi .
 \end{split}
 \tag{PX46}\label{px:pillar-open-bound}
\]For the sum, insert \(\xi_s=\eta L^s\); its geometric factor is bounded by a constant times \((\xi_r/\xi)^2\), which is no larger than \(\xi_r/\xi\). The constants do not multiply with the number of levels.

The internal mean on level \(r\) is a finite smooth mean of paths of \(O(L)\) edges. On its regular near-identity chart it differs from one by at most \(C_La_r\). The hierarchical transport has one such factor at each level below \(j\). All are unitary, hence\[
 |\mathsf T_j(U^g;y,x)-1|
       \le C_L\sum_{r<j}a_r
       \le C_L\epsilon
                    \quad(x\in B^j(y)).
 \tag{PX47}\label{px:uniform-pillar-transport}
\]The last sum uses \(\sum_{r<j}\xi_r/\xi \le1/(L-1)\). The fixed smallness of \(\epsilon\) also admits every mean used in this argument. Thus \eqref{px:uniform-pillar-transport} is a uniform-depth property of the actual iterated transports, not a hypothesis of such a property.

\subsubsection{A real weighted site gap}

We first prove the flat block inequality with its scale. Let \(\bar\lambda\) be the arithmetic mean on the \(n^4\) sites of a block. The exact variance identity is\[
 \sum_x|\lambda(x)-\bar\lambda|^2
       =\frac1{2n^4}\sum_{x,z}|\lambda(x)-\lambda(z)|^2 .
 \tag{PX48}\label{px:block-variance}
\]Connect each pair by the coordinate-order path. Its length is at most \(4n\). A fixed oriented edge lies in at most \(n^5\) of these paths: its two longitudinal endpoint coordinates give at most \(n^2\) choices, and the three other coordinates give at most \(n^3\). Cauchy on the path and then \eqref{px:block-variance} give a constant times \(n^2\) times the sum of squared edge differences. Multiplying by \(\eta^4\xi^{-2}\) proves\[
 \xi^{-2}\|\lambda\|_{2,B}^2
  \le C\big\{
         \|\partial_{\rm int}\lambda\|_{2,B}^2
                              +\xi^2|\bar\lambda|^2\big\}.
 \tag{PX49}\label{px:flat-block-poincare}
\]Here the fine norm uses \(\eta^4\), and \(\partial\) divides a bond difference by \(\eta\). The mean part is \(\eta^4n^4\xi^{-2} |\bar\lambda|^2 =\xi^2|\bar\lambda|^2\). For \(n=1\) the variance vanishes and the same inequality holds.

Apply the real comb gauge of \eqref{px:fine-comb-strip} in each pillar. Only internal bonds are used in the block inequality, so the gauges of neighboring pillars may be different. They define a unitary site gauge globally, and all covariant norms are unchanged. In this frame\[
 \|(\partial_{\rm int}-D_{U,{\rm int}})\lambda\|_{2,B}
                       \le C_L\epsilon\,\xi^{-1}\|\lambda\|_{2,B}.
 \tag{PX50}\label{px:block-gradient-comparison}
\]The estimate follows from \(|\operatorname{Ad}(U_b)-1| \le C_L\epsilon\eta/\xi\) and the at-most-eight incidences at a site.

Let \(Q'_0\) be arithmetic block average. Equation \eqref{px:uniform-pillar-transport}, Cauchy for the finite average and \(\xi^4=n^4\eta^4\) give\[
 \xi^4|(Q'_{\rm pill}-Q'_0)\lambda|^2
                        \le C_L\epsilon^2\|\lambda\|_{2,B}^2 .
 \tag{PX51}\label{px:block-average-comparison}
\]Use \eqref{px:block-gradient-comparison}--\eqref{px:block-average-comparison} on the right of \eqref{px:flat-block-poincare}. The result is the same two terms with \(D_U\) and \(Q'_{\rm pill}\), plus \(C_L\epsilon^2\xi^{-2} \|\lambda\|_{2,B}^2\). Choose the fixed curvature ceiling so that this coefficient is at most one half of the left coefficient. Absorption proves\[
 \xi^{-2}\|\lambda\|_{2,B}^2
 \le C_L\big\{
       \|D_{U,{\rm int}}\lambda\|_{2,B}^2
                         +\xi^2|Q'_{\rm pill}\lambda|^2\big\}.
 \tag{PX52}\label{px:covariant-block-poincare}
\]

Sum over the disjoint pillars. Every internal fine bond is counted at most once; all other gradient terms in the full energy are nonnegative. Suppose \(a(y)=c_y\xi_y^{-2}\) with \(c_y\ge a_*>0\). Then\[
 \begin{split}
 \sum_x\eta^4\xi(x)^{-2}|\lambda(x)|^2
 &\le C_{L,a_*}\big\{
        \|D_U\lambda\|_2^2
             +\langle Q'_{\rm pill}\lambda,
                              aQ'_{\rm pill}\lambda\rangle_{\rm tops}\big\}\\
 &=C_{L,a_*}\langle\lambda,L_{\rm site}\lambda\rangle .
 \end{split}
 \tag{PX53}\label{px:real-weighted-site-gap}
\]This is the actual site weighted gap, uniform in depth, volume, region shape and real holonomy. It uses the small real curvature and actual transport construction; no physical current or criticality is used. Crossing-edge energy is kept, not replaced by a clipped block mean.

For a complex site vector split into Hermitian real and imaginary parts. Since the real unitary connection is a real matrix in an orthonormal Lie-algebra basis, the same inequality holds for its Hermitian energy. If the top scale is at most one, it in particular gives an ordinary finite \(L^2\) gap independent of the bare mesh. In other physical units its smallest unweighted gap is scaled by the inverse square of the largest local scale.

\subsubsection{The complex site domain uses the actual iterated means}

The preceding transport argument has a complex version with a fixed physical radius, independent of the number of fine sites. We give it to avoid assuming a uniform radius for the iterated site average.

Consider one step on a fixed two-block neighborhood. Suppose all its edge matrices satisfy \(|W_b-1|\le a\), and let \(\varpi=\sup|\partial W-1|\). There are constants \(r_L,C_L\) for which the actual means and logarithms are admitted if \(a,\varpi<r_L\), and\[
 \begin{split}
 a_{\rm next}&\le La+C_La^2+C_L\varpi,\\
 \varpi_{\rm next}&\le
       L^2(1+C_La)\varpi+C_L\varpi^2 .
 \end{split}
 \tag{PX54}\label{px:complex-one-step-recurrences}
\]For the first line, a central straight product differs from one by at most \(La+C_La^2\). Plaquette sliding bounds each closed comparison logarithm by \(C_L\varpi\): all prefix transports have condition number at most \(1+C_La\) on this finite neighborhood. The defining exponential average therefore changes the central product by at most that amount.

For the second line, take the local complex comb gauge. Its values and inverses are \(1+O(C_La)\). It sets the tree bonds to one; each remaining bond is then \(1+O(C_L\varpi)\) by the finite plaquette sliding identity. In that gauge the Taylor/flux argument for the actual one-step map gives the coefficient \(L^2\) and a quadratic error \(C_L\varpi^2\). Conjugating the fine and coarse plaquettes through the small complex gauge changes their norms by \(1+C_La\). This proves \eqref{px:complex-one-step-recurrences}. Complex covariance here follows directly from the conjugation identity for the small logarithm, or from uniqueness of the admitted finite mean. No absolute-value formula is continued holomorphically.

These constants can be bounded before an iteration is run. The finite mean construction has a joint raw-edge holomorphic ball of radius \(c/L\), with universal supremum bounds; two scalar Cauchy circles bound its second derivative by \(CL^2\). A comb on the one-step patch sweeps at most \(CL^2\) plaquettes, so its non-tree logarithms have size \(CL^2\varpi\). The quadratic Taylor error in the coarse plaquette is therefore at most \(CL^2(L^2\varpi)^2 =CL^6\varpi^2\). The fine logarithmic curl error times its \(L^2\) flux factor has the same upper order. Prefix conjugations cost \(1+CLa\). For example \(r_L=cL^{-3}\) and \(C_L=CL^8\), with fixed universal shrinkages and enlargements, suffice in \eqref{px:complex-one-step-recurrences}. This count uses directional Cauchy on the proved full supremum ball, not one coordinate estimate summed over all field coordinates.

On a pillar let \(t_r=\xi_r/\xi\), and normalize \(A_r=a_r/t_r\), \(V_r=\varpi_r/t_r^2\). Equation \eqref{px:complex-one-step-recurrences} implies, with an increased fixed constant,\[
 Z_{r+1}\le Z_r+C_Lt_r(Z_r+Z_r^2),
            \qquad Z_r=A_r+V_r,\qquad
                 \sum_{r<j}t_r\le\frac1{L-1}.
 \tag{PX55}\label{px:complex-transport-gronwall}
\]For \(Z_r\le1\), the factor is at most \(1+2C_Lt_r\). Its product is at most \(\exp(2C_L/(L-1))\). Choose \(Z_0\) below one half the minimum one-step chart radius times the inverse of this number. Induction then gives \(Z_r\le C_LZ_0\) and keeps every step in its chart. The constants depend on \(L\), not \(j\). With the preceding finite choices a sufficient common small input radius has the explicit form \(cL^{-3}\exp(-CL^8)\). The precise numerical constants are not needed below, but this gives a positive lower radius rather than an asserted upper cap. The product of the internal means and its inverse have norm at most \(\exp(C_L\sum_ra_r)\). Hence\[
 |\mathsf T_j(W;y,x)-1|\le C_LZ_0,\qquad
 |\mathsf T_j(W;y,x)|+
                  |\mathsf T_j(W;y,x)^{-1}|\le C_L .
 \tag{PX56}\label{px:complex-pillar-transport}
\]This proves the existence and boundedness of the whole finite hierarchical formula on that domain.

Now start with a real field satisfying \eqref{px:pillar-curvature}, put it in the real comb gauge \eqref{px:fine-comb-strip}, and let \(W=e^{i\eta E}U\). Assume the actual physical first-order input satisfies\[
 h=\max\{\|\xi E\|_\infty,
                    \|\xi^2\nabla_UE\|_\infty\}\le h_L .
 \tag{PX57}\label{px:site-complex-input}
\]In the comb frame the ordinary first difference of \(E\) differs from its covariant one by \(O(\epsilon\xi^{-1}|E|)\). The elementary plaquette product then gives \(a_0\le C(\epsilon+h)\eta/\xi\) and \(\varpi_0\le C(\epsilon+h+h^2)(\eta/\xi)^2\). Thus the starting \(Z_0\) in \eqref{px:complex-transport-gronwall} is at most \(C(\epsilon+h)\).

Choose a fixed \(h_L>0\) so that the same argument holds with \(E\) replaced by \(sE\) for \(|s|<h_L/h\), and assume \(h\le h_L/4\). Reduce the real curvature ceiling at the same time so that \(\epsilon\le h_L\) and the resulting starting \(Z_0\) satisfies \eqref{px:complex-transport-gronwall}. The transport is holomorphic in \(s\) on this disc, with a uniform bound. Cauchy's formula on a smaller concentric disc, or the geometric sum of its Taylor coefficients, yields\[
 |\mathsf T_j(e^{i\eta E}U;y,x)
                       -\mathsf T_j(U;y,x)|\le C_Lh .
 \tag{PX58}\label{px:site-transport-lipschitz}
\]The bound in \eqref{px:complex-pillar-transport} is uniform in \(j\) and in the number of coordinates. This is a single directional Cauchy argument on an actual admitted physical disc. It does not assert an entrywise absolute Hessian sum for an arbitrary many-variable functional.

For the complex site average write \(Q'_E\) and for the real reference write \(Q'_U\). The transported finite-average calculation used in \eqref{px:block-average-comparison} gives\[
 \|a^{1/2}(Q'_E-Q'_U)\lambda\|_{2,{\rm tops}}
                   \le C_Lh\,\|\xi^{-1}\lambda\|_2 .
 \tag{PX59}\label{px:complex-site-average-difference}
\]Here the real row weights satisfy \(a(y)\le a^*\xi_y^{-2}\). The proof squares the average difference in each disjoint pillar, uses \eqref{px:site-transport-lipschitz}, and sums its fine volume. There is no factor equal to the number of fine sites. The same argument applies to the outer endpoint extension \eqref{px:outward-site-extension}: its neighborhood uses one adjacent scale, so the input constants change by fixed powers of \(L\).

\subsubsection{Relative forms give the complex site inverse}

It remains to compare the actual site operators, not their unscaled matrix coefficients. Let \(L_U\) be the real site operator and let \(L_E=D_{e^{i\eta E}U}^{\mathsf T} D_{e^{i\eta E}U}+Q_E'^{\mathsf T} aQ'_E\). The transposes in \(L_E\) are bilinear analytic transposes. At a real reference they are the usual Hermitian adjoints.

For a fine bond the difference of the two site derivatives is bounded by \(C|E_b||\lambda(b_+)|\). Adjacent pillar scales differ by at most the fixed factor \(L\); the region gaps exclude a bare bond jumping over another layer. The finite incidence count therefore gives\[
 \|(D_{e^{i\eta E}U}-D_U)\lambda\|_2
                        \le C_Lh\|\xi^{-1}\lambda\|_2 .
 \tag{PX60}\label{px:complex-site-gradient-difference}
\]Keeping derivatives on the two arguments, expand the kinetic difference. Use \eqref{px:complex-site-gradient-difference}, then \eqref{px:complex-site-average-difference} for the mass difference, and finally the real weighted site gap \eqref{px:real-weighted-site-gap}. For arbitrary complex vectors \(v,w\), this gives\[
 |\langle v,(L_E-L_U)w\rangle_{\rm H}|
 \le C_Lh\,
      \langle v,L_Uv\rangle_{\rm H}^{1/2}
      \langle w,L_Uw\rangle_{\rm H}^{1/2}.
 \tag{PX61}\label{px:relative-complex-site-form}
\]For example the quadratic mass term is \((Q'_E-Q'_U)^{\mathsf T} a(Q'_E-Q'_U)\); the two cross terms contain only one difference. On replacing the left argument by its complex conjugate, the same absolute bounds apply to the analytic bilinear transpose. No Hermitian conjugation of the complex coefficients is inserted into their definition.

Let \(B_E=L_U^{-1/2}(L_E-L_U)L_U^{-1/2}\). Equation \eqref{px:relative-complex-site-form} proves \(\|B_E\|_{2\to2}\le C_Lh\). For \(C_Lh\le1/2\) one has\[
 G'_E=L_U^{-1/2}(1+B_E)^{-1}L_U^{-1/2},\qquad
 \operatorname{Re}\langle\lambda,L_E\lambda\rangle_{\rm H}
                  \ge\tfrac12\langle\lambda,L_U\lambda\rangle_{\rm H}.
 \tag{PX62}\label{px:complex-site-inverse}
\]The inverse is holomorphic on the whole admitted physical ball. It has the same weighted coercivity, up to a fixed factor, uniformly in depth and volume. This is a relative-form argument. The unscaled operator norm of a coefficient perturbation need not be bounded as \(\eta\downarrow0\).



\subsubsection{Exponential decay from the weighted site form}

Let \(V\) denote multiplication by \(\xi(x)^{-2}\) on sites. The real gap \eqref{px:real-weighted-site-gap} is \(L_U\ge cV\). We prove a weighted inverse bound before converting it to pointwise regularity. This avoids estimating a bare kinetic operator by its unscaled norm.

Let \(\psi\) be bounded and one-Lipschitz for \eqref{px:layer-metric}, and let \(M_s=e^{s\psi}\) on sites. Use the value \(\psi(b_-)\) on an edge and \(\psi(y)\) on a top. Put\[
 D_s=M_s^{\rm edge}D_UM_s^{-1},\qquad
 Q'_s=M_s^{\rm top}Q'_UM_s^{-1}.
 \tag{PX63}\label{px:site-weighted-rows}
\]On an edge the two endpoint values of \(\psi\) differ by at most \(C_L\eta/\xi_b\). Inside a pillar their difference from the top value is at most a fixed constant. The difference from the unweighted row therefore satisfies\[
 \|(D_s-D_U)\lambda\|_2
  +\|a^{1/2}(Q'_s-Q'_U)\lambda\|_{2,{\rm tops}}
               \le C_L|s|\|V^{1/2}\lambda\|_2
 \tag{PX64}\label{px:site-weighted-row-difference}
\]for \(|s|\le s_L\). The edge estimate follows from \(|e^{st}-1|\le C|s||t|\) on the stated bounded interval. For the site row use the same estimate in its finite transported average, then Cauchy with the actual pillar volume. Its constant is independent of \(\|\psi\|_\infty\).

The conjugated form has the exact factorization\[
 M_sL_UM_s^{-1}
              =D_{-s}^{\mathsf T} D_s+(Q'_{-s})^{\mathsf T} aQ'_s .
 \tag{PX65}\label{px:conjugated-site-form}
\]Expanding it and using \eqref{px:site-weighted-row-difference}, \eqref{px:real-weighted-site-gap} proves a relative form error at most \(C_L|s|\). With a complex field perturbation in the ball of \eqref{px:complex-site-inverse}, the same calculation gives error \(C_L(h+|s|)\) relative to the real \(L_U\). Choose both small enough for that number to be at most one half. All these forms are invertible.

For any such inverse \(\mathcal G_s\), testing the equation \(L_s\lambda=V^{1/2}f\) against \(\lambda\) gives\[
 \|V^{1/2}\mathcal G_sV^{1/2}\|_{2\to2}
                                        \le 2/c .
 \tag{PX66}\label{px:weighted-site-inverse}
\]Indeed its real part is at least \((c/2)\|V^{1/2}\lambda\|_2^2\), whereas the right pairing is at most \(\|V^{1/2}\lambda\|_2\|f\|_2\). This proof uses neither a normal matrix nor a pointwise positive complex kernel.

Choose \(\psi(x)=\min\{d(x,\Delta(z)),N\}\). It is one-Lipschitz. The diameter of a pillar in \eqref{px:layer-metric} is bounded, so on \(\Delta(y)\) its value is at least \(d(y,z)-C\) after \(N\) is large enough. Equation \eqref{px:weighted-site-inverse} then yields\[
 \|P_yG'_EP_z\|_{2\to2}
       \le C_L\xi_y\xi_z e^{-m d(y,z)}
 \tag{PX67}\label{px:site-cell-l2-decay}
\]for a fixed \(m=m(L)>0\). On a finite torus take the periodic distance. On the cover first use bounded \(N\) and then increase it. The constants are independent of this truncation.

In particular the normalized operator \(V^{1/2}G'_EV^{1/2}\) has exponentially decaying cell \(L^2\) entries. Its two-sided cell Schur norm follows from \eqref{px:metric-ball-count}--\eqref{px:summability-rate} when the separation is chosen with\[
             \frac{8\log L}{\rho}<m/2 .
 \tag{PX68}\label{px:site-separation-choice}
\]The rate \(m\) above depends only on the fixed one-step constants and the site gap; it does not depend on \(\rho\). Thus this is a genuine order of choice: increase the separation after those constants, or verify that the already chosen separation satisfies this finite inequality. It is not an appeal to unspecified polynomial volume growth.

\subsubsection{A scalar subsolution estimate at the fine mesh}

We use a local estimate whose proof also explains why a fine Coulomb singularity causes no loss in \eqref{px:site-cell-l2-decay}. On the infinite scalar fine lattice in four dimensions, let \(\mathcal L=-\Delta_\eta\). The positive heat kernel gives a positive Green kernel\[
 G_0(x)=\int_0^\infty p_t^\eta(x)\,dt,\qquad
 |G_0(x)|\le C(|x|+\eta)^{-2},\quad
 |\nabla G_0(x)|\le C(|x|+\eta)^{-3}.
 \tag{PX69}\label{px:scalar-massless-kernel}
\]Here is the pointwise heat-kernel bound used in this integration.
Put $s_t=\max(\eta,\sqrt t)$. The Fourier formula (BGH2)--(BGH4)
gives $\|D^r p_t^\eta\|_\infty\le C_rs_t^{-4-r}$ for $r=0,1$:
for $t\ge\eta^2$ integrate $|\theta|^r e^{-ct|\theta|^2/\eta^2}$
and rescale, and for smaller $t$ use the bounded Fourier multiplier.
Write the half-time convolution with the derivative on one fixed
factor, $D^rp_t^\eta(x)=\sum_y\eta^4
(D^rp_{t/2}^\eta)(y)p_{t/2}^\eta(x-y)$.
Split the sum into $|y|\ge |x|/2$ and its complement.
On the first set apply (BGH5) of order $r$ to the first factor
and the order-zero supremum bound to the second.
On the complement $|x-y|\ge |x|/2$; apply the order-zero tail
bound to the second factor and the order-$r$ supremum bound
to the first. Since $s_{t/2}$ and $s_t$ are comparable, for
every fixed $N$ this gives
\[
 |D^r p_t^\eta(x)|\le C_{r,N}s_t^{-4-r}
                    (1+|x|/s_t)^{-N},\qquad r=0,1.
\]
Integrating these inequalities gives the asserted Green bounds. For example split at \(t=|x|^2\), use a polynomial off-diagonal tail before the split, and \(Ct^{-2}\), respectively \(Ct^{-5/2}\), afterwards. At \(|x|\le\eta\) the bounds are the corresponding lattice scaling bounds. Also \(\mathcal LG_0\) is the normalized lattice delta, by integrating \(-\partial_tp_t^\eta\).

Suppose \(w\ge0\) and \(\mathcal Lw\le F_0\) on a ball of radius \(2r\), \(r\ge8\eta\). Choose a nonnegative cutoff \(\chi\), equal to one on the \(r\)-ball, supported in the \(2r\)-ball, with its first two differences bounded by \(Cr^{-1}\), \(Cr^{-2}\). At the center \(x\), summation by parts gives\[
 w(x)=
       \langle\chi G_0(x-\cdot),\mathcal Lw\rangle
             -\langle[\mathcal L,\chi]G_0(x-\cdot),w\rangle .
 \tag{PX70}\label{px:subsolution-cutoff-identity}
\]All sums use \(\eta^4\). This identity follows by moving \(\mathcal L\) off \(w\) and using \(\mathcal LG_0=\delta_x\). There is no derivative of \(w\) in its last term. Since the first test is nonnegative, its contribution is at most \(CF_0r^2\). Indeed the lattice sum of \(G_0\) on that ball is at most \(Cr^2\), by dyadic shells.

The commutator is supported on the cutoff annulus. Writing it as \((\mathcal L\chi)G_0\) plus products of first differences shows that its absolute value is at most \(Cr^{-4}\) there, by \eqref{px:scalar-massless-kernel}. Cauchy and the annulus volume therefore prove\[
 w(x)\le C\{r^{-2}\|w\|_{2,B_{2r}}+r^2F_0\}.
 \tag{PX71}\label{px:scalar-subsolution-estimate}
\]If \(r/\eta\) is bounded, the point-value estimate \(|w(x)|\le\eta^{-2}\|w\|_2\) gives the same inequality with a changed constant. Thus the local estimate is uniform at coarse meshes as well.

For a real unitary connection, let \(\lambda\) be Lie-algebra valued and put \(w=|\lambda|\). For \(\lambda(x)\ne0\), pair each term \(\lambda(x)-\operatorname{Ad}(U_{x,x'})\lambda(x')\) with \(\lambda(x)/|\lambda(x)|\). Unitarity gives\[
 -\Delta_\eta|\lambda|(x)
       \le |D_U^{\mathsf T} D_U\lambda(x)|.
 \tag{PX72}\label{px:discrete-kato}
\]At a zero value of \(\lambda\) the left side is nonpositive, so the same inequality holds. This proves the covariant scalar subsolution inequality without differentiating a gauge. It is used only on a patch with the full covariant Laplacian stencil. Bare boundary cells will use the point-value fallback instead.

\subsubsection{Pointwise site Green bounds}

First take a real \(U\) and \(\lambda=G'_Uf\). At a site in pillar \(\Delta(y)\), the mass part of its equation obeys\[
 |(Q_U'^{\mathsf T} aQ'_U\lambda)(x)|
                 \le C_L\xi_y^{-4}\|\lambda\|_{2,\Delta(y)} .
 \tag{PX73}\label{px:site-mass-pointwise}
\]In fact it is the transported vector \(a(y)Q'_U\lambda(y)\), whose norm is at most \(a^*\xi_y^{-2}|Q'_U\lambda(y)|\). Cauchy in the average gives the additional \(\xi_y^{-2}\).

Take a local ball of radius comparable to \(\xi_y\). It meets a bounded \(L\)-dependent number of pillars, all with scales in \([\xi_y/L,L\xi_y]\). The gaps exclude a second intervening layer in that ball once \(\rho>C_LL\). Call the corresponding whole-pillar neighborhood \(\mathcal N(y)\). Its cells are at bounded \(d\)-distance from \(y\). Apply \eqref{px:scalar-subsolution-estimate}--\eqref{px:site-mass-pointwise} on a lift of this ball. This proves\[
 \|\lambda\|_{\infty,\Delta(y)}
   \le C_L\xi_y^{-2}
                 \sum_{w\in\mathcal N(y)}\|\lambda\|_{2,\Delta(w)}
         +C_L\xi_y^2\|f\|_{\infty,\mathcal N(y)} .
 \tag{PX74}\label{px:site-local-l2-to-infinity}
\]Near a boundary of the action region, every pillar lacking this full patch is bare: the nonbare regions have the prescribed gap from that boundary. The single-cell fallback in \eqref{px:scalar-subsolution-estimate} proves \eqref{px:site-local-l2-to-infinity} there directly. On a torus the local ball is lifted to the periodic cover. Its fixed radius can meet only a fixed number of copies of a cell, since every period contains at least two top blocks. This changes only the constant.

If \(f\) is supported in \(\Delta(z)\), use \eqref{px:site-cell-l2-decay} and \(\|f\|_{2,\Delta(z)} \le\xi_z^2\|f\|_\infty\) on the right side. The far contribution is \(C\xi_z^3\xi_y^{-1} e^{-m d(y,z)}\|f\|_\infty\). The local forcing term has \(d(y,z)\le C_L\). Equation \eqref{px:level-barrier} absorbs the scale ratio:\[
 \|G'_Uf\|_{\infty,\Delta(y)}
       \le C_L\xi_y^2
                e^{-m_1d(y,z)}\|f\|_\infty,
 \qquad 0<m_1<m-\frac{3\log L}{\rho}.
 \tag{PX75}\label{px:site-pointwise-green}
\]Choosing the separation large enough leaves a positive fixed \(m_1\). This bound concerns a bounded cell source.

The actual equation also gives\[
 \|D_U^{\mathsf T} D_UG'_Uf\|_{\infty,\Delta(y)}
                  \le C_Le^{-m_1d(y,z)}\|f\|_\infty .
 \tag{PX76}\label{px:site-laplacian-green}
\]Use \(D_U^{\mathsf T} D_U\lambda =f-Q_U'^{\mathsf T} aQ'_U\lambda\), and bound its mass term by the supremum on its same pillar. The cell diameter in the layer metric is bounded.

\subsubsection{Site gradients and their complex continuation}

For differentiated pointwise bounds use the actual local gauge clause of the regular-pair class. On the required patch around a cell it supplies a unitary gauge with\[
 U=e^{i\eta A},\qquad
 \xi|A|+\xi^2|\nabla A|\le a_{\rm loc},
                         \qquad a_{\rm loc}\le a_L .
 \tag{PX77}\label{px:site-local-coefficient-class}
\]The constant \(a_L\) is fixed small. Near one interface the patch may be described in the finer adjacent scale; converting to \(\xi\) changes these bounds by fixed powers of \(L\). The region gaps and the fixed patch collar are chosen to include no third layer. This gauge clause is an input defining the regular-pair class. Deriving it for an actual background from weaker tests is the separate background estimate.

In this local gauge the site equation is a scalar componentwise Laplace equation with first- and zeroth-order errors. Expanding the transported differences gives\[
 |\Delta_\eta\lambda|
 \le |f-Q'^{\mathsf T} aQ'\lambda|
       +C_La_{\rm loc}
          \{\xi^{-1}|\nabla\lambda|
                                      +\xi^{-2}|\lambda|\}.
 \tag{PX78}\label{px:site-scalar-equation}
\]Each difference of a link exponential costs \(\eta|\nabla A|\); each undifferentiated link error costs \(\eta|A|\). After division by the two powers of \(\eta\) in the Laplacian the terms have exactly the displayed dimensions.

Here is the interior absorption that supplies the gradient. On a slightly larger ball of radius \(C\xi\), put \(M=\|\lambda\|_\infty\) and \(F=\|f-Q'^{\mathsf T} aQ'\lambda\|_\infty\). Let \(d_b\) be distance of a fine edge to the ball boundary and take the finite maximum of \(d_b|\nabla\lambda(b)|\). If \(d_b\) is only a fixed multiple of \(\eta\), the point-value bound on a difference controls it by \(CM\). Otherwise apply the scalar interior gradient estimate to a ball of radius \(d_b/4\) around an edge where this maximum is approached. Distances in that smaller ball are comparable to \(d_b\). Its gradient supremum is therefore at most the weighted maximum divided by \(c d_b\). Using \eqref{px:site-scalar-equation} in the scalar estimate gives\[
 \sup_b d_b|\nabla\lambda(b)|
       \le C_L(M+\xi^2F)
                       +C_La_{\rm loc}
                           \sup_b d_b|\nabla\lambda(b)|.
 \tag{PX79}\label{px:site-gradient-absorption}
\]Choose \(C_La_{\rm loc}\le1/2\). Absorption, \eqref{px:site-pointwise-green}--\eqref{px:site-laplacian-green} and the scalar interior Hölder estimate then give, for a bounded source in \(\Delta(z)\),\[
 \begin{split}
 \|\nabla_UG'_Uf\|_{\infty,\Delta(y)}
       &\le C_L\xi_y e^{-\delta d(y,z)}\|f\|_\infty,\\
 [\nabla_UG'_Uf]_{\alpha,\Delta(y)}
       &\le C_{L,\alpha}\xi_y^{1-\alpha}
                           e^{-\delta d(y,z)}\|f\|_\infty,
                       \qquad 0<\alpha<1 .
 \end{split}
 \tag{PX80}\label{px:site-gradient-green}
\]The Hölder comparison is made in the assigned local frame, or with its declared transport to a common base. The bounded potential and pointwise gradient give the transported version. The constant is not asserted bounded as \(\alpha\uparrow1\). Cells with a bounded number of fine sites use the direct difference bound; near the action boundary they are bare, as already noted.

These pointwise estimates continue to the actual complex site operator by a stronger norm argument than a pointwise matrix coefficient estimate. For a fixed source cell \(z\), let\[
 \begin{split}
 \|\lambda\|_{\mathcal Y_z}
   &=\sup_y e^{\delta d(y,z)}
       \max\{\xi_y^{-2}\|\lambda\|_{\infty,\Delta(y)},
              \xi_y^{-1}\|\nabla_U\lambda\|_{\infty,\Delta(y)}\},\\
 \|f\|_{\mathcal F_z}
   &=\sup_y e^{\delta d(y,z)}\|f\|_{\infty,\Delta(y)} .
 \end{split}
 \tag{PX81}\label{px:site-profile-spaces}
\]Use a rate \(\delta\) strictly below the cell Green rate, with its spare rate larger than the ball-growth rate in \eqref{px:metric-ball-count}. Summing the bounded-source cell estimates then proves \(\|G'_U\|_{\mathcal F_z\to\mathcal Y_z} \le C_L\), uniformly in \(z\).

For the actual left perturbation \(U_E=e^{i\eta E}U\), expand the two covariant site differences and the two site averages. The kinetic difference has a coefficient bounded by \(|E|\) on a first difference of \(\lambda\), and a coefficient bounded by \(|\nabla_UE|+|E|^2\) on \(\lambda\). The mass difference uses \eqref{px:site-transport-lipschitz}, has size \(C_Lh\xi^{-2}\), and reads only the same pillar. Hence\[
 \|(L_E-L_U)\lambda\|_{\mathcal F_z}
                         \le C_Lh\|\lambda\|_{\mathcal Y_z}.
 \tag{PX82}\label{px:strong-site-perturbation}
\]No second difference of \(\lambda\) occurs: the difference of two connections in a site Laplacian is lower order. The first-order input on \(E\) is used here; its supremum alone would not bound the zeroth-order coefficient. The metric diameter of a pillar is bounded, so the profile weight passes through the mass average with a fixed factor.

For \(h\) below the resulting fixed threshold, the series\[
 G'_E=\sum_{n\ge0}
            [-G'_U(L_E-L_U)]^nG'_U
 \tag{PX83}\label{px:strong-site-neumann}
\]converges from \(\mathcal F_z\) to \(\mathcal Y_z\). At finite volume it coincides with the inverse already constructed in \eqref{px:complex-site-inverse}, by uniqueness. It therefore proves the complex versions of the pointwise potential and gradient estimates. The equation supplies the complex covariant Laplacian bound. In the same real local frame, \(e^{i\eta E}e^{i\eta A}\) has a physical logarithm with the first-order bounds \(C(a_{\rm loc}+h)\). The scalar interior estimate then gives the complex Hölder member as well.

The argument establishes the site estimates for bounded cell sources. A mixed gradient with a divergence source has a Calderón--Zygmund singularity and requires its specified cancellation or Hölder input. No uniform bound for that mixed operator on all bounded raw sources has been inferred.

\subsubsection{The site Gram inverse has its own lower-bound proof}

Invertibility of \(G'\) does not by itself give a uniform inverse for \(Q'G'^2Q'^{\mathsf T}\). We prove that estimate using actual site bumps, rather than assuming it as part of the site Green theorem. Let \(\Xi\) multiply a top value by \(\xi_y\).

Choose a fixed smooth scalar bump \(\chi_y\), supported in the middle portion of its pillar and equal to one on a smaller middle portion. Its first two differences are bounded by \(C\xi_y^{-1}\) and \(C\xi_y^{-2}\). The arithmetic average \(\beta_y=Q'_0\chi_y\) has a positive lower bound independent of the fine mesh: count the grid points in that smaller cube. The finitely coarse cases use the center point.

Use the real local gauge \eqref{px:site-local-coefficient-class} on a neighborhood of the pillar and its one-fine-cell collar. In this gauge the initial open bond bound is \(C a_{\rm loc}\eta/\xi_y\). The recurrence \eqref{px:pillar-open-recurrence}, with the same real curvature bound, therefore puts the hierarchical transports within \(C_L(a_{\rm loc}+\epsilon)\) of one. Consequently the three by three matrix \(m_y\psi=Q'_U(\chi_y\psi)\) equals \(\beta_y I+ O(C_L(a_{\rm loc}+\epsilon))\). For the fixed small class ceiling it is invertible with uniform bound.

Define a site lift on that pillar by \(\chi_y m_y^{-1}\psi(y)\) in the local gauge and undo the gauge. At the node use its transformed frame as well. The resulting operator \(E_{\rm sm}\) satisfies\[
 Q'_UE_{\rm sm}=1,\qquad
 \|E_{\rm sm}\psi\|_2\le C_L\|\psi\|_{2,{\rm tops}},\qquad
 \|L_UE_{\rm sm}\psi\|_2
                         \le C_L\|\Xi^{-2}\psi\|_{2,{\rm tops}} .
 \tag{PX84}\label{px:site-smooth-lift}
\]The first equality holds pillar by pillar, since the site averages have disjoint supports. For the last inequality, the covariant Laplacian of one bump is bounded by \(C_L\xi_y^{-2}|\psi(y)|\). Expand it in the local gauge: the second cutoff differences, the connection times first differences, and the differentiated connection all have that bound by \eqref{px:site-local-coefficient-class}. Its mass term is the transported vector \(a(y)\psi(y)\) on the same pillar. The volume is \(\xi_y^4\), giving the last \(L^2\) scale in \eqref{px:site-smooth-lift}. One fine collar of different bumps has bounded overlap, since each fine site has a fixed number of neighbors. For a bare pillar, take the single site bump; its Laplacian is bounded by \(C\eta^{-2}|\psi|\), which is exactly the same scale. Thus no missing local gauge at the outer bare boundary is used.

For a real top vector \(\psi\), use \eqref{px:site-smooth-lift} with \(\Xi^4\psi\) and the equation \(L_UG'_U=1\):\[
 \begin{split}
 \|\Xi^2\psi\|_{2,{\rm tops}}^2
   &=\langle(Q'_U)^{\mathsf T}\psi,
                         E_{\rm sm}\Xi^4\psi\rangle\\
   &=\langle G'_U(Q'_U)^{\mathsf T}\psi,
                           L_UE_{\rm sm}\Xi^4\psi\rangle\\
   &\le C_L\|G'_U(Q'_U)^{\mathsf T}\psi\|_2
                            \|\Xi^2\psi\|_{2,{\rm tops}} .
 \end{split}
 \tag{PX85}\label{px:site-gram-duality}
\]Divide by the last node norm, with the zero case immediate. This proves the operator lower bound\[
 Q'_UG_U'^2(Q'_U)^{\mathsf T}\ge c_L\Xi^4 .
 \tag{PX86}\label{px:site-gram-lower}
\]The \(\Xi^4\) test in \eqref{px:site-gram-duality} is essential. Testing only with \(E_{\rm sm}\psi\) would give a weaker ratio of two norms and would not prove the stated weighted lower bound at different local scales.

For the upper bound and locality, set\[
 \mathcal F_U=G'_U(Q'_U)^{\mathsf T}\Xi^{-2},\qquad
 M_U=\Xi^{-2}Q'_UG_U'^2(Q'_U)^{\mathsf T}\Xi^{-2}
                                      =\mathcal F_U^{\mathsf T}\mathcal F_U .
 \tag{PX87}\label{px:normalized-site-gram}
\]The site injection \((Q'_U)^{\mathsf T}\) is an isometry between node \(L^2\) and its transported constant pillar fields. Equation \eqref{px:site-cell-l2-decay} therefore bounds the normalized cell entry of \(\mathcal F_U\) by\[
          C_L(\xi_y/\xi_z)e^{-m d(y,z)} .
 \tag{PX88}\label{px:site-column-cell-bound}
\]The same estimate holds for the complex \(\mathcal F_E\), with the bounded transport condition numbers already proved. The cell summability calculation with \(p=1\) gives a uniform exponential two-sided Schur norm of these operators, after increasing the fixed separation if needed. Their products therefore have the same type of norm. In particular \(c_LI\le M_U\le C_LI\) at a real pair.

The complex perturbation is uniform as well. The finite matrix \(\mathcal F_E\) is holomorphic on the actual physical first-order ball constructed above. Its operator norm is uniformly bounded there by \eqref{px:site-column-cell-bound}. Apply one scalar Cauchy circle to the line through \(E\). This gives \(\|\mathcal F_E-\mathcal F_U\| \le C_Lh\), independent of the number of coordinates. Hence\[
 \|M_E-M_U\|_{2\to2}\le C_Lh,\qquad
                         \|M_E^{-1}\|_{2\to2}\le 2/c_L
 \tag{PX89}\label{px:complex-site-gram}
\]on a fixed smaller ball. Here \(M_E=\mathcal F_E^{\mathsf T}\mathcal F_E\) uses the analytic bilinear transpose. Its norm difference is bounded by \((\|\mathcal F_E\|+\|\mathcal F_U\|) \|\mathcal F_E-\mathcal F_U\|\). The transpose has the same Hilbert operator norm in the real weighted orthonormal bases. No Hermitian positivity of the complex Gram matrix is assumed.

An exponentially weighted Schur norm also controls its conjugation. For a bounded node function \(\psi\) that is one-Lipschitz in \(d\), its cell entries change by the factor \(e^{s(\psi(y)-\psi(z))}-1\). The inequality\[
 |e^{st}-1|\le |s|\,|t|e^{|s||t|}
 \tag{PX90}\label{px:nonlocal-conjugation}
\]and a spare fraction of the exponential rate bound this perturbation by \(C_L|s|\) in operator norm. Indeed \(\sup_{r\ge0}r e^{-\delta r/2}<\infty\). Choose \(s\) small and use \eqref{px:complex-site-gram}. The inverse of the conjugated \(M_E\) is uniformly bounded. The distance-weight argument used for \eqref{px:site-cell-l2-decay} gives exponential entries of \(M_E^{-1}\).

In the physical top pairing the weighted orthonormal basis at \(y\) is \(\xi_y^{-2}\) times the color basis. If \(B_{yz}\) are the entries of \(M_E^{-1}\) in those orthonormal coordinates, then \(\|B_{yz}\|\le C_Le^{-m_2d(y,z)}\). Undoing both \(\Xi^{-2}\) factors and the orthonormal bases gives the counting kernel\[
 \left|
 (Q'_EG_E'^2(Q'_E)^{\mathsf T})^{-1}(y,z)
 \right|
                   \le C_L\xi_y^{-4}e^{-m_2d(y,z)} .
 \tag{PX91}\label{px:site-gram-counting-kernel}
\]Explicitly the unnormalized entries of \(M_E^{-1}\) are \((\xi_z^2/\xi_y^2)B_{yz}\). Multiplying on the left and right by \(\xi_y^{-2}\), \(\xi_z^{-2}\) leaves \(\xi_y^{-4}B_{yz}\). This explains the asymmetric fourth power in a counting-kernel formula, despite symmetry in the physical weighted pairing.

Finally the actual site projector has the representation\[
 R_E=1-\mathcal F_EM_E^{-1}\mathcal F_E^{\mathsf T} .
 \tag{PX92}\label{px:site-projector-kernel}
\]It agrees with its original Gram formula by \eqref{px:normalized-site-gram}. The smoothing part has an exponential cell \(L^2\) Schur norm, and the matrix is holomorphic on the whole smaller physical ball. At a real field it is the orthogonal projector already identified. This completes the site Gram and projector part of the package. Fine vector coercivity, mixed derivative source norms and the current constraint terms are not consequences of \eqref{px:site-projector-kernel}.


\subsubsection{The Hessian read from the constrained minimum}

We next identify the three quadratic forms in the physical calculation. This is a finite-dimensional identity on any admitted chart; it precedes a uniform inverse estimate. Let \(U\) be the reference critical field, \(j_U\) its action current, \(Q\) the derivative of the actual straightened constraint, and \(\omega\) its multiplier. Take a fixed reference slice, so that its linear condition is \(RD_U^{\mathsf T} a=0\).

Criticality on that slice already gives the full stationarity equation\[
                         j_U=Q^{\mathsf T}\omega .
 \tag{PX93}\label{px:physical-stationarity}
\]To see this, first note that \(Q\mathsf S=Q\), \(\mathsf S a\) is in the slice, and \(a-\mathsf S a=D_UG'RD_U^{\mathsf T} a\). For \(a\in\ker Q\), criticality annihilates \(\mathsf Sa\); the Ward identity annihilates its remaining gauge part. Thus \(j_U\) annihilates \(\ker Q\). The finite-dimensional annihilator identity gives \eqref{px:physical-stationarity}. When the row space is the actual independent datum space, \(Q\) is onto and \(\omega\) is unique. Equivalently it is unique whenever the row Gram inverse exists.

Because \(Q\mathsf H=1\), \eqref{px:physical-stationarity} implies \(\mathsf H^{\mathsf T} j_U=\omega\). Define the actual Lagrangian Hessian\[
             K_{\rm L}=\Delta_U^{\rm W}-\Delta^{(2)}[\omega].
 \tag{PX94}\label{px:physical-lagrangian-hessian}
\]The quadratic constraint term here is the second derivative of the same straightened constraint as \(Q\), with the same per-length pairing. It cannot be replaced by the quadratic part of a different row map.

For the invariant extension of the Wilson Hessian, the physical second variant is therefore\[
 \mathsf G_1[j_U]^{-1}
    =\mathsf S^{\mathsf T} K_{\rm L}\mathsf S
                      +D_URD_U^{\mathsf T}+Q^{\mathsf T} aQ .
 \tag{PX95}\label{px:physical-full-form}
\]Let \(W_1=Q\mathsf G_1[j_U]Q^{\mathsf T}\) and \(H_1=\mathsf G_1[j_U]Q^{\mathsf T} W_1^{-1}\). The physical gauge identity gives \(RD_U^{\mathsf T} H_1=0\). Consequently \(\mathsf SH_1=H_1\), and direct multiplication of \eqref{px:physical-full-form} gives\[
                   H_1^{\mathsf T} K_{\rm L}H_1=W_1^{-1}-a .
 \tag{PX96}\label{px:physical-row-readout}
\]For example, \(H_1^{\mathsf T}\mathsf G_1^{-1}H_1=W_1^{-1}\), the gauge contribution is zero, and the row mass contributes \(a\), since \(QH_1=1\). This proves all signs in \eqref{px:physical-row-readout}. It does not prove the existence of either inverse.

Here and in \eqref{px:coarse-action-hessian}--\eqref{px:active-datum-flatness},
$B$ and $b$ are per-length row coordinates. With
$\ell=\operatorname{diag}(\xi_c)$, they are related to the raw
data of \eqref{ym4:exp500-bg:mb:reduced-datum} by
$b(B)=\ell^{-1}b_{\rm raw}(\ell B)$.
The multiplier $\omega$ is dual to this per-length convention.
Let the actual local critical branch be parametrized by a coarse variation \(B\), with reduced datum \(b(B)\). Its derivative \(h_B\) at zero is taken in the fixed reference slice. Differentiated stationarity gives the constrained linear equations for \(K_{\rm L}\), with row value \(Db(0)\dot B\). They imply\[
        h_B=H_1Db(0),\qquad
 D^2\mathcal A_{\rm eff}(0)
      =Db(0)^{\mathsf T}(W_1^{-1}-a)Db(0)
                            +\langle\omega,D^2b(0)\rangle .
 \tag{PX97}\label{px:coarse-action-hessian}
\]Here \(\mathcal A_{\rm eff}\) is the action evaluated on that branch. For the first assertion, the differentiated equation tested against slice vectors in \(\ker Q\) says \(\mathsf S^{\mathsf T} K_{\rm L}h_B\) lies in the range of \(Q^{\mathsf T}\). Adding \(Q^{\mathsf T} aQh_B\) and using \eqref{px:physical-full-form} then solves the row problem by its defining inverse. For the second, differentiate the value function twice. The second derivative of the branch contracts with \(j_U=Q^{\mathsf T}\omega\); differentiating its constraint replaces that contraction by \(\omega D^2b- \omega D^2q[h_B,h_B]\). The latter term is exactly the subtraction in \(K_{\rm L}\).

\subsubsection{Which coarse coordinates remove the datum curvature}

The statement \(b(B)=B\) has a specific coordinate meaning. In the present chart the actual coarse group field is\[
                             V(B)(c)=e^{i\xi_cB(c)}V_0(c).
 \tag{PX98}\label{px:left-coarse-chart}
\]The fine straightened row logarithm is relative to the same reference \(V_0\). On an interior row the exact node identity then gives \(b(B)(c)=B(c)\). It holds for nontrivial \(V_0(c)\) as well. No assertion about the absolute logarithm of \(V(B)(c)\) is involved.

An outward row instead has its prescribed factor at the finer outer endpoint. In group form its equation is\[
       e^{i\xi_cb(B)(c)}
           =e^{i\xi_cB(c)}\operatorname{Ad}_{V_0(c)}
                                      v_z((\ell B)_{\rm int}) .
 \tag{PX99}\label{px:outward-datum-factor}
\]Here $v_z$ is the raw endpoint factor in
\eqref{ym4:exp500-bg:mb:reduced-datum}. It can have both linear and quadratic datum corrections. If only the highest active level is varied, the lower-level data \(B_{\rm int}\) in this factor are fixed to zero. Then \(v_z(0)=1\) and the factor is constant. More generally the next-step support mask must be checked against every outward carrier. Deep interior rows give an immediate instance; a mask by itself does not change a logarithmic coordinate.

For contrast, if an absolute coordinate \(W_c=\log(V_0(c))/i\) is used and the variable is still a relative increment \(D\), its new absolute logarithm is \(\log(e^{iD}e^{iW_c})/i\). Its second derivative in \(D\) generally does not vanish. Equation \eqref{px:coarse-action-hessian} then retains its datum-curvature term. This distinction is independent of the support mask.

Let \(\iota\) inject the actual active next-step coarse bond variables into the row space. When the literal map \eqref{px:left-coarse-chart}--\eqref{px:outward-datum-factor} satisfies\[
 Db(0)\iota=\iota,\qquad
                     D^2b(0)[\iota\,\cdot,\iota\,\cdot]=0,
 \tag{PX100}\label{px:active-datum-flatness}
\]its coarse action precision on those variables is \(\iota^{\mathsf T}(W_1^{-1}-a)\iota\). Its coarse current is \(\iota^{\mathsf T} H_1^{\mathsf T} j_U\). The restriction \(\iota^{\mathsf T}\) and the physical dual weights are part of the readout, even when notation suppresses them.

\subsubsection{The next curved fibre inserts a different multiplier}

At the next lattice, write raw bond increments as \(D\), so that the group field is \(e^{iD}V\). The actual plain constraint logarithm is\[
 \mathfrak q_V(D)
    =\log\!\left(\mathcal M(e^{iD}V)\mathcal M(V)^{-1}\right)/i .
 \tag{PX101}\label{px:next-raw-constraint}
\]Its local pivot graph has derivative \(C_V\), and the pivot injection \(h_V\) satisfies \(D\mathfrak q_V(0)h_V=1\). Twice differentiating the graph equation gives\[
 D^2\Phi_V(0)[u,v]
                =-h_VD^2\mathfrak q_V(0)[C_Vu,C_Vv].
 \tag{PX102}\label{px:next-graph-second-derivative}
\]This is the actual pivot graph, not a free extension chosen after the Hessian.

Let \(\widehat K\) and \(\widehat J\) be the coarse effective action Hessian and gradient in these raw counting coordinates. The chain rule on that graph is\[
 \begin{split}
 D^2(\mathcal A_{\rm eff}\circ\Phi_V)(0)[u,v]
   ={}&\langle C_Vu,\widehat K C_Vv\rangle\\
     &-\langle h_V^{\mathsf T}\widehat J,
                   D^2\mathfrak q_V(0)[C_Vu,C_Vv]\rangle .
 \end{split}
 \tag{PX103}\label{px:next-constrained-hessian}
\]Thus the next constraint multiplier is
$h_V^{\mathsf T}\widehat J$. It is not $\widehat J$ on the
new row space unless the notation already includes this transpose.
Define $E_V^{(2)}[\widehat J]$ by the last contraction in
\eqref{px:next-constrained-hessian}. Let
$\widehat J_{\rm GF}$ and $\widehat K_{\rm GF}$ be the gradient
and Hessian at $D=0$ of the actual function
$D\mapsto\mathrm{GF}(e^{iD}V)$, in the same raw counting coordinates.
The chain rule for the complete next free precision gives
\[
 C_V^{\mathsf T}
 \bigl(\widehat K+\widehat K_{\rm GF}
       -E_V^{(2)}[\widehat J+\widehat J_{\rm GF}]\bigr)C_V .
 \tag{PX104}\label{px:complete-next-precision}
\]
The physical step chooses $V=V_k^{\rm bg}=\mathcal M^kU$
in block axial gauge. The actual nonnegative $\mathrm{GF}_k$
vanishes there, so its full first derivative is zero and
$\widehat J_{\rm GF}=0$. This recovers the shorter physical
formula with $E_V^{(2)}[\widehat J]$. At an arbitrary reference,
both gauge derivatives in \eqref{px:complete-next-precision}
must be retained. An absolute contour gauge and its relative version
can have different derivatives away from the axial locus; the formula
uses the actual gauge definition of the model being considered.

For a uniform next-lattice spacing \(\xi\), the raw and physical conversions are\[
       \widehat K=\xi^2 K_{\rm phys,count},\qquad
                              \widehat J=\xi^3J_{\rm phys}.
 \tag{PX105}\label{px:coarse-raw-duality}
\]Indeed \(\xi^4\sum B(KB)= \xi^2\sum D(KD)\) when \(D=\xi B\), and \(\xi^4\sum J_{\rm phys}B= \sum \xi^3J_{\rm phys}D\). At nonuniform row lengths the corresponding diagonal length matrices must be kept on both sides; their commutation is not assumed. The coupling coefficients in the original action remain inside \(K\) and \(J\).

Combining \eqref{px:physical-row-readout}, \eqref{px:active-datum-flatness} and \eqref{px:next-constrained-hessian} gives the physical diagram: fine Lagrangian Hessian, coarse value Hessian, then next-fibre Hessian. Each arrow has its own row map, its own dual current, and its own gauge term. The physical identity proves none of those three objects interchangeable.

\subsubsection{A flat Feynman comparison on all the layers}

There is a useful vector estimate which does not involve the Landau projection or an independent current. It uses the actual flat straightened averages. We prove it to specify exactly what is available before those two additional operations.

First take the full periodic fine lattice as the ambient graph. Let \(a_\mu(x)\) be a fine vector one-form, with physical \(L^2\) measure \(\eta^4\). For a positive level-\(j\) row of length \(\xi_j=n\eta\), \(n=L^j\), its flat per-length average is\[
 (Q_j^0a)_\mu(y)
    =n^{-5}\sum_{u\in\mathcal B_n}
                      \sum_{t=0}^{n-1}a_\mu(y+\eta u+\eta t e_\mu).
 \tag{PX107}\label{px:flat-field-average}
\]Here \(\mathcal B_n\) is the centered set of \(n^4\) fine block sites. The formula is the flat straightened map, not the plain path map. At \(j=0\) it is the identity. The coefficient of each fine bond is nonnegative, the coefficients sum to one, and each is at most \(n^{-4}\). The last fact follows because at most \(n\) pairs \((u_\mu,t)\) give a fixed longitudinal position, while the three other coordinates are fixed.

Its support lies in the union of the two endpoint blocks. Consequently\[
 |Q_j^0a(c)|\le
               C\xi_j^{-2}\|a_\mu\|_{2,\mathcal D_c} .
 \tag{PX108}\label{px:flat-row-ltwo}
\]For example Cauchy gives \(\sum_b w_b^2\le n^{-4}\) from \(\sum_bw_b=1\) and \(w_b\le n^{-4}\); converting counting \(L^2\) to physical \(L^2\) contributes \(\eta^{-2}\). Thus the resulting coefficient is \(\eta^{-2}n^{-2}=\xi_j^{-2}\). This bound is specific to this flat row formula.

For a fine bond \(b\), put \(\xi_b\) equal to the larger of its endpoint pillar scales and define\[
 \|a\|_{\rm mass}^2
          =\sum_b\eta^4\xi_b^{-2}|a(b)|^2,\qquad
 \mathcal E_{\rm F}^0(a)
     =\|\partial a\|_2^2+
                  \sum_{c\in\Gamma}c_c\xi_c^2|Q^0a(c)|^2 ,
 \quad 0<c_*\le c_c\le c^* .
 \tag{PX109}\label{px:flat-feynman-energy}
\]The component gradient contains all fine coordinate differences on the ambient torus. The second term is exactly the physical row pairing with \(a(c)=c_c\xi_c^{-2}\). In particular it has neither an extra fine volume factor nor an unweighted counting transpose.

For each nonbare top \(y\) and each direction \(\mu\), one of the adjacent level-\(j\) rows has \(y\) as an endpoint and is a determining row. Indeed the excluded higher region is a union of aligned blocks of side at least \(L\xi_j\). Along that coordinate line, the excluded block containing an exterior point has at least \(L\ge3\) level-\(j\) sites. Both neighboring sites cannot belong to the higher region. Choose one that does not. If the other endpoint is outside \(\Omega_j\), this is precisely an outward row, which is retained. At the highest level there is no higher region to exclude.

Let \(m_{y,\mu}\) be the arithmetic mean of \(a_\mu\) on the pillar enlarged by one fine cell. A rectangular carrier containing that enlarged pillar and the chosen row's two endpoint blocks has side length at most \(C\xi_j\). For \(j\ge1\) it lies inside the ambient fine graph by the nesting gap. The ordinary box Poincare inequality, proved by the coordinate-path count used in the scalar block inequality, gives\[
 |Q_j^0a(c(y,\mu))-m_{y,\mu}|
           \le C\xi_j^{-1}\|\partial a_\mu\|_{2,\mathcal D_{y,\mu}} .
 \tag{PX110}\label{px:flat-row-mean-comparison}
\]To verify this without a pointwise trace estimate, subtract the carrier's arithmetic mean. Both functionals annihilate that constant in their difference and have \(L^2\) norm at most \(C\xi_j^{-2}\). Poincare then supplies one factor \(\xi_j\).

A second use of Poincare on the enlarged pillar gives\[
 \xi_j^{-2}\|a\|_{2,\mathcal B_y^+}^2
       \le C\sum_\mu\|\partial a_\mu\|_{2,\mathcal D_{y,\mu}}^2
                    +C\xi_j^2\sum_\mu|Q_j^0a(c(y,\mu))|^2 .
 \tag{PX111}\label{px:flat-field-local-poincare}
\]Every fine bond with a nonbare endpoint can be assigned to a pillar of its larger endpoint scale and is included in that pillar's enlargement. If both endpoints are bare, its level-zero row is the identity and controls it directly.

The overlap in the gradient sum is uniform. For a fixed level, all carrier centers are on its grid and their diameters are \(C\xi_j\), so a fine point belongs to a fixed number of them. A carrier of an outward level-\(j\) row is inside \(\Omega_{j-1}\), and its endpoint blocks are outside \(\Omega_{j+1}\). Its one-fine-cell collar can enter only an adjacent level. Thus a fixed fine point is counted at only a fixed number of levels, not at every level of the sequence. Each chosen row has at most its two endpoints as selecting tops. Summing \eqref{px:flat-field-local-poincare} therefore proves\[
                         \|a\|_{\rm mass}^2
                                    \le C_L\mathcal E_{\rm F}^0(a).
 \tag{PX112}\label{px:flat-feynman-gap}
\]The constants are independent of depth, fine mesh, volume and the shapes of the separated regions. Periodic carriers which wrap are lifted to their rectangles; the compatible torus has at least two top blocks, so the multiplicity of that lift is bounded.

There is also the upper row estimate\[
              \sum_{c\in\Gamma}\xi_c^2|Q^0a(c)|^2
                                      \le C_L\|a\|_{\rm mass}^2 .
 \tag{PX113}\label{px:flat-field-row-upper}
\]Use \eqref{px:flat-row-ltwo} on each row carrier. All fine endpoint scales there are comparable to \(\xi_c\), within the fixed adjacent scale factor. The preceding overlap count then applies. Equations \eqref{px:flat-feynman-gap} and \eqref{px:flat-field-row-upper} are estimates for the flat component Feynman form.

For a proper active region, extend the sequence by bare cells on the ambient torus and compress this full positive form by zero extension of its active bond coordinates. Its lower bound survives exactly: apply \eqref{px:flat-feynman-gap} to the extended vector. This is the Dirichlet compression of the component Laplacian. An inner-plaquette Maxwell Hessian has a different boundary form. Nor has the nonlocal subtraction \(D_U(1-R)D_U^{\mathsf T}\) been included. Those comparisons are separate in the vector precision problem.

\subsubsection{The Wilson Hessian at small curvature}

The Wilson part admits a relative form estimate without assuming a matrix norm bound. We derive it directly from the plaquette product. Use \(\langle X,Y\rangle_{\rm col} =\tfrac12\operatorname{Tr}(XY)\) on Hermitian traceless matrices. For a plaquette with oriented bonds \(b_1,\ldots,b_4\), put \(P=U(b_1)\cdots U(b_4)\) and \(T_{r-1}=U(b_1)\cdots U(b_{r-1})\). A left variation \(e^{i\eta v(b)}U(b)\) inserts \(X_r=\eta\operatorname{Ad}(T_{r-1})v(b_r)\). On reversed bonds the variation is \(v(\bar b)=-\operatorname{Ad}(U(b)^{-1})v(b)\), as follows by inverting the group field.

The varied plaquette is \(\prod_{r=1}^4e^{iX_r}P\). At \(P=1\), the trace of the mixed second derivative is the trace of the product of the two sums of insertions. Cyclicity of the trace cancels their commutators. For general real \(P\), subtract that calculation and use the factor \(P-1\) in every remaining term. Thus the plaquette action \(1-\tfrac12\operatorname{Tr}P\) satisfies\[
 \begin{split}
 D^2\mathcal A_p(U)[v,w]
     &=\left\langle\sum_rX_r,\sum_rY_r\right\rangle_{\rm col}
                                      +\mathcal R_p(v,w),\\
 |\mathcal R_p(v,w)|
     &\le C\|P-1\|
           \left(\sum_r|X_r|^2\right)^{1/2}
           \left(\sum_r|Y_r|^2\right)^{1/2}.
 \end{split}
 \tag{PX114}\label{px:plaquette-hessian-exact}
\]There are only sixteen mixed insertion terms. Unitarity of the prefixes makes \(\sum|X_r|^2=\eta^2\sum|v(b_r)|^2\). The same polynomial trace expression, with no real-part operation, is its holomorphic continuation.

Define the plaquette linearized curl by\[
 (\mathfrak d_Uv)(p)
          =\eta^{-1}\sum_{r=1}^4
                              \operatorname{Ad}(T_{r-1})v(b_r).
 \tag{PX115}\label{px:plaquette-linearized-curl}
\]Then \(\sum_rX_r=\eta^2\mathfrak d_Uv(p)\). If \(\|P-1\|\le\epsilon\eta^2\xi_p^{-2}\), with \(\xi_p\) the largest corner scale, summing \eqref{px:plaquette-hessian-exact} over the actual plaquettes gives\[
 \left|\langle v,\Delta_U^{\rm W}w\rangle
                       -\langle\mathfrak d_Uv,\mathfrak d_Uw\rangle_2\right|
              \le C_L\epsilon\|v\|_{\rm mass}\|w\|_{\rm mass}.
 \tag{PX116}\label{px:wilson-form-remainder}
\]Every bond is in at most six unoriented plaquettes. Each error has the factor \(\eta^2\epsilon\xi_p^{-2}\) from \(P-1\) and another \(\eta^2\) from the two insertions. Their product is the physical \(\eta^4\) measure times the natural inverse scale squared. There is no loss \(\epsilon\eta^{-2}\) in this form estimate.

This curl is close to the ordinary forward covariant curl \(d_Uv=\nabla_{U,\mu}v_\nu- \nabla_{U,\nu}v_\mu\). For example on a \(\mu\nu\) plaquette the difference of their numerator sums is\[
 (1-\operatorname{Ad}P)
           \bigl(\operatorname{Ad}U_\nu(x)v_\mu(x+\eta e_\nu)
                                                    +v_\nu(x)\bigr).
 \tag{PX117}\label{px:curl-convention-difference}
\]Direct multiplication of the four prefixes proves this identity. It gives \(\|\mathfrak d_Uv-d_Uv\|_2 \le C_L\epsilon\|v\|_{\rm mass}\), since \(\eta\le\xi_p\). In particular\[
 \langle v,\Delta_U^{\rm W}v\rangle
     \ge(1-C_L\epsilon)\|d_Uv\|_2^2
                                     -C_L\epsilon\|v\|_{\rm mass}^2
 \tag{PX118}\label{px:wilson-lower-form}
\]for a real field, after decreasing the fixed curvature ceiling. Positive bounded plaquette weights give the same statement with their lower and upper bounds. Signed layer profiles or additional current subtractions must be estimated separately.

On the full periodic graph, the covariant Hodge calculation is also explicit:\[
 \left|
   \|d_Uv\|_2^2+\|D_U^{\mathsf T} v\|_2^2
                              -\sum_\mu\|\nabla_{U,\mu}v\|_2^2
 \right|
                          \le C_L\epsilon\|v\|_{\rm mass}^2 .
 \tag{PX119}\label{px:covariant-hodge-error}
\]Expand both squares. The diagonal terms agree because each covariant shift is unitary. For different directions the uncancelled terms are commutators of the two shifts, or of a shift and an inverse shift. Their coefficients are conjugates of \(\operatorname{Ad}P-1\). The two difference denominators contribute \(\eta^{-2}\); the plaquette bound leaves \(C\epsilon\xi^{-2}\). There are a fixed number of neighboring terms at each bond. Adjacent scales are comparable, so Cauchy proves \eqref{px:covariant-hodge-error}. On a proper inner graph the summation by parts has additional boundary terms.

The full Feynman gauge uses \(D_UD_U^{\mathsf T}\). The actual fine Landau form instead uses \(D_URD_U^{\mathsf T}\). Put $\mathcal F_U=G'Q'^{\mathsf T}$ and
$M_U=Q'G'^2Q'^{\mathsf T}$. The difference of the two gauge terms
is the nonnegative real smoothing form\[
          D_U(1-R)D_U^{\mathsf T}
             =(D_U\mathcal F_U)M_U^{-1}
                                      (D_U\mathcal F_U)^{\mathsf T} .
 \tag{PX120}\label{px:landau-feynman-difference}
\]The projected estimate uses the constrained local models and their
parametrix, as explained after
\eqref{px:feynman-coefficient-comparison}. This retains the
subtraction together with the row mass. Its boundary model is the
all-bare local block, where $Q_\square=Q'_\square=1$ and
$R_\square=0$; no vanishing of the global $R$ at an interface
is asserted.

For the admitted complex perturbation \(E\), the actual plaquette Taylor formula gives \(\|P(U_E)-P(U)\|\le C_Lh\eta^2\xi_p^{-2}\). Indeed the first order boundary sum of \(E\) is its covariant curl; the remaining terms contain either two increments or one increment times a real connection on a local gauge patch. The first-order core bounds control each by that scale. Also \(\|d_{U_E}v-d_Uv\|_2 \le C_Lh\|v\|_{\rm mass}\). Subtracting \eqref{px:plaquette-hessian-exact} therefore proves\[
 |\langle v,(\Delta^{\rm W}_{U_E}-\Delta_U^{\rm W})w\rangle|
   \le C_Lh
       (\|d_Uv\|_2+\|v\|_{\rm mass})
       (\|d_Uw\|_2+\|w\|_{\rm mass}) .
 \tag{PX121}\label{px:complex-wilson-relative-form}
\]This uses a physical first-order perturbation radius, not an unscaled matrix norm for the fine precision. It is a genuine part of the complex comparison. The field-row comparison below supplies the corresponding row estimate. We must join it to the projected-current and boundary comparisons for the full precision.


\subsubsection{Compact dual bumps for the actual flat field rows}

We can also prove a uniform right inverse for the flat layered field rows. This is a different statement from a right inverse on site fields. The proof uses two longitudinal moments and a strictly triangular correction between adjacent levels.

Fix a top \(y\) of level \(j\ge1\), put \(n=L^j=2m+1\), and work on its fine block \(\eta\mathcal B_n+y\). Choose a field with just its \(\mu\)-component nonzero there:\[
 a_\mu(y+\eta u)
       =\psi\,\phi(u_\mu)
                         \prod_{\nu\ne\mu}\chi(u_\nu).
 \tag{PX125}\label{px:row-dual-bump}
\]It is zero outside the block and on bonds whose other endpoint is outside. For sufficiently large \(n\), take \(\chi(u)\) from a fixed even smooth profile of \(u/n\), supported in \((-1/4,1/4)\) and positive on a smaller interval. Write\[
 \beta=n^{-1}\sum_u\chi(u),\qquad
       \sigma=n^{-1}\sum_u(u/n)^2\chi(u).
 \tag{PX126}\label{px:row-profile-moments}
\]Both have positive lower and finite upper bounds independent of \(n\).

For a fine longitudinal position \(u\) in this block, the right coarse row in \eqref{px:flat-field-average} counts that bond \(u+m+1\) times. The left row counts it \(m-u\) times. No other coarse row sees it. After the transverse sums, their two response functionals are\[
 \ell_+(\phi)=\beta^3n^{-2}\sum_u(u+m+1)\phi(u),
 \qquad
 \ell_-(\phi)=\beta^3n^{-2}\sum_u(m-u)\phi(u).
 \tag{PX127}\label{px:two-row-moments}
\]Define\[
 \begin{split}
 \phi_+(u)&=\beta^{-4}\chi(u)
           +\frac{m}{n}\frac{\beta^{-3}}{\sigma}
                                             \frac{u}{n}\chi(u),\\
 \phi_-(u)&=\beta^{-4}\chi(u)
           -\frac{m+1}{n}\frac{\beta^{-3}}{\sigma}
                                             \frac{u}{n}\chi(u).
 \end{split}
 \tag{PX128}\label{px:explicit-dual-profiles}
\]Evenness gives \(\sum u\chi(u)=0\). Substitution into \eqref{px:two-row-moments} now gives\[
 (\ell_+(\phi_+),\ell_-(\phi_+))=(1,0),\qquad
       (\ell_+(\phi_-),\ell_-(\phi_-))=(0,1).
 \tag{PX129}\label{px:dual-profile-exactness}
\]For example the sum of the two functionals is \(\beta^3n^{-1}\sum\phi\). For either profile this is one; the odd term in \(\phi_+\) sets the left functional to zero. The right one is then one. The calculation for \(\phi_-\) is the same with the other moment. The signed profiles are uniformly bounded and their first differences are \(O(n^{-1})\).

The finitely many odd values below this smooth-profile range need no limiting argument. Put the three transverse profiles at their center point and use longitudinal bonds with positions \(-1,0\). Both endpoints of these bonds are in the block for \(n\ge3\). Their two response columns form\[
             n^{-5}
                    \begin{pmatrix}m&m+1\\m+1&m\end{pmatrix}.
 \tag{PX130}\label{px:finite-row-profile-matrix}
\]Its determinant is \(-n^{-9}\). Its inverse therefore gives the required two exact dual profiles, with a bound depending only on the fixed cutoff between this finite set and the smooth range. That is a universal constant, not a depth dependence.

Assign a determining row \(c\) to one of its level-\(j\) top endpoints \(y\). An outward row has its interior endpoint as that top. Use \(\phi_+\) or \(\phi_-\) according to the chosen endpoint. Its seed \(E_c^0\psi\) is supported on fine bonds strictly inside that pillar and satisfies, on the full regular level-\(j\) grid,\[
                       Q_j^0E_c^0\psi=\delta_c\psi .
 \tag{PX131}\label{px:one-row-seed}
\]The transverse support lies in a single coarse transverse block, and \eqref{px:dual-profile-exactness} removes the only other longitudinal row. For a bare row take its single fine bond value. It has both endpoints in the bare layer.

Let \(E^0\psi\) be the sum of these seeds at every level. There are at most eight assigned seeds at a top. Their supports belong to distinct pillars, and the collars in their gradient norm have bounded overlap. The amplitude and first-difference bounds just proved give\[
 \|E^0\psi\|_{\rm mass}^2+
                \|\partial E^0\psi\|_2^2
                     \le C\sum_{c\in\Gamma}\xi_c^2|\psi(c)|^2 .
 \tag{PX132}\label{px:seed-energy}
\]The factor on the right is forced by the dimensions: the supporting volume is \(\xi_c^4\), and a first difference or mass factor is \(\xi_c^{-1}\). For a bare seed the same calculation is \(\eta^4\eta^{-2}|\psi|^2\).

\subsubsection{The correction between adjacent levels}

A seed of level \(j\) has zero value under every lower determining row. Indeed its fine support is in \(\Omega_j\), whereas a lower row's entire carrier is outside the next region of that lower level. For \(i<j\), that next region contains \(\Omega_j\). It also has zero value under every row of level \(i\ge j+2\): such a carrier is inside \(\Omega_{i-1} \subset\Omega_{j+1}\), whereas the seed pillar is outside \(\Omega_{j+1}\). The same argument uses the two bare endpoints for \(j=0\). Thus only the next level can receive an unwanted value.

The full flat averages compose exactly:\[
                  Q_{j+1}^{0,{\rm full}}
                         =Q_L^0Q_j^{0,{\rm full}} .
 \tag{PX133}\label{px:flat-average-composition}
\]In the defining sums, the transverse digits are \(nr+u\), where \(r\) is in the centered \(L\)-range and \(u\) in the centered \(n\)-range. They tile the centered \(Ln\)-range. The longitudinal digits are \(nt+s\), with \(0\le t<L\), \(0\le s<n\). They tile \(0,\ldots,Ln-1\). The product normalization is \(L^{-5}n^{-5}\). This proves \eqref{px:flat-average-composition}, including the per-length normalization.

Let the row Hilbert norm be \(\|\psi\|_{\rm row}^2 =\sum_c\xi_c^2|\psi(c)|^2\). A regular \(L\)-step average has row coefficient sum one and column coefficient sum \(L^{-4}\). The column identity can be checked directly: the two longitudinal counts add to \(L\), while the three transverse coordinates select one block. The ordinary counting \(\ell^2\) operator norm is therefore at most \(L^{-2}\) by Schur. The row norm at the next level has one extra factor \(L\) in its square root. Hence the map between these row norms has norm at most \(L^{-1}\).

Writing \(T\) for the next-level leakage, \eqref{px:one-row-seed}--\eqref{px:flat-average-composition} give\[
              Q^0E^0=1+T,\qquad
                  T_{ij}=0\ (i\ne j+1),\qquad
                         \|T\|_{{\rm row}\to{\rm row}}\le L^{-1}.
 \tag{PX134}\label{px:row-triangular-leakage}
\]Restriction to actual determining rows and injection from their subset of the regular grid cannot increase the norm. Different level blocks of \(T\) have different input and output levels, so the whole direct-sum norm is the supremum of those block norms. No sum over the number of levels is introduced.

We have proved the explicit flat row lift\[
 E_{\rm row}^0=E^0(1+T)^{-1},\qquad
 Q^0E_{\rm row}^0=1,\qquad
 \|E_{\rm row}^0\psi\|_{\rm mass}^2+
       \|\partial E_{\rm row}^0\psi\|_2^2
       \le \frac{C}{(1-L^{-1})^2}\|\psi\|_{\rm row}^2 .
 \tag{PX135}\label{px:flat-layered-row-lift}
\]At finite depth the inverse is the finite triangular series; its norm is bounded by the displayed geometric series. The same series converges in the row Hilbert space for an infinite compatible sequence. The assertion is a right inverse with an energy bound, not an estimate of a nonlinear background map.

At a real flat connection the same construction uses a flat unitary frame on each pillar. The response at the other coarse endpoint differs by its fixed unitary transport. The two scalar moments still annihilate the unwanted row; the desired row is normalized in its own color frame. Schur uses the same positive scalar coefficient majorants, so its constants are unchanged. All profile supports lie inside their pillars. This handles periodic flat holonomies without choosing a single globally small logarithm. Under this gauge transfer the gradient in the energy bound is \(\nabla_U\), or equivalently the ordinary component gradient in its flat parallel frame. An ordinary gradient in an arbitrary original gauge is not bounded by that statement. No analogous bare \(L^2\) field-row estimate at a general nonflat field has been assumed.

\subsubsection{The field-row Gram bounds}

The last proof gives the flat instance of a general useful duality reduction. Let \(K\) be a real symmetric positive fine precision and \(G=K^{-1}\). Let \(\Xi\) multiply each field row by its length. Suppose the actual row map has an extension \(E_{\rm row}\) such that\[
 QE_{\rm row}=1,\qquad
 \langle E_{\rm row}\psi,K E_{\rm row}\psi\rangle
                            \le C\|\Xi^{-1}\psi\|_{2,{\rm rows}}^2 .
 \tag{PX136}\label{px:field-row-extension-premise}
\]The node pairing here has volume \(\xi_c^4\); its last norm is just \(\sum_c\xi_c^2|\psi(c)|^2\). Testing with \(E_{\rm row}\Xi^2\psi\) gives\[
 \|\Xi\psi\|_{2,{\rm rows}}^2
   =\langle Q^{\mathsf T}\psi,E_{\rm row}\Xi^2\psi\rangle
   \le C^{1/2}\langle\psi,QGQ^{\mathsf T}\psi\rangle^{1/2}
                                       \|\Xi\psi\|_{2,{\rm rows}} .
 \tag{PX137}\label{px:field-row-gram-lower}
\]Thus \(QGQ^{\mathsf T}\ge C^{-1}\Xi^2\). No pointwise delta-function test in the fine field was used.

An upper bound needs the separate actual trace estimate\[
               \|\Xi^{-1}Qa\|_{2,{\rm rows}}^2
                                      \le C\langle a,Ka\rangle .
 \tag{PX138}\label{px:field-row-trace-premise}
\]Indeed apply it to \(a=G^{1/2}f\). It proves \(\|\Xi^{-1}QG^{1/2}\|\le C^{1/2}\) and therefore \(QGQ^{\mathsf T}\le C\Xi^2\). For the flat Feynman form, \eqref{px:flat-field-row-upper}, \eqref{px:flat-feynman-gap} and \eqref{px:flat-layered-row-lift} prove both premises. For the flat Landau form, \eqref{px:flat-layered-row-lift} proves the upper energy cost in \eqref{px:field-row-extension-premise}, since \(0\le R\le1\); the needed fine coercivity is still a distinct step.

Normalize the row covariance as \(N=\Xi^{-1}QGQ^{\mathsf T}\Xi^{-1}\). If its exponential cell Schur bound has also been proved, the exponentially weighted conjugation argument applies to \(N\). Conjugating $N$ by $e^{\delta d(\cdot,c_0)}$ changes it by
$O(\delta)$ in the two-sided Schur norm: use the exponential
weight already present in its kernel bound and
$|e^{\delta t}-1|\le\delta|t|e^{\delta|t|}$. For fixed sufficiently
small $\delta$, the real lower bound keeps the conjugated inverse
bounded. Undoing the conjugation gives exponential entries of $N^{-1}$. In physical weighted orthonormal row coordinates, undoing normalization gives the counting kernel estimate\[
                   |(QGQ^{\mathsf T})^{-1}(c,d)|
                \le C\,\xi_c^{-3}\xi_d\,e^{-m_3d(c,d)} .
 \tag{PX139}\label{px:field-row-inverse-counting}
\]To check the powers, an orthonormal entry of \(N^{-1}\) becomes \((\xi_d^2/\xi_c^2)\) times that entry on unnormalized values. The two exterior \(\Xi^{-1}\) factors leave \(\xi_c^{-3}\xi_d\). This differs from the site-Gram counting weight \(\xi_y^{-4}\): the site Gram contains \(G'^2\), whereas this row covariance contains \(G\).

Finally a nonflat row extension is reduced to a definite error estimate. Starting from a specified flat or locally flat extension \(E_0\), set \(D_Q=QE_0-1\). If the actual primitive comparison proves\[
          \|D_Q\|_{{\rm row}\to{\rm row}}\le\vartheta<1,
 \tag{PX140}\label{px:nonflat-row-error}
\]then \(E_{\rm row}=E_0(1+D_Q)^{-1}\) is an actual right inverse, with its energy cost multiplied by \((1-\vartheta)^{-2}\), provided that cost for \(E_0\) has been proved in the actual precision. We establish the row error and row trace estimates for the admitted domain below.

\subsubsection{The actual straightened row is a finite analytic recursion}

We now obtain the needed row coefficient bounds from the actual formula, rather than from a proposed inverse. Use raw one-step link logarithms on a lattice whose spacing has been set to one. Let \(T_{yx}(W)\) be the internal mean transport of the block chart, with its specified coordinate-order path multiset. For a reference \(W\) and a raw increment \(a\), put\[
 \begin{split}
 u_W(a)(y)
   &=\exp\!\left[L^{-4}\sum_{x\in B_L(y)}
       \log\!\left(T_{yx}(e^{ia}W)T_{yx}(W)^{-1}\right)\right],\\
 F_W(a)(c)
   &=\frac1i\log\!\left[
        u_W(a)(c_-)^{-1}\mathcal M(e^{ia}W)(c)
                    u_W(a)(c_+)\mathcal M(W)(c)^{-1}\right].
 \end{split}
 \tag{PX141}\label{px:actual-one-step-straightening}
\]All logarithms are the regular branch. The first line is an exponential of averaged relative logarithms; it is not another unspecified mean. The internal transports themselves use the actual Federbush mean defined in the block chart. This distinction fixes both levels of averaging.

For \(\|\log W/i\|_\infty\) and \(\|a\|_\infty\) below \(c/L\), these are holomorphic finite functions. Every path has at most \(CL\) bonds. Products of their exponentials remain in a fixed logarithmic neighborhood; the mean contraction proved in the block chart is uniform in the number of paths. After a fixed decrease of \(c\), the expression inside the second logarithm remains in that same neighborhood. Thus \(\|F_W(a)\|_\infty\le C\) on this joint raw ball, with universal \(C\).

Differentiating at \(W=1,a=0\) cancels the two internal path sums from the plain average against the two factors \(u\). The remaining sum is exactly the straight segment average \(S_L\), with raw row sum \(L\) and raw column sum \(L^{-3}\). It is essential to use both factors in \eqref{px:actual-one-step-straightening}; the plain average has a different derivative.

Let \(A_W(a)=D_aF_W(a)\). Two scalar Cauchy circles, in arbitrary sup-norm directions of the reference logarithm and of \(a\), give\[
 \begin{aligned}
 \|A_W(a)\|_{\rm row}
    &\le L(1+C_L\delta),&
 \|A_W(a)\|_{\rm col}
    &\le L^{-3}(1+C_L\delta),\\
 \sup_b\sum_{c,b'}\|D_bD_{b'}F_W(a)(c)\|
    &\le C_L,&
 \delta&=\|\log W/i\|_\infty+\|a\|_\infty .
 \end{aligned}
 \tag{PX142}\label{px:finite-one-step-schur}
\]Here and below the norm of a color entry is its finite matrix norm. We explain why the derivative estimate is independent of a field-coordinate count. The two circles bound a second directional derivative by \(C L^2\). For a fixed output color row, its operator norm from bounded input colors controls the sum of its entry norms, up to the fixed color dimension: sum the dual norms of its three scalar rows. A fixed input bond affects only the coarse rows at the endpoints of the one or two blocks containing it. Their number is fixed. Apply the same argument with one input bond fixed to obtain the last line. Dividing the row and column errors by their flat values costs only powers of \(L\). For definiteness, all one-step constants in \eqref{px:finite-one-step-schur} may be enlarged to \(C L^{12}\). This deliberately loose polynomial bound is before any number of levels is introduced.

The multilevel recursion is literal. Write \(W_r=\mathcal M^r(W_0)\), \(a_0=\eta A\), and\[
                 a_{r+1}=F_{W_r}(a_r),\qquad
                   v_{r+1}=v_r|_{\rm next}\,u_{W_r}(a_r).
 \tag{PX143}\label{px:actual-straightened-recursion}
\]Equivariance of the plain mean gives, by induction,\[
                \mathcal M^r(e^{i\eta A}W_0)
                                  =(e^{ia_r}W_r)^{v_r}.
 \tag{PX144}\label{px:actual-factorization}
\]Indeed apply \(M\) to the previous factorization, restrict its gauge to the next nodes, and then use \eqref{px:actual-one-step-straightening}. The gauge factors multiply in the order in \eqref{px:actual-straightened-recursion}. They do not alter the recursion for \(a_r\).

\subsubsection{Summability of the raw radii gives a depth-uniform row bound}

For the row estimates from this point on, \(\epsilon\) is a common ceiling for the real curvature and the local real-gauge bounds. If the curvature and local-gauge ceilings are supplied separately, take their maximum for \(\epsilon\). Existence of those local gauges remains the primitive regular-pair hypothesis. No bound for \(a_{\rm loc}\) is inferred from curvature alone.

Fix a row of length \(\xi=n\eta\), \(n=L^j\). On its determining carrier, the regular-pair gauge clause gives a real local potential and its first difference at scale \(\xi\), with at most the fixed adjacent-layer factor at a seam. In the local gauge write $\alpha_r=\|\log W_r/i\|_\infty$ and
$\omega_r=\|\partial W_r-1\|_\infty$. The one-step product,
plaquette-sliding and Taylor estimates give
\[
 \alpha_{r+1}\le L\alpha_r+C_L(\alpha_r^2+\omega_r),\qquad
 \omega_{r+1}\le L^2(1+C_L\alpha_r)\omega_r+C_L\omega_r^2.
\]
For the first estimate a straight product contributes $L\alpha_r$
and its quadratic error, while each closed contour comparison
contributes $C_L\omega_r$. For the second, put the one-step patch
in its complex comb gauge. Its non-tree bonds have size
$C_L\omega_r$; the linear flux through a coarse plaquette is
$L^2$ times the fine flux, and the Taylor error is quadratic.
Conjugating back costs $1+C_L\alpha_r$. These are finite one-step
estimates, with $L$ fixed.

At the bottom of the carrier,
$\alpha_0\le C_L(\epsilon+h)/n$ and
$\omega_0\le C_L(\epsilon+h)/n^2$. Divide the two recurrences by
$t_r$ and $t_r^2$, respectively. Since
$\sum_{r<j}t_r\le(L-1)^{-1}$, a first-exit argument bounds their
normalized values at a fixed multiple of the initial ceiling.
After the fixed weak-field shrinkage this proves\[
 \|\log W_r/i\|_\infty
                    \le C_L(\epsilon+h)t_r,\qquad
                    t_r=L^r/n,\qquad r<j .
 \tag{PX145}\label{px:reference-raw-radii}
\]These estimates are used only on the dependence cone of that row. They also follow by extending its local small potential with a cutoff on a fixed larger cube and then using locality. The cutoff keeps its first difference at the same scale. The torus room required for nonflat local classes permits this cube; the flat periodic formulas need no such extension. No gauge bound growing with \(M_1\) is used.

Allow an arbitrary fine increment with \(\xi\|A\|_\infty\le r\). It need not have small fine derivatives. On the recursion ball, \eqref{px:finite-one-step-schur} and \(F_W(0)=0\) give\[
 \|a_{r_0}\|_\infty\le 2r\,t_{r_0}
                  \quad(0\le r_0\le j)
 \tag{PX146}\label{px:fluctuation-raw-radii}
\]when \(\epsilon+h+r\) is below a fixed threshold depending on \(L\). To avoid confusing the index with the radius, write \(p_s=\|a_s\|_\infty/t_s\). As long as \(p_s\le2r\), Taylor's formula and \eqref{px:finite-one-step-schur} imply\[
 p_{s+1}\le p_s
       [1+C_Lt_s(\epsilon+h+p_s)] .
 \tag{PX147}\label{px:radius-discrete-gronwall}
\]The sum of \(t_s\) is below \((L-1)^{-1}\). Choose the fixed ceiling so that \(C_L(\epsilon+h+2r)/(L-1)<\log2\). Multiplying the factors proves \eqref{px:fluctuation-raw-radii} by first exit. The raw radii also stay strictly inside the \(c/L\) one-step balls. Finally \(\|u_{W_s}(a_s)-1\|\le C_L\|a_s\|_\infty\); their condition numbers have product at most \(\exp(C_L\sum_s\|a_s\|_\infty)\). A further fixed decrease makes this at most two, supplying the room for the complex covariance in \eqref{px:actual-factorization}.

Put \(A_s=D_aF_{W_s}(a_s)\) and \(P_{b,a}=A_{b-1}\cdots A_a\). The same summable radii in \eqref{px:finite-one-step-schur} give\[
 \|P_{b,a}\|_{\rm row}\le C L^{b-a},\qquad
                   \|P_{b,a}\|_{\rm col}\le C L^{-3(b-a)} .
 \tag{PX148}\label{px:multilevel-row-column-products}
\]The constant is bounded by the product of \(1+C_L(\|\log W_s/i\|_\infty+\|a_s\|_\infty)\). It is independent of \(a,b,j\). Multiplying only uniform bounds larger than one at each level would miss this point.

The per-length row logarithm is \(q_j(A)=a_j/\xi\). Its derivative with respect to the physical fine coordinates is \((\eta/\xi)P_{j,0}=n^{-1}P_{j,0}\). Consequently\[
 \|Dq_j(A)\|_{\rm row}\le C,\qquad
 \|Dq_j(A)\|_{\rm col}\le C n^{-4},\qquad
                   |D_bq_j(A)(c)|\le C(\eta/\xi)^4 .
 \tag{PX149}\label{px:actual-row-density}
\]These are counting-coordinate derivatives. A functional derivative uses the additional \(\eta^{-4}\). Telescoping the products against their flat factors also gives row error \(C_L(\epsilon+h+r)\) and column error \(C_L(\epsilon+h+r)n^{-4}\), in the common local unitary frames.

An actual row carrier has only its two endpoint blocks. A fixed fine bond in a layer can occur in carriers of only adjacent levels: the outer block is inside the previous region and both endpoint blocks are outside the next. At each of those levels the number of rows is fixed. Thus \eqref{px:actual-row-density} implies for the full layered field map\[
       \|\Xi^{-1}Qa\|_{2,{\rm rows}}
                                  \le C_L\|a\|_{\rm mass}.
 \tag{PX150}\label{px:actual-row-mass-trace}
\]For a single row, its coefficient sum is bounded and each entry is at most \(C(\eta/\xi_c)^4\). Cauchy therefore gives the analogue of \eqref{px:flat-row-ltwo}; the carrier overlap argument of \eqref{px:flat-field-row-upper} then proves \eqref{px:actual-row-mass-trace}. This estimate has been obtained from the actual finite recursion, not by treating an infinity norm bound alone as an \(L^2\) bound.

\subsubsection{The two-legged constraint coefficient has a uniform moment}

The same computation gives the coefficient estimate needed for the quadratic current term. Let \(n_s=L^s\) and \(m_s=L^{j-s-1}\). Twice differentiating the recursion gives\[
 D^2_{a_0}a_j
      =\sum_{s<j}P_{j,s+1}
                      D^2F_{W_s}(a_s)[P_{s,0}\,\cdot,P_{s,0}\,\cdot].
 \tag{PX151}\label{px:actual-hessian-chain}
\]Fix the first fine bond \(b\) and sum absolutely over the second. For the term indexed by \(s\), the largest entry of the outer product is at most \(C m_s^{-3}\), by its column sum. The sum over the first inner leg is at most \(C n_s^{-3}\). The second inner row sum is at most \(C n_s\). The remaining sum over the one-step Hessian, with its first input fixed, is bounded by \eqref{px:finite-one-step-schur}. Thus\[
 \sum_{b'}\|D_{a_0(b)}D_{a_0(b')}a_j(c)\|
                 \le C_L\sum_{s<j}m_s^{-3}n_s^{-2}.
 \tag{PX152}\label{px:raw-hessian-leg-sum}
\]This step uses a column sum on the fixed leg and a row sum on the summed leg. Using two row sums would lose the fine density.

Physical input derivatives contribute \(\eta^2\); the output per-length normalization contributes \(\xi^{-1}\). Since \(m_s^{-3}n_s^{-2} =L^3n_s/n^3\) and \(\sum_{s<j}n_s\le n/(L-1)\), we obtain\[
 \sup_b\sum_{b'}\|D_bD_{b'}q_j(A)(c)\|
                  \le C_L\,\xi\,(\eta/\xi)^4 .
 \tag{PX153}\label{px:constraint-hessian-leg-moment}
\]There is no number-of-levels factor. The geometric sum is dominated by its final scales. Symmetry of the second derivative gives the corresponding other-leg sum.

Let \(\omega(c)\) be an arbitrary row current with \(|\omega(c)|\le\gamma_\omega\xi_c^{-3}\). Pair the Hessian at zero with the physical row volume. Equation \eqref{px:constraint-hessian-leg-moment}, Schur and the same bounded carrier overlap give\[
 \left|\sum_c\xi_c^4
              \langle\omega(c),D^2q_c(0)[u,v]\rangle_{\rm col}\right|
                  \le C_L\gamma_\omega
                                      \|u\|_{\rm mass}\|v\|_{\rm mass}.
 \tag{PX154}\label{px:quadratic-current-form-bound}
\]For clarity, the row prefactor \(\xi_c^4|\omega(c)|\) is bounded by \(\gamma_\omega\xi_c\). Multiplying it by the leg moment \(C_L\xi_c(\eta/\xi_c)^4\) leaves \(C_L\gamma_\omega\eta^4\xi_c^{-2}\) on each fine leg. Adjacent scales replace \(\xi_c\) by the fine bond scale at only a fixed cost.

The holomorphic ball was used to extract these coefficients. The quadratic form in \eqref{px:quadratic-current-form-bound} then extends to every finite fine vector by its coefficient formula. It does not require a Gaussian sample to lie in that holomorphic ball. For the second variant one must still bound the actual induced current \(\mathsf H[J]^{\mathsf T} J\) in this row norm, and bound \(\mathsf S\) in the weighted fine energy spaces. Those are separate operator readouts.

\subsubsection{The nonflat field-row lift follows from the row comparison}

The coefficient estimates allow one to complete \eqref{px:nonflat-row-error} at a regular pair. Keep the real local gauges of its center. Construct the profiles \eqref{px:row-dual-bump} in each pillar in such a gauge, and undo the gauge. For a seed assigned to the far endpoint of a row, normalize its color by the fixed unitary transport of that row. Denote this operator again by \(E^0\). It is fixed at the real center while the complex increment varies.

Two local real gauges used on the same pillar differ there by a nearly constant unitary matrix. To prove this, write their connection equation on each fine bond and telescope along a coordinate path of length \(C\xi_j\). Both gauge potentials have size at most \(C_L\epsilon/\xi_j\). The transition therefore differs from its root value by at most \(C_L\epsilon\). The first-order class on a fixed larger cube supplies both gauges. Only a path of size \(C_L\xi_j\) is used, not the full \(M_1\xi_j\) cube.

It follows from \eqref{px:actual-row-density} and this transition bound that the same-level seed response is the identity with an error of row and column norm \(C_L(\epsilon+h)\) in the row energy norm. Only a fixed number of same-level rows meet a seed pillar. The next-level response has a leading part with the positive scalar majorant of \(Q_L^0\), multiplied by fixed unitary color matrices, and an error. There are at most \(C L^4\) lower-level seeds in one such carrier; a seed affects at most a fixed number of next-level rows. An error entry of size \(C_L(\epsilon+h)\) therefore has row sum \(C_LL^5(\epsilon+h)\) and column sum \(C_LL(\epsilon+h)\) after the row-energy normalization. Schur bounds this error by \(C_LL^3(\epsilon+h)\). Absorb this fixed power into \(C_L\).

All other level blocks are zero by the carrier argument preceding \eqref{px:flat-average-composition}. The leading next-level part still has norm at most \(L^{-1}\); Schur uses only the scalar majorants, not a compatible global choice of flat color frames. We have proved\[
 \|Q_{\mathbf U}E^0-1\|_{{\rm row}\to{\rm row}}
                 \le L^{-1}+C_L(\epsilon+h)=:\vartheta<1 .
 \tag{PX155}\label{px:actual-row-lift-error}
\]The ceiling needed here depends on \(L\) and not on depth, volume or \(M_1\). The actual right inverse is\[
             E_{\rm row}=E^0(Q_{\mathbf U}E^0)^{-1}.
 \tag{PX156}\label{px:actual-layered-row-lift}
\]Both its squared mass norm and its covariant gradient energy are bounded by \(C_L(1-\vartheta)^{-2}\|\psi\|_{\rm row}^2\). For the gradient of a seed, expand in its local gauge: the profile derivative is of size \(C\xi_j^{-1}|\psi|\) and the connection term has the same bound. Supports and their fine collars have bounded overlap. The small complex perturbation adds at most \(C_Lh\) times the mass norm. This proves the stated energy bound, including bare seeds.

There is a corresponding actual Feynman Poincare inequality. Repeat \eqref{px:flat-row-mean-comparison} in a common local real gauge for a pillar and its chosen row carriers. The row comparison after \eqref{px:actual-row-density} has both row and column error bounds. Schur therefore bounds its local \(L^2\) functional norm by \(C_L\epsilon\xi_j^{-2}\). Ordinary gradients and covariant gradients differ there by \(C_L\epsilon\xi_j^{-1}\) times the field. Consequently \eqref{px:flat-field-local-poincare}, summed over the carriers, gives\[
 \|a\|_{\rm mass}^2
       \le C_L\bigl(\|\nabla_Ua\|_2^2+
                         \|a^{1/2}Q_Ua\|_{2,{\rm rows}}^2\bigr)
                         +C_L\epsilon^2\|a\|_{\rm mass}^2 .
 \tag{PX157}\label{px:actual-feynman-poincare}
\]The repeated letter \(a\) in \(a^{1/2}\) is the fixed positive row weight, whereas the arguments are fine fields. Equivalently that term is \(\sum_cc_c\xi_c^2|Q_Ua(c)|^2\). Absorb the last term by the fixed weak-field ceiling. Only carrier differences in the interior are used; the bare boundary fields are controlled by their identity rows. Thus the same argument applies to an active component-gradient form containing all those interior differences.

The full complex Feynman form is stable on a fixed physical ball. The actual row comparison gives\[
 \|a^{1/2}(Q_{\mathbf U}-Q_U)v\|_2
                       +\|(\nabla_{\mathbf U}-\nabla_U)v\|_2
                                \le C_Lh\|v\|_{\rm mass}.
 \tag{PX158}\label{px:feynman-coefficient-comparison}
\]For the first term use the row and column differences, the physical row volumes and the bounded carrier overlap. For the second use \(\operatorname{Ad}(e^{i\eta E})-1=O(\eta E)\). Expanding the two bilinear squares and using \eqref{px:actual-feynman-poincare} gives a relative form perturbation \(C_Lh\). Its inverse exists by the normalized Neumann series. This completes the component Feynman estimate. The full projected
estimate can be completed by the same local construction.
For the Landau precision use the constrained one- and two-scale
models of (c), rather than subtracting
\eqref{px:landau-feynman-difference} from a Feynman lower bound.
The decomposition $A=A_2+H_jB-\partial H'_j\vartheta$ splits
the flat form orthogonally with its row mass. The three terms in
(d) comparing $\partial(1-R)\partial^{\mathsf T}$ with its
local model cost $CM_1^{-1}$ for commutators and $Ce^{-cM_1}$
for changes of domain and separated tails. Their parametrix
therefore has a convergent error series. Near the outer boundary
use the massive all-bare local blocks, whose flat inverses satisfy
$a_0\le\eta^2G_\square^{-1}\le a_0+16$ on arbitrary bond sets.
The small local perturbation is treated in the same way; no
periodic inner-Hodge identity at the boundary is required.

Include the current corrections in their defining order:
\[
 \begin{aligned}
 \mathsf G
   &=\mathsf G_0\sum_{n\ge0}
                    (\Delta'_\pi[\mathbf J]\mathsf G_0)^n,\\
 \mathsf G_1
   &=\mathsf G\sum_{n\ge0}
      (\mathsf S^{\mathsf T}
       \Delta^{(2)}[\mathsf H^{\mathsf T}\mathbf J]
                                  \mathsf S\mathsf G)^n.
 \end{aligned}
\]
These are the expansions of
\cite[(3.130)--(3.138)]{B85c}, with the assigned row and site
maps. Transfer each differential factor to the neighboring inverse
on the left or right; no derivative of the independent current
is needed. The row and column estimates and
\eqref{px:constraint-hessian-leg-moment} give the constraint
term its scale factors. Expand $\mathsf H$ inside the induced
multiplier: its local part stays in the local inverse and its
remaining walks give the small tree vertices of (g). The local
coefficients are small with the regularity and current radii,
so the same construction gives the projected inverse and its
potential and gradient bounds. We must check the species weights,
the derivative allocation and the adjacent-level endpoint equations
when identifying these products with the actual blocks. The
covariant Laplacian member excluded in
\cite[Theorem~3.12]{B85c} is not asserted here.


\subsection*{Differentiating the physical response}
\label{ym4:exp500-bg:response}
We give the calculations behind the response and current readouts in
the preceding construction. Every matrix below is the derivative of
the assigned nonlinear equation. The constraint-curvature term and
the outward datum factors remain present throughout.

\subsubsection*{The actual equations and their differentiated right-hand sides}

Use the real reference field $U_0$ and write the varied fine field as
$e^{i\eta A}U_0$.  For the multilevel determining set $\mathbf B$, the
straightened average is denoted by $\mathcal Q^{\rm st}(A)$ and its
derivative by $K_A$.  This is the nonlinear straightening of the group
average, rather than the straight linear average.  In (BGR1)--(BGR9), $\mathcal Q^{\rm st}$, $B$ and
$\mathsf b$ use raw row logarithms, as in
\eqref{ym4:exp500-bg:mb:reduced-datum}; the fine field $A$
still has the physical normalization $e^{i\eta A}U_0$.
The row multiplier is dual to this raw convention.
With the reduced datum
$\mathsf b(B)$, the physical critical pair satisfies
\begin{align*}
 F'(A)&=K_A^{\mathsf T}\omega,
       &\mathcal Q^{\rm st}(A)&=\mathsf b(B),\\
 RD^{\mathsf T}A&=0,
       &F(A)&=A_{\Omega_0}(e^{i\eta A}U_0).
 \tag{BGR1}
\end{align*}
Here $R$ and $D$ are formed at the reference field.  The first equation is the physical stationarity identity proved
above by testing the remaining gauge directions; its restriction
to the slice suffices for the following linear algebra.
When the reference is changed with the datum, differentiating the slice
also contributes $(\dot R D^{\mathsf T}+R\dot D^{\mathsf T})A$.
We keep the reference fixed in (BGR1).  A calculation that moves the
reference and then omits this contribution does not differentiate (BGR1).

The datum forcing is $Db(B)\dot B$, where $b$ is the exact reduced
datum in \eqref{ym4:exp500-bg:mb:reduced-datum}. Its diagonal
logarithmic differentials and its finite two-step inverse were computed
in \eqref{ym4:exp500-bg:mb:seam-datum-nilpotence}. The outward
block factors therefore enter this forcing with their actual
interior-row derivatives.

Let $\mathcal V$ be the fixed linear slice, after the residual node gauge
has been fixed.  Its fine pairing and the datum pairing are the actual
cell-weighted pairings.  Write $q=K_A|_{\mathcal V}$ and let $q^\sharp$
be its transpose for these pairings.  The constrained second variation is
\[
 \mathcal K=F''(A)
       -\sum_a \mu_a\omega_a
                       (\mathcal Q^{\rm st}_a)''(A)
 \quad\hbox{on }\mathcal V.
 \tag{BGR4}
\]
The weights $\mu_a$ in this coordinate expression are the datum cell
weights.  Equivalently the second term is the Hessian of the scalar
pairing $\langle\omega,\mathcal Q^{\rm st}(A)\rangle$.  This latter
description fixes the convention if a different coordinate basis is used.
The curvature term in (BGR4) cannot be replaced by the Hessian of $F$
alone.  It is the usual second fundamental form of the constraint fibre.

Differentiating (BGR1) at a fixed reference gives the exact system
\[
 \mathcal K\dot A-q^\sharp\dot\omega=0,
       \qquad q\dot A=D\mathsf b(B)\dot B.
 \tag{BGR5}
\]
More generally, for explicit parameters $p$ in the action or
constraints, with the slice fixed, put $\mathcal C(A,p)=0$ and
$\mathcal E(A,\omega,p)=F'_p(A)-\mathcal C_A^\sharp\omega$.
The right-hand sides are then
\[
 f=-D_p\mathcal E\,\dot p,\qquad
 b=-D_p\mathcal C\,\dot p,\qquad
 \mathcal K\dot A-q^\sharp\dot\omega=f,\quad q\dot A=b.
 \tag{BGR6}
\]
Every differentiation in $D_p$ is made before substituting the solution
$A(p),\omega(p)$.  This distinction prevents an implicit extra response
term on the right-hand side.

An independent-current extension is an additional specified operator.
Formula (BGR4) identifies the physical constrained Hessian.  It does not
identify every separately defined $\Delta'_\pi[J]$ with that Hessian.
Likewise, an absolute contour gauge functional and its relative version
have the same quadratic gauge block at an axial field, but generally have
different blocks away from that field.  The complete gauge and current
normal form must be matched before an independently named propagator can
be inserted into (BGR5).

\subsubsection*{A complete block inverse, including the multiplier}

The following identity recovers the multiplier along with the field.
It is finite-dimensional and exact. It also applies to
bounded operators when the displayed inverses and compositions exist on
the stated spaces.  Every transpose in the complex calculation is
bilinear, with no conjugation.

\begin{proposition}[Response from a primitive constrained inverse]
Suppose $q:\mathcal V\to\mathcal D$ is onto, $E:\mathcal D\to\mathcal V$
satisfies $qE=I$, and $P:\mathcal X\to\mathcal V$ is an isomorphism onto
$\ker q$.  Assume
\[
 H=P^\sharp\mathcal K P\quad\hbox{is invertible},\qquad
 L=P H^{-1}P^\sharp,\qquad T=E-L\mathcal K E.
 \tag{BGR7}
\]
Then the unique solution of (BGR6) is
\begin{align*}
 \dot A&=Lf+Tb,\\
 \dot\omega&=E^\sharp(\mathcal K L-I)f
                          +E^\sharp\mathcal K T b.
 \tag{BGR8}
\end{align*}
If $\mathcal K^\sharp=\mathcal K$, this simplifies to
\[
 \dot\omega=-T^\sharp f+M b,\qquad
 M=E^\sharp\mathcal K E
                     -E^\sharp\mathcal K L\mathcal K E,
       \qquad M^\sharp=M.
 \tag{BGR9}
\]
\end{proposition}

\begin{proof}
Write $\dot A=Eb+Pv$.  Applying $P^\sharp$ to the first equation of
(BGR6), and using $qP=0$, gives
$Hv=P^\sharp(f-\mathcal K Eb)$.  This determines $v$ and gives the first
formula of (BGR8).  The residual $\mathcal K\dot A-f$ annihilates $\ker q$.
It therefore belongs to the range of $q^\sharp$.  Since
$E^\sharp q^\sharp=I$, its unique multiplier is the second formula of
(BGR8).  Conversely these formulas satisfy both equations, so there is
no lost multiplier equation.  For a homogeneous solution, the same
argument gives $v=0$, then $\dot A=0$ and $\dot\omega=0$.

For the symmetric assertion, $H^\sharp=H$ and $L^\sharp=L$.
Thus $T^\sharp=E^\sharp-E^\sharp\mathcal K L$, which gives (BGR9).
Finally $E^\sharp\mathcal K T$ equals the displayed expression for $M$;
both its terms are symmetric.  Positivity of the full $\mathcal K$ was
not used.  Invertibility was needed on the actual constrained slice.
\end{proof}

The inverse $L$ acts on fine forces; $T$ lifts the constraint datum after
the energetic correction; $M$ is the multiplier response.  Bounding $L$
alone leaves the products $L\mathcal K E$ and
$E^\sharp\mathcal K T$ to be controlled.  The size and locality of $E$
also matter.  A right inverse with a norm growing in the number of layers
would not prove a uniform response theorem.

Before applying per-length row estimates to these formulas, convert
the datum and its dual together. Put $\ell=\operatorname{diag}(\xi_c)$. With the cell pairings unchanged,
\[
 \begin{gathered}
 q_{\rm len}=\ell^{-1}q_{\rm raw},\qquad
 \omega_{\rm len}=\ell\omega_{\rm raw},\qquad
 E_{\rm len}=E_{\rm raw}\ell,\\
 \mathsf b_{\rm len}(B_{\rm len})
       =\ell^{-1}\mathsf b_{\rm raw}(\ell B_{\rm len}),\qquad
 D\mathsf b_{\rm len}(B_{\rm len})\dot B_{\rm len}
       =\ell^{-1}D\mathsf b_{\rm raw}(\ell B_{\rm len})
                                      \ell\dot B_{\rm len}.
 \end{gathered}
\]
These conversions preserve both $qE=I$ and the scalar Lagrangian,
so they preserve $\mathcal K$. At the identity in unit-spacing one-step coordinates,
division of the raw coarse logarithm by $L$ changes its pivot
derivative with respect to a raw fine-bond angle from $L^{-3}$
to $L^{-4}$; the right inverse and multiplier acquire the reciprocal
factors above. For fine spacing $\eta$ and a physical field $A$,
the raw derivative instead contains $\eta L^{-3}$, and division
by the coarse length $L\eta$ again gives $L^{-4}$. The cell weights and amplitudes below are inserted
after this coordinate choice.

\subsubsection*{Cell weights and a kernel algebra}

We describe the norm in which the preceding formulas can be estimated
without a bare-mesh counting loss.  Let each component $i$ have a positive
cell weight $\mu_i$ and an amplitude $s_i>0$.  A vector is measured by
$\|v\|_s=\sup_i |v_i|/s_i$.  A kernel is written as
\[
 (Av)_i=\sum_j\mu_j A_{ij}v_j,
       \qquad \widehat A_{ij}=s_i^{-1}A_{ij}s_j.
 \tag{BGR10}
\]
For different input and output species their respective amplitudes and
cell weights are used.  The identity kernel is $\delta_{ij}/\mu_j$.
Thus even the identity is not represented by $\delta_{ij}$ when it is
used with a weighted summation convention.

For a metric $d$ common to the component locations define
\[\begin{aligned}
 \|A\|_m=\max\bigg\{&\sup_i\sum_j\mu_j
                 |\widehat A_{ij}|e^{m d(i,j)},\\
                 &\sup_j\sum_i\mu_i
                 |\widehat A_{ij}|e^{m d(i,j)}\bigg\}.
\end{aligned}\tag{BGR11}\]
Matrix norms in a colour or local finite-dimensional fibre are operator
norms.  The two sums use the correct weights for their two species.
The triangle inequality for $d$ and absolute convergence give
\[
 \|AB\|_m\le\|A\|_m\|B\|_m.
 \tag{BGR12}
\]
Indeed, insert the intermediate sum with weight $\mu_k$, use
$d(i,j)\le d(i,k)+d(k,j)$ and cancel the intermediate amplitudes.
For the row bound take the supremum of the inner $B$ row sum; for the
column bound use the corresponding $A$ column sum.  These are distinct
calculations and neither counts the number of fine bonds.

The first row sum bounds $A$ on the normalized supremum norm.  The two
sums together bound it on the normalized, cell-weighted $\ell^2$ space
by the Schur estimate.  To see the latter directly, apply weighted
Cauchy--Schwarz to each row with weights
$\mu_j|\widehat A_{ij}|$, then sum in $i$ and use the column bound.
The result is $\|Av\|_{2,s}\le\|A\|_0\|v\|_{2,s}$.

For the natural response components above a top of scale $\xi_y$, the
amplitudes are
\[
 s_H(y)=\xi_y^{-1},\qquad
 s_{\nabla H}(y)=\xi_y^{-2},\qquad
 s_J(y)=\xi_y^{-3}.
 \tag{BGR13}
\]
Equivalently the output is
$(\xi_yH,\xi_y^2\nabla H,\xi_y^3J)$.  The cell weights are fixed by
the actual fine, coarse and pillar pairings, rather than chosen to make
a desired estimate true.  Each constraint length factor belongs in the
matrix before (BGR10) is formed.  The transpose is determined by the two declared pairings.
The input and output amplitudes must be inserted before applying (BGR12)
to a product that changes species at a seam.

For a raw logarithm $a=g y$, a physical radius $\alpha$ becomes
$\alpha/g$ in the canonical Gaussian variable $y$.  A derivative with
respect to $y$ contributes $g$ times the derivative with respect to $a$.
Thus the second derivative of a scalar activity in Gaussian coordinates
is $g^2$ times its raw-logarithm Hessian.  This conversion occurs in
addition to the length and cell-weight conversions in (BGR10)--(BGR13).

\begin{corollary}[Uniform response after the primitive bounds]
Suppose the actual normalized kernels in (BGR7) satisfy
\[
 \|L\|_m\le a,\quad \|E\|_m\le e,\quad
 \|\mathcal K E\|_m\le k,
       \quad \|E^\sharp\mathcal K\|_m\le k'.
 \tag{BGR14}
\]
Then
\[
 \|T\|_m\le e+ak,\qquad
 \|M\|_m\le k'(e+ak).
 \tag{BGR15}
\]
In the nonsymmetric case the force-to-multiplier block has norm at most
$k'a+\|E^\sharp\|_m$.  In the symmetric case it is $-T^\sharp$.
Its norm equals the norm of $T$ with reversed species if the dual
amplitudes are chosen reciprocally for the same cell-weighted pairings.
For other force and current amplitudes its weighted transpose bound
must be checked separately.
\end{corollary}

The last qualification is necessary.  A raw transpose and a transpose
after conjugation by amplitude weights are different matrices.  If
$\widehat T_{ij}=s_i^{-1}T_{ij}s_j$, its transposed normalized kernel
is obtained by giving the dual components amplitudes $s_i^{-1}$ and
$s_j^{-1}$.  A natural force weight of a different physical dimension
does not satisfy this rule automatically.  Its conversion belongs in
the kernel before a Schur norm is claimed.

For the actual row map, \eqref{px:actual-layered-row-lift} and
\eqref{px:actual-row-mass-trace} give the lift and trace estimates.
The other bounds in (BGR14) come from the field and multiplier
blocks of the projected construction after
\eqref{px:feynman-coefficient-comparison}. Retain the row equation
and transfer derivatives to neighboring inverse or lift factors
before estimating the energetic products. The conversions after
(BGR9) give their datum and multiplier the same physical pairings.
The remaining identification consists of recording each block's
species weight and outward datum derivative. Then (BGR15) follows
by the displayed kernel products.

\subsubsection*{Invertibility on the whole complex ball}

On the fixed slice $\mathcal V$, let
\[
 \mathcal S(h,\nu)
      =(\mathcal K h-q^\sharp\nu,\,qh).
\]
Its domain consists of slice fields and row multipliers; its range
consists of the paired field forces and row data in (BGR6)--(BGR9).
Fix a reference operator $\mathcal S_0$ on these same spaces,
with the same pairings and datum convention.  If the primitive estimates give
\[
 \|\mathcal S_0^{-1}\|_m\le C_0,
 \qquad
 \|\mathcal S(B)-\mathcal S_0\|_m\le\varepsilon,
       \qquad C_0\varepsilon<1,
 \tag{BGR16}
\]
then the Neumann series converges in the kernel algebra and proves
\[
 \|\mathcal S(B)^{-1}\|_m
          \le {C_0\over1-C_0\varepsilon}.
 \tag{BGR17}
\]
It is essential that the second bound in (BGR16) holds on the whole
retained complex ball.  A bound only along the real solutions does not
give the analytic disc needed later for Cauchy estimates.

If $\mathcal S_0$ is unbounded in the chosen norm, the correct statement
uses the relative perturbation
$\mathcal S_0^{-1}(\mathcal S-\mathcal S_0)$ between its actual graph
spaces.  The same series applies there.  In particular an elliptic
second-order operator cannot be treated as a bounded supremum-norm
matrix with a mesh-independent norm merely because its inverse is
bounded.

If the coefficients of $\mathcal S(B)$ are holomorphic in the indicated
operator norm, holomorphy follows term by term on a strictly smaller
disc, or by the resolvent identity
\[
 \mathcal S(B')^{-1}-\mathcal S(B)^{-1}
   =-\mathcal S(B')^{-1}
              [\mathcal S(B')-\mathcal S(B)]\mathcal S(B)^{-1}.
 \tag{BGR18}
\]
The identity also proves continuity and the derivative formula wherever
the coefficient derivative exists in the indicated operator norm.
No complex positivity assertion is used.  Real coercivity may construct
$\mathcal S_0^{-1}$; the complex inverse then requires a genuine norm
comparison or another explicitly convergent construction.

For the actual response make this comparison on the common complex
ball for (BGR1), with its reference and slice fixed. Use the local
analytic coefficients of $K_A$, the full constraint-curvature and
prescribed gauge blocks, and the outward factors of
\eqref{ym4:exp500-bg:mb:seam-datum-nilpotence}.
Equations \eqref{px:complex-wilson-relative-form} and
\eqref{px:feynman-coefficient-comparison}, together with the
projected-current expansions and cutoff commutators above, give
the relative comparison for the assembled system. To identify it
with the actual $\mathcal S$, we must match the auxiliary equations,
including each seam's datum and multiplier normalization. These are
local coefficient checks of the same construction. After this
identification, the estimates verify (BGR16) in its graph-space
form on a smaller ball and supply the uniform continuation radius.

\subsubsection*{From local forcing to first response}

Suppose a change in one datum coordinate $B(c)$ gives a forcing column
$F_c$ in the enlarged system.  The column consists of both the field
force and the constraint force in (BGR6).  If it is supported within
distance $r_0$ of the top $y_c$ and has normalized supremum norm at most
$b_0$, then (BGR17) gives
\[
 |D_{B(c)}Z_y[e]|/s_y
   \le {C_0b_0 e^{mr_0}\over1-C_0\varepsilon}
                                  e^{-m d(y,y_c)}|e|.
 \tag{BGR19}
\]
Here $Z$ contains the output species actually controlled by the inverse,
with amplitude $s_y$ from (BGR10). Equivalently apply the bound to
the normalized readout $Z_y/s_y$; subsequent Cauchy estimates use
normalized readouts.
The proof is the row estimate: for each contributing component $j$,
$d(y,y_c)\le d(y,j)+r_0$, so multiply by $e^{m d(y,y_c)}$ and sum the
kernel against $F_c$.  The reduced-datum derivative in
\eqref{ym4:exp500-bg:mb:seam-datum-nilpotence} changes
$b_0,r_0$ by its actual block bounds.  It must be included in $F_c$.

For an exponentially local forcing column, a bound
$|F_c(j)|/s_j\le b_0e^{-a d(j,y_c)}$ gives (BGR19) with any
$0<m'<\min(m,a)$ and $m'$ in place of $m$.  Indeed the triangle
inequality leaves a factor $e^{-(a-m')d(j,y_c)}\le1$, and the remaining
sum is a weighted row sum at rate $m'$.  No uncontrolled number of
sites appears.

This proves response decay from three smaller tasks: construct the
actual primitive inverse, verify its weighted row norm, and compute the
actual datum forcing.  If the output inverse controls only $H$ and the multiplier, the estimate
has not yet been proved for $\nabla H$ or the current.  Those outputs
require their own readout maps.

For example the physical current is a specified nonlinear function
$J(U)$.  In a fixed logarithmic chart its derivative is
\[
 \dot J=D_A[J(e^{i\eta A}U_0)]\dot A.
 \tag{BGR20}
\]
One may express this derivative by the multiplier and
lower-order terms through the differentiated field equation.  Such a
replacement must be derived from (BGR1) and the definition of $J$.
  On an independent-current core the variable $\mathbf J$
and the computed current $J(\mathbf U)$ are distinct until the physical
pair is imposed.

\subsubsection*{Spatial differences of the response}

The differentiated equation supplies spatial regularity in addition
to the inverse-kernel bound.

There is a useful exact identity for that estimate.  On a translation
patch let $T$ shift by one mesh step, $\nabla=(T-I)/\eta$ and
$\mathcal K_T=T\mathcal K T^{-1}$.  If $\mathcal K v=f$, then
\[
 \mathcal K_T\nabla v
    =\nabla f-{\mathcal K_T-\mathcal K\over\eta}v.
 \tag{BGR21}
\]
To prove it, write $T(\mathcal K v)=\mathcal K_TTv$, subtract
$\mathcal K v$ and divide by $\eta$.  At a boundary this calculation
must be made on the common interior patch; it does not shift an
inadmissible boundary row into the equation.

For a scalar divergence-form example
$\mathcal K=D^\sharp aD+V$, with constant mesh difference $D$,
the coefficient commutator in (BGR21) is
\[
 {\mathcal K_T-\mathcal K\over\eta}
      =D^\sharp(\nabla a)D+\nabla V.
 \tag{BGR22}
\]
If the energy norm is
$\|v\|_{\mathcal E}^2=\|Dv\|_2^2+\|v\|_2^2$, the bilinear pairing
and Cauchy--Schwarz imply
\[
 \big\|[(\mathcal K_T-\mathcal K)/\eta]v\big\|_{\mathcal E^*}
 \le (\|\nabla a\|_\infty+\|\nabla V\|_\infty)
                                               \|v\|_{\mathcal E}.
 \tag{BGR23}
\]
Thus a bounded inverse from $\mathcal E^*$ to $\mathcal E$, together
with coefficient regularity, bounds $\nabla v$ in energy norm.  This
does not assert a supremum bound; passing to a supremum or H\"older
bound requires the corresponding local elliptic estimate.

In the gauge system, shifting the covariant differences adds transport
commutators, the average rows add contour terms, and a seam adds its
own boundary comparison.  The required finite list is: the translated
interior equation, the gauge-row commutator, the constraint-row
commutator, and each seam or exterior boundary row.  Each term must
be bounded between its actual force and output norms.  The uniform
primitive regularity estimate is precisely this list of bounds and
the associated local elliptic estimate.  It is not supplied by
exponential decay of $H^{-1}$ alone.

\subsubsection*{Joint Cauchy estimates with joining distances}

Once (BGR19) has been established uniformly on a complex product ball,
several useful higher response bounds require no further model
estimate.  We give the argument with all radius losses visible.

Let $\mathcal B_r$ be a ball in the supremum of the dimensionless
bond logarithms, with a fixed norm on the three colour coordinates.
Suppose a holomorphic map $Z_y$ on $\mathcal B_r$ satisfies
\[
 |\partial_c Z_y(B)[e]|
        \le C e^{-m d(y,y_c)}|e|,
                       \qquad B\in\mathcal B_r.
 \tag{BGR24}
\]
Fix $0<\theta<1$.  There is a colour constant $c_*$ such that the
coordinate polydisc of radius $c_*(1-\theta)r$ about every
$B\in\mathcal B_{\theta r}$ lies in $\mathcal B_r$.  The number of
bonds does not enter this assertion, since the norm is a supremum over
bonds; simultaneous variation of three colour components changes only
$c_*$.  Cauchy's formula applied to the first derivative gives
\[
 |\partial_{c_1}\cdots\partial_{c_q} Z_y(B)|
 \le C(q-1)!
       [c_*(1-\theta)r]^{-(q-1)}
                       e^{-m\max_j d(y,y_{c_j})}.
 \tag{BGR25}
\]
Repeated coordinates are allowed.  The relevant multi-index factorial
is at most $(q-1)!$.  Holomorphic derivatives commute, so each one of
the $q$ coordinates can be selected as the first derivative in
(BGR24); taking the minimum of the resulting exponential bounds proves
the maximum in (BGR25).

Let $\mathcal T(y,c_1,\ldots,c_q)$ be the minimum spanning-tree length
in the underlying metric.  The star through $y$ is a competing tree,
and therefore
\[
 \mathcal T\le\sum_{j=1}^q d(y,y_{c_j})
                \le q\max_jd(y,y_{c_j}).
 \tag{BGR26}
\]
Consequently (BGR25) gives a joining-tree exponential with rate $m/q$.
For the second derivative it also gives the more explicit bound
\[
 |\partial_c\partial_{c'}Z_y(B)|
 \le {C\over(1-\theta)r}
       e^{-\frac m2[d(y,y_c)+d(y,y_{c'})]}.
 \tag{BGR27}
\]
This is the estimate used in the pointed Taylor remainder.  It costs
one datum radius.  Formula (BGR25) is not a uniform-in-$q$ tree bound
at a fixed rate: its rate deteriorates with $q$.  If an arbitrary-order
polymer estimate needs a fixed rate, one must use the local nonlinear
equations and their tree expansion, with the resulting factorial and
entropy costs.  A fixed-order application may use (BGR25) directly.

The same Cauchy argument applies to an output H\"older seminorm if
the first response is already bounded in that seminorm on the larger
ball.  Cauchy's formula can be applied to every normalized difference
$[Z_y-Z_{y'}]/|y-y'|^\gamma$ and then a supremum taken.  It cannot
create spatial H\"older regularity from a pointwise bound.

\subsubsection*{Finite boxes and compatible boundary forcing}

Let $\mathcal S$ be the actual linearized saddle operator and
$\mathcal S_\Lambda$ its prescribed finite-box version.  To compare
solutions, their equations must be put on common variables and rows
in an interior collar.  Write the resulting exact difference equation
as
\[
 \mathcal S_\Lambda\,\delta X=f_\Lambda.
 \tag{BGR28}
\]
The force includes the difference of boundary rows, the datum
restriction, the gauge-restoration term and any cutoff commutator.
The identity (BGR28) is checked before its right-hand side is bounded.
Replacing an actual box boundary by a convenient Dirichlet boundary
changes this force and requires a new comparison.

Suppose the inverse has weighted row bound $C$ at rate $m$ and the
force is supported in a collar $\mathcal C$. Let $R$ be the physical size
of the box and $p\ge0$ a fixed exponent. If
$\|f_\Lambda\|_s\le A R^p$, then for a component at $y$
\[
 |\delta X_y|/s_y
       \le CA R^p e^{-m\operatorname{dist}(y,\mathcal C)}.
 \tag{BGR29}
\]
This follows by extracting the exponential distance from the weighted
row sum.  In particular no factor equal to the number of bare sites
in the collar is necessary.  Such a factor would generally destroy
uniformity as the mesh tends to zero.

The same conclusion, at any rate $m'<\min(m,a)$, holds if
\[
 |f_\Lambda(j)|/s_j
             \le A R^p e^{-a\operatorname{dist}(j,\mathcal C)}.
 \tag{BGR30}
\]
For each $j$, the distance from $y$ to the collar is at most
$d(y,j)+\operatorname{dist}(j,\mathcal C)$.  Multiplying the summand
by $e^{m'\operatorname{dist}(y,\mathcal C)}$ leaves a weighted row
sum and a remaining exponential at most one.  This proves the claim.

A cutoff must be made in physical coordinates.  For a smooth function
$\chi(x/R)$ on a mesh $\eta$, its first difference divided by $\eta$
is $O(R^{-1})$ and its second difference divided by $\eta^2$ is
$O(R^{-2})$, by the integral form of the mean-value formula.
The unnormalized difference itself has an extra factor $\eta$.
These are different quantities.  A cutoff of an affine potential of
size $O(R)$ can therefore have bounded first derivatives and
polynomially bounded forcing, rather than a bare-mesh singularity.

For a nonlinear equation a cutoff of a solution is not a solution even
where the cutoff is a constant strictly between zero and one.  All
such points belong to the transition collar.  On that collar the
force is computed from the actual nonlinear formula or its integral
Taylor formula.  The admissibility of the cutoff field, its averages,
and its tower backgrounds is a separate domain check.  A bounded
formal commutator cannot replace it.

For two nonlinear backgrounds, subtracting their equations gives an
averaged Jacobian along a common admissible coordinate segment.  The
inverse estimate for that averaged operator must be justified, for
example by a uniform comparison of every Jacobian with the same
reference as in (BGR16).  An arbitrary average of invertible complex
matrices need not be invertible.  The first-step background section
uses just such a common comparison; a multilevel application needs
it on the actual multilevel ball.

\subsubsection*{The full substitution and the test for staying inside a core}

An old-origin source is evaluated on data of the form
\[
 V_j(\nu)=\mathcal M^j
        \big(e^{i\eta(t+\nu\zeta)H}U\big)|_{\Gamma_j}.
 \tag{BGR31}
\]
The parameters on the fine bonds in $\zeta$ may vary independently.
The inner analytic argument uses $|\nu|<c/s$, not merely the physical
segment between zero and one.  There are also outer and removed
substitutions.  Each must satisfy the same list of field, independent
current, node, gauge and tower conditions that defines its assigned
core.  A proof only of the curvature clause is incomplete.

Here is a precise useful criterion.  Express the core as inequalities
$\|\mathcal R_a(X)\|<r_a$ for its actual clauses, including all
hereditary tower clauses.  Let $X_0$ be the complete reference pair and
assume margins $\|\mathcal R_a(X_0)\|\le(1-\epsilon_a)r_a$.
If a common analytic curve satisfies
\[
 \sup_{\nu\ {\rm on\ the\ retained\ disc}}
 \|\mathcal R_a(X_\nu)-\mathcal R_a(X_0)\|
                         <\epsilon_a r_a
                         \quad\hbox{for every }a,
 \tag{BGR32}
\]
then it stays in the core.  The proof is the triangle inequality for
every defining clause.  Its substantive content in an application is
the uniform bound for the infinite hereditary list, not that triangle
inequality.

If the actual response and average derivatives imply
\[
 {\|\mathcal R_a(X_\nu)-\mathcal R_a(X_0)\|\over r_a}
      \le C_a|\nu|\,\delta_k\sigma_j/\alpha_j,
 \tag{BGR33}
\]
then the inner disc is admitted whenever
$c\sup_a\{C_a\delta_k\sigma_j/(\epsilon_a s\alpha_j)\}<1$.
For the window maps, Section~\ref{sec:substitution-room} gives the
normalized displacement in (BGR33). Its average derivative controls
the datum clauses, the reduced-current response controls the current
clause, and the real-to-complex curvature calculation controls the
entire tower list. For a layer gap $\epsilon_a$ decreasing as
$2^{-n}$, use the quotient $C_a/\epsilon_a$ rather than a
positive lower bound on all the gaps. The response estimate supplies
the factor $e^{-c\delta_0ML(n-1)}$; thus the relevant ratio is
bounded by
\[
 C2^n\sqrt{1+n}e^{-c\delta_0ML(n-1)},
\]
which is uniformly bounded after $M$ is chosen. With
$s=\delta_k/\alpha_k$, the other quotient in the admission condition is
$\sigma_j\alpha_k/\alpha_j$. The slow-radius inequality gives
$\sigma_j\alpha_k/\alpha_j\le C\sqrt{1+k-j}L^{-(k-j)}$,
whose supremum is finite. Choose the common amplification constant
$c$ from these two suprema and the retained margin. Together with
the common continuation bound below, these strict inequalities admit
the disc simultaneously for all old levels and tower clauses. The powers of $\sigma_j$ here come from the physical
length conversion in the nonlinear average.

To use the displacement estimate for admission, continue the common
analytic solution branch up to a possible first exit. Before exit,
the local inverse construction following (BGR18) applies on the
same scaled ball and gives a uniform continuation radius. Equations
(BGR32)--(BGR33), with the tower estimate of
Section~\ref{sec:substitution-room}, keep every clause strictly
inside its allotted margin. The branch therefore continues past
any proposed first exit. The finite address check is that the
outer, inner and removed substitutions use their assigned tower,
current and node maps; the normalized margin estimate already
handles their full hereditary lists.

The unamplified reference value must also belong to the core.
At the finitely many young ages where the critical half-datum
chart is replaced by the whole-core estimate, use that estimate
for the complete reference value before applying the displacement
bound. The amplification argument then spends its remaining
margin. This separates admission of $X_0$ from admission of
its increment.

For removed substitutions the current and the background may change
together.  The path used in the comparison must then remain in the
joint core.  Separate coordinate Cauchy estimates at the two endpoints
do not imply that the segment connecting them is admitted, especially
when a hereditary condition is not convex.  One can instead use the
actual common analytic chart and the strict clause margins in (BGR32).



\subsection*{The mesh-uniform H\"older estimate used by the response}
\label{ym4:exp500-bg:holder}

The response system needs spatial regularity as well as decay.  We prove
the particular scalar estimate that converts a bounded potential and a
bounded componentwise Laplacian into a H\"older bound for its first
differences.  The estimate is uniform in the mesh.  It also applies in a
parallel frame for an unweighted flat unitary connection, including
arbitrary flat holonomies on a compatible torus.  We then apply the estimate in the local real gauges of the background
construction, with the actual covariant Laplacian as its source.

Fix $0<\alpha<1$.  On $\eta\mathbb Z^4$ define
\[
 D_i^+u(x)={u(x+\eta e_i)-u(x)\over\eta},\quad
 D_i^-u(x)={u(x)-u(x-\eta e_i)\over\eta},\quad
 \Delta_\eta=\sum_{i=1}^4D_i^-D_i^+.
 \tag{BGH1}
\]
The Laplacian has nonpositive Fourier symbol.  The counting measure in
all convolution kernels is $\eta^4$ times counting measure.  We use the
Euclidean distance, allowing dimension-dependent constants when a
coordinate path is used.  A derivative of order $r$ means an arbitrary
product of $r$ of the forward or backward differences in (BGH1).

\subsubsection*{Heat kernels with enough derivative and tail control}

Start on the unit lattice.  For $T\ge0$ set
\[
 q_T(n)={1\over(2\pi)^4}\int_{[-\pi,\pi]^4}
 e^{-T\omega(\theta)}e^{in\cdot\theta}\,d\theta,
 \qquad\omega(\theta)=4\sum_i\sin^2(\theta_i/2).
 \tag{BGH2}
\]
This is a nonnegative probability kernel.  Indeed, if $\mathsf A$ is
the sum of the eight nearest-neighbour translations, then
$\Delta_1=\mathsf A-8I$ and
\[
 e^{T\Delta_1}=e^{-8T}\sum_{k=0}^{\infty}{T^k\mathsf A^k\over k!}.
 \tag{BGH3}
\]
Every coefficient is nonnegative, and applying the operator to the
constant function gives one.  Fourier transformation gives (BGH2), so
the two constructions agree.  The physical heat-kernel density is
\[
 p_t^\eta(\eta n)=\eta^{-4}q_{t/\eta^2}(n).
 \tag{BGH4}
\]
It follows directly that convolution with $p_t^\eta$ is a contraction
on bounded functions.

We need polynomial tails rather than a sharp Gaussian bound.  The
following estimates suffice for the entire argument.

\begin{lemma}[Uniform differentiated heat-kernel tails]
For $\eta>0$, $a\ge0$ and integers $r,N\ge0$ there is $C_{r,N}$ independent of $\eta,t,a$
such that
\begin{align*}
 \sum_{|x|\ge a}\eta^4|D^r p_t^\eta(x)|
 &\le C_{r,N}t^{-r/2}(1+a/\sqrt t)^{-N},
                                      &&t\ge\eta^2,\\
 \sum_{|x|\ge a}\eta^4|D^r p_t^\eta(x)|
 &\le C_{r,N}\eta^{-r}(1+a/\eta)^{-N},
                                      &&0\le t\le\eta^2.
 \tag{BGH5}
\end{align*}
\end{lemma}

\begin{proof}
A product of $r$ unit differences has Fourier multiplier $m_r(\theta)$,
a product of factors $e^{\pm i\theta_i}-1$.  For every multi-index $b$,
\[
 |\partial^b m_r(\theta)|
           \le C_{r,b}|\theta|^{(r-|b|)_+}.
 \tag{BGH6}
\]
The bound follows by assigning the derivatives to factors.  A factor
that has not been differentiated vanishes to first order at zero;
all its positive-order derivatives are bounded.  If more than $r$
derivatives occur the asserted power is zero.

For $\theta\in[-\pi,\pi]^4$, $\omega(\theta)\ge c|\theta|^2$.
Its first derivatives are $O(|\theta|)$ and its higher derivatives
are bounded.  Differentiating $e^{-T\omega}$ for $T\ge1$ gives
\[
 |\partial^b e^{-T\omega(\theta)}|
 \le C_b T^{|b|/2}
           (1+|\sqrt T\theta|)^{M_b}e^{-c_bT|\theta|^2}.
 \tag{BGH7}
\]
For completeness, each term of the repeated product rule consists of
$e^{-T\omega}$ times factors $T\partial^{b_j}\omega$, with
$b_j\ne0$ and $\sum b_j=b$.  A first derivative contributes
$\sqrt T\,O(|\sqrt T\theta|)$.  A factor of order at least two
is at most a constant times $T^{|b_j|/2}$ because $T\ge1$.
Multiplying proves (BGH7); the finite polynomial can be absorbed by
a smaller Gaussian exponent when integrating.

Leibniz' rule, (BGH6), and (BGH7) now imply
\[
 \int_{[-\pi,\pi]^4}
  |\partial^b(m_r e^{-T\omega})|^2\,d\theta
                    \le C_{r,b}T^{|b|-r-2},\qquad T\ge1.
 \tag{BGH8}
\]
When at most $r$ derivatives hit the multiplier, scale
$\theta=T^{-1/2}v$ in that term.  When more than $r$ hit it, its
bound is smaller than the displayed power since $T\ge1$.  The
four-dimensional volume factor is $T^{-2}$.  This explains every
power in (BGH8).

Let $\delta^r$ denote the corresponding unit differences.  Parseval's
identity and differentiation of the periodic Fourier multiplier give,
for every integer $M\ge0$,
\[
 \sum_n (1+|n|^2/T)^M|\delta^r q_T(n)|^2
                    \le C_{r,M}T^{-r-2},\qquad T\ge1.
 \tag{BGH9}
\]
Expand the weight into coordinate monomials.  The term $|n^b|^2T^{-|b|}$
is bounded by (BGH8) times $T^{-|b|}$.  All multiplier functions are
smooth and periodic, so there are no boundary terms at $\pm\pi$.

Choose $M>N+2$, increasing it by one if necessary.  Comparing the
lattice sum with the radial integral in four dimensions gives
\[
 \sum_{|n|\ge A}(1+|n|^2/T)^{-M}
       \le C_{M,N}T^2(1+A/\sqrt T)^{-2N},\qquad T\ge1.
 \tag{BGH10}
\]
One way to check the comparison is to cover the lattice by unit cubes.
For $T\ge1$ the values of $1+|x|/\sqrt T$ at any two points of a cube
are comparable by a fixed dimensional factor.  The integral of
$(1+|x|^2/T)^{-M}$ outside a ball of radius comparable with $A$ is
$T^2$ times the corresponding integral in scaled coordinates.  That
integral decays at least as $(1+A/\sqrt T)^{-2N}$ for this choice of $M$.
Small values of $A$ are absorbed by the full integral bound.

Cauchy--Schwarz with (BGH9) and (BGH10) proves
$\sum_{|n|\ge A}|\delta^rq_T(n)|
\le C T^{-r/2}(1+A/\sqrt T)^{-N}$.
For $0\le T\le1$, all derivatives of
$m_re^{-T\omega}$ of any fixed order are uniformly bounded.
The same Parseval argument, with weight $(1+|n|^2)^M$, gives instead
$C(1+A)^{-N}$.  Finally (BGH4) contributes $\eta^{-r}$ to the
weighted $\ell^1$ norm.  Substitution of $T=t/\eta^2$ proves (BGH5).
\end{proof}

The constants depend only on the fixed derivative order and tail power. The small-time estimate is needed for the discrete second and third differences.

\subsubsection*{A massive resolvent and its first-difference modulus}

For $\mu>0$ define
\[
 G_\mu^\eta(x)=\int_0^\infty e^{-\mu^2t}p_t^\eta(x)\,dt,
       \qquad R_\mu f=G_\mu^\eta*f.
 \tag{BGH11}
\]
Its kernel is nonnegative and has total mass $\mu^{-2}$.  For fixed
$\eta$, $\Delta_\eta$ is a bounded operator on bounded functions.
Integrating the derivative of
$e^{-\mu^2t}e^{t\Delta_\eta}$ in operator norm shows
\[
 (-\Delta_\eta+\mu^2)R_\mu f=f,
       \qquad \|R_\mu f\|_\infty\le\mu^{-2}\|f\|_\infty.
 \tag{BGH12}
\]
Conversely, any bounded solution of this equation equals $R_\mu f$.
Indeed integrate the same identity applied to the bounded solution;
the terminal term tends to zero because of $e^{-\mu^2t}$ and the
heat contraction.  This also proves uniqueness without a maximum
attainment argument on an infinite lattice.

Assume $\mu\eta\le1$.  Integrating (BGH5) for $r=1$, first on
$[0,\eta^2]$ and then on $[\eta^2,\infty)$, gives
\[
 \sum_x\eta^4|DG_\mu^\eta(x)|\le C\mu^{-1}.
 \tag{BGH13}
\]
The first interval contributes $C\eta$ and the second is bounded
by $C\int_0^\infty e^{-\mu^2t}t^{-1/2}\,dt=C'\mu^{-1}$.

For a lattice displacement $h$, telescoping along a coordinate path
and using (BGH5) with $r=1,2$ gives, for $t\ge\eta^2$,
\[
 \|\tau_h Dp_t^\eta-Dp_t^\eta\|_1
       \le C\min(t^{-1/2},|h|t^{-1}).
 \tag{BGH14}
\]
The second bound follows by summing the second difference over the
path of physical length at most $2|h|$ in four dimensions; the first
is twice the first-difference norm.  For $t\le\eta^2$ use
$C\eta^{-1}$ instead.  Every nonzero lattice displacement has
$|h|\ge\eta$.

Suppose first that $\eta\le|h|\le\mu^{-1}$.  Integrate (BGH14) on
$[\eta^2,|h|^2]$, on $[|h|^2,\mu^{-2}]$ and on the remaining
tail.  Together with the small-time interval, this gives
\[
 \|\tau_h DG_\mu^\eta-DG_\mu^\eta\|_1
       \le C|h|\big(1+\log(1/(\mu|h|))\big).
 \tag{BGH15}
\]
On the middle interval the integral is
$|h|\int_{|h|^2}^{\mu^{-2}}t^{-1}dt$.
The first interval contributes $O(|h|)$, and the exponentially
damped last interval does as well.  If $|h|>\mu^{-1}$ use twice
(BGH13).  Since for $0<x\le1$
\[
 x^{1-\alpha}(1+\log(1/x))
        \le 1+{1\over e(1-\alpha)},
 \tag{BGH16}
\]
both ranges imply
\[
 \|\tau_h DG_\mu^\eta-DG_\mu^\eta\|_1
          \le C_\alpha\mu^{\alpha-1}|h|^\alpha.
 \tag{BGH17}
\]
The value at $h=0$ is zero.  It follows for bounded $f$ that
\[
 \|DR_\mu f\|_\infty\le C\mu^{-1}\|f\|_\infty,
 \qquad [DR_\mu f]_\alpha
                \le C_\alpha\mu^{\alpha-1}\|f\|_\infty.
 \tag{BGH18}
\]
The seminorm compares lattice points, using their physical distance.
The logarithmic modulus (BGH15) explains why these estimates do not
give a mesh-uniform Lipschitz bound for $DR_\mu f$ from bounded $f$
alone.  The constants are permitted to diverge as $\alpha$ tends to one.

\subsubsection*{Off-diagonal resolvent derivatives}

Let $R\ge\eta$ and $\mu=R^{-1}$.  For every fixed $c>0$ and
$r=0,1,2,3$,
\[
 \sum_{|x|\ge cR}\eta^4|D^rG_{1/R}^\eta(x)|
                           \le C_{r,c}R^{2-r}.
 \tag{BGH19}
\]
Choose $N>\max(0,r-2)$.  On the time interval $[\eta^2,R^2]$,
the large-time part of (BGH5) is bounded by
$Ct^{-r/2}(\sqrt t/R)^N$.  Its integral is $O(R^{2-r})$.
On $[R^2,\infty)$ the exponential $e^{-t/R^2}$ gives the same
bound by scaling.  The interval $[0,\eta^2]$ contributes at most
$C\eta^{2-r}(\eta/R)^N$, which is also $O(R^{2-r})$.
This proves (BGH19), including the second and third derivatives for
which the uncut global norms would have a worse small-scale behaviour.

There is no assertion here that the second resolvent derivative has
a uniformly bounded global absolute sum.  Its off-diagonal bound is
what makes the interior cutoff argument work.

\subsubsection*{An interior estimate with no unknown cutoff gradient}

Let $Q_R$ be the cube of side $2R$ centred at a lattice point.  The
following formulation uses a fixed ratio of cubes; changing that ratio
changes only the constants.

\begin{theorem}[Interior first-difference H\"older bound]
Let $R\ge64\eta$.  Suppose $u$ is defined on $Q_{2R}$ and
$\Delta_\eta u=f$ at every point needed by the equation there.
Put $M=\|u\|_{\infty,Q_{2R}}$ and
$F=\|f\|_{\infty,Q_{2R}}$.  On $Q_{R/2}$, with all first differences
and comparisons taken inside its one-mesh enlargement,
\begin{align*}
 \|Du\|_\infty&\le C(R^{-1}M+RF),\\
 [Du]_\alpha&\le C_\alpha
                    (R^{-1-\alpha}M+R^{1-\alpha}F).
 \tag{BGH20}
\end{align*}
The constants are independent of $R,\eta,u$.
\end{theorem}

\begin{proof}
Choose the restriction of a smooth cutoff $\chi(x/R)$, equal to one
on $Q_R$ and supported in $Q_{3R/2}$.  It satisfies
$|D\chi|\le C/R$ and $|\Delta_\eta\chi|\le C/R^2$.
Its difference supports lie in a collar at distance at least $cR$
from the evaluation cube and its one-mesh enlargement.  The condition
$R\ge64\eta$ makes this separation uniform after every difference
used below.  Extend $u$ arbitrarily outside $Q_{2R}$; the value of
$v=\chi u$ and all forces below are independent of that extension.

With $\mu=1/R$ the exact product identity is
\[
 (-\Delta_\eta+\mu^2)v
      =a-[\Delta_\eta,\chi]u,
       \qquad a=\chi(-f+\mu^2u).
 \tag{BGH21}
\]
Both sides are extended by their actual compactly supported formulas.
The equation for $u$ is used only where $\chi$ or its differences are
nonzero, inside the given domain.  By (BGH12),
$v=R_\mu a-R_\mu[\Delta_\eta,\chi]u$.
The first term has force norm at most $F+R^{-2}M$, so (BGH18)
gives exactly its contributions to (BGH20).

To treat the commutator one must avoid estimating the unknown
$Du$ in the cutoff collar.  Move the commutator onto the kernel.
For scalar multiplication by $\chi$ and the constant cell-weighted
pairing,
\[
 [\Delta_\eta,\chi]^{\mathsf T}
       =[\chi,\Delta_\eta]=-[\Delta_\eta,\chi].
 \tag{BGH22}
\]
For $j=1,2$, apply this identity to
$K_j(y)=D_x^jG_\mu^\eta(x-y)$ in the convolution.  The discrete
product rule shows that its absolute value is bounded by
\[
 C M\left(R^{-2}\|1_{|\cdot|\ge cR}D^jG_\mu^\eta\|_1
       +R^{-1}\|1_{|\cdot|\ge cR}D^{j+1}G_\mu^\eta\|_1\right).
 \tag{BGH23}
\]
To check the rule explicitly,
$\Delta(\chi K)-\chi\Delta K-K\Delta\chi$
is the sum over $i$ of
$(D_i^+\chi)(D_i^+K)+(D_i^-\chi)(D_i^-K)$.
All factors containing a difference of $\chi$ are supported in the
separated collar.  Changes of the evaluation point by at most the
fixed number of mesh steps in a derivative merely change $c$.
The kernel is absolutely summable on that collar by (BGH19), which
also justifies the summation by parts.

For $j=1$, (BGH19) bounds (BGH23) by $CM/R$.
For $j=2$ it bounds it by $CM/R^2$.  Let
$w=R_\mu[\Delta_\eta,\chi]u$.  A coordinate path between two points
of $Q_{R/2}$ stays in that cube; telescoping the first differences
of $w$ gives
$|Dw(x)-Dw(x')|\le CM R^{-2}|x-x'|$.
Its supremum bound is $CM/R$.  Since $|x-x'|\le C R$ in the cube,
this gives $[Dw]_\alpha\le C_\alpha M R^{-1-\alpha}$.
Adding the contribution of $R_\mu a$ proves both estimates.
On the evaluation cube $v=u$, including all the differences used,
so the estimates apply to $u$ itself.
\end{proof}



The result holds componentwise for vectors in any fixed
finite-dimensional normed space.  Equivalent coordinate norms
contribute a constant depending on that fixed dimension.  
\subsubsection*{Tori and flat unitary holonomies}

On a torus with periods $\ell_i$, the scalar heat kernel and massive
resolvent are periodizations of their infinite-lattice kernels.
Their absolute derivative sums are bounded by the corresponding
infinite sums.  For a set at periodic distance at least $cR$ from
the origin, every lift has Euclidean distance at least $cR$.
Thus (BGH19) holds for periodic distance as well.  If the cubes and
their collars in the interior theorem embed without wrapping, the
same cutoff proof gives (BGH20) with the same constants.  For example
$\min_i\ell_i\ge 8R$ suffices for the stated cubes.

A flat unitary connection can be made trivial on the lifted lattice.
The deck transformations act by its commuting unitary holonomies.
The kernel on the torus is the image sum with the corresponding
unitary factors.  On a common parallel frame in the embedded cube,
differences act on the lifted kernel.  Taking the norm and using
unitarity bounds the image sum by the scalar sum already estimated.
Hence all the preceding local bounds remain uniform in the
holonomies.  Their distance from the identity is irrelevant.

The argument requires flatness.  Curved parallel transports have
additional commutators.  Nor does it give the same estimate for an
arbitrary nonlocal or layered positive operator: the proof used the
componentwise constant-coefficient Laplacian in (BGH1).

\subsubsection*{Covariant differences and the response equation}

Here is how (BGH20) enters the background and response estimates.
Suppose a local real gauge has supplied a potential on a cube
enlarged by a fixed multiple of its local scale $\xi$, with
\[
 \|A\|_\infty\le C_0m\varepsilon\xi^{-1},\qquad
 \|\Delta_\eta A\|_\infty
                       \le C_1m\varepsilon\xi^{-3}.
 \tag{BGH24}
\]
Apply the theorem on cubes of radius a fixed fraction of $\xi$,
covering the interior.  It proves
\[
 [DA]_\alpha\le C_\alpha m\varepsilon\xi^{-2-\alpha},
       \qquad \|DA\|_\infty\le C m\varepsilon\xi^{-2}.
 \tag{BGH25}
\]
For pairs farther apart than a fixed fraction of $\xi$, the
supremum bound gives the same H\"older estimate.  Thus local cubes
give the full indicated interior seminorm.  This proves the
H\"older implication once (BGH24) is known; it does not construct the
gauge or prove the nonlinear Laplacian estimate in (BGH24).

There is a finite coarse-mesh case in this application.  If the
chosen fraction of $\xi$ is less than $64\eta$, the ratio
$\xi/\eta$ is bounded by a fixed constant.  Directly,
$|DA|\le2\|A\|_\infty/\eta$ and
$[DA]_\alpha\le4\|A\|_\infty\eta^{-1-\alpha}$.
They give (BGH25) with the same powers of $\xi$ and a larger
fixed constant.  The same fallback applies to the response
application below.  Thus no unmentioned lower bound on
$\xi/\eta$ is needed.

For the physical constrained-background response we use the differentiated
Euler--Lagrange equation to estimate the positive covariant
componentwise Laplacian
$\mathcal L_{U_0}=\sum_i\nabla_{U_0,i}^{\mathsf T}\nabla_{U_0,i}$.
Thus $\mathcal L_{\mathbf 1}=-\Delta_\eta$ with the heat-generator
sign in (BGH1).  Its conversion can also be made quantitative.
In a local gauge write $U_0(b)=e^{i\eta A_0(b)}$
and let $T_b=\operatorname{Ad}(U_0(b))$ in the chosen unitary colour
representation.  If $\eta|A_0|$ stays in a fixed small exponential
chart, then
\[
 \|(T_b-I)/\eta\|\le C|A_0(b)|.
 \tag{BGH26}
\]
Differencing $T_b$ in a transverse direction uses the integral
derivative of the exponential.  It gives
$|D T_b|\le C\eta|D A_0|$, with the adjacent bond values interpreted
in the same local frame.  The inverse transport satisfies the same
bound.  Expanding the forward and backward covariant differences
therefore gives
\[
 |(\mathcal L_{U_0}+\Delta_\eta)h|
 \le C\big(|A_0|\,|Dh|
                     +( |DA_0|+|A_0|^2)|h|\big).
 \tag{BGH27}
\]
One can verify this by first writing each covariant first difference
as the ordinary difference plus $(T_b-I)/\eta$ times a neighbouring
value.  In the second difference the apparent $\eta^{-1}$ terms
combine into a first difference of $h$, a difference of the
coefficient $(T_b-I)/\eta$, and a product of two such coefficients.
These are exactly the three terms on the right of (BGH27); the factors there denote maxima on the fixed neighboring stencil.
There is no estimate of an individual uncombined $\eta^{-1}A_0h$
term.

Assume on a fixed enlargement of the local-scale evaluation cube that
\begin{align*}
 \xi|A_0|+\xi^2|DA_0|&\le a_0,\\
 \xi|h|+\xi^2|\nabla_{U_0}h|
                  +\xi^3|\mathcal L_{U_0}h|&\le B.
 \tag{BGH28}
\end{align*}
Formula (BGH26) first gives
$\xi^2|Dh|\le C(1+a_0)B$.
Then (BGH27) implies
\[
 \xi^3|\Delta_\eta h|
                   \le C(1+a_0+a_0^2)B.
 \tag{BGH29}
\]
Applying (BGH20) at scale $\xi$ yields
$[Dh]_\alpha\le C_{\alpha,a_0}B\xi^{-2-\alpha}$.
The difference between $Dh$ and $\nabla_{U_0}h$ is a bounded
coefficient times a neighbouring $h$.  Its H\"older seminorm is
controlled by the same local bounds, using the Lipschitz coefficient
bound from $DA_0$ and the first-difference bound for $h$.
Thus
\[
 [\nabla_{U_0}h]_\alpha
                         \le C_{\alpha,a_0}B\xi^{-2-\alpha}.
 \tag{BGH30}
\]

This initially compares gradient components in the chosen local
gauge.  If the assigned H\"older seminorm transports one component
to the other along paths of length at most $C|x-x'|$, unitary
product telescoping and (BGH28) give
$\|\operatorname{Ad}U_0(\Gamma)-I\|
\le C a_0|x-x'|/\xi$.
Multiplying by $|\nabla_{U_0}h|\le B/\xi^2$ adds at most
$C a_0B\xi^{-2-\alpha}|x-x'|^\alpha$ on the scale-$\xi$
cube.  Hence (BGH30) also holds in that transported seminorm.
Its premise still includes the pointwise covariant-gradient
bound in (BGH28); the argument removes an additional H\"older
premise, rather than that gradient bound.

If the right-hand side of (BGH28) has the response factor
$B=C e^{-m d(y,y_c)}|e|$, that factor persists on a fixed-size
local-scale cube up to the fixed exponential comparison of distances
within the cube.  This gives the corresponding H\"older member of
the response estimate. For the physical background response at
$s\equiv1$, we apply this implication after obtaining the local real gauge
and the three readouts in (BGH28) from the two physical response systems
above and their identification in mechanism (k). At a seam we must verify these readouts and compare
the local frame and equation on their common two-layer determining
neighborhood.

For arbitrary decoupled data $s$, a bound on the componentwise
Laplacian is not among the declared operator norms. We do not use
(BGH28) for that response. Its gradient H\"older estimate instead
uses the H\"older member already retained in the norms of
$\mathfrak G(s)$, $\widetilde G(s)$ and the datum lift: apply
those operator bounds in the same contraction for the potential
and its first differences. The force and reduced-current terms
in that contraction involve only the declared first-order norms.
This use of the physical Laplacian estimate therefore imposes no
stronger second-order bound on the decoupled operator data.

% END MODULE chapters/04_propagators.tex


% BEGIN MODULE chapters/pointed_hessian_cancellation_v6.tex
\section{The Hessian of a pointed member and its own action counterterm}
\label{ch:pointed-hessian-cancellation}

The object in this section is the actual pointed total composed with
the actual free response. It need not be a convolution. Its bulk sum
will be a convolution, and the distinction is retained throughout.
The result identifies its marginal quadratic term and its own action
counterterm. The nonlinear remainder uses the localized response and
field-domain hypotheses of Section~\ref{ch:pointed-window-remainder}.

\subsection*{The pointed object and the analytic hypotheses}

Work in coarse lattice units on $\mathbb Z^4$. An independently
oriented bond $(x,\mu)$ starts at $x$ and ends at $x+e_\mu$.
A real field $B_\mu(x)$ is a traceless Hermitian matrix. Choose a
color basis $t_a$, $1\le a\le3$, orthonormal for
$\langle A,C\rangle=\frac12\operatorname{Tr}(AC)$. Thus
$B=\sum_a B^a t_a$. Complexification is bilinear, not sesquilinear.
The normalization gives the quadratic Wilson action
\[
 A(e^{itB})=\frac{t^2}{2}\sum_{p}|dB(p)|^2+O(t^3)
 \tag{PH1}\label{ph:action-normalization}
\]
on every finite set of plaquettes. Here $d$ is the oriented lattice
curl and each coordinate-plane plaquette is counted once.

An individual member $G$ has the pointed pair total, in its
exponential chart,
\[
 \mathcal G_z(e^{i\eta A},J)=\sum_{X\ni z}
                  G(X;(e^{i\eta A},J),z),\qquad z\in\mathbb Z^4.
 \tag{PH2}\label{ph:pointed-total}
\]
The bare spacing $\eta$ can differ from the coarse unit spacing.
There are no stored spectators in a regular member. Assume the
following analytic and geometric properties of this actual family.

The sum in \eqref{ph:pointed-total} and its first two derivatives
converge normally in the neighborhoods where they are used. Every
summand is local in its pair variables and invariant under all local
gauge transformations for which the arguments remain in its declared
domain. In particular it is invariant under constant adjoint rotations.
Let $\mathcal P(B)$ be the actual background pair
formed from a sufficiently small finitely supported coarse field,
with $\mathcal P(0)=(1,0)$. It is twice differentiable, is gauge
covariant, and has first response, expressed in this potential chart,
\[
 hB=D\mathcal P(0)B=(HB,D_1^*D_1HB).
 \tag{PH3}\label{ph:response}
\]
The precise linear map $H$ is part of the background construction;
it is not replaced by a convolution chosen for this proof. Put
$F_z(B)=\mathcal G_z(\mathcal P(B))$. Assume the pointed totals
are covariant under unit translations and the full hypercubic group,
including reflections. These are properties of one individual member
at all its base points. They do not require two different members
chosen at different sites to be related.

Finally assume an absolute exponential Hessian moment bound. With
$\delta_{x,\mu}^a$ the field of color $a$ at one coarse bond, let
\[
 \mathcal K_z^{\mu\nu;ab}(x,y)
  =\sum_{X\ni z}D^2G(X,(1,0),z)
                      [h\delta_{x,\mu}^a,h\delta_{y,\nu}^b].
 \tag{PH4}\label{ph:actual-kernel}
\]
All derivatives of a pair function are taken in the $(A,J)$ chart
of \eqref{ph:pointed-total}, including its factor $\eta$.
For some $m>0$ assume
\[
 \sum_{\mu,\nu,a,b}\sum_{x,y\in\mathbb Z^4}
 e^{m(|x-z|+|y-z|)}|\mathcal K_z^{\mu\nu;ab}(x,y)|\le C_m.
 \tag{PH5}\label{ph:exponential-moment}
\]
The corresponding derivative representation on finitely supported
fields is assumed, as is usual when the normally convergent
localized Hessians are composed with the free response. This bound
and the actual background/core hypotheses are model inputs here.
In particular, this section does not infer a uniformly bounded
coefficient extractor from a potential norm on a shrinking field ball.

We first justify that \eqref{ph:actual-kernel} really is the Hessian
of $F_z$. Constant adjoint invariance makes the first derivative of
each summand of \eqref{ph:pointed-total} zero at $(1,0)$: each
coefficient of a potential or current coordinate would be an
invariant linear functional on the adjoint color space, which has
no fixed vector. This applies to its finitely many local coordinates
and hence to every direction supplied by $h$. The normal convergence
permits summation. In the second-order chain rule the term containing
$D\mathcal G_z(1,0)D^2\mathcal P(0)$ is therefore zero.
Consequently \eqref{ph:actual-kernel}, with the displayed actual
$h$, is exactly $D^2F_z(0)$.

The exponential moment hypothesis has a concrete sufficient
localized form. Suppose the analytic core contains a usable pair
norm ball of radius $a_{\rm core}$, and on it
$|G(X,\cdot,z)|\le g e^{-\kappa d(X)}$. Suppose also
\[
 \|h\delta_{x,\mu}^a\|_{1;X}
       \le C e^{-\delta\operatorname{dist}(x,X)},\qquad
 \#\{X\ni z:q\le d(X)<q+1\}\le C e^{\kappa_D q},
 \qquad \kappa>\kappa_D,
\]
and $\operatorname{diam}(X)\le C_M(d(X)+1)$, with fixed
localization cube size. These are the core, response and entropy
conditions used in quadratic extraction. Two-variable Cauchy gives
\[
 |D^2G(X,0,z)[h\delta_x,h\delta_y]|
 \le Cg a_{\rm core}^{-2} e^{-\kappa d(X)}
 e^{-\delta(\operatorname{dist}(x,X)+
                 \operatorname{dist}(y,X))}.
\]
Choose $m<\delta/2$ so small that
$2mC_M<(\kappa-\kappa_D)/2$. Since $z\in X$, inserting the
weight in \eqref{ph:exponential-moment} costs at most
$e^{2mC_M(d(X)+1)}$ after retaining the two response distances.
The remaining sums over $x$ and $y$ cost at most
$C(d(X)+1)^8$: enclose $X$ in a box of diameter
$C_M(d(X)+1)$ and sum the exponentially decreasing distances
outside it. The rooted domain count now gives
\[
 C_m\le Cg a_{\rm core}^{-2}
 \sum_{q\ge0}(q+1)^8
 e^{-(\kappa-\kappa_D-2mC_M)q}
 \le C' g a_{\rm core}^{-2}.
 \tag{PH5a}\label{ph:moment-from-core}
\]
Thus \eqref{ph:exponential-moment} need not be introduced as a
new physical estimate when these localized hypotheses have already
been supplied. The factor $a_{\rm core}^{-2}$ is retained. The
coordinate normalization and near-zero core inclusions needed to
use this Cauchy estimate are treated in
Section~\ref{ch:extraction-radius-normalization}.

\subsection*{Color, symmetry, and the local gauge Ward identity}

Constant adjoint invariance implies
\[
 \mathcal K_z^{\mu\nu;ab}(x,y)=\delta_{ab}K_z^{\mu\nu}(x,y),
 \qquad
 K_z^{\mu\nu}(x,y)=\frac13\sum_{a=1}^3
                         \mathcal K_z^{\mu\nu;aa}(x,y).
 \tag{PH6}\label{ph:color}
\]
For a direct proof, fix the spatial indices and view the color
coefficients as a $3\times3$ complex matrix. Adjoint rotations
through $\pi$ about the color axes force its off-diagonal entries
to zero, and rotations interchanging the axes make its diagonal
entries equal. The adjoint action of $\mathrm{SU}(2)$ contains
these rotations. This also proves the assertion for complex
members. Equality of mixed derivatives gives
\[
 K_z^{\mu\nu}(x,y)=K_z^{\nu\mu}(y,x).
 \tag{PH7}\label{ph:Hessian-symmetry}
\]
For a real member all these coefficients are real. Thus its bulk
kernel is self-adjoint in addition to being bilinearly symmetric.
No positivity of a Hessian is assumed.

For a finitely supported Lie-algebra-valued site function $\chi$,
write $(d\chi)_\mu(x)=\chi(x+e_\mu)-\chi(x)$. Gauge covariance
of $\mathcal P$ and gauge invariance of the member give
$F_z(B^g)=F_z(B)$, where $B^g$ is the logarithmic coarse bond
field after a small gauge transformation. Choose the gauge
convention $U(x,\mu)^g=g(x)^{-1}U(x,\mu)g(x+e_\mu)$, so
the derivative at the identity of $g=e^{it\chi}$ is $d\chi$.
Differentiate this identity once in $t$ and once in a finitely
supported field direction $C$. The mixed derivative of the gauge
coordinate itself is multiplied by $DF_z(0)=0$. Therefore
\[
 D^2F_z(0)[d\chi,C]=0.
 \tag{PH8}\label{ph:compact-Ward}
\]
The other sign convention for gauge transformations gives the same
zero identity. The argument initially uses real directions and
extends complex bilinearly. It uses neither a noncompact gauge
transformation nor a polynomial field in a nonlinear field domain.

\subsection*{Polynomial extensions justified by absolute moments}

For fields of polynomial growth define the extended Hessian by
\[
 \mathcal H_z(B,C)=
 \sum_{\mu,\nu,a}\sum_{x,y}
          K_z^{\mu\nu}(x,y)B_\mu^a(x)C_\nu^a(y).
 \tag{PH9}\label{ph:polynomial-Hessian}
\]
Every such sum converges absolutely by
\eqref{ph:exponential-moment}. Compact cutoffs converge to it by
dominated convergence, independently of the cutoffs. This defines
an extension of the quadratic jet, not a value of the nonlinear
member at an unbounded polynomial field.

The Ward identity extends to polynomial gauge parameters and
polynomial second arguments. To prove it, let $\chi$ be a
polynomial of degree $q$ on sites and let
$\chi_R(x)=\chi(x)\vartheta(x/R)$ with a fixed smooth compact
cutoff equal to one near zero. Its forward differences converge
pointwise to $d\chi$, and satisfy a bound
\[
 |d\chi_R(x)|\le C_\chi(1+|x|)^{\max(q-1,0)}
                                      \quad(R\ge1).
\]
Indeed away from the cutoff annulus the usual polynomial difference
has this bound; on the annulus the extra difference of the cutoff
costs $R^{-1}$ and $|x|$ is comparable to $R$. Apply
\eqref{ph:compact-Ward} with a single-bond second argument and
pass to the limit using \eqref{ph:exponential-moment}. Summing
against any polynomial second argument is then legitimate by
absolute convergence. The radius of the nonlinear gauge chart
may depend on $R$: only its derivative at zero was used for each
fixed $R$.

In particular the linear polynomial $\chi(x)=x_\mu t_a$ gives
the separate row and column zeroth moments
\[
 \sum_x K_z^{\mu\nu}(x,y)=0,\qquad
 \sum_y K_z^{\mu\nu}(x,y)=0.
 \tag{PH10}\label{ph:row-column-zero}
\]
The second identity follows also from
\eqref{ph:Hessian-symmetry}. Both hold for each fixed remaining
edge and pair of directions, not only after a full trace.

For the moment at $z=0$ set, initially in lower bond-start
coordinates,
\[
 T_{\mu\nu,\kappa\lambda}
       =\sum_{x,y}K_0^{\mu\nu}(x,y)x_\kappa y_\lambda.
 \tag{PH11}\label{ph:T}
\]
For a symmetric matrix $S$, the quadratic gauge polynomial
$\chi(x)=\frac12\sum_{\mu,\kappa}S_{\mu\kappa}
x_\mu x_\kappa t_a$ has forward difference
\[
 (d\chi)_\mu(x)=
       \sum_\kappa S_{\mu\kappa}x_\kappa+\tfrac12S_{\mu\mu}.
\]
Use the polynomial Ward identity with second argument
$C_\nu(y)=y_\lambda t_a$ at one direction $\nu$. The constant
term vanishes by \eqref{ph:row-column-zero}. Since $S$ is
arbitrary symmetric, one obtains
\[
 T_{\mu\nu,\kappa\lambda}
       =-T_{\kappa\nu,\mu\lambda},\qquad
 T_{\mu\nu,\kappa\lambda}
       =-T_{\mu\lambda,\kappa\nu}.
 \tag{PH12}\label{ph:T-Ward}
\]
The second identity follows by interchanging the two Hessian
arguments and repeating the first. These are exactly the two
antisymmetries which remove symmetric gradients.

\subsection*{The bulk sum and its second moment}

Unit translation covariance of this individual member gives
$K_z^{\mu\nu}(x,y)=K_0^{\mu\nu}(x-z,y-z)$. The bulk quadratic
kernel is
\[
 \Pi^{\mu\nu}(x,y)=\sum_zK_z^{\mu\nu}(x,y)
                         =\pi_{\mu\nu}(y-x),\qquad
 \pi_{\mu\nu}(u)=\sum_xK_0^{\mu\nu}(x,x+u).
 \tag{PH13}\label{ph:bulk}
\]
All sums converge absolutely. In fact, for $0<m'\le m$,
\[
 \sum_u e^{m'|u|}|\pi_{\mu\nu}(u)|
 \le\sum_{x,y}e^{m'(|x|+|y|)}|K_0^{\mu\nu}(x,y)|<\infty.
 \tag{PH14}\label{ph:bulk-moment-bound}
\]
For fixed $x,y$, translation of $z$ in
\eqref{ph:bulk} is a sub-sum of the same absolute double sum.
This proves both the bulk formula and every rearrangement below.
Only this bulk object is a convolution. The original $K_0(x,y)$
continues to have two independent spatial arguments.

The zeroth and first bulk moments vanish by
\eqref{ph:row-column-zero}. Its second moment is
\[
 \begin{split}
 P_{\mu\nu,\kappa\lambda}
  &:=\sum_u\pi_{\mu\nu}(u)u_\kappa u_\lambda\\
  &=\sum_{x,y}K_0^{\mu\nu}(x,y)
                    (y_\kappa-x_\kappa)(y_\lambda-x_\lambda)\\
  &=-T_{\mu\nu,\kappa\lambda}
                         -T_{\mu\nu,\lambda\kappa}.
 \end{split}
 \tag{PH15}\label{ph:P-versus-T}
\]
The pure $y_\kappa y_\lambda$ term is zero after summing $x$
by its row identity, and the pure $x_\kappa x_\lambda$ term is
zero after summing $y$. This is the reason for both minus signs;
no pointwise convolution identity for $K_0$ has been used.

The Fourier series
$\widehat\pi_{\mu\nu}(k)=\sum_u\pi_{\mu\nu}(u)e^{ik\cdot u}$
is holomorphic for sufficiently small $|\operatorname{Im}k|$
by \eqref{ph:bulk-moment-bound}. In particular
$\partial_\kappa\partial_\lambda\widehat\pi_{\mu\nu}(0)
=-P_{\mu\nu,\kappa\lambda}$. This optional Fourier description
concerns the bulk sum alone.

\subsection*{Bond midpoints and the hypercubic tensor}

Under reflections, lower bond starts undergo a direction-dependent
integer displacement. To apply tensor symmetry without dropping
that displacement, move first to bond-midpoint coordinates
$\widetilde x=x+e_\mu/2$ and
$\widetilde y=y+e_\nu/2$. The value of $T$ does not change:
expanding $\widetilde x_\kappa\widetilde y_\lambda$ adds only
terms of the forms $\sum_{x,y}K_0y_\lambda$,
$\sum_{x,y}K_0x_\kappa$ and $\sum_{x,y}K_0$, all zero by
\eqref{ph:row-column-zero}. The displacement $u$ changes by the
constant $(e_\nu-e_\mu)/2$. Because the bulk zeroth and first
moments are zero, its second moment $P$ is also unchanged.
Thus the lower-start extractor and the midpoint extractor agree
exactly.

The midpoint bond sets are
$\mathbb Z^4+e_\mu/2$. A signed permutation $w$ with
$we_\mu=s_\mu e_{\sigma(\mu)}$ maps this set to the midpoint
set of direction $\sigma(\mu)$; a negative orientation contributes
the sign $s_\mu$ to a one-form. Pointed covariance at the fixed
base point zero and the ordinary transformation of the midpoint
coordinates therefore give
\[
 T_{\sigma(\mu)\sigma(\nu),\sigma(\kappa)\sigma(\lambda)}
       =s_\mu s_\nu s_\kappa s_\lambda
                                  T_{\mu\nu,\kappa\lambda}.
 \tag{PH16}\label{ph:W4}
\]
There is no affine correction left in this formula; the preceding
moment calculation proved that it vanished.

We now classify this tensor directly. By
\eqref{ph:T-Ward} it is zero if $\mu=\kappa$ or
$\nu=\lambda$. Reflecting one coordinate in
\eqref{ph:W4} shows that each coordinate label must occur an
even number of times among the four indices of a nonzero entry.
The two antisymmetric pairs must consequently be the same
unordered pair. The only possibilities are
$T_{\mu\mu,\kappa\kappa}$ and
$T_{\mu\kappa,\kappa\mu}$ for $\mu\ne\kappa$.
Permutations make all the first values equal to a number $b$;
\eqref{ph:T-Ward} makes the second values $-b$. Hence
\[
 T_{\mu\nu,\kappa\lambda}
      =b\bigl(\delta_{\mu\nu}\delta_{\kappa\lambda}
                     -\delta_{\mu\lambda}\delta_{\kappa\nu}\bigr).
 \tag{PH17}\label{ph:Maxwell-T}
\]
Full reflections are used here: four distinct indices, including a
possible alternating tensor, are excluded by a single reflection.
This is not merely an argument using proper rotations.

Combining \eqref{ph:P-versus-T} with
\eqref{ph:Maxwell-T} gives
\[
 P_{\mu\nu,\kappa\lambda}
    =b\bigl(\delta_{\mu\lambda}\delta_{\kappa\nu}
            +\delta_{\mu\kappa}\delta_{\lambda\nu}
            -2\delta_{\mu\nu}\delta_{\kappa\lambda}\bigr),
 \qquad b=P_{\mu\nu,\mu\nu}\quad(\mu\ne\nu).
 \tag{PH18}\label{ph:extractor}
\]
Thus the extractor is explicitly
$b[G]=\sum_u\pi_{\mu\nu}(u)u_\mu u_\nu$ for any distinct
$\mu,\nu$. It has no extra color factor, because
\eqref{ph:color} used the average of the three diagonal color
Hessians. The quadratic term of the bulk symbol is
$b(|k|^2\delta_{\mu\nu}-k_\mu k_\nu)$; higher lattice
derivative terms have not been set to zero.

For fixed actual $h$, the extractor is complex linear in the member
and ignores additive constants. It is continuous in the exponential
Hessian moment norm: the absolute second moment is bounded by that
norm times $\sup_{r\ge0}r^2e^{-mr}$, with a harmless choice of a
smaller exponential weight. Normal convergence in that norm also
permits termwise extraction and preserves holomorphy in a member
parameter. These are Hessian-norm statements, and make no assertion
that the Hessian norm is uniformly bounded by a potential norm on
coupling-dependent shrinking cores.

\subsection*{The affine quadratic jet and the tent action}

Let $B_\mu(x)=a_\mu+\sum_\kappa C_{\mu\kappa}(x-z)_\kappa$,
using midpoint coordinates or lower starts. The difference is a
constant one-form, which the Ward identity annihilates. Equation
\eqref{ph:row-column-zero} removes every occurrence of $a$, and
\eqref{ph:Maxwell-T} proves
\[
 \frac12\mathcal H_z(B,B)
       =\frac b2\sum_{\mu<\nu}
                                  |C_{\nu\mu}-C_{\mu\nu}|^2.
 \tag{PH19}\label{ph:affine-jet}
\]
For complex fields the squared norm means the bilinear color
pairing; for real fields it is the ordinary squared norm. The
identity follows by expanding
$\sum_{\mu,\kappa}|C_{\mu\kappa}|^2
-\sum_{\mu,\nu}\langle C_{\mu\nu},C_{\nu\mu}\rangle$
into unordered pairs. In particular the axial affine choice
$\ell_\mu(x)=\sum_{\kappa<\mu}(x-z)_\kappa F_{\kappa\mu}$
has curl $F$ and quadratic jet $\frac b2\sum_{\mu<\nu}|F_{\mu\nu}|^2$.

To identify this with the member's own action counterterm, one more
normalization of the actual free response is needed. Assume that on
an affine coarse field $\ell$ of constant curl it reproduces the
same affine bare potential modulo a lattice gradient, and its
current component is zero. This is an affine-response statement
about the actual averaging/background problem, not a consequence
of the tensor classification. For this normalization argument
assume also that, at each bare bond in a bounded coarse region,
the matrix elements of $H$ in the coarse input bond decay
exponentially in its distance from that region. Their constants
must be uniform in its finitely many bare bonds at fixed $\eta$.
This is a precise sufficient free-response input. It makes
$H(\ell\vartheta(\cdot/R))$ converge at every such bond to an
absolutely convergent polynomially weighted matrix sum, by
dominated convergence. Thus its affine reproduction statement is
an equality of well-defined extended first responses, with no
nonlinear affine field inserted outside its analytic domain.

This input also gives the moment bound needed to apply the tensor
theorem to the pointed action itself. Its Hessian is a finite sum
over the bare plaquettes in a unit tent of products of first
differences of $H$. The exponential bounds at their finitely many
bonds give a product bound
$C_\eta e^{-m_H(|x-z|+|y-z|)}$. Summing in the two input bonds
with any smaller exponential weight proves
\eqref{ph:exponential-moment} for this action. This absolute
convergence statement does not assert that $C_\eta$ has a uniform
cutoff bound; any such bound is a further quantitative input.

There is a useful stronger scope. Write
$d_\eta H(p)=\eta^{-1}\sum_{b\in\partial p}H(b)$ with oriented
signs. If the actual response satisfies the physical first-difference
bound
\[
 |d_\eta H\delta_{x,\mu}^a(p)|
      \le C e^{-m_H|x-z|}\qquad (h_z(p)>0),
\]
uniformly in $\eta$, then the moment constant for the pointed
action is uniform in $\eta$ as well. Indeed, its kernel contains
\[
 \sum_p \eta^4 h_z(p)
 (d_\eta H\delta_x)(p)\mathbin{\cdot}
 (d_\eta H\delta_y)(p).
\]
The exact quadrature proved below makes the weighted plaquette
sum one in each orientation. Consequently this kernel is bounded
by $C e^{-m_H(|x-z|+|y-z|)}$, and any smaller exponential moment
is summable uniformly. This explains precisely why physical first
differences, rather than identity-only bounds on bare values, remove
the possible $C_\eta$ loss.

Let $\eta=1/N_0$ for an integer $N_0\ge1$, and use the midpoint
tents
\[
 h_z(p)=\prod_{r=1}^4(1-|p_r-z_r|)_+,
 \qquad A_z(U)=\sum_{\mu<\nu}\sum_{p_{\mu\nu}}h_z(p)A_p(U).
 \tag{PH20}\label{ph:tent-action}
\]
For each fixed plaquette orientation,
\[
                 \eta^4\sum_{p_{\mu\nu}}h_z(p)=1.
 \tag{PH21}\label{ph:quadrature}
\]
Here its barycenters have half-mesh shifts in directions $\mu,\nu$
and zero shifts in the other directions. The one-dimensional
unshifted identity is
$N_0^{-1}[1+2\sum_{r=1}^{N_0-1}(1-r/N_0)]=1$.
For the half-mesh shift it is
$2N_0^{-1}\sum_{r=0}^{N_0-1}(1-(r+1/2)/N_0)=1$.
Multiplication of these four identities proves
\eqref{ph:quadrature}, including the coarse case $N_0=1$.
Also $\sum_z h_z(p)=1$, since the two neighboring one-dimensional
tents sum to one in each coordinate.

For a bare affine potential of curl $F$, the first variation of a
plaquette holonomy in the chart $U_b=e^{it\eta A_b}$ is
$it\eta^2F_{\mu\nu}$. Equation
\eqref{ph:action-normalization}, gradient invariance and
\eqref{ph:quadrature} therefore give
\[
       \frac12D^2(A_z\circ\mathcal P)(0)[\ell,\ell]
                      =\frac12\sum_{\mu<\nu}|F_{\mu\nu}|^2.
 \tag{PH22}\label{ph:action-b-one}
\]
The affine-response hypothesis was used exactly in this line.
Using the polynomial Ward identity to remove symmetric gradients and constants,
polarization of \eqref{ph:action-b-one} determines its affine tensor directly
as \eqref{ph:Maxwell-T} with $b=1$. Writing $K^A_z$ for its
color-averaged composed Hessian, the same gauge Ward calculation
gives its separate row and column zeroth moments. Applying the algebra
of \eqref{ph:P-versus-T} to its absolutely convergent pointed
contraction $\sum_x K^A_z(x,x+u)$ therefore gives
$P^A_{\mu\nu,\mu\nu}=1$ for $\mu\ne\nu$, the moment
extractor $b[A_z]=1$. Identifying this contraction with the full bulk
sum over base points additionally requires translation covariance of
$A_z\circ\mathcal P$; applying the full tensor theorem also requires
its stated hypercubic covariance. Gauge covariance alone does not
supply those spacetime properties. The direct identities
\eqref{ph:affine-jet} and \eqref{ph:action-b-one} already give
\[
 D^2\bigl(F_z-b[G](A_z\circ\mathcal P)\bigr)(0)
                                    [\ell,\ell]=0
 \tag{PH23}\label{ph:own-counterterm}
\]
for every affine constant-curvature field. Polarization gives the
corresponding mixed affine identity. This proves the exact
extractor-to-counterterm match with the normalization
$A(e^{iB})=\frac12|dB|^2+O(B^3)$.

\subsection*{Selecting different members at different sites}

The proof can be applied separately to each member $G^{[z]}$
chosen at a site. Within that member the family at all hypothetical
base points is translation covariant, so its own bulk extractor
$b[G^{[z]}]$ is defined by \eqref{ph:bulk} and
\eqref{ph:extractor}. Its pointed total at the actual site $z$
is canceled by that number times $A_z$, by
\eqref{ph:own-counterterm}. No translation covariance of the
array $z\mapsto G^{[z]}$ has been used.

This distinction is necessary for an inhomogeneous scalar source:
one may freeze it to a different constant-source member at each
site, but must use the coefficient of that same member at that
site. Equality of two unpointed totals alone does not identify
their separate site weights. The nonlinear third-order and H\"older remainder uses the actual
localized response, uniform field-domain inclusion and analytic derivative
bounds stated for the window.

% END MODULE chapters/pointed_hessian_cancellation_v6.tex


% BEGIN MODULE chapters/pointed_window_remainder_v6.tex
\section{The nonlinear remainder in a pointed window}
\label{ch:pointed-window-remainder}

The algebraic cancellation of a pointed Hessian concerns the second derivative
at the identity.  Its use at a nonlinear background requires a local Taylor
estimate.  We prove that estimate from the background hypotheses, retaining
the domain assertions that make the compositions meaningful.  In particular,
the datum is large compared with its value close to the base point.  A
Taylor bound using only its global supremum would lose the required power
of the scale ratio.

\subsection*{The background obligations and their coordinate radii}

Work first in units in which the coarse spacing of a single-scale window is
one.  The background package A1--A3 is used in the following explicit form.
On a box with collars there is a holomorphic potential/current representative
\(\Phi_\Lambda(B)\), with \(\Phi_\Lambda(0)=0\), for the actual constrained
critical orbit.  Its datum domain contains
\[
 \mathcal B_{r_B}
   =\{B:\hbox{determining bonds}\longrightarrow\mathfrak{su}(2)^\mathbb C:
                        \sup_c|B(c)|<r_B\}.
 \tag{PW1}
\]
Here \(B(c)\) is a dimensionless logarithm of a group-valued coarse bond.
This is a product of balls at bonds.  It has no coupled derivative
restriction.  The orbit representative, its current, and the response
estimates hold throughout this complex domain.  Core and tower memberships
under restriction and substitution are also part of A1; they are not inferred
from the formal existence of a logarithm.  All differentiations and coordinate
sums below are first carried out on finite boxes.  Their uniform bounds permit
the stated coherent infinite-volume limits; no representation of a derivative
on an infinite \(\ell^\infty\) space by coordinate derivatives is assumed.

For clarity, the output norm over a top of local spacing \(\xi_y\) is
\[
 Z_y=(\xi_yH,\xi_y^2\nabla H,\xi_y^3J),\qquad
 |Z_y|=\max\{\xi_y|H|,\ \xi_y^2|\nabla H|,\ \xi_y^3|J|\}.
 \tag{PW2}
\]
The first-response bound supplied by A3 is
\[
 |\partial_{B(c)}Z_y(B)[e]|
       \le C e^{-m d(y,y_c)}|e|,
       \qquad B\in\mathcal B_{r_B},
 \tag{PW3}
\]
together with the stated spatial H\"older bound for \(\nabla H\).
The distance is the local-scale path distance.  For the window below all
tops have the same scale; after rescaling, (PW3) uses ordinary lattice
distance and all the output weights are one.  Constants may depend on the
fixed window size, the decay rate and the H\"older exponent, but not on the
bare mesh, the torus or the number of preceding steps.

For the finite-box comparison we also use the background package's
compatibility of the actual linearised equations and their boundary rows.
In the natural force norms, the A3 inverse bound from one coarse pillar
\(y'\) to \(y\) is \(C e^{-m d(y,y')}\), including the output differences
and current used here.  A cutoff of an affine solution leaves a forcing
of polynomial size supported near the boundary, or with an exponentially
decaying boundary tail if a gauge row is exponentially local.
This boundary-locality assertion must hold for the actual gauge and
constraint operators.  It follows by commuting the cutoff through their
local enlarged saddle equations, or through their compatible exponentially
local rows.  It is not a consequence of (PW1)--(PW3) for unrelated maps.
The single-pillar inverse estimate includes the integration or cell-volume
weights in its natural norms; a bare entrywise bound without those weights
would not control a fine-mesh sum.

Choose a fixed basis of the three-dimensional complex colour space.  If
\(\sup|B|\le r_B/2\), varying one basis coordinate by
\(|s|<c_{\rm col}r_B/2\) stays in (PW1), where \(c_{\rm col}>0\) depends
only on this choice of norm and basis.  Thus there is an admissible
coordinate Cauchy radius independent of the number of bonds.  The factors
\(\xi_y\) in (PW2) are output weights and commute with this Cauchy estimate.
If instead a physical potential datum is written \(a(c)=B(c)/\xi_c\),
its coordinate radius is \(c_{\rm col}r_B/(2\xi_c)\), and
\(\partial_{a(c)}=\xi_c\partial_{B(c)}\).  Contracting a second derivative
against \(a(c)a(c')\) cancels both input factors \(\xi_c,\xi_{c'}\).
No inverse local spacing is lost by using dimensionless bond coordinates.

\begin{lemma}[Product decay for the second response]
On \(\mathcal B_{r_B/2}\), the first-response hypothesis implies
\[
 |\partial_{B(c)}\partial_{B(c')}Z_y(B)[e,e']|
 \le \frac{C}{r_B}
       e^{-\frac m2(d(y,y_c)+d(y,y_{c'}))}|e|\,|e'|.
 \tag{PW4}
\]
The same conclusion holds for the local first-order pair norm on a domain
\(X\), with \(d(y,y_c)\) replaced by \(\operatorname{dist}(c,X)\).
\end{lemma}

\begin{proof}
Apply the coordinate Cauchy estimate to
\(\partial_{B(c)}Z_y(B)[e]\), using a slightly smaller fixed fraction of
the radius just specified.  This gives
\(Cr_B^{-1}e^{-m d(y,y_c)}|e|\,|e'|\).
The finite-dimensional holomorphic map has symmetric second derivatives,
so the same argument interchanging \(c,c'\) gives the corresponding bound
with \(c'\).  The smaller of two positive bounds is at most their geometric
mean, which proves (PW4).  Taking the supremum over the output bonds,
first differences and current components gives the stated local norm
version.  Colour indices cost only a fixed factor.
\end{proof}

This argument uses response decay at every complex datum in a ball.
Decay of the linear response at \(B=0\) alone is insufficient.

\subsection*{Weighted polynomial data and the local expansion}

Fix a base point \(z\), a scale ratio \(0<\sigma\le1\), and a window whose
distance from \(z\) to its boundary is comparable with \(\sigma^{-1}\).
Write \(a\) for a radius of a regular member, \(a_{\rm out}\) for the
output field radius and
\[
 \rho=a_{\rm out}/a\ge1,\qquad
     a\le C_0r_B,\qquad a_{\rm out}\le c_0,\qquad
     0<\alpha\le1 .
 \tag{PW5}
\]
The constant \(c_0\) is chosen small enough for the matrix charts.
A fixed constant \(C_{\rm ax}\) in the axial datum estimate below is understood.
For the local Taylor argument we require
\(C_{\rm ax}a_{\rm out}\sigma\le r_B/2\).
This is a global input-domain condition, distinct from the local radius-ratio
dichotomy used later.  The ordinary age bound will deal with its failure at
the finitely many young ages.
A harmless enlargement of the radius ratio to at least one is allowed.
For a domain \(X\) containing \(z\), let \(D=2+\operatorname{diam}X\), in
the present units.

The local regularity conclusion of A1--A3 for the combined background and
fluctuation gives, on the box before choosing its axial gauge, a logarithm
\(B\) with
\[
 |B|\le C a_{\rm out}\sigma,\quad
 |\nabla B|\le C a_{\rm out}\sigma^2,\quad
 |\nabla B(x)-\nabla B(x')|
       \le C a_{\rm out}\sigma^{2+\alpha}|x-x'|^\alpha.
 \tag{PW6}
\]
These are spatial regularity bounds, not consequences of holomorphy alone.
All statements in this section require the actual pair and its restrictions
to remain in the member's prescribed core.  The small chart ball used for
Taylor's formula must be contained in that core.  The radius \(a\) is a
fixed fraction of its curvature radius.  Throughout the admitted core
\(|\partial\mathbf U(p)-1|\le C a\eta_j^2\), so its own tent action obeys
\(|A(h_z,\mathbf U)|\le C a^2\), by the determinant identity and exact
tent plaquette count.  This is a bound on the whole core, not just its
contained chart ball.

In the axial gauge about \(z\), the datum has a local form
\[
 \widehat B(c)=\ell_F(c)+r(c),\qquad
 \ell_{F,\mu}(x)=\sum_{\kappa<\mu}(x-z)_\kappa F_{\kappa\mu},
 \tag{PW7}
\]
where \(|F|\le C a_{\rm out}\sigma^2\), and
\[
 |\widehat B(c)|\le C a_{\rm out}\sigma^2(1+|c-z|),\qquad
 |r(c)|\le C a_{\rm out}\sigma^{2+\alpha}(1+|c-z|)^2.
 \tag{PW8}
\]
The first bound also gives
\(\sup|\widehat B|\le C_{\rm ax} a_{\rm out}\sigma\), so the stated
global input-domain condition places \(\widehat B\) and its segment from zero
in \(\mathcal B_{r_B/2}\).

Here are details of (PW7)--(PW8).  Fixing the transports along a coordinate
comb to the identity expresses a remaining bond as a product of at most
\(C(1+|c-z|)\) transported plaquette holonomies.  With constant plaquette
strengths, the linear part of this product is exactly \(\ell_F\).
Equation (PW6) shows that a plaquette logarithm differs from its value at
\(z\) by at most
\[
 C\big[a_{\rm out}\sigma^{2+\alpha}(1+|c-z|)^\alpha
              +a_{\rm out}^2\sigma^3(1+|c-z|)\big].
\]
The second term includes differentiation of the non-Abelian commutator.
Conjugating a plaquette logarithm by a comb transport adds at most
\(C a_{\rm out}^2\sigma^3(1+|c-z|)\).
Since \(a_{\rm out}\) is bounded and \(\sigma^{1-\alpha}\le1\), both
errors are bounded by
\(C a_{\rm out}\sigma^{2+\alpha}(1+|c-z|)\).
Summing along the strip proves the second bound in (PW8).

For completeness, the finite-product estimate used here follows by
differentiating
\(P(t)=\prod_i e^{tX_i}\) twice: if \(\sum_i|X_i|\le c\), then
\(\sup_{0\le t\le1}|P''(t)|\le C(\sum_i|X_i|)^2\).
The analytic logarithm in this fixed matrix ball has bounded first two
derivatives, so
\[
 \left|\log\prod_i e^{X_i}-\sum_iX_i\right|
                         \le C\left(\sum_i|X_i|\right)^2.
\]
Its error here is at most
\(C a_{\rm out}^2\sigma^4(1+|c-z|)^2\), which is absorbed in (PW8).
The unexpanded product bound gives the first bound of (PW8).
All transports stay in a fixed matrix ball, because the original comb has
length \(O(\sigma^{-1})\) and its original bond logarithms are
\(O(a_{\rm out}\sigma)\).

We shall use the elementary weighted sum
\[
 \sum_{c\in\mathbb Z^4}
       e^{-\mu\operatorname{dist}(c,X)}(1+|c-z|)^k
          \le C_{\mu,k}D^{k+4},\qquad k\ge0.
 \tag{PW9}
\]
Indeed \(|c-z|\le D+\operatorname{dist}(c,X)+C\).
The shell at distance between \(q\) and \(q+1\) is contained in a box
with at most \(C(D+q+1)^4\) bonds.  Bounding each shell by this larger
number and summing \(e^{-\mu q}(D+q+1)^{k+4}\) proves (PW9).
Only a polynomial in \(D\) is spent.  Every power needed below is fixed.

Integrate the first response on the segment \(t\widehat B\), and then use
(PW3), (PW8) and (PW9).  With a fixed integer \(s\), increased below when
necessary, this gives
\[
 \|\Phi_\Lambda(\widehat B)\|_X\le C D^s a_{\rm out}\sigma^2,\qquad
 \|D\Phi_\Lambda(0)r\|_X
              \le C D^s a_{\rm out}\sigma^{2+\alpha}.
 \tag{PW10}
\]
Taylor's formula with integral remainder for \(\Phi_\Lambda\), followed by
(PW4) and the product of the two weighted sums, gives
\[
 \|\Phi_\Lambda(\widehat B)-D\Phi_\Lambda(0)\widehat B\|_X
       \le \frac{C D^s}{r_B}a_{\rm out}^2\sigma^4.
 \tag{PW11}
\]
The use of (PW9) in (PW11) is essential.  Using only
\(\sup|\widehat B|=O(a_{\rm out}\sigma)\) would produce the wrong power
\(\sigma^2\) in its right-hand side.

\subsection*{The finite-box linear response}

The free response reproduces an affine constant-curvature field, modulo a
gauge direction, with zero linear current.  This follows from the actual
linearised equations: the affine field has constant curl, zero current
and zero divergence; centred block averages reproduce it with the correct
bond-length normalization.  It therefore solves the stationarity,
constraint and gauge equations.  For the coherent response, the finite-box comparison below identifies its affine extension by exhaustion and coordinatewise dominated convergence; no uniqueness assertion in a class of unbounded fields is required.
The response is applied only inside gauge-invariant Hessians, so its gauge
direction has no effect.

We also explain the finite-box error instead of identifying a box response
with an infinite response.  The two responses solve the same local
linearised equations in the interior.  Multiply the affine solution by a
cutoff equal to one on the interior and zero near the box boundary.
Its first two finite differences are bounded in the present scale units;
it can be obtained by sampling a fixed smooth profile across a boundary
layer of fixed coarse width.  The cut-off affine pair satisfies the imposed
homogeneous boundary conditions.  Its discrepancy from the prescribed
equations and constraints is supported in that boundary layer.
The equations and their boundary conditions give a bound
\(C|F|(R+1)^s\) for that forcing, where \(R\) is the box radius.
Here the constraint rows, possibly exponentially local gauge rows and their differences
are measured in the local norms of A2--A3.  Apply the inverse bound of A3
from this boundary layer to \(X\).  Decompose the forcing into coarse
pillars, so the sum has unit-scale density rather than bare-mesh density.
For an exponentially local boundary tail use
\(d(x,\partial\Lambda)\le d(x,y)+d(y,\partial\Lambda)\);
splitting the two decay exponents in half and summing a polynomial
unit-scale shell gives the same bound with a smaller fixed rate.
When
\(\operatorname{diam}X\le\min\{R/4,\operatorname{dist}(z,\partial\Lambda)/2\}\),
every point of \(X\) is at distance at least
\(\operatorname{dist}(z,\partial\Lambda)/2\) from the boundary.
For \(R\asymp\sigma^{-1}\), this distance is comparable with \(R\),
and the inverse estimate gives
\[
 \|D\Phi_\Lambda(0)\ell_F-L_F\|_X
       \le C |F|(R+1)^s e^{-cR},
 \tag{PW12}
\]
where \(L_F\) is the infinite affine response pair.  Polynomial boundary
counts, and the finitely many local differences used to form the current,
are absorbed in \(s\).  This argument uses agreement of the actual local
equations; bounds for unrelated kernels on two regions would not imply
(PW12).

For \(R\asymp\sigma^{-1}\),
\((R+1)^s e^{-cR}\le C_\alpha\sigma^\alpha\).  Combining
(PW7), (PW10)--(PW12), we obtain
\[
 \begin{gathered}
 \Phi_\Lambda(\widehat B)=L_F+E_X,\qquad
 \|L_F\|_X\le C D^s a_{\rm out}\sigma^2,\\
 \|E_X\|_X
 \le C D^s\left[
 a_{\rm out}\sigma^{2+\alpha}
       +r_B^{-1}a_{\rm out}^2\sigma^4\right].
 \end{gathered}
 \tag{PW13}
\]
Domains too large for the interior condition in (PW12) are put into the
exponentially long tail.  In applications their tree length is at least
\(c/\sigma-C\).

\subsection*{Taylor's formula for the actual member}

Let \(G_X\) denote a regular pointed member on \(X\), in its exponential
potential/current chart.  Assume its supremum on that core is
\(g_X=\mathsf g e^{-\kappa d_j(X)}\), and that its chart contains the
ball of radius \(a\).  A fixed change of constants allows the use of a
ball of radius \(2a\) and estimates on its half ball; this is the specified
room in the field domains.  Constant gauge transformations act by the
adjoint representation on both chart components.  There is no invariant
linear functional on \(\mathfrak{su}(2)\), so
\(DG_X(0)=0\).  In particular the second derivative of the background
does not enter the Hessian of this composition:
\[
 D^2(G_X\circ\Phi_\Lambda)(0)[B,B']
       =D^2G_X(0)[D\Phi_\Lambda(0)B,D\Phi_\Lambda(0)B'].
 \tag{PW14}
\]
This assertion concerns regular members without stored colour spectators.

\begin{theorem}[Pointed nonlinear window remainder]
\label{thm:pointed-window-remainder}
Under the background and domain hypotheses above, for the near domains
there is a fixed integer \(s_1\) such that
\[
 \begin{split}
 G_X(\Phi_\Lambda(\widehat B))-G_X(0)
    &=\tfrac12D^2G_X(0)[L_F,L_F]+\mathcal R_X,\\
 |\mathcal R_X|
    &\le C\mathsf g e^{-\kappa d_j(X)}D^{s_1}
       \big[\rho^2\sigma^{4+\alpha}+\rho^3\sigma^6\big].
 \end{split}
 \tag{PW15}
\]
The same estimate holds after absorbing the member's own tent-action
counterterm, with its own Hessian subtracted.
\end{theorem}

\begin{proof}
Put \(t=D^s\rho\sigma^2\), increasing \(s\) to cover every polynomial
factor in (PW10)--(PW13).  First suppose \(t\) is less than a sufficiently
small fixed constant.  Equation (PW10) then places the composed pair in
the half chart ball of \(G_X\).  Cauchy's formula and Taylor's theorem give
\[
 \left|G_X(P)-G_X(0)-\tfrac12D^2G_X(0)[P,P]\right|
                  \le Cg_X(\|P\|_X/a)^3,\qquad
 \|D^2G_X(0)\|\le Cg_X/a^2.
 \tag{PW16}
\]
In (PW13), divide by \(a\), use \(a/r_B\le C_0\), and expand the
difference of the two quadratic forms.  The product of \(L_F/a\) and
the first error term is bounded by
\(CD^{2s}\rho^2\sigma^{4+\alpha}\).
The product with the second error term is bounded by
\(CD^{2s}\rho^3\sigma^6\).  The square of the first error is smaller
than the first of these bounds.  The square of the second is bounded
by \(C t^4\), which is at most a constant times \(t^3\) in this case.
The mixed product is bounded by the sum of the squares.  Finally the
Taylor remainder in (PW16) is at most \(C g_Xt^3\).
Increasing \(s_1\) proves (PW15).

If \(t\) exceeds that fixed small constant, use the actual core bound
\(|G_X(P)-G_X(0)|\le2g_X\), rather than Taylor's formula at \(P\).
The Hessian estimate still bounds its polynomial value at \(L_F\) by
\(C g_Xt^2\).  For \(t\ge c>0\),
\(1+t^2\le C_c t^3\), and this proves (PW15) after increasing \(s_1\)
again.  Thus large radius ratios do not impose an upper bound on the age.

The coefficient of a member obeys
\(|b(G)|\le C\mathsf g/a^2\) by its Hessian moments.  The tent action
on its fixed local support has norm at most \(Ca^2\).
Consequently its counterterm is another holomorphic family of size
\(C\mathsf g\), with bounded support.  Apply the same proof to this
family and subtract.  The affine field and its curvature in the two
quadratic terms are the same.
\end{proof}

The exact own-counterterm identity \eqref{ph:own-counterterm} of the
pointed Hessian cancellation theorem now applies to this last,
normalised family: its complete sum of pointed curvature tensors is zero.
A sum cut to the window is the complete sum minus the omitted domains.
The omitted domains have \(d_j(X)\ge c/\sigma-C\).
If the number through a cube of length \(q\) is at most \(Ce^{\kappa_Dq}\)
with \(\kappa>\kappa_D\), polynomial factors in \(D\le C M(d_j(X)+2)\)
are summable.  Thus (PW15), summed over near domains, gives
\[
 C\mathsf g\big[\rho^2\sigma^{4+\alpha}+\rho^3\sigma^6\big],
 \tag{PW17}
\]
and completing the tensor sum costs at most
\(C\mathsf g\rho^2\sigma^4 e^{-c/\sigma}\).
If a far domain is not contained in any non-wrapping chart, its value has
only the norm bound and contributes an additional
\(C\mathsf g e^{-c/\sigma}\).  Such a term has no justified factor
proportional to the actual curvature: a flat field around a hole or a torus
cycle need not be a pure gauge.  It must be retained when a field-sensitive
tail estimate is claimed.

For the radius estimate used in a window, \(\rho\ge1\) permits both
exponentially small tails to be absorbed in (PW17).  Moreover
\[
 \rho^2\sigma^{4+\alpha}+\rho^3\sigma^6
       \le2\rho^3\sigma^{4+\alpha},
       \qquad \rho\ge1,\quad \sigma\le1,\quad \alpha\le1.
 \tag{PW18}
\]
There are \(O(\sigma^{-4})\) level-\(j\) base points in one window of the
later scale.  Its remaining factor is therefore
\(C\mathsf g\rho^3\sigma^\alpha\).  With
\(\rho^2\le C(1+n-j)\) and \(\sigma=L^{-(n-j)}\), this is the summable
ordinary weight
\(C\mathsf g(1+n-j)^{3/2}L^{-\alpha(n-j)}\).
The global input-domain condition also holds at every sufficiently old age.
Indeed
\[
 \frac{C_{\rm ax}a_{\rm out}\sigma}{r_B}
       \le C_{\rm ax}C_0\rho\sigma
       \le C(1+n-j)^{1/2}L^{-(n-j)}.
\]
Choose a fixed age \(a_\star\) beyond which this is at most \(1/2\).
At each of the finitely many younger ages, the given core membership bounds
the member and its own counterterm by
\(C\mathsf g(1+\rho^2\sigma^4)\).
Here \(\rho\) is uniformly bounded and
\(\sigma\ge L^{-a_\star}\), so this is bounded by a fixed constant times
\(\mathsf g\rho^3\sigma^{4+\alpha}\).
This last argument uses the ordinary age bound; it is not valid for an
arbitrary unrelated pair \((\rho,\sigma)\).
Each site may carry a different Euclidean member: cancellation and the
estimate are performed at that site with that member's own coefficient.

\begin{proposition}[Vacuum-only window bound]
\label{cor:vacuum-only-window}
For an analytic gauge-invariant member without stored colour spectators,
no lattice covariance is needed to obtain
\[
 |G_X(\Phi_\Lambda(\widehat B))-G_X(0)|
   \le C\mathsf g e^{-\kappa d_j(X)}D^{s_2}
                                      \rho^2\sigma^4.
 \tag{PW19}
\]
After the domains in a window have been summed, the bound is
\(C\mathsf g\rho^2\sigma^4\), with the same qualifications on remote,
non-wrapping charts and young ages as above.
\end{proposition}

\begin{proof}
Gauge invariance under constant colour rotations again gives \(DG_X(0)=0\),
independently of any spatial symmetry.  By (PW10), with
\(t=D^s\rho\sigma^2\), the local pair lies in the half chart ball whenever
\(t<c\).  Cauchy's estimate for the second-order Taylor remainder then gives
\(C g_Xt^2\).  For \(t\ge c\), use the core bound \(2g_X\le C_cg_Xt^2\).
This proves (PW19), and the domain entropy absorbs its fixed polynomial
in \(D\).  A remote domain not admitting a chart contributes only
\(C\mathsf g e^{-c/\sigma}\); it is bounded by
\(C\mathsf g\rho^2\sigma^4\) for \(\rho\ge1\).
For the finitely many young ages at which the global input-disc condition
is unavailable, the core bound is absorbed since
\(\rho\ge1\) and \(\sigma\ge L^{-a_\star}\).
Finally the \(O(\sigma^{-4})\) sites in the window cancel \(\sigma^4\),
leaving the ordinary weight \(\rho^2\) for a defect or a remainder.
\end{proof}

% END MODULE chapters/pointed_window_remainder_v6.tex


% BEGIN MODULE chapters/localized_gaussian_comparison_v7.tex
\section{Comparison of the actual localized complex Gaussian functional}
\label{ch:localized-gaussian-comparison}

The finite regular expansion uses a normalized Gaussian functional on a
local set of fluctuation coordinates. Its auxiliary Gaussian can have
as many coordinates as the surrounding volume. We prove its comparison
bound without paying for that auxiliary dimension. The argument permits
sharp, measurable cutoffs; it never shifts their arguments into the
complex domain.

\subsection*{The local matrices and the normalized functional}

All matrix norms below are operator norms for the Euclidean complex
space. Transposition is bilinear transposition, not conjugate
transposition. Let \(C=C^{\mathsf T}\) be an invertible complex
\(m\)-by-\(m\) matrix and let \(\Gamma\) be a complex
\(m\)-by-\(n\) matrix, where \(m\ge1\) and \(n\ge0\).
Assume real reference matrices \(C_0=C_0^{\mathsf T}>0\) and
\(\Gamma_0\), and one constant \(K\ge1\), such that
\begin{equation}
 \|C_0\|+\|C_0^{-1}\|+\|\Gamma_0\|\le K,
 \qquad
 \|C-C_0\|+\|\Gamma-\Gamma_0\|\le\varepsilon.
 \tag{FG1}
\end{equation}
The number \(\varepsilon\) will be smaller than a constant depending
only on \(K\). No dimension enters that choice. Define
\begin{equation}
 D=C+C\Gamma\Gamma^{\mathsf T}C,
 \qquad D_0=C_0+C_0\Gamma_0\Gamma_0^{\mathsf T}C_0.
 \tag{FG2}
\end{equation}
The real matrix \(D_0\) is positive definite. In the small neighborhood
in (FG1), normalized Gaussian densities are defined by
\begin{equation}
 d\mu_D(b)=(2\pi)^{-m/2}(\det D)^{-1/2}
             e^{-b^{\mathsf T}D^{-1}b/2}\,db,
 \qquad b\in\mathbb R^m.
 \tag{FG3}
\end{equation}
The square root is the holomorphic continuation of the positive root
at the reference. We justify both convergence and normalization below.

The auxiliary representation of the localized functional is
\begin{equation}
 \begin{split}
 {\cal J}_{C,\Gamma}[F]
  =\int d\mu_{I_n}(\xi)\,
       e^{-\xi^{\mathsf T}\Gamma^{\mathsf T}C\Gamma\xi/2}
       \int d\mu_C(b)\,e^{b^{\mathsf T}\Gamma\xi}F(b).
 \end{split}
 \tag{FG4}
\end{equation}
For \(n=0\) the outer integral is absent. An integrand \(F\) need
only be Borel measurable on \(\mathbb R^m\). Suppose
\begin{equation}
 |F(b)|\le F_0 e^{\rho|b|^2},\qquad F_0\ge0,\quad\rho\ge0.
 \tag{FG5}
\end{equation}

\begin{theorem}[Local Gaussian comparison]
\label{thm:localized-complex-Gaussian}
There are positive constants \(\varepsilon_0(K),\rho_0(K)\) and
\(A_K\), independent of \(m,n\), such that (FG1),
\(\varepsilon\le\varepsilon_0\) and \(\rho\le\rho_0\) imply
absolute convergence of (FG4),
\begin{equation}
 {\cal J}_{C,\Gamma}[F]=\int F\,d\mu_D,
 \qquad {\cal J}_{C,\Gamma}[1]=1,
 \qquad
 |{\cal J}_{C,\Gamma}[F]|
       \le F_0e^{A_K(\varepsilon+\rho)m}.
 \tag{FG6}
\end{equation}
The last estimate also bounds \(\int |F|\,|d\mu_D|\).
In particular its exponential cost uses the local dimension \(m\),
even when the auxiliary dimension \(n\) grows with the full volume.
\end{theorem}

\subsection*{Accretivity, integration and the determinant branch}

\begin{proof}
The joint precision in (FG4) is the symmetric block matrix
\begin{equation}
 Q=\begin{pmatrix}
 C^{-1}&-\Gamma\\
 -\Gamma^{\mathsf T}&I_n+\Gamma^{\mathsf T}C\Gamma
 \end{pmatrix}.
 \tag{FG7}
\end{equation}
At the real reference it factors as
\begin{equation}
 Q_0=L_0^{\mathsf T}
       \begin{pmatrix}C_0^{-1}&0\\0&I_n\end{pmatrix}L_0,
 \qquad
 L_0=\begin{pmatrix}I_m&-C_0\Gamma_0\\0&I_n\end{pmatrix}.
 \tag{FG8}
\end{equation}
Both \(\|L_0\|\) and \(\|L_0^{-1}\|\) are at most
\(1+K^2\). Thus \(Q_0\ge c_K I\) for a positive \(c_K\).
The resolvent identity for \(C^{-1}-C_0^{-1}\), followed by
expanding the products in (FG7), gives
\(\|Q-Q_0\|\le A_K\varepsilon\) for small
\(\varepsilon\); the constants depend only on \(K\).
For real \((b,\xi)\), consequently,
\[
 \operatorname{Re}(b,\xi)^{\mathsf T}Q(b,\xi)
                   \ge(c_K/2)(|b|^2+|\xi|^2).
\]
Choose \(\rho_0<c_K/8\). This proves absolute convergence of the
joint integral with (FG5), hence Fubini's theorem applies. It does
not yet give the desired dimension bound; taking its absolute value
before integrating \(\xi\) would pay unnecessarily for \(n\).

Set \(A=I_n+\Gamma^{\mathsf T}C\Gamma\). It is a small perturbation
of the real positive matrix \(I_n+\Gamma_0^{\mathsf T}C_0\Gamma_0\).
Integrate \(\xi\) first. The normalized Gaussian exponential formula
gives the factor
\[
 (\det A)^{-1/2}
       \exp\{b^{\mathsf T}\Gamma A^{-1}\Gamma^{\mathsf T}b/2\}.
\]
This formula for a complex linear force does not require a contour
change. One proof normalizes the complex Gaussian first and
differentiates its exponential transform in the force. Integration
by parts, justified by accretivity, gives
\(\nabla_\ell G=A^{-1}\ell G\), with \(G(0)=1\), and hence
\(G(\ell)=e^{\ell^{\mathsf T}A^{-1}\ell/2}\).

For completeness, the normalization used in that proof follows
without a complex change of variables. On an accretive symmetric
precision neighborhood the real-space integral is holomorphic by a
fixed Gaussian majorant. It equals \((\det A)^{-1/2}\) at every
real positive symmetric matrix there by real orthogonal
diagonalization and the one-dimensional real Gaussian integral.
Successive one-variable identity theorems in the independent
symmetric matrix entries extend the equality to the complex
neighborhood, with the root continued from its positive value.
The real symmetric matrices have a full set of real coordinates
for this argument; no identity is being extended from a
lower-dimensional background graph.

The elementary inverse and determinant identities are
\begin{equation}
 \begin{split}
 D^{-1}&=C^{-1}-\Gamma
             (I_n+\Gamma^{\mathsf T}C\Gamma)^{-1}\Gamma^{\mathsf T},\\
 \det D&=\det C\,\det(I_n+\Gamma^{\mathsf T}C\Gamma).
 \end{split}
 \tag{FG9}
\end{equation}
The inverse identity follows by multiplication with (FG2).
For the determinant identity, eliminate either off-diagonal block in
\(\bigl(\begin{smallmatrix}I_m&-C\Gamma\\
\Gamma^{\mathsf T}&I_n\end{smallmatrix}\bigr)\); this proves
\(\det(I_m+C\Gamma\Gamma^{\mathsf T})
=\det(I_n+\Gamma^{\mathsf T}C\Gamma)\).
The roots in (FG9) agree at the real reference, and agree throughout
the connected small perturbation neighborhood by continuity and
nonvanishing. The remaining \(b\) integral is precisely (FG3).
This proves the first identity in (FG6), including for Borel \(F\),
since no complex argument of \(F\) was introduced. The same
normalization argument gives \(\int d\mu_D=1\).

\subsection*{The local-dimension estimate}

The reference bounds imply
\(D_0\ge C_0\), \(\|D_0\|\le A_K\) and
\(\|D_0^{-1}\|\le K\). Expanding (FG2) yields
\begin{equation}
 T=D_0^{-1/2}(D-D_0)D_0^{-1/2},\qquad
 \|T\|\le A_K\varepsilon=:\eta.
 \tag{FG10}
\end{equation}
Shrink \(\varepsilon_0\) so that \(\eta\le1/4\).
Then
\(D^{-1}=D_0^{-1/2}(I+T)^{-1}D_0^{-1/2}\) and
\(\|(I+T)^{-1}-I\|\le\eta/(1-\eta)\).
For the real symmetric matrix
\[
 R=D_0^{1/2}\bigl(\operatorname{Re}D^{-1}-2\rho I_m\bigr)D_0^{1/2}
\]
we obtain \(\|R-I_m\|\le A_K(\varepsilon+\rho)=:u\).
Shrink the constants further so that \(u\le1/2\).
Taking the absolute value of the \emph{marginal} density, and then
integrating the real positive Gaussian, gives
\begin{equation}
 \int e^{\rho|b|^2}|d\mu_D(b)|
   =|\det(I+T)|^{-1/2}(\det R)^{-1/2}
   \le(1-\eta)^{-m/2}(1-u)^{-m/2}.
 \tag{FG11}
\end{equation}
Every singular value of \(I+T\) is at least \(1-\eta\), and
every eigenvalue of \(R\) is at least \(1-u\). Finally
\(-\log(1-v)\le2v\) for \(0\le v\le1/2\) proves (FG6).
\end{proof}

\subsection*{The matrices supplied by the regular step}

We now identify the matrices to which the comparison applies. Let
$P=(\mathbf U,\mathbf J)$ be a regular complex pair. In the notation
of Section~\ref{ch:rg-objects-v2}, put $V=\mathcal M^k\mathbf U$ and
\begin{align}
 \Delta_k(P)&=(Q\mathsf G_1Q^{\mathsf T})^{-1}-a
       -\mathsf E^{(2)}[\mathsf H_1^{\mathsf T}\mathbf J],
       \label{eq:actual-regular-Delta}\\
 \mathsf A(P)&=\mathsf C_V^{\mathsf T}
       \bigl(\mathrm{GF}^{(2)}_V+\Delta_k(P)\bigr)\mathsf C_V.
       \label{eq:actual-regular-precision}
\end{align}
Here $\mathsf C_V$ eliminates the pivot bonds. The many-scale rows
$Q$ come from the straightened average. The quadratic form
$\mathsf E^{(2)}$ comes from the plain one-step constraint on $V$.
These two averages have different roles.

At a background pair $P_W=(U_+,J(U_+))$, the form in
\eqref{eq:actual-regular-precision} is the Hessian on the fibre of
$A\circ U_k+\mathrm{GF}_k$. Indeed, the multiplier in the
background equation is the derivative of the critical action with
respect to its datum. Differentiating that equation gives
$(Q\mathsf G_1Q^{\mathsf T})^{-1}-a$. A curve in the fibre of the
next average has a second derivative in its constraint; pairing it
with the multiplier gives the last term in
\eqref{eq:actual-regular-Delta}. Thus the Gaussian of the step is
$C^{(k)}=\mathsf A(P_W)^{-1}$. Soft gauge fixing supplies the
positive quadratic form on the orbit directions, on which the
critical action has zero Hessian.

The comparison starts with the flat bound
$\mathsf A(\mathbf1,0)\ge\gamma_0>0$. Here and below the
constants are independent of the cutoff and the volume. Fix
$\ell_0$, a sufficiently large fixed multiple of $M_1$.
On each cube of side $O(\ell_0)$ the real factor has a real gauge
in which its potential and first difference are bounded by
$a_{\rm w}$. The core supplies this gauge with
$a_{\rm w}=CLM a_0$, where $a_0,a_1$ denote its two radii.
Schwarz estimates first compare $\mathsf A(P)$ with
$\mathsf A(U,0)$ and then, in this gauge, with the flat operator.
The walks that leave the cube cost $Ce^{-c\ell_0}$. For a
partition $\sum_i h_i^2=1$ on these cubes, the localization formula
is
\[
 \mathsf A=\sum_i h_i\mathsf A h_i
       +\frac12\sum_i[h_i,[h_i,\mathsf A]],\qquad
 \left\|\sum_i[h_i,[h_i,\mathsf A]]\right\|
       \le C\ell_0^{-2}.
\]
The double commutator estimate uses the exponential moment of the
walk kernels and the slopes of the cutoffs. Consequently
\[
 \operatorname{Re}\mathsf A(P)\ge
 \gamma_0-C\ell_0^{-2}-Ce^{-c\ell_0}
      -C(a_{\rm w}+a_0+a_1)\ge\tfrac12\gamma_0.
\]
Choose $\ell_0$, then $a_{\rm w}$, then $M\ge C\ell_0$,
then the radii with $CLM a_0\le a_{\rm w}$. Every compression
inherits this gap. The same exponential kernel bounds, by
conjugation, give uniform decay of the compressed shifted
resolvents. Integrating
\[
 \mathsf A^{-1/2}=\frac1\pi\int_0^\infty
        y^{-1/2}(y+\mathsf A)^{-1}\,dy
\]
produces the walk-chain expansion of the inverse root.

Let $Z_0$ be the fine support of an integrand and set
$\chi=\chi_{Z_0}$. Only coarse cubes outside the collar
$Z_0^\sharp$ carry parameters. Decouple the precision and its
inverse root by their respective expansions, and write
\begin{equation}
 C(Z_0;s)=(\chi\mathsf A(s)\chi)^{-1},\qquad
 \Gamma(Z_0;s)=-\chi\mathsf A(s)(1-\chi)\mathsf r(s).
 \label{eq:actual-collared-Gaussian}
\end{equation}
The symbol $\mathsf r(s)$ denotes the decoupled inverse-root
series; it does not denote $\mathsf A(s)^{-1/2}$. The minus
sign in $\Gamma$ matches (FG4). Equivalently one can use the
opposite sign and replace the auxiliary white noise by its
negative. This change leaves (FG2) and every comparison bound
unchanged.

The reference matrices are obtained by setting $s=0$ and replacing
$P$ by $(U,0)$. They are real, and $C_0^{-1}\ge c\gamma_0$.
For independent $|s_\Delta|\le e^{\bar\kappa}$, a carrier that
meets $m$ parameter cubes contributes at most $e^{m\bar\kappa}$.
The first $2^4$ cubes may meet at a corner. Every further cube
requires length proportional to $M$, whose exponential decay
pays for its parameter once $cM\ge C\bar\kappa$.
A carrier connecting a controlled row in $Z_0$ to a parameter
must also cross the collar. Its bound contains
$e^{C\bar\kappa-cM}$, which is $O(e^{-c'M})$ after $M$ is
chosen. The weighted Schur bounds and the resolvent identity
therefore give
\[
 \|C-C_0\|+\|\Gamma-\Gamma_0\|
       \le C(e^{-cM}+LM a_0+a_1),
 \qquad
 \|C_0\|+\|C_0^{-1}\|+\|\Gamma_0\|\le K.
\]
The corresponding comparison also holds for $C^{-1}$ and
$\Gamma^{\mathsf T}C\Gamma$. Taking the radii small after $M$
makes the perturbation as small as required in (FG1), uniformly
on these independent discs. A complex coefficient ratio $\tau$
changes the structure to $(\tau^{-1}C,\tau^{1/2}\Gamma)$ and
adds $C|\tau-1|$ to the perturbation. The choice
$|\tau-1|\le c(LM)^{-4}$ suffices for the Gaussian comparison;
the free quadratic potential may impose a smaller radius.

Only rows near $Z_0$ have been compared. Far from $Z_0$ the
mutilated form need not have a gap. The collar is essential even
for the reference comparison: without it, the reference rows
along a wall of $Z_0$ can vanish.

Finally set $s=1$. The Schur complement gives
\[
 C+C\Gamma\Gamma^{\mathsf T}C
       =\chi\mathsf A(P)^{-1}\chi.
\]
At $P=P_W$ this is the marginal covariance of the step Gaussian.
The local comparison theorem consequently supplies its normalized
functional, its holomorphy with fixed real cutoff sets, and the
local-coordinate estimate used in (UM4)--(UM6). This construction
uses the regular operator and core bounds; it does not assert a
Gaussian law for the full nonlinear averaging disintegration.

\subsection*{Sharp failures, holomorphic parameters, and scope}

Partition the real coordinates into disjoint bond blocks. If a set
\(P_0\) of distinct blocks fails the tests \(|b_\beta|<\mathfrak b\),
then \(\sum_{\beta\in P_0}|b_\beta|^2
\ge|P_0|\mathfrak b^2\). Therefore, for \(\gamma\ge0\) with
\(\rho+\gamma/2\le\rho_0\), (FG11) gives
\begin{equation}
 \int \prod_{\beta\in P_0}\mathbf1_{\{|b_\beta|\ge\mathfrak b\}}
       e^{\rho|b|^2}|d\mu_D|
 \le \exp\{A_K(\varepsilon+\rho+\gamma/2)m
                     -\gamma\mathfrak b^2|P_0|/2\}.
 \tag{FG12}
\end{equation}
This is a product-of-failures estimate, with one Gaussian tail per
distinct tested bond. Multiplication by additional indicators only
improves the bound.

If the matrices and the potential part of an integrand depend
holomorphically on external parameters, and the sharp cutoff sets
remain fixed, the same estimates give a parameter-independent
Gaussian majorant on each smaller parameter neighborhood.
Fubini and Morera then prove holomorphy of the integrals. No
holomorphy of a characteristic function in its real field argument
is asserted or required.

The rescaled structure \((\tau^{-1}C,\tau^{1/2}\Gamma)\) has
effective covariance \(\tau^{-1}D\). For real \(\tau\) in a fixed
compact subinterval of \((0,\infty)\), rescale the real reference
in the same way; all constants remain uniform. A small complex
displacement of \(\tau\) is included in the matrix perturbation.

In the regular step, let \(N_c(Z)\) be the number of coarse
localization cubes. The local free-bond color-coordinate count obeys
\(m\le c(LM)^4N_c(Z)\). Consequently (FG6) supplies the factor
\(\exp\{A_K(LM)^4(\varepsilon+\rho)N_c(Z)\}\), and (FG12)
supplies the failed-cutoff tail. If a Gaussian comparison instead uses
the symbol \(N(Z)\) for the local coordinate count itself, its bound is
\(\exp\{A_K(\varepsilon+\rho)N(Z)\}\), with no second factor
\((LM)^4\). These are two conventions for the same local cost. The matrix gap and kernel closeness
are the actual A2--A3 obligations, which we must prove; this section proves the resulting
finite-dimensional integration estimate.

Completeness is separate. Equality of this local functional with a
marginal of one whole-torus Gaussian requires the actual completeness
identity for its effective covariance. In the source construction
that identity is supplied on complete physical pairs and their
physical-data continuation. Arbitrary independent-current cores
need only their own local holomorphic functionals and the bounds
proved here. This theorem does not manufacture a global covariance
on those cores, verify nonlinear background/core estimates, or
establish a Gaussian law for actual averaging disintegrations.

% END MODULE chapters/localized_gaussian_comparison_v7.tex


% BEGIN MODULE chapters/unmarked_polymer_majorant_v8.tex
\section{The unmarked majorant for the actual regular polymer gas}
\label{ch:unmarked-polymer-majorant}

We derive the activity and local logarithm bounds from the finite
Mayer, cutoff and coarse-decoupling formulas.  The Gaussian estimate
is used at a potential budget larger than the physical one.  No
logarithm or nonvanishing assertion is needed at that larger budget.
Only after the activities have been estimated do we apply the
proved hard-core polymer theorem.

All bounds below are local on the specified complex potential/current
cores.  At complete physical pairs the completeness identity also
identifies their gas with the regular fluctuation integral.  No
global Gaussian on an arbitrary independent-current pair is assumed.

\subsection*{The actual formulas and their analytic inputs}

Fix a finite four-dimensional periodic cube lattice, with fine
localization cubes of side \(M\) and coarse cubes of side \(LM\).
Both grids tile the torus.  Lengths \(d_f\) and \(d_c\) are tree
lengths in fine and coarse cube units.  There are a fixed finite
number of potential types: the free, intrinsic, already formed
older, and optional inserted potential groups may each be
combined within their type before the Mayer expansion.
Thus this type count is independent of the number of older
levels or source profiles.
A tagged fine domain \(Y\) includes
its type in its name; geometry ignores this tag.
For the actual local potentials assume
\[
 \begin{aligned}
 |V(Y;P,B)|&\le a_Ye^{-\lambda_Vd_f(Y)}
       \left(1+M^{-4}\sum_{b\subset Y}|B(b)|^2\right),
 \qquad a_Y\le a,\\
 a&=\sum_{\text{types}}\|\mathbf V^{\text{type}}\|_{\lambda_V}.
 \end{aligned}
 \tag{UM1}
\]
The potentials are holomorphic on their declared complex fluctuation
boxes and external cores.  If a potential is defined only on a box,
the corresponding sharp cutoff is kept wherever that potential occurs.
The free potential, when present, is defined on all real fluctuation
coordinates.

Let \(\mathbf D\) be a selected set of tagged domains.
Write \(Y_0\) for the union of the free-bond sets requiring a field
box among those selections, and let \(P_0\) be a set of distinct
failed-cutoff bonds outside \(Y_0\).
The exact integrand is
\[
 F_{\mathbf D,P_0}
   =(-1)^{|P_0|}\chi_{Y_0}
       \prod_{b\in P_0}(1-\chi_b)
       \prod_{Y\in\mathbf D}(e^{V(Y)}-1),
 \qquad \chi_b=\mathbf1_{\{|B(b)|<\mathfrak b\}}.
 \tag{UM2}
\]
The cutoff product on \(Y_0\) makes every selected factor
well defined.  Outside that product its integrand is set to zero.
This convention introduces no complex argument into a characteristic
function.

For any set \(S\) of fine cubes, let \(\mathcal N_f(S)\)
consist of \(S\) and its fixed first layer of neighbors in the
given cubical adjacency.  Define \(S^\sharp\) to be the union
of the coarse cubes containing at least one fine cube of
\(\mathcal N_f(S)\).  This is a local cube union, not a
rectangular bounding box.  A bond uses its incident fine cubes
in this definition.  The collar preserves unions.  A fixed
number of additional neighboring cube layers is permitted only
if it is part of the assigned collar in every structural
identity; its fixed width is then included in the constants.

For the index \(\iota=(\mathbf D,P_0)\), let \(Z_0(\iota)\)
be its actual fine support, including the globally defined selected
potentials and the failed bonds.
Its coarse collar is \(Z_0^\sharp\).
For \(Z\supset Z_0^\sharp\), the actual terms and activities are
\[
 \begin{split}
 T(Z,\iota)&=
 \prod_{\Delta\in Z\setminus Z_0^\sharp}
       \int_0^1ds_\Delta\,\partial_{s_\Delta}\,
       {\cal J}(Z_0;s_Z)[F_{\mathbf D,P_0}],\\
 H(Z)&=\sum_{\iota}T(Z,\iota).
 \end{split}\tag{UM3}
\]
Here \(s_Z=1\) on \(Z_0^\sharp\), has the displayed parameters
on the remaining cubes of \(Z\), and is zero elsewhere.
Polymers \(Z\) are nonempty wall-connected coarse domains.
The empty selection contributes only the gas's leading one.
It is omitted from \(H\): on a nonempty \(Z\), its term has a
derivative of the normalized constant functional
\({\cal J}[1]=1\), and is zero.
Disconnected regions factor into their wall-components rather than
being counted again as single polymers.

The local functional \({\cal J}\) has the actual matrix
representation and reference bounds of the local Gaussian
comparison theorem.  Uniformly on independent coarse parameter
discs
\[
 |s_\Delta|<R_s,\qquad R_s>2,
 \tag{UM4}
\]
its matrix perturbation is at most \(\varepsilon\), its gap
constant is \(K\), and its local real scalar-coordinate dimension
obeys
\[
 m(Z_0;s_Z)\le c_m(LM)^4N,\qquad N=\#Z.
 \tag{UM5}
\]
The auxiliary Gaussian dimension is unrestricted.
For an integrand bounded by \(F_*e^{\rho|B|^2}\), including
distinct failed tests \(P_0\), the proved comparison gives
\[
 |{\cal J}[F]|
 \le F_*
 \exp\!\left\{A_K(\varepsilon+\rho+\gamma/2)m
                    -\gamma\mathfrak b^2|P_0|/2\right\}
 \tag{UM6}
\]
when \(\gamma>0\), \(\varepsilon\le\varepsilon_0(K)\), and
\(\rho+\gamma/2\le\rho_0(K)\).
This is the marginal local-coordinate estimate; no factor using
the exterior auxiliary dimension is allowed.
The potentials, matrices and integrals are jointly holomorphic
in the scalar potential couplings and \(s\), with the fixed
cutoff sets just described.

The remaining structural hypotheses are the actual heredity,
normalized-functional and wall-factorization identities.
They identify (UM3) with the finite expansion on complete
physical pairs and define coherent local terms on every local
complex core.  These hypotheses concern the assigned operators
and supports.  They are not consequences of a numerical
Gaussian bound.

\subsection*{The two geometric counts}

Let \(Y^\sharp\) be the locally defined coarse collar above.
All neighborhood widths are fixed independently of \(L\)
and of the diameter or shape of \(Y\).
There is a fixed geometric constant \(C_h\) such that
\[
 \#Y^\sharp\le C_h\left(1+\frac{d_f(Y)}L\right).
 \tag{UM7}
\]
To prove this, take a fine covering tree for \(Y\), with length
within one of its infimum, and traverse every edge twice.
After adding a bounded initial segment, this gives a closed
walk of length at most \(C(d_f(Y)+1)\) passing within a fixed
fine-cube distance of every cube of \(Y\).
Split the walk into segments of length at most \(L\).
Each segment and its fixed fine collar meet a bounded number
of coarse cubes, independent of \(L\ge2\).
The number of segments is at most \(C(1+d_f(Y)/L)\).
This proves (UM7), also on the periodic graph.
A failed bond has a coarse collar with at most a fixed number
of cubes.  Therefore, writing
\[
 r=\#\mathbf D,\qquad v=\#P_0,\qquad
 h=\#(Z\setminus Z_0^\sharp),
\]
and increasing \(C_h\) once,
\[
 N\le C_h\left[
       r+\frac1L\sum_{Y\in\mathbf D}d_f(Y)+v+h\right].
 \tag{UM8}
\]
There is no term proportional to the surrounding volume.

We also use the rooted fine-domain entropy bound.  For every
\(\xi>\kappa_f\),
\[
 S_\xi:=
  \sup_{\text{fine cube }q}
       \sum_{\text{tagged }Y\ni q}e^{-\xi d_f(Y)}<\infty.
 \tag{UM9}
\]
One direct counting proof rounds a covering tree to the
bounded-degree graph of fine cubes, traverses its edges twice
starting at \(q\), and records which neighboring cubes are
selected in \(Y\).
The walk has length at most \(C(d_f(Y)+1)\).
Its choices and the selection flags are bounded by a fixed
constant to that power.
This proves exponential rooted entropy, uniformly on the torus.
The fixed number of type tags changes that constant only.
In particular, if \({\cal Y}(Z)=\{Y:Y^\sharp\subset Z\}\),
\[
 \sum_{Y\in{\cal Y}(Z)}e^{-\xi d_f(Y)}
        \le c_fL^4S_\xi N .
 \tag{UM10}
\]
This follows by charging a domain to any one of its fine
cubes and summing over all fine cubes in \(Z\); overcounting
only enlarges the bound.
The number of free bonds whose collars fit in \(Z\) is at
most \(n_bN\), where
\[
 n_b\le c_b(LM)^4 .
 \tag{UM11}
\]
The bond count in (UM11) and the scalar-coordinate dimension
in (UM5) are distinct counts.  Neither is counted twice.

\subsection*{A Gaussian-only budget and repeated Schwarz}

Choose a positive number \(W_0\), independent of cutoff and volume,
small enough for a Gaussian tilt, and choose
\(\lambda_0>\kappa_f\).
This is not yet the physical potential budget.
For every selected tagged domain introduce its own coupling
\(\zeta_Y\).  For \(a>0\) use the radii
\[
 {\cal R}_Y=\frac{W_0}{a}
             e^{(\lambda_V-\lambda_0)d_f(Y)} .
 \tag{UM12}
\]
We will require \(a\le W_0/2\), so every radius exceeds one.
On the product of these discs, the local density bound (UM9)
gives
\[
 \begin{split}
 \sum_{Y\in\mathbf D}|\zeta_YV(Y;P,B)|
 &\le W_0\sum_{Y\in\mathbf D}e^{-\lambda_0d_f(Y)}
             \left(1+M^{-4}\sum_{b\subset Y}|B(b)|^2\right)\\
 &\le C_pW_0L^4N+C_pW_0M^{-4}|B|^2 .
 \end{split}\tag{UM13}
\]
The constant \(C_p\) includes the finite tags and their
rooted sums, but no number of selected domains or global
volume.
The quadratic term uses the rooted sum at each bond.
The constant term uses the fine-cube count (UM10).

Choose \(W_0,\gamma\) so that
\[
 C_pW_0M^{-4}+\gamma/2\le\rho_0(K).
 \tag{UM14}
\]
For complex \(\zeta_Y\) the Mayer product is bounded by
\(2^r\exp(\sum_Y|\zeta_YV(Y)|)\).
Apply (UM6) with the tilt in (UM13).
Uniformly on both parameter polydiscs its integral is at most
\[
 2^r e^{C_0N}\,p^v,\qquad
 p=e^{-\gamma\mathfrak b^2/2},
 \quad
 C_0=C_pW_0L^4+
 A_Kc_m(LM)^4
       \big(\varepsilon+C_pW_0M^{-4}+\gamma/2\big).
 \tag{UM15}
\]
The failed tests are distinct bond blocks, so the tail is \(p^v\).
The Gaussian theorem was applied only to a fixed real integration
space with measurable cutoffs.
There is no zero-freeness assertion on this larger coupling
polydisc.

The integral vanishes whenever any selected \(\zeta_Y=0\).
Repeated Schwarz on the polydisc consequently divides the
bound by \(\prod_Y{\cal R}_Y\) at \(\zeta_Y=1\).
For clarity, division by \(\prod_Y\zeta_Y\) has removable
singularities; apply the maximum principle on smaller
polydiscs and let their radii increase to (UM12).
This proves the product factor without a factorial.
Mixed coarse Cauchy on circles centered at
\(s_\Delta\in[0,1]\), of radii tending to \(R_s-1\),
then gives the literal term bound
\[
 |T(Z,\mathbf D,P_0)|
 \le e^{C_0N}
       \left(\frac{2a}{W_0}\right)^r
       e^{-(\lambda_V-\lambda_0)
                           \sum_{Y\in\mathbf D}d_f(Y)}
       p^v (R_s-1)^{-h}.
 \tag{UM16}
\]
The coupling derivatives and the coarse derivatives commute
by joint holomorphy and the same Gaussian domination.
If \(a=0\), all nonempty Mayer selections are zero.
Apply the same estimate to the cutoff-only terms, with
\(r=0\); this is also the limiting interpretation of (UM16).

\subsection*{The activity bound, including its linear amplitude}

Fix a desired activity rate \(\Lambda>0\).
For wall-connected coarse polymers there is a fixed \(c_\ell\)
such that
\[
 \ell(Z):=d_c(Z)+5\le c_\ell N .
 \tag{UM17}
\]
A cube-adjacency spanning tree proves \(d_c(Z)\le CN\);
the added five is absorbed because \(N\ge1\).
Choose
\[
 K_*=2+c_\ell\Lambda,\qquad
 \Theta=C_h(C_0+K_*).
 \tag{UM18}
\]
Choose one \(\sigma_0>\kappa_f\), and impose the explicit
rate conditions
\[
 \lambda_V-\lambda_0\ge \Theta/L+\sigma_0,
 \qquad \log(R_s-1)\ge\Theta .
 \tag{UM19}
\]
Using (UM8) to pay for
\(e^{(C_0+K_*)N}\), (UM16) becomes
\[
 |T(Z,\mathbf D,P_0)|
 \le e^{-K_*N}
       \prod_{Y\in\mathbf D}
          \left(\frac{2a}{W_0}e^\Theta
                         e^{-\sigma_0d_f(Y)}\right)
       \prod_{b\in P_0}(pe^\Theta).
 \tag{UM20}
\]
The unused factor
\(((R_s-1)^{-1}e^\Theta)^h\) is at most one.
Every selection has thereby paid its own part of the hull
cost.  This is the place where the coarse hull estimate
retains the favorable \(d_f(Y)/L\).

At fixed \(Z\), dropping the restriction
\(P_0\cap Y_0=\emptyset\) only enlarges the nonnegative sum.
All selected domains and failed bonds still have collars
inside \(Z\), so (UM10)--(UM11) give
\[
 \begin{split}
 \sum_{\substack{\mathbf D,P_0\\
                     (\mathbf D,P_0)\ne(\emptyset,\emptyset)}}
   \prod_{Y\in\mathbf D}u_Y
   \prod_{b\in P_0}q
 &\le
  \prod_{Y\in{\cal Y}(Z)}(1+u_Y)
       \prod_{b:b^\sharp\subset Z}(1+q)-1\\
 &\le e^{\tau N}-1 ,
 \end{split}\tag{UM21}
\]
where
\[
 \begin{split}
 u_Y&=(2a/W_0)e^\Theta e^{-\sigma_0d_f(Y)},\qquad
 q=pe^\Theta,\\
 p_\star&=n_bp,\qquad
 \tau=e^\Theta\left[
        \frac{2c_fL^4S_{\sigma_0}}{W_0}a+p_\star\right].
 \end{split}\tag{UM22}
\]
The products are finite in finite volume.  In particular no
factor from bonds outside \(Z\) has been inserted.
The subtraction of one is essential: it retains a selected
potential or failed cutoff, rather than creating a unit
amplitude for every \(Z\).

Take the physical smallness conditions
\[
 a\le W_0/2,\qquad \tau\le1 .
 \tag{UM23}
\]
For \(N\ge1\),
\(e^{\tau N}-1\le\tau N e^{\tau N}\) and \(N\le e^N\).
Equations (UM20)--(UM23) prove
\[
 \begin{aligned}
 |H(Z)|&\le e^{-K_*N}(e^{\tau N}-1)
   \le\tau e^{-(K_*-2)N}\\
   &\le C_\tau(a+p_\star)e^{-\Lambda\ell(Z)} ,\\
 C_\tau&=e^\Theta
       \max\!\left\{\frac{2c_fL^4S_{\sigma_0}}{W_0},1\right\}.
 \end{aligned}
 \tag{UM24}
\]
This is an estimate for the actual activities (UM3).
Their smallness has not been assumed.
The larger Gaussian budget \(W_0\) was used only to bound
the analytic integrals; the physical budget in (UM23) is
smaller and is what makes the activity amplitude at most one.

\subsection*{The actual union filing and the local logarithm}

Choose a desired logarithmic output rate \(\mu>0\), and set
\(g_*=\Lambda-\mu\).
Let \(N\le C_v(d_c(Z)+1)\), and let the number of coarse
polymers containing a fixed cube with
\(m\le d_c(Z)<m+1\) be at most \(C_De^{\kappa_Dm}\).
These geometric bounds follow from bounded-degree cube
covering walks as above.
There is a fixed \(c_{\rm adj}\) such that polymers
incompatible with \(Z\) contain at least one cube in a
set of at most \(c_{\rm adj}N\) cubes: use \(Z\) and
its wall neighbors.

For these counting estimates use the positive norm majorants
\[
 h_Z=\sup_{P_Z}|H(Z;P_Z)|,\qquad
 \widehat w_Z=h_Ze^{\mu\ell(Z)},
\]
where the supremum is over the declared local core of \(Z\),
not over restrictions of one global field.
At a fixed core pair \(P_X\) on \(X\), actual activities are
used only for polymers \(Z\subset X\), with their values
obtained by the assigned hereditary restrictions.
Their weighted absolute values are bounded by
\(\widehat w_Z\).
The bound with amplitude \(\tau\le1\) in (UM24) gives
\[
 \sup_q\sum_{Z\ni q}\widehat w_Ze^{\#Z}
 \le
 \tau\,C_De^{2C_v}e^{-5g_*}
       \sum_{m\ge0}
               e^{-(g_*-\kappa_D-C_v)m}
 =:\tau B_* .
 \tag{UM25}
\]
Choose the rates so that
\[
 g_*>\kappa_D+C_v,\qquad c_{\rm adj}B_*\le1 .
 \tag{UM26}
\]
The explicit factor \(e^{-5g_*}\) allows the last
inequality by increasing the rate gap.
Only the upper bound \(\tau\le1\) is needed here.
It follows that
\[
 \sum_{Z':Z'\not\sim Z}\widehat w_{Z'}e^{\#Z'}
       \le c_{\rm adj}N\tau B_*\le N .
 \tag{UM27}
\]
This verifies, for the actual activities, the criterion
of the already proved hard-core tree theorem, with
the root weight \(\#Z\).
In particular it verifies that criterion on every hereditary
subsystem \(Z\subset X\), without assigning actual fields
outside \(X\).

If two polymers are incompatible, their covering trees
can be joined at a common cube or wall with length at
most five in four dimensions.
For a connected incompatibility list this gives
\[
 \ell\!\left(\bigcup_{i=1}^nZ_i\right)
                         \le\sum_{i=1}^n\ell(Z_i).
 \tag{UM28}
\]
Indeed a spanning tree of the incompatibility graph
has \(n-1\) joins, and
\(d_c(\cup_iZ_i)\le\sum_i d_c(Z_i)+5(n-1)\).
Thus the union filing has the same exponential weight.

Define the actual local logarithmic term by
\[
 E(X)=\sum_{n\ge1}\frac1{n!}
       \sum_{\cup_iZ_i=X}
          \phi(Z_1,\ldots,Z_n)\prod_iH(Z_i),
 \tag{UM29}
\]
where \(\phi\) is the connected hard-core graph coefficient.
For a fixed cube \(q\in X\), mark a vertex whose polymer
contains \(q\).  This only enlarges the absolute series.
Symmetry of the ordered list changes its factor
\(1/n!\) into \(1/(n-1)!\).
When dropping the union restriction, retain \(Z_i\subset X\)
for the actual terms at \(P_X\).
Use (UM28), and then the rooted tree theorem for that
subsystem, dominated by the numerical majorants
\(\widehat w\).
Only the subsequent nonnegative majorant sum may be enlarged
to all geometric polymers containing \(q\).
The result is
\[
 \begin{split}
 e^{\mu\ell(X)}|E(X)|
 &\le\sum_{\substack{Z\subset X\\Z\ni q}}
                       \widehat w_Ze^{\#Z}
 \le\tau B_*\\
 &\le C(a+p_\star).
 \end{split}\tag{UM30}
\]
There is no factor \(\#X\) or ambient volume in this
local estimate.  It proves the unmarked potential-ball
majorant used in the finite insertion theorem.

Uniform convergence on smaller external and coupling
domains follows from these same positive majorants.
The actual local functions are holomorphic, so their
union-filed sum is holomorphic on each declared core.
At complete physical pairs, the complete-Gaussian and
wall-factorization identities identify the gas with the
finite regular integral; the hard-core theorem then gives
\(\exp(\sum_XE(X))=\Xi\ne0\).
On general local cores the activity and logarithm bounds
do not require that global Gaussian interpretation.

\subsection*{Parameter choices and the available insertion slot}

All constants above concern fixed geometry and the actual
matrix/core package.  One must verify (UM4)--(UM6) and
(UM14), (UM19), (UM23), (UM26) together.
For example, it suffices that the matrix package control
\(A_Kc_m(LM)^4\varepsilon\) by a fixed constant on the
chosen decoupling discs.
Choose \(W_0\) so its two contributions to \(C_0\) are
bounded, and choose positive \(\gamma\) of order at most
\((LM)^{-4}\) to bound its contribution.
Then \(C_0\) is bounded independently of exterior volume.
One can choose the target rate gap, the coarse parameter
radius and \(\lambda_V\) by (UM18)--(UM19), and finally
make \(a\) small and \(\mathfrak b\) large to obtain (UM23).
This lists explicit sufficient inequalities; it does not
assert them for a matrix family whose constants have not
been controlled on the required enlarged discs.

For an explicit insertion budget, put
\[
 A_a=\frac{2c_fL^4S_{\sigma_0}}{W_0},\qquad
 a_*=\min\!\left\{\frac{W_0}{2},
                           \frac1{4e^\Theta A_a}\right\},
 \qquad p_\star\le\frac1{4e^\Theta}.
\]
Then every total norm \(a\le a_*\) has \(\tau\le1/2\)
and satisfies (UM23), uniformly above this cutoff threshold.
If the base total potential norm is at most \(a_*/2\), reserve
one additional potential slot of norm \(w_\star\le a_*/2\).
The entire slot stays within (UM23) and hence within
(UM24)--(UM30).
This supplies the local uninserted majorant uniformly on
that slot, rather than presupposing a zero-free activity
domain.  The constants may depend on the fixed
localization parameters and rate choices, but not on
cutoff, torus size, old-level count or the auxiliary
Gaussian dimension.

% END MODULE chapters/unmarked_polymer_majorant_v8.tex


% BEGIN MODULE chapters/finite_source_insertion_v7.tex
\section{A source insertion in the finite regular fluctuation integral}
\label{ch:finite-source-insertion}

We keep the calibration, the real action weight, the Gaussian structure,
the cutoffs and the background maps fixed.  The result below concerns the
finite regular fluctuation integral.  It does not enlarge the source domain
of a large-field history operation.  This distinction permits a small
potential direction to be used at its actual norm.
Localized terms and their insertion estimates are defined on the
specified complex local cores, with potential and current independent
there.  Their sum is identified with a complete Gaussian integral
only at physical complete background pairs \(P=p_W\).
An identity on that graph is not extended to arbitrary independent
currents.

\subsection*{The fixed-reference integral}

Let \(\theta\ge\theta_{\min}\), \(g=\theta^{-1/2}\), and
\(\tau=x/\theta\), with \(|\tau-1|<\sigma_\star\).
The fixed small number \(\sigma_\star\), the allowed external pairs
and their Gaussian gap are restricted by the ordinary Gaussian
comparison and potential allowance below.  In particular the
decoupled precisions stay in their coercive regime and the free
potential leaves the declared allowance; no assertion is made
for arbitrary positive \(x/\theta\).  The sourced run considered
here keeps \(x>0\) fixed.
At a physical complete background pair \(P=p_W\), in the
fluctuation coordinates \(B\) on the non-pivot bonds, the reference
measure is the normalized Gaussian with covariance \(\tau^{-1}C(P)\).
The same identity holds on a holomorphic continuation in physical
background data wherever that continuation and the ordinary Gaussian
bounds have been supplied.  It is not an identity on every complex
core pair.  Put
\[
 \Xi(P,z)=\int d\mu_{\tau^{-1}C(P)}(B)\,
       \prod_b\chi_b(B)\,
       \exp\left\{\sum_Y V(Y;P,B)+z\sum_Yv(Y;P,B)\right\}.
 \tag{SI1}
\]
The parameter \(z\) enters only the potentials.  In particular the covariance,
the cutoff thresholds and the Gaussian determinant do not depend on \(z\).
Here the cutoffs are sharp indicators
\(\chi_b(B)=\mathbf1\{|B(b)|<\mathfrak b\}\) on free bonds.
For a bond set \(A\), write
\(\chi_A=\prod_{b\in A}\chi_b\) and
\(\chi_A^c=\prod_{b\in A}(1-\chi_b)\).
Thus the latter denotes a product of complements.
After the old action at the background has been carried separately, the
new fluctuation increment is
\[
 {\cal E}(P,z)={\cal N}[\theta;x](P)+\log\Xi(P,z).
 \tag{SI2}
\]
The field-independent Gaussian number includes
\(-\frac32|\mathfrak X|\log\tau\), as well as the Haar-coordinate,
scaling and gauge-fixing counts.  The localised normalizer
\(\mathbf N^\circ\) and the free-potential family \(\mathbf G_\tau\)
also use the fixed \((\theta,x)\).  All these quantities have zero
\(z\)-derivative.  Moving \(z\) into \(x\), and differentiating the resulting
Gaussian determinant, would describe a different presentation.

For an inserted exponent \(H\), the potential actually inserted is its
fluctuation difference.  Thus, when the fibre field is \(U_P(B)\) and
\(U_P(0)\) is its background, at a physical complete pair,
\[
 \sum_Yv(Y;P,B)=H(U_P(B))-H(U_P(0)).
 \tag{SI3}
\]
Consequently the background subtraction in the step functional is already
present in (SI1); it must not be subtracted a second time.
Equivalently, the raw fibre integral with the additional exponent
\(zH(U_P(B))\) has the extra factor \(e^{zH(U_P(0))}\) outside (SI1).
That factor belongs to the separately carried old action, rather than to
the newly generated cluster activities.

\subsection*{The ordinary Gaussian expansion used here}

We state the analytic inputs at the level at which they are applied.
On an arbitrary complex local core pair \(P\), the actual localized
Gaussian functionals \({\cal J}(Z_0;s;P)\), potentials and
union-filed terms below are defined by their local rules.
We do not assume that the complete values for different \(Z_0\)
all arise as marginals of one global Gaussian at that \(P\).
A potential has norm
\[
 |V(Y;P,B)|\le\|\mathbf V\|_\lambda e^{-\lambda d(Y)}
       \left(1+M^{-4}\sum_{b\subset Y}|B(b)|^2\right)
 \tag{SI4}
\]
on its specified complex field box.  The potential may be defined only
there: the cutoff product on its domain is kept in every term in which
it occurs.  The ordinary Gaussian comparison estimate bounds a localized
decoupled functional by \(F_0e^{CuN(Z)}\) when its integrand is bounded by
\(F_0e^{\rho|B|^2}\), with \(u=\widehat\varepsilon+\rho\).
Here \(N(Z)\) counts the affected real scalar Gaussian coordinates,
of order \((LM)^4N_c(Z)\), where \(N_c(Z)\) is the number of
coarse cubes.  In the cube-count convention the same bound is
\(F_0e^{C(LM)^4uN_c(Z)}\); the dimension factor is counted once.
The Gaussian reference gap, radii and tilt are chosen so these factors are
bounded per coarse cube.  The decoupling parameters have a fixed enlarged
complex disc.  These are the ordinary background/Gaussian inputs; no
source derivative is contained in them.

Choose a potential allowance \(w_\star>0\), independent of the cutoff and
volume, small enough for the following ordinary expansion.  The base
potentials use at most a fixed number of slots, and leave one slot available
for an additional potential of norm at most \(w_\star\).  The same
Gaussian comparison, rooted domain count and decoupling majorant are valid
throughout this allowance.  The convergent hard-core expansion gives
local logarithmic terms \(E(X)\) and bounds
\[
 |E(X)|\le C_\star e^{-\lambda_{\rm out}d(X)}.
 \tag{SI5}
\]
The constant is uniform in all allowed external pairs, decoupling
parameters and potentials.  This is the ordinary uninserted polymer
majorant, not an assumed differentiability bound for a renormalization map.

Here is the explicit finite expansion to which that majorant applies.
Let \(\mathbf D\) be a set of selected interaction domains.  Write
\(Y_0\) for the union of their free-bond sets on which a field box is
required; a globally defined free potential does not contribute to this
union.  Let \(P_0\subset\mathfrak X\setminus Y_0\) be a set of failed
fluctuation cutoffs.  The indicator identities \(\chi_b^2=\chi_b\)
and \(\chi_b(1-\chi_b)=0\) give exactly
\(\prod_b\chi_b=\chi_{Y_0}
\sum_{P_0\subset\mathfrak X\setminus Y_0}
(-1)^{|P_0|}\chi_{P_0}^c\).
Before
decoupling, its integrand is
\[
 F_{\mathbf D,P_0}(z)
    =(-1)^{|P_0|}\chi_{Y_0}\chi^c_{P_0}
             \prod_{Y\in\mathbf D}(e^{V(Y)+zv(Y)}-1).
 \tag{SI6}
\]
The empty choice \((\mathbf D,P_0)=(\emptyset,\emptyset)\)
supplies the leading one in the gas and is omitted from its
activities.  All polymer domains \(Z\) are nonempty and
wall-connected, so no \(H(\emptyset)\) is defined.
Disconnected regions are products of their wall-component
activities, rather than additional single polymers.
Equivalently, for a nonempty \(Z\) the empty choice has at least
one decoupling derivative of the normalized constant functional
\({\cal J}[1]=1\) and contributes zero.
When several potential origins have the same geometric support, their
origin is included in the index \(Y\); the product selects tagged local
terms.  The support and its tree length ignore this tag.
For each support index \(\iota=(\mathbf D,P_0)\) let \(Z_0(\iota)\)
be its support.  The decoupled activities are formed by
\[
 \begin{aligned}
 T(Z,\iota;z)&=
 \prod_{\Delta\subset Z\setminus Z_0(\iota)^\sharp}
          \int_0^1ds_\Delta\,\partial_{s_\Delta}\,
       {\cal J}(Z_0(\iota);s_Z)[F_{\mathbf D,P_0}(z)],\\
 H(Z,z)&=\sum_\iota T(Z,\iota;z).
 \end{aligned}
 \tag{SI7}
\]
In this formula \(s_Z\) is one on \(Z_0(\iota)^\sharp\), is the
displayed free parameter on \(Z\setminus Z_0(\iota)^\sharp\), and
is zero outside \(Z\).  The normalized Gaussian functional
\({\cal J}\) is independent of \(z\).  At a physical complete
pair \(P=p_W\), its complete-structure identity is
\[
 {\cal J}(Z_0;1)[F]
       =\int d\mu_{\tau^{-1}C(P)}(B)\,F(B).
\]
This physical identity, together with interpolation of every remaining
coarse-cube parameter and the ordinary decoupling property, is what
identifies the gas with (SI1).
Two polymers are compatible when they have neither a cube nor a wall
in common.  Multiplicativity over their wall-components yields
\(\Xi=\sum_{\Gamma\ {\rm compatible}}\prod_{Z\in\Gamma}H(Z)\),
with empty product one, at a physical complete pair.
On general local cores, the same local formulas define activities and
their union-filed terms without asserting one global Gaussian
representation.  Those logarithmic terms are
\[
 E(X,z)=\sum_{n\ge1}\frac1{n!}
       \sum_{\cup_{i=1}^n Z_i=X}
            \rho^T(Z_1,\ldots,Z_n)\prod_{i=1}^n H(Z_i,z).
 \tag{SI8}
\]
Here \(\rho^T(Z_1,\ldots,Z_n)\) is the sum over connected graphs on
\(\{1,\ldots,n\}\) of the products of their edge weights; an edge has
weight \(-1\) for an incompatible pair and zero otherwise.
For \(n=1\), \(\rho^T=1\).
Every sum is first formed in finite volume; the logarithmic series and
its differentiated series are justified by the same absolute majorants.

\begin{proposition}[The inserted-potential estimate]
\label{prop:finite-inserted-potential}
Let \(\varepsilon=\|\mathbf v\|_\lambda\).  Under the ordinary
local Gaussian/polymer hypotheses above, on every specified
complex local core pair,
\[
 \left|\left.\partial_zE(X,z)\right|_{z=0}\right|
       \le C\varepsilon e^{-\lambda_{\rm out}d(X)}.
 \tag{SI9}
\]
At a complete physical pair \(P=p_W\), its sum over \(X\)
is the normalized expectation
\[
 \left.\partial_z\log\Xi(P,z)\right|_{z=0}
   =\frac{ \int d\mu\,\prod_b\chi_b\,e^{\sum_YV(Y)}
                            \sum_Yv(Y)}{
        \int d\mu\,\prod_b\chi_b\,e^{\sum_YV(Y)}}.
 \tag{SI10}
\]
There is no volume factor in the local estimate (SI9).
\end{proposition}

\begin{proof}
If \(\varepsilon=0\), every inserted potential vanishes on its domain
and the assertion is immediate.  Otherwise take the actual auxiliary
potential \(z\mathbf v\) on
\(|z|<w_\star/\varepsilon\).  This is a legitimate potential in the
actual local Mayer and decoupling expressions,
regardless of whether it would be an admissible physical source for a
history construction.  The ordinary Gaussian comparison and the polymer
criterion hold on this disc.  Every integrand and every activity is
holomorphic there; Gaussian domination and the parameter-independent
majorant give holomorphy of the union-filed series (SI8).
This local argument does not require one global Gaussian at a
general independent-current core pair.
Apply Cauchy's estimate on
\(|z|=w_\star/(2\varepsilon)\) to (SI5).  It gives (SI9) with
constant \(2C_\star/w_\star\), which depends only on fixed ordinary
allowances.

The same argument gives a norm derivative, which is useful when the
actual source family of potentials is nonlinear.  On the potential
space with its fixed declared domains, joint holomorphy on every
two-dimensional affine plane makes the directional derivative
complex linear: the degree-one term of
\(E_X(\mathbf V+z_1\mathbf h_1+z_2\mathbf h_2)\) is the sum of
its two degree-one terms.  Cauchy's estimate makes this linear map
bounded.  For \(\|\mathbf h\|_\lambda<w_\star/2\), apply
one-variable Cauchy to \(t\mapsto E_X(\mathbf V+t\mathbf h)\)
on the disc \(|t|<w_\star/\|\mathbf h\|_\lambda\).
Summing its Taylor coefficients of degree at least two gives
\[
 \left|E_X(\mathbf V+\mathbf h)-E_X(\mathbf V)
                         -DE_X(\mathbf V)[\mathbf h]\right|
 \le C C_\star\frac{\|\mathbf h\|_\lambda^2}{w_\star^2}
                         e^{-\lambda_{\rm out}d(X)}.
\]
All constants are uniform in the external core pair, so this is
also the remainder bound in the unpointed supremum norm
\(\|F\|_{\rm out}=\sup_{X,P}
e^{\lambda_{\rm out}d(X)}|F_X(P)|\).
It does not assert a pointed or stronger rooted norm bound.
Uniform local Lipschitz bounds on smaller potential balls give the
unique continuous extension to the potential-norm completion.
The displayed remainder passes to that extension and proves
Fr\'echet differentiability there as well.
If an actual ordinary source family satisfies
\(\mathbf V(w)=\mathbf V(0)+w\dot{\mathbf V}+o(|w|)\)
in that norm, the remainder bound proves its chain rule
\(\partial_wE_X(\mathbf V(w))|_0=
DE_X(\mathbf V(0))[\dot{\mathbf V}]\).
Pointwise holomorphy of \(\mathbf V(w)\) on its ordinary source
disc, with a uniform potential-norm bound there, gives its
norm Taylor expansion by the same uniform Cauchy estimates.
Thus the auxiliary linear direction computes the actual first
source jet; the nonlinear physical source family need not exist
on the large auxiliary-potential disc.

This is also the marked-activity derivative of the explicit expansion.
The derivative of (SI6) replaces precisely one selected factor by
\[
 v(Y)e^{V(Y)}
       \prod_{\substack{Y'\in\mathbf D\\Y'\ne Y}}
                                     (e^{V(Y')}-1).
 \tag{SI11}
\]
Domains with \(V(Y)=0\) are included: their unmarked term is zero,
but their derivative need not be zero.  In a term with a marked potential
the cutoff on its domain is retained.  Derivatives commute with (SI7)
by the same Gaussian comparison and with (SI8) by absolute convergence.
Thus this proof differentiates the actual finite interaction formula,
not a hypothetical normed map without a specified integral.
The pinned cluster estimate is what keeps the local Cauchy bound (SI5)
independent of volume.

At a complete physical pair, the sum rule for the ordinary
cluster expansion gives
\(\exp(\sum_XE(X,z))=\Xi(P,z)\ne0\).
Differentiating it and the finite integral proves (SI10).
The denominator is this one regular Gaussian integral; no division by
the full sourced partition function or by a history mass is used.
\end{proof}

\subsection*{The actual scalar action direction}

Let \(H_u(U)=A(u,U)\), where \(u\) is a bounded function on bare
plaquettes.  It is filed by assigning a plaquette to the union of the
fixed cubes containing its corners.  In the background field domain,
\(|\partial U(p)-1|\le C\alpha\eta^2\).
The identity
\(a_p(U)=\frac12\det(1-\partial U(p))\) gives a raw local action
member of norm at most \(CM^4\alpha^2\sup|u|\).
It has no coupling counterterm at this stage.

More precisely, at a complete physical pair the potential in (SI3)
is a first difference of this action.
Suppose the actual background response and its decoupled
substitutions stay in their stated field domains and shift a plaquette
by at most \(C\delta\eta^2\), where \(0<\delta\ll\alpha\).
Along this shift, the plaquette remains at distance at most
\(C\alpha\eta^2\) from the identity.  The derivative of the determinant
of a \(2\times2\) matrix is bounded by a constant times its size, so
\[
 |a_p(U_P(B))-a_p(U_P(0))|
                          \le C\alpha\delta\eta^4.
 \tag{SI12}
\]
This proof works at complex pairs.  It uses the actual curvature and
response bounds, not positivity.  On a local core the same determinant
bound applies to the two actual substituted plaquette matrices whenever
the declared response bounds give these curvature and difference
estimates; it does not require a global fibre map there.

We give the quantitative fine-response localization separately from
the coarse Gaussian decoupling in (SI7).
For a plaquette \(p\), let \(X_0(p)\) be the union of its fixed
corner cubes.  Its number of cubes is at most a dimensional constant
\(C_{\rm base}\), and its tree length is zero.
Write \(\mathcal U(P')\) for the bare configuration used by the
actual background evaluation of a pair \(P'\), and set
\(F_p(P')=u(p)a_p(\mathcal U(P'))\).
Thus this local function is explicitly a function of a pair.
The ordinary response-bracket hypotheses used here are the following.
On each \(Y\supset X_0(p)\) the same response rule gives
\[
 \Phi_{p,Y}(s,\nu)=
 F_p\big(P|_{X_0(p)}\dotplus\nu\,\mathfrak h_Y(s;B)\big),
 \qquad
 |s_\Delta|<R=e^{\bar\kappa},\quad |\nu|<\rho,
\]
holomorphic jointly on these independent discs, uniformly in the
allowed complex external fields and fluctuation box, with
\[
 \rho\ge c_R\alpha/\delta>1,\qquad
 \sup|\Phi_{p,Y}|\le C\alpha^2\eta^4|u(p)|.
\]
Here \(c_R>0\) can contain the fixed loss
\(e^{-C\bar\kappa}\).
In particular all substituted pairs must stay in their actual cores
on the whole joint polydisc, not just for real interpolation parameters.
The operation \(\dotplus\) shifts the potential and current by the
actual response.  These are quantitative background and membership
hypotheses; the following Cauchy argument does not prove them.

For disjoint fine-cube sets \(K,N\), define the mixed bracket by
\[
 \Psi[K|N]=\int_{[0,1]^K}
       \prod_{\Delta\in K}\partial_{s_\Delta}
                        \Psi(s_K+1_N)\,ds_K ,
\]
with all other parameters zero.  Its scalar-action potential is
the bracket of
\(\Psi(s)=\Phi_{p,Y}(s,1)-\Phi_{p,Y}(s,0)\)
with \(K=Y\setminus X_0(p)\), \(N=X_0(p)\).
At \(\nu=0\) the field is the unshifted restricted background.
Schwarz's lemma in \(\nu\) first gives
\(\sup_s|\Psi(s)|\le
2C\rho^{-1}\alpha^2\eta^4|u(p)|\).
For every \(s_\Delta\in[0,1]\), a coordinate Cauchy circle has
radius arbitrarily close to \(R-1\).
Iterated Cauchy therefore bounds the bracket by
\[
 2C\rho^{-1}\alpha^2\eta^4|u(p)|
             (R-1)^{-|K|}
 \le C_R\alpha\delta\eta^4|u(p)|
              e^{-(\bar\kappa-1)|K|},
\]
where \(\bar\kappa\ge1\).
There is no factorial in this estimate, because each independent
cube parameter is differentiated once.
The ordinary rooting property says that a component containing a
differenced cube and missing \(X_0(p)\) makes the bracket zero.
For the present choice \(N=X_0(p)\), every component missing
the base has such a differenced cube.  No assertion about an
arbitrary remote component consisting only of undifferenced
parameters is being made.
The local response rule agrees with the response built on the
active wall-component.  Between local ambient domains on which
the substitutions are defined it has hereditary consistency,
\(\Phi_{p,Y}(s,\nu)=
\Phi_{p,\Lambda}(s\,\mathbf1_Y,\nu)\), with the base parameters
set to one.
At a complete physical pair on the full lattice its endpoints are
\[
 \Phi_{p,\Lambda}(1,1)=u(p)a_p(U_P(B)),\qquad
 \Phi_{p,\Lambda}(1,0)=u(p)a_p(U_P(0)).
\]
Apply the finite fundamental theorem of calculus in every fine-cube
parameter outside \(X_0(p)\), expanding evaluation at one as evaluation
at zero plus the integral of the derivative.
The resulting sum of brackets equals
\(\Phi_{p,\Lambda}(1,1)-\Phi_{p,\Lambda}(1,0)\).
At a complete physical pair this is exactly
\(u(p)[a_p(U_P(B))-a_p(U_P(0))]\).
Wall-component consistency identifies each bracket with the
corresponding \(\Phi_{p,Y}\); rooting removes the disconnected
terms.  Thus the construction gives a uniformly bounded local
extension on every declared core and localizes the actual
difference (SI3) at complete physical pairs.

For a surviving bracket, the ordinary domain geometry gives
\(d(Y)\le C_{\rm geo}(|K|+1)\).
At fixed \(Y\) there are at most
\(CM^4\eta^{-4}(|K|+C_{\rm base})\) possible plaquette anchors.
Choose the fixed surplus \(\bar\kappa\) so that
\(\bar\kappa-1>C_{\rm geo}\lambda+2\), increasing it further
if a subsequent rooted-domain entropy sum is included.
Sum the displayed bracket bound over its plaquette anchors and
use that a linear polynomial times \(e^{-2|K|}\) is bounded.
This proves
\[
 \|\mathbf v_u\|_\lambda
       \le C M^4\alpha\delta\,\sup|u|
       =C M^4\alpha^2 s\,\sup|u|,\qquad s=\delta/\alpha.
 \tag{SI13}
\]
The factors \(\eta^4\) and \(\eta^{-4}\) cancel exactly.
The required membership and response estimates on all decoupling
substitutions are background obligations that we must prove.  The bracket argument neither
creates those inclusions nor replaces them by bounds at real fields.

Combining (SI9) with (SI13) proves a one-step unpointed source response
of size \(C\alpha^2s\sup|u|\) for the actual local regular family
on every declared complex core.  Its physical sum is the source
response of the fixed-reference regular integral.
The Gaussian normalizer contributes zero derivative.
Its interaction vacuum value contributes
\(\partial_zE(X,(1,0),0)\), which is the source vacuum energy and obeys the
same bound.  Subtracting it leaves a factor at most two in the estimate.

\subsection*{Own counterterms and the remaining iteration}

For a constant source, once the actual bond-density pointing is applied,
the derivative member is Euclidean.  Its own coefficient is
\(\partial_z b[E_z]|_0=b[\partial_zE_z|_0]\), because the free response
and the coefficient extractor are fixed and complex linear.
The counterterm to remove from the derivative action at site \(y\) is
therefore this number times \(A(h_y,U)\).  A coefficient belonging to
another member, or a derivative of a Gaussian presentation with a
different weight, would not supply that identity.

There are two further tasks before the one-step estimate is a complete
multilevel source-response bound.  First, renewing the actual pointing
uses three bond-indexed interpolation legs.  Differentiating a normalized
leg expectation produces both the derivative of its leg insertion and a
connected expectation of that insertion with \(\sum_Yv(Y)\).
A root-plus-insertion cluster bound and the path's admissible field room
must control those terms before the same estimate is asserted for each
pointed term.

Second, the old regular derivative members and the nonconstant source
defects have to be inserted again with their own coefficients.  A
constant-source derivative of size \(\alpha_j^2s_j\) can be normalized
by the pointed Hessian theorem.  Write \(\alpha_H>0\) for the fixed
spatial H\"older exponent, as distinct from an analytic field radius.
On a non-wrapping local chart, retaining
the \emph{true} radius ratio
\(\rho_j=\delta_{\rm out}/\alpha_j\) in the near-domain Taylor proof,
before any enlargement to \(\max\{1,\rho_j\}\), would give after the
fine-site count the contribution
\[
 C\alpha_j^2s_j
   \left[\left(\frac{\delta_{\rm out}}{\alpha_j}\right)^2
                       L^{-\alpha_H(n-j)}
       +\left(\frac{\delta_{\rm out}}{\alpha_j}\right)^3
                       L^{-2(n-j)}\right].
 \tag{SI14}
\]
Its quadratic part sums to \(C\delta_{\rm out}^2s_{\rm out}\) when
the ordinary flow and age weights apply.  For
\(\alpha_{\rm out}\asymp x_{\rm out}^{-1/2}(\log x_{\rm out})^q\)
and
\(\delta_{\rm out}\asymp x_{\rm out}^{-1/2}
(\log x_{\rm out})^{p_0}\), this is
\[
 Cx_{\rm out}^{-1}(\log x_{\rm out})^{3p_0-q}.
 \tag{SI15}
\]
Under this homogeneous local refinement, the arithmetic supplies the
indicated scale bound. The derivative recurrence also uses the regular
member and defect hypotheses of the linear bootstrap. In
particular the defects must retain the factor \(\alpha_j^2s_j\) as
well as the scaled Lipschitz gain.  A fixed-disc source bound alone
does not provide it.  Marked and remote terms need their own uniform
fixed-disc bounds.  Field-independent tails around holes or torus cycles
cannot be assigned an extra \(\rho_j^2\) when \(\rho_j<1\); they require
a non-wrapping chart or a separately small class estimate.

% END MODULE chapters/finite_source_insertion_v7.tex


% BEGIN MODULE chapters/renewed_source_pointing_v7.tex
\section{Renewing the actual bond pointing after one source insertion}
\label{ch:renewed-source-pointing}

An estimate for an unpointed localized logarithm is not by itself an
estimate for a particular pointing of that logarithm. We retain the
actual three interpolation legs and differentiate their bond densities.
For a source filed as an older weight member, the Gaussian structure,
free potential, cutoffs and numerical normalizers are fixed. Thus the
first, free-potential leg and the second, cutoff leg are independent of
this source. Only the interaction leg contributes below. This choice
of presentation is essential to the statement.

\subsection*{The finite local logarithm and its actual locality}

Fix a finite regular fluctuation integral, its Gaussian reference, its
cutoffs and all geometric and external data. Fix nested full complex
fluctuation boxes
\[
 \mathcal B_0=\{B:\sup_b|B(b)|<\mathfrak b_0\},\qquad
 \mathcal B_1=\{B:\sup_b|B(b)|<(1+\nu)\mathfrak b_0\},
 \qquad \mathfrak b_0\ge1,\quad \nu>0.
\]
The ordinary spare-potential estimate below accepts potentials
holomorphic and bounded in its growth norm on $\mathcal B_0$;
all its actual Gaussian cutoff supports lie there. The larger box
$\mathcal B_1$ is reserved for the one coordinate Cauchy estimate
on an inserted potential. Both boxes and the external-pair and
decoupling domains are fixed once for this step. Ordinary inputs
whose cutoffs fill half their analytic box leave this room by
choosing an intermediate $\mathcal B_0$. This uses that existing
room once, and is not a new shrinking at each source derivative.
A local potential family
has the Gaussian growth norm
\[
 |V(Y;P,B)|\le \|\mathbf V\|_\lambda e^{-\lambda d(Y)}Q_Y(B),
 \qquad Q_Y(B)=1+M^{-4}\sum_{b\subset Y}|B(b)|^2.
 \tag{SP1}\label{sp:potential-norm}
\]
There are finitely many bond variables but arbitrarily many of them
as the volume grows. The norm and all constants used below are those
of the ordinary regular expansion, uniform in that number. A fixed
finite number of potential types can be combined by addition of their
norms.

Let $E_X(\mathbf V)$ denote the actual union-filed logarithmic term of
the local Gaussian functionals, Mayer expansion and connected-cluster
expansion. These local objects are defined on their independent-current
pair cores. Away from complete physical data, their wall-component
functionals need not be margins of one global Gaussian measure.
Normalized global Gaussian expectations will be used only at a complete
physical pair $P=p_W$, or by holomorphic continuation in the physical
datum $W$ where the same complete construction is defined. We use
the following consequence of its uninserted estimate, as established
for the fixed-reference integral in
Section~\ref{ch:finite-source-insertion}. At every base system used
below there is an available potential allowance $w_\star>0$ such that
\[
 |E_X(\mathbf V+\mathbf U)|
       \le C_\star e^{-\lambda_{\rm out}d(X)},
       \qquad \|\mathbf U\|_\lambda<w_\star.
 \tag{SP2}\label{sp:unmarked-allowance}
\]
The actual finite local Gaussian functionals and the absolutely
convergent logarithmic series are holomorphic in every finite collection of potential
parameters on this allowance. A single additional family is enough:
several auxiliary directions are added together in $\mathbf U$, with
the sum of their norm loads less than $w_\star$.

The locality needed here is part of this actual expansion, not a
choice made after taking its logarithm. A contribution filed at $X$
uses only selected interaction domains whose enlarged union is in
$X$. In particular, if a family $\mathbf a_b$ vanishes except on
domains containing $b^\sharp$, then
\[
 DE_X(\mathbf V)[\mathbf a_b]=0,
 \qquad
 D^2E_X(\mathbf V)[\mathbf a_b,\mathbf v]=0
                  \quad\hbox{if }b^\sharp\not\subset X.
 \tag{SP3}\label{sp:root-locality}
\]
The assertions include domains with a zero unmarked potential. In a
Mayer derivative their zero factor can be replaced by the nonzero
inserted factor, and the marked domain must then be included in the
union. Omitting such terms would invalidate both locality and the
insertion identity.

\begin{lemma}[Two derivatives of the actual localized term]
\label{lem:root-source-local-Cauchy}
For potential families $\mathbf a,\mathbf v,\mathbf c$, the allowance
\eqref{sp:unmarked-allowance} gives
\[
 \begin{split}
 |D^2E_X(\mathbf V)[\mathbf a,\mathbf v]|
 &\le \frac{16C_\star}{w_\star^2}
       \|\mathbf a\|_\lambda\|\mathbf v\|_\lambda
                      e^{-\lambda_{\rm out}d(X)},\\
 |DE_X(\mathbf V)[\mathbf c]|
 &\le \frac{2C_\star}{w_\star}
       \|\mathbf c\|_\lambda e^{-\lambda_{\rm out}d(X)}.
 \end{split}
 \tag{SP4}\label{sp:mixed-Cauchy}
\]
These derivatives preserve \eqref{sp:root-locality}.
\end{lemma}

\begin{proof}
If either direction in the first estimate has zero norm, the claim
is immediate. Otherwise put the single auxiliary family
$\zeta\mathbf a+z\mathbf v$ in the available slot. On the bidisc
with radii $w_\star/(4\|\mathbf a\|_\lambda)$ and
$w_\star/(4\|\mathbf v\|_\lambda)$ its total norm is at most
$w_\star/2$. Apply the two-variable Cauchy formula to the actual
holomorphic term in \eqref{sp:unmarked-allowance}. For the second
estimate use radius $w_\star/(2\|\mathbf c\|_\lambda)$.
The support assertion follows from the union filing before
differentiation. No radius for a full history construction has been
enlarged in this argument.
\end{proof}

Thus a separate second-mark estimate is not needed if the actual
ordinary bound is already uniform on the full potential allowance
\eqref{sp:unmarked-allowance}. If only a one-direction derivative
bound were known, without this allowance and locality, the lemma
would not follow from it.

\subsection*{The interpolation paths and the differentiated density}

Let $\mathbf V_t$, $0\le t\le1$, be the actual interaction leg,
with local bond directions $\mathbf a_{t,b}$ satisfying
\[
 \partial_t\mathbf V_t=\sum_b\mathbf a_{t,b}.
 \tag{SP5}\label{sp:actual-leg}
\]
Every summand is supported on domains containing $b^\sharp$.
At a finite volume the sum is finite. We must prove that the ordinary
interpolation construction supplies the uniform bound
\[
 \sup_b\int_0^1\|\mathbf a_{t,b}\|_\lambda\,dt\le A.
 \tag{SP6}\label{sp:ordinary-leg-norm}
\]
This is a property of the actual unmarked legs and their admissible
field domains. It is not inferred merely from an endpoint norm.
Uniform ordinary bounds and field room are required throughout the
path, with the same allowance \eqref{sp:unmarked-allowance}.
The third-order vanishing of the intrinsic $V^{G,A}$ makes its
$t^{-2}$ radial path regular at zero. It does not by itself identify
the actual bond split. We do not replace that split by an Euler
decomposition or any other pointing to infer
\eqref{sp:ordinary-leg-norm}.

For the inserted source let
\[
 \mathbf V_t(s)=\mathbf V_t+s\mathbf v_t,
 \qquad
 \mathbf a_{t,b}(s)=\mathbf a_{t,b}+s\mathbf c_{t,b},
 \qquad
 \partial_t\mathbf v_t=\sum_b\mathbf c_{t,b},
 \tag{SP7}\label{sp:source-leg}
\]
with the same bond support and with $\mathbf v_0=0$. It is enough
that these are the first source derivatives of holomorphic paths;
the affine notation in $s$ makes the formula easier to read. Suppose
\[
 \sup_{0\le t\le1}\|\mathbf v_t\|_\lambda\le C_v\varepsilon,
 \qquad
 \sup_b\int_0^1\|\mathbf c_{t,b}\|_\lambda\,dt
                                      \le C_c\varepsilon.
 \tag{SP8}\label{sp:source-leg-norm}
\]
Choose $|s|C_v\varepsilon<w_\star/4$. The remaining allowance
in \eqref{sp:unmarked-allowance} bounds the first derivative in
direction $\mathbf a_{t,b}+s\mathbf c_{t,b}$ by a constant times
$\|\mathbf a_{t,b}\|_\lambda+|s|\|\mathbf c_{t,b}\|_\lambda$.
This is integrable by \eqref{sp:ordinary-leg-norm} and
\eqref{sp:source-leg-norm}. On a smaller source disc the mixed
Cauchy estimate similarly dominates its source derivative.
Consequently the actual path integral is holomorphic there and
may be differentiated under the integral.

The actual interaction-leg term carrying bond $b$ is
\[
 E^\bullet_s(X,b)=\int_0^1
 DE_X(\mathbf V_t(s))[\mathbf a_{t,b}(s)]\,dt.
 \tag{SP9}\label{sp:actual-pointed-leg}
\]
Its source derivative contains two terms:
\[
 \dot E^\bullet(X,b)=\int_0^1
 \left\{D^2E_X(\mathbf V_t)[\mathbf a_{t,b},\mathbf v_t]
             +DE_X(\mathbf V_t)[\mathbf c_{t,b}]\right\}dt.
 \tag{SP10}\label{sp:differentiated-leg}
\]
The first is the connected root-plus-source term. The second is
the direct change of the leg insertion. Neither can be dropped.
Equations \eqref{sp:mixed-Cauchy}, \eqref{sp:ordinary-leg-norm}
and \eqref{sp:source-leg-norm} prove
\[
 |\dot E^\bullet(X,b)|
 \le C_\star\left(\frac{16AC_v}{w_\star^2}
                     +\frac{2C_c}{w_\star}\right)
                \varepsilon e^{-\lambda_{\rm out}d(X)}.
 \tag{SP11}\label{sp:bond-bound}
\]
It is zero unless $b^\sharp\subset X$.

Here is the exact normalized-expectation interpretation on a complete
physical pair $P=p_W$. Write
$\langle\cdot\rangle_t$ for the normalized, possibly complex,
fixed-reference integral with interaction $\sum_YV_t(Y)$.
The denominator is nonzero by the same ordinary expansion.
With $a_{t,b}=\sum_Ya_{t,b}(Y)$,
$v_t=\sum_Yv_t(Y)$ and $c_{t,b}=\sum_Yc_{t,b}(Y)$,
\[
 \partial_s\langle a_{t,b}(s)\rangle_{t,s}\big|_{s=0}
   =\langle c_{t,b}\rangle_t
           +\langle a_{t,b}v_t\rangle_t
           -\langle a_{t,b}\rangle_t\langle v_t\rangle_t.
 \tag{SP12}\label{sp:Duhamel}
\]
This follows by differentiating its numerator and denominator.
It requires no positivity. On these complete data, summing
\eqref{sp:differentiated-leg} over $X$ gives \eqref{sp:Duhamel}
by the physical sum rule for the actual connected expansion.
The disconnected product cancels; an unnormalized insertion would
give the wrong formula and could carry a volume factor. On an
arbitrary independent-current core,
\eqref{sp:differentiated-leg} remains the local derivative identity
and bound. It is not identified there with a covariance under an
unavailable global Gaussian measure.

\subsection*{The older-weight path and its field room}

An older weight enters through its fluctuation difference
$v(Y;P,B)$, so $v(Y;P,0)=0$. Its interaction path is
$v_t(Y;P,B)=v(Y;P,tB)$, and
\[
 c_{t,b}(Y;P,B)=\sum_{a=1}^3
                   B_b^a\partial_{B_b^a}v(Y;P,tB).
 \tag{SP13}\label{sp:older-weight-root}
\]
The intrinsic path $t^{-2}V^{G,A}(tB)$ belongs to the terms with
the action or gauge-fixing inverse-coupling coefficient. It must
not be applied to this older weight. In particular, no artificial
vanishing of its linear and quadratic fluctuation terms is assumed.

We spell out sufficient field room for \eqref{sp:source-leg-norm}.
Suppose $v$ is holomorphic and obeys
\eqref{sp:potential-norm} with norm $\varepsilon$ on
$\mathcal B_1$, for the actual external pairs and decoupling
parameters. We prove the bound on the entire $\mathcal B_0$
required by \eqref{sp:unmarked-allowance}, not merely on real
integration points. For real
$t\in[0,1]$, vary one coordinate at $tB$ by a complex number of
radius
\[
 r_b=c_\nu(1+|B_b^a|),
\]
where $c_\nu>0$ is a sufficiently small fixed fraction of the
available room and the fixed color norm. This disc lies in
$\mathcal B_1$, uniformly for $B\in\mathcal B_0$ and in
$t,b$ and the number of variables. Moreover its growth factor is
bounded by $C_\nu Q_Y(B)$: the one changed coordinate adds at
most $C_\nu(1+|B_b^a|^2)$ to the sum defining $Q_Y$, with $M$
fixed. Cauchy's formula therefore gives
\[
 |B_b^a\partial_{B_b^a}v(Y;P,tB)|
       \le C_\nu\varepsilon e^{-\lambda d(Y)}Q_Y(B).
 \tag{SP14}\label{sp:coordinate-source-root}
\]
Also $Q_Y(tB)\le Q_Y(B)$, so
$\|\mathbf v_t\|_\lambda\le\varepsilon$ on the actual cutoff
box. Summing the three color coordinates in
\eqref{sp:older-weight-root} proves \eqref{sp:source-leg-norm}.
The original $B$ is used in this majorant, so no $1/t$ or
$\mathfrak b_0^2$ loss occurs as $t$ tends to zero. Thus the
differentiated root is a potential on exactly $\mathcal B_0$,
while the undifferentiated source used its one fixed larger box.

A complex disc for a common radial parameter alone would not
justify \eqref{sp:coordinate-source-root}. The individual bond
legs require the independent-coordinate field room just used.
Nor does an endpoint inserted-potential norm bound the variation
of an arbitrary interpolation path. The prescribed path and its
ordinary bounds \eqref{sp:ordinary-leg-norm} remain necessary.

\subsection*{Midpoint tents, entropy, and the exact sum rule}

Use the actual geometric tents at bond midpoints. The collars in
the union filing are assumed large enough that
$h_z(b)\ne0$ implies $z\in b^\sharp\subset X$ for a nonzero
bond-root term. The fixed block ratio $L$ gives
\[
 \sum_{z\in X}h_z(b)=1,
 \qquad \sum_bh_z(b)\le4L^4.
 \tag{SP15}\label{sp:tents}
\]
These geometric identities are source independent. The renewed
pointed source family is therefore
\[
 \dot F(X,z)=\sum_bh_z(b)\dot E^\bullet(X,b),
 \qquad
 \sup_{X,z}e^{\lambda_{\rm out}d(X)}|\dot F(X,z)|
                                      \le C\varepsilon.
 \tag{SP16}\label{sp:pointed-source-bound}
\]
The constant depends on $L$, the fixed ordinary allowances, $A$
and the field room, and is independent of volume and field-variable
count. No additional source-position factor is spent in this
supremum norm.

For completeness, the geometric root cost also remains harmless
when an application sums domains through $z$. If that count at
length $m$ is at most $C_D e^{\kappa_Dm}$ and
$\lambda_{\rm out}>\kappa_D$, then a smaller output rate costs
only
\[
 4L^4 C\varepsilon\sum_{m\ge0}(m+1)^q
                e^{-(\lambda_{\rm out}-\kappa_D)m}<\infty
\]
for every fixed polynomial power $q$. The two-variable Cauchy
estimate has already summed all cluster multiplicities and choices
of source domains through the ordinary rooted majorant. No second
unsummed count of them is justified or needed. Additional geometric
position sums cost a polynomial only when the actual domain
geometry supplies that bound. This does not permit unrecorded extra
domain types or an exponent loss larger than the available gap.

Finally sum \eqref{sp:actual-pointed-leg} over $b$ at a fixed
localization domain $X$. The chain rule
and \eqref{sp:actual-leg} give
$E_X(\mathbf V_1(s))-E_X(\mathbf V_0(s))$.
Because $\mathbf v_0=0$, its source derivative is
$DE_X(\mathbf V_1)[\mathbf v_1]$. On every declared local core,
\eqref{sp:tents} therefore proves the algebraic identity
\[
 \sum_{z\in X}\dot F(X,z)
                         =DE_X(\mathbf V_1)[\mathbf v_1].
\]
At a complete physical pair $p_W$, the physical Gaussian sum rule
then gives
\[
 \sum_{X,z}\dot F(X,z)
       =\sum_XDE_X(\mathbf V_1)[\mathbf v_1]
       =\partial_s\log\Xi(\mathbf V_1+s\mathbf v_1)\big|_0.
 \tag{SP17}\label{sp:pointed-source-sum}
\]
The same formula is valid on the permitted analytic continuation
of the physical datum $W$. This is not continuation from the
physical graph to an arbitrary independent current. All sums are
finite-volume sums followed by absolutely convergent cluster sums.
Coherent infinite-volume extensions of the local families use the
same uniform local bounds and companion limits, when supplied.

If the actual reference, union filing and interpolation paths are
covariant under lattice isometries and a constant source has that
same covariance, these identities also hold after transporting
$X,b,P$. The midpoint tent then gives
$\dot F(\mathsf rX,\mathsf rz;\mathsf rP)
 =\dot F(X,z;P)$. The same reasoning preserves gauge invariance
and reality. These properties are inherited from the actual
construction, not inferred from a scalar norm estimate. A
nonconstant source does not acquire Euclidean covariance merely
from \eqref{sp:pointed-source-bound}.

If the source entered a Gaussian structure or the cutoff leg,
additional derivatives would be required. With a fixed reference,
a cutoff root is in general a boundary functional $N_{t,b}$, not the
bounded function $\partial_t\chi/\chi$. Its differentiated
normalized density is
$N_{t,b}(v e^V)/\Xi-(N_{t,b}(e^V)/\Xi)\langle v\rangle$,
plus any direct variation of $N_{t,b}$. If the reference measure
also varies, its score contributes to the denominator derivative
as well. The retained Gaussian tail would have to be proved under
insertion. The fixed-reference
older-weight presentation avoids those additional terms here.

The conclusion is one finite-step pointed insertion estimate.
It does not differentiate a full history class, prove closure of
the response under iteration, supply the spatial defect gain of
a nonconstant profile, or establish the clean-history Gaussian
disintegration needed for an interacting two-point law.

% END MODULE chapters/renewed_source_pointing_v7.tex


% BEGIN MODULE chapters/window_potential_formation_v7.tex
\section{From ordinary windows to local fluctuation potentials}
\label{ch:window-potential-formation}

The estimate in this section is an ordinary formation estimate.
It concerns the actual functions obtained by rewriting old terms on
box backgrounds, before their fine-cube bracket localization.
The outer potential/current pair may have independent components.
This causes no substitution of an arbitrary pair into a theorem
about critical backgrounds: the inner window function first
constructs its own constrained box background and its own current.
The hypotheses below retain the domain and response assertions needed
to perform that construction on the whole complex parameter set.
We fix the window-first rule below as part of the localization:
at one old base point, its near pointed sum is one window
member before any fine bracket is applied.

\subsection*{The maps and geometric inputs}

Fix a step from level \(k\) to level \(k+1\).
For an old level \(j\le k\), put
\[
 a=k-j,\qquad \sigma_j=L^{-a},\qquad
 \rho_j=\max\{1,\alpha_k/\alpha_j\},\qquad
 s_k=\delta_k/\alpha_k .
 \tag{WF1}
\]
The radii and the ordinary flow satisfy
\(\rho_j\le C(1+a)^{1/2}\).
In particular
\[
 C_{\rm age}:=\sup_{a\ge0}C(1+a)^{1/2}L^{-a}<\infty,
 \qquad \rho_j\sigma_j\le C_{\rm age}.
 \tag{WF2}
\]
Here \(\delta_k\) is the tested fluctuation amplitude, whereas
\(\alpha_k\) is the analytic field radius.  The ratio in (WF1)
is an ordinary full-radius ratio.  No small-field gain is being
deduced by replacing a smaller ratio with one.

We spell out the ordinary background and geometry inputs.  They
must concern the actual assigned maps, not arbitrary functions
with similar notation.

First, let \(P\) range over the local complex core used to form the
new potentials.  Let \(\mathbf H_k(P;s,B')\) be the actual response
to the tested free fluctuation coordinates, with the indicated
fine-cube parameters.  The plain nonlinear iterated block average
is denoted \(\mathcal M^j\).  On a window \(\square\), with a fixed
partition function \(\zeta_\square\), the datum used for an inner
term is the restriction of
\[
 {\cal V}_{j,\square}(P;s,t,\nu)=
 \mathcal M^j\!\left(
 e^{\,i\eta_k(t+\nu\zeta_\square)
                  \mathbf H_k(P;s,B')}\,U_+(P)\right),
 \qquad 0\le t\le1 .
 \tag{WF3}
\]
The notation \(U_+(P)\) denotes the configuration in the assigned
background evaluation and \(\eta_k\) the conversion from its
potential to a bare bond logarithm.  Formula (WF3) uses the
nonlinear average.  Replacing it by its differential would not
give the window identity.
The box-background map applied to this datum returns a pair
\[
 P^{\rm box}_{j,\square}(P;s,t,\nu)
   =\big(U^{\rm box}_{j,\square}({\cal V}_{j,\square}),
                    J^{\rm box}_{j,\square}({\cal V}_{j,\square})\big)
 \tag{WF4}
\]
in which \(J^{\rm box}\) is the current of that constrained critical
configuration.  The ordinary maps include their prescribed node
and gauge representatives.

Second, all these maps are holomorphic on the full independent
fine-parameter discs \(|s_\Delta|<R=e^{\bar\kappa}\).
The unamplified substitutions have a fixed positive amount of
datum, chart and core room.  If \(r_{B,j}\) is the dimensionless
datum radius, the background package supplies
\(\alpha_j/r_{B,j}\le C\).
Its response and nonlinear-average estimates bound the change
caused by the \(\nu\) term, in every relevant normalized datum,
chart and core norm, by
\[
 C_R|\nu|\,\delta_k\sigma_j/\alpha_j
            \le C_R|\nu|s_k\rho_j\sigma_j .
 \tag{WF5}
\]
The constant \(C_R\) includes the fixed loss from the enlarged
fine-parameter discs.  These are estimates for the actual
nonlinear average and response on their assigned analytic
domains.  For example, in a datum norm the estimate follows by
the uniform first-response bound, composed with the derivative
of the actual average, and integration along the perturbation.
It must also hold for the other core inequalities and tower
substitutions used in the inner, outer and removed terms.
Existence and holomorphy on the corresponding datum balls are
part of the background package.

It follows from (WF2)--(WF5) and the fixed room that there is a
constant \(c_\nu>0\), independent of \(j,k\), such that all
substitutions are admissible on the common disc
\[
 |\nu|<{\cal R}_\nu=c_\nu/s_k .
 \tag{WF6}
\]
Indeed choose \(c_\nu\) so that
\(C_Rc_\nu C_{\rm age}\) is smaller than half of every
fixed room.  Equation (WF5) then keeps each substitution inside
its assigned domain.  This argument verifies an inclusion
using the actual input-domain estimates; it does not extend
functions to unions of complex gauge orbits.
We take \(s_k<c_\nu/2\), so that \({\cal R}_\nu>2\).

Third, the critical box-background maps in (WF4) obey the
weighted first-response, H\"older and compatible boundary-comparison
estimates used in the nonlinear pointed window theorem.
The complete fields on (WF6), including the unperturbed field,
have output size at most \(C\alpha_k\).
For old ages the axial datum is bounded by
\(C\alpha_k\sigma_j\).  Consequently its global half-datum-ball
condition holds when
\(C\rho_j\sigma_j\) is sufficiently small.
By (WF2) this fails at most at a fixed finite set of young ages.
At those ages the actual core memberships and the whole-core
bound for a member's own tent action remain available.
Finite boxes are treated first; their linearized boundary
comparison is uniform in the bare mesh.  Periodic or distant
tails are retained with their stated exponential lengths.

Finally, assume the actual window identities and elementary
localization geometry.  The partition \(\sum_\square
\zeta_\square=1\) is inserted along the direction of
differentiation.  For complete physical pairs the old action
difference is the sum over windows of
\(\int_0^1\partial_\nu|_0\) of its window evaluations (WF3)--(WF4),
plus the prescribed outer and removed terms.
The same local formulas define their extensions at every local
complex core.  Equality with a global physical fibre integral
is asserted only at complete physical pairs.

An inner window contains at most \(C_M\sigma_j^{-4}\) old base
points.  At a fixed base point the number of old domains of
tree length between \(m\) and \(m+1\) is at most \(C_De^{\kappa_Dm}\).
Diameters are bounded by \(C_M(1+m)\), and the input decay rate
has a fixed surplus over \(\kappa_D\).
Thus every fixed polynomial diameter factor is summable.
An outer field has the additional decay from its distance to
the window, while a removed field has the factor
\(e^{-cM/\sigma_j}\).
In the member's local pair norm in level-\(j\) physical units,
whose analytic radius is \(\alpha_j\), their changes are bounded by
\[
 C_R|\nu|\delta_k\sigma_j e^{-cM/\sigma_j}.
 \tag{WF7}
\]
The corresponding statement with the actual distance is used
when summing outer terms.
Both endpoints and the interpolation segment stay in the same
member core, with a fixed fraction of its analytic radius
available for Cauchy's inequality.
Long input domains retain their original tree decay: one
fraction supplies \(e^{-c_\kappa/\sigma_j}\), and the remaining
fraction supplies the declared output spatial rate.
Here \(c_\kappa>0\) comes from tree length in \(M\)-cubes and
the allocated input rate.  When these long-domain tails are
written together with removed-response tails as
\(e^{-cM/\sigma_j}\), the fixed \(c\) is allowed to depend on
\(M\) and the input rate; it is not asserted independent of them.
These are geometric and response hypotheses, not bounds on
the value of an old member.

\subsection*{The room used by the substitutions}
\label{sec:substitution-room}

The displacement estimate must include the hereditary tower
clauses. Their control uses the curvature of both factors,
not a bound on the complex potential alone. Suppose that
$\mathbf U=e^{i\eta A'}U$ and that $U,\mathbf U$ both have
plaquettes bounded by $\alpha'_0\eta^2$, while
$|A'|,|\nabla_UA'|\le\alpha'_1$. The local expansion is
\[
 \partial(e^{i\eta A'}U)
 =\bigl(1+i\eta^2D_UA'
      +O(\eta^2|A'|^2+
          \eta|A'|\,|\partial U-1|)\bigr)\partial U.
\]
It follows that
$|D_UA'|\le C(\alpha'_0+(\alpha'_1)^2)$.
Since $U$ is real, its adjoint transports commute with the real
conjugation. The real and imaginary parts of $A'$ separately
have this bound. Their real perturbations therefore have
curvature of size $C(\alpha'_0+(\alpha'_1)^2)$.

Let $\phi$ be a rescaled tower curvature or current. It is
holomorphic in $A'$ on a fixed ball of radius $r$. Its linear
Taylor part is determined by the values at $0,A'_{\rm R}$ and
$A'_{\rm I}$. Consequently
\[
 |\phi(A')|\le C\bigl(|\phi(0)|+
       |\phi(A'_{\rm R})|+|\phi(A'_{\rm I})|\bigr)
       +C\sup|\phi|\,\|A'\|^2/r^2.
\]
The real regularity estimate controls the three real values.
Only the quadratic remainder uses the crude tower bound.
At an old level, the complex potential gains one power of the
scale ratio and its square gains the two powers required by
the curvature normalization. Thus, uniformly in the depth,
the tower bound is
$K^\sharp(\alpha'_0+(\alpha'_1)^2)$ in its own units.
The relations $\alpha_1^2\le C\alpha_0$ and
$K^\sharp(\alpha'_0+(\alpha'_1)^2)\le\alpha_0$
are the corresponding restrictions on the radii.

This also explains why a fixed core margin gives actual
analytic room. Apply Schwarz's lemma to a tower quantity on a
larger potential disc on which the preceding a priori bound
holds. Its increment is bounded by a constant times the
potential displacement measured in the allowed radius. The
constant is uniform over the tower and its depth. A sufficiently
small displacement therefore consumes at most a prescribed
fraction of every tower margin. Clauses involving only the
real factor are unchanged by this shift. The independent
current is shifted by its reduced-current response, whose
bound requires only first differences of the potential.

The same argument is used on layered spaces with their
geometric gaps. If $\nu$ layers separate a clause from the
integrated field, the quotient of its response bound by its
margin contains
\[
 2^\nu(1+\nu)^{1/2}
       e^{-c\delta_0ML(\nu-1)}.
\]
This is uniformly bounded after $M$ is fixed. The factor $2^\nu$
is the reciprocal gap, and $(1+\nu)^{1/2}$ is the largest
running-radius loss. No sum over the old levels is introduced.

We apply this room estimate to the three substitutions as
follows. The outer field is the response on the original
domain. The inner field is the background of the enlarged
window, evaluated at the datum obtained by the actual
nonlinear average in (WF3). The two removed fields are the
differences produced by replacing the original response and
the original background by their box versions. Their bounds
have the additional factor $e^{-cM/\sigma_j}$. Their analytic
discs are chosen separately from the inner disc. All these
fields use the assigned node factors. The inner pair in (WF4)
carries the actual current of its generated box background, also
for complex data on its analytic domain. Outer and removed
substitutions use their prescribed reduced-current assignments
away from the critical physical locus. At physical parameters
these assignments agree with the currents of the corresponding
critical configurations.

For the inner substitution, the datum has the form
$\varsigma B+B'$ with $0\le\varsigma\le1$. Here $B$ is the
nonlinear average of the shifted box background and
$|B'|\le c\min\{\alpha_{0,j-1},\alpha_{0,j}\}$.
The box background is taken in Landau form. Its lower-level
restrictions are backgrounds over their own averages. The
running-radius inequality and the two powers of scale in the
curvature estimate put its triple in the core of every level
at which it is used. The box gauge has bounded condition
number; its distance from the identity is not required to be
small when the real factor is an arbitrary pure gauge.

The window multiplier can therefore be amplified until its
displacement uses a fixed fraction of this room. For a term
of level $j$ one obtains the radius
\[
 \varrho_F\asymp
       \sigma_j^{-1}\alpha_j\delta_k^{-1}
       e^{-C\bar\kappa}
       \ge c\sigma_j^{-1}(\rho_j s_k)^{-1}
       e^{-C\bar\kappa}.
\]
Because $\rho_j\le C(1+k-j)^{1/2}$ and
$\sigma_j=L^{-(k-j)}$, their product is uniformly bounded.
It follows that $\varrho_F^{-1}\le Ce^{C\bar\kappa}s_k$,
and that a common disc $|\nu|<c_\nu/s_k$ is available.
This supplies the scale arithmetic in (WF5)--(WF6) and the
corresponding normalized displacement estimate in (OA7).
In applying it, the three substituted maps and their node and
current conventions are kept fixed throughout the disc.
The estimate is a statement about those maps, including all
their tower clauses; it is not a consequence of a bound on an
old activity's Hessian.

\subsection*{Renormalize the critical inner function first}

Let an old regular member have norm \(g_j\), be gauge invariant
and Euclidean as a pointed total, and have its own coefficient
\(b_j\).  At every base point remove its vacuum value and
\(b_jA(h_z,\cdot)\) before bracket interpolation.
The pointed Hessian theorem fixes this coefficient by the
unit normalization of the same tent action.
Site-frozen members may differ at different base points: the
argument uses each member's own hypothetical Euclidean family
and its own coefficient.  No translation invariance of that
array is required.

For \(j<k\), here is the quantitative window estimate on (WF4).
The axial field near a base point has linear constant-curvature
part of size \(C\alpha_k\sigma_j^2\).
The H\"older and nonlinear response errors have sizes
\(C\alpha_k\sigma_j^{2+\alpha_H}\) and
\(C r_{B,j}^{-1}\alpha_k^2\sigma_j^4\), respectively.
Cauchy's second and third derivatives of a member cost
\(Cg_j/\alpha_j^2\) and \(Cg_j/\alpha_j^3\).
The linear derivative vanishes by constant-color invariance.
The quadratic value on the affine part is canceled on the
whole pointed total at that base point by the member's own
tent action.  No coefficient is assigned to an individual
domain \(X\).  The individual Taylor remainders have a fixed
polynomial in \(D=2+\operatorname{diam}_jX\) times their
original spatial decay; the domain count above sums these
polynomials.  Multiplying the remaining errors therefore
gives, at one base point after this domain sum,
\[
 Cg_j\left[
       \rho_j^2\sigma_j^{4+\alpha_H}
                     +\rho_j^3\sigma_j^6\right].
 \tag{WF8}
\]
This is the local Taylor argument on the constrained critical
box pair, with its uniform weighted response sums.
When its local Cauchy-size parameter is large, the whole-core
bound and the own-action bound give the same estimate after
increasing that summable polynomial, since \(1+T^2\le CT^3\)
for \(T\ge1\).
The finite set of young ages failing the global half-datum-ball
condition is handled by the same whole-core bounds; their
scale ratios have a fixed positive lower bound.
The resulting constant is independent of the number of old
levels.

Now sum the \(C_M\sigma_j^{-4}\) base points.
Because \(\rho_j\ge1\), \(0<\alpha_H\le1\), and
\(\sigma_j\le1\), the renormalized inner window function has
bound
\[
 Cg_j\rho_j^3\sigma_j^{\alpha_H}.
 \tag{WF9}
\]
Truncating the pointed total at the enlarged window contributes
only the long-domain tail already set aside in the geometry.
It does not replace the pointed total by an unnormalized
individual term.

A defect or an aged remainder has no coupling counterterm.
After its value at the identity is removed, constant-color
invariance still eliminates the linear derivative, provided
there are no fixed color spectators.  The vacuum-only window
argument gives at one base point, after the same polynomial
domain sum, \(Cd_j\rho_j^2\sigma_j^4\), where \(d_j\) is that origin's norm.
The large local-size case uses \(1\le CT^2\) for \(T\ge1\).
After the base-point sum its inner bound is
\[
 Cd_j\rho_j^2 .
 \tag{WF10}
\]
Stored-field boundary origins are not included in this
statement.  Their spectators and their hull geometry need the
separate boundary formation estimate.

At \(j=k\) no age improvement is needed.
The vacuum-normalized old member together with its own action
has whole-core norm at most \(Cg_k\), by the coefficient bound
\(|b_k|\le Cg_k/\alpha_k^2\) and the core bound on its tent action.
The actual fluctuation change has size \(C\delta_k\) within
a core with fixed analytic room \(c\alpha_k\).
Cauchy's inequality therefore gives \(Cs_kg_k\), with its
original spatial decay.  There is no fine-volume age factor
at this level.  This is the term with weight \(\omega_1\);
no cancellation between different current-level domains is
used for it.

\subsection*{The one small factor and the outer terms}

Fix \(t\in[0,1]\) and the fine parameters.
Subtract the window function's value at \(\nu=0\).
It vanishes there and is holomorphic on (WF6).
The bounds (WF9)--(WF10) hold on that whole disc.
Schwarz's lemma, or Cauchy's formula for its derivative at zero,
therefore supplies precisely one factor \(1/{\cal R}_\nu\le Cs_k\).
The actual window formula uses that derivative and its integral
over \(t\); the interval has length one.
Thus an old regular inner term contributes
\(Cs_kg_j\rho_j^3\sigma_j^{\alpha_H}\), and a vacuum-only
origin contributes \(Cs_kd_j\rho_j^2\).
No second amplification factor is extracted later.

Outer and removed terms need no marginal cancellation on a
general independent-current pair.
Use Cauchy's inequality in the actual member core and (WF7).
Before its \(t\)-integration and the fine-site count, a difference
is bounded by
\(C g_j|\nu|(\delta_k/\alpha_j)
\sigma_j e^{-cM/\sigma_j}\).
Equivalently, use the uniform disc (WF6) and then its one
Schwarz or Cauchy factor.
After the fine-site count the bound is
\[
 Cs_kg_j\rho_j\sigma_j^{-3}e^{-cM/\sigma_j}.
 \tag{WF11}
\]
The same reasoning applies to a vacuum-only origin with \(g_j\)
replaced by \(d_j\).
For \(0<\sigma\le1\) and \(\rho\ge1\),
\[
 \rho\sigma^{-3}e^{-cM/\sigma}
       \le C\rho^3\sigma^{\alpha_H},
 \qquad
 \rho\sigma^{-3}e^{-cM/\sigma}\le C\rho^2.
 \tag{WF12}
\]
Indeed the suprema of
\(\sigma^{-3-\alpha_H}e^{-cM/\sigma}\) and
\(\sigma^{-3}e^{-cM/\sigma}\) are finite.
Any fixed additional polynomial from spatial summation is
absorbed in the same way, by decreasing the exponential rate
once.  Long-domain tails are treated with their retained tree
decay and the same elementary inequality.
This proves the ordinary outer and removed bounds without
applying the pointed cancellation to an arbitrary pair
\((U,J)\).

\subsection*{Fine brackets and the resulting formation estimate}

We now make the base of an inner bracket explicit.
For \(j<k\) and one old base point \(z\), let
\(\square=\square_k(z)\) be its assigned current window and
\(R_z=X'_k(z)\) its fixed enlarged box.
There are only a fixed number of direction-partition windows
near this base window.  Group their directions by linearity
before applying the box-background replacement, and choose
\(R_z\) to contain their fixed enlargements.  In (WF3) their
sum is the inner cutoff \(\zeta_\square\); the remaining
directions are outer terms.  This is part of the fixed
window-first rule.  The derivative at zero is linear in
these directions, so this grouping preserves the partition
identity without changing any incoming pointed member.
Let \(\mathsf A_{j,z}(X,\cdot)\) be the fixed action-family
filing with total \(A(h_z^{(j)},\cdot)\), and put
\(\widetilde G_{j,z}=G_{j,z}-b_{j,z}\mathsf A_{j,z}\).
The one inner window member is the function
\[
 \begin{split}
 {\cal W}_{j,z}(P')=
 \sum_{\substack{X\ni z\\X\subset\square^{\sim2}}}
 \big[&
 \widetilde G_{j,z}\big(X,P^{\rm box}_{j,z}(P'),z\big)\\
 &-\widetilde G_{j,z}\big(X,(1,0),z\big)\big].
 \end{split}\tag{WF13a}
\]
Here \(P'\) is a local potential/current pair of the current
resolved level \(k\), on \(R_z\).
The map \(P^{\rm box}_{j,z}\) first forms its actual nonlinear
averaged datum and then takes the old-level constrained
critical box pair of that datum.
The same box-background map is used for every \(X\) in this
sum.  The own tent action is included in the near sum; its
fixed support lies inside the assigned enlarged window.
The declared domain of this entire member is \(R_z\), and its
fine-cube base is
\(B_z=\operatorname{Hull}_{\pi_k}R_z\).
Its number of current fine cubes is bounded by a fixed
constant \(C_{\rm win}\), independent of \(j,k,z\).
This is a presentation of an inner member before localization;
we do not apply different old-\(X\) bracket operators to its
summands.

Let \({\cal W}_{j,z,Y}(s,t,\nu)\) mean this one function
evaluated by the assigned local response and the datum
(WF3) on \(Y\supset B_z\).
The inner potential contributed by this member is
\[
 V^{\rm in}_{j,z}(Y)=
 \int_0^1
 \left[
   \left.\partial_\nu
       {\cal W}_{j,z,Y}(\,\cdot\,,t,\nu)\right|_{\nu=0}
 \right][\,Y\setminus B_z\mid B_z\,]\,dt .
 \tag{WF13b}
\]
The brackets in this formula are linear in the whole function
(WF13a).  Thus its cancellation and bound (WF8) are established
before this single bracket is estimated.
Different \(z\)'s need not have a common base: there is no
cancellation between them, and their bounds can be added.
Hereditary restriction identifies the bracket on \(Y\) with
the bracket of the same ambient interpolation function whose
parameters outside \(Y\) are zero.  Finite fundamental theorem
of calculus therefore gives, at a complete physical pair,
\[
 \sum_{Y\supset B_z}V^{\rm in}_{j,z}(Y;p_W)
   =\int_0^1
       \left.\partial_\nu{\cal W}_{j,z,\Lambda}
                             (1,t,\nu;p_W)\right|_{\nu=0}\,dt .
\]
Every local term uses only the pair on \(Y\) and the input
members on its declared domain \(B_z\subset Y\).
This proves chart-coherent locality and the physical sum
identity for this fixed window-first rule.
Summing the window identities and the retained outer and
removed origins recovers the old action difference.
No equality of individual output domains with a different
filing is required.
The current-level origins \(j=k\), and the far and removed
origins, retain their own declared bases and the bounds
already proved.  For vacuum-only near origins the same
window-first grouping may be used without a marginal
subtraction.

The assigned maps are hereditary on active wall-components,
and complete at all parameters equal to one.
For an interpolation function \(\Psi\), its bracket is
\[
 \Psi[K|N]=\int_{[0,1]^K}
        \prod_{\Delta\in K}\partial_{s_\Delta}
                            \Psi(s_K+1_N)\,ds_K .
 \tag{WF13}
\]
For (WF13b), \(K=Y\setminus B_z\), \(N=B_z\).
For an origin retaining its own base \(X_0\), use
\(K=Y\setminus X_0\), \(N=X_0\).
A component carrying a differenced cube and missing the base
makes this bracket zero.
Each coordinate has the disc \(R=e^{\bar\kappa}\), so mixed
Cauchy costs \((R-1)^{-|K|}
\le e^{-(\bar\kappa-1)|K|}\), with no factorial.
The common \(\nu\)-disc and its one derivative commute with
these brackets by joint holomorphy.
The preceding amplitude bounds therefore retain this additional
spatial factor.

Here is a quantitative rate allocation.
For a retained old domain \(X\) and its current fine hull
\(X_0\), put \(m=\#(Y\setminus X_0)\).
The actual cube geometry and rooting give a fixed constant
\(C_h\) such that
\[
 d_k(Y)\le C_h\big[d_j(X)+m+1\big].
 \tag{WF13c}
\]
To see this, view a tree covering \(X\) on the old lattice
in current units and cover it by current fine cubes.
Its hull has a covering tree of length at most
\(C(d_j(X)+1)\).  Each component of the additional cubes
meets that hull, or its bracket vanishes.
A rooted spanning forest adds at most a fixed number of
edges per additional cube.  This proves (WF13c).
For an inner window member the same argument and
\(\#B_z\le C_{\rm win}\) give
\(d_k(Y)\le C_h(m+1)\), increasing \(C_h\) by a fixed
factor if necessary.

Split the old input rate into positive parts
\(\kappa_{\rm sum}+\kappa_{\rm far}+\kappa_{\rm ret}\le\kappa\).
Choose \(\kappa_{\rm sum}>\kappa_D\) with enough strict
gap to sum the fixed polynomial diameter factors, and
use \(\kappa_{\rm far}\) for the long-domain tail in
(WF7)--(WF12).
The retained old-domain factor is
\(e^{-\kappa_{\rm ret}d_j(X)}\).
Put \(\beta=\bar\kappa-1\).
Choose an output rate \(\lambda\) and a number \(\epsilon>0\)
with
\[
 C_h(\lambda+\epsilon)
                <\min\{\kappa_{\rm ret},\beta\}.
 \tag{WF13d}
\]
All rates are chosen once, sufficiently large for the
subsequent ordinary polymer criterion.
Combining the retained old factor with \(e^{-\beta m}\)
and (WF13c) gives
\(C e^{-(\lambda+\epsilon)d_k(Y)}\).
For inner members the bounded base needs only the
\(\beta\) part of this inequality.
The number of possible current-window anchors at fixed
\(Y\) is polynomial in \(\#Y\); the old bases within one
such window have already been counted by
\(C_M\sigma_j^{-4}\).
The bound \(\#Y\le C(d_k(Y)+1)\) and the remaining
\(\epsilon\) absorb this polynomial.
This is a pointwise output-domain estimate; any later
sum over output domains spends its separately declared
entropy reserve.

The function
\(\rho_j^3\sigma_j^{\alpha_H}\) is bounded by
\[
 \omega_{k+1-j}
       :=C(1+k-j)^{3/2}L^{-\alpha_H(k-j)},
 \qquad \sum_{n\ge1}\omega_n<\infty .
 \tag{WF14}
\]
All age functions have the same enlarged fine-parameter domains
and the same \(\nu\)-disc.
At fixed \(j,z\), (WF13a)--(WF13b) provide the common-root
function before any bracket.  The resulting bounds for distinct
base points and levels are then added absolutely.
No bound for a sum of differently based bracket operators has
been inferred from a sup norm of their unbracketed sum.
The uniform bounds justify passage through the bracket
integrals and to limits of the old-level sum.

We have proved the following formation estimate.  Let \(g_j\)
be regular-member norms, \(d_j\) defect norms and \(r_j\)
aged-remainder norms.  Under the explicit background, window
and geometry inputs above, the actual old-origin potential
satisfies
\[
 \|\mathbf V_k^{\rm old}\|_\lambda
 \le Cs_k\left[
       \sum_{j\le k}\omega_{k+1-j}g_j
       +\sum_{j\le k}\rho_j^2d_j
       +\sum_{j\in{\cal J}_k}\rho_j^2r_j\right].
 \tag{WF15}\label{wf:formation-bound}
\]
These are local bounds on every specified complex core.
At complete physical pairs, finite fundamental theorem of
calculus, the window partition, and hereditary resummation of
(WF13) identify their sum with the actual old action difference.
There is no assertion of a single global off-graph Gaussian
or a global torus background for an arbitrary local pair.

Adding a scalar Wilson-weight direction contributes at most
\(CM^4s_k\alpha_k^2\|u\|_\infty\), by the separate
determinant/response-bracket argument.
The formation rules are linear in their origins, so (WF15)
applies to their first source derivatives.
For a frozen regular member that direction is
\(f(p_z){\cal L}_j[1]\) with its own \(b_j'(0)\);
the remaining direction is the source defect.
The finite potential-ball Taylor remainder justifies the
chain rule even when the actual ordinary source families are
nonlinear.  This gives the formation inequality used in the
linear source bootstrap, with no additional unexplained
source-response norm assumption.

Assume the joint complex window/core room, nonlinear averaged-datum
estimates, H\"older response and compatible boundary comparison,
and the outer/removed-field locality of the assigned maps.
These are the analytic inputs used in the formation proof.
The large-field history operation uses its separate history-induction
hypotheses.

% END MODULE chapters/window_potential_formation_v7.tex


% BEGIN MODULE chapters/linear_source_bootstrap_v7.tex
\section{A first-source-response bootstrap on the ordinary domains}
\label{ch:linear-source-bootstrap}

This section isolates a linear consequence of the ordinary local formation
bounds.  It uses the inserted-potential and renewed-pointing estimates at
one step, rather than a larger domain of physical sources.  In particular,
it distinguishes the derivative of a newly generated activity from the
old background action that is carried to the next scale.
The formation and domain hypotheses below concern the actual local
maps.  Their statement is not an assertion that a nonlinear history
construction satisfies them.

\subsection*{Meshes, norms, and the fixed ordinary run}

Work first on a finite periodic lattice.  Let \(h_j=L^jh_0\) be the
physical mesh of level \(j\), with \(L>1\).  A source profile is a
bounded Lipschitz function on the physical torus; complex profiles are
allowed by complexification.  Put
\[
 \|f\|_{\sharp,j}=\max\{\|f\|_\infty,C_Mh_j\operatorname{Lip}f\}.
 \tag{LB1}
\]
At level \(j\), a pointed localization domain \(X\) contains its base
point \(z\).  Its enlarged source region \(X^\sim\) and the fixed
plaquette point \(p_z\) satisfy
\[
 \sup_{p\in X^\sim}|p-p_z|_{\rm tor}
       \le C_Mh_j(1+d_j(X)).
 \tag{LB2}
\]
Distances and diameters in this section use the torus metric.
The family norm with rate \(\gamma\) is the supremum of
\(e^{\gamma d_j(X)}|F(X,P,z)|\) over the actual specified complex
field domain and all \(X,z\).  Companion families, when present,
are included in the same norm.  A bound with rate \(\kappa+1\)
therefore leaves one unit of rate for a diameter.

The ordinary run fixes all Gaussian structures, background maps,
thresholds, label spaces and reference operations.  Sources change
local exponent members only.  Let
\[
 \alpha_j=\theta_j^{-1/2}(\log\theta_j)^q,\qquad
 \delta_j=\theta_j^{-1/2}(\log\theta_j)^{p_0},\qquad
 s_j=\delta_j/\alpha_j,\qquad q>p_0\ge0.
 \tag{LB3}
\]
Fixed multiplicative constants in these definitions can be included.
Assume the ordinary flow bounds, with constants independent of the
number of levels,
\[
 c(k-j)-C_0\le\theta_j-\theta_k
                       \le C(k-j)+C_0\qquad(j\le k).
 \tag{LB4}
\]
All \(\theta_j\) are sufficiently large.  These bounds imply, with a
uniform constant \(C_F\), that for \(1\le j\le k\),
\[
 \begin{aligned}
 \alpha_{j-1}^2s_{j-1}&\le C_F\alpha_k^2s_k,\\
 \varrho_{k,j}^2\alpha_{j-1}^2s_{j-1}
                         &\le C_F\alpha_k^2s_k,\qquad
 \varrho_{k,j}=\alpha_k/\alpha_j,\\
 \varrho_{k,j}^2&\le C_F(1+k-j).
 \end{aligned}\tag{LB5}
\]
Indeed the functions \(t^{-1}(\log t)^{2q}\) and
\((\log t)^{p_0-q}\) are decreasing for sufficiently large \(t\).
The bounded displacement \(C_0\) in (LB4) changes their ratios by
at most a constant.  The second line also uses the bounded
consecutive ratio \(\alpha_{j-1}/\alpha_j\).  The upper flow bound
and \(\theta_k\ge1\) prove the last line.  These observations also
give \(s_j\le C_Fs_k\) when \(j\le k\).

\subsection*{Locality supplies the source-gradient gain}

Let \(E_j^w\) be an ordinary sourced regular family, holomorphic in a
neighborhood of \(w=0\).  Define its first derivative by
\[
 {\cal L}_j[f]=\left.\frac{d}{dt}E_j^{tf}\right|_{t=0},
 \qquad
 A_j=\sup_{f\ne0}
        \frac{\|{\cal L}_j[f]\|_{\kappa+1}}{\|f\|_{\sharp,j}}.
 \tag{LB6}
\]
The supremum is taken over the permitted profiles and the tori in
the run.  At a fixed finite lattice all differentiations are in
finite-dimensional source coordinates.  Locality means that the
value indexed by \(X\) depends only on the values of \(f\) on
\(X^\sim\).  To apply the following argument it suffices that the
source class admit bounded Lipschitz truncations of a profile and
constants, as does the class defined by (LB1).  No density claim
about an infinite-dimensional algebraic dual is used.

The derivative defect at a base point is exactly
\[
 {\cal D}_j[f](X,P,z)
       ={\cal L}_j[f](X,P,z)
                   -f(p_z){\cal L}_j[1](X,P,z).
 \tag{LB7}
\]
We claim
\[
 \|{\cal D}_j[f]\|_\kappa
                  \le C_M A_jh_j\operatorname{Lip}f.
 \tag{LB8}
\]
To prove it, fix \(X,z\), and subtract \(f(p_z)\) from \(f\).
On \(X^\sim\) the resulting profile has modulus at most
\(C_Mh_j(1+d_j(X))\operatorname{Lip}f\).
For real \(f\), clip its values outside an interval twice this
large, using a Lipschitz scalar map that is the identity on that
interval and has Lipschitz constant at most a fixed constant.
For a complex profile clip the real and imaginary parts separately.
The resulting global profile \(g\) agrees with \(f-f(p_z)\) on
\(X^\sim\) and satisfies
\(\|g\|_{\sharp,j}\le
C_Mh_j(1+d_j(X))\operatorname{Lip}f\).
Locality and linearity give
\({\cal D}_j[f](X,P,z)={\cal L}_j[g](X,P,z)\).
Apply (LB6) and
\((1+d)e^{-d}\le1\) for \(d\ge0\) to obtain (LB8).
Smooth profiles may instead use a smooth clipping map with the same
fixed bounds.  This proof applies uniformly at every external
complex field point in the family norm.
For a source measured at level \(m\ge j\), it gives in particular
\[
 \|{\cal D}_j[f]\|_\kappa
       \le C A_jL^{-(m-j)}\|f\|_{\sharp,m}.
 \tag{LB9}
\]
Thus the derivative defect retains the response amplitude \(A_j\)
and the scale-gradient gain together.

\subsection*{The precise one-step inputs}

For clarity we list the ordinary inputs needed for the bootstrap.
The first two are estimates for local formation and the finite
regular gas, at fixed background maps; they are not estimates for
complex history masses.

The formation of potentials is linear in each origin.  When a
source is differentiated at zero, the regular-origin direction
at a site is \(f(p_z){\cal L}_j[1]\), accompanied by its own
coefficient, and the remaining direction is (LB7).
We require the following ordinary formation estimate on the
actual localized complex pairs:
\[
 \begin{split}
 \|\dot{\mathbf V}_k[f]\|_\lambda
 \le C s_k\bigg[
   &\alpha_k^2\|\dot u_k[f]\|_\infty
       +\sum_{j\le k}\omega_{k+1-j}
                          A_j\|f\|_\infty\\
   &+\sum_{j\le k}\varrho_{k,j}^2
                            \|{\cal D}_j[f]\|_\kappa
       +\sum_{j\in{\cal J}_k}\varrho_{k,j}^2
                            \|\dot R_j[f]\|_\kappa\bigg].
 \end{split}\tag{LB10}
\]
Here \(\omega_n\le C(1+n)^CL^{-\alpha_H n}\) with
\(\alpha_H>0\), and \({\cal J}_k\) contains the aged remainder
origins that actually enter the regular part.  Younger remainders
remain in the remainder class.  The factor \(s_k\) is the
fluctuation-amplitude/field-radius ratio.  In a formulation with
additional positive powers of \(\theta_k^{-1}\) in this factor,
they are bounded by a constant times \(s_k\) at weak coupling.
The current scalar-origin term in (LB10) follows more sharply
from its fluctuation-difference estimate
\(C\alpha_k\delta_k\|\dot u_k[f]\|_\infty\).
All substitutions in (LB10) must lie in their actual specified
cores; a norm inequality does not create those inclusions.
In particular, a proof of the pointed Hessian cancellation and
nonlinear window estimate on constrained critical backgrounds
alone does not prove (LB10) for independent potential/current pairs
used in localization.  Passing from those background estimates
to all the actual localized pairs remains part of this ordinary
formation input.

The inserted-potential estimate for the actual finite
Mayer/decoupling series, followed by its actual renewed bond
pointing, gives
\[
 \|{\cal L}_{k+1}[f]\|_{\kappa+1}
                  \le C_P\|\dot{\mathbf V}_k[f]\|_\lambda.
 \tag{LB11}
\]
Its proof uses the ordinary spare-potential slot.  The auxiliary
potential parameter is not a physical source in a history box.
The normalized Gaussian, all reference measures, and
\(\mathbf N^\circ,\mathbf G_\tau\) have zero source derivative.
These facts are essential to (LB10)--(LB11).

The ordinary constant-source construction supplies numbers \(b_j'(0)\)
with \(\sum_j|b_j'(0)|\le B_0\), uniformly in the cutoff.
Its fixed tents and stage factors satisfy
\[
 \dot u_k[f](p)=f(p)+\sum_{j\le k}b_j'(0)\varphi_j(p)
                    \sum_yh_y^{(j)}(p)f(p_y),\qquad
 \|\dot u_k[f]\|_\infty\le(1+B_0)\|f\|_\infty.
 \tag{LB12}
\]
The last bound uses \(0\le\varphi_j\le1\), \(h_y^{(j)}\ge0\),
and \(\sum_yh_y^{(j)}=1\).
The coefficient of \({\cal L}_j[1]\) is exactly \(b_j'(0)\):
the free response and coefficient extractor are fixed and
complex linear.  Thus no coefficient of another member enters
(LB10), even when \(f(p_z)\) varies with the site.

Finally suppose the ordinary remainder families are bounded by
\(C_R\theta_j^{-p}e^{-\kappa d_j(X)}\), with \(p>3\), on the fixed
physical source disc \(|t|<r/\|f\|_{\sharp,j}\) around zero.
Cauchy's inequality at half that radius proves
\[
 \|\dot R_j[f]\|_\kappa
       \le (2C_R/r)\theta_j^{-p}\|f\|_{\sharp,j}.
 \tag{LB13}
\]
This is the only source-disc estimate needed for the small
remainder class.  It is not used to produce an \(\alpha_j^2\)
factor for regular activities.

\subsection*{The induction and its constants}

\begin{proposition}[Linear regular-response bootstrap]
\label{prop:linear-source-bootstrap}
Under (LB1)--(LB13), suppose the initial regular list is empty.
There are \(M_0<\infty\) and \(s_\star>0\), independent of the
number of levels, cutoff and volume, such that \(s_j\le s_\star\)
for all \(j\) implies
\[
 A_j\le M_0\alpha_{j-1}^2s_{j-1}\qquad(j\ge1).
 \tag{LB14}
\]
\end{proposition}

\begin{proof}
Divide (LB10) by \(\|f\|_{\sharp,k+1}\).
Use (LB9) with \(m=k+1\), (LB11)--(LB13), and the monotonicity
\(\|f\|_{\sharp,j}\le\|f\|_{\sharp,k+1}\).  This gives
\[
 A_{k+1}\le C_1s_k\left[
 (1+B_0)\alpha_k^2+
 \sum_{j\le k}\omega_{k+1-j}A_j+
 C\sum_{j\le k}\varrho_{k,j}^2A_jL^{-(k+1-j)}
 +C_r\sum_{j\le k}\varrho_{k,j}^2\theta_j^{-p}\right].
 \tag{LB15}
\]
There is no sum over the number of source labels in this inequality.
The norms of linear directions have already combined them into \(f\).
The flow bound and (LB5) imply
\[
 \sum_{j\le k}\varrho_{k,j}^2\theta_j^{-p}
       \le C\sum_{a\ge0}(1+a)(\theta_k+ca)^{-p}
       \le C_p\theta_k^{2-p}\le C_p\alpha_k^2 .
 \tag{LB16}
\]
To check the middle bound, compare the sum with its first term and
the integrals of \((\theta+ct)^{-p}\) and
\(t(\theta+ct)^{-p}\) on \([0,\infty)\).
They are fixed constants times
\(\theta^{1-p}\) and \(\theta^{2-p}\), respectively; the latter
integral is finite exactly when \(p>2\).
The final bound uses \(p>3\), \(q\ge0\), and large \(\theta_k\).

Assume (LB14) through level \(k\).  Both sums involving \(A_j\)
in (LB15) are at most
\[
 C M_0\alpha_k^2s_k
       \left(\sum_{n\ge1}\omega_n+\sum_{n\ge1}L^{-n}\right)
\]
by (LB5).  Consequently, for constants \(C_2,C_3\) independent
of \(k\),
\[
 A_{k+1}\le\alpha_k^2s_k(C_2+C_3M_0s_k).
\]
Choose \(M_0\ge2C_2\) and
\(s_\star\le(2C_3)^{-1}\).  The right-hand side is at most
\(M_0\alpha_k^2s_k\), as required.  At the first step there are
no old regular members, so the same argument starts the induction.
The direct bottom-level scalar potential, if it occupies a separate
initial slot, must satisfy the same fixed-reference fluctuation
estimate used in (LB10).
\end{proof}

\subsection*{A consequence on the near, non-wrapping window class}

Let the final fluctuation size be \(\delta_n=\alpha_ns_n\), and
consider only terms for which the actual local window chart, the
pointed affine comparison and its complex domain inclusions hold
without a field-independent wrapping or holonomy tail.
In addition, the axial and H\"older bounds of the nonlinear window
proof must hold for the \emph{complete} near background-plus-fluctuation
pair with output size \(a_{\rm out}\le C\delta_n\).
A bound on its fluctuation increment alone is insufficient:
an unrestricted background of size \(\alpha_n\) would require that
larger size in the radius ratio.
The intended near class is the one on which the actual small-field
tests impose a fixed multiple of \(\delta_n\) on the complete field.
Keep the true ratio \(\rho_j=\delta_n/\alpha_j\); it need not
exceed one.  The global half-datum-disc condition in the near-domain
Taylor proof holds uniformly at all ages.  The actual window-domain
input includes \(\alpha_j/r_{B,j}\le C\), where \(r_{B,j}\)
is the independent-coordinate datum radius.  Hence, with
\(\sigma=L^{-(n-j)}\), the flow gives
\[
 \frac{C_{\rm ax}\delta_n\sigma}{r_{B,j}}
 \le C\rho_j\sigma
 \le C s_n(1+n-j)^{1/2}L^{-(n-j)}.
\]
Choose \(s_\star\) smaller by a fixed factor so that the right-hand
side stays below the required half-disc constant.
This uses the near-domain proof for \(\rho_j>0\) itself and avoids
the young-age core-only fallback, which would lose a \(\rho_j^2\)
gain when \(\rho_j<1\).
After the four-dimensional fine-site count, the
renormalized constant-member contribution at age \(a=n-j\)
is bounded by
\[
 C A_j\|f\|_\infty
       \left[\left(\frac{\delta_n}{\alpha_j}\right)^2
                              L^{-\alpha_Ha}
            +\left(\frac{\delta_n}{\alpha_j}\right)^3
                              L^{-2a}\right].
 \tag{LB17}
\]
The value at the identity and this member's own
\(b_j'(0)A(h_z,\cdot)\) are removed here.
By (LB14)--(LB5),
\(A_j/\alpha_j^2\le Cs_n\) and
\(\delta_n/\alpha_j\le Cs_n(1+a)^{1/2}\).
The sum of (LB17) is therefore at most
\(C\delta_n^2s_n\|f\|_\infty\).

For the defect, the vacuum-only window estimate instead gives
\[
 C\|{\cal D}_j[f]\|_\kappa
                 \left(\frac{\delta_n}{\alpha_j}\right)^2
 \le C\delta_n^2s_nh_j\operatorname{Lip}f.
 \tag{LB18}
\]
Since \(h_j=h_nL^{-(n-j)}\), its sum is at most
\(C\delta_n^2s_n\|f\|_{\sharp,n}\).
Thus these near regular and source-gradient terms together have
size \(C\delta_n^2s_n\|f\|_{\sharp,n}\) per final window unit.
With the radii (LB3), the prefactor is
\(\theta_n^{-1}(\log\theta_n)^{3p_0-q}\).
The gain is obtained because the newly generated derivative member
already has amplitude \(\alpha_j^2s_{j-1}\).

No field-independent far tail is included in (LB17)--(LB18).
Such a tail cannot acquire a factor \(\rho_j^2\) merely from a
source derivative.  Recent large-field terms, marked remainder
windows and torus or hole windings still require their own
ordinary source-disc bounds, domain comparisons or separately
small-class estimates.  Nor do the local implications in this
section verify the nonlinear formation hypotheses for every
history.  They show exactly how the \(\alpha^2\) normalization and
the source-gradient gain propagate once those ordinary local
hypotheses and the finite pointed estimate hold.

% END MODULE chapters/linear_source_bootstrap_v7.tex


% BEGIN MODULE chapters/old_activity_hessian_v7.tex
\section{The Hessian of old activities in a run window}
\label{ch:old-activity-hessian}

This section proves a conditional ordinary estimate.  Its inputs are
the complex domains and first-response estimates of the actual background
maps, the pointed quadratic cancellation, and the window formation rule.
The desired Hessian estimate is not one of the inputs.  In particular, an
estimate on the Gaussian cutoff box alone is not used as an analytic
continuation theorem.  We construct the larger fluctuation domain from
the unscaled constraint chart and the datum domains of the background
maps.

The conclusion concerns the regular old-origin part of a step exponent.
The conditional exterior-force estimate, boundary origins with stored
fields, and the identification of a clean-class conditional measure are
separate issues.

\subsection*{Physical and Gaussian coordinates}

Fix a step from level \(k\) to level \(k+1\).  Write
\[
 g_k=\theta_k^{-1/2},\qquad
 \alpha_i\asymp\theta_i^{-1/2}(\log\theta_i)^q,
 \qquad \mathfrak b_k\asymp(\log\theta_k)^p,
 \qquad \delta_k=g_k\mathfrak b_k .
 \tag{OA1}
\]
Here \(q>p\), and \(\alpha_i\) is the smaller of the usable physical
radii in the datum, potential, curvature and current conditions.
Their fixed ratios are included in the constants.  A finite alteration
of an index in \(\alpha_i\), as needed for adjacent input layers, has the
same meaning.  We must prove that the ordinary flow gives
\[
 \rho_i:=\max\{1,C\alpha_k/\alpha_i\}
       \le C_\rho(1+k-i)^{1/2},\qquad
 \sigma_i=L^{-(k-i)},\qquad i\le k .
 \tag{OA2}
\]
Constants depend on the fixed block and collar choices, but not on the
number of old levels, the cutoff or the volume.  In particular,
\(C_{\rm a}:=\sup_{i\le k}\rho_i\sigma_i<\infty\).

The free Gaussian variable \(y_b\) has three real color coordinates at
each nonpivot bond of the level-\(k\) lattice.  The physical bond
logarithm is of size \(g_k y_b\).  Thus a physical analytic radius
\(\alpha_k\) corresponds to a Gaussian radius of order
\(\alpha_k/g_k\), rather than to \(\alpha_k\).
For complex variables we use the Euclidean norm of the three complex
components at a bond and its supremum over bonds.

Let \(I\) be the finite box of free bonds under consideration.  Its
outside free coordinates and other exterior data are fixed.  Every
real free coordinate appearing in the regular fields is bounded by
\(C_{\rm w}\mathfrak b_k\).  The number of coordinates satisfies
\[
 m_I=3\#I\le C(rL^{-a_*}+D)^4\le C(r+D)^4.
 \tag{OA3}
\]
This count will not enter the localized estimate below.  It is relevant
to the weaker estimate available without spatial locality.

\subsection*{The primitive analytic inputs}

We specify the domain inputs used in the proof.  These are statements
about assigned maps and their first responses, not bounds on the
Hessian of an old activity.

For the actual one-step average, at its actual background
\(V^{\rm bg}(W)\), put
\[
 q_W(a)=\frac1i\log\!\left[
       \mathcal M(e^{ia}V^{\rm bg}(W))W^{-1}\right]
        =LQ_Wa+\mathcal R_W(a).
 \tag{OA4}
\]
The logarithm is taken on its assigned local branch.  The fixed-pivot
construction supplies linear maps \(\mathsf C_W\) and \(h_W\) with
\[
 LQ_W\mathsf C_W=0,\qquad LQ_Wh_W=1,
 \qquad\|\mathsf C_W\|_{\infty\to\infty}
              +\|h_W\|_{\infty\to\infty}\le C_L .
 \tag{OA5}
\]
The map \(\mathsf C_W\) is the pivot elimination map, not a Gaussian
covariance.  The maps in (OA4)--(OA5) are the ones used in the run.
Uniformly on an unscaled physical ball \(\|a\|_\infty<r_0\), the
remainder is holomorphic and obeys
\[
 \|\mathcal R_W(a)\|_\infty\le C_0\|a\|_\infty^2,
 \qquad
 \|D\mathcal R_W(a)\|_{\infty\to\infty}
                                      \le C_0\|a\|_\infty .
 \tag{OA6}
\]
Here \(r_0,C_0,C_L\) are independent of the box size.  The same bounds,
with fixed enlarged constants, apply to the assigned local charts on
the independent fine-parameter discs \(|s_\Delta|<e^{\bar\kappa}\).
At physical parameters the local charts agree with the run chart by
their restriction and completeness identities.

The ordinary background package supplies holomorphic maps on their
declared datum balls and their declared core slices.  The datum radii
satisfy \(r_{B,i}\ge c\alpha_i\).  At zero fluctuation all substituted
data and pairs have a fixed positive fraction \(\epsilon_0\) of room
in these domains.  This room includes the tower conditions, node
representatives, current bounds, and the inner, outer and removed
substitutions.  These are local complex assertions; they do not assert
the existence of a global Gaussian measure for independent currents.

We use the following precise first-response form of that package.
Measure each defining datum or core displacement at old level \(i\)
in units of its allowed radius.  Along an admitted physical perturbation
of supremum size \(u\), its derivative in any of these normalized
quantities is bounded by
\[
 C_R\sigma_i\alpha_i^{-1}u.
 \tag{OA7}
\]
The bound holds throughout the declared complex domains, for all the
fine parameters above.  It is the integrated first-response and
nonlinear-average estimate: the actual average, not its differential,
is used to construct the datum.  The constant \(C_R\) includes the
fixed enlarged parameter discs and the chart derivative.  For outer
and removed derivatives there are additionally their assigned
exponential distance factors.  We must also prove a compatible boundary
comparison and the weighted H\"older response estimates for the actual
box backgrounds.

Finally, the field maps are hereditary on active wall components.
After the parameters outside a domain \(Y\) are set to zero, a local
potential uses only the free coordinates in \(Y\).  A fixed collar can
instead be included in \(Y\) throughout; it changes the constants
below.  The number of domains of tree length in \([m,m+1)\) containing
a fixed cube is at most \(C_De^{\kappa_Dm}\), and
\[
 \operatorname{diam}Y\le C_M(1+d_k(Y)).
 \tag{OA8}
\]
The spatial rates have the strict reserves needed to sum these
domains and the fixed polynomial diameter factors.

\subsection*{The actual chart on the larger ball}

For clarity we give the contraction argument behind the radius claim.
Put \(z=g_k\mathsf C_Wy\).  The nonlinear correction is the solution
of
\[
 d=\mathcal R_W(z-h_Wd),\qquad
 a_W(y)=z-h_Wd,
 \qquad V_W(y)=e^{ia_W(y)}V^{\rm bg}(W).
 \tag{OA9}
\]
Choose \(r_1>0\) so small that
\(2r_1<r_0\) and \(4C_0C_Lr_1<1/2\).
For \(\|z\|_\infty\le r_1\), the right side maps the ball
\(\|d\|_\infty\le4C_0\|z\|_\infty^2\) to itself.
Indeed \(\|z-h_Wd\|_\infty\le2\|z\|_\infty\), and (OA6)
then gives that bound.  Its Lipschitz constant in \(d\) is at most
\(2C_0C_Lr_1<1/2\).  Iteration from zero therefore converges uniformly
to the unique solution, including locally uniformly in all complex
parameters.  Uniform limits of the holomorphic iterates are
holomorphic.  Differentiating the fixed-point equation, or its
uniformly convergent derivative iteration on a smaller ball, gives
\[
 \|a_W(y)\|_\infty\le Cg_k\|y\|_\infty,
 \qquad \|D_ya_W(y)\|_{\infty\to\infty}\le Cg_k .
 \tag{OA10}
\]
Moreover, \(q_W(a_W(y))=0\), because \(LQ_Wz=0\) and
\(LQ_Wh_Wd=d\).  This is the actual constraint, with its nonlinear
average.  Uniqueness shows that this chart agrees on the operational
cutoff box with the chart already used in the run.

There is consequently a fixed \(c_*>0\) for which all the old
substitutions are defined and holomorphic on
\[
 \|y\|_\infty<R_k,
 \qquad R_k=c_*\alpha_k/g_k .
 \tag{OA11}
\]
Here and below fixed exterior coordinates may be retained with norm
at most \(C_{\rm w}\mathfrak b_k\).  To prove the assertion, first make
\(\alpha_k\) small enough that (OA10) stays in its physical chart.
For any fixed complex \(y\) in (OA11), follow the segment \(ty\),
\(0\le t\le1\).  On the portion where a composition of assigned
background maps is admitted, (OA7) and (OA10) bound each normalized
displacement by
\[
 C_R g_k\|y\|_\infty\sigma_i/\alpha_i
       \le C_R c_*\rho_i\sigma_i
       \le C_R c_*C_{\rm a}.
 \tag{OA12}
\]
Choose \(c_*\) so that the last expression is less than
\(\epsilon_0/4\).  The segment cannot reach a domain boundary:
every defining inequality still has at least half of its initial
room.  Atomic background maps are already given on those domains, so
this is an inclusion argument for a composition of assigned
holomorphic maps, not an assertion that a branch continues through
an unknown singularity.  The same argument gives the required joint
holomorphy with the fine parameters.  It also applies to local
independent-current formulas without postulating an off-graph global
Gaussian integral.

The window multiplier used to differentiate an inner direction can
range on a fixed disc \(|\nu|<c_1/c_*\), where \(c_1\) is chosen from
the same room estimates and \(c_*<c_1/4\).  In fact its normalized
increment is bounded by \(C_R|\nu|c_*\rho_i\sigma_i\).
The unamplified part costs only \(C_Rc_*C_{\rm a}\); choose \(c_1\)
so that their sum still spends less than half the room.  Thus this
disc contains \(\nu=1\) and permits the window derivative at zero.

Since \(q>p\), increasing the coupling floor gives
\(C_{\rm w}\mathfrak b_k<R_k/4\).  Two coordinate Cauchy circles
with radii \(cR_k\), centered anywhere on the real product window,
therefore fit in (OA11).  There is no factor depending on \(m_I\)
in this inclusion.  No indicator function or cutoff is analytically
continued.  Only the activity functions are evaluated on the larger
ball; the real measure, Gaussian scaling and all numerical choices
of the run remain fixed.

\subsection*{An amplitude bound before the Gaussian-growth norm}

Let the regular member at old level \(i\) have norm \(e_i\), with
\(e_i\le E_0\).  The norms include its infinite companion if one is
required.  Each pointed total has local gauge invariance, constant
color invariance and the full site symmetry needed for the pointed
Hessian theorem.  At every old base point subtract its vacuum and
its own extracted coefficient times its own tent action.  The tent
has unit quadrature; the action convention is
\(A(e^{iB})=\tfrac12|dB|^2+O(B^3)\).
This is the regular part of the actual ledger decomposition: its
counterterms have been transferred to the running Wilson coefficient.
A profile deficit is retained as a lag or boundary origin, rather
than silently subtracted from the measure.  The clean interior has
no such deficit in the canonical redrawn-profile ledger.
Thus the affine quadratic term cancels with coefficient one in that
normalization.  We do not replace the pointed object by a bulk
convolution or assign an action coefficient to an individual domain.

We recall explicitly the value estimate supplied by this cancellation.
Write \(D_X=2+\operatorname{diam}_iX\).  On an inner box background
the affine part and its error have sizes bounded, with fixed powers
of \(D_X\), by
\[
 \alpha_k\sigma_i^2,
 \qquad
 \alpha_k\sigma_i^{2+\alpha_H}
                +r_{B,i}^{-1}\alpha_k^2\sigma_i^4,
 \qquad 0<\alpha_H\le1.
 \tag{OA13}
\]
These estimates concern the complete substituted field, not just its
fluctuation increment.  They hold on the whole disc just constructed.
Second and third derivatives of a member cost
\(Ce_i\alpha_i^{-2}\) and \(Ce_i\alpha_i^{-3}\), with its
original tree decay.  Its linear derivative at the trivial pair is
zero by constant color invariance.

Set \(T=C D_X^s\rho_i\sigma_i^2\).  When \(T\) is small, Taylor's
formula, the cancellation of the affine quadratic value on the
pointed total, and \(r_{B,i}\ge c\alpha_i\) bound the remaining
terms by
\(Ce_iD_X^{s'}[\rho_i^2\sigma_i^{4+\alpha_H}
 +\rho_i^3\sigma_i^6]\).
For example, the product of the affine part and its error, divided
by \(\alpha_i^2\), gives exactly these two powers.  Squares of the
error and cubic terms are bounded by the same expression, since
their normalized sizes are bounded by \(C(T\sigma_i^{\alpha_H}+T^2)\)
and \(CT^3\).  If \(T\) is not small, whole-core bounds for the
member and its own action, and the comparison affine quadratic
value, give \(Ce_i(1+T^2)\le Ce_i C_TT^3\).  The fixed powers
of \(D_X\) are summable against the tree decay.  The finite set of
young ages for which the global half-datum-ball condition fails is
also covered by the whole-core bound: their \(\sigma_i\) have a
fixed positive lower bound.  This is precisely the reason to use
the ordinary full-radius ratio \(\rho_i\ge1\) here.

At one base point, after its domain sum, the bound is therefore
\[
 Ce_i\left[\rho_i^2\sigma_i^{4+\alpha_H}
                         +\rho_i^3\sigma_i^6\right].
 \tag{OA14}
\]
A current window contains at most \(C\sigma_i^{-4}\) old base
points.  Its value bound is at most
\(Ce_i\rho_i^3\sigma_i^{\alpha_H}\).
Terms extending beyond that window retain an exponentially small
tree-length tail.  Outer and removed terms retain their actual
response factors.  Even the cruder amplitude
\(Ce_i\sigma_i^{-4}e^{-c/\sigma_i}\), when only a long-domain
value bound is used, is bounded by
\(Ce_i\rho_i^3\sigma_i^{\alpha_H}\).
This does not give such a tail a square of the actual small
curvature; it is a full analytic-radius estimate.

It is important where the inner estimate is localized.  At each
old base point form one window function from the near pointed sum,
its own action counterterm, and one common critical box map.  This
entire function has one fixed current-scale base, of uniformly
bounded size.  Only then apply the fine brackets of the window
formation rule.  Thus the cancellation in (OA14) occurs before a
single common bracket.  No sup norm of a sum is used to bound a
sum of differently based bracket operators.

On the independent fine discs the bracket with \(m\) differenced
cubes costs at most
\((e^{\bar\kappa}-1)^{-m}\), with no factorial.  The derivative
at \(\nu=0\) costs the reciprocal of its fixed disc, hence only
a fixed constant.  For outer origins retaining their original
base, split the old input rate into the domain-sum, long-tail
and retained parts.  If \(Y\) is the resulting current domain,
the hull and rooting geometry gives
\(d_k(Y)\le C_h(d_i(X)+m+1)\); for a bounded window base it gives
\(d_k(Y)\le C_h(m+1)\).
Choose the retained rates and \(\bar\kappa-1\) larger than
\(C_h(\lambda+\varepsilon)\), after spending the old rooted
entropy and diameter powers.  The remaining \(\varepsilon\)
absorbs the polynomial number of window anchors at a fixed \(Y\).
This proves the following pointwise analytic bound for the
actual chosen window filing:
\[
 |V_{k,i}^{\rm old}(Y;y)|
       \le Ce_i\omega_{k+1-i}e^{-\lambda d_k(Y)},
 \qquad \|y\|_\infty<R_k,
 \quad
 \omega_{a+1}=C(1+a)^{3/2}L^{-\alpha_Ha}.
 \tag{OA15}
\]
The current origin \(i=k\) uses its direct whole-core Cauchy
bound and \(\omega_1=O(1)\); it needs no age improvement.
The proof was carried out on the larger ball using a physical
fluctuation amplitude at most \(c_*\alpha_k\), rather than the
tested amplitude \(\delta_k\).  It does not use the final
potential bound with a Gaussian-growth factor
\(1+M^{-4}\sum_{b\subset Y}|y_b|^2\).  Using that weaker
envelope on a radius \(R_k\) would introduce \(R_k^2\) and
erase the desired Cauchy gain.

Let
\[
 \mathcal E_k=\sum_{i\le k}\omega_{k+1-i}e_i
           \le C E_0 .
 \tag{OA16}
\]
Absolute convergence is uniform on every smaller ball in (OA11),
by (OA15).  It justifies summing old levels before taking two
coordinate derivatives.  All estimates can first be made with
finitely many old levels and on a finite torus.  Their constants
do not change in either limit.

\subsection*{Cauchy, rooted counting and the weighted precision gap}

Write
\(P_{k,I}^{\rm old}(y)=\sum_Y\sum_{i\le k}V_{k,i}^{\rm old}(Y;y)\),
with exterior variables fixed.  Additive terms depending only on
the exterior can be omitted.  By two coordinate Cauchy formulas,
including the factor \(2!\) for a repeated coordinate, (OA15) gives
\[
 \left|\partial_{y_b^a}\partial_{y_{b'}^{a'}}
            \sum_{i\le k}V_{k,i}^{\rm old}(Y;y)\right|
       \le C\mathcal E_k R_k^{-2}e^{-\lambda d_k(Y)}
 \tag{OA17}
\]
throughout the real product window.  The derivative is zero unless
both bonds belong to the declared domain, including its fixed
collar.  For such a domain,
\(d(b,b')\le C_M(1+d_k(Y))\).
Choose \(\mu>0\) so small that
\(\lambda-\mu C_M>\kappa_D\), retaining a strict gap.  Root the
domains at a cube containing \(b\); the fixed number of possible
root cubes costs only a constant.  Then
\[
 \sum_{Y\ni b,b'}e^{-\lambda d_k(Y)}
 \le C e^{-\mu d(b,b')}
       \sum_{m\ge0}e^{-(\lambda-\mu C_M-\kappa_D)m}
 \le C'e^{-\mu d(b,b')}.
 \tag{OA18}
\]
Consequently the full old-activity Hessian, in the actual Gaussian
coordinates, satisfies
\[
 \big\|D_bD_{b'}P_{k,I}^{\rm old}(y)\big\|
 \le C\mathcal E_k\left(\frac{g_k}{\alpha_k}\right)^2
                   e^{-\mu d(b,b')}
 \le C E_0(\log\theta_k)^{-2q}e^{-\mu d(b,b')} .
 \tag{OA19}\label{oa:localized-hessian}
\]
Color matrix norms cost a fixed factor only.  The count of bonds at
distance in \([r,r+1)\) is at most \(C(1+r)^3\).  Thus for every
\(0<\mu'<\mu\), both weighted row and column sums of the Hessian
are at most \(C\mathcal E_k(g_k/\alpha_k)^2\).  In particular,
Schur's inequality gives
\[
 \|T_\phi D^2P_{k,I}^{\rm old}(y)T_\phi^{-1}\|_{2\to2}
       \le C\mathcal E_k(g_k/\alpha_k)^2,
 \qquad
 |\phi(b)-\phi(b')|\le\mu' d(b,b'),
 \tag{OA20}
\]
where \(T_\phi\) multiplies the three coordinates at \(b\) by
\(e^{\phi(b)}\).  This estimate is uniform in \(I\).

For completeness, suppose the actual free precision \(A\), after
the prescribed constraint and gauge fixing, has a gap
\(A\ge a_-\mathbf1\) and uniform exponentially summable entries.
This is a separate free-background input, concerning the constrained
pullback Hessian in the coordinates \(\mathsf C_W\), not the ambient
curl Hessian.  Its Dirichlet compression \(A_{OO}\) has the same
gap.  For small Lipschitz weights its conjugation difference has
norm at most \(C\mu'\).  Hence the conjugated inverse of \(A_{OO}\)
has norm at most \(2/a_-\) when \(C\mu'<a_-/2\), by the Neumann
series.  Applying this to
\(\phi(b)=\mu' d(b,b')\) bounds each inverse entry by
\((2/a_-)e^{-\mu'd(b,b')}\).  This proof is uniform in \(O\).
It follows by two exponentially summable convolutions that
\[
 K_I=((A^{-1})_{II})^{-1}
       =A_{II}-A_{IO}A_{OO}^{-1}A_{OI}
\]
has exponentially summable entries at a possibly smaller rate.
Furthermore,
\(\langle v,K_Iv\rangle
 =\inf_w\langle(v,w),A(v,w)\rangle\ge a_-\|v\|^2\).
Thus the inverse marginal covariance, with its Schur
correction, retains the free gap. When exterior coordinates are fixed,
the free conditional precision is instead \(A_{II}\), which also
retains that gap by compression. The Schur matrix \(K_I\) describes
the free marginal obtained by integrating the exterior. Neither
statement identifies an interacting marginal with the pointwise
potential \(P_{k,I}^{\rm old}\) at fixed exterior data. The conjugation difference of \(K_I\)
also has norm at most \(C\mu'\): use
\(|e^{\phi(b)-\phi(b')}-1|
 \le\mu'd(b,b')e^{\mu'd(b,b')}\) and the row and column sums.
Choose \(\mu'\) so that this norm is at most \(a_-/4\).
Then
\[
 \operatorname{Re}\langle v,
      T_\phi(K_I-D^2P_{k,I}^{\rm old})T_\phi^{-1}v\rangle
 \ge\left(\frac{3a_-}{4}
       -C\mathcal E_k(g_k/\alpha_k)^2\right)\|v\|^2.
 \tag{OA21}
\]
For a sufficiently large coupling floor the old activities cost
at most \(a_-/4\), leaving a weighted gap of at least \(a_-/2\).
Other intrinsic or remainder Hessians must be charged separately.
The reference precision is held fixed while differentiating in
\(y\); no derivative of a covariance with respect to \(y\) is
introduced.

\subsection*{Remainders and the weaker finite-box alternative}

A holomorphic aged remainder with no fixed color spectators has
zero linear term at the trivial pair, but has no action
counterterm.  Its vacuum-only window bound has weight \(\rho_i^2\)
instead of \(\omega_{k+1-i}\).  The same proof gives (OA19) with
\(\mathcal E_k\) augmented by \(\sum_i\rho_i^2 r_i\), provided its
actual maps have the same domains and locality.  If
\(r_i\le C\theta_i^{-s}\), \(s>2\), and
\(\theta_i\ge\theta_k+c(k-i)-C_0\), this sum is at most
\[
 C\sum_{a\ge0}(1+a)(\theta_k+ca)^{-s}
                  \le C_s\theta_k^{2-s} .
 \tag{OA22}
\]
To see the last inequality, compare the sum with its integral and
put \(a=\theta_k t\).  The remaining integral of
\((1+t)(1+ct)^{-s}\) is finite.  This optional conclusion does not
cover a boundary term retaining stored color variables or a term
for which only an integrated norm is known.

If only a bound on the whole finite-box analytic function is
available, a useful but weaker statement remains.  Suppose
\(|P(y)|\le C E_0(r+D)^A\) on a Gaussian polydisc of radius
\(c\alpha_k/[g_k(r+D)^B]\). Then, on the concentric polydisc of half that radius,
the same coordinate Cauchy estimate and \(m_I\le C(r+D)^4\) give
\[
 \|D^2P\|_{2\to2}
       \le C E_0(r+D)^{A+2B+4}(g_k/\alpha_k)^2 .
 \tag{OA23}
\]
Applying (OA23) throughout the operational real window requires that
this window lie in that smaller polydisc; this inclusion is an
additional hypothesis, especially when \(B>0\).
Indeed each matrix entry has the bound without the last power
\((r+D)^4\), and the row and column contain at most \(m_I\)
entries.  For a direct volume bound \(A=4,B=0\), this is
\(CE_0(r+D)^8(\log\theta_k)^{-2q}\).  When
\(r+D\le C(\log\theta_k)^{2p}\), it tends to zero if \(q>8p\).
It preserves an unweighted gap.  It does not imply (OA20) at a
fixed positive spatial weight: conjugation could cost
\(e^{\mu'\operatorname{diam}I}\).  The rooted localization in
(OA15)--(OA18), rather than a count of window coordinates, is
what proves the stronger weighted conclusion.

\subsection*{What the estimate does and does not settle}

Under the stated ordinary analytic map hypotheses and the reviewed
pointed and window inputs, the old regular-activity contribution to
the localized Hessian is bounded by
\(CE_0(\log\theta_k)^{-2q}\), uniformly over old ages and box sizes.
The analytic radius is in the physical-to-Gaussian coordinate
conversion (OA11).  A bare operational cutoff domain, or only a
Gaussian-growth potential bound on that domain, would not prove it.
The two required primitive verifications are the uniform unscaled
pivot estimates (OA5)--(OA6) and the all-domain first-response and
membership estimates (OA7), including towers and outer substitutions.
They belong to the actual ordinary background analysis and have not
been proved for the Yang--Mills maps merely by this deduction.

The statement has no automatic consequence for the change in force
when exterior data are changed.  For that one needs, separately,
uniform estimates such as
\(\|D_b(P_e-P_{e_0})\|\le C\mathfrak b_k e^{-c d(b,O)}\)
on the actual compatible exterior and stored data.  A Hessian bound
in the variables of \(I\) does not bound derivatives with respect
to the exterior parameters.  Nor does it identify a clean-class
restriction with a Gaussian fibre kernel.  Field-independent
remote or holonomy bounds without the analytic domain assertion
used above remain separate tails; they receive neither an invented
curvature factor nor an unjustified Cauchy estimate.

% END MODULE chapters/old_activity_hessian_v7.tex


% BEGIN MODULE chapters/05_large_field.tex
\section{Large-field histories and the induction}
\label{candidate:large_field}
\candidatebanner{C}

\subsection*{2. Where the fields are wild: large-field regions, the factor paying for them, and the closed induction}

Convergence of the expansion of Step 1 requires that the plaquette variables of the backgrounds lie within $\varepsilon_k\eta_k^2$ of $\mathbf{1}$, with $\varepsilon_k=g_kp_0(g_k)$, and that the fluctuation fields stay below $\delta_k=(A_1/A_0)\varepsilon_k$. The two thresholds carry the same power of the logarithm and differ by a small fixed ratio. Nothing guarantees these conditions. For this reason each step inserts decompositions of unity $1=\sum\prod\chi\prod\chi^c$ built from local tests. At a place where a test fails, no expansion is made and the fields there are stored. The weights present at that place then supply a factor $e^{-c\,x_k\varepsilon_k^2}=e^{-c\,p_0(g_k)^2}$. For a failed plaquette test this weight is the Wilson weight. When an approximate fluctuation test fails, the weights present are the Wilson weight together with the kept gauge-fixing exponents. When a fluctuation test, a collar test or a rim test fails, the weight is only the Gaussian form of the stored fluctuation field. (The failure of one of the sharper tests of the large-field operation $\mathbf{R}$ of (2b) yields only $e^{-p_0(g_k)}$: 2a, 2c.) In this way the density becomes a sum over histories $h$, collected according to their labels $\mathcal{S}=(\Omega_j,\Lambda_j,S_j)_{j\le k}$, which consist of the regions of all levels. Besides the regions, a label carries the marks left by the operations $\mathbf{R}$, one for each island and for each drop-out that stands (2a); the labels are thus labels with marks. The exponential form given below holds for sums of histories and not for an individual history (2a, 2d).

\textbf{Theorem 2.1 [inductive description; ultraviolet stability].} Suppose that the constants are chosen in the prescribed order, beginning with the exponents and a large $L$ and ending with a small $\gamma$ (2e). Suppose also that the torus is admissible up to level $k$, and up to $K$ for the stability bounds; that is, $\ell L^{K-j}\ge2L^{11}M\check{R}_j$ for $j\le k$, respectively for $j\le K$. This is a lower bound on $\ell$, which is chosen after $M$ and $g$. When $\ell\in\mathfrak{L}(g)$ and the datum is a first-passage datum, the condition follows from the upper flow bound up to $K$ (3b). Assume further that $g_j\le\gamma$ and $\lambda_\sharp(j-i)-c_5\le x_i-x_j\le B_\sharp(j-i)$ for $i\le j\le k$ (up to $K$ for the stability bounds). When $g$ is small, (3a) provides these conditions for any first-passage datum; they are conditions on how the couplings flow: $\beta$ has no upper bound, and $K=K(\beta,g)$ is determined together with it. Then $\rho_k=(\mathbf{R}\mathbf{T})^k\rho_0$ is defined, and

\begin{gather*}
\FitMath{\rho_k(V_k)=\sum_{\mathcal{S}}\mathbf{T}'_k(\mathcal{S})\big[\chi_{k,\mathcal{S}}\exp A_{k,\mathcal{S}}(U_k(\mathbf{V}))\big],} \\[3pt]
\FitMath{A_{k,\mathcal{S}}=-A(x_k(\cdot),U)+\mathbf{E}_k+\mathbf{R}_k+\mathbf{B}_k-\mathcal{E}_{k,\mathcal{S}},}
\end{gather*}


\begin{gather*}
\FitMath{\|E^{(j)}\|_{j,\kappa}\le E_0,} \\[3pt]
\FitMath{\|R^{(j)}\|_{j,\kappa}\le g_j^{\kappa_0},} \\[3pt]
\FitMath{|B^{(j)}(X,\cdot)|\le\mathsf{B}\,e^{-\kappa d^\natural_j(\mathcal{S};X)},} \\[3pt]
\FitMath{\mathsf{M}(i,X)\le e^{-(2-c')p_0(g_i)},}
\end{gather*}

where $-\mathcal{E}_{k,\mathcal{S}}$ is the constant term of the action of the label and $\mathsf{B}$ is a constant fixed before $\gamma$ (2d, 2e); $d^\natural_j$ is the hull length of (2d), and $\mathsf{M}(i,X)$ is the mass of a component $X$ that, at some level $i\le k$, has passed all those tests of the operation $\mathbf{R}$ that come before the last one. Such a component is either healed at that level or an island, which remains in the label. The mass is defined as follows. Let $e^{\Phi_\omega}$ be the integrand of the operation $\mathbf{T}'_i(X)$ belonging to the component. This integrand contains the Wilson factor of the core (2b, 2c), and its exponent has the form $\Phi_\omega(\xi;\mathbf{p})=\Phi^\circ_\omega(\xi)+\sum_f\Phi_{\omega,f}(\xi;\mathbf{p})$, which is Borel in $\xi$, analytic in $\mathbf{p}$ and real at real pairs; each factor $\Phi_{\omega,f}$ depends on $\xi$ and on $\mathbf{p}$ only on a support $R_f\subset X^\sim$ (2d). Let $\omega$ run over the local histories of $X$, each with a frame in the hypercubic group $W_4$. Let $\xi$ denote the stored variables and $\mu_\omega$ their measure. Let $\mathbf{p}=(\mathbf{U},\mathbf{J})$ denote the complex pairs of the space $U^c_i(X^\sim,\alpha_{0,i},\alpha_{1,i})$, which contains the pairs at which a term evaluates the factor. Let $\mathsf{v}(\omega)$ be the space-time volume of the patch, which counts $\log x_j$ for each stored site of level $j$ (2c). Then

\begin{gather*}
\FitMath{\mathsf{M}(i,X):=\frac1{|W_4|}\sum_\omega e^{c\,\mathsf{v}(\omega)}\int\sup_{\mathbf{p}}\big|e^{\Phi_\omega(\xi;\mathbf{p})}\big|\,\mu_\omega(d\xi).}
\end{gather*}

The operation of a drop-out contains no Wilson factor of its own. For a drop-out the same bound holds in a modified form, with two changes. First, the Wilson weight over $X$ of the carrying term is included with the operation; the carrying term is a term of some level $k\ge i$ in which $X$ stands. Second, the supremum is taken over the evaluations at $\xi$. These are the real configurations that are compatible with $\xi$, regular layer by layer, and at which every test listed by $\omega$, with the exception of the test whose failure made $X$ a drop-out, takes its recorded value (the full list of these conditions is given in (2c)). This form follows from the induction and is not one of its clauses. The mass bound holds for every $c\ge0$ in the weight $e^{c\,\mathsf{v}}$ and for every small $c'>0$ in its exponent, provided the ceiling on the couplings is chosen after both $c$ and $c'$; the reserve constant $c_{\mathrm{r}}$ of (2e) is one such $c$. Apart from the core factors of the islands in $Z$ (2b), the operation $\mathbf{T}_k(Z)$ contains no Wilson weight. The factor coming from a failed plaquette test sits in $e^{-A(x_k(\cdot),U_k)}$, which is the function the operation acts on. In the sum over the labels, every component $C$ of $Z_k$ keeps a factor $e^{-c\,p_0(g_k)-\kappa_1|C|_k}$. The constant $\kappa_1>0$ outweighs the number $C^n$ of connected sets of $n$ cells that contain a given cell (2c). Each term is a non-negative measure. Neither $E^{(j)}$ nor $R^{(j)}$ carries a label index. On non-wrapping domains they are the same functions on every torus and at every cutoff: they are the values, at the lifts of the domains, of a single family on $\mathbb{Z}^4$. (The constant $\mathcal{E}_{k,\mathcal{S}}$, by contrast, depends on the label and on the torus.) Finally, let $\chi_k=\chi_k(T_\ell)$ denote the product, taken over every cube of the torus, of the tests $\chi_k(\square)$. Then, uniformly, there is a constant $E_-$ such that

\begin{gather*}
\FitMath{\chi_k\,e^{-x_kA(U_k(V_k))-E_-|\mathbb{T}^{(k)}|}\le\bar\rho_k,} \\[3pt]
\FitMath{\int\bar\rho_k\,dV_k\le e^{C(1+\log x_k)|\mathbb{T}^{(k)}|},} \\[3pt]
\FitMath{\int\rho_k\,dV_k=Z_{K,\ell,\beta}.}
\end{gather*}

The theorem is the case $\theta=x(\beta)$ of an induction carried out along an arbitrary frame $\theta$ (the opening of Step 3). That induction runs on all admissible tori simultaneously. The tube is a clause of the class $\mathfrak{C}_k(\theta)$ of (2e), and the step by itself does not yield this clause again (Theorem 2.11). For complex field arguments, each term depends on a triple $(U,A',\mathbf{J})$, in which $U$ is $G$-valued and $A',\mathbf{J}$ are complex. For a term of level $j$, the pair of a triple is $\mathbf{p}=(\mathbf{U},\mathbf{J})$ with $\mathbf{U}=e^{i\eta_jA'}U$, and the triple is called a splitting of its pair. The domains of analyticity are defined by conditions on triples. One of the statements propagated by the induction is that a term takes the same value at any two splittings of one pair. The splitting is therefore a coordinate and not an argument, and a term is a function of the pair. (In Step 1 the values coincide on a complex gauge orbit, with a margin left in the radii. Here what is asserted concerns two splittings of a single pair, and it holds on the domain itself.)

The core of the theorem is due to Ba{\l}aban: Comm. Math. Phys. 119 (1988) 243, Theorem 1 and Corollary 3, supplemented by 122 (1989) 175 and 122 (1989) 355, Theorem 1. For the idea of $\mathbf{R}$, he cites Gallavotti and co-workers as the closest earlier work. He states, for one torus, that neither the regular terms nor the part of the remainders produced by the fluctuation integral depend on the label (119, pp. 278--279). He also names a second choice of the thresholds and sets it aside; in that choice only the logarithm is evaluated on a bound of the flow known a priori (p. 246). The upper bound he gives is pointwise, $\varrho_k(V_k)\le e^{E_+|T_\eta|}$, stated for his density and his volume ((2.50)), and Theorem 2.1 gives the integrated form in its place. The extension adds the following, which belong to the Package C statements that we must prove:
\begin{itemize}
\item that independence with a range covering the remainders in their entirety, simultaneously on all admissible tori and cutoffs, by means of one family on $\mathbb{Z}^4$ (for the counterterm, the paper 119, p. 281, uses ``$\mathbf{E}^{(j)}(U_j,z)$ defined on the whole lattice $L^{-j}Z^d$'');
\item the positivity of each term, as a clause of the statement (in 122 (1989) 175, p. 176, Ba{\l}aban announces that the densities $\rho(Z'',V)$ are positive);
\item measure and formula as two separate layers;
\item the class taken as a set of presentations along a frame;
\item the flow inequalities, hypotheses of Theorem 2.1, obtained as conclusions of the induction, carried out jointly with the flow (3a);
\item covariance of the $R$-terms under translations only, term by term, and covariance of their totals by way of measures;
\item the mass, with the complex supremum inside the integral and with a reserve weight;
\item the hull length;
\item the three possible fates of a component, two of which are recorded as marks of the label (in 122 (1989) 175, pp. 181--183, Ba{\l}aban removes from the region on which $\mathbf{R}$ acts any component in which a test of some stage fails, and in 122 (1989) 355, (1.102), he retains the operation $\mathbf{T}_k(Y)$ where the last test fails), and the two sequences of regions attached to a label;
\item an order of the constants that contains no circle;
\item exactness for each partition of the histories.
\end{itemize}

The route also makes use of the following:
\begin{itemize}
\item the collar and the rim of (2a), which ensure that in the interior the mean of the conditional Gaussian measure stays under the thresholds set by the tests;
\item the factor depending on two levels in the sum over the scales of the $R$-terms (2e);
\item the terms $V(Y)$, which are not of small amplitude, with the consequent removal of the empty family before the exponentiation (2d);
\item in (2c), a single count over the genealogy of a component, which replaces the double induction of 122 (1989) 355, pp. 384--387.
\end{itemize}

In that induction, the case in which two regions join rests on a sentence on p. 387. That sentence states that after $n_1$ further hull operations, where $n_1$ is a fairly small number, for instance $n_1<10$, the domain obtained from the region whose wait is the longer of the two contains the domain obtained from the other region. The count never compares two such regions. A joining of regions costs it no separate payment, because the count charges the joining to a created cell (2c). The later steps rely on the theorem and on its sourced form (4a), which contains the insertion formula.

\subsubsection*{(2a) The tests, the nested regions, and the content of a history}

\textbf{The tests of step $t=k+1$} are applied to a term of $\rho_k$ whose live region is $Y=\mathcal{Z}_k$. First, outside a margin of $Y$, the new field $V_t$ undergoes the plaquette test $\chi_t(\square)$, and the cubes where this test fails make up $P_t$. The factor $\mathfrak{G}_k$ is inserted away from $P_t\cup Y$. By the Faddeev--Popov argument, its integral over the gauge group of the blocks equals $1$. Next comes the approximate fluctuation test $\chi'_k(\square)=\mathbf{1}\{|V_k(\mathcal{M}^kU_{t,\square}(V_t))^{-1}-\mathbf{1}|<2\delta_k\}$. This test is evaluated at the bonds of $\square^{\sim2}$ other than the pivot bonds (1d), and the third test below is evaluated at the same bonds. Here one takes from the orbit $U_{t,\square}(V_t)$ an element whose $k$-th average lies in the block axial gauge, and every such element gives the same tested number. The cubes where this test fails make up $Q_t$. These sets determine $\Omega_t$, the region on which $U_t(\mathbf{V})$ and $\mathcal{A}_k$ exist. The backgrounds and the propagators on sequences of regions are taken from (1b), (1c), together with their regularity and their decay in the local scale. A label's sequence starts at $\Omega_0=\Omega_1^\sim$, the first region enlarged by one layer of level-$1$ cells, and not at $T_\ell$; Theorem 4.13 and Lemma 4.15 require this. After that the test $\chi^{(k)}(\square)$ is made, and its failures make up $R_t$. Hulls ($'$, $\widehat{\cdot}$) and layers are taken in level-$t$ cells. The boxing map $\mathsf{bx}$ replaces each component contained in a cube of $O(1)$ cells by its bounding box, and repeats this until nothing changes. With these conventions,

\begin{gather*}
\FitMath{\Omega_t^c=Q_t'^{\sim2}\cup P_t'^{\sim4}\cup\widehat{Y}^{\sim8},} \\[3pt]
\FitMath{Z^{[0]}=\mathsf{bx}\big((\Omega_t^c)^{\sim2}\cup R_t'^{\sim}\big),} \\[3pt]
\FitMath{Z^{[n+1]}=\mathsf{bx}\big(Z^{[n]}\cup(\text{cells of the failures of round }n)^{\sim}\big),}
\end{gather*}

where the rounds are those of the recursion described below. The set $Z_t=\Lambda_t^c$ is the final $Z^{[n]}$, and $S_t=(\Omega_t\setminus\Omega_t^{\sim-1})\cup R_t$. The numbers of layers are those of Ba{\l}aban. If the boxing map is omitted, $Z^{[0]}$ coincides with his $\Lambda_t^c$ (Comm. Math. Phys. 119 (1988) 243, (3.16), (3.20)). The order of these numbers matters: an older set receives more layers, and each move is made one layer inside the region of the move before it. The creation set $\mathcal{W}^\wedge_t$ is the union of the hulls of every failure of the step: $P_t'$, $Q_t'$, $R_t'$, the cells of the collar bonds that fail, and the rim cells that fail. If nothing is created, the step reduces to the free evolution

\begin{gather*}
\FitMath{\Phi_t(Y)=\mathsf{bx}\big(\widehat{Y}^{\sim10}\big);} \\[3pt]
\FitMath{\text{in general}} \\[3pt]
\FitMath{Z_t\subset\widehat{Y}^{\sim10}\cup(\mathcal{W}^\wedge_t)^{\sim O(1)}} \\[3pt]
\FitMath{\text{up to boxing}}
\end{gather*}

That is, one coarsens, adds a margin and boxes; this is the $S(Z)$ of (2c). The size $\check{R}_k$ is the smallest power of $L$ that is not less than $(\log\min_{i\le k}x_i)^r$, with $r\ge2$. Because it is a power of $L$, all partitions are nested. Because it is formed from the running minimum, it is non-increasing (Ba{\l}aban forms it with $g_k$ itself), and therefore $t\mapsto t-\mathsf{N}_t$ is strictly increasing. Moreover $\check{R}_k\le L\check{R}_{k+1}$, and two consecutive drops never occur, because $x$ changes by $O(1)$ per step and $\gamma$ is small. A test cube $\square$ has side of order $M_2\check{R}_t$. Here $M_2$ is a power of $L$ with $L^2M_2\le M$ and $c\,\delta_0M_2\ge1$. The last condition is a floor on $M_2$, and $\delta_0$ is fixed together with $L$, which comes earlier. Hence a test cube lies well inside a cell of its own level.

\textbf{What factorises, and why.} The tests $\chi_t$ and $\chi'_k$ are built from localised backgrounds. They depend on $V_t,V_k$ only within $O(1)$ test cubes. The third test, and also the collar and rim tests, depend on $\mathcal{A}_k=\tfrac1i\log(V_k(V^{\mathrm{bg}}_k)^{-1})$, which is formed with the true averaged background on $\Omega_t$. These tests are therefore local in the new variable $\mathcal{A}_k$, but not in $(V_k,V_t)$. Once $\mathcal{A}_k$ has become the variable, the surviving $\chi'_k$ (on $S_t$) compare $e^{i\mathcal{A}_k}V^{\mathrm{bg}}_k$ with $\mathcal{M}^kU_{t,\square}$. They therefore depend on $V_t$ over the whole of $\Lambda_t$. Four facts together ensure that everything attached to $Z_t$ factorises over its components:
\begin{itemize}
\item The backgrounds are localised.
\item The substitution on the strip (2b) is made. After it, every indicator kept in $\Omega_t\cap\Lambda_t^c$ depends on $V_t$ in $\Lambda_t^c$ alone, apart from the factors that equal $1$ on the supports.
\item The conditional quadratic form is cut (2b).
\item The pivot values are dealt with as follows. An old test $\chi_k(\square)$ of a cube $\square\subset\Omega_t^c$ adjacent to $\Omega_t$ also depends on $V_k$ at pivot bonds (1d) inside $\Omega_t$. There $V_k$ is not a stored variable, and through $e^{i\mathcal{A}_k}V^{\mathrm{bg}}_k$ a test of one component depends on the fields stored in another. At such a bond, $V_k$ is given the value that solves, for the pivot bond, the averaging constraint $\mathcal{M}(V_k)(c)=V_t(c)$ of the single coarse bond $c$. Ba{\l}aban (122 (1989) 175, (1.3)) retains only those old tests whose cubes lie outside $\Omega_t^{\sim}$; the tests that he leaves out equal $1$ on the supports by Lemma 2.2 (a), whose margin is this layer.
\end{itemize}

Inside $\mathbf{R}$, the tests are built from backgrounds localised in the component, never from the background of the torus. The induction carries the following property: each indicator stored in $\mathbf{T}_k(Z)$ depends only on those field variables that are localised in $Z$.

\textbf{Totality.} By definition, a test is the indicator function of the set where all objects occurring in it are defined and where its inequalities hold. These objects are the localised minimal orbit, which must exist and be a single orbit, the contour variables, and the logarithms. Consequently $\chi^c=1$ wherever some object is undefined. A factor of $\mathfrak{G}_k$ equals $1$ where $V(y,x)$ is undefined, by a separate convention for that factor. Hence $\int\mathfrak{G}_k(V^u)\,du=1$ for every $V$, and each insertion $1=\sum\prod\chi\prod\chi^c$ is an identity that needs no assumption. The tests are not empty conditions. The averaging constraint transports smallness: $|\partial V_k-\mathbf{1}|<2\varepsilon_k$ implies $|\partial V_t-\mathbf{1}|<CL^2\varepsilon_k$. Therefore $U_{t,\square}(V_t)$ exists on the carried supports, and the new test further restricts a field that is regular already (Comm. Math. Phys. 119 (1988) 243, p. 246). The cost is that $\chi^c$ contains configurations that are not large in any quantitative sense. These configurations must also produce the small factor (2c).

\textbf{Collar and rim}: this device is an addition to the step of Comm. Math. Phys. 119 (1988) 243, Sect. 3. The stored part of the fluctuation field moves the mean of the conditional Gaussian measure of (2b). The thresholds $\delta_k$ are stated in the unscaled normalisation, and in that normalisation $\mathcal{A}_k|_{\Lambda_t}$ has the mean

\begin{gather*}
\FitMath{\mathbf{m}=-\,C^{(k)}(\Lambda_t)\,\Delta^{(k)}\,\mathcal{A}_k\big|_{\Omega_t\setminus\Lambda_t}}
\end{gather*}

This map is linear, so in terms of the scaled fields $A$, $A'$ of (2b) it is the same map, with each threshold divided by $g_k$; only $|\mathbf{m}|$ enters. This shift must remain well below the least threshold among the inner tests; that threshold is $\theta_c\delta_k$ on bonds that cross $\partial\Lambda_t$, and $\delta_k$ further in. A recursion therefore keeps enlarging the region until, along its boundary, $\mathcal{A}_k$ satisfies two sharper tests:
\begin{itemize}
\item $|\mathcal{A}_k|<\delta_k/(4\Theta_1)$ on the collar, the set of stored bonds whose distance from $\Lambda_t$, measured in level-$k$ units, is at most $\varpi_T$. With a qualifying word added, the same name serves for three further objects. In (2b), the collar of a stage means the shell over which that stage integrates. In (2c), the waiting collar is the set that keeps $X\cap\Omega_{m_\flat}$ apart from the floor. In (2d), (2e), the collar of a component $X$ means the enlarged component $X^\sim$, that is, $X$ together with one layer of $M$-cubes of its level.
\item $|\mathcal{A}_k|<\theta_c\delta_k$ on the rim, the set of cells of level $t$ of $Z_t$ that touch $\Lambda_t$.
\end{itemize}

At greater distance, all that is known is $|\mathcal{A}_k|<C\delta_k$. Every round examines bonds and cubes that no earlier round examined. Hence the recursion terminates, and the admissible outcomes it produces constitute a decomposition of unity. The kernel of $\mathbf{m}$ is bounded, with exponential decay. The contributions of the collar, the rim and the remaining part give

\begin{gather*}
\FitMath{|\mathbf{m}|\le C\big(\Theta_1^{-1}+\theta_ce^{-c\varpi_T}+e^{-cM\check{R}_t}\big)\delta_k<\tfrac12\theta_c\delta_k} \\[3pt]
\FitMath{\text{if }\theta_c\Theta_1\ge C,\,\varpi_T\ge C,\,\gamma\text{ small}.}
\end{gather*}

The rim fraction $\theta_c\in(0,1)$, the collar width $\varpi_T$ and the collar factor $\Theta_1$ serve this purpose.

\textbf{Lemma 2.2 [the new tests imply the old; the form reproduces itself].} (a) On $\Omega_t^\sim$, that is, on $\Omega_t$ enlarged by a margin of one layer of cells, the indicators in $\mathfrak{G}_k$ and the tests $\chi_k(\square)$ that the term carries all equal $1$. This holds on the support of the new tests together with the tests and averaging constraints that the term still carries alongside them. (b) The test $\chi^{(k)}$ is made only on cubes $\square\subset\Omega_t^{\sim-1}$. For such a cube, if $\square\not\subset R_t$, then $\chi^{(k)}(\square)$ implies $\chi'_k(\square)$. Consequently, within $\Omega_t$, the test $\chi'_k$ remains only on $S_t$. (c) After the step, the tests that stand outside the integral operations are exactly $\chi_{t,\mathcal{S}}:=\chi_t(\Omega_t)=\prod_{\square\subset\Omega_t}\chi_t(\square)$. They are formed with the action sequence defined below, and with $\chi^c_t$ on an island less a margin.

For (a): up to a gauge transformation, 

\begin{gather*}
\FitMath{U_{k,\square}=\exp\big[i\eta_k\mathbf{H}_{k,\square}\big(\tfrac1i\log[V_k(\mathcal{M}^kU_{t,\square'})^{-1}]\big)\big]U_{t,\square'}}
\end{gather*}

. Here $\mathbf{H}_{k,\square}$ is the localised response $\mathbf{H}_k$ of (1b), so the argument has the form displayed only off a boundary layer of the cube; inside that layer the argument is built from the lift of $V_k$ and is bounded by $C\varepsilon_k$. It follows that

\begin{gather*}
\FitMath{|\partial U_{k,\square}-\mathbf{1}|\le(1+c)\varepsilon_t\eta_t^2+C\big(\delta_k+e^{-c\,\delta_0M_2\check{R}_k}\varepsilon_k\big)\eta_k^2\le\big(CL^{-2}+CA_1/A_0+Ce^{-\check{R}_k}\big)\varepsilon_k\eta_k^2<\varepsilon_k\eta_k^2 .}
\end{gather*}

Three mechanisms appear repeatedly:
\begin{itemize}
\item Curvature scales by the factor $\eta_t^2/\eta_k^2=L^{-2}$.
\item The ratio $\delta_k/\varepsilon_k=A_1/A_0$ is small.
\item The artificial boundary conditions of localised minimisers act only through $e^{-c\delta_0\operatorname{dist}}$ (1b, 1c). On cubes whose side is proportional to $\check{R}_k\approx(\log x_k)^r$, this quantity is smaller than every power of $g_k$, because $r>1$. In the present estimate it needs to beat only constants; on the strip (2b) it must beat a negative power of $g_k$ (Ba{\l}aban, Comm. Math. Phys. 119 (1988) 243, (3.2)--(3.8)).
\end{itemize}

For (b): 

\begin{gather*}
\FitMath{|V^{\mathrm{bg}}_k(\mathcal{M}^kU_{t,\square})^{-1}-\mathbf{1}|\le C\tfrac{A_0}{A_1}e^{-c\check{R}_k}\delta_k<\delta_k}
\end{gather*}

 ((3.17)--(3.19) of the same paper). The factor $2$ between the two thresholds gives exactly the room needed for this. For (c): a test cube of step $t$ serves also as a test cube of the density of level $t$, carrying the same background and the same threshold. The old tests inside $\Omega_t^c$ pass into the integral operation. The same ladder of implications holds inside $\mathbf{R}$ (2b). On the support of the tests of a stage, the following tests equal $1$: those of the preceding stage, and those of the single stored step whose fields the stage integrates (the plaquette tests of its two levels, the gauge indicators, the approximate fluctuation tests). Further, consider any point at which a term is evaluated: the triple $(U,0,J(U))$, where $U$ is the background of the label, at fields where no characteristic function of the term vanishes. At such a point, $\chi_q(\square)(V_q)=1$ holds on every cube $\square$ of level $q$ that lies in the layer $\Omega_{q+1}^\sim\setminus\Omega_{q+1}$ and not in a drop-out. Also at such a point, the current of the bare field, formed from the plaquettes inside $\Omega_0$ only, stays within the bound that the domains of analyticity impose on their current (the size $|\mathbf{J}|$ is smaller than the first radius of the domain, measured in the units of the local scale). This bound holds also at the bonds of $\Omega_0$ that have a plaquette leaving $\Omega_0$.

\textbf{Proposition 2.3 [geometry].} The regions are nested, $T_\ell\supset\Omega_1\supset\Lambda_1\supset\dots\supset\Lambda_k$, with gaps of at least $cM\check{R}_j$ in level-$j$ units. Every large-field region and every creation set is a union of cells of its own level. Both statements hold for the regions that the steps produce (119, p. 256). Each operation $\mathbf{R}$ has stage sequences and prepared sequences (namely the regions that the label has after $m$ of its stages, and after all $\mathsf{N}_k$ of its stages; see (2b)), on which the backgrounds of (2b), (2d) are formed, and drop-outs keep nests. In these sequences and nests the boxes are only unions of $M$-cubes, and the only lower bound on the gaps is one $M$-cube of the later level:

\begin{gather*}
\FitMath{\operatorname{dist}_\infty(Z_j,\Omega_{j+1})\ge LM,} \\[3pt]
\FitMath{\Omega_{j+1}\subset\Lambda_j^{\sim-1},} \\[3pt]
\FitMath{Z_j^\sim\cap\Omega_{j+2}=\emptyset} \\[3pt]
\FitMath{(\text{level-}j\text{ units; layers of }M\text{-cubes of level }j),}
\end{gather*}

the first of these inequalities being an equality in the upper part of a nest, where $\Lambda_j=\Omega_j$; the reason is that a step places eight layers around the old region, and the preparation removes most of them again. Consequently the estimates of (1b), (1c) must hold with a separation of only one $M$-cube. A smallness $e^{-c\check{R}_j}$ from a separation is available only at regions produced by steps. In each decomposition exactly one outcome equals $1$, and hence $\rho_k\,dV_k=\sum_h\tau_h$ holds exactly. A history records every outcome: $P_t$, $Q_t$, $R_t$, the collars and rims, and the outcomes inside $\mathbf{R}$. It determines the label, the creations and the ages. There are two kinds of creation:
\begin{itemize}
\item Creations by step $t$ are the cells $\mathfrak{c}$ of $\mathcal{W}^\wedge_t$, which are $M\check{R}_t$-cubes of level $t$. They are counted one cell at a time, and each carries one factor $e^{-c\,p_0(g_t)^2}$.
\item Creations by an operation are its marked components (next paragraph); the tests of an operation are decompositions of unity component by component. Each marked component counts as one creation and carries one factor $e^{-p_0(g_t)}$, which must pay for $\mathsf{N}_t+1$ levels of one small box.
\end{itemize}

All later statements of locality rest on this.

\textbf{Three fates, marked components, two sequences of regions} (Ba{\l}aban declines ``a systematic description of all possible cases'': 122 (1989) 175, p. 183). We set $t_\flat:=t-\mathsf{N}_t$. In each component $X$ of its class (2b), the operation $\mathbf{R}$ of level $t$ prepares a nest $Z''_l(X)$, $t_\flat<l\le t$. The nest consists of increasing boxes in $\Omega_l^c\cap X$. Each box is what remains of the large-field region of scale $l$ after the stage that integrates the stored fields of level $l-1$. A component has one of three fates:
\begin{itemize}
\item \emph{Healed} (every test passed): $\Omega_l,\Lambda_l,S_l\mapsto\Omega_l\cup X,\Lambda_l\cup X,S_l\setminus X$ for $l\le t$, and $X$ leaves the label.
\item \emph{Island}, also called second class (only the last test fails, namely $\chi_t$ of the new field on $X$ less a margin): the same replacement is made for $l<t$ and for $\Omega_t$, $S_t$, but $\Lambda_t$ is left unchanged. Thus $X\subset Z_t$ remains, as a component without inner regions. On it, $\mathbf{T}'_t(X)$ is kept together with its variables, its measure and its factors (2d).
\item \emph{Drop-out} (some test fails after $n\le\mathsf{N}_t$ stages): $X$ remains in $Z_t$ and retains, as its integral operation (2b), the operations carried out up to that point. Moreover

\begin{gather*}
\FitMath{\Omega_l\cap X=X\setminus Z''_l(X)} \\[3pt]
\FitMath{(t_\flat<l\le t_\flat+n),} \\[3pt]
\FitMath{\Lambda_l\cap X=X\setminus Z''_l(X),} \\[3pt]
\FitMath{S_l\cap X=\emptyset} \\[3pt]
\FitMath{(\text{the same }l,\,l<t),} \\[3pt]
\FitMath{\text{all other entries unchanged, so }\Lambda_t\cap X=\emptyset .}
\end{gather*}

(For $n=\mathsf{N}_t$, then, the entry $S_t\cap X$ keeps its previous value; on the support of the tests of the last stage, the tests retained on it equal $1$, by the ladder of implications inside $\mathbf{R}$ that is stated in the argument for Lemma 2.2.)
\end{itemize}

Islands and drop-outs are the \emph{marked components}. A mark persists until a later healed component or island swallows it. The word also denotes two other objects: the mark $\Theta^{(k)}$ below, which is a weight on histories, and a marked term of a cluster expansion (2e). A drop-out's mark keeps the regions $\Omega^\circ_l=\Omega_l\cap X$, and likewise $\Lambda^\circ_l,S^\circ_l$, which $X$ possessed before the operation. For $0\le\boldsymbol{\nu}(X)\le n(X)$ we set

\begin{gather*}
\FitMath{\Omega^{\boldsymbol{\nu}}_l\cap X=X\setminus Z''_l(X)\,\,(t_\flat<l\le t_\flat+\boldsymbol{\nu}(X)),} \\[3pt]
\FitMath{\Omega^{\boldsymbol{\nu}}_l\cap X=\Omega^\circ_l\,\,(t_\flat+\boldsymbol{\nu}(X)<l\le t_\flat+n(X)).}
\end{gather*}

The choice $\boldsymbol{\nu}=n$ gives the \emph{determining sequence}, which describes the field in its present state. It governs the determining sets, $U_k(\mathbf{V})$ and the propagators; the analyticity spaces involve all the $\Omega^{\boldsymbol{\nu}}$. The choice $\boldsymbol{\nu}=0$ gives the \emph{action sequence}, the regions for which the terms were constructed. It governs the index sets of $\mathbf{E}_k,\mathbf{R}_k,\mathbf{B}_k$, the profiles $\varphi_j$ and hence $x_k(\cdot)$, the carried tests, and the constants $\mathcal{E}_{k,\mathcal{S}}$. No single sequence can play both roles. A stored region in a determining position requires variables that have already been integrated. A stage region in the sums re-sorts terms that were made for other regions. Three facts make the device workable. \emph{Marks wait}: a mark of level $s$ lying inside a class component of level $t$ satisfies $s<t_\flat$. Hence standing marks that share a part are nested, their windows $(s_\flat,s]$ are disjoint, and at any one place and level at most one drop-out mark is active. Next, each $\Omega^{\boldsymbol{\nu}}$ is the label of an actual history, the one in which the operations stop after $\boldsymbol{\nu}$ stages. Its steps, classes and nests coincide with those of the history at hand, and so Proposition 2.3 applies to it: \emph{the later geometry does not depend on how far an interrupted operation had progressed}. Finally, a term created at a moment $\mu$ (following the step, or the operation, of a level $j\le t$) retains permanently the space in which it was made:

\begin{gather*}
\FitMath{\Omega^{\boldsymbol{\nu}_\mu}_l(\mathcal{S}_t)=\Omega_l(\mathcal{S}^\mu)\cup Z_{[j,t]}} \\[3pt]
\FitMath{(l\le j),}
\end{gather*}

where $\mathcal{S}_t$ denotes the level-$t$ label as it stands once the operation $\mathbf{R}$ has been performed, $\boldsymbol{\nu}_\mu=n$ on each mark older than $\mu$ and is $0$ on each later mark, $\mathcal{S}^\mu$ is the label at that moment, and $Z_{[j,t]}$ is the union of all components that were healed or became islands at levels $j,\dots,t$.

\textbf{Locality in the history.} If the large-field regions of two histories never touch, the two histories combine into a single history $h^1\cup h^2$. There are two reasons for this. (i) \emph{A margin beats a reach}: two regions made of level-$t$ cells that do not touch are separated by at least one cell. One cell exceeds the largest range of a bond test, of a pivot bond and of a collar (in the case of the collar, this is the floor $M\ge\varpi_T$ of (2e)). Taking hulls and adding layers commute with forming unions, and boxing depends on a single component only. A step therefore acts on each component separately. (ii) The class of $\mathbf{R}$ depends on the past: it depends on every object whose creation falls in the window $(t_\flat,t]$, or which carries a mark at a level in $[t_\flat,t)$. Moreover, the recursion is not monotone: a later region may cover a healed place once more. \emph{Window lemma}: no creation within the window of $t$ has been healed. Indeed, healing at a level $i<t$ something created at step $s$ requires $s\le i-\mathsf{N}_i$, and for something marked at level $s$ it requires $s<i-\mathsf{N}_i$. On the other hand $i-\mathsf{N}_i<t-\mathsf{N}_t\le s$, since $\mathsf{N}_t=\check{R}_t$ is non-increasing. Within the window, therefore, the creations and the marks of each history still lie in that history's own region. It follows that each history is the superposition of its clusters, and that clusters may be interchanged between histories. Each cluster has a root. The root is of one of two kinds: a component that remains live at the last level, or a healed component, considered at the level where its healing took place. The cluster is made up of the components of earlier levels that flow into the root, together with everything the history records in them. Each component, at whatever level, lies in the cluster of precisely one root, namely the last of its successors; thus nothing recorded about the healed components is lost.

\textbf{A hard-core gas; the standing part.} Erasing clusters of healed components from a history yields a history with the same label. The \emph{standing part} $h^\flat$ of $h$ is what remains after the clusters of all its healed components are removed. For a string $s$ of outcomes of the next step one has $(h,s)^\flat=(h^\flat,s)^\flat$. Let $\bar h$ be a history of level $k$ in which no component is healed at any level. In the rest of this paragraph each history is taken without the clusters healed below level $k$. It stands for the class of histories that yield it once these clusters are erased. Consider the histories that reduce to $\bar h$ once one erases the clusters belonging to those components that were healed at level $k$. They are in bijection with the pairs $(\mathfrak{X},(\varsigma_X)_{X\in\mathfrak{X}})$. Here $\mathfrak{X}$ is a finite collection of possible healed components which touch neither one another nor $Z_k(\bar h)$. Further, $\varsigma_X\in\mathcal{H}^1(X)$ ranges over the histories of $X$ alone in which $X$ is healed and nothing is healed before it, and the choices for different components are independent. All these histories share one label. An individual $\tau_h$ is a measure produced by the moves, and it has no exponent of the form in Theorem 2.1. The part of $R^{(k)}$ produced by $\mathbf{R}$ arises only through the sum over $(\mathfrak{X},\varsigma)$, a sum which $\mathbf{R}$ itself performs (Proposition 2.10 [outlined below; Section~\ref{sec:actual-history-closure}]). This sum is a polymer sum. The reasons are the gas property and two facts of locality in the history: $\mathbf{T}'_k(X)$ is the same in all the summands containing $X$, and a healed patch leaves no trace outside itself (2d).

\textbf{Traces and charts.} Let $W$ be a region, and let $W^+_i$, the \emph{window} of level $i$, denote the enlargement of this region by a single layer of level-$i$ cells. The trace of a label is

\begin{gather*}
\FitMath{\operatorname{tr}_W\mathcal{S}:=\big(\Omega_i^c\cap W^+_i,\,Z_i\cap W^+_i,\,S_i\cap W^+_i\big)_{i\le k},}
\end{gather*}

with two additions. The first addition is the marks of the label. They enter collectively rather than individually. For every level, every kind of mark and, for drop-outs, every number of stages carried out, the trace records the union of all components carrying such a mark, together with the unions of their stored regions and of their nests. Hence the trace does not show which pieces in $W^+_i$ belong to the same component. The second addition consists of three separate sets of stored bonds of level $i-1$ contained in $W^+_i$: those in the collar, those at a rim cube, and those in the cut-off zone of the strip. The cut-off zone consists of those bonds of the strip that lie near the set where the fluctuation field is large and carry no kept small-field condition of that field; there the substitution of (2b) uses the cut-off $\vartheta_{\mathrm{co}}$. Where a bond carries a kept condition, the variables are left unchanged, while on the remaining bonds of the strip the substitution proceeds with no cut-off. The terms depend on these three sets, and the regions in $W^+_i$ do not fix them: to decide if a cube is a rim cube one has to look at its neighbours, and $\Omega_i^c$ enters up to two layers away. The trace $\operatorname{tr}_Wh$ of a history is formed in the same way from everything that $h$ records, including the labels of its earlier moments. It is finer: knowing the trace of a label does not fix the traces of the histories that carry it. A \emph{chart} $\iota:W\to W'$ is a translation whose vector belongs to the lattice of level-$t$ cells. It acts either within one torus or, through a common lift of a non-wrapping region, from one torus and cutoff to another. It leaves invariant the lattices $\mathbb{T}^{(i)}$, $i\le t$, and also the partitions of all levels up to $t$ into cubes and cells. Neither a shift by a single site nor a mere isometry of the windows does so. Suppose that $\iota\operatorname{tr}_Wh=\operatorname{tr}_{W'}h'$, that the increments and the earlier terms in $W$ correspond, and that the bare weight corresponds on the window $W^+_t$ (this window contains $W^{\sim O(1)}$, since the cells of level $t$ are $M\check{R}_t$-cubes). Then the two histories produce corresponding terms with domains in $W$, including their natural domains of definition. Here a term is regarded as a partial map on the fields on $W$, and one of two corresponding terms is defined exactly where the other is. The integration of $V_k$ over the box is the move (3) of (2b) and forms part of the main step of $\mathbf{R}$. Consider the part of the regular terms on a domain of level $m$ that this integration takes. It depends only on the trace through the levels $i\ge m$, whatever the values of $k$, of $i-m$, of the size of the domain and of the volume. The standing operations, that is the $\mathbf{T}'_i(X)$ of the islands carried by a term, are transported likewise, measures included and one factor at a time (2e). For the analyticity spaces more is needed. The spaces of two corresponding terms are mapped onto one another if two further conditions hold. First, the bonds of the bare lattice carrying their fields must form corresponding sets. This is needed because whether a bond leaving the domain ends in $\Omega_1^\sim$ is decided outside the window; for a domain in $\Omega_1$ there is no condition. Second, at each level, two unit cubes of the bare lattice lying in the domain must belong to a common drop-out mark of that level if and only if their images do. This is needed because a domain that is not convex may cut a mark into two pieces; where no mark meets the domain there is no condition. On the infinite lattice no sum over labels is formed. The only terms formed there are those of the never-large label, the label at which no region stands; domain by domain, these terms contain the sums over healed components that $\mathbf{R}$ performs (Proposition 2.10 [outlined below; Section~\ref{sec:actual-history-closure}]). A term whose window $W^+_t$ lies in $\Lambda_t$ coincides with the corresponding term of that label (the terms belong to the standing part $h^\flat$, and everything that this part records lies in $Z_t$: 2e). If the domain of such a term does not wrap, it is one and the same function on every torus and at every cutoff, and in this way $E^{(j)}$ and $b_j$ are computed (3b). It is this that links the geometry to (3b) and (4b).

\textbf{The net form of the exponent.} Each function $\mathfrak{t}$ of the exponent has a fixed place, given by its level, its domain and the way it is made. The label enters only through the sorting into $\mathbf{E}_k,\mathbf{R}_k,\mathbf{B}_k$, the profile and one number:

\begin{gather*}
\FitMath{A_{k,\mathcal{S}}(U)=-\beta A(U)+\sum_{\mathfrak{t}}\mathfrak{t}(U)+\mathrm{const}(\mathcal{S}),} \\[3pt]
\FitMath{\text{because}} \\[3pt]
\FitMath{x_k(\cdot)+\sum_{j\le k}b_j\varphi_j(\cdot)=\beta}
\end{gather*}

and $c\mapsto A(c,U)$ is linear. Consequently a change of label changes no local term.

\textbf{The law of the histories.} The set $\mathfrak{H}_k$ of histories of level $k$ is finite. A run has the datum $\mathfrak{d}=(\beta,m,g)$ and takes place on $T_{\ell_m}$, $\ell_m=2L^m$. The mass of a history $h\in\mathfrak{H}_k$ is $\mathsf{t}(h):=\tau_h(\mathcal{V}_k)\ge0$. Since the steps are exact (2e), $\sum_{h\in\mathfrak{H}_k}\mathsf{t}(h)=Z_{K,\ell,\beta}$ at each level. The law of the histories of the last level, and in the same way the law of the histories of any level, is

\begin{gather*}
\FitMath{\tilde\nu_{K,\ell,\beta}(\mathsf{E}):=Z_{K,\ell,\beta}^{-1}\sum_{h\in\mathsf{E}}\mathsf{t}(h),}
\end{gather*}

where events of histories are written in sans-serif letters, $\mathsf{E}_k\subset\mathfrak{H}_k$. The creation set of step $j$ is $\mathcal{W}^\wedge_j(h)$, as before. Without the superscript, $\mathcal{W}_K(g)$ denotes the first-passage window. The live region $\mathcal{Z}_k(h)$ is $Z_k(h)$ with the components healed at step $k$ removed; it is the input region of step $k+1$. The event that the cell $\mathfrak{c}$ is live at level $k$, that is $\mathfrak{c}\subset\mathcal{Z}_k$, is written $\mathfrak{lv}_k(\mathfrak{c})$. For a cell $z$ of level $n$, the survival event $\mathcal{E}^{(n)}_z(k)$ is the event that $z\subset\mathcal{Z}_i$ for $n\le i\le k$. The statements in which a run is continued past $K$ are (3d), (4d) and (7e). For these statements we fix an anchor coupling $\gamma_\star\in(0,\gamma_J]$, which is determined by the constants up to $c_a$ and by $\gamma_J$, and put $X_\star=\gamma_\star^{-2}$. The anchor level is $K_\star=K(\beta,\gamma_\star)$; it goes to infinity as $\beta$ does. In these statements the reference coupling obeys the ceiling $g\le\gamma_\star$, and the bounds on the flow hold up to $K_\star$. Also, at a given $\beta$, the run passes the window of $g$ no later than the window of $\gamma_\star$, $K(\beta,g)\le K_\star$; this follows from $g\le\gamma_\star$ and the monotonicity of the first passage.

\textbf{Lineages and marks.} Let $X$ be a component of $Z_k(h)$, at any level, that is either live at level $k$ or healed at step $k$; a component live at $k=K$ is a terminal component. Its lineage regions are $\mathsf{Z}_k(X):=X$ and, going down the levels, $\mathsf{Z}_j(X):=$ the union of all components that are live at level $j$ and contained in $\mathsf{Z}_{j+1}(X)$. Its lineage consists of these components $(j,C)$ together with the creations $(j,W)$ inside the $\mathsf{Z}_j(X)$. The weight of a level-$j$ cell is $\mathsf{r}_j:=M^4\check{R}_j^5$. We write $|C|_j$ for the number of level-$j$ cells in $C$. For comparison, $\#_j(Y)$ counts the sites of $\mathbb{T}^{(j)}$ in $Y$, and a cell contains $(M\check{R}_j)^4$ of them. Finally, for a component $W$ of the set that step $j$ creates in the lineage, $d'_j(W)$ is its tree length measured in these cells; the marks set by the operations $\mathbf{R}$ do not enter $\operatorname{Dep}$. Then

\begin{gather*}
\FitMath{\operatorname{Vol}_k(X;h):=\sum_{(j,C)\text{ in the lineage}}2\,\mathsf{r}_j\,|C|_j,} \\[3pt]
\FitMath{\operatorname{Dep}_k(X;h):=\sum_{(j,W)\text{ created in the lineage}}c\,p_0(g_j)^2\big(d'_j(W)+1\big).}
\end{gather*}

The mark at level $k$ has two dials, $\mathsf{w}\in[0,1)$ and $c_a\ge0$, the second also called the age constant. The mark and the end mark $\Theta:=\Theta^{(K)}$ are

\begin{gather*}
\FitMath{\Theta^{(k)}(h):=\sum_{X\text{ live at level }k}\big[\mathsf{w}\operatorname{Dep}_k(X;h)+c_a\operatorname{Vol}_k(X;h)\big]\ge0 .}
\end{gather*}

The mark keeps back the fraction $\mathsf{w}$ of the small factor of each creation in such a lineage. It also charges $2c_a$ for every cell, multiplied by the weight of the cell, and for every level. The constant $c$ in $\operatorname{Dep}$ is a single number. It is chosen after the fractions of the fluctuation threshold, namely those at which a step evaluates its tests: $1/(4\Theta_1)$ on the collar and $\theta_c$ on the rim. It is subject to $c\le c'\min\{\Theta_1^{-2},\theta_c^2\}$, where $c'$ is chosen after $(A_1,A_0)$ and contains the ratio $(A_1/A_0)^2$ of $\delta_k^2/g_k^2$ to $p_0(g_k)^2$. No ceiling on $g$ can take the place of this choice, since such a test has a deposit and a gain that are both of degree $2\mathfrak{p}_0$ in $\log x_j$.

\textbf{Three ages} are attached to a component $C$ that is live at level $k$. (1) The age in the label is determined by the oldest region of the label within $C$: $\operatorname{age}^{\mathrm{lab}}(C)=k-\min\{t:C\cap Z_t\ne\emptyset\}$, where $Z_t$ denotes the large-field region of the label $\mathcal{S}(h)$ at scale $t$. (2) The age of the lineage is $\operatorname{age}_k(C;h)=k-j_0$, with $j_0$ the earliest level that has a non-empty lineage region $\mathsf{Z}_j(C)$. The events $\mathfrak{lv}^{=a}_k(\mathfrak{c})$ and $\mathfrak{lv}^{\ge a}_k(\mathfrak{c})$ say that at level $k$ the cell lies in a live component whose lineage has age $a$, respectively at least $a$. (3) The seed age is $k$ minus the seed level. The seed level is the least $t$ such that $C$ contains, as a point set, some region created at step $t$. The corresponding events are written $\mathfrak{sd}^{=a}_k(\mathfrak{c})$. The first age is at most the second, and the second is at most the third. The second and the third can differ. Suppose that a region was healed at some earlier step and that a later live region again covers its place. Such a region counts for the seed level but is not part of the lineage. If the seed level is defined using the regions that no step before $k$ has healed, the two ages coincide. Payment is made for the age of the lineage alone (2e). Age has to be paid for with creations. In a lineage with $n$ creations, those made by operations included, a creation made at step $t$ satisfies $k-t\le Cn(\mathsf{N}_t+\log n)$. Here a creation is a created cell or a mark. By the waiting bound and the count of (2c), this bound holds for every lineage; hence an old lineage contains many creations. These objects enter (2e) and (4d). Through the age series of (4d) they underlie both the two-point law (7a)--(7c) and the four-point law (7d), namely clauses (c)(3), (c)(4), (c)(5).


\subsubsection*{(2b) The large-field operation: which variables are frozen, which are integrated out, and what a patch retains}

\textbf{Measure and formula: two layers.} At the level of measures, a step acts as a finite sequence of moves on states $\sigma=(h,\xi)$, where $h$ records the outcomes obtained up to that point and $\xi$ the variables stored up to that point. Each move is given by positive kernels $\Pi^{(o)}$, finitely many and one for each outcome $o$, such that $\sum_o\Pi^{(o)}(\sigma,\cdot)$ has total mass one. A move is one of four kinds. It may be a test, that is, an indicator whose outcome is added to $h$. It may be a map. It may be a draw, that is, the image of a probability law, as for gauge fixings, hypercubic frames $r\in W_4$ and reference laws. Or it may delete a stored variable, which amounts to an integration. Everything expressed by a formula, such as Gaussian weights, Jacobians, expansions, resummations or the replacement of a test by $1$, only recasts an integrand in another form. A gauge fixing is a draw as a move. For it to be the factor that the formula writes, the measure must be gauge invariant. Let $\nu_0$ be any complex bare measure, let $\mu_0\ge0$, and let $\mathfrak{w}$ be any bounded Borel function of the bare field. Denote by $\mathsf{K}_h$ the kernel of the history $h$, obtained by composing the kernels $\Pi^{(o)}$ of its moves along its outcomes. Then

\begin{gather*}
\FitMath{\tau_h[\nu_0]=\int\nu_0(dU)\,\mathsf{K}_h(U,\cdot),} \\[3pt]
\FitMath{\sum_h\mathsf{K}_h(U,\mathcal{V}_k)=1\,\,\forall U,} \\[3pt]
\FitMath{\sum_h|\tau_h|(\mathcal{V}_k)\le|\nu_0|(\mathcal{V}_0),} \\[3pt]
\FitMath{\big|\tau_h[\mathfrak{w}\mu_0]\big|\le\sup|\mathfrak{w}|\,\tau_h[\mu_0].}
\end{gather*}

By totality (2a; a reference law whose normalisation does not lie in $(0,\infty)$ is not drawn, and in that case the component drops out; this case does not arise on the supports, see (2e), ``Reality and positivity''), every kernel is defined at every field, and no smallness condition is needed. They depend on the lattices and the constants. Of the frame $\theta$ (2e: the real sequence from which the choices are formed; in the present case $\theta=x$, and it is only through this frame that the kernels then vary with $\beta$; it is unrelated to the hypercubic frames $r$), they use only the gauge-fixing law, the coupling of the reference law, the cube sizes and the thresholds. Apart from their dependence through the frame, the kernels do not depend on any measure, on $\beta$, on a source, on any term of an effective action, or on whether some series converges. Exactness and positivity therefore come at no cost, because a move does nothing but split mass. The sourced densities come from $\nu_0=\rho^w_0\,dU$ with the same kernels (4a), and runs whose choices are frozen are runs of that same procedure (3c). The estimates of (2c)--(2e) add only one further statement: the formulas are formed on the supports and there coincide with the measures. \emph{Passengers}: divide the stored variables into those, $\mathfrak{v}$, that a move uses and the remaining ones, $\mathfrak{q}$. In step $t=k+1$ the operation $\mathbf{T}$ uses the fields of level $k$ and, from the lower levels, only the data of its backgrounds; the operation $\mathbf{R}$ uses variables only in the components of its class and next to them. Any variable not used by a move of the step is carried along unchanged, 

\begin{gather*}
\FitMath{\Pi^{(o)}((h,(\mathfrak{v},\mathfrak{q})),\cdot)=\Pi^{\flat(o)}((h,\mathfrak{v}),\cdot)\otimes\delta_{\mathfrak{q}}}
\end{gather*}

, in which $\Pi^{\flat(o)}((h,\mathfrak{v}),\cdot)$ is a measure on the states $((h,o),\mathfrak{v}')$. This is the locality of a step, in measure-theoretic form.

\textbf{$\mathbf{T}$ on a label} (Ba{\l}aban, Comm. Math. Phys. 119 (1988) 243, Sect. 3) comprises the following, carried out on the regions fixed by the tests of (2a). First the multi-scale background $U_{j+1}(\mathbf{V})$ is formed, on the determining set $\mathbf{B}[\Omega_\bullet]$ (1b); the third test of (2a) is formed with this background. Next, a change of variables is made on the strip. Then a conditional Gaussian integral is taken over $\Lambda_{j+1}$. Finally, the logarithm of this integral is split into $E$, $R$ and $B$. The operation sits between (1d), in which no regions occur, and the operation $\mathbf{R}$.

\textbf{The strip.} The tests kept on $S_{j+1}\subset\Omega_{j+1}\setminus\Lambda_{j+1}$ depend on $V_{j+1}$ non-locally, by way of $V^{\mathrm{bg}}_j$. Take a background localised to $\Lambda_{j+1}^c\cap\Lambda_j$ and denote its average by $V^{\mathrm{loc}}_j$. Then

\begin{gather*}
\FitMath{V^{\mathrm{bg}}_j(V^{\mathrm{loc}}_j)^{-1}=e^{i\mathbf{H}^{(j)}},} \\[3pt]
\FitMath{|\mathbf{H}^{(j)}|\le C\varepsilon_je^{-\check{R}_j};} \\[3pt]
\FitMath{V_j(b)=W'e^{-i\vartheta_{\mathrm{co}}(W')\mathbf{H}^{(j)}(b)}e^{i\mathbf{H}^{(j)}(b)}\,V^{\mathrm{loc}}_j(b),}
\end{gather*}

where the second relation is a substitution made bond by bond, with new variable $W'=e^{i\mathcal{A}'(b)}$. In it, $\vartheta_{\mathrm{co}}$ is a smooth cut-off in $|\log W'|$. It equals $0$ up to slightly above $\delta_j$ and $1$ from slightly below $2\delta_j$ onwards. This cut-off is applied on the cut-off zone of the strip ((2a)); at the remaining bonds where the substitution is made, $\vartheta_{\mathrm{co}}\equiv1$. Consequently, on the cut-off zone, a test whose threshold does not exceed $\delta_j$ is the same function of the new variable as of the old one, while a test at $2\delta_j$ is now made against the localised background. Because $|\mathbf{H}^{(j)}|\ll\delta_j$, a contraction argument shows that the map on a bond is a diffeomorphism of $G$. Its Jacobian equals $1$ wherever $\vartheta_{\mathrm{co}}$ is locally constant, and it contributes to the action as a boundary term. Once the substitution has been made, the factors of the kept tests at those bonds next to $\Lambda_{j+1}$ where no substitution takes place are omitted. On the supports these factors equal $1$, and where $\vartheta_{\mathrm{co}}<1$, a test at $2\delta_j$ has a factor equal to $1$ whichever background is used. Moreover, at a pivot bond of $\Omega_{j+1}$ the value of $V_j$ is taken from the constraint $\mathcal{M}(V_j)=V_{j+1}$. With this, the tests depend on $V_{j+1}$ only in $\Lambda_{j+1}^c$, and the integral operations therefore factorise ((1.21)--(1.25), (3.22) there).

\textbf{What is frozen and what is integrated.} On $\Omega_{j+1}^c$, step $j+1$ stores $V_j$ together with its constraint. On $\Omega_{j+1}$ the action is expanded around $U_{j+1}(\mathbf{V})$. The only variable integrated there is $A=\mathcal{A}^{\mathrm{s}}_j|_{\Lambda_{j+1}}$, conditionally on the stored field $A'=\mathcal{A}^{\mathrm{s}}_j|_{\Omega_{j+1}\setminus\Lambda_{j+1}}$. The measure is Gaussian with covariance $C^{(j)}(\Lambda_{j+1})$, which is the inverse of the restriction of the form $\Delta^{(j)}$ to the bonds of $\Lambda_{j+1}$, and its linear term is $-\langle A,\Delta^{(j)}A'\rangle$. The coupling flows only where this integration is carried out, $\varphi_{j+1}\in C_0^\infty(\Lambda_{j+1})$, and it stops on the whole of $Z_{j+1}$, including $\Omega_{j+1}\setminus\Lambda_{j+1}$. Once the square is completed, the conditional weight appears as a factor inside $\mathbf{T}_k(Z)=\prod_{j<k}\mathbf{T}^{(j)}$, with the regions cut by the component $Z$ ((2.21), (3.23) there; this operation is the product shown for a component on which $\mathbf{R}$ has never acted, and in general the operation is defined by a recursion over the marks):

\begin{gather*}
\FitMath{\mathbf{T}^{(j)}=\int dV_j\big|_{\Omega_{j+1}^c}\delta\big(\mathcal{M}(V_j)V_{j+1}^{-1}\big)\,\mathfrak{G}_j\,\chi\chi^c\,z^{(j)}\!\int dA'\,\chi\chi^c\,e^{-\frac12\langle A',[\Delta^{(j)}-\Delta^{(j)}C^{(j)}(\Lambda_{j+1})\Delta^{(j)}]A'\rangle}.}
\end{gather*}

This operation contains no Wilson weight. That weight belongs to $\exp A_{k,\mathcal{S}}$, the function on which the operation acts; $z^{(j)}$ is explained below. The old $R$-terms, $\mathbf{R}'_j=\sum_{i\le j_1}\mathbf{R}^{(i)}$, appear alongside $\mathbf{E}_j$ in the regular action $\widetilde{\mathbf{E}}_j$ of the step (2e). Here $\mathbf{R}^{(i)}=\sum_X[R^{(i)}(X,\cdot)-R^{(i)}(X,\mathbf{1})]$ is the summand of level $i$ of $\mathbf{R}_j$, and $j_1$ is the index $k_1$ of (2d) at level $j$. The young ones, $\mathbf{R}''_j=\sum_{j_1<i\le j}\mathbf{R}^{(i)}$, are handled separately. Put $\delta F:=F(U_j)-F(U_{j+1})$. Attach a parameter $t$ to the terms present only next to $Z_j$, and a parameter $s$ to $\mathbf{R}''_j$ ((3.23)--(3.28) and pp. 270--271 there). Let $z^{(j)}=\prod_cz^{(j)}(c)$ be the factors that come from eliminating the delta functions. In $\mathbf{T}^{(j)}$ their product runs over those coarse bonds $c$ of $\Omega_{j+1}$ that lie outside $\Lambda_{j+1}$; under the logarithm it runs over those of $\Lambda_{j+1}$. Then

\begin{gather*}
\FitMath{\log\Big[z^{(j)}\!\!\int\!dA\,\chi^{(j)}e^{-\frac12\langle A,\Delta^{(j)}A\rangle-\langle A,\Delta^{(j)}A'\rangle+\cdots}\Big]-\tfrac12\big\langle A',\Delta^{(j)}C^{(j)}(\Lambda_{j+1})\Delta^{(j)}A'\big\rangle} \\[3pt]
\FitMath{=\int_0^1\!dt\,\big\langle\delta\mathbf{B}_j-\delta A(x_j(\cdot)-x_j,\cdot)+\cdots\big\rangle_{s=1,t}+\int_0^1\!ds\,\langle\delta\mathbf{R}''_j\rangle_{s,t=0}+\Big(\log[\cdots]-\tfrac12\big\langle A',\Delta^{(j)}C^{(j)}(\Lambda_{j+1})\Delta^{(j)}A'\big\rangle\Big)_{s=t=0}.}
\end{gather*}

In this identity $\langle\,\cdot\,\rangle_{s,t}$ denotes the normalised mean with respect to the measure that the parameters $s$, $t$ interpolate. The left-hand side is the unnormalised integral. Its second term is present since the square, once completed, was moved into $\mathbf{T}^{(j)}$. Inside the integral over $t$, the dots stand for two kinds of terms. Those of the first kind are the terms of the strip, namely the Jacobians together with the differences of all quantities that depend on a shifted variable. Those of the second kind are the differences, taken between the two backgrounds, of the exponents of the integral operations carried by the term. The coupling difference carries the sign it has in $A_{j,\mathcal{S}}=-x_jA-A(x_j(\cdot)-x_j,\cdot)+\dots$. A cluster expansion as in (1d) is applied to each part. The Gaussian normaliser is contained in the last part; it is that part of $E^{(j+1)}$ which carries $b^0_{j+1}$. On a label it is a determinant: that of a multi-scale operator, compressed to $\Lambda_{j+1}$ and taken in the background of that label; so by definition it is non-local. It is converted into local terms,

\begin{gather*}
\FitMath{\sum_c\log z^{(j)}(c)-\tfrac12\log\det\big(\mathsf{C}^{\mathsf{T}}\Delta^{(j)}\mathsf{C}\big|_{\Lambda_{j+1}}\big)=\text{number}+\sum_X\mathbf{N}(X,U_{j+1}),} \\[3pt]
\FitMath{|\mathbf{N}(X)|\le Ce^{-cMd_{j+1}(X)},}
\end{gather*}

which Ba{\l}aban does by means of a resolvent integral and walks (1d), whereas here it is done by sums with alternating signs, taken over the sub-regions $Y\subset X$, whose terms are the diagonal entries of the logarithm of the operator built on $Y$. M{\"o}bius inversion yields the sum rule, and the decay is due to a sign-reversing involution, which cancels everything except the chains of walks that visit all of $X$. Only the last two parts give terms with domain $X\subset\Lambda_{j+1}$. These terms are built from walks and fields inside $X$, so the regions do not affect them. They are of two kinds. The first is $R^{(j+1)}$; it arises from the part with $s$, to which what $\mathbf{R}$ leaves is added (2d). The second is $E^{(j+1)}$, which comes from the last part, where $\mathbf{E}_j$ stands with the old $R$-terms. These terms are the same for every label. For the regular terms this holds for every such $X$. For the remainders it holds for $X$ one cell inside $\Lambda_{j+1}$, the only place where a label shows them, and this is the margin of (2e). Thus the young $R$-terms go into $R^{(j+1)}$, and the old ones go into $E^{(j+1)}$ and $b_{j+1}$ (2d, 2e). Everything else makes up $B^{(j+1)}$; among it is the part of the quadratic form that is rejected (next paragraph). It amounts to $O(1)$ for each $M$-cube of level $j+1$ next to $Z_{j+1}$; this cost is covered by (2c), and for that reason no renormalisation of it is required.

\textbf{What a one-step operation keeps.} The conditional form is a Schur complement. It is a sum over chains of operators with connected carriers, where a chain is either a walk term of $\Delta^{(j)}$ or a sequence of a walk, a resolvent walk and a walk. A chain is kept if the cells of level $j+1$ that it meets intersect $\Omega_{j+1}\cap Z_{j+1}$ in a non-empty connected set and remain within one layer of that set; a kept chain is then counted with the component $\mathfrak{k}$ of $\Omega_{j+1}\cap Z_{j+1}$ that contains this set:

\begin{gather*}
\FitMath{\big\langle A',[\Delta^{(j)}-\Delta^{(j)}C^{(j)}(\Lambda_{j+1})\Delta^{(j)}]A'\big\rangle=\sum_{\mathfrak{k}}\langle A',\Delta^{\mathrm{loc}}_{\mathfrak{k}}A'\rangle-2\sum_X\mathsf{B}^\times(X),} \\[3pt]
\FitMath{|\mathsf{B}^\times(X)|\le Cp_0(g_j)^2e^{-cM\check{R}_{j+1}}e^{-\kappa d_{j+1}(X)},} \\[3pt]
\FitMath{\langle A',\Delta^{\mathrm{loc}}_{\mathfrak{k}}A'\rangle\ge c\sum_{b\text{ in }\mathfrak{k}}|A'(b)|^2:}
\end{gather*}

(the bound on $\mathsf{B}^\times$ is for stored fields in range, $|A'(b)|\le Cp_0(g_j)$, which they are on the support of the kept tests) a rejected chain has length greater than the size of a cell, and the size $p_0(g_j)^2$ of the stored field squared is outweighed by $e^{-cM\check{R}_{j+1}}$. The rule of Ba{\l}aban, by which the form ``couples only the fields $A_j$ in the same component'' (119, p. 258), fails to be local: whether two bonds belong to a common component is settled at a distance from them, and two histories having the same trace can differ on it. The kept form is not a Schur complement but a truncated series. Its lower bound is therefore not inherited and has to be proved. The proof localises to small cubes (Ismagilov--Morgan--Simon) and compares with a genuine Schur complement. \emph{Carriage}: in an integral operation that a term carries, the tests, the gauge factors and the delta functions depend on no background, whereas the forms $\Delta^{\mathrm{loc}}_{\mathfrak{k}}$ and the factors $z^{(j)}$ do. These forms and factors are given as functions of a pair $(\mathbf{U},\mathbf{J})$ near $\mathfrak{k}$, and they are constructed once and for all. They are evaluated at the current background, which in a term of level $k$ is $U_k(\mathbf{V})$. The background changes with every later step and with every stage of $\mathbf{R}$. Each change produces differences 

\begin{gather*}
\FitMath{-\tfrac12\langle A',[\Delta^{\mathrm{loc}}_{\mathfrak{k}}(U^\sharp)-\Delta^{\mathrm{loc}}_{\mathfrak{k}}(U^\flat)]A'\rangle}
\end{gather*}

, where $U^\flat$ is the background after the move and $U^\sharp$ the one before it, both written in the new variables. Localisation turns these differences into boundary terms. For a piece of a kept form whose carrier is $R$, the oscillation with respect to the pair is bounded by $Cp_0(g_j)^2e^{-\kappa d^\natural_{j+1}(R)}$, measured in the hull length of the carrier; this bound is large, being the square of a field at its threshold. The difference of such a piece at the two backgrounds of a stage of level $i$ is localised as well, and each of its localised pieces, on a domain $Y$, has modulus at most $C(\log x_i)^{-(q-3\mathfrak{p}_0-Cr)}e^{-c\,d^\natural_k(Y)}$. This is a negative power, because $q$ is chosen after $\mathfrak{p}_0$ and $r$, with $q>3\mathfrak{p}_0+Cr$.

\textbf{The idea of $\mathbf{R}$} (Comm. Math. Phys. 122 (1989) 175, (0.3)--(0.4)). We split $Z=Z'\cup Z''$. Following that paper (p. 176), $Z''$ is the union of the components ``for which the small factors connected with large field control some number of next steps'', and $Z'$ is the union of all the other components (the rule by which the present construction sorts the components into these two sets is given under ``Where it acts'', below). The density is a sum $\rho(V)=\sum_Z\rho(Z,V)$ over the large-field regions $Z$ ((0.2) of that paper), and $\rho(Z'',V)$ denotes the density belonging to the region $Z''$. Then

\begin{gather*}
\FitMath{(\mathbf{R}\rho)(V)=\sum_Z\rho(Z'',V)\,\frac{\int dV|_{Z'}\,\rho(Z,V)}{\int dV|_{Z'}\,\rho(Z'',V)},} \\[3pt]
\FitMath{\int\mathbf{R}\rho=\int\rho .}
\end{gather*}

In $Z'$ the law of the field is replaced by the small-field law, and the mass is preserved. This idealised form is not linear in $\rho$. The form actually used is linear, because its reference law contains the Wilson weight and nothing else.

\textbf{Where it acts.} Let $k_\flat=k-\mathsf{N}_k\ge1$ (\emph{depth}); this restriction excludes nothing at $k=\mathsf{N}_k$, since by the second condition below a member of the class is then a component of $Z_k$ in which no step has created anything, and no such component exists. A component $X$ of $Z_k$ belongs to the class when three conditions hold. First, for $k_\flat\le i\le k$ the set $X\cap Z_i$ is a single component of $Z_i$ and is small, that is, contained in a cube of $O(1)$ level-$i$ cells (\emph{single boxes}, and at $i=k$ \emph{size}). Second, for $k_\flat<i\le k$, step $i$ creates nothing in it, and it does not absorb a second component, $X\cap Z_i=\Phi_i(X\cap Z_{i-1})$ (\emph{quiet}). Third, it has not been acted on by any operation whose level lies in $[k_\flat,k)$, the lower end point included (\emph{quarantine}). Ba{\l}aban imposes two conditions, that the component be small and quiet (ibid., p. 177); to these two, the exclusion of a joining is added here as a separate condition, named explicitly. Two small quiet boxes that join within the window form a small component without creation, yet its trace one scale lower has two components: a joining does not count as a creation. Let $H_i$ denote the hull, in level-$i$ cells, of $X\cap Z_{i-1}$, and let $B^{[r]}$ be the box $B$ with each face pushed outward by $r$ such cells. Then

\begin{gather*}
\FitMath{\Omega_i^c\cap X=H_i^{[8]},} \\[3pt]
\FitMath{Z_i\cap X=H_i^{[10]}} \\[3pt]
\FitMath{(k_\flat<i\le k):}
\end{gather*}

the boxes are nested, the two of each level being concentric, and the number of cells along a side of $X$ is comparable to the fixed point of the free evolution (2c). As long as $k\le\mathsf{N}_k$ the class is empty, so $\mathbf{R}$ does nothing at the first levels. For $k>\mathsf{N}_k$ the class is non-empty: a single cell created at step $k-\mathsf{N}_k$, with nothing else, belongs to it. Up to level $k$ the torus is admissible, $\tfrac12\ell L^{K-i}\ge CM\check{R}_i$ for $i\le k$, as the opening of Theorem 2.1 requires. For a first-passage datum, $x_K\le g^{-2}+B_\sharp$, this condition follows from $\ell\in\mathfrak{L}(g)$ by way of $\check{R}_i\le L^{K-i}\check{R}_K$, a consequence of the upper flow bound; it is a lower bound on $\ell$. Hence no cell, no box and no neighbourhood of a box wraps around the torus.

\textbf{The operation.} (0) \emph{The frame sum}, described first although it belongs inside (3). In each component the identity $1=|W_4|^{-1}\sum_{r\in W_4}1$ is inserted after the integration over the box and before the tree gauge of (3). Consequently no frame sum is attached to a component that leaves at (2) or during the stages. In the summand of $r$, every object of (3), (4) and (2d) that belongs to $X$ is formed in the frame $r$: the reference law with its denominator, the two trees, the test on $B$, and the last-shell test. If a component leaves after the tree gauge, its frame sum stays with it, and an island keeps, in each summand, the cut-off and the delta function belonging to the frame of that summand. Formed instead in a single frame $r_0$, these make the bracket of $X$ integrate to $|W_4|^{-1}\sum_r\mathfrak{a}_r\,\mathsf{D}^{(r_0)}_X/\mathsf{D}^{(r)}_X$ instead of $|W_4|^{-1}\sum_r\mathfrak{a}_r$. Here $\mathfrak{a}_r$ denotes that part of the summand of $r$ which is independent of the box, and $\mathsf{D}^{(r)}_X$ denotes the denominator of $\mathfrak{s}_X$ formed in the frame $r$. The reason for keeping the frame sum is that neither the test on $B$ nor the last-shell test is gauge invariant. Each of them enters only by way of decompositions of unity performed within a single gauge-fixed measure. Let $\mathsf{m}$ be the measure before the tree gauge, and let $\mathsf{a}_r,\mathsf{b}_r$ be its parts that pass, respectively fail, these tests in the frame $r$. Then $\mathsf{a}_r+\mathsf{b}_r=\mathsf{m}$ for every $r$, and hence $|W_4|^{-1}\sum_r(\mathsf{a}_r+\mathsf{b}_r)=\mathsf{m}$. With the drop-outs kept instead in a single fixed frame $r_0$, the total is $\mathsf{m}+(|W_4|^{-1}\sum_r\mathsf{a}_r-\mathsf{a}_{r_0})$. The average is taken so as to obtain covariance under the lattice isometries; trees break this covariance. (1) \emph{Stages}: for $j=k_\flat,\dots,k-1$, one value after another, the stored variables of the outer shell of $Z_{j+1}\cap X$ are integrated out as in a small-field step. Before this, tests at the sharper threshold $\delta'_j:=g_jA_1p_1(g_j)$ are made, formed with backgrounds localised in $X$. What remains is a box $\Lambda^{\mathrm{pr}}$ of side at most $CM$: it is the set $Z''_k(X)$ of the nest of (2a) enlarged by one more layer of $M$-cubes, and this layer is the shell of level $k$. In it the levels $k_\flat,\dots,k$ are stacked in thin shells, the shell of level $j$ carrying the field $V_j$. Together these fields form the datum $V''$ attached to the prepared determining set $\mathbf{B}''$; its background is $U''_k$. At the innermost position lies the core $C_{k_\flat}=X\cap Z_{k_\flat}$, whose side, measured in the unit of level $k$, is at most $CM\check{R}_{k_\flat}L^{-\mathsf{N}_k}$. (2) The test asks whether the exterior field $V_k$ admits an extension into $\Lambda^{\mathrm{pr}}$ whose curvature is less than $2\varepsilon_k$. If an exterior field satisfies $|\partial V_k-\mathbf{1}|<\varepsilon$, then it does admit an extension, one with $|\partial V_k-\mathbf{1}|<CM^2\varepsilon$. To see this, set the field equal to $\mathbf{1}$ inside after choosing an axial gauge on the surface $\partial^+\Lambda^{\mathrm{pr}}$. Hence ``no regular extension'' means that some exterior plaquette has $|\partial V_k-\mathbf{1}|\ge cM^{-2}\varepsilon_k$, and this pays as in (2c) (p. 193 there). (3) $V_k$ is integrated over $\Lambda^{\mathrm{pr}}$. Taken together, the parts from the test (2) up to the division into the two classes (2d) form the \emph{main step} of the operation; the stages that precede it and the exponentiation that follows it are counted separately. This integration is the first move to alter the density. Integrals against any bounded $F$ that is independent of $V_k$ on the box are left unchanged by it. Level by level, from the top down, it resolves the bare chain $\prod_j\int dV_j\,\delta(\mathcal{M}(V_j)V_{j+1}^{-1})$, $\int dV_i(c)\,\delta(\mathcal{M}(V_{i-1})(c)V_i(c)^{-1})=1$. What remains is a free Haar integral,

\begin{gather*}
\FitMath{\chi_k\,\chi_{k,\Lambda}\int dV_{k_\flat}\big|_{C_{k_\flat}}\chi\;\mathbf{T}_{k_\flat}(C_{k_\flat})\int dV''\big|_{\mathbf{B}''\cap(\Lambda^{\mathrm{pr}}\setminus C_{k_\flat})}\,\chi''_k\,e^{A''_k(U''_k)}}
\end{gather*}

Here $\chi_{k,\Lambda}$ is the test of (2); $\chi$ stands for the plaquette tests of level $k_\flat$ carried by the term on $\Omega_{k_\flat}\cap C_{k_\flat}$; and $\chi''_k$, $A''_k$ are the tests and the action that (1) leaves. The integral is regarded as a density that is constant in $V_k$ on the box. The field that is expanded is $V''$. It is expanded around the averages of the minimiser $U_0$ below, each average taken at its own level:

\begin{gather*}
\FitMath{V''=e^{ig_kB}V_0\,\text{ on the fluctuation bonds }\mathbb{B}_0=\mathbf{B}''\cap(\text{shells}),} \\[3pt]
\FitMath{V_0=\big(\mathcal{M}^jU_0\text{ on the shell of level }j\big)_{k_\flat\le j\le k}.}
\end{gather*}

Thus $B$ is a multi-scale field on bonds of every level from $k_\flat$ to $k$, with a single scaling $g_k$ and a single test $|B|<\delta'_k/g_k$. This is the field $B$ that occurs in Proposition 2.8. The gauge of this field is fixed along a tree $T_0$ in $\mathbb{B}_0$, whose bonds come from all these lattices, one lattice for each shell. For two nested boxes $P_1\supset P_2$, the tree is made up of axial contours which leave from a single initial site and go around $P_2$. The family of trees has two properties (pp. 195--196 there). First, the family is stable. Suppose one enlarges $P_2$ to a box $P'_2\subset P_1$ and deletes the contour of each site of $P'_2\setminus P_2$; then what is left is the tree of $P_1\setminus P'_2$. Ba{\l}aban describes this property as important, and boxes are needed because of it. Second, the tree of a shell uses up the gauge freedom at every site of the shell except one. For this reason, for each inner box on which a gauge is fixed, one external bond crossing its boundary is added, and one more is added for the outermost box. Through these extra bonds the trees of the shells become a single tree, and it fixes the gauge entirely, with Jacobian one. The core operation $\mathbf{T}_{k_\flat}(C_{k_\flat})$, which contains the failed tests, is at no point unfolded. One further test, the \emph{last-shell test}, compares with the vacuum and constrains $V_{k_\flat}$ between the core and $\mathbb{B}_0$. It requires 

\begin{gather*}
\FitMath{|\partial U_{k_\flat,\square}((\mathbf{1},V_{k_\flat}))-\mathbf{1}|<\tfrac12\varepsilon_{k_\flat}\eta_{k_\flat}^2}
\end{gather*}

 on the test cubes of that region, where $(\mathbf{1},V_{k_\flat})$ equals $\mathbf{1}$ on the side of $\mathbb{B}_0$. (4) For each term $\tau$, its mass $m_\tau=\int\tau\,dV_k|_{\Lambda^{\mathrm{pr}}}$ multiplies the reference law. The level-$k$ variable $V'_k:=V_kV_\Lambda^{-1}$, where $V_\Lambda$ is the element of the critical orbit of Proposition 2.4 that equals $\mathbf{1}$ on the tree of $\delta_{\mathrm{ax}}$, occurs only in this law and in its denominator. Let $U_{k,X}$ denote the background localised in $X$, and let $\chi(\Lambda^{\mathrm{pr}})$ be the cut-off $\mathbf{1}\{|\tfrac1i\log V'(b)|<M_0\varepsilon_k$ on the bonds of the box$\}$, in which $M_0$ is a constant chosen large after $M$. Let $\delta_{\mathrm{ax}}$ be the tree gauge of $\mathbb{T}^{(k)}\cap\Lambda^{\mathrm{pr}}$, taken on a tree that contains all the sites of this set together with the external bond of $T_0$ into the box; that bond is part of the tree because, without it, a constant gauge transformation of the box is left unfixed. Let $\phi_X$ be a cut-off function that equals $1$ well inside $X$. Then

\begin{gather*}
\FitMath{\mathfrak{s}_X(V_k)=\frac{\delta_{\mathrm{ax}}(V'_k)\,\chi(\Lambda^{\mathrm{pr}})(V'_k)\,e^{-x_kA(\phi_X,U_{k,X}(V'_kV_\Lambda))}}{\int dV'|_{\Lambda^{\mathrm{pr}}}\,\delta_{\mathrm{ax}}(V')\,\chi(\Lambda^{\mathrm{pr}})(V')\,e^{-x_kA(\phi_X,U_{k,X}(V'V_\Lambda))}}.}
\end{gather*}

If one of the tests of (1)--(3) fails in a component, that component drops out (2a): it remains live, and the failure counts as a creation. This can happen to the same component repeatedly. Each occurrence counts as a creation and is followed by a wait of at most $\mathsf{N}_k+1$ steps.

\textbf{The purpose of the stages, and the reason for a sharper threshold.} Continuing the field by least action costs a constant factor: for exterior data with $|\partial V_k-\mathbf{1}|<\varepsilon$ one has $|\partial V_\Lambda-\mathbf{1}|\le C(M)\varepsilon$, and $|\partial U_0-\mathbf{1}|\le C(M)\varepsilon_k\eta_k^2$ near the box. For every value of the coupling, the stages reduce the component from side $M\check{R}_k$ to a box of side $O(M)$. Hence $C(M)$ is a constant fixed before $\gamma$. On the shell of level $j$ it is overcome by the factor $L^{-(k-j)}$ by which a level-$k$ field scales when viewed from level $j$. On the outermost shells it is overcome by $\check{R}_k^{-c}$, because $1+\log_L\check{R}_i$ levels are compressed into the box, namely the levels from the last level $i$ with $L^{k-i}=\check{R}_i$ up to level $k$. Consequently $U_0$ belongs to the domain of the carried tests. It also belongs to the analyticity domains of the terms, insofar as the domain of a term lies inside the component or at a distance from the box; the remaining terms are treated in (2d). In the words of 122 (1989) 175, p. 195, ``it is the reason why we have done the preliminary integrations''. As in Lemma 2.2, the tests made at a stage have to imply the earlier ones. Here this must happen $\mathsf{N}_k=\check{R}_k$ times in succession, and the losses together amount to a fixed fraction. If two backgrounds satisfy the same constraints on a region, their difference is the response to the difference of their data, and this response decays. Let $U^{(n)}_{k,X}$ denote the background after $n$ stages, localised in $X$. Put $j=k_\flat+n$, and let $\beta_{\mathrm{s}}$ be a fixed small fraction, written $\beta$ by Ba{\l}aban. Then, on the support of the tests of stage $n+1$ and on the shell of level $i\le j$,

\begin{gather*}
\FitMath{|\partial U^{(n)}_{k,X}-\mathbf{1}|\le|\partial U^{(n+1)}_{k,X}-\mathbf{1}|\big(1+c\beta_{\mathrm{s}}2^{-(j-i)}\big)+c\beta_{\mathrm{s}}2^{-(j-i)}\varepsilon_i\eta_i^2,}
\end{gather*}

where the error term is 

\begin{gather*}
\FitMath{C\delta'_je^{-\delta_0\operatorname{dist}}\eta_i^2\le C\,\tfrac{\delta'_j}{\varepsilon_j}(1+(j-i)^{1/2})e^{-c(M/M_1)(j-i)}\varepsilon_i\eta_i^2}
\end{gather*}

. For $M$ large, the decay $e^{-c(M/M_1)(j-i)}$ across the shells dominates $1+(j-i)^{1/2}$, and $2^{-(j-i)}$ remains. The ratio $\delta'_j/\varepsilon_j=(A_1/A_0)(\log x_j)^{-(\mathfrak{p}_0-\mathfrak{p}_1)}$ dominates the fixed constants, and at this point smallness of $A_1/A_0$ suffices as well. After $n$ stages the test requires $|\partial U^{(n)}_{k,X}-\mathbf{1}|<\big(1-\beta_{\mathrm{s}}(1-2^{-(j-i+1)})\big)\varepsilon_i\eta_i^2$ on that shell. On the outermost shells, where the collars of the stages overlap, this requirement is relaxed by a factor $L_\flat^2$ for each overlap, where $L_\flat$ is an integer with $2\le L_\flat<\frac12L$ that is chosen together with $L$; such a choice requires $L\ge5$. Let $k_\sharp+1$ be the last of the levels at which an $M\check{R}$-cube, measured in units of level $k$, has side precisely $M$; in formulas, $M\check{R}_{k_\sharp+1}L^{k_\sharp+1-k}=M$. The collars of the stages with $j\le k_\sharp$ are disjoint, and the collars of the later stages overlap. On the shell $i$ the factor is $L_\flat^{2(i-k_\sharp-1)_+}$. It depends on the shell only, not on the stage, and it becomes larger with each successive shell until it reaches $\check{R}_k^{\,c}$. In the region not covered by the shells, the bound is expressed in the scale $j$ of the stage and equals $(1-\tfrac12\beta_{\mathrm{s}})\,c_l\,L_\flat^{2e}\varepsilon_j\eta_j^2$, where $c_l=3$ adjacent to $\Omega_k$ at the last stage and $c_l=1$ elsewhere. Far from the box, the curvature of the least-action continuation $U_0$ can be as large as $(2+c)\varepsilon_k\eta_k^2$, and the test of the last stage must be implied there as well. Here $e=0$ for $j\le k_\sharp$. For $j>k_\sharp$ one takes $e=(j-l-1)_+$ on $\Omega_l\setminus\Omega_{l+1}$, $k_\sharp<l\le j$, and $e=j-k_\sharp-1$ between the shells and $\Omega_{k_\sharp+1}$. This is the form of 122 (1989) 175, (1.24), where the level is denoted $k_0$ and the base $L_0$. On a layer $l<j$ this bound remains below $\varepsilon_l\eta_l^2$, the scale of that layer. In passing from the scale $j+1$ to the scale $j$ one gains $L^{-2}$ and pays $L_\flat^2$. Hence the bound of each stage implies the bound of the preceding stage. The margin between consecutive stages pays for the error, which decays geometrically. The sharper threshold enters in two places. In the first, the new fluctuation test has to imply the old fluctuation tests of step $j+1$, at $\delta_j$ and $2\delta_j$. This holds because $CL\delta'_j=CL(\log x_j)^{-(\mathfrak{p}_0-\mathfrak{p}_1)}\delta_j\le\tfrac12\delta_j$ for $\gamma$ small, and it puts no new condition on the constants chosen earlier. In the second, at the main step, $\delta'_k$ has to remain a small fraction of $\varepsilon_{k_\flat}$, which is the curvature threshold of the core level. In this estimate the thresholds of two levels may differ by $\varepsilon_j/\varepsilon_i\le C(1+(j-i)^{1/2})$, the bound valid for every pair of levels (2e). Restrict to the common support of the carried tests $\chi_k$, the tests of (2) and (3), and the tests and delta functions that stand under the core operation; the latter make the data stored in the core regular. On this set, and provided $C(M)\check{R}_k^{-c}\le c$ on the outermost shells, every test that the stages leave behind equals $1$ when read on $X$ minus one layer of level-$k$ cells. The equality does not reach $\partial X$: adjacent to $X^c$, the background localised in $X$ is constrained bond by bond on the bare lattice. There the test therefore demands $|\partial V_k-\mathbf{1}|<C\varepsilon_k\eta_k^2$ of the block field itself, while the carried tests provide only $2\varepsilon_k$. All of this holds under the condition

\begin{gather*}
\FitMath{C\,\mathsf{N}_k^{1/2}\,\frac{p_1(g_{k_\flat})}{p_0(g_{k_\flat})}+C(M)\,L^{-\mathsf{N}_k}\mathsf{N}_k^{1/2}\le c:}
\end{gather*}

the first term is $\delta'_k$ measured relative to $\varepsilon_{k_\flat}$, including the allowance $\mathsf{N}_k^{1/2}$. Since $\mathsf{N}_k=\check{R}_k\approx(\log x_k)^r$, it is a condition on the exponents, $\mathfrak{p}_1<\mathfrak{p}_0-Cr$, together with a ceiling on $g_k$. The second term is a level-$k$ field viewed from level $k_\flat$, so $\mathsf{N}_k$ is taken large.

\textbf{How a patch recovers its coupling.} Each stage gives back to the collar it integrates the increments that this collar did not receive. After $m$ stages the Wilson term is formed with

\begin{gather*}
\FitMath{x^{(m)}(\cdot):=\beta-\sum_{i\le k}b_i\,\varphi\big[\Lambda^{(m)}_i\big](\cdot),} \\[3pt]
\FitMath{\Lambda^{(m)}_i=\Lambda_i\cup(X\setminus Z''_i)\,\,(k_\flat<i\le k_\flat+m),} \\[3pt]
\FitMath{\Lambda^{(m)}_i=\Lambda_i\,\text{otherwise}}
\end{gather*}

where $\varphi[\Lambda]$ is the profile of the region $\Lambda$. The difference $A(x^{(m)}-x^{(0)},\cdot)$ is carried separately, as a Wilson piece (2d). In this way, at each bare plaquette the weights of all the Wilson terms in the exponent still sum to $-\beta$ (2a). Hence $x''_k=x^{(\mathsf{N}_k)}$ coincides with $x_k$ on the part of $X$ outside the box, and equals $x_k+O(1+k-j)$ on the shell of level $j$. This is what makes the three Wilson terms of Proposition 2.8 [outlined below; Section~\ref{sec:actual-history-closure}] cancel; two of them carry the constant $x_k$.

\textbf{What a patch retains.} A component that has passed every test carries the operation below. It is (1.71) of 122 (1989) 355, in which the summation extends over $r$ and over all sequences of regions contained in $X$, with the prefactor $1/(2^dd!)$ in front. The notation is that of (2d), and $\chi'$ denotes the test $|B|<\delta'_k/g_k$.

\begin{gather*}
\FitMath{\mathbf{T}'_k(X)F=\frac1{|W_4|}\sum_{\text{histories in }X,\,r}\int dB\,\delta_{rT_0}(B)\,\chi'\!\int dV_{k_\flat}\,\chi\,\mathbf{T}_{k_\flat}(C_{k_\flat})\,e^{-\frac12\|\phi^{1/2}D\mathsf{H}B\|^2-A(x''_k\phi_{\mathrm{c}},U''(e^{ig_kB},V''))+V(X^\sim,U_k)+\text{numbers}}\,F .}
\end{gather*}

A component is healed when the new $V_k$ is regular in $X$. Then nothing is retained once the sum over the families has been exponentiated (2d), and $X$ is removed from the label. For an island, $X$ retains 

\begin{gather*}
\FitMath{\chi^c_k\,\chi_{k,\Lambda}\,\chi(\Lambda^{\mathrm{pr}})\,\delta_{\mathrm{ax}}(V_kV_\Lambda^{-1})\,\mathbf{T}'_k(X)}
\end{gather*}

 as a single block under its frame sum, for as long as it stands ((1.72), (1.104)). The cut-off $\chi(\Lambda^{\mathrm{pr}})$ is absent from (1.102). A drop-out after $n$ stages retains the following: the one-step operations of the levels that were not reached; the bare integrals $\int dV_l\,\delta(\mathcal{M}(V_l)V_{l+1}^{-1})$ of the levels already treated; the core; and the functions present together with the failing factor $1-\chi$. The terms produced up to that point are added to the boundary terms.

\textbf{Proposition 2.4 [minimiser in a box; Comm. Math. Phys. 122 (1989) 175, Proposition 1].} There exist a radius $a_\Lambda>0$ and a constant $C(M)$, both chosen after $M$ and before $\gamma$, with the following property. Suppose the exterior data satisfy $|\partial V_k-\mathbf{1}|<\varepsilon$ with $C(M)\varepsilon\le a_\Lambda$. In the application $\varepsilon\le C\varepsilon_k$, so this is a ceiling on $\gamma$. Take the class of continuations with $|\partial V_k-\mathbf{1}|<a_\Lambda$ at those level-$k$ plaquettes whose corners all lie within distance $1$ of $\Lambda^{\mathrm{pr}}$. On this class the background $U_{k,X}(V_k)$ exists. Within the class, the map $V_k|_{\Lambda^{\mathrm{pr}}}\mapsto A(U_{k,X}(V_k))$, which is a function of $V_k$ on the level-$k$ bonds having an end point in $\Lambda^{\mathrm{pr}}$, has exactly one critical gauge orbit $V_\Lambda$. This orbit is the set of its minima in the class, it is analytic in the data, and it satisfies $|\partial V_\Lambda-\mathbf{1}|\le C(M)\varepsilon$. Ba{\l}aban states it for the function on fields of mild regularity and for $\varepsilon$ sufficiently small, with the bound $CM^5\varepsilon$. The orbit of this class is the one used in (3) and (4). Moreover, $U_0:=U_{k,X}(V_\Lambda)$ satisfies $J(U_0)=0$ in $\Lambda^{\mathrm{pr}}$. The proof is a contraction argument in a ball whose radius depends on $M$, which is why a class of continuations is fixed. In axial gauge the Hessian is bounded below by $c(M)>0$; Ba{\l}aban obtains the lower bound $cM^{-5}$. Indeed, the Hessian is at least $c\|\partial B\|^2$ less a small multiple of $\|B\|^2$. In a tree gauge one has the Poincar{\'e} inequality

\begin{gather*}
\FitMath{\sum_b|B(b)|^2\le C(M)\sum_p|\partial B(p)|^2}
\end{gather*}

(122 (1989) 355, (1.7)--(1.9), with $CM^5$). The corresponding inequality for a single bond of the nest, $|B(b)|^2\le C(M\check{R}_k)^C\sum_p|\partial B(p)|^2$, yields the factor of the test in (3). It is (1.78) there, with every bond and every plaquette taken in its own lattice. The power of $\check{R}_k$ has two origins: the side of a shell in its own lattice, which is at most $CML\check{R}_k$ sites, and the number $\mathsf{N}_k+1$ of shells, where $\mathsf{N}_k=\check{R}_k$. The first form serves the level-$k$ variable $V'_k$ and the uniqueness of $V_\Lambda$. The form for one bond serves the multi-scale field $B$.

\textbf{A gauge for the configuration generated in a healing.} The configuration generated in a healing must belong to the analyticity spaces. To this end it is put in the form $e^{i\eta_kA^{\mathrm{g}}}U$, where $A^{\mathrm{g}}$ and its covariant differences $\nabla_UA^{\mathrm{g}}$ ($\nabla^U_\mu$ as in the Notation) are smaller than $(1+c\beta_{\mathrm{s}})\alpha_{1,k}$ (2d). On boxes of side $O(1)$ around the sites $y$ there are local gauges $\hat{u}_y$, with potentials $A^{(y)}$. They are glued together by a partition of unity $(\chi_y)$ taken in the quaternions,

\begin{gather*}
\FitMath{u^{\mathrm{g}}(x)=\mathcal{N}\Big(\sum_y\chi_y(x)\hat{u}_y(x)\Big),} \\[3pt]
\FitMath{\mathcal{N}(m)=(\det m)^{-1/2}m,}
\end{gather*}

and then $A^{\mathrm{g}}=\sum_y\chi_yA^{(y)}$ is a mean whose weights sum to $1$. This holds apart from quadratic terms and from terms comparable to the differences $\hat{u}_y\hat{u}_{y'}^{-1}-\mathbf{1}$ between local gauges at neighbouring sites. In $\nabla_UA^{\mathrm{g}}$, the term $\sum_y(\partial\chi_y)A^{(y)}$ has weights summing to $0$. It is therefore comparable in size to the differences between the potentials of neighbouring sites, which have a leading part in common. The square root is the branch on $|\det m-1|<1$ that takes the value $1$ at $\det m=1$. The mean in the display lies in this domain. It is the product of one local gauge $\hat{u}_y$ with $\mathbf{1}$ plus an average of the small differences $\hat{u}_{y'}\hat{u}_y^{-1}-\mathbf{1}$ between local gauges at neighbouring sites. The real linear span of $\mathrm{SU}(2)$ is $\mathbb{H}$. There $\det m=|m|^2$, and $\mathcal{N}(m)=m/|m|\in G$ for $0<|m|<\sqrt2$. Wherever the root is defined, $\mathcal{N}$ is holomorphic and bi-covariant: $\mathcal{N}(u_1mu_2)=u_1\mathcal{N}(m)u_2$ whenever $u_1,u_2\in G^c$. Hence the same construction applies to complex fields.


\subsubsection*{(2c) The rarity of wild patches: a small factor for each creation set against a cost for each site}

\textbf{Lemma 2.5 [one factor per creation].} The following holds on the support of a term, at step $t=k+1$. (i) When the plaquette test of step $t$ fails at $\square$, each real $U$ that has the stored averages satisfies $\sum_{p\text{ near }\square}a_p(U)\ge c\,\varepsilon_t^2$. (ii) When the approximate fluctuation test fails, this sum, increased by the part of $\mathrm{GF}_k$ located at the blocks near $\square$, is bounded below by $c\,\delta_k^2$. (In what follows, $\mathrm{GF}_j(X;V_j)$ denotes the part of $\mathrm{GF}_j(V_j)$ located at the blocks in $X$.) (iii) One has $\langle B,\Delta^{(k)}B\rangle\ge c\|B\|^2$, and the bound extends to the kept conditional forms (2b). Consequently a failed fluctuation test yields 

\begin{gather*}
\FitMath{\tfrac12\langle\mathcal{A}^{\mathrm{s}}_k,\Delta^{(k)}\mathcal{A}^{\mathrm{s}}_k\rangle\ge c\,\delta_k^2/g_k^2}
\end{gather*}

. The same is true of a failed collar test or rim test, because the threshold of either differs from $\delta_k$ by a constant factor, namely $1/(4\Theta_1)$ or $\theta_c$ (2a). Each of the three is proved in contrapositive form, and this form also covers the configurations at which some object entering a test is not defined. Suppose that the quantity which pays lies below the bound and below a constant. Then the data are regular, so (1b) provides the localised orbit and the contour variables, and the inequalities of the test are satisfied. In (iii), if the logarithm of the tested field is undefined, the variable of integration must satisfy $|\mathcal{A}'|\ge\delta_k$ (on the cut-off zone of the strip, the shift vanishes when $\mathcal{A}'$ is small). Since the coupling at that place is $x_{t-1}$, every created cell produces the factor $e^{-c\,p_0(g_t)^2}$. Inside $\mathbf{R}$ the tests work with the threshold $\delta'_j$, and the gauge trees they use have length at most $C\check{R}_j^{C}$. When a bond variable is large in a tree gauge, all that follows is that some plaquette exceeds $c\check{R}_j^{-C}\delta'_j$. For this reason these tests give no more than $e^{-c\check{R}_j^{-C}p_1(g_j)^2}\le e^{-p_0(g_j)}$, and this inequality requires $2\mathfrak{p}_1-Cr>\mathfrak{p}_0$.

\textbf{The energy of the averages.} To prove (i), apply (1b) in the reverse direction: it gives $|\partial V_t-\mathbf{1}|\ge c\,\varepsilon_t$ somewhere nearby. Moreover, a configuration of small action has flat averages, $\sum|\partial\mathcal{M}^tU-\mathbf{1}|^2\le C\sum a_p(U)$, and this inequality is scale invariant in dimension four. Stated separately, the bound reads as follows. Let $U$ be any real configuration, with no regularity required, so that the whole constraint set is covered. Let $S$ be any set of plaquettes of $\mathbb{T}^{(t)}$, and let $\mathcal{N}_0(S)$ denote the bare plaquettes lying beneath it, namely those all of whose corners lie in the union of the blocks $B^t(y)$, taken over the corners $y$ of a single $p'\in S$. Then

\begin{gather*}
\FitMath{\sum_{p\in\mathcal{N}_0(S)}a_p(U)\,\ge\,\min\Big\{\tfrac14\sum_{p'\in S}\big|\partial\mathcal{M}^tU(p')-\mathbf{1}\big|^2,\,\tfrac12e_*\Big\},} \\[3pt]
\FitMath{e_*=e_*(L)>0.}
\end{gather*}

Put $f_l=|\partial\mathcal{M}^lU-\mathbf{1}|$. A single averaging step yields $f_{l+1}\le\mathsf{W}f_l+C\sum f_l^2$, in which $\mathsf{W}$ is an averaging matrix whose row sums are $L^2$ and whose column sums are $L^{-2}$. By Schur's test, $\|\mathsf{W}\|_{2\to2}\le L^{(4-d)/2}=1$ in dimension four. The bound $\|\mathsf{W}\|_{1\to1}\le L^{-2}$ makes the quadratic errors sum geometrically, so no constant builds up as $t$ grows. (Ba{\l}aban gives (i) at the background and under regularity: in 122 (1989) 355, p. 380, the factor of a failed plaquette test is derived from Theorem 1 of 102 (1985) 277 and from (71) of 102 (1985) 255. Here, by contrast, the bound is valid for all real $U$, with constants independent of how many steps have been taken.)

\textbf{No integration is undone.} For the superrenormalisable model (Comm. Math. Phys. 86 (1982) 555, Sect. 3.A and Proposition 3.1), the conditional integrations are reversed at the end. This recovers a single positive quadratic form over many scales, and the factors of all the characteristic functions are then extracted from it simultaneously. Here the operations $\mathbf{T}_k(Z)$ keep the integrals, and every factor is obtained at the place where it sits. It comes from the Wilson weight, through the display above; or from the Wilson weight together with the kept gauge-fixing exponents (Lemma 2.5 (ii)); or from the conditional form, whose lower bound is $c$. (Taken as a whole, the conditional form inherits the bound because it is a Schur complement; the kept part of this form is treated in (2b).) With this, the cost of the stored Gaussian integrals is $e^{C\#_j(Z_j\cap\Omega_j)}$ (122 (1989) 355, p. 383).

\textbf{One action and many claimants.} Gains are added up in the ledger that follows, and this addition is legitimate: nothing is spent twice and the integrals that remain converge. A single Wilson term, a single set of gauge-fixing factors and a single set of stored Gaussians pay for every failed test at every level. A \emph{piece} is a pair $(t,D)$ in which $D$ is a component of the region produced by step $t$ on which no $\mathbf{R}$ has acted up to level $k$. We write $\mathsf{w}_{t,D}$ for the weight of its one-step operation. This weight consists of the gauge-fixing factors with their normalisers, the Jacobians that arise when the averaging constraint is solved for the pivot bonds, and the conditional Gaussian. We write $\xi_{t,D}$ for the variables that belong to the piece. Let the shares of the action be as defined in (1), and extend $U_k$ beyond $\Omega_0$ by means of the stored bare field. Then at each point of the support of the term of a history one has $A(x_k(\cdot),U_k)\ge A_{\mathrm{cr}}:=\sum_{(t,D)}x_{t-1}\sum_{p\in\text{share of }(t,D)}a_p(U_k)\ge0$, so that the only part of the action spent is the Wilson sums over the shares. Moreover

\begin{gather*}
\FitMath{e^{-A_{\mathrm{cr}}}\prod_{(t,D)}\mathsf{w}_{t,D}\le\exp\Big[-c\sum_{\text{created cells }(t,\mathfrak{c})\text{ in the pieces}}p_0(g_t)^2\Big]\prod_{(t,D)}\bar{\mathsf{w}}_{t,D}(\xi_{t,D}),} \\[3pt]
\FitMath{\int\bar{\mathsf{w}}_{t,D}\le e^{C\#_t(D)},}
\end{gather*}

where each $\bar{\mathsf{w}}_{t,D}$ is a Borel function of the variables of its own piece only (the background no longer enters it), and $C=C(L)$ carries no $\log x$ per site. (1) \emph{The action is divided among the pieces, not held in common.} The share of a piece consists of the bare plaquettes that lie within one cell of those of its cells at which a plaquette test or an approximate fluctuation test has failed. Distinct pieces have disjoint shares, because the tests of a step are carried out several cells away from the hull of the older region. On a share one has $x_k(\cdot)=x_{t-1}$, since the older profiles have completed their transition and the newer ones have not yet begun it. Every other plaquette contributes $\ge0$. (2) \emph{The tests concern averages of the configuration that appears in the action.} The tests were built from localised backgrounds of step $t$, whereas the Wilson term contains $U_k$. The two are connected as follows. The identity $\mathcal{M}^{t-1}U_k=V_{t-1}$ on $\Gamma_{t-1}$ is the constraint that defines the background, and on the same set $V_t=\mathcal{M}(V_{t-1})$ is imposed by the corresponding delta function. Lemma 2.5 holds for every $U$ with these averages, so it applies to $U_k$. (3) \emph{Each exponent is split into two halves}: one half is used for payment and the other is kept for the integration. The only normaliser that depends on the coupling is integrated against half of its own exponent,

\begin{gather*}
\FitMath{\int du\,e^{-\frac12x(1-\frac12\operatorname{Re}\operatorname{tr}u)}\le C\int du\,\mathbf{1}\{|u-\mathbf{1}|<\epsilon\}\,e^{-x(1-\frac12\operatorname{Re}\operatorname{tr}u)}} \\[3pt]
\FitMath{(\epsilon\le1,\,x\epsilon^2\ge1),}
\end{gather*}

and both sides are of order $x^{-3/2}$. On the cut-off zone of the strip the shift (2b) vanishes for small $\mathcal{A}'$; consequently a large tested field is a large variable of integration. (The approximate test requires $\mathrm{GF}$ for the following reason. Take a gauge transformation at a site of $\mathbb{T}^{(t-1)}\setminus\mathbb{T}^{(t)}$. It moves $V_{t-1}$ away from $\mathcal{M}^{t-1}U_{t,\square}(V_t)$, the average of the localised background (2a), which is in the block axial gauge (1a). It does so at no cost in action, since it alters neither the action nor $V_t$.) (4) \emph{Within a patch on which the $\mathbf{R}$ of level $m$ acts}, a test that fails may not be paid for by the large field from which the patch arose. Whatever is older than the wait was paid for by plaquettes that lie in the \emph{floor}. The floor is made up of those cells of level $m_\flat$ of the old core, $\Omega^c_{m_\flat}\cap X$, which do not touch $\Omega_{m_\flat}$; this old core is a part of the core $C_{m_\flat}=X\cap Z_{m_\flat}$ of (2b), because $\Lambda_i\subset\Omega_i$. The gauge-fixing factors and stored Gaussians that paid for it are all of levels below $m_\flat$. Let $X$ be a drop-out, and consider stored variables with two properties: each test carried by $X$ takes its recorded value, and the plaquette variables of the datum in the old core, taken layer by layer ($V_l$ on $\Gamma_l$), are within $C\check{R}_l^2\varepsilon_l$ of $\mathbf{1}$. (These two properties hold on the supports. They also hold at the points of evaluation that occur below in the mass of a later drop-out that contains $X$. Within the mass of $X$ itself, no use is made of the test whose failure turned it into a drop-out.) For every real $U$ compatible with stored variables of this kind, one has

\begin{gather*}
\FitMath{c\,x_m\!\!\sum_{p\subset X^+,\,\text{no corner in the floor}}\!\!a_p(U)+\sum_{j=m_\flat}^{m-1}x_j\mathrm{GF}_j(X;V_j)+c\sum_{j=m_\flat}^{m-1}\|\mathcal{A}^{\mathrm{s}}_j\|_X^2\,\ge\,(\log x_m)^{2\mathfrak{p}_1-Cr}\,\ge\,p_0(g_m),}
\end{gather*}

where $X^+$ denotes $X$ enlarged by two layers of unit cubes of level $m$ and $U$ is compatible on $X^+$. This follows from a single dichotomy: if the left-hand side is the smaller, then each test of the operation takes the value $1$. In every case the bridge runs along the same chain. It starts from the action of a compatible configuration and passes to the plaquettes of the stored fields (by the energy of the averages). From there it passes to the regularity of the background used by the test, and then to the bond variables. That last link comes either from the rigidity of a single averaging step or from loops in the shells. In a shell of side $s\le CM\check{R}_m$, a loop composed of $\mathsf{n}$ straight pieces has a deviation from $\mathbf{1}$ that is bounded by $C\mathsf{n}s^2$ times the plaquette bound, and this is where the loss $\check{R}_m^{-C}$ comes from. The waiting collar screens off the unknown large field in the floor; this collar is what separates the floor from the region $X\cap\Omega_{m_\flat}$ in which the operation acts, and its width is at least one cell in the local units. Where the operation acts, the background of any stage has curvature at most $C(\varsigma+\check{R}_m^2\varepsilon_me^{-\delta_0M\check{R}_m})\eta_j^2$ on the layer $j$, which is the local scale. Here $\varsigma$ is the size of the plaquettes of the stored fields that follows from the smallness of the left-hand side. A drop-out uses no plaquette that has a corner in its floor. By contrast, each plaquette used beneath it lies inside the floor, and each gauge-fixing factor and stored Gaussian used beneath it is of a level below $m_\flat$. The windows of nested operations are disjoint.

\textbf{The cost.} Carrying $Z_n$ through step $n$ incurs three costs. The first is a factor $e^{C\log x_{n-1}}$ for each site of $\mathbb{T}^{(n)}$ lying in $T_\ell\setminus\Omega_n$. The logarithm in question is $-\tfrac32\log x_{n-1}$ times the number of free bonds less the number of gauge-fixed sites. It is produced where the fluctuation field is introduced and rescaled, $\mathcal{A}_{n-1}=g_{n-1}\mathcal{A}^{\mathrm{s}}_{n-1}$, and where the gauge is fixed, that is on $\Omega_n$, and not where the field is integrated; the normalisation removes it everywhere. The second is a factor $e^C$ for each site of $\mathbb{T}^{(n)}$ lying in the rest of $Z_n$ enlarged by a layer. It comes from $\mathbf{E}_n,\mathbf{R}_n,\mathbf{B}_n$ and from the stored Gaussians, which are written in the scaled variable. The third is the entropy of the subregions. Altogether, for each component $Z$ of $Z_n$, the cost is $\exp[CM^4\check{R}_n^5(d'_n(Z)+1)]$, where $d'_n$ is the tree length measured in cells. Of the five powers of $\check{R}_n$, four account for volume and one for $\log x_{n-1}\le C\check{R}_n$. In the count below, the number $|Z_n|_n$ of level-$n$ cells of the region takes the place of Ba{\l}aban's linear size $d'_n$, because a region that consists of a single cell has $d'_n=0$ and still incurs a cost.

\textbf{Proposition 2.6 [the count, the ledger, the credit].} Let a history and one of its components be given. Then the count and the ledger displayed below are valid. If moreover the component is neither healed nor marked, the credit inequality displayed below is valid for it as well. In all of these the constants depend neither on the history nor on the number of steps. The ledger and the credit are statements of Ba{\l}aban (Comm. Math. Phys. 122 (1989) 355, (1.79)--(1.81), (1.89)), here with volumes counted in cells. The proof does not follow his induction. It rests instead on counting once, globally, along the whole genealogy of a component.

\textbf{Shrinking and waiting.} We denote by $S(Z)$ the free evolution of (2a), that is, the hull of $Z$ in the next partition together with $O(1)$ further layers, boxed. Every step that is not a drop of the cube size $\check{R}$ divides linear sizes, counted in cells, by the factor $L$. When a step is a drop, the partitions before and after it coincide, so that nothing shrinks; and of two consecutive steps, at most one is a drop. Consequently

\begin{gather*}
\FitMath{d'_n(S^{n-j}Z)\le CL^{-\nu}d'_j(Z)+C,} \\[3pt]
\FitMath{\nu\ge\lfloor(n-j)/2\rfloor\,\text{ the number of the steps }j+1,\dots,n\text{ that are not drops}}
\end{gather*}

(this display replaces the scaling of the tree length in Ba{\l}aban, Comm. Math. Phys. 122 (1989) 355, p. 384, which in the present notation, with $Z^{(n-j)}$ denoting the hull of $Z$ in level-$n$ cells, reads: $d'_n(Z^{(n-j)})\le L^{-\frac12(n-j)}d'_j(Z)\le2^{-(n-j)}d'_j(Z)$ over $n-j>1$ steps, and $d'_n(S^{n-j}Z)\le C\,d'_n(Z^{(n-j)})$, or $\le C$ when the right-hand side vanishes). The added layers do not pile up, because every change of scale contracts what the earlier levels added. It follows that within $\mathsf{k}\le\mathsf{N}_j+C\log(d'_j(Z)+2)$ steps a component $Z$ of $Z_j$ does one of three things, provided the history continues up to level $j+\mathsf{k}$: it meets a creation, it joins another component, or it enters the class of $\mathbf{R}$.

\textbf{Why a quiet component is bound to enter the class.} Shrinking makes the component small, but smallness alone does not place it in the class. Consider a quiet line, that is, the sequence of components $v_i$, of the levels $i$, into which a single component is carried as long as nothing is created in it and nothing joins it. Such a line develops deterministically, $v_i=\Phi_i(v_{i-1})$, where $\Phi_i$ is the rule of the step for a lone component, and its side in cells obeys $\alpha_{i+1}\le\alpha_i+C$, $\alpha_{i+2}\le\alpha_i/L+2C$. This is an affine contraction over two levels; its fixed point is less than half of the threshold of smallness, so that room is left for a drop. Apart from the depth $k_\flat\ge1$, which is a matter of the arithmetic of the cube sizes, a line of this kind acquires on its own precisely the conditions of the class (2b). Moreover, a condition that fails on account of an event in the window also fails at every intermediate level, because $n\mapsto n-\mathsf{N}_n$ is strictly increasing. Hence the carrier of the event cannot have been healed in between. The waiting bound is founded on this property of the class (in Ba{\l}aban, a single sentence, p. 385).

\textbf{The count.} An unhealed component of level $n-1$ is a predecessor of the component $(n,C)$ of $Z_n$ that contains its hull. For a component $u_*$, the genealogy $\operatorname{Gen}(u_*)$ is the tree formed by all $u=(n(u),C(u))$ that equal $u_*$ or are iterated predecessors of it. It is the same object as the two of (2a). The cluster of a root is its genealogy together with everything that the history records within it. The lineage of a live component is made up of the members of its genealogy, regarded as regions $\mathsf{Z}_j$, together with the creations in them. None of these three contains a component that was healed earlier and whose place a later region covers again; such a component is the root of a cluster of its own. Write $\mathsf{n}^{\mathrm{cr}}_t(u_*)$ for the number of created cells of level $t$ in $\operatorname{Gen}(u_*)$. Write $\mathsf{n}^{\mathrm{mk}}_t(u_*)$ for the number of its marked components of level $t$ other than $u_*$ itself. These are the drop-outs and islands, that is, the creations by an operation of (2a), and they are counted separately here because, for each of them, the count uses only the factor $e^{-p_0}$, which is all that a drop-out contributes. Set $\mathsf{T}_t:=C\check{R}_t$; this number bounds the length of a wait and does not increase with $t$. Then for non-increasing costs $\hat{\mathsf{r}}_1\ge\hat{\mathsf{r}}_2\ge\dots\ge0$,

\begin{gather*}
\FitMath{\sum_{u\in\operatorname{Gen}(u_*)}\hat{\mathsf{r}}_{n(u)}\,|C(u)|_{n(u)}\,\le\,C\sum_t\mathsf{n}^{\mathrm{cr}}_t(u_*)\,\mathsf{T}_t\hat{\mathsf{r}}_t+C\sum_t\mathsf{n}^{\mathrm{mk}}_t(u_*)\,\mathsf{T}_t\hat{\mathsf{r}}_t .}
\end{gather*}

In particular, for each history and without any bound on the number of steps,

\begin{gather*}
\FitMath{\sum_{n\le k}\#\{\text{cells of }Z_n\}\le C\sum_{\text{creations of both kinds, at steps }t}(\mathsf{N}_t+1).}
\end{gather*}

The proof consists of three moves, and nothing is carried through an induction over the levels. (1) Contraction. The cells are counted by way of the occupied cells of a component, which are the hulls of the created cells beneath it. These form a skeleton that commutes with taking hulls, and the component never lies more than $O(1)$ cells away from it. Over a fixed number $s_0$ of levels, at least half of the occupied cells of a component that has many of them disappear, whereas a light component (one with at most $C$ of them) retains at most $C$. Let $\mathsf{o}_n$ be the number of occupied cells of level $n$ in $\operatorname{Gen}(u_*)$, and let $l_n$ be the number of its light components of level $n$. Then

\begin{gather*}
\FitMath{\mathsf{o}_n\le\tfrac12\mathsf{o}_{n-s_0}+C\,l_{n-s_0}+\sum_{n-s_0<t\le n}\mathsf{n}^{\mathrm{cr}}_t} \\[3pt]
\FitMath{(\text{quantities of levels}\le0\text{ are }0),} \\[3pt]
\FitMath{\text{whence}} \\[3pt]
\FitMath{\sum_n\hat{\mathsf{r}}_n\mathsf{o}_n\le C\sum_n\hat{\mathsf{r}}_nl_n+C\sum_t\hat{\mathsf{r}}_t\mathsf{n}^{\mathrm{cr}}_t,}
\end{gather*}

where the second inequality is obtained from the first by multiplying by $\hat{\mathsf{r}}_n$ and summing, since the costs are non-increasing. (2) Stretches. A member of $\operatorname{Gen}(u_*)$ is called an event if it contains a created cell of its own level, or comes after a marked component, or has two or more predecessors (in which case it is a joining). The genealogy is a disjoint union of stretches. Each stretch starts at an event and continues through components that have a single predecessor and contain nothing new. A stretch that starts at level $t$ has at most $\mathsf{T}_t$ light members. Indeed, after $O(1)$ levels a light component is contained in a cube of $O(1)$ cells, and once a further $\mathsf{N}_t$ quiet levels have passed it has entered the class, in which it is healed or else marked. (3) For each of the three kinds, the events are paid for through an injective assignment, and no created cell is charged more than twice. An event that contains a created cell is paid by that cell, the successor of a mark is paid by that mark, and a joining is paid by a created cell at a leaf of the tree. For joinings, select at each member that has predecessors one of them as its first predecessor. To a joining $v$, assign the leaf reached by starting from a predecessor of $v$ other than the first and always moving to the first predecessor. No two joinings receive the same leaf, since on the way up from that leaf the first member entered otherwise than from its first predecessor is $v$ (the leaves of a tree outnumber its branchings). Moreover, the leaf lies at a lower level, at which $\mathsf{T}_t\hat{\mathsf{r}}_t$ is no smaller (because $\mathsf{T}_t=C\check{R}_t$ is non-increasing in $t$), and each leaf contains a created cell. Thus a joining requires no payment of its own.

\textbf{The ledger.} Each created cell brings the gain $\mathsf{g}_t=c\,p_0(g_t)^2$ and each mark the gain $\mathsf{p}_t=p_0(g_t)$ (Lemma 2.5). Against these stands, for every cell and every level, the cost $\mathsf{r}^{\mathrm{c}}_n=C\log x_n\,(M\check{R}_n)^4$ found above, which is bounded by the non-increasing sequence $\hat{\mathsf{r}}_n=C\mathsf{r}_n=CM^4\check{R}_n^5$; here $\mathsf{r}_n$ is the weight of a cell, (2a). With these choices,

\begin{gather*}
\FitMath{\mathsf{F}(u_*):=\sum_t\mathsf{g}_t\mathsf{n}^{\mathrm{cr}}_t(u_*)+\sum_t\mathsf{p}_t\mathsf{n}^{\mathrm{mk}}_t(u_*)-\sum_{u\in\operatorname{Gen}(u_*)}\mathsf{r}^{\mathrm{c}}_{n(u)}|C(u)|_{n(u)}\,} \\[3pt]
\FitMath{\ge\,2\sum_t\mathsf{p}_t\mathsf{n}^{\mathrm{cr}}_t(u_*)\,\ge\,(2-c)\,p_0(g_{n(u_*)}),}
\end{gather*}

provided that $\mathsf{g}_n\ge2\mathsf{p}_n+C\mathsf{T}_n\hat{\mathsf{r}}_n$ and $\mathsf{p}_n\ge C\mathsf{T}_n\hat{\mathsf{r}}_n$. An event costs $\mathsf{T}_n\hat{\mathsf{r}}_n\le CM^4\check{R}_n^6\approx(\log x_n)^{6r}$. This cost is made up of $\log x_n\le C\check{R}_n$ for each site, $(M\check{R}_n)^4$ sites in a cell and a wait of length $\check{R}_n$. Since $\mathfrak{p}_0\ge6r+1$, the quantity $p_0(g_n)=A_0(\log x_n)^{\mathfrak{p}_0}$ exceeds this cost, and the extra power of the logarithm absorbs the constants. The final inequality holds for the following reason. The genealogy has a leaf, and so it has a created cell at some level $t\le n(u_*)$. Moreover $p_0(g_t)\ge(1-c)p_0(g_{n(u_*)})$, because, up to a bounded error, the flow is monotone: $x_t\ge x_{n(u_*)}-C$. When volumes are counted in cells, $e^{-\mathsf{F}(u_*)}$ coincides, up to two modifications, with the quantity whose supremum is taken on the right-hand side of (1.79). First, the $2p_0(g_j)$ that Ba{\l}aban reserves for each created region is taken from $c\,p_0(g_j)^2$ (the first of the two conditions). Second, Ba{\l}aban's cost of a cell, which in the notation of this text reads $CM^4\check{R}_j^5$, is replaced by the finer form $C\log x_j\,(M\check{R}_j)^4$. Components that consist of a single cell are included: such a component costs $\mathsf{r}^{\mathrm{c}}_n$ at every level, and this cost is paid through the light members.

\textbf{The credit, as a corollary.} Let $z=(j,Z)$ be a component that is neither healed nor marked. We write its factor as $e^{-\mathsf{F}(z)}=e^{-\kappa_j(Z)-2p_0(g_{j(Z)})}$, where $j(Z)$ is the level of the first creation occurring in its genealogy. We extend the history of $z$ by the fictitious continuation $S(Z),\dots,S^{\mathsf{k}}(Z)$. Nothing is created in this continuation, and it runs at the levels $n\le K$, where the flow bounds are available; beyond $K$ it is cut off. We mark the last member of this continuation and apply the ledger to the extended history. The extended history has the same created cells as $z$, and apart from that last member it has the same marked components. The ledger therefore gives 

\begin{gather*}
\FitMath{\mathsf{F}(z)-\sum_n\mathsf{r}^{\mathrm{c}}_n|S^{n-j}Z|_n\ge2\sum_t\mathsf{p}_t\mathsf{n}^{\mathrm{cr}}_t(z)\ge2\mathsf{p}_{j(Z)},}
\end{gather*}

which is to say

\begin{gather*}
\FitMath{\kappa_j(Z)\,\ge\,\sum_{n=j+1}^{\min\{j+\mathsf{k},K\}}\mathsf{r}^{\mathrm{c}}_n\,\big|S^{n-j}Z\big|_n ,}
\end{gather*}

that is, (1.80), in which Ba{\l}aban has $K$ in place of $\mathsf{k}$ and the summand 
\begin{gather*}
\FitMath{O(1)M^dR_n^{d+1}d'_n(S^{n-j}(Z)),}
\end{gather*}
 where $R_n$ is his size, built from $g_n$ itself. Thus the credit $\kappa_j(Z)$ pays for the entire wait. By the shrinking, the sum is bounded by $CM^4L^5\check{R}_j^6(d'_j(Z)+1)$ ((1.81)).

\textbf{Who pays.} At a creation, the deposit $c\,p_0(g_j)^2(d'_j+1)$ has degree $2\mathfrak{p}_0$ in $\log x_j$. It exceeds the total cost $CM^4L^5\check{R}_j^6(d'_j+1)$, whose degree is $6r$, so that per cell $p_0^2$ covers $\check{R}_j^6$. Now suppose a test fails inside $\mathbf{R}$. All that is known of the old credit is $\kappa\ge0$, so the renewed wait of a small region, $CM^4\check{R}_j^6$, must be paid by the new factor $e^{-p_0(g_j)}$ on its own. For this reason that factor has to be at least as good as $e^{-p_0}$. A joining neither pays anything nor releases anything.

\textbf{Joining: Ba{\l}aban's induction and the count.} Ba{\l}aban proves (1.80) by induction on $j$ and, within a step, by induction on the number of regions that join (pp. 384--387). When $X$ and $Y$ join to form $Z$, he sets free the reserve $2p_0(g_{j(X)})+2p_0(g_{j(Y)})-2p_0(g_{j(Z)})\ge c\,p_0(g_j)$ to pay for the renewed wait. This step relies on the claim of p. 387 that, once $S$ has been applied $n_1$ further times, the domain built from $Y$ includes the domain built from $X$; there $n_1$ is called a rather small number, with $n_1<10$ given as an instance. The count uses no containment of this kind. For the hull operation of this part, whether such a containment holds depends on where the two regions lie relative to the grids. In a single coordinate the hull map is $\mathsf{h}(y)=\lfloor(y+\tfrac12(L-1))/L\rfloor$, and $S$, taken with ten layers, maps an upper end $y$ to $\mathsf{h}(y)+10$. Set $y_{m_*}=10$ and $y_m=L(y_{m+1}-10)+\tfrac12(L-1)$. Then

\begin{gather*}
\FitMath{\mathsf{h}(y_m)=y_{m+1}-10,} \\[3pt]
\FitMath{\mathsf{h}(y_m+l)=y_{m+1}-9} \\[3pt]
\FitMath{(0\le m<m_*,\,1\le l\le L).}
\end{gather*}

Suppose that $Y$ ends at $y_0$ and that $X$, a block of side $5$, begins at $y_0+1$. The two regions touch. The upper end of $S^m(Y)$ is $y_m$ and the upper end of $S^m(X)$ is $y_m+1$ (for $m=0$ it is $y_0+5$), so $S^m(X)\not\subset S^m(Y)$ for every $m\le m_*$, where $m_*$ is arbitrary. At each of these steps the two partitions have ratio $L$, so that none of the steps is a drop. The count never compares $S^m(X\cup Y)$ with $S^m(Y)$.

\textbf{Proposition 2.7 [the mass bound; where the small factor stands].} Consider a component $X$ that has passed every test of the operation $\mathbf{R}$ of some level $i\le k$ before the last decomposition $1=\chi_i+\chi_i^c$ (a healed component or an island). The operation of (2b) on it is written

\begin{gather*}
\FitMath{\mathbf{T}'_i(X)F\,(\mathbf{p})=\frac1{|W_4|}\sum_\omega\int e^{\Phi_\omega(\xi;\mathbf{p})}F\,\mu_\omega(d\xi),} \\[3pt]
\FitMath{\mathbf{p}=(\mathbf{U},\mathbf{J})\,\text{a pair on }X^\sim,}
\end{gather*}

where $\omega$ ranges over the local histories of $X$, that is, the results of every decomposition of unity inside $X$ together with a frame in $W_4$ of the gauge tree, $\xi$ ranges over the stored variables, and $e^{\Phi_\omega}$ contains the core Wilson factor $\exp[-A(x''_i\phi_{\mathrm{c}},U''(\xi))]$. (a) For a healed component or an island, in which case $\mathsf{k}=0$,

\begin{gather*}
\FitMath{\mathsf{M}(i,X)=\frac1{|W_4|}\sum_\omega e^{c\,\mathsf{v}(\omega)}\int\sup_{\mathbf{p}\in U^c_i(X^\sim,\alpha_{0,i},\alpha_{1,i})}\big|e^{\Phi_\omega(\xi;\mathbf{p})}\big|\,\mu_\omega(d\xi)\le e^{-(2-c')p_0(g_i)},} \\[3pt]
\FitMath{\mathsf{v}(\omega)=\sum_{u\text{ listed by }\omega}\log x_{n(u)}\,\#_{n(u)}(C(u)):}
\end{gather*}

here the coupling is taken at the level of the action, and the current coupling then follows from $p_0(g_i)\ge(1-c')p_0(g_k)$; the supremum stands inside the integral; and a weight growing exponentially with the space-time volume of the patch is paid for in advance. It is this bound that the induction carries forward. Ba{\l}aban has (1.89) instead, at each fixed pair ((1.73)). (b) Take a single history and a single component $Z$ of $Z_k$ on which no $\mathbf{R}$ has acted. For every value of the exterior variables (in total variation),

\begin{gather*}
\FitMath{\Big\|\mathbf{T}_k(Z)\big[\chi_{k,\mathcal{S}}\,e^{-A(x_k(\cdot),U_k)}\big]\Big\|\le\prod_{\text{created cells }(t,\mathfrak{c})\subset Z}e^{-c\,p_0(g_t)^2}\cdot\exp\Big[C\sum_{t\le k}\#_t(Z_t\cap Z)\Big],}
\end{gather*}

with each factor taken at the coupling of its own step. Bounds of this kind do not multiply over the components, since the background couples the components. The pointwise majorant given above does multiply (Tonelli). Consequently, after integration and for an entire label, there exists a history $h$ carrying the label $\mathcal{S}$ for which the bound below holds, where $(t,\mathfrak{c})$ ranges over the created cells of $h$ lying outside the marked $X$. (The created cells are not determined by the label: a row of three failing cells and that row with its middle cell removed yield the same regions.)

\begin{gather*}
\FitMath{\int\mathbf{T}'_k(\mathcal{S})\big[\chi_{k,\mathcal{S}}\,e^{-A(x_k(\cdot),U_k)}\big]dV_k\le e^{C\sum_j\#_j(Z_j)}\prod_{(t,\mathfrak{c})\text{ outside the marked }X}e^{-c\,p_0(g_t)^2}\prod_{(i,X)\text{ marked, outermost}}\mathsf{M}(i,X).}
\end{gather*}

(The bound holds a fortiori when each site carries $C(1+\log x_j)$ and each created cell carries $e^{-c\,p_0(g_t)}$. This is the form used in Proposition 2.13 [outlined below; Section~\ref{sec:actual-history-closure}], the statement that carries the logarithm, and in (4b).) (c) Consider the sum over the labels, with the costs paid by Proposition 2.6 and the nested sequences inside each component summed. In this sum every component $C$ of $Z_k$ keeps the factor $\exp[-c\,p_0(g_k)-\kappa_1\#_k(C)]$, in which the volume is counted in sites (and so, a fortiori, the factor with $\kappa_1|C|_k$, in which it is counted in cells; this is the form in the introduction to this step). Since $d_k(C)\le\#_k(C)$, this per-site factor includes the term $-\kappa_1d_k(C)$, which is what p. 387 asserts for an arbitrary region. The decay is one further cost $O(1)$ per site; like the other costs, it is paid for by the events (Proposition 2.6). If failed plaquette tests have produced a component, then $\mathbf{T}_k(Z)1\le\exp[C\sum_{t\le k}\#_t(Z_t\cap Z)]$ (119, (2.21), and p. 264, where it is said of the Wilson term that it ``yields the small factors for large field characteristic functions'').

\textbf{The bound has to reproduce itself.} (i) \emph{Supremum under the integral.} In a term of level $k\ge i$, the factor remaining on $X$ is evaluated at $(U_k,J(U_k))|_{X^\sim}$, and $U_k$ depends on the integration variables of every component. Such an integral is not controlled by a bound at each fixed pair, and bounds that are merely uniform in the other variables do not multiply. For example, take $\mathbf{1}\{|\xi_1-\xi_2|<\epsilon\}$ on $[0,T]^2$: each integral in one variable is $\le2\epsilon$, whereas the double integral is about $2\epsilon T$. What does multiply is a majorant given by a Borel function that depends only on the variables of $X$. The integral of the supremum itself is understood as an upper integral. Such a majorant can be obtained because only two factors depend on the pair: $\Phi_\omega=\Phi^\circ_\omega(\xi)+[-\tfrac12q_X+V(X^\sim)](\xi;\mathbf{p})$, where $q_X=\|\phi^{1/2}D\mathsf{H}B\|^2\ge0$ is formed at a real background, and $|V(X^\sim)|\le C\,\#\{\text{stored sites}\}$, with every site counted at its own level. The real background here is the critical configuration of the prepared sequence, taken with $B=0$ and at the datum of the main step, and it is built with the $G$-valued factor of the triple. Through that factor, $q_X$ depends on the splitting, whereas the exponent $\Phi_\omega$, taken in its entirety, is determined by the pair (see clause (vi) of the class of (2e)). (ii) \emph{A fixed set.} If the supremum is taken over the arguments that occur at level $k$, the bound is not stable, since the next step replaces $U_k$ by $U_{k+1}$. The supremum is therefore taken over a larger fixed set, built only from the data of the levels $\le i$. This set contains every later argument, since on the supports the following bounds hold, with nothing lost in the age $k-j$ of the layer $\Omega_j\setminus\Omega_{j+1}$,

\begin{gather*}
\FitMath{|\partial U_k-\mathbf{1}|<C\check{R}_j^2\varepsilon_j\eta_j^2,} \\[3pt]
\FitMath{|J(U_k)|<C\check{R}_j^2\varepsilon_j} \\[3pt]
\FitMath{\text{on }\Omega_j\setminus\Omega_{j+1}\,(\text{the current in level-}j\text{ units}),} \\[3pt]
\FitMath{C\check{R}_j^2\varepsilon_j<\alpha_{0,j}\,\text{ for all }j\,\,(q\text{ is chosen after }\mathfrak{p}_0\text{ and }r,\text{ the ceiling on the couplings after }L)}
\end{gather*}

(the last inequality holds since the quotient is a negative power of $\log x_j$, since $q>\mathfrak{p}_0+2r$, and since the couplings stay below a ceiling; moreover, a layer at index distance $a$ has an influence that decays like $e^{-(a-1)_+}$, while its threshold grows like $(1+a)^{1/2}$). Hence the mass seen at level $k$ does not exceed the mass of level $i$, and a single step $\mathbf{R}\mathbf{T}$ gives back the bound with unchanged constants. (iii) \emph{The weight} $e^{c\,\mathsf{v}(\omega)}$ has as its only cost a larger constant in the ledger. Since it carries $\log x_j$ per site, it absorbs every perturbation of size $O(1)$ per site of the patch, such as the factor $e^\Psi$ of the sources, (4b). In Proposition 2.13 [outlined below; Section~\ref{sec:actual-history-closure}] the weight covers whatever the label still exhibits inside a drop-out, since the small factors produced there have already been consumed by the mass; this is why the masses are required with the weight raised. The argument requires a volume that counts each pair (level, cell) only a bounded number of times. A volume that counts the past again at every drop-out may exceed the total of all the gains; an example is a row of cells all created at the first step, with a drop-out at each completed wait. (iv) \emph{Drop-outs.} The factors kept by a drop-out contain no Wilson action, and no inner background has been formed; the kept factors are evaluated at the background of whichever later term contains $X$. On the other hand, the configuration enters the Wilson term of the carrying term, and the failed plaquette tests listed by $\omega$, which are the tests of the creations in $X$, bound plaquettes formed from averages of this configuration. That weight, over the bare plaquettes all of whose corners lie in $X$, is therefore counted with the operation. The pair is replaced by an \emph{evaluation} at $\xi$. By this we mean a real $U$ on $X$ with four properties. First, it is compatible with $\xi$, that is, $\mathcal{M}^jU=V_j(\xi)$ on the bonds of $\Gamma_j$ in $X$. Second, it is regular layer by layer, in the same sense as on the supports: on the layer $j$ one has $|\partial U-\mathbf{1}|<C\check{R}_j^2\varepsilon_j\eta_j^2$, while the plaquette variables of the stored fields differ from $\mathbf{1}$ by less than $C\check{R}_j^2\varepsilon_j$. Third, at this configuration each test listed by $\omega$ takes its recorded value, with the exception of the test whose failure made $X$ a drop-out; here the fields not formed by $\xi$, including $V_i$ on $X$, are read as averages of $U$. Fourth, $(U,J(U))|_{Y^\sim}\in U^c_m(Y^\sim,\alpha_{0,m},\alpha_{1,m})$ for each island $(m,Y)$ in $X$. Together with this configuration, the supremum also runs over further free parameters subject to one-sided bounds. These are the weight $\upsilon$ of the Wilson term, with $\upsilon\ge c\,x_t$ on $Z_t$, and the kept forms and Jacobians, which are evaluated at a background (2b), with form $\ge c$ and Jacobian $\le C$:

\begin{gather*}
\FitMath{\frac1{|W_4|}\sum_\omega e^{c\,\mathsf{v}(\omega)}\int\sup_{\text{evaluations at }\xi}e^{-A(\upsilon\mathbf{1}_X,U)}\,(\text{kept factors})\,\mu_\omega(d\xi)\le e^{-(2-c')p_0(g_i)}} \\[3pt]
\FitMath{(\upsilon\ge c\,x_t\text{ on }Z_t).}
\end{gather*}

(Over the empty set the supremum is $0$.) Take a term of some level $k\ge i$ in which $X$ occurs, and a point of its support. Put $\upsilon=x_k(\cdot)$ and take the forms and Jacobians at the background of that term. Then the configuration that the Wilson term of that term carries on $X$ is an evaluation at those coordinates of the point that belong to $X$. Hence, as in (ii), the set contains all later arguments. For a drop-out, the left-hand side is its mass $\mathsf{M}(i,X)$. Where $X$ has no sum over frames, nothing depends on the frame listed by $\omega$, by the identity $1=|W_4|^{-1}\sum_r1$. This bound does not use the test whose failure made $X$ a drop-out. The gain $p_0(g_i)$ of that test is reserved for the renewed wait, and it is drawn on once $X$ occurs as a mark within a subsequent component (Proposition 2.6). The supremum is taken over a set fibred over $\xi$. Over a product set it selects a flat $U$, at which $e^{-A}=1$, and the factor is lost. In a toy case, take $\Phi=-T(1-\cos p)$ together with the test $|\xi|\ge\epsilon$: its mass is $e^{-cT\epsilon^2}$ on the graph $p=\xi$, but approximately $1$ over the product set. For this reason each inequality of Lemma 2.5 is required to hold at every configuration of this set, and not merely at backgrounds. These inequalities do hold for all real configurations having the prescribed averages. Nothing is assumed inductively at a drop-out; the inductive bound is the one for the islands. (v) \emph{Islands inside $X$.} The count and the ledger are carried out on the outer genealogy of $X$, that is, on the genealogy cut off at the islands $(m,Y)$ in $X$. Each such island is a minimal node, and all of its cells, of which there are $O(1)$, are counted as created cells of level $m$. These cells are paid for in the same way as created cells. The factor $\chi^c_m$ of the new field on $Y$ less a margin is present in the term. Consequently there is a test cube of level $m$ in which the test fails at $V_m(\xi)$; this is the test $\chi_m$ of a step, and not one of the sharper tests of the stages. Lemma 2.5 (i) then extracts the gain $c\,p_0(g_m)^2$ from the core Wilson factor of $X$, on plaquettes that are used by no other gain; when $X$ is a drop-out, the gain is extracted instead from the Wilson weight that (iv) counts together with its operation, and this holds at every evaluation. This test does not enter the mass of $Y$, since that mass is formed before the last decomposition. One evaluates the factor of $Y$ at a real pair of $U^c_m(Y^\sim,\alpha_{0,m},\alpha_{1,m})$, where it is bounded by its majorant. Since $\mathsf{v}$ is additive, the sum over the local histories of $Y$ separates and is at most $\mathsf{M}(m,Y)\le1$. Every leaf of the cut genealogy contains a created cell, and the ledger leaves $2\sum_t\mathsf{p}_t\mathsf{n}^{\mathrm{cr}}_t$, just as it does on the full genealogy.


\subsubsection*{(2d) Cancellation, localization and the sum over healed components}
\label{sec:actual-history-closure}

The main operation has two tasks. It must replace the prepared background
by the healed background without leaving a large Wilson action on the
shells. It must then sum the histories inside the healed components while
keeping the exterior data fixed. The local terms produced by these two
tasks have different bounds. A term confined to the collar of a component
has a bound by its stored volume. A term extending beyond that collar
decays in the boundary hull length. The volume term stays inside the
component integral, where the large-field mass pays for it.

Write $U''$ for the prepared background, $\bar U$ for the healed
background, and $U_0$ for the real critical configuration at the centre
of the reference law. Let $\phi_X$ cut off inside the treated component
$X$, let $\phi_{\rm c}$ cut off its old core, and put
$\phi=\phi_X(1-\phi_{\rm c})$. At a real point of the tests the
Wilson part satisfies
\[
\begin{split}
&-x_kA(\phi_X,U_{k,X}(e^{ig_kB}V_\Lambda))
 -A(x''_k(\cdot),U'')+x_kA(\phi_X,U_0)\\
&\qquad=-A(x_k(\cdot),\bar U)
 -A(x''_k\phi_{\rm c},U'')
 -\tfrac12\|\phi^{1/2}D\mathsf H B\|^2+\mathsf R_X .
\end{split}
\]
The positive term in the first line comes from the denominator of the
reference law. The linear term in $B$ vanishes because $U_0$ is critical
on the fibre. The second term on the right retains the Wilson weight
which pays for old large plaquettes. The quadratic form is nonnegative.

To estimate $\mathsf R_X$, insert the localized backgrounds of the
successive prepared sequences. The ten resulting terms are differences
of neighboring Wilson actions, Wilson actions with small weights, and
the remainder beyond the quadratic Taylor term. For each comparison the
two determining sets and their data agree in a collared region. The
plaquettes being compared lie at weighted distance at least
$cM\check R_k$ from its boundary. Exponential locality therefore gives
an error bounded by
\[
 C(1+\bar c)(\log\theta_k)^{2q+Cr}
 e^{-c\delta_0M(\log\theta_k)^r}.
\]
This tends to zero uniformly in the level and in the torus. The cubic
Taylor terms contain a positive power of $g_k$; the terms from the
shrinking core contain $L^{-c\mathsf N_k}$. The weight difference on
a shell is $O(k-j)$ rather than $O(x_k(k-j))$. Its growth is absorbed
by the weighted-distance decay across the layers.

The inner-core term has a useful complex continuation. Move the complex
part into the response from data supported outside the core. At the
stored real variables $\xi$ one obtains
\[
 A(x''_k\phi_{\rm c},U''(\mathbf U))
 =A(x''_k\phi_{\rm c},U''(\xi))+O(1).
\]
The first term on the right is the action of a $G$-valued configuration.
It is retained in the component measure. The response error has the
factor $L^{-\mathsf N_k}$ and is bounded even though the old shell
weights increase with the depth.

The remainders and the changes in earlier activities are localized by
one interpolation. Each walk in the response carries a parameter
$s(\Delta)$, $|s(\Delta)|<e^{\bar\kappa}$, for every cube it visits.
Setting a parameter to zero prevents a walk from crossing that cube.
For a term $f$ of bound $b_f$, interpolate its argument by
$\lambda(\mathbf H,\mathsf j)(s)$. The response and hereditary-domain
estimates give an analytic interpolation on
$|\lambda|<\mathsf a^{-1}$, where $\mathsf a$ is the size of the
response divided by the room in the domain of $f$. Its values at
$(s,\lambda)=(\mathbf1,1)$ and at $s=0$ are the old and new
evaluations. Brackets in the $s$ variables and Schwarz's lemma in
$\lambda$ then give
\[
 f(U^{\rm old})-f(U^{\rm new})=\sum_{Y\supset X}f_Y,
 \qquad
 |f_Y|\le2\mathsf a b_f
 e^{-(\bar\kappa-1)\mathfrak m(X,Y)}.
\]
Here $\mathfrak m$ counts the new cubes and the cubes needed to join
them to both the original term and the integrated field. The new
argument, its gauge factor and its stored data are formed together.
In particular, a stored field is not transported to the top of a tower
whose depth increases at every later step.

At a stage the amplitude is at most
$C\mathsf N_k\delta_j/\min_\iota\alpha_{\iota,j}$; at the main
step away from the box it is at most
$Cg_k^{-1}e^{-cM\check R_k}$. Inside a healed box the new field is
evaluated through its own generated background at its averages. The
critical-configuration identity identifies the arrival with the old
background. The corresponding gauge has leading potential given by an
average of the argument's potential. Thus it obeys the potential bound
required by the old term; an axial representative without this control
would not suffice.

Near a sharp test no analyticity in the stored field is claimed.
Instead the term is the value at $\mu_b=1$ of a function analytic for
$|\mu_b|<\lambda^{\rm st}$, with the switch properties that it does
not read the field when $\mu_b=0$ and does not read $\mu_b$ when the
field is zero. Block differences therefore gain
$(\lambda^{\rm st}-1)^{-1}$ per group of variables. Choosing
$\lambda^{\rm st}=(\log\theta_j)^{6r+1}$ pays for the
$O(\check R_k^4)$ cubes in a component and the $\mathsf N_k$ stages.
Once a stage integrates a stored field, its carried terms are evaluated
at field zero and its switch is removed. This preserves their bounds.

The sum of the stage logarithms is consequently bounded by
\[
 \sum_{\nu<\mathsf N_k}\sup|\mathbb C^{(\nu+1)}_k(Y)|
 \le C_{\mathbf S}
 e^{-(1+3\beta_{\rm s})\kappa d_k^\natural(\mathcal S;Y)}.
\]
The constant does not grow with the number of stages. Every older term
in a stage occurs through a difference, whose amplitude beats
$\mathsf N_k$; the stage normalizers have summable amplitudes.

Before summing healed components, complete the action at the output
label. If an old slot disappears, use the term formed from the history
with the whole treated cluster erased. Superposition makes this a
history of the construction. The erased history agrees with the original
one wherever the latter is carried unchanged. In this way the completed
action depends on the exterior data alone. The correction $V_{\mathfrak X}(Y)$
depends on the healed family $\mathfrak X$ only through its components
inside $Y$, their histories and the label near $Y$.

Fix these exterior data. The remaining expression is
\[
 \sum_{\mathfrak X\in\mathfrak F}
 \prod_{C\in\mathfrak X}\mathbf T'_k(C)
 \exp\sum_YV_{\mathfrak X}(Y).
\]
The operations act also on the exponential. Each integrates the stored
variables of its component, and $V(Y)$ reads those variables whenever
the component lies in $Y$. The component measure is the same in every
summand containing that component. If
$\mathfrak X\subset\mathfrak X'$ and
$\mathfrak X'(Y)=\mathfrak X(Y)$, then
$V_{\mathfrak X'}(Y)=V_{\mathfrak X}(Y)$.

Factor out the empty family. The Mayer expansion then has a healed
component in every nonempty polymer. Its mass supplies
$e^{-(2-c)p_0(g_k)}$, while the remaining local factors cost at most
exponentials of the component volume. The mass estimate below has
already prepaid those costs. The connected expansion converges and
gives one family, consistently on all restricted domains,
\[
 \log\mathsf K=\sum_Y\bigl(V_\emptyset(Y)+\mathbf R'_\sharp(Y)\bigr),
 \qquad
 |\mathbf R'_\sharp(Y)|
 \le C e^{-(1+c)p_0(g_k)}
 e^{-(1+\beta_{\rm s}/2)\kappa d_k^\natural(\mathcal S';Y)}.
\]
The same bound holds near standing islands. The part with no standing
object in its domain agrees with the term formed at the never-large
label and enters $R^{(k)}$; the rest is a boundary term.

The successive rates are
\[
 (1+4\beta_{\rm s})\kappa>
 (1+3\beta_{\rm s})\kappa>
 (1+2\beta_{\rm s})\kappa>
 (1+\tfrac12\beta_{\rm s})\kappa>\kappa.
\]
Each resummation spends only its indicated gap. An older boundary
constant enters a difference, a long domain, or the collar-volume
estimate under the component integral. It therefore changes the final
coupling ceiling, rather than the returned boundary constant.

\paragraph*{Stored operations and the norm propagated by the induction.}
We give the mass argument with its integration variables visible. This
also specifies the bound used in the preceding sum over healed
components. Fix a frame $(\theta_j)$ and put
\begin{equation}
 \mathsf l_j=\log\theta_j,\qquad
 g_j=\theta_j^{-1/2},\qquad R_j=\check R_j,
 \qquad P_j=p_0(g_j)=A_0\mathsf l_j^{\mathfrak p_0}.
 \label{ym4:exp500-history:scales}
\end{equation}
The constants below are independent of the torus, cutoff and height of
the history. They can depend on the geometric constants chosen earlier.
Whenever a prescribed extra volume weight is introduced, its coefficient
is fixed before the coupling ceiling is lowered.

For a component $X$ treated at level $i$, write the operation after
using its surviving averaging delta functions as
\begin{equation}
 \mathbf T'_i(X)F(\mathbf p)
 =\frac1{|W_4|}\sum_{\omega}
   \int_{S_\omega}e^{\Phi_\omega(\xi;\mathbf p)}
                            F\,d\mu_\omega(\xi).
 \label{ym4:exp500-history:operation}
\end{equation}
Here $\omega$ records the gauge-tree frame, all outcomes in the outer
genealogy, and recursively a local history for each inner island.
The variables $\xi$ are the fields still free after a surviving delta
function has been read as the definition
$V_{j+1}(c)=\mathcal M(V_j)(c)$. Their measure is a product of Haar
and Lebesgue measures. A Gaussian normalizer, gauge normalizer or
elimination determinant occurring in the operation remains a factor
of its integrand. It is not absorbed into a new probability measure.

The pair $\mathbf p=(\mathbf U,\mathbf J)$ belongs to the complex
domain on $X^\sim$ formed at level $i$. To read its exponent one may
use an admitted splitting $\mathbf U=e^{i\eta_j A'}U$, with $U$
real on the layer $j$. The action calculation above gives
\begin{equation}
 \Phi_\omega(\xi;\mathbf p)
   =\Phi^\circ_\omega(\xi)
      -\tfrac12q_X(\xi;U,A',\mathbf J)
      +V_X(\xi;U,A',\mathbf J),
 \qquad q_X\ge0,
 \label{ym4:exp500-history:exponent}
\end{equation}
where
$q_X=\|\phi^{1/2}D\mathsf H B\|^2$ is formed at the real glued
background. The two last summands can depend separately on the
splitting. Their sum, and hence the full exponent, depends only on the pair. The defining formula of the operation makes that exponent holomorphic on the pair domain. We apply Cauchy estimates to the complete integrand; the splitting-dependent terms are used for its modulus bound.

The part $\Phi^\circ_\omega$ contains the unexpanded inner Wilson
factor
\begin{equation}
 \exp\{-A(x''_i\phi_{\rm c},U''(\xi))\},
 \label{ym4:exp500-history:inner-wilson}
\end{equation}
the stored Gaussian and gauge weights, their normalizers, chart
densities, elimination factors, constants and the inner operations at
their real reads. The background $U''(\xi)$ in this formula is the
real inner background, obtained from the truncated prepared sequence
and its stored innermost datum. It is distinct from the glued
background used in $q_X$. Only the part outside the old core was
expanded in obtaining \eqref{ym4:exp500-history:exponent}. Thus the
old large plaquettes remain in
\eqref{ym4:exp500-history:inner-wilson} with their full positive
action.

The localization estimates for the main step and the earlier
activities give, on every admitted splitting,
\begin{equation}
 |V_X|\le C N_0(\omega),
 \qquad
 \sup_{\mathbf p}|e^{\Phi_\omega(\xi;\mathbf p)}|
       \le e^{\Phi^\circ_\omega(\xi)+C N_0(\omega)}.
 \label{ym4:exp500-history:pointwise-pair}
\end{equation}
Here $N_0(\omega)=\sum_j\#_j(\omega)$ counts the stored sites at
their own levels, including the assigned collars and the operation's
fixed coordinate multiplicities. Write $N_0^\circ(\omega)$ for the same
count in the outer record, before inserting its inner histories; thus
$N_0(\omega)=N_0^\circ(\omega)+\sum_yN_0(\omega'_y)$.
The bound in \eqref{ym4:exp500-history:pointwise-pair} follows by sorting the localized
terms into those confined to $X^\sim$, for which the rooted sum is
a constant per stored site, and those that leave that collar, for
which the boundary hull decay is summable. The response differences
in the inner operations are included in $V_X$ before this estimate
is taken. In particular an inner operation is read at a real
admissible pair in $\Phi^\circ_\omega$.

The right side of \eqref{ym4:exp500-history:pointwise-pair} depends
only on the component's own stored variables. That fact permits
composition at much later backgrounds. Taking an integral at a fixed
pair first would lose it: a pair at a later stage depends on the
stored fields of other components.
For an ordinary component put
\[
\begin{split}
 m_\omega(\xi)&=1_{S_\omega}(\xi)
       \sup_{\mathbf p}|e^{\Phi_\omega(\xi;\mathbf p)}|,\\
 m_\omega^{\mathbb R}(\xi)&=1_{S_\omega}(\xi)
       \sup_{\mathbf p\ \mathrm{real}}e^{\Phi_\omega(\xi;\mathbf p)},
\end{split}
\]
where both suprema use the admitted pair domain.
The mass is
\begin{equation}
 \mathsf M_c(i,X)
 =\frac1{|W_4|}\sum_\omega e^{c\mathsf v^\natural(\omega)}
    \int^*m_\omega(\xi)\,d\mu_\omega(\xi).
 \label{ym4:exp500-history:mass}
\end{equation}
The star denotes the infimum of the integrals of nonnegative Borel
majorants. For an ordinary fixed finite-dimensional core the
supremum is itself Borel: use continuity in the pair and a countable
dense set, and then take a countable supremum of Borel functions of
$\xi$. The upper integral also covers the varying evaluation fibres
below. All products are estimated by explicitly constructed Borel
majorants before Tonelli's theorem is applied.

\paragraph*{The evaluations used by a drop-out.}
A drop-out has no new inner background and no new core Wilson
factor. Its retained operation contains the tests, integrations,
gauge factors, complete kept forms and elimination factors left
when it stopped. The required action is supplied by the term
carrying it. Transfer to its factor the Wilson weight on the bare
plaquettes all of whose corners lie in $X$. This is an exact
redistribution of factors in the same integrand.

For fixed $\omega,\xi$ the admissible evaluations form a set
$\mathfrak K^\sharp(\omega,\xi)$. An evaluation includes a real
field $U$ on $X$ satisfying
\begin{equation}
 \mathcal M^jU(b)=V_j(\xi)(b)
 \quad\text{on every determining bond whose blocks lie in }X.
 \label{ym4:exp500-history:compatibility}
\end{equation}
At the bare level these equations use the stored bare field. The
field is regular layer by layer, as on the support of a history;
each recorded test has its recorded value, except for the last
test whose failure stopped this particular operation. Fields not
supplied by $\xi$ are read as the specified averages of $U$.
On each inner island $(m,Y)$ the pair
$(U,J(U))|_{Y^\sim}$ belongs to the real domain formed at its own
level. The evaluation also includes the kept complete forms and
Jacobian factors, with their proved one-sided bounds, and a real
Wilson weight $\upsilon$ with $\upsilon\ge c\theta_t$ on the
action layer owning a plaquette of level $t$.

The drop-out version of \eqref{ym4:exp500-history:mass} uses
\begin{equation}
 \mathsf E_\omega(\xi;\mathfrak e)
   =e^{-A(\upsilon1_X,U)}\mathsf K_\omega(\xi;\mathfrak e),
 \qquad
 \sup_{\mathfrak e\in\mathfrak K^\sharp(\omega,\xi)}
                                    \mathsf E_\omega(\xi;\mathfrak e)
 \label{ym4:exp500-history:dropout-factor}
\end{equation}
in place of the pair supremum: for a drop-out both $m_\omega$ and
$m_\omega^{\mathbb R}$ denote this fibre supremum multiplied by
$1_{S_\omega}$. Its value is zero when the fibre is empty.
The fibre depends on $\xi$ through
\eqref{ym4:exp500-history:compatibility} and its tests. In particular
a flat field cannot be chosen independently of stored nontrivial
averages. The Wilson field in
\eqref{ym4:exp500-history:dropout-factor} stays real. Complexifying
it would destroy the positive energy comparison used here.

There are two uses of the last failed test. It is excluded from the
mass of the operation that it stops. When this drop-out occurs
inside a later component, its record includes that test, and its
failure supplies the gain for the renewed wait. The latter
evaluation fibre is a restriction of the former one. No root gain
is charged in the earlier mass and then charged again in the later
one. If a particular operation carries no gauge-tree frame, insert
$1=|W_4|^{-1}\sum_{r\in W_4}1$; this keeps the same normalization
in \eqref{ym4:exp500-history:mass} for both kinds of mark.

These domains are fixed at the birth level. To see that every later
read lies in them, use the support estimates, in local level-$j$
units,
\begin{equation}
 |\partial U_k-1|\le C R_j^2\varepsilon_j\eta_j^2,
 \qquad |J(U_k)|\le C R_j^2\varepsilon_j,
 \qquad C R_j^2\varepsilon_j<\alpha_{0,j}.
 \label{ym4:exp500-history:later-reads}
\end{equation}
The last inequality follows from $q>\mathfrak p_0+2r$ after the
coupling ceiling is lowered. It has no loss in $k-j$. Influences
across layers at index distance $a$ carry $e^{-(a-1)_+}$, whereas
the threshold comparison costs at most $(1+a)^{1/2}$. Thus their
sum is bounded. The action sequence gives the retained Wilson
weight, and the determining sequence gives its compatible field.
Using either sequence in both places would not prove
\eqref{ym4:exp500-history:later-reads} for the actual operation.

\paragraph*{The action seen by the tests.}
We recall why failed tests of averages can be paid for by the
action of any compatible real field. This is needed for the whole
drop-out fibre, not only for one minimizing background.
For a real field $W$ put $f_W(p)=|\partial W(p)-1|$, using the
matrix operator norm. In $\mathrm{SU}(2)$,
\begin{equation}
 f_W(p)^2=2a_p(W).
 \label{ym4:exp500-history:su2-energy}
\end{equation}
On the four blocks at the corners of a coarse plaquette, a comb
gauge and the discrete Stokes formula give
\begin{equation}
 f_{\mathcal M W}\le \mathsf P f_W+\mathsf r(W),
 \qquad \|\mathsf r(W)\|_1\le C\|f_W\|_2^2.
 \label{ym4:exp500-history:average-one-step}
\end{equation}
This formula is used while the plaquettes in those blocks lie in
the fixed averaging chart. The matrix $\mathsf P$ is nonnegative.
Its row sums are $L^2$ and its column sums are $L^{-2}$ in four
dimensions. The row count is the area of a fine spanning surface;
the column count includes the $L^{-4}$ average over block
origins. For $j$ successive levels write
$\mathsf P^{[j]}=\mathsf P_{j-1}\cdots\mathsf P_0$, with each
factor acting from its fine plaquette grid to the next coarse
grid. Each factor has these same row and column bounds.
Thus Schur's inequality gives
\begin{equation}
 \|\mathsf P\|_{2\to2}\le1,
 \qquad \|\mathsf P\|_{1\to1}\le L^{-2},
 \qquad \|\mathsf P^{[j]}h\|_2\le L^{-2j}\|h\|_1.
 \label{ym4:exp500-history:average-schur}
\end{equation}
In the last inequality use $\|v\|_2\le\|v\|_1$ after applying the
column bound $j$ times. Restrictions to the relevant ancestor
blocks can only decrease these positive row and column sums.

Here is the uniform bootstrap. Let $E_0$ be the sum of $f_U^2$
over those ancestor blocks. Iterating
\eqref{ym4:exp500-history:average-one-step} along them and using
\eqref{ym4:exp500-history:average-schur} gives
\begin{equation}
 \|f_{\mathcal M^tU}\|_2
 \le E_0^{1/2}
       +C\sum_{l<t}L^{-2(t-1-l)}E_l,
 \qquad E_l=\|f_{\mathcal M^lU}\|_2^2.
 \label{ym4:exp500-history:average-bootstrap}
\end{equation}
Suppose successively that $E_l\le4E_0$ for $l<t$. The sum in
\eqref{ym4:exp500-history:average-bootstrap} is then at most
$4CE_0/(1-L^{-2})$. For $E_0$ below a fixed threshold it is at
most $E_0^{1/2}/2$. The bound $E_t\le4E_0$ follows with room
to spare. The same threshold keeps every intervening plaquette
inside its averaging chart, since its size is at most
$E_l^{1/2}$. This starts at level zero and proves the assertion
at every finite depth. Retaining the sharper right side also gives
$E_t^{1/2}\le E_0^{1/2}+C'E_0$.

Consequently, for a set $S$ of tested plaquettes of level $t$,
\begin{equation}
 \sum_{p\in\mathcal N_0(S)}a_p(U)
 \ge \min\left\{
       c\sum_{p'\in S}|\partial\mathcal M^tU(p')-1|^2,
       e_*\right\}.
 \label{ym4:exp500-history:energy-averages}
\end{equation}
Here $\mathcal N_0(S)$ is the union of the ancestor corner blocks
and $c,e_*>0$ depend on $L$, not on $t$. If the left side is
below $e_*$, the bootstrap proves the first bound; otherwise the
second alternative applies. No regularity of $U$ was assumed.
The small-energy alternative proves the regularity required to
apply the averaging charts. This is what makes the bound usable
at every compatible evaluation in
\eqref{ym4:exp500-history:dropout-factor}.

For a failed plaquette test at step $t$, linear regularity first
shows that some nearby plaquette of $V_t$ has size
$c\varepsilon_t$. Otherwise the localized background exists and
passes the test. Equation
\eqref{ym4:exp500-history:energy-averages} therefore supplies
$c\varepsilon_t^2$ from the compatible field's action. If the
approximate fluctuation test fails, the corresponding estimate is
\begin{equation}
 A_{\mathfrak c^\sim}(U)
       +\mathrm{GF}_{t-1}(\mathfrak c^\sim;V_{t-1})
                       \ge c\delta_{t-1}^2.
 \label{ym4:exp500-history:approximate-test}
\end{equation}
Indeed, smallness of its left side makes the coarse plaquettes and
the contour gauge variables $O(\sqrt\sigma)$. One-step rigidity,
using the average and that gauge, then bounds the difference from
the localized background by $C\sqrt\sigma$. This contradicts
failure when $\sigma<c\delta_{t-1}^2$. The gauge exponent is
essential here: a gauge transformation inside a block can move its
fine field while preserving both its Wilson action and its block
average.

For a failed fluctuation, collar or rim test, the stored Gaussian
supplies the energy. Choose $\kappa_{\rm pos}>0$ below both the
uniform real lower bound for the complete precision
\eqref{eq:actual-regular-precision} and the lower bound for its
kept conditional form in (2b), in the stored scaled coordinates.
If $a$ is the stored part and $B$ the part already integrated,
the complete conditional form satisfies
\begin{equation}
 \langle a,\Delta_{\mathsf Y'}a\rangle
 =\min_B\langle a+B,\mathsf A(a+B)\rangle
 \ge\kappa_{\rm pos}\|a\|^2.
 \label{ym4:exp500-history:stored-schur}
\end{equation}
The free coordinate groups are orthogonal, so the last inequality
follows from $\mathsf A\ge\kappa_{\rm pos}$. The actual kept
form has the stated lower bound by the separate cut-form
comparison in (2b). A failed
test forces a stored scaled coordinate of size at least
$c\delta_{t-1}/g_{t-1}$, with the stated fixed collar or rim
fraction included in $c$. On the cut-off strip the shift vanishes
when the integration variable is small. Thus the same implication
holds where the tested logarithm is not defined: a small integration
variable would put that logarithm in its chart and make the test
pass. Multiplication by $\theta_{t-1}$ in
\eqref{ym4:exp500-history:approximate-test}, or insertion in
\eqref{ym4:exp500-history:stored-schur}, gives a gain
$cP_t^2$ for a created cell. The energy pieces used by distinct
certificates are assigned on the cell neighborhoods fixed in (2a).
Their bounded overlaps are included in $c$ before the gains are
summed.

\paragraph*{A new mark pays outside the old floor.}
Let a component be treated at level $m$ with preparation bottom
$m_\flat$. Its old floor consists of the level-$m_\flat$ cells in
$\Omega_{m_\flat}^c\cap X$ that do not touch
$\Omega_{m_\flat}$. The plaquettes spent before the wait lie in
this floor. The gauge and stored Gaussian factors spent then
have levels below $m_\flat$. The waiting collar separates the
floor from the part on which the new operation acts.

We check a renewed test by contraposition. Suppose the sum of the
unspent Wilson, gauge and stored Gaussian energies in the new
window is below $c\mathsf l_m^{2\mathfrak p_1-8r}$. The energy
of averages and the gauge bounds make all corresponding stored
plaquettes and contour variables at most
\begin{equation}
 \varsigma=Cg_m\mathsf l_m^{\mathfrak p_1-4r}.
 \label{ym4:exp500-history:mark-smallness}
\end{equation}
The unknown field in the floor does not enter this estimate.
For a stage background $\widehat U$, the background locality
estimate through the waiting collar gives, outside the floor,
\begin{equation}
 |\partial\widehat U(p)-1|
 \le C\bigl(\varsigma+
       R_m^2\varepsilon_m e^{-\delta_0MR_m}\bigr)
                                      L^{-2j(p)}.
 \label{ym4:exp500-history:floor-screening}
\end{equation}
The index $j(p)$ is the local layer of $p$. The second term is
smaller than the first after the coupling ceiling is lowered.
This uses a collar of fixed positive width in its own local
units and the weighted-distance background estimate. It does
not require a small field inside the floor.

The most restrictive bond test is the main-step test on $V'$.
On a shell between two boxes, with side at most $CMR_m$, a loop
of $n$ straight pieces has holonomy within
$Cns^2\varpi$ of the identity if the shell plaquettes are within
$\varpi$. To obtain this estimate choose the shell reference
paths, move the loop across plaquettes in the shell, and use
telescoping products of unitary matrices. The required spanning
surfaces have at most $Cns^2$ plaquettes. The inner box sits
strictly inside the outer one in every direction, so the paths
and surfaces remain in the shell. The same reference gauge
extends a one-site-thick shell into the inner box by setting its
new bonds to the identity and transforming back; the resulting
plaquettes have size at most $Cs^2\varpi$.

Apply this to the averages of the compatible field and the basic
background. They agree on the gauge tree and outside the central
box. The tested $V'(b)$ is conjugate to the quotient of their two
loop holonomies, and hence
\begin{equation}
 |\tfrac1i\log V'(b)|
       \le CM^2R_m^2\varsigma<\delta'_m.
 \label{ym4:exp500-history:mark-bond}
\end{equation}
Two powers of $R_m$ pay for the shell area, and the other two in
\eqref{ym4:exp500-history:mark-smallness} absorb its fixed
constant. The other stage tests use one-step rigidity or the
same shell comparison with a smaller loss. Existence and extension
tests also pass: small exterior plaquettes give the extension just
described, and the localized background theorem then supplies
their background and logarithm charts. Thus failure of any test
provides, outside the floor, an energy witness of size at least
\begin{equation}
 c\mathsf l_m^{2\mathfrak p_1-8r}\ge P_m.
 \label{ym4:exp500-history:mark-gain}
\end{equation}
The exponent hierarchy $2\mathfrak p_1-8r>\mathfrak p_0$ gives
the last inequality at sufficiently small coupling.

For this root witness the compatible field may be needed on
$X^+$, two unit layers beyond $X$. We do not claim the witness
in the mass whose evaluation only supplies a field on $X$.
It is used in the later containing component, where that field
and the last failed test are present. Nested operation windows
are disjoint. Their gauge and Gaussian witnesses have different
levels, and their Wilson witnesses avoid the earlier floors.
This verifies the allocation needed to add the gains without
spending the same action twice.

\paragraph*{A volume with bounded multiplicity.}
A node $u=(n,C)$ is a component of the region made at level $n$.
An unhealed node has one successor, and the predecessors of a
node are disjoint. Let $\mathcal G(x)$ be the genealogy below
$x=(i,X)$, and cut it at the inner islands to form
$\mathcal G^\circ(x)$. The cut nodes are leaves of the outer
tree and roots of their own inner trees. Define
\begin{equation}
\begin{split}
 \operatorname{vol}(u)&=\mathsf l_{n(u)}(MR_{n(u)})^4
                                      |C(u)|_{n(u)},\\
 \mathsf V(x)&=\sum_{u\in\mathcal G(x)}\operatorname{vol}(u),
 \qquad
 \mathsf V^\circ(x)=\sum_{u\in\mathcal G^\circ(x)}
                                             \operatorname{vol}(u),\\
 \mathsf v^\natural(x)&=\mathsf V^\circ(x)
                  +\sum_{y\text{ inner}}\mathsf v^\natural(y).
\end{split}
 \label{ym4:exp500-history:volume}
\end{equation}
Expanding this finite recursion counts every uncut node once.
A non-root clearing node occurs twice, as the leaf of the outer
tree and as the root of its inner tree. No third outer tree
meets it. Consequently
\begin{equation}
             \mathsf V(x)\le\mathsf v^\natural(x)
                                      \le2\mathsf V(x).
 \label{ym4:exp500-history:volume-multiplicity}
\end{equation}
Keeping a drop-out does not make another copy of all its older
nodes. A component healed earlier and covered again later has
its own root; its old genealogy is not inserted in the new one
merely because their physical locations overlap.

The site count without logarithms, with these same multiplicities,
is at most $\mathsf v^\natural/\mathsf l_{\min}$. A fixed
coarse-$M$ read collar multiplies the site count of an owning
$MR_j$-cell by at most $(1+C/R_j)^4$. At large $R_j$ this is
at most two. Overlaps only decrease the union count. Thus the
number $n^\sharp$ of stored and owned read sites satisfies
\begin{equation}
 n^\sharp(\omega)\le
       \frac{2}{\mathsf l_{\min}}\mathsf v^\natural(\omega).
 \label{ym4:exp500-history:read-volume}
\end{equation}
Bond directions and color multiplicities are included in the
constant of a per-site energy estimate. The same ownership
shows that the part of a label still visible below a drop-out
obeys
\begin{equation}
 \mathsf V_{\rm in}(i,X;\mathcal S)
   :=\sum_{j<i}\mathsf l_j\#_j(Z_j^{\rm a}\cap X)
                          \le\mathsf v^\natural(\omega_X).
 \label{ym4:exp500-history:label-volume}
\end{equation}
The regions on the left are action regions. The determining
regions may have redone some of these levels and cannot replace
them in this count.

\paragraph*{Contraction and the count over generations.}
The geometric count is the step that converts the individual
event gains into a uniform mass. Let $O(u)$ consist of the
images, in the cell grid of $u$, of its earlier creation cells.
The actual hull, enlargement and bounded-size boxing rules give
\begin{equation}
 O(u)\subset C(u)\subset O(u)^{[e]},\qquad
 O(u)\text{ is }h_0\text{-connected},\qquad
 |C(u)|_{n(u)}\le C_e|O(u)|_{n(u)}.
 \label{ym4:exp500-history:skeleton}
\end{equation}
All constants are fixed. The layers added at earlier levels have
a bounded total reach in current cell units: every two levels
contract the earlier reach by $L^{-1}$, since size drops are
not consecutive. The sum is a geometric series. Boxing affects
only components of bounded side, and adds a further fixed reach.
Mergers take unions of these sets with touching collars. These
facts prove \eqref{ym4:exp500-history:skeleton}; a large
rectangular hull of an arbitrary thin component is not taken.

For an $h_0$-connected set $A$ with $N$ cells, its image under a
coarsening of ratio $q$ has at most
\begin{equation}
                  C_0(1+h_0N/q)
 \label{ym4:exp500-history:coarse-count}
\end{equation}
cells, including a fixed coarse collar. Traverse a spanning tree
of $A$ in both directions, joining neighboring cell centers.
The walk has length at most $2h_0(N-1)$. Divide it into pieces
of length at most $q$. Each piece meets a fixed number of
coarse cells in four dimensions. Their number is at most
$1+2h_0N/q$, proving
\eqref{ym4:exp500-history:coarse-count} with its collar.

Choose a fixed number $s_0$ of levels so that their net
coarsening makes the coefficient of $N$ at most $1/4$, and
choose $\nu_*\ge4C_0$. A node with at most $\nu_*$ occupied
cells is called light. A heavy component propagated for $s_0$
levels without new material loses at least half its occupied
cells. A light component contributes only a fixed number.
If $o_n$ is the total occupied count of a genealogy at level
$n$, $l_n$ its number of light nodes, and $n_t^{\rm cr}$ its
number of created cells, we obtain
\begin{equation}
 o_n\le\tfrac12o_{n-s_0}+C_1l_{n-s_0}
                   +C_1\sum_{n-s_0<t\le n}n_t^{\rm cr}.
 \label{ym4:exp500-history:occupied-recurrence}
\end{equation}
Unions at a merger only decrease this count. At a cut through an inner island $(m,Y)$, the last ordinary
plaquette test $\chi_m$ fails at threshold $\varepsilon_m$.
The energy-of-averages estimate supplies $cP_m^2$ from the outer
component's retained core Wilson action or, for a containing drop-out,
from its carrying-term Wilson weight at every compatible evaluation.
The plaquettes have not been spent on another witness. This test is
excluded from the mass of $Y$.
The uniformly bounded number $|Y|_m\le N_*$ of cut cells is counted
as creations. Choose $c_G$ small and the floor large so that
$N_*G_m+P_m\le c_{\rm w}P_m^2$, where $c_{\rm w}$ is the reserved
witness constant. This pays the cut cells and any mark credit from
this single witness. Keeping drop-outs are not cut: their renewed
marks retain the credit $P_m$ already assigned above.

For any nonnegative non-increasing sequence $r_n$, multiply
\eqref{ym4:exp500-history:occupied-recurrence} by $r_n$ and
sum over the finite levels. The shifted first sum is at most
half of $\sum_nr_no_n$. A creation occurs in at most $s_0$
of the last summands, with weights at most its birth weight.
Thus
\begin{equation}
 \sum_nr_no_n\le2C_1\sum_nr_nl_n
                   +2C_1s_0\sum_tr_tn_t^{\rm cr}.
 \label{ym4:exp500-history:weighted-contraction}
\end{equation}
No term proportional to the starting height occurs.

It remains to count light nodes. A quiet chain has one
predecessor at each step, no new creation, and no preceding
mark. Its side lengths satisfy
\begin{equation}
 \alpha_{n+1}\le\alpha_n+A,
 \qquad\alpha_{n+2}\le L^{-1}\alpha_n+2A.
 \label{ym4:exp500-history:quiet-contraction}
\end{equation}
The fixed point $2A/(1-L^{-1})$ lies strictly below the
small-component threshold. A light node has bounded diameter
by \eqref{ym4:exp500-history:skeleton}; after a fixed
number $b_0$ of quiet levels it is therefore a single small
box. After another preparation depth, its window lies wholly
after the event that began the chain. The depth, quietness,
quarantine and single-box trace conditions now hold, and the
dynamic class includes the component. It cannot be withheld.
The number of light nodes on a chain starting at level $t$
is consequently at most
\begin{equation}
              T_t=b_0+\overline N_t+3\le C_TR_t,
 \label{ym4:exp500-history:light-wait}
\end{equation}
where $T_t$ has a non-increasing choice. Heavy nodes before
this light segment have already been paid for in
\eqref{ym4:exp500-history:weighted-contraction}.

A quiet chain starts at a creation, after a mark, or at a
merger. The first two are charged to their creation or mark.
For a merger choose a non-first predecessor and descend through
first predecessors until reaching a leaf. Different mergers
receive different leaves. For incomparable mergers this follows
from disjointness of their predecessor trees. For comparable
mergers, a first-predecessor descent from the higher merger
cannot take the non-first branch selected at the lower merger.
Each leaf contains a created cell. A creation pays once for
its own chain and at most once for an assigned merger.

Since $T_tr_t$ is non-increasing, payment at the earlier
assigned event only increases the available weight. Combining
this assignment with
\eqref{ym4:exp500-history:weighted-contraction} gives
\begin{equation}
 \sum_{u\in\mathcal G(x)}r_{n(u)}|C(u)|_{n(u)}
  \le C_2\sum_tT_tr_t
                  \bigl(n_t^{\rm cr}+n_t^{\rm mk}\bigr).
 \label{ym4:exp500-history:genealogy-count}
\end{equation}
The root's own mark is excluded from $n_t^{\rm mk}$, since
it has not yet begun a renewed wait. On a torus too small
for a required small-box collar to embed, a bounded periodic
cell side leaves only a bounded number of admissible further
coarsenings. Increasing the fixed part of $T_t$ includes
these terminal levels. This alternative uses no wrapped-box
background.

\paragraph*{The scalar reserve.}
Let $c_t$ be the complete cost per level-$t$ cell, including
the normalizer, local energy, counting and prescribed volume
weight prices. A non-increasing majorant is
\begin{equation}
 c_t\le r_t=C_3c_1M^4R_t^5.
 \label{ym4:exp500-history:cell-price}
\end{equation}
Indeed the cost before this enlargement is
$c_1\mathsf l_t(MR_t)^4$, and
$\mathsf l_t\le CR_t$. Use a creation credit
$G_t=c_GP_t^2$ and a conservative mark credit $P_t$.
Equation \eqref{ym4:exp500-history:genealogy-count} shows
that the balance
\begin{equation}
 \mathfrak F(x)
 =\sum_tG_tn_t^{\rm cr}+\sum_tP_tn_t^{\rm mk}
              -\sum_{u\in\mathcal G(x)}c_{n(u)}|C(u)|_{n(u)}
 \label{ym4:exp500-history:balance}
\end{equation}
satisfies
\begin{equation}
 \mathfrak F(x)\ge2\sum_tP_tn_t^{\rm cr}
 \quad\text{if}\quad
 G_t\ge2P_t+C_2T_tr_t,
 \qquad P_t\ge C_2T_tr_t.
 \label{ym4:exp500-history:reserve}
\end{equation}
The mark can spend its conservative credit on the new wait;
the remaining creation credit supplies the small mass.

Here the hierarchy is effective. Since $T_t\le C R_t$ and
$R_t\asymp\mathsf l_t^r$, the cost on the right is at most
$C c_1M^4\mathsf l_t^{6r}$. Take
$\mathfrak p_0\ge6r+1$, as permitted by the exponent order.
Dividing by $P_t$ gives a bound
$C'c_1\mathsf l_t^{-(\mathfrak p_0-6r)}$. Therefore a
floor
\begin{equation}
 \mathsf l_t\ge K(1+c_1),\qquad c_GA_0\mathsf l_t\ge4
 \label{ym4:exp500-history:log-floor}
\end{equation}
with sufficiently large fixed $K$ implies both inequalities
in \eqref{ym4:exp500-history:reserve}. These conditions are
one ceiling on the coupling, imposed after $c_1$. They are
uniform over all levels. In particular the coefficient of
an additional fixed volume weight only changes the floor
by an affine bound; no height-dependent ceiling is needed.

Every nonempty genealogy has a creation leaf. The bounded-error
monotonicity of the frame gives, for $t\le i$,
$P_t\ge(1+\beta_0)^{-1}P_i$ at a sufficiently high common
floor. Thus \eqref{ym4:exp500-history:reserve} leaves at least
$2(1+\beta_0)^{-1}P_i$. For an unfinished unmarked root $z=(j,Z)$ let $S$ denote one step
of deterministic quiet propagation of its cell set. Continue this
geometry for $q$ levels, stopping strictly before a new creation,
merger or completed wait, and at the terminal level if necessary.
Only the cell costs change. Applying
\eqref{ym4:exp500-history:reserve} to the extended genealogy,
whose first creation level is $t_0$, gives
\[
 \mathfrak F(z)-2P_{t_0}
   \ge\sum_{n=j+1}^{j+q}c_n|S^{n-j}Z|_n\ge0.
\]
This is the credit for the future quiet evolution of a live component.
It uses neither a new test nor the admission of a new event.

\paragraph*{The entropy of the literal histories.}
The count must include the level and position of an earlier
event. Make a graph of cell addresses together with their
levels. Join neighboring cells at the same level and each
cell to its parent at the next. Its degree is at most
\begin{equation}
                         D=L^4+3^4+2.
 \label{ym4:exp500-history:address-degree}
\end{equation}
At a size drop a parent has one child; otherwise it has at
most $L^4$. The active addresses of an outer component form
a connected set rooted at a fixed cell of $X$. A connected
set of $N$ vertices through that root has at most
$D^{2(N-1)}$ choices: order the neighbors, choose its
canonical spanning tree and follow the depth-first walk
that traverses every tree edge twice. The visited set
recovers the original set. Downward paths record earlier
birth levels and positions, so neither the height nor the
torus volume supplies another factor.

For a fixed active set the remaining code has at most
$e^{b\mathsf V^\circ}$ choices. At level $j$ there are
$O((MR_j)^4)$ local sites and bonds per cell, with a fixed
finite list of test outcomes at each. The stage position
and failure type have at most a fixed power of $R_j+1$
choices; their logarithm is $O(\mathsf l_j)$. Stored
variable types, seam factors and class entrances use the
same local addresses. The normalized frame sums contribute
one, including the root frame and the frames inside the
outer record. Since $\mathsf V^\circ\ge N$, choosing
\begin{equation}
 c_{\rm e}\ge b+2\log D+\log2
 \label{ym4:exp500-history:entropy-price}
\end{equation}
gives
\begin{equation}
 \sum_{\omega^\circ}|W_4|^{-1-f(\omega^\circ)}
                       e^{-c_{\rm e}\mathsf V^\circ(\omega^\circ)}
 \le\sum_{N\ge1}D^{2(N-1)}e^{-(c_{\rm e}-b)N}\le1.
 \label{ym4:exp500-history:entropy}
\end{equation}
The integer $f$ counts the additional outer frames. The
component is nonempty, which explains the starting value
$N=1$. The identity operation for an empty component has
no rare-mass assertion. The bound by one, rather than by
an unspecified constant, is useful in the induction.

\paragraph*{The common envelope and the mass induction.}
We now combine the field estimates with this count. First
bound every nonquadratic term of the complete integrand
on the support by its proved cost per owned site. This
includes the boundary terms involving stored fluctuation
fields. Divide the positive action, gauge and Gaussian exponents in fixed
proportions. After decreasing the fixed $c_G$ and raising the coupling
floor, the witness share supplies a credit
$\Gamma_e=G_t=c_GP_t^2$ for each outer creation at level $t$, and
$\Gamma_e=P_t$ for each outer renewed mark. The root's own mark is
excluded, as in \eqref{ym4:exp500-history:balance}. The creation and
mark certificates above prove these bounds on their owned pieces.
A separate fixed share supplies the integrable reference envelope;
it is not spent again on the event gains.

The scalar normalizers have the appropriate local prices.
For the truncated gauge factor,
\begin{equation}
 z(\theta,\epsilon)
   =\int_G1_{\{|u-1|<\epsilon\}}
              e^{-\theta(1-\frac12\operatorname{tr}u)}\,du,
 \quad
 c\theta^{-3/2}\le z(\theta,\epsilon)\le C\theta^{-3/2}
 \label{ym4:exp500-history:gauge-normalizer}
\end{equation}
when $\epsilon\le1$ and $\theta\epsilon^2\ge1$.
Integrating a Lie-algebra ball of radius $c\theta^{-1/2}$
proves the lower bound; the radial Haar density and the
quadratic lower bound on $1-\frac12\operatorname{tr}u$
prove the upper bound. The same upper bound with
$\theta/2$ shows that half the exponent integrates
against the original normalizer at a bounded cost.
An elimination determinant and its chart Jacobian cost
a fixed factor per owning coordinate; their inverse
determinants are those of the actual pivot charts.
Rescaling $g_j$ and the remaining normalizers cost at
most $C\mathsf l_j$ per level-$j$ site.

For stored Gaussian blocks, use the complete kept forms
and the actual free coordinates of the constraint chart.
Their coercivity and the local reference-coordinate
comparison in (2b) leave an envelope $g_{\omega^\circ}$
whose conditional integral satisfies
\begin{equation}
 \sup_{\xi_{\rm in}}
       \int g_{\omega^\circ}(\xi^\circ;\xi_{\rm in})
                                    \,d\mu^\circ(\xi^\circ)
                         \le e^{C_5\mathsf V^\circ}.
 \label{ym4:exp500-history:reference-integral}
\end{equation}
One can see the local cost directly in the triangular
stored-coordinate representation. At fixed later
variables replace each earlier free block $b_j$ by
its actual shifted coordinate $A_j$. The inverse
Jacobian is bounded by $J_0^{d_j}$ on that chart.
A translation has $J_0=1$. Dropping restrictions
then gives the full Euclidean Gaussian integral,
\begin{equation}
 \int e^{-\beta\sum_j|A_j|^2}\,d\mu^\circ
 \le J_0^{\sum d_j}(\pi/\beta)^{\frac12\sum d_j}
 \le e^{C_5\mathsf V^\circ}.
 \label{ym4:exp500-history:triangular-gaussian}
\end{equation}
The dimension count is the count of stored scalar
coordinates, not the number of all exterior coordinates
used to construct their covariance. This retains
uniformity in the exterior volume. For a localized
marginal the same estimate follows from its local
Gaussian comparison. In both forms it is the actual
complete conditional form, including its compensating
gauge potential, that is bounded. A flat covariance
with the same formal name would not establish
\eqref{ym4:exp500-history:reference-integral}.

At every inner island replace the real admissible
read by its strong real majorant. With the scalar
and local energy prices just listed, the complete
outer factor is bounded by
\begin{equation}
\begin{split}
 m_\omega(\xi)
 &\le |W_4|^{-f(\omega^\circ)}
       e^{C_7\mathsf V^\circ-\sum_e\Gamma_e}
       g_{\omega^\circ}(\xi^\circ;\xi_{\rm in})
           \prod_y\frac{m^{\mathbb R}_{\omega'_y}(\xi'_y)}{|W_4|},
 \\
 \sup_{\xi_{\rm in}}\int g_{\omega^\circ}\,d\mu^\circ
 &\le e^{C_5\mathsf V^\circ}.
\end{split}
 \label{ym4:exp500-history:complete-envelope}
\end{equation}
For a drop-out the identical envelope holds at every
evaluation in its fibre. Its compatible fields obey
the same certificates
\eqref{ym4:exp500-history:energy-averages}--
\eqref{ym4:exp500-history:mark-gain}; its complete kept
forms and Jacobians obey the same one-sided bounds.
The envelope is therefore formed before the fibre
supremum. Replacing the inner suprema by nonnegative Borel majorants makes
its right side Borel on the full product of stored coordinates.
The upper integrals permit these majorants to be chosen with
arbitrarily small integral excess. Tonelli is applied to those
majorants, without asserting that a varying-fibre supremum is Borel.
The integral excesses are sent to zero after the outer integration. This is stronger than separate
integral bounds at each exterior argument.

Fix $\omega^\circ$. Its inner islands, domains and
levels are now fixed, and their local histories can
be summed independently. Multiply
\eqref{ym4:exp500-history:complete-envelope} by
$e^{c\mathsf v^\natural}$. Additivity in
\eqref{ym4:exp500-history:volume} and Tonelli allow
the outer variables to be integrated first. The
remaining integrals and history sums are the
earlier strong real masses, including their own
frame normalizations. If these masses are at most
one, the result is
\begin{equation}
 \frac1{|W_4|}\sum_{\omega\text{ over }\omega^\circ}
  e^{c\mathsf v^\natural(\omega)}\int^*m_\omega\,d\mu_\omega
 \le |W_4|^{-1-f(\omega^\circ)}
     e^{(C_5+C_7+c)\mathsf V^\circ-\sum_e\Gamma_e}.
 \label{ym4:exp500-history:mass-step}
\end{equation}
No measurable maximizing pair or maximizing evaluation
has been selected. The supremum has not been moved
through an integral.

Apply the scalar reserve with
$c_1=C_5+C_7+c+c_{\rm e}$, enlarged for its fixed
coordinate multiplicities. The cut leaves' new
tests supply their outer creation witnesses, while
their old histories have just been integrated in
their own masses. The reserve therefore gives
\begin{equation}
 \sum_e\Gamma_e-(C_5+C_7+c+c_{\rm e})\mathsf V^\circ
                      \ge2\sum_tP_tn_t^{\rm cr}.
 \label{ym4:exp500-history:outer-reserve}
\end{equation}
There is a creation leaf. Comparing its level with
$i$ leaves $\mathsf a_*P_i$, where
$\mathsf a_*=2(1+\beta_0)^{-1}=2-c'$. Sum
\eqref{ym4:exp500-history:mass-step} over its outer
records and use \eqref{ym4:exp500-history:entropy}.
We obtain
\begin{equation}
                  \mathsf M_c(i,X)\le e^{-\mathsf a_*P_i}.
 \label{ym4:exp500-history:mass-conclusion}
\end{equation}
This proves the complex-pair bound from the earlier
strong real bounds. Restricting the supremum to
real pairs returns the induction hypothesis at
level $i$. At the first level the product of
inner masses is empty. The drop-out estimate
uses the same proof and needs no additional
inductive drop-out hypothesis. Constants are
returned unchanged at every level.

\paragraph*{Restriction and the healed histories.}
We also justify the history product used before
the mass estimate. Include in the read region
of a move every stored variable, determining
datum, action plaquette, gauge factor and
reference normalizer that it uses. A localized
background formed on a box reads that box;
the short path of one of its kernel terms
does not reduce this read region. Class
membership reads the entire preparation
window and its quarantine marks.

If two clusters never touch, their cell margin
exceeds the fixed reach of every test, pivot
and collar. Hulls and enlargement commute
with their union, and boxing acts on its
separate components. For the preparation
window there is one further point. If an
event created at level $s$ lies in the window
of a later level $t$, it has not healed at
an intermediate $i<t$. Such healing would require $s\le i-\mathsf N_i$. Since
$\mathsf N_t=\check R_t$ is non-increasing, however,
\begin{equation}
            i-\mathsf N_i<t-\mathsf N_t\le s.
 \label{ym4:exp500-history:window-no-healing}
\end{equation}
For a mark the first inequality is strict at
the lower endpoint as well. Thus the reads
that decide the class are still present in
their own cluster. The actual class contains
all eligible components; there is no global
choice withholding one of two disjoint
eligible components.

Order the moves by their production, reading
and deletion of stored slots, including the
just described class reads. Independent
moves have disjoint read regions, product
reference measures and disjoint output
substitutions. Positive kernels commute by
Tonelli; their admitted complex continuations
commute by Fubini using
\eqref{ym4:exp500-history:complete-envelope}.
A test is unchanged by a move that changes
none of its inputs. Swapping adjacent
independent moves therefore preserves the
whole operation. Induction along this order,
using the hereditary background identities,
shows that restriction to a union of full
clusters produces those same local moves
and their same normalizers.

Erasing healed clusters is consequently an
actual history, and inserting finitely many
compatible clusters recovers an actual
history with the same standing part.
Connected components of the erased record
give the inverse correspondence. At fixed
exterior data its local operations are
exactly \eqref{ym4:exp500-history:operation}.
This proves the product and sum used above
to form $\mathbf R'_\sharp$, including
their measures. It supplies no global
pointwise gas over arbitrary final fields;
the remaining standing operations must
still be integrated and counted.

\paragraph*{Action shares and the loaded label sum.}
For the latter count first cut a history
at its outermost marks. Above the cuts,
a piece $(t,D)$ is a component made at
step $t$ on which no $\mathbf R$ has
acted. It retains its one-step variables
$\xi_{t,D}$ and their gauge, Gaussian
and elimination factors
$\mathsf w_{t,D}$. Assign to it the
bare plaquettes within one cell of
its failed plaquette and approximate
fluctuation tests. Write this set as
$\mathcal P(N_{t,D})$.

These sets are disjoint. A piece
untouched by $\mathbf R$ stays in
the live region at every later
step, and new tests are made only
beyond the prescribed margin of
that older region. No outermost
mark meets these action shares.
On a share the older profiles have
completed their transition and
the newer ones have not begun
theirs, so its action weight is
exactly $x_{t-1}$. Extending $U_k$
by the stored bare field gives
\begin{equation}
 A(x_k(\cdot),U_k)\ge
 \sum_{(t,D)}x_{t-1}
       \sum_{p\in\mathcal P(N_{t,D})}a_p(U_k)\ge0.
 \label{ym4:exp500-history:action-shares}
\end{equation}
At these plaquettes the background
constraints and the surviving
averaging delta functions give
the same averages as the tested
fields. Hence the all-compatible
field estimates apply to this
$U_k$ as well. After reserving
the event payments and integrating
the remaining half-exponents,
one obtains Borel functions of
each piece's own variables with
\begin{equation}
\begin{split}
 &e^{-A_{\rm cr}}\prod_{(t,D)}\mathsf w_{t,D}
 \le e^{-c\sum_{(t,\mathfrak c)\text{ created}}P_t^2}
                       \prod_{(t,D)}\overline{\mathsf w}_{t,D}
                                                   (\xi_{t,D}),\\
 &\hspace{18mm}
 \int\overline{\mathsf w}_{t,D}\,dm_{t,D}
                                  \le e^{C\#_t(D)}.
\end{split}
 \label{ym4:exp500-history:piece-majorants}
\end{equation}
The Gaussian integration uses the
Schur minimum in
\eqref{ym4:exp500-history:stored-schur}.
The gauge normalizer uses half
its own exponent as in
\eqref{ym4:exp500-history:gauge-normalizer}.
The resulting $C$ has no
$\log\theta_t$ factor. All
logarithmic prices will enter
with the vacuum constants.

For a marked component replace
its actual later read by the
majorant in
\eqref{ym4:exp500-history:mass}.
The groups $\xi_{t,D}$ and
$\xi_X$ are disjoint. Their
majorants therefore multiply
by Tonelli, even though $U_k$
depends on all of them. Summing
the local histories at the marks
produces their masses. The
restrictions of histories that
give the same label have at
most a fixed exponential number
of choices per stored site;
that is the local code already
counted in
\eqref{ym4:exp500-history:entropy}.
We do not identify a label with
one prescribed set of failed
cells.

Put $N_k=|\mathbb T^{(k)}|$ and
define the required load
\begin{equation}
 \mathscr L_c(\mathcal S)
 =c\sum_{n\le k}|\Gamma_n|
  +\sum_{j\le k}
      (C_{\log}\log\theta_{j-1}+C_{\rm loc}+c)
                                            \#_j(Z_j^+).
 \label{ym4:exp500-history:label-load}
\end{equation}
The first term counts determining bonds; $Z_j^+$ is the
collared action defect. Assign each bond to an endpoint in its
determining layer, and then to that site's owned cell. There are
only finitely many incident bond directions. For an unchanged
layer, charge the resulting site address to its action cell;
for a redone layer, charge it to the local history of the mark
that owns the redone slot, at its original level. A one-step
conversion costs at most $L^4$, and collars have bounded
multiplicity. For nested marks the inner record already
accounts for the earlier history. Thus each owned level-site
address receives a bounded number of charges, independently of
the number of subsequent reads. The fixed incidence and collar
multiplicities are included in the cell price and in
$c^\sharp(c)$ below. Equations
\eqref{ym4:exp500-history:volume-multiplicity}--
\eqref{ym4:exp500-history:label-volume}
show that the load inside
a mark is paid by
$e^{(c^\sharp-c)\mathsf v^\natural}$
for a fixed affine
$c^\sharp(c)>c$.

The vacuum constants have
the same assignment. The
regular one-step number on
the full torus is a sum of
one common local number
over its sites. Subtracting
the label's number cancels
the sites where that label
still uses it. Only the
missing or redone sites
remain. The scalar gauge
and Gaussian normalizers
bound each by
$C_{\log}\log\theta_{j-1}+C_{\rm loc}$.
Stage numbers are assigned
to their own integration
sites. The cancellation
uses the actual same local
number, with its actual
sign, on both lists; bounds
for unrelated volume
constants would not suffice.

After this normalization
the total unmarked local
energy is at most
$C_{\rm sf}N_k$ plus the
assigned defect prices.
The summation over old
levels uses the regular
member's own vacuum and
quadratic subtraction,
as in the earlier
pointed cancellation;
it is not a sum of a
constant times every
fine-lattice volume.
The unspent real coarse
Wilson factor is at most
one and final Haar measure
is normalized. Inside
outermost marks use
\eqref{ym4:exp500-history:mass-conclusion}
at weight $c^\sharp$.
Above them use
\eqref{ym4:exp500-history:piece-majorants}
and the scalar reserve.
For unfinished live trees
use their future-wait
credit after
\eqref{ym4:exp500-history:log-floor}.

It remains to sum the
positive records. Root
every carried component
at a cell of the fixed
$M$-partition of the
last lattice. There
are at most $N_k/M^4$
such addresses. Earlier
healed components are
already in the local
energies and supply no
additional final roots.
After adding the load
price to $c_1$, the
connected-set code and
\eqref{ym4:exp500-history:entropy}
give integrated weights
$a_\gamma\ge0$ with
\begin{equation}
              \sum_{\gamma\text{ rooted at }z}a_\gamma\le1.
 \label{ym4:exp500-history:rooted-label-sum}
\end{equation}
Dropping compatibility
from a positive sum
then gives
\begin{equation}
 \sum_{\mathfrak F\text{ allowed}}
                  \prod_{\gamma\in\mathfrak F}a_\gamma
 \le\prod_\gamma(1+a_\gamma)
 \le\exp\!\left(\sum_\gamma a_\gamma\right)
 \le e^{N_k/M^4}.
 \label{ym4:exp500-history:forest-sum}
\end{equation}
This also includes all
local outcome codes
leading to a given
label. For any finite
cutoff the histories
are finite; increasing
finite partial sums
give the same positive
bound if a countable
local representation
is used.

Combining the estimates gives
\begin{equation}
 \sum_{\mathcal S}
   |e^{\mathcal E_k(\mathcal S)}|
   \tau^{\rm abs}_{\mathcal S,c}(\mathcal V_k)
                    e^{\mathscr L_c(\mathcal S)}
                         \le e^{C(c)N_k}.
 \label{ym4:exp500-history:weighted-label-conclusion}
\end{equation}
Here each outer retained
operation in $\tau^{\rm abs}$
has had its integrand
replaced by its modulus
times $e^{c\mathsf v^\natural}$.
The constant depends on
$c$ and the earlier
construction parameters,
and not on the cutoff,
torus or last coupling.
The surplus mass weight
pays the lower levels
still visible inside a
drop-out after their
creation gains have
already been spent in
its mass.

\paragraph*{The source excess uses the same reserve.}
The ordinary complex
source keeps all sharp
tests, reference measures
and background choices
at their unsourced
values. It changes the
complete exponent by
an excess $\Psi_\omega$.
At a main operation its
new part is
\begin{equation}
 \Psi_\omega^{\rm new}
 =A(\zeta''\phi_{\rm c},U''(\xi))
                 +V_X^{,w}-V_X,
 \qquad
 \Psi_\omega=\Psi_\omega^{\rm new}
       +\sum_y\Psi_{\omega'_y}(\xi'_y;\mathbf q_y(\xi)).
 \label{ym4:exp500-history:source-recursion}
\end{equation}
The running source profile
$\zeta''$ is distinct
from the core cutoff
$\phi_{\rm c}$. The
first action is at the
real inner background,
and the second difference
is formed by the same
localized rule as the
unsourced potential.
If an interpolation cuts
an inner island, that
rule forms its excess
recursively on the
admitted data box.
It does not keep an
unrelated partial sum
of its carrier factors.
For a drop-out there is
no new inner background
or new potential from
a completed main step.
Its retained excess is
the sum of its inner
excesses. The source
in the Wilson factor
supplied by its carrier
stays in that carrier's
action comparison. This
keeps the two types
of operation in the
same fixed unsourced
representation.

For completeness, the
per-site estimate is
obtained before summing
the nesting. Assign a
new local member $\psi_f$
to its owned carrier
$R_f$, and bound both
its value at a fixed
reference pair and its
oscillation on the
admitted carrier core.
Rooted decay with rate
strictly above the
incidence entropy gives
\begin{equation}
 \sum_{f:\,x\in R_f}
   \frac{\sup|\psi_f|}{|R_f|}\le C(\varrho).
 \label{ym4:exp500-history:source-site}
\end{equation}
Indeed, sum by tree
length $m$; the count
is $C_0C_1^m$, and
the bound is
$C e^{-\kappa m}$
with $\kappa>\log C_1$.
Any polynomial root
multiplicity is absorbed
by a smaller retained
decay rate. The reference
value is controlled by
the same vacuum and
scalar-action readouts;
an oscillation bound
alone would miss an
arbitrarily large
constant.

For the Wilson member
in
\eqref{ym4:exp500-history:source-recursion}
use $0\le a_p(U'')\le2$
where no small-field
information is present,
and the curvature radius
of its own layer where
it is present. Each
bare plaquette is charged
to its actual owned
level; it is not
counted as just one
last-level plaquette.
More explicitly, on
the layer $n$,
$|\partial U''-1|\le
 C R_n^2\varepsilon_n\eta_n^2$.
The factor $R_n^2\varepsilon_n$
is bounded at small
coupling. Equation
\eqref{ym4:exp500-history:su2-energy}
then bounds the bare
action density by
$C\eta_n^4$.
There are at most
$C\eta_n^{-4}$ such
bare plaquettes per
site of that layer.
The two powers cancel,
leaving $O(\varrho)$
per owned level-$n$
site, uniformly in
its depth. On the
stored bare layer
the compact-group
bound supplies the
same conclusion
without a curvature
estimate.
Together with
\eqref{ym4:exp500-history:source-site},
the formed estimates
for $V_X^{,w}-V_X$
give
\begin{equation}
 |\Psi_\omega^{\rm new}|
          \le f(\varrho)N_0^\circ(\omega),
 \qquad
 |\Psi_\omega|\le f(\varrho)n^\sharp(\omega).
 \label{ym4:exp500-history:source-bound}
\end{equation}
For the second inequality
expand the recursion in
\eqref{ym4:exp500-history:source-recursion}.
Each new part uses its
own sites and each
inner history occurs
once. The same constant
therefore returns at
every depth; it is not
added once per ancestor.

Choose the mass weight
so that $c\mathsf l_j\ge2f(\varrho)$
on every relevant level.
Then
\eqref{ym4:exp500-history:read-volume}
gives
\begin{equation}
 \sup|e^{\Phi_\omega+\Psi_\omega}|
 \le e^{c\mathsf v^\natural(\omega)}
                             \sup|e^{\Phi_\omega}|.
 \label{ym4:exp500-history:source-absorption}
\end{equation}
This bound is taken
inside every integral.
The same fixed Borel
envelope works uniformly
on a compact source
polydisc. Holomorphy
then follows from
Fubini and Morera,
with no parameter
dependence of the
test support or
reference measure.
The source contributes
an ordinary mass price;
no comparison of
derivatives of two
different histories
is involved.

Finally remove from
each sourced term its
own vacuum constant.
The change of the
remaining small-field
exponent is bounded
by $C(1+\varrho)$ per
determining-layer site,
and the differences
between its vacuum
number and the common
one have the defect
price in
\eqref{ym4:exp500-history:label-load}.
Apply
\eqref{ym4:exp500-history:weighted-label-conclusion}
at the weight enlarged
for
\eqref{ym4:exp500-history:source-absorption}.
It supplies the actual
total-variation sum
needed in
\eqref{oh:term-bound}--\eqref{oh:label-sum},
at the volume $N_k$
of the last lattice.
The partition floor
remains a separate
positive-label argument.
No location of the
complex zeros of the
full partition function
is used in this mass
or label calculation.

\subsubsection*{(2e) The induction closed: the order of the constants, the sum over histories, exactness}

\textbf{The class.} Let a frame $\theta_0,\dots,\theta_k$ be fixed (see the opening of Step 3). Each choice of level $j$ that depends on the coupling is made with $\theta_j$. These choices are the thresholds, the radii, the scaling $\mathcal{A}_j=\theta_j^{-1/2}\mathcal{A}^{\mathrm{s}}_j$, the gauge-fixing weight, the Gaussian forms of the stages of $\mathbf{R}$ and the reference laws. The cube size and the waiting depth $\check{R}_j=\mathsf{N}_j$ are made instead with the running minimum $\min_{i\le j}\theta_i$ (a frame can go up by as much as $c_5$, and taking the minimum ensures that $\check{R}_j$ does not increase). Consequently the tests, the regions, the histories, the domains of analyticity and the measures of the integral operations depend on $\theta$ only, and on no entry of the density. In the bounds of this part we put $g_j:=\theta_j^{-1/2}$, so that $g_j\le\gamma$ is the floor $\theta_j\ge\gamma^{-2}$ of a frame (at $\theta=x$ this is the $g_j=x_j^{-1/2}$ of Theorem 2.1). In this way the circle ``regions depend on couplings, couplings on activities, activities on regions'' is broken. The elements of $\mathfrak{C}_k(\theta)$ are presentations rather than densities: the families are not determined by the density, and locality in the history is a property of how the terms are placed. Each element is an array, specified simultaneously on every admissible torus and cutoff,

\begin{gather*}
\FitMath{\mathfrak{P}=\Big((x_j);\,(E^{(j)},\widehat{E}^{(j)},b_j);\,(R^{(j)},\widehat{R}^{(j)});\,\big(B^{(j)}_{\mathcal{S}}(X)\big);\,(\Phi_\omega)\Big)_{j\le k},} \\[3pt]
\FitMath{b_j=x_{j-1}-x_j=b[\widehat{E}^{(j)}]} \\[3pt]
\FitMath{(1\le j\le k;\,\text{at }j=0\text{ the one entry }x_0),}
\end{gather*}

whose density $\rho[\mathfrak{P}]$ has the form given in Theorem 2.1. In it, $\widehat{E}^{(j)},\widehat{R}^{(j)}$ denote single families on $\mathbb{Z}^4$ that coincide with the families on the torus on lifts of non-wrapping domains; $b[\cdot]$ denotes the coefficient of (1e); and $\Phi_\omega=\Phi^\circ_\omega(\xi)+\sum_f\Phi_{\omega,f}(\xi;\mathbf{p}|_{R_f})$ are the exponents of the integral operations attached to the islands. Membership in the class requires the following.
(i) \emph{the tube}: $|\mathsf{z}_j|\le\mathsf{t}_j\approx c\,\theta_j(\log\theta_j)^{-C}$, where $\mathsf{z}_j=\theta_j-x_j$. This relates the frame, which is given, to the weights of the array, and the choice of the frame ensures it. The tube is wide in absolute terms, but narrow enough that the discrepancy between the weight and the choices, regarded as a potential, is small: $|\mathsf{z}_j|\alpha_{0,j}\delta_j\le c(\log\theta_j)^{-c}$.
(ii) The bounds of Theorem 2.1 on $E^{(j)},R^{(j)},B^{(j)}$ hold, together with analyticity, reality, and invariance under a single fixed gauge action on the stored data. A boundary term is analytic as well in the switches $\mu_b$ that belong to those of its stored bonds lying close to a test, and in those of its stored fields that lie far from every test. Any two splittings of one pair give each term the same value. If the tests of a term take the value $1$, then each boundary term and each exponent has in its domain the triple $(U,0,J(U))$ formed from the background $U$.
(iii) \emph{identification}: $E^{(j)}$ coincides with the family produced from the earlier entries, in each of its two chains (see below), by the small-field step at the pair $(\theta_{j-1},x_{j-1})$, carried out on the whole lattice.
(iv) \emph{lattice covariance} holds in the forms stated below. Above all, the pointed totals are covariant under every isometry of $\mathbb{T}^{(j)}$, unit translations included (119, (2.29)).
(v) \emph{coherence}: whenever a chart (2a) identifies, scale by scale, what two histories record within one cell of $X$, it identifies the terms at $X$. This single clause contains locality in the history, covariance under translations and independence of the torus. For a standing operation, the chart identifies $\mu_\omega$ together with the Borel set carrying it, as well as its set of factors and $\Phi^\circ_\omega$. It also identifies the factors one by one: $\Phi_{\iota\omega,\iota f}(\iota\xi;\iota\mathbf{p})=\Phi_{\omega,f}(\xi;\mathbf{p})$, for arbitrary components and histories, provided that the chart $\iota$ transports both the record of the history at $R_f$ and $\xi$ on $R_f$.
(vi) For an island, $\mu_\omega$ is a Borel measure, positive and defined on a Borel set; $\Phi_\omega$ depends in a Borel way on $\xi$, is analytic in $\mathbf{p}$, and takes real values whenever the pair is real. Moreover, gauge invariance holds pointwise in $\xi$ and for each factor separately. The bounds of \emph{mass} (Proposition 2.7, with a reserve constant $c_{\mathrm{r}}\ge0$ in the weight) and of \emph{oscillation} (2d) hold, the coupling in them being taken at the level at which the island was made. The class imposes no bound on the exponent itself, which is large and negative, and no clause concerns a drop-out.
(vii) Each term, taken together with its stored variables, is a non-negative measure $\hat\tau_h$.
The form and the bounds alone do not give a class that a step preserves. The gain $L^{-c(n-j)}$ of the regular terms does not follow from $\|E^{(j)}\|_{j,\kappa}\le E_0$, and neither do $|\mathbf{E}_k|\le C\sum_n|\Gamma_n|+C\sum_j\#_j(Z_j^+)$ (below) and the input of the next step, which depend on that gain. Norm bounds alone do not imply age decay, because they allow marginal action contributions; the example below exhibits one before its own correct subtraction. Take the tents $h^{(j)}_z$ of (1d); then $\mathsf{Q}_z:=c\,\alpha_{0,j}^{-2}A(h^{(j)}_z,\cdot)$ defines an activity whose norm does not exceed $1$, and

\begin{gather*}
\FitMath{\sum_z\mathsf{Q}_z=c\,\alpha_{0,j}^{-2}A(U),} \\[3pt]
\FitMath{\alpha_{0,j}^{-2}=c\,\theta_j(\log\theta_j)^{-2q}:}
\end{gather*}

this is the inverse coupling $\theta_j$ itself, up to the factor $c(\log\theta_j)^{-2q}$. For a regular family satisfying the full hypotheses of Section~\ref{ch:pointed-hessian-cancellation}, its own extracted coefficient $b_j$ and its own tent action remove this marginal part. The Ward identities, full hypercubic tensor symmetry and response hypotheses identify its affine quadratic jet at each base point with that multiple of the tent action. Take the format of the complex class below, in which each site has its own member. In that format one can find arrays which obey all the bounds of clause (ii) and which, at each level, add no less than the same multiple of the action, with no part of it decaying as the age grows. Uncancelled marginal contributions of this kind are excluded by those structural hypotheses together with each member's own correct subtraction. Unit-translation covariance alone neither forces isotropy nor removes a marginal action term.

\textbf{Theorem 2.11 [the step preserves the class].} Suppose the following:
\begin{itemize}
\item $\mathfrak{P}\in\mathfrak{C}_k(\theta)$;
\item $\theta_{k+1}$ prolongs the frame;
\item all $g_j\le\gamma$, and they lie below a ceiling that depends on the reserve constant $c_{\mathrm{r}}$ and does not increase with it;
\item the next weight $x_{k+1}:=x_k-b[\widehat{E}^{(k+1)}]$ lies in the tube (here $\widehat{E}^{(k+1)}$ is formed first, by clause (iii), at the pair $(\theta_k,x_k)$, and its formation involves no choice of level $k+1$).
\end{itemize}

Then each move of step $k+1$ is defined. The identity $\mathbf{R}\mathbf{T}\rho[\mathfrak{P}]=\rho[\mathfrak{P}']$ holds separately for each standing part (2a), every part keeping its stored variables, and $\mathfrak{P}'\in\mathfrak{C}_{k+1}(\theta)$, the constants being unchanged. Moreover, $\mathfrak{P}'$ prolongs $\mathfrak{P}$. An array belongs to the class of level $k+1$ exactly when its cut at $k$ belongs to the class of level $k$ and, in addition, it satisfies the conditions of level $k+1$. Finally, the array whose single entry is $x_0=\beta$, with $|\beta-\theta_0|\le\mathsf{t}_0\le c\,\theta_0(\log\theta_0)^{-C}$, belongs to $\mathfrak{C}_0$. \emph{Not returned}: the hypothesis that the next weight belongs to the tube. At every step this hypothesis is provided by how the frame is chosen. For $\theta=x$ the deviation is $0$, and the requirement is that $x_{k+1}$ prolong the frame, which is the content of the flow inequalities (3a) and of the nesting of the cube sizes, a consequence of the upper flow inequality (3b). For the reference line it is supplied by the drift bound of (3c), below. (In Ba{\l}aban, Comm. Math. Phys. 119 (1988) 243, the second remark after Theorem 1 concerns the space of ``all densities satisfying the conditions of the inductive assumption'', and these conditions include (2.29), (2.32), (2.33).)

\textbf{Why the induction closes.} (a) Room: fix a fraction $\beta_{\mathrm{s}}$. The terms produced by $\mathbf{T}$ (2b) are defined on the radii $(1+\beta_{\mathrm{s}})\alpha_{i,k}$ and decay at the rate $(1+4\beta_{\mathrm{s}})\kappa$. Moreover, $e^{-\kappa d_j(X)}\le e^{-c\kappa d_j(X)}e^{-(1+3\beta_{\mathrm{s}})\kappa d_k(\widehat{X})}$ for $j<k$, where $\widehat{X}$ denotes the union of all level-$k$ cubes meeting $X$, taken without any added layer. The operation $\mathbf{R}$ uses up this room and nothing beyond it. (b) Drift: for $m<n$ one has $g_m\le(1+c)g_n\le C(n-m)^{1/2}g_m$. This does no harm over the $\mathsf{N}_k$ levels that $\mathbf{R}$ spans, because every threshold changes only by powers of $\log x_k$. (c) Order. Every constant is of one of four kinds:
\begin{itemize}
\item free upward: its range has no upper bound (the odd integers for $L$, the powers of $L$ for $M$);
\item free downward: its range is some $(0,c]$;
\item confined to a window;
\item determined by earlier constants: a function of those alone.
\end{itemize}

The text states finitely many conditions on the constants. Each condition is assigned to the last constant on which it depends, after every determined constant in it has been expressed through the earlier ones. For every admissible choice of the earlier constants, the condition is then, at that constant, of one of three kinds:
\begin{itemize}
\item a floor on a constant free upward, which holds from some value onward;
\item a ceiling on a constant free downward, which holds on some $(0,c]$;
\item an endpoint of a non-empty window.
\end{itemize}

A condition that follows from such conditions is not counted separately, with one exception: the requirement that a window be non-empty. That requirement is a floor or a ceiling on some earlier constant, and it counts at that earlier constant. Two further stipulations apply:
\begin{itemize}
\item Every required comparison of thresholds must retain its stated margin. Lowering a threshold changes both its passing set and the gain available from its failure, so it cannot be treated as preserving all the estimates automatically.
\item Some families of conditions are indexed by a quantified number, for instance the floor $c_a\ge C(b)\log L$, $C(b)\to\infty$, with an integer $b\ge1$. No single $c_a$ satisfies this for all $b$. Such a family belongs to the statement in which the number is quantified, and it is imposed only at finitely many values of that number, which are fixed before the constant.
\end{itemize}

Thus, apart from the constant it bounds, each of these conditions depends only on earlier constants, in the order

\begin{gather*}
\FitMath{(r,\mathfrak{p}_0,\mathfrak{p}_1,q,\kappa_0,\beta_{\mathrm{s}},\theta_{\mathrm{s}},\theta_c)\prec L\prec(M_1,R_1)\prec\kappa\prec(\bar\kappa,\kappa_1)\prec M_1'\prec M_2\prec(\Theta_1,\varpi_T)\prec M} \\[3pt]
\FitMath{\prec\sigma_\star\prec a_\Lambda\prec(A_1,A_0,C_0,C_1,\alpha_1,\alpha_0,E_0,M_0,\alpha',\mathsf{B})\prec(\underline\lambda,c_5,\bar B)\prec\gamma\prec(\varrho,c_a)\prec(g,\ell,\beta,K)}
\end{gather*}

Some of these constants require comment:
\begin{itemize}
\item $\underline\lambda=\tfrac14c_0$ and $c_5=2c_2$, where $c_2$ is the constant of Proposition 3.3, are fixed together with $L$; in the display they appear as members of the triple that defines the frames.
\item $\bar\kappa$ is the exponent of the decoupling parameters of (1d). The factor $e^{-\bar\kappa}$ of a link of the cluster expansion comes from the decay of the covariance across an $M$-cube.
\item $\delta_0M\gg\bar\kappa\gg L\kappa$: rates are chosen before the cube sizes $M_1'$, $M_2$, $M$. These are floors on $\bar\kappa$ and then on $M$, both chosen after $L$ and $\kappa$.
\item $M_1$ depends on $L$ alone, and $R_1$ depends on $L$ and $M_1$.
\item $\Theta_1\ge C/\theta_c$ and $\varpi_T\ge C$, and both are free upward.
\item $M\ge\max\{\varpi_T,L^2M_2\}$ and $A_0\gg A_1$.
\item The later ceilings on $g$ may depend on $c_a$ and on the dial $\mathsf{w}\in[0,1)$. The dial is a fraction that depends on no constant and is fixed together with the exponents.
\item The reserve constant $c_{\mathrm{r}}\ge0$ of (vi) is free upward and comes before the ceiling that depends on it.
\end{itemize}

Induction along this order now gives two facts, because finitely many floors hold together from some value onward, finitely many ceilings hold together on some $(0,c]$, and the finitely many endpoints of a window leave it non-empty owing to the conditions that count at earlier constants. Every admissible choice made up to a given position can be extended, and no later condition places a restriction on a constant chosen earlier. Adding one more floor or ceiling of the same kind leaves both facts true. The choice is therefore not circular, and a condition arising in the construction fits the order if, at the last of the constants on which it depends, it is of the kind named. A bound may contain a number that a statement returns, or a number given as a function of a whole group of constants. Any such number must be fixed before the constant that it bounds; the cycle cut below is a case of this. A representative condition on $M$ or on $\gamma$ asserts that $e^{-cM}$, $e^{-c\check{R}_k}$ or a power of $\log x_k$ dominates constants fixed earlier. \emph{The supremum form}: a condition imposed on all couplings of all runs reduces to a single ceiling on $\gamma$, namely $\sup_{x\ge\gamma^{-2}}h(x)\le c$ with $h\to0$, because the scale appears in it solely by way of $x_k$ (122 (1989) 175, p. 186: the conditions ``do not depend on any scale''). \emph{The cycle that is cut}: $\lambda_\sharp$ and $B_\sharp$ depend on $\gamma$ through $\sup|\mathsf{q}|$ (here $\mathsf{q}_i$ is the part of the increment $b_i$ contributed by the remainders; see below), whereas the ceiling $\gamma$ depends on the flow constants. Frames are therefore defined with constants $(\underline\lambda,c_5,\bar B)$ fixed before $\gamma$: $\underline\lambda(j-i)-c_5\le\theta_i-\theta_j\le\bar B(j-i)$ for $i\le j$. Hence a frame with first member $\theta_0$ has length at most $(\theta_0-\gamma^{-2}+c_5)/\underline\lambda$, in the same way that $K$ is bounded in terms of $\beta$ in (3a). In addition, $\underline\lambda\le\lambda_\sharp$ and $B_\sharp\le\bar B$. \emph{No caps}: there is no upper bound on $L,M,\kappa,q,\mathfrak{p}_0,c_a$, on the cutoff, the volume, an age or a number of steps, and no lower bound on a coupling. \emph{The role of each exponent}:
\begin{itemize}
\item $\mathfrak{p}_0-\mathfrak{p}_1>Cr$, for comparing the ratio $\delta'_j/\delta_j$ with the square root of the waiting depth, and for the accumulation of the stages ((2b), (2d));
\item $2\mathfrak{p}_1-Cr>\mathfrak{p}_0$, so that a failed stage test still gives $e^{-p_0}$;
\item $\mathfrak{p}_0>3Cr$, so that $p_0$ dominates $\kappa\check{R}^4$ and the cost $\check{R}^6$ of a wait, and so that the window left for $\mathfrak{p}_1$ by the two preceding conditions is non-empty;
\item $q-\mathfrak{p}_0>Cr$, so that $\mathsf{N}_k\mathsf{s}_k\to0$;
\item $q>C\mathfrak{p}_0$, so that the product of an amplitude and an oscillation of degree $2\mathfrak{p}_0$ remains bounded;
\item $\kappa_0>C$, for the sum over the scales.
\end{itemize}

(d) \emph{No circle.} New terms are built from old ones, and the old ones are bounded with the same constants that the new ones must satisfy. The circle is avoided because an old term enters again in one of three ways:
\begin{itemize}
\item unchanged;
\item through a difference, which carries a small factor. In $\mathbf{T}$ this factor is $\mathsf{s}_k$, the factor that the Cauchy estimate gives for a first difference in the fluctuation field (2d).
\item as a whole, in the main step of $\mathbf{R}$. In that case there are two possibilities. Either the domain of the term must stretch across $c\check{R}_kL^{k-j}$ cubes of the term's own level $j<k$, in which case its decay makes the term small. Alternatively, the domain lies inside an enlarged component (a collar $X^\sim$ of (2d)). There only a volume bound $C\sum_{j\le k}\#_j(Z_j\cap X)$ is available, each scale counted in its own unit, and its constant is chosen later. The list is given below under \emph{Boundary terms}.
\end{itemize}

\emph{Uniformity of the input.} A term of level $j$ is defined on the radii at which it was created, and step $n+1$ evaluates it up to the output radius of that step. Let $w_{n,j}$ be the square of the ratio of these two radii. Let $\mathsf{f}_\nu\approx\nu^CL^{-c\nu}$ be forgetting weights; in (1d) they are written $\omega_\nu$. Then

\begin{gather*}
\FitMath{w_{n,j}\approx\frac{\theta_j}{\theta_n}\Big(\frac{\log\theta_n}{\log\theta_j}\Big)^{2q}\le C(n+2-j),} \\[3pt]
\FitMath{\sum_{j\le n}\mathsf{f}_{n+1-j}\|E^{(j)}\|_{j,\kappa}+\sum_{j\le n}w_{n,j}\|R^{(j)}\|_{j,\kappa}\le N_\star}
\end{gather*}

uniformly in $n$ and in the frame. This is everything that the small-field step requires of its input. The two factors in the approximation above are never unfavourable at the same time: the upper flow bound gives $\theta_j/\theta_n\le1+(n-j)$, and the lower one gives $\theta_n\le\theta_j+c_5$. The regular terms are of size $O(1)$ at each level and gain $L^{-c(n-j)}$ through their two counterterms. The remainders gain nothing; they are summed by means of the coupling. \emph{Regular terms}: let $\tau_k=x_k/\theta_k$, and let $E_{\mathsf{N}}$ be the bound on the localised Gaussian normaliser. Then

\begin{gather*}
\FitMath{\|E^{(k+1)}\|_{k+1,\kappa}\le E_{\mathsf{N}}+C\big(|1-\tau_k|+\delta_k+\mathsf{s}_kN_\star+e^{-c\,p_0(g_k)^2}\big)\le E_0 ,}
\end{gather*}

so $E_0$ is determined by the one-loop normaliser alone, and $C$ depends neither on $E_0$ nor on $N_\star$. \emph{Boundary terms}: $\mathsf{B}$ is the maximum of the bounds coming from their origins (2d), and none of these bounds depends on $\mathsf{B}$. This holds because data bounded by $\mathsf{B}$ enter the step in four ways:
\begin{itemize}
\item In $\mathbf{T}$, they enter as differences in the fluctuation field, with $C\mathsf{B}\mathsf{s}_k\le1$.
\item In the stages, they enter as differences in the collar field, taken one block of bonds at a time. Each block that is switched off contributes a negative power of $\log\theta_k$, with an exponent larger than the one belonging to (number of stages)$\times$($M$-cubes of level $k$ of a component).
\item In the main step, outside the enlarged components, they enter either through the difference under the substitution, or as undifferenced functions of a level $j<k$. In the second case the domain must extend over $c\check{R}_kL^{k-j}$ cubes, and $\mathsf{B}e^{-c\check{R}_k}\le1$.
\item Inside an enlarged component there is no smallness, and they enter only through the volume bound on $V(X^\sim)$. The constant of this bound is chosen after $\mathsf{B}$, and the bound is paid by the ledger of (2c).
\end{itemize}

(e) \emph{No growth with the level.} A carried term is literally unchanged: it has the same function, the same domain of analyticity and the same tree length. The bounds of a standing island refer to the level at which it was made and do not involve $k$. Indeed, the supremum in the mass bound is taken over a fixed set, built from the data of the levels up to that level, and this set contains the arguments of all later levels. Hence the mass at level $k$ does not exceed the mass at the level at which the island was made ((2c), ``The bound has to reproduce itself'', (ii)). The sizes that are returned are maxima rather than sums.

% END MODULE chapters/05_large_field.tex


% BEGIN MODULE chapters/ordinary_source_formation_audit_v6.tex
\section{Ordinary source formation and volume bounds}
\label{ch:ordinary-source-formation-audit}

This section separates three estimates that have different mathematical
content: forming local interactions, exponentiating a polymer gas, and
integrating the resulting expression.  A convergent polymer expansion supplies
the second estimate.  To apply it to a renormalisation construction one must
first verify its activity bounds for the interactions that the construction
actually forms.  The elementary implications below specify the ordinary source
estimates needed for that verification.

\subsection*{Site-dependent members and the summation over old levels}

Fix an odd integer \(L\ge3\).  Let \((\theta_j)_{0\le j\le n}\) satisfy
\[
 \theta_j\ge Y,\qquad
 \lambda(n-j)-c_5\le\theta_j-\theta_n\le B(n-j),
 \qquad Y\ge2(c_5+1),
 \tag{OSF1}
\]
where \(\lambda>0\) and \(c_5\ge0\).  A pointed family at level \(j\) is a collection
\(G(X,P,z)\), where \(X\) is a localisation domain, \(z\in X\) is a base
point and \(P\) belongs to a specified complex field domain.  Its norm is
\[
 \|G\|_{j,\kappa}
   =\sup_{X,z,P}e^{\kappa d_j(X)}|G(X,P,z)|.
 \tag{OSF2}
\]
Any companion families used to define its quadratic coefficient are included
in this supremum.  The field domains and restrictions are part of the data;
holomorphy does not assert that a function extends beyond these domains.

At every site \(z\), permit a different family \(G_{j,z}\).  Each such family
comes with its own quadratic coefficient \(b(G_{j,z})\).  The cancellation invoked here uses the full hypotheses of
Section~\ref{ch:pointed-hessian-cancellation}, including gauge invariance,
the actual normalized response, absolute Hessian moments and covariance
of its \emph{pointed} total under unit translations and the full hypercubic
group. Under those hypotheses, subtracting its own coefficient times its
own tent action cancels its marginal affine quadratic term at that base point.  No covariance relating the chosen members at different
sites is required.  This distinction permits a scalar source to be frozen at
the value belonging to each site.

Here is a precise quantitative formulation of the input used in this
argument.  Put \(\epsilon_{jn}=L^{-(n-j)}\), and let \(\rho_{jn}\) be the
ratio between the field radius required at level \(n\) and the radius on
which the member of level \(j\) is given, enlarged to be at least one.  Suppose
\(\rho_{jn}^2\le C_\rho(1+n-j)\).  Assume a linear localisation rule with the
following estimates, including summability over its rooted domains.
For one base point of level \(j\), a normalised regular member contributes at
most
\[
 C\,\|G_{j,z}\|_{j,\kappa}\,
       \rho_{jn}^{\,3}\epsilon_{jn}^{\,4+\alpha}
 \qquad(0<\alpha<1)
 \tag{OSF3}
\]
to a window of level \(n\), after its vacuum and quadratic subtractions.
A member with its vacuum value alone subtracted contributes at most
\[
 C\,\|F_{j,z}\|_{j,\kappa}\,
       \rho_{jn}^{\,2}\epsilon_{jn}^{\,4}.
 \tag{OSF4}
\]
The constants include the permitted loss of spatial decay and the count of
domains through the base point.  The rule must evaluate every member within
its stated field domain for every allowed decoupling and window parameter.
This domain inclusion is a separate hypothesis of (OSF3)--(OSF4).

The factor \(\epsilon_{jn}^4\) has a direct dimensional meaning.  A window of
level \(n\) contains at most \(C\epsilon_{jn}^{-4}\) base points of level
\(j\).  Thus vacuum subtraction balances the number of base points in four
dimensions, whereas the quadratic subtraction must supply the additional
\(\epsilon_{jn}^{\alpha}\) to make old regular members summable.

\begin{proposition}[Uniform ordinary input]
Suppose the regular members obey
\[\sup_z\|G_{j,z}\|_{j,\kappa}\le G_\star.\]
Suppose defects and remainders
obey, respectively,
\[
 \|D_j\|_{j,\kappa}\le L^{-(n+1-j)},\qquad
 \|H_j\|_{j,\kappa}\le C_H\theta_j^{-p},\qquad p>2.
 \tag{OSF5}
\]
Under (OSF3)--(OSF4), their combined ordinary input at level \(n\) is bounded
by
\[
 C G_\star
 +C\sum_{a=0}^{n}(1+a)L^{-(a+1)}
 +C C_H\sum_{a=0}^{n}(1+a)(\theta_n+\lambda a-c_5)^{-p}
 \le C G_\star+C+C C_H\theta_n^{\,2-p}.
 \tag{OSF6}
\]
The constant is independent of \(n\), the number of sites and the choice of
the member at each site.
\end{proposition}

\begin{proof}
Sum (OSF3) over the base points in a window.  The inverse volume factor
cancels \(\epsilon_{jn}^4\), leaving
\[
 C G_\star\sum_{a\ge0}(1+a)^{3/2}L^{-\alpha a}<\infty.
\]
This uses only the supremum of the member norms, so choosing a different
covariant member at each site has no effect on the bound.  The same
cancellation in (OSF4) gives the second sum in (OSF6) for defects.  That sum
is bounded by a convergent geometric series and its first derivative.

For remainders use (OSF1) and put \(y=\theta_n-c_5\ge\theta_n/2\).
For \(a\ge0\), \(1+a\le C(y+\lambda a)\), since \(y\ge1\).
Comparison with the integral of the decreasing function
\((y+\lambda t)^{1-p}\) gives
\[
 \sum_{a\ge0}(1+a)(y+\lambda a)^{-p}
 \le C\left(y^{1-p}+\int_0^\infty(y+\lambda t)^{1-p}\,dt\right)
 \le C y^{2-p}.
 \tag{OSF7}
\]
This proves the stated, deliberately nonoptimal, bound.
\end{proof}

For \(SU(2)\), a bounded scalar action weight supplies another admissible
ordinary input.  If \(W\in SL(2,\mathbb C)\), then
\[
 1-\tfrac12\operatorname{tr}W=\tfrac12\det(1-W).
 \tag{OSF8}
\]
On a field domain where each plaquette satisfies
\(|W_p-1|\le\alpha_n\eta_n^2\), its action density is bounded by
\(C\alpha_n^2\eta_n^4\).  A local cube contains at most
\(C M^4\eta_n^{-4}\) plaquettes.  Filing each plaquette in a domain containing
its corners therefore gives a weight member of norm at most
\(C M^4\alpha_n^2\sup|u|\).  This proof uses the determinant identity and the
count of plaquettes; it does not use covariance of \(u\).  With
\(\alpha_n=O(\theta_n^{-1/2}(\log\theta_n)^q)\), any fixed bound on
\(\sup|u|\) is eventually admissible.

These summations do not prove (OSF3).  That estimate contains the
renormalisation assertion: covariance of the pointed Hessian, the correct
counterterm, regularity of the background within the window and uniform
membership after substitution must be proved for the actual formation rule.
They cannot be replaced by a bound on the full total or by holomorphy in the
source alone.

\subsection*{Exponential reserves and generated local perturbations}

Let a local history \(\omega\) have a positive measure \(\mu_\omega\) on
\(S_\omega\), a reference exponent \(\Phi_\omega(\xi;P)\), and a site count
\(N_\omega\).  Suppose its reserve volume \(V_\omega\) is additive over
the older histories swallowed by this history and satisfies
\[
 V_\omega\ge A N_\omega,\qquad
 \sum_\omega\int_{S_\omega}
      e^{c_\star V_\omega}\sup_P|e^{\Phi_\omega(\xi;P)}|
             \,d\mu_\omega(\xi)\le\varepsilon,\qquad c_\star\ge1.
 \tag{OSF9}
\]
The supremum in (OSF9) is assumed measurable, or is replaced by a measurable
majorant throughout.  If a holomorphic perturbation satisfies
\(|\Psi_\omega|\le fN_\omega\) and \(A\ge f\), then for bounded \(F\)
\[
 \left|\sum_\omega\int F e^{\Phi_\omega+\Psi_\omega}\,d\mu_\omega\right|
 \le\varepsilon\|F\|_\infty.
 \tag{OSF10}
\]
Indeed, \(|e^{\Psi_\omega}|\le e^{V_\omega}\), so take absolute values
inside every integral and apply (OSF9).  If each integrand is holomorphic
in the parameter and all reference measures and integration sets are
parameter independent, a common integrable majorant on compact parameter
sets also proves holomorphy by Fubini and Morera.
This is an ordinary complex mass estimate.  It requires no derivative
with respect to the number of source zones.

Suppose a formation rule produces a new excess exponent in the explicit form
\[
 \Psi_\omega^{\rm new}
   = A(\widehat c\,\upsilon1_X,U''(\xi))
       +V(X^\sim;\mathfrak L)-V(X^\sim;\mathfrak L^0),
 \tag{OSF11}
\]
where \(U''(\xi)\) is real, \(0\le\upsilon\le1\) and
\(\sup|\widehat c|\le\bar c\).  To apply (OSF10) it is sufficient to prove
the actual formation bounds
\[
 |A(\widehat c\,\upsilon1_X,U''(\xi))|\le C_W\bar c\,N_\omega^0,\qquad
 |V(X^\sim;\mathfrak L)|+|V(X^\sim;\mathfrak L^0)|
          \le2c_{\rm int}(1+\bar c)N_\omega^0.
 \tag{OSF12}
\]
If \(\Psi_\omega=\Psi_\omega^{\rm new}
+\sum_{\omega'}\Psi_{\omega'}\) and
\(N_\omega=N_\omega^0+\sum_{\omega'}N_{\omega'}\), these inequalities give
\[
 |\Psi_\omega|
 \le\max\{C_W\bar c+2c_{\rm int}(1+\bar c),f\}\,N_\omega.
 \tag{OSF13}
\]
The maximum, rather than a sum of the old and new constants, is what permits
unbounded nesting.  Preserving the carrier oscillation norm needed in the
next formation step is an additional estimate; it does not follow from
(OSF13).

There is a genuine reason to generate local perturbations through this
formation rule.  A bound on an exponent and on the oscillation of each
carrier does not bound arbitrary partial sums of carrier factors.  On a
one-point integration space, write
\(\Psi^\circ(\lambda)=-N\lambda\) and
\(\psi_1(\lambda;P)=N\lambda\), with a second, unused carrier if desired.
Then \(\Psi^\circ+\psi_1=0\), both carrier oscillations in \(P\) are zero,
and the full mass is unchanged for every \(\lambda\).  Interpolating only
\(\psi_1\), while leaving \(\Psi^\circ\) at its value at zero, instead gives
the mass \(e^{N\lambda}\).  No bound uniform in \(N\) survives on a fixed
complex disc.  This example refutes that partial-interpolation argument,
not a construction which defines its localized exponents recursively by
(OSF11) and proves (OSF12) on the whole complex data box.

\subsection*{What a pointwise polymer logarithm proves}

Consider a finite hard-core polymer gas \(\Xi_y\), depending on an exterior
field \(y\).  Suppose its activities \(h_y(C)\) satisfy a uniform
Koteck\'y--Preiss criterion, including absolute convergence of the logarithmic
series and the bound
\[
 |\log\Xi_y|\le\sum_C|h_y(C)|e^{a(C)}\le c\,N.
 \tag{OSF14}
\]
If the activities are real for real \(y\), then \(\Xi_y>0\).
One proof is that the absolutely convergent logarithmic series is real.
Equivalently, replace every activity by \(t h_y(C)\), \(0\le t\le1\).
The same criterion applies, so \(\Xi_y(t)\) is real, continuous and nonzero.
Its value at zero is one, proving positivity at one.  Individual activities
need not be positive.  Moreover \(e^{-cN}\le\Xi_y\le e^{cN}\).

If a final representation really has the form
\(Z=\int q(y)\Xi_y\,d\nu(y)\), with \(q\ge0\), then
\[
 e^{-cN}\int q\,d\nu\le Z\le e^{cN}\int q\,d\nu.
 \tag{OSF15}
\]
Thus even a pointwise polymer logarithm requires control of the remaining
integral.  It does not turn a restricted Haar integral into a probability
integral.  For example, take the empty gas and let \(q\) be the indicator of
a Haar set of measure \(e^{-N^2}\).  The polymer criterion is automatic,
but a lower bound \(e^{-CN}\) for \(Z\) is false.

In a representation that sums over standing labels, the useful upper
estimate instead has the form
\[
 \sum_{\mathcal S}
       \big|e^{\mathcal E_{\mathcal S}(w)}
                  \tau_{\mathcal S}^{\,w}\big|(\mathcal V_n)
       e^{W(\mathcal S)}
 \le e^{cN_n}.
 \tag{OSF16}
\]
Here the absolute value is total variation, \(W\) pays for missing vacuum
contributions and the measures inside the retained integral operations are
also majorised before they are integrated.  A pointwise bound for one gas
does not supply this sum: arbitrarily many positive labels can have the
same individually bounded gas.  The counting of labels and the exponential
reserve in their operations are essential.

For the lower bound at zero source, positivity permits retaining the
all-small-field label.  Suppose its density, after vacuum normalisation,
is at least
\(\chi(y)\exp[-x A(U(y))-C N_n]\).  Assume the ball in which every one of
the \(4N_n\) link variables has distance less than \(s\) from the identity
is contained in \(\{\chi=1\}\), and that
\(A(U(y))\le C s^2N_n\) there.  Normalized Haar measure on \(SU(2)\) assigns
this link ball mass at least \((c s^3)^{4N_n}\).  It follows directly that
\[
 e^{-E}Z(0)\ge
  \exp[-(C+4\log(1/(c s^3))+Cxs^2)N_n].
\]
Taking \(s=c_0x^{-1/2}\) gives
\(e^{-E}Z(0)\ge\exp[-(C+6\log x)N_n]\).
The nontrivial model input here is that the link ball passes every background
test and that the action of its background has the asserted bound.
No zero-free statement for the sourced full partition function is needed.

% END MODULE chapters/ordinary_source_formation_audit_v6.tex


% BEGIN MODULE chapters/flow_arithmetic_v2.tex
\section{First passage and the arithmetic of a marginal flow}
\label{ch:flow-arithmetic}

Assume the coupling recursion and uniform increment bounds below.
We derive their first-passage and scale-summation consequences,
allowing individual increments to be negative. The RG realization,
torus and cutoff compatibility, and nonconstant-source estimates
are included in the separate flow and source hypotheses.

\subsection{Hypotheses and first passage}

Fix numbers $c_0>0$, $C_0\ge0$, $h>0$, and
$0\le\epsilon<c_0/4$. For each initial value $\beta\ge h+B$ consider a real
run with $x_0=\beta$. Assume that a next step exists whenever $x_j>h$, and
that every such step has the form

\begin{gather*}
\FitMath{x_j=x_{j-1}-b_j,} \\[3pt]
\FitMath{b_j=b_j^0+r_j.}
\end{gather*}

For every interval of defined steps assume

\begin{gather*}
\FitMath{\left|\sum_{l=i+1}^j(b_l^0-c_0)\right|\le C_0,} \\[3pt]
\FitMath{|r_j|\le\epsilon} \\[3pt]
\FitMath{\hbox{if }x_{j-1}\ge h.}
\end{gather*}

All constants are independent of $\beta$ and of the length of the run. The
existence of the next step is an assumption: an inequality about hypothetical
increments cannot establish that an underlying analytic construction is
defined. Set

\begin{gather*}
\FitMath{\lambda=c_0-\epsilon>0,} \\[3pt]
\FitMath{\Lambda=c_0+\epsilon,} \\[3pt]
\FitMath{B\ge c_0+C_0+\epsilon.}
\end{gather*}

The one-step case of the interval estimate gives $b_j\le B$. This is an upper
bound, not an assertion $b_j\ge0$.

Define

\begin{gather*}
\FitMath{K(\beta)=\min\{k\ge0:x_k\le h+B\},} \\[3pt]
\FitMath{\mathcal W_k=\{\beta\ge h+B:K(\beta)=k\}.}
\end{gather*}

The minimum is initially a possible stopping index; its finiteness will be
proved. At the endpoint $\beta=h+B$ it equals zero. When $\beta>h+B$ it is
positive, so the last increment is an actual step of the run.

\begin{proposition}[Block drift and a protected terminal window]
The first-passage time is finite. For defined segments up to that time,

\begin{gather*}
\FitMath{\lambda(j-i)-C_0\le x_i-x_j
                    \le\Lambda(j-i)+C_0} \\[3pt]
\FitMath{(0\le i\le j\le K).}
\end{gather*}

If $K>0$, then $h<x_K\le h+B$. Moreover

\begin{gather*}
\FitMath{\max\left\{0,\frac{\beta-h-B-C_0}{\Lambda}\right\}
 \le K\le\frac{\beta-h+C_0}{\lambda}.}
\end{gather*}

\end{proposition}

\begin{proof}
Summing the recursion over a segment gives

\begin{gather*}
\FitMath{x_i-x_j=c_0(j-i)
       +\sum_{l=i+1}^j(b_l^0-c_0)+\sum_{l=i+1}^jr_l.}
\end{gather*}

The two assumed bounds give the block drift estimate. If the run never reached
$h+B$, all its values would remain above $h$, so it could be extended for
every finite number of steps. The lower drift bound would then imply
$x_j\le\beta-\lambda j+C_0$, contradicting $x_j>h+B$ for sufficiently large
$j$. Thus first passage occurs.

For $K>0$, minimality gives $x_{K-1}>h+B$. Since $b_K\le B$,
$x_K=x_{K-1}-b_K>h$. The upper endpoint follows from the definition of $K$.
The extra width $B$ is exactly what prevents the final step from overshooting
the lower reference value $h$.

Finally $\beta-x_K\le\Lambda K+C_0$ and $x_K\le h+B$ give the lower bound
on $K$. Conversely $\lambda K-C_0\le\beta-x_K<\beta-h$ gives the stated
upper bound for $K>0$. The endpoint case $K=0$ also satisfies both inequalities.
\end{proof}

In particular $\beta\to\infty$ forces $K(\beta)\to\infty$. Conversely a
sequence with $K\to\infty$ has $\beta\to\infty$, by the lower drift estimate
and the terminal lower bound. Since these stopping times are integers, the
sets $\mathcal W_k$ are nonempty for infinitely many different $k$. Nothing
here says that every one is nonempty, that one is an interval, or that the
terminal values in it coincide. Such assertions would require additional
information about the dependence of the analytic recursion on $\beta$.

For example, $b_j^0=c_0+(-1)^jq$, with $q>c_0$ and $r_j=0$, satisfies the
interval hypothesis with $C_0=q$. Every other increment is negative.
Nevertheless the block drift is positive and the preceding proof applies.
Replacing block drift by a silently assumed monotonicity would discard this
perfectly admissible numerical example.

\subsection{Vanishing mean error and the leading logarithm}

Assume in addition that

\begin{gather*}
\FitMath{|r_j|\le\omega(x_{j-1}),} \\[3pt]
\FitMath{\omega_* (M):=\sup_{x\ge M}\omega(x)\longrightarrow0} \\[3pt]
\FitMath{(M\to\infty),}
\end{gather*}

where $\omega$ is nonnegative on $[h,\infty)$. The earlier global bound by
$\epsilon$ remains in force. No monotonicity of $\omega$ is needed.

\begin{lemma}[Time spent below a fixed inverse coupling]
For $M\ge h$, the number $N_M$ of steps $1\le j\le K$ with $x_{j-1}\le M$
is bounded by

\begin{gather*}
\FitMath{N_M\le1+\frac{M-h+C_0}{\lambda}.}
\end{gather*}

This bound does not depend on the initial value.
\end{lemma}

\begin{proof}
There is nothing to prove if no such step occurs. Otherwise let $i$ be the
smallest index in $\{0,\ldots,K-1\}$ with $x_i\le M$. Every counted step
has its preceding index between $i$ and $K-1$, so $N_M\le K-i$. The block
drift from $i$ to $K$ gives
$\lambda(K-i)-C_0\le x_i-x_K<M-h$. The displayed, slightly weaker bound
therefore follows. Excursions back above $M$ cause no problem: all subsequent
indices are included in this bound on the entire remaining time.
\end{proof}

\begin{theorem}[Leading first-passage asymptotics]
With $h$ and the other constants fixed,

\begin{gather*}
\FitMath{\frac1{K(\beta)}\sum_{j=1}^{K(\beta)}|r_j|\longrightarrow0,} \\[3pt]
\FitMath{\frac{\beta}{K(\beta)}\longrightarrow c_0} \\[3pt]
\FitMath{(\beta\to\infty).}
\end{gather*}

If $a_K=L^{-K}$ for a fixed $L>1$, then
$\beta/\log(a_K^{-1})\to c_0/\log L$.
\end{theorem}

\begin{proof}
Split the steps according to whether $x_{j-1}\le M$. The lemma gives

\begin{gather*}
\FitMath{\frac1K\sum_{j=1}^K|r_j|
 \le\frac{\epsilon}{K}
             \left(1+\frac{M-h+C_0}{\lambda}\right)+\omega_*(M).}
\end{gather*}

First let $\beta\to\infty$ at fixed $M$, using $K\to\infty$, and then let
$M\to\infty$. This proves the first limit. Summing the whole recursion gives

\begin{gather*}
\FitMath{\beta=x_K+c_0K+e_K+\sum_{j=1}^Kr_j,} \\[3pt]
\FitMath{|e_K|\le C_0.}
\end{gather*}

The terminal value belongs to a fixed bounded interval, so division by $K$
proves the second limit. Finally $\log(a_K^{-1})=K\log L$ exactly.
\end{proof}

If the analytic one-loop calculation establishes
$c_0=11N^2\log L/(12\pi^2)$ in the convention $\beta=g_0^{-2}$, this
arithmetic yields $\beta/\log(a_K^{-1})\to11N^2/(12\pi^2)$.
The one-loop coefficient is an input to this implication, not a consequence
of the stopping argument. The result is uniform over any family of runs
satisfying the same constants and error envelope.

\subsection{A quantitative optional error bound}

Suppose $h\ge e$, $p\ge0$, and
$\omega(x)\le Cx^{-1/2}(\log x)^p$. Then

\begin{gather*}
\FitMath{|\beta-c_0K|\le h+B+C_0+
 C'\sqrt{\beta+C_0+2}\,[\log(\beta+C_0+2)]^p.}
\end{gather*}

To prove it, partition $[h,\infty)$ into unit intervals. If two preceding
indices $i<j<K$ have values in the same interval, then $x_i-x_j<1$ and
the lower drift gives $j-i<(1+C_0)/\lambda$. A unit interval is therefore
visited at most $J=2+\lceil(1+C_0)/\lambda\rceil$ times. Also
$x_i\le\beta+C_0$ by the lower drift from zero. Summing the supremum of the
error envelope in each visited interval consequently gives

\begin{gather*}
\FitMath{\sum_{j=1}^K|r_j|
 \le CJ[\log(\beta+C_0+2)]^p
        \sum_{n=0}^{\lceil\beta+C_0-h\rceil}(h+n)^{-1/2}.}
\end{gather*}

The last sum is at most a constant times $\sqrt{\beta+C_0+2}$, by comparison
with the integral of $x^{-1/2}$. Substitute this bound in the summed recursion.
For fixed $h$ this is an $O(\sqrt\beta(\log\beta)^p)$ error. It supplies no
second logarithmic coefficient, convergence rate for observables, or
summable comparison between neighboring cutoffs.

\subsection{Depth envelopes and sums over scales}

The same block estimate gives useful bounds measured from the terminal
scale. For $0\le s\le K$, apply it between $K-s$ and $K$ to obtain

\begin{gather*}
\FitMath{h+\lambda s-C_0<x_{K-s}
             \le h+B+\Lambda s+C_0.}
\end{gather*}

This estimate records the possible bounded backward excursions without
requiring the inverse couplings to be monotone. It also distinguishes the
fixed reference value $h$ from the terminal value $x_K$, which is allowed to
range over its window.

For an explicit scale-summation corollary suppose
$h\ge\max\{1,2C_0\}$. Then the lower bound becomes
$x_{K-s}\ge h/2+\lambda s$. Write $g_j=x_j^{-1/2}$. For every $q>2$,
comparison with a decreasing integral gives

\begin{gather*}
\FitMath{\sum_{j=0}^K g_j^q
 \le\sum_{s=0}^\infty(h/2+\lambda s)^{-q/2}
 \le(h/2)^{-q/2}
   +\frac{(h/2)^{1-q/2}}{\lambda(q/2-1)}
 \le C(q,\lambda)h^{1-q/2}.}
\end{gather*}

The bound is uniform in the number of steps. The first term of the sum is
bounded separately, and the integral accounts for all later terms; this
explains both the strict condition $q>2$ and the displayed dependence on the
reference inverse coupling.

At the marginal exponent $q=2$ the integral instead grows logarithmically.
Using the two depth envelopes gives, for example,

\begin{gather*}
\FitMath{\frac1\Lambda
 \log\left(1+\frac{\Lambda(K+1)}{h+B+C_0}\right)
 \le\sum_{j=0}^K g_j^2
 \le\frac2h+\frac1\lambda
                     \log\left(1+\frac{2\lambda K}{h}\right).}
\end{gather*}

For the lower bound integrate the decreasing function
$(h+B+C_0+\Lambda s)^{-1}$ from zero to $K+1$. For the upper bound separate
its counterpart's first term and integrate from zero to $K$. These elementary
estimates are often the numerical part of a renormalization argument:
sufficiently high powers of the running coupling are summable, whereas the
marginal power is not. They cannot replace the analytic cancellation or
renormalization that produces those powers in a particular model.

The extra lower bound on $h$ belongs only to this corollary. For a smaller
fixed $h$, separate the finitely many terminal depths $s<2C_0/\lambda$, bound
them using $x_j>h$, and apply the same integral comparison to the remaining
depths. The resulting constants may depend on $h$ and $C_0$.

\subsection{An exact change of lattice units}

There is also an identity that requires no coupling-flow assumption. Let
$\Lambda_{K,\ell}=(L^{-K}\mathbb Z/\ell\mathbb Z)^4$ be a compatible finite
lattice and let $b$ be an integer with $K-b\ge0$. Dilation $D_bx=L^bx$ maps
it bijectively to $\Lambda_{K-b,L^b\ell}$. Transfer a bond variable by
$U'_{D_b e}=U_e$. The map bijects positively oriented bonds and plaquettes,
preserves their incidence, and satisfies $U'_{D_bp}=U_p$. Product normalized
Haar measure is therefore unchanged, as is the Wilson action
$A(U)=\sum_p[1-\tfrac12\operatorname{Re}\operatorname{tr}U_p]$.

For a test function on $T_\ell$, put $f^{[b]}(x')=f(L^{-b}x')$ on
$T_{L^b\ell}$. Plaquette barycenters obey $x(D_bp)=D_bx(p)$, so

\begin{gather*}
\FitMath{\mathcal O_{K-b,L^b\ell}(f^{[b]})(U')
 =\sum_p f(L^{-b}D_bx(p))a_p(U)
 =\mathcal O_{K,\ell}(f)(U).}
\end{gather*}

No spacing prefactor appears in this observable. Consequently the two Wilson
partition functions at the same $\beta$ are equal, and all joint moments and
cumulants satisfy the exact identity

\begin{gather*}
\FitMath{S^T_{n;K,\ell,\beta}(f_1,\ldots,f_n)
 =S^T_{n;K-b,L^b\ell,\beta}(f_1^{[b]},\ldots,f_n^{[b]}).}
\end{gather*}

This is one finite-dimensional integral written in two units, including the
same number of sites in each direction. It is not a comparison between
different bare measures.

For $\|f\|_\sharp=\max(\|f\|_\infty,c_M\operatorname{Lip}f)$ the exact
norm transformations are

\begin{gather*}
\FitMath{\|f^{[b]}\|_\infty=\|f\|_\infty,} \\[3pt]
\FitMath{\operatorname{Lip}f^{[b]}=L^{-b}\operatorname{Lip}f.}
\end{gather*}

Thus $\|f^{[b]}\|_\sharp\le\|f\|_\sharp$ for $b\ge0$, and it is at most
$L^{|b|}\|f\|_\sharp$ for $b<0$. All support distances and support radii
are multiplied by $L^b$. The same bond bijection carries a fixed-spacing DLR
specification to its dilated specification, since it carries each finite
integration set and all plaquettes meeting that set to their counterparts.
It therefore also transports infinite-lattice Gibbs states and their local
moment identities.

To identify this unit change with two initial portions of a renormalization
run requires more: the blocking rules, gauges, auxiliary choices and
extraction of increments must commute with the bond bijection. Likewise,
torus-independent increments require a separately defined common local
construction. Neither property follows from the scalar recursion or the
exact dilation identity proved here.

% END MODULE chapters/flow_arithmetic_v2.tex


% BEGIN MODULE chapters/06_flow.tex
\section{The running coupling and first-passage arithmetic}
\label{candidate:flow}
\candidatebanner{B, C, and E}

\subsection*{(3a) The run, its coefficient, and the analytic requirements}

The numerical running coupling and the auxiliary frame have different
roles. This distinction is needed even for a real bare coupling. A frame
fixes the tests, cubes, Gaussian references, coordinate choices and
large-field operations. The weight of the Wilson action is another entry
of the construction. A derivative with the frame fixed differentiates
that weight; it does not differentiate a characteristic function whose
threshold is moving. We first state the actual analytic requirements on
these objects. The arithmetic consequences below are conditional
theorems under those requirements.

Fix odd \(L>1\), the block parameter \(M\), and the remaining geometric
constants. In the ordinary-source application use its actual
\(M=L^{m_0}\), fixed collar layers and compatible cofinal tori;
admissibility below refers to those assigned lattices.
The group is \(\mathrm{SU}(2)\); in the coefficient formulas
we retain \(N=2\). The pointed families, hereditary complex
potential/current cores, midpoint tents, companions and Euclidean list
norms are those of Section~\ref{ch:rg-objects-v2}. In particular a pointed
family of level \(j\) has terms \(E(X,P,z)\), where \(X\) is a connected
localization domain and \(z\) is a unit lattice site in \(X\). Its norm
is a supremum in \(z\), not a sum in \(z\). The inequality
\[
 |E(X,P,z)|\le \mathcal A e^{-\kappa d_j(X)}
 \tag{RF1}
\]
is required on its declared core. Summing domains requires the separate
rooted-domain entropy estimate. A companion agrees on every common
nonwrapping lift, with the required collar, for all admissible tori
with the same frame prefix and initial weight. Local analyticity does
not establish that compatibility.

Write \(P_\circ=(\mathbf1,0)\), and let \(\mathfrak h_\infty\) be the
infinite-volume first response of the background potential/current pair
at \(P_\circ\). Its existence, gauge covariance, affine reproduction
modulo gauge and exponential decay belong to the background hypotheses.
For the composed physical Hessian, average over an orthonormal color
basis \(t_a\), \(1\le a\le3\), and set
\[
 \begin{aligned}
 \Pi(x,y)&=\frac13\sum_{a=1}^3\sum_z\sum_{X\ni z}
  D^2E(X,P_\circ,z)
       [\mathfrak h_\infty\delta_x^a,\mathfrak h_\infty\delta_y^a],\\
 b(E)&=\sum_{u\in\mathbb Z^4}\Pi_{\mu\nu}(u)u_\mu u_\nu,
                  \qquad \mu\ne\nu .
 \end{aligned}\tag{RF2}
\]
The second line is used only after translation covariance and independence
of the choice \(\mu\ne\nu\) have been justified. Hessians are in the pair
exponential chart, not in an unspecified coordinate system. Their
directions have norm
\(\max\{\|A\|_\infty,\|\nabla A\|_\infty,\|J\|_\infty\}\).
For \(a'=\min\{a_1,c,a_0/C\}\), a ball of radius \(ca'\) in this
norm must lie in the core. A ball controlled by \(\|A\|_\infty\) alone
does not supply the uniform mesh normalization.

The coefficient and pointed-cancellation theorem in
Section~\ref{ch:pointed-hessian-cancellation} applies under its full
response, summability and symmetry hypotheses. Its consequences used
here are absolute convergence, vanishing zeroth and first moments, and
\[
 \begin{aligned}
 \sum_u\Pi_{\mu\nu}(u)u_\rho u_\sigma
 &=b(E)(\delta_{\mu\rho}\delta_{\nu\sigma}
       +\delta_{\mu\sigma}\delta_{\nu\rho}
       -2\delta_{\mu\nu}\delta_{\rho\sigma}),\\
 |\Pi(x,y)|&\le C(a')^{-2}\mathcal A
                      e^{-c|y_--x_-|},
 &|b(E)|&\le C(a')^{-2}\mathcal A .
 \end{aligned}\tag{RF3}
\]
The extraction is linear, holomorphic in admitted parameters and real
on real families. It depends on the total, so two coherent localizations
of the same physical functional have the same coefficient. A constant
total has coefficient zero. At absolute radii the factor \((a')^{-2}\)
is a constant depending on \(M\); at running radii it can cost
\(C x_j(\log x_j)^{-2q}\). Neither the kernel nor its coefficient
depends on which admitted radii were used to estimate it.

The pointed Wilson family assigns the action to the same tents as the
regular member. Its norm is \(Ca_0^2\), so subtracting its own extracted
coefficient preserves a constant multiple of the family norm. This is
the subtraction used below. The stronger statement that the remaining
translation-invariant quadratic kernel is divisible by three lattice
derivatives, with exponentially decaying quotient kernels on a fixed
polyannulus, is a separate analytic target when that representation is
used. The moment identity alone is not presented as a proof of such
division.

\paragraph*{Frame, weight, tube and box.}
The functions used for the auxiliary choices are
\[
 \begin{aligned}
 \varepsilon(y)&=A_0y^{-1/2}(\log y)^{\mathfrak p_0},
 &\delta(y)&=A_1y^{-1/2}(\log y)^{\mathfrak p_0},\\
 \delta'(y)&=A_1y^{-1/2}(\log y)^{\mathfrak p_1},
 &\alpha_\iota(y)&=C_\iota y^{-1/2}(\log y)^q,\\
 R(y)&=\min\{L^b:L^b\ge(\log y)^r\},
 &q&>\mathfrak p_0+2r,\quad \mathfrak p_1<\mathfrak p_0 .
 \end{aligned}\tag{RF4}
\]
The fluctuation scaling is \(y^{-1/2}\). Gauge-fixing weights,
large-field stage scalings, stage Wilson constants and reference laws
also use the frame. Every occurrence of a coupling is assigned either
to such a choice or to the action weight. No occurrence is assigned
to both roles. Cube sizes and waiting depths use the running minimum
of the frame.

A finite reference sequence \(\theta_0,\ldots,\theta_n\) is real,
above its floor, and obeys
\[
 \lambda_\flat(j-i)-c_5\le\theta_i-\theta_j\le\bar B(j-i),
 \qquad
 \check R_j=R(\min_{i\le j}\theta_i),\quad
 \check R_j\le L\check R_{j+1}.
 \tag{RF5}
\]
In the run along this sequence,
\[
 x_0^\theta=\beta,\qquad
 x_{j+1}^\theta=x_j^\theta-b_{j+1}^\theta,\qquad
 z_j=\theta_j-x_j^\theta,\qquad \tau_j=x_j^\theta/\theta_j .
 \tag{RF6}
\]
The frame and weights are separate lists. Put
\(\mathsf t^\sharp(y)=\varpi y(\log y)^{-2q_{\rm t}}\).
A tube has positive, nonincreasing radii \(\mathsf t_j\), with
\(\mathsf t_j\le\mathsf t^\sharp(\theta_{j+1})\) for \(j<n\).
For a one-member sequence require
\(\mathsf t_0\le2\mathsf t^\sharp(\theta_0)\). The constants
\(\varpi,q_{\rm t}\) are fixed before the coupling ceiling, with
\(2\varpi\le\min\{\varpi_0,\sigma_\star\}\) and \(q_{\rm t}\ge q_0\).
The pair is alive through \(k\) if its coefficients through \(k\) are
defined and \(|z_j|<\mathsf t_j\). It has room \(r'\) when
\(|z_j|+r'\le\mathsf t_j\) at every required level.

The box of level \(k\) is the full inductive array: weights, regular
families with companions and coefficients, remainders, boundary terms
and standing large-field operations. It includes the bounds
\(|E^{(j)}|\le E_0e^{-\kappa d_j}\),
\(|R^{(j)}|\le\theta_j^{-\kappa_0/2}e^{-\kappa d_j}\),
the covariance and positivity clauses, identification with the two
chains below, and the drift bounds. Its reserve constant \(c_{\rm m}\),
also denoted \(c_{\rm r}\), enters the large-field mass bound; increasing
it shrinks the box and does not change the step. An admissible torus
must satisfy
\[
 \tfrac12\ell L^{K-i}\ge C(L)M\check R_i\qquad(i\le k).
 \tag{RF7}
\]
All box statements are simultaneous statements on these tori and on
their common companions. The positivity of a label measure
\(\tau_{\mathcal S}=\sum_{\mathcal S(h)=\mathcal S}\tau_h\)
must not be confused with the numerical ratio \(\tau_j\).

\paragraph*{The Gaussian and one-loop requirements; statement 3.3.}
Let \(\Phi_l=\log Z^{(l-1)}-\log Z^{(l-1)}(\mathbf1)\).
The localized Gaussian family is required to satisfy
\[
 \Phi_l(V)=\sum_X\sum_{z\in X}N^{(l)}(X,P_V,z),\qquad
 |N^{(l)}|\le E_{\mathsf N}e^{-\delta_{\mathsf N}Md_l(X)},
 \qquad b(N^{(l)})=b_l^0 .
 \tag{RF8}
\]
The constants depend only on \(L\). No coupling occurs in these
terms. Nonwrapping terms agree on all common lifts, and their pointed
totals are covariant under every unit lattice isometry.

For precision, the localization meant here is the following actual
one. Set \(k=l-1\), let \(A_Y=C^{(k)}(Y)^{-1}\) be the supplied
fluctuation form constructed inside \(Y\) and compressed to free bonds
at depth \(cM\), and define
\[
 \begin{aligned}
 f_b(Y)&=\operatorname{tr}(\log A_Y)(b,b),\\
 E_{\mathsf N}(X;b)
   &=\sum_{\substack{Y\subset X\\b\ {\rm at\ depth}\ cM\ {\rm in}\ Y}}
                      (-1)^{|X|-|Y|}f_b(Y).
 \end{aligned}\tag{RF9}
\]
The free-bond contribution to \(N^{(l)}\) is
\(-[E_{\mathsf N}(X;P,b)-E_{\mathsf N}(X;P_\circ,b)]/2\),
smeared by the midpoint tents. The pivot contribution is the logarithm
of its actual elimination factor minus its value at \(P_\circ\), smeared
by the same tents. The required spectral gap and complex continuation
of \(A_Y\), the locality of its walk expansion and the \(c/M\) radius
ceiling are part of this Gaussian-localization input. Finite Gaussian
integration estimates alone do not prove them.

The distinction between an on-shell pair and a critical pair matters.
On shell, the explicit current form expresses gauge covariance.
Identifying it with the Hessian on a constrained fibre also requires
the multiplier equation and actual criticality. The four tempered
localization properties used in the one-loop analysis are its sum
formula, common lifts, exponential domain decay on the shrinking
supremum polydisc of radius \(c\eta_l^2\), and invariance under constant
gauge transformations. Uniform Hessian Cauchy estimates instead use
the pair-direction norm preceding (RF3).

For obligation B3 we must prove, with the fixed momentum convention
\((2\pi)^{-4}dq\),
\[
 \begin{aligned}
 c_0&=\frac{11N^2}{12\pi^2}\log L,\qquad
 \left|\sum_{l=1}^k b_l^0-kc_0\right|\le c_2,\\
 |b_{l+1}^0-b_l^0|&\le C_\alpha L^{-\alpha l}
                         \quad(0<\alpha<1).
 \end{aligned}\tag{RF10}
\]
Thus \(b_l^0\to c_0\). The standard convention is
\(\beta=2N/g_{\rm std}^2\), giving the increment
\((11N/(24\pi^2))\log L\) of \(g_{\rm std}^{-2}\) per factor \(L\).
The numerical value in (RF10) requires the actual lattice trace
calculation. It is not deduced by quoting a continuum perturbative
coefficient.

We retain the precise auxiliary targets of that calculation. A tempered
functional is real analytic near the trivial field on arbitrarily
large tori, gauge and lattice-isometry invariant, normalized at that
field, and has a plane-wave Hessian per unit volume converging uniformly
in momentum to a bounded holomorphic symbol on a fixed strip.
For such functionals the background substitution is required to preserve
the coefficient. If \(\widehat\pi^F\) is the limiting bond Hessian symbol
with lower-endpoint Fourier convention, its midpoint form obeys
\[
 \begin{aligned}
 e^{i(p_\mu-p_\nu)/2}\widehat\pi^F_{\mu\nu}(p)
  &=b[F](|p|^2\delta_{\mu\nu}-p_\mu p_\nu)+O(|p|^4),\\
 \widehat\pi^{F\circ p_k}(p)
  &=L^2\widehat\pi^F(p/L)+O(|p|^3).
 \end{aligned}\tag{RF11a}
\]
The second formula is an actual background/average requirement;
its coefficient \(L^2=L^{-1}L^{-1}L^4\) records two form scalings
and the volume ratio. Uniform analytic control of the aliased
modes is part of that requirement. The first formula follows
from the moment identity and midpoint lattice symmetries.
The constrained Gaussian logarithms
\(\Psi_k\) then obey
\(\Psi_k=\Phi_k+\Psi_{k-1}\circ p_k\), where \(p_k\) is the preceding
background map, and
\(\sum_{l\le k}b_l^0=b[\Psi_k]\).
The finite Schur representation uses
\[
 \mathbb S(X,Q)=
 \begin{pmatrix}X&Q^{\mathsf T}\\Q&0\end{pmatrix}.
 \tag{RF11}
\]
Here \(Q\) is the covariant form average, \(R\) is the covariant site
average, \(D\) is the covariant gradient, \(\Delta=D^{\mathsf T}D\)
is the site Laplacian, \(\Delta^{\rm W}\) is the Wilson-action
Hessian on forms, and
\(\mathcal K=\Delta^{\rm W}+DD^{\mathsf T}\),
\(\Xi=\mathcal K-D(1-P)D^{\mathsf T}+Q^{\mathsf T}Q\),
where \(P\) is the orthogonal projection onto \(\ker R\).
On real backgrounds the representation to be justified is
\[
 \begin{aligned}
 \Psi_k={}&-\tfrac12\log|\det\mathbb S(
       \Xi-Q^{\mathsf T}Q+\mathfrak j[J,\mu],Q)|\\
 &+\log|\det\mathbb S(\Delta,R)|
          -(\text{the same expression at }\mathbf1).
 \end{aligned}\tag{RF12}
\]
Here \(U=U_k(V)\) is the actual critical background,
\(J=\nabla_{\rm fine}A(U)=J_{\eta_k}(U)\), and
\(\mu=\nabla_{\rm block}(A\circ U_k)(V)\).
The convention is the bilinear color pairing
\(\langle B,C\rangle=\frac12\operatorname{tr}(BC)\);
fine-bond inner products have weight \(\eta_k^4\), and
block-bond inner products have unit weight. Explicitly,
\[
 \begin{aligned}
 \left.\frac d{dt}A(e^{i\eta_k tB}U)\right|_0
     &=\eta_k^4\sum_b\langle J(b),B(b)\rangle,\\
 \left.\frac d{dt}(A\circ U_k)(e^{itC}V)\right|_0
     &=\sum_c\langle\mu(c),C(c)\rangle .
 \end{aligned}\tag{RF12a}
\]
All gradients and adjoints in (RF12) use those fixed conventions.
The form \(\mathfrak j[J,\mu]\) is the current/multiplier
correction obtained by second variation of the constrained
background problem. Its explicit form and the cancellation of
its quadratic momentum coefficient are established in
\eqref{ym4:actual-current-form}--\eqref{ym4:current-correction-coefficient}.
For complex continuation the absolute
determinants in this real identity must be replaced by coherent
logarithmic branches; an absolute value is not holomorphic.

The mass calculation uses finite determinants before any
thermodynamic limit. Work here on the unit bare lattice. For an
expanding box \(\Lambda_R=[-R,R]^4\cap\mathbb Z^4\), let
\(\iota_R^0\) and \(\iota_R^1\) be zero extension from its sites
and its positive bonds with both endpoints in the box. Put
\[
 \begin{aligned}
 U_f(x,1)&=\exp(-ifx_2t_3),&
 U_f(x,\nu)&=I\quad(\nu\ne1),&
 \tfrac12\operatorname{tr}(t_3^2)&=1,\\
 \Delta_{f,R}&=\iota_R^{0*}\Delta_{U_f}\iota_R^0,&
 \mathcal K_{f,R}&=\iota_R^{1*}\mathcal K_{U_f}\iota_R^1 .
 \end{aligned}\tag{RF13a}
\]
These are Dirichlet compressions of the scalar and vector
operators appearing in (RF12), before its averaged-constraint
corrections. The vector Wilson Hessian is defined by second
variation against finitely supported bond fields, so no
infinite total action is differentiated. The oriented
\(12\)-plaquette is \(\exp(if t_3)\). Adjoint charged colors
therefore carry fluxes \(2f\) and \(-2f\), while the third
color is neutral.

For each fixed \(z>0\) and finite box, take the real logarithms
near \(f=0\), where the massive matrices are positive, and set
\[
 \begin{aligned}
 \mathcal D_R(f,z)
   &=\log\det(\Delta_{f,R}+zI)
           -\tfrac12\log\det(\mathcal K_{f,R}+zI),\\
 \widetilde\beta(z)
   &=\lim_{R\to\infty}\frac1{|\Lambda_R|}
           \left.\frac{\partial^2}{\partial f^2}
                         \mathcal D_R(f,z)\right|_{f=0}.
 \end{aligned}\tag{RF13b}
\]
These identifications can be obtained directly from the fixed-mass
decay used in the calculation below. Centre a gauge at the point where
the differentiated trace density is read. The two field derivatives
then cost polynomial factors in relative distances, while the massive
kernels decay exponentially. Terms affected by the boundary are
therefore summable and their contribution per unit volume tends to
zero along the stated box exhaustions. This gives the bulk coefficient
and its independence of the exhaustion. The order is the derivative
at fixed positive mass, then volume, then \(z\downarrow0\).

For the finite-volume normalized Hessian kernels the Bloch-symbol
calculation below gives the translation-invariant form
\(\mathcal L(z)\). Decay makes the Bloch integrands holomorphic on
a common momentum strip; their torus Riemann sums converge to the
bulk integrals. The comparison concerns
\(\mathcal L^\sharp=\mathcal L(1)-\mathcal L'(1)+
\mathcal L''(1)/2\), not \(\mathcal L(1)\) alone, uniformly
on the prescribed external-momentum strip.
The required ultraviolet formulas for \(N=2\) are
\[
 \begin{aligned}
 \widetilde\beta(z)&=\frac{11}{6\pi^2}\log(1/z)+O(1),\\
 |z\widetilde\beta'(z)|+|z^2\widetilde\beta''(z)|&\le C,\\
 \sum_{l=1}^k b_l^0
 &=\widetilde\beta(\eta_k^2)
   -\eta_k^2\widetilde\beta'(\eta_k^2)
   +\tfrac12\eta_k^4\widetilde\beta''(\eta_k^2)+O(1).
 \end{aligned}\tag{RF13}
\]
All field expansions are at fixed positive mass before taking that
mass to zero. The vector form need not be nonnegative without its
mass; the small-field hypothesis must make the massive form coercive.
The comparison with the averaged constraints and gauge determinants
can be carried out by the stopped resolvent expansion below. Its two
local checks are the mixed energy bounds for the differentiated factors
and the grouping of the finite magnetic remainder before mass
differentiation. We give the estimates that make both checks uniform.

The geometric-convergence target can equivalently be organized through
the Taylor triple
\(\mathfrak J_k=(\Delta_k,S_3^{(k)},S_4^{(k)})\) of
\(\mathcal S_k=A\circ U_k\). With
\(\Lambda(B,\upsilon)=S(B)+\langle\upsilon,m(B)\rangle\),
\(\mathbb G=(D^2\Lambda(0))^{-1}\),
\(V_i=D^i\Lambda(0)\) and
\(\dot z=\mathbb G(0,C)=(HC,-\Delta_{k+1}C)\), the identity is
\[
 \begin{aligned}
 \Phi_{k+1}''(C,C')={}&\tfrac23\langle C,C'\rangle
                 -\tfrac23\langle HC,HC'\rangle\\
 &-\tfrac12\operatorname{Tr}(\mathbb G V_4[\dot z,\dot z'])
 +\tfrac12\operatorname{Tr}
       (\mathbb G V_3[\dot z]\mathbb G V_3[\dot z']).
\end{aligned}\tag{RF14}
\]
Here \(\eta_k=L^{-k}\) in the mass comparison,
\(S(B)=\mathcal S_k(e^{iB})+\mathrm{GF}_k(e^{iB})\),
\(m(B)=i^{-1}\log\mathcal M(e^{iB})\),
\(\Delta_k=\mathcal S_k''(\mathbf1)\), and
\(S_j^{(k)}\) is its degree-\(j\) Taylor form in the same chart.
The vector \(H C\) is the linear constrained background response.
Its uniform weighted-norm continuity and
\(\|\mathfrak J_{k+1}-\mathfrak J_k\|\le C_\alpha L^{-\alpha k}\)
are the additional actual-background comparison hypotheses used here. These
statements preserve the geometric comparison used at equal depth in
the stronger matched-cutoff route.

\paragraph*{The two-coupling Gaussian factor.}
At level \(k\), write the fibre form as
\(\Delta=\Delta_A+\Delta_{\rm GF}\). The frame fixes the field scaling
and gauge-fixing weight; the action weight supplies
\(\tau\Delta_A+\Delta_{\rm GF}\). The actual finite-dimensional
hypothesis is that \(\Delta_A\) vanishes on the block gauge directions
and \(\Delta_{\rm GF}\) acts only on those directions, with a positive
nondegenerate form on each complementary integration space. Under this
splitting the two ordinary Gaussian integrals factor, giving
\[
 \begin{aligned}
 \int dB\,\delta(QB)
 e^{-\langle B,(\tau\Delta_A+\Delta_{\rm GF})B\rangle/2}
     &=\tau^{-d_l/2}Z^{(k)}(V),\\
 \int d\mu_{\tau^{-1}C^{(k)}}(B)
 e^{-(1-\tau)\langle B,\widehat G B\rangle/2}
     &=\tau^{n_{\mathcal O}/2},\\
 n_{\mathcal O}&=3\#(\mathbb T^{(k)}\setminus\mathbb T^{(k+1)}),
 &d_l&=3|\mathfrak X_k|-n_{\mathcal O},\quad l=k+1 .
 \end{aligned}\tag{RF15}
\]
Here \(Q\) is the one-step average on the fibre, not the \(k\)-fold
average, \(\mathfrak X_k\) is the free-bond set, and
\(\widehat G\) is the gauge-fixing form in those coordinates.
For positive real \(\tau\), the formula follows by rescaling only the
non-gauge coordinates. Analytic continuation to the right half-plane
uses the branch of \(\log\tau\) real on the positive axis. Both
dimensions are independent of the background.

Consequently the total of the extra Gaussian family \(G_\tau\) is
\(n_{\mathcal O}\log\tau/2\). Its physical second derivatives vanish
and \(b(G_\tau)=0\). Its norm may nevertheless be
\(K_{\mathbf M}|1-\tau|e^{-(1+c)\kappa d}\), with no small-coupling
gain. This cancellation is needed to avoid a spurious accumulation of
slopes of order \(\theta_k^{-1}\). Formula (RF15) presupposes the actual
gauge splitting; it is not a consequence of a bound for a general
complex Gaussian.

The complete regular family has the form
\(E^{(k+1)}=N^{(k+1)}+G_{\tau_k}+F\).
Its physical step functional, up to field-independent normalizations,
is
\[
 \begin{aligned}
 \mathcal E^{(k+1)}[\theta_k;x;S](V)
 ={}&\log\!\left(
  \mathbf T_k^{\rm sf}[\theta_k]\,
       e^{-xA(U_k(\cdot))+S(U_k(\cdot))}\right)(V)\\
 &+xA(U_{k+1}(V))-S(U_{k+1}(V)).
 \end{aligned}\tag{RF16}
\]
The operator uses the whole torus as its small-field region, and \(S\)
is the action of the incoming list. Equality between this physical
functional and localized totals is asserted on their admitted
backgrounds. It is not an identity of a global Gaussian at arbitrary
independent-current pairs. The local formulas on such pairs have their
own complex cores.

\paragraph*{Lists, the two chains, and the increment as a function.}
A list contains regular families \(G^{(j)}\) with their own action
counterterms and unpointed remainders \(H^{(j)}\). At the output radii
\(\bar a_\iota\) its required load is
\[
 \begin{aligned}
 \mathcal N(\mathfrak f)
  &=\sum_{j=1}^k\omega_{k+1-j}\|G^{(j)}\|
    +\sum_{j\ {\rm aged}}[\bar a/a_j]^2\|H^{(j)}\|,\\
 \omega_n&=(n+1)^{3/2}L^{-\alpha_{\rm H}(n-1)},\\
 |F(X)|&\le K_{\mathbf M}
   \{\delta(\theta_k)+\mathsf s\,\mathcal N(\mathfrak f)
                 +e^{-c\theta_k\delta(\theta_k)^2}\}
                  e^{-(1+c)\kappa d(X)} .
 \end{aligned}\tag{RF17}
\]
The radius ratio is the maximum of the two coordinate-radius ratios;
\(\mathsf s=\delta(\theta_k)/\bar a_0\).
The first, second and third terms pay respectively for the interaction
beyond quadratic order, old-family differences, and failed Gaussian
tests. The parameter choices used to infer a power bound for the last
term must actually satisfy that power bound.

The regular part has a summable age weight because its own affine
quadratic action is removed before localization. A remainder has no
such counterterm, so its squared-radius load is not given an age decay.
Sections~\ref{ch:pointed-window-remainder},
\ref{ch:finite-source-insertion}, \ref{ch:renewed-source-pointing},
\ref{ch:window-potential-formation} and
\ref{ch:linear-source-bootstrap} prove the corresponding finite
implications under their actual background, bond-split, potential,
core and boundary hypotheses. Their general polynomial age factors
do not automatically give the exact exponent \(3/2\) in (RF17).
One must either verify that sharper estimate or allocate a smaller
positive exponential rate and prove the resulting comparison.

A Euclidean list has one companion, the correct coefficients of its
regular members, covariant pointed totals, and invariant remainder
totals on the declared covariance class. Its radii of covariance obey
\[
 e_{j'}\le L^{2(j'-j)}e_j\qquad(j\le j').
 \tag{RF18}
\]
Closure means that the actual output again has these properties.
Restriction to smaller cores does not prove covariance. The
backgrounds of fluctuating fields must first lie in a common
covariance class on an open real fluctuation ball; continuation then
requires a connected analytic domain on which both expressions exist.
The actual domain and naturality hypotheses remain part of closure.

The unmarked chain uses absolute radii and starts with the empty list.
The marked chain uses running radii
\(\check\alpha_{\iota,j}=\min\{\alpha_\iota(\theta_{j-1}),
\alpha_\iota(\theta_j)\}\) and inserts a remainder only when
\(M\check R_j\le L^{m+1-j}\) at step \(m+1\).
This age is distinct from the waiting depth of a large-field operation.
Before that age the remainder lies outside the regular exponent and
contributes to the remainder and boundary outputs.

Restriction from absolute to running radii must restrict the same
actual functions. Under that compatibility the marked difference
cancels \(N\) and \(G_\tau\). The required recurrence is
\[
 \begin{aligned}
 u_{m+1}&\le K_{\mathbf M}\mathsf s(\theta_m)
       \left(\sum_{i=1}^m\omega_{m+1-i}u_i+
                         C_{\mathfrak h}\theta_m^{-p}\right),\\
 p&=(\kappa_0-2)/2>1,\\
 \sum_{j\le m}\frac{\alpha_0(\theta_m)^2}
                       {\check\alpha_{0,j}^2}
              \theta_j^{-\kappa_0/2}
       &\le C_{\mathfrak h}\theta_m^{-p},\\
 C_{\mathfrak h}&\le C\left(1+
                         \frac1{\lambda_\flat(p-1)}\right).
 \end{aligned}\tag{RF19}
\]
The lower flow bound pays one inverse-coupling power in this sum.
With a constant frame its load can grow with the number of levels.
Smallness of \(K_{\mathbf M}\mathsf s\), relative to the summable
kernel and \(C_{\mathfrak h}\), closes a weighted supremum induction.
This is an actual marked-chain requirement, not a conclusion of an
unmarked finite source derivative.

The resulting increment requirement, with independent deviations
before imposing their recursion, is
\[
 \begin{aligned}
 b_{k+1}&=b_{k+1}^0+f_{k+1}(z_0,\ldots,z_k)+q_{k+1},\\
 \sup|f_{k+1}|&\le C\theta_k^{-1/2}
                            (\log\theta_k)^{\mathfrak p_0}
                           =:\epsilon_k,\\
 \sup|q_{k+1}|&\le C_v\theta_k^{-p_R}=:\epsilon'_k,
 &p_R&=(\kappa_0-4)/2 .
 \end{aligned}\tag{RF20}
\]
For \(f\), the supremum is on the whole product
\(|z_i|<\sigma_\star\theta_i\), with
\(0<\sigma_\star\le1/2\), and requires neither a lower flow bound
nor remainders. For \(q\), the frame obeys the weak flow bounds,
\(|z_i|<\varpi_0\theta_i(\log\theta_i)^{-2q_0}\) with
\(\varpi_0\le\sigma_\star/2\), and the admitted aged remainder data
range over their stated holomorphic parameter domain. Their bounds are
\(4\theta_j^{-\kappa_0/2}e^{-\kappa d_j}\), with (RF18).
The functions are real on real data. Constants are independent of
the frame length, torus, cutoff and admitted parameter; the weak
triple affects their threshold.

There is also a weaker sufficient marked-multiplier estimate. If the
marked exponent has norm \(\mathfrak n_k\le Cg_k^{\kappa_0-5}\)
and the actual step is holomorphic in its multiplier on a disc
of radius comparable to \(1/\mathfrak n_k\), Schwarz yields
\[
 |E^{(k+1),\mathfrak R}|
 \le C\mathsf s(x_k)\mathfrak n_k e^{-(1+c)\kappa d},
 \qquad |b_{k+1}^{\mathfrak R}|\le Cg_k^{\kappa_0-7}.
 \tag{RF21}
\]
This preserves an alternative sufficient error for the real flow;
it does not prove the stronger exponent in (RF20).

\paragraph*{Well-founded induction and statement 3.1.}
The substantive construction input says that, for an admitted box
and frame prefix through \(k\), \(E^{(k+1)}\) and its coefficient
are formed before any choices using \(\theta_{k+1}\). They use the
creation radii \((1+\beta_{\rm s})\alpha_\iota(\theta_k)\).
The terminal covariance radius is
\(e_{k+1}=c\varepsilon(\theta_k)\), so it also needs no next frame
entry. Once the computed next weight is admitted, the actual
\(\mathbf R\mathbf T\) operation closes the next box, consistently
on every admissible torus and its companion. The standing-history
kernel identity is
\((m(\kappa\otimes1))|_{\mathsf C}=(m|_{\mathsf A})\kappa\)
for a measure on \(\mathsf A\times\mathsf B\) and a kernel from
\(\mathsf A\) to \(\mathsf C\). Here \(\kappa\) denotes a kernel,
not a decay rate. This identity preserves the history sum; it does
not supply its quantitative mass or locality bounds.

Choose \(c_5=2c_2\), \(C_\beta\ge c_0+2c_2+1\), and first fix
\[
 \lambda_\flat=c_0/4,\quad
 \bar B=C_\beta+1,\qquad
 \bar\varsigma=C'\gamma^{\kappa_0-8},\quad
 \lambda_\sharp=c_0/2-\bar\varsigma,\quad
 B_\sharp=C_\beta+\bar\varsigma .
 \tag{RF22}
\]
The coupling ceiling is chosen so
\(\bar\varsigma\le\min\{c_0/4,1\}\).
The error exponent and other parameters must permit this ceiling;
\(\kappa_0>8\) is a sufficient choice for the displayed power of
\(\gamma\). An error estimate making that ceiling impossible
cannot be used in this parameter selection.
The marked-chain constants are computed at the weak triple before
this choice. Then choose \(Y_H=\gamma_H^{-2}\) above every required
floor, with
\(\epsilon(y)+\epsilon'(y)\le\min\{c_0/4,1\}\) and
\(\epsilon'(y)\le\bar\varsigma\) for \(y\ge Y_H\).
In particular \(\gamma_H\le
(\gamma^{-2}+c_5+1)^{-1/2}\). Increasing a reserve constant beyond
\([0,1]\) is allowed only after its stated floor is increased.

Put \(s_k=kc_0-\sum_{l\le k}b_l^0\); then \(|s_k|\le c_2\).
For any formed prefix in the weak class,
\[
 \begin{aligned}
 x_i-x_j-(j-i)c_0
 &=s_i-s_j+\sum_{i<l\le j}(b_l-b_l^0),\\
 \tfrac34c_0(j-i)-c_5
 &\le x_i-x_j\le B_\sharp(j-i).
 \end{aligned}\tag{RF23}
\]
The last upper bound uses the individual coefficient bound
\(|b^{\mathfrak A}|\le C_\beta\),
\(|b^{\mathfrak R}|\le\bar\varsigma\).
The lower bound is stronger than the strong triple, which in turn
implies the weak triple. Above a further floor, a bounded one-step
drop changes the smallest power \(R(y)\) by at most one factor \(L\).
Thus the computed next weight extends the frame while it remains
above the floor. The telescoping estimate supplies the numerical part of the induction.
Each next step is formed under the separate analytic formation and
closure hypotheses.

For \(h=g^{-2}\ge Y_H\), define
\[
 K(\beta,g)=\min\{k:x_k(\beta)\le h+B_\sharp\},\qquad
 \mathcal W_k(g)=\{\beta\ge h+B_\sharp:K(\beta,g)=k\}.
 \tag{RF24}
\]
The conditional flow statement 3.1 is that the run exists through
this index, has the same coefficients on all admissible tori, and
\[
 \begin{aligned}
 \lambda_\sharp(j-i)-c_5&\le x_i-x_j\le B_\sharp(j-i)
                         &&(i\le j\le K),\\
 h<x_K&\le h+B_\sharp,\\
 (\beta-h)/B_\sharp-1&\le K
 \le1+(\beta-h-B_\sharp+c_5)/\lambda_\sharp .
 \end{aligned}\tag{RF25}
\]
At \(\beta=h+B_\sharp\), \(K=0\); the terminal inequality still
holds. For \(K>0\), the previous value exceeds \(h+B_\sharp\)
and the last decrement is at most \(B_\sharp\), proving its strict
lower endpoint. The lower drift forces finite exit. The two bounds
for \(K\) follow by applying the upper drift to \(K\) and the lower
drift to \(K-1\). Arbitrarily large \(\beta\) give arbitrarily large
\(K\), hence infinitely many nonempty \(\mathcal W_k\).
No stepwise positivity of \(b_k\) was used.

At depth \(s\), write \(x_{[s]}=x_{K-s}\). The exact envelopes are
\[
 h+\max\{0,\lambda_\sharp s-c_5\}
 <x_{[s]}\le h+B_\sharp(s+1).
 \tag{RF26}
\]
Put \(\underline x_s=h+\lambda_\sharp s-c_5\),
\(\bar x_s=h+B_\sharp(s+1)\), and
\(g_-(g)=(g^{-2}+B_\sharp)^{-1/2}\).
A diagonal subsequence of any first-passage sequence with \(K\to\infty\)
therefore converges at every fixed depth, to values in the closed
versions of these intervals. Its couplings converge as well, with
terminal value in \([(h+B_\sharp)^{-1/2},g]\).
A whole-sequence limit along tuned runs requires the later
matched-cutoff input; compactness alone gives no such conclusion.

The arithmetic section, Section~\ref{ch:flow-arithmetic}, proves the
leading logarithm from vanishing mean errors, including bounded
backward excursions. In the present notation its uniform form is
\[
 \begin{aligned}
 |x_{[s]}-x_{[0]}-sc_0|&\le\Delta_g(s),&
 \Delta_g(s)/s&\longrightarrow0,\\
 \left|\frac{\beta}{K\log L}
       -\frac{11N^2}{12\pi^2}\right|
 &\le\frac{\Delta_g(K)+h+B_\sharp}{K\log L}
                       \longrightarrow0 .
 \end{aligned}\tag{RF27}
\]
Here \(K\ge1\), \(a_K=L^{-K}\). Denote the displayed relative
error bound by \(\delta_K\), distinct from the fluctuation
threshold \(\delta_j\). The function \(\Delta_g\) can be chosen
nondecreasing. Under the stated power/log error bound, with the
marked exponent sufficient for the same estimate, one may take
\(\Delta_g(s)\le C+C(h+c_0s/2)^{1/2}
(\log(h+c_0s/2))^{\mathfrak p_0}\).
The smaller ceiling \(\gamma_{\mathrm d}\le\gamma_H\)
imposes the monotonicity and envelope floors for this quantitative
bound; its constant is independent of \(g\) and \(s\).
For fixed \(g\) the relative error is
\(O_g(K^{-1/2}(\log K)^{\mathfrak p_0})\).
For example \(p_R>1\), used below, suffices for the marked sum.
The qualitative limit needs only an error tending to zero.

For the uniform depth statement, apply the arithmetic argument to the
segment from \(K-s\) to \(K\). Its shifted one-loop partial sums have
the same bound \(2c_2\). If \(\omega_*(y)\) is a nonincreasing envelope
for the error bound on \(x\ge y\), one may explicitly choose
\[
 \Delta_g(s)=2c_2+
 \sum_{u=0}^{s-1}
 \omega_*\!\left(h+\max\{0,\lambda_\sharp u-c_5\}\right).
 \tag{RF27a}
\]
This is independent of the initial coupling and cutoff, nondecreasing,
and sublinear by Cesaro averaging. Integral comparison for the
square-root/log envelope gives the quantitative bound above.

The sharper two-sided logarithmic comparison retained in statement
3.1 is
\[
 \begin{aligned}
 h-2c_2+\tfrac12c_0(K-k)
 &<x_k\le h+B_\sharp+2c_2+\tfrac32c_0(K-k),\\
 |x_k-x_K-c_0(K-k)|&\le\Delta_g(K-k).
 \end{aligned}\tag{RF28}
\]
These are the logarithmic inequalities since
\(\log(L^ka_K)^{-1}=(K-k)\log L\). The upper slope in the first
line uses the uniform error bound \( |b_l-b_l^0|\le c_0/2\)
and the partial-sum bound, not an assumption on the sign of each
one-loop increment.

\paragraph*{Scale comparisons and layered evaluation.}
These estimates retain the distinction between a numerical flow sum
and an estimate on an actual field. Write
\(t_p(x)=x^{-1/2}(\log x)^p\), \(p\ge0\). If
\(-c\le x_m-x_n\le B(n-m)\) for \(m<n\), then, above a floor
depending on a freely chosen \(\beta_0>0\),
\[
 \begin{aligned}
 t_p(x_n)&\le(1+\beta_0)(n-m)^{1/2}t_p(x_m),\\
 t_p(x_m)&\le(1+\beta_0)^2(n-m)^{\beta_0}t_p(x_n),\\
 \check R_j/\check R_{j+1}&\in\{1,L\},\\
 \sum_{j=0}^kg_j^p
 &\le C_p\left(1+\frac1{\lambda_\sharp(p-2)}\right)g_k^{p-2}
                       \qquad(p>2).
 \end{aligned}\tag{RF29}
\]
The last line also assumes the lower drift bound. The first two
follow by comparing the square roots and logarithms above their
monotonicity floor; a bounded backward excursion is absorbed by the
chosen \(1+\beta_0\). The third uses the running minimum and the
bounded one-step drop. For the last, use
\(x_j\ge x_k-c_5+\lambda_\sharp(k-j)\), absorb \(c_5\) by the
floor, and compare the decreasing power with its integral.
The exponent \(p=2\) is excluded: that integral is logarithmic.

For the field estimate fix an admitted label through \(k\), with
layers \(\Gamma_n=\Omega_n\setminus\Omega_{n+1}\),
\(\Omega_0=T_\ell\) and \(\Omega_{k+1}=\emptyset\).
The determining sequence starts at \(\Omega_1^\sim\).
Let \(\Lambda_j^+\) be \(\Lambda_j\) together with its level-\(j\)
islands. Base points are the unit sites lying one layer of
\(M\check R_j\)-cells inside this set. Let \(\mathsf Y\) be any
set of these points and \(\mathsf Y_n=\mathsf Y\cap\Gamma_n\).
The notation \(|D|_n\) counts sites of \(\mathbb T^{(n)}\) in \(D\);
it never counts cells.

Assume the actual window/core and boundary hypotheses for these
backgrounds, Euclidean regular members of norm at most \(E_0\),
remainders of norm at most \(g_j^{\kappa_0}\), and the stated
boundary-domain decay. In particular the actual boundary terms,
indexed by their recorded origins \(o\), must satisfy
\[
 |B_{\mathcal S}(j,X,o)|
       \le C_Be^{-\kappa d_j^\natural(\mathcal S;X)} .
 \tag{RF30a}
\]
The metric is the one in (R23) of
Section~\ref{ch:rg-objects-v2}: the minimizing tree represents
each chosen standing-island hull by one member, rather than
filling it with all of its cubes. It is not ordinary
\(d_j\)-decay across an island.
The retained layered requirements are
\[
 \begin{aligned}
 &\left|\sum_{z\in\mathsf Y}
   \left\{\sum_{\substack{X\ni z\\X\subset\Lambda_j^+}}
       [E^{(j)}(X,U_k,z)-E^{(j)}(X,\mathbf1,z)]
                      -b_jA(h_z,U_k)\right\}\right|\\
 &\hspace{12mm}\le E_1\sum_{n=j}^k
      (1+n-j)L^{-\beta_T(n-j)}L^{-4(n-j)}\#\mathsf Y_n,\\
 &\sum_{\substack{X\in\mathcal R_j\\X\cap\Omega\ne\emptyset}}
       |R^{(j)}(X,U_k)-R^{(j)}(X,\mathbf1)|\\
 &\hspace{12mm}\le R_1g_j^{\kappa_0}
             \sum_{n=j}^k(1+n-j)|\Gamma_n\cap\Omega^\sim|_n .
 \end{aligned}\tag{RF30}
\]
Here \(0<\beta_T\le\alpha_{\rm H}<1\) is fixed,
\(\mathcal R_j\) consists of domains outside \(Z_j^+\), and
\(Z_j^+\) is the stored large-field region enlarged by one cell
layer. The set \(\Omega\) is a union of level-\(k\) localization
cubes and \(\Omega^\sim\) has one level-\(j\) localization layer
added. The constant \(R_1\) in this display is not the geometrical
gap parameter. Constants may depend on \(\beta_T,L,N\) and the
construction parameters, but not on \(j,k,K,\ell\), the label or
the positions.

For clarity, the local estimate behind the first line, when its
actual background hypotheses have been verified, is
\[
 \begin{aligned}
 &\left|\sum_{X\in\mathcal X}
    [G(X,U,z)-G(X,\mathbf1,z)]-b(G)A(h_z,U)\right|\\
 &\qquad\le C\mathsf g
  \left\{(\varepsilon_z/a_0)^2
                \sigma^{4+\beta_T}
        +e^{-\kappa\min\{\mathsf D,c/\sigma\}/2}\right\},
 \qquad \sigma=L^{-(n-j)} .
 \end{aligned}\tag{RF31}
\]
The member has norm \(\mathsf g\) on radii \(a_0\le a_1\), with
\(a_0\le1\); the datum is \(\varepsilon_z\)-regular near \(z\),
\(\varepsilon_z\le c\), and
\(\varepsilon_z\sigma\le ca_0\).
Every term is evaluated in its actual core. The collection
\(\mathcal X\) contains all domains within \(\mathsf D\) layers
of \(z\). With complete near coverage to distance \(c/\sigma\),
the right side is bounded by
\(C\mathsf g(1+(\varepsilon_z/a_0)^2)\sigma^{4+\beta_T}\).
The local cancellation and nonlinear remainder sections supply
the analytic implication under their full hypotheses, including
the boundary-response compatibility and the own tent counterterm.
Checking those hypotheses on this layered background remains
part of the actual construction.

The tail in (RF31) has no factor
\((\varepsilon_z/a_0)^2\). For example a wrapping holonomy term,
multiplied by its allowed exponentially small domain weight, can
have zero nonwrapping companion and coefficient zero while its
normalized value is nonzero on a flat torus field. A curvature
amplitude cannot control that topological contribution.

To pass to (RF30) one also needs the actual geometry:
\(\varepsilon_z\le C\check R_n^2\varepsilon_n\le
c\alpha_{0,n}\), the comparison
\(\alpha_{0,n}\le C(2+n-j)^{1/2}\check\alpha_{0,j}\), and the
available collar depth. The condition \(q>\mathfrak p_0+2r\)
pays the first comparison. The factor \(L^{-4(n-j)}\) is used
entirely by converting level-\(j\) sites to level-\(n\) sites.
For remainders the corresponding near-domain Cauchy bound is
\[
 |R^{(j)}(X,U_k)-R^{(j)}(X,\mathbf1)|
 \le C(d_j(X)+3)^2L^{-4(n-j)}
  \left(\frac{\check\alpha_{0,n-1}}{\alpha_{0,j}}\right)^2
                    g_j^{\kappa_0}e^{-\kappa d_j(X)} .
 \tag{RF32}
\]
Far domains use their norm, not this amplitude estimate.
The second line of (RF30) also holds with \(1+n-j\) replaced
by \(\varpi_{jn}=\max\{1,(x_j/x_n)^2\}\).
Both dominate the required squared-radius ratio up to fixed
constants. The two resulting summations give
\(|\mathbf R_k(U_k)|\le C\sum_n
g_n^{\kappa_0-6}|\Gamma_n|_n\).

The regular counterterm in (RF30) is the sum of its own tents.
An arbitrary smooth weight equal to one on \(\Omega\) differs
from that sum on its collar; it produces a new action term
without an age gain. For the actual profile \(\varphi_j\),
\(\varphi_jh_z=h_z\) at every base point. The remaining rim
terms obey \(0\le A(\varphi_jh_y,U_k)\le C\alpha_{0,j}^2\).
The needed geometrical count is
\(\sum_j\#_j(Z_j^+)\le C\sum_n|\Gamma_n|_n\).
It includes the stored collars of healed components, not just
the re-done small-field regions. This window and history count is a separate geometric hypothesis.

Here \(\mathbf R_k\) denotes the sum of the normalized remainder
differences on the label. The net regular action
\(\mathbf E_k\) includes its profile counterterms and rim terms.
With those counts and boundary sums,
\(|\mathbf E_k(U_k)|\le C\sum_n|\Gamma_n|_n\).
For admitted complex perturbations of total coefficient size
\(\bar c\), regular norms at most \(E_0+1\), remainder norms
at most \(4g_j^{\kappa_0}\), and the same boundary hull bounds,
the sum of the elementary moduli is at most
\(C(1+\bar c)\sum_n|\Gamma_n|_n\).
This is a pointwise estimate on the admitted arrays; it is
not an assertion that arbitrary such arrays are produced by
the RG map. Its constants and evaluation floor need no new
dependence on \(\bar c\). The variable action coefficient is
\[
 x_k(y)=\beta-\sum_{j\le k}b_j\varphi_j(y),\qquad
 |x_k(y)-x_k|\le B_\sharp
                  \#\{j\le k:\varphi_j(y)<1\}.
 \tag{RF33}
\]
It equals \(x_k\) one coarse cell inside all the small-field
regions. Let \(A_{k,\mathcal S}\) be the ACTUAL label exponent
and \(\mathcal E_k\) its once-counted field-independent number.
The sign convention is \(\mathcal E_k=-E_{(k)}\), where
\(E_{(k)}\) is the accumulated log vacuum normalizer of (CP1).
Thus the constant part of the exponent is \(E_{(k)}\).
The exact density format, which is a construction input,
identifies the sum of the elementary regular, remainder, rim
and boundary terms with the expression on the left of
\[
 |A_{k,\mathcal S}(U_k)+\mathcal E_k+
                       A(x_k(\cdot),U_k)|
 \le C(1+\bar c)\sum_n|\Gamma_n|_n .
 \tag{RF33a}
\]
Their modulus sum proves this consequence on the admitted
backgrounds. Defining an exponent by a bound would not prove
the required exact density format.

\subsection{The constrained determinant and the logarithmic coefficient}
\label{ym4:new-one-loop-reduction}

The one-loop sum is computed before choosing a localisation of its
logarithm. This separates the coefficient from the individual polymer
terms and keeps the averaging and gauge determinants in the
calculation.

Let $\Psi_k$ be the normalized logarithm of the Gaussian integral
over the linearized fibre of the $k$-fold average at the actual
critical background. With every delta-function Jacobian retained,
Gaussian integration in stages is the Schur-complement identity
\[
 \Psi_k=\Phi_k+\Psi_{k-1}\circ p_k.
\]
Here $p_k$ is the preceding critical block field. The linearized
average annihilates the aliased modes at zero momentum, while its
unaliased branch reproduces a constant form. Its two responses and
the volume ratio give
\[
 \widehat\pi^{F\circ p_k}(p)
       =L^2\widehat\pi^F(p/L)+O(|p|^3),\qquad
 L^2=L^{-1}L^{-1}L^4.
\]
Taking the quadratic coefficient therefore gives
$b[F\circ p_k]=b[F]$ and
$\sum_{l\le k}b_l^0=b[\Psi_k]$. The uniform strip estimates below
justify the Hessians, the infinite-volume limit and the alias
expansion in this calculation.

The constrained Gaussian integral of a quadratic form $T$ with
linear constraint $M$ is, up to a numerical power of $2\pi$,
\[
 \left|\det\begin{pmatrix}T&M^{\mathsf T}\\M&0\end{pmatrix}
                                                    \right|^{-1/2}.
\]
The formula is valid when $M$ is onto and $T$ is positive on its
kernel. It follows by splitting off $\ker M$ and changing variables
on its orthogonal complement; the determinant of that change is
exactly the delta-function Jacobian. The same splitting, with the
residual gauge directions separated first, proves the linear
Faddeev--Popov identity. In this identity the constrained integral,
the determinant along residual gauge directions and the determinant
compensating a coarse gauge displacement of the constraint must
be changed together. The axial gauge is triangular along the
levels and its gauge determinant is one. Changing to the covariant
background gauge and straightening the constraint consequently
gives the saddle determinant in (RF12).

There are two current terms to retain. Gauge invariance at a field
with nonzero current relates the Wilson Hessian on a gauge
direction to a commutator paired with that current. Curvature of
the averaging fibre contributes the second differential of the
average paired with its multiplier. Equivalently, one first takes
the Hessian of the actual constrained Lagrangian, rather than
restricting the Wilson Hessian to a linear space, and then changes
gauge. This is the origin of $\mathfrak j[J,\mu]$ in (RF12).
Its value must be the one produced by these operations.

At the trivial field the residual gauge space is $\ker R$.
The projection in its gauge term is the orthogonal projection $P$
onto this space; it is not the Landau projection onto
$\Delta\ker R$ used in the propagator construction. The flat
operator
\[
 \Xi_0=\Delta^{\rm W}+DPD^{\mathsf T}+Q^{\mathsf T}Q
\]
has a uniform comparison with
$\Delta^{\rm F}+1$, where
$\Delta^{\rm F}=\Delta^{\rm W}+DD^{\mathsf T}$:
\[
 c(\Delta^{\rm F}+1)\le\Xi_0\le C(\Delta^{\rm F}+1).
\]
The curl controls the transverse modes. On gradients the identity
$QD=D_{\mathrm{block}}R$ and the projected gauge term control,
respectively, the block mean and the zero-mean part. The only
remaining constant mode is detected by $Q$. This is also the
comparison needed for the site-field saddle determinant.

In level-$k$ units take Bloch symbols, whose fine-cell dimension is
$O(\eta_k^{-4})$. A complex momentum shift is conjugation by an
exponential weight. The averages are bounded under this conjugation,
and the sandwiched operator
$(\Delta^{\rm F}+1)^{-1/2}\Xi_0
 (\Delta^{\rm F}+1)^{-1/2}$ is bounded on a fixed strip. Its
variation under an imaginary shift is $O(|\mathrm{Im}\,p|)$.
The real coercivity therefore gives invertibility on a smaller
strip of width independent of $k$. The symbols of the linearized
background have summable Fourier modes. Covariant averages at a
small complex potential obey
\[
 \|Q_k(U)\|\le\prod_{i<k}(1+CL^{i-k}\|A\|_\infty)\le e^C.
\]
Thus differentiating their finite jet does not introduce a factor
growing with the number of levels. Fourier summation on a torus
then differs from the Bloch integral by $O(e^{-c\mathfrak s})$,
where $\mathfrak s$ is its side in these units.

\paragraph{The current and multiplier correction.}
We compute its coefficient at each fixed depth $k$. Write
$(\cdot,\cdot)_{\rm f}$ and $(\cdot,\cdot)_{\rm c}$ for the
fine and coarse bond pairings in (RF12a), and retain $R_U$ for
the site average. The Landau projection and the projection of
forms to its slice are
\[
 \begin{aligned}
 G'_U&=(\Delta_U+R_U^{\mathsf T}R_U)^{-1},\\
 \mathcal R_U&=1-G'_UR_U^{\mathsf T}
       (R_UG_U'^2R_U^{\mathsf T})^{-1}R_UG'_U,\\
 \lambda_B&=G'_U\mathcal R_UD_U^{\mathsf T}B,
       &\Pi_UB&=B-D_U\lambda_B .
 \end{aligned}
\]
Here $\mathcal R_U$ projects onto $\Delta_U\ker R_U$;
it is different from $P_U$, which projects onto $\ker R_U$
in (RF12). In particular, $R_U\lambda_B=0$ and
$Q_U\Pi_U=Q_U$. Let $C_U^{(2)}$ be the quadratic homogeneous
part of the actual $k$-fold straightened average, normalized
so that its linear part is $Q_U$. All its accumulated node
factors are removed before taking this Taylor expansion.

For a site field $\lambda$, let $\lambda_+(b)$ denote its
value at the positive endpoint transported to the negative
endpoint by $U(b)$. With $\lambda=\lambda_B$, put
\[
 E_U(B)=i[\lambda,B]-i[B,\lambda_+]
                         -i[\lambda,D_U\lambda],
\]
where every bracket is taken at the negative endpoint.
The second variation of gauge invariance gives the exact identity
\[
 (\Pi_UB,\Delta_U^{\rm W}\Pi_UB)_{\rm f}
 -(B,\Delta_U^{\rm W}B)_{\rm f}
                  =-(E_U(B),J(U))_{\rm f}.
\]
The covariant difference contains the fine spacing: the
transported endpoint difference is
$\lambda_+-\lambda=\eta_kD_U\lambda$.
Thus the same identity follows from the mixed gauge variation
\[
 (B,\Delta_U^{\rm W}D_U\lambda)_{\rm f}
       =\tfrac i2(J(U),[\lambda+\lambda_+,B])_{\rm f};
\]
applying it to $B$ and $D_U\lambda$ gives the displayed
difference of Hessians because
$[\lambda_+-\lambda,D_U\lambda]=0$.

The curved constraint contributes
$-2(\mathsf H_U C_U^{(2)}(\Pi_UB),J(U))_{\rm f}$.
On the critical branch,
$J(U)=Q_U^{\mathsf T}\mu$ and $Q_U\mathsf H_U=1$,
so this term is $-2(C_U^{(2)}(\Pi_UB),\mu)_{\rm c}$.
The first identity also holds for the straightened average:
its linear part differs from the plain one by a coarse
gradient, which pairs to zero with $\mu$ since
$D_{\mathcal M^kU}^{\mathsf T}\mu=0$.
Consequently the actual correction, extended to independent
$J$ and $\mu$, is the symmetric form defined by polarization of
\[
 (B,\mathfrak j_U[J,\mu]B)_{\rm f}
   =-(E_U(B),J)_{\rm f}
        -2(C_U^{(2)}(\Pi_UB),\mu)_{\rm c}.
 \tag{RF12b}\label{ym4:actual-current-form}
\]
This is the invariant extension of the Wilson Hessian followed
by the second variation of the averaging constraint. The exact
change of gauge then allows the soft gauge term $D_UP_UD_U^{\mathsf T}$
in (RF12); it does not replace $\mathcal R_U$ by $P_U$ in
$\Pi_U$.

For the coefficient, it is essential to keep the two currents
independent until after differentiating. Set
$K_U=\Delta_U^{\rm W}+D_UP_UD_U^{\mathsf T}$ and
$z=(J,\mu)$. On each finite torus define, near $(\mathbf1,0)$,
\[
 F_k(U,z)=-\tfrac12\log
  \frac{\det\mathbb S(K_U+\mathfrak j_U[z],Q_U)}
       {\det\mathbb S(K_U,Q_U)} .
\]
The logarithm is its analytic branch with $F_k(U,0)=0$.
The site determinant cancels in this difference. At fixed $k$,
invertibility at the trivial field gives analyticity on each
finite torus. The kernels of local variations have thermodynamic
limits and decay exponentially. These kernels define the
normalized plane-wave jets; no total infinite-volume logarithmic
determinant is needed.

Every term in \eqref{ym4:actual-current-form} is covariant under
a background gauge transformation, also for independent $J,\mu$.
Indeed, the parallel transports, the removed node factors,
$G'_U$, $\mathcal R_U$ and $\Pi_U$ conjugate under that
transformation; taking the quadratic part preserves the same
rule for $C_U^{(2)}$. Both saddle matrices therefore conjugate
on their fine and coarse spaces. Their determinants, and hence
$F_k$, are invariant. This uses background gauge covariance,
which also holds for $K_U$ without the current correction.

Define $w_J(U)$ and $w_\mu(U)$ by
\[
 D_zF_k(U,0)[j,\nu]
      =(w_J(U),j)_{\rm f}+(w_\mu(U),\nu)_{\rm c}.
\]
These derivatives are the actual determinant traces:
\[
 D_zF_k(U,0)[\zeta]
 =-\tfrac12\operatorname{Tr}\left[
    \mathbb S(K_U,Q_U)^{-1}
    \begin{pmatrix}\mathfrak j_U[\zeta]&0\\0&0\end{pmatrix}
                                      \right].
\]
They transform as adjoint fields at the corresponding endpoints.
At $U=\mathbf1$, constant gauge transformations fix the
background. There is no nonzero adjoint-invariant vector in
$\mathfrak{su}(2)$, so $w_J(\mathbf1)=w_\mu(\mathbf1)=0$.
Differentiating their covariance under a local gauge transformation
now gives, for either field $w$,
\[
              D_Uw(\mathbf1)[d\lambda]=0.
\]
This is the gauge identity on the other leg of the mixed
current--background derivative. The gauge variation of the
current variable contributes only a commutator with $w(\mathbf1)$,
which is zero.

At fixed $k$, exponential decay allows a linear site function
$\lambda$ in this identity by cutoff approximation. Its gradient
is an arbitrary constant fine form. The mixed kernel thus
annihilates constant forms as well as gradients. The zero-momentum
linear background of a constant coarse form is constant modulo
a gradient. Its mixed symbol consequently vanishes at $p=0$;
analyticity makes the composed symbol $O_k(|p|)$.

The first variations of the physical current and multiplier are
$O_k(|p|^2)$. This holds in every fine cell, including the aliased
modes. In the linearized problem at $\mathbf1$, the symmetric
averages reproduce affine potentials, whose minimizers have
zero current and multiplier. Their
linearized current kernels therefore annihilate affine coarse
potentials, so both their constant and linear moments vanish.
A change of the fine gauge does not alter this conclusion, since
$J(\mathbf1)=0$ and $DJ(\mathbf1)d\lambda=0$.

Put $\Psi_k^0$ equal to (RF12) with $\mathfrak j$ omitted and
substitute the physical currents only now. In the second-variation
chain rule for $\Psi_k-\Psi_k^0$, every second variation of a
current is paired with $w(\mathbf1)=0$. Terms with two first
current variations are $O_k(|p|^4)$; mixed terms have orders
$O_k(|p|)O_k(|p|^2)=O_k(|p|^3)$. There is no term with two
background variations because $F_k(U,0)=0$ for every $U$. Hence
\[
 \widehat\pi^{\Psi_k-\Psi_k^0}(p)=O_k(|p|^3),\qquad
                     b[\Psi_k]=b[\Psi_k^0].
 \tag{RF12c}\label{ym4:current-correction-coefficient}
\]
The thermodynamic limit and the coefficient are taken at fixed
$k$. The equality of coefficients is exact for every $k$;
it requires no bound uniform in $k$ on the discarded cubic
remainder. The uniform massive comparison is applied to
$\Psi_k^0$ after this cancellation.


To compare the current-free saddle determinant with the massive
determinant pair, interpolate in the averaging perturbation and in
the mass. The perturbation is
\[
 Q^{\mathsf T}Q-D(1-P)D^{\mathsf T},
\]
with the corresponding site-field perturbation. Put
$\mathcal W_{\rm F}=(\Delta^{\rm F}+1)^{1/2}$. The flat comparison above controls
the interpolated inverse between its two energy weights. Its
continuation to the common smaller momentum strip follows from
the same bounded-conjugation argument.

After shrinking the common strip, the linearized background has
uniformly summable Fourier modes there. Its plane-wave variations
therefore enter the fixed potential ball with a bound independent
of the bare spacing. The average factors must be differentiated
before estimating the trace. Their bound on the fixed complex potential ball above gives,
by Cauchy's formula on a smaller ball of radius $h>0$,
\[
 \|D_A^jQ_k\|\le C_jh^{-j},\qquad j=0,1,2,
\]
uniformly in $k$; the site average is treated in the same way.
This differentiates one bounded holomorphic map, so it avoids a
loss for each of its $k$ factors. The factors from the nonlinear
averages carry $L^{i-k}$, whose sum is bounded. For the inverse use
$D(A^{-1})=-A^{-1}(DA)A^{-1}$ and its second derivative.
In these products every differentiated precision is paired with
the adjacent $\mathcal W_{\rm F}^{-1}$ factors. We must check the resulting finite
list of mixed energy bounds on this same strip, with the actual
average and constraint jets. This is the norm check needed in
addition to the coarse-rank count.

A trace containing an average can then be taken over the coarse
cell, whose dimension is bounded independently of $k$. The
remaining resolvent expansion is stopped when it has three free
propagators. Its four-dimensional Schatten majorant is
\[
 \int\|\widehat G(q)\|_{\mathfrak S^3}^3dq
       \le C\int_{\mathbb R^4}(1+c|q|^2)^{-3}dq<\infty,
 \qquad G=(\Delta^{\rm F}+1)^{-1}.
\]
The lower powers are not uniformly trace class as the fine spacing
tends to zero. Retain them exactly as the first two mass
derivatives of the determinant pair. Thus the explicit part is
\[
 \mathcal L^\sharp=\mathcal L(1)-\mathcal L'(1)
                         +\tfrac12\mathcal L''(1).
\]
The coarse traces and the stopped remainder give a bounded
holomorphic difference on the common external-momentum strip,
once the indicated mixed bounds are inserted. Cauchy's estimate
then bounds its quadratic coefficient uniformly in $k$.

The divergent coefficient is computed at fixed positive mass on a
constant Abelian field. Write
$\psi(q)=\sum_\mu(2-2\cos q_\mu)$ and
$\Upsilon=(\psi+z)^{-1}$. For sufficiently small plaquette flux
relative to $z$ the vector operator plus $z$ is positive: its
possibly negative part is bounded by $C|\sin f|$. Differentiation
at $f=0$ is therefore legitimate at each finite volume. Massive
decay controls the linearly growing vector potential and removes
boundary contributions after division by the volume.

The neutral color contributes no flux dependence. The charged
colors carry fluxes $\pm\varphi$, with $\varphi=2f$, and their
covariant shifts satisfy $u_1u_2=e^{i\varphi}u_2u_1$.
The second-order magnetic product is
\[
 a\star b=ab-\tfrac i2\varphi\{a,b\}
                  -\tfrac18\varphi^2\{\!\{a,b\}\!\}
                  +O(\varphi^3).
\]
At zero flux the vector symbol is $\psi I$. Solving for its
resolvent through flux degree two uses only $\Upsilon$, momentum
derivatives of $\psi$, and the first two flux coefficients of the
finite-range Wilson symbol. These coefficients are smooth
trigonometric matrices independent of $z$. The cross terms from
the first magnetic bracket cancel because $\Upsilon$ is a scalar
function of $\psi$. Integrating the mass derivative of the
logarithm and then integrating by parts in the periodic momentum
groups the second bracket into terms with a smooth numerator times
$\Upsilon^2$; the other terms have a bounded numerator times
$\Upsilon$, or are the vector spin term with $\Upsilon^2$.

Isolate the numerator at $q=0$ in every $\Upsilon^2$ term. The
remaining numerator is $O(|q|)$, and hence its absolute integral
is bounded uniformly for $0<z\le1$, since
\[
 \sup_{0<z\le1}\int_{[-\pi,\pi]^4}
       \bigl\{|q|(\psi+z)^{-2}+(\psi+z)^{-1}\bigr\}\,dq<\infty.
\]
Here $\psi(q)$ is comparable to $|q|^2$ near its only zero.
The constant numerators give
\[
 \widetilde\beta(z)
  =8(4-\tfrac13)\int_{[-\pi,\pi]^4}
             \frac{dq}{(2\pi)^4}(\psi(q)+z)^{-2}+r(z)
  =\frac{11}{6\pi^2}\log(1/z)+O(1).
\]
The $4$ is the spin contribution from the vector Hessian, read
from its intertwining with the covariant gradient on the critical
constant field. The $1/3$ is the orbital contribution of the
second magnetic bracket at $q=0$. To finish this finite
calculation, we must match the actual vector numerators to the
grouping just described, retaining its absolute integral bounds.

The same integral bounds control the mass derivatives of $r$. Indeed,
for a mass-independent numerator $F(q)$,
\[
 \begin{split}
 z\partial_z\{F(\psi+z)^{-m}\}
   &=-m\frac z{\psi+z}\,F(\psi+z)^{-m},\\
 z^2\partial_z^2\{F(\psi+z)^{-m}\}
   &=m(m+1)\left(\frac z{\psi+z}\right)^2F(\psi+z)^{-m}.
 \end{split}
\]
Both extra ratios are bounded. If a rearrangement introduces
explicit polynomial powers of $z$, the finite Leibniz rule gives
the same absolute integral bounds. Differentiating the grouped
integrals therefore gives
$|zr'(z)|+|z^2r''(z)|\le C$ as well as $|r(z)|\le C$.
For the leading integral each mass derivative adds a resolvent,
and $\int\Upsilon^j dq\le Cz^{2-j}$ for $j=3,4$. Consequently
$|z\widetilde\beta'(z)|+|z^2\widetilde\beta''(z)|\le C$.
For $z=\eta_k^2$, these bounds, the exact cancellation
\eqref{ym4:current-correction-coefficient} and the uniform
current-free remainder give (RF13), and hence the bounded
telescoping error in (RF10).

\subsection*{(3b) A common companion and the torus floor}

\paragraph*{Statement 3.4 and its compatibility hypotheses.}
For each fixed frame and weight prefix there must be a single
family \(\widehat E^{(j)}(\widehat X,P,z)\) on the infinite
level-\(j\) lattice, and likewise a companion for \(R^{(j)}\).
For every admissible \((\ell,K')\), \(K'\ge j\), and every
label, it agrees with the torus term when the projection of
\(\widehat X\) is injective with its margin and lies one full
level-\(j\) cell layer inside \(\Lambda_j\). A sufficient margin
places \(\widehat X\) in a rectangle with sides less than half
the corresponding torus side. The admissibility used here is
\[
 \ell L^{K'-i}\ge2L^{11}M\check R_i\qquad(i\le j).
 \tag{RF34}
\]
The family has the same analytic bounds and cube-translation
covariance on every such lift. Its coefficient \(b_j\) is a
function only of the frame and weights through level \(j-1\).

Uniqueness follows because every bounded domain, collar included,
embeds in arbitrarily large admissible tori. Existence requires
more. First, the walk, propagator, localized background and gauge
choices must be made once by a rule natural under covering maps,
cube translations and the later dilations. Independently chosen
objects on separate tori cannot be compared. Second, each local
formation formula must use only its supplied neighborhood,
earlier terms restricted there and the common numbers \(x_i,b_i\).
The induction is simultaneous over tori of unbounded side; it
does not define \(b_j\) using only domains that happen to fit a
particular torus.

Third, formation must be local in the history trace on the
prescribed one-layer windows \(X_i^+\). A chart identifying
these traces, the bare weight, prefix coefficients and incoming
local terms identifies the newly formed term on the common
local core. In the one-layer interior the standing history has
empty trace: healed components have already been absorbed into
local terms. This actual standing-history property is what
makes that interior term independent of the label. It is not
implied by an inequality for the modulus of an activity.

The order \(E^{(j)},b_j,x_j,R^{(j)}\) avoids a circular
definition. In particular the new large-field weight \(x_j\)
does not enter the preceding extraction of \(b_j\).
We must prove this compatibility for the full interacting
construction. The finite formulas proved elsewhere can verify
pieces of it; their holomorphy alone does not prove it.

Under these conditions, the coefficient of the physical totals,
taken in the large-torus limit, equals the companion
coefficient in (RF2). Use the all-small-field totals for their
unit-lattice symmetries; a truncated label total has no such
invariance in general. On nonwrapping domains the Hessians
agree by compatibility and the torus response converges to
\(\mathfrak h_\infty\). Domains escaping a large fixed box
carry an exponential tree-length tail. The rooted entropy is
absorbed by the chosen \(\kappa\); response decay controls
remote differentiations. Dominated convergence therefore
passes the differentiated invariant totals to the companion
kernel. This proves the sum-rule identification, including
that a constant total has coefficient zero. It requires the
uniform response convergence and domain bounds just stated.

Wrapping domains have
\(d_j(X)\ge c\ell L^{K-j}/M\). Their exponential decay is used
directly and does not define an additional renormalized
coefficient. For a self-calibrated run, the cube ratio shows
that (RF34) at the terminal level implies it at all preceding
levels. Since \(x_K\le h+B_\sharp\), it suffices that
\[
 \ell\ge\ell_\star(g):=
               2L^{11}M R(g^{-2}+B_\sharp).
 \tag{RF35}
\]
Write \(\mathfrak L(g)\) for the compatible torus sides above
this floor. The same floor works on the reference line,
whose terminal value is \(h\). The side is selected after
\(g\); the cutoff \(K\) remains unrestricted. This statement
is a uniform construction requirement, not an assertion of
infinite-volume uniqueness.

\subsection*{(3c) Frozen choices, slopes, and tuned bare couplings}

The reference line for a run of length \(K\) is
\[
 \theta_k=h+c_0(K-k),\qquad h=g^{-2}.
 \tag{RF36}
\]
Equal depth in runs of lengths \(K\) and \(K+1\) gives equal frame
numbers, lattices, partitions, radii and cube sizes. It does not give
equal tests, since their backgrounds live over different bare lattices,
and it identifies no labels or histories between those runs.
The line beginning at depth \(s\) is the reference line of
\(g_{[s]}=(h+c_0s)^{-1/2}\).

For fixed choices, every inserted characteristic function is the same
real Borel function of its integration fields. The actual operations
act linearly on the initial density. If their required dominated
analytic estimates hold, applying them to
\(e^{-(\beta-\zeta)A}\) gives holomorphic measures and local families
in \(\zeta\). No derivative of a cutoff is taken. Continuity of the
real run follows from the actual fixed-frame input-list Lipschitz
estimate while the arrays retain their tube room. It is continuity,
not a contraction. When each bare coupling chooses its own frame,
the tests and domains can change discontinuously and this argument
does not apply.

The exact decompositions of unity and normalized gauge factors imply
that the bare measure, partition function and observable integrals
are independent of admitted auxiliary choices. Intermediate
densities after a large-field operation need not coincide.
Comparing their extracted couplings requires E4. The retained
self-calibrated versus reference-run target is
\[
 \max_{k\le K}|x_k(\beta)-x_k^\natural(\beta)|
 \le Cg(\log g^{-2})^{3\mathfrak p_0+1}\le1 .
 \tag{RF37}
\]
It is required for first-passage data with
\(\varrho\ge B_\sharp+1\) and
\(0<g\le\gamma_{U'}\le\gamma_U\). Threshold and large-field
changes contribute at most \(C\theta_k^{-(\kappa_0-4)/2}\)
per step, and gauge-fixing changes at most
\(C\theta_k^{-3/2}(\log\theta_k)^{3\mathfrak p_0+1}\).
Proving those actual comparisons remains necessary. Equality of
bare integrals does not prove (RF37).

The small-field coupling polydisc has relative radius
\(\sigma_\star\theta_j\). The large-field tube is narrower.
On a level-\(j\) unit cube, an action-weight change has size
\[
 |x-\theta_j|\alpha_{0,j}^2
 \le C\mathsf t_j\theta_j^{-1}(\log\theta_j)^{2q}.
 \tag{RF38}
\]
A logarithmic tube factor at least \((\log\theta_j)^{-2q}\)
keeps this bounded. A fluctuation difference gains a further
\(\delta_j/\alpha_{0,j}\). This explains the different roles
of the full small-field polydisc and the large-field tube.
It does not justify extending the entire interacting history
construction to the larger polydisc. Complex bare couplings in
the latter construction are obtained near admitted real runs
through the fixed-source analytic domain.

\paragraph*{Statement 3.5: the weaker marked-disc slope route.}
For this paragraph require \(p_R>1\). The unmarked term has the
independent-variable bound in (RF20). Along the actual run the
marked term is allowed a different analytic input: on a disc in
the bare variable about each admitted real centre, of radius at
least \(c_\beta\varrho\), it is holomorphic and bounded by
\(\epsilon'_k\). Bounds are uniform on centres whose required
prefixes retain the prescribed tube room. On a smaller concentric
disc, Cauchy therefore gives \(C\epsilon'_k/\varrho\).
This hypothesis includes the actual remainder dependence on that
bare variable. It does not assert that this dependence is a
function of independent deviations alone.

The dimension-free polydisc lemma,
Lemma~\ref{lem:polydisc-dimension-free-v7}, gives on the half
polydisc
\[
 \sum_i r_i|\partial_iF|\le\tfrac43\sup|F|,
 \qquad
 |F(z)-F(z')|\le\tfrac43\sup|F|
                         \max_i|z_i-z_i'|/r_i .
 \tag{RF39}
\]
One proves it by sending a single unit disc through the point
with coordinate disc automorphisms, choosing their phases to align
the derivative terms, and applying the one-variable Schwarz
inequality. Integrating the derivative bound on the segment gives
the second formula. Neither bound contains the number of earlier
levels. Neither implies age forgetting: the bounded function
\(F(z)=Bz_0/(\sigma_\star\theta_0)\) already has a nondecaying
weighted derivative in its first coordinate.

Here is the needed forced derivative estimate, including the
marked term. Let \(b_i'=d b_i/d\beta\),
\(M_i=\max_{0\le j\le i}|x_j'|\), with \(x_0'=1\).
The frame is fixed, so \(z_j'=-x_j'\).
The tube constants and coupling floor put these evaluation points
in the half polydisc required in (RF39).
Since \(\min_{j<i}\theta_j\ge\theta_{i-1}-c_5
\ge\theta_{i-1}/2\) above the floor, (RF39) and the separate
bare-disc Cauchy estimate imply
\[
 \begin{aligned}
 |b_i'|&\le a_iM_{i-1}+d_i,\\
 a_i&\le C\theta_{i-1}^{-3/2}
                      (\log\theta_{i-1})^{\mathfrak p_0},
 &d_i&\le C\varrho^{-1}\theta_{i-1}^{-p_R}.
 \end{aligned}\tag{RF40}
\]
The \(\theta^{-1}\) factor in \(a_i\) is the unmarked
independent-polydisc gain. The factor \(1/\varrho\) in \(d_i\)
comes from the actual bare-variable disc and must not be
replaced by \(1/\theta_i\).

Put \(V_i=1+\sum_{l\le i}|b_l'|\). The differentiated recursion
gives \(M_i\le V_i\), and hence
\[
 V_i\le(1+a_i)V_{i-1}+d_i,\qquad
 \sum_{i=1}^K|b_i'|
 \le e^{A_K}(1+D_K)-1 ,
 \quad A_K=\sum_{i=1}^Ka_i,\quad D_K=\sum_{i=1}^Kd_i .
 \tag{RF41}\label{fc:forced-gronwall}
\]
Indeed iterating the scalar inequality gives
\(V_i\le\prod_{l\le i}(1+a_l)+
\sum_{l\le i}d_l\prod_{m=l+1}^i(1+a_m)\);
each product is at most \(e^{A_i}\).
No triangular estimate for the physical field response is hidden
in this one-dimensional calculation.

On the reference line, integral comparison yields, uniformly in
\(K\),
\[
 A_K\le Ch^{-1/2}(\log h)^{\mathfrak p_0},
 \qquad D_K\le C\varrho^{-1}h^{1-p_R}.
 \tag{RF42}
\]
The condition \(p_R>1\) is exactly what makes the second bound
uniform over all depths for a fixed bare-disc radius.
For the original full Lipschitz numbers one may now take
\(\mathsf a_{[K-i]}=(4/3)a_i+d_i\), once the uniform slope
bound below is imposed. Thus
\[
 \sum_{s<K}\mathsf a_{[s]}
 \le C\{h^{-1/2}(\log h)^{\mathfrak p_0}
                   +\varrho^{-1}h^{1-p_R}\}.
 \tag{RF42a}
\]
The sharper original \(O(g(\log g^{-2})^{\mathfrak p_0})\)
rate follows when \(p_R\ge3/2\), in particular in the
intended large-\(\kappa_0\) regime. For merely \(p_R>1\),
retain the general forced rate displayed here.
Choose \(h\) large enough that
\(e^{A_K}(1+D_K)-1\le1/3\). Then
\(|x_j'-1|\le1/3\) at every admitted prefix. On real intervals,
\(2/3\le x_j'\le4/3\).
The statement is uniform only where the analytic inputs and
room conditions hold. It does not identify different frames.

The stronger independent marked-coupling polydisc route in the
comparison appendix remains available separately. Under its
actual hypotheses the marked row has order
\(\theta^{-p_R-1}\log^C\theta\); it is not the premise used
in (RF40)--(RF42). Both routes require their actual domains.
In particular a fixed physical-source disc alone does not
supply the unmarked independent-variable gain.

\paragraph*{Canonical tube and tuning theorem 3.2.}
For a general finite reference sequence of length \(n\ge1\),
the largest permitted tube is
\[
 \mathsf t_j^*=
 \min_{1\le t\le\min\{j+1,n\}}\mathsf t^\sharp(\theta_t).
 \tag{RF43}
\]
This is just the combination of all preceding upper bounds and
the requirement that tube radii do not increase with the level.
The one-member case retains its separate convention.
A constant-tube criterion follows by setting
\(\mathsf K(a)=\sup_{y\ge e}y^{-1/4}(\log y)^a<\infty\).
If the frame floor is \(x_\star\ge e\) and
\(\mathsf K(2q_{\rm t})t\le\varpi x_\star^{3/4}\), then
\[
 t\le\frac{\varpi y^{3/4}}{\mathsf K(2q_{\rm t})}
       \le\varpi y(\log y)^{-2q_{\rm t}}
       \qquad(y\ge x_\star).
 \tag{RF43a}
\]
Thus the constant sequence of radius \(t\) is a tube; its
criterion uses the floor rather than the frame length.
On the line (RF36), set
\[
 \begin{aligned}
 \iota_k&=\epsilon_k+\epsilon'_k,\qquad
 T_K=\varrho+c_2,\qquad
 T_k=T_K+\sum_{j=k}^{K-1}\iota_j,\\
 F_k(\beta)&=\theta_k-x_k^\natural(\beta)+s_k,\qquad
 \widehat{\mathsf t}_k=T_k+c_2+2R',
 \qquad R'=4\varrho .
 \end{aligned}\tag{RF44}
\]
This is the canonical tube: in depth notation
\(T_{[s]}=T_{[0]}+\sum_{u<s}\iota_{[u]}\).
The correction by \(s_k\) removes the bounded one-loop
oscillation, giving
\(|F_{k+1}-F_k|\le\iota_k\).
The accumulated drift satisfies
\(\sum_{j=k}^{K-1}\iota_j
\le C\theta_k^{1/2}(\log\theta_k)^{\mathfrak p_0}\)
in the stated power regime. For sufficiently large \(h\),
uniformly over \(K\) and \(k\), this and the fixed additive
terms are smaller than
\(\varpi\theta_{k+1}(\log\theta_{k+1})^{-2q_{\rm t}}\).
The same holds at the terminal convention. Thus
\(\widehat{\mathsf t}\) is an admitted tube.
The choice is made after \(\varrho\), not by shrinking the
prescribed offset range as \(K\) grows.

We spell out the domain premise needed by the interval argument.
At a prefix satisfying \(|F_j|\le T_j\), the actual array is
defined with room \(2R'=8\varrho\). The next coefficient is
formed using that prefix before the new large-field operation.
The uniform bounds and the marked bare-disc hypothesis above
apply on a neighborhood of every such prefix. Whenever the
computed next value satisfies \(|F_{k+1}|\le T_{k+1}\), the
next array closes with the same room. These are the actual
prefix formation/closure hypotheses, not consequences of
the scalar intermediate value theorem.

Define \(I_0=\{\beta:|F_0(\beta)|\le T_0\}\).
Since \(F_0=\theta_0-\beta\), this is a compact interval
mapped decreasingly onto \([-T_0,T_0]\). Choose the floor so
all its initial weights are admitted. Inductively suppose
\(I_{k-1}\) has this property for \(F_{k-1}\), and every
earlier prefix lies in its required room. The preceding
domain premise defines \(b_k\), hence \(F_k\), on a
neighborhood of \(I_{k-1}\), even before restricting its
new output to the target window. Equations (RF40)--(RF42)
make \(F_k'=-x_k'\) strictly negative there.

At the left endpoint,
\(F_{k-1}=T_{k-1}=T_k+\iota_{k-1}\), so
\(F_k\ge T_k\). At the right endpoint the corresponding
inequality is \(F_k\le-T_k\). Continuity supplies unique
points with values \(T_k\) and \(-T_k\), and strict
monotonicity orders them. Their intervening compact interval
\[
 I_k=\{\beta\in I_{k-1}:|F_k(\beta)|\le T_k\}
 \tag{RF45}
\]
is mapped decreasingly onto the whole target interval.
On it
\(|\theta_j-x_j|\le T_j+c_2\) for \(j\le k\), leaving
room \(2R'\) in (RF44). The actual closure premise therefore
permits the induction to continue. This proves existence
and interval connectedness rather than presuming them from
an implicit tuning equation.

At \(K\), \(x_K=h+s_K-F_K\). Since \(|s_K|\le c_2\),
its image contains every \(x\) with \(|x-h|\le\varrho\).
Each such value has exactly one preimage in \(I_K\).
Let \(\gamma_U\le\gamma_J\le\gamma_H\) impose the preceding
conditions for the fixed \(\varrho\), with \(\gamma_J\) also
imposing the admitted source analytic domain. For every \(K\)
and \(0<g\le\gamma_U\), the tuning theorem under these hypotheses is
\[
 x_K^\natural(\beta_K^\natural)=g^{-2},\qquad
 \tfrac12c_0K-c_5\le\beta_K^\natural-g^{-2}
                       \le\tfrac32c_0K+c_5 .
 \tag{RF46}\label{fc:tuned-coupling}
\]
The common companion makes this the same number on every
admissible torus. Uniqueness has been proved in the admitted
canonical interval, not among arbitrary disconnected
presentations at all bare couplings.
Also
\(|\beta_K^\natural-h-c_0K|\le2c_2+\sum_k\iota_k\).
The single-run spacing estimate is
\[
 \tfrac12c_0(j-i)-c_5
 \le x_i^\natural-x_j^\natural
 \le\tfrac32c_0(j-i)+c_5\qquad(0\le i\le j\le K).
 \tag{RF46a}
\]
It follows by summing \(|b_l-b_l^0|\le c_0/2\).
It gives (RF46), but gives no
bound between the tuned numbers for two different cutoffs.

The bare offset \(t\) and terminal coordinate \(x\) are
distinct:
\(x_K^\natural(\beta_K^\natural-t)=h-t+O(c|t|)\)
when the slope error is at most \(c\). The reference frame
is independent of \(\beta\); the self-calibrated frame
\(\theta_j=x_j(\beta)\) and its frozen shifted run are
other applications of the uniform prefix induction.
For a constant source \(\zeta\), the frozen shifted
weights are \(x_j(\beta)-\zeta_j\), and their slope
normalization is \(x_j'(\beta)\). Identifying this derivative
with a particular F3 witness still requires that witness's
admitted run and comparison hypotheses.

\paragraph*{A finite nontriangular comparison lemma.}
The stronger two-run route uses an estimate different from
(RF41). Fix integers \(n\ge1\), \(0\le r_0<n\), numbers
\(0<\rho<1\), \(\Pi,\mathsf M\ge0\), nonnegative
\(y_w\) for \(r_0<w\le n\), and nonnegative \(D_u\),
\(k_{u,w}\), \(h_u,h'\) for the finite indicated ranges.
Put
\[
 S_w=\frac{\rho^{\,n-w+1}}{1-\rho}+(w-r_0)\Pi .
 \tag{RF47}
\]
Suppose, for \(r_0\le s<n\) and \(r_0\le u<n\),
\[
 \begin{aligned}
 y_{s+1}&\le\sum_{u=r_0}^sD_u,\\
 D_u&\le\mathsf M(\rho^{\,n-u}+\Pi)
                           +\sum_{w=u+1}^n k_{u,w}y_w,\\
 \sum_{w=u+1}^nk_{u,w}S_w
       &\le h_uS_{u+1}+h'\Pi,\qquad
 \sum_{u=r_0}^{n-1}h_u+h'\le\tfrac12 .
 \end{aligned}\tag{RF48}
\]
Then \(y_{s+1}\le2\mathsf M S_{s+1}\).
Indeed \(S_w>0\) and increases with \(w\). Let
\(Z=\max_{r_0<w\le n}y_w/S_w\), a finite number.
The forcing sum is at most \(\mathsf M S_{s+1}\).
The kernel sum is at most
\[
 Z\left(\sum_{u=r_0}^s h_uS_{u+1}
                 +(s-r_0+1)h'\Pi\right)
 \le\tfrac12 Z S_{s+1}.
 \tag{RF49}
\]
Taking the maximum gives \(Z\le\mathsf M+Z/2\), as
claimed. The sums couple a difference to later depths,
so this is weighted-max absorption, not a triangular
recursion. Its actual profile-kernel bound in (RF48)
remains a substantive two-run input. Bounds on increments
alone do not establish that input.

\paragraph*{Statement 3.6: real presentations of the same measure.}
A real presentation over a frame is an actual inductive array
whose weights remain in the tube. At each level it carries one
Euclidean member, with one scalar coefficient common to all
sites and no defect, and one shift \(\mathsf c_j\). These shifts
obey the admitted \(\bar c\) bound, with the frame floor chosen
in terms of \(\bar c\). Members, companions and coefficients
are independent of the admissible torus and real at real fields.
The all-small-field label is the label whose regions fill the
torus. In its exponent the action coefficient is
\(x_j-\mathsf c_j\). The weight actually used in the Gaussian
step is \(x_j\), and \(\mathsf c_jA\) is carried in the local
action. A spatially varying source array with defects is not
a presentation of this special kind.

Write \(b_j=x_{j-1}-x_j\), let \(\widehat b_j\) be the
coefficient actually extracted at step \(j\), and set
\(u=\mathsf c_0\). Coefficient bookkeeping gives
\[
 \mathsf c_j=u+\sum_{l=1}^j(\widehat b_l-b_l),\qquad
 x_j-\mathsf c_j=x_{j-1}-\mathsf c_{j-1}-\widehat b_j .
 \tag{RF50}
\]
In particular \(\widehat b_j\) need not equal the planned
decrement \(b_j\) when the shifts vary with the level.

The comparison hypothesis is equality of the all-small-field
measures at every level, on every admissible torus, for two
real presentations over the same frame. Require hereditary
open windows \(\mathfrak O_m\ni\mathbf1\) of block fields:
their backgrounds lie in every needed core, all the relevant
tests equal one, and
\(\mathcal M^iU_m(V)\in\mathfrak O_i\) when \(V\in\mathfrak O_m\)
and \(i<m\). Also require the common connected analytic
fluctuation domain on which continuation from such a window
to the next actual integral is used. These are properties of
the actual maps and tests.

Under these hypotheses the conclusions are
\[
 \begin{aligned}
 x_j-\mathsf c_j&=x'_j-\mathsf c'_j,\\
 \widehat b_j&=\widehat b'_j,\\
 \sum_X(E^{(j)}-E'^{(j)})(X,U)
   &=\tfrac32|\mathfrak X_{j-1}|
                         \log(x_{j-1}/x'_{j-1}).
 \end{aligned}\tag{RF51}
\]
The first line holds through level \(n\), and the coefficient
equality through \(n+1\) only when both presentations admit
that next step. The total identity is evaluated at backgrounds
\(U=U_m(V)\), \(j\le m\le n\), of the common windows.
On those backgrounds the remainder totals agree, and the
vacuum constants agree. When the total suppresses \(z\),
\(E^{(j)}(X,U)=\sum_{z\in X}E^{(j)}(X,P_U,z)\).
No assertion identifies individual localized terms of the two
presentations.

Here is the induction. At the bare level equality identifies
the effective weight \(\beta-u\), since the action is not a
constant function. On an open small-field window the positive
continuous densities agree everywhere if they agree almost
everywhere, because product Haar measure has full support on the
open window. Their logarithmic exponents are real, so their
equality involves no integer ambiguity. Analytic continuation
on the admitted fluctuation domain therefore identifies the
next unnormalized integral. The normalized Gaussian in each
presentation uses that presentation's own weight. Its
normalization has the term
\(-3|\mathfrak X_{j-1}|\log(x_{j-1}/\theta_{j-1})/2\).
It follows that the next regular totals differ by the
field-independent constant displayed in (RF51).
The common-companion sum rule makes their extracted
coefficients equal. Equation (RF50) then identifies the next
effective weights. The same physical formation identities on
the hereditary windows identify remainder totals and the
once-counted vacuum constants. This completes the induction.
It uses all-torus compatibility to extract a coefficient
from a constant total; one finite torus alone is insufficient.

Apply the result to a real \(\beta\) whose prefix along
\(\theta\) has room \(2\varrho\), and a real
\(|\zeta|<\varrho\) for which both constructed presentations
are admitted. One presentation starts at \(\beta-\zeta\)
with no shift; the other starts at \(\beta\) with constant
source \(\zeta\) and running shift \(\zeta_k\). The exact
identity is
\[
 x_k^\theta(\beta-\zeta)
                  =x_k^\theta(\beta)-\zeta_k .
 \tag{RF52}
\]
The slope bound keeps the running shift at most
\(4|\zeta|/3\) in its stated range. For
\(\theta=x(\beta)\), this describes the frozen shifted run,
not the self-calibrated run of \(\beta-\zeta\).
For the reference line it identifies the two presentations
of \(x_k^\natural(\beta-\zeta)\).

Complex gluing also uses one frame. First restrict the real
centres to the canonical interval \(I_K\) constructed above,
or to one of its admitted prefix intervals. Each source disc
is supplied by the actual analytic source construction.
Where a real point in that disc also belongs to the canonical
interval, (RF52) identifies its continuation with the real run.
If two discs about such centres overlap, their real overlap
contains an open interval in \(I_K\): the segment between
the centres lies in \(I_K\), and the open real disc intervals
overlap along that segment. For coincident centres at an
endpoint, a one-sided open interval inside the nondegenerate
\(I_K\) suffices. The identity theorem therefore identifies
the coefficient functions on the full complex overlap.

The union carries a single holomorphic bare-coupling
function. Remainder totals on common windows glue likewise;
regular totals glue after accounting for their explicit
Gaussian constant. This proves uniqueness of the supplied
continuations, not their original bounds or radii.

The broader target for arbitrary admitted real centres is
retained conditionally: every overlap used for gluing must
contain a common open interval of ADMITTED real runs, or the
actual prefix/source-disc construction must prove real
admission there. That extra obligation is substantive if
the real admitted set has disconnected components.
Geometric intersection of two discs alone does not settle
it. With the stated admission property the same identity
theorem proof applies. Nothing here identifies runs made
with different frames, individual localized terms about
different centres, or functions on independent-current
cores from equality only on a physical graph.

\subsection*{(3d) Units, common prefixes and change of reference}

\paragraph*{Statement 3.7: the finite dilation identity.}
For an integer \(b\) with \(0\le b\le K\), the map
\(D_bx=L^bx\) identifies the bare lattices at
\((K,\ell)\) and \((K-b,L^b\ell)\). The exact result of
Section~\ref{ch:flow-arithmetic} is
\[
 \begin{aligned}
 Z_{K,\ell,\beta}(w)
     &=Z_{K-b,L^b\ell,\beta}(w^{[b]}),\\
 S^T_{n;K,\ell,\beta}(\underline f)
     &=S^T_{n;K-b,L^b\ell,\beta}(\underline f^{[b]}),
 &f^{[b]}(x)&=f(L^{-b}x).
 \end{aligned}\tag{RF53}
\]
There is no spacing prefactor in
\(\mathcal O(f)=\sum_pf(x(p))a_p\).
The map transports bonds, plaquettes and Haar variables
bijectively; both sides are the same finite integral.
The partition identity is entire in its finite source
variables and bare coupling. Logarithmic derivatives use
its nonzero domain. The same transport identifies periodic
Gibbs states on the two infinite lattices.
Radii and distances multiply by \(L^b\),
\(\operatorname{Lip}f^{[b]}=L^{-b}\operatorname{Lip}f\),
and \(\|f^{[b]}\|_\sharp\le\|f\|_\sharp\).
The free bilinear form is invariant. This proves neither
convergence nor agreement of two RG constructions.

\paragraph*{Statement 3.8: what the step reads.}
Construction naturality is the additional assertion that
step \(k\) reads only its actual tower through \(k\),
partitions, frame and weight prefixes, earlier outputs
and fixed supply rules. Those rules are made for the
lattice, not for its unit of length. For two presentations
of the same bare lattice, their images under \(D_b\)
therefore identify the histories and their measures,
regions, healing data and local nonwrapping terms through
level \(K-b\), provided the prefixes agree. This is required
for both fixed-frame and self-calibrated runs. Separate
arbitrary choices on the two lattices would not establish it.

For a self-calibrated run this permits continuation from
first passage at \(g\) to first passage at \(g'\), provided
the prolonged run remains below the anchor ceiling and
the rescaled torus is admissible. Its history law through
the earlier level is unchanged. For reference runs,
\(\theta_{[r+s]}(g)=\theta_{[r]}(g_{[s]})\);
the prefix at \(g\), down to depth \(s\), is the
reference run at \(g_{[s]}\) of length \(K-s\), on
\(T_{L^s\ell}\). Its deviations, increments and tube
are the same. The implication
\(\ell\in\mathfrak L(g)\Rightarrow
L^s\ell\in\mathfrak L(g_{[s]})\) uses the stated torus
floor and parameter choices.

Keep the beyond-first-passage scope separate. If
\(\beta\in\mathcal W_{K'}(g)\) and \(K=K'+q\), \(q\ge0\),
then (RF53) identifies the cutoff-\(K\) measure with the
first-passage measure on \(T_{L^q\ell}\).
Estimates for it require \(L^q\ell\in\mathfrak L(g)\).
A torus volume factor becomes
\(\exp[nc_V(L^q\ell)^4]\), whereas a per-ball estimate
has its radius and separation changed to \(L^qR,L^qd\)
and no side dependence in its constant:
\[
 |S^T_{n;K,\ell,\beta}(\underline f)|
 \le A(L^qd;L^qR)B(L^qR)^n(n!)^{\varkappa_\star}
                         \prod_i\|f_i\|_\sharp .
 \tag{RF54}
\]
The stated lower torus floor for that bound also applies.
The factorial finite-torus estimate (RF54) is an additional
conditional target: its use requires that estimate for the actual
finite-torus source construction. Theorem~\ref{thm:graded-rooted-source}
proves its stated graded bound for periodic Wilson DLR states;
it does not by itself supply (RF54), nor a polynomial separation
prefactor of degree independent of the number of tests.
The finite dilation identity transports a supplied bound and
does not prove this additional target.
For \(q<0\), use the identity in reverse, require
\(L^q\ell\in\mathfrak L(g)\), and allow a norm cost
at most \(L^{|q|}\) per test. Applying this to tuned
offsets requires the separate comparison
\(|K(\beta_K^\natural-t,g)-K|\le C(\varrho)\).
It does not follow from tuning at a frozen frame.

\paragraph*{Statement 3.9: all four reference-change targets.}
Let
\[
 S_\ell^{[t]}[g]=
   \lim_KS^T_{\cdot;K,\ell,\beta_K^\natural(g)-t},
 \qquad \mathsf J=[-\varrho/4,\varrho/4].
 \tag{RF55}
\]
Existence along the whole sequence, its common complex
offset strip and moving-offset continuity require E3 and
the terminal-coordinate comparison in
Appendix~\ref{ch:comparison-capsule-v7}.
The strip width and the cutoff floor may depend on the
fixed torus. The coordinates are
\[
 \Xi_K(t)=x_K^\natural(\beta_K^\natural-t),\qquad
 \Xi_{\infty,s}(t)=
   \lim_Kx_{K-s}^\natural(\beta_K^\natural-t).
 \tag{RF55a}
\]
The depth limit requires the matched comparison. Terminal
matching uses \(n_K=\Xi_{K+1}^{-1}\circ\Xi_K\) on its
admitted common strip. Its near-identity and summable-error
properties are part of obligation E3, which we must prove; they are not consequences of the
existence of individual tuning inverses. First choose \(\varrho\) large compared with
\(c_2,B_\sharp+1\) and a bound for the tuned increments;
then choose \(g\) small. The bound for those increments
is a two-run comparison requirement fixed before
\(\varrho\), not a consequence of one-run spacing.

(i) At fixed depth \(s\), the same bare coupling has two
offset labels:
\[
 \begin{aligned}
 \widehat t_{K,s}
   &=\beta^\natural_{K-s}(g_{[s]})-\beta^\natural_K(g),\\
 \beta^\natural_K(g)-t
   &=\beta^\natural_{K-s}(g_{[s]})
                              -(t+\widehat t_{K,s}).
 \end{aligned}\tag{RF56}
\]
Require the actual two-run estimate
\(|\widehat t_{K,s}|\le
C(c_2+\sum_{K-s\le j<K}\iota_j)\), its summable
cutoff increments, and the comparison domain on which
both presentations exist. The initial comparison requires
\(2c_2+\sum_{K-s\le j<K}\iota_j\le\varrho\); a smaller
fixed fraction of that bound keeps the offsets in the inner window.
For \(g\) small in terms of
the fixed \(s\), these offsets lie in
\(\frac12\mathsf J\) and converge to \(\widehat t_s\).
The whole-sequence moving-offset limit then gives
\[
 \begin{aligned}
 S_\ell^{[t]}[g](\underline f)
 &=S_{L^s\ell}^{[t+\widehat t_s]}[g_{[s]}]
                                  (\underline f^{[s]}),\\
 \widehat t_{s+s'}(g)
 &=\widehat t_s(g)+\widehat t_{s'}(g_{[s]}).
 \end{aligned}\tag{RF57}
\]
Here \(t\in\frac12\mathsf J\); each depth condition
must hold at its displayed reference coupling.
The condition on \(g\) is an upper bound, never a
lower bound. The terminal coordinate on the right is
the limiting depth coordinate \(\Xi_{\infty,s}(t)\).
Its existence belongs to the matched comparison.

(ii) For arbitrary sufficiently small \(g,g'\) and
\(t\in\mathsf J\), retain the stronger target that some
\(t'\in\mathsf J\) and integer \(b\) satisfy
\[
 S_\ell^{[t]}[g](\underline f)
       =S_{L^b\ell}^{[t']}[g'](\underline f^{[b]}).
 \tag{RF58}
\]
Its torus conditions are both
\(\ell\ge\ell_\star(g)\) and
\(\ell\ge
L^{(|g'^{-2}-g^{-2}|+C)/\lambda_\sharp}
\ell_\star(g')\).
The sign of \(b\) is that of
\(g'^{-2}-g^{-2}\), up to a bounded ambiguity.
The number of intervening passage levels lies between
\(|g'^{-2}-g^{-2}|/B_\sharp-1\) and
\((|g'^{-2}-g^{-2}|+C)/\lambda_\sharp\);
\(|b|\) differs from it by \(O(1)\).
This assertion needs recognition of the same bare run
in both calibrated families, not merely (RF53).

(iii) Choose once an anchor
\(0<\gamma_\star\le\gamma_J\), sufficiently small in
terms of the fixed large-field window dials \(\mathsf w,c_a\). For any proposed
\(\mathbb R^4\) limit at \(g\le\gamma_\star\), retain
the prolonged-run target: after a subsequence,
\(K(\beta_j,\gamma_\star)-K_j=s_\star\) is constant,
of order \(g^{-2}-\gamma_\star^{-2}\), and
\[
 S_n(\underline f)=S'_n(\underline f^{[-s_\star]}).
 \tag{RF59}
\]
The anchor limit uses the same bare-coupling sequence.
This preserves the target for arbitrary defining
sequences carrying the required construction and
recognition hypotheses. It does not give an automatic
mirror/cone theorem for unrestricted growing-torus
limits. Dilation-invariant properties need be checked
at only one reference coupling once this target is
established.

(iv) The bare tuned increments
\(\alpha_K=\beta^\natural_{K+1}-\beta^\natural_K\)
are a new notation, unrelated to the field radii.
The retained conclusion is
\[
 \alpha_K\longrightarrow c_0,\qquad
 S_\ell^{[t]}(\underline f)
     =S_{L^b\ell}^{[t-bc_0]}(\underline f^{[b]})
 \quad(b\ge1,\ t,t-bc_0\in\mathsf J).
 \tag{RF60}
\]
For completeness the shape convention is
\(\mathfrak T_{\mathbb R^4}(d;\vartheta)\): both profiles are
nonnegative nonzero Lipschitz functions supported in radius-\(d\)
balls, with support distance in \([d,6d]\), and each obeys
\(\int f\ge\vartheta d^4\|f\|_\infty\) and
\(\operatorname{Lip}f\le\|f\|_\infty/(\vartheta d)\).
Clause (iv) retains the condition
\(\mathfrak T_{\mathbb R^4}(1;\vartheta)\ne\emptyset\) and
\(g\le\gamma_W(\vartheta)\). The proof also needs the actual
uniform smooth-profile envelopes on the pairs used in its
iteration; class nonemptiness by itself is not an envelope estimate.
The comparison appendix proves this implication with
the actual two-run depth/profile estimates, summable
comparison errors and uniform two-point envelopes.
Those envelopes are uniform in the offsets and tori
needed for repeated blocking; the fixed shape class
must be nonempty. Two distinct cluster increments
would give an offset period by analyticity, and
repeated blocking would contradict the uniform
short-distance envelopes. The unique increment
limit is then \(c_0\) by the averaged one-loop law.
The same argument yields
\(\beta^\natural_{K+b}-\beta^\natural_K\to bc_0\)
for every fixed integer \(b\).
The exact input hypotheses remain in force; finite
dilation and weak numerical inequalities do not
establish them.

\subsection*{Vacuum constants and the interfaces used later}

The vacuum recursion is (CP1) of
Appendix~\ref{ch:comparison-capsule-v7}.
Each assigned step number, Gaussian local number and
regular/remainder value at \(P_\circ\) is counted once.
Using the normalized \(N\) requires adding its removed
local numbers; using \(N^\circ\) does not. At fixed
calibration these numbers have zero source derivative.
They are not additional freely chosen normalizations.

The further vacuum targets retained here are
\[
 \begin{aligned}
 \log(\mathsf c_k(\theta_k)Z^{(k)}(\mathbf1))
      &=\mathsf N^{(k)}
               +\sum_X\sum_{z\in X}\mathsf n^{(k+1)}(X,z),\\
 \mathsf n_k^{\rm tot}
      &=36N_{k+1}\log L
        -\tfrac12\log\det(\mathsf C^{\mathsf T}
                             \Delta^{(k)}\mathsf C)|_{\mathbf1},
 &|\mathsf n_k^{\rm tot}|&\le CN_{k+1},\\
 \mathsf n_0^{\rm tot}
      &\le-3(L-1)^4N_1\log2<0 .
 \end{aligned}\tag{RF61}
\]
Here \(N_j\) counts sites, \(\mathsf c_k\) is the
actual change-of-variable normalization and
\(\mathsf N^{(k)}\) is the free-bond/coarse-site
count, with \(-3\log\theta_k/2\) per free bond.
The map \(\mathsf C\) is the actual linear embedding
of nonpivot coordinates into full constrained
fluctuations; this determinant is on the free color
coordinates. It is not the covariance \(C^{(k)}\).
The local numbers are Gaussian values before
normalization, real and coupling independent,
with common nonwrapping lifts and the Gaussian
domain bound. Their coefficients vanish.
The count alone is not the entire step number.
The determinant estimates and first-step sign in
(RF61) remain actual Gaussian targets, beyond
the once-only accounting identity.

The explicit leading vacuum requirement is
\[
 \left|E_{K,\ell,\beta}
   +\tfrac32(N^2-1)(L^4-1)
        \sum_{j=0}^{K-1}(\log x_j)N_{j+1}\right|
                  \le CN_0 ,
 \tag{RF62}
\]
uniform in the admitted runs, tori and cutoffs.
The estimate (RF62) is the quantitative hypothesis on the vacuum term
used here.

The order of choices is retained. Fix the geometric
exponents and \(L\), with the Gaussian constants and
one-loop data; then choose \(\kappa\) above the domain
entropy, the required lower block scales and \(M\).
Choose \(\sigma_\star\), absolute radii below \(c/M\),
\(K_{\mathbf M}\), and
\(E_0\ge\max\{2E_{\mathsf N},4K_{\mathbf M}\sigma_\star\}\).
Next choose \(N_\star\), the ceiling on \(\mathsf s\),
the unmarked constants, \(C_\beta\) and the weak
triple. Compute the marked-load and tube constants
at that triple before choosing \(\gamma\), the
strong triple, \(\gamma_H\) and the smaller drift
ceiling. The source radius and comparison windows
have their stated place before the final smallness
of \(g\). Only then are \(\ell,\beta,K\) varied.
Each claimed uniformity depends on respecting this
order.

The later uses retain their distinct scopes.
Beyond-first-passage data use (RF54); prolonged runs
use statement 3.8; fixed-depth reference prefixes
and terminal coordinates use (RF57).
The germ/rigidity route in the comparison appendix
also carries its own shape and second-difference
hypotheses. In the anchor unit its correlation
constants acquire \(L^{-(4-n)s_\star}\).
Blocking and distinct-member comparisons use
(RF60), not just a relabeling of \(K\).
The sub-unit cells of the angular route keep
their separate moment, mirror and ultraviolet
interfaces. None of these preserved stronger
targets changes the conditional status of the
ordinary interacting source construction,
the actual clean-class two-point extraction,
the external reconstruction theorem or the
identification with the full gauge theory.

% END MODULE chapters/06_flow.tex


% BEGIN MODULE chapters/07_sources.tex
\section{Sources, local energies, and separated correlation bounds}
\label{candidate:sources}
\candidatebanner{C and D}


% BEGIN MODULE chapters/source_domains_v2.tex
\subsection*{Source domains and the exact running-shift identity}

We define the source variables and prove their finite-dimensional
identities. Their use for effective energies takes the multiscale
estimates from the source propositions stated separately.

Fix $K$, a torus $T_\ell$, an odd $L>1$, and positive constants $M,c_M$ with $c_M\ge M$.
The depth $s$ corresponds to physical spacing $L^{-s}$; level $j$ of a
run of length $K$ corresponds to depth $K-j$. A source triple is

\begin{gather*}
\FitMath{w=(c,(f_1,\ldots,f_n),(z_1,\ldots,z_n)),} \\[3pt]
\FitMath{w(x)=c+\sum_i z_if_i(x),}
\end{gather*}

where $f_i$ are real Lipschitz functions, $c,z_i$ are complex, and the
family may be empty. Put

\begin{gather*}
\FitMath{\|f\|_{\sharp,s}=\max\{\|f\|_\infty,
                         c_ML^{-s}\operatorname{Lip}f\},} \\[3pt]
\FitMath{\|w\|_{\flat,s}=|c|+\sum_i|z_i|\|f_i\|_{\sharp,s}.}
\end{gather*}

The class $\mathcal W_r^{(s)}$ means the triples with
$\|w\|_{\flat,s}<r$. It is a union of finite-dimensional source balls,
one for each fixed real family. Holomorphy on this class means holomorphy
in $(c,z)$ on each such ball. The number $\|w\|_{\flat,s}$ belongs to the
chosen triple and can change when its function is represented differently.
The physical partition function and every source-local energy are required
to agree for two triples that have the same values on all bare plaquettes
on which that quantity depends. This representation-independence is a
separate part of the source format.

Always $\|w\|_\infty\le\|w\|_{\flat,s}$ and
$c_ML^{-s}\operatorname{Lip}w\le\|w\|_{\flat,s}$. Thus, if
$m=K-s$ and $j\le m$,

\begin{gather*}
\FitMath{c_Ma_{K-j}\operatorname{Lip}w
   =L^{-(m-j)}c_Ma_{K-m}\operatorname{Lip}w
   \le L^{-(m-j)}\|w\|_{\flat,s},} \\[3pt]
\FitMath{a_t=L^{-t}.}
\end{gather*}

This is the exact scale gain used when a source is frozen on an earlier
block. It follows from the norm definition, not from an analyticity theorem.

Here is an explicit algebraic version of the running profiles. Let $P_K$
be the finite set of bare plaquettes, with barycentres $x(p)$. For each
level $j$ choose a finite set $Y_j$ of base points, anchors $p_y\in P_K$,
and nonnegative functions $h_y^{(j)}$ on $P_K$ such that

\begin{gather*}
\FitMath{\sum_{y\in Y_j}h_y^{(j)}(p)=1,} \\[3pt]
\FitMath{h_y^{(j)}(p)\ne0\,\Longrightarrow\
       |x(p)-x(p_y)|\le C_hMa_{K-j}.}
\end{gather*}

Choose also profiles $0\le\varphi_j(p)\le1$. The grid, anchors, partition
and profiles are fixed independently of the source throughout the analytic
family. They are supplied by the geometric construction; the identities
below hold for every finite choice satisfying these conditions.

Let $b_j$ be holomorphic on $D(0,\rho)$ and define

\begin{gather*}
\FitMath{d_j=\sup_{0<|z|<\rho}\frac{|b_j(z)-b_j(0)|}{|z|},} \\[3pt]
\FitMath{\sum_{j\le k}d_j\le\frac13.}
\end{gather*}

For triples whose function has $\|w\|_\infty<\rho$, set

\begin{gather*}
\FitMath{x_k(p)=\beta-\sum_{j\le k}b_j(0)\varphi_j(p),}
\end{gather*}


\begin{gather*}
\FitMath{c_j(p;w)=\varphi_j(p)\sum_{y\in Y_j}
                   h_y^{(j)}(p)b_j(w(x(p_y))),}
\end{gather*}


\begin{gather*}
\FitMath{\zeta_k(p;w)=w(x(p))+
   \sum_{j\le k}\varphi_j(p)\sum_{y\in Y_j}h_y^{(j)}(p)
             [b_j(w(x(p_y)))-b_j(0)].}
\end{gather*}

Then the exact bookkeeping identity is

\begin{gather*}
\FitMath{x_k(p)-\zeta_k(p;w)
      =\beta-w(x(p))-\sum_{j\le k}c_j(p;w).}
\end{gather*}

To prove it, expand both sides and use $\sum_yh_y^{(j)}=1$ in each
constant term $b_j(0)\varphi_j$. Consequently the Wilson weights in a
sourced exponent can be added or subtracted without changing their total
bare coefficient. This assertion concerns that total coefficient; it does
not assert equality of two effective measures unless the exact integral
transport has also been proved.

The same definitions give

\begin{gather*}
\FitMath{|\zeta_k(p;w)-w(x(p))|\le\tfrac13\|w\|_\infty,} \\[3pt]
\FitMath{|\zeta_k(p;w)|\le\tfrac43\|w\|_\infty.}
\end{gather*}

Each value is holomorphic on every source ball with supremum norm below
$\rho$, because only finitely many evaluations and compositions occur.
It depends on the source only inside distance $C_hMa_{K-k}$ from $p$,
including the value at $p$ itself; here $a_{K-j}\le a_{K-k}$ for $j\le k$.

For a constant source $w\equiv z$, in a region where every
$\varphi_j(p)=1$, this reduces to

\begin{gather*}
\FitMath{\zeta_k(z)=z+\sum_{j\le k}[b_j(z)-b_j(0)],} \\[3pt]
\FitMath{\zeta_0(z)=z,} \\[3pt]
\FitMath{\zeta_{k+1}(z)=\zeta_k(z)+b_{k+1}(z)-b_{k+1}(0).}
\end{gather*}

Differentiating at zero gives
$\zeta_k'(0)=1+\sum_{j\le k}b_j'(0)$ and
$|\zeta_k'(0)-1|\le1/3$. The derivative bound follows from the definition
of $d_j$ by taking the limit at zero. For spatially varying sources, freezing
the function $w$ itself costs at most $\operatorname{Lip}w$ times the anchor
distance. To freeze $b_j\circ w$ one must work on a smaller source disc.
Indeed $h_j(z)=[b_j(z)-b_j(0)]/z$ extends holomorphically to zero and is
bounded by $d_j$. For $|z|\le\rho'<\rho$, Cauchy's estimate gives

\begin{gather*}
\FitMath{|b_j'(z)|=|h_j(z)+zh_j'(z)|
       \le d_j\frac{\rho}{\rho-\rho'}.}
\end{gather*}

Thus on the half-disc the additional factor is at most $2d_j$.
Apply the mean value integral on the segment between the two source values
to bound their composed difference. The displayed depth identity then supplies
the factor $L^{-(m-j)}$, with the fixed anchor constant $C_hM/c_M$.
The radial bound $d_j$ alone is not a Lipschitz bound on the whole open disc.

All smallness, radius and representation-independence conditions used here
must be retained when invoking the sourced RG format. They do not follow
from the entire dependence of the bare lattice partition function alone.

% END MODULE chapters/source_domains_v2.tex

\subsection*{4. Transporting the observables: sources, analyticity, and bounds independent of the cutoff and of the volume}

The measure $\mu_{K,\ell,\beta}$ is described at the unit scale in Steps 1--3, and the present step gives the corresponding description of its correlations. Because $\mathcal{O}(f)=A(f,U)$, a source $w=c+\sum_iz_if_i$ acts as a complex shift of the bare inverse coupling that depends on position:

\begin{gather*}
\FitMath{\rho^w_0=e^{-A(\beta-w,U)},} \\[3pt]
\FitMath{Z(w)=\int\rho^w_0\,dU,} \\[3pt]
\FitMath{S^T_{n;K,\ell,\beta}(f_1,\dots,f_n)=\partial_{z_1}\cdots\partial_{z_n}\log Z(w)\big|_{w=0}.}
\end{gather*}

The renormalisation steps are applied to $\rho^w_0$. Whatever the source leaves behind is gathered into local numbers, the source energies $e_j(X;w)$ of (4c). We set $\mathcal{E}(w):=\sum_{1\le j\le K}\sum_{X\in\mathcal{D}_j}e_j(X;w)$, $\widetilde{Z}(w):=e^{-E-\mathcal{E}(w)}Z(w)$ and $C_g:=C(1+\log g^{-2})$.

\textbf{Theorem 4.1 (The construction with sources).} Let $0<g\le\gamma_J$. For each first-passage datum $(K,\beta)$, each $\ell\in\mathfrak{L}(g)$ and each source satisfying $\|w\|_\flat<2\varrho$, the pair $(E,(e_j(X;\cdot)))$ forms a system of localised source energies. That is, each $e_j(X;w)$ is holomorphic in $(c,z)$ (analyticity). If the $f_i$ are real, it takes real values whenever $(c,z)$ is real. It depends on $w$ only through its values on $X^{\sim1}$, and it is zero at $w=0$ (locality). Moreover, uniformly,

\begin{gather*}
\FitMath{|e_j(X;w)|\le Ce^{-\kappa_ed_j(X)},} \\[3pt]
\FitMath{|\widetilde{Z}(w)|\le e^{C_g\ell^4},} \\[3pt]
\FitMath{\widetilde{Z}(0)\ge e^{-C_g\ell^4}}
\end{gather*}

(these are, in order: the exponential bound, whose rate $\kappa_e:=c\kappa$, $c<1$, exceeds the entropy rate of the localisation domains through a cube; upper stability; the floor). The same conclusions hold along a reference run, with every choice frozen at $\theta_k$, as long as $x^\natural_k$ remains in a tube around $\theta_k$ with a margin of $C\varrho$. For each $K$ this requirement is a condition on $\beta$ only. When $g\le\gamma_U$, it is met on the interval $I_K$ of (3c), an interval that contains the tuned $\beta^\natural_K$. The conclusions also hold for the run stopped at any level $k$, where the sums are stopped at $j\le k$, the vacuum constant $E_{(k)}$ of the first $k$ steps replaces $E$, and $C(1+\log x_k)$ per site of $\mathbb{T}^{(k)}$ replaces $C_g\ell^4$ (4b).

\textbf{Theorem 4.2 (Correlation bounds: Theorem 0 (c)(1)--(2)).} Take $g,K,\beta,\ell$ as above and $n\ge2$. Let the $f_i$ be real Lipschitz functions whose supports lie at mutual distances $\ge\mathsf{d}>0$. For the second and the third bound, let each support also be
contained in a ball of radius $R$. The third bound below concerns
periodic Wilson DLR states. It assumes real $C^1$ tests and the
ordinary source hypotheses on the compatible cofinal torus family
with the actual rooted geometry specified by D1.
These are hypotheses of the graded deduction used below. Then the following estimates hold in their stated scopes:

\begin{gather*}
\FitMath{|S^T_n|\le C_n(\mathsf{d})e^{nc_V\ell^4}\prod_i\|f_i\|_\sharp;} \\[3pt]
\FitMath{|S^T_n|\le C^\infty_n(\mathsf{d};R)\prod_i\|f_i\|_\sharp\,\,(\ell\ge Cn(R+1));} \\[3pt]
\FitMath{|S^T_n[\nu]|\le [2B(R)]^n(n!)^{\varkappa_\star}\max\{1,\mathsf{d}^{-5n}\}\prod_i\|f_i\|_\sharp,}
\end{gather*}

uniformly, where every constant may also depend on $g$ and
$C^\infty_n(\mathsf{d};R)$ does not increase with $\mathsf{d}$.
In the third bound, $\varkappa_\star$ is one finite integer and
$B(R)=C(R+1)^4$, as in Theorem~\ref{thm:graded-rooted-source}
under the actual rooted geometry and analytic hypotheses of D1 and G1.
The second bound also holds for $S^T_n[\nu]$ in each periodic
Gibbs state $\nu$ at $\beta$ on $\mathbb{L}_K$, with no condition
on the volume. The third bound is the stated graded theorem for
such periodic Wilson DLR states; it is not a finite-torus factorial
conclusion. A uniform finite-torus factorial estimate remains an
additional target. A stronger bound with a single polynomial
separation prefactor $A(\mathsf d;R)$ of degree independent of $n$
remains an additional target. It does not follow from the graded
estimate displayed here and is not used as a proved consequence of D2.

\ensuremath{\dagger} Theorem 4.13 is the third main result. Under the law of the histories, the probability that a cell is live, or that it is created, is $\le Ce^{-cp_0(g_k)}$, uniformly in the volume. The short-distance laws rest on this result. (7b) treats the two-point law, clauses (c)(3), (c)(4). It uses these bounds on a single cell and, from their age-resolved form (Proposition 4.14), the age series at a rate $b\ge8$. (7d) treats the four-point law, clause (c)(5), and uses the series at a rate $b>16$.

The same tools also give two further results that are used later. The first is the one-point transfer (Proposition 4.12 [proved here; Section~\ref{ch:all-state-source-transfer}]), $|\nu(\mathcal{O}(f))-\nu'(\mathcal{O}(f))|\le C_1(R)\|f\|_\sharp$, which compares periodic Gibbs states taken at the same first-passage datum $(K,\beta)$. With it, one can centre a pure phase at the means of that periodic Gibbs state whose decomposition contains the phase, the coordinates of (10c) become compact, and the one-point numbers $\mathfrak{c}_\xi$ of Theorem 0'(3)(ii) are bounded. It also enters the growth constants of (b)(v), so that a single set of constants covers these components as well. The second consists of the insertion bound and the plate bound (Propositions 4.17 [proved here; Section~\ref{ch:uniform-insertion}], 4.18), which are the a priori estimates on which (10b) rests.

(5b) uses only the bounds of Theorem 4.2. (5c) (for the re-centring), (6a), (6c) and (10c)--(10e) also use the one-point transfer of Proposition 4.12. The tuned couplings need not be first-passage data. (5a) and Step 8 apply Theorem 4.1 along reference runs, and this gives the tuned limits the torus bound only. The format underlying Theorem 4.1 is the starting point of (7a)--(7b). The difficulty lies in two places. The first is Theorem 4.1: the source term is marginal, since a step leaves it no smaller, just as it leaves the action no smaller; it takes complex values, so positivity is no longer available; and it must be carried through the large-field operations. The second is Theorem 4.13, in which the ratio of the stability bounds loses $e^{C_g\ell^4}$. The finite Cauchy and Gibbs-state transfer implications are proved
under their explicit analytic inputs in the corresponding sections.
The history witnesses, their locality and replication, and their
age-resolved weights in Lemma 4.15 remain construction requirements;
reflection positivity by itself does not verify them.

\subsubsection*{(4a) The curvature observable as a source, and the reason the choices of the construction stay fixed}

\textbf{Objects.} The norm $\|f\|_\sharp=\max\{\|f\|_\infty,CM\operatorname{Lip}f\}$ controls how much $f$ can vary within any $M$-cube of the unit lattice. The sourced densities $\rho^w_k$ arise from $\rho^w_0$ through the operations of Steps 1--2, and every choice in them is the one made in the run at $w=0$. These choices are the tests $\chi$, the thresholds $\varepsilon_k,\delta_k$, the radii, the sizes $\check{R}_k$, the backgrounds $U_k(\mathbf{V})$ (which minimise the action with no source), the scaling $\mathcal{A}_k=g_k\mathcal{A}^{\mathrm{s}}_k$, the covariances $C^{(k)}$, the gauge-fixing weight and the reference laws of $\mathbf{R}$.

\textbf{Proposition 4.3 (Kernels; same histories; exactness).} To each history $h$ of level $k$ there corresponds a positive kernel $\mathsf{K}_h(U,dV_k)$, built from the choices of the run and from nothing else. For $\Phi$ bounded and gauge invariant, the operations carry $\Phi\rho_0$ to

\begin{gather*}
\FitMath{\sum_h\tau_h[\Phi],} \\[3pt]
\FitMath{\tau_h[\Phi]=\int\Phi(U)\,\rho_0(U)\,\mathsf{K}_h(U,\cdot)\,dU,} \\[3pt]
\FitMath{\sum_h\tau_h[\Phi](\mathcal{V}_k)=\int\Phi\rho_0\,dU.}
\end{gather*}

For the choice $\Phi=e^{A(w,U)}$ one obtains $\rho^w_k\,dV_k=\sum_h\tau^w_h$, where the sum runs over the histories of $\rho_k$. Each $\tau^w_h$ is a complex measure which, as a function of $(c,z)$, is entire, and $\tau^0_h=\tau_h\ge0$. Moreover $\int\rho^w_k\,dV_k=Z(w)$ for $k\le K$.

Each step consists of finitely many moves. The first is the insertion of $1=\int dV\,\delta(\mathcal{M}(U)V^{-1})$. The second is a family of decompositions of unity $1=\sum\prod\chi\prod\chi^c$. The third is a gauge-fixing factor, of total integral $1$ over the gauge group; its insertion is harmless since $\Phi\rho_0$, like each $a_p$, is gauge invariant. The fourth is a set of changes of variables. The last occurs inside $\mathbf{R}$: a term is replaced by the product of a fixed probability law and its integral over a component. Every move is linear in the density and leaves the total integral unchanged, and no move requires positivity. The tests depend on the fields only, so the regions, labels and histories are the same as at $w=0$.

\textbf{Proposition 4.4 (Exact insertion of block observables).} Given a history $h$ of level $k$, denote by $Y(h)$ the union of all components that $\mathbf{R}$ has acted upon at levels not exceeding $k$. A bond $c$ of $\mathbb{T}^{(k)}$ is said to be far from $Y(h)$ when neither the two blocks $B^k(c_\pm)$ nor the bare bonds that join them meet $Y(h)^\sim$. There exists a probability measure $\mathsf{P}_k$ on triples $(U,h,V)$, each made of a bare field, a history and a block field, with four properties. First, its marginal on $U$ is $\mu_{K,\ell,\beta}$. Second, under its marginal on $(h,V)$ each bounded gauge-invariant $\Phi(h,V)$ has integral $Z(0)^{-1}\sum_h\int\Phi(h,V)\,\tau_h(dV)$. Third, it is invariant under $(U,h,V)\mapsto(U^u,h,V^{u|})$, where $u$ is a gauge transformation of the bare lattice and $u|$ is its restriction to the sites of $\mathbb{T}^{(k)}$. Fourth, under it $V=\mathcal{M}^kU$ holds on every bond far from $Y(h)$. Consequently, setting $F^{(h)}(V):=\mathbb{E}_{\mathsf{P}_k}[F(\mathcal{M}^kU)\mid h,V]$, which is the conditional expectation with respect to $\mathsf{P}_k$, taken in a version that is invariant in $V$, one has for each bounded gauge-invariant $F$ on $\mathcal{V}_k$

\begin{gather*}
\FitMath{\big\langle F(\mathcal{M}^kU)\big\rangle_{K,\ell,\beta}=\frac{\sum_h\int F^{(h)}\,d\tau_h}{\sum_h\int d\tau_h},} \\[3pt]
\FitMath{|F^{(h)}|\le\sup|F|,} \\[3pt]
\FitMath{F^{(h)}=F\,\text{ if }F\text{ reads only bonds far from }Y(h);}
\end{gather*}

in particular, $F^{(h)}=F$ for any history in which $\mathbf{R}$ has acted on none of the components. In precisely this sense $\rho_k$ represents the block averages of the bare field. The mechanism has three parts. First, $\mathbf{T}$ integrates $F\circ\mathcal{M}$ exactly against the density one level lower. Second, on each component on which it acts, $\mathbf{R}$ puts a reference law in place of the law of the block field; this alters the joint distribution of $\mathcal{M}^kU$ and $V_k$ only in a neighbourhood of that component, and it leaves every mass as it was. Third, the steps are glued in succession, each being regarded as a transport of the gauge-averaged terms. All the assertions are identities between measures and involve neither a constant nor a smallness condition. The role of the proposition is to fix what $\rho_k$ means; nothing is derived from it. In the two-point law (7b), a single scale is not isolated by this coupling, because, inside the region it heals, the law of the block field is replaced by $\mathbf{R}$. It is isolated by the exactness with the sources read at depth $s$, that is, by $\sum_h\tau^w_h(\mathcal{V}_{K-s})=Z(w)$ (Proposition 4.3, together with (4b)).

\textbf{Why the frozen choices still suffice.} Ba{\l}aban's theorems concern a single real coupling at each step (from these numbers the profiles build the function $g_k^2(x)$ of 119 (1988), (2.24), which coincides with $g_k^2$ on the last domain), a background that minimises the action appearing in the exponent, and a single coupling-constant counterterm in each step. For this reason the source is not built into that framework but is carried alongside it. In terms of the running shift $\zeta_k(\cdot;w)$ of (4b),

\begin{gather*}
\FitMath{-A\big(x_k(\cdot)-\zeta_k(\cdot;w),U\big)=-A\big(x_k(\cdot),U\big)+\sum_{y\in\mathbb{T}^{(k)}}A\big(h^{(k)}_y\zeta_k(\cdot;w),U\big),}
\end{gather*}

where $h^{(k)}_y$ denote Ba{\l}aban's tent functions (Comm. Math. Phys. 119 (1988), (3.40), (3.65)). On the set where $|\partial\mathbf{U}-\mathbf{1}|<\alpha_{0,k}\eta_k^2$ one has

\begin{gather*}
\FitMath{\big|A(h^{(k)}_y\zeta_k(\cdot;w),\mathbf{U})\big|\le C\alpha_{0,k}^2\sup|\zeta_k|\le C\varrho\,g_k^2(\log x_k)^{2q}\ll1:}
\end{gather*}

it is merely an additional local term, small and analytic, once $g$ is chosen small after $\varrho$. When the fields are real, $|e^{-A(x_k(\cdot)-\zeta_k,U)}|\le e^{-(1-C\varrho g^2)A(x_k(\cdot),U)}$; hence a failed test still costs $e^{-cp_0(g_k)^2}$.

\textbf{Proposition 4.5 (Reality and positivity at real sources).} All the choices are real, and so are the $f_i$; consequently each move takes real data to real data. Hence, at real $(c,z)$, each of the following quantities is real: the boundary terms, the increments, the activities $E^{(j),w},R^{(j),w}$ (at real fields), the vacuum constants, the running shifts, and also the $e_j(X;w)$. The exponents are therefore real at real block fields. Suppose moreover that the source is a real constant with $|\zeta|<2\varrho$. Then, in the operations $\mathbf{T}_k(Z)$ and $\mathbf{T}'_i(X)$ of (4b), every exponential factor is strictly positive at real pairs, and

\begin{gather*}
\FitMath{x_k(\cdot)-\zeta_k(\cdot;\zeta)>0,} \\[3pt]
\FitMath{\tau^\zeta_h\ge0:}
\end{gather*}

thus the stability bounds of the run rely on a positivity that is lost at complex $w$ and is still available at real shifts. Furthermore, $\rho^\zeta_0$ coincides with the bare density at $\beta-\zeta$, written with the choices belonging to $\beta$: these are two presentations of a single measure (3c). The choices of a reference run do not depend on $\beta$. Take a real $\beta$ whose run has room $2\varrho$ in the tube, and a real $\zeta$ with $|\zeta|<\varrho$; along such a run one then obtains $x^\natural_j(\beta-\zeta)=x^\natural_j(\beta)-\zeta_j(\zeta)$ and $b^\natural_j(\beta-\zeta)=b_j(\zeta)$. These facts enter the runs with real shifts in (3c) and in Step 8, as well as the monotone map that sends the bare coupling to the renormalised coupling (8b). The conjugated copies of (4d) and (4e) make no use of these facts; holomorphy is all they require (Lemma 4.11 [to be proved; Package D]).

\textbf{The scope of the mechanism.} Because $\mathcal O(f)=A(f,U)$ is a weighted action, its source is a coupling shift. A different bounded gauge-invariant local observable $\Phi$ need not be a weighted plaquette action. The exact insertion identity can transport $e^{z\Phi}$, but does not put it in the variable-Wilson format or supply its counterterms, source stability or uniform correlation estimates. Those require a separately closed source construction.


\paragraph*{Exact insertion of block observables: a coupling of bare field, history and block field}

\textbf{(i) The setting and the objects.}

\emph{Purpose of the block.} The effective density of level $k$ decomposes into a finite sum of positive measures indexed by the histories $h$ of level $k$: $\rho_k\,dV_k=\sum_h\tau_h$. A single application of $\mathbf T$ integrates $F\circ\mathcal M$ exactly against the full density one level lower, so after $k$ applications one has $\int F\,d(\mathbf T^k\rho_0)=\int F(\mathcal M^kU)\,\rho_0(U)\,dU$. In the case where $\mathbf R$ acts nowhere, $\rho_k\,dV_k$ is the image of $\rho_0\,dU$ under $\mathcal M^k$, and the assertion of this block then amounts to the definition of an image measure. Where $\mathbf R$ acts on a component, a reference law takes the place of the law of the block field inside that component. This modifies how $\mathcal M^kU$ and $V_k$ are jointly distributed on bonds close to that component, and no mass changes. What this block asserts is that the replacements fit together into a single probability measure, defined on triples (bare field, history, block field), with respect to which, away from the components acted upon, the block field coincides with the block average of the bare field. No small factor and no constant enters Statements 1 to 6.

\emph{Spaces.} Each $\mathcal V_j$ is a finite product of copies of the compact group $G$, and is therefore a compact metrisable space. Finite sets are given the discrete topology, and all measures are Borel measures.

\emph{Terms as the measures produced by moves.} In parts (4a) and (4c) the construction is regarded as a chain of moves, a move being a finite family of positive kernels with one kernel for each outcome. The outcomes of the moves of the steps $1,\dots,j$, listed in order, form a string, and such a string is called a history of level $j$. We write $\mathsf K_h(U,\cdot)$ for the kernel obtained by composing along the string $h$; it runs from $\mathcal V_0$ to the states reached after step $j$. The admissible histories of level $j$, namely the strings for which $\mathsf K_h\ne0$, form the finite set $\mathfrak H_j$, and $\mathfrak H_0=\{\varnothing\}$. The term of a history $h\in\mathfrak H_j$ is defined as the measure into which its string transforms the bare measure, taken as a measure in the block field. Its mass is $\mathsf t(h)$:

\begin{gather*}
\FitMath{\tau_h(\mathcal B):=\int\rho_0(U)\,\mathsf K_h\big(U,\{V_j\in\mathcal B\}\big)\,dU,} \\[3pt]
\FitMath{\tau_\varnothing=\rho_0\,dU,} \\[3pt]
\FitMath{\mathsf t(h):=\tau_h(\mathcal V_j).}
\end{gather*}

The statements and the proofs below use the terms only through this definition.

\emph{The two stages of a history.} Given $h'\in\mathfrak H_{j+1}$, we write $h'|_j\in\mathfrak H_j$ for the sub-list of those entries of $h'$ that come from the steps and from the large-field operations of the levels $\le j$. The strings at the stage reached once the moves of $\mathbf T$ of step $j+1$ have been made form a finite set, $\mathfrak H^{\mathbf T}_{j+1}$. Its elements are the lists $a=(h;P,Q,\Xi)$. Each is formed by some $h\in\mathfrak H_j$ together with the outcomes of those three moves of $\mathbf T$ which are built from tests:
\begin{itemize}
\item $P=P_{j+1}$, the outcome of the plaquette tests (at step $1$, the pair $(P_0,P_1)$);
\item $Q=Q_{j+1}$, the outcome of the approximate fluctuation tests;
\item $\Xi=\Xi_{j+1}$, the outcome of the fluctuation tests together with the recursion along the boundary, which determines $\Omega_{j+1},\Lambda_{j+1},S_{j+1}$.
\end{itemize}

For such an $a$ we set $a|_j:=h$. We let $\tau^{\mathbf T}_a$ be the measure produced by the string $a$, taken in the variable $V_{j+1}$, so that it is a measure on $\mathcal V_{j+1}$. The class of the large-field operation of level $j+1$ is denoted $\mathcal Q_{j+1}(a)$. It is determined by the entries of $a$ and consists of every component of $Z_{j+1}$ that satisfies the class conditions (part (2b)). The union of its members is

\begin{gather*}
\FitMath{\mathsf Z(a):=\bigcup_{X\in\mathcal Q_{j+1}(a)}X\subset T_\ell .}
\end{gather*}

The components that leave the class in the course of the operation (the drop-outs) are contained in $\mathsf Z(a)$, and so are the second-class components, although both kinds remain live. For $h'\in\mathfrak H_{j+1}$, the sub-list $h'^{\,\mathbf T}\in\mathfrak H^{\mathbf T}_{j+1}$ consists of $h'|_j$ together with the outcomes $P,Q,\Xi$ of step $j+1$. Every other entry of $h'$ is an outcome or a class produced within the large-field operation of level $j+1$.

\emph{The region of action.} Let $h\in\mathfrak H_j$. We denote by $Y(h)\subset T_\ell$ the union of those components of the large-field regions of $h$ which have been acted on by the large-field operation at one of the levels $\le j$; equivalently, it is the union of the class components of these levels. In particular $Y(\varnothing)=\emptyset$, and

\begin{gather*}
\FitMath{Y(h')=Y(h'|_j)\cup\mathsf Z\big(h'^{\,\mathbf T}\big)} \\[3pt]
\FitMath{(h'\in\mathfrak H_{j+1}).}
\end{gather*}


\emph{Clear and far bonds.} To a bond $c$ of $\mathbb T^{(j)}$ we attach the set of bare sites $\mathrm{Bl}_j(c):=B^j(c_-)\cup B^j(c_+)$, where $B^0(y)=\{y\}$. We say that the bond $c$ is clear of a set $Y$ when $\mathrm{Bl}_j(c)\cap Y=\emptyset$. A bond $c$ of $\mathbb T^{(k)}$ is called far from $Y(h)$ when neither $\mathrm{Bl}_k(c)$ nor any bare bond joining its two blocks meets $Y(h)^\sim$. This set is the union of the enlarged sets $X^\sim$, one for each of the class components $X$ whose union is $Y(h)$. The enlargement adds to such a component one layer of cubes of $\pi_{j(X)}$, and $j(X)$ is the level at which the operation acted on $X$.

\emph{Gauge averages.} The $G$-valued gauge transformations of $\mathbb T^{(j)}$ form the group $\mathcal G_j$. It is a finite product of copies of $G$ and carries the Haar probability measure $dv$. It acts through $V\mapsto V^v$, with $(V^v)^{v_0}=V^{v_0v}$. Given $u\in\mathcal G_i$ with $i\le j$, we write $u|_j\in\mathcal G_j$ for the gauge transformation obtained from it by restriction to the sites of $\mathbb T^{(j)}$. A function on $\mathcal V_j$, or a measure on it, is called invariant if this action leaves it invariant. The same convention applies on the configurations over a set of bonds closed under reversal, and on $\mathfrak H\times\mathcal V_j$, where the group acts trivially on the finite set $\mathfrak H$. Given a bounded Borel function $F$ and a finite positive measure $\tau$ on $\mathcal V_j$, we define

\begin{gather*}
\FitMath{\bar F(V):=\int F(V^v)\,dv,} \\[3pt]
\FitMath{\int F\,d\bar\tau:=\int\bar F\,d\tau .}
\end{gather*}


\textbf{(ii) The statements.} Throughout, $\beta$ is real, $0\le s\le K$ and $k=K-s$. The run is assumed to satisfy, through level $k$, the inductive description of the effective densities (part (2e)), with $x_0,\dots,x_k\ge\gamma^{-2}$, together with the two-sided flow bound (part (3a)). Of all these data, the proofs use only the fact that the moves of the steps $1,\dots,k$ are formed, that is, they use nothing beyond a reference sequence and a torus that is large compared with the cubes of the levels $\le k$.

\textbf{1. The step, taken term by term.} Let $0\le j<k$ and $h\in\mathfrak H_j$. Then each $\tau^{\mathbf T}_a$ with $a|_j=h$ is a finite non-negative measure on $\mathcal V_{j+1}$. Moreover, every bounded Borel invariant $F$ on $\mathcal V_{j+1}$ satisfies the following identity (the transport identity of the stage $\mathbf T$):

\begin{gather*}
\FitMath{\sum_{a\in\mathfrak H^{\mathbf T}_{j+1}:\,a|_j=h}\int F\,d\tau^{\mathbf T}_a=\int(F\circ\mathcal M)\,d\tau_h .}
\end{gather*}

Now let $a$ be such that $a|_j=h$, and let $F$ be a bounded Borel invariant function on $\mathcal V_{j+1}$ which does not depend on the variables $V(b)$ for which $b$ is a bond of $\mathbb T^{(j+1)}$ having at least one end point in $\mathsf Z(a)$. Then (the transport identity of the stage $\mathbf R$)

\begin{gather*}
\FitMath{\sum_{h'\in\mathfrak H_{j+1}:\,h'^{\,\mathbf T}=a}\int F\,d\tau_{h'}=\int F\,d\tau^{\mathbf T}_a .}
\end{gather*}

After summation over $h\in\mathfrak H_j$, the first identity expresses the exactness of $\mathbf T$ on the whole density. Taking $F\equiv1$ in both identities gives

\begin{gather*}
\FitMath{\sum_{h'\in\mathfrak H_{j+1}:\,h'|_j=h}\mathsf t(h')=\mathsf t(h):}
\end{gather*}

thus the step divides the mass of a history among its continuations. The law of the histories (part (4d)) is based on this identity.

\textbf{2. Gauge averages.} Suppose a compact group $\mathcal G$ with Haar probability $dv$ acts continuously on a compact metrisable space $\mathcal X$. Define $\bar F(x)=\int F(v\cdot x)\,dv$ and $\int F\,d\bar\tau=\int\bar F\,d\tau$.
(a) The function $\bar F$ is bounded, Borel and invariant, with $\sup|\bar F|\le\sup|F|$, and $\bar F=F$ when $F$ is invariant. The measure $\bar\tau$ is finite, positive and invariant. It satisfies $\int F\,d\bar\tau=\int F\,d\tau$ for invariant $F$, and $\bar\tau(N)=\tau(N)$ for invariant Borel sets $N$.
(b) Every invariant finite positive measure $\nu$ satisfies $\int\bar F\,d\nu=\int F\,d\nu$. Hence, if two invariant finite positive measures give the same integral to every bounded Borel invariant function, they are equal.
(c) Let $r:\mathcal G\to\mathcal G'$ be surjective, and let $q:\mathcal X\to\mathcal X'$ be a Borel map with $q(v\cdot x)=r(v)\cdot q(x)$. Then the image under $q$ of any invariant measure is invariant, and $f\circ q$ is invariant whenever $f$ is invariant.

\textbf{3. Gluing.} Let $\mathcal X_1,\mathcal X_2,\mathcal Y$ be compact metrisable spaces and let $p_i:\mathcal X_i\to\mathcal Y$ be Borel maps. Let $\lambda_i$ be finite positive measures on $\mathcal X_i$ such that $(p_1)_*\lambda_1=(p_2)_*\lambda_2$. Then $\mathcal X_1\times\mathcal X_2$ carries a finite positive measure $\lambda$ with the following two properties: its marginals are $\lambda_1$ and $\lambda_2$, and $p_1(x_1)=p_2(x_2)$ for $\lambda$-almost every $(x_1,x_2)$.

\textbf{4. Clear bonds.} (a) One has $c_-,c_+\in\mathrm{Bl}_j(c)$. (b) Let $c$ be a bond of $\mathbb T^{(j+1)}$, and let $b$ be a bond of $\mathbb T^{(j)}$ whose two end points lie in the double block $B(c_-)\cup B(c_+)$. Then $\mathrm{Bl}_j(b)\subset\mathrm{Bl}_{j+1}(c)$. (c) Every bond of $\mathbb T^{(k)}$ that is far from $Y(h)$ is clear of $Y(h)$.

\textbf{5. A single step as a transport of the averaged terms.} Let $0\le j<k$. Then $(\mathfrak H_j\times\mathcal V_j)\times(\mathfrak H_{j+1}\times\mathcal V_{j+1})$ carries a finite positive measure $\mathsf m_{j+1}$ with the following properties:
(a) the marginal on $\mathfrak H_j\times\mathcal V_j$ equals $\sum_{h\in\mathfrak H_j}\delta_h\otimes\bar\tau_h$;
(b) the marginal on $\mathfrak H_{j+1}\times\mathcal V_{j+1}$ equals $\sum_{h'\in\mathfrak H_{j+1}}\delta_{h'}\otimes\bar\tau_{h'}$;
(c) for $\mathsf m_{j+1}$-almost every $((h,V),(h',V'))$ both of the following hold: $h'|_j=h$, and $V'(c)=\mathcal M(V)(c)$ for every bond $c$ of $\mathbb T^{(j+1)}$ that has no end point in $\mathsf Z(h'^{\,\mathbf T})$.

\textbf{6. Exact insertion of block observables.} We write $Z=Z_{K,\ell,\beta}=Z(0)$ and $\mu=\mu_{K,\ell,\beta}$, and $\mathbb E_{\mathsf P}$ denotes the expectation with respect to a probability measure $\mathsf P$.
(a) On $\mathcal V_0\times\mathfrak H_k\times\mathcal V_k$ there exists a probability measure $\mathsf P$ with the following four properties.
Its marginal on $\mathcal V_0$ equals $\mu$.
Its marginal on $\mathfrak H_k\times\mathcal V_k$ is built from the averaged terms:

\begin{gather*}
\FitMath{\bar\nu_k:=Z^{-1}\sum_{h\in\mathfrak H_k}\delta_h\otimes\bar\tau_h .}
\end{gather*}

Hence this marginal assigns to every bounded Borel function $\Phi(h,V)$ that is invariant in $V$ the integral $Z^{-1}\sum_h\int\Phi(h,V)\,\tau_h(dV)$, and the law of the histories of level $k$ is the marginal of $\mathsf P$ on $\mathfrak H_k$: for $\mathsf E_k\subset\mathfrak H_k$, $\mathsf P(h\in\mathsf E_k)=Z^{-1}\sum_{h\in\mathsf E_k}\mathsf t(h)=\tilde\nu(\mathsf E_k)$.
It is invariant under the maps $(U,h,V)\mapsto(U^u,h,V^{u|_k})$, $u\in\mathcal G_0$.
Conditionally on $(U,h)$, the third coordinate agrees with $\mathcal M^kU$ on every bond that is clear of $Y(h)$, and in particular on every bond far from $Y(h)$.
(b) For each bounded Borel invariant function $F$ on $\mathcal V_k$, let

\begin{gather*}
\FitMath{F^{(h)}(V):=\mathbb E_{\mathsf P}\big[F(\mathcal M^kU)\,\big|\,h,V\big]}
\end{gather*}

be chosen among the versions of the conditional expectation as one that is invariant in $V$. Versions with this invariance exist, and for each $h$ any two of them coincide $\tau_h$-almost everywhere. Each of them satisfies

\begin{gather*}
\FitMath{\big\langle F(\mathcal M^kU)\big\rangle_{K,\ell,\beta}=\frac{\sum_h\int F^{(h)}\,d\tau_h}{\sum_h\int d\tau_h},} \\[3pt]
\FitMath{|F^{(h)}|\le\sup|F|\,\,\tau_h\text{-almost everywhere},} \\[3pt]
\FitMath{\sum_h\int d\tau_h=Z .}
\end{gather*}

(c) Suppose $F$ depends only on the variables of bonds that are clear of $Y(h)$; this is the case in particular when its dependence is confined to bonds far from $Y(h)$. Then one has $F^{(h)}=F$ almost everywhere with respect to $\tau_h$. In particular, for every history in which the large-field operation has acted on none of the components, $F^{(h)}=F$. Without this restriction one has, in general, $F^{(h)}\ne F$.
Here the depth $s$ is arbitrary, $0\le s\le K$, and $k=K-s$.

\textbf{7. Bounded block observables.} Take $F$ as in Statement 6 (b) and level-$k$ cells $\mathfrak c_1,\dots,\mathfrak c_n$, and denote by $\mathsf E_k:=\bigcup_{i\le n}\mathfrak{lv}_k(\mathfrak c_i)\subset\mathfrak H_k$ the event that at least one of these cells is contained in the live region of level $k$. Then

\begin{gather*}
\FitMath{\Big|\big\langle F(\mathcal M^kU)\big\rangle_{K,\ell,\beta}-\frac{\sum_{h\notin\mathsf E_k}\int F^{(h)}\,d\tau_h}{\sum_h\int d\tau_h}\Big|=\Big|\frac{\sum_{h\in\mathsf E_k}\int F^{(h)}\,d\tau_h}{\sum_h\int d\tau_h}\Big|} \\[3pt]
\FitMath{\le\sup|F|\;\tilde\nu(\mathsf E_k)\le n\,C\,\sup|F|\,e^{-c_{\rm lv}p_0(g_k)} .}
\end{gather*}

The final inequality is valid whenever the data satisfy the hypotheses of the one-cell bound of part (4d), with $C$ and $c_{\rm lv}>0$ the constants of that bound (written $C$ and $c$ in Theorem 4.13). It holds uniformly in $K$, $\beta$, $\ell$ and the positions of the cells. In this form the bound applies to an observable of bounded modulus that depends on the block variables of a bounded number of cells at its own scale. In the expectation of such an observable, the part contributed by the histories of $\mathsf E_k$ is smaller than any power of $g_k$.

\textbf{(iii) Proofs.}

\textbf{On 1.} The two transport identities and the specified
locality of each move are properties of the actual history operations
that we must prove: obligation C1 of Section~\ref{ch:dependency-catalogue}.
The following argument explains their consequences for specified
kernels and observables. It does not replace their verification for
the Yang--Mills construction by a reference to (4c).
Positivity. All kernels of all moves are non-negative. Hence the measure into which a string turns $\rho_0\,dU\ge0$ is finite and non-negative.
The stage $\mathbf T$. In its first move, $\mathbf T$ sets $V_{j+1}:=\mathcal M(V_j)$, and no later move of $\mathbf T$ alters $V_{j+1}$. The reason is that the Faddeev--Popov draw uses gauge transformations supported on the sites $B(y)\setminus\{y\}$, and these act on $V_j$ without touching $V_{j+1}$. Hence at this stage the identity is valid for every bounded Borel $F$. As an equality of measures, $\sum_{a|_j=h}\tau^{\mathbf T}_a=\mathcal M_*\tau_h$.
The stage $\mathbf R$. Given a class component $X$, denote by $\Lambda^{\rm pr}\cap X$ its box. This box lies well inside $X$, and $\mathbf R$ integrates the block field over it. Denote by $\mathfrak b_X$ the set of bonds of $\mathbb T^{(j+1)}$ that have an end point in that box. The moves of $\mathbf R$ alter $V_{j+1}$ on the bonds of the sets $\mathfrak b_X$ and nowhere else. Hence at this stage the identity is valid for every bounded Borel $F$ that does not depend on these bonds. Each bond of $\mathfrak b_X$ has both end points in $X$. Therefore every function that depends on no bond with an end point in $\mathsf Z(a)$ belongs to this class. The property that matters is that $F$ does not depend on $\mathfrak b_X$. Gauge invariance by itself is not enough: the plaquette function of (i) is a counterexample.
Why history by history. Each move, once summed over its outcomes, is a probability kernel at each state. For a move consisting of tests, exactly one outcome takes the value $1$. For the gauge fixing the property holds because the gauge fixing is carried out as a procedure, namely as a draw. Let $\mathfrak G\ge0$ be normalised by $\int\mathfrak G(V^u)\,du=1$. Then the measure

\begin{gather*}
\FitMath{(\tau\kappa_{\mathfrak G})(\mathcal B)=\int\tau(dV)\int du\,\mathfrak G(V^u)\,\mathbf 1_{\mathcal B}(V^u)}
\end{gather*}

has the same mass as $\tau$, for every measure $\tau$. The product $\mathfrak G\,\tau$, by contrast, has mass $\int\mathfrak G\,d\tau$. If $\tau$ is invariant under the group with respect to which $\mathfrak G$ is normalised, this coincides with the mass of $\tau$, and then $\tau\kappa_{\mathfrak G}=\mathfrak G\,\tau$. Suppose that the bare weight is invariant under gauge transformations and that the representatives transform covariantly. Then, history by history, the term on which a gauge-fixing move acts has this invariance (see the treatment of the steps as stochastic kernels at the beginning), and for such a term the draw coincides with the product. The invariance meant here is invariance under the group of the move. A term need not be invariant under the full gauge group of the level. For a measure without this invariance, replacing the draw by the product does not conserve the mass in general; this already happens in exactly computable models over finite groups. The masses considered from here on are the ones produced by the draw, which preserves them for any bounded Borel weight. Likewise, the axial frames are drawn before the tests that use them. The test that follows the gauge fixing along the tree lacks gauge invariance. Therefore the mass $a_r$ that passes it depends on the frame $r$. Suppose the components that leave kept a single frame $e$ while the survivors were averaged over the frames. Then the mass after the reference laws would be off by $\bar a-a_e$, with $\bar a$ the average of the $a_r$, and already for $F\equiv1$ the identity of the stage $\mathbf R$ would not hold. The actual history construction must verify both points: the frame draw precedes its tests, and the product formula represents that draw on the required invariant terms.
Why it does not suffice to have the identities in their form summed over the histories. Let $\mathcal V_0=\mathcal V_1=\mathcal V_2=\{0,1\}^2$, where bond $1$ is exterior and bond $2$ is the component. Let $\mathcal M=\mathrm{id}$, let the gauge group be trivial, and let $\mu$ be the uniform law on the diagonal. Level $1$ has two outcomes $a_1,a_2$, with $\tau^{\mathbf T}_{a_1}=\tau^{\mathbf T}_{a_2}=\mu/2$, and $\mathbf R$ acts on bond $2$ in $a_2$ alone. Thus $\tau_{a_1}=\mu/2$, while $\tau_{a_2}$ is the product of the marginal of $\mu/2$ on bond $1$ with the uniform law on bond $2$. At level $2$ no $\mathbf R$ acts, and every history has a single continuation, $a_i'$ for $a_i$. We set $\tau_{a_1'}:=\tau_{a_2}$, $\tau_{a_2'}:=\tau_{a_1}$. Every term is non-negative. The totals are transported exactly, and the masses are conserved history by history, $\mathsf t(a_i')=\mathsf t(a_i)$. Moreover, $\mathbf R$ replaces a term by the product of its exterior marginal with a probability law. On $\{h=a_1'\}$, however, $Y(a_1')=\emptyset$ forces $V=U$, which is a point of the diagonal. Yet $\tau_{a_1'}$ assigns mass $\frac14$ to the complement of the diagonal. Hence no $\mathsf P$ exists. The failure lies in the identity of the stage $\mathbf T$ applied to the one history $a_1$ by itself.
Below, the two identities are applied to invariant $F$ alone.

\textbf{Proof of 2.} (a) The map $(v,x)\mapsto F(v\cdot x)$ is bounded and Borel, so Fubini's theorem shows that $\bar F$ is Borel. Right invariance of Haar measure gives $\bar F(v_0\cdot x)=\int F(vv_0\cdot x)\,dv=\bar F(x)$. The functional $F\mapsto\int\bar F\,d\tau$ is positive and continuous under monotone limits. It therefore defines a measure $\bar\tau$. With $F_{v_0}(x)=F(v_0\cdot x)$, left invariance yields $\overline{F_{v_0}}=\bar F$, and this is the invariance of $\bar\tau$. The remaining assertions follow from $\bar F=F$ for invariant $F$, applied to indicators as well. (b) If $\nu$ is invariant, Fubini's theorem gives $\int\bar F\,d\nu=\int dv\int F(v\cdot x)\,\nu(dx)=\int F\,d\nu$. Hence $\nu$ is determined by the values it assigns to the integrals of the invariant functions $\bar F$. (c) This follows by a change of variables.
By this device, the proof converts the two identities of Statement 1 into equalities of measures, namely equalities between the gauge averages of the terms. The proof applies these identities to invariant functions only.

\textbf{Proof of 3.} We apply the disintegration theorem in the following form. Let $\mathcal X$, $\mathcal Y$ be compact metrisable, let $p:\mathcal X\to\mathcal Y$ be Borel, let $\lambda$ be a finite positive measure on $\mathcal X$, and let $\nu=p_*\lambda$. Then there exist probability measures $\kappa(y;\cdot)$ on $\mathcal X$, Borel in $y$, such that $\lambda(B)=\int\nu(dy)\,\kappa(y;B)$ and $\kappa(y;p^{-1}\{y\})=1$ for $\nu$-almost every $y$. Set $\nu=(p_i)_*\lambda_i$, and let $\kappa_i(y;\cdot)$ be a disintegration of this kind of $\lambda_i$ along $p_i$. Define

\begin{gather*}
\FitMath{\lambda(B):=\int\nu(dy)\,\big(\kappa_1(y;\cdot)\otimes\kappa_2(y;\cdot)\big)(B).}
\end{gather*}

The integrand is Borel in $y$, because the sets $B$ with this property form a Dynkin class that contains the products of Borel sets. We have $\lambda(B_1\times\mathcal X_2)=\int\nu(dy)\kappa_1(y;B_1)=\lambda_1(B_1)$, and the second marginal is treated in the same way. The set $\{p_1(x_1)=p_2(x_2)\}$ is Borel and contains $p_1^{-1}\{y\}\times p_2^{-1}\{y\}$. For $\nu$-almost every $y$, the latter set has full measure under $\kappa_1(y;\cdot)\otimes\kappa_2(y;\cdot)$.

\textbf{Proof of 4.} (a) Every site lies in its own block. (b) For a site $y$ of $\mathbb T^{(j+1)}$, its cube is the union, over the sites $x\in B(y)$ of $\mathbb T^{(j)}$, of their cubes. After intersection with the bare lattice this becomes $B^{j+1}(y)=\bigcup_{x\in B(y)}B^j(x)$. Consequently $B^j(b_\pm)\subset B^{j+1}(c_-)\cup B^{j+1}(c_+)$. (c) This follows from $Y(h)\subset Y(h)^\sim$.

\textbf{Proof of 5.} The measures $\tau_h$, $\tau_{h'}$, $\tau^{\mathbf T}_a$ are all finite and non-negative. Consequently the gauge average of each of them is a finite, positive and invariant measure.
Step 1 (the stage $\mathbf T$: Radon--Nikodym). Let $h\in\mathfrak H_j$ be fixed. The map $\mathcal M$ is covariant, $\mathcal M(V^v)=\mathcal M(V)^{v|_{j+1}}$, as well as Borel measurable, and $v\mapsto v|_{j+1}$ is a surjection of $\mathcal G_j$ onto $\mathcal G_{j+1}$. Therefore $\mathcal M_*\bar\tau_h$ is invariant, and $F\circ\mathcal M$ is invariant whenever $F$ is. Let $F$ be any bounded Borel function on $\mathcal V_{j+1}$. Applying the transport identity of the stage $\mathbf T$ to the invariant function $\bar F$, we obtain

\begin{gather*}
\FitMath{\sum_{a|_j=h}\int F\,d\bar\tau^{\mathbf T}_a=\sum_{a|_j=h}\int\bar F\,d\tau^{\mathbf T}_a=\int(\bar F\circ\mathcal M)\,d\tau_h=\int(\bar F\circ\mathcal M)\,d\bar\tau_h=\int\bar F\,d(\mathcal M_*\bar\tau_h)=\int F\,d(\mathcal M_*\bar\tau_h).}
\end{gather*}

Hence $\sum_{a|_j=h}\bar\tau^{\mathbf T}_a=\mathcal M_*\bar\tau_h$. In particular $0\le\bar\tau^{\mathbf T}_a\le\mathcal M_*\bar\tau_h$. Therefore $\bar\tau^{\mathbf T}_a=f_a\cdot\mathcal M_*\bar\tau_h$ for Borel functions $0\le f_a\le1$ with $\sum_{a|_j=h}f_a=1$ almost everywhere. Define $\mathsf m^{\mathbf T}_a$ as the image of $(f_a\circ\mathcal M)\,\bar\tau_h$ under $V\mapsto(V,\mathcal M(V))$. It is a measure on pairs $(V,W)\in\mathcal V_j\times\mathcal V_{j+1}$. The marginals of this measure are $(f_a\circ\mathcal M)\bar\tau_h$ and $\bar\tau^{\mathbf T}_a$, and almost surely $W=\mathcal M(V)$. What the identity transports is the sum over the outcomes and nothing more. It does not fix which portion of $\tau_h$ goes to which $a$, and the only splitting it supplies is the one by $f_a\circ\mathcal M$. Accordingly, under the coupling just constructed, the outcome $a$ is drawn with probability $f_a(\mathcal M(V))$, and it is not determined by the tests.
Step 2 (the stage $\mathbf R$: restriction off the class). Fix $a$. Let $\mathrm{Bd}(a)$ denote the set of bonds of $\mathbb T^{(j+1)}$ having no end point in $\mathsf Z(a)$. Let $\mathcal V(a)$ be the space of configurations on $\mathrm{Bd}(a)$, and let $\mathrm{res}_a:\mathcal V_{j+1}\to\mathcal V(a)$ be the restriction map, which commutes with the gauge action. Set $\lambda_a:=\sum_{h'^{\,\mathbf T}=a}\delta_{h'}\otimes\bar\tau_{h'}$. Let $f$ be a bounded Borel invariant function on $\mathcal V(a)$. Then $f\circ\mathrm{res}_a$ is invariant and depends on no bond with an end point in $\mathsf Z(a)$. The transport identity of the stage $\mathbf R$ therefore gives

\begin{gather*}
\FitMath{\sum_{h'^{\,\mathbf T}=a}\int(f\circ\mathrm{res}_a)\,d\bar\tau_{h'}=\sum_{h'^{\,\mathbf T}=a}\int(f\circ\mathrm{res}_a)\,d\tau_{h'}=\int(f\circ\mathrm{res}_a)\,d\tau^{\mathbf T}_a=\int(f\circ\mathrm{res}_a)\,d\bar\tau^{\mathbf T}_a .}
\end{gather*}

The measure $(\mathrm{res}_a)_*\bar\tau^{\mathbf T}_a$ and the image of $\lambda_a$ under $(h',V')\mapsto\mathrm{res}_a(V')$ are both invariant, so Statement 2 (b) shows that they are equal. The identity says nothing about the bonds with an end point in $\mathsf Z(a)$. For this reason the two measures are glued over $\mathcal V(a)$. The result is a measure $\mathsf m^{\mathbf R}_a$ on $\mathcal V_{j+1}\times(\mathfrak H_{j+1}\times\mathcal V_{j+1})$ with marginals $\bar\tau^{\mathbf T}_a$ and $\lambda_a$. Almost surely under it, $\mathrm{res}_a(V')=\mathrm{res}_a(W)$ and $h'^{\,\mathbf T}=a$.
Step 3 (composition). On the middle field $W$, the marginal of $\mathsf m^{\mathbf T}_a$ and the marginal of $\mathsf m^{\mathbf R}_a$ are both equal to $\bar\tau^{\mathbf T}_a$. We glue over $\mathcal V_{j+1}$ and discard the middle field. We then attach to $V$ the label $h=a|_j$ and sum over $a$. The resulting measure is $\mathsf m_{j+1}$. (a): $\sum_{a|_j=h}(f_a\circ\mathcal M)\bar\tau_h=\bar\tau_h$. (b): $\sum_a\lambda_a=\sum_{h'}\delta_{h'}\otimes\bar\tau_{h'}$, because each $h'$ has exactly one first stage. (c): almost surely $W=\mathcal M(V)$, $\mathrm{res}_a(V')=\mathrm{res}_a(W)$ and $h'^{\,\mathbf T}=a$, so that $h'|_j=a|_j=h$.

\textbf{Proof of 6.} For $0\le j\le k$ we define

\begin{gather*}
\FitMath{\mathsf A_j:=\big\{(U,h,V)\in\mathcal V_0\times\mathfrak H_j\times\mathcal V_j:\,V(b)=(\mathcal M^jU)(b)\text{ for every bond }b\text{ of }\mathbb T^{(j)}\text{ clear of }Y(h)\big\},}
\end{gather*}

which is a Borel set.
Step 1 (induction on the level). We show that, for each $0\le j\le k$, the space $\mathcal V_0\times\mathfrak H_j\times\mathcal V_j$ carries a probability measure $\mathsf P_j$ with the following three properties. Its projection to $\mathcal V_0$ is $\mu$. Its marginal on $\mathfrak H_j\times\mathcal V_j$ equals $Z^{-1}\sum_{h\in\mathfrak H_j}\delta_h\otimes\bar\tau_h$, which is the family of averaged terms. Finally, $\mathsf P_j(\mathsf A_j)=1$.
To begin the induction: the measure $\rho_0\,dU$ is invariant, so $\bar\tau_\varnothing=\tau_\varnothing$, and we take for $\mathsf P_0$ the push-forward of $\mu$ under $U\mapsto(U,\varnothing,U)$.
Passage from $j$ to $j+1$: by the second property together with Statement 5 (a), the measures $\mathsf P_j$ and $Z^{-1}\mathsf m_{j+1}$ induce the same marginal on $\mathfrak H_j\times\mathcal V_j$. We glue the two over this space and project the result onto $(U,h',V')$, which gives $\mathsf P_{j+1}$. The first two properties carry over, and Statement 5 (b) is what carries over the second. As for the third, with probability one we have $(U,h,V)\in\mathsf A_j$, $h'|_j=h$, and $V'(c)=\mathcal M(V)(c)$ outside $\mathsf Z(h'^{\,\mathbf T})$. Take a bond $c$ clear of $Y(h')=Y(h)\cup\mathsf Z(h'^{\,\mathbf T})$. Both of its end points belong to $\mathrm{Bl}_{j+1}(c)$, so neither lies in $\mathsf Z(h'^{\,\mathbf T})$, and therefore $V'(c)=\mathcal M(V)(c)$. For each bond $b$ of the double block of $c$ we have $\mathrm{Bl}_j(b)\subset\mathrm{Bl}_{j+1}(c)$. Hence such a bond is clear of $Y(h)$, and $V(b)=(\mathcal M^jU)(b)$. Because $\mathcal M(V)(c)$ depends on $V$ only through these bonds, we obtain

\begin{gather*}
\FitMath{V'(c)=\mathcal M(V)(c)=\mathcal M(\mathcal M^jU)(c)=(\mathcal M^{j+1}U)(c).}
\end{gather*}

Step 2 (symmetrisation). The group $\mathcal G_0$ acts on triples by $(U,h,V)\mapsto(U^u,h,V^{u|_k})$. Averaging $\mathsf P_k$ over this action gives a measure $\mathsf P$, and it is invariant under the action. Since $\mu$ and $\bar\nu_k=Z^{-1}\sum_h\delta_h\otimes\bar\tau_h$ are both invariant, $\mathsf P$ still has the first two properties. Covariance gives $\mathcal M^k(U^u)=(\mathcal M^kU)^{u|_k}$, so the condition $V(b)=(\mathcal M^kU)(b)$ holds exactly when $V^{u|_k}(b)=(\mathcal M^k(U^u))(b)$ does. Thus the set $\mathsf A_k$ is invariant, and $\mathsf P(\mathsf A_k)=\mathsf P_k(\mathsf A_k)=1$. Let $\Phi$ be a function and $N$ a Borel set, both invariant in $V$, with sections $N_h$. Then Statement 2 (a) gives

\begin{gather*}
\FitMath{\int\Phi\,d\bar\nu_k=Z^{-1}\sum_h\int\Phi(h,V)\,\tau_h(dV),} \\[3pt]
\FitMath{\bar\nu_k(N)=Z^{-1}\sum_h\tau_h(N_h).}
\end{gather*}

This is why the statement is restricted to invariant $\Phi$ and $F$. The marginal of $\mathsf P$ consists of the averaged terms, and the averaged terms coincide with the terms only on invariant functions and invariant sets.
Step 3 (assertion (a)). The marginals and the invariance were obtained in Step 2. A function of $h$ alone is invariant in $V$, and this yields the law of the histories. Because $\mathsf P(\mathsf A_k)=1$, for almost every $(U,h)$ a regular conditional distribution of $V$ given $(U,h)$ is concentrated on the configurations that agree with $\mathcal M^kU$ on every bond clear of $Y(h)$, and so in particular on every far bond.
Step 4 (invariant versions). Suppose that $F$ is invariant, and set $\mathsf f(U):=F(\mathcal M^kU)$. This function is invariant under $\mathcal G_0$. Take an arbitrary version $\Psi$ of $\mathbb E_{\mathsf P}[\mathsf f(U)|h,V]$ with $|\Psi|\le\sup|F|$, so that $\mathbb E_{\mathsf P}[\mathsf f(U)\phi(h,V)]=\int\Psi\phi\,d\bar\nu_k$ for every bounded Borel $\phi$. Each $v\in\mathcal G_k$ is of the form $u|_k$, where $u\in\mathcal G_0$ takes the value $v(y)$ on the block $B^k(y)$. Using the invariance of $\mathsf P$, of $\mathsf f$ and of $\bar\nu_k$, we get

\begin{gather*}
\FitMath{\mathbb E_{\mathsf P}[\mathsf f(U)\phi(h,V)]=\mathbb E_{\mathsf P}[\mathsf f(U^u)\phi(h,V^v)]=\mathbb E_{\mathsf P}[\mathsf f(U)\phi(h,V^v)]=\int\Psi(h,V)\phi(h,V^v)\,d\bar\nu_k=\int\Psi(h,V^{v^{-1}})\phi(h,V)\,d\bar\nu_k .}
\end{gather*}

The first equality holds because $\mathsf P$ is jointly invariant. We integrate over $v$, apply Fubini's theorem and use the fact that Haar measure is invariant under inversion. It follows that $F^{(h)}(V):=\int\Psi(h,V^v)\,dv$ is also a version, and that it is Borel, invariant in $V$ and bounded by $\sup|F|$. Any two invariant versions differ only on an invariant Borel set $N$ with $\bar\nu_k(N)=0$. Since the terms are non-negative, this gives $\tau_h(N_h)=0$ for every $h$.
Step 5 (assertion (b)). Take an invariant version. Then

\begin{gather*}
\FitMath{\big\langle F(\mathcal M^kU)\big\rangle=\int\mathsf f\,d\mu=\mathbb E_{\mathsf P}[\mathsf f(U)]=\int F^{(h)}(V)\,\bar\nu_k(dh,dV)=Z^{-1}\sum_h\int F^{(h)}\,d\tau_h .}
\end{gather*}

The last equality uses the invariance of $F^{(h)}$ in $V$. Taken by itself, the defining property of a conditional expectation yields the integral against $\bar\nu_k$ and nothing more. No use can be made at this point of a non-invariant version. Consider a term in which some component has been put into the second class by the large-field operation at level $k$. By part (4c), such a term carries, in the level-$k$ block field, a Dirac distribution on the gauge tree of the box. Consequently $\tau_h$ and $\bar\tau_h$ are mutually singular for such a term. A version is determined $\bar\tau_h$-almost everywhere, and in this case the integral against $\tau_h$ of a version that is not invariant is then arbitrary. Taking $\Phi\equiv1$ gives $\sum_h\int d\tau_h=Z$. For every $h$, any invariant version coincides $\tau_h$-almost everywhere with the version constructed in Step 4, and that version is bounded by $\sup|F|$.
Step 6 (assertion (c)). Fix $h_0$, and suppose that $F$ depends only on bonds clear of $Y(h_0)$. On $\mathsf A_k\cap\{h=h_0\}$ we then have $F(\mathcal M^kU)=F(V)$. It follows that 

\begin{gather*}
\FitMath{\mathbb E_{\mathsf P}[\mathsf f(U)\mathbf 1_{\{h=h_0\}}\phi(V)]=\mathbb E_{\mathsf P}[F(V)\mathbf 1_{\{h=h_0\}}\phi(V)]}
\end{gather*}

 for every bounded Borel $\phi$. In other words, $F^{(h_0)}=F$ $\bar\nu_k$-almost everywhere on $\{h_0\}\times\mathcal V_k$. The set on which this equality fails is invariant in $V$ and is therefore a $\tau_{h_0}$-null set. When the operation has acted on no component of $h_0$, we have $Y(h_0)=\emptyset$, and then every bond is clear.

\textbf{Proof of 7.} The equality comes from splitting the numerator of Statement 6 (b) into the part over $\mathsf E_k$ and the part over its complement. We have $\tau_h\ge0$, and $|F^{(h)}|\le\sup|F|$ holds $\tau_h$-almost everywhere. Hence $|\int F^{(h)}\,d\tau_h|\le\sup|F|\,\mathsf t(h)$. Summing this over $h\in\mathsf E_k$ and dividing by $Z=\sum_h\mathsf t(h)$ gives $\sup|F|\,\tilde\nu(\mathsf E_k)$. Subadditivity together with the one-cell bound $\tilde\nu(\mathfrak{lv}_k(\mathfrak c))\le C\,e^{-c_{\rm lv}p_0(g_k)}$ of part (4d) yields $\tilde\nu(\mathsf E_k)\le n\,C\,e^{-c_{\rm lv}p_0(g_k)}$. For a history outside $\mathsf E_k$, no one of the cells lies in a large-field region that is live at level $k$.

The proofs above use gauge averaging, the disintegration of measures and their gluing, the Radon--Nikodym theorem, an induction on the level, and the symmetrisation of a coupling.

\textbf{(iv) What is used.}

\emph{From other parts.} The following results of part (1a) are used. The map $\mathcal M$ is Borel. It is gauge covariant, $\mathcal M(V^v)=\mathcal M(V)^{v|}$. It is local: $\mathcal M(V)(c)$ depends on $V$ only through the bonds of the double block of $c$. From part (4c) we use three things. The first is the description of the construction as a chain of moves. The second is that every kernel is non-negative, so that every $\tau_h$ and every $\tau^{\mathbf T}_a$ is a finite non-negative measure. The third is the pair of transport identities of Statement 1. From part (2b) we use the class of the large-field operation, which enters through the recursion $Y(h')=Y(h'|_j)\cup\mathsf Z(h'^{\,\mathbf T})$. Statements 1 to 6 use nothing further from the construction: no bound, no analyticity, and no form of a term. Only Statement 7 draws in addition on part (4d), through the one-cell bound there on the law of the histories.

\emph{What reads it.} No argument in another part draws its conclusion from Statement 6. Within this part Statement 6 has a single consequence, Statement 7, and the proof of Statement 7 combines it with the one-cell bound of part (4d). Of the results of this block, the other parts rely only on the last line of Statement 1, $\sum_{h'|_j=h}\mathsf t(h')=\mathsf t(h)$. This identity makes the laws of the histories of the levels $j\le K$ a consistent family of probability measures. That family is the law $\tilde\nu$ of part (4d). The bounds on large fields rest on it, and through them so do the laws at one scale. For $k\ge1$ the observables of the theorem do not have the form $F(\mathcal M^kU)$. Indeed, $\mathcal O(f)=\sum_pf(x(p))\,a_p(U)$ is a function of the bare plaquettes, and Statement 6 cannot be applied to it. These observables are carried as sources (first block of this part). A single scale is isolated by exactness with sources at depth $s$:

\begin{gather*}
\FitMath{\sum_{h\in\mathfrak H_{K-s}}\tau^w_h(\mathcal V_{K-s})=Z(w),} \\[3pt]
\FitMath{\text{whence, at }c=0,} \\[3pt]
\FitMath{\Big\langle\exp\sum_iz_i\,\mathcal O(f_i)\Big\rangle_{K,\ell,\beta}=\frac{Z(w)}{Z(0)}}
\end{gather*}

and this quantity is an expectation, taken over the histories, the block fields and the stored variables at level $K-s$, of the exponential of the sourced exponent minus the unsourced one.

\emph{Against Ba{\l}aban's papers.} Ba{\l}aban presents the operation $\mathbf R$ in a simplified form in the Introduction of Comm. Math. Phys. 122 (1989) 175, formula (0.3) on p. 176: 
\begin{gather*}
\FitMath{(\mathbf R\rho)(V)=\sum_Z\rho(Z'',V)\int dV|_{Z'}\,\rho(Z,V)\big/\int dV|_{Z'}\,\rho(Z'',V).}
\end{gather*}
 This form has the normalisation property $\int dV\,(\mathbf R\rho)(V)=\int dV\,\rho(V)$ ((0.4)). In the procedure as actually carried out, the denominators integrate parts of the densities that are localised near $Z'$ (p. 177), and the normalisation property takes the form (1.102), p. 201. The objects treated in Ba{\l}aban's papers are densities. On observables, he states that his Theorem 1 ``is the basis for many other applications, for example, for an analysis of expectation values of physical observables, like loop variables, averaged loop variables'', and calls these problems that ``deserve detailed analysis and further publication'' (Comm. Math. Phys. 122 (1989) 355, p. 356).


\subsubsection*{(4b) Analytic dependence on the source; locality and decay in tame and in wild regions}

\textbf{Objects.} We write $w(p):=w(x(p))$. When the source is a constant shift $w\equiv\zeta$, the increments $b_j(\zeta)$ are computed on the infinite lattice in the same way as $b_j$. They therefore do not depend on the torus, and they are real for real $\zeta$. The running shift is $\zeta_k=\zeta+\sum_{j\le k}[b_j(\zeta)-b_j(0)]$, and the inverse couplings are $x_k-\zeta_k$. For a general source, let $p_y$ denote a bare plaquette at $y$. Then

\begin{gather*}
\FitMath{\zeta_k(p;w)=w(p)+\sum_{j\le k}\varphi_j(p)\sum_{y\in\mathbb{T}^{(j)}}h^{(j)}_y(p)\,\big[b_j(w(p_y))-b_j(0)\big]:}
\end{gather*}

in other words, the coupling counterterm of level $j$ is applied locally, and its coefficient is taken at the local value of the source.

\textbf{Proposition 4.6 (The format with sources).} The density $\rho^w_k$ has the form of (2e), with $x_k(\cdot)$ replaced by $x_k(\cdot)-\zeta_k(\cdot;w)$. Its activities $E^{(j),w},R^{(j),w},B^{(j),w}$ are holomorphic in $(c,z)$, depend on $w$ only through its values on $X^{\sim1}$, and obey the bounds valid at $w=0$ up to constants. Moreover $|\zeta_k(p;w)|\le\frac43\sup|w|$, and for every constant $|\zeta_*|<2\varrho$

\begin{gather*}
\FitMath{\big|E^{(j),w}(X,\cdot)-E^{(j),\zeta_*}(X,\cdot)\big|\le E_0\,\varrho^{-1}\mathsf{m}_j\,e^{-\kappa d_j(X)},} \\[3pt]
\FitMath{\mathsf{m}_j:=\max\Big\{\sup_{X^{\sim1}}|w-\zeta_*|,\,CML^{j-K}\operatorname{Lip}w\Big\}.}
\end{gather*}


\textbf{Strategy: induction on the levels, on a complex polydisc of data.} On a complex polydisc of data (activities, boundary terms, shifts) centred at the real run, a single step (the cluster expansion of (1d) together with $\mathbf{R}$ of (2b)--(2d)) is local, analytic and uniformly bounded. This is the step on the complex class of (2e). Its small-field core is Proposition 1.6 [to be proved; Package A--H] of (1d), which is the one-step theorem stated as a holomorphic map on a class of families of activities. The sourced data remain inside the polydisc, so holomorphy follows by composition. Locality follows because an activity at $X$ depends on the data, and on the outcomes of the history, only near $X$ (locality in the history, (2a)).

\emph{The marginal term.} When the source is constant, the step preserves the symmetries of (1e), so once $b_{k+1}(\zeta)A$ has been extracted only irrelevant terms remain. The estimate is the one of (3c). There is a remainder $\le Cg_k^{\kappa_0-4}$, coming from the $R^{(j)}$, which is holomorphic in $\zeta$ on a disc of radius $C\varrho$. Apart from this remainder, the increment $b_{k+1}$ is holomorphic in the shifts $\zeta_j$ of the earlier levels on discs of radius $cx_j$, where $c\le\sigma_\star$. These discs make up the polydisc of (1d), taken along the run's own couplings. On these discs its non-constant part, less the remainder, is bounded by $Cg_k(\log x_k)^{\mathfrak{p}_0}$. Cauchy's estimate together with $x_k\ge g^{-2}+c(K-k)$ yields

\begin{gather*}
\FitMath{|b_{k+1}(\zeta)-b_{k+1}(0)|\le\mathsf{a}_{[K-k-1]}|\zeta|,} \\[3pt]
\FitMath{\sum_{s<K}\mathsf{a}_{[s]}\le C\sum_k\big[g_k^3(\log x_k)^{\mathfrak{p}_0}+\varrho^{-1}g_k^{\kappa_0-4}\big]\le\tfrac13,}
\end{gather*}

where the quantity $\mathsf{a}_{[s]}$ carries a depth index and denotes, for the step that arrives at depth $s$, the Lipschitz number of its increment (the same letter is used in (3c) and (8b)). At zero source, reality and the displayed radial bound give $\zeta'_k(0)\in[\frac23,\frac43]$. Ordered derivative bounds on other admitted real constant-source intervals use the full slope estimate of (3c). On $X$ a general source is compared with $\zeta_*=w(p_*)$, where $p_*$ is a bare plaquette of $X^{\sim1}$. When $\mathsf{m}_j<\varrho$, the slope displayed above comes from the Schwarz lemma, applied along a line of data and not along a line of sources. This line is a family, holomorphic in $t$, of data of the level below; at $t=0$ it coincides with the data of $\zeta_*$, and at $t=1$ it coincides on $X^{\sim1}$ with those of $w$. For $|t|<\varrho/\mathsf{m}_j$ the family remains in the convex polydisc. When $\mathsf{m}_j\ge\varrho$, the slope is simply the sup bound. Subtract from each of the two activities the normaliser part $\mathbf{N}^{(j)}$ of (1d), which is common to both; what is left of each is at most $\frac12E_0e^{-\kappa d_j(X)}$. The slope is thus a clause that the induction carries along. The induction carries similar slopes for $R^{(j)}$ and $B^{(j)}_{\mathcal S}$, and also, among the operations treated below, for the excess exponents $\Psi_\omega$. The slope does not follow, up to the boundary of the class of sources $\{\|w\|_\flat<2\varrho\}$, from holomorphy in $w$, from the bound, from locality and from the agreement at constant sources. Moreover, $\mathsf{m}_j\le e^{d_j(X)}L^{-(K-j)}\|w\|_\flat$. The part that no counterterm renormalises therefore carries a factor $L^{-(K-j)}$, which is summable over the scales. In short, a gradient of the source is irrelevant.

\emph{Wild regions.} Two kinds of operations act on the sourced exponent. The first kind consists of the operations $\mathbf{T}_k(Z)$ on the components on which $\mathbf{R}$ has not acted. These are built from the choices of the run only, namely delta functions, tests, gauge-fixing factors and conditional Gaussian weights, and they contain no Wilson weight (2b). On such components the source appears in the function on which these operations act, that is, in the Wilson weight and in the boundary terms. Taken on its own, $\mathbf{T}_k(Z)1$ carries no small factor: for a component produced by failed plaquette tests, this quantity is $\le\exp[C\sum_{t\le k}\#_t(Z_t\cap Z)]$. Three places carry the small factor (2c): the Gaussian forms of the fluctuation fields, the kept gauge-fixing exponents and the Wilson weight. Only the last of these depends on the source; the first two are built from the choices of the run. The pointwise majorant of (2c) gives one factor $e^{-cp_0(g_t)^2}$ for each creation, at the coupling of that creation's own step, and one constant for each site. It is used in the form it has at $w=0$; the sourced exponent and the unsourced one are compared further below, on the supports of the tests.

The second kind concerns a component $X$ on which $\mathbf{R}$ has acted at some level $i$ (whether it was healed, remained as an island, or dropped out because a test of a stage or of the main step failed). For such a component the source enters inside its operation $\mathbf{T}'_i(X)$. Except for a drop-out, this operation contains the Wilson factor of the core. When the component is healed or is an island, the operation coincides with that of Ba{\l}aban, Comm. Math. Phys. 122 (1989) 355, (1.71). Let $\omega$ run over the local histories of $X$, let $\xi$ be the stored variables and $\mu_\omega\ge0$ their measures, and let $\mathbf{p}=(\mathbf{U},\mathbf{J})$ be the pairs on $X^\sim$ at which a term evaluates the factor of $X$ (2c). For each $\omega$, $e^{\Phi_\omega}$ is replaced by $e^{\Phi_\omega+\Psi_\omega}$, with an \emph{excess exponent}

\begin{gather*}
\FitMath{\Psi_\omega=A\big(\hat\zeta\,\phi_{\mathrm{c}}\mathbf{1}_X,U''(\xi)\big)+\big[V^w-V^0\big](X^\sim)+\sum_{\omega'}\Psi_{\omega'}:}
\end{gather*}

its summands are the Wilson term, weighted by the shift $\hat\zeta$ of the prepared regions ($|\hat\zeta|\le C\varrho$, and $0\le\phi_{\mathrm{c}}\le1$ is the cut-off on the old core, (2d)) and evaluated at the background $U''(\xi)$ of the inner region, which is constructed from the stored fields; the difference between the interaction terms $V$ of $\mathbf{R}$ with and without the source; and the exponents of the inner histories $\omega'$, namely those of the swallowed islands and drop-outs. The operation modified in this way is written $\mathbf{T}'_i(X;w)$. (In the case of a drop-out, the operation (2c) is taken to include the Wilson weight that the carrying term has over $X$. The operations that are kept are built from the choices of the run only, and none of their factors depends on the source, which appears in that Wilson weight and in the boundary terms. The exponent is then $\sum_{\omega'}\Psi_{\omega'}$ alone, the contribution of the swallowed islands, and it obeys the bound given below.) The exponent $\Psi_\omega$ is holomorphic in $(c,z)$, is real for real pairs and real sources, vanishes at $w=0$, and depends on $w$ only on $X^{\sim2}$. The background $U''$ is real, and its regularity holds only at the scale on which the respective layer was stored. Where $\phi_{\mathrm{c}}\ne0$, one has $|\partial U''-\mathbf{1}|\le C\check{R}_n^2\varepsilon_n\eta_n^2\le C\eta_n^2$ on $\Omega_n\setminus\Omega_{n+1}$, so the bound $a_p\le C\eta_n^4$ balances the $C\eta_n^{-4}$ plaquettes per site of level $n$ (on the bare layer $a_p\le2$, since $G$ is compact). Each layer therefore contributes $O(\varrho)$ for each site of its own scale:

\begin{gather*}
\FitMath{|\Psi_\omega|\le\bar{f}\,\#(\omega),} \\[3pt]
\FitMath{\bar{f}=C(1+\varrho),}
\end{gather*}

where, as in (2e), $\#(\omega)$ is the number of stored sites, each being counted at its own level. Since $\#(\omega)$ is additive over the inner histories, the constant after a step is the maximum of $\bar{f}$ and the constant of the new term, not their sum. Hence the constant does not increase with the level. The perturbation is extensive and must therefore be paid in advance. This is done by the mass bound of (2c), (2e), which the induction maintains with a reserve weight $e^{c\mathsf{v}(\omega)}$, where $c=c_{\mathrm{r}}$ is the reserve constant of (2e), one fixed number (the same constant stands in every weight $e^{c\mathsf{v}}$ of this part except the raised one in the bound of the whole label, whereas in the exponents that contain $p_0$ the letter $c$ denotes a generic constant). Here the space-time volume $\mathsf{v}(\omega)$ counts $\log x_j$ for each stored site of level $j$:

\begin{gather*}
\FitMath{\mathsf{M}(i,X):=\frac1{|W_4|}\sum_\omega e^{c\mathsf{v}(\omega)}\int\sup_{\mathbf{p}}\big|e^{\Phi_\omega(\xi;\mathbf{p})}\big|\,\mu_\omega(d\xi)\le e^{-(2-c')p_0(g_i)},}
\end{gather*}

with the supremum over the complex pairs taken inside the integral and the sum. This form applies to a healed component and to an island. In the case of a drop-out, the integral contains the Wilson weight of the carrying term over $X$. The supremum is then taken over the evaluations at $\xi$, that is, over the real configurations which are compatible with the stored variables, are regular layer by layer, and give the listed tests their recorded values, with the exception of the test whose failure made $X$ a drop-out; these evaluations form a set fibred over $\xi$ (2c). A product set contains a flat $\mathbf{U}$, at which $e^{-A}=1$, and the supremum over such a set removes the factor. Ba{\l}aban's bound is $\mathbf{T}'_k(X)1\le e^{-(2-c')p_0(g_k)}$ (ibid., (1.89), where Ba{\l}aban has $2(1+\beta_0)^{-1}$ in place of $2-c'$), for each fixed pair and without the weight. If $g$ is small enough once $\varrho$ has been fixed, then $c\log x_j\ge2\bar{f}$. Consequently $|e^{\Psi_\omega}|\le e^{c\mathsf{v}(\omega)}$, and for every bounded function $F$ of the stored variables and every evaluation of this kind

\begin{gather*}
\FitMath{\big|\mathbf{T}'_i(X;w)F\big|=\Big|\frac1{|W_4|}\sum_\omega\int F\,e^{\Phi_\omega+\Psi_\omega}\,d\mu_\omega\Big|\le\sup|F|\,\mathsf{M}(i,X)\le\sup|F|\,e^{-(2-c')p_0(g_i)}.}
\end{gather*}

Under a positive measure, the modulus of a complex factor serves as a bound for it; in the case of the small complex term produced by a complex pair, this is the device of Ba{\l}aban, Comm. Math. Phys. 122 (1989) 355, (1.73). The estimate just displayed is made on the moduli, and a sourced component retains the small factor of the unsourced one. That $\mathbf{T}'_i(X;w)F$ is holomorphic follows from dominated convergence together with the theorems of Morera and Fubini. For a healed component, the exponentiation of (2d) produces its local terms from this bound, each at most $e^{-p_0(g_i)}e^{-\kappa d_i(Y)}$ on a domain $Y$, just as at $w=0$.

\emph{The whole label.} Bounds on the individual components cannot be multiplied together, even though each is uniform in the variables outside its own component, because the background couples the components. (For $\mathbf{1}\{|\xi_1-\xi_2|<\epsilon\}$ on $[0,T]^2$, every integral in one variable is $\le2\epsilon$, while the double integral equals $2\epsilon T-\epsilon^2$.) By Tonelli, what does multiply is a product of functions of disjoint groups of variables; here the factors of this one product are the pieces of the pointwise majorant together with the integrands of the masses. For this reason the supremum is placed inside the integral. The result is the integrated bound of (2c), with one factor for each creation and one mass for each outermost component on which $\mathbf{R}$ has acted. If every $e^{\Phi_\omega}$ in $\mathbf{T}'_k(\mathcal{S})$ is replaced by $e^{c\mathsf{v}(\omega)}|e^{\Phi_\omega}|$, then

\begin{gather*}
\FitMath{\int\mathbf{T}'_k(\mathcal{S})\Big[\chi_{k,\mathcal{S}}\,e^{-A(x_k(\cdot),U_k(\mathbf{V}))}\Big]dV_k\le\exp\Big[C\sum_{j\le k}(1+\log x_j)\,\#_j(Z_j)\Big]\prod_{(t,\mathfrak{c})}e^{-cp_0(g_t)}\prod_{(i,X)}\mathsf{M}(i,X),}
\end{gather*}

where the first product runs over those creation cells, outside these components, that belong to a single history with the label $\mathcal{S}$, namely the history for which this product is largest (the creation cells are not determined by the label, (2c)). The second runs only over the outermost ones among the components on which $\mathbf{R}$ has acted and which remain in the label. (In (2c) the bound carries $p_0(g_t)^2$ and has no logarithm. The sum over the labels needs only the present weaker form, because $CM^4\check{R}_t^5\log x_t\le\tfrac12c\,p_0(g_t)$ (Proposition 2.13 [outlined below; Section~\ref{sec:actual-history-closure}]): half of the exponent pays the cost and the other half is kept. Inside a drop-out the mass has consumed the factors of the creations, and whatever the label still shows there is paid for by the weight. The masses are required at a weight that exceeds the $e^{c\mathsf{v}(\omega)}$ spent by the sources by a constant; this constant is determined by the per-site cost of the part of the label visible inside the drop-out, and its choice comes after that of $\varrho$ and before the ceiling on $g$ is set (the mass bound is valid at every weight $c\ge0$, (2c), and the masses $\mathsf{M}(i,X)$ in the display are understood at the raised weight). The surplus comes from that part of $\mathsf{M}(i,X)$ which is made up of the local histories admitted by the label; on these, $\mathsf{v}(\omega)$ is bounded below by the volume shown (2c), (2e).) Once its vacuum constant is set aside, the sourced term of the label is dominated by the left-hand side multiplied by $e^{C(1+\varrho)}$ for each bond of the determining set, for three reasons. The moduli are taken inside the integrals. In the operations, $|e^{\Psi_\omega}|\le e^{c\mathsf{v}(\omega)}$. On the supports of the tests, what remains of the sourced exponent once the Wilson term $-A(x_k(\cdot),U_k)$ and the vacuum constant are set aside is at most $C(1+\varrho)$ for each bond of the determining set (4c). Consider the weighted sum over the labels of (2e), that is, the total masses summed over $\mathcal{S}$, where the credit pays the costs. In it every component $C$ of $Z_k$ keeps a factor

\begin{gather*}
\FitMath{\exp\big[-cp_0(g_k)-\kappa_1|C|_k\big],}
\end{gather*}

where $|C|_k$ is measured in cells. In (2c) the factor is taken per site, with $\#_k(C)\ge d_k(C)$ in place of $|C|_k$. This contains the form ``with the additional term $-\kappa_1d_k(X)$'', which Ba{\l}aban gives for a general region (ibid., p. 387).

\emph{The cost of the reserve.} A factor $e^{c\log x_j}$ for each site amounts to at most $e^{cM^4\check{R}_j^5}$ for each cell. Its form is that of the volume costs $CM^4\check{R}_n^5$ per cell and per step, which, until $\mathbf{R}$ acts on a component, are covered by the credit $\kappa_j(Z)$ of that component: an amount $p_0(g_j)^2$ for each cube, deposited at its creation, pays for $\check{R}_j^6$ for each cube (ibid., (1.79)--(1.81) and (1.89); these bound $\mathbf{T}'_k(X)1$ for the operation of (1.71), which contains the Wilson factor). Ba{\l}aban already pays $O(1)\log g_j^{-2}$ for each site of $Z_j$, on account of the normalisation constants (ibid., p. 380). The spare weight, whose constant $c$ is free, holds once $g$ is small enough, the constant $c$ being chosen first (2c).

\textbf{Proposition 4.7 (Sources along a frame).} Let $(\theta_j)_{j\le n}$ be a frame such that $\theta_j\ge\gamma_J^{-2}$ throughout, and let $(x_j)$ be the run along it starting from a real $\beta$, as in the opening of Step 3. The cube sizes are $\check R_j=R(\min_{i\le j}\theta_i)$. All the other auxiliary choices of level $j$ (the reference laws, the thresholds, the gauge-fixing weight and the scaling of the fluctuation field) are made with $\theta_j$, whereas the weight of the action obeys its own recursion. The reference line is a frame. A first-passage run, that is, the run starting from a first-passage datum $(K,\beta)$, is the run taken along its own couplings. The run is said to stay in the tube if $|\theta_j-x_j|\le\mathsf{t}_j$ for all $j$, and to have room $r$ if $|\theta_j-x_j|+r\le\mathsf{t}_j$. Suppose that the run has room $2r$ for $r=\varrho$ or for $r=2\varrho$. Then, for all sources with $\|w\|_\flat<r$ and for the steps of the frame, Propositions 4.3, 4.5, 4.6 and Theorem 4.1 are valid, with upper stability and the floor holding with $C(1+\log\theta_n)$ for each site of $\mathbb{T}^{(n)}$. They hold on every admissible pair $(\ell,K)$ of a torus and a cutoff, admissibility meaning $K\ge n$ and $\ell L^{K-i}\ge CM\check{R}_i$ for $i\le n$, as in (3b); all of this is uniform in $\ell$, in $K$, in $n$, in $\beta$ and in the frame. Moreover

\begin{gather*}
\FitMath{\mathsf{z}_j(\zeta):=(\theta_j-x_j)+\zeta_j(\zeta)\in D(0,\mathsf{t}_j),} \\[3pt]
\FitMath{\sum_{j\le n}\,\sup_{0<|\zeta|<r}\frac{|b_j(\zeta)-b_j(0)|}{|\zeta|}\le\tfrac13.}
\end{gather*}

The mechanism is as follows. A step depends separately on the choice $\theta_j$ and on the weight $x_j-\zeta_j$. For given remainders, its small-field part is holomorphic in their difference $\mathsf{z}_j(\zeta)$ on $D(0,\sigma_\star\theta_j)\supset D(0,\mathsf{t}_j)$ (Proposition 1.6 [to be proved; Package A--H]); this difference is the deviation $\mathsf{z}_j=\theta_j-x_j$ displaced by the running shift. The full step is formed at a real deviation in the tube, together with a complex shift about it of modulus at most $\bar c$ (the complex class of (2e)). Consequently, as far as the increments are concerned and with the remainders $R^{(i)}$ held fixed, a deviation from the frame and a constant source are the same variable. The remainders depend on the source as well, and the activities in the two presentations of Proposition 4.5 need not coincide term by term. The room is what keeps the complex excursion $|\zeta_j|\le\frac43|\zeta|$ inside the disc. The objects constructed do not depend on the tube. A single statement therefore serves the first-passage runs, the comparison of two cutoffs in (8a), and the runs with frozen choices that occur in (3c) and (5a); it is used with $r=\varrho$, where the room is $2\varrho$, and with $r=2\varrho$, where the room is $4\varrho$ (which is the room $C\varrho$ of Theorem 4.1).

\emph{Reading at depth $s$.} For the run stopped at level $K-s$, everything above holds with the following changes. The norm $\|f\|_{\sharp,s}=\max\{\|f\|_\infty,CML^{-s}\operatorname{Lip}f\}$ replaces $\|f\|_\sharp$, the Lipschitz constant being measured at scale $L^{-s}$. The sums are restricted to $j\le K-s$. Either normaliser may be used, each with exponents of its own: $E$ of the whole run, or the vacuum constant $E_{(K-s)}$ of the first $K-s$ steps. The number of sites of $\mathbb{T}^{(K-s)}$ replaces $\ell^4$, and $L^{-(K-s-j)}$ replaces $L^{-(K-j)}$. The exponents of the two global bounds also depend on $s$ (7b).


\subsubsection*{(4c) Finite ordinary sources and the positive-label floor}
\label{sec:actual-ordinary-source-proof}

\paragraph*{A finite source neighbourhood without loss of radius.}
Fix the choices of the real run and put $r=2\varrho$. At depth $s$ use
\[
 \|f\|_{\sharp,s}=\max\{\|f\|_\infty,
 CM L^{-s}\operatorname{Lip}f\},\qquad
 \|w\|_{\flat,s}=|c|+\sum_i|z_i|\|f_i\|_{\sharp,s}.
\]
The constant-source disc is doubled. Constant sources are constructed
first. The increment is holomorphic on the coupling polydisc of radii
proportional to $\theta_j$. A single Schwarz estimate moving every
coordinate simultaneously bounds its slope by the maximum of the
coordinate displacements divided by these radii. Thus the slope gains
$\theta_j^{-1}$; it does not contain a sum of reciprocals over all
earlier levels. The lower flow bound makes the resulting slopes summable.
For a holomorphic increment put
\[
 \mathsf d_\rho[b]
 =\sup_{0<|\zeta|<\rho}\frac{|b(\zeta)-b(0)|}{|\zeta|}.
\]
For a sufficiently small coupling,
\[
 \sum_j\mathsf d_{2r}[b_j]\le\tfrac13,
 \qquad |\zeta_j(\zeta)|\le\tfrac43|\zeta|.
\]
This leaves the same fixed margin at every level.

For a general source, freeze its value near a base point. The difference
between the actual regular activity and the one at the frozen constant
is a defect with no coupling counterterm. At an old level $j\le m$ it
has size $L^{-(m-j)}$. This gain must hold on the whole source class.
Holomorphy on that class alone would give a derivative estimate only
after shrinking the class, which cannot be repeated indefinitely.

Instead interpolate the data of a step. We specify the distance
used for this interpolation. Identify a region $Y$ with the bare
plaquettes having a corner in it, and use the torus distance between
their barycentres to set
\[
 \operatorname{Lip}(Y;w)
 =\sup_{\substack{p,p'\in Y\\x(p)\ne x(p')}}
       \frac{|w(p)-w(p')|}{d_{T_\ell}(x(p),x(p'))},\qquad
 \mathsf l_j(Y;w)=CM L^{-(K-j)}\operatorname{Lip}(Y;w),
\]
\[
 \mathsf m_j(Y;w,\zeta_0)
 =\max\left\{\sup_{p\in Y}|w(p)-\zeta_0|,
                     \mathsf l_j(Y;w)\right\}.
\]
An empty supremum is zero. These quantities are monotone in $Y$,
and $\mathsf m_j\le\mathsf m_{j+1}$ on a fixed region.

For a union $W$ of cubes of the output partition of step $k+1$,
take $\|w\|_{\flat,K-k-1}<r$, fix one constant
$\zeta_0\in D(0,r)$, and put
$\mathsf m=\mathsf m_{k+1}(W;w,\zeta_0)$.
Keep the unsourced tests, backgrounds, coordinate maps and reference
measures fixed. At sites in $W$, interpolate the local weight and each
Euclidean member, together with its own coefficient, from the data of
$\zeta_0$ to those of $w$. Interpolate defects, remainders and boundary
terms only when their complete read regions lie in $W$; treat a
boundary sum site by site. Outside $W$ use the constant-source data.
Write $\mathfrak L_t^W$ for this line of data. Its complex parameter
satisfies $|t|\mathsf m<r$. It need not be the system of one physical
source function. If $\mathsf m=0$, locality gives the desired zero
differences directly; the disc below is used for $\mathsf m>0$.

The input slopes admit the entire line to the larger data box.
Its weight size is at most $3r$ and its regular norm at most $E_0+1$.
For every older read $Y\subset W$ of level $j\le k+1$,
\[
 \mathsf l_j(Y;w)
 \le L^{-(k+1-j)}\mathsf l_{k+1}(W;w)
 \le L^{-(k+1-j)}\mathsf m.
\]
Together with vanishing at the constant source, this admits the defects
with their required older-level factors. The other members retain
their own majorants. The one-step formation theorem is used on this
whole box and on the original common field cores.

For the regular output indexed by $X$, take $W=X^\sim$.
This enlargement contains its complete input read region, including
the finer read enlargements, so locality identifies the line's
endpoints with the two source outputs below when $t=1$ is admitted.
Subtract the common Gaussian normalizer from the regular output.
The cluster contribution and the thin-tube weight discrepancy bound
what remains by $\frac12e^{-\kappa'd(X)}$, with
$\kappa'\ge\kappa+1$, throughout $|t|<r/\mathsf m$.
Schwarz's lemma therefore gives
\[
 |E^{(k+1),w}(X,\cdot,z)-E^{(k+1),\zeta_0}(X,\cdot,z)|
 \le\frac{\mathsf m}{r}e^{-(\kappa+1)d_{k+1}(X)}.
\]
For $r/\mathsf m\le1$, use the two ordinary endpoint bounds instead
of evaluating the data line at $t=1$. The factor $1/2$ pays for the
factor two in Schwarz's estimate and returns the defect constant $1$.
Taking $\zeta_0=w(p_z)$ and using the spare decay to sum the diameter
of $X$ gives
\[
 |D^{(j)}(X,\cdot,z)|\le L^{-(m-j)}e^{-\kappa d_j(X)}
\]
at the terminal level $m$. The slopes of the other regular activities
and boundary terms follow with their own majorants. The same
$\zeta_0$ is used for every object in the enclosing region.

For an operation cut by $W$, its exponent on $\mathfrak L_t^W$ is
formed by the earlier steps on that line. Affine interpolation of a
partial sum of its factors is not used. This recursion is well-defined:
an output partition region is also a region for the preceding partition,
the next source class is contained in the preceding class, and
$\mathsf m_k\le\mathsf m_{k+1}$. A line needed at level $k+1$ has
therefore already been admitted at level $k$. At $t=0$ the formed
exponents are those of the constant source. At $t=1$, when this point
is admitted, a factor satisfies $\psi_f(1)=\psi_f^w$ if
$R_f^\sim\subset W$; agreement of the scalar part
$\Psi^\circ(1)=\Psi^{\circ,w}$ requires $X^{\sim2}\subset W$.
Agreement on one carrier does not identify the whole scalar exponent.

Let $N_0^\circ(\omega)$ count the operation's own sites, at their
respective levels, and define its total source size recursively by
\[
 N_0^\circ(\omega)=\sum_j\#_j(Z_j(\omega)\cap X),\qquad
 \mathsf n(\omega)=N_0^\circ(\omega)
       +\sum_{\omega'\in\operatorname{In}(\omega)}\mathsf n(\omega').
\]
Each inner history enters once. At a new completed operation,
\[
 \begin{split}
 \Psi_\omega(t)&=\Psi^{\rm new}_\omega(t)
                      +\sum_{\omega'}\Psi_{\omega'}(t),\\
 \Psi^{\rm new}_\omega(t)
 &=A(\hat c(t)\upsilon\mathbf1_X,U''(\xi))
       +V(X^\sim;\mathfrak L_t^W)-V(X^\sim;\mathfrak L^0).
 \end{split}
\]
Here $\mathfrak L^0$ is the fixed unsourced system,
$|\hat c(t)|\le3r$, $0\le\upsilon\le1$, and $U''(\xi)$ is real.
Its layer bounds cost $3C_Wr$ per newly owned site for the Wilson
term; the bare layer is counted in bare sites. The two collar
potentials cost $2c_{\rm int}(1+3r)$ per such site. Put
\[
 c_{\rm new}=3C_Wr+2c_{\rm int}(1+3r),\qquad
 c_\Phi=2c_{\rm new}.
\]
Additivity of $\mathsf n$ returns the ordinary exponent constant as
$\max(c_{\rm new},c_\Phi)=c_\Phi$ and its formed-line bound as
$\max(c_{\rm new},2c_\Phi)=2c_\Phi$.
On $W=X^{\sim2}$, Schwarz gives slope
$2c_{\rm new}\mathsf m/r=c_\Phi\mathsf m/r$ per newly owned site.
Each inner exponent retains its previous slope toward this same
$\zeta_0$, and its local distance is at most the enclosing distance.
The size recursion thus returns $c_\Phi$, with no new price per
ancestor. A drop-out has only the sum of its inner excesses.

For the carrier oscillations this argument must be made in the sum
norm. For a carrier $R$ born at level $i$, put
\[
 M_R=(\log\theta_i)^{2\mathfrak p_0}
          e^{-\kappa d_i^\natural(\mathcal S;R)},\qquad
 \operatorname{osc}_{\mathbf p}h
       =\sup_{\mathbf p,\mathbf p'}|h(\mathbf p)-h(\mathbf p')|.
\]
At fixed stored variables the local formation bound is
\[
 \sum_{f:R_f=R}\operatorname{osc}_{\mathbf p}\psi_f(t)
       \le c_\times(3r)M_R,
       \qquad |t|\mathsf m<r.
\]
The pair cores in this bound are fixed common admitted cores for the
entire line. A carried mark retains the bound; at a new mark the same
one-step estimate returns it.

Apply the data-line construction at the carrier's formation level
$i$, with $W=R^\sim$ and now
$\mathsf m=\mathsf m_i(R^\sim;w,\zeta_0)$.
The case $\mathsf m=0$ again follows by locality.
First suppose $T=r/\mathsf m>1$.
For finitely many factors choose fixed admitted pairs
$\mathbf p_f,\mathbf p'_f$ and fixed unimodular numbers $\sigma_f$
so that the increments between zero and one of
$\sigma_f[\psi_f(t;\mathbf p_f)-\psi_f(t;\mathbf p'_f)]$
are nonnegative real numbers. Their sum $H(t)$ is holomorphic and
bounded by $c_\times(3r)M_R$ on $D(0,T)$. Schwarz gives
$|H(1)-H(0)|\le2c_\times(3r)M_R/T$.
Take the pair suprema after this estimate, then increase the finite
factor sets. It follows that
\[
 \sum_{f:R_f=R}\operatorname{osc}_{\mathbf p}
          [\psi_f^w-\psi_f^{\zeta_0}]
 \le2c_\times(3r)
     \frac{\mathsf m_i(R^\sim;w,\zeta_0)}r M_R.
\]
For $T\le1$ the two ordinary oscillation bounds give the same result.
The phases test already formed functions of one parameter; they do
not introduce independent carrier coefficients or local sources.
This proves the slope bound for the sum of suprema without treating
that sum itself as holomorphic.

The formed-line bounds propagate with the same ordinary radius $r$.
Their inputs are the actual one-step formation estimate, hereditary
core inclusions and collar-potential bound used above. The mass
reserve absorbs the resulting exponent; it does not prove those
formation estimates. This argument does not assert the stronger
weighted discs for independently marked zone uses.

The coupling ceiling is chosen after $r$ and these exponent constants.
On each component, $|e^{\Psi_\omega}|\le e^{C(r)\mathsf n(\omega)}$
is absorbed by its weighted unsourced mass. All estimates are made in
total variation before the operations are composed. With
$c_{\rm L}=C(r)$ this gives the actual term comparison
\[
 |e^{\mathcal E_k(\mathcal S;w)}\tau^w_{\mathcal S}|(\mathcal V_k)
 \le e^{\mathcal E_k(\mathcal S)}
 \tau^{\rm abs}_{\mathcal S,c_{\rm L}}(\mathcal V_k)
 e^{c_{\rm L}\sum_l|\Gamma_l|}.
\]
The weighted label sum applies at the raised mass weight required by
$c_{\rm L}$. The source-energy share missing from a label is bounded by
$\sum_j(C\log\theta_{j-1}+C)\#_jZ_j^+$. The same weight absorbs it.
Thus
\[
 \left|e^{-E_{(k)}-\sum_{j\le k,X}e_j(X;w)}Z(w)\right|
 \le e^{[C\log\theta_k+C(r)]N_k}.
\]
No source-dependent denominator has entered this argument.

At zero source all terms are positive measures. Keep the never-large
label and restrict every last-level bond to a ball of radius
$s=c\theta_k^{-1/2}$ about $\mathbf1$. To check the tests, lift the
last block field to each earlier lattice by taking the field equal
to $\mathbf1$ inside its blocks and putting the block bonds on the
crossings. A lifted plaquette is $\mathbf1$, a cancelling product
$VV^{-1}$, or a plaquette of $V$; its curvature is at most $4s$.
The local variational background estimate on each test cube then
gives the required small-field tests and
$A(U_k)\le Cs^2N_k$. Thus $x_kA(U_k)\le CN_k$ on this ball.
If $D=\dim\mathrm{SU}(2)=3$, its product
Haar measure costs at most
$\frac12dD\log\theta_k+C$ per site. Hence
\[
 e^{-E_{(k)}}Z(0)
 \ge \exp[-(C+\tfrac12dD\log\theta_k)N_k].
\]
At $k=K$, $N_K=\ell^4$ and first passage confines $\theta_K$ to a
fixed window above $g^{-2}$. The upper bound and the floor therefore
have constants depending on $g$ alone. They verify the two history
criteria and the positive-label criterion in
Section~\ref{ch:ordinary-history-stability} for this construction.

All formations read a source only in their enlarged domain and begin
with the function $w(p)$ at bare plaquettes. The resulting energies
therefore depend on that function, independently of its presentation as
a finite linear combination. Their exponential rate exceeds the domain
entropy. At a chessboard continuation the copies vanish near cell
faces; joining them increases neither their sup norm nor their
Lipschitz constant. Splitting even and odd copies costs only the fixed
factor two already allowed by the doubled source radius.

The energies just formed also have the rooted geometry used in
Section~\ref{ch:graded-rooted-source}. The sum over a regular
point $z\in X$ and the remainder at $X$ use the same domains of
$\pi_j$; summing histories introduces no new independent domain
index. At depth $s=K-j$ these are connected unions of cubes of
side $Ma_s$. With $M=L^{m_0}$ the number of such cubes along a
torus side is $2L^{m-m_0+s}$, hence even on the admitted tori.
Their complete source-dependence region is a fixed-layer
enlargement. For a domain of cubical edge-tree length $m(X)$,
the sum over $z$ costs at most $C_M(m(X)+1)$. Combining the fixed
origin groups and using the strict entropy-rate surplus absorbs
this polynomial into the exponential bound. The rooted energy
majorant therefore has precisely the geometry used in (GS1)--(GS3),
with fixed constants independent of the cutoff.

\paragraph*{The common vacuum density.}
The vacuum constants are sums of elementary Gaussian count constants
and the values of the regular and remainder terms at $\mathbf1$.
The Gaussian count constants per physical volume agree on two tori:
their couplings and choices agree, and each count is proportional to
that volume. Distribute $R^{(j)}(X,\mathbf1)$ equally among its sites.
The remaining density at a site is a sum over domains through it.
Domains contained in a box of side less than one quarter of the smaller
torus have a unique common lift and give identical densities.

At depth $u=K-j$, every remaining domain has
$d_j(X)\ge c\ell_\wedge L^u/M-C$. Its exponential decay, summed
against the domain entropy, costs $Ce^{-c\ell_\wedge L^u/M}$ per
site. A unit physical volume contains $L^{4u}$ sites. Consequently
\[
 \left|\frac{E_{K,\ell,\beta}}{\ell^4}
       -\frac{E_{K,\ell',\beta}}{\ell'^4}\right|
 \le C\sum_{u\ge0}L^{4u}
 e^{-c\ell_\wedge L^u/M}\le C.
\]
The bound is uniform in the cutoff. It uses the exact companion
formation on nonwrapping domains, not a comparison of unspecified
families at fixed norm. Reflection-positive subadditivity then moves a
mirror-periodized source to the fixed reference torus. Upper stability
is used there, while the larger torus contributes its unsourced floor
and the displayed vacuum-density comparison. Thus the common
normalizer can be formed on the reference torus itself.

\subsubsection*{(4d) Rarity, witnesses and ages}

The history propositions below rest on obligation D3 of
Section~\ref{ch:dependency-catalogue}. In particular, we must prove
witness locality, narrow-window estimates and replication with all
required weights for the actual history operations. The argument below
explains how these estimates enter the rarity and age bounds.

\textbf{Theorem 4.13 (Large fields are rare uniformly in the volume).} (Live regions, creation sets, lineages, marks and ages, as well as the law $\tilde\nu$ of the histories, are defined in (2a), and the prepaid bounds are stated in (2e). They have here the meaning given there.) For $g$ small, for every level $0\le k\le K$ (with $\mathcal Z_0=\emptyset$) and every level-$k$ cell $\mathfrak{c}$, one has uniformly

\begin{gather*}
\FitMath{\tilde\nu\big(\mathfrak{lv}_k(\mathfrak{c})\big)=\tilde\nu\big(\mathfrak{c}\subset\mathcal{Z}_k\big)\le Ce^{-cp_0(g_k)},} \\[3pt]
\FitMath{\tilde\nu\big(\mathfrak{c}\text{ is created at step }k\big)=\tilde\nu\big(\mathfrak{c}\subset\mathcal{W}^\wedge_k\big)\le Ce^{-cp_0(g_k)}.}
\end{gather*}

The second bound ($k\ge1$) is a consequence of the first, because creation cells are live. Indeed, $\mathbf{R}$ touches a component only if no creation has occurred in it during the last $\mathsf{N}_k$ steps. Hence a component of $Z_k$ containing a creation cell of step $k$ is not healed at step $k$, which gives $\mathcal{W}^\wedge_k\subset\mathcal{Z}_k$. The same inclusion also bounds the event that $\mathfrak{c}$ lies in a live component in which a creation occurs at step $k$. The theorem relies on the following point. A label's determining set starts from a neighbourhood of its first region, $\Omega_0=\Omega_1^\sim$, and not from the whole torus. Its background therefore requires the bare field to be regular only on that neighbourhood (see under Lemma 4.15). If regularity is required at every bare plaquette of $T_\ell$, then a single rough plaquette, wherever it lies, leaves the step-$1$ background undefined. By totality (2a), every test formed with that background then fails at every place, so that $Z_1$ equals the whole torus and is never healed. Moreover, at fixed $(\beta,g)$, the probability that such a plaquette is present tends to $1$ as $\ell\to\infty$. The theorem concerns only the law of the histories of the run; it gives no bound on the law that $\mu_{K,\ell,\beta}$ induces on the bare plaquettes inside a box of fixed physical size.

The proof is a first-moment argument. For a level-$n$ cell $\mathfrak{z}$, let $\mathcal{E}^{(n)}_{\mathfrak{z}}(k)$ be the survival event of (2a), namely $\mathfrak{z}\subset\mathcal{Z}_i$ for $n\le i\le k$. The event $\mathcal{E}^{(k)}_{\mathfrak{c}}(K)\subset\mathfrak{lv}_k(\mathfrak{c})$ is the never-healed event. (i) If healing takes place at a level $j>k$, the factor of $\mathbf{R}$, $Ce^{-cp_0(g_k)-cp_0(g_j)}$, appears, and it is summable in $j$. (ii) We use $\operatorname{Vol}_i$, $\operatorname{Dep}_i$, the mark $\Theta^{(i)}$ with its dials $\mathsf{w}$, $c_a$, and the end mark $\Theta=\Theta^{(K)}$, all as in (2a). For a component live at level $i$, along its lineage $\operatorname{Vol}_i$ charges every live cell, at every level $j\le i$, with twice the weight $\mathsf{r}_j=M^4\check{R}_j^5$, while $\operatorname{Dep}_i$ contributes $cp_0(g_j)^2(d'_j(W)+1)$ for every component $W$ that is created. By (2e), the upper bound of (4c) stays valid once the mark is prepaid: for the run stopped at an arbitrary level when the dial is $\mathsf{w}=0$, and at the last level when the dial is $\mathsf{w}>0$. In the form needed here, the prepaid bound takes up the large-field bound one lineage at a time and consists of two statements, independent of each other. \emph{Volumes are paid by events.} Let $\mathsf r_j=M^4\check R_j^5$. For every history, and for every component $X$ of $Z_k$, whether live or healed at step $k$, one has

\begin{gather*}
\FitMath{\sum_{(j,C)}\mathsf r_j\,|C|_j\le c_G\sum_{(j,C)}\mathsf r_j\check R_j\,e(j,C),} \\[3pt]
\FitMath{e(j,C):=\big(|\operatorname{pred}|-1\big)^++\big|\operatorname{pred}\cap\operatorname{Act}_{j-1}\big|+\sum_W\big(d'_j(W)+1\big),}
\end{gather*}

where both sums run over the components of the lineage of $X$. Here $\operatorname{pred}$ denotes the set of components live at level $j-1$ that are contained in $C$; $\operatorname{Act}_{j-1}$ is the set of those components that are live even though $\mathbf R$ of level $j-1$ has acted on them; $W$ ranges over those components of the step-$j$ creation set that lie in $C$; and the constant $c_G$ is determined by $L$, $M$ and the remaining constants fixed ahead of the ceiling $\gamma$ imposed on the couplings, and by nothing else. The coupling $\gamma_\star$ is taken small. The inequality is a statement on the geometry of the regions alone; neither a field nor a measure enters. At levels where nothing happens, a component, measured in cells of its own level, shrinks geometrically until it becomes small; and once a component has stayed small over $\check R_j+1$ levels without any event, it belongs to the class of $\mathbf R$. \emph{Masses by lineage traces.} The lineage trace $\mathfrak l$ of $X$ is the list of the regions $\mathsf Z_j(X)$, of the sets created in them, and of the sets $\operatorname{Act}_j$. We set 

\begin{gather*}
\FitMath{\operatorname{wt}(\mathfrak l):=\exp\big[C_{\rm chg}\operatorname{Vol}_k-\operatorname{Dep}_k-\sum_{(j,W)}2p_0(g_j)-\sum_j|\operatorname{Act}_j|\,p_0(g_j)\big]}
\end{gather*}

, where $C_{\rm chg}$ is the charge constant; it is defined by the requirement that, for each cell of level $j$, the sums over the admissible regions, the Gaussian integrals, the normalisation constants and the vacuum-energy counterterms contribute a factor no larger than $e^{2C_{\rm chg}\mathsf r_j}$. This constant is determined by the parameters of the construction alone, and it is fixed before $\gamma_\star$. Let $\mathscr L$ be any set of lineage traces of level $k$, and sum over those histories of level $k$ whose live components carry the lineage traces $\mathscr L$. Then

\begin{gather*}
\FitMath{e^{-E_{(k)}}\sum_h\mathsf t(h)\le e^{C|\mathbb T^{(k)}|}\prod_{\mathfrak l\in\mathscr L}\operatorname{wt}(\mathfrak l):}
\end{gather*}

indeed, one unfolds the majorants of (2c), passing through the components that $\mathbf R$ has acted on, as far as the creations, and one sums the entropy with the lineage trace held fixed. No factor arises from a joining, since a creation at a leaf of the lineage pays for it (Proposition 2.6, move (3)). The relation $\mathsf r_j\check R_j=M^4\check R_j^6\ll p_0(g_j)$ is valid for all couplings not exceeding the anchor coupling $\gamma_\star$. This anchor coupling is fixed small only once $L$, $M$, the dials $\mathsf w$ and $c_a$, and the rate $c'$ introduced below (a constant of this proof only, different from the small $c'$ of the mass bound) have been fixed, and it does not exceed $\gamma_J$; $g$ is subsequently taken small depending on $\gamma_\star$. Consequently the volumes, the entropy of the lineage traces and the mark are all paid for by the events. Precisely, take the rate $c'\ge0$, which is fixed before $\gamma_\star$ and may be given any value; the value $c'=1$ already outweighs the number $C^{|X|_k}$ of connected sets $X$ of cells containing a given cell, because $\mathsf r_k\ge M^4$. Then for every set of lineage traces ending in $X$, 

\begin{gather*}
\FitMath{\sum_{\mathfrak l}\operatorname{wt}(\mathfrak l)\,e^{\mathsf w\operatorname{Dep}_k+c_a\operatorname{Vol}_k}\le e^{-c'\mathsf r_k|X|_k}\,e^{-2\min p_0(g_{j_0})}}
\end{gather*}

, where $j_0$ denotes the first level of a lineage. It is this bound that prepays the mark. There are two reasons for carrying out the argument at the last level, with the end mark. The first reason applies already at $\mathsf{w}=0$. The gain obtained from the first moment in (iii) is the quotient of two quantities: the exponent in the moment bound, equal to $\mathsf{B}(y)=C(1+\log y)$ for every site of the lattice on which the run stops, and the product of $c_a\mathsf{r}_k$ with the count of level-$k$ cells. When the run is stopped at level $k$, the quotient equals $\mathsf{B}(x_k)(M\check{R}_k)^4/(c_a\mathsf{r}_k)\approx\log x_k/(c_a\check{R}_k)$, and this carries no smallness of order $e^{-cp_0(g_k)}$. At the last level there are fewer sites, by the factor $L^{-4(K-k)}$. The second reason concerns the alternative of (iv) in which the torus shrinks. There a single deposit is weighed against the volume, and this requires $\mathsf{w}>0$. Whatever value the dial takes, the end mark registers only lineages of components still live at the last level, that is, only the never-healed lineages. This is why the healed part is split off in (i). Together with the floor one obtains $\mathbb{E}_{\tilde\nu}e^\Theta\le e^{C_g\ell^4}$, hence $\mathbb{E}_{\tilde\nu}\Theta\le C_g\ell^4$ by Jensen's inequality. On the other hand $\operatorname{Vol}_K$ charges every never-healed cell, so $\sum_{\mathfrak{c}}\mathbf{1}_{\mathcal{E}^{(k)}_{\mathfrak{c}}(K)}\le\Theta/(c_a\mathsf{r}_k)$, where $c_a>0$. (iii) Translations by level-$k$ cells commute with the steps up to level $k$; therefore $\tilde\nu(\mathfrak{lv}_k(\mathfrak{c}))$ is the same for every $\mathfrak{c}$. The event $\mathcal{E}^{(k)}_{\mathfrak{c}}(K)$, by contrast, depends on coarser partitions and is not covariant. Averaging gives

\begin{gather*}
\FitMath{\tilde\nu\big(\mathfrak{lv}_k(\mathfrak{c})\big)\le\max_{\mathfrak{c}'}\tilde\nu\big(\mathfrak{lv}_k(\mathfrak{c}')\setminus\mathcal{E}^{(k)}_{\mathfrak{c}'}(K)\big)+\frac{C_g\ell^4}{c_a\mathsf{r}_k\,\#\{\mathfrak{c}'\}}\le Ce^{-cp_0(g_k)}+C_gL^{-4(K-k)}.}
\end{gather*}

(iv) To make $K-k$ large, the run is extended on the same bare lattice up to a larger reference coupling, not exceeding the anchor coupling $\gamma_\star$. This is a fixed upper limit for the reference couplings. It is chosen small once the dials $\mathsf{w}$ and $c_a$ are fixed, in such a way that (i), (ii) and the comparison made in the first alternative below are valid at each reference coupling up to $\gamma_\star$; only afterwards is $g$ chosen small, depending on $\gamma_\star$. This continuation leaves the law up to level $k$ unchanged (3d). Steps (i)--(iii) are applied to the continued run, with $\Theta$ the mark at its last level. There are two alternatives. In the first, the torus shrinks to polylogarithmic size. Every never-healed lineage then contains a creation, so $\Theta\ge\mathsf{w}\,cp_0(g_k)^2$ on $\mathcal{E}^{(k)}_{\mathfrak{c}}(K)$, and by Markov's inequality this single deposit outweighs $C_g\ell^4$. In the second, at least $c\,g^{-2}$ steps are gained, and $L^{-4(K-k)}$ outweighs $e^{cp_0(g_k)}$. Between them, these two alternatives cover all data.

\textbf{Proposition 4.14 (Ages).} Denote by $C$ the component live at level $k$ that contains $\mathfrak{c}$, by $\mathsf{Z}_j(C)$ its lineage regions, and by $j_0$ the least level at which these regions are non-empty. The three ages of $C$ are the age read in the label, the age of the lineage, $\operatorname{age}_k(C;h)=k-j_0$, and the seed age. These ages and the events $\mathfrak{lv}^{=a}_k(\mathfrak{c})$, $\mathfrak{lv}^{\ge a}_k(\mathfrak{c})$, $\mathfrak{sd}^{=a}_k(\mathfrak{c})$ are defined as in (2a). The age read in the label refers to the oldest region of the label that appears in $C$: $\operatorname{age}^{\mathrm{lab}}(C)=k-\min\{t:C\cap Z_t\ne\emptyset\}$, where $Z_t=\Lambda_t^c$ is the large-field region of scale $t$ in the label $\mathcal{S}(h)$ of the history. The proposition states that the three ages are ordered,

\begin{gather*}
\FitMath{\operatorname{age}^{\mathrm{lab}}(C)\,\le\,\operatorname{age}_k(C;h)\,\le\,\text{seed age},}
\end{gather*}

and it gives a decay in the second of them. For the first inequality, note that every point of $C\cap Z_t$ belongs to the live region at each level between $t$ and $k$. It therefore belongs to $\mathsf{Z}_t(C)$, and so $t\ge j_0$. For the second inequality, no component of $\mathsf{Z}_{j_0}(C)$ contains a component that is live at the preceding level. Each such component is therefore created at step $j_0$, and it is contained in $C$. The second inequality may be strict. Suppose a region created at a step $t<j_0$ is contained in $C$. Then it was healed at some step before $k$, and the live region of some later level covers its place anew. As a point set it counts for the seed level. It contributes no factor, however, since no cell of it stays live through all the levels from $t$ to $k$. In terms of events, $\mathfrak{sd}^{=a}_k(\mathfrak{c})\subset\bigcup_{a'\le a}\mathfrak{lv}^{=a'}_k(\mathfrak{c})$. Let $b\ge1$, let $c_a\ge C(b)\log L$ (a floor on $c_a$, which is chosen after $L$), and let $g$ be small depending on $b$ and $c_a$. Then, uniformly, with constants $C$ on the right-hand sides that may depend on $b$ and $c_a$,

\begin{gather*}
\FitMath{\tilde\nu\big(\mathfrak{lv}^{=a}_k(\mathfrak{c})\big)\le Ce^{-cp_0(g_k)}(2L^b)^{-a},} \\[3pt]
\FitMath{\tilde\nu\big(\mathfrak{lv}^{=a}_k(\mathfrak{c})\cap\mathcal{E}^{(k)}_{\mathfrak{c}}(K)\big)\le C(2L^b)^{-a-(K-k)}.}
\end{gather*}

Both bounds decay in the age of the lineage, and so does the age series given further below. Suppose the seed level is formed only from regions that are healed at no step before $k$. Then the two bounds and the age series hold in the seed age as well, that is, for $\mathfrak{sd}^{=a}_k(\mathfrak{c})$: the seed level equals $j_0$, the two ages coincide, and $\mathfrak{sd}^{=a}_k(\mathfrak{c})=\mathfrak{lv}^{=a}_k(\mathfrak{c})$. When the seed level is formed from all regions, the price of survival yields no decay. Two facts show this. First, a level-$k$ cell contains of the order of $L^{4(k-t)}$ cells of level $t$, whereas the bound per creation has an exponent that is merely polylogarithmic in $k-t$. Second, take $k=K$ and $a=K-1$. Applied there, the bound gives at most $C(2L^b)^{-(K-1)}$, at fixed $g$ and for every $K$, for the probability that a fixed cell at the final level is live and occupies the place of a region created at step $1$. Yet the tests of the last step depend only on $x_{K-1},x_K\approx g^{-2}$. A first-moment argument yields only a density, and never a factor per level for a single cell. For this reason the chessboard estimate is now applied to the large fields, which is a second use of reflection positivity.

\textbf{Lemma 4.15 (Chessboard for the law of the histories).} Put $q_a=e^{-c_a}$, where $c_a\ge c^{-1}\log(2L^4)$, equivalently $L^4q_a^c\le\frac12$. In this condition, in the two bounds below and wherever these bounds are quoted, the constant is $c=\frac1{16}$; it is fixed before $c_a$, and it arises as the root, taken over the tiles, of the bound for a pattern that is repeated in one tile out of $2^4$. Fix $\epsilon>0$ and take $g$ small depending on $c_a,\epsilon$. A \emph{tiling} is a translate, by a multiple of the side length of a level-$k$ cell, of the partition of the torus into cubes of side $\mathsf{l}_\square$, called \emph{tiles}, whose faces lie in mirror planes that pass through midpoints of bonds. In each tile, the cube with the same centre and with side $\frac34\mathsf{l}_\square$ is called its \emph{core}. A cell, of whatever level, is called \emph{deep} in a tile when the level-$k$ cell containing it lies in the core of that tile. Then for each level $k\le K$ there is a tile side $\mathsf{l}_\square\approx M\check{R}_k$, measured in level-$k$ units, which is a power of $L$ and spans at least sixteen level-$k$ cells, with the following property. Let $n\le k$, let $\mathfrak{z}$ be a level-$n$ cell situated deep inside a tile $\square_0$ of a tiling, and let $\square_1,\dots,\square_r$ be further tiles, all distinct. Then

\begin{gather*}
\FitMath{\tilde\nu\big(\mathcal{E}^{(n)}_{\mathfrak{z}}(k)\big)\le\mathsf{l}_\square^{-4}\,q_a^{c(k-n+1)},} \\[3pt]
\FitMath{\sum_{\mathfrak{c}_1,\dots,\mathfrak{c}_r}\tilde\nu\Big(\mathcal{E}^{(n)}_{\mathfrak{z}}(k)\cap\bigcap_{i\le r}\mathfrak{lv}_k(\mathfrak{c}_i)\Big)\le\mathsf{l}_\square^{-4}\,q_a^{c(k-n+1)}\,\epsilon^r,}
\end{gather*}

where the sum runs over the level-$k$ cells $\mathfrak{c}_i$ lying deep inside $\square_i$. The first bound is valid for every level-$n$ cell on every torus, including tori too small to carry a tiling with tiles of this side. The bounds are independent of how many tiles the torus contains. Survival costs a factor $q_a^c$ per level, and every further tile that contains a deep live cell costs a factor $\epsilon$. The proof has three steps. \emph{Price of survival:} take the dial $\mathsf{w}=0$, so that only the volumes are prepaid. For the run stopped at level $k$, the moment bound of (2e) then applies. Its exponent per site is $\mathsf{B}(x_k)=C(1+\log x_k)$, which has degree one in the logarithm. Markov's inequality gives, for each event $\mathsf{E}_k\subset\mathfrak{H}_k$,

\begin{gather*}
\FitMath{\tilde\nu(\mathsf{E}_k)\le\exp\Big[\mathsf{B}(x_k)|\mathbb{T}^{(k)}|-\inf_{\mathsf{E}_k}\Theta^{(k)}\Big].}
\end{gather*}

Consider a cell that is live at each of the levels $n,\dots,k$. At every one of these levels it belongs to the lineage of the component that is live at level $k$ and contains it. Thus $\operatorname{Vol}_k$ charges it once for each level. Consequently the joint survival of several cells is bounded by $e^{\mathsf{B}(x_k)|\mathbb{T}^{(k)}|}$, multiplied by a factor $q_a^{\mathsf{r}_j}$ for each level $j$ and each distinct cell encountered at that level. The other cells of a pattern, that is, of a tuple of cells of one level in a tile as defined below, are paid for by other payers. Take a cell that has no large-field cell within two cells of it and at one of whose bonds a lower threshold is passed by the scaled fluctuation field of the step. The price of such a cell is borne by the conditional law of the fluctuation field in its box. If the component of a cell is absorbed into a component healed at a step $\le k$, the lineage of that component pays for the cell. \emph{Witnesses:} the chessboard estimate concerns functions of the bare field, whereas $\tilde\nu$ is a law on histories. A history is determined by the bare field together with independent auxiliary variables, namely one uniform variable for every possible application of a probability kernel in the steps up to $k$, each of them attached to a location. By the tile data of a run we mean the restriction of the bare field to the tile, together with the auxiliary variables attached to locations in it. The estimate rests on a locality property, the \emph{agreement of two runs inside a tile}. Consider two runs with identical tile data. Assume that in the first run, up to some level $J'\le k$, two things hold in the tile. First, no tested scaled fluctuation value lies in a narrow window around its threshold (at a test of step $j$ of a value scaled by $g_i$, $i\le j$, the window has width $2\varepsilon^{\rm dec}_j/g_i$, with $\varepsilon^{\rm dec}_j=Ce^{-\delta D_j}\varepsilon_j$; here $\delta$ is the rate at which two backgrounds that have the same data in a region approach each other within that region, and $D_j=\delta^{-1}[C+4\log\mathsf{l}_\square+(4\log L+\log2)(k-j)+\dots]$, where the dots denote terms of order $\mathsf{N}_j\log L+\log p_0(g_j)$. Thus the width decreases like $\mathsf{l}_\square^{-4}(2L^4)^{-(k-j)}$, which is set against the $\mathsf{l}_\square^4L^{4(k-j)}$ sites of level $j$ in a tile. For the decoupling, a separation of $CD_j$ lattice units of level $j$ is needed, which amounts to $\lceil CD_j/(M\check{R}_j)\rceil$ layers of level-$j$ cells. Measured in units of level $k$, these cells shrink like $L^{-(k-j)}$, so the layers consumed over all the levels together amount to a fraction of the tile side). Second, no long cluster of cells of a level $j\le J'$ that lie in $\Omega_j^c$ or carry a tested bond at which a lower threshold is passed by the scaled fluctuation field of the step (a trigger) appears. Then, at the end of step $J'$, the two runs coincide in the inner layers of the tile: their block fields, both current and stored, agree there, and so do the regions of the label. Moreover, every live component that meets these layers is contained in the tile and is common to both runs, and the same is true of its label, of its inner history and of its stored block fields. The recorded collar test is centred at the background localized in the component, as in the stages above. Its comparison with the true averaged background supplies the collar bound, rather than an equality of the two sharp indicators. The common data do not include the stored fluctuation fields. On the bonds examined by a passed fluctuation test (2a) they are built from the averaged background, which is not local. Among the later moves, a single test of $\mathbf{R}$ depends on them, and on those bonds that test is passed in every run. The hypotheses concern only one of the two runs; no condition is placed on the other run outside the tile. On the runs that meet the two hypotheses, this property makes every survival event read deep inside a tile a function of the tile data only. As a result, products over distinct tiles can be integrated against the Wilson measure. Two cases occur for the survival of a cell deep in a tile. Either the tile data force it, or, at the first level with a trigger, it is accompanied by the event of a pattern, an event that is likewise forced by the tile data. By a pattern we mean a tuple of $m=1+2(k-J)$ cells of that level $J$ in the tile, any two of which are at least $2c_{\rm adj}(J)+3$ cells apart and each of which lies within $m_0(J)/2$ cells of the first cell of the tuple. Every cell of the tuple either lies within two cells of some cell of the large-field region of that level, or carries a bond at which a lower threshold is passed by the scaled fluctuation field of the step. Here $c_{\rm adj}(J):=\lceil CD_J/(M\check{R}_J)\rceil+C$, the number of layers that the decoupling takes at level $J$ plus a constant, is the largest step, measured in cells, by which a cluster advances; and $m_0(J):=6(c_{\rm adj}(J)+1)m$ is the diameter, measured in cells, beyond which a cluster counts as long, the factor $6$ leaving room for the $m$ cells. The widths $m_0(J)$ belong to the layers consumed at the levels $J\le k$, and these layers together amount to at most $\mathsf{l}_\square/32$. The proof is an induction on the level and, within one level, over the stages of $\mathbf{R}$. Locality requires that the determining set of a label start from a neighbourhood of the first region ($\Omega_0=\Omega_1^\sim$), not from the whole torus. In Ba{\l}aban, Comm. Math. Phys. 119 (1988), the determining set has as its support a small neighbourhood of $\Omega_1$ (p. 255), and the region $\Omega_0=\Omega_1^\sim$ appears in a localisation property of the boundary terms (pp. 261--262). If regularity is required of all the bare plaquettes of the torus, then a single rough plaquette anywhere turns off the averaged background on the whole torus, and nothing is local. Hence every survival event read deep inside a tile is bounded above by a finite family of witness functions $f_\iota\in[0,1]$, each depending only on the bare field in that tile. The cost is $C$ per tile, which pays for the windows: under the conditional law of a tested fluctuation value, a narrow window has small probability. Thus

\begin{gather*}
\FitMath{\tilde\nu\Big(\bigcap_{i\le r}\mathcal{E}^{(n_i)}_{\mathfrak{z}_i}(k)\Big)\le C^{r+1}\sum_{\iota_0,\dots,\iota_r}\Big\langle\prod_if_{\iota_i}\Big\rangle_{K,\ell,\beta}\le\prod_{i\le r}C\sum_{\iota_i}\|f_{\iota_i}\|_{\mathrm{ch}},}
\end{gather*}

where the second inequality follows from reflection positivity of the Wilson measure in the faces of the tiles. \emph{Root and surplus:} a witness has norm bounded by a root of the probability that its pattern occurs, copied across the whole torus by reflections; among these copies only the translates are retained, and they lie in one tile out of $2^4$, which gives $c=\frac1{16}$. There are three payers for a replicated pattern. Consider first a cell such that no large-field cell lies within two cells of it and that has a bond at which a lower threshold is passed by the scaled fluctuation field of the step; it is paid for by the conditional law of the fluctuation field in its box. A cell whose component remains live up to level $k$ is paid for by the mark $\Theta^{(k)}$ of the run stopped at level $k$. If the component of a cell is absorbed into a component healed at a step $\le k$, the cell is paid for by the lineage of that component. Taking the root turns the volume factor into $e^{\mathsf{B}(x_k)\mathsf{l}_\square^4}$ per tile and turns $q_a$ into $q_a^c$. Since $\mathsf{l}_\square^4\log x_k\ll\mathsf{r}_k=M^4\check{R}_k^5$, the surplus $q_a^{c\mathsf{r}_k}$ of one level absorbs this factor, together with the number of level-$k$ cells in a tile and $\epsilon^{-1}$. The patterns are summed one level at a time. A pattern of level $J$ pays $q_a^{\mathsf{r}_J}$ for each of its $m=1+2(k-J)$ cells, and survival from level $n$ up to level $J-1$ pays $q_a$ for each of those levels. A tile contains at most $\mathsf{l}_\square^4L^{4(k-J)}(m_0(J)+1)^{4(m-1)}$ patterns of level $J$, where $m_0(J)\le C(x_k+k-J)^3$ is the length from which onward a cluster counts as long. Half the weight of each cell beats this number, except for the positions of the first cell: $Cm\log(x_k+k-J)$ is set against $c\,c_a\mathsf{r}_Jm$, where $\mathsf{r}_J\ge M^4(\log(x_k+c(k-J)-C))^{5r}$ and $r\ge2$ is the exponent of $\check{R}_k$. This is a logarithm against a power of a logarithm, uniformly in $k-J$, because $x_J$ grows linearly in $k-J$. Finally, the inequality $(J-n)^++1+2(k-J)\ge(k-n+1)+(k-J)$ retains $q_a^{c(k-n+1)}$ and leaves a further factor $q_a^{c(k-J)}$. This factor is set against the $L^{4(k-J)}$ positions, and together they give $(L^4q_a^c)^{k-J}\le2^{-(k-J)}$. At this point $\ell$ disappears.

\emph{Proof of Proposition 4.14.} \emph{Anchor:} choose any cell $\mathfrak{z}$ of the lineage region $\mathsf{Z}_n(C)$ at the first level $n=j_0=k-a$ of the lineage. The lineage regions are nested and contained in the live regions. Hence $\mathfrak{z}$ stays live throughout the levels $n,\dots,k$, so that the event $\mathcal{E}^{(n)}_{\mathfrak{z}}(k)$ occurs, and the cell belongs to the component of $\mathfrak{c}$. This is the point at which the age of the lineage enters: a region healed in the meantime provides no such cell. \emph{Sum over anchors:} a box of side $\mathsf{l}_\square$ contains at most $\mathsf{l}_\square^4L^{4(k-n)}$ anchors, and each costs $\mathsf{l}_\square^{-4}q_a^{c(k-n+1)}$. For an anchor at a distance of $i\ge i_0$ boxes, the live component is connected and therefore passes through a chain of at least $i$ distinct adjacent boxes, each of which contains a live cell. These cells need not be deep. Consider the $2^4$ tilings whose face planes are displaced by $\pm\lfloor n_\square/4\rfloor$ cells from a corner of the level-$k$ cell containing $\mathfrak{z}$, where $n_\square\ge16$ is the number of level-$k$ cells along a side of a tile. Every level-$k$ cell is deep in a tile of one of these tilings, and $\mathfrak{z}$ is deep in all of them. By the pigeonhole principle a single tiling serves at least $i/16$ of the cells. These cells lie in at least $i/16^2$ distinct tiles, since a tile meets at most $2^4$ boxes, and at least $i/i_0$ of these tiles, where $i_0=2\cdot16^2$, are different from the tile of $\mathfrak{z}$. By the second bound of the lemma and the choice $\epsilon=(2C_{\rm ent})^{-i_0}$, where $C_{\rm ent}$ is the entropy per link of the chain, such an anchor costs an additional factor $16\,C_{\rm ent}^{\,i}\epsilon^{i/i_0}=16\cdot2^{-i}$, which is summable against the $(2i+3)^4$ boxes of a shell. Finally $L^4q_a^c\le(2L^b)^{-2}$, and this is exactly the floor $c_a\ge C(b)\log L$ imposed in the statement of the proposition. Therefore $\tilde\nu(\mathfrak{lv}^{=a}_k(\mathfrak{c}))\le C(2L^b)^{-2a}$. The geometric mean of this bound and that of Theorem 4.13 is the first bound. Summing it gives the age series

\begin{gather*}
\FitMath{\sum_{a\ge0}L^{ba}\,\tilde\nu\big(\mathfrak{lv}^{\ge a}_k(\mathfrak{c})\big)\le Ce^{-cp_0(g_k)},}
\end{gather*}

which two later parts use. In (7b) the series is needed with rate $b\ge8$, because for two points a live component of age $a$ carries a weight of up to $L^{8a}$; there the series underlies the two-point law, Theorem 0 (c)(3), (c)(4). In (7d) the series enters with rate $b>16$, on each torus, since a weight $L^{16a}$ arises for four points; this concerns clause (c)(5). The expansions of (7b) and (7d) require the age up to which a live component stores fields. This is the age read in the label. For a component holding an island, the block of which retains the fields of its core (2b), it is at most the age of the lineage. Here, as in (7d), the series is formulated with the age of the lineage, and this age is never smaller. The series therefore covers what these expansions require: 

\begin{gather*}
\FitMath{\{\mathfrak{c}\subset C,\,\operatorname{age}^{\mathrm{lab}}(C)\ge a\}\subset\mathfrak{lv}^{\ge a}_k(\mathfrak{c})}
\end{gather*}

. Applying the lemma at level $K$ in place of $k$ yields the second bound.

\emph{A second route to decay in the age: the two-term age bound.} This bound complements Lemma 4.15 and does not replace it. The prepaid bounds of (2e) alone give decay in the age of the lineage. They need no tiles, no auxiliary variables and no second run, and therefore no windows and no triggers. As in (i) and (ii) of the proof of Theorem 4.13, the lineage of the live component containing $\mathfrak{c}$ is classified by its fate. Suppose it is healed at a level $j>k$. Then it contributes the factor of $\mathbf{R}$, and the sum over lineage traces is restricted to traces of age at least $a$. Suppose instead it is never healed. Then the moment bound is applied at the last level, at the dial $\mathsf{w}=0$, to the histories that contain such a component. Throughout this event the mark collects at least $\bar c_a:=2c_aM^4$ for each level at which the lineage has a member, and these contributions add up over the levels. Consequently no sum over positions of an anchor arises, and the linear gain of the oldest creation remains. Let $c_a>0$ and take $\gamma_\star$ small depending on $c_a$. Then on each torus of the family, for each level $k$, each level-$k$ cell $\mathfrak{c}$ and each $a\ge0$,

\begin{gather*}
\FitMath{\tilde\nu\big(\mathfrak{lv}^{\ge a}_k(\mathfrak{c})\big)\le Ce^{-\bar c_aa}\,e^{-cp_0(g_k)}+C\exp\Big[\mathsf{B}(x_{K'})\,|\mathbb{T}^{(K')}|-2p_0(g_k)-\bar c_a\,(a+K'-k+1)\Big],}
\end{gather*}

where $K'\ge K$ is the last level of an arbitrary prolongation of the run (3d), the run itself counting as one of its prolongations. The first term contains no volume, and no condition ties $k$, $a$, $K$ and $\ell$ together. The consequences of the bound follow from its second term, case by case; none of the cases is a hypothesis of any statement. (1) There is a side $\ell_V(g)$, tending to infinity as $g\to0$, with the following property. For every torus with $\ell\le\ell_V(g)$, the prolongation of (iv) terminates in the alternative where the torus is reduced to polylogarithmic size. In that case the gain $2p_0(g_k)$, of degree $\mathfrak{p}_0$ in the logarithm, dominates the volume term, of degree $4r+1$. The bound is then $Ce^{-\bar c_aa}e^{-cp_0(g_k)}$ at every level. (2) The last lattice of the reference torus $T_{\ell_\star(g)}$ of Proposition 7.3 has polylogarithmic size even without prolongation. Hence that torus satisfies the same bound at all levels. (3) If $\ell>\ell_V(g)$, the torus satisfies the same bound at those levels at which $\bar c_a(K'-k+1)\ge\mathsf{B}(x_{K'})|\mathbb{T}^{(K')}|$. This is a floor on the depth that depends on the torus; once it is met at one level, it is met at every lower level. At all of these levels, if $\bar c_a\ge b\log L+\log2$, summing the bound gives the age series at the rate $b$. Thus this road alone supplies what (7b) uses, namely the age series at the rate $b\ge8$, at all levels of the reference torus, and hence Theorem 0 (c)(3), (c)(4). For (7d) it supplies the series at $b>16$ on each torus. On the tori of (1) it holds at all levels; on the tori of (3) it holds at the levels below one that depends on $\ell$. Theorem 7.5 [to be proved; Package F] states clause (c)(5) in exactly this form: $\mathsf{d}\le\mathsf{d}_4$ if $\ell\le\ell_4(g)$, and $\mathsf{d}\le\mathsf{d}_4(\ell)$ otherwise, where $\ell_4(g):=\max\{\ell_V(g),\ell_\star(g)\}$. The contribution of Lemma 4.15, through its tiles, windows and triggers, is decay at all levels and on all tori, with constants independent of $\ell$. This is uniformity as $\ell\to\infty$. Because of it, one threshold suffices for every torus, and the four-point law carries over to the limits on $\mathbb{R}^4$ (7d).

\emph{The age constant.} In (2e) two roles of $c_a$ are described; the one at work here is role (A), the dial of the mark. Its effect on the law of the histories is $C(2L^b)^{-a-(K-k)}$. Role (B), the age rate $q_a$ of the aged load of Proposition 2.12 [to be proved; Package C], concerns the pairs of lists belonging to a class of the kind given there (2e), and it is used in (8a). The floors on $c_a$ listed in (2e) concern histories, rotations and age-resolved rarity. Proposition 4.14 imposes the third of these.


\subsubsection*{(4e) The available source-transfer consequences}

Section~\ref{ch:all-state-source-transfer} proves the transfer of
the ordinary analytic source input to periodic Gibbs states, including
the one-point difference bound. Theorem~\ref{thm:graded-rooted-source}
gives the graded bound used in Theorem 4.2, with its
$\mathsf d^{-5n}$ separation loss. Section~\ref{ch:uniform-insertion}
proves the insertion implication under its stated source and mirror
hypotheses. These results do not verify the ordinary interacting
source input or the history-persistence requirements of (4d).

\clearpage
\ensuremath{\dagger} \textbf{Cutoffs beyond first passage.} The dilation identity of (3d) gives

\begin{gather*}
\FitMath{S^T_{n;K,\ell,\beta}(f_1,\dots,f_n)=S^T_{n;K-q,\,L^q\ell,\,\beta}\big(f^{[q]}_1,\dots,f^{[q]}_n\big),} \\[3pt]
\FitMath{\|f^{[q]}\|_\sharp\le\|f\|_\sharp,}
\end{gather*}

since $\mathcal{O}$ carries no power of the spacing. The same relation holds between $\mathbb{L}_K$ and $\mathbb{L}_{K-q}$. Let $\beta\ge g^{-2}+B_\sharp$ and $q\ge0$, and suppose that $K-q=K(\beta,g)$. The bounds of Theorem 4.2 are then valid at the cutoff $K$, with $(L^q\mathsf{d};L^qR;L^q\ell)$ in place of $(\mathsf{d};R;\ell)$.
% END MODULE chapters/07_sources.tex


% BEGIN MODULE chapters/ordinary_stability_deduction_v6.tex
\section{Ordinary stability from a normalized history representation}
\label{ch:ordinary-history-stability}

Assume the normalized positive history representation and weighted
label bounds stated below. We derive upper stability and the unsourced
partition floor, retaining each label-dependent integral operation.
The proof uses the total-variation label sum and the regular-label
lower bound as separate estimates.

Let the final four-dimensional periodic lattice have $N$ sites and
$4N$ independently oriented bonds. Haar measure on each
$\mathrm{SU}(2)$ bond is normalized. Let $w$ range over a specified
ordinary complex source neighborhood containing zero. Assume there
are complex measures $\tau_S(w)$, indexed by a finite set of labels,
such that
\[
 Z(w)=\sum_S\tau_S(w)(\mathcal V),\qquad
                      \tau_S(0)\ge0.
 \tag{OH1}\label{oh:exactness}
\]
Assume all normalizers and scalar quantities below are well defined;
$E$ is the real vacuum constant of the unsourced all-small-field
label. Put
\[
 \mathcal E(w)=\sum_{j,X}e_j(X;w),\qquad
 \mathcal E(0)=0,\qquad
                  \widetilde Z(w)=e^{-E-\mathcal E(w)}Z(w).
\]
At each finite cutoff the sums are finite. Their locality and uniform
estimates, rather than convergence of a sum over cutoffs, are what
matter here.

\subsection*{A total-variation criterion for upper stability}

Let $\mathcal E_S(w)$ denote the vacuum constant in the exponent of
the term carrying label $S$. Suppose there are nonnegative weights
$W_S$ and real label costs $q(S)$ satisfying
\[
 \|e^{\mathcal E_S(w)}\tau_S(w)\|_{\mathrm{TV}}\le W_S,
 \qquad
 \operatorname{Re}[-E-\mathcal E(w)-\mathcal E_S(w)]
                                      \le C_0N+q(S),
 \tag{OH2}\label{oh:term-bound}
\]
and the weighted label bound
\[
                     \sum_S e^{q(S)}W_S\le e^{C_1N}.
 \tag{OH3}\label{oh:label-sum}
\]
The constants are independent of cutoff, volume and $w$ in the
specified source neighborhood. Then
\[
                         |\widetilde Z(w)|\le e^{(C_0+C_1)N}.
 \tag{OH4}\label{oh:upper}
\]
Indeed multiply each measure in \eqref{oh:exactness} by the common
normalization, insert and remove its own $e^{\mathcal E_S(w)}$,
take total variations, apply \eqref{oh:term-bound}, and sum by
\eqref{oh:label-sum}. This proof takes no logarithm of $Z(w)$ and
asserts no absence of its complex zeros.

In a history construction, $W_S$ is obtained by bounding the
unsourced integrations after inserting a reserve factor such as
$e^{c v(\omega)}$ on every local history. The ordinary source changes
an island exponent by an extensive amount,
$|\Psi_\omega(w)|\le C_r\,n(\omega)$, where sites are counted
in their own scales. A sufficient reserve is
$c v(\omega)\ge C_r n(\omega)$, uniformly in the history and
the exterior arguments. This is an ordinary mass reserve for one
presentation, and is distinct from estimates comparing marked or
weighted derivatives of two different presentations.

It is essential that the majorant used for each component be a
function of its own integration variables. A bound on its integral
at each fixed exterior field is insufficient in a coupled product.
For example, let both component factors be
$f_1(x,y)=f_2(x,y)=\mathbf1_{|x-y|<\epsilon}$ on $[0,T]^2$.
Each integral of a factor in its own variable is at most
$2\epsilon$, whereas the double integral of their product has
order $\epsilon T$, not $\epsilon^2$.
A bound with the supremum over the admissible exterior pair
\emph{inside} the local integral supplies a component majorant;
Tonelli then permits multiplication when the integration-variable
groups are disjoint. The construction must also specify which
Wilson factors are allotted to those majorants. An operation with
no Wilson factor of its own has no independent plaquette-tail gain.

Neither \eqref{oh:term-bound} nor \eqref{oh:label-sum} is a
consequence of the last triangle inequality. They encode ordinary
renormalization estimates: suppression of creation events,
control of their genealogies, the allocation of Wilson and Gaussian
factors without reuse, and payment of the missing vacuum normalizers
of every large-field label. In particular, estimates with the volume
of each fine lattice cannot simply be summed to obtain a bound
with the volume $N$ of the final lattice.

\subsection*{The floor from one positive label}

Let $x\ge1$ be the last inverse coupling and set
$r=cx^{-1/2}$ with a fixed $0<c\le1$. Write
\[
 \mathcal B_r=\{V:\|V_b-1\|_{\mathrm{op}}<r
                                      \hbox{ for every bond }b\}.
\]
Suppose the all-small-field label $S_0$ has a density with respect
to final product Haar measure satisfying
\[
 d\tau_{S_0}(0)(V)
   \ge e^E\chi(V)\exp[-x\mathcal A(V)-C_IN]\,dV,
 \quad
 \chi=1\text{ on }\mathcal B_r,
 \quad \mathcal A(V)\le C_Ar^2N\text{ on }\mathcal B_r.
 \tag{OH5}\label{oh:small-field-floor-input}
\]
Here $\mathcal A$ may be the action of a minimizing background.
If so, inclusion in the characteristic set and the displayed action
bound are background estimates which must be verified; they do not
follow merely by renaming $\mathcal A$ as a coarse action.

For normalized Haar measure there is an absolute $c_H>0$ such that
the operator-norm ball of radius $r\le1$ about the identity has
measure at least $c_Hr^3$. To prove this write
$U=\cos\theta\,1+i\sin\theta\,n\cdot\sigma$, so
$\|U-1\|_{\mathrm{op}}=2\sin(\theta/2)$ and normalized Haar
measure is $(2\pi^2)^{-1}\sin^2\theta\,d\theta\,d\Omega(n)$.
The interval $0\le\theta\le r$ lies in this ball. Since
$\sin\theta\ge2\theta/\pi$ on $[0,1]$, its measure is at least
$8r^3/(3\pi^3)$. Decreasing the constant handles the irrelevant
choice of an open instead of a closed ball.

Positivity in \eqref{oh:exactness} permits restriction to $S_0$,
then to $\mathcal B_r$. Independence of the $4N$ Haar bonds gives
\[
 \begin{split}
 e^{-E}Z(0)
  &\ge e^{-(C_I+C_Ac^2)N}(c_Hc^3x^{-3/2})^{4N}\\
  &\ge\exp[-(C+6\log x)N].
 \end{split}
 \tag{OH6}\label{oh:floor}
\]
Since $\mathcal E(0)=0$, this is the required floor for
$\widetilde Z(0)$. The constant depends on the fixed choices and
on $C_I,C_A$, but not on cutoff or volume.

For the Wilson action of the coarse bond field itself, the action
bound in \eqref{oh:small-field-floor-input} is elementary.
Telescoping the product of four unitary matrices gives
$\|V_p-1\|_{\mathrm{op}}\le4r$. For $\mathrm{SU}(2)$,
$1-\frac12\operatorname{Re}\operatorname{Tr}V_p
=\frac12\|V_p-1\|_{\mathrm{op}}^2\le8r^2$.
There are $6N$ plaquettes, so one may take $C_A=48$.
This calculation does not prove the corresponding assertion for
a non-Abelian minimizing background or for the tests at all its
earlier levels; those are the additional parts of
\eqref{oh:small-field-floor-input} in a multiscale construction.

\subsection*{Local energy bounds and the first-passage window}

Suppose the source families are actual functions of the source
values in a fixed-layer enlargement $X^+$, holomorphic in its
finite-dimensional source parameters and real at real arguments.
Suppose their differences at the trivial field obey
\[
 |\Delta E_j(X,z;w)|\le C_Ee^{-\kappa d_j(X)},\qquad
 |\Delta R_j(X;w)|\le C_Re^{-\kappa d_j(X)}.
\]
These are source-minus-unsourced differences: assume explicitly
$\Delta E_j(X,z;0)=\Delta R_j(X;0)=0$.
Let the number of possible base points $z\in X$ be at most
$C_M(d_j(X)+1)$, as follows for domains covered by the vertices
of a cubical spanning tree. For
$e_j(X;w)=\sum_{z\in X}\Delta E_j(X,z;w)+\Delta R_j(X;w)$,
\[
 |e_j(X;w)|\le C_\delta
                     e^{-(\kappa-\delta)d_j(X)}\qquad(\delta>0).
 \tag{OH7}\label{oh:energy-bound}
\]
This uses only $(m+1)e^{-\delta m}\le C_\delta$. If the rooted
entropy rate is $\kappa_D<\kappa$, choose
$0<\delta<\kappa-\kappa_D$. The resulting energy decay still
exceeds entropy. Locality, holomorphy, realness and vanishing at
zero pass through the finite base-point sum. Dependence on the
source function itself is a structural hypothesis here; it is not
implied by a bound on a chosen list of source coefficients.

If a first-passage argument supplies
$g^{-2}<x\le g^{-2}+B_\sharp$, then the constants in
\eqref{oh:upper} and \eqref{oh:floor}, when bounded by
$C(1+\log x)$, are uniform over the cutoff data at this fixed
$g$. This is the final elementary step producing constants of
the form $C_g$. It supplies no proof that an admissible run reaches
that window.

\subsection*{What a pointwise polymer criterion would supply}

There is a different sufficient mechanism when a construction
actually has a pointwise hard-core gas. Let its polymers be nonempty
subsets of a set of $N$ sites. Suppose, uniformly in a real coarse
field $V$, its real activities obey
\[
 \sum_{\gamma'\not\sim\gamma}|z_{\gamma'}(V)|e^{a(\gamma')}
       \le a(\gamma),\qquad
 \sup_x\sum_{\gamma\ni x}|z_\gamma(V)|e^{a(\gamma)}\le b.
 \tag{OH8}\label{oh:KP}
\]
These inequalities remain valid after multiplication of all
activities by any $t\in[0,1]$. The convergent cluster
expansion makes its partition function $\Xi_V(t)$ nonzero and
continuous on this interval. It is real there and
$\Xi_V(0)=1$; hence $\Xi_V(1)>0$. The absolute cluster theorem
bounds its logarithm by
$\sum_\gamma|z_\gamma|e^{a(\gamma)}\le bN$: sum the rooted
cluster bound and then use the nonempty site supports in
\eqref{oh:KP}. The branch is the one continued from $t=0$.
It is real on the interpolation, and therefore gives
\[
                     e^{-bN}\le\Xi_V(1)\le e^{bN}.
\]
For complex source activities, the corresponding uniform criterion
gives $|\Xi_V(w)|\le e^{bN}$. With a verified common Haar/Wilson
integral representation and its action estimates, these facts
could be used to prove stability and a floor.

This lemma is not a license to replace a label/history representation
by that integral. Local polymer exponentiations on small-field
regions, or exponentiation of the sum over healed components, do
not by themselves provide a pointwise gas over all final coarse
fields. For the history representation the concrete sufficient
conditions are \eqref{oh:term-bound}, \eqref{oh:label-sum}, and
\eqref{oh:small-field-floor-input}; verification of those ordinary
conditions remains the upstream work.

% END MODULE chapters/ordinary_stability_deduction_v6.tex


% BEGIN MODULE chapters/all_state_source_transfer_v3.tex
\section{Transferring ordinary source bounds to all periodic states}
\label{ch:all-state-source-transfer}

The transfer in this section uses reflection positivity and ordinary
complex source bounds on finite tori. It does not use geometric
forgetting estimates for derivatives of inherited activities. Nor does
it assume moment bounds for Gibbs states of a perturbed interaction.
The original arbitrary periodic state enters through its DLR equation;
the only measures to which a source estimate is applied are finite
Wilson torus measures.

Fix an odd integer $L\ge3$, $a=L^{-K}\le1$, a coupling $\beta>0$, and
$\mathcal O_a(f)=\sum_p f(x(p))a_p$. All sums defining this observable
are finite when $f$ has compact support. Although its supremum norm
may grow as $a^{-4}$, it is bounded at each fixed spacing. Every
periodic Wilson DLR state is invariant under all lattice isometries
and positive at all lattice mirrors by
Theorem~\ref{thm:amrp-all-periodic}.

\subsection*{An exact domination lemma, before source estimates}

Use a cubic board of cells of side $l=L^j\ge a$, whose face
planes cut links in half. Take cubic tori of side $\ell=2L^m$ with
$m\ge\max(0,j)$, having an even number of cells in every direction.
Every sequence of sides tending to infinity in this family eventually
accommodates any fixed coarser board of allowed side. All sequences in
this section are taken in this family. When localization cubes have
side $Ma_s$ with the construction convention $M=L^{m_0}$, restrict
to the cofinal subfamily $m\ge m_0$. Then
$\ell/(Ma_s)=2L^{m-m_0+s}$ is an even integer at every scale, and all
localization partitions are retained. The cell algebra consists of
bounded functions of bonds lying wholly inside the open cell; crossing
bonds belong to no cell algebra. Reflected copies are conjugated when
the reflecting isometry has negative determinant. For a cell function
$F$ define

\begin{gather*}
\FitMath{p_{l,\ell}(F)=
 \left\langle\prod_{C} F^{[C]}\right\rangle_{a,\ell,\beta}^{1/N},} \\[3pt]
\FitMath{N=(\ell/l)^4.}\tag{ST1}
\end{gather*}

The number inside the root is nonnegative. To see this, split the
board at one face and the opposite face. The disseminated product is
$\Theta\Phi\,\Phi$ for the product of its copies in one half.
Finite-torus reflection positivity applies.

We will use, including for complex functions,

\begin{gather*}
\FitMath{\left|\left\langle\prod_C F_C\right\rangle\right|
 \le\prod_Cp_{l,\ell}(F_C),} \\[3pt]
\FitMath{p(F+G)\le p(F)+p(G),} \\[3pt]
\FitMath{p(\bar F)=p(F),} \\[3pt]
\FitMath{p(1)=1.}\tag{ST2}
\end{gather*}

Here are details extending the finite chessboard argument to this
form. In one direction transport all slab profiles to one reference
slab, using conjugated reflection for odd copies. For an arbitrary
filling of the $n$ slabs, $n$ even, reflection Cauchy--Schwarz bounds
the squared absolute partition integral by the two integrals obtained
by retaining either half and reflecting it. Normalize label $j$ by
$(Z_j^{\rm dis}+\epsilon)^{1/n}$, where $Z_j^{\rm dis}\ge0$ is its
fully disseminated integral and $\epsilon>0$. The normalized filling
of maximum absolute value exists. Both reflected fillings also
maximize whenever that maximum is positive: each is nonnegative,
bounded by the same maximum, and their product bounds its square.
A constant run in a maximizing filling can be doubled by retaining
the half starting at that run. Once its length reaches $n/2$, a
constant filling results. Every constant normalized filling is at
most one. Thus the maximum is at most one; the zero maximum is
trivial. Letting $\epsilon\downarrow0$ proves the one-direction bound
even when some disseminated integral is zero. Applying this argument
successively to the four coordinate directions proves the first
inequality in (ST2): coordinate reflections commute, and each
successive dissemination copies the preceding pattern literally in
the other coordinates. The resulting exponent for each original cell
is $1/N$.

Expanding the disseminated product for $F+G$ and applying the first
inequality term by term bounds it by
$\sum_{k=0}^N\binom Nk p(F)^kp(G)^{N-k}$. Its nonnegative $N$th root
proves the triangle inequality. Exactly half the copies are
conjugated, giving $p(zF)=|z|p(F)$; conjugating the entire nonnegative
disseminated integral proves $p(\bar F)=p(F)$. In particular
$p(\operatorname{Re}(e^{-i\phi}F))\le p(F)$.

\paragraph{Domination.}
For any periodic Wilson DLR state $\rho$, any finite collection of
occupied cells, and any sequence of compatible tori
$\ell_k\to\infty$,

\begin{gather*}
\FitMath{\left|\rho\left(\prod_C F_C\right)\right|
 \le\liminf_{k\to\infty}\prod_Cp_{l,\ell_k}(F_C).}\tag{ST3}
\end{gather*}

Only finitely many factors differ from $1$.

First consider a bounded real function $H$ of one cell. Sum its mirror
copies over a cube of side $\ell$. Reflection invariance of $\rho$
makes the expectation of the sum equal to $N\rho(H)$. Conditional
Jensen bounds that expectation by the supremum, over exterior
configurations, of the logarithm of its conditional exponential
moment. Replace the exterior condition by periodic bonds. After
adjoining dummy independent Haar variables if necessary, the two
integrals use the same variables. Only boundary plaquettes change;
each changes the Hamiltonian by at most $2\beta$, and there are
$O((\ell/a)^3)$ such plaquettes. Source factors do not cross the outer
cell faces and therefore agree in the two integrals. It follows that

\begin{gather*}
\FitMath{\rho(H)\le\log p_{l,\ell}(e^H)
                  +C\beta l^4a^{-3}\ell^{-1}.}\tag{ST4}
\end{gather*}

This estimate is uniform in the exterior configuration, not in $a$;
the volume limit is taken with $a$ fixed.

For bounded real $G$, put $H=\log(1+tG)$ with
$0<t<(2\|G\|_\infty)^{-1}$. The elementary inequalities
$\log(1+tG)\ge tG-Ct^2\|G\|_\infty^2$ and
$\log p(1+tG)\le t p(G)$ turn (ST4), first as
$\ell_k\to\infty$ and then as $t\downarrow0$, into
$\rho(G)\le\liminf_kp_{l,\ell_k}(G)$. For complex $F$ choose $\phi$
so that $\rho(\operatorname{Re}(e^{-i\phi}F))=|\rho(F)|$ and use
the last consequence of (ST2).

Finally enclose the finitely many occupied fine cells strictly inside
one larger coarse board cell. The one-cell inequality applies to
$G=\prod_CF_C$. On every compatible torus the fine-board chessboard
inequality gives
$p_{\rm coarse,\ell}(G)\le\prod_Cp_{l,\ell}(F_C)$: disseminate the
coarse cell, apply (ST2) on the fine board, and take the coarse root.
Choose the larger cell as a sufficiently large allowed integral power
of the fine side, so that both boards are compatible along the chosen
torus sequence. This proves (ST3). The lower limit along
\emph{any} sequence is essential below.

The perturbation $\log(1+tG)$ in this proof is arbitrary and need not
be an admissible linear source. No renormalization-group bound was
applied to it or to its Gibbs states. It enters only an exact Jensen
and boundary-comparison argument and the elementary seminorm
inequality.

\subsection*{The precise ordinary torus hypothesis}

Here is one sufficient source hypothesis, stated to expose the actual
input to the transfer. Fix constants $r,A,C_\pm>0$, an entropy
rate $\kappa_0$, and $\kappa>\kappa_0$, uniformly over the selected
cutoff--coupling data $(a,\beta)$ and allowed tori. Let
$\|f\|_\sharp=\max(\|f\|_\infty,c_M\|\nabla f\|_\infty)$.
For every finite family of real Lipschitz torus functions require the
ordinary complex source bound, on
$\sum_i|z_i|\|f_i\|_\sharp<2r$,

\begin{gather*}
\FitMath{|Z_{a,\ell}(W)|\le
 \exp\left[V_{a,\ell}+\operatorname{Re}U_{a,\ell}(W)
       +\operatorname{Re}\sum_{s=0}^K\sum_D e_{a,\ell,s,D}(W)
       +C_+\ell^4\right],} \\[3pt]
\FitMath{Z_{a,\ell}(0)\ge e^{V_{a,\ell}-C_-\ell^4}.}\tag{ST5}
\end{gather*}

The terms in this expression have the following meaning and bounds.
The source-independent vacuum number $V_{a,\ell}$ is real.

The point counterterm is
$U_{a,\ell}(W)=\sum_p u_{a,\ell,p}(W(x(p)))$, where each $u$ is
holomorphic and zero at zero. For each fixed cutoff its norm on
smaller discs is bounded uniformly in $p$ and $\ell$; this bound may
depend on the cutoff. The domain term $e_{a,\ell,s,D}$ is holomorphic
in the ordinary source variables, is zero at zero, depends only on
$W|_{D^+}$, and satisfies
$|e_{a,\ell,s,D}|\le Ae^{-\kappa m(D)}$. Put $a_s=L^{-s}$ and require
$\operatorname{diam}D^+\le C_0a_s(m(D)+1)$.
The domain entropy bound, in a convenient localized form, is

\begin{gather*}
\FitMath{\sum_{D:\,D^+\cap B\ne\varnothing,\,m(D)\ge t}
                     e^{-\kappa m(D)}
 \le C(R+1)^4a_s^{-4}e^{-ct}}\tag{ST6}
\end{gather*}

for every radius-$R$ ball $B$ and $t\ge0$, uniformly in its position.
A count per unit cube bounded by $e^{\kappa_0m}$ implies this form
after absorbing a fixed polynomial in $m$ into a smaller entropy gap.
All domain sums used are normally convergent on smaller source balls.
Source representation independence is required: local terms depend on
the source function on their indicated region, not on the particular
list representing that function.

Condition (ST5) is an upper bound and a floor, not a claim that
$Z(W)$ is zero-free on the whole source ball. Its ordinary source
class includes the global real Lipschitz profiles constructed by
mirror periodization. This last inclusion matters. For
$w=\sum_i z_if_i$ in one cell, the periodized source is

\begin{gather*}
\FitMath{W^{\rm ch}=\sum_i z_iH_i^++\sum_i\bar z_iH_i^- .}\tag{ST7}
\end{gather*}

The even and odd copy profiles $H_i^\pm$ are real. If each support is
at distance $b> a$ from the faces, the profiles vanish near the faces
and their Lipschitz norms are at most $\|f_i\|_\sharp$. Hence
$\|W^{\rm ch}\|_\flat\le2\sum_i|z_i|\|f_i\|_\sharp$, independently
of the number of cells. The ordinary radius $2r$ in (ST5) therefore
supplies the radius $r$ used below. One first regards the two lists of
coefficients as independent complex variables; evaluating one at the
conjugate of the other does not require a new analytic extension.

\subsection*{Holomorphic cell normalizers from the ordinary hypothesis}

Suppose $m$ radius-$R$ support balls lie at distance at least a fixed
$b>1$ from the faces of a cell $C$ of side $l$. For each torus form a
normalizer $E_{C,\ell}(z)$ as follows. Retain in $U+\sum e$ only the
point terms and domains whose complete dependence region lies in one
open board cell. Evaluate each retained term at the source copy in
that cell, conjugate its value if that cell has odd parity, sum over
cells, and divide by $N$. The result is holomorphic in $z$, even
though (ST7) itself involves $\bar z$: on an odd cell the retained
term has the form $g(\bar z)$, and
$\overline{g(\bar z)}$ is holomorphic. Its real part is unchanged by
this outer conjugation. Also $E_{C,\ell}(0)=0$.

Any discarded source-dependent domain meets a support ball and
reaches a cell face. Its diameter is at least $b$, so
$m(D)\ge b/(C_0a_s)-1$. Counting once per ball, (ST6) shows that the
discarded contribution divided by $N$ is at most

\begin{gather*}
\FitMath{B_{\rm face}(m,R,b)
 \le CAm(R+1)^4\sum_{s\ge0}a_s^{-4}e^{-c'b/a_s}<\infty.}\tag{ST8}
\end{gather*}

The sum is independent of the cutoff. Overcounting domains meeting
several balls only enlarges this upper bound. Point counterterms
whose source is nonzero lie strictly inside cells, so none are
discarded. Combining (ST1), (ST5) and (ST8) proves

\begin{gather*}
\FitMath{p_{l,\ell}(e^{\mathcal O_a(w)})
 \le\exp[\operatorname{Re}E_{C,\ell}(z)+B_C],} \\[3pt]
\FitMath{B_C=(C_++C_-)l^4+B_{\rm face}(m,R,b).}\tag{ST9}
\end{gather*}


At a fixed cutoff the normalizers are locally bounded uniformly in
$\ell$. For the point terms this follows from their fixed-cutoff
bound and $O(l^4a^{-4})$ plaquettes per cell. For the domain terms use
(ST6) with $t=0$ and the finite sum of $a_s^{-4}$ over $s\le K$.
These bounds may grow with $K$, which is harmless for the next step.
Montel's theorem therefore gives a subsequence on which
$E_{C,\ell}$ converges locally uniformly to a holomorphic $E_C$.
Explicitly, Cauchy estimates on smaller source balls give
equicontinuity there; Arzela--Ascoli and a diagonal selection over a
countable compact exhaustion give the subsequence, and Cauchy's
formula proves holomorphy of its limit.
For finitely many occupied cells take one common subsequence. It is
chosen from the torus normalizers alone, independently of every
infinite-volume state.

There is a second uniform bound that does not grow with $K$. If
$n\ge2$ sources in the same cell have pairwise support distance at
least $d$, a point counterterm has zero mixed derivative, because a
point cannot belong to two of the supports. A contributing domain
must meet all supports; in particular its diameter is at least $d$.
Cauchy's formula and (ST6), with the first support ball as anchor,
give

\begin{gather*}
\FitMath{|\partial_1\cdots\partial_nE_C(0)|
 \le D_n(d;R)\prod_i\|f_i\|_\sharp,} \\[3pt]
\FitMath{D_n(d;R)=CA(R+1)^4\max(1,d^{-4})(n/r)^n.}\tag{ST10}
\end{gather*}

Indeed the remaining scale sum is
$\sum_sa_s^{-4}e^{-c'd/a_s}\le C\max(1,d^{-4})$.
This follows by splitting at $a_s\approx d$; the tail becomes a
convergent sum $Cd^{-4}\sum_{j\ge0}L^{4j}e^{-c''L^j}$.
The same estimate holds before the Montel limit, and hence after it.

\subsection*{All-state generating functions, cumulants and mean differences}

\begin{theorem}[Ordinary torus source estimates imply all-state bounds]
\label{thm:all-state-source-transfer}
Assume the ordinary torus source hypothesis (ST5)--(ST6), including
the stated locality, point counterterms and uniformity, and the proved
lattice symmetries. Then (ST13), (ST14) and (ST15) hold for every
periodic Wilson DLR state, uniformly in the cutoff, the states and the
support positions. Their constants depend only on the displayed
source constants, the number of entries, the support radius and,
for higher cumulants, the support separation.
\end{theorem}

Let $f_1,\ldots,f_n$ have radius-$R$ supports at arbitrary positions.
A zero test makes every corresponding multilinear assertion trivial,
so such tests may be omitted.
Choose an allowed side $l\le C_Ln(R+b+1)$ and a board offset that
puts each support ball at distance at least $b$ from every face.
For each coordinate, the forbidden offsets occupy at most
$2n(R+b+a)$ of a period; choosing $l$ larger than a fixed multiple of
this number leaves an allowed link-plane offset. Thus at most $n$
cells are occupied. Put $w_C=\sum_{i:\,\operatorname{supp}f_i\subset C}z_if_i$.
Apply (ST3) to $F_C=e^{\mathcal O_a(w_C)}$ along the common Montel
subsequence in (ST9). For every periodic DLR state $\rho$ this gives

\begin{gather*}
\FitMath{\left|e^{-E(z)}\rho\left(e^{\sum_i z_i\mathcal O_a(f_i)}\right)\right|
 \le Q,} \\[3pt]
\FitMath{E=\sum_CE_C,} \\[3pt]
\FitMath{Q=e^{\sum_CB_C},} \\[3pt]
\FitMath{\sum_i|z_i|\|f_i\|_\sharp<r.}\tag{ST11}
\end{gather*}

The same $E$ and $Q$ serve all states at the datum. The number $Q$
depends on $n,R$ and the fixed source constants, but not on the
cutoff, state, or support positions. This conclusion uses the lower
limit along the selected sequence, rather than an upper limit over
all tori.

Write the normalized holomorphic function in (ST11) as $h_\rho$.
It has $h_\rho(0)=1$ and $|h_\rho|\le Q$. A logarithm is needed only
near zero, where it exists by continuity. For any nonempty subset
$B\subset\{1,\ldots,n\}$, Cauchy's formula, with all other variables
zero, gives

\begin{gather*}
\FitMath{|\partial_Bh_\rho(0)|
 \le Q(|B|/r)^{|B|}\prod_{i\in B}\|f_i\|_\sharp.}
\end{gather*}

The finite derivative identity for a logarithm now implies

\begin{gather*}
\FitMath{|\partial_1\cdots\partial_n\log h_\rho(0)|
 \le r^{-n}\!\sum_{\pi}(|\pi|-1)!Q^{|\pi|}
                 \prod_{B\in\pi}|B|^{|B|}
                 \prod_i\|f_i\|_\sharp.}\tag{ST12}
\end{gather*}

The sum is over set partitions. No uniform zero-free polydisc is
needed for this identity at zero.

For $n\ge2$, the cumulant is the sum of the mixed derivatives of $E$
and $\log h_\rho$. If more than one cell is occupied, the mixed
derivative of $E=\sum_CE_C$ is zero because every summand lacks a
variable. If one cell is occupied, use (ST10). Equations (ST10)--(ST12)
prove the all-state separated cumulant bound

\begin{gather*}
\FitMath{|\kappa_n^\rho(\mathcal O_a(f_1),\ldots,\mathcal O_a(f_n))|
 \le C_n(d;R)\prod_i\|f_i\|_\sharp .}\tag{ST13}
\end{gather*}

Here $C_n$ is explicit as (ST12), with the uniform choice of $Q$ above,
plus $D_n$. This proves a finite constant for each $n$; no factorial
growth assertion beyond these formulas is being made.

For one radius-$R$ test choose one cell with $l\le C_L(R+b+1)$ and
the corresponding common normalizer $E_f$. Cauchy directly gives

\begin{gather*}
\FitMath{|\rho(\mathcal O_a(f))-E_f'(0)|
       =|h_\rho'(0)|\le(Q_R/r)\|f\|_\sharp.}
\end{gather*}

Consequently, for every pair of periodic states at the same datum,

\begin{gather*}
\FitMath{|\rho(\mathcal O_a(f))-\nu(\mathcal O_a(f))|
 \le(2Q_R/r)\|f\|_\sharp .}\tag{ST14}
\end{gather*}

The common derivative $E_f'(0)$ may diverge with the cutoff. The
estimate controls its differences from the means and hence differences
between the means. The argument also covers separated reflected
copies: their support radii, distances and test norms are unchanged.

If desired the one-point constant can be chosen to grow only as
$(R+1)^4$. Choose a fixed smooth partition of unity by translates of
compactly supported functions of uniformly bounded derivative. Its
members have supports in balls of a fixed radius $R_0$, only
$C(R+1)^4$ meet a radius-$R$ support, and
$\|f\psi_j\|_\sharp\le C\|f\|_\sharp$. Apply (ST14) to these pieces
with the fixed constant $Q_{R_0}$ and add. Such a partition is obtained
by dividing a nonnegative compact bump, positive on a unit cube, by
the positive periodic sum of its integer translates.

Finally put $b_1=C_1(R)$ and $b_m=C_m(d;R)$ for $m\ge2$. The
moment--cumulant formula and invariance of higher cumulants under
constant shifts give

\begin{gather*}
\FitMath{\left|\rho\prod_{i=1}^k
     [\mathcal O_a(f_i)-\nu(\mathcal O_a(f_i))]\right|
 \le\left(\sum_{\pi\in\operatorname{Part}\{1,\ldots,k\}}
                     \prod_{B\in\pi}b_{|B|}\right)
                     \prod_i\|f_i\|_\sharp .}\tag{ST15}
\end{gather*}

Singleton blocks use (ST14); the other blocks use (ST13). This is the
centered moment estimate required when a component is centered at a
different state or mixture, with the same uniformity in both states.

Thus (ST5)--(ST6), source locality, the stated ordinary holomorphy,
and the proved lattice symmetries imply the bounds used under G4.
The transfer itself introduces no weighted analytic-extension
obligation. The analytic construction must still prove the actual
ordinary source bound, entropy gap, upper stability and floor. A
torus statement controlling only separated mixed derivatives of a
logarithm does not state the cell estimate (ST9) or a common one-point
normalizer. Those concrete conditions must be retained when replacing
an all-state input by the transfer theorem.

% END MODULE chapters/all_state_source_transfer_v3.tex


% BEGIN MODULE chapters/graded_source_bound_v4.tex
\section{Factorial order bounds from rooted domain counting}
\label{ch:graded-rooted-source}

We strengthen the spatial counting used in the ordinary-source transfer.
The resulting estimates have factorial growth in the number of entries
and a separation singularity of order proportional to that number.
Only the ordinary finite-torus source estimates are used. In particular,
neither weighted forgetting nor source bounds for perturbed Gibbs states
are added.

Retain the ordinary source upper bound and floor (ST5), the point-local
counterterms, and the proved chessboard domination (ST3). Let $L\ge3$
be odd, $a=L^{-K}$, and $a_s=L^{-s}$, $0\le s\le K$. All constants
below are uniform over the chosen cutoff--coupling data and all large
allowed tori. The physical source norm is
$\|f\|_\sharp=\max(\|f\|_\infty,c_M\|\nabla f\|_\infty)$.

\subsection*{The stronger count and its geometric origin}

The localization parameter obeys $M=L^{m_0}$ for a fixed integer
$m_0\ge0$. Use the cofinal family of original tori
$\ell=2L^m$, $m\ge m_0$. At every actual scale,
$\ell/(Ma_s)=2L^{m-m_0+s}$ is an even integer, so all localization
partitions tile the torus; no scale terms are omitted. Every fixed
power-of-$L$ board also tiles all sufficiently large tori in this
family, with even cell counts. At scale $s$ use a cubical partition
of side $\rho_s=M a_s$. A localization domain $X$ is a union of these cubes. It admits
a finite connected cubical edge tree meeting every cube of $X$, and
$m(X)$ is the minimum number of edges of such a tree. The source
dependence region is contained in a fixed number of neighbouring cube
layers around $X$. This includes source values at plaquette barycentres
whose plaquette has a corner in a one-layer enlargement of $X$,
because a bare plaquette has diameter $O(a)\le O(a_s)$.

For a fixed cube $Q$ the elementary rooted entropy estimate is

\begin{gather*}
\FitMath{\#\{X:Q\subset X,\,m(X)=m\}\le C e^{\kappa_Dm}.}\tag{GS1}
\end{gather*}

Here is a proof for the stated domain indexing. A minimizing tree with
$m$ edges has at most $m+1$ vertices. Every cube meeting an edge of
the cubical grid contains an endpoint of that edge; hence every cube
of $X$ contains a tree vertex. Choose a tree vertex at a corner of $Q$
in at most $16$ ways. Traversing each tree edge twice gives a closed
walk of length $2m$ from that corner. There are at most $8^{2m}$ such
walks. At most $16(m+1)$ cubes contain a vertex of the walk, and a
subset of these specifies $X$. Thus

\begin{gather*}
\FitMath{\#\le16\,8^{2m}2^{16(m+1)}.}
\end{gather*}

This covers all admissible domains; restricting the domain class can
only reduce the count. The same upper bound applies on the finite
cubical quotient.

Suppose $|e_{s,X}(w)|\le A e^{-\kappa m(X)}$, where
$\kappa>\kappa_D$. Summing (GS1) gives a positive entropy gap $c>0$.
If the dependence region sees a radius-$u$ ball, some cube of $X$ lies
within $C\rho_s$ of that ball. There are at most
$C(1+u/a_s)^4$ choices of this root cube. Consequently

\begin{gather*}
\FitMath{\sum_{X:\,X^+\text{ sees }B,\,m(X)\ge t}
                       e^{-\kappa m(X)}
 \le C(1+u/a_s)^4e^{-ct}.}\tag{GS2}
\end{gather*}

There is no factor $a_s^{-4}$ when $u\ll a_s$. This improvement uses
the fixed-layer dependence geometry and rooted indexing; a bound only
by $C(u+1)^4a_s^{-4}$ would not suffice for the argument below.
Also, summing over every possible root cube, with harmless overcounting,
gives the global estimate

\begin{gather*}
\FitMath{\sum_X e^{-\kappa m(X)}\le C\ell^4a_s^{-4}.}\tag{GS3}
\end{gather*}

The tree representation implies
$\operatorname{diam}X^+\le C a_s(m(X)+1)$, as needed in these estimates.

\subsection*{A polynomial bound for one finely supported cell source}

Let a board have side $\eta\le1$ and one real test $v$ supported in a
ball of radius $\delta\le\eta$, lying at distance at least $\delta$
from the faces. Assume $\delta\ge16a$. Periodize $zv$ by reflected
copies, conjugating odd copies. Its ordinary source norm is at most
$2|z|\|v\|_\sharp$. Thus ordinary source radius $2r$ allows
$|z|<r/\|v\|_\sharp$, independently of the number of cells.

Construct the interior-cell holomorphic normalizer $E_{v,\ell}(z)$
as in (ST9). The discarded terms are the source-dependent domains
whose full dependence region is not contained in one cell. Such a
domain sees a copy of the support ball and reaches a face, so
$m(X)\ge c_0\delta/a_s-C_0$. Let $N=(\ell/\eta)^4$.
The discarded sum divided by $N$ has three scale ranges.

For $a_s\ge\eta$, discard even the source restriction and use (GS3).
The cost is at most $CA(\eta/a_s)^4$ per scale and per cell; its sum
is bounded by a geometric series. For
$\delta\le a_s<\eta$, (GS2) gives a bounded cost per scale, and there
are at most $C(1+\log(\eta/\delta))$ such scales. For $a_s<\delta$,
(GS2) and the forced domain size give a cost at most

\begin{gather*}
\FitMath{CA(1+\delta/a_s)^4e^{-c'\delta/a_s}.}
\end{gather*}

The sum is bounded uniformly in $\delta$ by comparison with
$\sum_{j\ge0}(1+L^j)^4e^{-c''L^j}$. All ranges may be truncated at
the actual cutoff, which only decreases the bounds. The point-local
counterterms are retained: their nonzero sources are strictly inside
cells. Their possibly large first derivatives cause no face error.

The stability and floor cost per cell is
$(C_++C_-)\eta^4\le C_++C_-$. We obtain

\begin{gather*}
\FitMath{p_{\eta,\ell}(e^{z\mathcal O_a(v)})
 \le \exp[\operatorname{Re}E_{v,\ell}(z)]\,
             Q_0(\eta/\delta)^p}\tag{GS4}
\end{gather*}

for constants $Q_0\ge1$, $p\ge0$ independent of $a,\eta,\delta$ and
the support position. Fixed-cutoff local boundedness and Montel
selection supply a limiting holomorphic univariate normalizer $E_v$.
For any finite list of pieces use one common subsequence. It is chosen
from the torus data and is independent of the infinite-volume states.

\subsection*{Board offsets and a genuine product source polydisc}

Consider $n$ tests supported in radius-$R$ balls, with distinct supports
at distance at least $d>0$. Put $\bar d=\min(d,1)$, and choose an
allowed board side, a power of $L$, such that

\begin{gather*}
\FitMath{\frac{\bar d}{16L}<\eta\le\frac{\bar d}{16},} \\[3pt]
\FitMath{\delta=\frac{\eta}{64n}.}\tag{GS5}
\end{gather*}

First suppose $a\le\delta/16$; thus this board side is a multiple of
the mesh. A fixed smooth partition of unity at mesh proportional to
$\delta$ splits each $f_i$ into pieces $v_{i,k}$ with

\begin{gather*}
\FitMath{\#\{k:v_{i,k}\ne0\}\le C(1+R/\delta)^4,} \\[3pt]
\FitMath{\operatorname{supp}v_{i,k}\subset B_{i,k}(\delta),} \\[3pt]
\FitMath{\|v_{i,k}\|_\sharp\le C\delta^{-1}\|f_i\|_\sharp.}\tag{GS6}
\end{gather*}

For example, scale a smooth compact partition obtained from translates
of a bump divided by their positive periodic sum. Scaling its bounded
first derivatives gives the last inequality. A zero piece contributes
nothing and is omitted.

Choose one piece from each original test. In any coordinate the
offsets for which one enclosing radius-$\delta$ ball comes within
$\delta$ of a face occupy an interval of length at most $4\delta$
modulo $\eta$. On the discrete link-plane offsets, one such interval
contains at most $4\delta/a+2$ candidates. The $n$ intervals exclude
at most $n(4\delta/a+2)<\eta/a$ candidates by (GS5) and the fine-mesh
assumption. There is an allowed offset in each coordinate. Hence all
chosen balls have the required margin. Different original tests
occupy different cells: a four-dimensional side-$\eta$ cube has
Euclidean diameter $2\eta<d$, whereas their nonempty supports are
$d$ apart.

Write these chosen pieces as $v_1,\ldots,v_n$ and their different cells
as $C_1,\ldots,C_n$. By (GS4) each has a univariate normalizer $E_i$
and a bound

\begin{gather*}
\FitMath{Q_i\le Cn^p,} \\[3pt]
\FitMath{r_i=r/\|v_i\|_\sharp.}\tag{GS7}
\end{gather*}

Product domination (ST3) is applied to the separate cell factors.
It therefore gives, for every subset $S$ and every periodic state
$\rho$,

\begin{gather*}
\FitMath{\left|\exp\left[-\sum_{i\in S}E_i(z_i)\right]
      \rho\left(e^{\sum_{i\in S}z_i\mathcal O_a(v_i)}\right)\right|
 \le\prod_{i\in S}Q_i} \\[3pt]
\FitMath{(|z_i|<r_i, i\in S).}\tag{GS8}
\end{gather*}

This is a product polydisc; no bound on its combined source norm is needed. Each
chessboard factor in (ST3) uses only one periodized piece and its own
ordinary source disc. Inactive factors are exactly $1$, with
$p(1)=1$, so (GS8) contains the product over the active subset,
not the product over all $n$ indices. The same normalizers work for
every state, including a second centering state $\nu$.

\subsection*{Common centering and factorial order}

Define $X_i=\mathcal O_a(v_i)-E_i'(0)$.
Differentiating the normalized function in (GS8) once in each active
variable, and using product Cauchy, gives

\begin{gather*}
\FitMath{|\rho\prod_{i\in S}X_i|\le\prod_{i\in S}(Q_i/r_i).}\tag{GS9}
\end{gather*}

Only the first derivative of each univariate $E_i$ occurs; no bound on
its higher derivatives is required. Applying the one-variable case
to $\nu$ gives $|\nu(\mathcal O_a(v_i))-E_i'(0)|\le Q_i/r_i$.
Expanding the product centered at $\nu$ over its active subsets and
using (GS9) therefore proves
\begin{equation}
 \left|\rho\prod_{i=1}^n
       [\mathcal O_a(v_i)-\nu(\mathcal O_a(v_i))]\right|
 \le2^n\prod_i(Q_i/r_i).
 \tag{GS10}
\end{equation}
This explains both the necessity of subset bounds in (GS8) and the
role of a common normalizer. Uniform moments for auxiliary perturbed
Gibbs states are not being presumed.

Sum (GS10) over the piece choices from (GS6), using linearity of
centering. The board may depend on that finite piece choice; every
constant is independent of that choice. Since $\delta\le1$,
\begin{equation}
\FitMath{\begin{aligned}
 \left|\rho\prod_{i=1}^n
       [\mathcal O_a(f_i)-\nu(\mathcal O_a(f_i))]\right|
 &\le\big[C(R+1)^4n^p\delta^{-5}\big]^n
                           \prod_i\|f_i\|_\sharp\\
 &\le B_0(R)^n n^{(p+5)n}\bar d^{-5n}
                           \prod_i\|f_i\|_\sharp.
 \end{aligned}}
 \tag{GS11}
\end{equation}

If $a>\delta/16$, dispense with the partition and the board. A
radius-$R$ ball contains at most $C(1+R/a)^4$ plaquette barycentres.
Since $0\le a_p\le2$,

\begin{gather*}
\FitMath{\|\mathcal O_a(f)-\nu(\mathcal O_a(f))\|_\infty
 \le C(R+1)^4a^{-4}\|f\|_\sharp
 \le C(R+1)^4\delta^{-4}\|f\|_\sharp.}
\end{gather*}

Its product bound is absorbed by (GS11), because $\delta$ is
comparable to $\bar d/n$, $p\ge0$ and $\bar d\le1$.
This handles every cutoff, including the case when the prospective
side $\eta$ is smaller than the mesh.

Choose the fixed integer $k=\lceil p\rceil+5$. The elementary inequality
$n!\ge(n/e)^n$, obtained by integrating $\log x$, converts (GS11) to

\begin{gather*}
\FitMath{\left|\rho\prod_{i=1}^n
       [\mathcal O_a(f_i)-\nu(\mathcal O_a(f_i))]\right|
 \le B(R)^n(n!)^k\max(1,d^{-5n})
                           \prod_i\|f_i\|_\sharp.}\tag{GS12}
\end{gather*}


\begin{theorem}[Graded moments and cumulants]
\label{thm:graded-rooted-source}
Under the ordinary torus source hypotheses and the rooted domain
geometry stated above, (GS12) holds for every pair of periodic Wilson
DLR states at the same datum. For $n\ge2$ their cumulants satisfy

\begin{gather*}
\FitMath{|\kappa_n^\rho(\mathcal O_a(f_1),\ldots,\mathcal O_a(f_n))|
 \le [2B(R)]^n(n!)^k\max(1,d^{-5n})
                              \prod_i\|f_i\|_\sharp.}\tag{GS13}
\end{gather*}

The constants are independent of the cutoff, the states and support
positions. The exponent $k$ is one fixed finite integer.
One may take $B(R)=C(R+1)^4$ with a constant fixed by the source and
domain parameters.
\end{theorem}

To verify the cumulant assertion without an unnecessary extra
factorial, center at $\rho$ and apply the moment--cumulant formula.
For a partition with block sizes $n_1,\ldots,n_m$, use (GS12) on each
block and $\prod_jn_j!\le n!$. The remaining combinatorial bound is

\begin{gather*}
\FitMath{\sum_\pi(|\pi|-1)!\prod_{D\in\pi}(|D|!)^k
 \le(n!)^{k-1}\sum_\pi(|\pi|-1)!\prod_{D\in\pi}|D|!
 \le2^{n-1}(n!)^k.}
\end{gather*}

For exactly $m$ blocks, the middle sum equals
$\frac{n!}{m}\binom{n-1}{m-1}$: order the blocks, sum over positive
compositions of $n$, and cancel their factorial denominators. Summing
over $m$ proves the inequality and (GS13).

% Keep the closing statement and argument together.
\clearpage
The separation singularity in this theorem is $d^{-5n}$, not a single
$d^{-5}$ prefactor independent of $n$. This is a factorial-growth
estimate at each fixed separation, with diagonal order proportional
to $n$. It is sufficient for extensions whose derivative order may
grow linearly with the number of entries. A stronger separation law
requires an additional argument. The ordinary analytic source and stability hypotheses are those
stated in D1.

One may also recover a presentation with a separation prefactor independent
of $n$, at the cost of one factorial power. Put $x=\max(1,d^{-5})$.
The exponential series gives $x^n\le n!e^x$, so (GS12) implies

\begin{gather*}
\FitMath{\left|\rho\prod_{i=1}^n
       [\mathcal O_a(f_i)-\nu(\mathcal O_a(f_i))]\right|
 \le e^{\max(1,d^{-5})}B(R)^n(n!)^{k+1}
                            \prod_i\|f_i\|_\sharp.}\tag{GS14}
\end{gather*}

Thus $A(d;R)=e^{\max(1,d^{-5})}$ is one possible prefactor in that
notation. Its separation behavior is much worse than a polynomial.
The quantitative extension uses the sharper graded form (GS12), rather
than attempting to infer diagonal order from this rewritten bound.

% END MODULE chapters/graded_source_bound_v4.tex


% BEGIN MODULE chapters/continuum_compactness_v5.tex
\section{Compactness on separated tests and the lattice Gibbs limit}
\label{ch:continuum-compactness}

The conclusion of this section is subsequential compactness. A sequence
of cutoff data with uniform source estimates is an input. Existence of
Wilson Gibbs states at each fixed datum is proved below, but does not
produce the cutoff data or prove those uniform estimates. No matching
of different cutoffs, tuning, or uniqueness of their continuum limits
is asserted.

\subsection*{The fixed-cutoff Gibbs measure argument}

Fix a spacing $a>0$, a finite inverse coupling $\beta\ge0$, and the
compact group $G=\mathrm{SU}(2)$. Choose one orientation of each nearest
neighbor bond of $a\mathbb Z^4$ and denote this countable set by $E$.
A configuration is $U\in\mathcal X=G^E$; the opposite bond has matrix
$U_e^{-1}$. The product topology is compact and metrizable. For example,
enumerating the bonds gives a metric
$\sum_{r\ge1}2^{-r}\min(1,\|U_{e_r}-V_{e_r}\|)$.
Continuous functions depending on finitely many bonds are uniformly
dense in $C(\mathcal X)$. Indeed, replacing all coordinates after the
first $N$ by a fixed configuration changes the metric by at most
$2^{-N}$, and uniform continuity proves this approximation.

Let $A_p(U)=1-\frac12\operatorname{Re}\operatorname{Tr}U_p\in[0,2]$
be the Wilson plaquette action. For a finite bond window $W$, let
$\mathcal P(W)$ be the plaquettes meeting $W$ and let $\partial W$
be their bonds outside $W$. Define
\[
 \gamma_W F(U)=\frac{\int_{G^W}F(VU)
       \exp[-\beta\sum_{p\in\mathcal P(W)}A_p(VU)]\,dV}
      {\int_{G^W}\exp[-\beta\sum_{p\in\mathcal P(W)}A_p(VU)]\,dV}.
 \tag{CC1}\label{cc:specification}
\]
Here $VU$ replaces the variables in $W$, and Haar measure is
normalized. The denominator lies between
$e^{-2\beta|\mathcal P(W)|}$ and one. The normalized density is
therefore a continuous function of the finite variables
$(V,U_{\partial W})$. If $F$ is continuous local, $\gamma_WF$ is
continuous local, depending only on the remaining variables of $F$
and on $\partial W$. This proves the Feller property directly;
integration of a uniformly continuous function over the compact
finite product preserves continuity. By uniform approximation it
also holds for $F\in C(\mathcal X)$.

Take Wilson probability measures on periodic tori with integer periods
$N\to\infty$, and lift them periodically to $\mathcal X$. The finite
partition function is finite and strictly positive because the
integrand is continuous and positive on a compact space. The lifted
measures have a weakly convergent subsequence. One direct compactness
argument is as follows. The space $C(\mathcal X)$ is separable: on a
compact metric space finite approximations by distance functions with
rational coefficients form a countable uniformly dense family.
Successively extract convergence of the integrals on this family.
Uniform approximation then defines a bounded positive functional on
all of $C(\mathcal X)$, with value one at the constant one. Its Riesz
probability measure is the weak limit. This argument simultaneously
gives convergence against every continuous function, not merely the
chosen countable family.

The lifted torus measure is not itself an infinite-volume DLR state:
remote periodic copies can determine an inside variable. The identity
needed for passage to the limit is a local one. Fix $W$ and a
continuous local $F$. For all sufficiently large periods, the finite
bonds involved in $W$, $\partial W$ and $F$ embed injectively into
the torus, and the plaquettes meeting $W$ are precisely their usual
torus counterparts. Integrating the variables in that torus window
first gives exactly
\[
       \widetilde\mu_N(\gamma_WF)=\widetilde\mu_N(F).
 \tag{CC2}\label{cc:torus-local-DLR}
\]
Both integrands are continuous local. Their weak limits satisfy
$\mu(\gamma_WF)=\mu(F)$. Cylinder approximation and
$\|\gamma_WF\|_\infty\le\|F\|_\infty$ extend this identity to all
continuous $F$. Equivalently the Borel probability measures
$\mu\gamma_W$ and $\mu$ agree. The identity consequently holds for
every bounded Borel $F$.

It is also the conditional-law formulation of DLR. For bounded
outside-measurable $g$, the kernel is proper and satisfies
$\gamma_W(gF)=g\gamma_WF$. Hence
$\mu(gF)=\mu(g\gamma_WF)$, and $\gamma_WF$ is outside measurable.
Thus \eqref{cc:specification} is the conditional expectation given
the bonds outside $W$. This holds for every finite window.

The torus lifts are invariant under lattice translations and the
hypercubic group. These transformations preserve Haar measure and the
Wilson action, including when a chosen bond orientation is reversed.
Weak convergence preserves their invariance, so $\mu$ is an invariant
Wilson DLR state. In particular it is periodic, with the whole lattice
as a period group. The same Feller argument makes the class of
translation-invariant DLR states weakly closed, hence compact, and
its defining identities make it convex. No compactness of the union
over arbitrary period groups is being asserted.

If a source argument uses a prescribed compatible cofinal family of
tori, take that family in this proof. No measure estimate beyond
compactness is needed. Reflection positivity and the stronger symmetry
theorems for every periodic DLR state are separate results; this
measure argument proves neither their all-state estimates nor any
uniformity as $a\downarrow0$.

\subsection*{The continuum hypothesis}

Fix an odd $L\ge3$ and a sequence $a_j=L^{-K_j}\downarrow0$.
At the same datum $(a_j,\beta_j)$ let $\rho_j,\nu_j$ be periodic
Wilson DLR states. The all-mirror lattice theorem supplies the following
properties used here: both states are invariant under unit lattice
translations and the hypercubic group, and $\rho_j$ is reflection
positive at coordinate and diagonal lattice mirrors. The centering
state $\nu_j$ is invariant at those mirrors as well. These properties
are lattice inputs to the compactness argument.

These are limits through states on an infinite lattice. They include
volume-first limits when the states are torus limit points. A direct
limit on growing tori, or a spacing-first limit on fixed tori followed
by removal of volume, is a different construction until a separate
recognition theorem identifies it with this class. In particular this
section does not assign diagonal reflection positivity to an arbitrary
growing cubic-torus limit. Tuned fixed-torus limits, when obtained by
matched-cutoff estimates, remain separate from the subsequence theorem.

For $f\in C_c^1(\mathbb R^4)$ write
\[
 \mathcal O_j(f)=\sum_p f(x_p)A_p,
 \qquad X_j(f)=\mathcal O_j(f)-\nu_j(\mathcal O_j(f)),
 \qquad m_{j,n}(\mathbf f)=\rho_j\prod_{i=1}^n X_j(f_i).
 \tag{CC3}\label{cc:moments}
\]
There is one plaquette for each coordinate plane and lattice site,
so six orientations are counted once each; $x_p$ is its barycenter.
There is no power of $a_j$ in the observable. At a fixed cutoff these
are bounded continuous random variables with finite bond dependence.
Set $m_{j,0}=1$ and complexify multilinearly.

Let $\mathfrak S_n$ be the tuples of compact $C^1$ tests with pairwise
disjoint supports. A zero entry gives value zero and can be set aside.
The separation is the minimum distance of the supports, not that of
balls containing them. Every support may lie in its own radius-$R$
ball; these balls may overlap. Use the fixed physical norm
\[
 \|f\|_\sharp=\max(\|f\|_\infty,c_M\|\nabla f\|_\infty),
                      \qquad c_M>0\hbox{ fixed}.
\]
The ordinary source hypothesis is used only through the following
proved implications of it. At each fixed order there is a bound
\[
 |m_{j,n}(\mathbf f)|\le C_n(d;R)\prod_i\|f_i\|_\sharp
 \tag{CC4}\label{cc:local-bound}
\]
when the supports have distance at least $d>0$ and the stated radius
bound. Bounds holding eventually on each fixed support neighborhood
suffice for extraction. The one-point version is included: for
different centering and integrating states it is a uniform bound on
their mean difference, not on either unsubtracted mean.

For the quantitative conclusion assume the graded bound, with one
finite integer $k\ge1$ and $B(R)<\infty$ independent of $j$ and the
two states,
\[
 |m_{j,n}(\mathbf f)|
  \le B(R)^n(n!)^k\max(1,d^{-5n})\prod_i\|f_i\|_\sharp.
 \tag{CC5}\label{cc:graded-bound}
\]
For $n=1$ take $d=1$. A harmless change of $B$ includes complex
entries. These are all-state bounds at the indicated data and also
apply to the pair $(\rho_j,\rho_j)$ used for canonical recentering.
The ordinary-source theorem obtains this bound from D1 and
its actual rooted domain geometry, including the fixed-layer
dependence and strict entropy gap. A weaker counting assertion alone
does not supply this hypothesis. Those source hypotheses, and
existence of the given sequence of cutoff data with common constants,
are not proved by this section.

\subsection*{A countable support-adapted family}

Let $\mathcal U$ be the bounded finite unions of open boxes with
rational endpoints. This is countable. Every finite collection of
separated compact sets $K_i$ has neighborhoods $U_i\in\mathcal U$
such that $K_i\subset U_i$ and the closures $\overline U_i$ are
pairwise separated. Cover each compact set by sufficiently small
rational boxes inside a neighborhood of radius less than one fourth
of its distance from the other sets, then take a finite subcover.
The neighborhoods can also stay in any prescribed open support
regions.

For each $U\in\mathcal U$ we now exhibit a countable family dense
in $C_c^1(U)$ for the $C^1$ norm, with every approximant supported
in $U$. Fix a nonnegative smooth unit-mass bump $\varrho$ in the unit ball.
For integers $m\ge1$ let
\[
 \chi_{U,m}=
  \mathbf1_{\{x:\operatorname{dist}(x,U^c)>1/(2m)\}}
                                      *\varrho_{1/(8m)}.
\]
It is smooth, compactly supported in $U$, takes values in $[0,1]$,
and equals one where $\operatorname{dist}(x,U^c)\ge1/m$.
Use finite rational linear combinations of
\[
 \chi_{U,m}(x)\cos(2\pi q\cdot x/N),\qquad
 \chi_{U,m}(x)\sin(2\pi q\cdot x/N),
                   \quad q\in\mathbb Z^4,\ N\in\mathbb N,
 \tag{CC6}\label{cc:dense-functions}
\]
and take their rational span over all indices. Call this real
countable family $D_U$.

To verify density, fix $f\in C_c^1(U)$ and choose $m$ such that
$\chi_{U,m}=1$ on a neighborhood of its support. Extend $f$ by
zero, smooth it on a sufficiently small scale, and view the resulting
smooth compact function as periodic on a large integer-side box.
Smoothing converges to $f$ in $C^1$ because $f$ and its gradient are
uniformly continuous. The Fourier partial sums of this smooth
periodic function converge with their first derivatives: repeated
integration by parts makes the differentiated Fourier series
absolutely summable. Approximate their finitely many real sine and
cosine coefficients by rationals and multiply by $\chi_{U,m}$.
The multiplication is bounded in $C^1$, and $\chi_{U,m}f=f$.
This produces the required approximants. Complexification gives the
countable dense family $D_U+iD_U$ for complex tests.

List all tuples $(g_1,\ldots,g_n)$, for every $n\ge1$ and every
ordered list of pairwise separated $\overline U_i$, with
$g_i\in D_{U_i}+iD_{U_i}$. This is one countable list. By
\eqref{cc:local-bound} the values of each coordinate form a bounded
sequence; finite initial values do not matter if a bound starts
eventually. Extract an infinite subsequence for convergence of the
first coordinate, then a subsequence of it for the second, and so
on. If $I_r$ is the $r$th nested infinite index set, choose increasing
$j_r\in I_r$. For each fixed coordinate its values along $(j_r)$
eventually lie in the convergent subsequence for that coordinate.
All listed values therefore converge along this single subsequence.

\subsection*{Extension to every separated tuple}

The following telescoping estimate explains why that one extraction
is enough. If $\operatorname{supp}f_i\cup\operatorname{supp}g_i$
lies in $\overline U_i$, let $d_U>0$ be a lower bound for their
pairwise distances, taking $d_U=1$ when $n=1$, and choose a common
finite radius bound $R_U$. Multilinearity
and \eqref{cc:local-bound} give, uniformly in the relevant tail,
\[
 \begin{split}
 |m_{j,n}(\mathbf f)-m_{j,n}(\mathbf g)|
 \le C_n(d_U;R_U)\sum_{i=1}^n
   \|f_i-g_i\|_\sharp
   \prod_{h\ne i}\max(\|f_h\|_\sharp,\|g_h\|_\sharp).
 \end{split}
 \tag{CC7}\label{cc:telescoping}
\]
Every mixed tuple in this telescoping sum is separated because the
neighborhoods, rather than just the two endpoint tuples, carry the
separation.

Given any $\mathbf f\in\mathfrak S_n$, choose such rational
neighborhoods and approximate each entry in $C^1$ by the corresponding
countable family. The right side of \eqref{cc:telescoping} tends
to zero uniformly in $j$. For a fixed approximating tuple its values
converge along $(j_r)$. Comparing two indices through this tuple
shows that $m_{j_r,n}(\mathbf f)$ is Cauchy. Define its limit by
$\sigma_n(\mathbf f)$. This definition is independent of all
neighborhoods and approximants because it is the limit of the
original sequence at that tuple. The subsequence has not been chosen
again for $\mathbf f$; convergence now holds on every separated
compact $C^1$ tuple simultaneously. Set $\sigma_0=1$.

Passing in \eqref{cc:telescoping} proves the same local continuity
estimate for $\sigma_n$. On compact subsets of each product of the
fixed support spaces, convergence is uniform: choose a finite
$C^1$ net in that compact set, apply pointwise convergence at its
centers, and use \eqref{cc:telescoping} at both ends. Its constants
are uniform on a norm-bounded set.

Multilinearity passes whenever all tuples in the proposed identity
are admissible. In particular the union of the alternative supports
in a single entry must stay separated from the other entries.
Permutation symmetry, conjugation and reality pass directly from
their exact lattice identities. The bounds
\eqref{cc:local-bound} and \eqref{cc:graded-bound} pass with their
original constants at every fixed tuple, since absolute-value
inequalities are closed under convergence. Their radius and
separation parameters concern the original supports, not the larger
neighborhoods used to construct the limit.

\subsection*{Moving arguments, real symmetries, and rounded mirrors}

Suppose $f_{i,j}\to f_i$ in $C^1$ and all their supports lie in
fixed pairwise separated compact neighborhoods. Equation
\eqref{cc:telescoping} and fixed-tuple convergence give
\[
 m_{j_r,n}(f_{1,j_r},\ldots,f_{n,j_r})
                   \longrightarrow\sigma_n(f_1,\ldots,f_n).
 \tag{CC8}\label{cc:moving-tests}
\]
This is also valid for every finite combination of tuples.
For translations the needed $C^1$ convergence has the explicit
modulus
\[
 \|\tau_v f-\tau_w f\|_\sharp
 \le\max\{\|\nabla f\|_\infty|v-w|,
                       c_M\omega_{\nabla f}(|v-w|)\}\longrightarrow0,
 \tag{CC9}\label{cc:translation-modulus}
\]
where $\omega_{\nabla f}$ is the uniform modulus of continuity of
the compactly supported continuous gradient. No second derivative
is required. The chain rule gives the corresponding convergence
for a convergent sequence of Euclidean isometries.

For any fixed real translation $v$ choose lattice vectors
$v_j\in a_j\mathbb Z^4$ with $v_j\to v$. Both states are invariant
under $v_j$, so $m_{j,n}(\tau_{v_j}\mathbf f)=m_{j,n}(\mathbf f)$.
The moving-test statement proves real translation invariance of
$\sigma_n$. Hypercubic invariance is exact at every spacing and
passes without rounding.

For clarity the reflection limit also works at arbitrary affine
positions. A coordinate mirror $x_\mu=c$ can be approximated by link
mirrors with $c_j\in a_j(\mathbb Z+1/2)$ and $c_j\to c$.
A diagonal mirror $x_\mu\pm x_\lambda=c$ is approximated by
$c_j\in a_j\mathbb Z$. In unit-normal coordinates its offset grid
is $a_j\mathbb Z/\sqrt2$. These are different conventions for
the same diagonal mirror; no factor $\sqrt2$ is discarded.
The reflected tests converge in $C^1$: the linear reflection part
is fixed and the translation parts differ by $2(c_j-c)e_\mu$
in the coordinate case, and by
$(c_j-c)(e_\mu\pm e_\lambda)$ in the diagonal case.

Let $F$ be a finite complex combination of separated tuples whose
entries lie strictly on one open side of the limiting mirror.
Its compact support has positive distance from the mirror. For all
fine cutoffs the barycenters sampled by those entries and all the
bonds of their plaquettes lie on the same side of the rounded
mirror: a plaquette has diameter $O(a_j)$, while that distance is
fixed. Form its centered lattice polynomial $F_j$.
Reflection invariance of $\nu_j$ makes centering commute with the
lattice reflection. Lattice reflection positivity gives
\[
       0\le\rho_j(\overline{\theta_jF_j}F_j)
              =m_j(\theta_j\overline F\cdot F).
\]
Every tuple on the right has a common positive support separation
for large $j$: internal separations are fixed and the two sides
are separated by the mirror gap. Apply \eqref{cc:moving-tests}
to get
\[
              \sigma(\theta\overline F\cdot F)\ge0.
 \tag{CC10}\label{cc:reflection-limit}
\]
Reflection invariance itself passes in the same way. This proves
coordinate and diagonal reflection positivity, including their
affine translates. It proves no new reflection direction.

\subsection*{The one-point limit and the connected functions}

The first moment is determined more precisely by the plaquette
means. Unit translations and hypercubic symmetry make
\[
 \delta_j=\rho_j(A_p)-\nu_j(A_p)
\]
independent of the plaquette position and its coordinate plane.
Define
\[
 R_j(f)=\frac{a_j^4}{6}\sum_p f(x_p),\qquad
                         c_j=6a_j^{-4}\delta_j.
 \tag{CC11}\label{cc:Riemann-mean}
\]
Then $m_{j,1}(f)=c_jR_j(f)$. Each orientation samples a translate
of $a_j\mathbb Z^4$. Comparing its sample value to the integral
on the corresponding side-$a_j$ cube gives
\[
 R_j(f)\longrightarrow\int f,
 \qquad |R_j(f)-\textstyle\int f|
             \le C_R a_j\|\nabla f\|_\infty
 \tag{CC12}\label{cc:Riemann-limit}
\]
for a test supported in a fixed radius-$R$ ball. Indeed the sample
and any point of its cube are at distance at most $2a_j$; only
$C(R+1)^4a_j^{-4}$ cubes can contribute. Average the six such
estimates. The same bound is independent of the center of that
ball.

Choose one real smooth compact bump $f_*$ with $\int f_*=1$.
Its limiting moment exists by the preceding extraction, and
$R_j(f_*)\to1$. Consequently along the extracted subsequence
\[
 c_j=\frac{m_{j,1}(f_*)}{R_j(f_*)}\longrightarrow c\in\mathbb R,
             \qquad \sigma_1(f)=c\int f.
 \tag{CC13}\label{cc:first-moment}
\]
In particular the coefficients are bounded on that tail. This
bound is uniform over both states:
$|c_j|\le2B(R_*)\|f_*\|_\sharp$ once $R_j(f_*)\ge1/2$,
and $|c|\le B(R_*)\|f_*\|_\sharp$.
This
argument uses the uniform one-point comparison and the common
plaquette means; it is not a deduction from formal continuum
translation invariance alone. If $\rho_j=\nu_j$, then $c=0$.

All connected limits are finite algebraic consequences, requiring
no common disc of analyticity for an exponential. For a fixed
$n$ work in the algebra with $z_i^2=0$ and let
\[
 M_j(z)=1+\sum_{\varnothing\ne A\subset[n]}
                    m_{j,|A|}(\mathbf f_A)\prod_{i\in A}z_i.
\]
The logarithm is the finite series for $\log(1+(M_j-1))$.
Its coefficient of $\prod_{i=1}^nz_i$ is
\[
 K_{j,n}(\mathbf f)=
 \sum_{\pi\in\operatorname{Part}[n]}
       (-1)^{|\pi|-1}(|\pi|-1)!
                            \prod_{A\in\pi}m_{j,|A|}(\mathbf f_A).
 \tag{CC14}\label{cc:cumulant}
\]
In fact the term of degree $r$ in the logarithm counts the $r!$
orders of each partition with $r$ blocks and divides by $r$,
giving $(r-1)!$. Every subtuple is separated. Hence these finite
polynomials converge to the same formula with $\sigma$ in place
of $m_j$; denote it by $S_n$. Conversely the finite exponential
series gives
$\sigma_n=\sum_\pi\prod_{A\in\pi}S_{|A|}$, since its factor
$1/r!$ cancels the orders of the blocks.

Adding constants to the $X_j(f_i)$ multiplies $M_j$ by
$\exp(\sum_i b_i z_i)$, so changes its logarithm only in degree
one. Thus for $n\ge2$, $K_{j,n}$ is the usual cumulant of the
raw variables $\mathcal O_j(f_i)$ under $\rho_j$, independent of
the centering state. The graded cumulant bound passes as well:
\[
 |S_n(\mathbf f)|\le[2B(R)]^n(n!)^k
                  \max(1,d^{-5n})\prod_i\|f_i\|_\sharp,
                         \qquad n\ge2.
 \tag{CC15}\label{cc:connected-bound}
\]
One may obtain it either from the uniform lattice cumulant bound,
or directly from \eqref{cc:cumulant}: if block sizes are $n_h$,
$\prod_h n_h!\le n!$, and
\[
 \sum_\pi(|\pi|-1)!\prod_{A\in\pi}(|A|!)^k
       \le(n!)^{k-1}\sum_\pi(|\pi|-1)!\prod_{A\in\pi}|A|!
       \le2^{n-1}(n!)^k.
\]
For exactly $r$ blocks the last unweighted sum is
$\frac{n!}{r}\binom{n-1}{r-1}$, by ordered compositions of $n$.

There is a canonical centered functional distinct from the original
two-state moments. Center at $\rho_j$ instead of $\nu_j$; expanding
each product and using convergence of $m_{j,1}$ gives along the
same subsequence
\[
 \sigma^0_n(\mathbf f)=
  \sum_{A\subset[n]}\sigma_{|A|}(\mathbf f_A)
                  \prod_{i\notin A}\bigl(-c\textstyle\int f_i\bigr)
  =\sum_{\pi\in\operatorname{Part}_{\ge2}[n]}
                                      \prod_{A\in\pi}S_{|A|}.
 \tag{CC16}\label{cc:centered-functional}
\]
It has empty value one, singleton zero, the same higher connected
functions, and all the bounds and positivity properties above.
For the bounds use the lattice pair $(\rho_j,\rho_j)$ before
passing to the limit, so no uncontrolled constant from recentering
is introduced.

\subsection*{The quantitative flat extension and closed inequalities}

For completeness the topology retained by this compactness statement
is explicit. Let $\Delta_n=\bigcup_{i<h}\{x_i=x_h\}$ and let
$\mathcal S_{\Delta,n}$ be the Schwartz functions whose every jet
vanishes on $\Delta_n$. Use
\[
 p_s(\Phi)=\max_{|\alpha|\le s}\sup_x
          (1+|x|^2)^{s/2}|\partial^\alpha\Phi(x)|.
\]
On unit-radius supports, \eqref{cc:graded-bound} and
$\|f\|_\sharp\le\max(1,c_M)\max(\|f\|_\infty,
\|\nabla f\|_\infty)$ give the hypothesis of the proved
quantitative flat-extension theorem, with separation exponent five.
That theorem supplies unique extensions $\Sigma_{j,n}$ and
$\Sigma_n$ of the lattice and limiting compact forms and a common
bound
\[
 |\Sigma_{j,n}(\Phi)|+|\Sigma_n(\Phi)|
            \le 2C^n(n!)^{k+32}p_{20n}(\Phi).
 \tag{CC17}\label{cc:flat-bound}
\]
Its separated smooth tensors are dense in this Schwartz subspace.
The same conclusions apply to $\sigma^0$ and, after changing $C$,
to the connected forms. The theorem is an extension on this
specified flat space, not a claim of a canonical extension across
the diagonals on all Schwartz functions.

There is convergence on the whole flat space along the already
chosen subsequence. Given $\Phi$, approximate it by a finite sum
$T$ of smooth separated tensors in the seminorm $p_{20n}$.
Equation \eqref{cc:flat-bound} controls both
$\Sigma_{j,n}(\Phi-T)$ and $\Sigma_n(\Phi-T)$ uniformly in $j$.
On $T$ convergence has already been proved. Taking first the cutoff
limit and then the approximation error to zero proves
\[
                \Sigma_{j_r,n}(\Phi)\longrightarrow\Sigma_n(\Phi).
 \tag{CC18}\label{cc:flat-convergence}
\]
The positive-time version of the density theorem uses tensors
strictly above the mirror. Their reflected products are flat on
all cross diagonals, including the time boundary. Thus the limit's
symmetry, reality, real translations and reflection positivity extend by this
density and the continuity of tensor products in Schwartz
seminorms. For a coordinate or diagonal mirror the corresponding
positive half-space is used. It follows by an orthogonal change
of variables in the flat-space density construction; it does not
assume invariance of the functional under that change of variables.

More generally any closed inequality involving finitely many fixed
tuple evaluations and continuous functions of them passes to the
limit. The statement allows moving tests as in
\eqref{cc:moving-tests}, convergent scalar coefficients and
eventual uniform bounds. Its proof is simply convergence of those
finitely many numbers followed by continuity and closedness of the
inequality. It does not turn strict inequalities into uniform
strict bounds or provide estimates whose constants diverge with
the cutoff.

\begin{theorem}[Subsequential continuum compactness]
\label{thm:continuum-compactness}
For every input sequence satisfying the lattice properties and
two-state uniform bounds above, there is one subsequence whose
moments converge on every separated compact $C^1$ tuple. The
limiting functional is admissibly multilinear, symmetric and real,
invariant under real translations and the hypercubic group, and
reflection positive at every coordinate and diagonal affine mirror.
It retains the graded bound and its quantitative diagonal-flat
extension. Its first moment is $c\int f$, its higher connected
functions are the limits of the integrating states' cumulants,
and \eqref{cc:centered-functional} is their canonical centered
functional. Moving arguments and fixed-tuple closed inequalities
pass as stated above.
\end{theorem}

The preceding sections prove every assertion. The fixed-cutoff
measure argument ensures that invariant Wilson DLR states are
available at each given datum. It supplies neither a sequence of
admissible renormalization data with $K_j\to\infty$ nor D1's
uniform source bounds. It also supplies no identification of two
different extracted limits, no matched-cutoff convergence rate,
and no claim of a full Yang--Mills observable theory.

% END MODULE chapters/continuum_compactness_v5.tex


% BEGIN MODULE chapters/quantitative_flat_extension_v4.tex
\section{A quantitative extension to diagonal-flat Schwartz tests}
\label{ch:quantitative-flat-extension}

Assume the graded compact separated-moment bound below. We construct
the diagonal-flat extension and control its constants and derivative
orders in the number of arguments. In the gauge-theory application
the graded-source theorem supplies this bound under D1 on its actual GS geometry and G1.
The proof uses explicit Fourier series on product boxes to retain
that quantitative control.

\subsection{The hypothesis and the quantitative conclusion}

Use $\|f\|_1=\max(\|f\|_\infty,\|\nabla f\|_\infty)$ on
$C_c^1(\mathbb R^4;\mathbb C)$. Suppose $S_n$, $n\ge1$, are complex
multilinear forms on tuples with pairwise separated compact supports.
Multilinearity is required whenever the supports in the identity remain
separated. Fix numbers $1\le b<\infty$, $\kappa\ge0$ and $\alpha\ge0$,
independent of $n$, such that

\begin{gather*}
\FitMath{|S_n(f_1,\ldots,f_n)|
 \le b^n(n!)^\kappa\max(1,d^{-\alpha n})
                         \prod_i\|f_i\|_1}\tag{QF1}
\end{gather*}

if every support lies in a ball of radius one, placed anywhere, and different
supports have Euclidean distance at least $d>0$. For $n=1$ only the radius
and norm conditions are imposed. A centered first moment equal to zero is
included. Real forms may first be complexified, changing $b$ by an absolute
factor. Neither invariance nor positivity is needed for the extension.

Put $A=\lceil\alpha\rceil$, $k_*=A+15$, and
$\Delta_n=\bigcup_{i<j}\{x_i=x_j\}$.
Let $\mathcal S_{\Delta,n}$ be the Schwartz functions on $\mathbb R^{4n}$
whose derivatives of every order vanish on $\Delta_n$. For $n=1$ this is
the ordinary Schwartz space. Write

\begin{gather*}
\FitMath{p_k(\varphi)=\max_{|\beta|\le k}\sup_x
       (1+|x|^2)^{k/2}|\partial^\beta\varphi(x)|.}\tag{QF2}
\end{gather*}


\begin{theorem}[Quantitative diagonal-flat extension]
\label{thm:quantitative-flat-extension}
There is $C<\infty$, depending on $b,\alpha$ and fixed numerical cutoff
choices but not on $n$, and a unique continuous linear functional
$\Sigma_n$ on $\mathcal S_{\Delta,n}$ agreeing with $S_n$ on smooth
separated tensors. It satisfies

\begin{gather*}
\FitMath{|\Sigma_n(\varphi)|
       \le C^n(n!)^{\kappa+32}p_{k_*n}(\varphi),} \\[3pt]
\FitMath{n\ge1.}\tag{QF3}
\end{gather*}

Finite sums of smooth compact separated tensors are dense in this
diagonal-flat space in its Schwartz topology. In its positive-time
subspace, the tensors may be chosen with supports compactly contained in
$\{x_4>0\}$. These are constrained density assertions, not density in
the unrestricted Schwartz space across the diagonals.
\end{theorem}

The conclusion is an H1 growth estimate with linear derivative order and
a fixed factorial exponent; its constants are intentionally not optimized.
For example, $\alpha=5$ gives the sufficient order $20n$.
The probability-theory normalization $\Sigma_0(1)=1$ may be added
separately.

The radius-one formulation loses no information about a hypothesis stated
as
\[
 B(R)^n(n!)^\kappa\max(1,d^{-\alpha n})
\]
with constants independent
of support positions. Choose a fixed smooth unit-scale partition in each
slot. A function supported in a radius-$R$ ball is a sum of at most
$C(1+R)^4$ supported pieces of radius at most one, each with $C^1$ norm
at most $C\|f\|_1$. Supports of pieces from different slots retain the
original separation. Multilinearity and (QF1) therefore imply the same
bound with

\begin{gather*}
\FitMath{B_{\rm eff}(R)\le Cb(1+R)^4.}\tag{QF4}
\end{gather*}

Thus only $B(1)$ is used, even if the supplied large-radius upper bound
was exponential. This assertion concerns a separate ball for each entry,
not one ball enclosing the entire configuration.

\subsection{Cutoffs with controlled derivatives}

We fix all cutoffs once. They can be chosen in Gevrey class two, with

\begin{gather*}
\FitMath{0\le\eta\le1,} \\[3pt]
\FitMath{|\partial^\beta\eta|\le C^{|\beta|+1}(|\beta|!)^2.}\tag{QF5}
\end{gather*}

Here is an explicit construction and a justification of this bound. Set
$a(t)=e^{-1/t}$ for $t>0$ and $a(t)=0$ for $t\le0$.
On the complex circle $|z-t|=t/2$ one has
$\operatorname{Re}(1/z)\ge2/(9t)$. Cauchy's formula bounds
$|a^{(m)}(t)|$ by $m!(2/t)^m e^{-2/(9t)}$.
Maximizing $t^{-m}e^{-c/t}$ and using $(m/e)^m\le m!$ gives
$C^{m+1}(m!)^2$. The step

\begin{gather*}
\FitMath{s(t)=\frac{a(t)}{a(t)+a(1-t)}}
\end{gather*}

is set equal to zero for $t\le0$ and one for $t\ge1$.
Its denominator on $[0,1]$ is at least $e^{-2}$.
The same derivative bound holds for its reciprocal and hence for $s$.
For precision, differentiate $D(1/D)=1$ inductively; after division by
$(m!)^2$, the convolution coefficients are
$1/\binom mr\le1$. If $|D^{(r)}|\le C_D^{r+1}(r!)^2$,
a derivative constant $C_D(1+2e^2C_D)$ for the reciprocal closes
the resulting geometric sum. Products
obey the same bound by Leibniz's formula. This proves (QF5), including
flatness at the join points.

Choose an even $\eta_0$ equal to one on $[-1,1]$, zero outside
$[-2,2]$, and nonincreasing on $[0,\infty)$, using $s(2-|t|)$ in
the transition. Put $\psi(t)=\eta_0(2t)$ and
$D(t)=\sum_{z\in\mathbb Z}\psi(t-z)$. At most three summands occur
locally and $D\ge1$. Then $\chi_1=\psi/D$ is supported in $[-1,1]$,
has (QF5), and $\sum_z\chi_1(t-z)=1$. Define

\begin{gather*}
\FitMath{\chi(u)=\prod_{a=1}^4\chi_1(u_a),} \\[3pt]
\FitMath{\chi^+(u)=\prod_{a=1}^4\eta_0(u_a),} \\[3pt]
\FitMath{\zeta(u)=1-\chi^+(u).}\tag{QF6}
\end{gather*}

Thus $\sum_{z\in\mathbb Z^4}\chi(u-z)=1$,
$\chi^+=1$ on the support of $\chi$, and $\zeta$ is zero for
$|u|_\infty\le1$ and one for $|u|_\infty\ge2$.
All these functions lie between zero and one and have (QF5), with
an absolute constant. In particular, undifferentiated factors cost one.

We will repeatedly use the following elementary product estimate.
For $N$ factors of this kind at scale $h$, every derivative of order
$q$ is bounded by

\begin{gather*}
\FitMath{(CN)^q(q!)^2h^{-q}.}\tag{QF7}
\end{gather*}

An absolute prefactor is allowed at $q=0$. Indeed, distribute the
multiindex among the factors by Leibniz. The sum of the multinomial
coefficients is $N^q$, and the product of the factorials of the
assigned derivative orders is at most $q!$. Derivatives of pair
differences $x_i-x_j$ introduce only signs. This argument produces no
factor $C^N$ when most factors are undifferentiated.

\subsection{An explicit local functional on a product box}

First, tensor relations cause no ambiguity. In a finite relation between
separated tensors, let $d$ be smaller than every separation. Localize all
entries by one finite smooth partition on sets of diameter less than
$d/8$. A product of slot sets with two sets closer than $d/2$ has a
zero factor in every original tensor: otherwise two original supports
would be closer than $d$. On every remaining product the slot sets are
separated, and the form is an ordinary multilinear map on their supported
function spaces. It factors through their algebraic tensor product.
The latter embeds injectively in functions on the Cartesian product:
for the finite-dimensional function spans in any relation, finitely many
point evaluations give invertible matrices in each slot. A zero tensor
therefore has zero value. Summing the localization proves independence
of the tensor presentation.

Here is a quantitative local extension. Let $0<h\le1/16$,
$\ell=h/64$, and let slot centers be $c_i\in\mathbb R^4$.
Suppose the cubes $Q_i^+=c_i+[-2\ell,2\ell]^4$ have pairwise distance
at least $h/2$. They lie in balls of radius one.
For a smooth function $u$ supported in
$\prod_i(c_i+[-\ell,\ell]^4)$, take its Fourier series on the product
of period-$8\ell$ cubes centered at these points. Write its coefficients
as $c_{k_1,\ldots,k_n}$, $k_i\in\mathbb Z^4$. Multiplying the series
by $\prod_i\chi^+((x_i-c_i)/\ell)$ gives a series of separated tensors
equal to $u$ everywhere. There is no boundary ambiguity: $u$ vanishes
near the boundaries of the period cubes.
Explicitly, with $Q=\prod_i(c_i+[-4\ell,4\ell]^4)$, the coefficients are

\begin{gather*}
\FitMath{c_{k_1,\ldots,k_n}=(8\ell)^{-4n}
    \int_Q u(x)\exp\!\left[-\frac{i\pi}{4\ell}
                  \sum_i k_i\cdot(x_i-c_i)\right]dx.}
\end{gather*}


In the scaled variables $y_i=(x_i-c_i)/(4\ell)$, apply
$(1-\pi^{-2}\Delta_i)^4$ in every
four-dimensional slot. Integration by parts gives

\begin{gather*}
\FitMath{|c_{k_1,\ldots,k_n}|
 \le C^n J_{8n,\ell}(u)
                    \prod_i(1+|k_i|^2)^{-4},} \\[3pt]
\FitMath{J_{q,\ell}(u)=\max_{|\beta|\le q}
                       \ell^{|\beta|}\|\partial^\beta u\|_\infty.}\tag{QF8}
\end{gather*}

The sum of absolute coefficients in the product differential operator
is at most $C^n$; no factor depending unspecifiedly on dimension is
hidden here. The localized exponential in slot $i$ has $C^1$ norm at
most $C\ell^{-1}(1+|k_i|)$. Since

\begin{gather*}
\FitMath{\sum_{k\in\mathbb Z^4}(1+|k|)(1+|k|^2)^{-4}<\infty,}
\end{gather*}

(QF1) makes the series of its moment values absolutely convergent and
gives a local functional $T(u)$ with

\begin{gather*}
\FitMath{|T(u)|\le (Cb)^n(n!)^\kappa
                 h^{-(\alpha+1)n}J_{8n,\ell}(u).}\tag{QF9}
\end{gather*}

For a tensor supported in the inner cubes it agrees with $S_n$:
its Fourier coefficients factor slot by slot, the partial sums converge
in $C^1$, and (QF1) permits passage through the multilinear form.

These local extensions glue to a unique functional $T$ on
$C_c^\infty(\mathbb R^{4n}\setminus\Delta_n)$. To check this explicitly,
in an overlap choose smaller product boxes enclosing the compact support
under consideration. Fourier expansions with cutoffs inside the overlap
converge with every fixed number of derivatives and consist of separated
tensors. Two local functionals agreeing on those tensors therefore agree
there, by their finite-order bounds (QF9). A finite partition of unity on
a compact off-diagonal set now defines $T$ independently of the cover.
For original separated tensors of larger radius, first partition each
entry into sufficiently small pieces and use multilinearity. This also
shows that the glued $T$ agrees with the given form on all smooth compact
separated tensors, not just on the boxes used in its construction.
This gluing argument is used for consistency, not for the numerical bound:
the latter uses only the fixed explicit grid and (QF9).

\subsection{Dyadic configuration layers and the factorial accounting}

Set $h_j=2^{-j}/16$ and define

\begin{gather*}
\FitMath{P_h(x)=\prod_{i<l}\zeta((x_i-x_l)/h),} \\[3pt]
\FitMath{D_0=P_{h_0},} \\[3pt]
\FitMath{D_j=P_{h_j}-P_{h_{j-1}}} \\[3pt]
\FitMath{(j\ge1).}\tag{QF10}
\end{gather*}

For $n=1$ the empty product is one and $D_j=0$ for $j\ge1$.
Off the diagonals $\sum_{j\ge0}D_j=1$. A fine layer has support

\begin{gather*}
\FitMath{h_j\le\min_{i<l}|x_i-x_l|_\infty\le4h_j.}\tag{QF11}
\end{gather*}

The statements about support also hold for its derivatives. Monotonicity
of the one-variable cutoffs makes $0\le D_j\le1$, although positivity
of these differences is not needed for the estimates.

At level $j$ set $\ell_j=h_j/64$ and use the product grid partition

\begin{gather*}
\FitMath{\chi_{j,z}(x)=\prod_{i=1}^n\chi(x_i/\ell_j-z_i),} \\[3pt]
\FitMath{z\in(\mathbb Z^4)^n.}\tag{QF12}
\end{gather*}

Retain a box if $\chi_{j,z}D_j$ is not identically zero. At one point
of that box every pair has distance greater than $h_j$; its centers
therefore have pair distance at least $h_j-2\ell_j$.
The expanded cubes $\ell_jz_i+[-2\ell_j,2\ell_j]^4$ have pair
distance at least $h_j-6\ell_j>h_j/2$. Thus (QF9) applies to
$u_{j,z}=\chi_{j,z}D_j\varphi$. A box with only boundary contact can
be ignored if its smooth localized function is zero.

There are at most $Cn^2$ pair and grid factors in such a cutoff.
Using (QF7), Leibniz's formula and $q=8n$, we obtain the uniform
scaled-derivative cost

\begin{gather*}
\FitMath{(Cn^2)^{8n}((8n)!)^2\le C^n(n!)^{32}.}\tag{QF13}
\end{gather*}

Here $n^{16n}\le e^{16n}(n!)^{16}$, while
$(8n)!\le8^{8n}(n!)^8$. This accounts for every dependence on the
number of pair cutoffs; the $\binom n2$ undifferentiated factors do
not create $C^{n^2}$.

On a fine layer choose a pair with $|x_i-x_l|_\infty\le4h_j$ and
replace $x_i$ by $x_l$. The resulting point $y$ lies on a diagonal
and $|x-y|\le8h_j$. Every derivative of $\varphi$ vanishes there.
Taylor's integral remainder along this segment gives, for $|\beta|\le8n$
and any positive integer $m$,

\begin{gather*}
\FitMath{|\partial^\beta\varphi(x)|
 \le\frac{(Ch_j)^m}{m!}
     C^{k_*n}(1+|x|^2)^{-k_*n/2}
                    p_{k_*n}(\varphi)}\tag{QF14}
\end{gather*}

when $8n+m\le k_*n$. The derivative of order $m$ along a vector
uses at most four coordinate directions; its multinomial sum is
bounded by the $m$-th power of that vector's $\ell^1$ norm.
Also $|x-y|\le1/2$, so the weight at every point of the segment is
comparable with that at $x$, at cost $C^{k_*n}$. These are exponential,
not factorial, costs in $n$.

Take

\begin{gather*}
\FitMath{m=(A+5)n+2,} \\[3pt]
\FitMath{8n+m=(A+13)n+2\le(A+15)n.}\tag{QF15}
\end{gather*}

In differentiating $u_{j,z}$, a derivative on a cutoff costs $h_j^{-1}$,
which its scaled-derivative factor in $J_{8n,\ell_j}$ cancels.
Derivatives remaining on $\varphi$ are bounded by (QF14).
Since $h_j\le1$, unused positive powers of $h_j$ can be discarded.
Equations (QF7), (QF13)--(QF15) consequently imply

\begin{gather*}
\FitMath{J_{8n,\ell_j}(u_{j,z})
 \le C_A^n(n!)^{32}h_j^m p_{k_*n}(\varphi)
       \sup_{x\in\operatorname{supp}\chi_{j,z}}
                         (1+|x|^2)^{-k_*n/2}.}\tag{QF16}
\end{gather*}

For $j=0$ use the same estimate without $h_j^m$, by the ordinary
Schwartz derivative bound; all powers of the fixed $h_0$ are absorbed
in $C_A^n$.

The grid sum has exactly the needed size. Since
$1+\sum_i|x_i|^2\ge\prod_i(1+|x_i|^2)^{1/n}$, and grid points in
one slot differ from its box by a fixed distance less than one,

\begin{gather*}
\FitMath{\begin{aligned}
 \sum_z\sup_{x\in\operatorname{supp}\chi_{j,z}}
                  (1+|x|^2)^{-k_*n/2}
 &\le C_A^n\prod_{i=1}^n
       \sum_{z_i\in\mathbb Z^4}(1+|\ell_jz_i|^2)^{-k_*/2}\\
 &\le C_A^n h_j^{-4n}.
 \end{aligned}}\tag{QF17}
\end{gather*}

The one-slot sum is bounded by $C_A\ell_j^{-4}$ by comparison with
the integral of $(1+|u|^2)^{-k_*/2}$ on $\mathbb R^4$; it is finite
because $k_*>4$. No number of configurations depending exponentially
on $n^2$ is being counted.

Combine (QF9), (QF16) and (QF17). Since $\alpha\le A$, the sum
of absolute values at every fine level is at most
\begin{equation}
\FitMath{C^n(n!)^{\kappa+32}p_{k_*n}(\varphi)
                 h_j^{m-(\alpha+5)n}
 \le C^n(n!)^{\kappa+32}p_{k_*n}(\varphi)h_j^2.}
 \tag{QF18}
\end{equation}
The coarse level has the same bound without $h_j^2$.
The dyadic sum $\sum_jh_j^2$ is finite. We may therefore define

\begin{gather*}
\FitMath{\Sigma_n(\varphi)=\sum_{j\ge0}\sum_zT(\chi_{j,z}D_j\varphi).}\tag{QF19}
\end{gather*}

This absolutely convergent sum proves (QF3) and continuity in its stated
seminorm. It agrees with $T$ on compact off-diagonal tests, because
there only finitely many layers and boxes occur and their partition
sums to one. In particular it has the required tensor values.

\subsection{Density, the time boundary, and independence of choices}

We give the density details also needed for uniqueness. First cut off
infinity by $Q_R(x)=\prod_i\chi^+(x_i/R)$.
For a fixed Schwartz seminorm $p_k$, an extra $L$ powers of decay give

\begin{gather*}
\FitMath{p_k((1-Q_R)\varphi)
       \le C_{n,k,L}R^{-L}p_{k+L}(\varphi),} \\[3pt]
\FitMath{R\ge1.}\tag{QF20}
\end{gather*}

This follows by Leibniz, since the nonzero terms have at least one
$|x_i|\ge R$ and cutoff derivatives cost $R^{-1}$.
Next cut out all pairwise diagonals with $P_h$.
Where a derivative of $1-P_h$ can contribute, some pair is within
$2h$. Applying the same Taylor argument with arbitrarily many extra
derivatives gives

\begin{gather*}
\FitMath{p_k((1-P_h)\varphi)
       \le C_{n,k,M}h^{M-k}p_{k+M}(\varphi).}\tag{QF21}
\end{gather*}

For each fixed $n,k$ choose $M>k$. Thus compact tests supported away
from every diagonal are dense in $\mathcal S_{\Delta,n}$.
For example both displayed density constants may be bounded by
$C^{k+M+L+1}(n^2+1)^k(k!)^2$, with unused parameters omitted.
This follows from (QF7) and the same directional Taylor remainder;
these constants are not used in place of the uniform estimate (QF18).
On each such compact set, use a finite product-box partition and the
Fourier expansions of (QF8). Their partial sums converge with every
derivative and have a common compact support. They consequently converge
in the Schwartz topology and are finite sums of separated smooth tensors.
This proves the asserted constrained tensor density.
Choosing the truncations and Fourier partial sums successively for the
first $1,2,\ldots$ seminorms gives a single convergent sequence.

For positive time define $\mathcal S_{\Delta,+,n}$ by requiring all
derivatives to vanish outside
$\{x_{1,4}>0,\ldots,x_{n,4}>0\}\setminus\Delta_n$.
In addition to the preceding cutoffs, set
\begin{equation}
\FitMath{B_h(x)=\prod_{i=1}^n s(x_{i,4}/h-1).}
 \tag{QF22}
\end{equation}
It is zero if some time is at most $h$ and one if all are at least
$2h$. Where $1-B_h$ contributes to a positive-time test, a time is
within $2h$ of zero. Taylor expansion to that forbidden hyperplane,
where every jet vanishes, proves the analogue of (QF21).
Hence $Q_RP_hB_h\varphi$ approximates every positive-time test in
the full Schwartz topology. Product boxes for this approximation can
be expanded while staying strictly above time zero. Their tensor
approximants are therefore strictly positive-time separated tensors.

There is also a uniform quantitative version: replace $P_h$ throughout
(QF10) by $P_hB_h$. A fine layer then has all pair distances and all
times at least $h$, while one pair distance or one time is at most
$4h$. The expanded grid boxes retain separation and positive time.
There are still at most $Cn^2$ cutoff factors. Taylor to the nearest
forbidden pair diagonal or time hyperplane supplies (QF14), so exactly
the same $k_*$ and factorial exponent apply. This verifies explicitly
that the time-boundary constraint introduces no hidden growth in $n$.

Uniqueness now follows without a choice of a particular partition.
Any two continuous extensions agreeing on separated tensors agree on
a dense subspace and hence on all of $\mathcal S_{\Delta,n}$.
Thus the Fourier boxes, cutoffs, dyadic grids and orders of their
absolutely convergent sums do not change the resulting functional.
The same statement holds on the positive-time subspace.

If the compact moments are symmetric, real or invariant under a specified
Euclidean motion, their extensions inherit that identity by density and
continuity of the test-function transformations. If the compact moments
are reflection positive, take finite positive-time Schwartz sequences
and approximate their components by the tensors just constructed.
For a negative-time and positive-time factor, their tensor product is
flat on all cross diagonals: away from the time boundary one factor
already vanishes, and at time zero its jets vanish. Internal diagonals
are handled by their own factors. The tensor product is continuous
in Schwartz seminorms. Thus the reflected quadratic forms converge,
and their nonnegativity passes to these full sequence domains.
This proves the distribution-space implication under the graded moment
hypothesis (QF1).

\paragraph{The graded-source limit and phase consequence.}
Under Theorem~\ref{thm:graded-rooted-source}, every pointwise limit on
separated tuples of its two-state centered moments retains (GS12): for
each fixed tuple the inequality is closed and its constants are uniform
in cutoff and in both states. The phase-interface theorem realizes each
of its component forms by such sequences of pure lattice states, centered
at the specified reference states. Therefore the same bound holds for
those components, including their possibly nonzero first moments; their
one-point shifts are covered by the all-state comparison bound.
The fixed equivalence between $\|\cdot\|_\sharp$ and $\|\cdot\|_1$
is absorbed in the radius-one constant $b$. Thus (QF1) applies with
$\alpha=5$ and the graded factorial exponent $k$. These limits and
components consequently have the unique diagonal-flat extensions with
order $20n$ and factorial exponent $k+32$, uniformly over the stated
phase realization. This uses that realization rather than assuming an
unrelated component decomposition has the same bounds. It supplies
neither contact terms across the diagonals nor full rotational covariance.

\subsection{Why the precise hypotheses matter}

Position uniformity cannot be replaced by unrestricted exponential
growth in an enclosing radius. For example,
$S_1(f)=\int e^{|x|^2}f(x)\,dx$ has finite bounds for compact tests
in every fixed ball about the origin, with an exponential radius cost,
but has no continuous extension to Schwartz space agreeing with those
values. Translated nonnegative compact bumps, damped by
$e^{-|x|^2/2}$, converge to zero in every Schwartz seminorm while their
values diverge. Their supports can be chosen in successively more distant
unit balls. This example violates the position-uniform radius-one
hypothesis. Under that actual hypothesis (QF4) rules out this obstruction.

Nor may a prefactor such as $e^{c n^2}$ be silently called factorial
growth of a fixed order. A lognormal random constant field has moments
$S_n(f_1,\ldots,f_n)=e^{n^2/2}\prod_i\int f_i$.
It is translation invariant and reflection positive, but violates every
estimate (QF3) with fixed linear seminorm order and fixed factorial
exponent. To see this, take unit-integral Gevrey-two bumps in $n$
separated unit balls with centers at distance $O(n)$ from the origin.
Their tensor is diagonal-flat. For every fixed $k$, its seminorm
$p_{kn}$ grows at most $C^n(n!)^{C_k}$, by (QF5) and the polynomial
location weights. Its moment grows like $e^{n^2/2}$, which is larger
than $C^n(n!)^q$ for every fixed $q$. Thus the graded, fixed-exponent
assumption in (QF1) has mathematical content; reflection positivity
alone does not supply it.

% END MODULE chapters/quantitative_flat_extension_v4.tex


% BEGIN MODULE chapters/euclidean_hilbert_bridge_v3.tex
\section{From separated moments to Euclidean evolution}
\label{ch:euclidean-hilbert-bridge}

This section proves a conditional reconstruction on compactly supported
separated test tuples. The time contraction follows from a bound which is
uniform in the positions of finitely many supports. No estimate on arbitrary
Schwartz sequences, and no bound uniform in the number of fields, enters
that argument. A further lattice approximation hypothesis gives a spectral
cone with a cubical spatial section. The round cone and continuous rotations
require additional input.

\subsection{The precise compact test domain}

Put $\mathcal D=C_c^1(\mathbb R^4;\mathbb C)$ and
\[
 \|f\|_1=\max\{\|f\|_\infty,\|\nabla f\|_\infty\}.
\]
A separated tuple $X=(f_1,\ldots,f_n)$ has compact supports at pairwise
positive distance. The empty tuple is denoted by $1$. Write $\mathscr V$
for the complex vector space of finite formal linear combinations of these
tuples. Multilinear and permutation relations may be imposed at the outset;
equivalently they will be null relations in the forms below. Concatenation
$X\cdot Y$ is used only if the supports of all its entries remain separated.
There is no claim that $\mathscr V$ is an algebra under unrestricted
concatenation.

Our translation convention is
$\tau_vf(x)=f(x-v)$, so its support is translated by $v$.
For the time reflection $\theta(x_1,x_2,x_3,x_4)
=(x_1,x_2,x_3,-x_4)$ put $\theta f=f\circ\theta$.
These maps act entrywise on tuples and linearly on formal combinations.
The bar conjugates both the coefficients and the test functions.
Let $\mathscr V_+$ consist of combinations whose entries have supports
compactly contained in $\mathbb E_+=\{x_4>0\}$.

Assume normalized symmetric multilinear moments $S_n$ on this separated
domain. Write $\sigma(X)=S_n(f_1,\ldots,f_n)$, $\sigma(1)=1$,
and extend linearly. They obey the following hypotheses.
They preserve conjugation, are translation invariant, and are invariant under
the coordinate reflections. Their reflection-positive form satisfies

\begin{gather*}
\FitMath{q(F,F):=\sigma(\overline{\theta F}\cdot F)\ge0} \\[3pt]
\FitMath{(F\in\mathscr V_+).}\tag{HB1}
\end{gather*}

This is positivity for every finite linear combination, including the empty
tuple, rather than just for individual monomials. The concatenations in
(HB1) and its polarized form are separated: reflected supports lie strictly
below time zero, and the others strictly above it.

For every integer $n\ge1$, $R<\infty$ and $d>0$, assume a finite number
$M_n(d,R)$ such that

\begin{gather*}
\FitMath{|S_n(f_1,\ldots,f_n)|
       \le M_n(d,R)\prod_{i=1}^n\|f_i\|_1,}\tag{HB2}
\end{gather*}

whenever each support is in a ball of radius $R$, centered anywhere, and
their pairwise distances are at least $d$. The centers need not lie in a
common bounded region. For one entry the separation condition is empty.
These are moment bounds, not merely cumulant bounds. Bounds on cumulants
give (HB2) after the finite moment--cumulant sum, if the needed one-point
bound is included. There is no restriction here on the growth of $M_n$ as
$n$ increases or $d$ decreases.

\begin{theorem}[Compact-tuple Euclidean reconstruction]
Under (HB1)--(HB2) and the stated invariances, the completion of
$\mathscr V_+$ by its reflection form is a separable Hilbert space
$\mathcal H$. The empty tuple has class $\Omega$ of norm one. The formulas

\begin{gather*}
\FitMath{T_t[F]=[\tau_{te_4}F]} \\[3pt]
\FitMath{(t\ge0),} \\[3pt]
\FitMath{U_x[F]=[\tau_{(x,0)}F]} \\[3pt]
\FitMath{(x\in\mathbb R^3)}\tag{HB3}
\end{gather*}

define, respectively, a strongly continuous semigroup of positive
self-adjoint contractions and a strongly continuous unitary representation.
They commute and fix $\Omega$. There is a unique nonnegative self-adjoint
operator $H$ with $T_t=e^{-tH}$. Each $T_t$, $t>0$, is injective.
The spatial generators $P_i$ and $H$ commute strongly.
\end{theorem}

\subsection{The quotient and the contraction proof}

Positivity of the quadratic form gives Hermiticity by polarization and the
Cauchy--Schwarz inequality
$|q(F,G)|^2\le q(F,F)q(G,G)$. Thus
$\mathscr N=\{F:q(F,F)=0\}$ is a linear subspace orthogonal to everything.
The quotient $\mathscr V_+/\mathscr N$, with inner product $q$, is
pre-Hilbert; its completion defines $\mathcal H$.
In particular $\|\Omega\|^2=\sigma(1)=1$ and
$\langle\Omega,[F]\rangle=\sigma(F)$.

Time and spatial translations in (HB3) preserve their positive-time domain.
Translation invariance and $\theta\tau_{te_4}=\tau_{-te_4}\theta$ give

\begin{gather*}
\FitMath{q(F,\tau_{te_4}G)=q(\tau_{te_4}F,G),} \\[3pt]
\FitMath{q(\tau_{(x,0)}F,\tau_{(x,0)}G)=q(F,G).}\tag{HB4}
\end{gather*}

The second identity proves that spatial translations preserve null vectors
and induce isometries with inverse $U_{-x}$.
For the first identity, boundedness of the induced time maps remains to be
proved; symmetry alone would not suffice.

Fix $F=\sum_\alpha c_\alpha X_\alpha\in\mathscr V_+$.
There are $\delta>0$, $d>0$ and $R<\infty$ such that every entry has
time coordinate at least $\delta$, every individual tuple has internal
separation at least $d$, and each entry has support in a ball of radius $R$.
For every $s\ge0$, every concatenation
$\overline{\theta\tau_{se_4}X_\alpha}\cdot
\tau_{se_4}X_\beta$ has internal separation at least
$d_*=\min\{d,2\delta\}>0$. Indeed its two groups lie below
$-\delta-s$ and above $\delta+s$. Moving the groups apart changes
their centers but not their radii, norms or internal separations.
Applying (HB2) to the finitely many pairs gives

\begin{gather*}
\FitMath{\sup_{s\ge0}\|\tau_{se_4}F\|_q\le C_F<\infty.}\tag{HB5}
\end{gather*}

The same assertion covers empty entries and constants by normalization.
Its essential hypothesis is that the bound is independent of support
positions. A bound with a growing enclosing-volume factor would not prove
(HB5).

Write $a(s)=\|\tau_{se_4}F\|_q$. Symmetry and Cauchy--Schwarz give

\begin{gather*}
\FitMath{a(t)^2=q(F,\tau_{2te_4}F)\le a(0)a(2t).}
\end{gather*}

Iteration yields

\begin{gather*}
\FitMath{a(t)\le a(0)^{1-2^{-k}}a(2^kt)^{2^{-k}}
       \le a(0)^{1-2^{-k}}C_F^{2^{-k}}.}\tag{HB6}
\end{gather*}

If $a(0)=0$, the first inequality already proves $a(t)=0$.
Otherwise let $k\to\infty$ to obtain $a(t)\le a(0)$.
Consequently time translations preserve $\mathscr N$ and extend to
contractions on $\mathcal H$. Their symmetry in (HB4) makes each bounded
extension self-adjoint. Moreover

\begin{gather*}
\FitMath{q(F,\tau_{te_4}F)=q(\tau_{(t/2)e_4}F,
                         \tau_{(t/2)e_4}F)\ge0,}
\end{gather*}

so every $T_t$ is positive. The semigroup law is inherited from the exact
translation law on the dense domain.

\subsection{Continuity, generators and elementary consequences}

For $f\in C_c^1$, uniform continuity of $\nabla f$ gives

\begin{gather*}
\FitMath{\|\tau_h f-f\|_1
 \le\max\left\{|h|\|\nabla f\|_\infty,
       \sup_{|x-y|\le|h|}|\nabla f(x)-\nabla f(y)|\right\}
       \longrightarrow0.}\tag{HB7}
\end{gather*}

This statement need not hold in the Lipschitz norm for a merely Lipschitz
function. To apply it to a moment, enclose the old and slightly moved
supports of each entry in one common neighborhood. For sufficiently small
shifts these neighborhoods retain half their original positive separation.
The multilinear telescoping identity and (HB2) then bound the change of a
moment by a constant times the sum of the $C^1$ changes of its entries.
Differences of functions are evaluated only on those fixed neighborhoods.

Expanding $\|[\tau_hF]-[F]\|^2$ into its finitely many moment terms and
using this observation proves convergence to zero for $h=te_4$, $t\downarrow0$,
and for spatial $h\to0$. The contraction and unitary bounds extend this
continuity from the dense quotient domain to all of $\mathcal H$.
The semigroup and group laws then give strong continuity at every parameter.

Here is a construction of the generator which also rules out a hidden
zero-transfer subspace. If $T_1v=0$, positivity and
$T_1=T_{1/2}^2$ give $T_{1/2}v=0$. Repeated halving and strong continuity
at zero imply $v=0$. Thus the spectral measure of $T_1$ has no atom at
zero. Define the densely defined self-adjoint operator
$H=-\log T_1$ by functional calculus; its spectrum is nonnegative.
Uniqueness of positive roots gives $T_{m/n}=T_1^{m/n}$ for nonnegative
rationals, and strong continuity gives $T_t=e^{-tH}$ for all $t\ge0$.
Functional calculus also proves uniqueness of $H$ and injectivity of every
$T_t$, $t>0$. Since translations fix the empty tuple,
$H\Omega=0$ and $U_x\Omega=\Omega$.

The commuting spatial unitaries have self-adjoint generators
$U_x=e^{i\sum_i x_iP_i}$. They commute with $T_1$ and hence with its
spectral projections and with $H$. The joint spectral theorem, equivalently
the spectral theorem for the strongly continuous abelian unitary group
$(s,x)\mapsto e^{isH}U_x$, supplies one projection-valued measure
$\mathsf E$ on $[0,\infty)\times\mathbb R^3$. Here and below strong
commutation means commutation of these spectral projections.
For $F,G\in\mathscr V_+$, $t\ge0$ and $x\in\mathbb R^3$ the dictionary is

\begin{gather*}
\FitMath{\begin{aligned}
 \sigma(\overline{\theta F}\,\tau_{(x,t)}G)
 &=\langle[F],e^{-tH}U_x[G]\rangle\\
 &=\int e^{-tE+ix\cdot p}
                  \langle[F],\mathsf E(dE\,dp)[G]\rangle.
 \end{aligned}}\tag{HB7a}
\end{gather*}


For completeness, the quotient is separable. Cover compact positive-time
supports by finite unions of rational boxes with closures still in
$\mathbb E_+$; choose these unions separately for each entry so that they
remain separated. There are countably many such finite systems. On each
fixed bounded union, the $C^1$ space is separable: restriction of a function
and its four derivatives embeds it isometrically in a finite product of
separable continuous-function spaces. Choose countable dense subsets inside
the supported $C^1$ spaces, and rational complex coefficients for formal
combinations. Telescoping with (HB2), on the fixed separated neighborhoods,
shows that the resulting countable set is total in $\mathcal H$.

The formal description does not create extra vectors from zero tensors.
For a finite homogeneous relation $\sum_\alpha c_\alpha
f_{\alpha,1}\otimes\cdots\otimes f_{\alpha,n}=0$ whose tuples are
separated, choose $d>0$ below all their separations. Localize every entry
with one finite smooth partition covering their supports, whose sets have
diameter less than $d/8$. In a product of partition sets, either two sets
have distance less than $d/2$, in which case every original tensor has a
zero factor there, or all the slot sets remain separated. In the latter
case the moment is multilinear on the product of their supported function
spaces and hence is a linear functional on their algebraic tensor product.
The map from this tensor product into functions on the Cartesian product
is injective: on the finitely many function spans involved, independent
functions can be distinguished by finitely many point evaluations, and
the corresponding evaluation matrices invert each slot. Thus the localized
zero tensor has zero moment. Summing proves the assertion. Applying it
also after concatenation with a reflected test shows that a zero tensor
is a reflection-null vector. This is an algebraic compact-support fact,
distinct from the Schwartz density input below.

Any realization with a dense image of the same test domain and the same
scalar products is therefore canonically unitarily equivalent, by sending
each test class to its other realization; the generators intertwine by
density. The construction can use any coordinate as time if its RP
hypothesis is supplied. For $c\ge0$ and $F$ supported strictly above time
$c$, the exact identity is $[F]=T_c[\tau_{-ce_4}F]$.

Any further linear symmetry of the test functions, acting entrywise on tuples,
which preserves the admissible positive-time domain and the moments and
commutes with $\theta$ and conjugation acts isometrically on this quotient. In particular,
if the moments also have the spatial signed-permutation symmetry $W_3$,
then $V_w[F]=[wF]$ is unitary,
$V_wT_t=T_tV_w$ and $V_wU_xV_w^{-1}=U_{wx}$.
Uniqueness of the joint spectral measure gives the corresponding covariance
under $(E,p)\mapsto(E,wp)$ and hence $W_3$ symmetry of the joint spectrum.
This optional finite symmetry is not continuous rotation invariance.

Spectral dominated convergence gives $T_t\to P_0=1_{\{0\}}(H)$ strongly
as $t\to\infty$. Hence the large-time reflected correlations equal the
matrix elements of $P_0$. In particular, $\ker H=\mathbb C\Omega$ exactly
when their limits, or their Ces\`aro limits, equal
$\overline{\sigma(F)}\sigma(G)$ for every positive-time $F,G$.
If $\sigma(f)=0$ and $S_2(\theta\bar f,f)>0$ for one positive-time test,
then $[(f)]$ is a nonzero vector orthogonal to $\Omega$, so
$\dim\mathcal H\ge2$; this requires an actual positive two-point input.

A shift of the fields does not change this representation. More precisely,
let $m$ be a complex linear, translation- and reflection-invariant functional
on $\mathcal D$, satisfying $m(\bar f)=\overline{m(f)}$ and bounded on
each fixed compact support. Define

\begin{gather*}
\FitMath{C_m(f_1,\ldots,f_n)
  =\sum_{A\subset\{1,\ldots,n\}}
       \prod_{i\notin A}[-m(f_i)]\,(f_i)_{i\in A},} \\[3pt]
\FitMath{\sigma^{\rm c}(F)=\sigma(C_mF).}
\end{gather*}

Then $C_{-m}$ is the inverse of $C_m$, and on every admissible
concatenation $C_m(FG)=C_m(F)C_m(G)$. It commutes with the bar, reflections
and translations. Therefore
$[F]_{\sigma^{\rm c}}\mapsto[C_mF]_\sigma$ is a unitary intertwining
the two evolutions. Its finite expansion also transfers (HB2), since local
boundedness and translation invariance bound $m$ uniformly on balls of each
fixed radius. This algebraic statement does not estimate the size of
the one-point difference between two cutoff states.

In the particular unweighted Wilson-density normalization, the one-point
calibration can be written explicitly. Let two states $\rho_j,\nu_j$ on
spacing $a_j$ be translation and hypercubic invariant, and put
$\mathcal O_j(f)=\sum_p f(x(p))a_p$. Their plaquette mean difference
$d_j=\rho_j(a_p)-\nu_j(a_p)$ is independent of $p$. Define

\begin{gather*}
\FitMath{R_j(f)=\frac{a_j^4}{6}\sum_p f(x(p)),} \\[3pt]
\FitMath{c_j=6a_j^{-4}d_j.}
\end{gather*}

Then
$\rho_j(\mathcal O_j(f))-\nu_j(\mathcal O_j(f))=c_jR_j(f)$
exactly, and the six shifted Riemann sums give $R_j(f)\to\int f$.
If these one-point differences converge to $m(f)$ on all compact $C^1$
tests, choose one $f_0$ with $\int f_0=1$. The ratio at $f_0$ shows
$c_j\to m(f_0)=c$, and consequently $m(f)=c\int f$ for every test.
A uniform one-point bound at $f_0$ bounds $|c|$; this bound is an input,
not a consequence of translation invariance alone. When the two states
coincide, $c_j=c=0$.

\subsection{Coarse lattice inequalities and the continuum cube cone}

The spectral cone needs additional information. Assume infinite-lattice
probability measures $\nu_j$ at spacings $a_j\downarrow0$, invariant
under lattice translations and reflection positive on the closed side of
all coordinate and diagonal lattice mirrors. Their moment forms will be
denoted by $S_n^{(j)}$. For periodic Wilson DLR states these hypotheses are proved in
Theorem~\ref{thm:amrp-all-periodic}. For other lattice measures they
remain explicit assumptions. The normal transfer $T_j$
and tangential translation $V_{j,i}$ in their reflection Hilbert spaces
then have the carrier established in the plateau argument:

\begin{gather*}
\FitMath{\sinh(E/2)\ge|\sin(p/2)|,} \\[3pt]
\FitMath{T_j=e^{-E},} \\[3pt]
\FitMath{V_{j,i}=e^{ip}.}\tag{HB8}
\end{gather*}

Here $E$ is dimensionless and may be infinity on $\ker T_j$, while $p$
is a circle variable. Only this lattice carrier, rather than a round
relativistic cone, is taken from that argument.

One lattice distinction is useful. For spacing $a$, the site mirror at
zero and the link mirror at $-a/2$ have the same positive-side lattice
algebra. Their forms obey
$q_{\rm link}(A,B)=q_{\rm site}(A,T B)$, by translation invariance.
Consequently $[A]_{\rm link}\mapsto T^{1/2}[A]_{\rm site}$ is an
isometry onto the closure of $\operatorname{Ran}T^{1/2}$, namely
$(\ker T)^{\perp}$. It intertwines the two normal transfers, since this
holds on the common local domain. Thus the link transfer is injective,
even if the site transfer has a kernel. This proves the relation between
the two pictures without discarding site-kernel contributions in (HB10).

For every integer $k\ge1$, convexity of $\sinh$ and the elementary
inequality $|\sin(kv)|\le k|\sin v|$ give

\begin{gather*}
\FitMath{\sinh(kE/2)\ge k\sinh(E/2)
       \ge k|\sin(p/2)|\ge|\sin(kp/2)|.}\tag{HB9}
\end{gather*}

The sine inequality follows by induction from the addition formula and
$|\cos v|\le1$. Thus functional calculus gives the bounded operator
inequality

\begin{gather*}
\FitMath{(I-T_j^k)^2\ge T_j^k(2I-V_{j,i}^k-V_{j,i}^{-k}).}\tag{HB10}
\end{gather*}

At $T_j=0$ this inequality holds directly, so no division by a transfer
eigenvalue is required. Its quadratic form is a five-point inequality.
If $s=ka_j$ and
$\Phi_{j,F}(v)=\nu_j(\overline{\theta F_j}\,\tau_vF_j)$,
it reads

\begin{gather*}
\FitMath{\Phi_{j,F}(0)-4\Phi_{j,F}(se_4)+\Phi_{j,F}(2se_4)
       +\Phi_{j,F}(se_4+se_i)+\Phi_{j,F}(se_4-se_i)\ge0.}\tag{HB11}
\end{gather*}


To pass this statement to the limit, specify the approximation rather than
identify Hilbert spaces at different cutoffs. Assume linear local lattice
bounded Borel representatives $\iota_j(f)$ for compact $C^1$ tests, preserving conjugation:
$\iota_j(\bar f)=\overline{\iota_j(f)}$. Their supports lie
within distance $\epsilon_j\to0$ of $\operatorname{supp}f$. They commute
with coordinate reflection and with translations by $a_j\mathbb Z^4$.
Let $F_j$ be the product representative of each tuple, extended linearly,
and let $S_n^{(j)}$ be its moments. Assume convergence on every fixed
separated compact tuple to $S_n$, and assume (HB2) for these lattice moments
with constants independent of $j$ once $j$ is sufficiently large for the
fixed $n,d,R$. For action-density representatives their supremum norms may
depend on $j$; the asserted uniform control is on these moments.

If $v_j\to v$, the supports of $\tau_{v_j}f$ and $\tau_vf$ lie in fixed
slightly enlarged sets. Whenever the limiting tuple is separated, these
sets retain a positive separation. Telescoping and the uniform moment
bound give

\begin{gather*}
\FitMath{S_n^{(j)}(f_{1,j},\ldots,f_{n,j})
       -S_n^{(j)}(f_1,\ldots,f_n)\longrightarrow0}\tag{HB12}
\end{gather*}

for such translated tests, by (HB7). Together with fixed-tuple convergence,
this proves convergence at the moving arguments.

The same argument also explains inheritance of the Hilbert hypotheses from
these approximants, when it is needed. A continuum translation is approached
by lattice vectors $v_j$; exact lattice covariance and (HB12) pass its
moment identity to the limit. Coordinate reflections through the origin
need no rounding. For any fixed $F\in\mathscr V_+$ the representatives
$F_j$ lie on the positive side for large $j$, so their finite quadratic
RP inequalities pass directly to (HB1). Normalization, conjugation,
multilinearity and the uniform moment bounds pass by their finite defining
relations. The lattice mirror hypotheses themselves are supplied by
Theorem~\ref{thm:amrp-all-periodic} in the periodic Wilson application.

There is also a finite-torus scalar-product statement, with a narrower
conclusion. If fixed separated moments converge from growing periodic tori
and satisfy uniform $C^1$ bounds on the fixed neighborhoods under
consideration, RP at the link mirror $x_4=a_j/2$ gives

\begin{gather*}
\FitMath{q(F,G)=\lim_j\nu_j(\overline{\theta_{a_j/2}F_j}\,G_j).}
\end{gather*}

Indeed $\theta_{a_j/2}f=\tau_{a_je_4}\theta f$, whose difference from
$\theta f$ tends to zero in $C^1$, and the same telescoping estimate
applies. This compares scalar products as numbers; it asserts no maps
between the cutoff Hilbert spaces. Finite-torus positivity of this kind
does not by itself supply the global lattice semigroup or (HB8).

Fix $t>0$ and choose integers $k_j\ge1$ with $k_ja_j\to t$.
For positive-time $F$, its lattice representative is supported on the
positive side of the fixed site mirror for large $j$.
All five moment expressions in (HB11) are separated uniformly: the negative
reflected group and every translated positive group retain their original
positive-time margin. Internal separations are unchanged. Therefore (HB12)
passes (HB11) to its continuum version. In $\mathcal H$ this is

\begin{gather*}
\FitMath{B_{t,i}:=(I-T_t)^2-T_t(2I-U_{te_i}-U_{-te_i})\ge0.}\tag{HB13}
\end{gather*}

Initially its form is nonnegative on the dense compact-tuple domain; since
$B_{t,i}$ is bounded and self-adjoint, the inequality holds on all of
$\mathcal H$. No differentiation of a cutoff correlation at $t=0$ has
been used. The fixed-time coarse inequality is passed first.

In the joint functional calculus its multiplier is

\begin{gather*}
\FitMath{b_{t,i}(E,p)=1-4e^{-tE}+e^{-2tE}
                     +2e^{-tE}\cos(tp_i)
       =2e^{-tE}[\cosh(tE)+\cos(tp_i)-2].}\tag{HB14}
\end{gather*}

Take a countable sequence $t_m\downarrow0$, for example $t_m=1/m$.
Positivity of (HB13) means that the joint projection-valued measure gives
zero mass to $\{b_{t_m,i}<0\}$. Remove the countable union of those
sets for the three spatial indices. At every remaining spectral point,

\begin{gather*}
\FitMath{0\le\lim_{m\to\infty}
       \frac{\cosh(t_mE)+\cos(t_mp_i)-2}{t_m^2}
       =\tfrac12(E^2-p_i^2).}
\end{gather*}

Since $E\ge0$, this proves

\begin{gather*}
\FitMath{\mathsf E\{(E,p):E<\max_i|p_i|\}=0.}\tag{HB15}
\end{gather*}

In particular $\operatorname{Dom}H\subset\operatorname{Dom}P_i$ and
$\|P_i\psi\|\le\|H\psi\|$ on that domain. The corresponding
form-domain statement follows by integrating $|p_i|\le E$ rather than
their squares. Zero-energy vectors are spatially invariant.
Consequently vectors fixed by all Euclidean time and spatial translations
are exactly $\ker H$. This statement concerns invariant Hilbert vectors;
identifying them with an ergodic decomposition of correlation states needs
the separate phase argument.
The cone (HB15) has a cubical spatial section.
Nothing here proves $E\ge|p|_2$ or supplies an $\mathrm O(4)$ action.
Finite periodic torus laws cannot simply be substituted for the infinite
lattice hypotheses at arbitrarily large separations; any required
infinite-volume reflection-positive state must first be obtained.

\subsection{What the full Schwartz sequence domain additionally needs}

The compact-tuple construction above does not establish the claimed
Schwartz-domain growth property. Here is a precise sufficient additional
input for that extension. For $n\ge2$ put
$\Delta_n=\bigcup_{i<j}\{x_i=x_j\}$ and let
$\mathcal S_{\Delta,n}$ be the Schwartz functions on $\mathbb R^{4n}$
whose derivatives of every order vanish on $\Delta_n$.
Use $\mathcal S_{\Delta,0}=\mathbb C$ and the usual Schwartz space for
$n=1$. Define the seminorms

\begin{gather*}
\FitMath{p_k(\varphi)=\max_{|\alpha|\le k}\sup_x
       (1+|x|^2)^{k/2}|\partial^\alpha\varphi(x)|.}
\end{gather*}

The positive-time subspace $\mathcal S_{\Delta,+,n}$ consists of functions
all of whose derivatives vanish outside
$\{x_{1,4}>0,\ldots,x_{n,4}>0\}\setminus\Delta_n$.
This includes flatness at the time boundary, as well as at the diagonals.

The separate continuity/growth input, called H1 here, supplies continuous
linear moment functionals $\Sigma_n$ on these full diagonal-flat spaces,
agreeing with the compact smooth separated tensors, and constants
$c,C<\infty$, $k\in\mathbb N$, $\varkappa<\infty$ such that

\begin{gather*}
\FitMath{|\Sigma_n(\varphi)|\le c C^n(n!)^{\varkappa}p_{kn}(\varphi).}\tag{HB16}
\end{gather*}

This is a condition on moments, not just connected functions. It is not a
consequence of having a finite $M_n(d,R)$ for every fixed triple of
parameters. The Density lemma of Section~\ref{ch:precise-reconstruction} proves that finite linear combinations of
smooth compact separated tensors are dense in $\mathcal S_{\Delta,n}$,
and those with strictly positive-time supports are dense in
$\mathcal S_{\Delta,+,n}$, in their Schwartz subspace topologies.
These assertions include the necessary cutoffs at infinity, at the time
boundary, and near the diagonals. They cannot be replaced by ordinary
tensor density without those constraints.

Using H1 and that density lemma, the finite Schwartz sequence space
$\bigoplus_{n\ge0}\mathcal S_{\Delta,+,n}$ has form
$\sum_{m,n}\Sigma_{m+n}(\overline{\theta F_m}\otimes G_n)$.
The tensor here lies in the diagonal-flat space: an internal coincidence
is killed by the corresponding factor, and a coincidence between a negative
and a positive time argument forces the common time to zero, where the
factors are flat. Tensor multiplication is continuous in the Schwartz
topologies. Thus H1 and tensor density transfer the compact reflection
positivity and show that its quotient completion is canonically the same
$\mathcal H$. Compact $C^1$ tuples themselves are approximated by compact
smooth ones inside the same separated positive-time neighborhoods, using
(HB2). The already proved contraction and unitary groups then extend by
this identification. We use H1 for continuity on the enlarged domain. For reconstruction
in a quantitative growth class, we also use the growth bound (HB16).
The constrained tensor-density argument and the quantitative extension
estimates are given in the cited extension section.

% END MODULE chapters/euclidean_hilbert_bridge_v3.tex


% BEGIN MODULE chapters/08_short_distance.tex
\section{The short-distance coefficient and nontriviality}
\label{candidate:short_distance}
\candidatebanner{F}

\subsection*{7. Non-triviality of the limit: the two-point law of asymptotic freedom, a positive four-point function on the tori, and short-distance universality}

At short distances, the connected two-point function behaves to leading order like that of $N^2-1$ free fields. On the tori the same holds for the connected four-point function, restricted to nearly collinear quadruples, which form the class $\mathfrak T^{(4)}$ introduced in (7d). The strength of these free fields equals the running coupling evaluated at the observation scale; it is non-zero and tends to zero only logarithmically.

\textbf{Objects.} The test pairs are $(f,h)\in\mathfrak{T}_{\mathbb{E}}(\mathsf{d};\vartheta)$, that is, pairs of non-negative Lipschitz bumps, each supported in a ball of radius $\mathsf{d}$, placed at a distance from one another that lies between $\mathsf{d}$ and $6\mathsf{d}$ ((7a)). The free kernel is $\mathcal{K}_{\mathrm{free}}(f,h)=\iint f(x)h(y)|x-y|^{-8}dx\,dy$. We take the coupling $g(\mathsf{d})=g_{K-s(\mathsf{d})}(\beta)$, as well as the normalisation $\zeta_K(\mathsf{d})=\zeta'_{K-s(\mathsf{d})}(\beta)$, at the level whose spacing is smaller than $\mathsf{d}$ by a power of $\log g(\mathsf{d})^{-2}$.

\textbf{Theorem 7.1 [non-triviality; Theorem 0(c)(3)-(5)].} Assume $g\le\gamma_W(\vartheta)$. Then there exist an integer $n_0\ge1$, a threshold $\mathsf{d}_0=\mathsf{d}_0(g,\vartheta)\in(0,1)$, a function $\bar\eta:(0,\mathsf d_0]\to[0,\frac13]$ that is non-decreasing with $\bar\eta(\mathsf d)\to0$ as $\mathsf d\to0$, and, for each $\ell\in\mathfrak L(g)$, a function $\bar\varepsilon_\ell:(0,\mathsf d_0]\to[0,\frac12]$ that is non-decreasing with $\bar\varepsilon_\ell(\mathsf d)\to0$ as $\mathsf d\to0$, with the following properties whenever $\mathsf{d}\le\mathsf{d}_0$. (i) Take any first-passage datum with $K-s(\mathsf{d})\ge n_0$, any $\ell\in\mathfrak{L}(g)$ and any pair in $\mathfrak{T}_{T_\ell}(\mathsf{d};\vartheta)$. Then

\begin{gather*}
\FitMath{\langle\mathcal{O}(f);\mathcal{O}(h)\rangle_{K,\ell,\beta}=\zeta_K(\mathsf{d})^2\,\tfrac{3(N^2-1)}{\pi^4}\,g(\mathsf{d})^4\,\mathcal{K}_{\mathrm{free}}(f,h)\,(1+\varepsilon),} \\[3pt]
\FitMath{|\varepsilon|\le\tfrac12,}
\end{gather*}

and moreover $|\varepsilon|\le\bar\varepsilon_\ell(\mathsf d)$, the inequality $|\tfrac{11N^2}{12\pi^2}g(\mathsf{d})^2\ln(1/\mathsf{d})-1|\le\bar\eta(\mathsf d)$ holds, and $|\zeta_K(\mathsf{d})-1|$ is bounded by the same right-hand side; all of these bounds are uniform in $(K,\beta)$. The same identity is also valid in every periodic Gibbs state at $\beta$ on $\mathbb{L}_K$, with the same $\zeta_K(\mathsf{d})$ and $g(\mathsf{d})$ and with an error satisfying $|\varepsilon|\le\tfrac12$ ((7a)). (ii) On pairs of class $C^1$, every continuum limit obtained from (5b) on the torus $T_\ell$, as well as every one obtained from (5c) on $\mathbb{R}^4$, obeys

\begin{gather*}
\FitMath{S_2(f,h)=\tfrac{3(N^2-1)}{\pi^4}\Big(\tfrac{12\pi^2}{11N^2}\Big)^2\frac{\mathcal{K}_{\mathrm{free}}(f,h)}{(\ln(1/\mathsf{d}))^2}\,(1+\eta),} \\[3pt]
\FitMath{c\le1+\eta\le C,}
\end{gather*}

and on each $T_\ell$ one has $\eta\to0$ as $\mathsf{d}\to0$. (iii) On the tori, let $\vartheta\le\vartheta_\square$, the number from Theorem 0(c)(5), which depends only on the fixed data. Then, by Theorem 7.5 [to be proved; Package F], there exist an integer $n_4\ge n_0$, a number $\mathsf d_4\in(0,\mathsf d_0]$, a side $\ell_4(g)\in\mathfrak L(g)$, together with, for each $\ell\in\mathfrak L(g)$ with $\ell>\ell_4(g)$, a threshold $\mathsf d_4(\ell)\in(0,\mathsf d_4]$, with the following property. Suppose $K-s(\mathsf{d})\ge n_4$. Suppose $\mathsf d\le\mathsf d_4$ when $\ell\le\ell_4(g)$, and $\mathsf d\le\mathsf d_4(\ell)$ otherwise. Then for every quadruple in $\mathfrak{T}^{(4)}_{T_\ell}(\mathsf{d};\vartheta)$ ((7d)), the following bounds hold uniformly in $(K,\beta)$. Here $c,C$ are constants that are allowed to depend on $g$ as well as on the shape $\vartheta$.

\begin{gather*}
\FitMath{c\,(N^2-1)\,g(\mathsf{d})^8\prod_i\|f_i\|_\infty\le S^T_{4;K,\ell,\beta}(f_1,\dots,f_4)\le C\,(N^2-1)\,g(\mathsf{d})^8\prod_i\|f_i\|_\infty.}
\end{gather*}

Clause (iii) is a statement about the lattice functions on the tori. For four points, Theorem 0 asserts nothing about the limits, and nothing about $\mathbb{R}^4$.

\ensuremath{\dagger} The sign of $S^T_4$ for quadruples of arbitrary shape at scale $\mathsf{d}$, at every place where the box kernel lies above a floor, and, for nearly collinear quadruples and for $g$ below a further threshold, the positivity of this function both in the limits on $\mathbb{R}^4$ and on a set of ergodic components having positive measure (7d); the rigidity of the short-distance germ and its relative form (7e).

\subsubsection*{(7a) The short-distance two-point law and its asymptotic-freedom normalisation}

\textbf{Objects.} The class $\mathfrak{T}_{\mathbb{E}}(\mathsf{d};\vartheta)$ consists of the pairs of non-zero Lipschitz functions $f,h\ge0$ satisfying

\begin{gather*}
\FitMath{\operatorname{supp}f\subset B(x,\mathsf{d}),} \\[3pt]
\FitMath{\operatorname{supp}h\subset B(y,\mathsf{d}),} \\[3pt]
\FitMath{\mathsf{d}\le\operatorname{dist}(\operatorname{supp}f,\operatorname{supp}h)\le6\mathsf{d},} \\[3pt]
\FitMath{\textstyle\int f\ge\vartheta\|f\|_\infty\mathsf{d}^4,} \\[3pt]
\FitMath{\operatorname{Lip}f\le\tfrac{\|f\|_\infty}{\vartheta\mathsf{d}},}
\end{gather*}

where the final two conditions are required of $h$ as well. It follows that $c\vartheta^2\|f\|_\infty\|h\|_\infty\le\mathcal{K}_{\mathrm{free}}(f,h)\le C\|f\|_\infty\|h\|_\infty$. On $\Lambda^2\mathbb{R}^4$ the free curvature propagator is given by

\begin{gather*}
\FitMath{\mathbb{K}(u)=\pi^{-2}|u|^{-4}\,\Lambda^2\mathsf{I}(\hat{u}),} \\[3pt]
\FitMath{\mathsf{I}(\hat{u})=1-2\hat{u}\hat{u}^{\mathsf{T}},} \\[3pt]
\FitMath{\hat u=u/|u|,}
\end{gather*}

which is the two-point tensor of $dA$ when $A$ is a one-form with covariance $\delta_{\mu\nu}(4\pi^2|x|^2)^{-1}$. Since $\Lambda^2\mathsf{I}$ is an involution, $\operatorname{tr}\mathbb{K}(u)^2=6\pi^{-4}|u|^{-8}$. On $\Lambda_{K,\ell}$, let $F=dA$ be the free lattice curvature, where $d$ denotes the lattice coboundary; this curvature is the standard Gaussian on exact $2$-cochains with a single colour and unit coupling, and $\mathbb{E}_{\mathrm{free}}$ denotes the expectation with respect to it. Its background part at depth $s$ is $F_B:=\mathbb{E}_{\mathrm{free}}[F\mid Q_{K-s}(\mathbf{1})A\text{ modulo closed forms}]$, and we set

\begin{gather*}
\FitMath{C^{\mathrm{free}}_K=C^{\mathrm{bg}}_{K,s}+C^{\mathrm{fl}}_{K,s},} \\[3pt]
\FitMath{C^{\mathrm{bg}}_{K,s}(p,q)=\mathbb{E}_{\mathrm{free}}\,F_B(p)F_B(q),} \\[3pt]
\FitMath{\mathcal{K}_{K,s}(f,h)=\sum_{p,q}f(x(p))h(x(q))\,C^{\mathrm{bg}}_{K,s}(p,q)^2.}
\end{gather*}

Let $Q^{\mathrm{pl}}_j$ be the average of plaquettes, that is, of $2$-cochains, over blocks of order $j$; it is the counterpart of the straight average of $1$-forms. Because $d\,Q_j(\mathbf{1})=d\,Q^{\mathrm{st}}_j(\mathbf{1})=Q^{\mathrm{pl}}_jd$, the conditioning above is conditioning on $Q^{\mathrm{pl}}_{K-s}F$. The $C^{\mathrm{bg}}_{K,K-j}$ are therefore kernels of orthogonal projections whose ranges are nested, and the slices

\begin{gather*}
\FitMath{\mathcal{C}_{(j)}:=C^{\mathrm{bg}}_{K,K-j+1}-C^{\mathrm{bg}}_{K,K-j}} \\[3pt]
\FitMath{(C^{\mathrm{bg}}_{K,K}=C^{\mathrm{free}}_K),} \\[3pt]
\FitMath{C^{\mathrm{fl}}_{K,s}=\sum_{1\le j\le K-s}\mathcal{C}_{(j)},}
\end{gather*}

are the covariances of mutually independent pieces of $F$. For the free field, the renormalisation group thus reduces to an orthogonal decomposition.

\textbf{Statement.} The error in Theorem 7.1(i) satisfies $|\varepsilon|\le\frac12$ and, in addition, $|\varepsilon|\le\bar\varepsilon_\ell(\mathsf{d})\to0$ on each torus of the family, uniformly in the cutoff, the bare coupling, the pair and its position. The same identity, with the same $\zeta_K(\mathsf{d})$ and $g(\mathsf{d})$ and with $|\varepsilon|\le\frac12$, holds for $\langle\mathcal{O}(f);\mathcal{O}(h)\rangle_\nu$ in every periodic Gibbs state $\nu$ at $\beta$ on $\mathbb{L}_K$. The constant arises as follows: $a_p\approx\frac12|F(p)|^2$; for each colour, Gaussian half-squares have covariance $\frac12(\text{covariance})^2$; and $\frac{N^2-1}2\cdot\frac6{\pi^4}=\frac{3(N^2-1)}{\pi^4}$.

\ensuremath{\dagger} \textbf{Two variants.} (i) \emph{The one-loop coupling.} Define $\mathsf{g}(\mathsf{d}):=\big(g^{-2}+\frac{11N^2}{12\pi^2}\ln\frac1{\mathsf{d}}\big)^{-1/2}$, which depends on $\mathsf{d}$ and $g$ only. Lemma 7.2, together with $0\le s\log L-\ln\frac1{\mathsf{d}}-\ln C_W(s)<\log L$, gives

\begin{gather*}
\FitMath{\big|g(\mathsf{d})^{-2}-\mathsf{g}(\mathsf{d})^{-2}\big|\le o(s)+C+\tfrac{c_0}{\log L}\ln C_W(s)=o(s(\mathsf{d})),}
\end{gather*}

uniformly in $(K,\beta)$. Hence $g(\mathsf{d})/\mathsf{g}(\mathsf{d})\to1$, and if $\mathsf{g}(\mathsf{d})$ is put in place of the running coupling $g(\mathsf{d})$ of the construction, the law still holds below a different threshold, with error functions that differ but are of the same type. (ii) \emph{Cutoffs off first passage.} Let $K'=K(\beta,g)$, and let $b=K-K'$, which may be positive, zero or negative. The dilation identity of (3d) then takes the form

\begin{gather*}
\FitMath{\langle\mathcal{O}(f);\mathcal{O}(h)\rangle_{K,\ell,\beta}=\langle\mathcal{O}(f^{[b]});\mathcal{O}(h^{[b]})\rangle_{K',L^b\ell,\beta},}
\end{gather*}

and the map $f\mapsto f^{[b]}$ sends $\mathfrak{T}_{T_\ell}(\mathsf{d};\vartheta)$ into $\mathfrak{T}_{T_{L^b\ell}}(L^b\mathsf{d};\vartheta)$ and leaves $\mathcal{K}_{\mathrm{free}}$ invariant. Consequently, whenever $L^b\ell\in\mathfrak L(g)$, $L^b\mathsf{d}\le\mathsf{d}_0$ and $K'-s(L^b\mathsf{d})\ge n_0$, the law at $(K',L^b\ell,\beta)$ with separation $L^b\mathsf{d}$ determines the behaviour of $\mu_{K,\ell,\beta}$ at separation $\mathsf{d}$. For $K$ large and under the conditions of clause (d)(1), (3d) and (8b) give $|b|\le C(\varrho)$ for the tuned data and for each of their offsets $t\in\mathsf J$. When $b<0$, the identity is to be read in the direction from $(K',L^b\ell)$ to $(K,\ell)$.

\textbf{Argument.} (1) The local analysis is carried out on a single torus, the smallest one, $T_{\ell_\star(g)}$. On that torus the law holds by Proposition 7.3 of (7b), and its kernel there is $\frac{N^2-1}2\mathcal{K}_{K,s}$. Reflection positivity then carries it from this torus to the remaining volumes and to the Gibbs states. For the Gibbs states this is the only route available, because the expansion deals with torus measures. The local analysis can also be carried out on the other tori. (7d) does this, with four sources, on every torus of the family, but only below a threshold on $\mathsf{d}$ which may decrease with $\ell$. The transfer, by contrast, yields a single threshold $\mathsf{d}_0$ common to all of them. (2) Comparison of lattice and continuum. Let $u$ be the shortest vector representing $x(p)-x(q)$ on $T_\ell$, with $|u|\approx\mathsf{d}$, let $\sigma_{pq}:=a_K^4\pi^{-2}|u|^{-4}$, and let $\mathbb{K}(\cdot)_{pq}$ be the entry belonging to the planes of $p,q$. Then

\begin{gather*}
\FitMath{\big|C^{\mathrm{free}}_K(p,q)-a_K^4\,\mathbb{K}(u)_{pq}\big|\le C\Big[\Big(\frac{a_K}{|u|}\Big)^{1/2}+\Big(\frac{|u|}\ell\Big)^4\Big]\sigma_{pq}} \\[3pt]
\FitMath{(10a_K\le|u|\le\ell/10),} \\[3pt]
\FitMath{\big|C^{\mathrm{fl}}_{K,s}(p,q)\big|\le o(1)\,\sigma_{pq}.}
\end{gather*}

The first inequality is the proved free comparison of Section~\ref{ch:lattice-curvature-comparison}; it compares second differences of the Green's function on the lattice torus with those of the continuum Green's function. Because $a_K/\mathsf{d}\le L^{-n_0}/C_W(s)$, the first term on its right is small by the choice of $n_0$, and it tends to $0$ with $\mathsf{d}$ uniformly in $K$, since the room increases. On the reference torus the second term is at most $C(\mathsf{d}/\ell_\star)^4$, which is small with $\mathsf{d}_0$ and tends to $0$ with $\mathsf{d}$. The second inequality is (7b)(c). Approximating by a Riemann sum, at a relative cost of $Ca_K/(\vartheta^3\mathsf{d})$, yields $\mathcal{K}_{K,s}=6\pi^{-4}\mathcal{K}_{\mathrm{free}}(1+o(1))$. (3) Lemma 7.6 [outlined here; Section~\ref{sec:finite-reference-normalizer}] of (7e) transfers the law to every $T_\ell$ and every $\nu$ at a cost of $C_2\mathsf{d}^2\|f\|_\infty\|h\|_\infty$, in which $C_2=C_2(g,\vartheta)$ denotes the constant of that lemma at $n=2$. Take $\mathsf{d}\le\mathsf{d}_0(g,\vartheta)$; there Lemma 7.2 gives $\tfrac{11N^2}{12\pi^2}g(\mathsf{d})^2\ln(1/\mathsf{d})\ge\tfrac23$. In this range the cost is a fraction of the main term. That fraction goes to $0$ as $\mathsf{d}$ does, and so do the $o(1)$ of (2) and the $R$ of Proposition 7.3. The threshold $\mathsf{d}_0$ is chosen at the point where these three errors combined give $|\varepsilon|\le\tfrac12$. The main term is $\ge c\vartheta^2(\ln(1/\mathsf{d}))^{-2}\|f\|_\infty\|h\|_\infty$. Consequently, on every torus, the error goes to $0$ as $\mathsf{d}$ does.

(\ensuremath{\dagger} The same convergence holds in every periodic Gibbs state and on $\mathbb{R}^4$.)

(4) The quantities $\zeta_K(\mathsf{d})$ and $g(\mathsf{d})$ are controlled by Lemma 7.2.




\subsubsection*{(7b) Isolating a single scale: coupling and normalisation at the scale $\mathsf d$}

\textbf{Objects.} We set $s=s(\mathsf{d})$ and $k_0:=K-s$, the witness level. We also write $x:=x_{k_0}(\beta)=g(\mathsf{d})^{-2}$. The separation measured in level-$k_0$ units is $\hat{\mathsf{d}}:=L^s\mathsf{d}\in[C_W(s),LC_W(s))$. This inclusion holds for $\mathsf d<1$, in which case $s(\mathsf d)\ge1$, and it follows from the minimality of $s(\mathsf d)$ and the monotonicity of $C_W$. We write $\underline x_s:=g^{-2}+\lambda_\sharp s-c_5$ and $\bar x_s:=g^{-2}+B_\sharp(s+1)$ for the ends of the envelope of (3a); they do not depend on $(K,\beta)$, and they are non-decreasing in $s$. The room $C_W(s)=C_{W,0}\,p_0(\bar{x}_s^{-1/2})^2$ contains a constant $C_{W,0}$, chosen large after $L$, $M$ and $\vartheta$ and before $g$. The lower bounds placed on this constant are $\vartheta\,C_W(s)\ge CM$, which is used in the Strategy, and those required in (b) and (c) below. The room tends to infinity, but $\log C_W(s)=o(s)$.

\textbf{Lemma 7.2 [coupling and normalisation].} Uniformly in $(K,\beta)$, and with no torus entering (3b), one has

\begin{gather*}
\FitMath{|x_{K-s}-x_K-c_0s|\le o(s),} \\[3pt]
\FitMath{\zeta'_k=1+\sum_{j\le k}b_j'(0)\in[\tfrac23,\tfrac43],} \\[3pt]
\FitMath{|\zeta'_{K-s}-1|\le Cg_{K-s}\,p_0(g_{K-s}).}
\end{gather*}

The first bound is the flow statement of (3a). Because $s(\mathsf{d})\log L/\ln(1/\mathsf{d})\to1$, it yields $g(\mathsf{d})^{-2}/\ln(1/\mathsf{d})\to c_0/\log L=\frac{11N^2}{12\pi^2}$. For the other two bounds, note that a constant source $w\equiv\zeta$ acts as a shift of $\beta$ in which every auxiliary choice is held frozen. The increment $b_{j+1}(\zeta)$ is holomorphic for $|\zeta|<2\varrho$, and its slope is $\le C\big[g_j^3(\log x_j)^{\mathfrak{p}_0}+\varrho^{-1}g_j^{\kappa_0-4}\big]$ (4b). Since $x_j\ge x_k+\lambda_\sharp(k-j)-C$, the slopes summed over the levels $j<k$ are at most $C\big[g_k(\log x_k)^{\mathfrak{p}_0}+\varrho^{-1}g_k^{\kappa_0-6}\big]$. This bound is a function of $x_k$ alone, it tends to $0$, and for $\kappa_0\ge7$ it is $\le Cg_k\,p_0(g_k)$.

\textbf{Proposition 7.3 [the law on the reference torus].} Let $g\le\gamma_W(\vartheta)$, let $(K,\beta)$ be a first-passage datum, and suppose that $\mathsf d\le\mathsf d_0$, $K-s(\mathsf d)\ge n_0$ and $(f,h)\in\mathfrak T_{T_{\ell_\star(g)}}(\mathsf d;\vartheta)$. Then on the torus $T_{\ell_\star(g)}$, uniformly in $(K,\beta)$ and in the pair, one has

\begin{gather*}
\FitMath{\langle\mathcal{O}(f);\mathcal{O}(h)\rangle=\zeta_{k_0}'^{\,2}\,\tfrac{N^2-1}2\,x^{-2}\,\mathcal{K}_{K,s}(f,h)\,(1+R),} \\[3pt]
\FitMath{|R|\le\tfrac14,} \\[3pt]
\FitMath{R\to0\,(s\to\infty).}
\end{gather*}


\textbf{Strategy.} In this proof the pair $(f,h)$ is also written $(f_1,f_2)$. We set $k_1:=k_0+\mathsf n$ with $\mathsf n:=\lceil2\log_L\bar x_s\rceil$. Here $g$ is fixed first. Then $\mathsf{d}_0$ is chosen, small enough that $\mathsf n\le s(\mathsf d)$ whenever $\mathsf d\le\mathsf d_0$. This gives $k_1\le K$. Moreover $x_j\approx x$ for $k_0\le j\le k_1$, because $B_\sharp\mathsf n\ll x$. The sourced construction (4a) is carried out up to a level $k\le k_1$, with the sources measured at depth $K-k$. At depth $s$ one has $\|f\|_{\sharp,s}=\|f\|_\infty$, because $\vartheta\,C_W(s)\ge CM$ for $\mathsf d\le\mathsf d_0(g,\vartheta)$. Above the witness level, $\|f\|_{\sharp,K-k}\le L^{k-k_0}\|f\|_\infty$, so the disc in $z$ becomes smaller. The identities for the derivatives at $z=0$ remain valid, and in the bounds each derivative in a source costs a factor $L^{k-k_0}$. One then has $Z(z)/Z(0)=\mathbb{E}_{\mathsf{P}^{(k)}}\,e^{\Phi^{(k)}(z)}$. In what follows we write $\mathsf{P}=\mathsf{P}^{(k_0)}$, which is a law distinct from the law $\mathsf{P}_{k_0}$ of Proposition 4.4 that appears further on, and $\Phi=\Phi^{(k_0)}$. Everything stated about these two objects in this paragraph holds at each level $k\le k_1$, with $k$ in place of $k_0$ and $K-k$ in place of $s$. $\mathsf{P}$ is a probability law on triples. Each triple consists of a class $[h]$ of histories of level $k_0$ sharing one standing part $h^\flat$ (2a), a block field $V\in\mathcal{V}_{k_0}$, and the stored fields. The law is determined by the unsourced sums $\tau_{[h]}:=\sum_{h'\in[h]}\tau_{h'}$. Here $\tau_{[h]}(dV)$ is an integral over the stored fields, which are the fields left unintegrated at level $k_0$. Its integrand is the exponential of the exponent of the sum, integrated against a positive measure composed of the gauge-fixing factors, the characteristic functions, the delta functions of the averaging constraints and Haar measures. Dividing this integrand by $Z(0)$ yields the law on triples. $\Phi$ is the difference between the sourced and the unsourced exponent of $\tau_{[h]}$. An exponent is attached to such a sum and not to an individual $\tau_h$ (2a), (2e). The marginal of $\mathsf{P}$ in $([h],V)$ equals $Z(0)^{-1}\tau_{[h]}(dV)$. Tested on gauge-invariant functions, this marginal agrees with the push-forward under $h\mapsto[h]$ of the $(h,V)$-marginal of $\mathsf{P}_{k_0}$, the law on triples (bare field, history, block field) of Proposition 4.4 (4a). Projected further onto $[h]$, the marginal is the image of $\tilde\nu$. Both laws therefore extend one and the same law of $[h]$ and of the gauge orbit of $V$, but they add different variables. $\mathsf{P}_{k_0}$ adds the bare field and $h$, the particular member of the class. $\mathsf{P}$ adds the stored fields, on which the part of $\Phi$ attached to the live components depends. With the notation $\Phi^{(k)}_B:=\prod_{i\in B}\partial_{z_i}\Phi^{(k)}|_{z=0}$ we have

\begin{gather*}
\FitMath{\langle\mathcal{O}(f_1);\mathcal{O}(f_2)\rangle=\langle\Phi^{(k)}_1;\Phi^{(k)}_2\rangle_{\mathsf{P}^{(k)}}+\mathbb{E}_{\mathsf{P}^{(k)}}\Phi^{(k)}_{12}} \\[3pt]
\FitMath{(1\le k\le k_1).}
\end{gather*}

All of these are expressions for one and the same quantity. At $k=k_0$ the first term on the right is the main term and the second term is an error ((a), (b)). At $k=k_1$ the roles are exchanged, and the law is obtained at that level ((c)). Both $\Phi^{(k)}_i$ and $\Phi^{(k)}_{12}$ are derivatives with respect to the sources. The argument uses conservation of mass alone, $\sum_h\tau^w_h(\mathcal{V}_{k})=Z(w)$. In other words, it uses exactness with the sources measured at depth $K-k$, which is the depth $s$ when $k=k_0$ (Proposition 4.3 of (4a), together with (4b)). This is what isolates the scale; no conclusion is drawn from $\mathsf{P}_{k_0}$, the coupling with the bare field. The law of $\mathcal{M}^{k_0}U$ is not identified with the marginal of $\mathsf{P}$ in $V$, because inside what $\mathbf{R}$ heals it replaces the law of the block field. Under $\mathsf{P}_{k_0}$, the equality $V=\mathcal{M}^{k_0}U$ is valid only on bonds at a large distance from the components that $\mathbf{R}$ has acted on. For this reason the sources are carried through the steps, and no observable is conditioned.

(a) \emph{Locality.} The source enters the exponent in three ways. \emph{Pointed pieces:} away from the live components, the coupling term is a sum of functions, each of which depends on a single value of $w$; this holds because the running shift at a plaquette is a finite sum of functions of this kind. As $f_1f_2=0$, these pieces give no mixed derivative. \emph{Regular local terms:} these are indexed by $j\le k_0$ and $X\in\mathcal{D}_j$, and they depend on $w$ only on $X^\sim$ (4b). Every domain that meets both supports satisfies $d_j(X)\ge c\,\hat{\mathsf{d}}/M$. By Cauchy bounds on the source disc, the contribution of these terms to $\Phi_{12}$ is bounded pointwise by $C\hat{\mathsf{d}}^4e^{-c\kappa\hat{\mathsf{d}}/M}\prod\|f_i\|_\infty$. This is of order $e^{-c\,p_0^2}$, where, in this step, $p_0$ without an argument stands for $p_0(g_{k_0})$. It is therefore smaller than any power of $g(\mathsf{d})$. \emph{The part attached to the live components:} only a cruder bound is used for this part. Consider a class $[h]$ in which a live component is in contact with a support, that is, lies within distance $c\,\mathsf d$ of it. The label is common to the class, and hence each live component, its lineage and the age of that lineage are common to the class as well; the members of the class differ only in healed components. Let $a(h)$ denote the largest lineage age among such components (Proposition 4.14 of (4d)). On such a class the mixed derivative of this part is bounded by

\begin{gather*}
\FitMath{(\log x)^C\Big[e^{-c\,\hat{\mathsf{d}}}+L^{C\check{R}_{k_0}}L^{8a(h)}\Big]\prod\|f_i\|_\infty.}
\end{gather*}

It is the mean that is small: $|\mathbb{E}_{\mathsf{P}}\Phi_{12}|\le o(1)\,x^{-2}\prod\|f_i\|_\infty$. This follows from the age series of (4d) (Proposition 4.14 and its proof), taken with rate $b\ge8$. The growth of the amplitude and the decay of the series are governed by one and the same age, namely the age of the lineage. The age recorded in the label restarts wherever $\mathbf{R}$ leaves a component standing as an island, and it is not larger. The number of cells of level $k_0$ in contact with the supports is $\le C\hat{\mathsf{d}}^4$, so

\begin{gather*}
\FitMath{\mathbb{E}_{\mathsf{P}}\Big[\mathbf{1}\{\text{contact}\}\,L^{8a(h)}\Big]\le C\hat{\mathsf{d}}^4\sum_{a\ge0}L^{8a}\max_{\mathfrak{c}}\tilde\nu\big(\mathfrak{lv}^{\ge a}_{k_0}(\mathfrak{c})\big)\le C\hat{\mathsf{d}}^4e^{-c\,p_0(g(\mathsf{d}))},}
\end{gather*}

and this bound overcomes $(\log x)^CL^{C\check{R}_{k_0}}x^2$ because $\mathfrak{p}_0\gg r$. \emph{At level $k_0$ the two-point function draws on this series twice}, both times at a rate $b\ge8$. At level $k_1$, in (c), it draws on the series once, at the same rate. The first use is the present one, the bound on the mean of the live part's mixed derivative. The second occurs in (b), in the covariances of the live first derivatives; there the product of two such derivatives carries $L^{4a}\cdot L^{4a}$. At level $k_0$ these two uses are the source of the weight $L^{8a}$ for two points. At level $k_1$ its source is (c1). Of the growth, only one property enters: as a function of the age, it is no faster than geometric. A different exponent, or an additional factor polynomial in $a$, is handled by a larger rate $b$, that is, by a larger $c_a$, a constant that may be increased freely. Hence the theorem on the ages (Proposition 4.14) underlies Theorem 0 (c)(3) and (c)(4). It also underlies (c)(5), for which (7d) employs this series with a larger rate.

(b) \emph{The marginal part.} Here the coupling carries the source. Let $h^{(j)}_y\ge0$, with $\sum_yh^{(j)}_y=1$, be the tent functions of (4a), (4b). These weights have range $CMa_{K-j}$ and serve to localise the counterterm of step $j$. They are written $h_z$ in (0b) and in Step 1, in which the level is held fixed. Let $p_y$ be a bare plaquette at $y$. At points where every profile $\varphi_j$, $j\le k_0$, takes the value $1$, that is, off the live components, the running shift satisfies, by its formula in (4b),

\begin{gather*}
\FitMath{\partial_{z_i}\zeta_{k_0}(p;w)\big|_0=f_i(x(p))+\sum_{j\le k_0}b_j'(0)\sum_{y\in\mathbb{T}^{(j)}}h^{(j)}_y(p)\,f_i(x(p_y))=:\zeta'_{k_0}\tilde{f}_i(p),} \\[3pt]
\FitMath{|\tilde{f}_i-f_i|\le\tfrac{CM}{\vartheta\hat{\mathsf{d}}}\sum_j|b_j'(0)|\,\|f_i\|_\infty;}
\end{gather*}

the factor $\zeta_K(\mathsf{d})^2$ arises here. Differentiating $-A(x_{k_0}(\cdot)-\zeta_{k_0}(\cdot\,;w),U_{k_0})$ gives

\begin{gather*}
\FitMath{\Phi_i=\text{const}+\zeta'_{k_0}X_i+(\text{irrelevant})+(\text{live})+(\text{far}),} \\[3pt]
\FitMath{X_i(V):=\sum_p\tilde{f}_i(p)\,\chi_{k_0}(\square(p))\,a_p\big(U_{k_0,\square(p)}(V)\big).}
\end{gather*}

In this formula $\square(p)$ is the test cube containing $x(p)$, and $U_{k_0,\square(p)}$ is the background localised about that cube, hence near $\operatorname{supp}f_i$. The test $\chi_{k_0}(\square(p))$ acts as a window: a plaquette contributes only if its localised background passes the plaquette test. Consequently $X_i$ depends on $V$ only near the support and is the same for every history. It also satisfies $|X_i|\le Cx^{-1}p_0^2\,\hat{\mathsf{d}}^4\|f_i\|_\infty$. Away from the large-field regions, the global background differs from the localised one by an amount exponentially small in the radius of localisation; this difference is the term (far). The size of the irrelevant terms is $\le C(\log x)^{3\mathfrak{p}_0-q}x^{-1}\hat{\mathsf{d}}^4\|f_i\|_\infty$. Per cube they are of size $\alpha^2(\varepsilon_{k_0}/\alpha)^3$, where $\alpha$ is of the order of the analyticity radii $\alpha_{0,k_0},\alpha_{1,k_0}$. This is a third-order Taylor remainder multiplied by the magnitude that the source term ($w$ times $a_p$) attains when the fields range over their analyticity domain. The factor $x^{-1}$ cannot be dispensed with. The law uses $p_0^2\,\hat{\mathsf{d}}^8(\log x)^{3\mathfrak{p}_0-q}\to0$ as the coefficient of $x^{-2}$, and Cauchy bounds over the source radius $\varrho$ yield no power of $x$. The factor comes from the bound on the regular family produced by a step ((2e), ``Dependence of the regular family on the array''). That bound also holds for shifts of the coupling that vary with position, and it weights a difference of shifts by $CM^4\alpha_{0,k}^2$. We apply it to the sourced and the unsourced arrays of (4b). The remainders enter through the slope towards the constants carried by the induction of (4b), and they are small because of their size. The result is a linear recursion for the slopes in the source, and the small factor in front of the bound closes this recursion over the levels. Its solution is $\approx x_j^{-1}(\log x_j)^{\mathfrak{p}_0+q}$, and this takes the place of $\varrho^{-1}$. The oscillating part of the source, namely what remains after subtracting from the source its value at a single point of the domain, has a slope bounded in the same way, with an extra factor $L^{-(k_0-j)}$. The live terms become larger as the age increases. If a live component has $\operatorname{age}^{\mathrm{lab}}=a$, the fields stored in it are those of every level from $k_0-a$ onwards, together with those of the next finer level. For a component holding an island, the block of which retains the fields of its core (2b), the stored fields begin at the first level of its lineage. Over one cube of level $k_0$, the plaquette sums of these fields contain at most $L^{4(a+1)}$ terms. For the live terms, only their first two moments on the events of contact enter, so that, as in (a), it suffices that their growth be geometric in the age. Let $h$ be a history in which a live component is in contact with $\operatorname{supp}f_i$. On such a history the live contribution to $\Phi_i$ is bounded by $(\log x)^C\big[e^{-c\,\hat{\mathsf{d}}}+L^{C\check{R}_{k_0}}(a(h)+2)^CL^{4a(h)}\big]\|f_i\|_\infty$. The power of the age in this bound is the one of (7d). In the bound $Ce^{-cp_0(g_j)}(2L^b)^{-a}$ of Proposition 4.14, the factor $2^{-a}$ absorbs this power, so the age series is used at the same rate. This live contribution has two covariances. One is with the live contribution of the other factor, with weight $L^{8a(h)}$. The other is with those functions of $V$ that occur in the other factor, with weight $L^{4a(h)}$, taken against their supremum bound. By the age series of (4d), as in (a), both of these covariances are $o(1)\,x^{-2}\prod\|f_i\|_\infty$.

(c) \emph{Gaussian extraction: the free template.} Denote by $\mathcal{H}_jB$ the exact bare $2$-cochain of smallest norm that satisfies $Q^{\mathrm{pl}}_j\mathcal{H}_jB=dB$. It is the curvature of the linear minimiser, namely $d\,\mathsf{H}_j$ at $U=\mathbf{1}$, a point at which $J=0$ and the three Green's functions of (1c) agree. We then have

\begin{gather*}
\FitMath{\mathcal{C}_{(j)}=\mathcal{H}_{j-1}C^{\mathrm{tr}}_j\mathcal{H}_{j-1}^*,} \\[3pt]
\FitMath{|\mathcal{H}_j(p,b)|\le C\eta_j^2e^{-\delta|x(p)-x_b|/a_{K-j}},} \\[3pt]
\FitMath{|\mathcal{C}_{(j)}(p,q)|\le C\eta_{j-1}^4e^{-\delta|x(p)-x(q)|/a_{K-j+1}}.}
\end{gather*}

The decay of $\mathcal{H}_j$ comes from Ba{\l}aban's work, apart from the passage to centred blocks and to level-$j$ units described below. In that work the linear minimiser is given by $H_k=GQ^*(QGQ^*)^{-1}$ (Comm. Math. Phys. 95 (1984) 17, (1.103); at $U=\mathbf{1}$ the straight average coincides with the $Q_k$ of that paper and of 96 (1984) 223), and both $G$ and $\nabla G$ decay exponentially, at a rate that depends on the dimension alone (95, Proposition 1.2 [cited from Ba{\l}aban; Package A--H], (1.110)). In 95 each block has $y$ as one of its corners ((1.6) there); from 109 onwards Ba{\l}aban's blocks are centred at $y$ ((0.3) there), like the blocks of the present construction. For centred blocks the same bounds follow from three devices, once one passes to level-$j$ units, and this passage is the origin of the factor $\eta_j^2$. (i) A spectral gap that is uniform in the fine spacing $\eta_j$, equivalently in the number $L^{4j}$ of fine sites in a block: the operator built from the mass of the averages (which depends on the block means alone), the Landau gauge term (a non-local projection) and the curl form is $\ge c(-\Delta+1)$. (ii) A Combes-Thomas conjugation with the exponential weights $e^{\delta|\cdot-x_b|}$, which yields the decay of the minimiser only in weighted $L^2$. (iii) A bootstrap that upgrades control in weighted $L^2$ to control of the first differences in weighted sup norms, as the pointwise bound on $\mathcal{H}_j$ requires. Uniformly in $\eta_j$, the kernel of $(-\Delta+1)^{-1}$ and its first differences belong to weighted $L^p$, the former for $p<2$ and the latter for $p<\frac43$ (in four dimensions the singularities are $|z|^{-2}$ and $|z|^{-3}$). Young's inequality is then applied twice, $L^2\to L^8\to L^\infty$. What results is elliptic regularity on the lattice with bounds independent of the fine spacing. The symbol $C^{\mathrm{tr}}_j$ denotes the single-step covariance in a tree gauge,

\begin{gather*}
\FitMath{C^{\mathrm{tr}}_j:=\iota\,\big(\iota^*\mathcal{H}_{j-1}^*\mathcal{H}_{j-1}\,\iota\big)^{-1}\iota^*,} \\[3pt]
\FitMath{d'\,C^{\mathrm{tr}}_j\,d'^*=C^{\mathrm{curv}}_j,}
\end{gather*}

where $\iota$ embeds into the $1$-forms $B$ of level $j-1$ the subspace of those that satisfy $B=0$ on one contour for each site (a tree in each block, imposed sharply) together with $Q^{\mathrm{st}}B=0$. In this formula $Q^{\mathrm{st}}$ denotes the straight average over the blocks of the single step from level $j-1$ to level $j$, at $U=\mathbf 1$. It is thus the covariance, on that subspace, of the quadratic form $\langle B,\mathcal{H}_{j-1}^*\mathcal{H}_{j-1}B\rangle=\|\mathcal{H}_{j-1}B\|^2$, which equals twice the free action of the minimiser with datum $B$. On the subspace one has $\|\mathcal{H}_{j-1}B\|^2\ge c\|d'B\|^2\ge c_L\|B\|^2$. It differs from the covariance $C^{(j-1)}$ of (1c)-(1d) at the trivial background in two ways. The gauge differs: there it is imposed softly, on the means over all contours. The constraint also differs: the straight average is not the linearisation of $\mathcal{M}$. The two agree in the curvature covariance $d'\cdot d'^*$, where $d'$ is the coboundary of level $j-1$. Indeed, given its block averages, the level-$(j-1)$ curvature $Q^{\mathrm{pl}}_{j-1}F$ has conditional covariance $C^{\mathrm{curv}}_j$. Since $\mathcal{H}_{j-1}B$ depends on $B$ only through $d'B$, it determines $\mathcal{C}_{(j)}$, and $a_p$ depends on nothing else. Summing over $j\le k_0$, where the coarsest slice gives the dominant contribution, one obtains the following bound. Here $t=|x(p)-x(q)|/a_s\ge\frac12C_W(s)$, and $\delta_{\mathrm{fl}}>0$ is a rate formed from the rates $\delta$ of the slices, depending only on $L$ and the dimension.

\begin{gather*}
\FitMath{\big|C^{\mathrm{fl}}_{K,s}(p,q)\big|\le C\big(\tfrac{a_K}{a_s}\big)^4e^{-\delta_{\mathrm{fl}}t}=C\pi^2t^4e^{-\delta_{\mathrm{fl}}t}\,\sigma_{pq}:}
\end{gather*}

the fluctuation part is small because the curvature of the linear minimiser decays. Because $t^4e^{-\delta_{\mathrm{fl}}t}$ is decreasing only for $t\ge4/\delta_{\mathrm{fl}}$, the room is given the floor $C_{W,0}A_0^2\ge8/\delta_{\mathrm{fl}}$. Over $j>k_0$ the sum of the exponentials reconstitutes the power law: $\sum_iL^{-4i}e^{-tL^{-i}}\le Ct^{-4}$.

\emph{The interacting measure.} We apply the formula of the Strategy at level $k_1$, whose backgrounds are flat on the scale of the pair. Only the classes $[h]$ occur, and no finer class of histories: the sourced exponent of level $k_1$ is the one that (4b) provides. (c1) \emph{The first member is an error term.} Consider the terms of the levels $k_0<j\le k_1$ and a source that is rough on their scale. Their slopes in such a source remain $\le Cx^{-1}(\log x)^{\mathfrak{p}_0+q}$, measured against $\|f_i\|_\infty$ (the norm $\|f_i\|_{\sharp,K-j}$ of the level of the term can exceed it by a factor of up to $L^{j-k_0}$). The recursion for these slopes closes since $\mathsf n\le C\log x$ (its size is set through $\bar x_s\le Cx$). Hence $\mathsf n^2(\log x)^{\mathfrak{p}_0-q}\le C(\log x)^{2+\mathfrak{p}_0-q}$ is small ($q-\mathfrak p_0>2$, $g$ small). Suppose first that no large-field region touches the supports at level $k_1$. Then $|\Phi^{(k_1)}_i-c_i|\le o(1)\,x^{-1}\|f_i\|_\infty$ for a constant $c_i$, where $o(1)$ refers to $\mathsf d\to0$, uniformly in the datum. In the other case $|\Phi^{(k_1)}_i-c_i|$ grows no faster than $L^{\mathsf n}L^{4a}\|f_i\|_\infty$ in the age $a$ of the touching region. The age series of (4d), taken at the single level $k_1$ with $b\ge8$, beats this growth: $L^{\mathsf n}\le L\bar x_s^2$ is a power of $x$, while the series lies below every power. Consequently $|\langle\Phi^{(k_1)}_1;\Phi^{(k_1)}_2\rangle_{\mathsf{P}^{(k_1)}}|\le o(1)\,x^{-2}\prod\|f_i\|_\infty$. (c2) \emph{The second member becomes one number.} The quantity $\Phi^{(k_1)}_{12}$ is the sum of two parts. The first is the mixed derivative of the sum $\sum_{j\le k_1}\sum_Xe_j(X;w)$ of localised source energies from (4c). The second is a part that depends on the fields and is bounded only crudely, by a quantity of order $L^{2\mathsf n}$ for each bond of the determining set; by (4c), the size of that set is counted by the unit volumes of level $k_1$ together with the sites $\sum_{j\le k_1}\#_j(Z_j)$ of the large-field regions (every source derivative costs a factor $L^{\mathsf n}$). We translate the pair by the points $v$ of a grid $\mathfrak P\subset\Lambda_{s,\ell}$ whose spacing is comparable to $\mathsf{d}$, and we write $f^v_i:=f_i(\cdot-v)$ for the translated sources. The left-hand side of the formula is unchanged by this. The sum over $v$ of the mixed derivatives equals the mixed derivative evaluated at the two array sources $\sum_v\epsilon_vf^v_i$, whose norms are those of $f_i$, averaged over the signs $\epsilon\in\{\pm1\}^{\mathfrak P}$ (this is the device of Lemma 7.6 [outlined here; Section~\ref{sec:finite-reference-normalizer}] of (7e)). The crude bound is therefore divided by the number $\approx L^{4\mathsf n}\hat{\mathsf{d}}^{-4}$ of pairs contained in a unit volume of level $k_1$. What survives is $L^{-2\mathsf n}\hat{\mathsf d}^4\le\bar x_s^{-4}\hat{\mathsf d}^4$, and this requirement fixes $\mathsf n$. There remains the number 

\begin{gather*}
\FitMath{c_{12}:=|\mathfrak P|^{-1}\sum_{v\in\mathfrak P}\partial_{z_1}\partial_{z_2}\sum_{j\le k_1}\sum_Xe_j(X;z_1f^v_1+z_2f^v_2)|_0}
\end{gather*}

. (c3) \emph{The number is evaluated.} The number $c_{12}$ is the value at a single point of the mean of $\Phi^{(k_1)}_{12}$ over $\mathfrak P$, namely of the term without large-field regions, taken at the trivial block field. It satisfies an exact single-step recursion at the flat background, leading from level $k$ to level $k+1$ for $k_0\le k<k_1$. The measures there are positive, so a remainder is carried through a step by its supremum alone. For each step, two inequalities on a single fluctuation integral (1d) are then required. That integral belongs to the small-field step taken over the whole lattice, with a background that is real and regular and with coupling $x_j\approx x$. The first inequality concerns the covariance of two sums of background plaquette energies. The second concerns the mean value of a bilinear form whose arguments are the background plaquette angles. The expansion of (1d), together with Cauchy bounds on a large polydisc, removes the interaction. A Gaussian computation remains, and Taylor's formula compares it with the flat background. At the flat background, the object that agrees with the free template is the curvature covariance $C^{\mathrm{curv}}_j$ introduced above. It determines $\mathcal{C}_{(j)}$, and $a_p$ depends on nothing else. Summed over $k_0<j\le k_1$, the $\mathcal{C}_{(j)}$ give $C^{\mathrm{bg}}_{K,s}$ minus $C^{\mathrm{bg}}_{K,s-\mathsf n}$, and the subtracted term is $\le CL^{-4k_1}$. Three facts are used here. The source arrives at level $k_0$ in the form $\zeta'_{k_0}\tilde{f}_i$ ((b)). Next, $a_p\approx\frac12|F(p)|^2$, where $F(p)$, an element of the Lie algebra, is the plaquette angle at $p$ of $U_{k_0,\square}(V)$ (the same angle is written $Y_p$ in (7d)). Finally, half-squares of Gaussian variables have covariance $\frac12(\text{covariance})^2$ for each colour. We obtain

\begin{gather*}
\FitMath{\Big|\langle\mathcal{O}(f_1);\mathcal{O}(f_2)\rangle-\zeta_{k_0}'^{\,2}\,\tfrac{N^2-1}2\,x^{-2}\,\mathcal{K}_{K,s}(\tilde{f}_1,\tilde{f}_2)\Big|\le o(1)\,x^{-2}\prod\|f_i\|_\infty.}
\end{gather*}

The formula at level $k_0$, combined with (a) and (b), gives the same estimate for $\zeta_{k_0}'^{\,2}\langle X_1;X_2\rangle_{\mathsf{P}}$ in the average over $\mathfrak P$.
\emph{Positions.} The cumulants are invariant under the translations of $\Lambda_{K,\ell}$. A translation $v\in\Lambda_{s,\ell}$ carries blocks onto blocks, which finer translations do not do. Hence

\begin{gather*}
\FitMath{C^{\mathrm{bg}}_{K,s}(p+v,q+v)=C^{\mathrm{bg}}_{K,s}(p,q),} \\[3pt]
\FitMath{\mathcal{K}_{K,s}(f,h)=\frac1{|\Lambda_{s,\ell}|}\sum_{v\in\Lambda_{s,\ell}}\mathcal{K}_{K,s}\big(f(\cdot-v),h(\cdot-v)\big):}
\end{gather*}

since both sides of the law are invariant, it is enough to make the comparison in the average over such positions.

(d) \emph{Wild histories.} For two points, the relevant event is that a large-field region touches the supports at the single level $k_1$ ((c1)). The weight of these histories is $\le C(M\check{R}_{k_1})^4e^{-c\,p_0(g_{k_1})}$, and no union over the levels is taken. In this bound the prefactor counts the sites of level $k_1$ in one cell of that level, and $O(1)$ cells touch the supports. For four points (7d), the unclean histories are the ones that matter. At each level $j$ of the segment of (7d), $k_0-O(L\check{R}_{k_0})\le j\le k_0+r_4$, the number of cells of level $j$ near the supports, that is, within $C\hat{\mathsf{d}}$ cells of level $k_0$, is $\le C\hat{\mathsf{d}}^4L^{4\max\{k_0-j,0\}}$; altogether there are $\le e^{CL\check{R}_{k_0}}$ such cells. A component on which $\mathbf{R}$ acts at a level $j$ of the segment contains a component that is live at level $j-1$. The one-cell bounds of (4d) give $\tilde\nu(\mathfrak{c}\subset\mathcal{Z}_j)\le Ce^{-c\,p_0(g_j)}$, and the same bound for creation at step $j$. Moreover $p_0(g_j)\ge c\,p_0(g(\mathsf{d}))$, because $x_j\ge x^{1/2}$. Hence $\tilde\nu(\text{unclean})\le e^{CL\check{R}_{k_0}-c\,p_0(g(\mathsf{d}))}$, and this is small because $\mathfrak{p}_0\gg r$. Even here these one-cell bounds must hold uniformly in the volume (prolonged run, translation covariance up to the level, first moment of the prepaid end mark: (4d)). The reason is that $T_{\ell_\star}$ contains $(\ell_\star/Ma_s)^4$ cubes of level $k_0$, and this number is unbounded.

\textbf{What must be overcome.} The natural size of a factor is $x^{-1}\|f_i\|_\infty$. The bound on $X_i$ is larger by $p_0^2\,\hat{\mathsf{d}}^4=(\log x)^{C\mathfrak{p}_0}$, and the irrelevant terms are controlled by sup bounds only. This is the reason for $q\gg\mathfrak{p}_0$ and for the super-polynomial smallness in (a), (d). The effective action is that of Ba{\l}aban (Comm. Math. Phys. 109 (1987) 249; 122 (1989) 175, 355); the sources and the extraction are added to it.

\paragraph*{The setting and the objects}

In (7a) the two-point law is stated for two bumps $f,h$ whose size and mutual separation are both $\approx\mathsf d$. They are placed on a lattice of spacing $a_K\ll\mathsf d$ inside a volume of side $\ell\gg\mathsf d$. Isolating the scale $\mathsf d$ has two aspects. \emph{From below}: $K-s(\mathsf d)$ renormalisation steps integrate out the scales between $a_K$ and $\mathsf d$. In the leading term these scales survive only through two numbers of the run: $g(\mathsf d)$, the coupling, and $\zeta_K(\mathsf d)$, the normalisation. \emph{From above}: comparing $T_\ell$ with a single fixed torus eliminates the scales from $\mathsf d$ up to $\ell$, together with the volume itself. This comparison costs $O(\mathsf d^2)$, which is of lower order than the main term $\approx(\ln\frac1{\mathsf d})^{-2}$. This part defines the two numbers, bounds them uniformly in the cutoff and in the bare coupling, and carries out the comparison of volumes. The Gaussian main term is extracted in (7a), at the level $K-s(\mathsf d)$ on the fixed torus. There the cumulant is written in terms of a law on resummed histories and block fields (Lemmas 7a.1.8 [to be proved; Package F], 7a.1.9). Resummed histories are classes of histories that become equal once their entries inside healed components are deleted. The large-field operation $\mathbf R$ sums over the outcomes inside the components it heals, so an exponent is attached to such a class rather than to an individual history; the class is the index $\mathsf h$ of Remark 7b.11\ensuremath{'} (ii).

\subsubsection*{(7c) The law passed to the limits}

\textbf{Corollary 7.4 [Theorem 0(c)(4)].} Let $S$ be a continuum limit, either on $T_\ell$ of (5b) or on $\mathbb{R}^4$ of (5c). Let $\mathsf{d}\le\mathsf{d}_0$, and let $(f,h)\in\mathfrak{T}_{\mathbb{E}}(\mathsf{d};\vartheta)$ be a pair of class $C^1$. Then the conclusion of Theorem 7.1(ii) is valid with

\begin{gather*}
\FitMath{(1-\bar\varepsilon)\big(1-\bar\eta(\mathsf{d})\big)^4\le1+\eta\le(1+\bar\varepsilon)\big(1+\bar\eta(\mathsf{d})\big)^4,}
\end{gather*}

in which $\bar\eta$ and $\bar\varepsilon_\ell$ denote the functions of Theorem 7.1. Their values lie in $[0,\frac13]$ and $[0,\frac12]$ respectively, and both tend to $0$ with $\mathsf d$. The quantity $\bar\eta(\mathsf{d})$ is a common bound for $|\zeta_K(\mathsf{d})-1|$ and for $|\frac{11N^2}{12\pi^2}g(\mathsf{d})^2\ln\frac1{\mathsf{d}}-1|$. On $T_\ell$ one has $\bar\varepsilon=\bar\varepsilon_\ell(\mathsf{d})$, and on $\mathbb{R}^4$ one has $\bar\varepsilon=\frac12$. Consequently $S_2(f,h)>0$, and on $\mathbb{R}^4$ we have $\dim\mathcal{H}(S)\ge2$.

\textbf{Argument.} The argument follows from the uniformity in Theorem 7.1(i). In the lattice law, the only quantities that depend on $(K,\ell,\beta)$ are $\zeta_K(\mathsf{d})=1+\eta''$, $g(\mathsf{d})^2\propto(1+\eta')/\ln(1/\mathsf{d})$ and $\varepsilon$. Therefore

\begin{gather*}
\FitMath{\frac{S^T_{2;K,\ell,\beta}(f,h)\,(\ln(1/\mathsf{d}))^2}{\frac{3(N^2-1)}{\pi^4}(\frac{12\pi^2}{11N^2})^2\mathcal{K}_{\mathrm{free}}(f,h)}=(1+\eta'')^2(1+\eta')^2(1+\varepsilon)}
\end{gather*}

takes its value in the closed interval of the display above, for every datum. For fixed $\mathsf{d}$, we have $K_j-s(\mathsf{d})\ge n_0$ once $j$ is large. Every pair on $\mathbb{R}^4$ is contained in some ball of radius $5\mathsf{d}$. Hence, once transported to $T_{\ell_j}$, it belongs to the class and keeps the same free kernel. A closed condition is preserved in the limit. The bound on $\varepsilon$ that clause (c)(4) records on $\mathbb{R}^4$ is the only one in Theorem 7.1(i) that is free of any side, namely $|\varepsilon|\le\frac12$.

\textbf{The Hilbert space.} Let $f_{\mathsf{d}}$ be a $C^1$ bump whose support is the ball $B(2\mathsf{d}e_4,\mathsf{d})$. Then $(\theta f_{\mathsf{d}},f_{\mathsf{d}})$ is a test pair. Its shape $\vartheta_1$ is fixed and does not depend on $\mathsf{d}$, because the bumps are dilates of a single profile. When $\vartheta>\vartheta_1$, the pair may fail to belong to $\mathfrak{T}_{\mathbb{R}^4}(\mathsf{d};\vartheta)$. We take the ceiling $\gamma_W(\vartheta)$ smaller than the ceiling under which Theorem 7.1(i) is valid at the shape $\vartheta_1$. That theorem is applied there with the witness depth $s_1(\cdot)$, which is built from the room constant $C_{W,0}$ of that shape, and with the thresholds of the same shape. Once $\mathsf{d}$ is below its threshold, the three inequalities $\zeta'_{K-s_1(\mathsf{d})}\ge\frac23$, $1+\varepsilon\ge\frac12$ and $g_{K-s_1(\mathsf{d})}^{-2}\le\bar{x}_{s_1(\mathsf{d})}$ bound $S^T_{2;K_j,\ell_j,\beta_j}(\theta f_{\mathsf{d}},f_{\mathsf{d}})$ from below by a positive number that does not depend on $j$. Consequently (6d) gives

\begin{gather*}
\FitMath{\big\|[(f_{\mathsf{d}})]\big\|^2=\sigma_S(\theta f_{\mathsf{d}},f_{\mathsf{d}})=S_2(\theta f_{\mathsf{d}},f_{\mathsf{d}})>0,} \\[3pt]
\FitMath{\langle\Omega,[(f_{\mathsf{d}})]\rangle=\sigma_S((f_{\mathsf{d}}))=0.}
\end{gather*}


\textbf{Periodic Gibbs states.} The expansion concerns torus measures. The law is carried over to these states by Lemma 7.6. That lemma relies on Lemma 4.10 [proved here; Section~\ref{ch:all-state-source-transfer}] (4d), according to which the chessboard dominates the periodic Gibbs states. Precisely: let $F$ be a function of a single board cell of side $\mathsf{l}=O(1)$, and on $T_\ell$ let the chessboard seminorm be denoted $\|\cdot\|_{\mathrm{ch},\ell}$. Then, along any compatible sequence of tori (ST3),

\begin{gather*}
\FitMath{|\nu(F)|\le\liminf_{\ell\to\infty}\|F\|_{\mathrm{ch},\ell}.}
\end{gather*}

This is a separate lemma. It is not the chessboard estimate, whose subject is torus measures. Its proof, taken in the form that involves all mirror copies, combines four ingredients: the reflections of the cells leave $\nu$ invariant (10a); Jensen's inequality is used together with the Gibbs property; the boundary condition is changed at the cost of a surface term, which is weighed against a volume; and finally the seminorm enters.

\textbf{Ergodic components.} At $(K_j,\beta_j)$, each periodic Gibbs state obeys the law. A closed condition that all of these states satisfy is inherited by the components (10c). This yields Theorem 0'(3)(ii): for a volume-first limit, each of its ergodic components satisfies the law in the form it takes on $\mathbb{R}^4$, with $\mathsf{d}_0$ and $\bar\eta$ unchanged.

\subsection{The conditional measure in a clean fluctuation box}
\label{ym4:new-clean-box}

We explain the conditional measure used in the Gaussian comparison.
Put $k_0=K-s$, $x=x_{k_0}$ and $k_1=k_0+n$, where
$n\asymp\log x$. On this segment $x_j\asymp x$. The supports are
separated by $\mathsf r$ spacings of level $k_0$, with
$\mathsf r\asymp(\log x)^{2\mathfrak p_0}$. At step $j$ choose a
read set $S_j$ containing every free bond read by the two insertions,
and a box $I_j$ with $d(S_j,I_j^c)\ge D$. Both insertions lie in the
same box, and
\[
 N_j:=\#I_j\le C(\mathsf rL^{-(j-1-k_0)}+D)^4
                  \le C(\mathsf r+D)^4.
\]
The box is contained in the tame region throughout the segment.
Let $\mathcal N$ be the neighbourhood of the two supports of radius
$C(\mathsf r+\check R_{k_0})a_{K-k_1}$, enlarged to contain the
boxes and all their determining bonds. It is chosen not to wrap.
Cleanliness means that the live region misses $\mathcal N$ at every
level in the finite window
\[
 \mathcal I_{\rm cl}
 =\{t:\max(0,k_0-C\check R_{k_0})\le t\le k_1\}.
\]
A new creation is live at its creation level. A component acted on
after its wait contains a component live at the preceding level.
Thus this condition excludes creations and large-field operations
near the boxes on the segment, including operations hidden by a merge.
The waiting factor is
$\Upsilon=L^{C\check R_{k_0}}=\exp[O((\log x)^r)]$.

Here we use the convention in which healing redraws every lower
small-field profile and compensates its change separately at each
plaquette. Thus the sum of the Wilson weights is still $-\beta$ at
each bare plaquette. A healed component leaves its deterministic
activity $R^{(i)}$; the discrepancy $x_j(\cdot)-x_j$ is supported on
the live region. It consequently vanishes in a clean box. This
convention is needed below: a bound on the average lag of individual
histories would not be a property of a resummed clean class.

The histories of a class are summed before taking a logarithm. The
following comparison is made with the same standing part, stored
variables and exterior fields on both sides. Mark the cubes near the
clean neighbourhood with a factor $e^u$. For a fixed cutoff the
identity of measures for one step is a polynomial identity in $e^u$.
Its coefficient of degree zero deletes precisely the inner histories
forbidden by cleanliness. The pointwise bracket estimate gives
\[
 d\tau_{\mathrm{cl}}=e^{-\Delta_j}d\tau_{\mathrm{res}},\qquad
 0\le\Delta_j\le\epsilon_j,
 \qquad
 \epsilon_j\le C\sum_{t\le j}c_t e^{-(1+c)p_0(g_t)}.
\]
Here the equality is on the common space of reference and stored
variables, before projection. Set $c_t=0$ for
$t\le\min\mathcal I_{\rm cl}$. For
$\min\mathcal I_{\rm cl}<t\le k_1$, $c_t$ counts the cubes of level
$t$ within $(200\check R_t+1)M$ level-$t$ units of $\mathcal N$.
It is bounded by $C(\mathsf r+\check R_{k_0})^4L^{4(k_1-t)}$.
Consequently the accumulated flag error is bounded by
$CL^{4n}\Upsilon(1+n)(\mathsf r+\check R_{k_0})^4e^{-cp_0(g)}$.
No sum over all earlier ultraviolet levels is taken. It is the pointwise bracket
estimate from the large-field construction that justifies this
display. A bound on integrated masses alone would not justify it.
Positivity preserves its two-sided comparison under all subsequent
operations. The conditional expectation at each clean step is
therefore a reference conditional expectation with a normalized
weight in $[e^{-\epsilon_{j-1}},e^{\epsilon_{j-1}}]$.

Their product dominates the reference segment kernel with weight
\[
 P_0=\exp\Big(-\sum_{k_0<j\le k_1}\epsilon_{j-1}\Big).
\]
Splitting off this common subkernel refines the resolved coupling by
one Bernoulli flag. On its reference branch no comparison weight
remains. The other branch has mass at most $n\epsilon_{k_1}$, in
addition to the mass of histories failing cleanliness. Both masses
are included in the unclean error. This construction uses the actual
resolved coupling of the run; an arbitrary disintegration having the
right first marginal would give no such conclusion.

\begin{proof}[The fluctuation law with an insertion]
At one tame step the shift to the critical background, the
linearization of the constraint, pivot elimination and rescaling
give, with a common prefactor $\mathfrak N(W)$,
\[
\begin{aligned}
&\int_{\mathcal M(V)=W}F(V)\chi(V)
 e^{-x_{j-1}[A(U_{j-1}(V))+\mathrm{GF}_{j-1}(V)]}\\
&\hspace{12mm}\times
 e^{\mathbf E_{j-1}(U_{j-1}(V))}\,dV\\
&\qquad=\mathfrak N(W)\int_{\mathcal W_j}
       F(\mathfrak v(y;W))e^{P_j(y;W)}d\gamma_j(y),
\end{aligned}
\]
where
\[
 \mathfrak v(y;W)=
 e^{i[g_{j-1}\mathsf Cy-\widetilde{\mathsf D}
                    (g_{j-1}\mathsf Cy)]}V^{\rm bg}_{j-1}(W),
 \qquad \mathcal W_j=\prod_{b\ \mathrm{free}}\{|y_b|<\mathsf b_j\},
 \qquad \mathsf b_j\asymp p_0(g_{j-1}).
\]
The map $\widetilde{\mathsf D}$ starts at degree two and changes only
pivot variables. Hence the window remains a product on the free
variables. The Gaussian $\gamma_j$ has covariance
$C^{(j-1)}(W)$. Its exponent $P_j$ contains the cubic action and
gauge terms, the Jacobian and Haar density, and the increment of the
old activities. All transformations retain the factor $F$ unchanged.
For gauge-invariant $F$, inserting or removing the gauge fixing also
retains its integral. Dividing by the formula for $F=1$ identifies
the normalized conditional expectation. Applied to a clean class,
this identity uses the common reference branch just constructed.
Operations outside the box and merges replace the exterior datum by
a mixture; their weights remain nonnegative.
\end{proof}

For completeness we give the estimate that removes that exterior.
It concerns the density on the read set, so it also applies to
bounded insertions without estimating their derivatives.

\begin{lemma}\label{ym4:new-window-comparison}
Let $I$ be finite, $\mathcal W=\prod_{b\in I}K_b$, with $K_b$ convex
bounded neighbourhoods of zero, and put
\[
 H_a(y)=\tfrac12\langle y,A_ay\rangle+\langle l_a,y\rangle-P_a(y),
 \qquad \nu_a=Z_a^{-1}\mathbf1_{\mathcal W}e^{-H_a}dy.
\]
For multiplication by $e^{\varphi(b)}$ write $T_\varphi$. Suppose
\[
 \langle v,T_\varphi A_0T_\varphi^{-1}v\rangle\ge\lambda_A|v|^2,
 \qquad
 \sup_{\mathcal W}\|T_\varphi P_0''T_\varphi^{-1}\|
                   \le\varepsilon_0<\lambda_A,
 \qquad\lambda=\lambda_A-\varepsilon_0.
\]
If $|\partial_b(H_1-H_0)|\le\delta_b$ and
$|G(y)-G(y')|\le\ell_G|T_\varphi(y-y')|$, then
\[
 |\nu_1G-\nu_0G|\le\frac{\ell_G}{\lambda}
       \Big(\sum_b e^{2\varphi(b)}\delta_b^2\Big)^{1/2}.
\]
No convexity assumption on $H_1$ is needed. Suppose the same
coercivity bounds hold for every $\varphi$ with
$\mathrm{Lip}\,\varphi\le\mu$. For $S\subset I$, $R=I\setminus S$,
let $m_a$ be the marginal density on $\prod_{s\in S}K_s$. Then
\[
 e^{-\Xi}\le\frac{m_1}{m_0}\le e^\Xi,
 \qquad
 \Xi=\sum_{s\in S}\mathrm{diam}(K_s)
       \left(\delta_s+\frac{L_sD_s}{\lambda}\right),
\]
where
\[
 L_s^2=\sup_{\mathcal W}\sum_{b\in R}
   e^{2\mu d(b,s)}\|(A_0-P_0'')(b,s)\|^2,
 \qquad D_s^2=\sum_{b\in R}e^{-2\mu d(b,s)}\delta_b^2.
\]
\end{lemma}

\begin{proof}
Shrink each $K_b$ slightly, extend $P_a$ smoothly by a retraction
that is the identity on the smaller product, and approximate the
walls by smooth convex one-bond potentials. Their Hessians commute
with $T_\varphi$ and are nonnegative. Choosing the walls sufficiently
steep preserves the displayed coercivity for the extended $H_0$.
The wall terms cancel in $\nabla(H_1-H_0)$. These approximations
have integrable Gaussian majorants, so the walls and then the
shrinking can be removed by dominated convergence.

For the smooth problem couple the Hamiltonian flows of $H_0,H_1$
using the same initial momentum, refreshing that momentum after
each interval of length $t$. Each marginal Gibbs law is preserved.
For $v=T_\varphi(q_1-q_0)$ the difference equation is
\[
 \ddot v=-M(s)v-T_\varphi\nabla(H_1-H_0)(q_1),\qquad\dot v(0)=0,
 \qquad\langle v,M(s)v\rangle\ge\lambda|v|^2.
\]
If $\|M\|\le\Lambda$, Taylor's formula bounds its homogeneous
contraction by $1-\lambda t^2/2+\Lambda^2t^4$. Iterating and then
letting $t\downarrow0$ gives
$\mathbb E|v|\le\nu_1(|T_\varphi\nabla(H_1-H_0)|)/\lambda$.
The Lipschitz bound proves the first assertion. The order of these
limits removes $\Lambda$ before the walls are removed.

For the second assertion fix the variables on $S$. The remaining
integration domain is independent of them because the window is a
product. Differentiate its marginal density:
\[
 \partial_s\log(m_1/m_0)
 =-\nu_1^u(\partial_s(H_1-H_0))
   -[\nu_1^u(\partial_sH_0)-\nu_0^u(\partial_sH_0)].
\]
Use the first assertion on $R$ with
$\varphi(b)=-\mu d(b,s)$ to bound the bracket by $L_sD_s/\lambda$.
Integration along segments bounds the oscillation by $\Xi$.
Both densities integrate to one, which fixes their ratio between
$e^{-\Xi}$ and $e^\Xi$.
\end{proof}

\subsection{The four exterior forces}
\label{ym4:new-exterior-forces}

Freeze an exterior datum $e$ for step $j$, including its stored
fields, and write $O=I_j^c$ for the remaining free bonds. On the
reference branch the conditional density on $I=I_j$ is proportional
to $\mathbf1_{\mathcal W_I}e^{-H_e}$, where
\[
 H_e(y)=\tfrac12\langle y,Q^e_{II}y\rangle
       +\langle Q^e_{IO}y_O(e),y\rangle-P_e(y).
\]
Choose $e_0$ with the same new field near the box, set
$\Gamma=(Q^{e_0})^{-1}$, and let $P_0$ be the exponent at $e_0$,
normalized at zero. The reference law on the box is the full
Gaussian marginal tilted by $e^{P_0}$, not the Gaussian obtained by
simply discarding its exterior. Its precision is
\[
 K=(\Gamma_{II})^{-1}=Q^{e_0}_{II}-Q^{\rm sc},\qquad
 Q^{\rm sc}=Q^{e_0}_{IO}(Q^{e_0}_{OO})^{-1}Q^{e_0}_{OI}.
\]
Consequently its discrepancy from $H_e$ has exactly four forces:
\[
 \nabla(H_e-H_0)=Q^{\rm sc}y+Q^e_{IO}y_O
       +(Q^e_{II}-Q^{e_0}_{II})y-\nabla(P_e-P_0).
\]
Uniform positivity and decay of the actual quadratic forms and
their inverses, the bound $|y_O|\le C\mathsf b_j$, and locality of
the background bound each force at $b$ by
$C\mathsf b_j e^{-c d(b,O)}$. For the last force one uses the
localized form of the old activities and boundary terms. A term
whose domain reaches $O$ has size at least $c\,d(b,O)/M$; half its
decay pays for that distance and the other half sums the domains.
For a domain contained in $I$, only the background changes, and
its difference is exponentially small in the distance to $O$.
The same argument applies to the actual localized quadratic form.
All these statements require the common cores for the entire
interpolation and uniform bounds on its exterior data.

The old-activity subtraction is essential for the small Hessian of
$P_0$. Before summation the second-order part is marginal. After
the coefficient under each tent is subtracted, it has the factor
$L^{-(4+\alpha)(j-1-i)}$ at level $i$, against
$L^{4(j-1-i)}$ sites. The remaining powers of the ratio of radii
grow only polynomially in $j-i$. The sum is therefore uniform.
For $R^{(i)}$ there is no coefficient subtraction; instead use
$\sum_{i<j}x_i^{1-\kappa_0/2}\le Cx^{2-\kappa_0/2}$.
On a box this yields
\[
 \sup|P_0|\le CN_j\left[g\mathsf b_j^3
       +(\log x)^{2(\mathfrak p_0-q)}
       +g^{\kappa_0-O(1)}\right],\qquad
 \|\partial_b\partial_{b'}P_0\|
       \le\varepsilon_P e^{-c d(b,b')},
\]
where $g=x^{-1/2}$ and $\varepsilon_P\to0$. There is no lag term
under the redrawn-profile convention above. The analyticity radii
in $y$ exceed $\mathsf b_j$ by a fixed margin, so Cauchy's estimate
preserves decay in this Hessian bound. Choosing $q$ sufficiently
large compared with $\mathfrak p_0$ makes the exponent and its
Hessian small despite $N_j=O((\log x)^{8\mathfrak p_0})$.

The preceding lemma now gives, on $S_j$, a density ratio in
$[e^{-\Xi_j},e^{\Xi_j}]$, with
\[
 \Xi_j\le CN_j\mathsf b_j^2e^{-cD}.
\]
It holds uniformly for the exterior data, hence also for their
positive mixtures. This proves the needed exterior removal for
bounded insertions. It uses the actual localized force; an internal
Hessian bound alone would not imply it.

\subsection{Completing the Gaussian segment comparison}
\label{ym4:new-gaussian-segment}

The remaining comparison can be made on a generating function.
For each step, first remove the exterior by the density estimate,
then the small tilt, then the Gaussian tails outside the product
window, and finally compare its covariance on the read set with
the flat covariance. Relative covariance comparison costs
$CN_j\varpi_j$, where
$\varpi_j\le C(\mathsf r g p_0+e^{-cD})$; the actual gap and the
analytic comparison of the quadratic form give this estimate.
The total one-step variation is bounded by
\[
 T_j=(e^{\Xi_j}-1)+C\theta_j+Ct_j+CN_j\varpi_j,
 \qquad\theta_j=\sup|P_0|,
\]
with $t_j$ its Gaussian-tail error. The flat covariances are
isotropic in color. Deterministic nesting and localization of the
minimizers, in one common box gauge, give the successive curvature
decomposition, with quadratic remainder
$d_{\rm st}\le C(\mathsf r g p_0)^2$.

We give the source normalization and the carried functions explicitly.
There are two insertions, so put $m=2$, keeping $n=k_1-k_0$ for
the length of the segment. Write $\phi_a$, $a=1,2$, for the
level-$k_0$ profiles $\widetilde f_a$ in (7b)(b), and
$\mathsf O_a$ for the corresponding windowed insertions $X_a$.
Set
\[
 |\lambda_a|\le\Lambda_a
 :=\frac{c x}{m\mathsf r^4p_0^2\|\phi_a\|_\infty},\qquad
 p_0=p_0(g),\qquad g=x^{-1/2}.
\]
The deterministic window bound gives
$\Lambda_a\le c'/(m\sup|\mathsf O_a|)$ after reducing $c$.
Thus the exponential of the two normalized insertions is uniformly
bounded, including on the unclean part of the measure.

For $k_0\le j\le k_1$, let $\mathsf X^{(j)}[A]$ be the
Landau representative of the localized minimizer with datum
$e^{iA}$ in the box around the pair, and put
\[
 \mathsf P_p^{(j)}(A)=\mathsf X^{(j)}[A](\partial p)-\mathbf1.
\]
The read sets contain these boxes and their determining bonds.
The actual linear curvature response $\mathcal H_j$ and the
localized minimizer estimates give
\[
 D\mathsf P_p^{(j)}(0)[A]=i(\mathcal H_j A)(p),\qquad
 |(\mathcal H_jA)(p)|\le C\eta_j^2\|A\|,
\]
\[
 |\mathsf P_p^{(j)}(A)-i(\mathcal H_jA)(p)|
       \le C\eta_j^2\|A\|^2.
\]
On the clean window,
$\mathsf O_a=\tfrac12\sum_p\phi_a(x(p))
|\mathsf P_p^{(k_0)}(A)|^2$.
Here $\eta_j$ is the fine spacing in level-$j$ units, as in the
notation appendix. In particular,
$\eta_{j-1}^2/\eta_{k_0}^2=L^{-2(j-1-k_0)}$.

Let $\gamma^0_j$ be the flat Gaussian law on the read set, with
covariance $C^{(j-1)}(\mathbf1)$, and let $\mathsf C^0$ eliminate
the flat pivot bonds. For independent flat copies $\xi_{j'}$ define
\[
 \mathsf m_{\le j}(p;\xi)
 =i\sum_{k_0<j'\le j}g_{j'-1}
       (\mathcal H_{j'-1}\mathsf C^0\xi_{j'})(p),
\]
\[
 \mathsf Z_j(\mathsf p)=
 \int\prod_{k_0<j'\le j}
       d\gamma^0_{j'}(\xi_{j'}\mid\mathcal W)
 \exp\left\{\frac12\sum_{a=1}^m\lambda_a
     \sum_p\phi_a(x(p))
       |\mathsf m_{\le j}(p;\xi)+\mathsf p(p)|^2\right\}.
\]
If $A^{\rm ax}(V)$ is a box axial potential for $V$, the carried
function is
\[
 Y_j(V)=\int_{\mathrm{SU}(2)}d\mathsf g\,
  \mathsf Z_j\bigl((\mathsf P_p^{(j)}
        (\operatorname{Ad}_{\mathsf g}A^{\rm ax}(V)))_p\bigr).
\]
The flat laws are isotropic in color and $\mathcal H_j$ acts on
each color separately. The constant gauge average therefore removes
the residual box frame. In particular $Y_j$ is a local Borel
gauge-invariant function, so the insertion identities apply to it.
At $j=k_0$ it is the original exponential. Integrating the next
flat copy transforms $\mathsf Z_{j-1}$ into $\mathsf Z_j$.

The polydisc just chosen bounds $\mathsf Z_j$ uniformly and gives
\[
 |\mathsf Z_j(\mathsf p')-\mathsf Z_j(\mathsf p)|
 \le\frac{C}{\eta_{k_0}^2gp_0}
                   \max_p|\mathsf p'(p)-\mathsf p(p)|.
\]
Apply this to the deterministic one-level curvature identity, whose
remainder is at most $\eta_{j-1}^2d_{\rm st}$, and then apply the
read-set variation bound. This yields
\[
 |\mathbb E_jY_{j-1}-Y_j|
 \le C\left[L^{-2(j-1-k_0)}d_{\rm st}/(gp_0)+T_j\right].
\]
The outer conditional expectations are positive and normalized.
Consequently they telescope this estimate by summing its suprema;
no derivative of an accumulated remainder is taken. At the top,
reflection of all the flat copies $\xi\mapsto-\xi$ removes the
linear term in the remaining background. Its contribution is
$O(L^{-4n})$.

Write $t$ for the Gaussian probability that at least one of the
flat copies leaves its window. The union bound gives
$t\le C\sum_jN_j e^{-cp_0^2}$. Removing those windows by
Cauchy--Schwarz costs $Ct^{1/2}$. Thus on the source polydisc the
conditional generating function differs from its flat Gaussian
value by at most
\[
 \delta_{\rm cl}\le C\left[
   \sum_{j=k_0+1}^{k_1}
      \left(L^{-2(j-1-k_0)}d_{\rm st}/(gp_0)+T_j\right)
       +L^{-4n}+t^{1/2}\right].
\]
In each color the resulting Gaussian covariance is
\[
 \mathcal C^x=\sum_{k_0<j\le k_1}x_{j-1}^{-1}\mathcal C_{(j)},
\]
and its generating function is
\[
 \mathsf Z^0(\lambda)=
 \det\left(1-\sum_{a=1}^m\lambda_a\phi_a\mathcal C^x\right)
                              ^{-(N^2-1)/2}.
\]
Its trace-log uses the branch equal to zero at $\lambda=0$.
The diagonal bound on $\mathcal C^x$ and the definition of
$\Lambda_a$ give
\[
 \left\|(\mathcal C^x)^{1/2}
       \left(\sum_a\lambda_a\phi_a\right)
                      (\mathcal C^x)^{1/2}\right\|_1
 \le Cc p_0^{-2}.
\]
Hence $\mathsf Z^0$ has no zeros and
$c\le|\mathsf Z^0|\le C$ uniformly on a fixed smaller polydisc.
It is independent of the clean class and of the top field.

Let $\epsilon_{\rm uc}$ include both the unclean histories and the
nonreference flag. Their masses are $e^{-cp_0}$ times powers of
$L^n$, the waiting factor $\Upsilon=L^{C\check R_{k_0}}$,
and polylogarithms of $x$. The same deterministic insertion
bound controls their generating functions. After averaging over the
classes and the top field, the true generating function $\mathsf Z$
therefore satisfies
\[
 |\mathsf Z/\mathsf Z^0-1|
                    \le C(\delta_{\rm cl}+\epsilon_{\rm uc}).
\]
When this bound is small, the logarithm of this ratio is bounded
by $C(\delta_{\rm cl}+\epsilon_{\rm uc})$. Cauchy's estimate on
the fixed smaller polydisc then gives
\[
 |\partial_{\lambda_1}\partial_{\lambda_2}
       (\log\mathsf Z-\log\mathsf Z^0)(0)|
 \le C(\delta_{\rm cl}+\epsilon_{\rm uc})
                             \Lambda_1^{-1}\Lambda_2^{-1}.
\]
The relative loss in units of
$x^{-2}\|\phi_1\|_\infty\|\phi_2\|_\infty$ is
\[
 \mathsf L_2
 :=x^2\prod_{a=1}^2(\Lambda_a\|\phi_a\|_\infty)^{-1}
 \le(C\mathsf r^4p_0^2)^2.
\]
It remains to pay this loss, rather than merely to show that
$\delta_{\rm cl}$ tends to zero. The only error with exclusively
logarithmic decay is the contribution of the old regular activities.
After multiplication by $\mathsf L_2$ it is at most
\[
 Cn\mathsf r^{12}p_0^4
                 (\log x)^{2(\mathfrak p_0-q)},
\]
which tends to zero when $q$ is chosen sufficiently large in the
stated order of constants. Terms carrying a positive power of $g$
beat every polylogarithmic loss. Choose the box margin $D$ and
$n\asymp\log x$ so that the exponentially small exterior and top
errors also survive $\mathsf L_2$. The window and exceptional masses
are negligible because $\mathfrak p_0>\max\{1,r\}$.
Consequently
\[
 \mathsf L_2(\delta_{\rm cl}+\epsilon_{\rm uc})\longrightarrow0,
\]
which gives the mixed logarithmic derivative with error
$o(x^{-2}\|\phi_1\|_\infty\|\phi_2\|_\infty)$.

The Gaussian covariance in each color is
$\sum_{k_0<j\le k_1}x_{j-1}^{-1}\mathcal C_{(j)}$.
Since $|x_{j-1}-x|\le Cn$, replacing its coefficients by $x^{-1}$
has relative error $O(n/x)$. Wick's formula then gives
$\tfrac12(N^2-1)x^{-2}\mathcal K_{K,s}$ up to the segment's top
tail. Transporting the source to level $k_0$ multiplies this by
$\zeta_{k_0}'{}^2$. Its profile error is
$CM\mathfrak b_{k_0}/(\vartheta\mathsf r)$, and the summable slope
bound makes $\zeta'_{k_0}=1+o(1)$. These are the interacting
comparisons needed before the free kernel can be substituted.

% END MODULE chapters/08_short_distance.tex


% BEGIN MODULE chapters/two_point_extraction_partial_v6.tex
\section{Covariance and transfer estimates for interacting extraction}
\label{ch:two-point-extraction-tests}

The finite insertion, renewed pointing, window formation and linear
bootstrap sections derive the ordinary regular first response under
their explicit unmarked, background, domain and slope inputs.
They retain the potential normalization through local integration
and give the near bound \(C\delta^2s\) on complete nonwrapping fields
of size \(\delta\). The remainder below must also include the
separately bounded marked, remote and boundary contributions.
None of these estimates proves an interacting Gaussian law.

\subsection{Covariance bookkeeping with the retained source normalization}

A source-disc bound alone cannot replace the potential estimate.
Indeed \(H(z,B)=(z/\rho)(B/\alpha)^4\) is bounded by one on its
product unit balls, but its source derivative at \(B=\epsilon\)
is \(\rho^{-1}(\epsilon/\alpha)^4\). Relative to
\(\alpha^2(\epsilon/\alpha)^3\), this is
\(x(\log x)^{p-3q}/\rho\) when
\(\alpha=x^{-1/2}(\log x)^q\) and
\(\epsilon=x^{-1/2}(\log x)^p\); it diverges at fixed \(\rho\).
The finite construction supplies its response by differentiating
the inserted potential, with the stated unmarked hypotheses.

Here is the corresponding covariance bookkeeping. Let $\lambda$ be any
probability measure, put $n_i=\|f_i\|_\infty$, fix $p>0$, and suppose, uniformly in
$x\to\infty$,
\[
 \begin{split}
 \|\mathsf O_i\|_\infty
       &\le C x^{-1}r^4(\log x)^{2p}n_i,\\
 \|\mathsf I_i\|_\infty
       &\le C x^{-1}r^4
                    [(\log x)^{3p-q}+x^{-a}]n_i,\qquad a>0,\\
 r&\le C_r(\log x)^{2p},\qquad q>21p.
 \end{split}
 \tag{ET3}
\]
Then
\[
 \begin{split}
 &|\operatorname{Cov}_\lambda(\mathsf O_1,\mathsf I_2)|
  +|\operatorname{Cov}_\lambda(\mathsf I_1,\mathsf O_2)|
  +|\operatorname{Cov}_\lambda(\mathsf I_1,\mathsf I_2)|\\
 &\hspace{1cm}\le
 C'x^{-2}\big[(\log x)^{21p-q}
                    +(\log x)^{18p}x^{-a}\big]n_1n_2
             =o(x^{-2}n_1n_2).
 \end{split}
 \tag{ET4}
\]
To prove it use $|\operatorname{Cov}(A,B)|\le2\|A\|_\infty\|B\|_\infty$.
The product of the two support counts costs $r^8\le C(\log x)^{16p}$,
and the a priori insertion costs another $(\log x)^{2p}$.
For large $x$ the bracket defining $\mathsf I_i$ is at most one, so its
square is absorbed by itself. This gives (ET4). The numerical inequality
$q>21p$ is a sufficient bound for this particular sup-norm step; it does
not claim to be the optimal requirement for the full extraction.


\subsection{A scalar transfer estimate with a genuinely common normalizer}

Let $\rho,\tau$ be two probability measures with real bounded observables
$\mathcal O(f)$ linear in a real compactly supported Lipschitz function
$f$. They may be measures on different configuration spaces. Write
\[
 D(f,h)=\operatorname{Cov}_\rho(\mathcal O(f),\mathcal O(h))
          -\operatorname{Cov}_\tau(\mathcal O(f),\mathcal O(h)).
\]
Assume that $D$ is invariant under joint translations by $a\mathbb Z^4$.
Set $\|f\|_\sharp=\max(\|f\|_\infty,c_M\operatorname{Lip}f)$ with
$c_M$ fixed.

Suppose there are $r_0>0$ and $Q\ge1$ such that, for every pair
$F_1,F_2$ supported in one common unit ball, there is one holomorphic
function $E(z_1,z_2)$, $E(0)=0$, serving both measures, for which
\[
 \left|e^{-E(z)}
       \omega\!\left(e^{z_1\mathcal O(F_1)+z_2\mathcal O(F_2)}\right)
       \right|\le Q,\qquad \omega\in\{\rho,\tau\},
 \tag{ET5}
\]
on $|z_1|\|F_1\|_\sharp+|z_2|\|F_2\|_\sharp<r_0$.
The constants are uniform in the ball position and in the measures.
The tests in this hypothesis may overlap and have either sign.

Put $c_i=\|F_i\|_\sharp/r_0$ and denote the normalized holomorphic
function in (ET5) by $h_\omega$. It has $h_\omega(0)=1$.
Schwarz's lemma along rays gives
$|h_\omega(z)-1|\le(Q+1)\sum_i c_i|z_i|$.
On the bidisc with radii $[4(Q+1)c_i]^{-1}$ its principal logarithm has
modulus at most $\log2$. Thus Cauchy's formula gives
\[
 |D(F_1,F_2)|
 \le 32\log2\,(Q+1)^2r_0^{-2}
                         \|F_1\|_\sharp\|F_2\|_\sharp.
 \tag{ET6}
\]
The derivatives of $E$ cancel because it is the same function for both
measures. If an entry is zero the conclusion is immediate. No global
zero-free statement for the original generating functions is needed.

Suppose now $a\le d$ is sufficiently small, and $f,h$ have supports
in one common ball of radius $5d$, with
$\operatorname{Lip}f\le\|f\|_\infty/(\vartheta d)$ and the analogous
condition for $h$. Choose a lattice multiple $r'\in[5d,10d)$ and
$N\ge c d^{-4}$ translates of that ball, with centers in $4r'\mathbb Z^4$,
all lying inside a common unit ball. For each choice of signs
$\varepsilon_v=\pm1$ let
$F^\varepsilon=\sum_v\varepsilon_v\tau_v f$ and
$H^\varepsilon=\sum_v\varepsilon_v\tau_v h$.
The translated balls are disjoint. Since the functions vanish on their
ball boundaries, these signed sums do not increase either supremum or
Lipschitz norm. Independent signs with the same sign used in the two
arguments give
\[
 \mathbb E_\varepsilon D(F^\varepsilon,H^\varepsilon)=N D(f,h).
\]
Apply (ET6) and divide by $N$. The result is
\[
 |D(f,h)|\le C(Q,r_0,c_M)\vartheta^{-2}
                         d^2\|f\|_\infty\|h\|_\infty.
 \tag{ET7}
\]
The exponent is $2=4-2$: the number of copies gains four powers and
the two Lipschitz norms lose one power each.

If a reference law has a main term bounded below by
$c_\vartheta x^{-2}\|f\|_\infty\|h\|_\infty$, its relative transfer
error is at most $C_\vartheta d^2x^2$. In the actual short-distance
shape class, each nonnegative entry also has the explicit mass lower
bound $\int f\ge\vartheta d^4\|f\|_\infty$, and the two supports are
at distance between $d$ and $6d$. These requirements give
$\iint f(x)h(y)|x-y|^{-8}\,dx\,dy
\ge10^{-8}\vartheta^2\|f\|_\infty\|h\|_\infty$.
The mass requirement is additional to the Lipschitz bound in (ET7).
At a witness depth with
$x\le\bar x_s=O(s)$ and
$d<L C_W(s)L^{-s}$, where $C_W$ is polylogarithmic in $s$, this tends
to zero uniformly. For a fixed $\vartheta=1/12000$ all constants remain
finite; their dependence on that fixed shape may affect the permitted
coupling and distance thresholds.

The common-normalizer hypothesis (ET5) is not a consequence of having
unrelated bounded normalizers for the two measures. There are two
different ways of supplying it. If the scalar stability construction
already provides one common cell normalizer for every periodic
infinite-volume Gibbs state at the same datum, then (ET5) holds for any
two such states. This includes a reference state obtained by taking
finite-torus limits, provided its Gaussian estimate was uniform along
arbitrarily large tori for each fixed cutoff and pair. The source
$z_1F_1+z_2F_2$ is a single cell source, even when the supports overlap.
This route needs no comparison with a fixed finite torus.

If, instead, the reference Gaussian estimate has only been proved on
one prescribed finite torus, a comparison of that torus with the
infinite-volume states remains necessary. Compatible local energies
and vacuum-density/floor bounds give one scalar chessboard route to
(ET5). Bounds valid on each torus with a distance threshold tending
to zero as its side grows do not by themselves furnish an infinite
reference with the needed fixed-pair estimate. In either route, a typed
Ward insertion estimate is not needed for (ET6)--(ET7).

\subsubsection*{A local companion criterion for vacuum-density compatibility}

For the finite-reference route just described, the following criterion
can remove a separate vacuum-pressure comparison assumption. It is not
needed for comparison of infinite-volume states when their common cell
normalizer is already available. It does not follow from upper stability
and a partition floor alone.
Fix $L>1$ and a cutoff $K$. In the depth convention $a_s=L^{-s}$,
level $j$ has spacing $a_{K-j}=L^{j-K}$, $0\le j\le K$.
Take compatible side lengths $\ell$ for which $\ell/a_{K-j}$ is integral.
Suppose the vacuum constant has the form
\[
 E_{\ell,K}=\sum_{j=0}^K
 \left[(\ell/a_{K-j})^4 c_{j,K}
       +\sum_{z\in(a_{K-j}\mathbb Z/\ell\mathbb Z)^4}
           \sum_{D\in\mathcal D_{\ell,j}(z)}v_{\ell,j}(D,z)\right].
 \tag{ET8}
\]
The numbers $c_{j,K}$, the coupling sequence and every auxiliary choice
are the same for all $\ell$. There is a translation-covariant infinite
companion $v_{\infty,j}$ and an integer length $m(D)\ge0$. For both
families, assume that the number of domains of length $m$ at a root is
at most $C_D e^{\kappa_D m}$ and
$|v_{\ell,j}(D,z)|,|v_{\infty,j}(D,z)|
 \le C_Ee^{-\kappa_E m}$ with $\kappa_E>\kappa_D$.
The constants are uniform in $j,K,\ell$. Assume that rooted domains
with $m<\ell/(C_0a_{K-j})$ correspond bijectively to companion domains,
with equal energy values. This is an explicit locality and nonwrapping
assumption, with $C_0$ fixed. It is stronger than simply asserting
that the two models use the same formula for their energies.

Absolute convergence defines
$\bar v_j=\sum_{D\in\mathcal D_{\infty,j}(0)}v_{\infty,j}(D,0)$ and
$A_K=\sum_{j=0}^K a_{K-j}^{-4}(c_{j,K}+\bar v_j)$.
At each root all terms below the stated threshold cancel in the
difference from the companion. Sum the two tails of the geometric
series with rate $\delta=\kappa_E-\kappa_D>0$. Division by the
physical volume therefore gives
\[
 \left|\ell^{-4}E_{\ell,K}-A_K\right|
       \le C\sum_{j=0}^K a_{K-j}^{-4}e^{-c\ell/a_{K-j}}.
 \tag{ET9}
\]
For every fixed $\ell_0>0$ and $\ell\ge\ell_0$, this is bounded by
$C_{\ell_0}e^{-c\ell/2}$, independently of $K$. Indeed, with $b=K-j$,
\[
 \sum_{b\ge0}L^{4b}e^{-c\ell L^b}
 \le e^{-c\ell/2}
                \sum_{b\ge0}L^{4b}e^{-(c\ell_0/2)L^b}
 \le C_{\ell_0}e^{-c\ell/2}.
 \tag{ET10}
\]
Thus two compatible tori obey the same bound with
$2C_{\ell_0}e^{-c\min(\ell,\ell')/2}$.
The divergent common number $A_K$ need not converge as the cutoff is
removed; it cancels exactly in the comparison.

In a localized source argument, equality of the corresponding
nonwrapping source energies also supplies the same interior-cell
counterterm for all tori. Bounds on the energies crossing the cell
faces and the appropriate partition floor are still needed to obtain
(ET5). This lemma addresses its vacuum-density datum only. If the
coupling sequence or local normalizer numbers differ between tori,
or if nonwrapping compatibility has not been established for the actual
formation rule, (ET9) has not been proved for that construction.

\subsection{A common normalizer for the finite reference torus and all states}
\label{sec:finite-reference-normalizer}

We now obtain (ET5) with the prescribed finite reference torus as one
of the two measures. Fix one first-passage datum and write
$T_*=T_{\ell_*}$ for the smallest admissible torus. Choose a fixed
cell side $l_c$ such that $n_*=\ell_*/l_c$ is even, and place the
support ball a fixed distance from the cell faces. All larger
admissible sides are integer multiples of $\ell_*$. The constants below
may depend on the reference coupling, on $l_c$ and on this fixed
distance. They are independent of the cutoff, the larger side and the
periodic Gibbs state.

The ordinary source construction gives the upper bound and the floor
in (ST5). Write $V_\ell$ for its source-independent vacuum number and
$\mathcal E_*(W)$ for the complete source energy on $T_*$, including
the point terms $U$ and the local domain terms. Thus
\[
 |Z_*(W)|\le
 \exp\{V_*+\operatorname{Re}\mathcal E_*(W)+C_+\ell_*^4\},
 \qquad Z_\ell(0)\ge\exp\{V_\ell-C_-\ell^4\}.
\]
The common companions of the regular and remainder families identify
the vacuum terms on every nonwrapping lifted domain. The explicit
Gaussian constants per site depend only on the common coupling run.
Consequently (ET8)--(ET10), applied to those families, give
\begin{equation}\label{eq:reference-vacuum-densities}
 \left|\ell^{-4}V_\ell-\ell_*^{-4}V_*\right|
 \le \delta_{\rm vac},\qquad \ell\ge\ell_*.
\end{equation}
Only a uniform bound is needed here. The common divergent vacuum
density cancels before the cutoff is removed.

Let $w_z=z_1F_1+z_2F_2$, where the real Lipschitz functions $F_i$
are supported in the chosen ball. Reflect this source into all the
cells, conjugating its coefficients in every odd cell, and denote the
result by $W_z$. The two real profiles of even and odd copies of
$F_i$ have the same supremum and Lipschitz bounds as $F_i$.
They vanish near the faces. Hence
\[
 \|W_z\|_{\flat}\le2\sum_i|z_i|\|F_i\|_\sharp.
\]
We use half the ordinary source radius to pay for this factor two.
No differentiability across a cell face is required; the joined
profiles are globally Lipschitz.

\paragraph*{Moving the repeated source to the reference torus.}
Put $M_\ell(z)=\langle e^{\mathcal O(w_z)}\rangle_\ell$ and
$n_\ell=\ell/l_c$. The finite chessboard inequality gives
\[
 |M_\ell(z)|\le
 Z_\ell(W_z)^{1/n_\ell^4}Z_\ell(0)^{-1/n_\ell^4}.
\]
The numerator is nonnegative. To compare it between volumes, cut in
one coordinate direction. The source is zero near the cutting faces,
so the crossing operator is the positive Wilson operator. The
operators through two reflected adjacent slabs are adjoints, including
the conjugation of the source. Their product is a positive trace-class
operator $Y$. Repeating the pair $q$ times replaces $\operatorname{Tr}Y$
by $\operatorname{Tr}Y^q$. Its eigenvalues are nonnegative, and
$\operatorname{Tr}Y^q\le(\operatorname{Tr}Y)^q$ for integer $q\ge1$.
Apply this successively in the four directions, allowing rectangular
tori at the intermediate stages. It follows that
\begin{equation}\label{eq:reference-pressure-comparison}
 0\le Z_\ell(W_z)\le Z_*(W_z)^{(\ell/\ell_*)^4}.
\end{equation}
This argument uses the exact finite lattice measure; the local
source-energy representation is required only on $T_*$.

Combine this inequality with stability on $T_*$, the floor on
$T_\ell$, and \eqref{eq:reference-vacuum-densities}. Since
$\ell^4/n_\ell^4=l_c^4$, both $\ell$ and $\ell_*$ satisfy
\begin{equation}\label{eq:reference-cell-upper-bound}
 |M_\ell(z)|\le
 \exp\{n_*^{-4}\operatorname{Re}\mathcal E_*(W_z)
              +(C_++C_-+\delta_{\rm vac})l_c^4\}.
\end{equation}
In particular the large-volume factor has become a fixed cell factor.

\paragraph*{The holomorphic part common to the two volumes.}
For each cell $rC$, retain the source-energy terms on $T_*$ whose
complete dependence region lies in its interior. Evaluate a retained
term at the single reflected copy $w_z^r$ in that cell. If the cell is
odd, conjugate this value once more. Define $E^\circ(z)$ to be the sum
of these values, divided by $n_*^4$. Include the point terms in the
same rule. In an even cell a summand is holomorphic in $z$; in an odd
cell it has the form $\overline{\phi(\bar z)}$, with $\phi$
holomorphic, and is holomorphic as well. We have $E^\circ(0)=0$.
Locality gives
\[
 \operatorname{Re}E^\circ(z)=n_*^{-4}
 \sum_{\text{interior terms}}
             \operatorname{Re}e_*(W_z).
\]
This definition averages over all cells. It requires no covariance
of the individual localized terms under the cell group.

The omitted domain terms meet a cell face and also reach the support
of a source copy. At depth $i=K-j$, such a domain has length at least
$[cL^i/M-C]_+$. It contains one of at most $(CL^i/M+C)^4$ anchor
cubes near the support ball. If $\kappa_e>\kappa_{\mathcal D}$ are
the energy and entropy rates, the sum per cell is bounded by
\begin{equation}\label{eq:reference-face-series}
 C_{\rm face}=C\sum_{i\ge0}(CL^i/M+C)^4
 e^{-(\kappa_e-\kappa_{\mathcal D})[cL^i/M-C]_+}<\infty.
\end{equation}
There are only finitely many coarsest levels without exponential
suppression. At finer levels the exponential in $L^i$ pays for the
anchor count. Thus this bound is uniform in the cutoff. Point terms
on the faces vanish, since the source is zero there. We obtain
\[
 \left|n_*^{-4}\operatorname{Re}\mathcal E_*(W_z)
                      -\operatorname{Re}E^\circ(z)\right|
 \le C_{\rm face}.
\]
Together with \eqref{eq:reference-cell-upper-bound}, this proves
\begin{equation}\label{eq:actual-common-reference-normalizer}
 |e^{-E^\circ(z)}M_\ell(z)|\le Q,
 \qquad
 Q=\exp\{C_{\rm face}+(C_++C_-+\delta_{\rm vac})l_c^4\},
 \quad \ell\ge\ell_*.
\end{equation}
The very same $E^\circ$ serves the finite reference torus and all
larger tori. Interior source energies on the larger tori have not
been compared term by term with those on $T_*$.

\paragraph*{Periodic states and the short-distance gain.}
For an arbitrary periodic DLR state, apply the exact domination lemma
of Section~\ref{ch:all-state-source-transfer} to the cell observable
$e^{\mathcal O(w_z)}$. That lemma bounds its expectation by the
limsup of the finite chessboard seminorms. The preceding calculation
bounds each such seminorm by $|e^{E^\circ(z)}|Q$. Hence
\eqref{eq:actual-common-reference-normalizer} also holds for every
periodic state, with the same normalizer and the same $Q$. At fixed
cutoff the observable is bounded and its generating function is
entire, so there is no additional interchange of an unbounded source
with an infinite-volume limit.

This supplies (ET5) for the finite reference and every periodic state.
The elementary logarithm estimate (ET6) then cancels $E^\circ$, and
the average over signed translates in (ET7) gives
\[
 |S_2^T[\nu](f,h)-S_2^T[T_*](f,h)|
       \le C_\vartheta d^2\|f\|_\infty\|h\|_\infty.
\]
The covariance difference is translation invariant under
$a\mathbb Z^4$, as required for that averaging. A reference main term
of size $x^{-2}\|f\|_\infty\|h\|_\infty$ therefore has relative
transfer error at most $C_\vartheta d^2x^2$, which tends to zero
at the witness depths in (ET7). The order is: finite reference
extraction, this uniform transfer at each cutoff, and then the
continuum and component limits. No assertion that every pure Gibbs
state is a limit of unperturbed torus measures is used.

% END MODULE chapters/two_point_extraction_partial_v6.tex


\section{Sharp cutoff boundaries in the classical Gaussian law}
\label{ch:gaussian-cutoff-clean-law}

The boundary term in the conditional integration-by-parts identity need
not be small merely because a Gaussian precision has a gap.  The event
that every bond is small can have a probability tending to zero with
the volume.  Dividing a whole-space tail estimate by that probability
would therefore give no uniform conditional estimate.  In the centered
Gaussian reference law the symmetry and convexity of all the other
cutoffs give a different argument: their conditional probability
decreases on every radial ray.  A radial boundary density can then be
compared with an inner radial integral in the same direction.  The
normalization cancels before directions or bonds are summed.

This section proves that comparison for possibly singular physical
bond covariances.  It derives the resulting second through eighth
moment errors, including the covariance and third and fourth connected
errors of localized quadratic curvature observables.  The actual flat
pivot chart and contour gauge supply a concrete fixed-step application.
The measure here is classical Gaussian.  Haar and coarea factors, old
activities, nonlinear action weights and conditional exterior forces
are additional terms in the complete conditional law of
Section~\ref{ch:two-point-extraction-tests}.

\subsection*{Physical blocks, actual covariance entries and the cutoff law}

Let \(Y\) be a centered Gaussian vector in a finite-dimensional real
space, with covariance \(G\ge0\), and let \(Z=CY\).  The blocks
\(Z_b\) belong to \(\mathbb R^3\); repeated or linearly dependent
blocks are allowed.  All covariances in this section are the actual
matrix entries
\[
 \Gamma_{bc}=\mathbb E[Z_bZ_c^{\mathsf T}]=C_bGC_c^{\mathsf T},
 \qquad \Gamma=CGC^{\mathsf T}.
 \tag{GCB1}\label{gcb:physical-covariance}
\]
For every tested block choose \(b_b>0\), put
\[
 \mathcal A=\bigcap_b\{|Z_b|<b_b\},\qquad
 \nu(P)={\mathbb E[P(Z)1_{\mathcal A}]\over
                       \mathbb P(\mathcal A)},
 \qquad d_b=\operatorname{rank}\Gamma_{bb}\le3,
 \tag{GCB2}\label{gcb:sharp-law}
\]
and choose an upper marginal scale \(\sigma_b>0\) with
\(\Gamma_{bb}\le\sigma_b^2 I\).  There is no lower eigenvalue
assumption.  A sufficiently small neighborhood of zero in the latent
Gaussian support is contained in \(\mathcal A\), so its probability
is positive.  If \(d_b=0\), then \(Z_b=0\) almost surely and that
constraint has no boundary contribution.  All blocks occurring in a
polynomial to be bounded below are included among the tested blocks.

For a finite volume no lower bound on \(\mathbb P(\mathcal A)\)
uniform in the volume is assumed or needed.  The estimates concern
these normalized finite-volume laws.  They do not define conditioning
on a simultaneous infinite collection of cutoff events, whose
intersection could have probability zero.

\subsection*{Marginal log-concavity from one-dimensional transport}

We first prove the integration fact needed for the conditional cutoff
probability.  Suppose nonnegative integrable functions \(f,g,h\)
on \(\mathbb R\), positive in the interiors of their interval
supports, satisfy, for \(0<t<1\),
\[
 h((1-t)x+ty)\ge f(x)^{1-t}g(y)^t.
 \tag{GCB3}\label{gcb:transport-hypothesis}
\]
It is enough initially that the functions be continuous and the
positive densities be smooth on their supports, as will hold in the
application.  Write \(F=\int f\), \(G_1=\int g\), and assume both
are positive.  The increasing quantile map \(T\) from density
\(f/F\) to \(g/G_1\) is continuously differentiable in the interior
and satisfies
\(T'(x)=f(x)G_1/[F g(T(x))]\).  The increasing map
\(T_t(x)=(1-t)x+tT(x)\) has derivative at least \(T'(x)^t\), by
the weighted arithmetic-geometric mean inequality.  Changing variables
on its image and using \eqref{gcb:transport-hypothesis} gives
\[
 \begin{aligned}
 \int h&\ge\int h(T_t(x))T_t'(x)\,dx\\
 &\ge\int f(x)^{1-t}g(T(x))^t T'(x)^t\,dx
       =F^{1-t}G_1^t.
 \end{aligned}
 \tag{GCB4}\label{gcb:transport-integral}
\]
One first integrates over compact subintervals of the positive support
and then lets them exhaust it.  All terms are nonnegative.  This also
handles unbounded interval endpoints and densities vanishing at the
boundary.  If either integral is zero the conclusion is immediate.

The same assertion holds in \(\mathbb R^q\).  Here is the induction
needed to avoid importing a stronger integration theorem.  Split
\(x=(x',u)\), \(y=(y',v)\).  For fixed \(x',y'\), apply the
one-dimensional assertion in \(u,v\), with the slice of \(h\)
at \((1-t)x'+ty'\).  The resulting marginals obey the same
inequality in \(q-1\) dimensions.  Apply the induction hypothesis
to those marginals and use Tonelli.  Slices of zero integral cause no
problem.  In the application all positive supports are open convex
sets and all slices have these interval supports.  Products of smooth
compact cutoffs with a Gaussian have integrable dominating functions
for their derivatives on every compact set of external parameters.
Their marginals are therefore continuous, smooth where required,
and positive on the interiors of their convex supports.  These facts
justify each induction step directly in this class of functions.

In particular, if \(k(z,w)\) is log-concave jointly in \(z,w\),
and its \(w\) slices are of the preceding integrable class, then
\[
 w_0(z)=\int k(z,w)\,dw
 \quad\hbox{is log-concave in }z.
 \tag{GCB5}\label{gcb:marginal-log-concavity}
\]
Indeed use the joint log-concavity inequality for the slices at
\(z,z'\) and \((1-t)z+tz'\), then integrate by
\eqref{gcb:transport-integral} and its induction.  The function
\(k\) need not be integrable over the external variable \(z\).
If \(k(-z,-w)=k(z,w)\), its marginal is even.  An even log-concave
function satisfies \(w_0(rn)\le w_0(sn)\) whenever \(0\le s\le r\):
write \(sn\) as a convex combination of \(rn\) and \(-rn\).
This ray monotonicity, rather than a volume estimate for the complete
cutoff probability, is the property used below.

\subsection*{Smooth compact approximants with an inner plateau}

For \(0<\delta\le1/2\), construct a radial cutoff as follows.  Put
\(\phi(t)=0\) for \(t\le0\), and on \(0<t<1\) set
\[
 \begin{aligned}
 \phi''(t)&={e^{-1/t^2}\over(1-t)^4},\qquad
 \phi(t)=\int_0^t(t-s){e^{-1/s^2}\over(1-s)^4}\,ds,
 \\
 \psi_{b,\delta}(z)&=
 \begin{cases}
 e^{-\phi((|z|/b-1+\delta)/\delta)},&|z|<b,\\
 0,&|z|\ge b.
 \end{cases}
 \end{aligned}
 \tag{GCB6}\label{gcb:smooth-cutoff}
\]
The function \(\phi\) is convex, nondecreasing, flat at zero and
diverges at least as a positive multiple of \((1-t)^{-2}\) near
one.  Consequently \(\psi_{b,\delta}\) is smooth on all of
\(\mathbb R^3\), including the sphere \(|z|=b\), is even and
log-concave, and equals one on \(|z|\le(1-\delta)b\).
The smoothness at zero follows because it is constant near zero;
the smoothness at the outer sphere follows from the exponential
decay of \(e^{-\phi}\) faster than every inverse power of
\(1-t\).  Convexity of \(z\mapsto\phi((|z|/b-1+\delta)/\delta)\)
follows from convexity of the norm and monotonic convexity of \(\phi\),
extended by \(+\infty\) at and beyond one.

As \(\delta\downarrow0\), these cutoffs increase pointwise to
\(1_{\{|z|<b\}}\).  To check monotonicity inside the ball, the
argument of \(\phi\) is \(1-(1-|z|/b)/\delta\), which decreases
as \(\delta\) decreases.  Define
\[
 \chi_\delta(Z)=\prod_b\psi_{b_b,\delta}(Z_b),\qquad
 \nu_\delta(P)={\mathbb E[P(Z)\chi_\delta(Z)]\over
                          \mathbb E\chi_\delta(Z)}.
 \tag{GCB7}\label{gcb:smooth-law}
\]
Every denominator is positive, and Gaussian dominated convergence
gives \(\nu_\delta(P)\to\nu(P)\) for every polynomial \(P\).
No independence of the tested coordinates enters this statement.

The plateau in \eqref{gcb:smooth-cutoff} is essential to the estimate:
the derivative is supported in the outer annulus
\([(1-\delta)b,b]\).  A smooth cutoff whose derivative remained
nonzero close to zero would not by itself yield the exponential
boundary tail established next.

\subsection*{Conditional Gaussian support and monotonicity on rays}

Fix a block \(b\) of positive rank and let
\(V_b=\operatorname{ran}\Gamma_{bb}\), of dimension \(d=d_b\).
Orthogonal decomposition of a whitened latent Gaussian gives, when
conditioning on \(Z_b=z\in V_b\),
\[
 Z_a=\Gamma_{ab}\Gamma_{bb}^{+}z+D_aW\quad(a\ne b),
 \qquad
 w_{b,\delta}(z)=
 \mathbb E_W\prod_{a\ne b}
 \psi_{b_a,\delta}(\Gamma_{ab}\Gamma_{bb}^{+}z+D_aW).
 \tag{GCB8}\label{gcb:conditional-weight}
\]
Here \(W\) is a standard Gaussian independent of \(Z_b\), possibly
in dimension zero, and \(\Gamma_{bb}^{+}\) is the inverse on
\(V_b\), extended by zero.  For completeness, write the latent
standard Gaussian as its orthogonal projection onto the row space
of the map defining \(Z_b\) and the perpendicular projection.
The first projection is determined by \(z\), and the second is
an independent standard Gaussian on its subspace.  This proves
\eqref{gcb:conditional-weight}, including singular full covariances.
There is no bound asserted on \(\Gamma_{ab}\Gamma_{bb}^{+}\).

The integrand defining \(w_{b,\delta}\), multiplied by the Gaussian
density of \(W\), is jointly log-concave in \((z,W)\) and unchanged
under \((z,W)\mapsto(-z,-W)\).  Formula
\eqref{gcb:marginal-log-concavity} therefore proves that
\(w_{b,\delta}\) is even and log-concave on \(V_b\).  Its value
decreases along every ray.  It is positive near zero, although it
may vanish elsewhere because of exact dependencies between blocks.
Such zeros strengthen the inequalities below and require no separate
boundary trace convention.

\subsection*{A normalized boundary flux with no lower marginal variance}

Define the normalized smooth flux
\[
 \beta_{b,\delta}=
 {\mathbb E[|\nabla\psi_{b_b,\delta}(Z_b)|
                 \prod_{a\ne b}\psi_{b_a,\delta}(Z_a)]
       \over\mathbb E\chi_\delta(Z)}.
 \tag{GCB9}\label{gcb:smooth-flux}
\]
For \(d_b=0\), set \(\beta_b^*=0\); the flux is zero because
the cutoff is constant near zero.  In positive rank the radial gradient
lies in \(V_b\), so the
ambient Euclidean norm in \eqref{gcb:smooth-flux} is also its intrinsic
gradient norm.  Set \(B=b_b\), \(a=(1-\delta)B\), and use polar
coordinates \(z=rn\) in \(V_b\).  Up to a constant independent
of \(r,n\), the intrinsic Gaussian density is
\(\exp[-q(n)r^2/2]\), with
\(q(n)=n^{\mathsf T}(\Gamma_{bb}|_{V_b})^{-1}n
\ge\sigma_b^{-2}\).

For a fixed direction let
\[
 \begin{aligned}
 D_n&=\int_0^B\psi_{B,\delta}(rn)w_{b,\delta}(rn)
                    r^{d-1}e^{-q(n)r^2/2}\,dr,\\
 I(q,a)&=\int_0^a s^{d-1}e^{-qs^2/2}\,ds.
 \end{aligned}
 \tag{GCB10}\label{gcb:radial-denominator}
\]
For every \(r\in[a,B]\), ray monotonicity and the inner plateau
give \(D_n\ge w_{b,\delta}(rn)I(q(n),a)\).
The integral of the nonnegative radial derivative
\(-\partial_r\psi_{B,\delta}\) on \([a,B]\) is one.  Thus the
corresponding numerator in direction \(n\) is at most
\[
 D_n\,{B^{d-1}e^{-q(n)a^2/2}\over I(q(n),a)}.
 \tag{GCB11}\label{gcb:directional-flux}
\]
This inequality remains true if \(D_n=0\), because the preceding
pointwise inequality then forces the numerator to vanish.  In
particular it does not divide by a possibly zero directional weight.

There are constants depending only on \(d\le3\) such that
\(I(q,a)\ge c_d\min(a^d,q^{-d/2})\).  If
\(a\le q^{-1/2}\), integrate over \([0,a]\), where the exponential
is at least \(e^{-1/2}\).  Otherwise integrate over
\([0,q^{-1/2}]\).  Since \(B/2\le a\le B\), the ratio in
\eqref{gcb:directional-flux} is bounded by
\(C_d B^{-1}[1+(qB^2)^{d/2}]e^{-qB^2/8}\).
Taking the supremum over all \(q\ge\sigma_b^{-2}\) gives
\[
 \boxed{\quad
 \beta_{b,\delta}\le\beta_b^*
 := {C_d\over b_b}
       \left[1+\left({b_b\over\sigma_b}\right)^d\right]
           e^{-c_d b_b^2/\sigma_b^2}\quad}
 \qquad(1\le d=d_b\le3).
 \tag{GCB12}\label{gcb:flux-bound}
\]
To justify the last maximization explicitly, put
\(x=qB^2\), \(x_0=B^2/\sigma_b^2\).  The product
\((1+x^{d/2})e^{-x/16}\) is bounded on \([0,\infty)\).
For \(x\ge x_0\), the remaining factor is at most
\(e^{-x_0/16}\); increasing the constant gives the displayed
bound, with for example a smaller fixed \(c_d\).  Now integrate
\eqref{gcb:directional-flux} over directions.  The determinant
normalization of the intrinsic Gaussian and the total integral of
\(D_n\) cancel.  This proves \eqref{gcb:flux-bound} uniformly in
\(\delta\), the volume, the number of other constraints, their
dependencies and every lower marginal eigenvalue.

For \(b_b\ge\sigma_b\) the useful simpler consequence is
\[
 \beta_b^*\le C_d b_b^{d-1}\sigma_b^{-d}
                  e^{-c_d b_b^2/\sigma_b^2}.
 \tag{GCB13}\label{gcb:large-radius-flux}
\]
The qualification is necessary.  At small radius the normalized
radial flux can be of order \(d/b_b\); an estimate proportional to
\(b_b^{d-1}\sigma_b^{-d}\) alone would be false there.  The proof
has retained that regime in \eqref{gcb:flux-bound}.

\subsection*{Gaussian integration by parts before the sharp limit}

For a polynomial \(P\) in the ambient physical blocks, ordinary
Gaussian integration by parts, applied to a whitened latent vector,
gives
\[
 \begin{aligned}
 \nu_\delta[Z_iP]&=
   \sum_j\Gamma_{ij}\nu_\delta[\nabla_jP]
       +\sum_b\Gamma_{ib}\mathcal L_{b,\delta}(P),\\
 \mathcal L_{b,\delta}(P)&=
 {\mathbb E[P(Z)\nabla\psi_{b_b,\delta}(Z_b)
                         \prod_{a\ne b}\psi_{b_a,\delta}(Z_a)]
        \over\mathbb E\chi_\delta(Z)}.
 \end{aligned}
 \tag{GCB14}\label{gcb:smooth-ibp}
\]
The gradients are column vectors and \(\Gamma_{ib}\) are the
corresponding \(3\) by \(3\) blocks.  Degeneracy of \(\Gamma\)
does not alter the formula: in the latent variables its two
chain-rule factors are exactly the factors producing
\(CGC^{\mathsf T}\).  Gaussian domination justifies integration
by parts even when the cutoffs have no compact support in null
directions of \(C\).

Let \(\|P\|_{\mathcal B}\) be the supremum on the product of
closed balls \(\mathcal B=\prod_b\{|z_b|\le b_b\}\).  Then
\[
 |\mathcal L_{b,\delta}(P)|
       \le\beta_b^*\|P\|_{\mathcal B}.
 \tag{GCB15}\label{gcb:individual-defect-bound}
\]
For \(L=\sum_i\ell_i^{\mathsf T}Z_i\), define its canonical
sharp defect by
\[
 \mathcal D_L(P)=\nu[LP]
       -\sum_{i,j}\ell_i^{\mathsf T}\Gamma_{ij}\nu[\nabla_jP].
 \quad
 |\mathcal D_L(P)|\le
 \left(\sum_{i,b}|\ell_i|\,\|\Gamma_{ib}\|\beta_b^*\right)
                   \|P\|_{\mathcal B}.
 \tag{GCB16}\label{gcb:sharp-defect}
\]
Indeed the left side is the limit of the left side of
\eqref{gcb:smooth-ibp} contracted with \(\ell\), while
\eqref{gcb:individual-defect-bound} bounds the right side uniformly.
This proves the sharp identity and its estimate without assigning
a meaning to a product of a boundary delta function with another
sharp indicator on the same surface.

For repeated or dependent constraints the separate sharp limits of
\(\mathcal L_{b,\delta}\) need not be intrinsic: the allocation
between coincident constraint surfaces can depend on the smoothing
schedule.  The total defect in \eqref{gcb:sharp-defect} is intrinsic,
because its definition uses only sharp conditional polynomial
expectations.  This is the only boundary object needed for the moment
estimates.  When a classical smooth-boundary free-coordinate trace
exists, \(\nabla1_D=-n\,dS\), so this defect is the negative of
the covariance contracted with the outward boundary functional in
the conditional-law identity.  This also fixes the sign convention.

\subsection*{The actual weighted Schur accounting}

Choose positive amplitude scales \(a_b\) and positive cell weights
\(\mu_b\).  Assume the following explicit two-sided bound, using
the covariance entries of \eqref{gcb:physical-covariance}:
\[
 S=\max\left\{
  \sup_i\sum_j\mu_j{\|\Gamma_{ij}\|\over a_i a_j},
  \sup_j\sum_i\mu_i{\|\Gamma_{ij}\|\over a_i a_j}
                         \right\}<\infty,
 \qquad
 \varepsilon_{\partial}
       =\sup_b{a_b\beta_b^*\over\mu_b}.
 \tag{GCB17}\label{gcb:weighted-schur}
\]
The two sums give the usual bounded operator on weighted counting
\(\ell^2\) after normalizing by \(a_b\): weighted Cauchy in each
row and then the column bound give norm at most \(S\).
For the boundary estimate the more specific computation is
\[
 {1\over a_i}\sum_b\|\Gamma_{ib}\|\beta_b^*
 =\sum_b\mu_b{\|\Gamma_{ib}\|\over a_i a_b}
                 {a_b\beta_b^*\over\mu_b}
 \le S\varepsilon_{\partial}.
 \tag{GCB18}\label{gcb:weighted-boundary-row}
\]
Thus, with \(U(L)=\sum_i a_i|\ell_i|\),
\[
 |\mathcal D_L(P)|\le
       \delta(L)\|P\|_{\mathcal B},\qquad
 \delta(L)=S\varepsilon_{\partial}U(L).
 \tag{GCB19}\label{gcb:linear-defect-norm}
\]
There is no number of ambient coordinates in this estimate.
The factor \(\mu_b^{-1}\) in
\(\varepsilon_{\partial}\) cannot be dropped: coordinate
Gaussian integration by parts uses the ordinary partial derivative
with respect to that coordinate.  A weighted kernel convention
does not change that derivative.

For example, if a precision is written as a self-adjoint operator
\(K\) in pairing \(\sum_b\mu_b y_b\cdot z_b\), its coordinate
density has Euclidean precision \(M K\), where
\(M=\operatorname{diag}(\mu_b I_3)\).  Its covariance matrix is
\((MK)^{-1}\).  If
\((K^{-1}f)_i=\sum_j\mu_j\mathsf G(i,j)f_j\), then
\(K^{-1}=\mathsf G M\), and \((MK)^{-1}=\mathsf G\).
Accordingly \(\mathsf G(i,j)\), rather than
\(\mu_j\mathsf G(i,j)\), is the covariance entry used here.
This calculation prevents either suppressing or inserting a second
cell-volume factor during the physical coordinate conversion.

\subsection*{All linear marginals remain Gaussian dominated}

There is a useful moment bound independent of the cutoff radii.
For a real linear statistic \(L\), put
\(v(L)^2=\mathbb E L^2=\ell^{\mathsf T}\Gamma\ell\).
If this is zero then \(L=0\) almost surely.  Otherwise condition
the latent Gaussian on \(L=x\).  The argument of
\eqref{gcb:conditional-weight}, now including every cutoff, gives
an even log-concave weight \(w(x)\) for its one-dimensional
Gaussian density, first for \(\nu_\delta\) and then by the
sharp limit.  Therefore \(w\) decreases as \(|x|\) increases.
For any nonnegative increasing function \(f\) of \(|x|\),
\[
 \nu[f(|L|)]\le
       \mathbb E f(|v(L)N|),\qquad N\sim N(0,1).
 \tag{GCB20}\label{gcb:linear-moment-domination}
\]
To see the normalization explicitly, let \(R,R'\) be independent
absolute Gaussian variables of this variance.  The product
\((f(R)-f(R'))(w(R)-w(R'))\) is nonpositive.  Its expectation
is twice \(\mathbb E[fw]-\mathbb E[f]\mathbb E[w]\).
Divide by \(\mathbb E[w]>0\), and truncate \(f\) followed by
monotone convergence if necessary.  This proves the assertion
without a lower bound on the cutoff probability.

In particular the law is centered and, for \(m\ge1\),
\[
 \nu[L^{2m}]\le(2m-1)!!\,v(L)^{2m},\qquad
 \nu[e^{tL}]\le e^{t^2v(L)^2/2},\qquad
 \operatorname{Cov}_\nu(Z)\le\Gamma.
 \tag{GCB21}\label{gcb:linear-subgaussian}
\]
The exponential assertion follows by symmetry and the increasing
function \(\cosh(t|L|)\).  The covariance ordering follows by
testing every linear functional.  Every joint moment of an odd
number of linear statistics is zero, because the entire law is
invariant under \(Z\mapsto-Z\).

\subsection*{Covariance and fourth-order defects of linear statistics}

For linear statistics \(L_r\), write
\[
 \begin{aligned}
 &\Gamma_{rs}^{L}=\mathbb E[L_rL_s],\quad
 \Sigma_{rs}=\nu[L_rL_s],\quad
 E_{rs}=\Sigma_{rs}-\Gamma_{rs}^{L},\\
 &B_r=\sum_i b_i|\ell_{r,i}|,\quad
 \delta_r=S\varepsilon_{\partial}U(L_r),\quad v_r=v(L_r).
 \end{aligned}
 \tag{GCB22}\label{gcb:word-data}
\]
Here \(B_r\) bounds \(|L_r|\) on the entire cutoff box, whereas
\(v_r\) is its untruncated Gaussian standard deviation.  These
two quantities have different uses.  Taking \(P=L_s\) in
\eqref{gcb:sharp-defect} gives
\[
 |E_{rs}|\le\min(\delta_r B_s,\delta_s B_r),\qquad
 |\Sigma_{rs}|+|\Gamma_{rs}^{L}|\le2v_rv_s.
 \tag{GCB23}\label{gcb:covariance-error}
\]
The second estimate follows separately for its two terms from
\eqref{gcb:linear-subgaussian} and Gaussian Cauchy--Schwarz.
Every third joint moment of linear statistics is exactly zero.

Apply the same identity to \(P=L_2L_3L_4\).  The derivative
of this product produces exactly three pairings, and hence
\[
 \nu[L_1L_2L_3L_4]
    =\sum_{j=2}^4\Gamma_{1j}^{L}\Sigma_{\{2,3,4\}\setminus\{j\}}
                    +\mathcal D_{L_1}(L_2L_3L_4).
 \tag{GCB24}\label{gcb:fourth-recursion}
\]
In this formula \(\Sigma_{\{r,s\}}=\Sigma_{rs}\).  Subtracting
the three Gaussian pairings gives an error at most
\(\delta_1B_2B_3B_4+
\sum_{j=2}^4|\Gamma_{1j}^{L}||E_{\{2,3,4\}\setminus\{j\}}|\).
Subtracting instead the three pairings with the actual covariance
\(\Sigma\) gives the fourth connected moment and the explicit bound
\[
 \begin{aligned}
 |\kappa_\nu(L_1,L_2,L_3,L_4)|
 \le\delta_1\big(&B_2B_3B_4+B_2v_3v_4\\
                         &+B_3v_2v_4+B_4v_2v_3\big).
 \end{aligned}
 \tag{GCB25}\label{gcb:fourth-connected-error}
\]
Indeed replace \(\Gamma_{1j}^{L}-\Sigma_{1j}\) by
\(-E_{1j}\), use \(|E_{1j}|\le\delta_1 B_j\), and bound each
remaining \(\Sigma\) by its two Gaussian standard deviations.
Any of the four indices can be chosen as the first one.  The
identity, rather than an assumed small fourth cumulant, proves
the connected estimate.

For example,
\(\operatorname{Cov}_\nu(L_1L_2,L_3L_4)
=\kappa_\nu(L_1,L_2,L_3,L_4)+\Sigma_{13}\Sigma_{24}
+\Sigma_{14}\Sigma_{23}\).
Together with \eqref{gcb:covariance-error} this gives the covariance
error of two factored quadratic observables relative to their
untruncated Gaussian Wick contractions.  The higher degree estimates
below also control three and four such observables.

\subsection*{The moment recursion through degree eight}

Let a collection of localized linear statistics obey
\[
 B_r\le R\,U_r,\qquad v_r\le v\,U_r,
 \qquad U_r=U(L_r),\qquad D=S\varepsilon_{\partial}.
 \tag{GCB26}\label{gcb:local-radius-data}
\]
One may take \(R=\max_{i\in\mathcal S}b_i/a_i\) and
\(v=\max_{i\in\mathcal S}\sigma_i/a_i\), where \(\mathcal S\)
is the union of the supports of these statistics.  The variance
bound follows from
\(\|\Gamma_{ij}\|\le\sigma_i\sigma_j\), itself a consequence of
Gaussian Cauchy--Schwarz.  Thus neither \(R\) nor \(v\) requires
a maximum over coordinates outside their supports.

Let \(M_\nu(1,\ldots,n)=\nu[\prod_{r=1}^nL_r]\) and let
\(M_0\) denote the corresponding untruncated Gaussian moment.
For even \(n\), integration by parts gives the exact recurrence
\[
 M_\nu(1,\ldots,n)-M_0(1,\ldots,n)
 =\sum_{j=2}^n\Gamma_{1j}^{L}
                  (M_\nu-M_0)(\{2,\ldots,n\}\setminus\{j\})
       +\mathcal D_{L_1}\Big(\prod_{r=2}^n L_r\Big).
 \tag{GCB27}\label{gcb:general-moment-recursion}
\]
For \(n=2\) the empty moment difference is zero.  Odd moments
vanish for both laws.  Consequently
\[
 \begin{aligned}
 |M_\nu(1,\ldots,2m)-M_0(1,\ldots,2m)|
             &\le d_{2m}\prod_{r=1}^{2m}U_r,\\
 d_2&=DR,\qquad d_n=DR^{n-1}+(n-1)v^2d_{n-2}.
 \end{aligned}
 \tag{GCB28}\label{gcb:moment-error-recurrence}
\]
This is a direct induction: each derivative pairing contributes
at most \(v^2U_1U_j\), and the boundary defect contributes at
most \(D U_1\prod_{r=2}^n B_r\).  In particular
\[
 \begin{aligned}
 d_2&=DR,\\
 d_4&=D(R^3+3v^2R),\\
 d_6&=D(R^5+5v^2R^3+15v^4R),\\
 d_8&=D(R^7+7v^2R^5+35v^4R^3+105v^6R).
 \end{aligned}
 \tag{GCB29}\label{gcb:degree-eight-errors}
\]
All constants are independent of the total number of Gaussian
coordinates.  The factors \(U_r\) retain the actual localization
and scale of each inserted statistic.

Absolute mixed moments have a second bound useful when subtracting
products of moments.  H\"older with \(2m\) equal exponents and
\eqref{gcb:linear-moment-domination} gives, for either law,
\[
 \left|\mathbb E_*\prod_{r=1}^{2m}L_r\right|
       \le(2m-1)!!\prod_{r=1}^{2m}v_r
       \le(2m-1)!!v^{2m}\prod_{r=1}^{2m}U_r,
 \qquad \mathbb E_*\in\{\nu,\mathbb E\}.
 \tag{GCB30}\label{gcb:mixed-moment-bound}
\]
No independence of the \(L_r\) is used.  The simultaneous use
of the support bound in the defect and the variance bound in
unaffected moment factors avoids an unnecessary cutoff radius
in every term.

\subsection*{Localized quadratic observables and their connected errors}

Write each real quadratic observable in a finite factored form
\[
 Q_j=\sum_\alpha c_{j,\alpha}
                      L_{j,\alpha,1}L_{j,\alpha,2},\qquad
 \mathcal A(Q_j)=\sum_\alpha |c_{j,\alpha}|
                        U(L_{j,\alpha,1})U(L_{j,\alpha,2}).
 \tag{GCB31}\label{gcb:quadratic-coefficient-norm}
\]
Choose common \(R,v\) for this finite list as in
\eqref{gcb:local-radius-data}.  Every product of \(p\le4\)
quadratic observables is a sum of words of degree \(2p\).
Equations \eqref{gcb:degree-eight-errors} and
\eqref{gcb:mixed-moment-bound} can therefore be summed with
absolute coefficients, with factor \(\prod_j\mathcal A(Q_j)\).
This step counts the inserted coefficients, not all coordinates
of the integration space.

For clarity define the joint cumulant by the finite partition formula
\[
 \kappa_*(Q_1,\ldots,Q_p)=
 \sum_{\pi\in\mathfrak P_p}(-1)^{|\pi|-1}(|\pi|-1)!
                  \prod_{A\in\pi}\mathbb E_*\prod_{j\in A}Q_j.
 \tag{GCB32}\label{gcb:quadratic-cumulant-definition}
\]
Write \(\Delta_p=\kappa_\nu(Q_1,\ldots,Q_p)
-\kappa_0(Q_1,\ldots,Q_p)\).  The following explicit estimates result:
\[
 \begin{aligned}
 |\nu Q_1-\mathbb E Q_1|&\le d_2\mathcal A(Q_1),\\
 |\Delta_2|
  &\le(d_4+2v^2d_2)\mathcal A(Q_1)\mathcal A(Q_2),\\
 |\Delta_3|
  &\le(d_6+3v^2d_4+15v^4d_2)
                              \prod_{j=1}^3\mathcal A(Q_j),\\
 |\Delta_4|
  &\le(d_8+4v^2d_6+30v^4d_4+156v^6d_2)
                              \prod_{j=1}^4\mathcal A(Q_j).
 \end{aligned}
 \tag{GCB33}\label{gcb:quadratic-connected-errors}
\]
Here the subscript zero always refers to the original untruncated
Gaussian with covariance \(\Gamma\).  In particular its quadratic
third and fourth cumulants are generally nonzero.  They are the
Gaussian Wick values, not quantities being set equal to zero.

We give the product bookkeeping in \eqref{gcb:quadratic-connected-errors}.
By multilinearity it is enough to use one factored monomial for each
\(Q_j\).  For covariance the full moment difference is bounded by
\(d_4\); replacing the product of its two means costs
\(2v^2d_2\).  At order three the three terms of partition type
\((2,1)\) cost at most
\(3(v^2d_4+3v^4d_2)\).  The product of three means, with
coefficient two, costs at most \(6v^4d_2\).  Together with
\(d_6\), these give the third line.  At order four, the four
partitions of type \((3,1)\) cost
\(4(v^2d_6+15v^6d_2)\); the three of type \((2,2)\) cost
\(18v^4d_4\).  The six of type \((2,1,1)\), with coefficient
two, cost \(12v^4d_4+72v^6d_2\).  The final product of four
means, with coefficient six, costs \(24v^6d_2\).
Their sum gives exactly the fourth line.  Each estimate uses the
telescoping difference of products, bounding unchanged factors by
\eqref{gcb:mixed-moment-bound}.  This establishes the constants
without an assumed connected-moment estimate.

Substituting \eqref{gcb:degree-eight-errors}, the three connected
error polynomials multiplying \(D\) are respectively
\[
 \begin{aligned}
 \mathcal P_2(R,v)&=R^3+5v^2R,\\
 \mathcal P_3(R,v)&=R^5+8v^2R^3+39v^4R,\\
 \mathcal P_4(R,v)&=R^7+11v^2R^5+85v^4R^3+411v^6R.
 \end{aligned}
 \tag{GCB34}\label{gcb:connected-error-polynomials}
\]
The same method covers mixed linear and quadratic observables:
expand each quadratic into its two factors, use
\eqref{gcb:general-moment-recursion} for each resulting word,
and then the finite partition formula.  An odd total linear degree
vanishes before taking absolute values.  No additional probabilistic
assumption is required for this mixed case.

The reference Gaussian Wick values themselves can be calculated
without a nonsingular physical covariance.  If
\(Q_j=Z^{\mathsf T}A_jZ\) with real symmetric \(A_j\), whitening
the latent Gaussian gives, for \(p\ge1\),
\[
 \kappa_0(Q_1,\ldots,Q_p)
   ={2^{p-1}\over p}\sum_{\pi\in S_p}
       \operatorname{tr}(\Gamma A_{\pi(1)}\cdots
                                      \Gamma A_{\pi(p)}).
 \tag{GCB35}\label{gcb:gaussian-quadratic-cycles}
\]
Indeed, near zero the logarithm of the Gaussian exponential
generating function is
\[
 \frac12\sum_{k\ge1}\frac{2^k}{k}
 \operatorname{tr}\left[\left(\Gamma\sum_jt_jA_j\right)^k\right].
\]
Taking the coefficient of \(t_1\cdots t_p\) gives
\eqref{gcb:gaussian-quadratic-cycles}.  This also identifies precisely
the contractions to which the sharp-cutoff cumulants are compared.

\subsection*{The fixed-step flat chart realizes the covariance hypothesis}

We now specialize to the actual first plain-average fibre at the
identity, with fixed odd block size \(L\), in current unit-spacing
coordinates.  Let \(P_L\) be the plain derivative of the nonlinear
average, \(\mathsf C_0\) its pivot inclusion with identity free
restriction, and \(\mathcal T\) the internal contour matrix.  The
definitions and the actual first-step precision are
\[
 \begin{aligned}
 &\operatorname{ran}\mathsf C_0=\ker P_L,
   \qquad (\mathsf C_0u)|_{\mathrm{free}}=u,\\
 &\mathsf K_L=\mathsf C_0^{\mathsf T}
                  (d_1^{\mathsf T}d_1+
                         \mathcal T^{\mathsf T}\mathcal T)\mathsf C_0,
 \quad Y\sim N(0,\mathsf K_L^{-1}),\quad Z=\mathsf C_0Y.
 \end{aligned}
 \tag{GCB36}\label{gcb:actual-flat-precision}
\]
The quadratic form here is exactly the first-step form, with no
previous finer-level minimization.  This is the classical reference with action weight
\(x=\theta\), so \(\tau=x/\theta=1\).  The canonical normalization
\(\theta g^2=1\) then gives coefficient one both to the Wilson curl
form and to the prescribed actual soft gauge.  In particular
\(\mathcal T\) is not a divergence gauge or a freely chosen
projector.

The pivot matrix has finite range with at most \(C_L\) nonzero
blocks in each row and column and entries bounded by \(C_L\).
The column assertion follows by the same double-block support as
the row assertion: a free bond belongs to only finitely many
double blocks when \(L\) is fixed.  The flat gap and the
exponential inverse estimate of
Section~\ref{ch:flat-exponential-kernels} give
\(\|\mathsf K_L^{-1}(b,c)\|\le C_L e^{-\mu_Ld(b,c)}\), uniformly
on the admitted tori.  Consequently
\[
 \Gamma=\mathsf C_0\mathsf K_L^{-1}\mathsf C_0^{\mathsf T},
 \quad
 \max\left(\sup_b\sum_c\|\Gamma_{bc}\|,
            \sup_c\sum_b\|\Gamma_{bc}\|\right)\le S_L,
 \quad \Gamma_{bb}\le\sigma_L^2 I.
 \tag{GCB37}\label{gcb:realized-flat-schur}
\]
To check volume uniformity, a four-dimensional shell of distance
\(m\) has at most \(C_L(m+1)^3\) bond indices in block
coordinates.  Summing \(e^{-\mu_Lm}\) over these shells is
finite; periodic image sums have the same upper bound.  Multiplying
the finite row and column sums of \(\mathsf C_0\), its transpose
and the inverse proves the two sums in
\eqref{gcb:realized-flat-schur}.  The diagonal bound follows, for
example, from either sum.  Increase \(\sigma_L\) to at least
one if convenient.

Thus \eqref{gcb:weighted-schur} is realized with
\(a_b=\mu_b=1\), \(S=S_L\), independent of the number of blocks.
This is a statement at a fixed current spacing and fixed \(L\).
Pointwise decay in a multilevel physical metric would not alone
give its unweighted analogue: the number and normalization of
coordinates at different scales must still be charged in the
actual weighted norm \eqref{gcb:weighted-schur}.

\subsection*{Why the actual contour gauge leaves no physical quadratic residue}

We prove the needed range statement rather than assuming that an
arbitrary soft gauge can be integrated away.  Let \(\mathcal S_0\)
be the real site fields vanishing at every block center.  For any
target internal contour data \(s=(s_{y,x})\), \(x\in B(y)\setminus
\{y\}\), define the site field \(R s\) by
\((Rs)(x)=s_{y,x}\) and \((Rs)(y)=0\).  Blocks partition the
sites, so this is well defined.  The plain-gradient identity and
the internal path telescoping identity give exactly
\[
 \begin{aligned}
 &P_Ld_0Rs=0,\qquad d_1d_0Rs=0,\qquad
 (\mathcal T d_0Rs)_{y,x}=s_{y,x},\\
 &\|Rs\|_2=\|s\|_2,\qquad \|d_0Rs\|_2\le4\|s\|_2.
 \end{aligned}
 \tag{GCB38}\label{gcb:actual-gauge-right-inverse}
\]
Here \(d_0\) is the forward site gradient and \(d_1\) the
plaquette circulation.  The last norm estimate follows by summing
\(|f(x+e_\mu)-f(x)|^2\le2|f(x+e_\mu)|^2+2|f(x)|^2\)
over four directions.  Internal contour means average paths, but
on a gradient every one of those paths gives the same difference
\(f(x)-f(y)\); hence no path count or \(L^{-4}\) factor occurs
in the third identity.  The constraint identity is
\(P_Ld_0f=d_L(f|_{\mathrm{centers}})\), not a site-average
gradient identity.

In particular the image of \(\mathcal T\) on
\(\ker d_1\cap\ker P_L\) is the entire internal target space.
It equals the image of \(\mathcal T\) on \(\ker P_L\).
This is the precise range condition that prevents a leftover
physical gauge penalty.  If an arbitrary alternative gauge failed
this condition, integrating its gauge directions could instead
leave a positive quadratic form in curvature.

Put \(\mathcal E_C=d_1\ker P_L\).  The linear map
\[
 \ker P_L\longrightarrow\mathcal E_C\oplus\mathcal S_0,
 \qquad a\longmapsto(d_1a,\mathcal T a)
 \quad\hbox{is a bijection}.
 \tag{GCB39}\label{gcb:curvature-gauge-coordinates}
\]
It is onto because any initial lift of a prescribed curvature in
\(\mathcal E_C\) can be adjusted to any prescribed contour data
by adding \(d_0Rs\) from
\eqref{gcb:actual-gauge-right-inverse}.  It is one-to-one because
the flat gap on \(\ker P_L\) says that both \(d_1a=0\) and
\(\mathcal T a=0\) force \(a=0\).  Equivalently use the actual
internal-bond and crossing-face zero-mode argument of
Section~\ref{ch:flat-contour-precision-gap}.

Apply the constant linear change of variables
\(Y\mapsto(F,s)=(d_1\mathsf C_0Y,\mathcal T\mathsf C_0Y)\).
Its Jacobian is a nonzero constant between the indicated Euclidean
spaces.  The quadratic exponent in
\eqref{gcb:actual-flat-precision} becomes exactly
\(\frac12\|F\|^2+\frac12\|s\|^2\).  The constant Jacobian and
the integral in \(s\) cancel from normalized curvature
expectations.  Therefore
\[
 F=d_1 Z\ \hbox{is a standard Gaussian on }\mathcal E_C,
 \qquad d_1\Gamma d_1^{\mathsf T}=P_{\mathcal E_C}.
 \tag{GCB40}\label{gcb:constrained-curvature-projection}
\]
The same calculation with any prescribed positive quadratic form
in \(s\) changes only the gauge normalization.

If the action weight is instead \(x=\tau\theta\), \(\tau>0\),
multiply the curl term in \eqref{gcb:actual-flat-precision} by
\(\tau\), leaving the reference soft gauge coefficient one.
The same bijection gives exponent
\(\frac{\tau}{2}\|F\|^2+\frac12\|s\|^2\), so the corresponding
covariance satisfies
\(d_1\Gamma_\tau d_1^{\mathsf T}=\tau^{-1}P_{\mathcal E_C}\).
All numerical reference applications below use \(\tau=1\).
No uniform covariance bound as \(\tau\downarrow0\) is inferred.

\subsection*{The precise relation with the full free-curvature comparison}

Let \(\mathcal E=\operatorname{ran}d_1\) on the same finite
periodic lattice, with each positive plaquette counted once.
Generally \(\mathcal E_C\) is a proper subspace of
\(\mathcal E\).  With \(\mathcal E_\perp=\mathcal E\cap
\mathcal E_C^\perp\), ordinary orthogonal decomposition gives
\[
 P_{\rm ex}=P_{\mathcal E}=P_{\mathcal E_C}+P_{\mathcal E_\perp},
 \qquad F_{\rm full}=F_C+F_\perp,
 \tag{GCB41}\label{gcb:free-curvature-decomposition}
\]
where the two summands on the right can be taken independent
standard Gaussians on their respective subspaces.  This proves
that the law of \(F_{\rm full}\) is the full exact-cochain
Gaussian used in the free-kernel comparison.  Conditional on
\(F_\perp\), its centered fluctuation has exactly the curvature
law \eqref{gcb:constrained-curvature-projection}.  Conversely,
integrating this complementary Gaussian supplies the missing
curvature covariance.

The decomposition is an exact classical Gaussian identity.  It
does not identify the law of a coarse mean in an interacting run.
In particular it does not permit replacing
\(P_{\mathcal E_C}\) by \(P_{\rm ex}\) in a single fixed-fibre
Gaussian estimate.  Fine fluctuations alone omit the second
term of \eqref{gcb:free-curvature-decomposition}; exterior or
coarse averaging is a separate covariance contribution in the
complete conditional-law identity.

There is a finite-dimensional expression for this distinction
which needs no chosen coarse-coordinate parametrization.  Let
\(Q_C\) be any matrix whose kernel inside \(\mathcal E\)
is \(\mathcal E_C\).  Then
\[
 P_{\mathcal E_C}=P_{\rm ex}
  -P_{\rm ex}Q_C^{\mathsf T}
       (Q_CP_{\rm ex}Q_C^{\mathsf T})^+
                         Q_CP_{\rm ex}.
 \tag{GCB42}\label{gcb:projection-correction}
\]
To prove the formula, set \(B=Q_CP_{\rm ex}\) restricted to
\(\mathcal E\).  Singular-value decomposition shows that
\(B^{\mathsf T}(BB^{\mathsf T})^+B\) is the orthogonal
projection onto \((\ker B)^\perp\) within \(\mathcal E\).
Subtract it from the identity on \(\mathcal E\), then extend
by zero on \(\mathcal E^\perp\).  This yields exactly
\eqref{gcb:projection-correction}.  The identity is not an
estimate that the correction is small.  Such an estimate on a
particular physical window would require its own quantitative proof.

At fine physical spacing \(\eta\), write the canonical raw
curvature as \(F\).  The physical curvature of this canonically
normalized field is \(\eta^{-2}F\).  The full exact-cochain covariance is therefore
\(\eta^{-4}P_{\rm ex}\), precisely the normalization of
the free-curvature comparison in Section~\ref{ch:lattice-curvature-comparison}.
Its plaquette components are sampled at their own barycentres.
For three independent color copies and
\(Q_f=\frac12\sum_p f(x_p)|F_p|^2\), Gaussian Wick contraction
gives
\[
 \operatorname{Cov}_0(Q_f,Q_h)
     ={1\over2}\sum_{p,q}f(x_p)h(x_q)
                    \|\mathbb E[F_pF_q^{\mathsf T}]\|_{\rm HS}^2.
 \tag{GCB43}\label{gcb:curvature-wick-readout}
\]
The physical Riemann weight \(\eta^4\) multiplying the squared
physical curvature cancels its \(\eta^{-4}\) scaling in
\(Q_f\).  Thus the $\operatorname{tr}(F^2)$ readout is the unweighted
plaquette sum displayed here; inserting an additional
\(\eta^4\) into that sum would change the observable.

For the full law in \eqref{gcb:free-curvature-decomposition},
\eqref{gcb:curvature-wick-readout} has the full free-curvature kernel and its
proved separated-smearing comparison.  For the first-step
fluctuation it has \(P_{\mathcal E_C}\) instead.  The cutoff
estimates compare each specified Gaussian reference with its own
Wick kernel.  They transfer that comparison only after the reference used in
the application has actually been identified with, or compared
to, the full exact-cochain reference.  The decomposition records
how this works in the free Gaussian theory without asserting it
for a conditional interacting measure.

\subsection*{Uniform cutoff errors for the curvature polynomial}

Take all physical tangent bond radii equal to \(B\ge\sigma_L\)
in \eqref{gcb:actual-flat-precision}, and retain counting weights
and unit amplitudes.  From
\eqref{gcb:large-radius-flux}, including all ranks from one
through three and enlarging the fixed constants,
\[
 \sup_b\beta_b^*\le\beta_L(B)
       :=C_L B^2e^{-c_LB^2},\qquad
 D\le S_L\beta_L(B),\qquad R=B,\quad v\le\sigma_L.
 \tag{GCB44}\label{gcb:flat-exponential-error}
\]
The bounds remain valid if some physical blocks are deterministic.
For a sampled test function let
\(W_f=\sum_p|f(x_p)|\), and define the canonical quadratic
curvature polynomial
\[
 Q_f(Z)={1\over2}\sum_p f(x_p)|d_1Z(p)|^2,
 \qquad \mathcal A(Q_f)\le24W_f.
 \tag{GCB45}\label{gcb:localized-curvature-norm}
\]
Indeed a fixed color component of a plaquette curl has four
coefficients of absolute value one, hence \(U\le4\).
There are three colors and the coefficient is \(f/2\), so
\((1/2)\cdot3\cdot4^2=24\).  Each positive direction pair is
already counted once in the sum over \(p\).  This calculation
requires no estimate on the number of unobserved bonds.

For \(p=2,3,4\), equations
\eqref{gcb:quadratic-connected-errors}--
\eqref{gcb:flat-exponential-error} now imply, with
\(\Delta_p=\kappa_\nu(Q_{f_1},\ldots,Q_{f_p})
-\kappa_0(Q_{f_1},\ldots,Q_{f_p})\),
\[
 \begin{aligned}
 |\nu Q_f-\mathbb E Q_f|&\le C_L\beta_L(B)B W_f,\\
 |\Delta_p|
  &\le C_{L,p}\beta_L(B)
       (B^{2p-1}+B^{2p-3}+\cdots+B)
                        \prod_{j=1}^pW_{f_j}.
 \end{aligned}
 \tag{GCB46}\label{gcb:curvature-connected-error}
\]
Here all expectations are normalized, including the sharp product
cutoff law.  The third and fourth lines implicit in the second
formula use sixth and eighth underlying linear moments.  The
vanishing of odd linear moments alone would not establish them.
The Gaussian values on the right are the actual cycle contractions
\eqref{gcb:gaussian-quadratic-cycles}, with the constrained curvature
projection of \eqref{gcb:constrained-curvature-projection} in this
particular application.

If each test is supported in a ball of current radius \(r\ge1\),
then \(W_f\le C(r+1)^4\|f\|_\infty\).  Thus the covariance
error has the expected two localized volume factors
\(C(r+1)^8\|f\|_\infty\|h\|_\infty\), rather than a factor
proportional to the volume of the torus.  At orders three and
four the factors are \((r+1)^{12}\) and \((r+1)^{16}\).


\subsection*{The nonlinear Wilson readout on the actual joint chart}

At the flat background let \(a=\Phi_0(u)\) be the actual nonlinear
fibre chart of Section~\ref{ch:block-average-pivot-chart}, so
\(u\) is the free raw logarithm and
\(D\Phi_0(0)=\mathsf C_0\).  Its domain contains
\(\|u\|_\infty<r_L\), with \(r_L=cL^{-9}\) after reducing
the fixed constant.  The finite local support and the uniform
holomorphic bounds on that joint domain give, on a smaller real ball,
\[
 \begin{aligned}
 \|\Phi_0(u)-\mathsf C_0u\|_\infty
                      &\le C_L\|u\|_\infty^2,\\
 A_p(e^{i\Phi_0(u)})
       &={1\over2}|d_1\mathsf C_0u(p)|^2+R_p(u),\\
 |R_p(u)|&\le C_L\|u\|_\infty^3.
 \end{aligned}
 \tag{GCB47}\label{gcb:actual-readout-taylor}
\]
For the second assertion use the holomorphic extension
\(1-\frac12\operatorname{tr}U_p\), which equals the real Wilson
readout on \(\mathrm{SU}(2)\).  Its constant and linear terms
at the identity are zero.  Differentiating the four bond products
shows its quadratic term is
\(\frac12|d_1\mathsf C_0u(p)|^2\), with Pauli pairing
\(|X|^2=\frac12\operatorname{tr}X^2\).  To bound the remainder,
apply one-variable Cauchy to the actual local function of \(tu\)
on a disc of radius proportional to
\(r_L/\|u\|_\infty\).  The joint chart bounds this function
uniformly; summing the Taylor terms of order at least three on
a smaller disc gives the displayed cubic estimate.  The first
assertion follows by the same argument at order two.  A plaquette
and its pivot dependencies occupy a fixed number of double
blocks, so neither remainder constant contains the lattice volume.

Set \(u=gY\).  Under the full physical tangent cutoff
\(|\mathsf C_0Y(b)|<B\), the identity free restriction implies
\(\|Y\|_\infty\le B\).  Thus \(gB<r_L/2\) places every
field of this cutoff event in the actual chart, and for
\(O_f=\sum_p f(x_p)A_p\),
\[
 O_f(e^{i\Phi_0(gY)})=g^2Q_f(Z)+\mathcal R_f(Y),
 \quad |\mathcal R_f|\le C_Lg^3B^3W_f,
 \quad |Q_f|\le C_LB^2W_f.
 \tag{GCB48}\label{gcb:nonlinear-readout-error}
\]
This uses the literal sampled $\operatorname{tr}(F^2)$ sum, with no extra
fine-cell volume factor.  The probability measure in this
identity is still the Gaussian cutoff measure, rather than the
nonlinear action measure.

The finite partition definition of a cumulant gives
\(|\kappa(X_1,\ldots,X_p)|\le C_p\prod_j\|X_j\|_\infty\).
Expand each slot of \eqref{gcb:nonlinear-readout-error}.  If
\(gB\le1\), the terms with at least one remainder have total
absolute value at most
\(C_{L,p}g^{2p+1}B^{2p+1}\prod_jW_{f_j}\).
Indeed a remainder replaces one factor \(g^2B^2\) by
\(g^3B^3\); the finite sum over subsets costs at most a
constant times \(gB(1+gB)^{p-1}\).
Together with \eqref{gcb:curvature-connected-error}, this proves
\[
 \begin{aligned}
 &|\kappa_\nu(O_{f_1},\ldots,O_{f_p})
                    -g^{2p}\kappa_0(Q_{f_1},\ldots,Q_{f_p})|\\
 &\quad\le C_{L,p}g^{2p}\left[
       \beta_L(B)(B^{2p-1}+B^{2p-3}+\cdots+B)
                          +gB^{2p+1}\right]\prod_{j=1}^pW_{f_j},
 \qquad p=2,3,4.
 \end{aligned}
 \tag{GCB49}\label{gcb:actual-readout-connected-bound}
\]
The same reasoning for one insertion gives mean error
\(C_Lg^2[\beta_L(B)B+gB^3]W_f\).


\subsection*{Literal radius conversions and comparison of cutoff sets}

At physical fine spacing \(\eta\), raw bond logarithm and physical
potential obey \(a_b=\eta A_b\).  In canonical tangent variables
\(a_b=g Z_b\).  Consequently the conversions are
\[
 |a_b|<\alpha_b\ \Longleftrightarrow\
       |Z_b|<\alpha_b/g,\qquad
 |\ell_b A_b|<r_b\ \Longleftrightarrow\
       |Z_b|<{\eta r_b\over g\ell_b}.
 \tag{GCB50}\label{gcb:literal-radius-conversion}
\]
The second relation refers to the tangent physical potential;
\(\ell_b\) is the actual scale multiplying it in that field
bound.  For the nonlinear chart use
\eqref{gcb:actual-readout-taylor}: its raw field differs from
\(gZ\) by at most \(C_Lg^2B^2\) on the event in question.
Any comparison with a cutoff on the nonlinear field must charge
this error and preserve a stated radius margin.

A free-coordinate box and the complete physical tangent product
box are not equal events.  With
\(c_L=\|\mathsf C_0\|_{\infty\to\infty}\ge1\), one has only
\[
 \{\|Y\|_\infty<B/c_L\}
   \subset\{\|\mathsf C_0Y\|_\infty<B\}
   \subset\{\|Y\|_\infty<B\}.
 \tag{GCB51}\label{gcb:free-physical-box-inclusions}
\]
The first inclusion uses the actual row norm of the pivot map;
the second uses its identity free restriction.  Inclusion alone
does not compare the corresponding normalized conditional
expectations.  The proof in this section works directly with
the middle, complete physical event, so requires no false
identification of these boxes.  If an actual integration rule
instead cuts only free coordinates, apply the abstract theorem
to those actual independent-coordinate blocks and carry the
observable coefficient map separately; it is a different cutoff
law with its own explicitly specified covariance and radii.

For growing canonical radius
\(B=c_*(\log\theta)^s\), \(s>0\), and
\(g=\theta^{-1/2}\), every fixed power of \(\log\theta\)
times \(\beta_L(B)\) tends to zero, as does every fixed such
power times \(g\).  The chart condition \(gB<r_L/2\) holds
eventually.  Thus for support radii bounded by any fixed power
of \(\log\theta\), the product of the bracket and localized volume factors
in \eqref{gcb:actual-readout-connected-bound} tends to zero for test
functions with bounded sup norms, uniformly in torus volume.  This is an absolute error relative
to the displayed insertion norms.  Converting it into a relative
error for a specified two-point coefficient also requires a
proved lower bound for that coefficient in the same test class.

\subsection*{Local limits and the conditional integration-by-parts identity}

The finite-volume cutoff laws have Gaussian-dominated moments
for every fixed local linear statistic by
\eqref{gcb:linear-moment-domination}.  In the fixed-step setting
their marginal variances are bounded uniformly by
\eqref{gcb:realized-flat-schur}.  These bounds give tightness on
each finite collection of bond coordinates.  A diagonal
subsequence on the countable lattice therefore has compatible
finite-dimensional limits.  Uniform higher local moments give
uniform integrability, so the local polynomial estimates above
pass to every such limit.  This conclusion does not assert
uniqueness of the limit, and does not identify it with conditioning
an infinite-volume Gaussian on a global event of positive
probability.  It only uses uniform finite-volume bounds.

The boundary proof also identifies its precise obstruction in
a general conditional law.  If in actual free coordinates, with
\(G>0\), its
complete density is
\(Z_e^{-1}\chi(CY)\exp[-\frac12Y^{\mathsf T}G^{-1}Y+R_e(Y)]\),
integration by parts has, in addition to the boundary defect,
the force term
\[
 \mathbb E_{\nu_e}\!\left[
       P(CY)\,(CG\nabla R_e(Y))_i\right].
 \tag{GCB52}\label{gcb:remaining-complete-force}
\]
This follows by differentiating the entire density, with its
actual normalizer; all integrability and chart-domain conditions
must hold.  An arbitrary exterior linear force can translate
the law toward the boundary, and an arbitrary old activity can
destroy the even log-concavity used in
\eqref{gcb:conditional-weight}.  Neither its small Hessian alone
nor the Gaussian boundary theorem bounds these two effects.

The same radial argument would remain valid for a further even
log-concave factor in the physical variables, because that factor
can be included in the joint log-concave integrand defining
\(w_{b,\delta}\).  Its new force term still has to be bounded
in the complete integration-by-parts identity.  No such convexity,
symmetry or force bound is asserted for the resummed interacting
conditional density.

The interacting application additionally requires the complete density,
exterior averaging and old-activity force estimates of
Section~\ref{ch:two-point-extraction-tests}, and comparison with the full
free-curvature reference.

% BEGIN MODULE chapters/uniform_insertion_bound_v4.tex
\section{A one-field insertion bound uniform in the outside fields}
\label{ch:uniform-insertion}

An estimate for each correlation function does not automatically give
an estimate in the Hilbert norm of an arbitrary polynomial of fields.
The argument below obtains that stronger norm by repeated reflections
at a fixed lattice spacing, before taking a continuum limit. Its
analytic input is the ordinary scalar source upper bound and floor,
not a typed insertion estimate.

Assume the all-mirror invariance and reflection positivity of every
periodic Wilson DLR state, the ordinary torus source hypothesis (ST5),
and the fixed-layer rooted domain geometry (GS1)--(GS3). The
localization size is $M=L^{m_0}$, with $L$ odd. Torus limits use the
cofinal compatible family $\ell=2L^m$, $m\ge m_0$. Every scale cube
$Ma_s$ then tiles the torus, since
$\ell/(Ma_s)=2L^{m-m_0+s}$. Every fixed board of power-of-$L$ side,
and every fixed coarser such board, tiles all sufficiently large tori
in this family. No scale terms are omitted.

\subsection*{A raw-algebra reflection argument}

Fix the spacing $a=L^{-K}$ and a periodic state $\rho$. Take a board
cell $C$ of side $l=L^j\ge a$, centered at a lattice site, with lower
time face $t$. Its face planes are link mirrors: $l/a$ is odd.
Its middle time plane is a site mirror. Both are supplied by the
all-mirror theorem. For bounded finite-support functions above $t$ put
\[
 \langle B,A\rangle_t=\rho(\overline{\theta_tB}A),
 \qquad \|A\|_t^2=\langle A,A\rangle_t.
 \tag{UI1}
\]
For a bounded complex cell function $D$, let
$q_D=\limsup_kp_{l,\ell_k}(D)$ along a compatible cofinal sequence.
Define, on the raw half-space algebra,
\[
 T_DA=D\,\tau_{le_4}A,
 \qquad D^*=\overline{\theta_{t+l/2}D}.
\]
Reflection and translation invariance give the formal adjoint identity
$\langle B,T_DA\rangle_t=\langle T_{D^*}B,A\rangle_t$.
No bounded operator on a quotient has been assumed.

Put $K_D=T_{D^*}T_D$ and $m_j=\langle A,K_D^jA\rangle_t$.
Formal adjoints and reflection Cauchy--Schwarz imply
\[
 m_{2r}=\|K_D^rA\|_t^2,\qquad
 m_{2r+1}=\|T_DK_D^rA\|_t^2,\qquad
 m_j^2\le m_{j-1}m_{j+1}.
 \tag{UI2}
\]
For odd $j=2r+1$ the last inequality uses the two vectors
$K_D^rA,K_D^{r+1}A$; for even $j=2r\ge2$ it uses
$T_DK_D^{r-1}A,T_DK_D^rA$. Thus it is a raw semidefinite-form
calculation, not an application of a spectral theorem to an operator
whose boundedness is still to be proved. If $m_0=0$, Cauchy--Schwarz
already gives $m_1=0$. Otherwise log-convexity gives
\[
 m_1^N\le m_0^{N-1}m_N. \tag{UI3}
\]

The expression for $m_N$ consists of the left endpoint
$\overline{\theta_tA}$, a column of $2N$ alternating translated copies
of $D^*$ and $D$, and the right endpoint $\tau_{2Nle_4}A$.
Choose a coarse side $l'=L^rl$ large enough to contain $A$ in a
coarse half-space cell above $t$, spatially aligned with $C$.
The oddness of $L^r$ makes all coarse faces fine-board faces.
Take $N$ through sufficiently large multiples of $L^r$.

On this coarse grid retain fixed endpoint zones containing both
endpoint supports and, if convenient, a fixed number $c_A$ of the
adjacent column factors. Bound the complete endpoint functions by
their fixed-cutoff supremum norms. This yields a finite constant
$C(a,A,D)$ independent of $N$. In the middle, every occupied coarse
cell is a product of fine cell functions only. Fine-to-coarse
chessboard refinement and product domination (ST3) apply there.
Grid-preserving translates and reflected conjugates of $D$ have the
same fine seminorm. Therefore, for every $\epsilon>0$,
\[
 m_N\le C(a,A,D)(q_D+\epsilon)^{2N-c_A}.
 \tag{UI4}
\]
The volume limit is taken first for each fixed $N$. The estimate uses
no assertion $p(AD)\le\|A\|_\infty p(D)$: a mixed endpoint block,
if present, is bounded in its entirety by its supremum norm.
The endpoint zones contain only a bounded number of coarse cells;
their factor count and supremum bound stay fixed as the column grows.

Taking the $N$th root in (UI3)--(UI4), then sending $N\to\infty$
and finally $\epsilon\downarrow0$, gives
\[
 \|T_DA\|_t\le q_D\|A\|_t,
 \qquad
 |\rho(\overline{\theta_tB}D\tau_{le_4}A)|
       \le q_D\|B\|_t\|A\|_t.
 \tag{UI5}
\]
The same order of limits proves this when $q_D=0$; no negative power
of $q_D$ is introduced. It also proves that null vectors are mapped
to null vectors. Only now may the map be regarded as a bounded
operator on the Hilbert quotient. Endpoint constants may depend on
the number of fields inside $A$, but they disappear before (UI5).
The final constant does not depend on that number, on $B$, or on
their polynomial coefficients.

\subsection*{One ordinary source and a common centering}

Let $h$ be real and supported in the ball of radius $l/4$ around the
cell center. For $a\le l/64$, apply the three-range count in (GS4)
with board side $\eta=l$, support radius $\delta=l/4$ and the same
face margin. Since $\eta/\delta=4$, its face cost is bounded
uniformly as $l\downarrow0$. Coarse scales use the global count
divided by the number of cells, not a separate root count at every
source copy. Upper stability and the floor cost only $Cl^4$.
Thus, for $l\le1$, there are holomorphic univariate normalizers with
\[
 E_{h,\ell}(0)=0,\qquad
 p_{l,\ell}\big(e^{z\mathcal O_a(h)-E_{h,\ell}(z)}\big)\le Q_0,
 \qquad |z|\|h\|_\sharp<r_0,
 \tag{UI6}
\]
where $Q_0,r_0>0$ are independent of the spacing and cell side.
They are also independent of the surrounding fields. At a fixed
cutoff these normalizers are locally bounded uniformly in the torus,
so a compatible Montel subsequence gives one $E_h$, independently of
every infinite-volume state.

The remaining coarse mesh ratios can also be included. If
$3\le l/a<64$, the ball and its contributing plaquettes lie wholly
inside the cell, and
$\|\mathcal O_a(h)\|_\infty\le C64^4\|h\|_\infty$.
Set $E_{h,\ell}=0$ and enlarge $Q_0$ in (UI6). The only smaller
allowed ratio is $l/a=1$; then a radius-$a/4$ ball about a lattice
site contains no plaquette barycentre, so $\mathcal O_a(h)=0$.
Consequently (UI6) covers every allowed ratio, with the same constants.

For fixed $z$ put $D_z=e^{z\mathcal O_a(h)-E_h(z)}$.
Homogeneity of the seminorm and the locally uniform convergence of
the normalizers give $q_{D_z}\le Q_0$. Apply (UI5) and then Cauchy's
formula on the one source disc. The expectation is holomorphic,
because at fixed cutoff all observables are bounded. We obtain
\[
 \left|\rho\left(\overline{\theta_tB}
       [\mathcal O_a(h)-E_h'(0)]\tau_{le_4}A\right)\right|
 \le(Q_0/r_0)\|h\|_\sharp\|B\|_t\|A\|_t.
\]
Taking $A=B=1$ in a second arbitrary periodic state $\nu$ bounds
$|\nu(\mathcal O_a(h))-E_h'(0)|$ by the same expression without
the outside norms. The case $D=1$ of (UI5) handles the scalar error.
Hence
\[
 \left|\rho\left(\overline{\theta_tB}
       [\mathcal O_a(h)-\nu(\mathcal O_a(h))]
                         \tau_{le_4}A\right)\right|
 \le C\|h\|_\sharp\|B\|_t\|A\|_t.
 \tag{UI7}
\]
Both states are arbitrary. Large cutoff-dependent normalizer
derivatives cancel; uniform moments of perturbed Gibbs states have
not been assumed. In using (UI5) for $D_z$, endpoint supremum bounds
may depend on $z$ and the cutoff. They disappear pointwise before the
uniform bound $q_{D_z}\le Q_0$ is used in Cauchy's formula.

\subsection*{Continuum passage with fixed support gaps}

Suppose $a_j\downarrow0$ and two periodic states $\rho_j,\nu_j$ at
each datum give convergent centered separated moments
\[
 \sigma_j(f_1,\ldots,f_n)
 =\rho_j\prod_i[\mathcal O_{a_j}(f_i)
                         -\nu_j(\mathcal O_{a_j}(f_i))]
 \longrightarrow\sigma(f_1,\ldots,f_n).
\]
Use the uniform all-state local moment bounds already proved to obtain
equicontinuity under small $C^1$ test changes. In particular, on each
fixed separated tuple, $\sigma_j(\tau_{v_j}X)\to\sigma(X)$ whenever
$v_j\to0$. This follows directly by multilinear telescoping and
$\|\tau_vf-f\|_\sharp\to0$ for compact $C^1$ functions. The constants
in this passage may depend on the fixed tuple and its length; they
do not enter the constant in (UI7).
Approximating a fixed real translation by lattice translations also
proves translation invariance of $\sigma$, since both states in each
datum are lattice-translation invariant.

Let $F,G$ be arbitrary finite complex linear combinations of field
products whose entries have compact supports strictly in positive
time and are separated within each product. Define
$\|[F]\|_\sigma^2=\sigma(\theta\overline F\cdot F)$.
This is nonnegative as a limit of the lattice reflection forms.
For a fixed power $l\le1$ and a real $h$ supported in
$\bar B((\vec y,0),l/4)$, choose the nearest spatial lattice point
$\vec y_j$ and put
\[
 v_j=(\vec y_j-\vec y,a_j/2),\quad t_j=a_j/2,
 \quad x_{*,j}=(\vec y_j,l/2+a_j/2).
\]
For all sufficiently fine cutoffs $l/a_j$ is odd. Therefore $x_{*,j}$
is a lattice site and the cell faces are link mirrors. The shifted
inserted test is explicitly
$h_j=\tau_{v_j+(l/2)e_4}h$, supported around that cell center.
The identity $\theta_{t_j}\tau_{v_j}=\tau_{v_j}\theta_0$ identifies
the lattice reflected norm tuples with the corresponding small
translations of the continuum norm tuples. No half-mesh translation
invariance of $\nu_j$ is being assumed.

Apply (UI7) to the centered lattice polynomials associated to
$\tau_{v_j}F$ and $\tau_{v_j}G$. All insertion and norm tuples have
fixed positive support gaps for large $j$. Their values converge by
the preceding equicontinuity, including when a limiting Hilbert norm
is zero. Thus
\[
 |\sigma(\theta\overline G\cdot(\tau_{(l/2)e_4}h)
                                      \cdot\tau_{le_4}F)|
 \le C\|h\|_\sharp\|[G]\|_\sigma\|[F]\|_\sigma.
 \tag{UI8}
\]
The constant is independent of the numbers of entries and the finite
polynomial coefficients of $F,G$. Arbitrary centering states $\nu_j$
are retained throughout.

For completeness time translations are contractions in these norms.
At fixed spacing, reflection Cauchy--Schwarz gives
$\|\tau_uA\|_t^2\le\|A\|_t\|\tau_{2u}A\|_t$ for every positive
lattice time $u$. Iteration against
$\|\tau_{2^Nu}A\|_t\le\|A\|_\infty$ proves contraction. The same
fixed-gap limit passage proves it for every positive real time in
the continuum. It uses no bounded insertion operator in advance.

\subsection*{A $C^1$ insertion estimate near the time mirror}

We record the form needed independently of any plate or angular
argument. Let $0<\eta\le1$, and let $f\in C_c^1$ be supported in
$\{\eta\le x_4\le3\eta\}$, with support diameter at most $R$.
Take the largest power $l$ of $L$ below $\eta$ and partition $f$
at mesh $cl$ into pieces in radius-$l/4$ balls. There are at most
$C_L(1+R)^3\eta^{-3}$ pieces: only the three spatial directions
have unbounded length in units of $\eta$. Since $f$ vanishes on both
time faces, $\|f\|_\infty\le\eta\|\partial_4f\|_\infty$.
This cancels the $1/l$ from the partition derivative, because
$\eta/l\le L$. Each piece has sharp norm at most
$C\|f\|_{C^1}$.

Translate a piece to the middle of its side-$l$ time slab and use
(UI8). Its center time lies between $\eta-l/4$ and $3\eta+l/4$;
both outside factors therefore retain a positive margin of at least
$\eta-3l/4$. Time contraction absorbs those margins. Summing gives
\[
 |\sigma(\theta\overline G\cdot(f)\cdot\tau_{4\eta e_4}F)|
 \le C(1+R)^3\eta^{-3}\|f\|_{C^1}
                                  \|[G]\|_\sigma\|[F]\|_\sigma.
\]
Taking $G=(f)\cdot\tau_{4\eta e_4}F$ and dividing its norm proves
the field-operator estimate; splitting real and imaginary parts handles
complex $f$. We may weaken $\eta^{-3}$ to $\eta^{-4}$ and write
\[
 \|[(f)\cdot\tau_{4\eta e_4}F]\|_\sigma
 \le C(R)\eta^{-4}\|f\|_{C^1}\|[F]\|_\sigma.
 \tag{UI9}
\]

\begin{theorem}[Uniform one-field insertion]
\label{thm:uniform-one-field-insertion}
Under the stated ordinary source, rooted-domain and lattice-symmetry
hypotheses, every limit just described satisfies (UI8)--(UI9).
Choose a fixed $0<c<1/64$ and let $0<\eta\le1$. If a real $k\in C_c^1$
is supported in a ball of radius $c\eta$ centered at $y$,
$|y_4|\le\eta/8$, and all
entries of $F,G$ lie strictly above $\eta/4$, then
\[
 |\sigma(\theta\overline F\cdot(k)\cdot G)|
 \le C\eta^{-4}\|k\|_\sharp
                  \|F\|_{(\eta/4)}\|G\|_{(\eta/4)},
 \quad
 \|F\|_{(\eta/4)}=\|[\tau_{-(\eta/4)e_4}F]\|_\sigma.
 \tag{UI10}
\]
The constant depends only on the ordinary source and localization
constants and the fixed geometric choices. It is independent of the
defining sequence, the states, support positions, and the numbers of
fields in the surrounding finite complex polynomials.
\end{theorem}

Indeed put $\xi=c\eta$ and translate the insertion in time by
$2\xi-y_4$, placing it in $[\xi,3\xi]$. After this translation the
left and right polynomials are, respectively,
$\tau_{y_4-2\xi}F$ and $\tau_{-y_4-2\xi}G$, the latter followed by
$\tau_{4\xi e_4}$. Each is a positive time translate of its
$\eta/4$-downshift by at least $\eta/8-2\xi>0$.
Apply (UI9), reflection Cauchy--Schwarz and contraction, absorbing
the fixed factor $c^{-4}$. This proves (UI10).
The theorem asserts the $C^1$ insertion bound only. It does not assert
a plate density theorem, an $L^1$ insertion bound, or full rotational
covariance.

% END MODULE chapters/uniform_insertion_bound_v4.tex


% BEGIN MODULE chapters/plate_density_from_insertion_v5.tex
\section{Analytic insertion densities and the plate bound}
\label{ch:plate-density-from-insertion}

This section derives an $L^1$ insertion bound from the proved $C^1$
Hilbert insertion bound. All operators are regularized by positive
Euclidean time. No unsandwiched point field is presumed to exist.
The exact spectral cube cone is used, but rotational covariance is not.

Let $\sigma$ be a symmetric, real, time-reflection-invariant,
translation-invariant
functional on admissible separated tests, with its reflection-positive
Hilbert reconstruction. Write $\mathcal D$ for the dense classes of
finite complex positive-time polynomials and $\Omega=[1]$.
Use the probability normalization $\sigma(1)=1$, so $\Omega$ is a
unit vector fixed by all translations.
Positive time translation is $e^{-tH}$, spatial translation is
$U(\mathbf x)=e^{i\mathbf x\cdot\mathbf P}$, and these operators
strongly commute. Assume their joint spectral measure is carried by
\[
 E\ge0,\qquad |q_j|\le E\quad(j=1,2,3).
 \tag{PD1}
\]
Also assume local $C^1$ continuity of the separated moments, as proved
from the ordinary source bounds. This continuity is used for fixed
finite polynomials only; none of its field-number constants enter the
uniform bounds below.

The insertion input has fixed constants $C_I,c_I>0$, $p\ge0$, with
$c_I<1/64$. For $0<\eta\le1$, let
\[
 I(F;h;G)=\sigma(\theta\overline F\cdot(h)\cdot G),\qquad
 \|F\|_{(\eta/4)}=\|[\tau_{-(\eta/4)e_4}F]\|.
\]
Whenever the entries of $F,G$ lie strictly above $\eta/4$ and $h$ is
real with support in a radius-$c_I\eta$ ball centered at $y$,
$|y_4|\le\eta/8$, the input is
\[
 |I(F;h;G)|\le C_I\eta^{-p}\|h\|_\sharp
                       \|F\|_{(\eta/4)}\|G\|_{(\eta/4)}.
 \tag{PD2}
\]
The norm is $\|h\|_\sharp=\max(\|h\|_\infty,c_M\|\nabla h\|_\infty)$.
Splitting real and imaginary parts permits complex $h$, at a fixed
factor two. The preceding insertion theorem supplies $p=4$; keeping
$p$ general makes the additional five powers transparent.

\subsection*{Bounded forms, real translations and holomorphy}

For $h$ supported in $\bar B(0,c_I\eta)$, define the form on raw
positive-time polynomials by
\[
 b_{h,\eta}(A,B)=I(\tau_{(\eta/4)e_4}A;h;
                                      \tau_{(\eta/4)e_4}B).
\]
By (PD2) it descends through null vectors and extends to a bounded
sesquilinear form on the Hilbert space. Riesz representation gives
a unique bounded operator $B_{h,\eta}$, with
\[
 \|B_{h,\eta}\|\le C\eta^{-p}\|h\|_\sharp,
 \quad
 B_{h,\eta}[A]=[(\tau_{(\eta/4)e_4}h)
                                \cdot\tau_{(\eta/2)e_4}A].
 \tag{PD3}
\]
The vector identity follows by testing against the dense polynomial
classes and translating the insertion form. Its displayed product
is admissible: the insertion is strictly positive in time and is
separated from the shifted entries of $A$. Thus this identity proves
the action of the bounded operator; it does not postulate a point
field on the quotient.

For complex $z=(z_4,\mathbf z)$ put
\[
 T(z)=\int e^{-z_4E+i\mathbf z\cdot\mathbf q}\,d\Pi(E,\mathbf q).
\]
The spectral cone gives $\|T(z)\|\le1$ if
$\operatorname{Re}z_4\ge\sum_j|\operatorname{Im}z_j|$.
It is operator-norm holomorphic in the strict interior. Indeed, on
a neighbourhood with margin $\epsilon>0$, derivatives of total order
$m$ of the integrand are bounded by $E^m e^{-\epsilon E}$.
Its supremum is at most $m!\epsilon^{-m}$. The exponential Taylor
series consequently converges uniformly in spectral variables on
smaller complex neighbourhoods, so its spectral integrals converge
in operator norm and may be differentiated there.

Now take $h$ supported in $\bar B(0,c_I\eta/2)$, and set
$\Phi=[\tau_{-(\eta/4)e_4}F]$, $\Psi=[\tau_{-(\eta/4)e_4}G]$.
For real $x$, $|x_4|<\eta/8$, direct use of (PD3) gives
\[
 I(F;\tau_xh;G)=\langle\Phi,M_{h,\eta}(x)\Psi\rangle,
 \quad M_{h,\eta}(z)=
 T(\tfrac\eta8e_4+z)B_{h,\eta/2}
 T(\tfrac\eta8e_4-z).
 \tag{PD4}
\]
For clarity, the real time factors in this formula are
$e^{-(\eta/8+x_4)H}$ and $e^{-(\eta/8-x_4)H}$; the spatial
factors are $U(\mathbf x)$ and $U(-\mathbf x)$.
Substitution in (PD3), followed by one common time translation,
recovers exactly $\theta\overline F\cdot(\tau_xh)\cdot G$.
Both factors in (PD4) are contractions on
\[
 \mathcal O_\eta=
 \{z:|\operatorname{Re}z_4|+
                     \textstyle\sum_{j=1}^3|\operatorname{Im}z_j|
                                                    <\eta/8\}.
 \tag{PD5}
\]
Hence $M_{h,\eta}$ is operator-norm holomorphic there and
$\|M_{h,\eta}(z)\|\le C\eta^{-p}\|h\|_\sharp$.
There is no restriction on real spatial translations.

\subsection*{A scaled compact parametrix, with its distribution proof}

In four dimensions let
$G(x)=-(64\pi^2)^{-1}|x|^2\log|x|$, with value zero at the origin.
It is $C^1$. For $x\ne0$,
\[
 \Delta(|x|^2\log|x|)=8\log|x|+6,
 \quad \Delta\log|x|=2|x|^{-2},
 \quad \Delta|x|^{-2}=0.
\]
The first two identities also hold as distributions: Green's formula
on $\{|x|>\epsilon\}$ has boundary errors tending to zero, respectively
of orders $O(\epsilon^4|\log\epsilon|)$ and $O(\epsilon^2)$,
up to smaller terms. In the third identity the boundary term is
nonzero. Against a compact smooth test $f$, the inner normal is
$-\partial_r$, so the term containing $f\partial_n|x|^{-2}$ tends
to $-2|S^3|f(0)=-4\pi^2f(0)$; the other boundary term is $O(\epsilon)$.
Thus
\[
 \Delta|x|^{-2}=-4\pi^2\delta_0,
 \qquad \Delta^3G=\delta_0.
 \tag{PD6}
\]

Choose real $\chi\in C_c^\infty(B(0,1))$ equal to one near zero.
Put $n=\chi G$ and $s=\delta_0-\Delta^3n$. Since $\chi=1$ near
the singularity, (PD6) cancels its delta term. Every remaining
Leibniz term contains a derivative of $\chi$ supported away from
zero, where $G$ is smooth. Therefore $s\in C_c^\infty(B(0,1))$.
For $0<r\le1$ define
\[
 n_r(x)=r^2n(x/r)
       =-\frac{\chi(x/r)|x|^2\log(|x|/r)}{64\pi^2},
 \qquad s_r(x)=r^{-4}s(x/r).
 \tag{PD7}
\]
Distributional scaling in four dimensions gives
$\Delta^3n_r=\delta_0-s_r$: six derivatives contribute $r^{-6}$,
the prefactor contributes $r^2$, and the scaled delta contributes
$r^4$. Ordinary differentiation gives
\[
 \|n_r\|_\sharp\le Cr,\qquad
 \|s_r\|_\sharp\le Cr^{-5}.
 \tag{PD8}
\]
Both supports lie in $\bar B(0,r)$. The logarithm is
$\log(|x|/r)$, not $\log|x|$ after simply cutting off $G$.
The subtracted quadratic polynomial removes a spurious $|\log r|$
in (PD8).

For smooth compact $k$, convolution with (PD6)--(PD7) gives
\[
 k=(\Delta^3k)*n_r+k*s_r.
 \tag{PD9}
\]
This can be verified directly against compact smooth tests by Fubini
and distributional integration by parts. Alternatively, the compact
$C^1$ kernel $n_r$ is integrable, so all derivatives of its convolution
with $k$ may be placed on $k$ under the integral; the distribution
identity then yields (PD9). Both sides are $C^1$, so their
distributional equality is also pointwise equality.

\subsection*{A bounded analytic density and the $L^1$ bound}

Fix $\epsilon=\min(c_I/8,1/1024)$ and put $r=\epsilon\eta$.
Both $n_r$ and $s_r$ fit the source ball used in (PD4).
Define the operator-valued holomorphic function
\[
 V_\eta(z)=\Delta_z^3M_{n_r,\eta}(z)+M_{s_r,\eta}(z).
 \tag{PD10}
\]
On the smaller domain
$|\operatorname{Re}z_4|+\sum_j|\operatorname{Im}z_j|<\eta/32$,
a polydisc of radius $\eta/128$ lies in $\mathcal O_\eta$.
Cauchy's inequalities bound the six derivatives by
$6!(128/\eta)^6$ times the bound in (PD4). The sum of coefficients
of $\Delta^3$ is $4^3$. Equations (PD8) therefore give
\[
 \|V_\eta(z)\|
 \le C\eta^{-p}(\eta^{-6}r+r^{-5})
 \le C\eta^{-(p+5)}.
 \tag{PD11}
\]
Every constant is independent of the surrounding polynomials and
their numbers of entries.

Let $k\in C_c^\infty$ have support in $|x_4|<\eta/64$.
For a continuous compact $\phi$ and a compact $C^1$ kernel $q$,
Riemann sums for $\phi*q=\int\phi(x)\tau_xq\,dx$ converge in
$C^1$, with common compact support. Local $C^1$ continuity of the
insertion form consequently implies
$I(F;\phi*q;G)=\int\phi(x)I(F;\tau_xq;G)\,dx$.
All convolution supports here stay strictly between the two outside
time supports. Applying this observation to (PD9), then integrating
six derivatives off $k$ onto the analytic matrix element in (PD4),
proves
\[
 I(F;k;G)=\int k(x)\langle\Phi,V_\eta(x)\Psi\rangle\,dx.
 \tag{PD12}
\]
The integrations by parts have no boundary terms, since $k$ is
compactly supported in the open strip. This proves the density
identity, not just an inequality.

Mollification carries (PD12) to compact $C^1$ tests in the strip:
the left side converges by local $C^1$ continuity and the right side
by $L^1$ convergence and (PD11). Thus it has the unique continuous
extension to every $L^1$ function on that strip, with
\[
 |I(F;k;G)|\le C\eta^{-(p+5)}\|k\|_{L^1}
                            \|F\|_{(\eta/4)}\|G\|_{(\eta/4)}.
 \tag{PD13}
\]
For an $L^1$ test this formula defines the extended insertion; no
pointwise field interpretation is required. The operator weak
integral $\int k(x)V_\eta(x)\,dx$ is well-defined by its bounded
sesquilinear forms and has the same norm bound. For compact $C^1$
tests it agrees with the regularized field product in (PD3).
If two choices of $\eta$ admit the same outside polynomials and
test strip, their forms agree first on compact smooth tests and then
on the common $L^1$ domain by density. Thus the extension is
compatible on intersecting pulled-down domains.

Different compact parametrices give the same density on the real
strip: (PD12) makes their continuous matrix-element differences zero
against every compact smooth test. They therefore agree pointwise.
Their holomorphic continuations agree on connected overlapping
domains, by uniqueness from an open set of real points, first in
scalar matrix elements and then as bounded operators.

\subsection*{Plates and compatible point-plate operators}

Split $x=(t,u,x_\perp)$, with $x_\perp\in\mathbb R^2$.
For $\varphi\in L^1(\mathbb R^2)$ set
\[
 V_{\varphi,\eta}(t,u)
   =\int\varphi(x_\perp)V_\eta(t,u,x_\perp)\,dx_\perp.
 \tag{PD14}
\]
The weak integral is a bounded operator, of norm at most
$C\eta^{-(p+5)}\|\varphi\|_1$. It is operator-norm holomorphic for
$|\operatorname{Re}t|+|\operatorname{Im}u|<\eta/32$.
Indeed the integrands are uniformly bounded and holomorphic on
smaller neighbourhoods, so their locally uniformly convergent Taylor
series may be integrated. For an $L^1$ plane profile $\alpha$ supported
in $|t|<\eta/64$, Tonelli and (PD12)--(PD13) give
\[
 |I(F;\alpha\otimes\varphi;G)|
 \le C\|\varphi\|_1\eta^{-\wp}\|\alpha\|_1
                 \|F\|_{(\eta/4)}\|G\|_{(\eta/4)},
 \qquad \wp=p+5.
 \tag{PD15}
\]
Profiles may extend arbitrarily far in the spatial coordinate $u$.
For smooth compact profiles and labels this is the original insertion;
all other stated $L^1$ cases are its continuous extension.

For $0<a\le1/4$ define the regularized point-plate operator
\[
 B_{\varphi;a}=V_{\varphi,4a}(0,0),\qquad
 \|B_{\varphi;a}\|\le C a^{-\wp}\|\varphi\|_1.
 \tag{PD16}
\]
It is the operator-norm limit of the regularized insertions of
$\alpha_\epsilon\otimes\varphi$, where $\alpha_\epsilon$ is a
unit-integral smooth plane approximate identity at zero. This follows
from the operator-norm continuity in (PD14). For $0<a<b\le1/4$,
\[
 B_{\varphi;b}=e^{-(b-a)H}B_{\varphi;a}e^{-(b-a)H}.
 \tag{PD17}
\]
To prove it, first use a compact smooth plate supported in the central
strips of both scales. The vector formula (PD3), or its dense
matrix-element identity, proves the relation for its two regularized
operators. Approximate the plane point by $\alpha_\epsilon$ and pass
in operator norm. Finally approximate an $L^1$ label by smooth compact
labels using (PD15). This proves compatibility on all the stated
domains rather than assuming an unbounded field common to the scales.

The same calculation gives real translation compatibility. For
$0<a'<a\le1/4$ and $|t|<a/16$,
\[
 V_{\varphi,4a}(t,u)=
 e^{-(a-a'+t)H}U(ue_1)B_{\varphi;a'}
 U(-ue_1)e^{-(a-a'-t)H}
 \tag{PD18}
\]
whenever $a-a'>|t|$. One proves this first for a sufficiently narrow
smooth approximate identity, using the real translation calculation
of (PD4), and then passes in operator norm. The right side also
supplies the corresponding holomorphic continuation. These identities
are suitable for composing independently regularized plate insertions.

\begin{theorem}[Analytic density and plate bound]
\label{thm:plate-density-from-insertion}
Under (PD1)--(PD2) and the stated local continuity, the insertion has
the analytic operator density (PD10)--(PD12), the uniform $L^1$ bound
(PD13), and the plate bound (PD15) with exponent $\wp=p+5$.
The point-plate operators (PD16) obey the compatibility and translation
identities (PD17)--(PD18). All constants are independent of the numbers
of outside fields. No additional typed source estimate is used.
\end{theorem}

As a useful spectral consequence, suppose $\varphi$ is real. Reflection
of real, time-even approximate identities proves
$B_{\varphi;a}^*=B_{\varphi;a}$.
Let $v_{\varphi,a}=B_{\varphi;a}\Omega$ and let $\mu_{\varphi,a}$
be its finite positive joint spectral measure for $(H,P_1)$.
Since $H\Omega=0$, (PD17) gives
$\mu_{\varphi,b}=e^{-2(b-a)E}\mu_{\varphi,a}$.
The positive measure $\mu_\varphi=e^{2aE}\mu_{\varphi,a}$ is thus
independent of $a$, is finite on bounded-energy sets, and may have
infinite total mass. For $t>0$ choose $0<a<\min(t/2,1/4)$ and put
\[
 K_\varphi(t,u)=
 \langle v_{\varphi,a},e^{-(t-2a)H+iuP_1}v_{\varphi,a}\rangle
 =\int e^{-tE+iuq}\,d\mu_\varphi(E,q).
 \tag{PD19}
\]
It is independent of $a$, finite, and its spectral measure is supported
on $E\ge|q|$. The integral is holomorphic on
$\operatorname{Re}t>|\operatorname{Im}u|$: on smaller neighbourhoods
all differentiated integrands are bounded by a polynomial in $E$
times a strictly decaying exponential, integrable by (PD19) at a
smaller positive real time.

To identify this as the original two-plate kernel, take real plane
approximate identities even in time, first with a smooth compact real
label $\varphi$. The vectors
$[(\tau_{ae_4}\alpha_\epsilon)\otimes\varphi]$
converge to $v_{\varphi,a}$ by (PD16). Their matrix elements in
(PD19) are precisely the two-plate moments at time separation $t$,
with the narrow profiles, by reflection and common translation.
For compact profiles with a fixed positive time gap, choose $a$
smaller than half that gap. Convolution approximates those profiles
in $C^1$, while the approximating matrix elements are uniformly
bounded by the two vector norm bounds. Fubini and dominated
convergence therefore recover the integral of $K_\varphi$ against
the original profiles. Smooth compact approximation of an $L^1$
label uses the uniform bounds above. The kernel here is the full
two-point moment kernel; for connected moments one applies the
construction to the centered functional.
This is a statement in the chosen coordinate frame. No diagonal-frame cone,
contact value at $t=0$, or rotational covariance is asserted here.

% END MODULE chapters/plate_density_from_insertion_v5.tex


% BEGIN MODULE chapters/angular_kernel_construction_v5.tex
\section{Analytic planar kernels and the finite angular band}
\label{ch:angular-kernel-construction}

This section constructs the analytic objects used by the angular
argument. Its quantitative input is a plate estimate uniform in the
numbers of outside fields. It proves neither that estimate nor the
short-distance two-point estimate. In particular, the conclusions
below do not by themselves establish rotation invariance of the
Euclidean functional.

\subsection*{The assumptions and the meaning of a product}

Write a point of Euclidean space as $(y,z)$, where
$y=(y^0,y^1)=(x_4,x_1)$ and $z=(x_2,x_3)$. Put
$\theta(y^0,y^1)=(-y^0,y^1)$. A tuple is a finite ordered list of
compactly supported $C^1$ test functions with pairwise disjoint
supports. Products of tuples mean concatenation, whenever the
concatenation still has disjoint supports. We use finite complex
linear combinations of such tuples, including the empty tuple $1$.
The notation $(h)$ denotes the tuple with its single entry $h$.

Assume the following properties of a centered functional $\sigma$.
It is multilinear in the entries, symmetric, real on real entries,
and invariant under simultaneous real translations and the
hypercubic group. Its values satisfy $\sigma(1)=1$ and
$\sigma((h))=0$. On every fixed collection of separated compact
support sets it is continuous for the $C^1$ norms, with a bound
independent of the positions of those sets when their diameters and
mutual distances are fixed. No bound on the growth of these constants
with the number of entries is assumed in this section.

Let $\mathscr A_+$ denote the combinations of tuples supported
strictly in $\{x_4>0\}$. Reflection positivity and the reconstruction
give a Hilbert space $\mathcal H$, a dense class map
$F\mapsto[F]$, and $\Omega=[1]$, such that
\[
 \langle[F],[G]\rangle=\sigma(\theta\overline F\cdot G),
 \qquad \|\Omega\|=1.
 \tag{AK1}\label{ak:hilbert}
\]
The inner product is antilinear in the first entry. The reconstructed
time and planar spatial translations have the actions
\[
 e^{-tH}[F]=[\tau_{te_4}F],\qquad
 e^{ixP}[F]=[\tau_{xe_1}F],\qquad t\ge0,
 \tag{AK2}\label{ak:translations}
\]
where $H,P$ are strongly commuting self-adjoint operators,
$H\ge0$, and
\[
 H\ge |P|,\qquad H\Omega=P\Omega=0.
 \tag{AK3}\label{ak:cone}
\]
Only this one planar consequence of the exact cube cone is used.

For $F$ supported above a height $a$ define the pulled-down norm by
\[
 \|F\|_{(a)}=\|[\tau_{-ae_4}F]\|.
\]
A label is a real $\varphi\in C_c^1(\mathbb R^2)$ in the transverse
variables. A profile is a complex $\alpha\in C_c^1(\mathbb R^2)$.
The corresponding plate is $\alpha\otimes\varphi$.
Fix constants $0<c<1/4$ and $\wp>0$. The plate input is that, for
each label, a constant $C_\varphi<\infty$ satisfies
\[
 \begin{split}
 &|\sigma(\theta\overline F\cdot
                    (\alpha\otimes\varphi)\cdot G)|\\
 &\hspace{8mm}\le C_\varphi\eta^{-\wp}\|\alpha\|_1
                 \|F\|_{(\eta/4)}\|G\|_{(\eta/4)}
 \end{split}
 \tag{AK4}\label{ak:plate-input}
\]
for $0<\eta\le1$, profiles supported in $|y^0|<c\eta$, and
$F,G$ supported strictly above $\eta/4$. This includes arbitrary
finite combinations and arbitrary numbers of entries in $F,G$.
The constant $C_\varphi$ is the same for all these choices and all
$\eta$. Every concatenation in this inequality has separated
supports. The plate theorem obtained from the ordinary-source
insertion estimate supplies precisely this input, with a fixed small
$c$ and $\wp=p+5$.

\subsection*{Bounded plates and their padding}

For a profile in the strip of \eqref{ak:plate-input} put $u=\eta/4$.
On the dense classes define
\[
 B_\varphi[\alpha;\eta][G]
   =\big[(\tau_{ue_4}(\alpha\otimes\varphi))
                          \cdot\tau_{2ue_4}G\big].
 \tag{AK5}\label{ak:plate-operator}
\]
The plate on the right lies at times strictly between $0$ and
$2u$; the other entries lie strictly above $2u$. Thus the right
side is an actual class in $\mathcal H$.

To prove that it depends only on $[G]$, test it against $[F]$ and
translate the full tuple down by $u$. The resulting expression is
\[
 \langle[F],B_\varphi[\alpha;\eta][G]\rangle
 =\sigma\big(\theta\overline{\tau_{ue_4}F}\cdot
          (\alpha\otimes\varphi)\cdot\tau_{ue_4}G\big).
\]
The two pulled-down norms on the right are $\|[F]\|$ and
$\|[G]\|$. By \eqref{ak:plate-input}, density and Hilbert duality,
\[
 \|B_\varphi[\alpha;\eta][G]\|
 \le C_\varphi\eta^{-\wp}\|\alpha\|_1\|[G]\|.
 \tag{AK6}\label{ak:operator-bound}
\]
In particular null classes map to null classes. There is a unique
bounded extension to the whole Hilbert space, with the same norm
bound. No domain containing an unspecified number of fields remains
in this statement.

The definitions also fix all padding conventions. If $\eta'\le\eta$
and the profile is supported in $|y^0|<c\eta'$, then, with
$d=(\eta-\eta')/4$,
\[
 B_\varphi[\alpha;\eta]
   =e^{-dH}B_\varphi[\alpha;\eta']e^{-dH}.
 \tag{AK7}\label{ak:padding}
\]
Indeed both sides applied to $[G]$ give the right side of
\eqref{ak:plate-operator}; density proves equality of the bounded
operators. Realness, symmetry, reflection invariance and a
translation by $2u$ similarly give
\[
 B_\varphi[\alpha;\eta]^*
       =B_\varphi[\overline\alpha\circ\theta;\eta].
 \tag{AK8}\label{ak:plate-adjoint}
\]
For example, conjugate the expression defining
$\langle[G],B_\varphi[\alpha;\eta][F]\rangle$, reflect all its
entries, translate by $2u$, and permute the entries. The resulting
expression is
$\langle[F],B_\varphi[\overline\alpha\circ\theta;\eta][G]\rangle$.
The label remains unchanged because it is real and the reflection
acts only on the plane.

\subsection*{The tube and its bounded translations}

For $v=(v^0,v^1)\in\mathbb C^2$ use the linear coordinates
$v_\pm=v^0\pm iv^1$, and set
\[
 \mathbb V=\{\operatorname{Re}v_+>0,
                   \operatorname{Re}v_->0\}
       =\{\operatorname{Re}v^0>|\operatorname{Im}v^1|\},
 \qquad v^*=(\overline{v^0},-\overline{v^1}).
 \tag{AK9}\label{ak:tube}
\]
This is an open convex cone. Its real points are the half-plane
$\mathbb H=\{v^0>0\}$. Both entrywise conjugation and the star
preserve it, and $\operatorname{Re}v\in\mathbb V$ whenever
$v\in\mathbb V$. These assertions follow directly from the two
strict real linear inequalities in \eqref{ak:tube}. Notice the
distinct identities
\[
 (v^*)_\pm=\overline{v_\pm},\qquad
 v^*=-\theta v\quad\hbox{for real }v.
\]

The joint spectral measure of $(H,P)$ is carried by
$\{(E,q):E\ge |q|\}$. Define on the closed tube
\[
 T(v)=\int e^{-v^0E+iv^1q}\,d\mathsf P(E,q).
 \tag{AK10}\label{ak:T}
\]
Writing $E_\pm=(E\pm q)/2\ge0$ gives
$-v^0E+iv^1q=-v_+E_--v_-E_+$. Consequently
\[
 \|T(v)\|\le1,\quad T(v)T(w)=T(v+w),\quad
 T(v)^*=T(v^*),\quad T(v)\Omega=\Omega.
 \tag{AK11}\label{ak:T-properties}
\]
The last assertion uses $H\Omega=P\Omega=0$. Dominated convergence
in the spectral measure of each vector proves strong continuity on
the closed tube.

The map $T$ is holomorphic in operator norm on $\mathbb V$.
Here is a direct proof, which also records why the exact cone is
appropriate. At a point with $\operatorname{Re}v_\pm\ge r>0$,
expand in increments $\Delta_\pm$. The coefficient of
$\Delta_+^j\Delta_-^k$ is the bounded functional-calculus operator
with multiplier
\[
 \frac{(-E_-)^j(-E_+)^k}{j!k!}
                    e^{-v_+E_--v_-E_+}.
\]
Its norm is at most $r^{-j-k}$, since
$x^je^{-x}/j!\le1$ for $x\ge0$. The resulting two power series
converge in operator norm when $|\Delta_\pm|<r$. For real $v$ with
$v^0\ge0$, the reconstruction gives
\[
 T(v)[F]=[\tau_{v^0e_4+v^1e_1}F].
 \tag{AK12}\label{ak:T-real}
\]

Let $R_z$ be the complex rotation with the usual matrix
\[
 R_z=\begin{pmatrix}\cos z&-\sin z\\ \sin z&\cos z\end{pmatrix}.
\]
In light coordinates $(R_zv)_\pm=e^{\pm iz}v_\pm$. In particular
\[
 R_{is}\mathbb V=\mathbb V,\qquad
 (R_{is}v)^*=R_{is}(v^*)\quad(s\in\mathbb R).
 \tag{AK13}\label{ak:boost-star}
\]
These are algebraic identities: the multipliers $e^{\mp s}$ are
positive real numbers. For a tuple $v=(v_i)$ in the tube define
\[
 d_{\mathbb V}(v)=
       \min\{1,\min_{i,\pm}\operatorname{Re}(v_i)_\pm\}.
\]
Then its exact transformation and a useful lower bound are
\[
 d_{\mathbb V}(R_{is}v)
 =\min\{1,\min_i e^{-s}\operatorname{Re}(v_i)_+,
                   \min_i e^s\operatorname{Re}(v_i)_-\}
 \ge e^{-|s|}d_{\mathbb V}(v).
 \tag{AK14}\label{ak:margin}
\]
The same assertions hold in every quarter-turn frame
$\mathbb V_\nu=R_\nu\mathbb V$, $\nu\in(\pi/2)\mathbb Z$.
The rotations commute; $R_\nu$ is real; and the real points of
$\mathbb V_\nu$ are $\{e_\nu\cdot v>0\}$, where
$e_\nu=R_\nu(1,0)$.

We will use the following elementary uniqueness principle. A
holomorphic function on a connected open subset of $\mathbb C^N$
which vanishes on a nonempty relatively open subset of $\mathbb R^N$
vanishes identically. At an interior real point all real derivatives
are zero, so its complex Taylor series is zero; overlapping balls
then propagate the equality along a path. The same proof applies to
Hilbert-valued holomorphic functions after testing against each
vector. An antiholomorphic variable is treated by conjugating that
variable first. This last use will always take place on domains
stable under conjugation.

\subsection*{Products with smooth plane profiles}

For $n\ge1$ define the open convex domain
\[
 \mathbb D^+_{n,\eta}=
 \{\mathbf x:x_1-\eta e_0/4\in\mathbb V,
       \ x_{i+1}-x_i-\eta e_0/2\in\mathbb V\ (i<n)\}.
 \tag{AK15}\label{ak:Deta}
\]
For profiles $\alpha_i$ in the strip $|y^0|<c\eta$ and labels
$\varphi_i$, abbreviate $B_i=B_{\varphi_i}[\alpha_i;\eta]$ and put
\[
 \begin{split}
 \Psi_\eta[\vec\alpha](\mathbf x)
   ={}&T(x_1-\eta e_0/4)B_1
        T(x_2-x_1-\eta e_0/2)B_2\cdots\\
      &\hspace{24mm}T(x_n-x_{n-1}-\eta e_0/2)B_n\Omega.
 \end{split}
 \tag{AK16}\label{ak:smooth-vector}
\]
It is Hilbert-valued holomorphic, and
\[
 \|\Psi_\eta[\vec\alpha](\mathbf x)\|
       \le\prod_{i=1}^n C_{\varphi_i}\eta^{-\wp}\|\alpha_i\|_1.
 \tag{AK17}\label{ak:smooth-vector-bound}
\]
The definitions remain valid on the closed domain using strong
continuity of $T$. When all $x_i$ are real and the displayed
non-strict time inequalities hold, direct substitution gives
\[
 \Psi_\eta[\vec\alpha](\mathbf x)
       =\Big[\prod_{i=1}^n\tau_{x_i}
                                (\alpha_i\otimes\varphi_i)\Big].
 \tag{AK18}\label{ak:real-vector}
\]
For clarity this is an induction from the right. The rightmost
operator first puts its plate at time $\eta/4$. The preceding
translation moves it by $x_n-x_{n-1}-\eta e_0/2$; the next plate
operator translates that entry forward by $\eta e_0/2$ and inserts
its own plate at $\eta e_0/4$. The two centers are therefore
$\eta e_0/4$ and $\eta e_0/4+x_n-x_{n-1}$. Iteration and the
leftmost translation put all centers at the $x_i$. Spatial shifts
follow the same addition rule. Time gaps and $c<1/4$ ensure that
every tuple in this induction is separated.

Define the smoothed scalar kernel on
$U_\eta=(\eta e_0/2+\mathbb V)^{n-1}$ by
\[
 K_{n,\eta}[\vec\alpha](v)
   =\langle\Omega,B_1T(v_1-\eta e_0/2)B_2\cdots
                   T(v_{n-1}-\eta e_0/2)B_n\Omega\rangle.
 \tag{AK19}\label{ak:smooth-kernel}
\]
It is holomorphic and bounded by the right side of
\eqref{ak:smooth-vector-bound}. If the $y_i$ are real and
$y_{i+1}^0-y_i^0\ge\eta/2$, then
\[
 \sigma\Big(\big(\tau_{y_i}
                      (\alpha_i\otimes\varphi_i)\big)_{i=1}^n\Big)
       =K_{n,\eta}[\vec\alpha](v(\mathbf y)),\qquad
 v(\mathbf y)_i=y_{i+1}-y_i.
 \tag{AK20}\label{ak:real-smooth-kernel}
\]
Translate the first center to $(\eta/4,0)$ and use
\eqref{ak:real-vector}; taking its scalar product with $\Omega$
proves the formula. Translation invariance then restores the
original centers. Set $K_0=1$ and $K_1=0$, consistently with the
empty and centered one-point moments.

\subsection*{Removing the plane smoothing}

Fix $\eta$ and choose real radial nonnegative smooth bumps
$\beta_{i,j}$, each of mass one, whose support radii tend to zero.
For all large $j$ the bumps fit in $D(0,c\eta)$.
The kernels $K_{n,\eta}[\vec\beta_j]$ have a common bound
\[
 C_\eta=\prod_i C_{\varphi_i}\eta^{-\wp n}
 \quad\hbox{on }U_\eta.
 \tag{AK21}\label{ak:normal-bound}
\]
The bumps need not have the same radii or shapes in different
entries. This freedom is needed later for mixed inner products.

Here are the details identifying every limit. Let $b_i$ be $C^1$
profiles whose successive support time gaps exceed $\eta$. For all
large $j$, their convolutions with $\beta_{i,j}$ have separated
supports in fixed compact sets. The identity
\[
 b_i*\beta_{i,j}=\int b_i(y)\tau_y\beta_{i,j}\,dy
\]
holds also as a $C^1$ Riemann integral. Indeed, for fixed $j$,
the integrand and its first derivatives in the test variable are
uniformly continuous on a common compact set. Riemann sums converge
there in $C^1$, with their supports in the same compact
neighborhood. Multilinearity, the assumed separated $C^1$
continuity, and successive Riemann summation in all entries give
\[
 \begin{split}
 \sigma\big(((b_i*\beta_{i,j})\otimes\varphi_i)_i\big)
   =\int\prod_i b_i(y_i)\,
       K_{n,\eta}[\vec\beta_j](v(\mathbf y))\,d\mathbf y.
 \end{split}
 \tag{AK22}\label{ak:convolution-identity}
\]
The integrands are continuous and compactly supported; thus Fubini
and the successive integrals are legitimate.

A locally bounded family of holomorphic functions has locally
uniformly convergent subsequences here. One can see this directly:
Cauchy estimates on smaller polydiscs bound first derivatives,
Arzel\`a--Ascoli gives convergence there, and a diagonal choice on
a countable exhaustion gives convergence on all compact subsets.
The Cauchy integral formula shows that the limit is holomorphic.
Apply this observation to \eqref{ak:normal-bound}. Along a
subsequence let the limit be $K'$. Since
$b_i*\beta_{i,j}\to b_i$ in $C^1$, passage in
\eqref{ak:convolution-identity} gives
\[
 \sigma\big((b_i\otimes\varphi_i)_i\big)
    =\int\prod_i b_i(y_i)K'(v(\mathbf y))\,d\mathbf y.
 \tag{AK23}\label{ak:eta-density}
\]
All differences in this integral lie in a compact subset of
$U_\eta$, so locally uniform convergence suffices.

There is at most one holomorphic $K'$ with this property. If two
continuous densities give the displayed integrals, choose the
$b_i$ to be unit-mass bumps in successively smaller discs around
any fixed real configuration with gaps $>\eta$. Their integrals
converge to the values of the two densities there. These values
coincide. Such differences contain the nonempty relatively open
set $\{v\hbox{ real}:v_i^0>\eta\}$ in the connected set $U_\eta$;
the uniqueness principle proves equality on $U_\eta$.
Consequently the entire smoothing family converges, not merely a
subsequence. Otherwise a sequence staying away on a compact set
would have a further subsequence with the same uniquely determined
limit. Denote it by $K_{n,\eta}$.

The argument proves a useful strengthened conclusion: for any
choice of unit-mass bumps in the different entries with maximal
support radius tending to zero, the corresponding kernels converge
to the same $K_{n,\eta}$, locally uniformly on $U_\eta$. No relation
between their rates of shrinkage is required.

If $\eta'<\eta$, the two limiting kernels agree on $U_\eta$.
Both satisfy \eqref{ak:eta-density} for gaps $>\eta$, and the same
uniqueness proof applies. Since $\bigcup_{\eta>0}U_\eta$
equals $\mathbb V^{n-1}$, this constructs a unique holomorphic
function
\[
 K_n(v;\vec\varphi)\quad\hbox{on }\mathbb V^{n-1}.
\]
It represents $\sigma$ on every time-ordered family of profiles:
\[
 \sigma\big((\alpha_i\otimes\varphi_i)_i\big)
   =\int\prod_i\alpha_i(y_i)
                       K_n(v(\mathbf y);\vec\varphi)\,d\mathbf y.
 \tag{AK24}\label{ak:unsmeared-density}
\]
Choose $\eta$ below all the support gaps to obtain this formula.
It also proves reality at every real point of the tube, by testing
against real profiles concentrated near the chosen points.

The bound retains exactly one factor per field:
\[
 |K_n(v;\vec\varphi)|
      \le\Big(\prod_i C_{\varphi_i}\Big)
                              d_{\mathbb V}(v)^{-\wp n}.
 \tag{AK25}\label{ak:kernel-bound}
\]
Indeed $v\in U_\eta$ whenever
$\eta<\min\{1,2\min_{i,\pm}\operatorname{Re}(v_i)_\pm\}$.
Take the limit up to this upper endpoint in
\eqref{ak:normal-bound}; that endpoint is at least
$d_{\mathbb V}(v)$. Neither a derivative of a shrinking bump nor
an estimate whose constant depends on outside field counts has
entered this step.

\subsection*{Real overlap identities and the four frames}

For a quarter turn $\nu$ define
$K_n^{(\nu)}(v;\vec\varphi)=K_n(R_{-\nu}v;\vec\varphi)$ on
$\mathbb V_\nu^{n-1}$. Hypercubic invariance and the change of
variables $y_i=R_\nu y_i'$ show that
\eqref{ak:unsmeared-density} holds with $K_n^{(\nu)}$ whenever
the profiles are ordered in the direction $e_\nu$. The rotations
act only on profiles and leave the labels unchanged.

It is useful to account for permutations as well. For a permutation
$\pi$ of $\{1,\ldots,n\}$ let
\[
 \mathcal U_{\nu,\pi}=
 \{\mathbf z\in(\mathbb C^2)^n:
       R_{-\nu}(z_{\pi(i+1)}-z_{\pi(i)})\in\mathbb V\ (i<n)\}.
 \tag{AK26}\label{ak:ordered-domain}
\]
On it use $K_n$ with these ordered differences and the correspondingly
permuted labels. Each domain is open, convex, nonempty, invariant
under simultaneous complex translations, and satisfies
$\operatorname{Re}\mathbf z\in\mathcal U_{\nu,\pi}$ whenever
$\mathbf z\in\mathcal U_{\nu,\pi}$.

Two of these kernels agree on every nonempty intersection. The
intersection is convex and hence connected. It contains the real
part of each of its points, so contains a relatively open real
ball. At any real configuration in the intersection take small
discs about its points which preserve both strict orders.
The moment of profiles supported in these discs has both continuous
kernel densities, by \eqref{ak:unsmeared-density}, the frame
identity and symmetry of $\sigma$. Testing with unit-mass profiles
shrinking to the points shows equality of the two densities at that
configuration. The holomorphic uniqueness principle now gives
equality on the whole complex intersection. This proves the
overlap identity, rather than presupposing an arbitrary-angle
rotation symmetry.

Thus the ordered kernels glue to a single holomorphic function
$\mathbf K_n(\mathbf z;\vec\varphi)$ on
\[
 \mathcal U_n^{\rm ord}=\bigcup_{\nu,\pi}\mathcal U_{\nu,\pi}.
 \tag{AK27}\label{ak:ordinary-domain}
\]
This union is connected. To check this explicitly, a real
configuration can be assigned any prescribed strict time ordering
and, independently, any prescribed strict abscissa ordering. Hence
each $0$-frame permutation domain intersects every $\pi/2$-frame
permutation domain. These convex domains form a connected union.
The other two frames give the same domains with their orderings
reversed, so add no disconnected components. Its real points form
\[
 \operatorname{Ord}_n=
 \{\mathbf y:\hbox{all times are distinct or all abscissae
                       are distinct}\}.
\]
The function is translation invariant, symmetric in the labelled
points, real on real configurations, and invariant under the
quarter turns. These properties follow first from the corresponding
moment densities on real open sets and then from uniqueness.
Imaginary rotations preserve each domain
$\mathcal U_{\nu,\pi}$ and hence preserve their union.

\subsection*{Point product vectors and their Gram kernel}

Put
\[
 \mathbb D_n^+=\{\mathbf x:x_1\in\mathbb V,
                           \ x_{i+1}-x_i\in\mathbb V\ (i<n)\}.
 \tag{AK28}\label{ak:Dplus}
\]
This is the union of the nested $\mathbb D^+_{n,\eta}$ as
$\eta$ decreases to zero. For $n=0$ there is one empty
configuration and its vector will be $\Omega$.

We prove strong convergence of the smooth product vectors, rather
than analytically continuing an inequality. Fix $\eta$, choose a
real radial unit-mass bump $\beta_\epsilon$ supported in
$D(0,\epsilon)$, and take $\epsilon<c\eta$. Use it in every
entry of \eqref{ak:smooth-vector}, writing the result as
$\Psi_{\eta,\epsilon}(\mathbf x;\vec\varphi)$.

If $\mathbf x\in\mathbb D^+_{n,\eta}$ and
$\mathbf y\in\mathbb D^+_{m,\eta}$, define the list of differences
\[
 J(\mathbf x,\mathbf y)=
 \big((x_n-x_{n-1})^*,\ldots,(x_2-x_1)^*,
             x_1^*+y_1,y_2-y_1,\ldots,y_m-y_{m-1}\big).
 \tag{AK29}\label{ak:J}
\]
If $n=0$, omit the left list and the central entry. If $m=0$,
omit the central entry and the right list. The labels of this list
are $(\varphi_n,\ldots,\varphi_1,\psi_1,\ldots,\psi_m)$.
Every entry belongs to $\eta e_0/2+\mathbb V$: for the central
one subtracting $\eta e_0/2$ gives the sum
$(x_1-\eta e_0/4)^*+(y_1-\eta e_0/4)$ in the open cone.
The other entries follow from star invariance.

For different radii $\epsilon,\epsilon'$ the exact identity is
\[
 \begin{split}
 &\langle\Psi_{\eta,\epsilon}(\mathbf x;\vec\varphi),
                  \Psi_{\eta,\epsilon'}(\mathbf y;\vec\psi)\rangle\\
 &\quad=K_{n+m,\eta}[
       \underbrace{\beta_\epsilon,\ldots,\beta_\epsilon}_{n},
       \underbrace{\beta_{\epsilon'},\ldots,\beta_{\epsilon'}}_{m}]
                    (J(\mathbf x,\mathbf y)),
 \end{split}
 \tag{AK30}\label{ak:mixed-inner-product}
\]
with the reversed left labels specified above. For real positions
this is just \eqref{ak:real-vector}, the Hilbert inner product,
symmetry, and \eqref{ak:real-smooth-kernel}. Radiality and realness
are important: $\theta\overline{\tau_x\beta_\epsilon}
=\tau_{\theta x}\beta_\epsilon$ for real $x$. The left side
is antiholomorphic in $\mathbf x$ and holomorphic in $\mathbf y$;
the right side has exactly the same properties because the star is
conjugate linear. The product domains are convex, conjugation
stable and have real interior points. Mixed holomorphic uniqueness
therefore proves \eqref{ak:mixed-inner-product} for complex
positions too. For total length one both sides are zero; for
length zero both equal one.

The strengthened smoothing conclusion after
\eqref{ak:eta-density} shows that the right side converges to
$K_{n+m}(J(\mathbf x,\mathbf y))$ when both radii tend to zero,
with no restriction on their relative rate. Convergence is uniform
on compact sets of the position domains. In particular, for a
compact subset of $\mathbb D^+_{n,\eta}$, expand
\[
 \|\Psi_{\eta,\epsilon}(\mathbf x)
              -\Psi_{\eta,\epsilon'}(\mathbf x)\|^2
\]
as its four mixed inner products. All four tend uniformly to the
same diagonal kernel value; their signed sum tends to zero.
This proves a locally uniform strong Cauchy property. The limit
$\Psi_\eta(\mathbf x;\vec\varphi)$ is Hilbert-valued holomorphic:
uniform convergence on smaller polydiscs passes the vector Cauchy
formula and its power series to the limit.

These vectors are independent of $\eta$ on overlaps. For any
fixed sufficiently small physical bump and real positions in an
overlap, both definitions give the same class in
\eqref{ak:real-vector}. Vector holomorphic uniqueness on the
connected overlap proves equality for complex positions, and then
the strong limits agree. We have therefore constructed
\[
 \Psi(\mathbf x;\vec\varphi)\in\mathcal H,
       \qquad\mathbf x\in\mathbb D_n^+,
 \tag{AK31}\label{ak:point-vector}
\]
holomorphically, with the bound
$\|\Psi(\mathbf x;\vec\varphi)\|
\le\prod_i C_{\varphi_i}\eta^{-\wp}$ for every admissible
$\eta$. On the domain where $v\in\overline{\mathbb V}$, the
semigroup rule also gives
\[
 \Psi(\mathbf x+v;\vec\varphi)=T(v)\Psi(\mathbf x;\vec\varphi).
 \tag{AK32}\label{ak:point-vector-translation}
\]
It holds before removing the bumps and hence afterward.

For labelled positive configurations define
\[
 \operatorname{Gr}(\mathsf x,\mathsf y)
   =K_{n+m}(J(\mathbf x,\mathbf y);
                       \varphi_n,\ldots,\varphi_1,\vec\psi).
\]
Use the empty and one-point conventions already specified. The
construction proves the exact equality
\[
 \operatorname{Gr}(\mathsf x,\mathsf y)
       =\langle\Psi(\mathsf x),\Psi(\mathsf y)\rangle.
 \tag{AK33}\label{ak:gram-vector}
\]
Thus this is a positive semidefinite kernel at complex positions,
for arbitrary finite families and different lengths, including the
empty configuration.

Finally $R_{is}\mathbb D_n^+=\mathbb D_n^+$, and
\eqref{ak:boost-star} gives
$J(R_{is}\mathbf x,R_{is}\mathbf y)=R_{is}J(\mathbf x,\mathbf y)$.
Consequently
\[
 \begin{split}
 \operatorname{Gr}_s(\mathsf x,\mathsf y)
   &:=K_{n+m}(R_{is}J(\mathbf x,\mathbf y);\hbox{the same labels})\\
   &=\langle\Psi(R_{is}\mathsf x),\Psi(R_{is}\mathsf y)\rangle
 \end{split}
 \tag{AK34}\label{ak:boost-gram}
\]
is positive semidefinite for every real $s$. This is an actual
Gram representation; positivity has not been inferred from analytic
continuation of a positive scalar function.

\subsection*{Narrow pencils and the entire angle function}

Fix $0<\phi_0<\pi/4$ and let
\[
 \operatorname{Sec}=\{v\in\mathbb R^2\setminus\{0\}:
                            |\arg(v^0+iv^1)|<\phi_0\}.
\]
For $n\ge2$, a pencil is a configuration with every consecutive
difference in this sector. The cases $n=0,1$ have their separate
constant kernels fixed above. Write $r_{\mathbf y}=\min_i|y_{i+1}-y_i|>0$.
For a pencil with differences $v_i$ and labels $\vec\varphi$ put,
on the vertical strip
\[
 S_\nu=\{z=t+is:|t-\nu|<\pi/2-\phi_0\},
       \qquad\nu\in(\pi/2)\mathbb Z,
\]
the function
\[
 a^{(\nu)}_{\mathbf y}(z)
     =K_n(R_{z-\nu}v;\vec\varphi).
 \tag{AK35}\label{ak:angle-patch}
\]
It is well defined and holomorphic. In fact, writing
$(v_i)_\pm=|v_i|e^{\pm i\phi_i}$ with $|\phi_i|<\phi_0$ gives
\[
 \operatorname{Re}(R_{t+is-\nu}v_i)_\pm
       =|v_i|e^{\mp s}\cos(\phi_i+t-\nu)>0.
 \tag{AK36}\label{ak:pencil-margin}
\]

Only neighboring strips intersect. Their intersection is exactly
\[
 S_\nu\cap S_{\nu+\pi/2}
   =\{t+is:\nu+\phi_0<t<\nu+\pi/2-\phi_0\},
\]
which is nonempty and connected. At a real angle there, all the
rotated differences lie strictly in the quadrant between the two
frame normals. Small discs about the rotated points are ordered in
both frames. The proved real overlap identity therefore identifies
the two values in \eqref{ak:angle-patch} on this real interval.
The one-variable identity theorem identifies the holomorphic
functions throughout the intersection.

The strip width $\pi-2\phi_0$ exceeds their spacing $\pi/2$.
They cover the complex plane and glue to a unique entire function
$a_{\mathbf y}(z;\vec\varphi)$. The construction immediately gives
\[
 a_{\mathbf y}(z+\pi/2)=a_{\mathbf y}(z).
 \tag{AK37}\label{ak:period}
\]
Indeed use the patch with center $\nu+\pi/2$ for the shifted
argument and the patch with center $\nu$ for the original argument;
the arguments of $K_n$ are identical. For real $t$ it equals
$\mathbf K_n(R_t\mathbf y;\vec\varphi)$, since a suitable frame
orders this rotated pencil. In particular it is real on the real
axis, and $a_{\mathbf y}(is)=K_n(R_{is}v;\vec\varphi)$.

For every real $t$ choose a nearest quarter turn $\nu$, so
$|t-\nu|\le\pi/4$. Set $c_0=\cos(\phi_0+\pi/4)>0$.
By \eqref{ak:pencil-margin}, every light-coordinate margin is at
least $c_0r_{\mathbf y}e^{-|s|}$. Hence
\[
 |a_{\mathbf y}(t+is)|
 \le\Big(\prod_iC_{\varphi_i}\Big)
       \min\{1,c_0r_{\mathbf y}e^{-|s|}\}^{-\wp n}
 \le A_{\mathbf y,\vec\varphi} e^{\wp n|s|},
 \tag{AK38}\label{ak:angle-type}
\]
where
$A_{\mathbf y,\vec\varphi}
=\prod_iC_{\varphi_i}\max\{1,(c_0r_{\mathbf y})^{-\wp n}\}$.
This proves the claimed exponential type with its linear dependence
on $n$; an uncontrolled prefactor depending on $n$ would not change
that conclusion.

\subsection*{The finite band and analytic harmonic coefficients}

Define
\[
 c_m(\mathbf y;\vec\varphi)=\frac2\pi
       \int_0^{\pi/2}a_{\mathbf y}(t;\vec\varphi)e^{-4imt}\,dt.
 \tag{AK39}\label{ak:Fourier}
\]
Cauchy's theorem on a rectangle of height $s$ and periodicity give
\[
 c_m=\frac2\pi\int_0^{\pi/2}
              a_{\mathbf y}(t+is)e^{-4im(t+is)}\,dt,
 \qquad |c_m|\le A_{\mathbf y,\vec\varphi}
                                      e^{4ms+\wp n|s|}.
\]
The vertical sides cancel exactly. Sending $s$ to $-\infty$ for
$4m>\wp n$, and to $+\infty$ for $4m<-\wp n$, proves
\[
 c_m=0\quad\hbox{if }4|m|>\wp n,\qquad
 a_{\mathbf y}(z)=\sum_{4|m|\le\wp n}c_m e^{4imz}.
 \tag{AK40}\label{ak:finite-band}
\]
For the equality, periodicity makes $a(z)$ a single-valued
holomorphic function of $w=e^{4iz}$ on $\mathbb C\setminus\{0\}$.
Its Laurent coefficients are the displayed $c_m$, so its Laurent
series is the finite sum after the vanishing just proved.
Equivalently one may use the uniformly convergent Fourier series
of the smooth periodic function on the real axis and then the
identity theorem. Reality implies $c_{-m}=\overline{c_m}$.

We next construct coefficients beyond pencils, without assuming
a continuity assumption across different ordering charts. Put
$N=\lfloor\wp n/4\rfloor$ and on planar configurations introduce
the angular vector field
\[
 \mathcal L=\sum_{i=1}^n
             (-y_i^1\partial_{y_i^0}+y_i^0\partial_{y_i^1}).
\]
On translation-invariant functions it is the same angular vector
field in the consecutive-difference coordinates. Let
\[
 P(X)=\prod_{m=-N}^N(X-4im),\qquad
 p_m(X)=\prod_{\substack{-N\le j\le N\\j\ne m}}
                         \frac{X-4ij}{4im-4ij}.
\]
The finite-band identity gives $P(\mathcal L)\mathbf K_n=0$
on the relatively open real set of pencils. The left side is
holomorphic on the connected domain $\mathcal U_n^{\rm ord}$.
Uniqueness therefore proves this equation everywhere on that
domain. Define there
\[
 c_m=p_m(\mathcal L)\mathbf K_n\quad(|m|\le N).
 \tag{AK41}\label{ak:projectors}
\]
On pencils these are the Fourier coefficients in
\eqref{ak:Fourier}: each polynomial $p_m$ selects its corresponding
distinct exponential. The polynomial identities
$\sum_m p_m=1$ and $(X-4im)p_m(X)=\text{constant}\,P(X)$ give
\[
 \sum_m c_m=\mathbf K_n,\qquad
                    \mathcal Lc_m=4imc_m.
 \tag{AK42}\label{ak:eigen-coefficients}
\]
The coefficients are holomorphic on the ordered domain, symmetric
in labelled points and invariant under simultaneous translations.
On its real points $c_{-m}=\overline{c_m}$: the angular vector
field has real coefficients, the density is real there, and
$\overline{p_m(X)}=p_{-m}(X)$ for real $X$.
For translation invariance, note that applying $\mathcal L$ to a
translation-invariant function again gives one: the extra term
under a common shift is a linear combination of the total
derivatives $\sum_i\partial_{y_i^\mu}$, which vanish.

Since imaginary rotations preserve the domain, the last equation
in \eqref{ak:eigen-coefficients} can be integrated along them:
\[
 c_m(R_{is}\mathbf z)=e^{-4ms}c_m(\mathbf z),\qquad
 \mathbf K_n(R_{is}\mathbf z)=\sum_m e^{-4ms}c_m(\mathbf z).
 \tag{AK43}\label{ak:boost-expansion}
\]
The same equation gives
$c_m(R_t\mathbf y)=e^{4imt}c_m(\mathbf y)$ on any real interval
whose rotated configurations remain in $\operatorname{Ord}_n$.

These coefficients extend uniquely and real-analytically to every
configuration of distinct real planar points. Here is the full
gluing argument. Let $\mathcal G_n$ consist of configurations with
distinct points and with no two nonzero difference vectors
perpendicular. This is open and dense: its complement, in addition
to the collisions, is a finite union of zero sets of nonzero real
polynomials. Each polynomial is nonzero, as can also be checked
by taking distinct points on one line. Every real rotation of a
configuration in $\mathcal G_n$ is ordinary. Failure of ordinariness
would produce a time tie and an abscissa tie, whose nonzero
difference vectors are perpendicular. Thus the preceding angular
ordinary differential equation proves covariance for every real
angle on $\mathcal G_n$.

For an arbitrary distinct configuration $\mathbf y$, choose an
angle $t$ at which all its rotated times are distinct. Only
finitely many angles are forbidden. In a neighborhood of
$\mathbf y$ define the extension by
\[
 \widetilde c_m(\mathbf u)
           =e^{-4imt}c_m(R_t\mathbf u).
 \tag{AK44}\label{ak:global-extension}
\]
This is real-analytic. Two choices agree on the dense set
$\mathcal G_n$ in every overlap, by the covariance already proved
there; continuity makes them agree everywhere in the overlap.
They consequently glue, agree with the original coefficients on
ordinary configurations, and obey exact covariance
\[
 \widetilde c_m(R_t\mathbf y)=e^{4imt}\widetilde c_m(\mathbf y)
 \quad(t\in\mathbb R).
 \tag{AK45}\label{ak:global-covariance}
\]
This last equality, permutation symmetry, translation invariance
and $\widetilde c_{-m}=\overline{\widetilde c_m}$ follow by
continuity from $\mathcal G_n$. Uniqueness follows because
$\operatorname{Ord}_n$ is dense in the collision-free set.
The collision-free set itself is connected: one may move its
points successively along polygonal paths avoiding the finitely
many positions occupied by the other points, with small preliminary
moves if necessary. Removing finitely many points does not
disconnect the plane.

Set $\overline m(0)=\overline m(1)=0$; for $n\ge2$ let
$\overline m(n)$ be the largest $|m|$ occurring for any pencil and
any real labels, or zero if none occurs. It is a finite integer and
\[
 \overline m(n)\le\wp n/4.
 \tag{AK46}\label{ak:top-frequency}
\]
Coefficients with $|m|>\overline m(n)$ vanish identically on the
ordered domain by uniqueness, and then on every collision-free
configuration by \eqref{ak:global-extension}. If
$\overline m(n)>0$, some pencil and labels have
$c_{-\overline m(n)}\ne0$, by reality.

\subsection*{Boost positivity without time ordering and the top kernel}

For arbitrary labelled configurations of $k$ and $l$ distinct
points in the real open half-plane, define
\[
 \widetilde{\operatorname{Gr}}_s(\mathsf x,\mathsf y)
  =\sum_m e^{-4ms}\widetilde c_m^{(k+l)}
       (\theta\mathbf x,\mathbf y;\vec\varphi,\vec\psi).
 \tag{AK47}\label{ak:global-gram}
\]
All points in this combined configuration are distinct: the two
groups lie on opposite sides of the mirror. Empty tuples have the
conventions $\widetilde c_m^{(0)}=\delta_{m0}$ and
$\widetilde c_m^{(1)}=0$.

This kernel is positive semidefinite for every real $s$.
First suppose each individual configuration has distinct times.
Sort those times. The combined reflected configuration is then
ordinary, and \eqref{ak:boost-expansion}, symmetry and
\eqref{ak:boost-gram} identify \eqref{ak:global-gram} with the
proved Gram kernel. For a finite family with time ties, perturb
the positions slightly inside the positive half-plane to remove
all such ties. The reflected combined configurations remain
collision free. Every coefficient is continuous there, so their
finite matrices converge entrywise to the desired matrix.
Positive semidefiniteness is closed under this limit. This proves
the assertion with time ties, arbitrary lengths and the empty
configuration included.

In particular, fix $k\ge1$, put $m_*=\overline m(2k)$, and set
$c_* =\widetilde c^{(2k)}_{-m_*}$. On $k$-point configurations
the kernel
\[
 (\mathsf x,\mathsf y)\longmapsto
                c_*(\theta\mathbf x,\mathbf y;\vec\varphi,\vec\psi)
 \tag{AK48}\label{ak:top-positive}
\]
is positive semidefinite: multiply the matrix in
\eqref{ak:global-gram} by $e^{-4m_*s}$ and send $s\to+\infty$.
All the other finitely many terms tend to zero.

There is also an exact formulation across an arbitrary line
$\mathfrak l$. Let $b\in\mathfrak l$ and let $R_\nu e_0$ be
its normal pointing into a chosen side $\mathfrak l_+$. The map
$r(y)=R_{-\nu}(y-b)$ sends that side to $\mathbb H$ and
intertwines its reflection $\theta_{\mathfrak l}$ with $\theta$.
Using translation invariance and \eqref{ak:global-covariance},
\eqref{ak:top-positive} becomes positivity of
\[
 (\mathsf x,\mathsf y)\longmapsto
       e^{4im_*\nu}c_*(\theta_{\mathfrak l}\mathbf x,
                                      \mathbf y;\vec\varphi,\vec\psi)
 \tag{AK49}\label{ak:top-all-lines}
\]
on $k$-point configurations in $\mathfrak l_+$. The compensating
unit phase is essential. This is a statement about the top
coefficient, and does not assume reflection positivity of the
original functional at every line.

\subsection*{The cut inequality, including the empty tuple}

The preceding constructions imply
\[
 2\overline m(k+l)\le\overline m(2k)+\overline m(2l)
                  \qquad(k,l\ge0).
 \tag{AK50}\label{ak:convexity}
\]
If the left top frequency is zero there is nothing to prove.
Otherwise choose a pencil and labels whose negative top coefficient
is nonzero. When $k,l\ge1$, translate the midpoint between
$y_k,y_{k+1}$ to the origin, and put
\[
 x_i=\theta y_{k+1-i}\ (1\le i\le k),\qquad
 z_j=y_{k+j}\ (1\le j\le l),
\]
with the correspondingly reversed left labels. Both are positive
configurations. Their first points and all consecutive differences
are in $\operatorname{Sec}$: reflection followed by a minus sign
changes the sign of a real spatial coordinate and preserves this
sector. The first points are the two reflected halves of the cut
difference.

The combined reflected configuration for $\mathsf x,\mathsf z$
is exactly the original pencil. The doubled configurations for
$\mathsf x,\mathsf x$ and $\mathsf z,\mathsf z$ are also pencils.
Their central differences are respectively $2x_1^0e_0$ and
$2z_1^0e_0$, and all other differences are sector vectors or their
stars. The Gram Cauchy--Schwarz inequality gives
\[
 |a_{\mathbf y}(is)|^2
   \le a_{(\theta x_k,\ldots,\theta x_1,x_1,\ldots,x_k)}(is)
       a_{(\theta z_l,\ldots,\theta z_1,z_1,\ldots,z_l)}(is).
\]
The left side has leading rate
$e^{8\overline m(k+l)s}$ as $s\to+\infty$, with nonzero
coefficient. Each right factor has rate at most
$e^{4\overline m(2k)s}$ or $e^{4\overline m(2l)s}$ by its finite
band. Comparison proves \eqref{ak:convexity}.

If $k=0$, translate a nonzero pencil so that its first point is in
$\operatorname{Sec}$. Apply Gram Cauchy--Schwarz to the empty
vector $\Omega$ and its product vector. The empty diagonal entry
is one, so the same comparison gives
$2\overline m(l)\le\overline m(2l)$. This is exactly the required
inequality with $\overline m(0)=0$. The case $l=0$ follows by
symmetry, and the centered one-point and fully empty cases are
immediate. Thus the initial convexity information has not been
lost by omitting the vacuum vector.

\subsection*{The two-point interface and its positive plane measure}

For two points, $\operatorname{Ord}_2$ is the entire real
collision-free set: two distinct points cannot have both coordinates
equal. Translation invariance therefore gives a real-analytic
function $G_{\varphi,\psi}(r)$ for every real $r\ne0$, equal to
$K_2(r;\varphi,\psi)$ when $r^0>0$. Its global coefficients obey
\[
 \begin{aligned}
 G_{\varphi,\psi}(r)&=\sum_m\widetilde c_m(r;\varphi,\psi),\\
 \widetilde c_m(R_t r)&=e^{4imt}\widetilde c_m(r),\\
 K_2(R_{is}r)&=\sum_m e^{-4ms}\widetilde c_m(r)
       \qquad(r^0>0).
 \end{aligned}
 \tag{AK51}\label{ak:two-point-interface}
\]
The last identity holds for every real $s$, not merely near zero;
it is \eqref{ak:boost-expansion}. Every sum is finite, with
$|m|\le\overline m(2)\le\wp/2$. The kernel in
\eqref{ak:global-gram} for one-point configurations is positive
semidefinite for every $s$. A quarter-turn frame with projection
at least $|r|/\sqrt2$ also gives the useful bound
\[
 |G_{\varphi,\psi}(r)|
 \le C_\varphi C_\psi
              \min\{1,|r|/\sqrt2\}^{-2\wp}.
 \tag{AK52}\label{ak:global-two-point-bound}
\]

The density formula for two points also holds for arbitrary
disjoint compact plane-profile supports, even if those whole
supports have no strict coordinate ordering. Let their distance
be $d>0$ and partition each profile into finitely many pieces
supported in sets of diameter less than $d/(10\sqrt2)$. For each
pair of pieces, two chosen points in their supports differ in at
least one coordinate by $d/\sqrt2$ or more. The small diameters
then make the two entire piece supports strictly ordered in that
coordinate, in one of the four frames. Apply
\eqref{ak:unsmeared-density} and the frame identity to this pair.
Summing the finitely many identities by bilinearity proves
\[
 \sigma(\alpha\otimes\varphi,\beta\otimes\psi)
  =\iint\alpha(y)\beta(y')G_{\varphi,\psi}(y'-y)\,dy\,dy'
\]
for every such pair of profiles. All partitions can be smooth;
their products with the original $C^1$ profiles remain $C^1$ and
have supports in the original separated sets.

There is a positive spectral plane measure for equal real labels.
Let $\zeta_s=\Psi((s,0);\varphi)$, $s>0$. By
\eqref{ak:point-vector-translation},
$\zeta_{s'}=e^{-(s'-s)H}\zeta_s$ if $s'>s$.
Let $\mu_s(B)=\langle\zeta_s,\mathsf P(B)\zeta_s\rangle$,
using the joint spectral measure of $(H,P)$. The locally finite
positive measure
\[
 d\mu_\varphi(E,q)=e^{2sE}\,d\mu_s(E,q)
 \tag{AK53}\label{ak:plane-measure}
\]
is independent of $s$, because the preceding vector identity
multiplies $\mu_s$ by $e^{-2(s'-s)E}$. It is carried by
$E\ge|q|$ and satisfies
\[
 K_2(v;\varphi,\varphi)
     =\int e^{-v^0E+iv^1q}\,d\mu_\varphi(E,q)
                       \qquad(v\in\mathbb V).
 \tag{AK54}\label{ak:Laplace-Fourier}
\]
For real $v=(t,x)$ choose $0<s<t/2$ and use
\eqref{ak:gram-vector} and \eqref{ak:point-vector-translation}:
the left side is
$\langle\zeta_s,T(t-2s,x)\zeta_s\rangle$, which is the
displayed integral. In particular
$\int e^{-tE}\,d\mu_\varphi<\infty$ for every $t>0$.
On a compact sub-tube the multiplier is bounded by $e^{-dE}$
for some $d>0$. Its derivatives are controlled by a smaller
exponential weight because $|q|\le E$. The integral is therefore
holomorphic there, and uniqueness extends its equality from real
points to all of $\mathbb V$.

No uniqueness of the vacuum or absence of a zero-energy mass is
asserted here. A positive plane measure in a diagonal-reflection
picture requires the reflection-positivity input for that picture;
it is not inferred from the coordinate cone. Likewise, eliminating
the nonzero two-point harmonics requires a separate ultraviolet
estimate and the diagonal-frame argument. The polygon and integer
arguments can use \eqref{ak:top-all-lines} and
\eqref{ak:convexity}, but are not contained in this section.
Finally, the analytic extension of coefficients to distinct planar
points does not identify the original moment functional on every
exceptional support arrangement. That distributional passage, in
particular arrangements with coincident planar projections and
separated transverse supports, is a separate obligation.

% END MODULE chapters/angular_kernel_construction_v5.tex


% BEGIN MODULE chapters/angular_two_point_elimination_v5.tex
\section{Ultraviolet elimination of the two-point angular harmonics}
\label{ch:angular-two-point-elimination}

This section proves an implication for the two-point function. Its inputs
are a continuous planar kernel with its axis-frame spectral realization,
the exact cone in that frame, one diagonal reflection, and a uniform
short-distance bound on a specified class of positive shapes. It constructs
the diagonal spectral measure and proves its cone support; a diagonal cone
and continuous rotation covariance are not assumed. The angular continuation
needed at order two is proved directly. In particular this argument does
not require an all-orders highest-harmonic theorem.

The kernel and Hilbert realization are supplied by the planar-kernel
construction under its insertion, continuity and reconstruction hypotheses.
We use these together with the interacting-source and short-distance
hypotheses stated below.

\subsection{The precise two-point interface}

Write a point of $\mathbb R^4$ as $(y,z)$, with
$y=(y^0,y^1)=(x_4,x_1)$ and $z=(x_2,x_3)$. On the plane let
$e_0=(1,0)$, let $R_t$ be the rotation
\[
 R_t(u,v)=(u\cos t-v\sin t,u\sin t+v\cos t),
 \qquad \theta(u,v)=(-u,v).
\]
The same notation denotes these maps on $\mathbb R^4$, leaving $z$ fixed.
Put $n=R_{\pi/4}e_0$. The reflection at $n^\perp$ is
$\theta_n=R_{\pi/4}\theta R_{-\pi/4}$.

Let $S_2$ be a real symmetric bilinear form on pairs of real
$C_c^1(\mathbb R^4)$ tests with disjoint supports, complexified bilinearly.
It is translation- and $W_4$-invariant and satisfies coordinate reflection
positivity. It also satisfies
\[
 S_2(\theta_n\overline f,f)\ge0
       \quad\bigl(f\in C_c^1(\{n\cdot y>0\};\mathbb C)\bigr).
 \tag{TU1}
\]
Positivity at any one diagonal normal in four dimensions implies this
choice by hypercubic conjugation. All scalar products below are conjugate
linear in the first variable.

The kernel interface is the following explicit assumption. For real
$\varphi,\psi\in C_c^1(\mathbb R^2)$ there are continuous functions
$G_{\varphi,\psi}$ on $\mathbb R^2\setminus\{0\}$ such that
\[
 S_2(\alpha\otimes\varphi,\beta\otimes\psi)
 =\iint\alpha(y)\beta(y')G_{\varphi,\psi}(y'-y)\,dy\,dy'
 \tag{TU2}
\]
whenever the compact supports of $\alpha,\beta\in C_c^1(\mathbb R^2;\mathbb C)$
are disjoint. There is a Hilbert space with strongly commuting
self-adjoint operators $H,P$, with $H\ge |P|$, and vectors
$\zeta_{\varphi,a}$, $a>0$, such that
\[
 \begin{split}
 \zeta_{\varphi,b}&=e^{-(b-a)H}\zeta_{\varphi,a}\quad(b>a),\\
 G_{\varphi,\psi}(t,x)
 &=\langle\zeta_{\varphi,a},
       e^{-(t-2a)H+ixP}\zeta_{\psi,a}\rangle\quad(t>2a).
 \end{split}
 \tag{TU3}
\]
This is a point-kernel interface, not a bound on unsmeared field
operators. At order two every distinct planar pair is in an ordinary
coordinate frame. Consequently the two-point specialization of the planar
kernel theorem supplies (TU2)--(TU3) without an exceptional planar set.

When $S_2$ is the connected function of a functional with unit vacuum vector
$\Omega$, apply the interface to
$\zeta_{\varphi,a}-\langle\Omega,\zeta_{\varphi,a}\rangle\Omega$.
The vacuum is translation invariant, so its orthogonal projection commutes
with $H,P$; this subtraction preserves (TU3), positivity and the cone. Thus
the interface here concerns the connected kernel when the ultraviolet
hypothesis concerns the connected function. No uniqueness of the vacuum
is required.

Kernel uniqueness follows by testing (TU2) with small nonnegative
approximate identities at two distinct points. The same test shows that
the kernels inherit reality, bilinearity in the labels, and all symmetries
of $S_2$ that preserve the plane. In particular, for
$G_\varphi=G_{\varphi,\varphi}$,
\[
 G_\varphi(R_{\pi/2}y)=G_\varphi(y),\qquad
 G_\varphi(y^0,-y^1)=G_\varphi(y),
 \qquad G_\varphi(y)\in\mathbb R.
 \tag{TU4}
\]

\subsection{The axis measure and the plate norm}

Fix a real label $\varphi$. Let $\mathsf P$ be the joint spectral measure
of $(H,P)$ and set
\[
 \mu_{\varphi,a}(B)=
    \langle\zeta_{\varphi,a},\mathsf P(B)\zeta_{\varphi,a}\rangle,
 \qquad d\mu_\varphi(E,q)=e^{2aE}\,d\mu_{\varphi,a}(E,q).
 \tag{TU5}
\]
The first measure is finite and positive. Equation (TU3) gives
$d\mu_{\varphi,b}=e^{-2(b-a)E}d\mu_{\varphi,a}$, so the second measure
is independent of $a$. It is positive and carried by
\[
 \mathcal C=\{(E,q)\in\mathbb R^2:E\ge |q|\}.
\]
For every $t>0$ choose $a<t/2$ in (TU3). Then
\[
 G_\varphi(t,x)=\int e^{-tE+ixq}\,d\mu_\varphi(E,q),
 \qquad
 \Gamma_\varphi(t)=\int e^{-tE}\,d\mu_\varphi=G_\varphi(t,0)<\infty.
 \tag{TU6}
\]
The measure is locally finite and $\sigma$-finite: on $\{E\le N\}$ its
mass is at most $e^{tN}\Gamma_\varphi(t)$. The function $\Gamma_\varphi$
is nonnegative and nonincreasing. By a quarter turn, any $y\ne0$ has a
positive first coordinate at least $|y|/\sqrt2$. Therefore
\[
 |G_\varphi(y)|\le\Gamma_\varphi(|y|/\sqrt2).
 \tag{TU7}
\]
This bound will give the diagonal reconstruction its bounded time orbits.
It uses no decay at large distance. Define also
\[
 \mathbb V=\{v\in\mathbb C^2:
                   \operatorname{Re}v^0>|\operatorname{Im}v^1|\},
 \qquad K_\varphi(v)=\int e^{-v^0E+iv^1q}\,d\mu_\varphi.
 \tag{TU7a}
\]
On a compact subset of this tube the absolute multiplier is at most
$e^{-cE}$ for some $c>0$. Its derivatives are dominated by a smaller
exponential because $|q|\le E$. Thus $K_\varphi$ is holomorphic and
agrees with $G_\varphi$ on the real positive-time half-plane.

On positive-time tests put
$\|[f]\|^2=S_2(\theta\overline f,f)$; reflection positivity and
polarization give the associated Hilbert seminorm and its triangle
inequality. For $\alpha$ supported in $\{y^0>0\}$, (TU2) and (TU6) give
\[
 \|[\alpha\otimes\varphi]\|^2
 =\int\left|\iint\alpha(t,x)e^{-tE+ixq}\,dt\,dx\right|^2
                                                 d\mu_\varphi(E,q).
 \tag{TU8}
\]
Indeed, if $t\ge\delta>0$ on the support, the absolute integral before
interchanging the variables is bounded by
$\|\alpha\|_{L^1}^2\Gamma_\varphi(2\delta)$. This proves Fubini's
applicability, including when the unweighted measure has infinite mass.

\subsection{The ultraviolet class and fixed shapes}

For $d>0$ and $\vartheta>0$, let $\mathfrak T(d;\vartheta)$ consist of
pairs of nonzero, nonnegative $C_c^1(\mathbb R^4)$ tests $f,h$, each
supported in a ball of radius $d$, satisfying
\[
 \begin{split}
 d&\le\operatorname{dist}(\operatorname{supp}f,
                          \operatorname{supp}h)\le6d,\\
 \int f&\ge\vartheta\|f\|_\infty d^4,\qquad
 \operatorname{Lip}f\le\|f\|_\infty/(\vartheta d),
 \end{split}
 \tag{TU9}
\]
and the analogous conditions for $h$. The balls may be placed anywhere.
The ultraviolet input is a finite nonnegative function $u$, with
$u(d)\to0$ as $d\downarrow0$, such that for $0<d\le d_0$,
\[
 0\le S_2(f,h)\le u(d)\|f\|_\infty\|h\|_\infty
               \quad\bigl((f,h)\in\mathfrak T(d;\vartheta_*)\bigr).
 \tag{TU10}
\]
Here the universal $\vartheta_*>0$ is chosen below, before $S_2$.
A hypothesis at a smaller shape parameter also suffices. A hypothesis
only on a narrower positive-shape class must be checked against this
choice; that inclusion is part of the ultraviolet premise.

The following explicit shapes fix all regularity constants. Choose
$a\in C_c^1((2,3))$, between zero and one, equal to one on
$[9/4,11/4]$, zero outside $[17/8,23/8]$, with cubic smoothstep ramps on
the two intervening intervals. Then
$\int a=5/8$ and $\|a'\|_\infty\le12$.
Let $I=[-1/4,1/4]$, $\tau=\mathbf1_I*\mathbf1_I$, and
\[
 b=12(\tau*\tau),\qquad c_b=\int b=3/4.
 \tag{TU11}
\]
This is a nonnegative $C_c^2$ spline, supported in $[-1,1]$, with maximum
one and $\|b'\|_\infty=2$. Explicitly it is
$1-6s^2+6|s|^3$ for $|s|\le1/2$, and $2(1-|s|)^3$ for
$1/2\le|s|\le1$. With the Fourier convention
$\widehat b(q)=\int b(s)e^{isq}\,ds$, convolution gives
\[
 \widehat b(q)=c_b\left(\frac{\sin(q/4)}{q/4}\right)^4,
 \tag{TU12}
\]
where the quotient has value one at zero.

For completeness a transverse partition can use
$q(s)=b(4s/3)$,
$D(s)=\sum_{j\in\mathbb Z}q(s-j)$ and $\chi(s)=q(s)/D(s)$.
There are at most two nonzero summands, $D\ge2/27$ by choosing a nearest
integer, and $\sum_j\chi(s-j)=1$.
The bounds $\|q'\|_\infty\le8/3$ and
$\|D'\|_\infty\le16/3$ give $\|\chi'\|_\infty\le1008$.
Put $\psi(z)=\chi(z_1)\chi(z_2)$. It takes values in $[0,1]$, has
support in $[-3/4,3/4]^2$, and its integer translates sum to one.
Let $\widetilde\psi(z)=\rho(z_1)\rho(z_2)$, where $\rho=1$ on
$[-1,1]$, $\rho=0$ outside $[-3/2,3/2]$, and on each intervening
interval use the cubic smoothstep. Thus
\[
 \int\widetilde\psi=25/4,\quad
 \|\nabla\widetilde\psi\|_\infty\le6,\quad
 \|\nabla\psi\|_\infty\le2016,\quad
 \psi\le\widetilde\psi.
\]
All these functions are globally $C^1$. The numerical choice
\[
 \vartheta_* =1/12000
 \tag{TU13}
\]
will suffice for all tests constructed next.

\subsection{Dilation removes the spline zeros}

For $h>0$ and $1\le\lambda\le2$ put
\[
 \alpha_h^\lambda(t,x)=a(t/h)b(4x/(\lambda h)),
 \qquad f_h^\lambda=\alpha_h^\lambda\otimes\varphi.
\]
The time transform and spatial transform are
\[
 A_h(E)=\int a(t/h)e^{-tE}\,dt,\qquad
 B_{h,\lambda}(q)=\frac{c_b\lambda h}{4}
       \left(\frac{\sin(\lambda hq/16)}{\lambda hq/16}\right)^4.
\]
For $E\ge0$, $A_h(E)\ge(h/2)e^{-3hE}$. There is an absolute $c>0$
such that
\[
 \int_1^2|B_{h,\lambda}(q)|^2\,d\lambda
                  \ge c h^2(1+h^2q^2)^{-4}.
 \tag{TU14}
\]
Here is a proof that includes every zero of the transform. Set
$x=h|q|/16$. For $0\le x\le1$, the sinc in the display is bounded
below on $0\le\lambda x\le2$, giving a fixed positive lower bound
after division by $h^2$. For $x\ge1$, that divided integral is at least
\[
 c x^{-8}\int_1^2\sin^8(\lambda x)\,d\lambda
       =c x^{-9}\int_x^{2x}\sin^8 v\,dv.
\]
The number
$c_1=\min_{r\in[0,\pi]}\int_r^{r+1}\sin^8v\,dv$ is positive:
the integral is continuous, periodic and never zero. An interval of length
$x\ge1$ contains $\lfloor x\rfloor\ge x/2$ successive intervals of
length one. The last display is therefore at least $c x^{-8}$.
Both regimes imply (TU14), with constants independent of $h,q$.

The finite number
$C_*=\sup_{r\ge0}e^{-r}(1+r^2)^4$ and $|q|\le E$ imply
\[
 e^{-6hE}(1+h^2q^2)^{-4}\ge C_*^{-1}e^{-7hE}.
\]
Using (TU8), (TU14) and nonnegative Fubini yields
\[
 \int_1^2\|[f_h^\lambda]\|^2\,d\lambda
    \ge c h^4\Gamma_\varphi(7h).
 \tag{TU15}
\]
The cone is essential here. Without it, arbitrarily large $|q|$ at small
$E$ could make the spline factor too small.

\subsection{Positive pedestals put the plate in the ultraviolet class}

Assume $\varphi\ne0$, choose $r_1>0$ with
$\operatorname{supp}\varphi\subset B(0,r_1)$, and take
$0<h\le\min(1,r_1,2d_0/5)$. Put
$A=\|\varphi\|_\infty+h\|\nabla\varphi\|_\infty>0$ and
\[
 \psi_i(z)=\psi(z/h-i),\quad
 \widetilde\psi_i(z)=\widetilde\psi(z/h-i),\quad
 \varphi_i^-=2A\widetilde\psi_i,\quad
 \varphi_i^+=2A\widetilde\psi_i+\psi_i\varphi.
\]
Only indices in
$\mathcal I=\{i:\psi_i\varphi\not\equiv0\}$ are needed. Its size is
at most $(2r_1/h+3)^2\le25r_1^2h^{-2}$. For a fixed $\lambda$ let
$h_i^\pm=\alpha_h^\lambda\otimes\varphi_i^\pm$. Then
\[
 f_h^\lambda=\sum_{i\in\mathcal I}(h_i^+-h_i^-),\qquad
 A\widetilde\psi_i\le\varphi_i^\pm\le3A\widetilde\psi_i.
 \tag{TU16}
\]
In particular these tests are nonnegative, nonzero, and have sup norm
between $A$ and $3A$.

We check (TU9), rather than invoking a positive decomposition without
shape control. Each support lies in the ball of radius $d=5h/2$ about
$(5h/2,0,hi)$: its four coordinate half-widths are at most
$h/2,h/2,3h/2,3h/2$, whose squared sum is $5h^2<d^2$.
The reflected ball also has radius $d$. The distance between the support
and its time reflection is at least $4h$ and at most $6h$, hence between
$d$ and $6d$.

Moreover, direct product integration gives
\[
 \int h_i^\pm\ge A h^4
    \frac{\int a\,\int b\,\int\widetilde\psi}{4}
       =\frac{375}{512}A h^4.
\]
Since $\|h_i^\pm\|_\infty\le3A$, the integral condition holds for
every $\vartheta\le1/160$. The product derivative rule gives
\[
 \operatorname{Lip}h_i^\pm\le\frac{A}{h}
 \left[3(12+4\cdot2)+2\cdot6+2016+1\right]
 =2089A/h.
\]
As $\|h_i^\pm\|_\infty\ge A$, (TU13) gives the Lipschitz condition:
$2089A/h\le\|h_i^\pm\|_\infty/(\vartheta_*d)$.
The reflected tests have the same integrals and norms. Thus
$(\theta h_i^\pm,h_i^\pm)\in\mathfrak T(5h/2;\vartheta_*)$.

By (TU10), reflection positivity and the Hilbert triangle inequality,
\[
 \|[f_h^\lambda]\|
 \le\sum_i\bigl(\|[h_i^+]\|+\|[h_i^-]\|\bigr)
 \le6\,\#\mathcal I\,A\sqrt{u(5h/2)}.
\]
Combining its square with (TU15) proves
\[
 \Gamma_\varphi(7h)
   \le C r_1^4
       (\|\varphi\|_\infty+\|\nabla\varphi\|_\infty)^2
       h^{-8}u(5h/2).
 \tag{TU17}
\]
Consequently
\[
                 t^8\Gamma_\varphi(t)\longrightarrow0
                         \quad(t\downarrow0).
 \tag{TU18}
\]
The zero label is immediate. The two factors of $h^{-4}$ in (TU17)
have different origins: normalization in (TU15), and the square of the
number of transverse cells. No uniform bound on a compact profile's
Fourier transform from below was used.

\subsection{Entire angular continuation and the critical Fourier modes}

For $r>0$, define on $|\operatorname{Re}z|<\pi/2$
\[
 g_r^{[0]}(z)=\int
       \exp[-rE\cos z+irq\sin z]\,d\mu_\varphi(E,q).
 \tag{TU19}
\]
Writing $z=u+is$, the absolute integrand is
$\exp[-r\cos u(E\cosh s+q\sinh s)]$.
On the cone this is at most $\exp[-r\cos u\,e^{-|s|}E]$.
On a compact substrip its derivatives are dominated by a polynomial in
$E$ times a smaller exponential. Thus the integral is holomorphic.
Its real restriction is $G_\varphi(r\cos u,r\sin u)$.

On $|\operatorname{Re}z-k\pi/2|<\pi/2$ define
$g_r^{[k]}(z)=g_r^{[0]}(z-k\pi/2)$. On real overlaps these functions
agree by (TU4). The complex overlaps are connected vertical strips with
real intervals, so the identity theorem glues them. Their union is the
complex plane. The resulting $g_r$ is entire and $\pi/2$-periodic.
For every real $u,s$ choose $k$ with $|u-k\pi/2|\le\pi/4$ to obtain
\[
 |g_r(u+is)|\le\Gamma_\varphi(r e^{-|s|}/\sqrt2)
                         =o(e^{8|s|}).
 \tag{TU20}
\]
The last little-o, uniform in $u$, is exactly (TU18).

Define
$a_m(r)=\frac2\pi\int_0^{\pi/2}g_r(u)e^{-4imu}\,du$.
Cauchy's theorem on a rectangle of height $s$, with cancellation of its
vertical sides by periodicity, gives
\[
 |a_m(r)|\le e^{4ms}\sup_{u\in[0,\pi/2]}|g_r(u+is)|.
\]
For $m\ge2$ take $s\to-\infty$; for $m\le-2$ take
$s\to+\infty$. Equation (TU20) gives $a_m(r)=0$ in both cases,
including $|m|=2$. A mere $O(t^{-8})$ bound on $\Gamma$ would not
eliminate that boundary mode. Fourier uniqueness for a continuous periodic
function now gives, first on the real line and then everywhere by the
identity theorem,
\[
 g_r(z)=a_0(r)+a_1(r)e^{4iz}+a_{-1}(r)e^{-4iz}.
 \tag{TU21}
\]
One elementary uniqueness proof uses the positive Fejer kernels: their
integrals are one, their mass outside every fixed neighbourhood of zero
tends to zero, and their convolution converges uniformly on continuous
periodic functions. A function with all Fourier coefficients zero has
all these convolutions zero. Thus no convergence assumption on a formal
Fourier series is implicit in (TU21).

\subsection{The positive measure in the diagonal frame}

Put $D(t,x)=G_\varphi(R_{\pi/4}(t,x))$ for $t>0$.
We construct its positive measure without assuming a diagonal cone.
For $y=(t,x)$, $t>0$, introduce a formal vector $e_y$ with
\[
 \langle e_y,e_{y'}\rangle_D=D(t+t',x'-x).
 \tag{TU22}
\]
This kernel is positive semidefinite. To see it, approximate each planar
point $R_{\pi/4}y$ by a nonnegative radial smooth mollifier supported
strictly on the positive side of the diagonal mirror, tensor it with
$\varphi$, and apply (TU1) to an arbitrary complex linear combination.
Formula (TU2) and continuity away from zero give (TU22) in the limit.
Thus its quotient and completion define a Hilbert space.

On finite sums define $T_s e_{(t,x)}=e_{(t+s,x)}$, $s\ge0$, and
$U_b e_{(t,x)}=e_{(t,x+b)}$, $b\in\mathbb R$. The $U_b$ preserve the
form and have inverse $U_{-b}$. The $T_s$ are symmetric and commute with
the $U_b$. If $v=\sum c_i e_{(t_i,x_i)}$ and
$\delta=2\min_i t_i>0$, (TU7) gives a finite bound on
$\|T_s v\|$ for every $s\ge0$, namely its square is at most
$\Gamma_\varphi(\delta/\sqrt2)(\sum|c_i|)^2$.
On the unquotiented vector space Cauchy--Schwarz gives
\[
 \|T_s v\|^2=\langle v,T_{2s}v\rangle_D
                         \le\|v\|\,\|T_{2s}v\|.
\]
If $\|v\|=0$ this preserves the null space. Otherwise iterate and use
the bounded orbit to get $\|T_s v\|\le\|v\|$. Therefore $T_s$ descend
to self-adjoint contractions. They are positive because
$T_s=T_{s/2}^2$. Continuity of $D$ implies strong continuity of both
families on point vectors, hence on their completion.

There are strongly commuting generators $H_D\ge0$ and $P_D$ with
$T_s=e^{-sH_D}$ and $U_b=e^{ibP_D}$. For specificity, $T_1$ is
injective: $T_1v=0$ implies $T_{1/2}v=0$, then all its dyadic roots
annihilate $v$, and strong continuity gives $v=0$.
Define $H_D=-\log T_1$ by spectral calculus. Uniqueness of positive
roots proves the semigroup identity first at rational times and then at
all times. The spectral theorem for the commuting unitary group supplies
$P_D$ and the joint spectral measure. This is the elementary Hilbert
reconstruction for (TU22), not an OS analytic-continuation theorem.

Let $v_a=e_{(a,0)}$. If $\mathsf P_D$ is that joint spectral measure,
define
\[
 \nu_a(B)=\langle v_a,\mathsf P_D(B)v_a\rangle_D,\qquad
 d\nu(E,q)=e^{2aE}\,d\nu_a(E,q).
 \tag{TU23}
\]
As $v_b=T_{b-a}v_a$ for $b>a$, $\nu$ is independent of $a$.
It is a positive, locally finite, $\sigma$-finite measure on
$[0,\infty)\times\mathbb R$, and for $t>0$
\[
 D(t,x)=\int e^{-tE+ixq}\,d\nu,\qquad
 \int e^{-tE}\,d\nu=D(t,0)<\infty.
 \tag{TU24}
\]
Indeed choose $a<t/2$ and evaluate
$\langle v_a,T_{t-2a}U_xv_a\rangle_D$.
Nothing so far asserts $|q|\le E$ for $\nu$.

\subsection{Fourier--Laplace uniqueness and the boost identity}

We shall use the following uniqueness fact. Positive Borel measures
$\rho_1,\rho_2$ on $[0,\infty)\times\mathbb R$ with finite
$\int e^{-tE}d\rho_i$ for every $t>0$ are equal if their transforms
$\int e^{-tE+ixq}d\rho_i$ agree for all $t>0,x\in\mathbb R$.
For each real $x$ these transforms are holomorphic in $t$ on
$\operatorname{Re}t>0$, by exponential domination, so the identity
theorem gives equality at $t=a+i\tau$. These are precisely the joint
Fourier transforms of the finite measures $e^{-aE}\rho_i$.
Finite Fourier transforms determine their measures: convolve their
difference with a Gaussian, apply the Gaussian Fourier integral and
Fubini to get zero, and let the Gaussian width tend to zero against
compactly supported continuous functions. Hence the weighted measures,
and then the original measures, agree. The same proof applies to finite
positive linear combinations on either side of an identity.

For $s\in\mathbb R$ define the Lorentz map of spectral variables
\[
 L_s(E,q)=(E\cosh s+q\sinh s,E\sinh s+q\cosh s),
 \qquad \mu_s=(L_s)_*\mu_\varphi.
 \tag{TU25}
\]
This preserves $\mathcal C$ and $L_sL_t=L_{s+t}$. Its energy satisfies
$E_s\ge e^{-|s|}E$, so $\int e^{-tE}d\mu_s\le
\Gamma_\varphi(te^{-|s|})$. Define
$F_s(t,x)=\int e^{-tE+ixq}d\mu_s$ for $t>0$.
Direct multiplication of the rotation matrix by $i s$ gives
\[
 F_s(v)=K_\varphi(R_{is}v)
       =g_r(\theta+is)
       \quad\bigl(v=R_\theta(re_0),\ v^0>0\bigr),
 \tag{TU26}
\]
where the middle expression uses (TU7a).
In fact $R_{is}(t,x)=(t\cosh s-ix\sinh s,
x\cosh s+it\sinh s)$, and its exponential multiplier is exactly
$e^{-tE_s+ixq_s}$. Thus (TU26) holds on the whole real half-plane.

Write $C=\cosh4>1$. The three terms in (TU21) obey the algebraic identity
\[
 (C-1)g_r(\theta+\pi/4)
       +g_r(\theta+i)+g_r(\theta-i)
                         =(C+1)g_r(\theta).
\]
Using (TU6), (TU24) and (TU26) and then transform uniqueness gives the
identity of positive measures
\[
 (C-1)\nu+\mu_1+\mu_{-1}=(C+1)\mu_\varphi.
 \tag{TU27}
\]
Every measure has the required finite Laplace weights. In particular
\[
       0\le\nu\le\frac{C+1}{C-1}\mu_\varphi.
 \tag{TU28}
\]
The diagonal measure is therefore carried by the exact cone. This
conclusion was deduced from (TU27), not inserted as an additional
reflection-positivity consequence.

For $r>0$ and real $s$ the function
$w_s(E,q)=\exp[-r(E\cosh s+q\sinh s)]$ is positive on the cone and
at most $e^{-r e^{-|s|}E}$. It is integrable against all four measures
in (TU27); for $\mu_{\pm1}$ use $L_sL_{\pm1}=L_{s\pm1}$.
Integrating (TU27) against $w_s$ yields
\[
 (C-1)\int w_s\,d\nu+F_{s+1}(r,0)+F_{s-1}(r,0)
                                  =(C+1)F_s(r,0).
\]
Substitute (TU21) into this identity to obtain
\[
 \begin{split}
 F_s(r,0)&=a_0(r)+a_1(r)e^{-4s}+a_{-1}(r)e^{4s}\ge0,\\
 \int w_s\,d\nu
       &=a_0(r)-a_1(r)e^{-4s}-a_{-1}(r)e^{4s}\ge0.
 \end{split}
 \tag{TU29}
\]
The first positivity comes from the axis measure, the second from the
diagonal measure. Divide by $e^{4s}$ and let $s\to+\infty$:
the two limits give $a_{-1}(r)\ge0$ and $a_{-1}(r)\le0$.
Then let $s\to-\infty$ to get both signs for $a_1(r)$.
The limits also show these coefficients are real, if that was not already
deduced from (TU4). Thus $a_1(r)=a_{-1}(r)=0$ for every $r>0$.

% Keep the closing statement and argument together.
\clearpage
\begin{theorem}[Two-point angular elimination]
\label{thm:angular-two-point-elimination}
Under (TU1)--(TU3) and the short-distance hypothesis (TU9)--(TU13),
every equal-real-label plane kernel is radial:
\[
             G_\varphi(R_t y)=G_\varphi(y)
                      \quad(y\ne0,\ t\in\mathbb R).
\]
Every nonzero angular coefficient of the two-point kernel vanishes,
including coefficients with mixed transverse labels.
\end{theorem}

\begin{proof}
The preceding argument proves the first assertion. Symmetry of $S_2$
gives $G_{\varphi,\psi}(y)=G_{\psi,\varphi}(-y)$; invariance under
the planar half turn gives $G_{\psi,\varphi}(-y)=G_{\psi,\varphi}(y)$.
Therefore the mixed kernel is symmetric in its real labels and
\[
 G_{\varphi,\psi}
       =\tfrac14(G_{\varphi+\psi}-G_{\varphi-\psi}).
\]
Both terms on the right are radial. Complex labels follow by bilinear
extension. Uniqueness of Fourier coefficients identifies this conclusion
with vanishing of every $m\ne0$ coefficient in any finite-angular-band
construction of the same two-point kernel. In its notation,
$\overline m(2)=0$.
\end{proof}

This order-two conclusion enters the all-orders harmonic-closure
argument. The passage to arbitrary separated four-dimensional tuples
uses its separately stated extension hypotheses. The premises here
are the stated ultraviolet law and kernel/spectral interface.
The Ward and infrared routes retain their separate hypotheses.

% END MODULE chapters/angular_two_point_elimination_v5.tex


% BEGIN MODULE chapters/angular_harmonic_closure_v5.tex
\section{From the two-point angular law to full Euclidean covariance}
\label{ch:angular-harmonic-closure}

This section closes the all-orders and distributional parts of the angular
argument. Its inputs are the analytic kernels and boost Gram identities of
Section~\ref{ch:angular-kernel-construction}, and the two-point elimination
of Section~\ref{ch:angular-two-point-elimination}. In particular, the
two-point ultraviolet estimate remains part of the hypotheses of the
combined implication. No rotational Ward identity is used here.

Work first with the centered scalar hierarchy. Compact separated moments
are symmetric, translation invariant and hypercubic invariant, have position-uniform local
$C^1$ bounds, and are reflection positive in every coordinate picture.
In each such picture the Hilbert generators have the exact cube cone
$H_j\ge |P_{j,r}|$ for every tangential coordinate $r\ne j$.
Their off-diagonal distributions are those constructed in
Section~\ref{ch:quantitative-flat-extension}.

\subsection{The harmonic data used below}

The chosen plane has coordinates $(x_4,x_1)$, and $R_u$ denotes its
rotation through $u$. A plate label is a real smooth compact function of
the two transverse coordinates. The analytic kernel construction supplies,
for each $n$, real-analytic functions
\[
 c_m(y_1,\ldots,y_n;\varphi_1,\ldots,\varphi_n),
 \qquad y_i\in\mathbb R^2,\quad y_i\ne y_j,
 \tag{HC1}
\]
multilinear in the labels, with only $|m|\le Cn$ possibly nonzero.
Their sum is the moment density on the ordinary ordering charts.
They are symmetric under simultaneous permutations of positions and labels,
translation invariant in the plane, and satisfy
\[
 c_m(R_u\mathbf y;\vec\varphi)
       =e^{4imu}c_m(\mathbf y;\vec\varphi),
 \qquad c_{-m}=\overline{c_m}
 \quad\hbox{for real labels.}
 \tag{HC2}
\]
An ordinary chart means that all time coordinates are distinct, or that
all abscissae are distinct; a fixed ordering is imposed on its connected
component. The functions (HC1) already extend over the nongeneric planar
configurations with distinct points. This is the analytic-projector
conclusion of the kernel construction, not a density assertion about
distributions on the exceptional set.

Let $b_n$ be the greatest absolute index of a harmonic that is not
identically zero for some choice of real labels. Set $b_n=0$ if there is
no nonzero nonconstant harmonic, and put $b_0=0$ for the vacuum kernel.
Thus $b_n$ are nonnegative integers and $b_n\le Cn$. The boost Gram
identity gives
\[
 2b_{k+l}\le b_{2k}+b_{2l},\qquad k,l\ge0.
 \tag{HC3}
\]
For clarity, this inequality compares global degrees. A nonzero coefficient
can be tested on an open configuration set split into $k$ negative and
$l$ positive times: real analyticity and connectedness prevent a nonzero
coefficient from vanishing on that whole open set. Cauchy--Schwarz for the
boosted product vectors bounds its largest exponential by the two diagonal
ones. The empty product is the vacuum vector, so the cases $k=0$ or $l=0$
are included.

For $m_*=b_{2k}$ the coefficient $c_{-m_*}$ is a positive kernel on
reflected pairs of $k$-point configurations in a half-plane. This follows
by multiplying the positive boost Gram matrix by $e^{-4m_*s}$ and taking
$s\to+\infty$. Covariance (HC2) transports it to every oriented line:
after multiplication by the fixed unit phase for that line it is again
positive. We use the zero-row consequence of this positivity, which is
unchanged by that phase. All these statements have explicit product-vector
proofs in the preceding kernel section; analytic continuation of a
nonnegative real-valued function alone would not prove them.

\subsection{The polygon divisibility argument, with label release}

Fix $k\ge1$, put $m_*=b_{2k}$, and abbreviate $H=c_{-m_*}$.
We prove
\[
                    k\mid 2m_* . \tag{HC4}
\]
Suppose otherwise. For any real plate label $\varphi$, place that same
label at every vertex of a regular $2k$-gon. Rotation through $\pi/k$
permutes the labelled vertices. Symmetry and (HC2) imply
\[
 H(\mathbf y;\varphi,\ldots,\varphi)
 =e^{-4\pi i m_*/k}H(\mathbf y;\varphi,\ldots,\varphi)=0.
 \tag{HC5}
\]
The polygon may be translated; rotation about its center has the same
effect by translation invariance. The assumption that (HC4) fails is
exactly what makes the phase in (HC5) different from one.

Call a labelled set $S$ a zero set if $H$ vanishes on every $2k$-point
extension of it, with all remaining real labels arbitrary. A full polygon
with the labels above is a zero set, since no completion is then needed.
Here is the reduction of such a set. Suppose it is contained, with its
labels, in $\theta Q\cup Q$, where $Q$ has $j\le k$ points strictly on
one side of a line, and the reflected labels agree. Fill $Q$ by any
$k-j$ additional points $A$ on that side, with arbitrary real labels.
The positive kernel has zero diagonal at $Q\cup A$, because its diagonal
configuration contains $S$. Its entire row is therefore zero. Indeed,
positivity of the $2\times2$ Gram matrix gives
$|K(X,Y)|^2\le K(X,X)K(Y,Y)$; a zero diagonal makes the right side zero.
Consequently
\[
 H\big(\theta(Q\cup A)\cup B\big)=0
 \tag{HC6}
\]
for every $k$-point $B$ on the other side of the reflected set, with
arbitrary real labels.

For each fixed choice of these free labels, (HC6) holds on a nonempty
open set of all the free positions: the reflected points of $A$ and the
points of $B$ can be varied independently in their respective half-planes.
The configuration space of these $2k-j$ free points, avoiding each other
and the fixed reflected $Q$, is connected. It is the complement of a
finite union of affine sets of codimension two. One elementary proof is
to join two configurations through a generic intermediate point; the
intermediate points whose segment meets one of these affine sets form a
finite union of measure-zero cones. Real analyticity and the identity
principle extend (HC6) to this whole configuration space. It follows that
$\theta Q$ itself is a zero set, with no restriction on the remaining
labels or half-planes. This is the label-release step; it is stronger than
just showing that a special diagonal configuration has zero value.

Apply this reduction to the polygon. Reflection in the perpendicular
bisector of an edge pairs its vertices and has no vertex on its line.
It first reduces the full polygon to a consecutive arc of $k$ vertices.
A consecutive arc of $j\le k$ vertices fits into the reflected union of
an arc of $\lceil j/2\rceil$ vertices. To check the odd case, index the
vertices by integers modulo $2k$. For the arc $0,\ldots,j-1$, reflect
indices by $i\mapsto j-i$ when $j$ is odd; the union has $j+1$ vertices,
paired without a fixed vertex. For even $j$, use $i\mapsto j-1-i$ and
the arc itself is paired. These are edge-bisector reflections, and the
selected half has the stated size. All its labels are still $\varphi$.
Repeated reduction reaches a single labelled point.

Translation covariance makes every point with label $\varphi$ a zero
set. Since $\varphi$ was arbitrary, every labelled $2k$-point
configuration has zero value: choose any one of its labels as that
single fixed label and use the unrestricted completion. Thus $H$ is
identically zero, contradicting the definition of $m_*$. This proves
(HC4). When $m_*=0$ the divisibility is automatic; no contradiction
argument is needed.

\subsection{Discrete convexity reduces every order to two points}

Set $a_k=b_{2k}$, $a_0=0$. Equation (HC3) with $k-1,k+1$ gives
\[
 2a_k\le a_{k-1}+a_{k+1},\qquad 0\le a_k\le C'k.
 \tag{HC7}
\]
The integer increments $d_k=a_k-a_{k-1}$ are nondecreasing, initially
nonnegative, and bounded. If one increment exceeded the linear-growth
constant, all later increments would do so, contradicting (HC7).
Hence they are eventually a fixed integer $d$. For large $k$ write
\[
 a_k=dk-D,\qquad D=\sum_{j\ge1}(d-d_j)\ge0,
 \tag{HC8}
\]
where only finitely many summands are nonzero. By (HC4), $k\mid2a_k$
for every $k$, so every sufficiently large $k$ divides $2D$. Thus
$D=0$. Each summand in (HC8) is nonnegative, and every increment equals
$d=a_1$. Therefore
\[
 b_{2k}=k b_2,\qquad 2b_n\le b_{2n}=n b_2.
 \tag{HC9}
\]
The second inequality is the empty-tuple case of (HC3). Convexity and
linear growth without (HC4) would not suffice: $a_k=\max(0,k-1)$ is
a counterexample. The polygon argument is what removes that defect.

The two-point elimination theorem gives $b_2=0$ under its ultraviolet
shape estimate. Equations (HC9) then give $b_n=0$ for every $n$.
Consequently the ordinary-chart kernel is invariant under rotations in
the chosen plane. The following argument upgrades this fact to the
actual off-diagonal distributions, including their exceptional supports.

\subsection{A local analytic translation lemma}

Fix a coordinate $j$ and divide the entries of a compact separated tuple
into two groups whose supports lie on opposite sides of a coordinate
hyperplane, with a positive gap. Their scalar moment is a matrix element
in the coordinate-$j$ Hilbert picture. After removing some unused gap,
it has the form
\[
 \sigma(F_-\cdot\tau_z F_+)
   =\langle\Phi,e^{-(d+z_j)H_j}U_{z_\parallel}\Psi\rangle,
 \qquad d>0,
 \tag{HC10}
\]
for real $z$ sufficiently small. Here the left group has been reflected
to define $\Phi$, both groups have been shifted away from their mirror,
and $d$ is the remaining normal separation. This is exactly the
compact-tuple translation dictionary, applied to products in each group.
Its vectors need not be point fields, and no coincident product is used.

Let $v$ be any real four-vector. The right side with $z=sv$ is
holomorphic for small complex $s$. The precise sufficient domain is
\[
 d+v_j\operatorname{Re}s>
          \sum_{r\ne j}|v_r|\,|\operatorname{Im}s|.
 \tag{HC11}
\]
Indeed the joint spectral cone $|P_{j,r}|\le H_j$ bounds the spectral
exponential by one there. On every smaller domain a positive damping
margin controls all powers of $H_j$ and $P_{j,r}$, so difference quotients
converge in operator norm. This proves holomorphy and local uniform
bounds, rather than formally translating test functions to complex
arguments. Translating the negative group instead gives the same
conclusion with the corresponding signs. In particular, any specified
half of a strict coordinate cut may be moved in any real direction,
with a holomorphic parameter near zero.

\subsection{Removing the exceptional distributional support}

For $n\ge2$ put
\[
 \mathcal C_n=\{(x_1,\ldots,x_n)\in(\mathbb R^4)^n:
                                      x_i\ne x_l\ (i\ne l)\}.
\]
Let $T_n$ be the off-diagonal moment distribution. For a fixed rotation
$R_u$ define its difference by
\[
 D_u(f_1\otimes\cdots\otimes f_n)
   =\sigma(f_1,\ldots,f_n)-\sigma(R_uf_1,\ldots,R_uf_n),
 \qquad (R_uf)(x)=f(R_u^{-1}x).
 \tag{HC12}
\]
Both terms are well-defined distributions on $\mathcal C_n$.

First suppose all original and all rotated time coordinates at a
configuration are distinct. Choose small product support boxes keeping
both time orderings strict. In those boxes approximate each smooth test
in $C^1$ by finite sums
$\alpha(x_4,x_1)\varphi(x_2,x_3)$, with supports staying inside the
boxes. For example, a product Fourier expansion with fixed compact
cutoffs gives these approximants with all first derivatives. The
rotation-invariant ordinary-chart kernel gives equality for the plate
approximants. The local moment bounds pass it to the original tests.
Local product-tensor density then says that $D_u$ vanishes on this
ordinary open set. This step alone would not rule out a distribution
concentrated at the remaining time ties.

Fix any $y\in\mathcal C_n$. For each coordinate take all threshold
hyperplanes halfway between successive distinct coordinate levels of
the finite set $\{y_1,\ldots,y_n\}$. Each threshold divides the indices
into two nonempty sets, and these finitely many cuts together separate
every pair of indices. No point lies on one of their hyperplanes.
There exist $u_0>0$ and a product neighborhood $W$ of $y$ such that,
for every $|u|<u_0$, each cut has the same index partition for $x$ and
for $R_ux$, for all $x\in W$, with a uniform positive gap in both
configurations. To see this, take the minimum distance of the finitely
many original points from these finitely many hyperplanes. Rotation is
continuous and the points are bounded, so their displacement and then
the size of $W$ can be made smaller than a quarter of that minimum.
Also take $u_0<\pi/3$. The cuts and their partitions are kept fixed
throughout the following local argument.

Fix $|u|<u_0$. At $x\in W$ let
\[
 q(x)=\#\{i<l:x_{i,4}=x_{l,4}\}
       +\#\{i<l:(R_ux_i)_4=(R_ux_l)_4\}.
 \tag{HC13}
\]
The set $\{q\le q_0\}$ is open: every inequality between distinct
time coordinates remains an inequality in a sufficiently small
neighborhood. We prove by induction that $D_u$ vanishes on this open
set for each integer $q_0$. The case $q_0=0$ is the ordinary-chart
argument above.

Take a point with $q=q_0>0$. Choose a tied pair in either frame and
one of the fixed common cuts that separates it. Let $A$ be the
positive side of that cut. Translate every entry in $A$ by $s e_4$.
Original cross-group time differences have slope $1$ or $-1$ in $s$;
rotated cross-group differences have slope $\cos u$ or $-\cos u$.
All slopes are nonzero. Internal differences in either group are
unchanged. For all sufficiently small nonzero real $s$, no previously
nonzero difference becomes zero, while every previously zero cross
difference becomes nonzero. At least the chosen tied pair is of this
kind, so the number in (HC13) decreases.

Choose one such $s_*$ and a small real interval $I$ about it. Shrink
the product neighborhood $V$ of the original point so that its translates
stay in $W$, all previously nonzero differences stay nonzero, and all
cross-group differences in both frames are nonzero for every $s\in I$.
Every configuration in the translated product boxes then has
$q\le q_0-1$. By induction, $D_u$ vanishes there as a distribution.

Now take arbitrary smooth slot tests supported in $V$. Their difference
(HC12), after the group translation, is a holomorphic function of $s$
near the real segment joining zero to $I$. For the first moment term
this is (HC10)--(HC11) with $v=e_4$. For the rotated term the same
common cut separates the same two groups, and the direction is
$v=R_ue_4$. Strict gaps persist on the chosen segment; shrinking the
segment and its complex neighborhood supplies (HC11) in both cases.
This is why both configurations were required to share the cut. The
holomorphic function is zero on $I$, hence is zero at zero. Thus $D_u$
vanishes on every such product tensor in $V$. The explicit local
Fourier tensor density for distributions makes it vanish on all tests
in $V$.

This proves the induction step. Since $q\le n(n-1)$, the induction
terminates and $D_u$ vanishes on $W$. It excludes a singular part on
any of the exceptional time-tie sets, including configurations with
coincident plane projections but separated transverse positions. At
no point was the kernel integral assumed analytic across those sets.

\subsection{All angles, all planes, and the Schwartz domain}

We have proved local invariance for every sufficiently small $u$ on a
neighborhood of each point of $\mathcal C_n$. Let a compactly supported
off-diagonal test be given, and fix a finite angle $v$. The union of
the supports of its rotations between zero and $v$ is a compact subset
of $\mathcal C_n$. Cover it by finitely many neighborhoods just
constructed, and take the minimum of their positive angle bounds.
Partition $v$ into finitely many smaller steps. On each step the
difference distribution is zero on the current support, so composition
gives invariance through $v$.

Thus all separated smooth moments are invariant in the chosen plane.
Compact $C^1$ tests follow by smooth approximation on fixed separated
neighborhoods. The hypercubic group conjugates this circle of rotations
to rotations in every coordinate plane. Givens elimination of the
columns of an orthogonal matrix expresses every member of
$\mathrm{SO}(4)$ as a finite product of such plane rotations: successively
zero the entries below each diagonal position, using a rotation in the
corresponding two-coordinate plane, and finish with diagonal signs.
A coordinate reflection supplies the other component of $\mathrm O(4)$.

For $n=1$, the translation-invariant one-point functional is
$c\int f$ and is already invariant. Restoring such a mean to a centered
field uses a finite subset expansion on each compact tuple; this
commutes with every rotation because $\int Rf=\int f$.
Finally rotations act continuously on Schwartz space, preserve every
pairwise coincidence diagonal, and preserve flatness there. The
diagonal-flat extension and its constrained tensor density therefore
transfer the same invariance to that full test domain.

\begin{theorem}[All-orders angular closure]
\label{thm:all-orders-angular-closure}
Under the moment, cone, analytic-kernel and two-point hypotheses stated
above, the full scalar hierarchy is $\mathrm O(4)$ invariant, both on
separated compact $C^1$ tuples and on its unique diagonal-flat Schwartz
extension. The polygon, discrete-convexity and exceptional-support steps
require no further analytic rotation assumption.
\end{theorem}

This closes a conditional Euclidean-covariance implication. It neither
proves the ordinary interacting source estimates nor the ultraviolet
two-point law, and it does not identify the scalar cyclic theory with a
complete family of Yang--Mills observables.

% END MODULE chapters/angular_harmonic_closure_v5.tex


% BEGIN MODULE chapters/compact_disintegration_v3.tex
\section{Conditional measures on a compact metric configuration space}
\label{ch:compact-disintegration}

The conditional laws used in the phase construction take values in the
Borel probability measures on a compact metric configuration space.
The conditioning field may be completed and need not be countably
generated as a literal collection of sets. It is the Borel field of the
\emph{target} space that is countably generated.

\begin{theorem}[Compact-metric conditional law]
\label{thm:compact-conditional-law}
Let $X$ be a nonempty compact metric space, let $\mu$ be a Borel
probability measure on $X$, and let $\overline{\mathcal B}_\mu$ be the
$\mu$-completion of its Borel field $\mathcal B(X)$. Suppose

\begin{gather*}
\FitMath{\mathcal N_\mu\subset\mathcal H\subset
       \overline{\mathcal B}_\mu,}\tag{CD1}
\end{gather*}

where $\mathcal H$ is a $\sigma$-field and $\mathcal N_\mu$ denotes
all subsets of Borel $\mu$-null sets. There exists a map
$K:X\to\mathcal P(X)$, where $\mathcal P(X)$ carries its weak topology,
with the following properties.

The map $K$ is both Borel measurable and $\mathcal H$-measurable. For
every bounded Borel $f:X\to\mathbb C$,

\begin{gather*}
\FitMath{Kf(x):=\int_X f(y)\,K(x,dy)} \\[3pt]
\FitMath{\text{is a version of }E_\mu[f\mid\mathcal H].}\tag{CD2}
\end{gather*}

In particular $x\mapsto K(x,B)$ is $\mathcal H$-measurable for every
$B\in\mathcal B(X)$, and

\begin{gather*}
\FitMath{\int_H K(x,B)\,\mu(dx)=\mu(H\cap B)} \\[3pt]
\FitMath{(H\in\mathcal H, B\in\mathcal B(X)).}\tag{CD3}
\end{gather*}

The measure on the right is understood on the completion. If $K'$ is
another probability-valued kernel satisfying (CD2), then there is a
single $\mu$-null set $N$ such that

\begin{gather*}
\FitMath{K(x,\cdot)=K'(x,\cdot)} \\[3pt]
\FitMath{(x\notin N)}\tag{CD4}
\end{gather*}

as Borel probability measures. One may take $N$ to be a Borel null set.
Consequently the law $K_*\mu$ on $\mathcal P(X)$ is independent of the
choice of conditional kernel.
\end{theorem}

\subsection*{Borel representatives and countably many tests}

We first record why completion causes no obstacle to the Borel version
in this statement. Every bounded
$\overline{\mathcal B}_\mu$-measurable real function $u$ has a Borel
representative. Indeed, for each rational $r$ replace $\{u>r\}$ by a
Borel set differing from it inside a Borel null set. Outside the union
of those countably many null sets, the rational level cuts reconstruct
$u$. Their supremum, clipped to a bound for $u$, is a Borel function
equal to $u$ there. If $u$ was $\mathcal H$-measurable, this
representative remains $\mathcal H$-measurable: the change is confined
to a Borel null set, all of whose subsets belong to $\mathcal H$ by
(CD1). This observation will be applied only countably many times.

Choose a countable dense set $\{x_j\}$ in $X$. Let $D\subset C(X;\mathbb R)$
be the collection generated by the constant $1$ and the functions
$y\mapsto d(y,x_j)$ under rational linear combinations, multiplication,
pointwise minimum and pointwise maximum. It is countable, is a rational
vector lattice and algebra, and separates points: if $x\ne y$, a dense
point sufficiently near $x$ gives different distances from $x$ and $y$.
Its uniform closure is a real algebra containing the constants and
separating points, so the Stone--Weierstrass theorem makes $D$ uniformly
dense in $C(X;\mathbb R)$.

For each $f\in D$ take a conditional expectation $u_f$ of $f$ given
$\mathcal H$. These conditional expectations can be obtained as the
orthogonal projections of $f$ onto the closed subspace
$L^2(\mathcal H,\mu)$ of $L^2(\mu)$. The projection identity against
$1_H$, $H\in\mathcal H$, gives

\begin{gather*}
\FitMath{\int_Hu_f\,d\mu=\int_H f\,d\mu.}\tag{CD5}
\end{gather*}

Testing the sets on which $u_f<0$, or $u_f$ exceeds an upper bound for
$f$, proves positivity and the usual $L^\infty$ bound. Linearity and
uniqueness follow from the same identity. Choose for each $u_f$ a
bounded representative that is both Borel and $\mathcal H$-measurable,
using the preceding observation.

There are only countably many assertions

\begin{gather*}
\FitMath{u_1=1,} \\[3pt]
\FitMath{u_{rf+sg}=ru_f+su_g} \\[3pt]
\FitMath{(f,g\in D, r,s\in\mathbb Q),} \\[3pt]
\FitMath{u_f\ge0} \\[3pt]
\FitMath{(f\in D, f\ge0\text{ on }X).}\tag{CD6}
\end{gather*}

Each holds almost surely. Their exceptional sets therefore have a
single Borel, $\mathcal H$-measurable null union $N_0$. This is the
simultaneous exceptional set needed for the construction. No
intersection over all continuous or all Borel functions is taken.

\subsection*{Pointwise Riesz representation}

For $x\notin N_0$ define $L_x(f)=u_f(x)$ on $D$. It is rational linear,
positive on nonnegative elements of $D$, and $L_x(1)=1$. If rational
$r>\|f\|_\infty$, then $r\pm f\in D$ are nonnegative, so

\begin{gather*}
\FitMath{|L_x(f)|\le r.}
\end{gather*}

Letting rational $r$ decrease to $\|f\|_\infty$ proves
$|L_x(f)|\le\|f\|_\infty$. Uniform density consequently gives a unique
continuous extension to $C(X;\mathbb R)$. The extension is real linear:
approximate a real scalar by rationals and use the norm bound. It is
positive: if $f\ge0$ is continuous, choose $f_j\in D$ with
$\|f_j-f\|_\infty<1/j$; then $f_j+1/j\ge0$ and
$L_x(f_j)\ge-1/j$. It still takes the value one at $1$.

The Riesz representation theorem supplies a unique regular Borel
probability measure $K(x,\cdot)$ representing $L_x$. For $x\in N_0$
set $K(x,\cdot)=\delta_{x_*}$ at one fixed $x_*\in X$. Thus $K(x,\cdot)$
is a probability measure at \emph{every} point. Its values on the null
set are a choice, with no conditional interpretation required.

We check the probability-valued measurability rather than inferring it
from uncountably many pointwise representations. Enumerate $D$ as
$\{d_j\}$ and define

\begin{gather*}
\FitMath{I:\mathcal P(X)\longrightarrow\mathbb R^{\mathbb N},} \\[3pt]
\FitMath{I(\rho)=(\rho(d_j))_j.}\tag{CD7}
\end{gather*}

The map is continuous and injective, by uniform density and uniqueness
in Riesz representation. The weak probability space $\mathcal P(X)$
is metrized by
$\sum_j2^{-j}|\rho(d_j)-\rho'(d_j)|/(1+2\|d_j\|_\infty)$:
uniform density identifies convergence in this metric with weak
convergence. It is compact. Indeed, compactness follows by taking a
diagonal subsequence of integrals on $D$, extending the resulting
positive normalized functional as above, and applying Riesz again.
Hence $I$ is a homeomorphism onto its compact image in the Hausdorff
space $\mathbb R^{\mathbb N}$.

Every coordinate of $I\circ K$ equals $u_{d_j}$ outside $N_0$ and the
constant $d_j(x_*)$ on $N_0$. It is both Borel and
$\mathcal H$-measurable. The countable product map is therefore
measurable for both fields; composing with the inverse of $I$ proves
the asserted measurability of $K$. This also supplies an explicit
countable determining family for uniqueness.

\subsection*{All bounded Borel tests and uniqueness}

For $f\in D$, the equality $Kf=u_f$ holds outside $N_0$; thus (CD5)
holds with $Kf$ in place of $u_f$. Uniform approximation extends it to
every real continuous $f$, and $Kf$ is measurable as a uniform limit
of the corresponding measurable integrals.

Here are the remaining monotone-class details. The collection of
bounded real Borel functions $f$ for which $Kf$ is measurable for both
fields and

\begin{gather*}
\FitMath{\int_H Kf\,d\mu=\int_H f\,d\mu} \\[3pt]
\FitMath{\text{for every }H\in\mathcal H}\tag{CD8}
\end{gather*}

is a vector space containing $C(X;\mathbb R)$. It is closed under
uniformly bounded increasing limits, by monotone convergence for each
probability $K(x,\cdot)$ and for $\mu$; subtracting a constant reduces
to nonnegative functions. The bounded functional monotone-class
theorem therefore makes it all bounded Borel functions, since the
continuous algebra generates $\mathcal B(X)$. One can see the latter
directly from a countable dense family of distance functions and the
countable base of metric balls. Complex functions follow by their real
and imaginary parts. This proves (CD2), and indicators give (CD3).
In particular it also verifies directly that $K$ is a probability
kernel, including measurability of every Borel-set evaluation.

For uniqueness, apply uniqueness of conditional expectation to $Kd_j$
and $K'd_j$ for the countably many $d_j\in D$. Outside a single null
set all these integrals agree. Uniform density then gives equality on
every continuous function, and Riesz uniqueness gives (CD4). The null
set belongs to the completed Borel field and hence is contained in a
Borel null set, as asserted. Changing a map on a null set does not
change its probability-valued law. The law has barycenter $\mu$,
because (CD8) with $H=X$ gives

\begin{gather*}
\FitMath{\mu(f)=\int_{\mathcal P(X)}\rho(f)\,(K_*\mu)(d\rho)} \\[3pt]
\FitMath{(f\text{ bounded Borel}).}\tag{CD9}
\end{gather*}

The integrand is Borel on $\mathcal P(X)$: this is immediate for
continuous $f$, and the same bounded monotone-class argument proves it
for every bounded Borel $f$.\hfill$\square$

\subsection*{The completed invariant field and the exact null-set scope}

For a countable group of Borel, $\mu$-preserving transformations,
the field of events in $\overline{\mathcal B}_\mu$ invariant modulo
$\mu$ under every group element is a $\sigma$-field containing
$\mathcal N_\mu$. Thus it satisfies (CD1). In particular the theorem
applies to the completed translation-invariant field of a compact
metric lattice configuration space. The conditional measures are Borel
probabilities on that configuration space; no countable generation of
the completed invariant field is needed.

The uniqueness assertion is simultaneous for the resulting probability
measures. The assertion that $Kf$ is a conditional-expectation version
holds separately for each bounded Borel $f$ and, of course, for any
preselected countable collection on one common conull set. It does not
assert equality outside one universal exceptional set to an arbitrary
previously selected version for \emph{every} Borel function.

Nor does the theorem assert pointwise properness
$K(x,H)=1_H(x)$ simultaneously for all conditioning events. Even when
those events are Borel, such a demand is false for completed fields.
For example let $\mu$ be nonatomic and let
$\mathcal H=\{H\in\overline{\mathcal B}_\mu:
\mu(H)\in\{0,1\}\}$ be the null/co-null field.
The conditional kernel is $K(x,\cdot)=\mu$ almost surely.
If simultaneous properness held outside one null set, then at any
point $x$ of that conull set the conditioning event $\{x\}\in\mathcal H$
would require $K(x,\{x\})=1$, whereas $\mu(\{x\})=0$. This does not
contradict (CD2): for a fixed Borel $H\in\mathcal H$, it gives
$K(x,H)=1_H(x)$ almost surely, with an exceptional set that may depend
on $H$. Equation (CD3) is the corresponding integral identity. The theorem concerns the countably generated Borel target
$\mathcal B(X)$; no regular-conditional-law assertion for an arbitrary
non-countably-generated target field is being made.

% END MODULE chapters/compact_disintegration_v3.tex


% BEGIN MODULE chapters/phase_interface_v2.tex
\section{A phase construction from upstream lattice estimates}
\label{ch:phase-interface}

Under the lattice estimates specified below, we construct the canonical
phase measure from explicit probabilistic objects. The finite-spacing law G3 is proved in
Section~\ref{ch:plateau-proof}; it is used before any continuum component
is constructed. The all-state moment and
one-point comparison estimate G4 is derived in
Section~\ref{ch:all-state-source-transfer} from the ordinary torus
source obligation D1. Conditional probability measures are
constructed in Section~\ref{ch:compact-disintegration}. The argument here
deduces the canonical phase measure from those precise ingredients.

In this section a pure state means an extreme point of the translation-invariant
DLR class $\mathcal G_j$, equivalently an ergodic state for lattice translations.
Extremality in the larger class of all, possibly noninvariant, DLR states is not
used and is not identified with this property without an additional argument.

Put $a_j=L^{-K_j}\downarrow0$, with $L$ an integer, and let
$\Omega_j=\mathrm{SU}(2)^{\mathcal B_j}$ be the space of bond configurations
on $a_j\mathbb Z^4$. The set $\mathcal B_j$ contains one orientation of each
bond. Its countability makes $\Omega_j$ compact and metrizable: enumerate the
bonds and use the sum of $2^{-m}$ times bounded metrics on the group factors.
The Wilson specification $\gamma_W$ integrates finitely many bonds $W$ with
density proportional to
$\exp[-\beta_j\sum_{p:p\cap W\ne\varnothing}a_p]$.
The denominator is positive, and compactness of the integration space shows
that $\gamma_W F$ is continuous and local whenever $F$ is. Define
$\mathcal G_j$ to consist of the translation-invariant probability measures
with $\rho(\gamma_WF)=\rho(F)$ for every finite $W$ and continuous local $F$.
The equations are weakly closed, so $\mathcal G_j$ is compact and convex.
Its nonemptiness follows from the fixed-spacing torus-limit construction,
and a centering state $\nu_j\in\mathcal G_j$ is fixed.
G1 supplies the lattice symmetries and mirror positivity used in G3.

Let $\mathfrak T$ be the set of finite tuples of real $C^1$ compactly supported
functions with pairwise disjoint supports, including the empty tuple.
Concatenation is used only when its resulting supports are disjoint. For
$X=(f_1,\ldots,f_k)$ define the bounded continuous local function

\begin{gather*}
\FitMath{X_j=\prod_{i=1}^k
       [\mathcal O_j(f_i)-\nu_j(\mathcal O_j(f_i))],} \\[3pt]
\FitMath{\varnothing_j=1.}
\end{gather*}

Assume $\nu_j(X_j)\to\sigma(X)$ on $\mathfrak T$.
For support radii $R$ and separation $d$, put $b_1=C_1(R)$ and
$b_n=C_n(d;R)$ for $n\ge2$, and define

\begin{gather*}
\FitMath{M_k(d;R)=\sum_{\pi\in\operatorname{Part}(\{1,\ldots,k\})}
                         \prod_{J\in\pi}b_{|J|},} \\[3pt]
\FitMath{M_0=1.}
\end{gather*}

The ordinary moment--cumulant identity, applied under any $\rho\in\mathcal G_j$,
gives

\begin{gather*}
\FitMath{|\rho(X_j)|\le M_k(d;R)\prod_i\|f_i\|_\sharp.}
\end{gather*}

Indeed, a singleton block contributes the difference between the means in
$\rho$ and $\nu_j$, and a larger block contributes the corresponding cumulant
under $\rho$. This is why a bound on one-point differences is needed.
Telescoping a difference of two products gives equicontinuity with the same
constants whenever all supports remain in fixed separated containers.

\subsection{Decomposition at one spacing}

Let $\mathcal I$ denote the completed invariant $\sigma$-field for a fixed
$\nu\in\mathcal G_j$. Theorem~\ref{thm:compact-conditional-law} gives
a probability kernel $\omega\mapsto\nu^\omega$ satisfying
$\nu^\omega(F)=E_\nu[F\mid\mathcal I](\omega)$ for bounded Borel $F$.
The following details ensure that its law is a decomposition into the required
states, rather than into arbitrary probability measures.

First, a bounded invariant density preserves the DLR equations. To prove it,
approximate such a density $h$ in $L^1(\nu)$ by bounded continuous local
functions $h_n$. This approximation is available because cylinder functions
are uniformly dense in continuous functions on the compact product space,
and continuous functions are dense in $L^1$ of a Borel probability measure.
Translate $h_n$ so its support misses $W$. Invariance of $h$ and $\nu$ leaves
the approximation error unchanged, while
$\gamma_W(h_nF)=h_n\gamma_WF$ for the translated function. Passing to the
limit proves $\nu(h\gamma_WF)=\nu(hF)$.

Consequently $\nu^\omega(\gamma_WF)=\nu^\omega(F)$ almost surely.
Likewise $\nu^\omega(\tau_vF)=\nu^\omega(F)$ almost surely for every lattice
translation $v$. These statements can be imposed on one common full-measure
set: finite bond sets and translations form countable families, and continuous
functions on the compact configuration space have a countable dense family
of continuous cylinder functions. Supremum-norm continuity then extends the
equations to all continuous local functions. A monotone-class argument extends
the integral identities to bounded Borel functions. Thus
$\nu^\omega\in\mathcal G_j$ almost surely.

The law G3 proves that almost every such state is ergodic. Fix a continuous
local $A$. Applied to $\nu^\omega$ and then integrated, its far-translation
limit gives

\begin{gather*}
\FitMath{\int\operatorname{Pl}_{\nu^\omega}(A,A)\,d\nu(\omega)
 =\operatorname{Pl}_\nu(A,A)
 =\int|\nu^\omega(A)|^2\,d\nu(\omega).}
\end{gather*}

The first passage to the limit uses the bound $\|A\|_\infty^2$ and dominated
convergence. The difference of the integrands on the two sides is the
nonnegative variance of the invariant conditional expectation of $A$ in
$\nu^\omega$. It is therefore zero almost surely. Impose this on a countable
dense family of continuous local $A$. Density in $L^2(\nu^\omega)$ then says
that every invariant $L^2$ function is constant, which is ergodicity.

Write $\mathcal E_j$ for these ergodic states. They are exactly the extreme
points of $\mathcal G_j$. A nontrivial invariant event decomposes a state
into two DLR states by the invariant-density argument. Conversely, the
Radon--Nikodym density in a convex decomposition of an ergodic invariant state
is invariant, and hence constant. Moreover $\mathcal E_j$ is closed. If
$\rho_n\in\mathcal E_j$ tends weakly to $\rho$, G3 gives a decay estimate
uniform in $n$. Passing first to this weak limit and then to far translations
identifies $\operatorname{Pl}_\rho(A,A)$ with $|\rho(A)|^2$. The same dense-set
argument proves that $\rho$ is ergodic. Thus $\mathcal E_j$ is compact.

Let $D_j$ be the law of $\nu_j^\omega$ under $\nu_j$. It is supported on
$\mathcal E_j$, has barycenter $\nu_j$, and does not depend on versions of the
kernel. For this last assertion choose one countable determining family of
continuous cylinder functions. Conditional expectations are unique almost
surely for that entire family simultaneously; equality on the family implies
equality of the probability measures. This also establishes measurability of
the probability-valued map and of every evaluation used below.

Uniqueness among all decompositions supported on $\mathcal E_j$ requires one
more observation. For any such law with barycenter $\nu_j$, the uniform law G3
identifies the integral of a product of finitely many state evaluations with
successive far-translation limits of $\nu_j$. These moments are therefore the
same as those of $D_j$. Evaluations of continuous local functions are continuous
and separate measures in the compact space $\mathcal E_j$. Their polynomial
algebra is dense by Stone--Weierstrass, so equal mixed moments imply equal laws.

\subsection{Continuum evaluation coordinates}

Choose countable dense test sets as follows. Let $\mathcal Q$ contain all finite
unions of closed boxes with rational corners. The space of $C^1$ functions
vanishing off $Q\in\mathcal Q$, with its $C^1$ norm, is separable: embed it
by its value and four first derivatives in a finite product of separable
continuous-function spaces. Choose a countable dense set in each space and
enlarge their union by the countable group of $L$-adic translations and by
$W_4$. Let $\mathfrak T^\circ$ consist of the separated tuples from this union.
It is countable. Every separated tuple can be approximated by such tuples
with supports in fixed separated containers: cover each original support by
a finite union of sufficiently small rational boxes inside its chosen open
container and use the dense set for that union.

For each $X\in\mathfrak T^\circ$ choose once and for all a moment bound $M_X$
from the preceding estimate. Define the compact metrizable product

\begin{gather*}
\FitMath{\mathsf K=\prod_{X\in\mathfrak T^\circ}[-M_X,M_X],} \\[3pt]
\FitMath{e_j(\rho)=(\rho(X_j))_{X\in\mathfrak T^\circ},} \\[3pt]
\FitMath{\lambda_j=(e_j)_*D_j.}
\end{gather*}

The evaluation map is continuous for each fixed $j$. In particular
$\lambda_j$ is a well-defined probability measure on $\mathsf K$.

\begin{theorem}[Conditional canonical decomposition]
Under G1, the finite-spacing law G3, the estimates G4, and the stated
convergence of $\nu_j(X_j)$, the measures $\lambda_j$ converge weakly to one
probability measure $\lambda$. Put $\mathsf X=\operatorname{supp}\lambda$.
Every $\xi\in\mathsf X$ defines a unique functional $\sigma_\xi$ on
$\mathfrak T$, obtained as a limit along pure states
$\rho_j\in\mathcal E_j$ centered at $\nu_j$. The maps
$\xi\mapsto\sigma_\xi(X)$ are continuous and separate points. Furthermore

\begin{gather*}
\FitMath{\sigma(X)=\int_{\mathsf X}\sigma_\xi(X)\,d\lambda(\xi).}
\end{gather*}

Their joint moments are the successive far-translation limits of $\sigma$.
More precisely, for $q\ge2$,

\begin{gather*}
\FitMath{\int_{\mathsf X}\prod_{i=1}^q\sigma_\xi(X^i)\,d\lambda(\xi)
 =\lim_{|v_q|\to\infty}\cdots\lim_{|v_2|\to\infty}
   \sigma(X^1\cdot\tau_{v_2}X^2\cdots\tau_{v_q}X^q).}
\end{gather*}

The innermost limit is taken first. At each stage the remaining groups must
have disjoint supports; every sufficiently distant next translation is then
admitted. Each of these intermediate limits exists.
These properties determine the compact space, measure and functionals uniquely
among compact metrizable full-support representations with continuous
point-separating evaluations and the displayed mixed-moment property, up to
the homeomorphism preserving all evaluations. Each component clusters
with a bound $C(X,Y)|v|^{-2}$ for sufficiently distant separated products,
where $C(X,Y)$ and the distance threshold are uniform in $\xi$.
\end{theorem}

\begin{proof}
The crucial uniformity comes from reflection sizes rather than supremum norms
of $X_j$. Choose an integer physical half-side $\varpi$ so that all supports
under consideration lie in $(-\varpi+1,\varpi-1)^4$. This also leaves room
for the corners of plaquettes whose barycenters are in those supports. Then
$r_j=\varpi/a_j$ is an integer. Each expression defining
$\mathcal N_{\rho,r_j}(X_j)$ is the moment of two separated reflected copies
of $X$. G4 bounds it uniformly in $j$ and $\rho$ by a constant determined by
$M_{2|X|}$ and the test norms. For a displacement $v$ common to the lattices,

\begin{gather*}
\FitMath{\frac{r_j}{|v|_\infty/a_j-2r_j}
 =\frac{\varpi}{|v|_\infty-2\varpi}.}
\end{gather*}

Thus G3 gives, in every pure state, factorization of the two separated products
with error at most $C(X,Y)|v|^{-2}$, uniformly in the cutoff. The norms of the
individual representatives, which can diverge, do not enter this bound.

For finitely many coordinates $X^1,\ldots,X^q$, integrate this factorization
over $D_j$. Choose $v_q$ first, then $v_{q-1}$, down to $v_2$, with each
new vector sufficiently far from the already fixed union. In the factorization
at that stage the enclosing support box and its constant may depend on the
fixed outer translations, but not on $j$. Repeating the estimate and multiplying
errors only by the uniform individual moment bounds makes the total error
smaller than any prescribed $\epsilon$, uniformly in $j$. At that one fixed
finite configuration the remaining single moment converges by the hypothesis
on $\nu_j$. The mixed coordinate moments of $\lambda_j$ are therefore Cauchy.
Their precise iterated-limit representation will follow after the components
and their uniform clustering have been constructed.
Polynomials in finitely many coordinates separate points of $\mathsf K$ and
contain constants. By Stone--Weierstrass they are uniformly dense in its
continuous functions. Compactness therefore supplies a weak limit, and the
convergent mixed moments make that limit unique along the full sequence.

If $\xi$ belongs to its support, every neighborhood of $\xi$ has positive
$\lambda$ measure. Portmanteau implies that its preimage has positive $D_j$
measure for all sufficiently large $j$. Use a nested countable neighborhood
base and a diagonal choice to obtain $\rho_j\in\mathcal E_j$ with
$e_j(\rho_j)\to\xi$. The uniform telescoping estimate extends convergence
from $\mathfrak T^\circ$ to every tuple, giving $\sigma_\xi$. It also proves
independence from the choice of $\rho_j$ and continuity in $\xi$, by uniformly
approximating a given evaluation with dense-set evaluations. Multilinearity,
symmetry and the moment bounds pass to this limit.

Barycenters and higher mixed moments pass to the limit first on the countable
set and then on all tuples by the same uniform approximation. The evaluations
separate points because their countable subfamily consists of the coordinates
of $\mathsf K$. If another compact full-support representation has the same
mixed moments, embed both spaces by \emph{all} their tuple evaluations into
$\mathbb R^{\mathfrak T}$. Their images have compact union. Polynomials in
finitely many of these coordinates are dense in the continuous functions on
that union, again by Stone--Weierstrass. Equality of mixed moments therefore
gives equality of the two image measures. Each compact image equals the
support of its image measure. The images consequently coincide, and
injectivity and compactness identify the two spaces by a homeomorphism
preserving the measure and every evaluation. This argument does not presume
that a countable subfamily separates an arbitrary competing representation.

Finally pass the pure-state factorization estimate to the limit for $L$-adic
translations. Continuity of $C^1$ tests in the norm $\|\cdot\|_\sharp$ extends
it to real translations, preserving the uniform constants after a fixed
enlargement of the support margin. This proves the asserted clustering.
Finally hold $v_3,\ldots,v_q$ fixed and let $|v_2|$ tend to infinity in the
barycenter identity. Uniform component clustering separates $X^2$ from the
fixed union of all the other groups. The uniform moment bounds justify passing
this limit through the integral. Then let $|v_3|$ tend to infinity, and continue
in that order. This proves existence of every intermediate limit and the
displayed mixed-moment formula, with the innermost limit taken first.
\end{proof}

% END MODULE chapters/phase_interface_v2.tex


% BEGIN MODULE chapters/phase_vacuum_v3.tex
\section{The invariant Hilbert space of a canonical phase measure}
\label{ch:phase-vacuum}

The canonical decomposition is a statement about correlation forms. Its
relation to the zero-energy Hilbert space is a further implication, proved
here with the required density argument. No assertion that a linear span
of phase evaluations is automatically an algebra is needed.

Let $X$ be a compact metrizable phase space, let $\lambda$ be a Borel probability
measure of full support on $X$, and let $\sigma_\xi$ be its component
correlation forms. They act on finite separated tuples of compactly supported
test functions and their complex
linear span, including the empty tuple $1$. Assume that

\begin{gather*}
\FitMath{f_F(\xi):=\sigma_\xi(F)\in C(X),} \\[3pt]
\FitMath{\sigma(F)=\int_X f_F\,d\lambda,} \\[3pt]
\FitMath{f_1=1.}\tag{PV1}
\end{gather*}

Assume that the functions $f_F$ separate points of $X$, and that the
components preserve conjugation, time reflection, and translations.
Concatenation is used only when all its supports are separated.

The quantitative phase construction supplies uniform clustering: for
any fixed separated tuples $F,G$, once their translated supports are
disjoint,

\begin{gather*}
\FitMath{\sup_{\xi\in X}
 |\sigma_\xi(F\,\tau_{te_4}G)
                   -\sigma_\xi(F)\sigma_\xi(G)|\longrightarrow0} \\[3pt]
\FitMath{(t\longrightarrow\infty).}\tag{PV2}
\end{gather*}

Finite linear combinations inherit this property by summation. Uniform
separated moment estimates ensure that all functions displayed in
(PV1)--(PV2) are bounded, for each fixed tuple. The phases need not have
the same one-point mean.

Assume that the reflection-positive quotient has been constructed and
completed to a Hilbert space $\mathcal H$, with dense classes $[F]$ of
separated tuples whose compact supports lie strictly in $x_4>0$, and their linear combinations. Suppose
positive time translations induce a strongly continuous contraction
semigroup $e^{-tH}$ with a nonnegative self-adjoint generator $H$ and

\begin{gather*}
\FitMath{\langle[F],e^{-tH}[G]\rangle
      =\sigma(\overline{\theta F}\,\tau_{te_4}G),} \\[3pt]
\FitMath{t>0.}\tag{PV3}
\end{gather*}

These are the explicit Hilbert-space hypotheses; this section does not
derive them from qualitative convergence alone. Write
$P_0=1_{\{0\}}(H)$ for the zero-energy projection.

\begin{theorem}[The zero-energy representation]
Under (PV1)--(PV3), there is a unique unitary

\begin{gather*}
\FitMath{U:\ker H\longrightarrow L^2(X,\lambda)} \\[3pt]
\FitMath{\text{such that}} \\[3pt]
\FitMath{U P_0[F]=f_F.}\tag{PV4}
\end{gather*}

The vacuum vector $[1]$ maps to the constant function $1$. If $X$ has
exactly $m<\infty$ points, then $\dim\ker H=m$. If $X$ is infinite,
then $\ker H$ has countably infinite Hilbert dimension.
\end{theorem}

\begin{proof}
The spectral theorem and dominated convergence give
$e^{-tH}\to P_0$ strongly as $t\to\infty$. For fixed $F,G$, use
(PV1), reflection invariance, conjugation, and (PV2) in (PV3) to obtain

\begin{gather*}
\FitMath{\begin{aligned}
 \langle P_0[F],P_0[G]\rangle
  &=\lim_{t\to\infty}\int_X
      \sigma_\xi(\overline{\theta F}\,\tau_{te_4}G)
                                                      \,\lambda(d\xi)\\
  &=\int_X\overline{f_F(\xi)}f_G(\xi)\,\lambda(d\xi).
 \end{aligned}}\tag{PV5}
\end{gather*}

The limit is justified by uniform clustering, applied to the finite
linear combinations as well as to individual tuples. In particular,
if $P_0[F]=0$, then $f_F=0$ in $L^2$. Thus the rule in (PV4) is
well defined and isometric on the span of the displayed vectors.
That span is dense in $\ker H$, since the classes $[F]$ are dense
in $\mathcal H$ and $P_0$ is a bounded projection. Its extension is
an isometry with closed range. It remains to prove surjectivity.

Let $\mathcal V$ be the complex linear span in $C(X)$ of all the
phase evaluations $f_F$ with $F$ supported at positive times.
Any separated tuple can be moved into this half-space by a positive
time translation, which leaves its phase evaluations unchanged.
Consequently $\mathcal V$ still separates points of $X$. It contains
$1$ and is closed under complex conjugation.

For positive-time tuples $F,G$, the concatenation
$F\,\tau_{te_4}G$ is a positive-time separated tuple for all sufficiently
large $t$. By (PV2),

\begin{gather*}
\FitMath{f_{F\,\tau_{te_4}G}\longrightarrow f_F f_G} \\[3pt]
\FitMath{\text{uniformly on }X.}
\end{gather*}

It follows first for generators, and then for their finite linear
combinations, that products of elements of $\mathcal V$ belong to
its uniform closure. The uniform closure is itself closed under
products: approximate each of two factors uniformly by elements of
$\mathcal V$ and use boundedness of the approximants. It is therefore
a unital conjugation-closed subalgebra separating points. The complex
Stone--Weierstrass theorem makes that closure equal to $C(X)$.
Since continuous functions are dense in $L^2(X,\lambda)$, the image
in (PV4) is dense and closed, hence is the whole space.

The stated vector correspondence determines the unitary on a dense
set, proving uniqueness. It sends $[1]=P_0[1]$ to $1$, since
translation fixes the empty tuple.

For a finite phase space, full support gives positive mass to
each of its $m$ points, so $L^2$ has dimension $m$. For an infinite
compact metric space, $L^2$ is separable. One way to see this is to
take a countable uniformly dense subset of $C(X)$ and then its
rational complex linear span, which is dense in $L^2$. Its dimension
cannot be finite: for every integer $m$ choose $m$ distinct points
and disjoint open neighbourhoods of them. Full support gives each
neighbourhood positive measure, so their indicator functions are
$m$ nonzero mutually orthogonal vectors. The dimension is therefore
countably infinite. This proves all assertions.
\end{proof}

In particular, a single phase gives a one-dimensional zero-energy
space under these hypotheses. Several phases can produce several
vacua, even when each component has a unique vacuum. An infinite
phase set is recorded by the extended natural count $\infty$ in
dimension statements; its set-theoretic cardinality is not equated
with Hilbert dimension. Nothing in this identification gives a
strictly positive lower spectral bound on $(\ker H)^\perp$.

% END MODULE chapters/phase_vacuum_v3.tex


% BEGIN MODULE chapters/reconstruction_audit_chapter.tex
\section{Reconstruction of the action-density sector}
\label{ch:precise-reconstruction}
\subsection*{Reconstruction: hypotheses, consequences, and scope}

All growth constants below are positive, the factorial exponents are
nonnegative, and derivative-order coefficients are positive integers.
The D2 and H1 identifiers are retained for the conditional graded-moment
and diagonal-flat growth conclusions. G2 now denotes the derived scalar
rotation requirement under the explicit ordinary-source, mirror, cone and
ultraviolet hypotheses described below.

The passage from Euclidean correlation functions to a quantum field
requires several logically separate inputs. Compactness of separated
correlations gives an off-diagonal distribution. Reflection positivity
gives a Hilbert space. Uniform growth in the number of arguments is
needed for the corrected reconstruction theorem. Full Euclidean
covariance and clustering have to be established in their own right.
Keeping these steps separate prevents a qualitative convergence
statement from carrying quantitative information it does not contain.

\subsubsection*{The Euclidean test domain}

Write

\begin{gather*}
\FitMath{D_n=\{(x_1,\ldots,x_n)\in(\mathbb R^4)^n:
             x_i\ne x_j\text{ whenever }i\ne j\}.}
\end{gather*}

Let $\mathcal S_{\ne,n}$ be the closed subspace of the Schwartz space
$\mathcal S(\mathbb R^{4n})$ consisting of functions whose every derivative
vanishes on the complement of $D_n$. These functions are flat on all
coincidence diagonals. Let $\mathcal S_{+,n}$ consist of those elements
that are also flat outside the region in which all time coordinates
are positive. Set $\mathcal S_{\ne,0}=\mathcal S_{+,0}=\mathbb C$.
The time reflection is $\theta(x_1,x_2,x_3,x_4)
=(x_1,x_2,x_3,-x_4)$, and

\begin{gather*}
\FitMath{(\Theta\overline\varphi)(x_1,\ldots,x_n)
       =\overline{\varphi(\theta x_1,\ldots,\theta x_n)}.}
\end{gather*}

Symmetry of the functions makes reversal of the order of the arguments
immaterial in this scalar setting.

These are valid test domains for the historical
Osterwalder--Schrader reconstruction problem. They do not assert that
the Euclidean functions have a preferred extension to every Schwartz
test. A distribution supported on a coincidence diagonal vanishes on
$\mathcal S_{\ne,n}$. Such a distribution therefore cannot be recovered
from separated correlations. Any later use of contact terms, composite
operator identities at coincident points, or local Ward identities has
to specify the additional extension that it uses. This observation
does not obstruct reconstruction on the stated domain.

\paragraph{Density lemma.}
Finite sums of tensors $f_1\otimes\cdots\otimes f_n$, with smooth
compactly supported $f_i$ having pairwise disjoint supports, are dense
in $\mathcal S_{\ne,n}$. The corresponding assertion holds in
$\mathcal S_{+,n}$ when all supports lie at positive times.

The proof, including the simultaneous infinity, diagonal and positive-time cutoffs and the product Fourier approximation, is given in Section~\ref{ch:quantitative-flat-extension}.

The lemma is a statement about the flat subspace. It gives no density
claim in the full Schwartz space, where evaluation on the diagonal
is a continuous functional that vanishes on every such tensor.

\subsubsection*{What convergence preserves}

Let $\sigma_j$ be centred Euclidean moment functionals at finite
cutoffs. Suppose that for every separated smooth tuple $X$ the
numbers $\sigma_j(X)$ converge to $\sigma(X)$. Suppose also that each
$\sigma_j$ is reflection positive for finite families of internally
separated positive-time tuples. Different tuples in a family may
overlap. This last qualification is essential: reflection positivity
is a matrix condition, not a condition on one tuple at a time.

\paragraph{Closed positivity lemma.}
For positive-time tuples $X_1,\ldots,X_r$ and complex coefficients
$c_1,\ldots,c_r$,

\begin{gather*}
\FitMath{\sum_{i,k=1}^r\overline c_i c_k\,
              \sigma(\theta\overline{X_i}\cdot X_k)\ge0.}
\end{gather*}


\emph{Proof.}
Each $X_i$ is internally separated. Its conjugated reflected entries have
negative times, so $\theta\overline{X_i}\cdot X_k$ is a separated tuple for
every pair $(i,k)$. The displayed expression is the limit of the
corresponding finite sum for $\sigma_j$, all of which are
nonnegative.\hfill$\square$

Suppose, in addition, that the limiting functionals extend
continuously to $\mathcal S_{\ne,n}$. Reflection positivity then holds
on finite Schwartz sequences $\Phi=(\varphi_0,\ldots,\varphi_N)$ with
$\varphi_n\in\mathcal S_{+,n}$:

\begin{gather*}
\FitMath{\sum_{n,m=0}^N
 S_{n+m}(\Theta\overline\varphi_n\otimes\varphi_m)\ge0.}
\end{gather*}

Indeed, the mixed tensor belongs to $\mathcal S_{\ne,n+m}$.
Coincidences within one factor are killed by its flatness; a
coincidence across the two factors has time coordinate zero or lies
outside the allowed half-space of one factor, so is killed there.
Approximate each $\varphi_n$ by the tensors in the density lemma,
using a seminorm strong enough for all the finitely many mixed
terms, and pass to the limit. Permutation symmetry and invariance
under any fixed Euclidean isometry pass in exactly the same way
once they hold on separated tensors.

There is one further step when translations are initially known
only for vectors in a dense subgroup, for example the union of
nested lattice translation groups. Given an arbitrary vector $v$,
choose subgroup vectors $v_k\to v$. Translation is continuous on
the smooth tests, and the local bounds permit passage from
$S_n(\tau_{v_k}\varphi)=S_n(\varphi)$ to
$S_n(\tau_v\varphi)=S_n(\varphi)$. This argument does not generate
rotations from the finite hypercubic group.

\subsubsection*{The quantitative reconstruction hypothesis}

Fix the seminorms

\begin{gather*}
\FitMath{|\varphi|_k=
 \max_{|\alpha|\le k}\sup_x
       (1+|x|^2)^{k/2}|\partial^\alpha\varphi(x)|.}
\end{gather*}

For reconstruction, assume a sequence of symmetric functionals
$S_n\in\mathcal S_{\ne,n}'$, with $S_0=1$, and constants
$A,B,k_0,p<\infty$, independent of $n$, such that

\begin{gather*}
\FitMath{|S_n(\varphi)|\le A B^n(n!)^p|\varphi|_{k_0n}.}\tag{RC1}
\end{gather*}

Assume full Euclidean invariance and reflection positivity on the
domains just described. Assume also spatial clustering on compact
positive-time sequences: for every such finite pair $F,G$ and every
nonzero spatial vector $v$,
\[
 \sigma(\theta\overline F\cdot\tau_{\lambda(v,0)}G)
       \longrightarrow\overline{\sigma(F)}\sigma(G)
                 \qquad(\lambda\longrightarrow\infty).
\]
Here each homogeneous term of a sequence is a finite sum of internally
separated smooth compact tuples. The Hilbert-space approximation below
extends this cluster identity to the full positive-time test domain.

The imported result is Theorem $E'\to R'$ of
\cite[\S IV.1, p.~287]{OS75}. Its condition $E0'$ is the bound by a
single Schwartz seminorm of order $sn$ and constants of factorial
growth, on the diagonal-flat space. This is exactly the form of
(RC1): take an integer $s\ge\max(1,k_0)$ and
$\sigma_n=A B^n(n!)^p\le A e^B(n!)^{p+1}$.
Euclidean invariance, the full-sequence reflection-positive form,
permutation symmetry and the cluster identity supply its remaining
conditions $E1$--$E4$, respectively. Thus the theorem gives a Hermitian
scalar Wightman field with positive energy, locality, Poincar\'e covariance
and a unique vacuum in its cyclic representation. Its analytic-continuation
proof is imported; the interacting estimates needed to obtain a hierarchy
with these properties remain separate premises.

Condition (RC1) must not be replaced by the weaker assertion that
every $S_n$ is separately continuous with unspecified constants.
The growth of the distribution order and the constants in $n$
matters to the existence of tempered boundary values in the
analytic continuation. Likewise, a local off-diagonal kernel
theorem alone does not establish (RC1).

\paragraph{The explicit rooted-source and dyadic extension route.}
There is now a quantitative route to (RC1), rather than an unspecified
fixed-$n$ Whitney estimate. Theorem~\ref{thm:graded-rooted-source} derives
moments directly from the ordinary source obligation D1 and the proved
G1 chessboard domination, using the actual domain geometry in (GS1)--(GS3).
At scale $a_s=L^{-s}$, domains are unions of $Ma_s$ cubes, indexed by
those unions, and their source dependence occupies only a fixed number of
neighboring cube layers. Their size is the minimum cubical edge-tree
length meeting all cubes, with source-energy decay above its rooted
entropy rate. Here $M=L^{m_0}$, and the cofinal tori $\ell=2L^m$,
$m\ge m_0$, tile every actual scale with even cube counts. No scale is
discarded for a failure to tile. The weak abstract (ST6) count alone does
not furnish the sharper rooted and global counts needed for this proof.

For real compact $C^1$ tests, the theorem gives one fixed finite integer
$k$, $B(R)=C(R+1)^4$, and the two-state centered bound

\begin{gather*}
\FitMath{\left|\rho\prod_{i=1}^n
   [\mathcal O_a(f_i)-\nu(\mathcal O_a(f_i))]\right|
 \le B(R)^n(n!)^k\max\{1,d^{-5n}\}
                                 \prod_i\|f_i\|_\sharp.}\tag{RC2}
\end{gather*}

Each support lies in its own radius-$R$ ball, placed anywhere; the
separation is $d$. The same factorial exponent bounds higher cumulants.
Uniformity lets (RC2) pass to the separated continuum limits and their
component limits. This is requirement D2 conditional on the actual D1
source representation and G1, not a new independent analytic estimate.

Theorem~\ref{thm:quantitative-flat-extension} uses only the unit-radius
version, with the fixed equivalence of the sharp and $C^1$ norms absorbed
in $B(1)$. Its product Fourier construction uses $8n$ derivatives;
Gevrey-two pair/grid cutoffs cost at most $C^n(n!)^{32}$.
Dyadic Taylor flatness to order $10n+2$ pays the local $h^{-6n}$ cost
and the $h^{-4n}$ configuration grid count, leaving a summable dyadic
square. It yields (RC1) with

\begin{gather*}
\FitMath{k_0=20,} \\[3pt]
\FitMath{p=k+32.}
\end{gather*}

The proof includes independence of tensor presentations, overlap gluing,
cutoffs at infinity and every pairwise diagonal, and the positive-time
boundary. Its constrained tensor density supplies the full-sequence
reflection-positive form. Thus H1 is a proved implication of D2 here;
the ordinary analytic RG source estimates in D1 are still assumptions.
Neither separate continuity with arbitrary order constants nor weak
(ST6) counting alone may be substituted for this GS/QF route.

\subsubsection*{Hilbert space and clustering}

Section~\ref{ch:euclidean-hilbert-bridge} proves the quotient construction
on compact positive-time tuples from their invariant reflection-positive
moments and position-uniform per-ball bounds. Fixed separation of the
reflected supports bounds the translated orbits before iterated
Cauchy--Schwarz proves contraction. The resulting time semigroup is
strongly continuous, positive and self-adjoint, and spatial translations
are strongly continuous unitaries. The stated lattice approximation also
gives the cube spectral cone. The GS/QF growth and density argument then
identifies its completion with the full positive-time Schwartz sequence
quotient. These deductions retain their explicit hypotheses; positivity
alone is not being substituted for the continuity and growth checks.

For completeness, clustering on a dense class of vectors extends
to the whole Hilbert space. Let $U(v)$ denote the spatial unitary
translations, and suppose that for dense vectors $\xi,\eta$,

\begin{gather*}
\FitMath{\langle\xi,U(v)\eta\rangle\longrightarrow
       \langle\xi,\Omega\rangle\langle\Omega,\eta\rangle} \\[3pt]
\FitMath{(|v|\to\infty).}
\end{gather*}

Choose dense approximants $\xi_k\to\xi$ and $\eta_k\to\eta$.
Unitarity bounds the difference of the matrix elements, uniformly
in $v$, by

\begin{gather*}
\FitMath{\|\xi-\xi_k\|\,\|\eta\|
       +\|\xi_k\|\,\|\eta-\eta_k\|.}
\end{gather*}

The vacuum matrix elements obey the same estimate. First let
$|v|\to\infty$ at fixed $k$, then let $k\to\infty$.
The argument proves the qualitative cluster limit. It does not
preserve a numerical decay rate on arbitrary vectors without
further quantitative control of the approximations.

\subsubsection*{Phases and full rotational covariance}

A symmetry of a mixture need not be a symmetry of each component.
For example, the uniform probability measure on a sphere is
rotation invariant, while its decomposition into point masses is
permuted by rotations. Therefore an invariant Euclidean limit
and an ergodic decomposition do not, by themselves, supply an
invariant component. A construction that reconstructs every phase
must prove covariance and positivity for those phases, or realize
each phase as a limit of cutoff states for which the properties
and uniform estimates have already been established. The phase-interface
construction provides this realization for the G1 mirror properties and
the D1/GS/QF moment bounds. The state-uniform version of the two-point
extraction must also hold for every approximating state. Under that
ultraviolet input, the angular chain proves covariance for each realized
component. One first subtracts its own mean $c\int f$; the finite subset
expansion and convergent cutoff one-point differences identify the
self-centered limit, and the same expansion restores the mean afterwards.
It is not inferred merely from a symmetric mixture.

Similarly, invariance under the hypercubic group, coordinate
reflection positivity, and translation invariance do not by
themselves imply invariance under $\mathrm O(4)$. The finite group
is not dense in the rotation group. Any route through diagonal
mirrors, a sharp energy--momentum cone, bounded insertions, or
the exclusion of nonconstant angular harmonics needs an additional
theorem. The following chain now supplies that theorem conditionally,
before the external reconstruction theorem is invoked.

Theorem~\ref{thm:all-orders-angular-closure} completes the scalar
covariance implication. Its preceding sections prove, in order, the
uniform Hilbert insertion bound, the analytic $L^1$ density, the product
kernels and continued Gram identities, and the ultraviolet elimination
of the two-point harmonics. Polygon divisibility and common coordinate
cuts then remove all higher harmonics and exceptional distributional
support. The proof works with finite support boxes; density of generic
point configurations alone would not exclude a singular part.

The ultraviolet premise is exactly (TU10) on the full positive-shape
class with $\vartheta_*=1/12000$, or a broader class, uniformly in the
approximating states and positions. The dependency catalogue proves its
implication from F2--F3 by the separated-kernel bound and logarithmic
coupling decay. A narrower shape class or a torus-only law cannot be
substituted. The resulting G2 deduction applies to periodic-state limits
and their realized components carrying the mirror and exact cone
interface; arbitrary growing cubic-torus limits require that interface
to be justified separately. Typed Ward, distant-shell flux and matched-run
estimates remain scoped to their separate stronger routes.

The actual Gaussian, history and finite-source estimates and the
interacting extraction are now developed as proof outlines in
Sections~\ref{sec:actual-history-closure},
\ref{sec:actual-ordinary-source-proof},
\ref{ym4:new-clean-box}--\ref{ym4:new-gaussian-segment} and
\ref{sec:finite-reference-normalizer}, with the Gaussian and
substitution estimates cited there. The OS
analytic-continuation theorem is still imported,
and the identification with full Yang--Mills observables remains separate.

\subsubsection*{What the reconstructed field establishes}

When the limiting observable is the centred scalar action density,
the theorem above constructs its scalar observable sector. A
nonzero two-point function on separated tests excludes a field
that is a multiple of the identity. A nonzero higher cumulant
excludes Gaussianity of that scalar field. Neither conclusion
alone identifies a full Yang--Mills observable algebra. In
particular, a Wick square of a free Gaussian field already has
nonzero higher cumulants. Establishing a Yang--Mills theory in a
stronger sense requires the stated additional observables,
identities, and compatibility relations that define that sense.

Nor does polynomial clustering establish a mass gap. The positive
measure $E e^{-E}\,dE$ on $[0,\infty)$ has Laplace transform

\begin{gather*}
\FitMath{\int_0^\infty E e^{-E}e^{-tE}\,dE=(1+t)^{-2}.}
\end{gather*}

Its support reaches zero, although its transform decays
quadratically. A positive spectral gap requires an estimate that
excludes spectral support in an entire interval above the vacuum.
That question remains separate from the existence and
reconstruction of the scalar sector described here.

% END MODULE chapters/reconstruction_audit_chapter.tex

\appendix

% BEGIN MODULE chapters/comparison_capsule_v7.tex
\section{Coupling coordinates and optional comparison implications}
\label{ch:comparison-capsule-v7}
\label{candidate:power_counting}
\label{candidate:uniqueness}
\label{candidate:rotations}

This section preserves the conditional targets of the coupling, matched-run
and Ward arguments. It supplies their elementary analytic implications
without treating a proposed interacting estimate as a proof. The
pointed families, companions, physical Hessian extractor, midpoint tents,
list action and radius weights are defined in Section~\ref{ch:rg-objects-v2}.
The pointed cancellation, nonlinear window estimate and finite formation
sections have their stated background and domain hypotheses. None computes
the one-loop trace or proves a matched-cutoff or typed Ward estimate.

\subsection*{Vacuum accounting and the stronger coupling inputs}

Write \(P_\circ=(\mathbf1,0)\). With \(E_{(0)}=0\), the accumulated
vacuum normalization is defined by
\begin{equation}
\begin{aligned}
 E_{(k+1)}&=E_{(k)}+\mathsf e_k,\qquad E=E_{(K)},\\
 \mathsf e_k&=\mathfrak n_k+
 \sum_X\left[\sum_{z\in X}E^{(k+1)}(X,P_\circ,z)
                         +R^{(k+1)}(X,P_\circ)\right].
\end{aligned}\tag{CP1}
\end{equation}
Here \(\mathfrak n_k\) is the assigned field-independent step number.
For \(\tau=x/\theta\), its Gaussian scaling part is
\(-3|\mathfrak X_k|\log\tau/2\); \(\mathfrak X_k\) is the set of
non-pivot bonds, with three real color coordinates per bond. Haar and
gauge-fixing counts are included in that assigned number.
The family in (CP1) contains the unnormalized local Gaussian numbers:
\(\mathbf N^\circ=\mathbf N+\mathbf N^\circ(P_\circ)\).
Those numbers must not also be added to \(\mathfrak n_k\).
Equivalently, using \(\mathbf N\) requires adding the removed local
numbers once separately. This convention counts every constant once.
At fixed calibration, \(\mathfrak n_k,\mathbf N^\circ,\mathbf G_\tau\)
have zero source derivative. Source energies are the differences of
the displayed \(E\)- and \(R\)-values between sourced and unsourced runs.
More explicitly, the localized source energy of level \(j\) is
\[
 e_j(X;w)=\sum_{z\in X}\{E^{(j),w}(X,P_\circ,z)-E^{(j),0}(X,P_\circ,z)\}
          +R^{(j),w}(X,P_\circ)-R^{(j),0}(X,P_\circ).
\]
The source vacuum exponent is the sum of these values over levels and
local domains. This definition introduces no additional arbitrary
normalizer. The zero source derivative of the Gaussian pieces refers
to the frozen potential presentation. Moving the source into the
Gaussian action weight changes the determinant and is another
presentation, whose numbers must be compared by its own sum rule.

Requirement B3 retains an actual one-loop input:
\begin{equation}
 b_i=b_i^0+f_i+q_i,\qquad
 c_0=\frac{11N^2}{12\pi^2}\log L,\qquad
 \left|\sum_{i=1}^r(b_i^0-c_0)\right|\le C_0 .
 \tag{CP2}
\end{equation}
Here \(N=2\), and \(b_i^0\) is independent of the bare coupling and
admissible torus. A summable estimate
\(|b_i^0-c_0|\le CL^{-\sigma i}\), \(\sigma>0\), suffices.
The asserted normalization requires the actual lattice determinant/trace
calculation and its zero-mode convention. Gaussian integration alone
does not establish it.

Some later slope conclusions require more than (CP2). Fix a calibration
\(\theta_0,\ldots,\theta_K\), independent of the varied bare coupling.
Put \(\xi_j=\theta_j-x_j\). Require the actual unmarked \(f_i\) and
marked \(q_i\), for \(1\le i\le K\), to be jointly holomorphic on the
independent earlier-variable polydiscs with radii
\begin{equation}
\begin{aligned}
 r_j^{\rm u}&=c\theta_j,&
 \sup|f_i|&\le C\theta_i^{-1/2}(\log\theta_i)^p,\\
 r_j^{\rm m}&=c\theta_j(\log\theta_j)^{-C_*},&
 \sup|q_i|&\le C\theta_i^{-p_R},
 \qquad 0\le j<i .
\end{aligned}\tag{CP3}
\end{equation}
The constants \(p,C_*\ge0\), \(p_R>0\), and \(c>0\) are fixed.
The actual runs must remain in the half polydiscs, with common
localization, counterterms and admissible domains.
These are stronger actual coupling inputs, including the marked tube.
They are not inferred from holomorphy of the bare partition function,
nor from a running-radius activity bound followed by two field
derivatives. The marked exponent is often \(p_R=(\kappa_0-4)/2\),
but any later use must retain the actual proved loss and tube radius.
All norms and source Taylor expansions here use fixed declared domains;
shrinking the field domains while differentiating would be a different
claim. Compatibility of the companion and earlier cutoff prefix
also remains part of the construction.

\subsection*{A dimension-free polydisc lemma and its precise consequence}

\begin{lemma}[The bounded-polydisc estimate, code \(1e.3.7\)]
\label{lem:polydisc-dimension-free-v7}
If \(F\) is holomorphic on \(\prod_{j=1}^dD(0,r_j)\) and
\(|F|\le B\), then
\begin{equation}
 \sum_jr_j\left(1-\frac{|\xi_j|^2}{r_j^2}\right)
                   |\partial_jF(\xi)|\le B .
 \tag{CP4}
\end{equation}
On the half polydisc this implies
\begin{equation}
\begin{aligned}
 \left|\sum_jv_j\partial_jF(\xi)\right|
       &\le\tfrac43 B\max_j|v_j|/r_j,\\
 |F(\xi)-F(\xi')|&\le\tfrac43B
                         \max_j|\xi_j-\xi'_j|/r_j .
\end{aligned}\tag{CP5}
\end{equation}
No constant depends on \(d\).
\end{lemma}
\begin{proof}
The case \(B=0\) is immediate. Set \(a_j=\xi_j/r_j\) and
\(\phi_a(z)=(a+z)/(1+\bar a z)\). Choose unit complex numbers
\(\epsilon_j\) so the terms in the derivative at zero of
\[
 H(z)=F\big(r_j\phi_{a_j}(\epsilon_jz)\big)_{j=1}^d
\]
have the same argument. Since \(\phi_a'(0)=1-|a|^2\), the modulus
of that derivative is the left side of (CP4).
The one-variable Schwarz lemma, applied after a disc automorphism
sending \(H(0)/B\) to zero, gives
\(|H'(0)|\le B(1-|H(0)/B|^2)\le B\).
The Schwarz lemma itself follows by applying the maximum principle
to a holomorphic function vanishing at zero divided by \(z\), first
on smaller circles and then letting their radius increase to one.
An attained boundary value makes \(H\) constant.
On the half polydisc, (CP4) gives
\(\sum_jr_j|\partial_jF|\le4B/3\).
This proves the first estimate of (CP5); integration along the segment
from \(\xi\) to \(\xi'\), which stays in that half polydisc, proves
the second.
\end{proof}

Assume \(\theta_j\ge\theta_i\) for \(j<i\), with all centers sufficiently
large that \(t/(\log t)^{C_*}\) is increasing.
Applying the lemma separately in (CP3) bounds the absolute row sum by
\begin{equation}
 a_i:=\sum_{j<i}|\partial_{\xi_j}(f_i+q_i)|
 \le C\left[
 \theta_i^{-3/2}(\log\theta_i)^p+
 \theta_i^{-p_R-1}(\log\theta_i)^{C_*}\right].
 \tag{CP6}
\end{equation}
Suppose also that the calibration has backward linear drift
\(\theta_i\ge h+c_1(K-i)\), \(c_1>0\).
Comparison with integrals gives
\[
 A:=\sum_{i=1}^Ka_i
 \le C\left[h^{-1/2}(\log h)^p+
                       h^{-p_R}(\log h)^{C_*}\right]\longrightarrow0.
\]
Indeed substitute \(t=hu\) in each integral; the integrals of
\(u^{-1-\delta}(1+|\log u|)^r\) are finite for \(\delta>0\).
A drift with a fixed additive error is equivalent after increasing \(h\).

For the actual differentiable recursion
\(x_0(\beta)=\beta\),
\(x_i(\beta)=\beta-\sum_{r\le i}b_r(\beta)\),
put \(M_i=\max_{j\le i}|x_j'|\). Calibration is held fixed, so
\(|b_i'|\le a_iM_{i-1}\) and
\(M_i\le(1+a_i)M_{i-1}\). Consequently
\begin{equation}
 M_K\le\prod_i(1+a_i)\le e^A,\qquad
 \sum_i|b_i'|\le\prod_i(1+a_i)-1\le e^A-1 .
 \tag{CP7}
\end{equation}
The middle estimate telescopes the successive products.
In particular
\[
 |x_K'(\beta)-1|\le\sum_i|b_i'|
 \le e^A-1
 =O\!\left(h^{-1/2}(\log h)^p+h^{-p_R}(\log h)^{C_*}\right).
\]
On witness prefixes with terminal \(h\to\infty\), under the same
uniformly admitted hypotheses, this makes the constant-source
normalization tend to one with the displayed error. Identifying
that prefix and derivative with the actual F3 witness remains
an input; the estimate does not establish that identification.
At weak coupling \(A\le\log(4/3)\), hence the last sum is at most
\(1/3\). Constant sources varying \(\beta\) to \(\beta-\zeta\)
have the same absolute slopes. This supplies the slope premise in
(LB12) under (CP3) and the admitted drift/domain conditions.
It does not prove those conditions.
Reality gives \(2/3\le x_K'(\beta)\le4/3\) on an admitted real
interval. Integrating the derivative proves the uniform terminal
bi-Lipschitz bound. If the same row estimates hold on a whole complex
offset disc, write \(\Xi_K(t)=h-t+R_K(t)\); then
\(R_K(0)=0\) and \(\sup|R_K'|\le1/3\).
One can construct the matching inverse without any global uniqueness
assertion. For \( |t|<\varrho\), solve
\[
 s=h-\Xi_K(t)+R_{K+1}(s)
\]
on the closed disc \( |s|\le2\varrho\), assuming the functions are
defined on a neighborhood of that disc. Its right side maps the disc
strictly inside itself and has Lipschitz constant at most \(1/3\).
Iteration converges geometrically to its unique solution, uniformly on
smaller \(t\)-discs. The iterates are holomorphic, so the solution
is holomorphic; reality gives a real solution at real \(t\).
The bounds \( |\Xi_K(t)-h|\le4|t|/3\) and
\( |\Xi_{K+1}(s)-h|\ge2|s|/3\) give
\( |\mathfrak n_K(t)|\le2|t|\).
This supplies the coordinate part of E2 when its actual full-disc
premises are available, rather than proving the existence of those
premises. Nor does a polydisc supremum
give forgetting in \(i-j\): \(F_i(\xi)=B\xi_0/r_0\) is bounded
by \(B\) for every \(i\), with unchanged first derivative.

\subsection*{The fixed-torus coordinate and the matched limit}

Fix a torus \(T_\ell\), \(h=g^{-2}\), and a separated tuple of
compact Lipschitz tests, with the fixed norm \(\|\cdot\|_\sharp\)
of the source convention. Every bound below includes their product
norm; its constant may depend on order and positive separation.
Let \(\beta_K^\natural\) be its torus-independent tuned bare center
and write depth as \([s]=K-k\). Define
\begin{equation}
\begin{aligned}
 \mathsf J&=[-\varrho/4,\varrho/4],&
 \Xi_K(t)&=x_K^\natural(\beta_K^\natural-t),\\
 F_{K,n}(t;\mathbf f)
   &=S^T_{n;K,\ell,\beta_K^\natural-t}(\mathbf f),&
 \mathfrak n_K&=\Xi_{K+1}^{-1}\circ\Xi_K .
\end{aligned}\tag{CP8}
\end{equation}
The inverse is the specified holomorphic matching branch.
For obligations E2--E3 we must prove its existence on a common
neighborhood, reality on the real axis, uniform upper and lower real
Lipschitz bounds for \(\Xi_K\), \(\Xi_K(0)=h\), and
\(|\Xi_K(t)-h+t|\le|t|/3\).
On a fixed smaller complex neighborhood of \(\mathsf J\), we must prove
\begin{equation}
\begin{aligned}
 |\mathfrak n_K-t|+|\mathfrak n_K'-1|
       +|\Xi_{K+1}-\Xi_K|&\le C\eta_K,\\
 |F_{K,n}(t)-\mathfrak n_K'(t)^n
                    F_{K+1,n}(\mathfrak n_K(t))|
       &\le C_{n,\ell,\mathbf f}\epsilon_K,\qquad
 \sum_K(\eta_K+\epsilon_K)<\infty .
\end{aligned}\tag{CP9}
\end{equation}
The \(F_{K,n}\) must be uniformly bounded and holomorphic on a
slightly wider strip, for example
\(\mathcal N_\ell=\{|\operatorname{Re}t|<\varrho/2,\
|\operatorname{Im}t|<\upsilon_\ell\}\), where
\(\upsilon_\ell=c\varrho e^{-C(g)\ell^4}>0\).
Matching paths stay in that strip after shrinking it.
Constants and cutoff floors may depend on this fixed \(\ell\).
Often \(\eta_K=\Upsilon_K\le
CK^{-3/2}(\log(2+K))^{p+1}\); a separate remainder error
\(\Pi_K\) is part of \(\epsilon_K\), not a polarization tensor.

\begin{proposition}[The conditional terminal-coupling limit]
\label{thm:terminal-coupling-limit-v7}
These inputs imply uniform convergence \(F_{K,n}\to S^{[t]}_{\ell,n}\)
on the smaller strip and \(\Xi_K\to\Xi_\infty\).
The latter is a real bi-Lipschitz homeomorphism from \(\mathsf J\)
onto \(I\supset[h-\varrho/6,h+\varrho/6]\). Set
\[
 S^{(\ell,x)}=S^{[\Xi_\infty^{-1}(x)]}_\ell,\qquad x\in I .
\]
If \(t_K\in\mathsf J\) and \(\Xi_K(t_K)\to x\), then
\(F_{K,n}(t_K)\to S^{(\ell,x)}_n\), without a subsequence.
\end{proposition}
\begin{proof}
Cauchy's estimate on the wider strip bounds \(F_{K,n}'\) uniformly.
Thus replacing \(\mathfrak n_K(t)\) by \(t\) in (CP9), and its
derivative factor by one, costs \(C\eta_K\). The resulting increments
are summable, uniformly on the smaller strip. Their limit is
holomorphic and continuous on \(\mathsf J\). The coordinate increments
are summable too. Its real Lipschitz bounds pass to the limit;
the endpoint bounds give the asserted interval.
Uniform coordinate convergence gives \(\Xi_\infty(t_K)\to x\);
the lower Lipschitz bound forces
\(t_K\to\Xi_\infty^{-1}(x)\). Uniform convergence of \(F_K\)
and continuity of its limit finish the proof.
\end{proof}

This preserves target \(0(d)\) under its comparison premises.
It proves no infinite-volume uniqueness.
The depth coordinate
\(\Xi_{\infty,s}(t)=\lim_Kx_{K-s}^\natural(\beta_K^\natural-t)\)
requires the corresponding fixed-depth comparison. If its envelope is
\[
 \lambda s-C\le\Xi_{\infty,s}(t)-\Xi_{\infty,0}(t)
                         \le\Lambda s+C,
\]
then \(\Xi_{\infty,s}(t)^{-1/2}\) decreases in order as
\(s^{-1/2}\); at length \(L^{-s}\) this is logarithmic running.
The reference value \(g_{[s]}=\theta_{[s]}^{-1/2}\) is a calibration,
not this arbitrary member's running coupling. Common-prefix matching
must also control
\(\hat t_{K,s}=\beta_{K-s}^\natural(g_{[s]})-\beta_K^\natural(g)\).
If it converges to \(\hat t_s\), exact finite dilation passes to
\(S^{[t]}_\ell[g]=S^{[t+\hat t_s]}_{L^s\ell}[g_{[s]}]
(\mathbf f^{[s]})\), wherever both offsets are in their windows.
Obligation E4 separately compares self-calibrated and frozen runs,
places their first-passage offsets in \(\mathsf J\), and supplies
\(|K(\beta_K^\natural-t,g)-K|\le C(\varrho)\).
Neither assertion follows from the new finite source derivative.
Exact lattice dilation uses \(f^{[b]}(y)=f(L^{-b}y)\), without
an observable prefactor, and \(T_\ell\mapsto T_{L^b\ell}\).
Conditional offset limits and depth shifts pass through that identity
only when the offsets converge and both torus floors hold.
The change-of-reference claims retain these matching premises.
The bare tuned-increment target must also be retained explicitly:
\begin{equation}
 \beta_{K+1}^\natural(g)-\beta_K^\natural(g)\longrightarrow c_0,
 \qquad
 \beta_{K+b}^\natural(g)-\beta_K^\natural(g)\longrightarrow bc_0
 \quad(b\in\mathbb Z\text{ fixed}).\tag{CP9a}
\end{equation}
The second assertion follows by finite telescoping from the first.
Here is a conditional proof preserving the stronger comparison route.
Write \(x^{[i]}_{[s]}\) for the tuned depth-\(s\) coupling in the run
of length \(i\), with \(x^{[i]}_{[0]}=h\) and
\(x^{[i]}_{[i]}=\beta_i^\natural\). Require
\(x^{[i]}_{[v]}-x^{[i]}_{[u]}\), for \(u<v\), to lie in
\([\lambda(v-u)-C,\Lambda(v-u)+C]\), and
\[
 |x^{[i+1]}_{[s]}-x^{[i]}_{[s]}|
 \le C\big(L^{-c(i-s)}+s\Pi_i\big),\qquad
 \sup_{s\ge1}s\sum_{i\ge s}\Pi_i<\infty .
\]
These constants and the last bound are fixed before \(\varrho\).
Comparing two bare centers at the same depth \(K'\), then telescoping
in the run length, bounds their extra discrepancy by
\(C\sum_{i\ge K'}L^{-c(i-K')}+CK'\sum_{i\ge K'}\Pi_i\le C'\).
The spacing bounds therefore give a uniform bound \(A\) on
\(\alpha_K=\beta_{K+1}^\natural-\beta_K^\natural\).
Choose \(\varrho>20A\).

Each cluster value \(\alpha\) of \(\alpha_K\) gives
\(S^{[t]}_\ell=S^{[t-\alpha]}_{L\ell}(\mathbf f^{[1]})\)
by exact finite dilation and uniform moving-offset convergence.
Two cluster values \(\alpha^-<\alpha^+\) give an offset period
\(\omega=\alpha^+-\alpha^-\) on an open interval. Holomorphy in the
offset extends that identity throughout the overlapping strips.
Repeated blocking can then restore every offset modulo \(\omega\)
into \([0,\omega)\); all intermediate offsets stay in \(\mathsf J\)
by the choice of \(\varrho\).
Require, additionally, one fixed positive shape \(\Psi\), a floor
\(\ell_0\), and \(r_0>0\), such that
\[
 0<\mathfrak m(r)\le S^{[t]}_{\ell,2}(\Psi_r)\le\mathfrak M(r),
 \quad \mathfrak M(r)\longrightarrow0,\qquad
 t\in\mathsf J,\ \ell\ge\ell_0,\ 0<r\le r_0 .
\]
Normalization by the two fixed supremum norms is understood.
The required uniformity in offset, shape placement and arbitrarily
large tori is part of the unverified tuned two-point comparison.
For \(r_m=L^{-m}r_0\), blocking gives
\(\mathfrak m(r_0)\le\mathfrak M(r_m)\to0\), a contradiction.
Thus \(\alpha_K\) has one cluster value and converges.
Finally (CP2)--(CP3), backward drift and \(x_K^\natural=h\) make the
average nonlinear increment error tend to zero, so
\(\beta_K^\natural/K\to c_0\). Cesaro averaging identifies the unique
cluster value with \(c_0\). Without these two-run and uniform-envelope
premises, (CP9a) remains an additional comparison obligation.
With (CP9a), exact finite dilation and the moving-offset limit give
\[
 S^{[t]}_\ell(\mathbf f)
       =S^{[t-bc_0]}_{L^b\ell}(\mathbf f^{[b]})\qquad(b\ge1),
\]
when both offsets lie in their windows and the torus floors hold.
Indeed the same bare coupling has the second offset
\(t-[\beta_{K+b}^\natural-\beta_K^\natural]\), which converges to
\(t-bc_0\). Thus the original blocking target is preserved without
claiming it follows from one-step potential estimates.

Distinct offsets need the additional \(7e\) germ-dilation premise.
Fix a nonempty class of nonnegative \(C^2\) pairs: each test is supported
in a radius-\(r\) ball, the two supports have distance in \([r,6r]\),
and each test satisfies
\(\int f\ge\vartheta r^4\|f\|_\infty\) and
\(\operatorname{Lip}f\le\|f\|_\infty/(\vartheta r)\).
Define \(\sigma\approx_2\sigma'\) by an error
\(Cr^2\|f\|_\infty\|h\|_\infty\), uniformly on this class at all
sufficiently small \(r\). Set
\((D_\mu\sigma)_2(f,h)=\sigma_2(f(\cdot/\mu),h(\cdot/\mu))\).
Require \(\sigma_t\approx_2D_\mu\sigma_{t'}\), with
\(\mu=\exp((t-t')\log L/c_0)\), and uniform two-point upper and
lower envelopes of order \((\log(1/r))^{-2}\).
A germ cannot be fixed by \(D_\mu\), \(\mu\ne1\). For \(\mu>1\),
apply its error bound to one fixed rescaled shape, with normalized
value \(Q(r)\), at radii \(r\mu^{-a}\).
The upper envelope gives \(Q(r\mu^{-a})\to0\), so
\[
 Q(r)\le C\sum_{a\ge1}(r\mu^{-a})^2
                     =\frac{Cr^2}{\mu^2-1}.
\]
This contradicts the logarithmic lower envelope for arbitrarily
small \(r\). For \(\mu<1\), invert the dilation. Thus equality of two
different-offset limits would contradict rigidity. The comparison
retains its shape, placement, offset and torus-floor quantifiers;
it is an additional premise, not a consequence of matching alone.

\subsection*{Why a geometric relabelling needs another construction}

Take three test tuples of the same fixed correlation order, for example
three pairs, with observable values
\[
 F_K(t)=(1,\,1+t,\,1+\Pi_K)
       =\alpha+\beta t+\gamma\Pi_K,\qquad
 \Pi_K=K^{-q},\quad q>1 .
\]
Here \(\alpha=(1,1,1)\), \(\beta=(0,1,0)\),
\(\gamma=(0,0,1)\), with determinant one.
On \(|t|\le1/4\) the values are positive and uniformly bounded.
With \(\Xi_K(t)=h-t\) and \(\mathfrak n_K=\mathrm{id}\), all
summable matching inequalities hold; their increments are
\(O(K^{-q-1})\).
Suppose a reparametrization \(t\mapsto\tau_K(t)\) and one common
normalization \(a_K(t)\), bounded above and away from zero, made
their increments geometrically small. At a fixed \(t\), geometric
summability forces \(a_K\), the first component, to converge
geometrically to a nonzero limit. The third component minus the
first is \(a_K\Pi_K\), tends to zero, and would also be geometrically
small. This contradicts its polynomial lower bound.
The same argument allows arbitrary reparametrizations: their value
does not enter either the first component or the third-minus-first.
Only a common, nonvanishing field normalization is required. This is
a counterexample to the proposed implication from numerical inequalities,
not an assertion that these three numbers already form a physical
Yang--Mills state. Thus those inequalities alone do not supply the
stronger second-coordinate claim.
A proposed coordinate
\(\hat x_K=x_K^\natural-\sum_{i:\,i+\log_L(M\check R_i)\le K}
\operatorname{coeff}_A(\operatorname{tot}R^{(i)})\)
needs an actual compatible remainder-transfer and pointing theorem.
The finite scalar insertion results do not establish that theorem.

\subsection*{The optional Ward target on arbitrary defining sequences}

In this section \(S_{j,n}\) and \(S_n\), \(n\ge2\), are
\emph{connected} forms. Require convergence on every smooth compact
separated tuple on \(\mathbb R^4\), along the specified defining
sequence. On every fixed collection of separated compact support
neighborhoods require uniform local multilinear \(C^1\) bounds.
Mollification inside slightly enlarged separated neighborhoods then
extends the forms and their convergence to compact \(C^1\) tuples:
first approximate each entry, use the uniform telescoping bound,
and then pass to the limit for the smooth approximants.
Assume full signed-permutation symmetry \(W_4\) of the limit.
The specified lattice probability laws realize these connected forms;
the mixed cumulants below use those laws.
For torus approximants a fixed \(R\) must fit in one common chart
for all sufficiently large \(j\).
For \(\omega\in\mathfrak{so}(4)\), set
\(\delta_\omega f=-(\omega x)\cdot\nabla f\) and
\[
 D^s_{j,\omega}f(x)=
 \sum_\rho(\omega x)_\rho
       \frac{f(x+a_je_\rho)-f(x-a_je_\rho)}{2a_j}.
\]
Here are actual additional Ward inputs, independent of the angular route.
Choose smooth \(\chi_R=1\) on \(B_R\), zero outside \(B_{2R}\),
with derivatives bounded by \(C_rR^{-r}\).
For every \(j,\omega,R\), supply fixed annular local functions
\[
 \mathfrak B_{j,\omega,R}=\sum_{\Delta\in\mathcal A_{j,R}}
                              \mathfrak B^\Delta_{j,\omega,R}.
\]
There are at most \(CR^4\) physical unit boxes; each summand is a
specified gauge-invariant Borel function of links in a fixed enlargement
of its box, linear in \(\omega\), supported near
\(R-Ca_j\le|x|\le2R+Ca_j\).
The functions may contain the divergent Wilson weight.
They are supplied by the actual change of variables, independently of
the tests, not defined as their desired rotation defect.
Require the exact finite identity
\begin{equation}
 \sum_i S_{j,n}(f_1,\ldots,D^s_{j,\omega}f_i,\ldots,f_n)
  =\mathfrak U_{j,\omega,R}(\mathbf f)
   +\langle\mathfrak B_{j,\omega,R};
                  \mathcal O_j(f_1);\ldots;\mathcal O_j(f_n)\rangle^T .
 \tag{CP10}
\end{equation}
In (CP10), \(\mathcal O_j(\mathbf f)\) denotes the separate list
of cumulant arguments, not their coincident product. The bulk functional
\(\mathfrak U\) must also be the specified correlation from that
same Ward construction; it is not an arbitrary error chosen after the
boundary functions. For every fixed sufficiently large \(R\), require
\(\mathfrak U_{j,\omega,R}\to0\) and the analogous sum with
\(D^s_{j,\omega}f_i+\delta_\omega f_i\) to tend to zero.
These include the actual typed anisotropy and volume-transfer estimates
where needed, and are not proved by scalar source bounds.

Multilinear telescoping and local \(C^1\) continuity differentiate the
limit's rotation orbit on smooth tuples. Applying the assumed convergence
at the fixed derivative tests in (CP10), without differentiating a
cutoff limit, gives
\begin{equation}
 \lim_j\langle\mathfrak B_{j,\omega,R};\mathcal O_j(\mathbf f)\rangle^T
     =-\delta_\omega S_n(\mathbf f).
 \tag{CP11}
\end{equation}
It is independent of \(R\).
Consequently target \(0(e)\)'s condition
\begin{equation}
 \liminf_{R\to\infty}\limsup_{j\to\infty}
 \left|\langle\mathfrak B_{j,\omega,R};
                      \mathcal O_j(\mathbf f)\rangle^T\right|=0
 \tag{CP12}
\end{equation}
is equivalent to vanishing of this derivative.
The inner limit always has \(R\) fixed.

For completeness, one nonzero generator suffices.
Assume also the invariant one-point convention \(c\int f\) and
zeroth moment one, or use the centered functional. Let \(G\) be the
closed subgroup of orthogonal transformations preserving the connected
forms. Closedness follows from \(C^1\) continuity on separated tuples.
Equivalently, define \(G\) from the centered moment forms obtained
by the finite sum over partitions of these cumulants; the inverse
partition formula gives the same group. The invariant one-point
function \(c\int f\) makes the uncentered moment group the same.
Let \(\mathcal A\) consist of skew matrices whose entire one-parameter
group lies in \(G\). It is closed under real scalar multiplication
and addition: the finite-matrix product formula
\[
 e^{t(A+B)}=\lim_{r\to\infty}(e^{tA/r}e^{tB/r})^r
\]
follows by expanding each factor to second order and telescoping the
product error, which is \(O(r^{-1})\). Closedness of \(G\) then applies.
It is also stable under conjugation by \(W_4\).
For \(a<b\), sign matrices \(D_\epsilon=\operatorname{diag}(\epsilon)\)
give the projection
\[
 2^{-4}\sum_{\epsilon\in\{\pm1\}^4}
       \epsilon_a\epsilon_b D_\epsilon A D_\epsilon
                         =A_{ab}E_{ab}.
\]
Here \(E_{ab}=e_ae_b^{\mathsf T}-e_be_a^{\mathsf T}\) is the rotation
generator of the \((a,b)\) coordinate plane; the identity uses that \(A\)
is skew. Thus any nonzero element of \(\mathcal A\) supplies a coordinate-plane
generator; permutations supply all of them.
Successive plane rotations reduce any determinant-one orthogonal
matrix to the identity column by column. Hence \(G\) contains
\(\mathrm{SO}(4)\), and its coordinate reflections give all
\(\mathrm O(4)\).
Condition (CP12) for a nonzero \(\omega\), on every smooth tuple,
gives its whole one-parameter group by applying the zero derivative
also at each rotated tuple. Conversely full invariance makes (CP11)
zero for every \(\omega\).
No Lie-subgroup theorem is needed. Indeed an infinite closed \(G\)
has nonidentity \(g_r\to I\); the normalized skew logarithms have a
nonzero subsequential limit \(A\).
Integer powers \(g_r^{\lfloor t/\|\log g_r\|\rfloor}\to e^{tA}\)
supply the required one-parameter group. Therefore such a
\(W_4\)-invariant group is finite or all of \(\mathrm O(4)\).

\clearpage
\subsection*{Infrared meaning and the Hilbert-space consequence}

A uniform mixed-insertion estimate
\(|\langle\mathfrak B^\Delta;\mathcal O_j(\mathbf f)\rangle^T|
\le C D^{-4-\delta}\), \(\delta>0\), at distance \(D\) from the
fixed supports, implies (CP12): its shell sum is at most
\(CR^4R^{-4-\delta}\).
Decay \(D^{-4}\) alone yields only a bounded sum.
It also suffices that the sum of absolute mixed insertions tend to
zero along one sequence \(R_l\to\infty\), after \(\limsup_j\).
Ordinary per-ball density bounds do not give either estimate for
these different, weight-bearing insertions.
If an older survival argument is used, its event
\(\mathcal L_{\mathfrak c,k,a,K}\) means that cell \(\mathfrak c\)
is live at level \(k\), its component lineage starts at \(k-a\),
and it remains live at every level \(k,\ldots,K\).
Its weighted law is an additional history input.

Assume the relevant centered or recentered moment functional also
satisfies the HB quotient hypotheses. Once full scalar rotation
invariance is supplied, spatial
\(R\in\mathrm O(3)\) acts on the HB quotient by
\(V_R[F]=[(R\oplus1)F]\).
It preserves the positive form and null space, has inverse \(V_{R^{-1}}\),
fixes \(\Omega\), commutes with time evolution and conjugates a spatial
translation by \(x\) to one by \(Rx\).
Local \(C^1\) continuity proves strong continuity on dense test vectors,
and unitarity extends it to the completion.
The joint spectral measure is consequently covariant under
\((E,p)\mapsto(E,Rp)\).
For limits for which HB's cube inequalities have been established,
they give \(E\ge|p|\): use a countable dense
set of rotations, remove their common null set, and approximate a
rotation carrying any given \(p\) to \(|p|e_1\).
Centering is the map \(C_m\) proved in
Section~\ref{ch:euclidean-hilbert-bridge} and changes none of these
conclusions. In the alternative notation,
\(\mathsf{sh}_\xi=C_{-m_\xi}\),
\(m_\xi(f)=c_\xi\int f\), so its inverse is \(C_{m_\xi}\).
No operator for a general time-mixing Euclidean rotation is asserted.
External OS reconstruction, F4 identification of the full theory and
any mass-gap estimate remain separate. The periodic-state angular
proof retains its own mirror, cone and ultraviolet interface; the
optional Ward premises above do not silently enlarge that state class.

% END MODULE chapters/comparison_capsule_v7.tex


% BEGIN MODULE chapters/notation_index_v2.tex
\section{Notation and units}
\label{ch:notation-index}

This index records the meanings and units of the principal symbols. Definitions
of analytic domains and history data are given in the technical sections.
When a construction requires an auxiliary map not yet instantiated, its type
and required estimates are inputs; a symbol in this index does not establish
their existence. Hand-numbered references to the notation in candidate
arguments refer to this index and the corresponding detailed definition.

\subsection*{Lattices, fields, and action}

\begin{description}
\item[$G,\mathfrak g,N,d_G$] The gauge group is $G=\mathrm{SU}(2)$;
$\mathfrak g$ is represented by Hermitian traceless matrices with
$\langle X,Y\rangle=\frac12\operatorname{tr}(XY)$. Here $N=2$ and
$d_G=N^2-1=3$. In statements retaining $N$, the corresponding matrix group
and trace convention must be specified.
\item[$L,M$] $L>1$ is the fixed odd blocking integer. $M$ is the fixed
side of a localization cube in the units of its own lattice, with
$M=L^{m_0}$ for a fixed integer $m_0\ge0$. The graded ordinary-source
argument uses the cofinal compatible tori $\ell=2L^m$, $m\ge m_0$, so
$\ell/(Ma_s)$ is an even integer at every actual scale. The other
geometric parameters $M_1,M_2$ and the buffer constants are fixed before
the cutoff, fields, and sources are chosen.
\item[$K,j,s,a_s,\eta_j$] $K$ is the number of ultraviolet blocking
steps. Level $j$ has depth $s=K-j$. Physical spacing at depth $s$ is
$a_s=L^{-s}$. The fine spacing measured in level-$j$ units is
$\eta_j=L^{-j}$; it is not the physical spacing $a_j$.
\item[$\Lambda_{K,\ell},\mathbb T^{(j)}$]
$\Lambda_{K,\ell}=(L^{-K}\mathbb Z/\ell\mathbb Z)^4$ is a compatible
finite torus lattice. $\mathbb T^{(j)}=\Lambda_{K-j,\ell}$ is the level-$j$
lattice in physical units. A fixed-spacing infinite lattice is
$a\mathbb Z^4$. The chosen torus family may restrict $\ell$ further.
\item[$U_b,U_p,\partial U$] $U_b\in G$ is a positively oriented bond
variable; $U_{\bar b}=U_b^{-1}$. $U_p$ or $\partial U(p)$ is the ordered
transport around a plaquette. $x(p)$ is its barycentre. Gauge transformations
act by $U_b\mapsto h_{b_-}U_bh_{b_+}^{-1}$.
\item[$a_p,A(w,U),\beta$]
$a_p=1-\frac12\operatorname{Re}\operatorname{tr}U_p$,
$A(w,U)=\sum_pw(x(p))a_p$, and $A(U)=A(1,U)$. Wilson's weight is
$e^{-\beta A(U)}$. The inverse coupling convention is $\beta=g_0^{-2}$;
with the usual Wilson convention it equals $2N/g_{\mathrm{std}}^2$.
\item[$\mathcal O(f),S_n^T$]
$\mathcal O(f)=\sum_pf(x(p))a_p$ has no additional spacing prefactor.
$S_n^T$ is its joint connected correlation, defined by differentiating the
logarithm of the sourced partition function. Centering changes only the
first cumulant.
\end{description}

\subsection*{Averages, Gaussian coordinates, and propagators}

\begin{description}
\item[$\mathcal M,Q_k,Q_k^\sharp$] $\mathcal M$ is the chosen Borel
group-valued block average, analytic near regular real configurations.
$Q_k$ is the linearization of its plain iterate. $Q_k^\sharp$ is the
straightened linearization after the coarse gauge factor is removed.
They are different operators at a nontrivial background.
\item[$Q_k^{\mathrm{st}},Q_k'$] $Q_k^{\mathrm{st}}$ is a straight
covariant bond average. $Q_k'$ averages site functions. An adjoint always
uses the inner products of both source and target lattices.
\item[$d_0,d_1,d_\eta$] $d_0$ and $d_1$ denote the lattice gradient and
circulation in the cochain convention of the relevant theorem.
$d_\eta=\eta^{-1}(\mathrm{shift}-1)$ is the physical scalar difference in
unit blocks. The unscaled combinatorial difference is $\eta d_\eta$.
The free curvature comparison is unchanged by a common scaling of $d_1$,
because it uses the projection onto its range.
\item[$Q,P,A_\eta,G'$] In the scalar block calculation $Q$ is the
arithmetic mean, the fine inner product is $\eta^4\sum_x$, and the coarse
inner product is the unweighted sum. Then $Q^*$ lifts constants,
$P=Q^*Q$, $A_\eta=d_\eta^*d_\eta+P$, and $G'=A_\eta^{-1}$.
With plain inner products on both lattices the arithmetic mean's adjoint
is different. The unitary conversion is proved in the block-normalization section.
\item[$R,\mathbf R,R^{(j)}$] In scalar Landau-gauge formulae $R$ is
an orthogonal projection, for example onto $\ker(QG')$. Bold $\mathbf R$
is the large-field operation. $R^{(j)}$ is a local remainder activity.
These three uses are distinguished by their types, not interchangeable.
\item[$P_{\mathrm{ex}},K,\mathsf K$]
$P_{\mathrm{ex}}$ is the orthogonal projection onto exact lattice two-forms.
Its physical kernel is $a^{-4}P_{\mathrm{ex}}(p,q)$.
$K$ or $\mathsf K$ denotes the six-by-six continuum curvature covariance
at unit coupling. Away from zero it equals
$(\pi^2|x|^4)^{-1}\Lambda^2(I-2\widehat x\otimes\widehat x)$.
The Fourier distributions also contain contact terms.
\item[$U_j(V),\mathsf G,\mathsf H$] $U_j(V)$ is a constrained regular
background with specified coarse averages. The symbols $\mathsf G$ and
$\mathsf H$ denote a constrained inverse and a response operator in that
background. Existence, analytic domains, gauges and uniform bounds are
the background obligations A1--A4, which we must prove; the proved scalar
model is only part of them.
\item[\(\Phi_\Lambda(B),r_B,Z_y\)]
In Section~\ref{ch:pointed-window-remainder}, \(B\) is a dimensionless
coarse bond logarithm and \(\Phi_\Lambda(B)\) the actual
potential/current background pair on a collared box.
The response bound holds throughout the product datum polydisc
\(\sup_c|B(c)|<r_B\), not merely at zero.
The local output norm uses
\(Z_y=(\xi_yH,\xi_y^2\nabla H,\xi_y^3J)\).
Finite-box comparison also uses the actual compatible equations and
boundary rows in natural force norms, including their cell weights.
\end{description}

\subsection*{Running quantities and source domains}

\begin{description}
\item[$x_j,g_j,b_j$] $x_0=\beta$,
$x_j=x_{j-1}-b_j$, and $g_j=x_j^{-1/2}$. The actual increment includes
all contributions assigned to the marginal action coefficient. Positivity
of every individual $b_j$ is not assumed in the arithmetic theorem.
\item[$c_0,b_j^0$] $b_j^0$ is a designated one-loop increment.
The proposed leading slope is $c_0=11N^2\log L/(12\pi^2)$ in this
normalization. The numerical first-passage theorem assumes its interval
sum law and an error envelope; it does not calculate this coefficient.
\item[$\theta_j,x_j^\natural$]
$\theta_j=g^{-2}+c_0(K-j)$ is a prescribed frozen reference sequence.
In the technical frame estimates $\bar g_j=\theta_j^{-1/2}$ denotes its
auxiliary coupling, distinct from the running coupling $g_j=x_j^{-1/2}$.
The superscript $\natural$ denotes quantities computed with the auxiliary
choices frozen on that sequence. Holomorphy of this calibrated coordinate
and its comparison across cutoffs are additional analytic assertions.
\item[$K(\beta,g),\mathcal W_K(g)$] $K(\beta,g)$ is the first scale
at which $x_j$ enters the reference window of width $B_\sharp$ above
$g^{-2}$. $\mathcal W_K(g)$ is the set of bare couplings with that first
passage. The terminal value can range over the window.
\item[$w,c,z,f_i$] A source is a triple consisting of a constant $c$,
a finite real test family $(f_i)$, and complex coefficients $(z_i)$,
with function $w=c+\sum_i z_if_i$. Analytic bookkeeping uses the triple;
physical quantities must depend only on its relevant plaquette values.
\item[$\|\cdot\|_{\sharp,s},\|\cdot\|_{\flat,s}$]
$\|f\|_{\sharp,s}=\max(\|f\|_\infty,c_ML^{-s}\operatorname{Lip}f)$,
where $c_M$ absorbs the fixed block constant. The source-triple norm is
$|c|+\sum_i|z_i|\|f_i\|_{\sharp,s}$. The index zero is usually omitted.
$\mathcal W_r^{(s)}$ is the corresponding open source ball on each slice.
\item[$h_y^{(j)},\varphi_j,\zeta_j$] $h_y^{(j)}$ is a partition over
base points, with its anchor support and normalization stated in the source
domain section. $\varphi_j$ is a fixed coupling profile. The running shift
$\zeta_j(p;w)$ is defined there by the exact counterterm identity; it is
not an additional freely chosen bare coupling.
\item[$\varrho,\rho,\alpha_{0,j},\alpha_{1,j}$]
$\varrho$ and $\rho$ are specified source or coupling-disc radii in their
respective statements. $\alpha_{0,j},\alpha_{1,j}$ are radii of complex
field domains. A radius used for coefficient extraction cannot silently
be replaced by a smaller running radius.
\item[$p_0,\check R_j,\varepsilon_j,\delta_j$]
$p_0(g)=A_0(\log g^{-2})^{\mathfrak p_0}$ is a deposit scale.
$\check R_j$ is the specified localization and waiting scale, rounded to
a power of $L$. The curvature and fluctuation tests use the frame thresholds
$\varepsilon_j,\delta_j$ defined in the presentation section; they are fixed
throughout a source family.
\end{description}

\subsection*{Local families, histories, and analytic requirements}

\begin{description}
\item[$E^{(j)},R^{(j)},B^{(j)}$] These are, respectively, regular local
activities, small remainder activities, and boundary activities. A pointed
family also has a base-point index. Its companion is a specified family on
the infinite lattice; a finite-torus total alone does not define it.
\item[$\mathsf b(E),e_j,\mathcal E$] $\mathsf b(E)$ is the scalar
marginal-coefficient extraction defined for an actual regular pointed
member by \eqref{ph:extractor}, under Requirement B2's conditional
hypotheses. Lowercase $e_j(X;w)$ denotes a localized source
energy, and $\mathcal E(w)$ their sum. The vacuum constant $E$ without a
domain index is a different scalar normalization.
\item[$d_j(X),\kappa,\kappa_{\mathcal D}$]
$d_j(X)$ measures the size or spanning length of the localization domain
in its level's cubes. $\kappa$ is a decay exponent; $\kappa_{\mathcal D}$
is a domain-counting exponent. Mixed derivatives spanning separated source
supports can be summed only when the declared decay beats that entropy.
\item[$\Omega_j,\Lambda_j,Z_j$] These are the nested regular region,
its buffered inner part, and its large-field complement in a history.
$\Lambda_j^+$ may also include islands left by the operation.
Exact inclusions, buffers, tower maps and their domains belong to the
declared geometric data; the names alone impose no such conditions.
\item[$h,\underline h,\omega$] $h$ records all outcomes of a history.
$\underline h$ records its standing data after absorbed inner outcomes
are summed out. $\omega$ often records a local history inside one component.
A measure associated with a standing history need not be the law of a
predetermined block variable after a healing operation.
\item[$\Xi_\omega,\mu_\omega,\Phi_\omega,\Psi_\omega$]
$\Xi_\omega$ is a stored integration space and $\mu_\omega$ its positive
reference measure. $\Phi_\omega$ is the exponent depending on stored and
external variables; $\Psi_\omega$ is an additional source exponent. Their
suprema and oscillations are taken on the stated domains under the mass
integral when that is what the induction requires.
\item[$\Theta(h),q_a,\Sigma_k$] $\Theta(h)$ is a prepaid history
weight for deposits and space-time volume. $q_a=e^{-c_a}$ is an age-survival
rate; here the subscript is not a lattice spacing. $\Sigma_k$ is the aged
load defined in the weighted-domain section. A survival probability belongs
in an output bound and does not by itself enlarge a complex input radius.
\item[$\mathfrak C_k,\mathcal U_j^c$] $\mathfrak C_k$ denotes a class
of complete presentations and $\mathcal U_j^c$ a complex pair domain.
Membership requires all declared structural clauses, maps and estimates.
An independent current is a coordinate in the pair, not automatically the
current evaluated on its field coordinate. No closure theorem follows from
writing these symbols without the formation data.
\item[\(\mathcal G_z,\mathcal K_z,A_z\)]
\(\mathcal G_z\) is one member's complete pointed pair total at
base point \(z\), and \(\mathcal K_z\) its Hessian after composition
with the actual free response. It need not be a convolution.
\(A_z\) is that base point's midpoint tent action in
\eqref{ph:tent-action}; its extractor is one. The identity
\eqref{ph:own-counterterm} uses the coefficient of that same member.
The full spatial array may choose different members at different sites.
\item[\(\rho,\sigma,a,a_{\rm out}\) in the window proof]
Here \(\rho=\max(1,a_{\rm out}/a)\) is a ratio of field radii and
\(\sigma=L^{-(n-j)}\) a scale ratio; this \(\rho\) is not a Gibbs
state. The local H\"older exponent called \(\alpha\) in that proof
is distinct from an analyticity radius \(\alpha_{i,j}\).
The ordinary age bound \(\rho^2\le C(1+n-j)\) gives the finite
young-age fallback. Actual whole-core membership is still required.
\item[OSF and OH]
Section~\ref{ch:ordinary-source-formation-audit} treats the given
regular members, defects, remainders and reserved histories.
Section~\ref{ch:ordinary-history-stability} treats a sum of complex
label measures \(\tau_S(w)\), positive at zero. Its site count \(N\)
is unrelated to the gauge-group parameter \(N=2\).
Its weighted label sum and all-small-field action bound are inputs
to its stability/floor implication, not supplied by the final sum.
\item[\(\mathscr T,W_f,\alpha^2\) in the response lemma]
\(\mathscr T\) is a holomorphic formation map on a potential-norm
ball; \(W_f\) is the inserted source potential, of norm
\(O(\alpha^2\|f\|_\infty)\) on the complex field domain.
This potential norm is used in the finite insertion and the response
bootstrap of Sections~\ref{ch:finite-source-insertion}
and~\ref{ch:linear-source-bootstrap}, with the actual analytic room
described in Section~\ref{sec:substitution-room}. The ordinary
coefficient bound on a bare weight \(u\) is a different norm and does
not imply this potential bound.
\item[\(E(z),E_{\ell,K},A_K\) in the transfer lemma]
\(E(z)\) is a single cell normalizer common to the states being
compared. \(E_{\ell,K}\) is a finite-torus vacuum constant;
\(A_K\) is its hypothetical common companion density in
(ET8)--(ET10), which may diverge with \(K\).
The companion comparison is one finite-reference route, unnecessary
when an infinite reference already shares D1's all-state normalizer.
\end{description}

\subsection*{Limits, phases, and reconstruction}

\begin{description}
\item[$\mathfrak S,\sigma$] $\mathfrak S$ is the set of separated
test tuples, including the empty tuple. Its span is a vector space;
concatenation is used only when the resulting supports are separated.
$\sigma$ is a moment functional, not the connected family itself.
\item[$\mathfrak V^1(A),\mathsf d(X),\mathfrak D_L$]
$\mathfrak V^1(A)$ consists of real compactly supported $C^1$ functions
whose supports lie in $A$. A separated-tuple chart is a product of these
spaces on sets at positive mutual distance. $\mathsf d(X)$ is the minimum
support distance; an empty support contributes infinite distance and a
zero entry makes a multilinear value zero. The dense translation subgroup
is $\mathfrak D_L^4$, with $\mathfrak D_L=\bigcup_Ka_K\mathbb Z$.
\item[$\sigma_S$, two-state centering] If $S=(S_n)_{n\ge2}$ is the
connected family, $\sigma_S$ is the partition sum with all blocks of size
at least two, with empty value one and singleton value zero.
The different expression
$\rho\prod_i[\mathcal O_K(f_i)-\nu(\mathcal O_K(f_i))]$ is centered in
$\nu$, not necessarily in its integrating state $\rho$. Its limiting
first moment can be $c\int f$; its cumulants of order at least two are
those of the integrating state.
\item[$\mathcal S_{\ne},|\varphi|_k$]
$\mathcal S_{\ne}(\mathbb R^{4n})$ consists of Schwartz functions flat
with all derivatives on every coincidence diagonal. One seminorm is the
supremum of $(1+|x|^2)^{k/2}|\partial^\alpha\varphi(x)|$ for $|\alpha|\le k$.
This space is different from all of $\mathcal S(\mathbb R^{4n})$.
\item[GS geometry, $m(X),X^+$] At scale $a_s$, localization domains
are indexed by unions of $Ma_s$ cubes. The size $m(X)$ is the least
number of cubical tree edges needed to meet all their cubes. The source
dependence $X^+$ lies in a fixed number of neighboring cube layers.
The rooted, local ball-root and global weighted counts are (GS1)--(GS3).
Their sharper geometry is not implied by the weak abstract (ST6) count.
\item[D2, $k,B(R)$] D2 is the graded all-state moment and cumulant
requirement derived in Theorem~\ref{thm:graded-rooted-source} under D1 in
that actual geometry and G1. One may take $B(R)=C(R+1)^4$ and one fixed
integer factorial exponent $k$. The separation factor is
$\max(1,d^{-5n})$, explicitly depending on the number of fields.
\item[H1, $k_{\rm E},\kappa_{\rm E}$] H1 is the diagonal-flat growth
requirement derived from D2 by Theorem~\ref{thm:quantitative-flat-extension}.
The sufficient constants are $k_{\rm E}=20$ and
$\kappa_{\rm E}=k+32$, with seminorm $|\varphi|_{20n}$.
The QF proof includes constrained tensor density at every diagonal,
infinity and the positive-time boundary. It does not prove the ordinary
analytic source representation D1 or extend tests across diagonals.
\item[UI, PD, $\mathsf p,\wp$]
UI is Theorem~\ref{thm:uniform-one-field-insertion}: an insertion into
a positive-time gap has a Hilbert operator bound with exponent
$\mathsf p$ independent of the number of outside fields. PD is
Theorem~\ref{thm:plate-density-from-insertion}; its analytic $L^1$
density has exponent $\wp=\mathsf p+5$ and compatible padded point
operators. These follow from the ordinary repeated-source input, not
from graded moment bounds alone.
\item[$G_{\varphi,\psi},\mu_\varphi,\Gamma_\varphi,L_s$]
$G_{\varphi,\psi}$ is the two-point kernel after smearing the two
transverse variables with real compact labels. For an equal label its
coordinate spectral measure is the positive measure $\mu_\varphi$ on
$\{E\ge|q|\}$, with
$G_\varphi(t,x)=\int e^{-tE+ixq}\,d\mu_\varphi$ and
$\Gamma_\varphi(t)=\int e^{-tE}\,d\mu_\varphi<\infty$ for $t>0$.
$L_s(E,q)=(E\cosh s+q\sinh s,E\sinh s+q\cosh s)$ is the spectral
boost used in the proof. The local letter $\nu$ in that proof denotes a
diagonal-frame spectral measure, rather than a Gibbs state; its cone
support is proved from the identity of measures (TU27).
\item[$\vartheta_*,u(d),c_m,b_n$]
The ultraviolet theorem fixes $\vartheta_*=1/12000$ and the full
positive-shape class (TU9). Its input is (TU10), with $u(d)\to0$,
uniform over that class, positions and approximating states.
The class includes the explicit mass condition
\(\int f\ge\vartheta_*d^4\|f\|_\infty\) for each nonnegative
radius-\(d\) test, in addition to the Lipschitz and separation bounds.
It is not inferred from the Lipschitz condition alone.
F2--F3 imply it only with those same quantifiers.
The angular coefficients transform as $c_m(R_t\mathbf y)=e^{4imt}c_m(\mathbf y)$.
$b_n$ is the largest absolute nonzero nonconstant frequency at order $n$,
with value zero when there is none. Polygon divisibility and convexity
give $b_{2k}=kb_2$; the ultraviolet argument gives $b_2=0$.
\item[G2 and its scope]
G2 is now the derived scalar rotation requirement supplied by
UI--PD--AK--UV--HC under D1 in its actual rooted geometry, G1, the
exact cube cone and the stated ultraviolet law. It applies to the
periodic-state limits and canonical components carrying that interface.
Arbitrary growing cubic-torus limits require their own mirror/cone
justification. Typed Ward and matched-cutoff estimates still belong to
the stronger retained routes that use them.
\item[$\mathcal G^{\rm per},\mathcal G^{\rm tor}$]
At fixed spacing and coupling the first is the union of Wilson DLR
classes with finite-index translation period groups; the second consists
of weak limits of periodically lifted torus measures as volume tends
to infinity. $\mathcal G^{\rm ex}$ denotes the translation-ergodic members
of the periodic class, not unrestricted DLR extremality.
A periodic lift before this limit is not asserted to be
an infinite-volume DLR measure. These classes, growing-torus limits and
conditionally tuned full sequences have distinct definitions in
Section~\ref{ch:continuum-compactness}.
\item[$f^{(\ell)}$, the common periodization chart]
If the union of all supports under consideration lies in one Euclidean
ball of radius $R<\ell/4$, periodization is defined by
$f^{(\ell)}\circ\pi_\ell=\sum_{k\in\mathbb Z^4}f(\,\cdot+\ell k)$.
On this common chart it preserves the relevant pairwise support distances
and norms. Independent small balls placed arbitrarily in $\mathbb R^4$
do not give this conclusion modulo $\ell$. Their separate radius bound
is the different, position-uniform convention in the moment estimates.
\item[$\mathcal G_j,\mathcal E_j,D_j$] $\mathcal G_j$ is the compact
translation-invariant DLR class at fixed spacing and coupling.
$\mathcal E_j$ is its ergodic, equivalently class-extreme, subset.
$D_j$ is the decomposition law obtained by conditioning the centering state
on the invariant sigma-field. Full DLR extremality without invariance is
not identified with this notion here.
\item[$\mathcal I,\operatorname{Pl}_\nu,\mathcal N_{\nu,r}$]
$\mathcal I$ is the full translation-invariant sigma-field.
$\operatorname{Pl}_\nu(A,B)$ is the inner product of their conditional
expectations on that field. $\mathcal N_{\nu,r}$ is the reflected size in
the upstream lattice law G3; it replaces divergent supremum norms of
cutoff field products in the phase limit.
\item[$\mathsf X,\lambda,\sigma_\xi$]
$\mathsf X$ is the compact full-support phase-coordinate space constructed
from the weak limits of lattice evaluation laws, $\lambda$ its probability
measure, and $\sigma_\xi$ its continuous component evaluations. These are
conclusions of the conditional phase theorem, not assumed continuum data.
\item[$\theta,\mathcal H,H,P,\Omega$] $\theta$ is the specified
Euclidean reflection. $\mathcal H$ is the completed quotient of its positive
form. $\Omega$ is the class of the empty tuple, $H$ the nonnegative time
generator, and $P$ the spatial momentum tuple when the corresponding
continuous groups have been constructed. $P_0$ projects onto $\ker H$.
\item[$\mathfrak U,\mathfrak B,W_4$] $\mathfrak U$ denotes a rotation
defect or limiting flux in the Ward argument; $\mathfrak B$ is its insertion.
$W_4$ is the group of signed coordinate permutations. Its invariance alone
does not imply invariance under $\mathrm O(4)$.
\end{description}

% END MODULE chapters/notation_index_v2.tex


% BEGIN MODULE chapters/argument_index.tex
\section{Index of arguments and proof obligations}

These code families organize the proof across the book. They locate
retained statements, detailed proofs of referenced results, and the
statements that we must still prove. Their exact forms are given at the
cited statement and in the dependency catalogue. Outlined supporting
arguments remain distinct from those obligations.

\subsection*{Named statements in the proof outline}

0, 1, 1a.1.1, 1a.1.10, 1a.1.11, 1a.1.12, 1a.1.13, 1a.1.14, 1a.1.2, 1a.1.3, 1a.1.4, 1a.1.5, 1a.1.6, 1a.1.7, 1a.1.9, 1a.2.1, 1a.2.10, 1a.2.2, 1a.2.3, 1a.2.4, 1a.2.7, 1a.2.8, 1a.2.9, 1a.3.1, 1a.3.11, 1a.3.12, 1a.3.13, 1a.3.3, 1a.3.4, 1a.3.6, 1a.3.7, 1a.3.8, 1a.3.9, 2, 2.1, 2.11, 2.2, 2.3, 2.4, 2.5, 2.6, 2.7, 3, 3.1, 3.2, 3.3, 3.4, 3.5, 3.6, 3.7, 3.8, 3.9, 4, 4.1, 4.13, 4.14, 4.15, 4.2, 4.3, 4.4, 4.5, 4.6, 4.7, 5, 6, 7.1, 7.2, 7.3, 7.4.

\subsection*{Detailed proofs of referenced statements}

\noindent 4.10: chessboard domination for arbitrary periodic states, Section~\ref{ch:all-state-source-transfer}.

\noindent 4.12: uniform one-point differences, Section~\ref{ch:all-state-source-transfer}.

\noindent 4.17: uniform one-field Hilbert insertion, Section~\ref{ch:uniform-insertion}.

\noindent 4.18: analytic insertion density and the plate-bound implication under its stated hypotheses, Section~\ref{ch:plate-density-from-insertion}.

\subsection*{Statements still to be proved, by package}

The ordinary history and source tasks 2.8, 2.10, 2.13 and 4.11 have
proof outlines in Sections~\ref{sec:actual-history-closure} and
\ref{sec:actual-ordinary-source-proof}. The extraction and transfer
tasks 7.5, 7.6 and 7a.1.8 are developed in
Sections~\ref{ym4:new-clean-box}--\ref{ym4:new-gaussian-segment}
and~\ref{sec:finite-reference-normalizer}. The stronger aged
comparison 2.12 remains separate. The one-loop calculation is in
Section~\ref{ym4:new-one-loop-reduction}.

In Package E, the code 8 refers both to Ba{\l}aban's Theorem 8, which he
states without proof and which we must prove, and to Proposition 8 of
Commun. Math. Phys. 102 (1985) 277--309, which is cited. In Package A--H,
the code 1.2 is Proposition 1.2 of Commun. Math. Phys. 95 (1984) 17--40,
which is cited. The following list retains the task identifiers;
the preceding paragraph gives the locations of the supplied outlines.

\paragraph{Package A.} 1a.3.2, 1a.3.5.

\paragraph{Package C.} 2.10, 2.12, 2.13, 2.8.

\paragraph{Package D.} 4.11.

\paragraph{Package E.} 8.

\paragraph{Package F.} 7.5, 7.6, 7a.1.8.

\paragraph{Package G.} 9.1, 9.10, 9.11, 9.3, 9.7.

\paragraph{Package A--H.} 1.2, 1.6.

\subsection*{Formula and statement navigation}

The first component of a hand-numbered code records its mathematical theme:
1a--1c backgrounds and propagators; 1d--1e local interactions; 2 histories;
3 flow; 4 sources; 5 limits; 6 Euclidean structure; 7 short distance;
8 comparison; 9 rotations; 10 phases and reconstruction. These codes retain
their meaning when a theme occupies separate book sections. A grouped
reference is read with the hypotheses and scope of its stated result.
% END MODULE chapters/argument_index.tex


% BEGIN MODULE chapters/bibliography.tex
\begin{thebibliography}{OSeiler78}\small

\bibitem[B82a]{B82a} T. Ba\l{}aban, $(\mathrm{Higgs})_{2,3}$ quantum fields in a finite volume. I. A lower bound, \emph{Communications in Mathematical Physics} \textbf{85} (1982), 603--626. \url{https://doi.org/10.1007/BF01403506}.

\bibitem[B82b]{B82b} T. Ba\l{}aban, $(\mathrm{Higgs})_{2,3}$ quantum fields in a finite volume. II. An upper bound, \emph{Communications in Mathematical Physics} \textbf{86} (1982), 555--594. \url{https://doi.org/10.1007/BF01214890}.

\bibitem[B83a]{B83a} T. Ba\l{}aban, $(\mathrm{Higgs})_{2,3}$ quantum fields in a finite volume. III. Renormalization, \emph{Communications in Mathematical Physics} \textbf{88} (1983), 411--445. \url{https://doi.org/10.1007/BF01213217}.

\bibitem[B83b]{B83b} T. Ba\l{}aban, Regularity and decay of lattice Green's functions, \emph{Communications in Mathematical Physics} \textbf{89} (1983), 571--597. \url{https://doi.org/10.1007/BF01214744}.

\bibitem[B84a]{B84a} T. Ba\l{}aban, Propagators and renormalization transformations for lattice gauge theories. I, \emph{Communications in Mathematical Physics} \textbf{95} (1984), 17--40. \url{https://doi.org/10.1007/BF01215753}.

\bibitem[B84b]{B84b} T. Ba\l{}aban, Propagators and renormalization transformations for lattice gauge theories. II, \emph{Communications in Mathematical Physics} \textbf{96} (1984), 223--250. \url{https://doi.org/10.1007/BF01240221}.

\bibitem[B85a]{B85a} T. Ba\l{}aban, Averaging operations for lattice gauge theories, \emph{Communications in Mathematical Physics} \textbf{98} (1985), 17--51. \url{https://doi.org/10.1007/BF01211042}.

\bibitem[B85b]{B85b} T. Ba\l{}aban, Spaces of regular gauge field configurations on a lattice and gauge fixing conditions, \emph{Communications in Mathematical Physics} \textbf{99} (1985), 75--102. \url{https://doi.org/10.1007/BF01466594}.

\bibitem[B85c]{B85c} T. Ba\l{}aban, Propagators for lattice gauge theories in a background field, \emph{Communications in Mathematical Physics} \textbf{99} (1985), 389--434. \url{https://doi.org/10.1007/BF01240355}.

\bibitem[B85d]{B85d} T. Ba\l{}aban, Ultraviolet stability of three-dimensional lattice pure gauge field theories, \emph{Communications in Mathematical Physics} \textbf{102} (1985), 255--275. \url{https://doi.org/10.1007/BF01229380}.

\bibitem[B85e]{B85e} T. Ba\l{}aban, The variational problem and background fields in renormalization group method for lattice gauge theories, \emph{Communications in Mathematical Physics} \textbf{102} (1985), 277--309. \url{https://doi.org/10.1007/BF01229381}.

\bibitem[B87]{B87} T. Ba\l{}aban, Renormalization group approach to lattice gauge field theories. I. Generation of effective actions in a small field approximation and a coupling constant renormalization in four dimensions, \emph{Communications in Mathematical Physics} \textbf{109} (1987), 249--301. \url{https://doi.org/10.1007/BF01215223}.

\bibitem[B88a]{B88a} T. Ba\l{}aban, Renormalization group approach to lattice gauge field theories. II. Cluster expansions, \emph{Communications in Mathematical Physics} \textbf{116} (1988), 1--22. \url{https://doi.org/10.1007/BF01239022}.

\bibitem[B88b]{B88b} T. Ba\l{}aban, Convergent renormalization expansions for lattice gauge theories, \emph{Communications in Mathematical Physics} \textbf{119} (1988), 243--285. \url{https://doi.org/10.1007/BF01217741}.

\bibitem[B89a]{B89a} T. Ba\l{}aban, Large field renormalization. I. The basic step of the $\mathbb R$ operation, \emph{Communications in Mathematical Physics} \textbf{122} (1989), 175--202. \url{https://doi.org/10.1007/BF01257412}.

\bibitem[B89b]{B89b} T. Ba\l{}aban, Large field renormalization. II. Localization, exponentiation, and bounds for the $R$ operation, \emph{Communications in Mathematical Physics} \textbf{122} (1989), 355--392. \url{https://doi.org/10.1007/BF01238433}.

\bibitem[Wilson74]{Wilson74} K. G. Wilson, Confinement of quarks, \emph{Physical Review D} \textbf{10} (1974), 2445--2459. \url{https://doi.org/10.1103/PhysRevD.10.2445}.

\bibitem[OS73]{OS73} K. Osterwalder and R. Schrader, Axioms for Euclidean Green's functions, \emph{Communications in Mathematical Physics} \textbf{31} (1973), 83--112. \url{https://doi.org/10.1007/BF01645738}.

\bibitem[OS75]{OS75} K. Osterwalder and R. Schrader, Axioms for Euclidean Green's functions II, \emph{Communications in Mathematical Physics} \textbf{42} (1975), 281--305. \url{https://doi.org/10.1007/BF01608978}.

\bibitem[OSeiler78]{OSeiler78} K. Osterwalder and E. Seiler, Gauge field theories on a lattice, \emph{Annals of Physics} \textbf{110} (1978), 440--471. \url{https://doi.org/10.1016/0003-4916(78)90039-8}.

\bibitem[FILS78]{FILS78} J. Fr\"ohlich, R. Israel, E. H. Lieb and B. Simon, Phase transitions and reflection positivity. I. General theory and long range lattice models, \emph{Communications in Mathematical Physics} \textbf{62} (1978), 1--34. \url{https://doi.org/10.1007/BF01940327}.

\bibitem[KP86]{KP86} R. Koteck\'y and D. Preiss, Cluster expansion for abstract polymer models, \emph{Communications in Mathematical Physics} \textbf{103} (1986), 491--498. \url{https://doi.org/10.1007/BF01211762}.

\bibitem[JW]{JW} A. Jaffe and E. Witten, \emph{Quantum Yang--Mills theory}, official problem description, Clay Mathematics Institute. \url{https://www.claymath.org/wp-content/uploads/2022/06/yangmills.pdf}.

\bibitem[Kallenberg02]{Kallenberg2002}
O. Kallenberg, \emph{Foundations of Modern Probability}, second edition,
Probability and Its Applications, Springer, New York, 2002.
\url{https://doi.org/10.1007/978-1-4757-4015-8}.

\end{thebibliography}

% END MODULE chapters/bibliography.tex

\end{document}
